%!TEX program = lualatex
\documentclass[letterpaper, 11pt]{report}
\usepackage{fontspec}
\usepackage{fancyhdr}

%%% MARGINS %%%
% use this for normal letter documents
% \usepackage[
%     letterpaper,
%     margin=1in,
%     includehead
% ]{geometry}

%use this for bookletized pdfs
\usepackage[
    paperwidth=6.875in,
    paperheight=10.625in,
    margin=0.75in,     % Now your margins will be exactly what you type
    includehead
]{geometry}

%%% FONT %%%
\setmainfont{Liberation Serif}

%%% SPEAKER \& TALK VARIABLES %%%
\newcommand{\speaker}{}
\newcommand{\talk}{}
\newcommand{\talkdatestr}{}
\newcommand{\setspeaker}[1]{\renewcommand{\speaker}{#1}}
\newcommand{\settalk}[1]{\renewcommand{\talk}{#1}}
\newcommand{\settalkdate}[1]{\renewcommand{\talkdatestr}{#1}}

%%% HEADER \& FOOTER SETUP %%%
\pagestyle{fancy}
\fancyhf{}

\fancyhead[L]{\textbf{\talk}}
\fancyhead[R]{\textbf{\speaker}}

\renewcommand{\headrulewidth}{0.4pt}

%%% CUSTOM HEADER CLASS FOR TALK TITLE %%%
\newcommand{\talkheading}[2]{%
    \vspace{1em}
    \noindent{\huge\bfseries #1}
    \par\nopagebreak\vspace{0.25em}
    \noindent{\normalsize\bfseries #2\ifx\talkdatestr\empty\else,\ \talkdatestr\fi}
    \par\nopagebreak\vspace{0.5em}
}

%%% IN-TALK HEADERS
\newcommand{\intalkheader}[1]{
    \vspace{0.75em}
    \noindent{\normalsize\bfseries #1}
    \par\nopagebreak % Prevents a break immediately after the text
    \vspace{0.25em}
    \nopagebreak     % Prevents a break after the whitespace
}

\begin{document}
\newpage
\settalk{Lost Battalions}
\setspeaker{Thomas S. Monson}
\settalkdate{April 1971}
\talkheading{Lost Battalions}{Thomas S. Monson}

This past November I stood on a very old bridge which spans the River Somme as it makes its steady but unhurried way through the heartland of France. Suddenly I realized that fifty-two years had come—then gone—since the signing of the Armistice of 1918 and the termination of the Great War. I tried to imagine what the River Somme looked like fifty-two years before. How many thousands of soldiers had crossed this same bridge? Some came back. For others, the Somme was truly a river of no return. For the battlefields of Vimy Ridge, Armentieres, and Nueve Chappelle took a hideous toll of human life. Acres of neat, white crosses serve as an unforgettable reminder.

\begin{verse}
“In Flanders fields the poppies blow\\
Between the crosses, row on row,\\
That mark our place; and in the sky\\
The larks, still bravely singing, fly\\
Scarce heard amid the guns below.
\end{verse}

—John McCrae

I found myself saying softly, “How strange that war brings forth the savagery of conflict, yet inspires brave deeds of courage—some prompted by love.”

As a boy, I enjoyed reading the account of the “lost battalion.” The “lost battalion” was a unit of the 77th Infantry Division in World War I. During the Meuse-Argonne offensive, a major led this battalion through a gap in the enemy lines, but the troops on the flanks were unable to advance. An entire battalion was surrounded. Food and water were short; casualties could not be evacuated. Hurled back were repeated attacks. Ignored were notes from the enemy requesting the battalion to surrender. Newspapers heralded the battalion’s tenacity. Men of vision pondered its fate. After a brief but desperate period of total isolation, other units of the 77th Division advanced and relieved the “lost battalion.” Correspondents noted in their dispatches that the relieving forces seemed bent on a crusade of love to rescue their comrades in arms. Men volunteered more readily, fought more gallantly, and died more bravely. A fitting tribute echoed from that ageless sermon preached on the Mount of Olives: “Greater love hath no man than this, that a man lay down his life for his friends.” (John 15:13.)

Forgotten is the plight of the “lost battalion.” Unremembered is the terrible price paid for its rescue. But let us turn from the past and survey the present. Are there “lost battalions” even today? If so, what is our responsibility to rescue them? Their members may not wear clothes of khaki brown nor march to the sound of drums. But they share the same doubt, feel the same despair, and know the same disillusionment that isolation brings.

There are the “lost battalions” of the handicapped, even the lame, the speechless, and the sightless. Have you experienced the frustration of wanting but not knowing how to help the individual who walks stiffly behind his Seeing Eye canine companion, or moves with measured step to the tap, tap, tap of a white cane? There are many who are lost in this trackless desert of darkness.

If you desire to see a rescue operation of a “lost battalion,” visit your city’s center for the blind and witness the selfless service of those who read to those who can’t. Observe the skills that are taught the handicapped. Be inspired by the efforts put forth in their behalf to enable them to secure meaningful employment.

Those who labor so willingly and give so generously to those who have lost so tragically find ample reward in the light that they bring into the lives of the sightless.

Do we appreciate the joy of a blind person as his nimble fingers pass quickly over the pages of the Braille edition of the New Testament? He pauses at the twelfth chapter of John and contemplates the depth of meaning in the promise of the Prince of Peace: “I am come a light into the world, that whosoever believeth on me should not abide in darkness.” (John 12:46.)

Consider the “lost battalions” of the aged, the widowed, the sick. All too often they are found in the parched and desolate wilderness of isolation called loneliness. When youth departs, when health declines, when vigor wanes, when the light of hope flickers ever so dimly, the members of these vast “lost battalions” can be succored and sustained by the hand that helps and the heart that knows compassion.

In Brooklyn, New York, there presides today in a branch of the Church a young man who, as a boy of thirteen, led a successful rescue of such persons in Salt Lake City. He and his companions lived in a neighborhood in which resided many elderly widows of limited means. All the year long, the boys had saved and planned for a glorious Christmas party. They were thinking of themselves, until the Christmas spirit prompted them to think of others. Frank, as their leader, suggested to his companions that the funds they had accumulated so carefully be used not for the planned party, but rather for the benefit of three elderly widows who resided together. The boys made their plans. As their bishop, I needed but to follow.

With the enthusiasm of a new adventure, the boys purchased a giant roasting chicken, the potatoes, the vegetables, the cranberries, and all that comprises the traditional Christmas feast. To the widows’ home they went carrying their gifts of treasure. Through the snow and up the path to the tumbledown porch they came. A knock at the door, the sound of slow footsteps, and then they met.

In the unmelodic voices characteristic of thirteen-year-olds, the boys sang “Silent night, holy night; all is calm, all is bright.” They then presented their gifts. Angels on that glorious night of long ago sang no more beautifully, nor did wise men present gifts of greater meaning.

I gazed at the faces of those wonderful women and thought to myself: “Somebody’s mother.” I then looked on the countenances of those noble boys and reflected: “Somebody’s son.” There then passed through my mind the words of the immortal poem by Mary Dow Brine:

\begin{verse}
“The woman was old and ragged and gray\\
And bent with the chill of the Winter’s day.\\
The street was wet with a recent snow,\\
And the woman’s feet were aged and slow.\\
She stood at the crossing and waited long,\\
Alone, uncared for, amid the throng\\
Of human beings who passed her by\\
Nor heeded the glance of her anxious eye.
\end{verse}

What was the message of the Master? “Inasmuch as ye have done it unto one of the least of these … ye have done it unto me.” (Matt. 25:40.)

There are other “lost battalions” comprised of mothers and fathers, sons and daughters, who have, through thoughtless comment, isolated themselves from one another. An account of how such a tragedy was narrowly averted is this occurrence in the life of a lad we shall call Jack.

Throughout Jack’s life, he and his father had many serious arguments. One day, when Jack was seventeen, they had a particularly violent one. Jack said to his father: “This is the straw that breaks the camel’s back. I’m leaving home, and I shall never return.” So saying, he went to the house and packed a bag. His mother begged him to stay, but he was too angry to listen. He left her crying at the doorway.

Leaving the yard, he was about to pass through the gate when he heard his father call to him: “Jack, I know that a large share of the blame for your leaving rests with me. For this I am truly sorry. I want you to know that if you should ever wish to return home, you’ll always be welcome. And I’ll try to be a better father to you. I want you to know that I’ll always love you.”

Jack said nothing but went to the bus station and bought a ticket to a distant point. As he sat in the bus watching the miles go by, he commenced to think about the words of his father. He began to realize how much love it had required for him to do what he had done. Dad had apologized. He had invited him back and had left the words ringing in the summer air, “I love you.”

It was then that Jack realized that the next move was up to him. He knew that the only way he could ever find peace with himself was to demonstrate to his father the same kind of maturity, goodness, and love that dad had shown toward him. Jack got off the bus. He bought a return ticket to home and went back.

He arrived shortly after midnight, entered the house, and turned on the light. There in the rocking chair sat his father, his head in his hands. As he looked up and saw Jack, he rose from the chair and they rushed into each other’s arms. Jack often said, “Those last years that I was home were among the happiest of my life.”

We could say here was a boy who overnight became a man. Here was a father who, suppressing passion and bridling pride, rescued his son before he became one of that vast “lost battalion” resulting from fractured families and shattered homes. Love was the binding band, the healing balm. Love—so often felt; so seldom expressed.

From Mt. Sinai there thunders in our ears, “Honour thy father and thy mother.” (Ex. 20:12.) And later, from that same God, the injunction, “… live together in love.” (D\&C 42:45.)

There are other “lost battalions.” Some struggle in the jungles of sin, some wander in the wilderness of ignorance. In reality, each one of us is numbered in what could well have been the lost battalion of mankind, even a battalion doomed to everlasting death.

“… by man came death. … For as in Adam all die.” (1 Cor. 15:21–22.) Each of us is a partaker of the experience called death. None escapes. Were we to remain unrescued, lost would be paradise sought. Lost would be family loved. Lost would be friends remembered. Realizing this truth, we begin to appreciate the supreme joy which accompanied the birth of the Savior of the world. How glorious the pronouncement of the angel: Behold, a virgin “shall bring forth a son, and thou shalt call his name JESUS: for he shall save his people from their sins.” (Matt. 1:21.)

While the rivers of France witnessed the advance of those who rescued the “lost battalion” in World War I, so did yet another river witness the commencement of the formal ministry of a universal rescuer, even a divine redeemer. The scripture records, “And there came a voice from heaven, saying, Thou art my beloved Son, in whom I am well pleased.” (Mark 1:11.)

Today, only ruins remain of Capernaum, that city by the lakeshore, heart of the Savior’s Galilean ministry. Here he preached in the synagogue, taught by the seaside, and healed in the homes.

On one significant occasion, Jesus (Luke 4:18) took a text from Isaiah: “The Spirit of the Lord God is upon me; because the Lord hath anointed me to preach good tidings unto the meek; he hath sent me to bind up the brokenhearted, to proclaim liberty to the captives, and the opening of the prison to them that are bound” (Isa. 61:1), a clear pronouncement of a divine plan to rescue the “lost battalion” to which we belong.

But Jesus’ preaching in Galilee had been merely prelude. The Son of Man had always had a dread rendezvous to keep on a hill called Golgotha.

Arrested in the Garden of Gethsemane after the Last Supper, deserted by his disciples, spat upon, tried, and humiliated, Jesus staggered under his great cross toward Calvary. He progressed from triumph—to betrayal—to torture—to death on the cross.

In the words of the hymn, “… the scene was changed; the morn was cold and chill, as the shadow of a cross arose upon a lonely hill.” For us our Heavenly Father gave his Son. For us our Elder Brother gave his life.

At the last moment the Master could have turned back. But he did not. He passed beneath all things that he might save all things—the human race, the earth, and all the life that ever inhabited it.

No words in Christendom mean more to me than those spoken by the angel to the weeping Mary Magdalene and the other Mary as they approached the tomb to care for the body of their Lord: “Why seek ye the living among the dead? He is not here, but is risen.” (Luke 24:5–6.)

With this pronouncement, the “lost battalion” of mankind—those who have lived and died, those who now live and one day will die, and those yet to be born and yet to die—this battalion of humanity lost had just been rescued.

Of him who delivered each of us from endless death, I testify he is a teacher of truth—but he is more than a teacher. He is the exemplar of the perfect life—but he is more than an exemplar. He is the great physician—but he is more than a physician. He who rescued the “lost battalion” of mankind is the literal Savior of the world, the Son of God, the Prince of Peace, the Holy One of Israel, even the risen Lord, who declared, “I am the first and the last; I am he who liveth, I am he who was slain; I am your advocate with the Father.” (D\&C 110:4.)

As his witness I testify to you that he lives. In the name of Jesus Christ. Amen.

\newpage
\settalk{With Hand and Heart}
\setspeaker{Thomas S. Monson}
\settalkdate{October 1971}
\talkheading{With Hand and Heart}{Thomas S. Monson}

Yesterday each person assembled in this historic Tabernacle was given the privilege to raise his right hand to sustain, in the positions to which they have been called, the leadership of the Church. The upraised hand is an outward expression of an inner feeling. As one raises his hand, he pledges his heart.

The Master frequently spoke of hand and heart. In a revelation given through the Prophet Joseph Smith at Hiram, Ohio, in March 1832, he counseled: “… be faithful; stand in the office which I have appointed unto you; succor the weak, lift up the hands which hang down, and strengthen the feeble knees.

“And if thou art faithful unto the end thou shalt have a crown of immortality, and eternal life in the mansions which I have prepared in the house of my Father.” (D\&C 81:5–6.)

As I ponder his words, I can almost hear the shuffle of sandaled feet, the murmurs of astonishment from listeners as they echo from Capernaum’s peaceful scene. Here multitudes crowded around Jesus, bringing the sick to be healed. A palsied man picked up his bed and walked, and a Roman centurion’s faith restored his servant’s health.

Not only by precept did Jesus teach, but also by example. He was faithful to his divine mission. He stretched forth his hand that others might be lifted toward God.

At Galilee there came to him a leper who pleaded: “Lord, if thou wilt, thou canst make me clean. And Jesus put forth his hand, and touched him, saying, I will; be thou clean. And immediately his leprosy was cleansed.” (Matt. 8:2–3.) The hand of Jesus was not polluted by touching the leper’s body, but the leper’s body was cleansed by the touch of that holy hand.

In Capernaum, at the house of Peter, yet another example was provided. The mother of Peter’s wife lay sick of a fever. The sacred record reveals that Jesus came “and took her by the hand, and lifted her up; and immediately the fever left her. …” (Mark 1:31.)

So it was with the daughter of Jairus, a ruler of the synagogue. Each parent can appreciate the feelings of Jairus as he sought the Lord, and, upon finding him, fell at his feet and pleaded, “My little daughter lieth at the point of death: I pray thee, come and lay thy hands on her, that she may be healed; and she shall live.” (Mark 5:23.)

“While he yet spake, there cometh one from the [ruler’s] house, saying to him, Thy daughter is dead; trouble not the Master.

“But when Jesus heard it, he answered him, saying, Fear not: believe only, and she shall be made whole.” Parents wept. Others mourned. Jesus declared: “Weep not; she is not dead, but sleepeth.

“[He] … took her by the hand, and called, saying: Maid, arise.

“And her spirit came again, and she arose straightway. …” (Luke 8:49–50, 52, 54–55.)

Once again, the Lord had stretched forth his hand to take the hand of another.

The beloved apostles noted well his example. He lived not so to be ministered unto, but to minister; not to receive, but to give; not to save his life, but to pour it out for others.

If they would see the star that should at once direct their feet and influence their destiny, they must look for it, not in the changing skies or outward circumstance, but each in the depth of his own heart and after the pattern provided by the Master.

Reflect for a moment on the experience of Peter at the gate Beautiful of the temple. One sympathizes with the plight of the man lame from birth who each day was carried to the temple gate that he might ask alms of all who entered. That he asked alms of Peter and John as these two brethren approached indicates that he regarded them no differently from scores of others who must have passed by him that day. Then Peter’s majestic yet gentle command: “Look on us.” (Acts 3:4.) The record states that the lame man gave heed unto them, expecting to receive something from them.

The stirring words Peter then spoke have lifted the hearts of honest believers down through the stream of time, even to this day: “Silver and gold have I none; but such as I have give I thee: In the name of Jesus Christ of Nazareth rise up and walk.” Frequently we conclude the citation at this point and fail to note the next verses: “And he took him by the right hand, and lifted him up: … he … stood, and walked, and entered with them into the temple. …” (Acts 3:6–8.)

A helping hand had been extended. A broken body had been healed. A precious soul had been lifted toward God.

Time passes. Circumstances change. Conditions vary. Unaltered is the divine command to succor the weak and lift up the hands which hang down and strengthen the feeble knees. Each of us has the charge to be not a doubter, but a doer; not a leaner, but a lifter. But our complacency tree has many branches, and each spring more buds come into bloom. Often we live side by side but do not communicate heart to heart. There are those within the sphere of our own influence who, with outstretched hands, cry out: “Is there no balm in Gilead … ?” (Jer. 8:22.) Each of us must answer.

Edwin Markham observed:

\begin{verse}
“There is a destiny that makes us brothers;\\
None goes his way alone:\\
All that we send into the lives of others\\
Comes back into our own.”
\end{verse}

—“A Creed”

One who lived much of his life ignoring his fellowmen and living for self alone was Dickens’ immortal character, Ebenezer Scrooge. But there came that wintry night when the ghost of Jacob Marley appeared to Scrooge and lamented:

“Not to know that any Christian spirit working kindly in its little sphere, whatever it may be, will find its mortal life too short for its vast means of usefulness. Not to know that no space of regret can make amends for one life’s opportunities misused! Yet such was I! Oh! Such was I!

“Why did I walk through crowds of fellow-beings with my eyes turned down, and never raise them to that blessed Star which led the Wise Men to a poor abode! Were there no poor homes to which its light would conduct me!”

In an effort to comfort Marley, Scrooge proffered, “But you were always a good man of business, Jacob.”

Lamented Marley, “Business! … Mankind was my business!” (A Christmas Carol.)

The change that then occurred in the life of Scrooge was miraculous indeed. He became overnight the most generous, the most lovable, the most kindhearted Christian soul. In his own words he described his condition: “I am not the man I was.” So it ever is when one inclines his heart to the example of the Christ.

“… he that loveth not his brother abideth in death,” wrote the apostle John 1900 years ago. (1 Jn. 3:14.)

Some point the accusing finger at the sinner or the unfortunate and in derision say, “He has brought his condition upon himself.” Others exclaim, “Oh, he will never change. He has always been a bad one.” A few see beyond the outward appearance and recognize the true worth of a human soul. When they do, miracles occur. The downtrodden, the discouraged, the helpless become “no more strangers and foreigners, but fellowcitizens with the saints, and of the household of God.” (Eph. 2:19.) True love can alter human lives and change human nature.

This truth was stated so beautifully on the stage in My Fair Lady. Eliza Doolittle, the flower girl, spoke to one for whom she cared and who later was to lift her from such mediocre status: “You see, really and truly, apart from the things anyone can pick up (the dressing and the proper way of speaking, and so on), the difference between a lady and a flower girl is not how she behaves, but how she’s treated. I shall always be a flower girl to Professor Higgins, because he always treats me as a flower girl, and always will; but I know I can be a lady to you, because you always treat me as a lady, and always will.” (Adapted from Pygmalion, in The Complete Plays of Bernard Shaw, p. 260.)

Eliza Doolittle was but expressing the profound truth: When we treat people merely as they are, they will remain as they are. When we treat them as if they were what they should be, they will become what they should be. (Adapted from a quotation by Johann Wolfgang Von Goethe.)

In reality, it was the Redeemer who best taught this principle. Jesus changed men. He changed their habits and opinions and ambitions. He changed their tempers, dispositions, and natures. He changed their hearts. He lifted! He loved! He forgave! He redeemed! Do we have the will to follow?

Prison warden Kenyon J. Scudder has related this experience: A friend of his happened to be sitting in a railroad coach next to a young man who was obviously depressed. Finally the man revealed that he was a paroled convict returning from a distant prison. His imprisonment had brought shame to his family, and they had neither visited him nor written often. He hoped, however, that this was only because they were too poor to travel and too uneducated to write. He hoped, despite the evidence, that they had forgiven him.

To make it easy for them, however, he had written them to put up a signal for him when the train passed their little farm on the outskirts of town. If his family had forgiven him, they were to put a white ribbon in the big apple tree which stood near the tracks. If they didn’t want him to return, they were to do nothing, and he would remain on the train as it traveled west.

As the train neared his home town, the suspense became so great he couldn’t bear to look out of his window. He exclaimed, “In just five minutes the engineer will sound the whistle, indicating our approach to the long bend which opens into the valley I know as home. Will you watch for the apple tree at the side of the track?” His companion changed places with him and said he would. The minutes seemed like hours, but then there came the shrill sound of the train whistle. The young man asked, “Can you see the tree? Is there a white ribbon?”

Came the reply: “I see the tree. I see not one white ribbon, but many. There must be a white ribbon on every branch. Son, someone surely does love you.”

In that instant he stood cleansed by Christ.

His friend said, “I felt as if I had witnessed a miracle.”

Indeed, he had witnessed a miracle appropriately described by the third verse of a favorite Christmas carol, “O Little Town of Bethlehem”:

\begin{verse}
“How silently, how silently, The wondrous gift is given!\\
So God imparts to human hearts The blessings of his heaven.\\
“No ear may hear his coming; But in this world of sin,\\
Where meek souls will receive him still, The dear Christ enters in.”
\end{verse}

—Hymns, no. 165

We, too, can experience this same miracle when we, with hand and heart, as did the Savior, lift and love our neighbor to a newness of life.

May we succor the weak, lift up the hands which hang down, and strengthen the feeble knees, thereby inheriting that eternal life promised by the Redeemer, I pray, in the name of Jesus Christ. Amen.

\newpage
\settalk{“Finishers Wanted”}
\setspeaker{Thomas S. Monson}
\settalkdate{April 1972}
\talkheading{“Finishers Wanted”}{Thomas S. Monson}

On sunlit days during the noon hour, the streets of Salt Lake City abound with men and women who for a moment leave the confines of the tall office buildings and engage in that universal delight called window shopping. On occasion I, too, am a participant.

One Wednesday I paused before the elegant show window of a prestigious furniture store. That which caught and held my attention was not the beautifully designed sofa nor the comfortable-appearing chair that stood at its side. Neither was it the beautiful chandelier positioned overhead. Rather, my eyes rested upon a small sign that had been placed at the bottom right-hand corner of the window. Its message was brief: “Finishers Wanted.”

The store had need of those persons who possessed the talent and the skill to make ready for final sale the expensive furniture that the firm manufactured and sold. “Finishers Wanted.” The words remained with me as I returned to the pressing activities of the day.

In life, as in business, there has always been a need for those persons who could be called finishers. Their ranks are few, their opportunities many, their contributions great.

From the very beginning to the present time, a fundamental question remains to be answered by each who runs the race of life. Shall I falter or shall I finish? On the answer await the blessings of joy and happiness here in mortality and eternal life in the world to come.

We are not left without guidance to make this momentous decision. The Holy Bible contains those accounts, even those lessons, which, if carefully learned, will serve us well and be as a beacon light to guide our thoughts and influence our actions. As we read, we sympathize with those who falter. We honor those who finish.

The apostle Paul likened life to a great race when he declared: “Know ye not that they which run in a race run all, but one receiveth the prize? So run, that ye may obtain.” (1 Cor. 9:24.)

And before the words of Paul fell upon the ears of his listeners, the counsel of the preacher, even the son of David, king in Jerusalem, cautioned: “… the race is not to the swift, nor the battle to the strong. …” (Eccl. 9:11.)

Could the son of David have been referring to his own father? Judged by any standards, the greatest king Israel ever had was David. Anointed by Samuel, he was honored by the Lord.

In the first flush of his incredible triumphs, David rode the crest of popularity. As he achieved fresh victories, the women greeted him with a new song: “… Saul hath slain his thousands, and David his ten thousands.” (1 Sam. 18:7.) In adoration the people exclaimed: “Behold, we are thy bone and thy flesh.” (2 Sam. 5:1.)

Power he won. Peace he lost.

It happened one evening when David was walking upon the roof of the king’s house that he saw from the roof a woman bathing, and the woman was very beautiful.

“And David sent and inquired after the woman. And one said, Is not this Bathsheba, … the wife of Uriah the Hittite?” So “David sent messengers, and took her. …” (2 Sam. 11:3–4.)

The gross sin of adultery was followed by yet another: “… Set ye Uriah in the forefront of the hottest battle, and retire ye from him, that he may be smitten, and die.” (2 Sam. 11:15.) Lust and power had triumphed.

David’s rebuke came from the Lord God of Israel: “… thou hast killed Uriah the Hittite with the sword, and hast taken his wife to be thy wife. … Now therefore the sword shall never depart from thine house. …” (2 Sam. 12:9–10.)

David commenced well the race, then faltered and failed to finish his course.

Lest we lull ourselves into thinking that only the gross sins of life cause us to falter, consider the experience of the rich young man who came running to the Savior and asked the question: “Good Master, what good thing shall I do, that I may have eternal life?”

Jesus answered him: “If thou wilt enter into life, keep the commandments.

“He saith unto him, Which?”

To Jesus’ enumeration of the commandments, “The young man saith … All these things have I kept from my youth up: what lack I yet?

“Jesus said unto him, If thou wilt be perfect, go and sell that thou hast, and give to the poor, … and come and follow me.

“But when the young man heard that saying, he went away sorrowful: for he had great possessions.” (Matt. 19:16–18, 20–22.)

He preferred the comforts of earth to the treasures of heaven. He would not purchase the things of eternity by abandoning those of time. He faltered. He failed to finish.

So it was with Judas Iscariot. He commenced his ministry as an apostle of the Lord. He ended it a traitor. For thirty paltry pieces of silver, he sold his soul. At last, realizing the enormity of his sin, Judas, to his patrons and tempters of his crime, shrieked: “I have sinned in that I have betrayed the innocent blood.” (Matt. 27:4.)

Remorse had led to despair, despair to madness, and madness to suicide. He had succeeded in betraying the Christ. He had failed to finish the apostolic ministry to which he had been divinely called.

Lust for power, greed of gold, and disdain for honor have ever appeared as faces of failure in the panorama of life. Captivated by their artificial attraction, many noble souls have stumbled and fallen, thus losing the crown of victory reserved for the finisher of life’s great race.

Concerning those who fall short, John Greenleaf Whittier’s words seem particularly fitting:

\begin{verse}
“For of all sad words of tongue or pen,\\
The saddest are these: ‘It might have been!’”
\end{verse}

—“Maude Muller”

May we turn from the lives of those who faltered and consider for a moment some who finished and won the prize.

There was a man in the land of Uz whose name was Job, and that man was perfect and upright and one that feared God and eschewed evil. Pious in his conduct, prosperous in his fortune, Job was to face a test that would tempt any man.

Shorn of his possessions, scorned by his friends, afflicted by his suffering, even tempted by his wife, Job was to declare from the depths of his noble soul: “… behold, my witness is in heaven, and my record is on high.” (Job 16:19.) “… I know that my redeemer liveth. …” (Job 19:25.)

Job did not falter. Job became a finisher.

Following the earthly ministry of the Lord, there were many who, rather than deny testimony of him, would forfeit their lives. Such was Paul the apostle. The impulse of his father to send him to Jerusalem opened the door to Paul’s destiny. He would pass through it and help to shape a new world.

Gifted in his capacity to stir, move, and manage groups of men, Paul was a peerless example of one who nobly made the transition from sinner to saint. Though disappointment, heartache, and trial were to beset him, yet Paul, at the conclusion of his ministry, could say: “I have fought a good fight, I have finished my course, I have kept the faith.” (2 Tim. 4:7.) Like Job, Paul was a finisher.

He admonished us to “lay aside … sin” and to “run with patience the race … , Looking [for an example] unto Jesus the author and finisher of our faith. …” (Heb. 12:1–2.)

Though Jesus was tempted by the evil one, yet he resisted. Though he was hated, yet he loved. Though he was betrayed, yet he triumphed. Not in a cloud of glory or chariot of fire was Jesus to depart mortality, but with arms outstretched in agony upon a cruel cross. The magnitude of his mission is depicted in the simplicity of his words.

To his Father he prayed, “… the hour is come. … I have glorified thee on the earth: I have finished the work which thou gavest me to do.” (John 17:1, 4.) “… into thy hands I commend my spirit. …” (Luke 23:46.)

Mortality ended. Immortality began.

Times change, circumstances vary, but the true marks of a finisher remain. Note them well, for they are vital to our success.

These, the marks of a true finisher, will be as a lamp to our feet in the journey through life. Ever beckoning us onward and lifting us upward is he who pleaded, “… come, follow me.” (Luke 18:22.)

Frequently his help comes silently—on occasion with dramatic impact. Such was my experience of some years ago when, as a mission president, I was afforded the privilege of guiding the activities of precious young men and women, even missionaries whom he had called.

Some had problems, others required motivation; but one came to me in utter despair. He had made his decision to leave the mission field when but at the halfway mark. His bags were packed, his return ticket purchased. He came by to bid me farewell. We talked; we listened; we prayed. There remained hidden the actual reason for his decision to quit.

As we arose from our knees in the quiet of my office, the missionary began to weep. Flexing the muscle in his strong right arm, he blurted out, “This is my problem. All through school my muscle power qualified me for honors in football and track, but my mental power was neglected. President Monson, I’m ashamed of my school record. It reveals that ‘with effort’ I have the capacity to read at but the level of the fourth grade. I can’t even read the Book of Mormon. How then can I understand its contents and teach others its truths?”

The silence of the room was broken by my young nine-year-old son, who, without knocking, opened the door and, with surprise, apologetically said, “Excuse me. I just wanted to put this book back on the shelf.”

He handed me the book. Its title: A Child’s Story of the Book of Mormon, by Dr. Deta P. Neeley. I turned to the foreword and read these words: “This book has been written with a scientifically controlled vocabulary to the level of the fourth grade.” A sincere prayer from an honest heart had been dramatically answered.

My missionary accepted the challenge to read the book. Half laughing, half crying, he declared, “It will be good to read something I can understand.” Clouds of despair were dispelled by the sunshine of hope. He completed an honorable mission. He became a finisher.

Today I think I shall once more walk by that furniture store in our city and again gaze at the small sign in the large show window, that I may indelibly impress upon my mind the true meaning of its words: “Finishers Wanted.”

I pray humbly that each one of us may be a finisher in the race of life and thus qualify for that precious prize: eternal life with our Heavenly Father in the celestial kingdom. I testify that God lives, that this is his work, and ask that each may follow the example of his Son, a true finisher, in the name of Jesus Christ. Amen.

\newpage
\settalk{Hands}
\setspeaker{Thomas S. Monson}
\settalkdate{October 1972}
\talkheading{Hands}{Thomas S. Monson}

When Jesus of Nazareth taught and ministered among men, he spoke not as did the scribes and scholars of the day but rather in language understood by all. Jesus taught through parables. His teachings moved men and motivated them to a newness of life. The shepherd on the hillside, the sower in the field, the fisherman at his net all became subjects whereby the Master taught eternal truths.

The divinely created human body, with its truly marvelous powers and intricate parts, acquired new meaning when the Lord spoke of eyes that were not blinded but did really see, ears that were not stopped but did truly hear, and hearts that were not hardened but did know and feel. In his teachings he referred to the foot, the nose, the face, the side, the back. Significant are those occasions when he spoke of yet another part—even the human hand. Considered by artists and sculptors the most difficult member of the human body to capture on canvas or form in clay, the hand is a wonder to behold. Neither color, size, shape, nor age distorts this miracle of creation.

First, let us consider the hand of a child. Who among us has not praised God and marveled at his powers when an infant is held in one’s arms. That tiny hand, so small yet so perfect, instantly becomes the topic of conversation. No one can resist placing his little finger in the clutching hand of an infant. A smile comes to the lips, a certain glow to the eyes, and one appreciates the tender feelings which prompted the poet to pen the lines: “A sweet new blossom of humanity, fresh fallen from God’s own home, to flower on earth.” (Gerald Massey.)

As the child grows, the tightly clutched hand opens in an expression of perfect trust. “Take me by the hand, Mother; then I won’t be afraid” bespeaks this confidence. The delightful song the little children sing so beautifully at once becomes a plea for patience, an invitation to teach—even an opportunity to serve:

\begin{verse}
“I have two little hands folded snugly and tight,\\
They are tiny and weak yet they know what is right,\\
During all the long hours till daylight is through,\\
There is plenty indeed for my two hands to do.
\end{verse}

—The Children Sing, no. 97

The sentiments such love and faith arouse should ever draw forth from each parent a pledge of fidelity—even a determination to do that which is right.

Should added emphasis be required, we need but refer to that account where the disciples came unto Jesus, saying, “Who is the greatest in the kingdom of heaven?

“And Jesus called a little child unto him, and set him in the midst of them,

“And said, Verily I say unto you, Except ye be converted, and become as little children, ye shall not enter into the kingdom of heaven.

“And whoso shall receive one such little child in my name receiveth me.

“But whoso shall offend one of these little ones which believe in me, it were better for him that a millstone were hanged about his neck, and that he were drowned in the depth of the sea.” (Matt. 18:1–3, 5–6.) Significant is the hand of a child.

Second, may we turn our attention to the hand of youth. This is the training period when busy hands learn to labor—and labor to learn. Honest effort and loving service become identifying features of the abundant life. Each was effectively taught the girls in the Mutual class when cookies were baked and taken by them to elderly women residing in a neighborhood rest home. The aged hand of a lonely grandmother clasped that of the thoughtful teenager. No word was spoken. Heart spoke to heart. The hand that baked the cookies was raised to wipe a tear. Such hands are clean hands. Such hearts are pure hearts.

Then comes that day when the hand of a boy takes the hand of a girl, and parents suddenly realize their children have grown. Never is the hand of a girl so delicately displayed as when there glistens on her finger a ring denoting a sacred pledge. Her step becomes quicker, her countenance brighter, and all is well with the world. Courtship has come. Marriage follows. And once again two hands are clasped, this time in a holy temple. Cares of the world are for a brief moment forgotten. Thoughts turn to eternal values. The clasped hands speak of promised hearts. Heaven is here.

Time passes. The hand of a bride becomes the hand of a mother. Ever so gently she cares for her precious child. Bathing, dressing, feeding, comforting—there is no hand like mother’s. Nor does its tender care diminish through the years. Ever shall I remember the hand of one mother—the mother of a missionary. Some years ago at a worldwide seminar for mission presidents, the parents of missionaries were invited to meet and visit briefly with each mission president. Forgotten are the names of each who extended a greeting and exchanged a friendly handshake. Remembered are the feelings which welled up within me as I took in my hand the calloused hand of one mother from Star Valley, Wyoming. “Please excuse the roughness of my hand,” she apologized. “Since my husband has been ill, the work of the farm has been mine to do, that our boy may, as a missionary, serve the Lord.”

Tears could not be restrained, nor should they have been. Such tears produce a certain cleansing of the soul. A mother’s labor sanctified a son’s service. Loved are the hands of a mother.

Not to be overlooked is the hand of a father. Whether he be a skilled surgeon, a master craftsman, or a talented teacher, his hands support his family. There is a definite dignity in honest labor and tireless toil.

During the period of the great depression I was a small boy. Fortunate were those men who had work. Jobs were few, hours long, pay scant. On our street was a father who, though old in years, supported with the labor of his hands his rather large family of girls. His firm was known as the Spring Canyon Coal Company. It consisted of one old truck, a pile of coal, one shovel, one man, and his own two hands. From early morning to late evening he struggled to survive. Yet during the monthly fast and testimony meeting, I specifically remember him expressing his thanks to the Lord for his family, for his work, and for his testimony. The fingers of those rough, red, chapped hands turned white as they gripped the back of the bench on which I sat as Brother James Farrell bore witness of a boy, even Joseph Smith, who, in a grove of trees near Palmyra, New York, knelt in prayer and beheld the heavenly vision of God the Father and Jesus Christ the Son. The memory of those hands of a father serve to remind me of his abiding faith, his honest conviction, and his testimony of truth. Honored are the hands of a father.

On Friday morning in this historic tabernacle, and in the homes of Church members viewing or listening to the conference session, hands were raised to sustain a prophet, a seer, and a revelator—even the President of The Church of Jesus Christ of Latter-day Saints. Our upraised hands were an outward expression of our inward feelings. As we raised our hands, we pledged our hearts. Could I for a moment mention the hands of that prophet, even President Harold B. Lee? I do so humbly and with his permission.

Some years ago, President Lee, directed by inspiration and revelation, called Dewitt J. Paul to serve as patriarch in one of the eastern stakes of the Church. The call humbled beyond words both Brother and Sister Paul. They wondered. They worried. They prayed for assurance and heavenly confirmation. Such did not come suddenly.

The vote of the people demonstrated their supporting approval. Then came the time for ordination. In a basement room situated two floors beneath the meeting hall in which the conference was held, Dewitt Paul nervously sat on a chair and, with a silent prayer in his heart, awaited his ordination. President Harold B. Lee then placed his hands upon the head of the newly called patriarch and began to speak. Peace replaced turmoil. Faith overcame doubt. Seated next to Sister Paul was a lifelong friend to whom Sister Paul had confided her concern. During the pronouncement of the blessing and ordination, she opened her eyes. As she did so she saw a ray of light shining upon President Lee as he placed his hands upon the head of Brother Paul. At the conclusion of the blessing, she hastened to tell Brother Lee of this confirmation of a call. She recounted how she saw the sunshine form the ray of light and how it brought a bright glow to the hands of President Lee. “Indeed, this is to you a confirmation of a sacred call,” said President Lee, “for as you look about this basement room, there is no window through which the sun might beam its rays.” Precious are the hands of a prophet.

Finally, may we speak of yet another hand—even the hand of the Lord. This was the hand which guided Moses, which strengthened Joshua—the hand promised to Jacob when the Lord declared: “Fear thou not; for I am with thee: be not dismayed; for I am thy God; … I will uphold thee with the right hand of my righteousness.” (Isa. 41:10.)

This was the determined hand which drove from the temple the money changers. This was the loving hand that blessed little children. This was the strong hand that opened deaf ears and restored vision to sightless eyes. By this hand was the leper cleansed, the lame man healed—even the dead Lazarus raised to life. With the finger of this hand there was written in the sand that message which the winds did erase but which honest hearts did retain. The hand of the carpenter. The hand of the teacher. The hand of the Christ. One called Pontius Pilate washed his hands of this man called King of the Jews. Oh foolish, spineless Pilate! Did you really believe that water could cleanse such guilt?

\begin{verse}
“I think of his hands pierced and bleeding to pay the debt!\\
Such mercy, such love, and devotion can I forget? …\\
Oh, it is wonderful that he should care for me enough to die for me!\\
Oh, it is wonderful, wonderful to me!”
\end{verse}

—Hymns, no. 80

Pitied is the hand that sins. Envied is the hand that paints. Honored is the hand that builds. Appreciated is the hand that helps. Respected is the hand that serves. Adored is the hand that saves—even the hand of Jesus Christ, the Son of God, the Redeemer of all mankind. With that hand he knocks upon the door of our understanding.

“Behold, I stand at the door, and knock: if any man hear my voice and open the door, I will come in to him. …” (Rev. 3:20.)

Shall we listen for his voice? Shall we open the doorway of our lives to his exalted presence? Each must answer for himself.

In this journey called mortality, clouds of gloom may appear on the horizon of our personal destiny. The way ahead may be uncertain, foreboding. In desperation we may be prompted to ask, as did another:

\begin{verse}
“… I said to the man who stood at the gate of the year:\\
‘Give me a light, that I may tread safely into the unknown.’\\
And he replied:\\
‘Go out into the darkness and put your hand into the hand of God.\\
That shall be to you better than a light and safer than a known way.’”
\end{verse}

Of this solemn truth I testify. I declare that our Lord and Savior does live and that he even today directs his church with his all powerful hand, in the name of Jesus Christ. Amen.

\newpage
\settalk{Yellow Canaries with Gray on Their Wings}
\setspeaker{Thomas S. Monson}
\settalkdate{April 1973}
\talkheading{Yellow Canaries with Gray on Their Wings}{Thomas S. Monson}

Some 23 years ago I was called as a young man to serve as the bishop of a large ward in Salt Lake City. The magnitude of the calling was overwhelming and the responsibility frightening. My inadequacy humbled me. But my Heavenly Father did not leave me to wander in darkness and in silence, uninstructed or uninspired. In his own way he revealed the lessons he would have me learn.

One evening at a late hour my telephone rang. I heard a voice say, “Bishop Monson, this is the hospital calling. Kathleen McKee, a member of your congregation, has just passed away. Our records reveal that she had no next of kin, but your name is listed as the one to be notified in the event of her death. Could you come to the hospital right away?”

Upon arriving there, I was presented with a sealed envelope which contained a key to the modest apartment in which Kathleen McKee had lived. A childless widow 73 years of age, she had enjoyed but few of life’s luxuries and possessed scarcely sufficient of its necessities. In the twilight of her life she had become a member of The Church of Jesus Christ of Latter-day Saints. Being a quiet and overly reserved person, little was known about her life.

That same night I entered her tidy basement apartment, turned the light switch, and in a moment discovered a letter written ever so meticulously in Kathleen McKee’s own hand. It rested face up on a small table and read:

“Bishop Monson,

“I think I shall not return from the hospital. In the dresser drawer is a small insurance policy which will cover funeral expenses. The furniture may be given to my neighbors.

“In the kitchen are my three precious canaries. Two of them are beautiful, yellow-gold in color, and are perfectly marked. On their cages I have noted the names of friends to whom they are to be given. In the third cage is ‘Billie.’ He is my favorite. Billie looks a bit scrubby, and his yellow hue is marred by gray on his wings. Will you and your family make a home for him? He isn’t the prettiest, but his song is the best.”

In the days that followed, I learned much more about Kathleen McKee. She had befriended many neighbors in need. She had given cheer and comfort almost daily to a cripple who lived down the street. Indeed, she had brightened each life she touched. Kathleen McKee was much like “Billie,” her prized yellow canary with gray on its wings. She was not blessed with beauty, gifted with poise, nor honored by posterity. Yet her song helped others to more willingly bear their burdens and more ably shoulder their tasks. She lived the message of the verse:

\begin{verse}
“Go, gladden the lonely, the dreary;\\
Go, comfort the weeping, the weary;\\
Go, scatter kind deeds on your way;\\
Oh, make the world brighter today!”
\end{verse}

—Deseret Sunday School Songs, 1909, No. 197

The world is filled with yellow canaries with gray on their wings. The pity is that so precious few of them have learned to sing. Perhaps the clear notes of proper example have not sounded in their ears or found lodgment in their hearts.

Some are young people who don’t know who they are, what they can be or even want to be. They are afraid, but they don’t know of what. They are angry, but they don’t know at whom. They are rejected, and they don’t know why. All they want is to be somebody.

Others are stooped with age, burdened with care, or filled with doubt—living lives far below the level of their capacities.

All of us are prone to excuse our own mediocre performance. We blame our misfortunes, our disfigurements, or our so-called handicaps. Victims of our own rationalization, we say silently to ourselves: “I’m just too weak,” or “I’m not cut out for better things.” Others soar beyond our meager accomplishments. Envy and discouragement then take their toll.

Can we not appreciate that our very business in life is not to get ahead of others, but to get ahead of ourselves? To break our own records, to outstrip our yesterdays by our todays, to bear our trials more beautifully than we ever dreamed we could, to give as we have never given, to do our work with more force and a finer finish than ever—this is the true idea: to get ahead of ourselves.

To live greatly, we must develop the capacity to face trouble with courage, disappointment with cheerfulness, and triumph with humility. You ask, “How might we achieve these goals?” I answer, “By getting a true perspective of who we really are!” We are sons and daughters of a living God, in whose image we have been created. Think of that truth: “Created in the image of God.” We cannot sincerely hold this conviction without experiencing a profound new sense of strength and power, even the strength to live the commandments of God, the power to resist the temptations of Satan.

True, we live in a world where moral character ofttimes is relegated to a position secondary to facial beauty or personal charm. We read and hear of local, national, and international beauty contests. Throngs pay tribute to Miss America, Miss World, and Miss Universe. Athletic prowess, too, has its following. The winter games, the world Olympics, the tournaments of international scope bring forth the adoring applause of the enthralled crowd. Such are the ways of men!

But what are the inspired words of God? From a time of long ago the counsel of Samuel the prophet echoes in our ears: “… the Lord seeth not as man seeth; for man looketh on the outward appearance, but the Lord looketh on the heart.” (1 Sam. 16:7.)

Sham and hypocrisy found no place with the King of kings and the Lord of lords. He denounced the scribes and Pharisees for their vanity and shallow lives, their pretense and feigned righteousness. He called them “whited sepulchres, which indeed appear beautiful outward, but are within full of dead men’s bones.” (Matt. 23:27.)

They, like the beautiful yellow canaries, were outwardly handsome, but a true song came not from their hearts.

To their counterparts on this continent God’s prophet declared: “For behold, ye do love money, and your substance, and your fine apparel, and the adorning of your churches, more than ye love the poor and the needy, the sick and the afflicted. …

“Why are ye ashamed to take upon you the name of Christ? …

“Why do ye adorn yourselves with that which hath no life, and yet suffer the hungry, and the needy, and the naked, and the sick and the afflicted to pass by you, and notice them not?” (Morm. 8:37–39.)

The Master could be found mingling with the poor, the downtrodden, the oppressed, and the afflicted. He brought hope to the hopeless, strength to the weak, and freedom to the captive. He taught of the better life to come—even eternal life. This knowledge ever directs those who receive the divine injunction: “Follow thou me.” It guided Peter. It motivated Paul. It can determine our personal destiny. Can we make the decision to follow in righteousness and truth the Redeemer of the world? With his help a rebellious boy can become an obedient man, a wayward girl can cast aside the old self and begin anew. Indeed, the gospel of Jesus Christ can change men’s lives.

In his epistle to the Corinthians, the apostle Paul taught: “… God hath chosen the weak things of the world to confound the things which are mighty.” (1 Cor. 1:27.)

When the Savior sought a man of faith, he did not select him from the throng of self-righteous who were found regularly in the synagogue. Rather, he called him from among the fishermen of Capernaum.

While teaching on the seashore, he saw two ships standing by the lake. He entered one and asked its owner to put it out a little from the land so he might not be pressed upon by the crowd. After teaching further, he said to Simon, “Launch out into the deep, and let down your nets.”

Simon answered: “Master, we have toiled all the night, and have taken nothing: nevertheless at thy word I will let down the net. And when they had this done, they inclosed a great multitude of fishes. … When Simon Peter saw it, he fell down at Jesus’ knees, saying, Depart from me; for I am a sinful man, O Lord.” (Luke 5:4–6, 8.)

Came the reply: “Follow me, and I will make you fishers of men.” (Matt. 4:19.) Simon the fisherman had received his call. Doubting, disbelieving, unschooled, untrained, impetuous Simon did not find the way of the Lord a highway of ease nor a path free from pain. He was to hear the rebuke: “O thou of little faith” (Matt. 14:31), and likewise the denunciation, “Get thee behind me, Satan: thou art an offence unto me” (Matt. 16:23). Yet, when the Master asked him, “… whom say ye that I am?” Peter answered: “Thou art the Christ, the Son of the living God.” (Matt. 16:15–16.)

Simon, man of doubt, had become Peter, apostle of faith. A yellow canary with gray on his wings qualified for the Master’s full confidence and abiding love.

When the Savior was to choose a missionary of zeal and power, he found him not among his advocates but amidst his adversaries. Saul of Tarsus made havoc of the church and breathed out threatenings and slaughter against the disciples of the Lord. But this was before the experience of Damascus Way. Of Saul, the Lord declared: “… he is a chosen vessel unto me, to bear my name before the Gentiles, and kings, and the children of Israel: … I will shew him how great things he must suffer for my name’s sake.” (Acts 9:15–16.)

Saul the persecutor became Paul the proselyter. Like the yellow canary with gray on his wings, Paul, too, had his blemishes. He himself said: “And lest I should be exalted above measure through the abundance of the revelations, there was given to me a thorn in the flesh, the messenger of Satan to buffet me. … For this thing I besought the Lord thrice, that it might depart from me. And he said unto me, My grace is sufficient for thee: for my strength is made perfect in weakness. …” (2 Cor. 12:7–9.)

Both Paul and Peter were to expend their strength and forfeit their lives in the cause of truth. The Redeemer chose imperfect men to teach the way to perfection. He did so then. He does so now—even yellow canaries with gray on their wings.

He calls you and me to serve him here below and sets us to the tasks he would have us fulfill. The commitment is total. There is no conflict of conscience. And in our struggle, should we stumble, then let us plead: “Lead us, oh lead us, great Molder of men; out of the darkness to strive once again.” (From the “Fight Song,” Yonkers High School.)

Our appointed task may appear insignificant, unnecessary, unnoticed. We may be tempted to question:

\begin{verse}
“‘Father, where shall I work today?’\\
And my love flowed warm and free.\\
Then he pointed out a tiny spot\\
And said, ‘Tend that for me.’\\
I answered quickly, ‘Oh no, not that!\\
Why, no one would ever see,\\
No matter how well my work was done.\\
Not that little place for me.’\\
And the word he spoke, it was not stern;\\
He answered me tenderly:\\
‘Ah, little one, search that heart of thine;\\
Art thou working for them or for me?\\
‘Nazareth was a little place,\\
And so was Galilee.’”
\end{verse}

—Meade MacGuire

My prayer today is that we indeed will follow the Man of Galilee. May we praise his name, and so order our lives as to reflect our love. May we ever remember that to us God our Father gave his Son, and that for us Jesus Christ gave his life. I testify that he lives and pray we may be worthy of such a divine gift, in the name of Jesus Christ the Lord. Amen.

\newpage
\settalk{“Behold Thy Mother”}
\setspeaker{Thomas S. Monson}
\settalkdate{October 1973}
\talkheading{“Behold Thy Mother”}{Thomas S. Monson}

One summer day I stood alone in the quiet of the American War Memorial Cemetery of the Philippines. A spirit of reverence filled the warm tropical air. Situated among the carefully mowed grass, acre upon acre, were markers identifying men, mostly young, who in battle gave their lives. As I let my eyes pass name by name along the many colonnades of honor, tears came easily and without embarrassment. As my eyes filled with tears, my heart swelled with pride. I contemplated the high price of liberty and the costly sacrifice many had been called upon to bear.

My thoughts turned from those who bravely served and gallantly died. There came to mind the grief-stricken mother of each fallen man as she held in her hand the news of her precious son’s supreme sacrifice. Who can measure a mother’s grief? Who can probe a mother’s love? Who can comprehend in its entirety the lofty role of a mother? With perfect trust in God, she walks, her hand in his, into the valley of the shadow of death that you and I might come forth unto life.

\begin{verse}
“The noblest thoughts my soul can claim,\\
The holiest words my tongue can frame,\\
Unworthy are to frame the name\\
More sacred than all other.\\
An infant when her love first came,\\
A man, I find it just the same:\\
Reverently I breathe her name—\\
The blessed name of mother.”
\end{verse}

—George Griffith Fetter

In this spirit, let us consider mother. Four mothers come to mind: first, mother forgotten; second, mother remembered; third, mother blessed; and finally, mother loved.

“Mother forgotten” is observed all too frequently. The nursing homes are crowded, the hospital beds are full, the days come and go—often the weeks and months pass—but mother is not visited. Can we not appreciate the pangs of loneliness, the yearnings of mother’s heart when hour after hour, alone in her age, she gazes out the window for the loved one who does not visit, the letter the postman does not bring. She listens for the knock that does not sound, the telephone that does not ring, the voice she does not hear. How does such a mother feel when her neighbor welcomes gladly the smile of a son, the hug of a daughter, the glad exclamation of a child, “Hello, Grandmother.”

There are yet other ways we forget mother. Whenever we fall, whenever we do less than we ought, in a very real way we forget mother.

Last Christmas I talked to the proprietress of a Salt Lake City nursing home. From the hallway where we stood, she pointed to several elderly women assembled in a peaceful living room. She observed, “There’s Mrs. Hansen. Her daughter visits her every week, right at 3:00 p.m. on Sunday. To her right is Mrs. Peek. Each Wednesday there is a letter in her hands from her son in New York. It is read, then reread, then saved as a precious piece of treasure. But see Mrs. Carroll; her family never telephones, never writes, never visits. Patiently she justifies this neglect with words which are heard but do not convince or excuse, ‘They are all so busy.’” Shame on all who thus make of a noble woman “mother forgotten.”

“Hearken unto thy father that begat thee,” wrote Solomon, “and despise not thy mother when she is old.” (Prov. 23:22.) Can we not make of a mother forgotten a “mother remembered”?

Men turn from evil and yield to their better natures when mother is remembered. A famed officer from the Civil War period, Colonel Higgenson, when asked to name the incident of the Civil War that he considered the most remarkable for bravery, said that there was in his regiment a man whom everybody liked, a man who was brave and noble, who was pure in his daily life, absolutely free from dissipations in which most of the other men indulged.

One night at a champagne supper, when many were becoming intoxicated, someone in jest called for a toast from this young man. Colonel Higgenson said that he arose, pale but with perfect self-control, and declared: “Gentlemen, I will give you a toast which you may drink as you will, but which I will drink in water. The toast that I have to give is, ‘Our mothers.’”

Instantly a strange spell seemed to come over all the tipsy men. They drank the toast in silence. There was no laughter, no more song, and one by one they left the room. The lamp of memory had begun to burn, and the name of “Mother” touched every man’s heart.

As a boy, I well remember Sunday School on Mother’s Day. We would hand to each mother present a small potted plant and sit in silent reverie as Melvin Watson, a blind member, would stand by the piano and sing, “That Wonderful Mother of Mine.” This was the first time I saw a blind man cry. Even today, in memory, I can see the moist tears move from those sightless eyes, then form tiny rivulets and course down his cheeks, falling finally upon the lapel of the suit he had never seen. In boyhood puzzlement I wondered why all of the grown men were silent, why so many handkerchiefs came forth. Now I know. You see, mother was remembered. Each boy, every girl, all fathers and husbands seemed to make a silent pledge: “I will remember that wonderful mother of mine.”

Some years ago I listened intently as a man well beyond middle age told me of an experience in his family history. The widowed mother who had given birth to him and his brothers and sisters had gone to her eternal and well-earned reward. The family assembled at the home and surrounded the large dining room table. The small metal box in which Mother had kept her earthly treasures was opened reverently. One by one each keepsake was brought forth. There was the wedding certificate from the Salt Lake Temple. “Oh, now Mother could be with Dad.” Then there was the deed to the humble home where each child had in turn entered upon the stage of life. The appraised value of the house had little resemblance to the worth Mother had attached to it.

Then there was discovered a yellowed envelope which bore the marks of time. Carefully the flap was opened and from inside was taken a homemade valentine. Its simple message, in the handwriting of a child, read, “I love you, Mother.” Though she was gone, by what she held sacred, Mother taught yet another lesson. A silence permeated the room, and every member of the family made a pledge not only to remember, but also to honor mother. For them it was not too little and too late, as in the classic poem of Rose Marinoni entitled “At Sunrise”:

\begin{verse}
“They pushed him straight against the wall,\\
The firing squad dropped in a row;\\
And why he stood on tiptoe,\\
Those men shall never know.\\
He wore a smile across his face\\
As he stood primly there,\\
The guns straight aiming at his heart,\\
The sun upon his hair.\\
For he remembered in a flash\\
Those days beyond recall,\\
When his proud mother took his height\\
Against the bedroom wall.”
\end{verse}

Now that we have considered “mother remembered,” let us turn to “mother blessed.” For one of the most beautiful and reverent examples, I refer to the holy scriptures.

In the New Testament of our Lord, perhaps we have no more moving account of “mother blessed” than the tender regard of the Master for the grieving widow at Nain.

“And it came to pass … that he went into a city called Nain; and many of his disciples went with him, and much people.

“Now when he came nigh to the gate of the city, behold, there was a dead man carried out, the only son of his mother, and she was a widow: and much people of the city was with her.

“And when the Lord saw her, he had compassion on her, and said unto her, Weep not.

“And he came and touched the bier: and they that bare him stood still. And he said, Young man, I say unto thee, Arise.

“And he that was dead sat up, and began to speak. And he delivered him to his mother.” (Luke 7:11–15.)

What power, what tenderness, what compassion did our Master and exemplar thus demonstrate. We, too, can bless if we will but follow his noble example. Opportunities are everywhere. Needed are eyes to see the pitiable plight, ears to hear the silent pleadings of a broken heart. Yes, and a soul filled with compassion that we might communicate not only eye to eye or voice to ear, but in the majestic style of the Savior, even heart to heart. Then every mother everywhere will be “mother blessed.”

Finally, let us contemplate “mother loved.” Universally applicable is the poem recalled from childhood and enjoyed by children even today, “Which Loved Best?”

\begin{verse}
“‘I love you, Mother,’ said little John;\\
Then, forgetting his work, his cap went on,\\
And he was off to the garden swing,\\
Leaving his mother the wood to bring.
\end{verse}

—Joy Allison

One certain way each can demonstrate genuine love for mother is to live the truths mother so patiently taught. Such a lofty goal is not new to our present generation. On this continent, in times described in the Book of Mormon, we read of a brave, a good, and noble leader named Helaman who did march in righteous battle at the head of 2,000 young men. Helaman described the activities of these young men: “… never had I seen so great courage, … as … they said unto me: … behold our God is with us, and he will not suffer that we should fall; then let us go forth; … Now they never had fought, yet they did not fear death; … yea, they had been taught by their mothers, that if they did not doubt, God would deliver them. And they rehearsed unto me the words of their mothers, saying: We do not doubt our mothers knew it.” (Alma 56:45–48.)

At the end of the battle, Helaman continued his description: “… behold, to my great joy, there had not one soul of them fallen to the earth; yea, and they had fought as if with the strength of God; yea, never were men known to have fought with such miraculous strength; and with such mighty power. …” (Alma 56:56.)

Miraculous strength, mighty power—mother’s love and love for mother had met and triumphed.

The holy scriptures, the pages of history are replete with tender, moving, convincing accounts of “mother loved.” One, however, stands out supreme, above and beyond any other. The place is Jerusalem, the period known as the Meridian of Time. Assembled is a throng of Roman soldiers. Their helmets signify their loyalty to Caesar, their shields bear his emblem, their spears are crowned by Roman eagles. Assembled also are natives to the land of Jerusalem. Faded into the still night, and gone forever are the militant and rowdy cries, “Crucify him, crucify him.”

The hour has come. The personal earthly ministry of the Son of God moves swiftly to its dramatic conclusion. A certain loneliness is here. Nowhere to be found are the lame beggars who, because of this man, walk; the deaf who, because of this man, hear; the blind who, because of this man, see; the dead who, because of this man, live.

There remained yet a few faithful followers. From his tortured position on the cruel cross, he sees his mother and the disciple whom he loved standing by. He speaks: “… woman, behold thy son! Then saith he to the disciple, Behold thy mother! …” (John 19:26–27.)

From that awful night when time stood still, when the earth did quake and great mountains were brought down—yes, through the annals of history, over the centuries of years and beyond the span of time, there echoes his simple yet divine words, “Behold thy mother!”

As we truly listen to that gentle command and with gladness obey its intent, gone forever will be the vast legions of “mothers forgotten.” Everywhere present will be “mothers remembered,” “mothers blessed,” and “mothers loved” and, as in the beginning, God will once again survey the workmanship of his own hand and be led to say, “It [is] very good.”

May each of us treasure this truth; one cannot forget mother and remember God. One cannot remember mother and forget God. Why? Because these two sacred persons, God and mother, partners in creation, in love, in sacrifice, in service, are as one.

May we, by our thoughts and our actions, honor God and mother, I pray humbly yet earnestly, in the name of Jesus Christ. Amen.

\newpage
\settalk{The Paths Jesus Walked}
\setspeaker{Thomas S. Monson}
\settalkdate{April 1974}
\talkheading{The Paths Jesus Walked}{Thomas S. Monson}

My beloved brothers and sisters, my heart is full to overflowing. You and I, this memorable day, have been partakers of the Spirit of the Lord Jesus Christ. This is his church. It bears his name. His prophet has lifted each of us beyond the shackles of this earth to the lofty heavens above. Our raised hands are supported by our pledged hearts. The kingdom of God moves forward in its undeviating and eternal course.

On a chill day last December, we gathered into this historic Tabernacle to pay honor and tribute to a man whom we loved, honored, and followed—even President Harold B. Lee. Prophetic in his utterance, powerful in his leadership, devoted in his service, President Lee inspired in all of us a desire to achieve perfection. He counseled us, “Keep the commandments of God. Follow the pathway of the Lord.”

One day later, in a very sacred room on an upper floor of the Salt Lake Temple, his successor was chosen, sustained, and set apart to his sacred calling. Untiring in his labor, humble in his manner, inspiring in his testimony, President Spencer W. Kimball invited us to continue the course charted by President Lee. He spoke the same penetrating words, “Keep the commandments of God. Follow the pathway of the Lord. Walk in his footsteps.”

Later that same evening, I happened to glance at a travel brochure which had arrived at my home several days earlier. It was printed in breathtaking color and written with persuasive skill. The reader was invited to visit the fjords of Norway and the Alps of Switzerland, all in one package tour. Yet another offering beckoned the reader to Bethlehem—even the Holy Land—cradle of Christianity. The closing lines of the brochure’s message contained the simple yet powerful appeal, “Come and walk where Jesus walked.”

My thoughts returned to the counsel God’s prophets—even President Lee and President Kimball—had provided: “Follow the pathway of the Lord. Walk in his footsteps.” I reflected on the words penned by the poet:

\begin{verse}
I walked today where Jesus walked,\\
In days of long ago;\\
I wandered down each path He knew,\\
With rev’rent step and slow.
\end{verse}

—Daniel S. Twohig

We need not visit the Holy Land to feel him close to us. We need not walk by the shores of Galilee or among the Judean hills to walk where Jesus walked.

In a very real sense, all can walk where Jesus walked when, with his words on our lips, his spirit in our hearts, and his teachings in our lives, we journey through mortality.

I would hope that we would walk as he walked—with confidence in the future, with an abiding faith in his Father, and a genuine love for others.

Jesus walked the path of disappointment.

Can one appreciate his lament over the Holy City? “O Jerusalem, Jerusalem, which killest the prophets, and stonest them that are sent unto thee; how often would I have gathered thy children together, as a hen doth gather her brood under her wings, and ye would not!” (Luke 13:34.)

Jesus walked the path of temptation.

That evil one, amassing his greatest strength, his most inviting sophistry, tempted him who had fasted for forty days and forty nights and was an hungred. Came the taunt: “… If thou be the Son of God, command that these stones be made bread.” The reply: “Man shall not live by bread alone. …” Again, “… If thou be the Son of God, cast thyself down: for it is written, He shall give his angels charge concerning thee. …” The answer: “… Thou shalt not tempt the Lord thy God.” Still again: “… the kingdoms of the world, and the glory of them … will I give thee, if thou wilt fall down and worship me. …” The Master replied: “Get thee hence, Satan: for it is written, Thou shalt worship the Lord thy God, and him only shalt thou serve.” (Matt. 4:3–4, 6–10.)

Jesus walked the path of pain.

Consider the agony of Gethsemane: “… Father, if thou be willing, remove this cup from me: nevertheless not my will, but thine, be done. … And being in an agony he prayed more earnestly: and his sweat was as it were great drops of blood falling down to the ground.” (Luke 22:42, 44.)

And who among us can forget the cruelty of the cross. His words: “… I thirst. … It is finished. …” (John 19:28, 30.)

Yes, each of us will walk the path of disappointment, perhaps due to an opportunity lost, a power misused, or a loved one not taught. The path of temptation, too, will be the path of each. “And it must needs be that the devil should tempt the children of men, or they could not be agents unto themselves. …” (D\&C 29:39.)

Likewise shall we walk the path of pain. We cannot go to heaven in a feather bed. The Savior of the world entered after great pain and suffering. We, as servants, can expect no more than the Master. Before Easter there must be a cross.

While we walk these paths which bring forth bitter sorrow, we can also walk those paths which yield eternal joy.

We, with Jesus, can walk the path of obedience.

It will not be easy. “Though he were a Son, yet learned he obedience by the things which he suffered.” (Heb. 5:8.) Let our watchword be the heritage bequeathed us by Samuel: “… Behold, to obey is better than sacrifice, and to hearken than the fat of rams.” (1 Sam. 15:22.) Let us remember that the end result of disobedience is captivity and death, while the reward for obedience is liberty and eternal life.

We, like Jesus, can walk the path of service.

Like a glowing searchlight of goodness is the life of Jesus as he ministered among men. He brought strength to the limbs of the cripple, sight to the eyes of the blind, hearing to the ears of the deaf, and life to the body of the dead.

His parables preach power. With the good Samaritan he taught: “… love … thy neighbour. …” (Luke 10:27.) Through his kindness to the woman taken in adultery, he taught compassionate understanding. In his parable of the talents, he taught each of us to improve himself and to strive for perfection. Well could he have been preparing us for our journey along his pathway. Why else would he counsel: “… Go, and do thou likewise.” (Luke 10:37.)

Finally, he walked the path of prayer.

Three great lessons from three timeless prayers. First, from his ministry: “… When ye pray, say, Our Father which art in heaven, Hallowed be thy name. …” (Luke 11:2.)

Second, from Gethsemane: “… not my will, but thine, be done.” (Luke 22:42.)

Third, from the cross: “… Father, forgive them; for they know not what they do. …” (Luke 23:34.)

It is by walking the path of prayer that we commune with the Father and become partakers of his power.

Shall we have the faith, even the desire, to walk these pathways which Jesus walked? God’s prophet, seer, and revelator has this day invited us to do so. All we need do is follow him, for this is the pathway he walks.

My first acquaintance with this prophet leader was 24 years ago when I served as a young bishop here in Salt Lake City. One morning, upon answering my telephone, a voice said, “This is Elder Spencer W. Kimball. I have a favor to ask of you. In your ward, hidden away behind a large building on Fifth South Street, is a tiny trailer home. Living there is Margaret Bird, a Navajo widow. She feels unwanted, unneeded, and lost. Could you and the Relief Society presidency seek her out, extend to her the hand of fellowship, and provide for her a special welcome?” This we did.

A miracle resulted. Margaret Bird blossomed in her newly found environment. Despair disappeared. The widow in her affliction had been visited. The lost sheep had been found. Each who participated in the simple human drama emerged a better person.

In reality, the true shepherd was the concerned apostle who, leaving the ninety and nine of his ministry, went in search of the precious soul who was lost. Spencer W. Kimball had walked the pathway Jesus walked. He did so then. He does so now.

As you and I walk the pathway Jesus walked, let us listen for the sound of sandaled feet. Let us reach out for the Carpenter’s hand. Then we shall come to know him. He may come to us as one unknown, without a name, as of old, by the lakeside he came to those men who knew him not. He speaks to us the same words, “… follow thou me …” (John 21:22), and sets us to the task which he has to fulfill for our time. He commands, and to those who obey him, whether they be wise or simple, he will reveal himself in the toils, the conflicts, the sufferings which they shall pass through in his fellowship; and they shall learn in their own experience who he is.

We discover he is more than the babe in Bethlehem, more than the carpenter’s son, more than the greatest teacher ever to live. We come to know him as the Son of God. He never fashioned a statue, painted a picture, wrote a poem, or led an army. He never wore a crown or held a scepter or threw around his shoulder a purple robe. His forgiveness was unbounded, his patience inexhaustible, his courage without limit. Jesus changed men. He changed their habits, their opinions, their ambitions. He changed their tempers, their dispositions, their natures. He changed men’s hearts.

One thinks of the fisherman called Simon, better known to you and to me as Peter, chief among the apostles. Doubting, disbelieving, impetuous Peter was to remember the night that Jesus was led away to the high priest. Present were the priests whose greed and selfishness the Master had reproved, the elders whose hypocrisy he had branded, the scribes whose ignorance he had exposed. And then there were the Sadducees, considered the most cruel and dangerous opponents. This was the night that the throng “began to spit on [the Savior], and to cover his face, … to buffet him, … and the servants did strike him with the palms of their hands.” (Mark 14:65.)

Where was Peter, who had promised to die with him and never to deny him? The sacred record reveals, “And Peter followed him afar off, even into the palace of the high priest: and he sat with the servants, and warmed himself at the fire.” (Mark 14:54.) This was the night that Peter, in fulfillment of the Master’s prophecy, indeed did deny him thrice. Amidst the pushing, the jeers, and the blows, the Lord, in the agony of his humiliation, in the majesty of his silence, turned and looked upon Peter.

As one chronologer described the change, “It was enough. Peter knew no more danger, he feared no more death. He rushed into the night to meet the morning dawn. This broken-hearted penitent stood before the tribunal of his own conscience, and there his old life, his old shame, his old weakness, his old self was doomed to that death of godly sorrow which was to issue in a new and a nobler birth.” (Frederic W. Farrar, The Life of Christ, Portland, Oregon: Farrar Publications, 1964, p. 604.)

Then there was Saul of Tarsus, a scholar, familiar with the rabbinical writings in which certain modern scholars find such stores of treasure. For some reason, these writings did not reach Paul’s need, and he kept on crying, “O wretched man that I am! who shall deliver me from the body of this death?” (Rom. 7:24.) And then one day he met Jesus, and behold, all things became new. From that day to the day of his death, Paul urged men to “put off … the old man …” and to “put on the new man, which after God is created in righteousness and true holiness.” (Eph. 4:22, 24.)

The passage of time has not altered the capacity of the Redeemer to change men’s lives. As he said to the dead Lazarus, so he says to you and me: “… come forth.” (John 11:43.) Come forth from the despair of doubt. Come forth from the sorrow of sin. Come forth from the death of disbelief. Come forth to a newness of life. Come forth.

As we do, and direct our footsteps along the paths which Jesus walked, let us remember the testimony Jesus gave: “Behold, I am Jesus Christ, whom the prophets testified shall come into the world. … I am the light and … life of the world. …” (3 Ne. 11:10–11.) “I am the first and the last; I am he who liveth, I am he who was slain; I am your advocate with the Father.” (D\&C 110:4.)

To his testimony I add my witness: He lives. His prophet this day has been sustained—even President Spencer W. Kimball. I so testify, in the name of Jesus Christ. Amen.

\newpage
\settalk{My Personal Hall of Fame}
\setspeaker{Thomas S. Monson}
\settalkdate{October 1974}
\talkheading{My Personal Hall of Fame}{Thomas S. Monson}

President Kimball, as this conference comes speedily to its close, the words of the apostle Peter seem to reflect the feelings of each person who has attended the conference or who has listened to or viewed the proceedings.

Following his experience on the Mount of Transfiguration, Peter said to Jesus, “Lord, it is good for us to be here.” (Matt. 17:4.) President Kimball, it is good for all of us to have been here.

I pray that the same sweet spirit which has prevailed will continue with me as I respond to this opportunity to address you.

On a clear winter day I was driving with a friend along the freeway which connects downtown Manhattan, New York, with suburban Westchester. He pointed out to me several of the historic sights which abound in this area where man has indiscriminately constructed his ribbon of highway through the pathway of history.

Suddenly, like an old friend, there came into view Yankee Stadium. Here it was—the stadium of champions, the home of my boyhood heroes. Indeed, what boy has not idolized those who, before cheering thousands, played superbly well the game of baseball.

Being winter, the parking lot surrounding the stadium was deserted. Gone were the crowds, the peanut vendors, the ticket clerks. Still present were the memories of Babe Ruth, Lou Gehrig, and Joe DiMaggio. The record of their prowess and skill is forever safe—they have been elected to the prestigious Baseball Hall of Fame.

As with baseball, so with life. In the interior of our consciousness, each of us has a private hall of fame reserved exclusively for the real leaders who have influenced the direction of our lives. Relatively few of the many men who exercise authority over us from childhood through adult life meet our test for entry to this roll of honor. That test has very little to do with the outward trappings of power or an abundance of this world’s goods. The leaders whom we admit into this private sanctuary of our reflective meditation are usually those who set our hearts afire with devotion to the truth, who make obedience to duty seem the essence of manhood, who transform some ordinary routine occurrence so that it becomes a vista whence we see the person we aspire to be.

For a moment, perhaps each of us could be the qualifying judge through whom each hall of fame entry must pass. Whom would you nominate for prominent position? Whom would I? Candidates are many—competition severe.

I nominate to the Hall of Fame the name of Adam, the first man to live upon the earth. His citation is from Moses: “And Adam was obedient unto the commandments of the Lord.” (Moses 5:5.) Adam qualifies.

For patient endurance there must be nominated a perfect and upright man whose name was Job. Though afflicted as no other, he declared: “My witness is in heaven, and my record is on high.

“My friends scorn me: but mine eye poureth out tears unto God.” (Job 16:19–20.) “I know that my redeemer liveth.” (Job 19:25.) Job qualifies.

Every Christian would nominate the man Saul, better known as Paul the apostle. His sermons are as manna to the spirit, his life of service an example to all. This fearless missionary declared to the world: “For I am not ashamed of the gospel of Christ: for it is the power of God unto salvation.” (Rom. 1:16.) Paul qualifies.

Then there is the man called Simon Peter. His testimony of the Christ stirs the heart:

“When Jesus came into the coasts of Caesarea Philippi, he asked his disciples, saying, Whom do men say that I the Son of man am?

“And they said, Some say that thou art John the Baptist; some, Elias; and others, Jeremias, or one of the prophets.

“He saith unto them, But whom say ye that I am?

“And Simon Peter answered and said, Thou art the Christ, the Son of the living God.” (Matt. 16:13–16.) Peter qualifies.

Of another time and place we recall the testimony of Nephi:

“I will go and do the things which the Lord hath commanded, for I know that the Lord giveth no commandments unto the children of men, save he shall prepare a way for them that they may accomplish the thing which he commandeth them.” (1 Ne. 3:7.) Surely Nephi is worthy of a place in the Hall of Fame.

There is yet another I choose to nominate—even the Prophet Joseph Smith. His faith, his trust, his testimony are reflected by his own words, spoken as he went to Carthage Jail and martyrdom: “I am going like a lamb to the slaughter; but I am calm as a summer’s morning; I have a conscience void of offense towards God, and towards all men.” (D\&C 135:4.) He sealed his testimony with his blood. Joseph Smith qualifies.

In our selection of heroes, let us nominate also heroines. First, that noble example of fidelity—even Ruth. Sensing the grief-stricken heart of her mother-in-law, who suffered the loss of each of her two fine sons, and feeling perhaps the pangs of despair and loneliness which plagued the very soul of Naomi, Ruth uttered what has become that classic statement of loyalty: “Intreat me not to leave thee, or to return from following after thee: for whither thou goest, I will go; and where thou lodgest, I will lodge: thy people shall be my people, and thy God my God.” (Ruth 1:16.) Ruth’s actions demonstrated the sincerity of her words. There is place for her name in the Hall of Fame.

Shall we not name yet another, a descendant of honored Ruth? I speak of Mary of Nazareth, espoused to Joseph, destined to become the mother of the only sinless man to walk the earth. Her acceptance of this sacred and historic role is a hallmark of humility. “And Mary said, Behold the handmaid of the Lord; be it unto me according to thy word.” (Luke 1:38.) Surely Mary qualifies.

Could we ask the question, “What makes of these men heroes and these women heroines?” I answer, unwavering trust in an all-wise Heavenly Father and an abiding testimony concerning the mission of a divine Savior. This knowledge is like a golden thread woven through the tapestry of their lives.

Who is that King of Glory, even the Redeemer, for whom such heroes and heroines faithfully served and valiantly died? He is Jesus Christ, the Son of God, even our Savior.

His birth was foretold by prophets; angels heralded the announcement of his earthly ministry. To shepherds abiding in their fields came the glorious proclamation:

“Fear not: for, behold, I bring you good tidings of great joy, which shall be to all people.

“For unto you is born this day in the city of David a Saviour, which is Christ the Lord.” (Luke 2:10–11.)

This same Jesus “grew, and waxed strong in spirit, filled with wisdom: and the grace of God was upon him.” (Luke 2:40.) Baptized of John in the river known as Jordan, he commenced his official ministry to men. To the sophistry of Satan, Jesus turned his back. To the duty designated by his Father, he turned his face, pledged his heart, and gave his life. And what a sinless, selfless, noble, and divine life it was. Jesus labored. Jesus loved. Jesus served. Jesus wept. Jesus healed. Jesus taught. Jesus testified. On a cruel cross, Jesus died. From a borrowed sepulchre, Jesus came forth to eternal life.

The name—Jesus of Nazareth—the only name under heaven given among men whereby we must be saved, has singular place and honored distinction in our Hall of Fame.

Some may question: “But what is the value of such an illustrious list of heroes, even a private Hall of Fame?” I answer. When we obey, as did Adam, endure as did Job, teach as did Paul, testify as did Peter, serve as did Nephi, give ourselves as did the prophet Joseph, respond as did Ruth, honor as did Mary, and live as did Christ, we are born anew. All power becomes ours. Cast off forever is the old self and with it defeat, despair, doubt, and disbelief. To a newness of life we come—a life of faith, hope, courage, and joy. No task looms too large. No responsibility weighs too heavily. No duty is a burden. All things become possible.

In our quest for an example, we need not necessarily look to years gone by or to lives lived long ago. Let me illustrate. Today Craig Sudbury presides over a ward here in Salt Lake City, but let me turn back the clock just a few years to the day he and his mother came to my office prior to Craig’s departure for the Australia Melbourne Mission. Fred, Craig’s father, was noticeably absent. Twenty-five years earlier, Craig’s mother had married Fred, who did not share her love for the Church and indeed did not belong to the Church.

Craig confided to me his deep and abiding love for his parents. He shared his innermost hope that somehow, in some way, his father would be touched by the Spirit and open his heart to the gospel of Jesus Christ. He pleaded earnestly with me for a suggestion. I prayed for inspiration concerning how such a desire might be rewarded. Such inspiration came, and I said to Craig, “Serve the Lord with all your heart. Be obedient to your sacred calling. Each week write a letter to your parents and, on occasion, write to Dad personally and let him know that you love him, and tell him why you’re grateful to be his son.”

He thanked me and, with his mother, departed the office. I was not to see Craig’s mother for some 18 months. She came to the office and, in sentences punctuated by tears, said to me, “It has been almost two years since Craig departed for his mission. His faithful service has qualified him for positions of responsibility in the mission field, and he has never failed in writing a letter to us each week. Recently my husband Fred stood for the first time in a testimony meeting and said, ‘All of you know that I am not a member of the Church, but something has happened to me since Craig left for his mission. His letters have touched my soul. May I share one with you?

“‘Dear Dad, Today we taught a choice family about the plan of salvation and the blessings of exaltation in the celestial kingdom. I thought of our family. More than anything in the world, I want to be with you and with Mother in that kingdom. For me it just wouldn’t be a celestial kingdom if you were not there. I’m grateful to be your son, Dad, and want you to know that I love you. Your missionary son, Craig.’

“Fred then announced, ‘My wife doesn’t know what I plan to say. I love her and I love our son, Craig. After 26 years of marriage I have made my decision to become a member of the Church, for I know the gospel message is the word of God. I suppose I have known this truth for a long time, but my son’s mission has moved me to action. I have made arrangements for my wife and me to meet Craig when he completes his mission. I will be his final baptism as a full-time missionary of the Lord.’”

A young missionary with unwavering faith had participated with God in a modern-day miracle. His challenge to communicate with one whom he loved had been made more difficult by the barrier of the thousands of miles which lay between him and his father. But the spirit of love spanned the vast expanse of the blue Pacific, and heart spoke to heart in divine dialogue.

No hero stood so tall as did Craig, when in far-off Australia he stood with his father in water waist deep and, raising his right arm to the square, repeated those sacred words: “Fred Sudbury, having been commissioned of Jesus Christ, I baptize you in the name of the Father, and of the Son, and of the Holy Ghost.”

The prayer of a mother, the faith of a father, the service of a son brought forth the miracle of God. Mother, father, son—each qualifies in a Hall of Fame.

May they and each of us so live as to merit the heavenly pronouncement:

“I, the Lord, am merciful and gracious unto those who fear me, and delight to honor those who serve me in righteousness and in truth unto the end.

“Great shall be their reward and eternal shall be their glory.” (D\&C 76:5–6.)

Our place in an everlasting and eternal Hall of Fame will thereby be assured. This is my earnest plea as I leave with you my witness that Jesus of Nazareth is our Savior and Redeemer, even our Advocate with the Father. In the name of Jesus Christ, the Lord. Amen.

\newpage
\settalk{The Way Home}
\setspeaker{Thomas S. Monson}
\settalkdate{April 1975}
\talkheading{The Way Home}{Thomas S. Monson}

Overlooking the azure blue waters of the famed Sea of Galilee is a historic landmark: the Mount of Beatitudes. Like a living sentinel with an eyewitness testimony, this silent friend seems to declare: “Here it was that the greatest person who ever lived delivered the greatest sermon ever given—the Sermon on the Mount.”

Instinctively the visitor turns to the Gospel of St. Matthew and reads: “And seeing the multitudes, he went up into a mountain: and when he was set, his disciples came unto him: And he opened his mouth, and taught them.” (Matt. 5:1–2.) Among the truths which he taught was this solemn statement: “Enter ye in at the strait gate: for wide is the gate, and broad is the way, that leadeth to destruction, and many there be which go in thereat:

“Because strait is the gate, and narrow is the way, which leadeth unto life, and few there be that find it.” (Matt. 7:13–14.)

Ageless in its application, wise men throughout the generations of time have sought to live by this simple statement.

When Jesus of Nazareth personally walked the rock-strewn pathways of the Holy Land, he, as the Good Shepherd, showed all who would believe how they might follow that narrow way and enter that strait gate to life eternal. “Come, follow me,” he invited. “I am the way.”

Little wonder that men did tarry for the outpouring of the Holy Ghost on the day of Pentecost. It was the gospel of Jesus Christ that was to be preached, his work that was to be done, and his apostles at the head of his church who were entrusted with the work.

History records that most men indeed did not come unto him, nor did they follow the way he taught. Crucified was the Lord, slain were the apostles, rejected was the truth. The bright daylight of enlightenment slipped away, and the lengthening shadows of a black night enshrouded the earth.

One word and one word alone describes the dismal state that prevailed: apostasy. Generations before, Isaiah had prophesied: “Darkness shall cover the earth, and gross darkness the people.” (Isa. 60:2.) Amos had foretold of a famine in the land: “Not a famine of bread, nor a thirst for water, but of hearing the words of the Lord.” (Amos 8:11.) Had not Peter warned of false teachers bringing damnable heresies, and Paul predicted that the time would come when sound doctrine would not be endured?

The dark ages of history seemed never to end. Was there to be no termination to this blasphemous night? Had a loving Father forgotten mankind? Would he send forth no heavenly messengers as in former days?

Honest men with yearning hearts, at the peril of their very lives, attempted to establish points of reference, that they might find the true way. The day of the reformation was dawning, but the path ahead was difficult. Persecutions would be severe, personal sacrifice overwhelming, and the cost beyond calculation. The reformers were like pioneers blazing wilderness trails in a desperate search for those lost points of reference which, they felt, when found would lead mankind back to the truth Jesus taught.

When John Wycliffe and others completed the first English translation of the entire Bible from the Latin Vulgate, the then church authorities did all they could to destroy it. Copies had to be written by hand and in secret. The Bible had been regarded as a closed book forbidden to be read by the common people. Many of the followers of Wycliffe were severely punished and some burned at the stake.

Martin Luther asserted the Bible’s supremacy. His study of the scriptures led him to compare the doctrines and practices of the church with the teachings of the scriptures. Luther stood for the responsibility of the individual and the rights of the individual conscience, and this he did at the imminent risk of his life. Though threatened and persecuted, yet he declared boldly: “Here I stand, I cannot do otherwise. God help me.”

John Huss, speaking out fearlessly against the corruption within the church, was taken outside the city to be burned. He was chained by the neck to a stake, and straw and wood were piled around his body to the chin and sprinkled with resin; and he was asked finally if he would recant. As the flames arose, he sang, but the wind blew the fire into his face, and his voice was stilled.

Zwingli of Switzerland attempted through his writings and teachings to rethink all Christian doctrine in consistently biblical terms. His most famous statement thrills the heart: “What does it matter? They can kill the body but not the soul.”

And who cannot today appreciate the words of John Knox? “A man with God is always in the majority.”

John Calvin, prematurely aged by sickness and by the incessant labors he had undertaken, summed up his personal philosophy with the statement: “Our wisdom … consists almost entirely of two parts: the knowledge of God and the knowledge of ourselves.”

Others could indeed be mentioned, but a comment concerning William Tyndale would perhaps suffice. Tyndale felt that the people had a right to know what was promised to them in the scriptures. To those who opposed his work of translation, he declared: “If God spare my life, … I will cause a boy that driveth the plough shall know more of the scripture than thou dost.”

Such were the teachings and lives of the great reformers. Their deeds were heroic, their contributions many, their sacrifices great—but they did not restore the gospel of Jesus Christ.

Of the reformers one could ask, “Was their sacrifice in vain? Was their struggle futile?” I answer with a resounding “No!” The Holy Bible was now within the grasp of the people. Each man could better find his way. Oh, if only all could read and all could understand. But some could read, and others could hear; and every man had access to God through prayer.

The long-awaited day of restoration did indeed come. But let us review that significant event in the history of the world by recalling the testimony of the plowboy who became a prophet, the witness who was there—even Joseph Smith.

Describing his experience, Joseph said: “There was in the place where we lived an unusual excitement on the subject of religion. … It … became general … [creating] division amongst the people, some crying, ‘Lo, here!’ and others, ‘Lo, there!’

“… I was one day reading the Epistle of James, first chapter and fifth verse, which reads: ‘If any of you lack wisdom, let him ask of God, that giveth to all men liberally, and upbraideth not; and it shall be given him.’

“Never did any passage of scripture come with more power to the heart of man than this did at this time to mine. It seemed to enter with great force into every feeling of my heart. I reflected on it again and again, knowing that if any person needed wisdom from God, I did; for how to act I did not know, and unless I could get more wisdom than I then had, I would never know; for the teachers of religion … understood the same passages of scripture so differently as to destroy all confidence in settling the question by an appeal to the Bible.

“At length I came to the conclusion that I must either remain in darkness and confusion, or else I must do as James directs, that is, ask of God. …

“So, in accordance with this, my determination to ask of God, I retired to the woods to make the attempt. It was on the morning of a beautiful, clear day, early in the spring of eighteen hundred and twenty.

“… I kneeled down and began to offer up the desire of my heart to God. …

“I saw a pillar of light exactly over my head, above the brightness of the sun, which descended gradually until it fell upon me.

“… When the light rested upon me I saw two Personages, whose brightness and glory defy all description, standing above me in the air. One of them spake unto me, calling me by name and said, pointing to the other—This is My Beloved Son. Hear Him!” (JS—H 1:5–17.)

The Father and the Son, Jesus Christ, had appeared to Joseph Smith. The morning of the dispensation of the fulness of times had come, dispelling the darkness of the long generations of spiritual night. As in the creation, light was to replace darkness; day was to follow night.

From then to now, truth has been and is available to us. Like the children of Israel in former times, endless days of wandering now can end with our entry to a personal promised land.

The restoration of the gospel dispels the gloom described in our time by the noted educator Robert Gordon Sproul. He had looked at the churches of America and declared:

“We have … the peculiar spectacle of a nation which, to some imperfect but nevertheless considerable extent, practices Christianity without actively believing in Christianity. We are asked to turn to the church for our enlightenment, but when we do so we find that the voice of the church is not inspired. The voice of the church today, we find, is the echo of our own voices. And the result of this experience, already manifest, is disillusionment. There is only one way out of the spiral. The way out is the sound of a voice, not our voice, but a voice coming from something not ourselves, in the existence of which we cannot disbelieve. It is the earthly task of the pastors to hear this voice, to cause us to hear it, and to tell us what it says. If they cannot hear it, or if they fail to tell us, we, as laymen, are utterly lost. Without it we are no more capable of saving the world than we were capable of creating it in the first place.” (Vital Speeches, Sept. 1, 1940, p. 701.)

Perhaps the famed Winston Churchill best declared the world’s pressing need. Said he: “I have lived perhaps longer experience than almost anyone, and I have never brooded over a situation which demanded more patience, composure, courage and perseverance than that which unfolds itself before us today: The need of a prophet.”

Today we have heard God’s prophet speak—even President Spencer W. Kimball. Today there goes forth from this pulpit an invitation to people throughout the world: Come from your wandering way, weary traveler. Come to the gospel of Jesus Christ. Come to that heavenly haven called home. Here you will discover the truth. Here you will learn the reality of the Godhead, the comfort of the plan of salvation, the sanctity of the marriage covenant, the power of personal prayer. Come home!

From our youth many of us may remember the story of a very young boy who was abducted from his parents and home and taken to a village situated far away. Under these conditions the small boy grew to young manhood without a knowledge of his actual parents or earthly home. Within his heart there came a yearning to return to that village called home.

But where was home to be found? Where were his mother and father to be discovered? Oh, if only he could remember even their names, his task would be less hopeless. Desperately he sought to recall even a glimpse of his childhood.

Like a flash of inspiration, he remembered the sound of a bell which, from the tower atop the village church, pealed its welcome each Sabbath morning. From village to village the young man wandered, ever listening for that familiar bell to chime. Some bells were similar, others far different from the sound he remembered.

At length the weary young man stood one Sunday morning before a church of a typical town. He listened carefully as the bell began to peal. The sound was familiar. It was unlike any other he had heard, save that bell which pealed in the memory of his childhood days. Yes, it was the same bell. Its ring was true. His eyes filled with tears. His heart rejoiced in gladness. His soul overflowed with gratitude. The young man dropped to his knees, looked upward beyond the bell tower—even toward heaven—and in a prayer of gratitude whispered, “Thanks be to God. I’m home.”

Like the peal of a remembered bell will be the truth of the gospel of Jesus Christ to the soul of him who earnestly seeks. Many of you have traveled long in a personal quest for that which rings true. The Church of Jesus Christ of Latter-day Saints sends forth to you an earnest appeal. Open your doors to the missionaries. Open your minds to the word of God. Open your hearts, even your very souls, to the sound of that still, small voice which testifies of truth. As the prophet Isaiah promised: “Thine ears shall hear a word … saying, This is the way, walk ye in it.” (Isa. 30:21.) Then, like the boy of whom I’ve spoken, you too will, on bended knee, say to your God and mine: “I’m home!”

May such be the blessing of all, I pray in the name of Jesus Christ. Amen.

\newpage
\settalk{The Faith of a Child}
\setspeaker{Thomas S. Monson}
\settalkdate{October 1975}
\talkheading{The Faith of a Child}{Thomas S. Monson}

What a truly glorious period of the year is conference time. Temple Square in Salt Lake City is the gathering place for tens of thousands who travel far, that they might hear the word of the Lord. Today the Tabernacle is filled to overflowing. Friendly conversation has been replaced by the music of the choir and the voices of those who pray and who speak. A sweet reverence fills the air.

It is a humbling experience to gaze into your faces and to appreciate your faith and devotion to truth. Patiently do you sit on those historic benches which the passing of time has somehow not made more comfortable.

Particularly am I grateful for the children who are here. In the balcony to my left I see a beautiful girl of perhaps ten years. Sweet little one, I do not know your name or whence you have come. This, however, I do know: the innocence of your smile and the tender expression of your eyes have persuaded me to place aside for a future time the message I had prepared for this occasion. Today, I am impressed to speak to you.

When I was a boy your age, I too had a teacher in Sunday School. From the Bible she would read to us of Jesus, the Redeemer and Savior of the world. One day she taught us how the little children were brought unto him, that he should put his hands on them and pray. His disciples rebuked those that brought the children. “But when Jesus saw it, he was much displeased, and said unto them, Suffer the little children to come unto me, and forbid them not: for of such is the kingdom of God.” (Mark 10:14.)

That lesson has never left me. Indeed, just a few months ago I relearned its meaning and partook of its power. My teacher was the Lord. May I share with you this experience.

Far away from Salt Lake City, and some eighty miles from Shreveport, Louisiana, lives the Jack Methvin family. Mother, dad, and the boys are members of The Church of Jesus Christ of Latter-day Saints. Until just recently there was a lovely daughter who, by her presence, graced that home. Her name was Christal. She was but ten years old when death ended her earthly sojourn.

Christal liked to run and play on the spacious ranch where her family lives. She could ride horses skillfully and excelled in 4-H work, winning awards in the local and state fairs. Her future was bright, and life was wonderful. Then there was discovered on her leg an unusual lump. The specialists in New Orleans completed their diagnosis and rendered their verdict: carcinoma. The leg must be removed.

She recovered well from the surgery, lived as buoyantly as ever and never complained. Then the doctors discovered that the cancer had spread to her tiny lungs. The Methvin family did not despair, but rather planned a flight to Salt Lake City. Christal could receive a blessing from one of the General Authorities. The Methvins knew none of the Brethren personally, so opening before Christal a picture of all the General Authorities, a chance selection was made. By sheer coincidence, my name was selected.

Christal never made the flight to Salt Lake City. Her condition deteriorated. The end drew nigh. But her faith did not waver. To her parents, she said, “Isn’t stake conference approaching? Isn’t a General Authority assigned? And why not Brother Monson? If I can’t go to him, the Lord can send him to me.”

Meanwhile in Salt Lake City, with no knowledge of the events transpiring in Shreveport, a most unusual situation developed. For the weekend of the Shreveport Louisiana Stake Conference, I had been assigned to El Paso, Texas. President Ezra Taft Benson called me to his office and explained that one of the other Brethren had done some preparatory work regarding the stake division in El Paso. He asked if I would mind were another to be assigned to El Paso and I assigned elsewhere. Of course there was no problem—anywhere would be fine with me. Then President Benson said, “Brother Monson, I feel impressed to have you visit the Shreveport Louisiana Stake.” The assignment was accepted. The day came. I arrived in Shreveport.

That Saturday afternoon was filled with meetings—one with the stake presidency, one with priesthood leaders, one with the patriarch, then yet another with the general leadership of the stake. Rather apologetically, Stake President Charles F. Cagle asked if my schedule would permit me time to provide a blessing to a ten-year-old girl afflicted with cancer. Her name: Christal Methvin. I responded that, if possible, I would do so, and then inquired if she would be at the conference, or was she in a Shreveport hospital? Knowing the time was tightly scheduled, President Cage almost whispered that Christal was confined to her home—more than eighty miles from Shreveport!

I examined the schedule of meetings for that evening and the next morning—even my return flight. There simply was no available time. An alternative suggestion came to mind. Could we not remember the little one in our public prayers at conference? Surely the Lord would understand. On this basis, we proceeded with the scheduled meetings.

When the word was communicated to the Methvin family, there was understanding but a trace of disappointment as well. Hadn’t the Lord heard their prayers? Hadn’t he provided that Brother Monson would come to Shreveport? Again the family prayed, asking for a final favor—that their precious Christal would realize her desire.

At the very moment the Methvin family knelt in prayer, the clock in the stake center showed the time to be 7:45. The leadership meeting had been inspirational. I was sorting my notes, preparing to step to the pulpit, when I heard a voice speak to my spirit. The message was brief, the words familiar: “Suffer the little children to come unto me, and forbid them not: for of such is the kingdom of God.” (Mark 10:14.) My notes became a blur. My thoughts turned to a tiny girl in need of a blessing. The decision was made. The meeting schedule was altered. After all, people are more important than meetings. I turned to Bishop James Serra and asked that he leave the meeting and advise the Methvins.

The Methvin family had just arisen from their knees when the telephone rang and the message was relayed that early Sunday morning—the Lord’s day—in a spirit of fasting and prayer, we would journey to Christal’s bedside.

I shall ever remember and never forget that early-morning journey to a heaven the Methvin family calls home. I have been in hallowed places—even holy houses—but never have I felt more strongly the presence of the Lord than in the Methvin home. Christal looked so tiny lying peacefully on such a large bed. The room was bright and cheerful. The sunshine from the east window filled the bedroom with light as the Lord filled our hearts with love.

The family surrounded Christal’s bedside. I gazed down at a child who was too ill to rise—almost too weak to speak. Her illness had now rendered her sightless. So strong was the spirit that I fell to my knees, took her frail hand in mine, and said simply, “Christal, I am here.” She parted her lips and whispered, “Brother Monson, I just knew you would come.” I looked around the room. No one was standing. Each was on bended knee. A blessing was given. A faint smile crossed Christal’s face. Her whispered “thank you” provided an appropriate benediction. Quietly, each filed from the room.

Four days later, on Thursday, as Church members in Shreveport joined their faith with the Methvin family and Christal’s name was remembered in a special prayer to a kind and loving Heavenly Father, the pure spirit of Christal Methvin left its disease-ravaged body and entered the paradise of God.

For those of us who knelt that Sabbath day in a sun-filled bedroom, and particularly for Christal’s mother and father as they enter daily that same room and remember how she left it, the immortal words of Eugene Field will bring back precious memories:

\begin{verse}
The little toy dog is covered with dust,\\
But sturdy and staunch he stands;\\
And the little toy soldier is red with rust,\\
And his musket moulds in his hands.\\
Time was when the little toy dog was new,\\
And the soldier was passing fair,\\
And that was the time when our Little Boy Blue\\
Kissed them and put them there.
\end{verse}

(“Little Boy Blue,” One Hundred and One Famous Poems, Chicago: Reilly \& Lee, 1958, p. 15.)

For us there is no need to wonder or to wait. Said the Master, “I am the resurrection, and the life: he that believeth in me, though he were dead, yet shall he live: And whosoever liveth and believeth in me shall never die.” (John 11:25–26.) To you, Jack and Nancy Methvin, he speaks: “Peace I leave with you, my peace I give unto you: not as the world giveth, give I unto you. Let not your heart be troubled, neither let it be afraid.” (John 14:27.) And from your sweet Christal could well come the comforting expression: “I go to prepare a place for you. … that where I am, there ye may be also.” (John 14:2–3.)

To you, my little friend in the upper balcony, and to believers everywhere, I bear witness that Jesus of Nazareth does love little children, that he listens to your prayers and responds to them. The Master did indeed utter those words, “Suffer the little children to come unto me, and forbid them not: for of such is the kingdom of God.” (Mark 10:14.)

I know these are the words he spoke to the throng gathered on the coast of Judea by the waters of Jordan—for I have read them.

I know these are the words he spoke to an apostle on assignment in Shreveport, Louisiana—for I heard them.

To these truths I bear record, in the name of Jesus Christ. Amen.

\newpage
\settalk{Hopeless Dawn—Joyful Morning}
\setspeaker{Thomas S. Monson}
\settalkdate{April 1976}
\talkheading{Hopeless Dawn—Joyful Morning}{Thomas S. Monson}

I am truly honored to follow at this pulpit the President of the Church, even the prophet of God, Spencer W. Kimball. My thoughts today have been centered on the land of his forebears, even Great Britain.

London, England, is steeped in history. Who has not heard of Trafalgar Square, Buckingham Palace, Big Ben, Westminster Abbey, or the River Thames? Of lesser renown, yet priceless in value, are the truly magnificent galleries of art situated in this city of culture.

One gray, wintry afternoon I visited the famed Tate Gallery. I marveled at the landscapes of Gainsborough, the portraits of Rembrandt, and the storm-laden clouds of Constable. Tucked away in a quiet corner of the third floor was a masterpiece which not only caught my attention but captured my heart. The artist, Frank Bramley, had painted a humble cottage facing a wind-swept sea. Kneeling at the side of an older woman was a young, grief-filled wife who mourned the loss of her seafaring husband. The spent candle at the window ledge told of her fruitless, night-long vigil. The huge gray clouds were all that remained of the tempest-torn night.

I sensed her loneliness. I felt her despair. The hauntingly vivid inscription which the artist gave to his work told the tragic story. It read: A Hopeless Dawn.

How the young widow longed for the comfort, even the reality, of Robert Louis Stevenson’s “Requiem”:

\begin{verse}
“Home is the sailor, home from the sea,\\
And the hunter home from the hill.”
\end{verse}

For her and many others who have loved and lost dear ones, each dawn is hopeless. Such is the experience of those who regard the grave as the end and immortality as but a dream.

The famed scientist, Madame Marie Curie, returned to her home the night of the funeral for her husband, Pierre Curie, who was killed in an accident in the streets of Paris, and made this entry in her diary:

“They filled the grave and put sheaves of flowers on it. Everything is over. Pierre is sleeping his last sleep beneath the earth; it is the end of everything, everything, everything.” (Vincent Sheehan, trans., Madame Curie: A Biography by Eve Curie, Garden City, New York: Garden City Publishing Co., 1943, p. 249.)

The atheist, Bertrand Russell, adds his testament: “No fire, no heroism, no integrity of thought and feeling can preserve an individual life beyond the grave.” And Schopenhauer, the German philosopher and pessimist, was even more bitter. He wrote: “To desire immortality is to desire the eternal perpetuation of a great mistake.”

In reality, every thoughtful person has asked himself the universal question, best phrased by the venerable, perfect, and upright man named Job, who, centuries ago, asked: “If a man die, shall he live again?” (Job 14:14.) Through inspiration from on high, Job answered his own question:

“Oh that my words were now written! oh that they were printed in a book!

“That they were graven with an iron pen and lead in the rock for ever!

“For I know that my redeemer liveth, and that he shall stand at the latter day upon the earth. …

“In my flesh shall I see God.” (Job 19:23–26.)

Few statements in scripture reveal so clearly a divine truth as does Paul’s epistle to the Corinthians: “For as in Adam all die, even so in Christ shall all be made alive.” (1 Cor. 15:22.)

Frequently, death comes as an intruder. It is an enemy that suddenly appears in the midst of life’s feast, putting out its lights and gaiety. It visits the aged as they walk on faltering feet. Its summons is heard by those who have scarcely reached midway in life’s journey, and often it hushes the laughter of little children. Death lays its heavy hand upon those dear to us and at times leaves us baffled and wondering. In certain situations, as in great suffering and illness, death comes as an angel of mercy. But for the most part, we think of it as the enemy of human happiness.

The plight of the widow, for instance, is a recurring theme throughout Holy Writ. Our hearts go out to the widow at Zarephath. Gone was her husband. Consumed was her scant supply of food. Starvation and death awaited. Then came Elijah, God’s prophet, who brought to her, through her faith, heavenly peace.

We remember also the widow of Nain. She grieved over the loss of her son. Her abiding faith, her earnest prayer, brought forth a divine gift. The Lord Jesus Christ returned to her and to life her precious son.

But what of today? Is there comfort for the grieving heart? Does God remember still the widow in her travail?

Not far from this tabernacle there lived two sisters. Each had two handsome sons. Each had a loving husband. Each lived in comfort, prosperity, and good health. Then the grim reaper visited their homes. First, each lost a son; then the husband and father. Friends visited; words brought a measure of comfort; but grief continued unrelieved.

The years passed. Hearts remained broken. The two sisters sought and achieved seclusion. They shut themselves off from the world which surrounded them. Alone they remained with their remorse. Then there came to a latter-day prophet of God, who knew well these two sisters, the inspiration of the Lord which directed him to their plight. Elder Harold B. Lee left his busy office and visited the penthouse home of the lonely widows. He listened to their pleadings. He felt the sorrow of their hearts. Then he called them to the service of God and to mankind. Each looked outward into the lives of others and upward into the face of God. Peace replaced turmoil. Confidence dispelled despair. God had once again remembered the widow and, through a prophet, brought divine comfort.

The darkness of death can ever be dispelled by the light of revealed truth. “I am the resurrection, and the life,” spoke the Master. “He that believeth in me, though he were dead, yet shall he live:

“And whosoever liveth and believeth in me shall never die.” (John 11:25–26.)

This reassurance, yes, even holy confirmation of life beyond the grave, could well be the peace promised by the Savior when he assured his disciples:

“Peace I leave with you, my peace I give unto you: not as the world giveth, give I unto you. Let not your heart be troubled, neither let it be afraid.” (John 14:27.)

“In my Father’s house are many mansions: if it were not so, I would have told you. I go to prepare a place for you …

“That where I am, there ye may be also.” (John 14:2–3.)

Out of the darkness and horror of Calvary came the voice of the Lamb, saying, “Father, into thy hands I commend my spirit.” (Luke 23:46.) And the dark was no longer dark, for he was with his Father. He had come from God and to God he had returned. So also those who walk with God in this earthly pilgrimage know from blessed experience that he will not abandon his children who trust in him. In the night of death his presence will be “better than a light and safer than a known way.” (From “God Knows,” by Minnie Louise Haskins.)

The reality of the resurrection was voiced by the martyr Stephen as he looked upward and cried, “I see the heavens opened, and the Son of man standing on the right hand of God.” (Acts 7:56.)

Saul, on the road to Damascus, had a vision of the risen, exalted Christ. Later, as Paul, defender of truth and fearless missionary in the service of the Master, he bore witness of the risen Lord as he declared to the saints at Corinth:

“Christ died for our sins according to the scriptures. …

“He was buried, and … he rose again the third day according to the scriptures. …

“He was seen of Cephas, then of the twelve:

“After that, he was seen of above five hundred brethren at once. …

“He was seen of James; then of all the apostles.

“And last of all he was seen of me.” (1 Cor. 15:3–8.)

In our dispensation, this same testimony was spoken boldly by the Prophet Joseph Smith, as he and Sidney Rigdon testified:

“And now, after the many testimonies which have been given of him, this is the testimony, last of all, which we give of him: That he lives!

“For we saw him, even on the right hand of God; and we heard the voice bearing record that he is the Only Begotten of the Father—

“That by him, and through him, and of him, the worlds are and were created, and the inhabitants thereof are begotten sons and daughters unto God.” (D\&C 76:22–24.)

This is the knowledge that sustains. This is the truth that comforts. This is the assurance that guides those bowed down with grief out of the shadows and into the light.

Such help is not restricted to the elderly, the well-educated, or a select few. It is available to all.

Several years ago, the Salt Lake City newspapers published an obituary notice of a close friend—a mother and wife taken by death in the prime of her life. I visited the mortuary and joined a host of persons gathered to express condolence to the distraught husband and motherless children. Suddenly the smallest child, Kelly, recognized me and took my hand in hers. “Come with me,” she said, and she led me to the casket in which rested the body of her beloved mother. “I’m not crying,” she said, “and neither must you. Many times my mommy told me about death and life with Heavenly Father. I belong to my mommy and my daddy. We’ll all be together again.” To my mind came the words of the Psalmist: “Out of the mouth of babes … hast thou ordained strength.” (Ps. 8:2.)

Through tear-moistened eyes, I saw my young friend’s beautiful and faith-filled smile. For her, whose tiny hand yet clasped mine, there would never be a hopeless dawn. Sustained by her unfailing testimony, knowing that life continues beyond the grave, she, her father, her brothers, her sisters, and indeed all who share this knowledge of divine truth can declare to the world: “Weeping may endure for a night, but joy cometh in the morning.” (Ps. 30:5.)

With all the strength of my soul, I testify that God lives, that his Beloved Son is the firstfruits of the resurrection, that the gospel of Jesus Christ is that penetrating light that makes of every hopeless dawn a joyful morning.

In the name of Jesus Christ. Amen.

\newpage
\settalk{Which Road Will You Travel?}
\setspeaker{Thomas S. Monson}
\settalkdate{October 1976}
\talkheading{Which Road Will You Travel?}{Thomas S. Monson}

A ribbon of black asphalt wends its way through the mountains of northern Utah into the valley of the Great Salt Lake, then meanders southward on its appointed course. Interstate 15 is its official name. This super freeway carries the output of factories, the products of commerce, and masses of humanity toward appointed destinations.

Several days ago, while driving to my home, I approached the entrance to Interstate 15. At the on-ramp I noticed three hitchhikers, each one of whom carried a homemade sign which announced his desired destination. One sign read “Los Angeles,” while a second carried the designation “Boise.” However, it was the third sign which not only caught my attention but caused me to reflect and ponder its message. The hitchhiker had lettered not Los Angeles, California, nor Boise, Idaho, on the cardboard sign which he held aloft. Rather, his sign consisted of but one word and read simply “ANYWHERE.”

Here was one who was content to travel in any direction, according to the whim of the driver who stopped to give him a free ride. What an enormous price to pay for such a ride. No plan. No objective. No goal. The road to anywhere is the road to nowhere, and the road to nowhere leads to dreams sacrificed, opportunities squandered, and a life unfulfilled.

Unlike the youthful hitchhiker, you and I have the God-given gift to choose the direction we go. Indeed, the apostle Paul likened life to a race with a clearly defined goal. To the saints at Corinth he urged: “Know ye not that they which run in a race run all, but one receiveth the prize? So run, that ye may obtain.” (1 Cor. 9:24.) In our zeal, let us not overlook the sage counsel from Ecclesiastes: “The race is not to the swift, nor the battle to the strong.” (Eccl. 9:11.) Actually, the prize belongs to him who endures to the end.

Each must ask himself the questions: “Where am I going?” “How do I intend to get there?” “Really, what is my divine destiny?”

When I reflect on the race of life, I remember another race, even from childhood days. Perhaps a shared experience from this period will assist in formulating answers to these significant and universally asked questions.

When I was about ten, my boyfriends and I would take pocketknives in hand and from the soft wood of a willow tree fashion small toy boats. With a triangular-shaped cotton sail in place, each would launch his crude craft in a race down the relatively turbulent waters of the Provo River. We would run along the river’s bank and watch the tiny vessels sometimes bobbing violently in the swift current and at other times sailing serenely as the water deepened.

During such a race, we noted that one boat led all the rest toward the appointed finish line. Suddenly, the current carried it too close to a large whirlpool, and the boat heaved to its side and capsized. Around and around it was carried, unable to make its way back into the main current. At last it came to rest at the end of the pool, amid the flotsam and jetsam which surrounded it, held fast by the fingerlike tentacles of the grasping green moss.

The toy boats of childhood had no keel for stability, no rudder to provide direction, and no source of power. Like the hitchhiker, their destination was “ANYWHERE,” but inevitably downstream.

We have been provided divine attributes to guide our destiny. We entered mortality not to float with the moving currents of life, but with the power to think, to reason, and to achieve.

Our Heavenly Father did not launch us on our eternal journey without providing the means whereby we could receive from Him God-given guidance to ensure our safe return at the end of life’s great race. Yes, I speak of prayer. I speak, too, of the whisperings from that still, small voice within each of us; and I do not overlook the holy scriptures, written by mariners who successfully sailed the seas we too must cross.

Individual effort will be required of us. What can we do to prepare? How can we assure a safe voyage?

First, we must visualize our objective. What is our purpose? The Prophet Joseph Smith counseled: “Happiness is the object and design of our existence; and will be the end thereof, if we pursue the path that leads to it; and this path is virtue, uprightness, faithfulness, holiness, and keeping all the commandments of God.” (Teachings of the Prophet Joseph Smith, pp. 255–56.) In this one sentence we are provided not only a well-defined goal, but also the way we might achieve it.

Second, we must make continuous effort. Have you noticed that many of the most cherished of God’s dealings with His children have been when they were engaged in a proper activity? The visit of the Master to His disciples on the way to Emmaus, the good Samaritan on the road to Jericho, even Nephi on his return to Jerusalem, and Father Lehi en route to the precious land of promise. Let us not overlook Joseph Smith on the way to Carthage, and Brigham Young on the vast plains to the valley home of the Saints.

Third, we must not detour from our determined course. In our journey we will encounter forks and turnings in the road. There will be the inevitable trials of our faith and the temptations of our times. We simply cannot afford the luxury of a detour, for certain detours lead to destruction and spiritual death. Let us avoid the moral quicksands that threaten on every side, the whirlpools of sin, and the crosscurrents of uninspired philosophies. That clever pied piper called Lucifer still plays his lilting melody and attracts the unsuspecting away from the safety of their chosen pathway, away from the counsel of loving parents, away from the security of God’s teachings. His tune is ever so old, his words ever so sweet. His prize is everlasting. He seeks not the refuse of humanity but the very elect of God. King David listened, then followed, then fell. But then so did Cain in an earlier era, and Judas Iscariot in a later one.

Fourth, to gain the prize we must be willing to pay the price. The apprentice does not become the master craftsman until he has qualified. The lawyer does not practice until he has passed the bar. The doctor does not attend our needs until internship has been completed.

\begin{verse}
You are the fellow that has to decide\\
Whether you’ll do it or toss it aside. …\\
Whether you’ll try for the goal that’s afar\\
Or just be contented to stay where you are.
\end{verse}

Edgar A. Guest, “You,” The Light of Faith, Chicago: Reilly and Lee, 1926, p. 133.

Let us remember how Saul the persecutor became Paul the proselyter, how Simon, the fisherman, became Peter, the apostle of spiritual power. And let us be mindful that before Easter there had to be a cross.

Our example in the race of life could well be our elder brother, even the Lord. As a small boy, he provided a watchword: “Wist ye not that I must be about my Father’s business?” (Luke 2:49.) As a grown man he taught by example compassion, love, obedience, sacrifice, and devotion. To you and to me his summons is still the same: “Come, follow me.”

One who listened and who followed was the Mormon missionary Randall Ellsworth, about whom you may have read in your daily newspaper or watched on the television set in your home.

Six months ago, while serving in Guatemala as a missionary for The Church of Jesus Christ of Latter-day Saints, Randall Ellsworth survived the devastating earthquake which hurled a beam down on his back, paralyzed his legs, and severely damaged his kidneys.

After receiving emergency medical treatment, Randall was flown to a large hospital near his home in Rockville, Maryland. While confined there, a television newscaster conducted with Randall an interview which I witnessed through the miracle of television. The reporter asked, “Can you walk?” The answer, “Not yet, but I will.” “Do you think you will be able to complete your mission?” Came the reply, “Others think not, but I will.”

With microphone in hand, the reporter continued: “I understand you have received a special letter containing a get-well message from none other than the president of the United States.” “Yes,” replied Randall, “I am very grateful to President Ford for his thoughtfulness; but I received another letter, not from the president of my country, but from the president of my church—The Church of Jesus Christ of Latter-day Saints—even President Spencer W. Kimball. This I cherish. With him praying for me, and the prayers of my family, my friends, and my missionary companions, I will return to Guatemala. The Lord wanted me to preach the gospel there for two years, and that’s what I intend to do.”

I turned to my wife and commented, “He surely must not know the extent of his injuries. Our official medical reports would not permit us to expect such a return to Guatemala.”

How grateful am I that the day of faith and the age of miracles are not past history but continue with us even now.

The newspapers and the television cameras directed their attention to more immediate news as the days turned to weeks and the weeks to months. The words of Rudyard Kipling described Randall Ellsworth’s situation:

\begin{verse}
The tumult and the shouting dies;\\
The Captains and the Kings depart:\\
Still stands Thine ancient sacrifice,\\
An humble and a contrite heart.\\
Lord God of Hosts, be with us yet,\\
Lest we forget—lest we forget!
\end{verse}

Rudyard Kipling’s Verse, Garden City, New York: Doubleday, 1946, p. 327.

And God did not forget him who possessed an humble and a contrite heart, even Elder Randall Ellsworth. Little by little the feeling in his legs began to return. In his own words, Randall described the recovery: “The thing I did was always to keep busy, always pushing myself. In the hospital I asked to do therapy twice a day instead of just once. I wanted to walk again on my own.”

When the Missionary Committee evaluated the amazing medical progress Randall Ellsworth had made, word was sent to him that his return to Guatemala was authorized. Said he, “At first I was so happy I didn’t know what to do. Then I went into my bedroom and I started to cry. Then I dropped to my knees and thanked my Heavenly Father.”

Two months ago Randall Ellsworth walked aboard the plane that carried him back to the mission to which he was called and back to the people whom he loved. Behind he left a trail of skeptics, a host of doubters, but also hundreds amazed at the power of God, the miracle of faith, and the reward of determination. Ahead lay honest, God-fearing, and earnestly seeking sons and daughters of our Heavenly Father. They shall hear His word. They shall learn His truth. They shall accept His ordinances. A modern-day Paul, who too overcame his “thorn in the flesh,” has returned to teach them the truth, to lead them to life eternal.

Like Randall Ellsworth, may each of us know where he is going, be willing to make the continuous effort required to get there, avoid any detour, and be ready to pay the often very high price of faith and determination to win life’s race. Then, as mortality ends, we shall hear the plaudit from our Eternal Judge, “Well done, thou good and faithful servant: thou hast been faithful over a few things, I will make thee ruler over many things: enter thou into the joy of thy lord.” (Matt. 25:21.)

Each will then have completed his journey, not to a nebulous “ANYWHERE,” but to his heavenly home—even eternal life in the celestial kingdom of God.

May such be our goal and our reward is my prayer, in the name of Jesus Christ. Amen.

\newpage
\settalk{Your Jericho Road}
\setspeaker{Thomas S. Monson}
\settalkdate{April 1977}
\talkheading{Your Jericho Road}{Thomas S. Monson}

My dear brothers and sisters, I seek the help of our Heavenly Father as I respond to the invitation to speak to you today. A great number of you have journeyed many miles to attend this conference. From the north, the south, the east, and the west you have traveled the roads to Salt Lake City.

The word road is most intriguing. A generation ago movie moguls featured Bob Hope, Bing Crosby, and Dorothy Lamour in films entitled The Road to Rio, The Road to Morocco, and The Road to Zanzibar. Earlier yet, Rudyard Kipling immortalized another road when he penned the lines of “On the Road to Mandalay.”

This afternoon my thoughts have returned to a road made famous by a parable Jesus told. I speak of the road to Jericho. May I open the Bible to the Gospel of St. Luke, that we might together relive the memorable event which made famous for all time the Jericho Road.

A certain lawyer stood and tempted the Master, saying, “What shall I do to inherit eternal life?

“He said unto him, What is written in the law? how readest thou?

“And he answering said, Thou shalt love the Lord thy God with all thy heart, and with all thy soul, and with all thy strength, and with all thy mind; and thy neighbour as thyself.

“And he said unto him, Thou hast answered right: this do, and thou shalt live.

“But he, willing to justify himself, said unto Jesus, And who is my neighbour?

“And Jesus answering said, A certain man went down from Jerusalem to Jericho, and fell among thieves, which stripped him of his raiment, and wounded him, and departed, leaving him half dead.

“And by chance there came down a certain priest that way: and when he saw him, he passed by on the other side.

“And likewise a Levite, when he was at the place, came and looked on him, and passed by on the other side.

“But a certain Samaritan, as he journeyed, came where he was: and when he saw him, he had compassion on him,

“And went to him, and bound up his wounds, pouring in oil and wine, and set him on his own beast, and brought him to an inn, and took care of him.

“And on the morrow when he departed, he took out two pence, and gave them to the host, and said unto him, Take care of him; and whatsoever thou spendest more, when I come again, I will repay thee.

“Which now of these three, thinkest thou, was neighbour unto him that fell among the thieves?

“And he said, He that shewed mercy on him. Then said Jesus unto him, Go, and do thou likewise.” (Luke 10:25–37.)

Each of us, in the journey through mortality, will travel his own Jericho Road. What will be your experience? What will be mine? Will I fail to notice him who has fallen among thieves and requires my help? Will you?

Will I be one who sees the injured and hears his plea, yet crosses to the other side? Will you?

Or will I be one who sees, who hears, who pauses, and who helps? Will you?

Jesus provided our watchword, “Go, and do thou likewise.” When we obey that declaration, there opens to our eternal view a vista of joy seldom equaled and never surpassed.

Now the Jericho Road may not be clearly marked. Neither may the injured cry out, that we may hear. But when we walk in the steps of that good Samaritan, we walk the pathway that leads to perfection.

Note the many examples provided by the Master: the crippled man at the pool of Bethesda; the woman taken in adultery; the woman at Jacob’s well; the daughter of Jairus; Lazarus, brother of Mary and Martha—each represented a casualty on the Jericho Road. Each needed help.

To the cripple at Bethesda, Jesus said, “Rise, take up thy bed, and walk.” (John 5:8.) To the sinful woman came the counsel, “Go, and sin no more.” (John 8:11.) To her who came to draw water, He provided a well of water springing up into everlasting life. To the dead daughter of Jairus came the command, “Damsel, I say unto thee, arise.” (Mark 5:41.) To the entombed Lazarus, the memorable words, “Lazarus, come forth.” (John 11:43.)

One may well ask the penetrating question, “These accounts pertained to the Redeemer of the world. Can there actually occur in my own life, on my Jericho Road, such a treasured experience?”

My answer is a resounding “yes.” Let me share with you two such examples—first, the account of one who was injured and was helped; second, the learning experience of one who traveled the Jericho Road.

Some years ago there went to his eternal reward one of the kindest and most loved men to grace the earth. I speak of Louis C. Jacobsen. He ministered to those in need, he assisted the immigrant to find employment, and he delivered more sermons at more funeral services than any other I have known.

One day while in a reflective mood, Louis Jacobsen told me of his boyhood. He was the son of a poor Danish widow. He was small in stature, not comely in appearance—easily the object of his classmates’ thoughtless jokes. In Sunday School one Sabbath morning, the children made light of his patched trousers and his worn shirt. Too proud to cry, tiny Louis fled from the chapel, stopping at last, out of breath, to sit and rest on the curb which ran along Second West in Salt Lake City. Clear water flowed along the gutter next to the curb where Louis sat. From his pocket he took a piece of paper which contained the outlined Sunday School lesson and skillfully shaped a paper boat, which he launched on the flowing water. From his hurt boyish heart came the determined words, “I’ll never go back.”

Suddenly, through his tears Louis saw reflected in the water the image of a large and well-dressed man. Louis turned his face upward and recognized George Burbidge, the Sunday School superintendent.

“May I sit down with you?” asked the kind leader.

Louis nodded affirmatively. There on the gutter’s curb sat a good Samaritan ministering to one who surely was in need. Several boats were formed and launched while the conversation continued. At last the leader stood and, with a boy’s hand tightly clutching his, they returned to Sunday School.

Later Louis himself presided over that same Sunday School. Throughout his long life of service, he never failed to acknowledge the traveler who rescued him along a Jericho Road.

When I first learned of that far-reaching experience, I reflected on the words:

\begin{verse}
He stood at the crossroads all alone,\\
The sunlight in his face.\\
He had no thought for the world unknown—\\
He was set for a manly race.\\
But the roads stretched east and the roads stretched west,\\
And the lad knew not which road was best.\\
So he chose the road that led him down,\\
And he lost the race and the victor’s crown.\\
He was caught at last in an angry snare\\
Because no one stood at the crossroads there\\
To show him the better road.
\end{verse}

May I relate to you my first journey along a personal Jericho Road. In about my tenth year, as Christmas approached, I yearned as only a boy can yearn for an electric train. My desire was not to receive the economical and everywhere-to-be-found wind-up model train, but rather one that operated through the miracle of electricity.

The times were those of economic depression, yet Mother and Dad, through some sacrifice, I am sure, presented to me on Christmas morning a beautiful electric train. For hours I operated the transformer, watching the engine first pull its cars forward, then push them backward around the track.

Mother entered the living room and said to me that she had purchased a wind-up train for Widow Hansen’s boy, Mark, who lived down the lane. I asked if I could see the train. The engine was short and blocky—not long and sleek like the expensive model I had received.

However, I did take notice of an oil tanker car which was part of his inexpensive set. My train had no such car, and pangs of envy began to be felt. I put up such a fuss that Mother succumbed to my pleadings and handed me the oil tanker car. She said, “If you need it more than Mark, you take it.” I put it with my train set and felt pleased with the result.

Mother and I took the remaining cars and the engine down to Mark Hansen. The young boy was a year or two older than I. He had never anticipated such a gift and was thrilled beyond words. He wound the key in his engine, it not being electric like mine, and was overjoyed as the engine and two cars, plus a caboose, went around the track.

Mother wisely asked, “What do you think of Mark’s train, Tommy?”

I felt a keen sense of guilt and became very much aware of my selfishness. I said to Mother, “Wait just a moment—I’ll be right back.”

As swiftly as my legs could carry me, I ran to our home, picked up the oil tanker car plus an additional car of my own, ran back down the lane to the Hansen home, and said joyfully to Mark, “We forgot to bring two cars which belong to your train.”

Mark coupled the two extra cars to his set. I watched the engine make its labored way around the track and felt a supreme joy difficult to describe and impossible to forget.

Mother and I left the Hansen home and slowly walked up the street. She, who with her hand in God’s had entered into the valley of the shadow of death to bring me, her son, across the bridge of life, now took me by the hand and together we returned homeward by way of our private Jericho Road.

Some remember mother for her rhymes recited, others for her music played, songs sung, favors bestowed, or stories told; but I remember best that day we together traveled our Jericho Road and, like the good Samaritan, found a cherished opportunity to help.

My brothers and sisters, today there are hearts to gladden, there are deeds to be done—even precious souls to save. The sick, the weary, the hungry, the cold, the injured, the lonely, the aged, the wanderer—all cry out for our help.

The road signs of life enticingly invite every traveler: This way to fame; this way to affluence; this way to popularity; this way to luxury. Pause at the crossroads before you continue your journey. Listen for that still, small voice which ever so gently beckons, “Come, follow me. This way to Jericho.”

May each of us follow Him along that Jericho Road, I pray in the name of Jesus Christ. Amen.

\newpage
\settalk{The Way of the Lord}
\setspeaker{Thomas S. Monson}
\settalkdate{October 1977}
\talkheading{The Way of the Lord}{Thomas S. Monson}

Frequently we sing the hymn, “Come, listen to a prophet’s voice and hear the word of God.” (Hymns, no. 46.) Today we have listened to the voice of a prophet, Spencer W. Kimball, proclaim the word of God.

Humbly and prayerfully I seek Divine help as I speak to you from the crossroads of the West. Salt Lake City is a mecca for tourists from all parts of the globe. Thousands throng to the beautiful ski slopes of Alta, Brighton, Park City, and Snowbird each winter. Each summer the canyons of Bryce and Zion’s host thousands more. An attraction for all seasons is Temple Square, with its historic Tabernacle, lofty, spired temple, and the beautiful Visitors Center which bids to one and all a friendly welcome.

Situated somewhat off the beaten path, away from the crowd, is yet another famous square. Here in a quiet fashion, motivated by a Christlike love, elderly and handicapped workers serve one another after the divine plan of the Master. I speak of Welfare Square, sometimes known as the Bishops Storehouse. At this central location and at numerous other sites throughout the world, fruits and vegetables are canned, commodities processed, labeled, stored, and distributed to those persons who are in need. There is no sign of government dole nor the exchange of currency here, since only the signed order from an ordained bishop is honored.

Journalists marvel at this unique welfare plan and write glowingly of a people who take justifiable pride in the independence of caring for their own. Most frequently the curious and pleasantly surprised visitor asks three fundamental questions: (1) How does this plan operate? (2) How is it financed? (3) What prompts such devotion on the part of every worker?

Over the years it has been my pleasant opportunity to supply many with the answers to these sincerely asked questions. To the question “How does this plan operate?” I usually respond by mentioning that I had the privilege during the period 1950 through 1955 to preside as a bishop over 1,000 members, situated in the central part of Salt Lake City. In the congregation were eighty-six widows and perhaps forty families who were judged to be in need, at varying times and to some extent, of welfare assistance. Each year, I, along with the thousands of other bishops, would prepare a commodity requirement budget estimating the needs of our people for the coming year. All such budgets were carefully reviewed and compiled and specific assignments given to units of the Church, that the requirements of the needy might be met.

In one ecclesiastical unit the Church members would produce beef, in another oranges, in another vegetables or wheat—even a variety of staples, that the storehouses might be filled and the elderly and needy supplied. The Lord provided the way when he declared, “And the storehouse shall be kept by the consecrations of the church; and widows and orphans shall be provided for, as also the poor.” (D\&C 83:6.) Then the reminder, “But it must needs be done in mine own way.” (D\&C 104:16.)

In the vicinity where I lived and served, we operated a poultry project. Most of the time it was an efficiently operated project supplying to the storehouse thousands of dozens of fresh eggs and hundreds of pounds of dressed poultry. On a few occasions, however, the experience of being volunteer city farmers provided not only blisters on the hands, but frustration of heart and mind. For instance, I shall ever remember the time we gathered together the teenaged Aaronic Priesthood young men to really give the poultry project a spring cleaning treatment. Our enthusiastic and energetic throng gathered at the project, and in a speedy fashion uprooted, gathered, and burned large quantities of weeds and debris. By the light of the glowing bonfires we ate hot dogs and congratulated ourselves on a job well done. The project was now neat and tidy. However, there was just one disastrous problem. The noise and the fires had so disturbed the fragile and temperamental population of 5,000 laying hens that most of them went into a sudden moult and ceased laying. Thereafter we tolerated a few weeds, that we might produce more eggs.

No member of The Church of Jesus Christ of Latter-day Saints who has canned peas, topped beets, hauled hay, or shoveled coal in such a cause ever forgets or regrets the experience of helping provide for those in need. Devoted men and women help to operate this vast and inspired program. In reality, the plan would never succeed on effort alone, for this program operates through faith after the way of the Lord.

Sharing with others that which we have is not new to our generation. We need but to turn to the account found in First Kings of the Holy Bible to appreciate anew the principle that when we follow the counsel of the Lord, when we care for those in need, the outcome benefits all. There we read that a most severe drought had gripped the land. Famine followed. Elijah the prophet received from the Lord what to him must have been an amazing instruction: “Get thee to Zarephath: … behold, I have commanded a widow woman there to sustain thee.” When he had found the widow, Elijah declared, “Fetch me, I pray thee, a little water in a vessel, that I may drink.

“And as she was going to fetch it, he called to her, and said, Bring me, I pray thee, a morsel of bread in thine hand.”

Her response described her pathetic situation as she explained that she was preparing a final and scanty meal for her son and for herself, and then they would die.

How implausible to her must have been Elijah’s response: “Fear not; go and do as thou hast said: but make me thereof a little cake first, and bring it unto me, and after make for thee and for thy son.

“For thus saith the Lord God of Israel, The barrel of meal shall not waste, neither shall the cruse of oil fail, until the day that the Lord sendeth rain upon the earth.

“And she went and did according to the saying of Elijah: and she, and he, and her house, did eat many days.

“And the barrel of meal wasted not, neither did the cruse of oil fail.” (1 Kgs. 17:9–11, 13–16.) This is the faith that has ever motivated and inspired the welfare plan of the Lord.

In response to the second question, “How is your welfare plan financed?” one needs but to describe the fast offering principle. The prophet Isaiah described the true fast by asking, “Is it not to deal thy bread to the hungry, and that thou bring the poor that are cast out to thy house? when thou seest the naked, that thou cover him; and that thou hide not thyself from thine own flesh?

“Then shall thy light break forth as the morning, and thine health shall spring forth speedily: and thy righteousness shall go before thee; the glory of the Lord shall be thy rereward.

“Then shalt thou call, and the Lord shall answer; thou shalt cry, and he shall say, Here I am. …

“And the Lord shall guide thee continually, and satisfy thy soul in drought: … and thou shalt be like a watered garden, and like a spring of water, whose waters fail not.” (Isa. 58:7–9, 11.)

Guided by this principle, in a plan outlined and taught by inspired prophets of God, Latter-day Saints fast one day each month and contribute generously to a fast offering fund at least the equivalent of the meals forfeited and usually many times more. Such sacred offerings finance the operation of storehouses, supply cash needs of the poor, and provide medical care for the sick who are without funds.

In many areas, the offerings are collected each month by the boys who are deacons as they visit each member’s home generally quite early on the Sabbath day. I recall that the boys in the congregation over which I presided had assembled one morning, sleepy-eyed, a bit disheveled, and mildly complaining about arising so early to fulfill their assignment. Not a word of reproof was spoken, but during the following week, we escorted the boys to Welfare Square for a guided tour. They saw firsthand a lame person operating the telephone switchboard, an older man stocking shelves, women arranging clothing to be distributed—even a blind person placing labels on cans. Here were individuals earning their sustenance through their contributed labors. A penetrating silence came over the boys as they witnessed how their efforts each month helped to collect the sacred fast offering funds which aided the needy and provided employment for those who otherwise would be idle.

From that hallowed day forward, there was no urging required by our deacons. On fast Sunday mornings they were present at 7:00, dressed in their Sunday best, anxious to do their duty as holders of the Aaronic Priesthood. No longer were they simply distributing and collecting envelopes. They were helping to provide food for the hungry and shelter for the homeless—all after the way of the Lord. Their smiles were more frequent, their pace more eager, their very souls more subdued. Perhaps now they were marching to the beat of a different drummer; perhaps now they better understood the classic passage, “Inasmuch as ye have done it unto one of the least of these my brethren, ye have done it unto me.” (Matt. 25:40.)

To the third and final question, “What prompts such devotion on the part of every worker?” the answer can be stated simply: An individual testimony of the gospel of the Lord Jesus Christ, even a heartfelt desire to love the Lord with all one’s heart, mind, and soul, and one’s neighbor as oneself.

This is what motivated a personal friend, now deceased, who was in the produce business, to telephone me during those days as a bishop and say, “I’m sending to the storehouse a semi-truck and trailer filled with citrus fruits for those who would otherwise go without. Let the storehouse management know the truck is coming, and there will be no charge; but Bishop, no one is to know who sent it.” Rarely have I seen the joy and appreciation this generous act brought forth. Never have I questioned the eternal reward to which that unnamed benefactor has now gone.

Such kind deeds of generosity are not a rarity, but are frequently found. Situated beneath the heavily traveled freeway which girds Salt Lake City is the home of a sixty-year-old single man who has, due to a crippling disease, never known a day without pain nor many days without loneliness. One winter’s day as I visited him, he was slow in answering the doorbell’s ring. I entered his well-kept home; the temperature in save but one room, the kitchen, was a chilly 40 degrees. The reason: not sufficient money to heat any other room. The walls needed papering, the ceilings to be lowered, the cupboards filled.

Troubled by the experience of visiting my friend, a bishop was consulted and a miracle of love, prompted by testimony, took place. The ward members were organized and the labor of love begun. A month later, my friend Lou called and asked if I would come and see what had happened to him. I did, and indeed beheld a miracle. The sidewalks which had been uprooted by large poplar trees had been replaced, the porch of the home rebuilt, a new door with glistening hardware installed, the ceilings lowered, the walls papered, the woodwork painted, the roof replaced, and the cupboards filled. No longer was the home chilly and uninviting. It now seemed to whisper a warm welcome. Lou saved until last showing me his pride and joy: there on his bed was a beautiful plaid quilt bearing the crest of his McDonald family clan. It had been made with loving care by the women of the Relief Society. Before leaving, I discovered that each week the Young Adults would bring in a hot dinner and share a home evening. Warmth had replaced the cold; repairs had transformed the wear of years; but more significantly, hope had dispelled despair and now love reigned triumphant.

All who participated in this moving drama of real life had discovered a new and personal appreciation of the Master’s teaching, “It is more blessed to give than to receive.” (Acts 20:35.)

To all within the sound of my voice I declare that the welfare plan of The Church of Jesus Christ of Latter-day Saints is inspired of Almighty God. Indeed, the Lord Jesus Christ is its Architect. To you I extend a heartfelt and sincere invitation: Come to Salt Lake City and visit Welfare Square. Your eyes will glow a little brighter, your heart will beat a little faster, and life itself will acquire a new depth of meaning. May such be your experience, I pray, in the name of Jesus Christ. Amen.

\newpage
\settalk{The Prayer of Faith}
\setspeaker{Thomas S. Monson}
\settalkdate{April 1978}
\talkheading{The Prayer of Faith}{Thomas S. Monson}

Our hearts are touched by the beautiful singing of these precious Primary boys and girls. All of the children participating here this afternoon enjoy the privilege of associating once each week with others of similar age in the meetings of the Primary. There are, however, other children, equally as sweet and precious, who are not so fortunate.

Some years ago while visiting the Australia Mission, I accompanied the mission president on a flight to Darwin to break ground for that city’s first Latter-day Saint chapel. We stopped for refueling at the small mining community of Mt. Isa. There we were met at the terminal by a mother and her two children of Primary age. She introduced herself as Judith Louden and mentioned that she and her two children were the only members of the Church in the town. Her husband, Richard, was not a member. We held a brief meeting, where I discussed the importance of holding a home Primary session each week. I promised to send from Church headquarters the home Primary materials to assist them. There was a commitment to pray, to meet, to persevere in faith.

Upon returning to Salt Lake City, I enlisted the help of then-President LaVern Parmley, and the home Primary materials were sent, along with a subscription to the Children’s Friend.

Years later, while attending the stake conference of the Brisbane Australia Stake, I happened to mention in a priesthood session the plight of this faithful woman and her children. I said, “Someday I hope to learn if that home Primary succeeded and meet the nonmember husband and father of that choice family.” One of the brethren in the meeting stood and said, “Brother Monson, I am Richard Louden, the husband of that good woman and the father of those precious children. Prayer and Primary brought me into the Church.”

The power of prayer again came to mind this past winter. I was on assignment many thousands of miles to the south in the beautiful city of Buenos Aires, Argentina. I paused by the historic Palermo Park, which graces the downtown area, and realized that this was sacred ground, for here on Christmas Day in 1925 Elder Melvin J. Ballard, an apostle of the Lord, dedicated all of South America for the preaching of the gospel. What a fulfillment to an inspired prayer is evident today as the growth of the Church in that land exceeds all expectations.

In that same park is a large statue of George Washington, the father of the United States and its first president. As I observed the statue, my thoughts returned to another historic place where prayer played a vital role—even Valley Forge. It was at Valley Forge that this same Washington led his badly battered, ill-fed, and scantily clad troops to winter quarters.

Today, in a quiet grove at Valley Forge, there is an heroic-sized monument to Washington. He is depicted not astride a charging horse nor overlooking a battlefield of glory, but kneeling in humble prayer, calling upon the God of Heaven for divine help. To gaze upon the statue prompts the mind to remember the oft-heard expression, “A man never stands taller than when upon his knees.”

Men and women of integrity, character, and purpose have ever recognized a power higher than themselves and have sought through prayer to be guided by that power. Such has it ever been. So shall it ever be.

In the very beginning, Father Adam was commanded, “Call upon God in the name of the Son forevermore.” (Moses 5:8.) Adam prayed. Abraham prayed. Isaac prayed. Moses prayed, and so did every prophet pray to that God from whence came his strength. Like the sands slipping through an hourglass, generations of mankind were born, lived, and then died. At long last came that glorious event for which prophets prayed, psalmists sang, martyrs died, and all mankind hoped.

The birth of the babe in Bethlehem was transcendent in its beauty and singular in its significance. Jesus of Nazareth brought prophecy to fulfillment. He cleansed lepers, He restored sight, He opened ears, He renewed life, He taught truth, He saved all. In so doing, He honored His Father and provided you and me with an example worthy of emulation.

More than any prophet or leader, He showed us how to pray. Who can fail to remember His agony in Gethsemane and that fervent prayer: “O my Father, if it be possible, let this cup pass from me: nevertheless not as I will, but as thou wilt.” (Matt. 26:39.) And His injunction: “Watch and pray, that ye enter not into temptation.” (Matt. 26:41.)

We remember His counsel: “When thou prayest, thou shalt not be as the hypocrites are: for they love to pray standing in the synagogues and in the corners of the streets, that they may be seen of men. …

“But thou, when thou prayest, … pray to thy Father which is in secret; and thy Father which seeth in secret shall reward thee openly.” (Matt. 6:5–6.)

This guiding instruction has helped troubled souls discover the peace for which they fervently yearn and earnestly hope.

Unfortunately, prosperity, abundance, honor, and praise lead some men to the false security of haughty self-assurance and the abandonment of the inclination to pray. Conversely, trial, tribulation, sickness, and death crumble the castles of men’s pride and bring them to their knees to petition for power from on high.

I suppose that during the holocaust of World War II more of the people living on this earth paused to pray than at any other time in our history. Who can calculate the concern of mothers, wives, and children who pleaded for Almighty God’s protecting care to be with absent sons, husbands, and fathers locked in mortal combat? Prayers are heard. Prayers are answered.

Heartwarming is the example of the mother in America who prayed for her son’s well-being as the ship on which he served sailed into the bloody cauldron known as the Pacific theater of war. Each morning she would arise from kneeling in prayer and serve as a volunteer on those production lines which became lifelines to men in battle. Could it be that a mother’s own handiwork might somehow directly affect the life of a loved one? All who knew her and her family cherished the actual account of her sailor son, Elgin Staples, whose ship went down off Guadalcanal. Staples was swept over the side; but he survived, thanks to a life belt that proved, on later examination, to have been inspected, packed, and stamped back home in Akron, Ohio, by his own mother!

\begin{verse}
I know not by what method rare,\\
But this I know, God answers prayer.\\
I know that He has given His Word,\\
Which tells me prayer is always heard,\\
And will be answered, soon or late.\\
And so I pray and calmly wait.
\end{verse}

(Eliza M. Hickock, “Prayer,” in The Best Loved Religious Poems, ed. James Gilchrist Lawson, New York: Fleming H. Revell Co., 1933, p. 160.)

Well might the younger generation ask the question: “But what about today? Does He still hear? Does He continue to answer?” To which I promptly reply: “There is no expiration date on the Lord’s injunction to pray. As we remember Him, He will remember us.”

Most of the time there are no flags waving nor bands playing when prayer is answered. His miracles frequently are performed in a quiet and natural manner.

Some years ago while I was attending the Grand Junction Colorado Stake conference, the stake president asked if I would meet with a grieving mother and father whose son had announced his decision to leave his mission field after having just arrived there. When the conference throng had left, we knelt quietly in a private place—mother, father, stake president, and I. As I prayed in behalf of all, I could hear the muffled sobs of a sorrowing mother and disappointed father.

When we arose, the father said, “Brother Monson, do you really think our Heavenly Father can alter our son’s announced decision to return home before completing his mission? Why is it that now, when I am trying so hard to do what is right, my prayers are not heard?”

I responded, “Where is your son serving?”

He replied, “In Duesseldorf, Germany.”

I placed my arms around that mother and father and said to them, “Your prayers have been heard and are already being answered. With more than twenty-eight stake conferences being held this day attended by the General Authorities, I was assigned to your stake. Of all the Brethren, I am the only one who has the assignment to meet with the missionaries in the Duesseldorf Germany Mission this very Thursday.”

Their petition had been honored by the Lord. I was able to meet with their son. He responded to their pleadings. He remained and completed a highly successful mission.

Some years later I again visited the Grand Junction Colorado Stake. Again I met the same parents. Still the father had not qualified to have his large and beautiful family join mother and father in a sacred sealing ceremony, that this family might be a forever family. I suggested that if the family members would earnestly pray, they could qualify. I indicated that I would be pleased to officiate on that sacred occasion in the temple of God.

Mother pleaded, father strived, children urged, all prayed. The result? Let me share with you a treasured letter that their young son, Todd, placed under Daddy’s pillow on Father’s Day morning.

“Dad,

“I love you for what you are and not for what you aren’t. Why don’t you stop smoking? Millions of people have … why can’t you? It’s harmful to your health, to your lungs, your heart. If you can’t keep the Word of Wisdom you can’t go to heaven with me, Skip, Brad, Marc, Jeff, Jeannie, Pam, and their families. Us kids keep the Word of Wisdom. Why can’t you? You are stronger and you are a man. Dad, I want to see you in heaven. We all do. We want to be a whole family in heaven … not half of one.

“Dad, you and Mom ought to get two old bikes and start riding around the park every night. You are probably laughing right now, but I wouldn’t be. You laugh at those old people, jogging around the park and riding bikes and walking, but they are going to outlive you. Because they are exercising their lungs, their hearts, their muscles. They are going to have the last laugh.

“Come on, Dad, be a good guy—don’t smoke, drink, or anything else against our religion. We want you at our graduation. If you do quit smoking and do good stuff like us, you and Mom can go with Brother Monson and get married and sealed to us in the temple.

“Come on, Dad—Mom and us kids are just waiting for you. We want to live with you forever. We love you. You’re the greatest, Dad.

Love,

Todd

“P.S. And if the rest of us wrote one of these, they’d say the same thing.

“P.P.S. Mr. Newton has quit smoking. So can you. You are closer to God than Mr. Newton!”

That plea, that prayer of faith, was heard and answered. A night I shall ever treasure and long remember was when this entire family assembled in a sacred room in the beautiful temple which graces this square. Father was there. Mother was there. Every child was there. Ordinances eternal in their significance were performed. A humble prayer of gratitude brought to a close this long-awaited evening.

May we ever remember …

\begin{verse}
Prayer is the soul’s sincere desire,\\
Uttered or unexpressed,\\
The motion of a hidden fire\\
That trembles in the breast.
\end{verse}

(Hymns, no. 220.)

He has taught us how to pray. That each of us will learn and live this lesson is my earnest plea and sincere prayer, in the name of Jesus Christ. Amen.

\newpage
\settalk{Profiles of Faith}
\setspeaker{Thomas S. Monson}
\settalkdate{October 1978}
\talkheading{Profiles of Faith}{Thomas S. Monson}

Who can help but be uplifted and inspired just to worship in this historic tabernacle and to listen to this glorious choir?

It has been said that “when Evan Stephens was conductor of the Tabernacle Choir, he was thrilled on one occasion by a sermon delivered by the late President Joseph F. Smith on the subject, [of “The Faith of Latter-day Saint Youth.”] At the close of the service Professor Stephens strolled alone up City Creek Canyon [to the north], pondering the inspired words of the President. Suddenly [the inspiration of heaven] came upon him and seated upon a rock which was standing firm under the pressure of the rushing water, he wrote with a pencil” these words:

\begin{verse}
Shall the youth of Zion falter\\
In defending truth and right?\\
While the enemy assaileth,\\
Shall [they] shrink or shun the fight?\\
No!
\end{verse}

(Hymns, no. 157; J. Spencer Cornwall, Stories of Our Mormon Hymns, Salt Lake City: Deseret Book, 1963, p. 173.)

In that early day, I am confident that youth were faced with difficult challenges to meet and vexing problems to solve. Youth is not a time of ease nor of freedom from perplexing questions. It wasn’t then, and it surely isn’t today. In fact, as time passes it seems that the difficulties of youth increase in size and scope. Temptation continues to loom large on life’s horizon. Accounts of violence, theft, drug abuse, and pornography blare forth from the television screen and peer constantly from most daily newspapers. Such examples blur our vision and fault our thinking. Soon assumptions become generally accepted opinions, and all youth everywhere are categorized as “not so good as yesteryear,” or “the worst generation yet.” How wrong are such opinions! How incorrect are such statements!

True, today is a new day with new trials, new troubles, and new temptations, but hundreds of thousands of Latter-day Saint youth strive constantly and serve diligently, true to the faith, as their counterparts of earlier years so nobly did. Because the contrast between good and evil is so stark, the exceptions to the prevailing trends are magnified, observed, and appreciated by decent persons throughout the world.

Let me share with you a pointed letter which came from a resident of Minnesota. It was addressed to Brigham Young University:

“Gentlemen:

“Beginning December 22, I made a bus trip from southern Minnesota to Florida via Des Moines and Chicago and points south.

“There was a large group of young men and women traveling the approximately same route from Des Moines. These fine young people were students from Brigham Young [University] going home for the holidays.

“They were all very polite, well-behaved, articulate young men and women. It was a pleasure to travel with them—to know them—and it gave me a new hope for the future.

“I realized that the university cannot do this. Young men and women of their caliber are the product of fine homes. The credit is due the parents. I cannot reach the parents, so my appreciation must go to the school.”

Such comments are not isolated, but rather typical, for which we are ever pleased. Our Latter-day Saint students are excellent examples of faith in action.

Another group which amazes the world and inspires faith is that army of Latter-day Saint missionaries, now more than 26,600 strong, currently serving throughout the world. All through their lives, these young men and women have prepared for and awaited that special day when a mission call is received. Fathers become justifiably proud and mothers somewhat anxious. Well do I remember the recommendation form of one missionary on which the bishop had written:

“This is the most outstanding young man I have ever recommended. He has excelled in all aspects of his life. He was president of his Aaronic Priesthood quorums and an officer at his high school. He lettered in track and football. I have never recommended a more outstanding candidate. I am proud to be his father.”

More generally, the bishop and stake president write, “John is a fine young man. He has prepared for his mission physically, mentally, financially, and spiritually. He will serve gladly and with distinction wherever he is called.”

One day I was with President Spencer W. Kimball as he signed these special calls to full-time missionary service. Suddenly he noticed the call of his own grandson. He signed his name as president of the Church and then penned a personal line at the bottom which read, “I’m proud of you. Love, Grandpa.”

When the call is received, the college text is closed and the scriptures opened. Family, friends, and often a special friend are left behind. Suspended are dating, dancing, and driving, as the three Ds are exchanged for the three Ts—tracting, teaching, and testifying.

Let us examine specifically several missionary profiles of faith, that we might better consider the question “Shall the youth of Zion falter?”

For a first profile, I mention Jose Garcia from Old Mexico. Born in poverty but nurtured in faith, Jose prepared for a mission call. I was present the day his recommendation was received. There appeared the statement: “Brother Garcia will serve at great sacrifice to his family, for he is the means of much of the family support. He has but one possession—a treasured stamp collection—which he is willing to sell, if necessary, to help finance his mission.”

President Kimball listened attentively as this statement was read to him, and then he responded: “Have him sell his stamp collection. Such sacrifice will be to him a blessing.” Then, with a twinkle in his eye and a smile on his face, this loving prophet said, “Each month at Church headquarters we receive thousands of letters from all parts of the world. See that we save these stamps and provide them to Jose at the conclusion of his mission. He will have, without cost, the finest stamp collection of any young man in Mexico.”

There seemed to echo from another place, another time, the experience of the Master:

“And he looked up, and saw the rich men casting their gifts into the treasury.

“And he saw also a certain poor widow casting in thither two mites.

“And he said, Of a truth I say unto you, that this poor widow hath cast in more than they all.” (Luke 12:1–3.)

“For all they did cast in of their abundance; but she of her want did cast in all that she had, even all her living” (Mark 12:44).

For a second profile, I turn from Mexico to a missionary at the Missionary Training Center in Provo, Utah, desperately struggling to become proficient in the German language, that he might be an effective missionary to the people of southern Germany. Each day as he opened his German grammar text, he noticed with interest and curiosity that the front cover displayed a picture of a most quaint and ancient house in Rothenburg, West Germany. Beneath the picture, the location was given. In his heart that young man determined, “I’ll visit that house and teach the truth to whoever lives within it.” This he did. The result was the conversion and baptism of Sister Helma Hahn. Today she devotes much of her time speaking to tourists who come from all over the world to see her house. She delights in telling them of the blessings which the gospel of Jesus Christ has brought to her. Her house is perhaps one of the most frequently photographed houses in the entire world. No visitor leaves without hearing in simple yet earnest words her testimony of praise and gratitude. That missionary who brought to Sister Hahn the gospel remembered the sacred charge: “Go ye therefore, and teach all nations, baptizing them in the name of the Father, and of the Son, and of the Holy Ghost” (Matt. 28:19).

Profile number three also relates to a missionary of unfaltering faith, Elder Mark Skidmore. When he received his call to Norway, he knew not one word of Norwegian—yet he realized that to teach and to testify he would need proficiency in the language of the Norwegian people. To himself he made a private vow: “I will not speak English until I have brought into the waters of baptism my first Norwegian family.” He plodded. He prayed. He pleaded. He worked. After the trial of his faith came the desired blessing. He taught and baptized a choice family. He then spoke in English for the first time in six months. I met with him that same week. His expression was one of thanksgiving and gratitude. I thought of the words of Moroni, that courageous captain: “I seek not for power. … I seek not for honor of the world, but for the glory of my God.” (Alma 60:36.)

For a final profile, I mention the mother of one noble missionary son. The family lived in the harsh climate of Star Valley, Wyoming. Summer there is brief and warm, while winter is long and cold. When a fine son of nineteen said farewell to home and family, he knew on whom the burden of work would fall. Father was ill and limited. To mother came the task of milking by hand the small dairy herd which sustained the family.

While serving as a mission president, I attended a seminar for all presidents held in Salt Lake City. My wife and I were privileged to devote an evening to meeting the parents of those missionaries who served with us. Some parents were wealthy and handsomely attired. They spoke in a gracious manner. Their faith was strong. Others were less affluent, of modest means and rather shy. They, too, were proud of their special missionary and prayed and sacrificed for his welfare.

Of all the parents whom I met that evening, the best remembered was that mother from Star Valley. As she took my hand in hers I felt the large calluses which revealed the manual labor she daily performed. Almost apologetically, she attempted to excuse her rough hands, her wind-whipped face. She whispered, “Tell our son Spencer that we love him, that we’re proud of him, and that we pray daily for him.”

Until that night I had never seen an angel nor heard an angel speak. I never again could make that statement, for that angel mother carried with her the Spirit of Christ. She, who with that same hand clasped in the hand of God had walked bravely into the valley of the shadow of death to bring to this mortal life her son, had indelibly impressed my life.

Nurtured and guided by such noble mothers, missionaries match the description of Helaman’s throng:

“And they were all young men, and they were exceedingly valiant for courage, and also for strength and activity; but behold, this was not all—they were men who were true at all times in whatsoever thing they were entrusted.

“Yea, they were men of truth and soberness, for they had been taught to keep the commandments of God and to walk uprightly before him.” (Alma 53:20–21.)

Such profiles prompt faith. They instill confidence. They teach truth. They testify of goodness. They help provide the answer to that question:

\begin{verse}
Shall the youth of Zion falter\\
In defending truth and right?\\
While the enemy assaileth,\\
Shall [they] shrink or shun the fight?\\
No!
\end{verse}

My sincere prayer is that we will stand with the youth of Zion, remain true to the faith, for which I pray in the name of Jesus Christ. Amen.

\newpage
\settalk{The Army of the Lord}
\setspeaker{Thomas S. Monson}
\settalkdate{April 1979}
\talkheading{The Army of the Lord}{Thomas S. Monson}

Tonight I am aware that you, my brethren, represent the largest gathering of the priesthood ever to assemble. I pray for the help of our Heavenly Father, that he may grant me inspiration coupled with courage.

Some twenty-four years ago I was seated in the choir seats of the Assembly Hall situated to the south of us here on Temple Square. The setting was stake conference. Elder Joseph Fielding Smith and Elder Alma Sonne had been assigned to reorganize our stake presidency. The Aaronic Priesthood, including members of bishoprics, were providing the music for the conference. Those of us who served as bishops were singing along with our young men. As we concluded singing our first selection, Brother Smith stepped to the pulpit and announced the names of the new stake presidency. I am confident the other members of the presidency had been made aware of their callings, but I had not. After reading my name, Brother Smith announced, “If Brother Monson is willing to respond to this call we shall be pleased to hear from him now.” As I stood at the pulpit and gazed out on that sea of faces, I remembered the song we had just sung. Its title was “Have Courage, My Boy, to Say No.” I selected as my acceptance theme “Have Courage, My Boy, to Say Yes.” Such is the courage I seek this evening.

The words of a better-known hymn describe you:

\begin{verse}
Behold! a royal army,\\
With banner, sword and shield,\\
Is marching forth to conquer,\\
On life’s great battlefield;\\
Its ranks are filled with soldiers,\\
United, bold and strong,\\
Who follow their Commander,\\
And sing their joyful song:
\end{verse}

(Hymns, no. 7.)

The priesthood represents a mighty army of righteousness—even a royal army. We are led by a prophet of God. In supreme command is our Lord and Savior, Jesus Christ. Our marching orders are clear. They are concise. Matthew describes our challenge in these words from the Master: “Go ye therefore, and teach all nations, baptizing them in the name of the Father, and of the Son, and of the Holy Ghost:

“Teaching them to observe all things whatsoever I have commanded you: and lo, I am with you alway, even unto the end of the world.” (Matt. 28:19–20.) Did those early disciples listen to this divine command? Mark records, “And they went forth, and preached every where, the Lord working with them” (Mark 16:20).

The command to go has not been rescinded. Rather, it has been reemphasized. Today twenty-eight thousand missionaries are serving in response to the call. Additional thousands will soon respond. Nine new missions will be created in July, bringing the total number of missions to 175. What a thrilling and challenging time in which to live!

You who hold the Aaronic Priesthood and honor it have been reserved for this special period in history. The harvest truly is great. Let there be no mistake about it; the opportunity of a lifetime is yours. The blessings of eternity await you. How might you best respond? May I suggest you cultivate three virtues, namely—

By so doing, you will ever be found part of that royal army of the Lord. Let us consider, individually, each of these three virtues.

First, a desire to serve. Remember the qualifying statement of the Master, “Behold, the Lord requireth the heart and a willing mind” (D\&C 64:34). A latter-day minister advised: “Until willingness overflows obligation, men fight as conscripts rather than following the flag as patriots. Duty is never worthily performed until it is performed by one who would gladly do more if only he could.”

Isn’t it appropriate that you do not call yourselves to this work? Isn’t it wise that your parents do not call you? Rather, you are called of God by prophecy and by revelation. Your call bears the signature of the President of the Church.

It was my privilege to serve for many years with President Spencer W. Kimball when he was chairman of the Missionary Executive Committee of the Church. Those never-to-be-forgotten missionary assignment meetings were filled with inspiration and occasionally interspersed with humor. Well do I remember the recommendation form for one prospective missionary on which the bishop had written: “This young man is very close to his mother. She wonders if he might be assigned to a mission close to home in California so that she can visit him on occasion and telephone him weekly.” As I read aloud this comment, I awaited from President Kimball the pronouncement of a designated assignment. I noticed a twinkle in his eye and a sweet smile cross his lips as he said, without additional comment, “Assign him to the Johannesburg South Africa Mission.”

Too numerous to mention are the many instances where a particular call proved providential. This I know—divine inspiration attends such sacred assignments. We, with you, acknowledge the truth stated so simply in the Doctrine and Covenants: “If ye have desires to serve God ye are called to the work” (D\&C 4:3).

Second, the patience to prepare. Preparation for a mission is not a spur-of-the-moment matter. It began before you can remember. Every class in Primary, Sunday School, seminary—each priesthood assignment—had a larger application. Silently, almost imperceptibly, a life was molded, a career commenced, a man made. Said the poet:

\begin{verse}
Who touches a boy by the Master’s plan\\
Is shaping the course of a future man.
\end{verse}

What a challenge is the calling to be an adviser to a quorum of boys. Advisers, do you really think about your opportunity? Do you ponder? Do you pray? Do you prepare? Do you prepare your boys?

As a boy of fifteen I was called to preside over a quorum of teachers. Our adviser was interested in us, and we knew it. One day he said to me, “Tom, you enjoy raising pigeons, don’t you?”

I responded with a warm “Yes.”

Then he proffered, “How would you like me to give you a pair of purebred Birmingham Roller pigeons?”

This time I answered, “Yes, sir!” You see, the pigeons I had were just the common variety trapped on the roof of the Grant Elementary School.

He invited me to come to his home the next evening. The next day was one of the longest in my young life. I was awaiting my adviser’s return from work an hour before he arrived. He took me to his loft, which was in a small barn at the rear of his yard. As I looked at the most beautiful pigeons I had yet seen, he said, “Select any male, and I will give you a female which is different from any other pigeon in the world.” I made my selection. He then placed in my hand a tiny hen. I asked what made her so different. He responded, “Look carefully, and you’ll notice that she has but one eye.” Sure enough, one eye was missing, a cat having done the damage. “Take them home to your loft,” he counseled. “Keep them in for about ten days and then turn them out to see if they will remain at your place.”

I followed his instructions. Upon releasing them, the male pigeon strutted about the roof of the loft, then returned inside to eat. But the one-eyed female was gone in an instant. I called Harold, my adviser, and asked: “Did that one-eyed pigeon return to your loft?”

“Come on over,” said he, “and we’ll have a look.”

As we walked from his kitchen door to the loft, my adviser commented, “Tom, you are the president of the teachers quorum.” This I already knew. Then he added, “What are you going to do to activate Bob?”

I answered, “I’ll have him at quorum meeting this week.”

Then he reached up to a special nest and handed to me the one-eyed pigeon. “Keep her in a few days and try again.” This I did, and once more she disappeared. Again the experience, “Come on over and we’ll see if she returned here.” Came the comment as we walked to the loft, “Congratulations on getting Bob to priesthood meeting. Now what are you and Bob going to do to activate Bill?”

“We’ll have him there this week,” I volunteered.

This experience was repeated over and over again. I was a grown man before I fully realized that, indeed, Harold, my adviser, had given me a special pigeon; the only bird in his loft he knew would return every time she was released. It was his inspired way of having an ideal personal priesthood interview with the teachers quorum president every two weeks. I owe a lot to that one-eyed pigeon. I owe more to that quorum adviser. He had the patience to help me prepare for opportunities which lay ahead.

Third, a willingness to labor. Missionary work is difficult. It will tax your energies. It will strain your capacity. It will demand your best effort—frequently, a second effort. Remember, the race goeth “not to the swift, nor the battle to the strong” (Eccl. 9:11)—but to him who endures to the end. Determine to—

\begin{verse}
Stick to your task till it sticks to you.\\
Beginners are many, but enders are few.\\
Honor, power, place and praise\\
Will always come to the one who stays.\\
Stick to your task till it sticks to you;\\
Bend at it, sweat at it, smile at it, too—\\
For out of the bend and the sweat and the smile\\
Will come life’s victories after awhile.
\end{verse}

During the final phases of World War II, I turned eighteen and was ordained an elder—one week before I departed for active duty with the Navy. A member of my ward bishopric was at the train station to bid me farewell. Just before train time, he placed in my hand a book which I hold before you tonight. Its title, the Missionary Handbook. I laughed and commented, “I’m not going on a mission.” He answered, “Take it anyway. It may come in handy.”

It did. During basic training our company commander instructed us concerning how we might best pack our clothing in a large sea bag. He advised, “If you have a hard, rectangular object you can place in the bottom of the bag, your clothes will stay more firm.” I suddenly remembered just the right rectangular object—the Missionary Handbook. Thus it served for twelve weeks.

The night preceding our Christmas leave our thoughts were, as always, on home. The barracks were quiet. Suddenly I became aware that my buddy in the adjoining bunk—a Mormon boy, Leland Merrill—was moaning with pain. I asked, “What’s the matter, Merrill?”

He replied, “I’m sick. I’m really sick.”

I advised him to go to the base dispensary, but he answered knowingly that such a course would prevent him from being home for Christmas.

The hours lengthened; his groans grew louder. Then, in desperation, he whispered, “Monson, Monson, aren’t you an elder?” I acknowledged this to be so; whereupon he asked, “Give me a blessing.”

I became very much aware that I had never given a blessing. I had never received such a blessing; I had never witnessed a blessing being given. My prayer to God was a plea for help. The answer came: “Look in the bottom of the sea bag.” Thus, at 2 a.m. I emptied on the deck the contents of the bag. I then took to the night-light that hard, rectangular object, the Missionary Handbook, and read how one blesses the sick. With about sixty curious sailors looking on, I proceeded with the blessing. Before I could stow my gear, Leland Merrill was sleeping like a child.

The next morning Merrill smilingly turned to me and said, “Monson, I’m glad you hold the priesthood.” His gladness was only surpassed by my gratitude.

Future missionaries, may our Heavenly Father bless you with a desire to serve, the patience to prepare, and a willingness to labor, that you and all who comprise this royal army of the Lord may merit his promise: “I will go before your face. I will be on your right hand and on your left, and my Spirit shall be in your hearts, and mine angels round about you, to bear you up.” (D\&C 84:88.)

This is my earnest and sincere prayer. I ask it humbly and in the name of Jesus Christ, amen.

\newpage
\settalk{Pornography—the Deadly Carrier}
\setspeaker{Thomas S. Monson}
\settalkdate{October 1979}
\talkheading{Pornography—the Deadly Carrier}{Thomas S. Monson}

This week, my brothers and sisters, the woodcutters are laying their massive axes and taking their power saws to the still stately and once mighty elm trees that grace the countryside surrounding London, England’s, Heathrow Airport.

It is said some of the majestic monarchs are over one hundred years old. One wonders how many persons have admired their beauty, how many picnics have been enjoyed in their welcome shade, how many generations of song birds have filled the air with music while capering among the outstretched and luxuriant branches.

The patriarchal elms are now dead. Their demise was not the result of old age, the recurring drought, nor the strong winds which occasionally lash the area. Their destroyer is much more harmless in appearance, yet deadly in result. We know the culprit as the bark beetle, carrier of the fatal Dutch elm disease. This malady has destroyed vast elm forests throughout Europe and America. Its march of death continues unabated. All efforts at control have failed.

Dutch elm disease usually begins with a wilting of the younger leaves in the upper part of the tree. Later the lower branches become infected. In about midsummer most of the leaves turn yellow, curl, and drop off. Life ebbs. Death approaches. A forest is consumed. The bark beetle has taken its terrible toll.

How like the elm is man. From a minute seed, and in accordance with a divine plan, we grow, are nurtured, and mature. The bright sunlight of heaven, the rich blessings of earth are ours. In our private forest of family and friends, life is richly rewarding and abundantly beautiful. Then suddenly, there appears before us in this generation a sinister and diabolical enemy—pornography. Like the bark beetle it too is the carrier of a deadly disease. I shall name it “pernicious permissiveness.”

At first we scarcely realize we have been infected. We laugh and make light-hearted comment concerning the off-color story or the clever cartoon. With evangelical zeal we protect the so-called rights of those who would contaminate with smut and destroy all that is precious and sacred. The beetle of pornography is doing his deadly task—undercutting our will, destroying our immunity, and stifling that upward reach within each of us.

Can this actually be true? Surely this matter of pernicious permissiveness is not so serious. What are the facts? Let’s look! Let’s listen! Then let’s act!

Pornography, the carrier, is big business. It is Mafia-spawned. It is contagious. It is addicting. In a study last year, the FBI estimated that Americans spent 2.4 billion dollars on hard-core pornography. Other estimates reach as high as 4 billion—a fortune siphoned away from noble use and diverted to a devilish purpose!

Apathy toward pornography stems mostly from a widespread public attitude that it is a victimless crime and that police resources are better used in other areas. Many state and local ordinances are ineffective, sentences are light, and the huge financial rewards far outweigh the risks.

The FBI points out that pornography may have a direct relationship to sex crimes. “In one large western city,” an agency report states, “the vice squad advised that 72 percent of the individuals arrested for rape and child-related sexual offenses had in their possession some type of pornographic material.”

Some publishers and printers prostitute their presses by printing millions of pieces of pornography each day. No expense is spared. The finest of paper, the spectrum of full color combine to produce a product certain to be read, then read again. Nor are the movie producer, the television programmer, or the entertainer free from taint. Gone are the restraints of yesteryear. So-called realism is the quest.

One of the leading box office stars of today lamented: “The boundaries of permissiveness have been extended to the limit. The last film I did was filthy. I thought it was filthy when I read the script, and I still think it’s filthy; but the studio tried it out at a Friday night sneak preview and the audience screamed its approval.”

Another star declared, “Movie makers, like publishers, are in the business to make money, and they make money by giving the public what it wants.”

Some persons struggle to differentiate between what they term “soft-core” and “hard-core” pornography. Actually, one leads to another. How applicable is Alexander Pope’s classic, “Essay on Man”:

\begin{verse}
Vice is a monster of so frightful mien\\
As to be hated needs but to be seen;\\
Yet seen too oft, familiar with her face,\\
We first endure, then pity, then embrace.
\end{verse}

(John Bartlett, Familiar Quotations, Boston: Little Brown and Co., 1968, p. 409).

The constant, consuming march of the pornography beetle blights neighborhoods just as it contaminates human lives. Some are particularly scarred by its insidious touch.

Come with me for a moment to a place portrayed in song—dear to the heart of America—New York City’s world-famous landmark of Broadway and Forty-fifth Street. There, standing so forlornly alone on a tiny island surrounded by bustling traffic, is a heroic-size statue of Father Francis P. Duffy, well-known chaplain of the Fighting Sixty-ninth of World War I fame. He wears the uniform of the battlefield. He carries a canteen to relieve the physical distress of the wounded and a Bible to bring spiritual comfort to the dying.

As we gaze at this splendid statue, there courses through memory’s corridors such melodies of the period as “Over There,” “Keep the Home Fires Burning,” and “Give My Regards to Broadway.” Were those fallen warriors who knew the song and remembered with affection Broadway and Forty-fifth Street to return and stand with us at the side of Father Duffy’s statue, what sight would meet their eyes and ours? On every hand are massage parlors, sex shops, X-rated movies—the neon-lighted signs flashing their facade of allure. The statue of Father Francis P. Duffy stands surrounded by sin, engulfed by evil. The pornography beetle has just about destroyed this area. He moves relentlessly closer to your city, your neighborhood, and your family.

An ominous warning was voiced by Laurence M. Gould, president emeritus of Carleton College:

“I do not believe the greatest threat to our future is from bombs or guided missiles. I don’t think our civilization will die that way. I think it will die when we no longer care. Arnold Toynbee has pointed out that 19 of 21 civilizations have died from within and not by conquest from without. There were no bands playing and flags waving when these civilizations decayed. It happened slowly, in the quiet and the dark when no one was aware.”

Just this month I read a review of a new movie. The leading actress told the reporter that she objected initially to the script and the part she was to play. The role portrayed her as the sexual companion of a fourteen-year-old boy. She commented: “At first I said, ‘No way will I agree to such a scene.’ Then I was given the assurance that the boy’s mother would be present during all intimate scenes, so I agreed.”

I ask: Would a mother stand by “watching,” were her son embraced by a cobra? Would she subject him to the taste of arsenic or strychnine? Mothers, would you? Fathers, would we?

From the past of long ago we hear the echo so relevant today:

“O Jerusalem, Jerusalem, which killest the prophets, and stonest them that are sent unto thee; how often would I have gathered thy children together, as a hen doth gather her brood under her wings, and ye would not!

“Behold, your house is left unto you desolate” (Luke 13:34–35).

Today we have a rebirth of ancient Sodom and Gomorrah. From seldom-read pages in dusty Bibles they come forth as real cities in a real world, depicting a real malady—pernicious permissiveness.

We have the capacity and the responsibility to stand as a bulwark between all we hold dear and the fatal contamination of the pornography beetle. May I suggest three specific steps in our battle plan:

First, a return to righteousness. An understanding of who we are and what God expects us to become will prompt us to pray—as individuals and as families. Such a return reveals the constant truth: “Wickedness never was happiness” (Alma 41:10). Let not the evil one dissuade. We can yet be guided by that still small voice—unerring in its direction and all-powerful in its influence.

Second, a quest for the good life. I speak not of the fun life, the sophisticated life, the popular life. Rather, I urge each to seek eternal life—life everlasting with mother, father, brothers, sisters, husband, wife, sons, and daughters, forever and forever together.

Third, a pledge to wage and win the war against pernicious permissiveness. As we encounter that evil carrier, the pornography beetle, let our battle standard and that of our communities be taken from that famous ensign of early America, “Don’t tread on me” (John Bartlett, Familiar Quotations, p. 1090).

Let us join in the fervent declaration of Joshua: “Choose you this day whom ye will serve … but as for me and my house, we will serve the Lord” (Josh. 24:15). Let our hearts be pure. Let our lives be clean. Let our voices be heard. Let our actions be felt.

Then the beetle of pornography will be halted in its deadly course. Pernicious permissiveness will have met its match. And we, with Joshua, will safely cross over Jordan into the promised land—even to eternal life in the celestial kingdom of our God.

That we may do so is my sincere prayer, in the name of Jesus Christ, amen.

\newpage
\settalk{Preparing the Way}
\setspeaker{Thomas S. Monson}
\settalkdate{April 1980}
\talkheading{Preparing the Way}{Thomas S. Monson}

We welcome to their new responsibilities and opportunities President Dwan J. Young and her counselors, Virginia B. Cannon and Michaelene P. Grassli. Certainly President Naomi Shumway and her counselors, Colleen B. Lemmon and Dorthea Lou C. Murdock, have established an enviable record of service on which to build.

Today I desire also to pay tribute to another Primary leader—a noble woman and personal friend. I speak of LaVern W. Parmley, former president of the Primary Association of The Church of Jesus Christ of Latter-day Saints and former member of the National Advisory Council, Boy Scouts of America. Sister Parmley, as she was affectionately addressed by those who knew her, completed her mission here on earth on Sunday, 27 January 1980. Her funeral services followed a few days later, where in word and song eulogies were expressed and comfort extended to those in attendance.

LaVern Parmley and Naomi Shumway, together with their counselors, gave freely of their time and talents over a span of many years, teaching little children to walk in the light of the gospel of Christ. They taught each child to sing with personal conviction:

\begin{verse}
I am a child of God, …\\
Lead me, guide me, walk beside me,\\
Help me find the way.\\
Teach me all that I must do\\
To live with him someday
\end{verse}

(“I Am a Child of God,” Sing with Me, B-76).

Part of their great love was teaching boys. Their inspired objective was to prepare them to receive the Aaronic Priesthood and to walk uprightly along the Scouting trail.

Under their direction, all eleven-year-old boys were asked to commit to memory the Articles of Faith of The Church of Jesus Christ of Latter-day Saints. You remember them. May I mention just two:

“We believe in God, the Eternal Father, and in His Son, Jesus Christ, and in the Holy Ghost” (A of F 1:1).

“We believe in being honest, true, chaste, benevolent, virtuous, and in doing good to all men; indeed, we may say that we follow the admonition of Paul—We believe all things, we hope all things, we have endured many things, and hope to be able to endure all things. If there is anything virtuous, lovely, or of good report or praiseworthy, we seek after these things” (A of F 1:13).

Can you think of a more firm foundation, a more basic philosophy to guide a boy, than the Articles of Faith? What a gift these noble women imparted when they expected each boy to know and indeed live by such a standard. They personally accepted the divine injunction, “Feed my sheep; feed my lambs” (see John 21:15–16).

Some may inquire: What is the significance of the Aaronic Priesthood for which such preparation takes place? Is it all that important in the life of a boy? The Priesthood of Aaron “is an appendage to the … Melchizedek Priesthood, and has power in administering outward ordinances” (D\&C 107:14). John the Baptist was a descendant of Aaron and held the keys of the Aaronic Priesthood. Perhaps we could review the life and mission of John, so that the importance of the Aaronic Priesthood might be more fully appreciated.

Long years ago and distant miles away, in the conquered country of Palestine, a marvelous miracle occurred. The setting was bleak, the time one of tumult. In these, the days of Herod, king of Judea, there lived a priest named Zacharias and his wife, Elisabeth. “They were … righteous before God” (Luke 1:6). However, long years of yearning had returned no reward—Zacharias and Elisabeth remained childless.

Then came that day of days ever to be remembered. There appeared to Zacharias the angel Gabriel, who proclaimed: “Fear not, Zacharias: for thy prayer is heard; and thy wife Elisabeth shall bear thee a son, and thou shalt call his name John.

“He shall be great in the sight of the Lord” (Luke 1:13–15).

Elisabeth did conceive. In due time a son was born, and according to the angel’s instruction he was named John.

As with the master, Jesus Christ, so with the servant, John—precious little is recorded of their years of youth. A single sentence contains all that we know of John’s history for a space of thirty years—the entire period which elapsed between his birth and his walk into the wilderness to commence his public ministry: “The child grew, and waxed strong in spirit, and was in the deserts till the day of his shewing unto Israel” (Luke 1:80).

His dress was that of the old prophets—a garment woven of camel’s hair. His food was such as the desert afforded—locusts and wild honey. His message was brief. He preached faith, repentance, baptism by immersion, and the bestowal of the Holy Ghost by an authority greater than that possessed by himself.

“I am not the Christ” (John 1:20), he told his band of faithful disciples. “But I am sent before him.” “I indeed baptize you with water … but one mightier than I cometh … : he shall baptize you with the Holy Ghost, and with fire (see Matt. 3:11).

Then there transpired the climactic scene of John’s mission—the baptism of Christ. Jesus came down from Galilee expressly “to be baptized” by John. Humbled of heart and contrite in spirit, John pleaded, “I have need to be baptized of thee, and comest thou to me?” The Master’s reply: “It becometh us to fulfil all righteousness.” (See Matt. 3:13–15.)

“And Jesus, when he was baptized, went up straightway out of the water: and, lo, the heavens were opened unto him, and he saw the Spirit of God descending like a dove, and lighting upon him:

“And lo a voice from heaven, saying, This is my beloved Son, in whom I am well pleased” (Matt. 3:16–17).

John’s testimony that Jesus was the redeemer of the world was declared boldly. Without fear, and with courage, John taught: “Behold the Lamb of God, which taketh away the sin of the world” (John 1:29).

Of John, the Savior later testified, “Among them that are born of women there hath not risen a greater than John the Baptist” (Matt. 11:11).

John’s public ministry moved toward its close. He had, at the beginning of it, condemned the hypocrisy and worldliness of the Pharisees and Sadducees; and he now had occasion to denounce the lust of a king. The result is well known. A king’s weakness and a woman’s fury combined to bring about the death of John.

The tomb in which his body was placed could not contain that body. Nor could the act of murder still that voice. To the world we declare that at Harmony, Pennsylvania, on 15 May 1829, an angel, “who announced himself as John, the same that is called John the Baptist in the New Testament” (D\&C 13, section heading), came as a resurrected personage to Joseph Smith and Oliver Cowdery. “The angelic visitant averred that he was acting under the direction of Peter, James, and John, the ancient Apostles, who held the keys of the higher Priesthood, which was called the Priesthood of Melchizedek” (D\&C 13, section heading). The Aaronic Priesthood was restored to the earth.

Thanks to that memorable event, I was given the privilege to bear the Aaronic Priesthood, as have millions of young men in these latter days. Its true significance was taught me by my former stake president, the late Paul C. Child.

When I was approaching my eighteenth birthday and preparing to enter military service in World War II, I was recommended to receive the Melchizedek Priesthood. Mine was the task of telephoning President Child for an appointment and interview. He was one who loved and understood the holy scriptures. It was his intent that all others should similarly love and understand them. Knowing from others of his rather detailed and searching interviews, our telephone conversation went something like this:

“Hello, President Child. This is Brother Monson. I have been asked by the bishop to seek an interview with you.”

“Fine, Brother Monson. When can you visit me?”

Knowing that his sacrament meeting time was six o’clock, and desiring minimum exposure of my scriptural knowledge to his review, I suggested, “How would five o’clock be?”

His response: “Oh, Brother Monson, that would not provide us sufficient time to peruse the scriptures. Could you please come at two o’clock, and bring with you your personally marked and referenced set of scriptures.”

Sunday finally arrived, and I visited President Child’s home on Indiana Avenue. I was greeted warmly, and then the interview began. He said: “Brother Monson, you hold the Aaronic Priesthood. Have you ever had angels minister to you?”

My reply was: “No, President Child.”

“Do you know,” said he, “that you are entitled to such?”

Again came my response: “No.”

Then he instructed, “Brother Monson, repeat from memory the thirteenth section of the Doctrine and Covenants.”

I began, “Upon you my fellow servants, in the name of Messiah I confer the Priesthood of Aaron, which holds the keys of the ministering of angels …” (D\&C 13:1).

“Stop,” President Child directed. Then in a calm, kindly tone he counseled: “Brother Monson, never forget that as a holder of the Aaronic Priesthood you are entitled to the ministering of angels.” It was almost as though an angel were in the room that day. I have never forgotten the interview. I yet feel the spirit of that solemn occasion. I revere the priesthood of Almighty God. I have witnessed its power. I have seen its strength. I have marveled at the miracles it has wrought.

Almost thirty years ago I knew a boy, even a priest, who held the authority of the Aaronic Priesthood. As the bishop, I was his quorum president. This boy, Robert, stuttered and stammered, void of control. Self-conscious, shy, fearful of himself and all others, he had an impediment of speech which was devastating to him. Never did he fulfill an assignment; never would he look another in the eye; always would he gaze downward. Then one day, through a set of unusual circumstances, he accepted an assignment to perform the priestly responsibility to baptize another.

I sat next to him in the baptistry of this sacred tabernacle. He was dressed in immaculate white, prepared for the ordinance he was to perform. I asked Robert how he felt. He gazed at the floor and stuttered almost incoherently that he felt terrible.

We both prayed fervently that he would be made equal to his task. Then the clerk read the words: “Nancy Ann McArthur will now be baptized by Robert Williams, a priest.” Robert left my side, stepped into the font, took little Nancy by the hand, and helped her into that water which cleanses human lives and provides a spiritual rebirth. He then gazed as though toward heaven and, with his right arm to the square, repeated the words “Nancy Ann McArthur, having been commissioned of Jesus Christ, I baptize you in the name of the Father, and of the Son, and of the Holy Ghost” (see D\&C 20:73). Not once did he stammer. Not once did he stutter. Not once did he falter. A modern miracle had been witnessed.

In the dressing room, as I congratulated Robert, I expected to hear this same uninterrupted flow of speech. I was wrong. He gazed downward and stammered his reply of gratitude.

To each of you this day, I testify that when Robert acted in the authority of the Aaronic Priesthood, he spoke with power, with conviction, and with heavenly help.

Such is the legacy of one called John, even John the Baptist. We hear his voice today. It teaches humility; it prompts courage; it inspires faith.

May we be motivated by his message. May we be inspired by his mission. May we be lifted by his life to a full appreciation of the Aaronic Priesthood and its divine power, I pray, in the name of Jesus Christ, amen.

\newpage
\settalk{The Bishop—Center Stage in Welfare}
\setspeaker{Thomas S. Monson}
\settalkdate{October 1980}
\talkheading{The Bishop—Center Stage in Welfare}{Thomas S. Monson}

Long years ago, the Apostle Paul wrote an epistle to his beloved associate Timothy in which he spoke of the qualifications a bishop should possess. He began:

“This is a true saying, If a man desire the office of a bishop, he desireth a good work.”

Today we might add, “and a good workout!” He continues:

“A bishop then must be blameless … vigilant, sober, of good behaviour, given to hospitality, apt to teach;

“… not greedy of filthy lucre; but patient. …

“Moreover he must have a good report of them which are without.” (1 Tim. 3:1–3, 7.)

These words burned into my soul when I read them as a newly called bishop thirty years ago. I was young—just twenty-two. The ward was large, with over 1,050 members, 87 of whom were widows. The welfare load was the heaviest of any ward in the entire Church.

The street addresses in that ward did not read like some: Colonial Hills Circle, Mountain View Estates, or Skyline Drive. Rather, they were noted as Orchard Place, Gale Street, Elford Court. The ward was not east of the tracks in Salt Lake City. It was not west of the tracks. That ward spanned the railroad tracks. Many of the widows and those in financial need were hidden away in basement apartments, upstairs back rooms, or dilapidated houses situated at the rear of little-known streets. I became the shepherd. This was my flock. I was reminded of God’s warning through Ezekiel: “Woe be to the shepherds of Israel that feed not my flock” (see Ezek. 34:2–3).

My teachers were heaven-sent. May I mention but a few: our former stake president, Harold B. Lee; President Marion G. Romney; and President J. Reuben Clark.

Brother Lee attended our stake conference the year I was appointed as a bishop. Looking over the priesthood leadership congregation on Saturday evening, he stepped from the pulpit, called for a chalkboard, came down from the stand, stood among us, and, as the master teacher, taught us our duty. He drew five circles under the heading “The Responsibilities of a Bishop.” He then gave to each circle a designation such as “The Father of the Ward,” “The President of the Aaronic Priesthood,” “The Common Judge in Israel,” and then placed emphasis on the bishop’s role in welfare. He cautioned us to seek after the poor, to care for them, and to do so in a spirit of love, kindness, and confidentiality.

Brother Romney was a frequent visitor to our stake and region. One evening he taught us the principle of faith by retelling the inspiring account of Elijah and the widow at Zarephath (see 1 Kgs. 17:8–16). He liken her circumstances to those of some widows in our area. As he taught welfare precepts from the handbook and responded to questions, one brother asked him, “Brother Romney, why do you seem to know whatever’s in that handbook?” to which Brother Romney, with that twinkle in his eye and smile upon his lips, responded, “I wrote it!”

President Clark too was a master teacher. It was my privilege during those years to assist him in the preparation of his manuscripts that they might find their way into printed volumes. What a unique and profitable experience to be with him frequently. Knowing that I was a new bishop presiding over a difficult ward, he emphasized the need for me to know my people, to understand their circumstances, and, in the spirit of tenderness, to minister to their needs. One day he recounted the example of the Savior as recorded in Luke, chapter seven, verses eleven through fifteen:

“And it came to pass … that he went into a city called Nain; and many of his disciples went with him. …

“When he came nigh to the gate of the city, behold, there was a dead man carried out, the only son of his mother, and she was a widow. …

“And when the Lord saw her, he had compassion on her, and said unto her, Weep not.

“And he came and touched the bier. … And he said, Young man, I say unto thee, Arise.

“And he that was dead sat up, and began to speak. And he delivered him to his mother.”

When President Clark closed the Bible, I noticed that he was weeping. In a quiet voice he said, “Tom, be kind to the widows, and look after the poor.”

Our bishops today need the same instruction and counsel. Many are new. They hear from every side that this program or that requires emphasis. Theirs is a sacred trust. Frequently that which counts most is recorded least. The visit to the elderly, the blessing to the sick, the comfort to the weary, the food to the hungry may not be recorded here, but I am convinced that they are known above and that we are guided in such ministrations.

The dimensions of the bishop’s welfare role are many. He is aided by his counselors, priesthood quorum leaders, and, of course, the ward Relief Society presidency. Perhaps a review would be helpful.

First, prevention. Paramount is the responsibility to coordinate personal and family preparedness efforts, including food storage. Also to be emphasized is the continuing need to insure that gainful employment is had by heads of families. Beyond this effort is the desirability to upgrade employment for those who may be underemployed. Such a duty involves the encouragement of wage earners to become skilled, that they will not be the last to be hired or the first to be fired.

Second, production. Participation on ward and stake welfare projects is a vital concern. Though times change, fields yet need to be plowed, crops to be thinned, buildings to be built, and storehouses to be filled.

I am grateful I learned to top sugar beets on our stake welfare farm. I am also grateful that we do not have to top beets in the same way today. That farm was not situated in a fertile belt of land but rather in the area of today’s industrial section of Salt Lake City. I testify, however, that when put to this sacred service, the soil was sanctified, the harvest blessed, and faith rewarded.

Third, processing. Oh, the joy of harvest time! Picture the scene of ward members canning peaches, sorting eggs, or cleaning vegetables, all for the use of those who are in need. Brows are sweat-lined, clothing is soiled, bodies are tired—but human souls are refreshed and lifted towards heaven.

Fourth, storage. The Lord in the revelations spoke frequently of his storehouses. On one occasion he counseled, “The storehouse shall be kept by the consecrations of the church; and widows and orphans shall be provided for, as also the poor” (D\&C 83:6). I am happy that over the entrance to our storehouses are the words Bishops’ Storehouse! Those who labor therein are recommended and sent by their respective bishops. Within such buildings there is found an atmosphere of love, of respect, and, indeed, of reverence. I am inspired each time I visit such a storehouse. There is no steeple or spire, no carpeted floors or stained-glass windows, but here is found the spirit of the Lord.

Fifth, distribution. This is where the bishop’s judgment is most severely tested. He cannot shirk this God-given responsibility. President J. Reuben Clark, Jr., summarized the bishop’s role in welfare services: He “is ‘to administer all temporal things’ … ; in his calling he is to ‘administer to the … poor and needy’; he is to search ‘after the poor to administer to their wants’ [see D\&C 107:68; 42:34; 84:112]. …

“Thus to the bishop is given all the powers, and responsibilities which the Lord has specifically prescribed in the Doctrine and Covenants for the caring of the poor. … No one else is charged with this duty and responsibility, no one else is endowed with the power and functions necessary for this work.

“Thus, ‘by the word of the Lord the sole mandate to care for and the sole discretion in caring for, the poor of the Church is lodged in the bishop.’ … ‘It is his duty and his only to determine to whom, when, how, and how much shall be given to any member of his ward from Church funds and as ward help.

“‘This is his high and solemn obligation, imposed by the Lord Himself. The bishop cannot escape this duty; he cannot shirk it; he cannot pass it on to someone else, and so relieve himself. Whatever help he calls in, he is still responsible.’” (Unpublished article, Church Historical Department, Salt Lake City, 9 July 1941, pp. 3–4.)

Every bishop needs a sacred grove to which he can retire to meditate and to pray for guidance. Mine was our old ward chapel. I could not begin to count the occasions when on a dark night at a late hour I would make my way to the stand of this building where I was blessed, confirmed, ordained, taught, and eventually called to preside. The chapel was dimly lighted by the street light in front; not a sound would be heard, no intruder to disturb. With my hand on the pulpit I would kneel and share with Him above my thoughts, my concerns, my problems.

On one occasion, a year of drought, the commodities at the storehouse had not been their usual quality, nor had they been found in abundance. Many products were missing, especially fresh fruit. My prayer that night is sacred to me. I pleaded that these widows were the finest women I knew in mortality, that their needs were simple and conservative, that they had no resources on which they might rely. The next morning I received a call from a ward member, a proprietor of a produce business. “Bishop,” he said, “I would like to send a semitrailer filled with oranges, grapefruit, and bananas to the bishops’ storehouse to be given to those in need. Could you make arrangements?” Could I make arrangements! The storehouse was alerted. Then each bishop was telephoned and the entire shipment distributed. Bishop Jesse M. Drury, that beloved welfare pioneer and storekeeper, said he had never witnessed a day like it before. He described the occasion with one word—“Wonderful!”

Other experiences may not be so dramatic but are nevertheless real and heartwarming. I recall an elderly couple whose frame home, situated at the end of a dirt lane, had not seen a coat of paint for too many years. These were neat and tidy people; they were concerned about the appearance of their small house. In a moment of inspiration I called, not upon the elders quorum or upon volunteers to wield paint brushes, but rather, following the welfare handbook, upon the family members who lived in other areas. Four sons-in-law and four daughters took brushes in hand and participated in the project. The paint had been provided by a dealer located in our area. The result was a transformation not only of the house but of the family. The children determined how they might best help mother and dad in their old age. They did so voluntarily and with gladness of heart. A house was painted, a family united, and respect preserved.

Fortunately, the blessings the welfare program provides are not received by the bishop alone. Rather, all who participate can share and share abundantly.

On a cold winter’s night in 1951 there was a knock at my door, and a German brother from Ogden, Utah, announced himself and said, “Are you Bishop Monson?” I answered in the affirmative. He began to weep and said, “My brother and his wife and family are coming here from Germany. They are going to live in your ward. Will you come with us to see the apartment we have rented for them?” On the way to the apartment, he told me he had not seen his brother for many years. Yet all through the holocaust of World War II, his brother had been faithful to the Church, serving as a branch president before the war took him to the Russian front.

I looked at the apartment. It was cold and dreary. The paint was peeling, the wallpaper soiled, the cupboards empty. A forty-watt bulb hanging from the living room ceiling revealed a linoleum floor covering with a large hole in the center. I was heartsick. I thought, “What a dismal welcome for a family which has endured so much.”

My thoughts were interrupted by the brother’s statement, “It isn’t much, but it’s better than they have in Germany.” With that, the key was left with me, along with the information that the family would arrive in Salt Lake City in three weeks—just two days before Christmas.

Sleep was slow in coming to me that night. The next morning was Sunday. In our ward welfare committee meeting, one of my counselors said, “Bishop, you look worried. Is something wrong?” I recounted to those present my experience of the night before, the details of the uninviting apartment. There were a few moments of silence. Then the group leader of the high priests said, “Bishop, did you say that apartment was inadequately lighted and that the kitchen appliances were in need of replacement?” I answered in the affirmative. He continued, “I am an electrical contractor. Would you permit the high priests of this ward to rewire that apartment? I would also like to invite my suppliers to contribute a new stove and a new refrigerator. Do I have your permission?” I answered with a glad “Certainly.”

Then the seventies president responded: “Bishop, as you know I’m in the carpet business. I would like to invite my suppliers to contribute some carpet, and the seventies can easily lay it and eliminate that worn linoleum.”

Then the president of the elders quorum spoke up. He was a painting contractor. He said, “I’ll furnish the paint. May the elders paint and wallpaper that apartment?”

The Relief Society president was next to speak: “We in the Relief Society cannot stand the thought of empty cupboards. May we fill them?”

The next three weeks are ever to be remembered. It seemed that the entire ward joined in the project. The days passed, and at the appointed time the family arrived from Germany. Again at my door stood the brother from Ogden. With an emotion-filled voice, he introduced to me his brother, wife, and their family. Then he asked, “Could we go visit the apartment?” As we walked up the staircase to the apartment, he repeated, “It isn’t much, but it’s more than they have had in Germany.” Little did he know what a transformation had taken place, that many who participated were inside waiting for our arrival.

The door opened to reveal a literal newness of life. We were greeted by the aroma of freshly painted woodwork and newly papered walls. Gone was the forty-watt bulb, along with the worn linoleum it had illuminated. We stepped on carpet deep and beautiful. A walk to the kitchen presented to our view a new stove and refrigerator. The cupboard doors were still open; however, they now revealed that every shelf was filled with food. The Relief Society as usual had done its work.

In the living room we began to sing Christmas hymns. We sang “Silent night! Holy night! All is calm; all is bright.” (Hymns, no. 160.) We sang in English; they sang in German. At the conclusion, the father, realizing that all of this was his, took me by the hand to express his thanks. His emotion was too great. He buried his head in my shoulder and repeated the words, “Mein Bruder, mein Bruder, mein Bruder.”

As we walked down the stairs and out into the night air, it was snowing. Not a word was spoken. Then a young girl asked, “Bishop, I feel better inside than I have ever felt before. Can you tell me why?”

I responded with the words of the Master: “Inasmuch as ye have done it unto one of the least of these my brethren, ye have done it unto me” (Matt. 25:40). Suddenly there came to mind the words from “O Little Town of Bethlehem”:

\begin{verse}
How silently, how silently,\\
The wondrous gift is given!\\
So God imparts to human hearts\\
The blessings of his heaven.
\end{verse}

(Hymns, no. 165.)

The poet said, “God gave his children memory, that in life’s garden there might be June roses in December” (C. Anketall Studdert-Kennedy, “Roses in December,” in The Best Loved Poems of the American People. sel. Hazel Felleman, Garden City, N.Y.: Garden City Publishing Co., 1936, p. 363). In my garden of memories no rose is more beautiful or fragrant than the rose brought to bloom by my participation in the welfare effort.

May our Heavenly Father ever bless our bishops in their sacred welfare responsibilities. Such duties are God-given. They were authored in heaven to bless in our day those who stand in need.

In the name of Jesus Christ, amen.

\newpage
\settalk{The Long Line of the Lonely}
\setspeaker{Thomas S. Monson}
\settalkdate{April 1981}
\talkheading{The Long Line of the Lonely}{Thomas S. Monson}

Today I desire to preach no sermon nor deliver a formal message. Rather, may I simply share with you my innermost thoughts. President David O. McKay referred to such as “heart petals.” I open to your view a window to my soul.

The Epistle of James has long been a favorite book of the Holy Bible. I find his brief message heart-warming and filled with life. Each of us can quote that well-known passage, “If any of you lack wisdom, let him ask of God, that giveth to all men liberally, and upbraideth not; and it shall be given him.” (James 1:5.) How many of us, however, remember his definition of religion? “Pure religion and undefiled before God and the Father is this, To visit the fatherless and widows in their affliction, and to keep himself unspotted from the world.” (James 1:27.)

The word widow appears to have had a most significant meaning to our Lord. He cautioned His disciples to beware the example of the scribes, who feigned righteousness by their long apparel and their lengthy prayers, but who devoured the houses of widows. (See Mark 12:38, 40.)

To the Nephites came the direct warning, “I will come near to you to judgment; and I will be a swift witness against … those that oppress the … widow.” (3 Ne. 24:5.)

To the Prophet Joseph Smith He directed, “The storehouse shall be kept by the consecrations of the church; and widows and orphans shall be provided for, as also the poor.” (D\&C 83:6.)

Such teachings were not new then. They are not new now. Consistently the Master has taught, by example, His concern for the widow. To the grieving widow at Nain, bereft of her only son, He came personally and to the dead son restored the breath of life—and to the astonished widow her son. To the widow at Zarephath, who with her son faced imminent starvation, He sent the prophet Elijah with the power to teach faith as well as provide food.

We may say to ourselves, “But that was long ago and ever so far away.” I respond: “Is there a city called Zarephath near your home? Is there a town known as Nain?” We may know our cities as Columbus or Coalville, Detroit or Denver. Whatever the name, there lives within each city the widow deprived of her companion and often her child. The need is the same. The affliction is real.

The widow’s home is generally not large or ornate. Frequently it is modest in size and humble in appearance. Often it is tucked away at the top of the stairs or the back of the hallway and consists of but one room. To such homes He sends you and me.

There may exist an actual need for food, clothing—even shelter. Such can be supplied. Almost always there remains the hope for that special hyacinth to feed the soul.

\begin{verse}
Go visit the lonely, the dreary;\\
Go comfort the weeping, the weary.\\
Oh, scatter kind deeds on your way\\
And make the world brighter today.
\end{verse}

The ranks of those in special need grow larger day by day. Note the obituary page of your newspaper. Here the drama of life unfolds to our view. Death comes to all mankind. It comes to the aged as they walk on faltering feet. Its summons are heard by those who have scarcely reached midway in life’s journey, and it often hushes the laughter of little children.

After the funeral flowers fade, the well wishes of friends become memories, the prayers offered and words spoken dim in the corridors of the mind. Those who grieve frequently join that vast throng I shall entitle “The Long Line of the Lonely.” Missed is the laughter of children, the commotion of teenagers, and the tender, loving concern of a departed companion. The clock ticks more loudly, time passes more slowly, and four walls do indeed a prison make.

Hopefully, all of us may again hear the echo of words spoken by the Master: “Inasmuch as ye have done it unto one of the least of these … , ye have done it unto me.” (Matt. 25:40.)

As we resolve to minister more diligently to those in need, let us remember to include our children in these learning lessons of life.

I have many memories of my boyhood days. Anticipating Sunday dinner was one of them. Just as we children hovered at our so-called starvation level and sat anxiously at the table with the aroma of roast beef filling the room, mother would say to me, “Tommy, before we eat, take this plate I’ve prepared down the street to Old Bob and hurry back.”

I could never understand why we couldn’t first eat and later deliver his plate of food. I never questioned aloud but would run down to his house and then wait anxiously as Bob’s aged feet brought him eventually to the door. Then I would hand him the plate of food. He would present to me the clean plate from the previous Sunday and offer me a dime as pay for my services. My answer was always the same: “I can’t accept the money. My mother would tan my hide.” He would then run his wrinkled hand through my blond hair and say, “My boy, you have a wonderful mother. Tell her thank you.”

You know, I think I never did tell her. I sort of felt mother didn’t need to be told. She seemed to sense his gratitude. I remember, too, that Sunday dinner always seemed to taste a bit better after I had returned from my errand.

Old Bob came into our lives in an interesting way. He was a widower in his eighties when the house in which he was living was to be demolished. I heard him tell my grandfather his plight as the three of us sat on the old front porch swing. With a plaintive voice, he said to grandfather, “Mr. Condie, I don’t know what to do. I have no family. I have no place to go. I have no money.” I wondered how grandfather would answer. Slowly grandfather reached into his pocket and took from it that old leather purse from which, in response to my hounding, he had produced many a penny or nickel for a special treat. This time he removed a key and handed it to Old Bob. Tenderly he said, “Bob, here is the key to that house I own next door. Take it. Move in your things. Stay as long as you like. There will be no rent to pay and nobody will ever put you out again.”

Tears welled up in the eyes of Old Bob, coursed down his cheeks, then disappeared in his long, white beard. Grandfather’s eyes were also moist. I spoke no word, but that day my grandfather stood ten feet tall. I was proud to bear his given name. Though I was but a boy, that lesson has influenced my life.

Each of us has his own way of remembering. At Christmas time I take delight in visiting the widows and widowers from the ward where I served as bishop. There were eighty-seven then—just nine today. On such visits, I never know what to expect; but this I do know: visits like these provide for me the Christmas spirit, which is, in reality, the Spirit of Christ.

Come with me, and we’ll together make a call or two. There’s the nursing home on West Temple where four widows reside. You never walk up the pathway but what you notice the parted curtain, as one inside waits hour after hour for the approaching step of a friend. What a welcome! Good times are remembered, perhaps a gift given, a blessing provided; but then it is time to leave. Never could I depart without first responding to the request of a widow almost one hundred years of age. Though she was blind, she would say, “Bishop, you’re to speak at my funeral and recite from memory Tennyson’s poem, ‘Crossing the Bar.’ Let’s hear you do it right now!” I would proceed:

\begin{verse}
Sunset and evening star,\\
And one clear call for me!\\
And may there be no moaning of the bar,\\
When I put out to sea, …
\end{verse}

(In Major British Writers, ed. G. B. Harrison, enl. ed., New York: Harcourt, Brace \& World, 1959, 2:466.)

Tears came easily, and then, with a smile, she would say, “Tommy, that was pretty good, but see that you do it a wee bit better at the funeral!” I later honored her request.

At another nursing home on First South, we might interrupt, as I did a few years ago, a professional football game. There, before the TV, were seated two widows. They were warmly and beautifully dressed, absorbed in the game. I asked, “Who’s winning?” They responded, “We don’t even know who’s playing, but at least it’s company.” I sat between those two angels and explained the game of football. I enjoyed the best contest I can remember. I may have missed a meeting, but I harvested a memory.

Let’s hurry along to Redwood Road. There is a much larger home here where many widows reside. Most are seated in the well-lighted living room. But in her bedroom, alone, is one on whom I must call. She hasn’t spoken a word since a devastating stroke some years ago. But then, who knows what she hears?—so I speak of good times together. There isn’t a flicker of recognition, not a word spoken. In fact, an attendant asks if I am aware that this patient hasn’t uttered a word for years. It made no difference. Not only had I enjoyed my one-sided conversation with her—I had communed with God.

When our beloved President Spencer W. Kimball met recently with those from a country where want is present, he asked not regarding statistics, but rather inquired: “Do our people have enough to eat? Are the widows cared for?” He was concerned.

During the administration of President George Albert Smith, there lived in our ward an impoverished widow who cared for her three mature daughters, each of whom was an invalid. They were large in size and almost totally helpless. To this dear woman fell the task to bathe, to feed, to dress, and to care for her girls. Means were limited. Outside help was nonexistent. Then came the blow that the house she rented was to be sold. What was she to do? Where would she go? The bishop came to the Church Office Building to inquire if there were some way the house could be purchased. It was so small, the price so reasonable. The request was considered, then denied.

A heartsick bishop was leaving the front door of the building when he met President George Albert Smith. After the exchange of greetings, President Smith inquired, “What brings you to the headquarters building?” He listened carefully as the bishop explained, but said nothing. He then excused himself for a few minutes. He returned wearing a smile and directed, “Go upstairs to the fourth floor. A check is waiting there for you. Buy the house!”

“But the request was denied.”

Again he smiled and said, “It has just been reconsidered and approved.” The home was purchased. That dear widow lived there and cared for her daughters until each of them had passed away. Then she, too, went home to God and to her heavenly reward.

The leadership of this Church is mindful of the widow, the widower, the lonely. Can we be less concerned? Emerson counseled that rings and jewels are not gifts, but substitutes for gifts. The only real gift is a portion of oneself. (See “Gifts,” by Ralph Waldo Emerson.)

We remember that during the meridian of time a bright, particular star shone in the heavens. Wise men followed it and found the Christ child. Today wise men still look heavenward and again see a bright, particular star. It will guide you and me to our opportunities. The burden of the downtrodden will be lifted, the cry of the hungry stilled, the lonely heart comforted. And souls will be saved—yours, theirs, and mine.

If we truly listen, we may hear that voice from far away say to us, as it spoke to another, “Well done, thou good and faithful servant.” (Matt. 25:21.)

May we see that special star, may we hear that same salutation, is my humble prayer, in the name of Jesus Christ, amen.

\newpage
\settalk{“He Is Risen”}
\setspeaker{Thomas S. Monson}
\settalkdate{October 1981}
\talkheading{“He Is Risen”}{Thomas S. Monson}

Not long ago a visitor asked, “What is there to see while I am in Salt Lake City?” Instinctively I suggested a tour of Temple Square, a drive to the nearby canyons, a visit to the Bingham copper mine, and perhaps a swim in the Great Salt Lake. A fear of being misunderstood kept me from expressing the thought, “Have you considered spending an hour or two at one of our cemeteries?” I never did reveal to him that wherever I travel I try to pay a visit to the town cemetery. It is a time of contemplation, of reflection on the meaning of life and the inevitability of death.

In the small cemetery in the equally small town of Santa Clara, Utah, I remember the preponderance of Swiss names which adorn the weathered tombstones. Many of those persons left home and family in verdant Switzerland and, in response to the call, “Come to Zion,” settled the communities where they now “rest in peace.” They endured spring floods, summer droughts, scant harvests, and back-breaking labors. They left a legacy of sacrifice.

The largest cemeteries, and in many respects those which evoke the most tender emotions, are honored as the resting places of men who died in the caldron of conflict known as war while wearing the uniform of their country. One reflects on shattered dreams, unfulfilled hopes, grief-filled hearts, and lives cut short by the sharp scythe of war.

Acres of neat, white crosses in the cities of France and Belgium accentuate the terrible toll of World War I. Verdun, France, is—in reality—a gigantic cemetery. Each spring, as farmers till the earth, they uncover a helmet here, a gun barrel there—grim reminders of the millions of men who literally soaked the soil with the blood of their lives.

A tour of Gettysburg, Pennsylvania, and other battlefields of the American Civil War marks that conflict, where brother fought against brother. Some families lost farms, others possessions. One family lost all. Let me share with you that memorable letter which President Abraham Lincoln wrote to Mrs. Lydia Bixby:

“Dear Madam:

“I have been shown in the files of the War Department a statement of the Adjutant General of Massachusetts that you are the mother of five sons who have died gloriously on the field of battle. I feel how weak and fruitless must be any words of mine which should attempt to beguile you from the grief of a loss so overwhelming. But I cannot refrain from tendering to you the consolation that may be found in the thanks of the republic they died to save. I pray that our Heavenly Father may assuage the anguish of your bereavement, and leave you only the cherished memory of the loved and lost, and the solemn pride that must be yours, to have laid so costly a sacrifice upon the altar of freedom.

“Yours very sincerely and respectfully, Abraham Lincoln.” (21 Nov. 1864; quoted in Selections from the Letters, Speeches, and State Papers of Abraham Lincoln, ed. Ida M. Tarbell, Boston: Ginn and Company, 1911, p. 109.)

A walk through Punchbowl Cemetery in Honolulu or the Memorial Cemetery of the Pacific at Manila reminds one that not all who died in World War II are buried in quiet fields of green. Many slipped beneath the waves of the oceans on which they sailed and on which they died.

Among the thousands of servicemen killed in the attack on Pearl Harbor was a sailor by the name of William Ball, from Fredericksburg, Iowa. What distinguished him from so many others who died on that day in 1941 was not any special act of heroism, but the tragic chain of events his death set in motion at home.

When William’s boyhood buddies, the five Sullivan brothers from the nearby town of Waterloo, received word of his death, they marched out together to enlist in the navy. The Sullivans, who wished to avenge their friend, insisted that they remain together, and the navy granted their wish. On November 14, 1942, the cruiser on which the brothers served, the U.S.S. Juneau, was hit and sunk in a battle off Guadalcanal in the Solomon Islands.

Almost two months went by before Mrs. Thomas Sullivan received the news, which arrived not by the usual telegram, but by special envoy: all five of her sons were reported missing in action in the South Pacific and presumed dead. Their bodies were never recovered.

One sentence only, spoken by one person only, provides a fitting epitaph: “Greater love hath no man than this, that a man lay down his life for his friends.” (John 15:13.)

Frequently the profound influence one life has on the lives of others is never spoken and, occasionally, little known. Such was the experience of a teacher of girls, even twelve-year-olds in the Beehive class of Mutual. She had no children of her own, though she and her husband dearly longed for children. Her love was expressed through the devotion to her special girls as she taught them eternal truths and lessons of life. Then came illness, followed by death. She was but twenty-seven.

Each year, on Memorial Day, her girls made a pilgrimage of prayer to the graveside of their teacher. First there were seven, then four, then two, and eventually just one, who continued the annual visit, always placing on the grave a bouquet of irises—a symbol of heartfelt gratitude. This year marked her twenty-fifth visit to the resting place of her teacher. Today she herself is a teacher of girls. Little wonder she is so successful. She mirrors the reflection of the teacher from whom came her inspiration. The life that teacher lived, the lessons that teacher taught, are not buried beneath the headstone which marks her grave, but live on in the personalities she helped to shape and the lives she so selflessly enriched. One is reminded of another master teacher, even the Lord. Once, with His finger, He wrote in the sand a message. (See John 8:6.) The winds of time erased forever the words He wrote, but not the life He lived.

“All that we can know about those we have loved and lost,” wrote Thornton Wilder, “is that they would wish us to remember them with a more intensified realization of their reality. … The highest tribute to the dead is not grief but gratitude.”

Two years ago, in beautiful Heber Valley just east of Salt Lake City, a loving mother and devoted father returned to that personal haven called home to discover that their three eldest sons lay dead. The night was bitter cold, and the fierce wind swept the falling snow, which covered the chimney, releasing deadly carbon monoxide fumes throughout the house.

The joint funeral service for the Keller boys was one of the most touching experiences of my life. The residents of the community had placed aside their daily tasks, children were excused from school, and all thronged to the chapel to express their deep feelings of condolence. So long as time and memory endure, I shall remember the scene of three shiny caskets, followed by grief-stricken parents and grandparents making their way to the front of the building.

The first speaker was the wrestling coach of the local high school. He paid tribute to Louis, the oldest boy. With an emotion-filled voice, and choking back the tears, he told how Louis was not necessarily the most gifted wrestler on the team, but added, “No one tried harder. What he lacked in athletic skill he made up with a determined heart.”

Then a youth leader spoke of Travis. He told how Travis had excelled in Scouting, in Aaronic Priesthood work, and was such a sterling example to his friends.

Finally, a distinguished appearing and obviously competent elementary school teacher told of Jason, the youngest of the three. She described him as quiet, even shy. Then, without embarrassment, she told how Jason had, in the scrawled penmanship of a boy, sent to her the sweetest and most welcome letter she had ever received. Its message was brief—just three words: “I love you.” She could barely complete her talk, so deep-felt were her emotions.

Through the tears and the sorrow of that special day, I observed eternal lessons that had been taught by those boys whose lives were honored and whose mortal missions concluded.

A coach expressed the determination to look beyond athletic prowess and into the heart of each boy. A youth leader made a solemn vow that every boy and girl would have the benefit which the program of the Church provided. An elementary school teacher looked at the small children, classmates of Jason. She said nothing, but her eyes revealed the determination of her soul. The message was unmistakably clear: “I will love each child. Each boy, each girl will be guided in the search for truth, in the development of talent, and be introduced to the wonderful world of service.”

And the audience, including Elders Marvin J. Ashton and Thomas S. Monson, will never again be the same. All will strive toward that perfection spoken of by the Master. Our inspiration? The lives of the boys who now rest from care and sorrow, and the fortitude of parents who trust in the Lord with all their hearts, who lean not to their own understanding, and who in all their ways acknowledge Him, knowing that He will direct their paths. (See Prov. 3:5–6.)

Let me share with you a portion of a letter sent to me by the noble mother of these three sons. It was written soon after their passing.

“We do have days and nights that right now seem so overwhelming. The change in our home life has been so drastic. With almost half our family gone now, the cooking, washing, and even shopping are different. We miss the noise and clutter, the teasing and playing together. Such are gone. Sunday is so quiet. We miss seeing the sacrament blessed and passed by our sons. Sunday was truly our family together day. We ponder the thought: no missions, no weddings, no grandchildren. We would not ask for their return, but we could not say we would ever have willingly given them up. We have returned to our Church duties and our family responsibilities. Our desire is to so live that the Keller family will be a forever family.”

To the Kellers, the Sullivans, and indeed to all who have loved and lost, let me share with you the conviction of my soul, the testimony of my heart, and the actual experiences of my life.

We know each one lived in the spirit world with Heavenly Father. We understand we have come to earth to learn, to live, to progress in our eternal journey toward perfection. Some remain on earth but for a moment, while others live long upon the land. The measure is not how long we live, but rather how well we live. Then come death and the beginning of a new chapter of life. Where does that chapter lead?

Many years ago I stood by the bedside of a young man, the father of two children, as he hovered between life and the great beyond. He took my hand in his, looked into my eyes and pleadingly asked, “Bishop, I know I am about to die. Tell me what happens to my spirit when I die.”

I prayed for heavenly guidance before attempting to respond. My attention was directed to the Book of Mormon, which rested on the table beside his bed. I held the book in my hand, and, as I stand before you here today, that book opened to the fortieth chapter of Alma. I began to read aloud:

“Now my son, here is somewhat more I would say unto thee; for I perceive that thy mind is worried concerning the resurrection of the dead. …

“Now, concerning the state of the soul between death and the resurrection—Behold, it has been made known unto me by an angel, that the spirits of all men, as soon as they are departed from this mortal body … are taken home to that God who gave them life.

“And then shall it come to pass, that the spirits of those who are righteous are received into a state of happiness, which is called paradise, a state of rest, a state of peace, where they shall rest from all their troubles and from all care, and sorrow.” (Alma 40:1, 11–12.)

My young friend closed his eyes, expressed a sincere thank-you, and silently slipped away to that paradise about which we had spoken.

Then comes that glorious day of resurrection, when spirit and body will be reunited, never again to be separated. “I am the resurrection, and the life,” said the Christ to the grieving Martha. “He that believeth in me, though he were dead, yet shall he live:

“And whosoever liveth and believeth in me shall never die.” (John 11:25–26.)

“Peace I leave with you, my peace I give unto you: not as the world giveth, give I unto you. Let not your heart be troubled, neither let it be afraid.” (John 14:27.)

“In my Father’s house are many mansions: if it were not so, I would have told you. I go to prepare a place for you … that where I am, there ye may be also.” (John 14:2–3.)

This transcendent promise became a reality when Mary and the other Mary approached the garden tomb—that cemetery which had but one occupant. Let Luke, the physician, describe their experience:

“Now upon the first day of the week, very early in the morning, they came unto the sepulchre. …

“And they found the stone rolled away. …

“… they entered in, and found not the body of the Lord Jesus.

“… as they were much perplexed thereabout, behold, two men stood by them in shining garments:

“And … said unto them, Why seek ye the living among the dead?” (Luke 24:1–5.)

“He is not here: for he is risen.” (Matt. 28:6.)

This is the clarion call of Christendom. The reality of the resurrection provides to one and all the peace that surpasses understanding. (See Philip. 4:7.) It comforts those whose loved ones lie in Flanders fields, who perished in the depths of the sea or rest in tiny Santa Clara or peaceful Heber Valley. It is a universal truth.

As the least of His disciples, I declare my personal witness that death has been conquered, victory over the tomb has been won. May the words made sacred by Him who fulfilled them become actual knowledge to all. Remember them. Cherish them. Honor them. He is risen. Such is my fervent prayer in the name of Jesus Christ, amen.

\newpage
\settalk{Sailing Safely the Seas of Life}
\setspeaker{Thomas S. Monson}
\settalkdate{April 1982}
\talkheading{Sailing Safely the Seas of Life}{Thomas S. Monson}

On February 14, 1939, Americans were celebrating Valentine’s Day. Postmen delivered sealed envelopes, and small children placed at the doorsteps of special friends folded papers containing brightly colored pictures. Each contained a greeting—a message of love. After all, Valentine’s Day is a day of love.

Far from America’s shores, in Hamburg, Germany, a public holiday also was being celebrated. However, a more somber mood prevailed. Amid fervent speeches, cheering throngs, and the playing of the national anthem, the new battleship Bismarck rumbled down into the River Elbe. This, the most powerful vessel afloat, carried not a message of love; rather, the Bismarck bristled with weapons of war.

The mighty colossus was a breathtaking spectacle of armor and machinery. Construction required more than fifty-seven thousand blueprints for the 406-millimeter, triple turret, radar-controlled guns. The vessel featured twenty-eight thousand miles of electrical circuits, and thirty-five thousand tons of armor-plate provided maximum safety. Majestic in appearance, gigantic in size, awesome in firepower, the Bismarck was considered unsinkable.

The Bismarck’s day of destiny dawned more than two years later, when on May 24, 1941, the two most powerful warships in the British navy, the Prince of Wales and the Hood, engaged in battle the Bismarck and the German cruiser Prinz Eugen. Within four minutes, the Bismarck had sent to the depths of the Atlantic the Hood and all but 3 men of a crew of 1,419. The other British battleship, the Prince of Wales, had suffered heavy damage and turned away.

Three days later, the Bismarck was engaged again, by four British warships. In all, the British concentrated the strength of eight battleships, two aircraft carriers, eleven cruisers, and twenty-one destroyers in an effort to seek and sink the mighty Bismarck.

Shell after shell inflicted but superficial damage. Was the Bismarck unsinkable after all? Then a torpedo scored a lucky hit which jammed the Bismarck’s rudder. Repair efforts proved fruitless. With guns primed, the crews at ready, the Bismarck could only steer a slow and stately circle. Just beyond reach was the powerful German air force. The safety of home port was ever so close. Neither could provide the needed haven, for the Bismarck had lost the ability to steer a charted course. No rudder; no help; no port. The end drew near. British guns blazed as the German crew scuttled and sank the once proud vessel. The hungry waves of the Atlantic first lapped at the sides, then swallowed the pride of the German navy. The Bismarck was no more. (See David Irving, Hitler’s War, New York: The Viking Press, 1977.)

Like the Bismarck, each of us is a miracle of engineering. Our creation, however, was not limited by human genius. Man can devise the most complex machines, but he cannot give them life or bestow upon them the powers of reason and judgment. Why? Because these are divine gifts, bestowed solely at God’s discretion. Our creator has provided us with a circulatory system to keep all channels constantly clean and serviceable, a digestive system to preserve strength and vigor, and a nervous system to keep all parts in constant communication and coordination. God gave man life, and with it, the power to think, to reason, to decide, and to love.

Like the vital rudder of a ship, we have been provided a way to determine the direction we travel. The lighthouse of the Lord beckons to all as we sail the seas of life. Our home port is the celestial kingdom of God. Our purpose is to steer an undeviating course in that direction. A man without a purpose is like a ship without a rudder—never likely to reach home port. To us comes the signal: Chart your course, set your sail, position your rudder, and proceed.

As with the ship, so it is with man. The thrust of the turbines, the power of the propellers are useless without that sense of direction, that harnessing of the energy, that directing of the power provided by the rudder, hidden from view, relatively small in size, but absolutely essential in function.

Our Father provided the sun, the moon, the stars—heavenly galaxies to guide mariners who sail the lanes of the sea. To all who walk the pathways of life, He cautions: Beware the detours, the pitfalls, the traps. Cunningly positioned are those clever pied pipers of sin beckoning here or there. Do not be deceived. Pause to pray. Listen to that still, small voice (see D\&C 85:6) which speaks to the depths of our souls the Master’s gentle invitation: “Come, follow me” (Luke 18:22). We turn from destruction, from death. We find happiness and life everlasting.

Yet, there are those who do not hear, who will not obey, who listen to the beat of a different drummer. Most prominent among their number was that son of Adam born of Eve, even Cain—a well-known name among men. Powerful in potential, but weak of will, Cain permitted greed, envy, disobedience, and even murder to jam that personal rudder which would have guided him to safety and exaltation. The downward gaze replaced the upward look; Cain fell. (See Moses 5:16–41.)

Less known, but more typical of our day, was that person of power, that cardinal of the cloth—even Wolsey. The prolific pen of William Shakespeare described the majestic heights, the pinnacle of power to which Cardinal Wolsey ascended. That same pen told how principle was eroded by vain ambition, by expediency, by a clamor for favor. Then came the tragic descent, the painful lament of one who had gained everything, then lost all. The words are beautiful; they border on scripture.

To Cromwell, his faithful servant, Cardinal Wolsey speaks:

\begin{verse}
When I am forgotten, as I shall be,\\
And sleep in dull cold marble, where no mention\\
Of me more must be heard of—say, I taught thee,\\
Say, Wolsey—that once trod the ways of glory,\\
And sounded all the depths and shoals of honour—\\
Found thee a way. …\\
A sure and safe one, though thy master mist it.\\
Mark but my fall, and that that ruin’d me.\\
… Fling away ambition:\\
By that sin fell the angels; how can man then,\\
The image of his Maker, hope to win by it?\\
Love thyself last; cherish those hearts that hate thee; …\\
Take an inventory of all I have,\\
To the last penny; ’tis the king’s: my robe,\\
And my integrity to heaven, is all\\
I dare now call my own.\\
O Cromwell, Cromwell!\\
Had I but serv’d my God with half the zeal\\
I served my king, He would not in mine age\\
Have left me naked to mine enemies.
\end{verse}

(King Henry the Eighth, act 3, sc. 2, lines 435–58.)

That heavenly rudder which would have ever been a guide to safety was ruined by the pursuit of power and quest for position. Like others before him and many more yet to follow, Cardinal Wolsey fell.

In an earlier time and by a wicked king, a servant of God was tested. Aided by the inspiration of heaven, Daniel, son of David, interpreted to the king the writing on the wall. Concerning the proffered rewards—even a royal robe and a necklace of gold—Daniel said: “Let thy gifts be to thyself, and give thy rewards to another.” (Dan. 5:17.)

Belshazzar’s successor, King Darius, also honored Daniel, elevating him to the highest position of prominence. There followed the envy of the crowd, the jealousy of princes, and the scheming of ambitious men.

Through trickery, aided by flattery, King Darius signed a proclamation that provided that anyone who made a request of any god or man, except the king, should be thrown into the lions’ den. (See Dan. 6:7.) The law was signed, the proclamation sent forth. Prayer was forbidden. In such matters, Daniel took direction not from an earthly king but from the king of heaven and earth, his God. Overtaken in his daily prayers, Daniel was brought before the king. Reluctantly, the penalty was pronounced. Daniel was to be thrown into the lions’ den. The sentence was carried out.

I love the biblical account which follows:

“The king went to his palace, and passed the night fasting … and his sleep went from him. …

“The king arose very early in the morning, and went in haste unto the den of lions.

“And when he came to the den, he cried with a lamentable voice unto Daniel. … O Daniel, servant of the living God, is thy God, whom thou servest continually, able to deliver thee from the lions?

“Then said Daniel unto the king, O king, live for ever.

“My God hath sent his angel, and hath shut the lions’ mouths, that they have not hurt me. …

“Then was the king exceeding glad. … Daniel was taken up out of the den, and no manner of hurt was found upon him, because he believed in his God.” (Dan. 6:18–23.)

In a time of critical need, Daniel’s determination to steer a steady course yielded divine protection and provided a sanctuary of safety.

The clock of history, like the sands of the hourglass, marks the passage of time. A new cast occupies the stage of life. The problems of our day loom ominously before us. Surrounded by the sophistication of modern living, we look heavenward for that unfailing sense of direction, that we might chart and follow a wise and proper course. He whom we call our Heavenly Father will not leave our sincere petition unanswered.

This lesson I learned anew some years ago as I received a rather unique and frightening assignment. Folkman D. Brown, then the Director of Mormon Relationships for the Boy Scouts of America, came to my office, having learned that I was about to depart for a lengthy assignment to New Zealand. He told me of his widowed sister, Belva Jones, who had been stricken with terminal cancer, who knew not how to tell her only son—a missionary in that far away country. Her wish, even her plea, was that he remain in the mission field and serve faithfully. She worried about his reaction; for the missionary, Elder Ryan Jones, had lost his father just a year earlier to the same dread disease.

I accepted the responsibility. Following a missionary meeting held adjacent to the majestically beautiful New Zealand Temple, I met privately with Elder Jones and, as gently as I could, explained the situation of his mother. Naturally there were tears—not all his—but then the handclasp of assurance and the pledge: “Tell my mother I will serve, I will pray, and I will see her again.”

I returned to Salt Lake City just in time to attend a conference of the Lost River Stake at Moore, Idaho. As I sat on the stand with the stake president, my attention was drawn almost instinctively to the east side of the chapel, where the morning sunlight bathed the lone occupant of a front bench. I said to the stake president, “Who is the sister upon whom the sunlight is resting? I feel I must speak to her today.” He replied, “Her name is Belva Jones. She has a missionary son in New Zealand. She is very ill and has requested a blessing.”

Prior to that moment, I had not known where Belva Jones lived. My assignment that weekend could have been to any one of fifty stakes. Yet the Lord, in His own way, had answered the prayer of faith of a concerned mother. We had a wonderful visit together. I reported word-for-word the reaction and the resolve of her son, Ryan. A blessing was provided, a prayer offered, a witness received. Belva Jones would live to see her son complete his mission. This privilege she enjoyed. Just one month prior to her passing, his mission completed, Ryan returned home.

As we venture forth on our individual voyages, may we sail safely the seas of life. With the never-failing rudder of faith guiding our passage, we too will find our way safely home. “Home is the sailor, home from sea.” Home to family, home to friends, home to heaven, home to God.

Of this truth I testify, in the name of Jesus Christ, amen.

\newpage
\settalk{“Run, Boy, Run!”}
\setspeaker{Thomas S. Monson}
\settalkdate{October 1982}
\talkheading{“Run, Boy, Run!”}{Thomas S. Monson}

Tuesday, June 8, 1982, dawned bright and clear in London, England. It was destined to be an historic day. A spirit of excitement permeated the very air and filled expectant hearts with keen anticipation. The President of the United States of America had arrived in Great Britain and soon would be addressing Parliament. Crowds gathered for the occasion, filled the streets and overflowed the nearby park. Uniformed policemen maintained order while famous Big Ben chimed its proud and clarion call which marked the appointed hour.

My wife, Frances, and I stood midst the milling crowd. Then, suddenly, Parliament’s doors swung open, the Prime Minister and the President greeted the throng, entered their limousines, and the motorcade drove slowly away. The crowd gave a mighty cheer, then began to disperse. Frances and I walked from the sunbathed street into the semi-dark, yet welcome, refuge of Westminster Abbey.

A reverence filled this world-famous edifice, as it should. For here, kings are crowned, royalty wedded, and rulers, whose mission of mortality has ended, are honored then buried. We walked along the aisleways, thoughtfully reading the inscriptions which marked the tombs of the famous. We remembered their achievements, recalled their deeds of valor, and marked their well-earned places in the world’s history. Then we paused before the Tomb of the Unknown Soldier, one of many who fell in France during the Great War. From an unmarked grave, the body of this fallen youth had been brought to London to forever lie in honor. I read aloud the inscriptions: “They buried him among the kings because he had done good toward God and toward His house.” “In Christ shall all be made alive.”

Toward the doorway we walked. Still visible in the park beyond were the remnants of the crowd. The immortal words of Rudyard Kipling coursed through my mind and spoke to my soul:

\begin{verse}
The tumult and the shouting dies,\\
The captains and the kings depart;\\
Still stands thine ancient sacrifice:\\
An humble and a contrite heart.\\
Lord God of hosts, be with us yet,\\
Lest we forget, lest we forget.
\end{verse}

(“Recessional”; see also Hymns, no. 77.)

One final marker to see, one more inscription to read. As a Scouter, I had come from America to view the plaque of honor dedicated to the memory of Scouting’s founder, Lord Baden-Powell. We stood before the magnificent marble memorial and noted the words:

On that day during this year which commemorates the 75th anniversary of Scouting and the 125th anniversary of its founder, I pondered the thought, “How many boys have had their lives blessed—even saved—by the Scout movement begun by Baden-Powell?” Unlike others memorialized within the walls of Westminster Abbey, Baden-Powell had neither sailed the stormy seas of glory, conquered in conflict the armies of men, nor founded empires of worldly wealth. Rather, he was a builder of boys, one who taught them well how to run and win the race of life.

Boys do become men.

\begin{verse}
Nobody knows what a boy is worth;\\
We’ll have to wait and see.\\
But every man in a noble place\\
A boy once used to be.
\end{verse}

(Quoted by Spencer W. Kimball, in Conference Report, Apr. 1977, p. 50.)

The reality of this thought is delightfully portrayed in the closing lines of the well-known musical Camelot. King Arthur’s Round Table has been destroyed by the jealousies of men, the infidelity of a queen, and the appearance in the present of a mistake from the past, even Mordred. Deprived of his dream, King Arthur and his forces prepare to meet the armies of Lancelot. All he held dear is gone; disillusionment has darkened into despair.

Suddenly, however, there appears a stowaway—the young boy Tom of Warwick. Filled with the hope of youth, he tells the king he has come to help him fight the mighty battle. He reveals his intention to become a knight. Under the questioning of Arthur, Tom declares his knowledge of the Round Table. He repeats the familiar goals: “Might for right! Right for right! Justice for all!”

A look of renewed confidence spreads across King Arthur’s face. All is not lost. To the boy he repeats the goals and glory of Camelot. Then he formally knights him “Sir Tom of Warwick.” Thus commissioned to depart the battlefield, to return to England, to renew the dream of Camelot, to grow up and to grow old, Sir Tom places aside the weapons of war; and armed with the tenets of truth, he hears his monarch command, “Run, boy, run!” A boy had been spared, an idea safeguarded, a hope renewed. (Alan J. Lerner, Camelot, New York: Random House, 1961, p. 115.)

Every boy blessed by Scouting learns in his youth far more than that envisioned by Sir Tom of Warwick. He adopts the motto “Be Prepared.” He subscribes to the code “Do a Good Turn Daily.” Scouting provides proficiency badges to encourage skills and personal endeavor. Scouting teaches boys how to live, not merely how to make a living. How pleased I am that The Church of Jesus Christ of Latter-day Saints in 1913 became the first partner to sponsor Scouting in the United States.

I love the inspired words of President Spencer W. Kimball when he spoke to Church members everywhere: “The Church of Jesus Christ of Latter-day Saints affirms the continued support of Scouting and will seek to provide leadership which will help boys keep close to their families and close to the Church as they develop the qualities of citizenship and character and fitness which Scouting represents. We’ve remained strong and firm in our support of this great movement for boys and of the Oath and the Law which are at its center.” (In Conference Report, Apr. 1977, pp. 50–51.)

What is the Scout Oath of which President Kimball spoke?

“On my honor I will do my best to do my duty to God and my country and to obey the Scout Law; to help other people at all times and to keep myself physically strong, mentally awake and morally straight.” (Boy Scout Handbook, North Brunswick, New Jersey: Boy Scouts of America, 1972, p. 34.)

A hero from war’s battlefield, General of the Army Douglas MacArthur, emphasized this same commitment when, in the twilight of his illustrious career, when the daylight of youth had departed and the shadows of age had descended, he declared in a message to young men: “In my dreams I hear again the crash of guns, the rattle of musketry, the strange, mournful mutter of the battlefield. But in the evening of my memory, I always come back to West Point. Always there echoes and re-echoes in my ears—Duty, Honor, Country.” (Address accepting Sylvanus Thayer Award, West Point, 12 May 1962.)

The Protestant minister Harry Emerson Fosdick phrased differently the same commitment: “Men will work hard for money. They will work harder for other men. But men will work hardest of all when they are dedicated to a cause. Until willingness overflows obligation, men fight as conscripts, rather than following the flag as patriots. Duty is never worthily performed until it is performed by one who would gladly do more, if only he could.” (Vital Quotations, comp. Emerson Roy West, Salt Lake City: Bookcraft, 1968, p. 38.)

And from the Confederate general, Robert E. Lee: “Duty is the sublimest word in the English language. Do your duty in all things. You cannot do more. You should never wish to do less.” (Inscription beneath his bust in the Hall of Fame.)

Let us consider the Scout Law referred to by President Kimball. When I think of the Scout Law, I reflect upon the life of one who knew the laws of God and who kept them—even the Lord, Jesus Christ. The twelve points of the Scout Law have their counterpart in the message of the Master.

Such inspired teachings, when taught by devoted leaders to precious boys of promise, influence not only the lives of the boys; they also affect eternity. “Cast thy bread upon the waters: for thou shalt find it after many days.” (Eccl. 11:1.) Such is Scouting.

Several years ago a group of men, leaders of Scouts, assembled in the mountains near Sacramento for Wood Badge training. This experience, where men camp out and live as do the Scouts they teach, is a most interesting one. They cook and then eat—burned eggs! They hike the rugged trails which age invariably makes more steep. They sleep on rocky ground. They gaze again at heaven’s galaxies.

This group provided its own reward. After days of being deprived, they feasted on a delicious meal prepared by a professional chef who joined them at the end of their endurance trail. Tired, hungry, a bit bruised after their renewal experience, one asked the chef why he was always smiling and why each year he returned at his own expense to cook the traditional meal for Scouting’s leaders in that area. He placed aside the skillet, wiped his hands on the white apron which graced his rotund figure, and told the men this experience. Dimitrious began:

“I was born and grew to boyhood in a small village in Greece. My life was a happy one until World War II. Then came the invasion and occupation of my country by the Nazis. The freedom-loving men of the village resented the invaders and engaged in acts of sabotage to show their resentment.

“One night, after the men had destroyed a hydroelectric dam, the villagers celebrated the achievement and then retired to their homes.”

Dimitrious continued: “Very early in the morning, as I lay upon my bed, I was awakened by the noise of many trucks entering the village. I heard the sound of soldiers’ boots, the rap at the door, and the command for every boy and man to assemble at once on the village square. I had time only to slip into my trousers, buckle my belt, and join the others. There, under the glaring lights of a dozen trucks, and before the muzzles of a hundred guns, we stood. The Nazis vented their wrath, told of the destruction of the dam, and announced a drastic penalty: every fifth man or boy was to be summarily shot. A sergeant made the fateful count, and the first group was designated and executed.”

Dimitrious spoke more deliberately to the Scouters as he said: “Then came the row in which I was standing. To my horror, I could see that I would be the final person designated for execution. The soldier stood before me, the angry headlights dimming my vision. He gazed intently at the buckle of my belt. It carried on it the Scout insignia. I had earned the belt buckle as a Boy Scout for knowing the Oath and the Law of Scouting. The tall soldier pointed at the belt buckle, then raised his right hand in the Scout sign. I shall never forget the words he spoke to me: ‘Run, boy, run!’ I ran. I lived. Today I serve Scouting, that boys may still dream dreams and live to fulfill them.” (As told by Peter W. Hummel.)

Dimitrious reached into his pocket and produced that same belt buckle. The emblem of Scouting still shone brightly. Not a word was spoken. Every man wept. A commitment to Scouting was renewed.

It has been said, “The greatest gift a man can give a boy is his willingness to share a part of his life with him.” To leaders who build bridges to the hearts of boys, to parents of Scouts, and to Scouts everywhere, on this our 75th anniversary, I salute you and pray our Heavenly Father’s blessings upon you. In the name of Jesus Christ, amen.

\newpage
\settalk{“Anonymous”}
\setspeaker{Thomas S. Monson}
\settalkdate{April 1983}
\talkheading{“Anonymous”}{Thomas S. Monson}

Recently, I approached the reception desk of a large hospital to learn the room number of a patient I had come to visit. This hospital, like almost every other in the land, was undergoing a massive expansion. Behind the desk where the receptionist sat was a magnificent plaque which bore an inscription of thanks to donors who had made possible the expansion. The name of each donor who had contributed \$100,000 appeared in a flowing script, etched on an individual brass placard suspended from the main plaque by a glittering chain.

The names of the benefactors were well known. Captains of commerce, giants of industry, professors of learning—all were there. I felt gratitude for their charitable benevolence. Then my eyes rested on a brass placard which was different—it contained no name. One word, and one word only, was inscribed: “Anonymous.” I smiled and wondered who the unnamed contributor could have been. Surely he or she experienced a quiet joy unknown to any other.

My thoughts turned backward in time—back to the Holy Land; back to Him whom we especially remember this Easter Sunday; back to Him who redeemed from the grave all mankind; back to Him who on that special mountain taught His disciples the true spirit of giving when He counseled, “Take heed that ye do not your alms before men, to be seen of them. …

“But when thou doest alms, let not thy left hand know what thy right hand doeth.” (Matt. 6:1, 3.)

Then, as though to indelibly impress on their souls the practical application of this sacred truth, He came down from the mountain with a great multitude following Him. “And, behold, there came a leper and worshipped him, saying, Lord, if thou wilt, thou canst make me clean.

“And Jesus put forth his hand, and touched him, saying, I will; be thou clean. And immediately his leprosy was cleansed.

“And Jesus saith unto him, See thou tell no man.” (Matt. 8:2–4.) The word anonymous had a precious meaning then. It still has.

The classics of literature, as well as the words from holy writ, teach us the endurability of anonymity. A favorite of mine is Charles Dickens’ “A Christmas Carol.” I can picture the trembling Ebenezer Scrooge seeing in vision the return of his former partner, Jacob Marley, though Jacob had been dead for seven years. The words of Marley penetrate my very soul, as he laments, “Not to know that any Christian spirit working kindly in its little sphere, whatever it may be, will find its mortal life too short for its vast means of usefulness. Not to know that no space of regret can make amends for one life’s opportunity misused! Yet such was I!” (“A Christmas Carol,” in The Best Short Stories of Charles Dickens, New York: Charles Scribner’s Sons, 1947, p. 435.)

After a fretful night—wherein Scrooge was shown by the Ghost of Christmas Past, the Ghost of Christmas Present, and the Ghost of Christmas Yet to Come the true meaning of living, loving, and giving—he awakened to discover anew the freshness of life, the power of love, and the spirit of a true gift. He remembered the plight of the Bob Cratchit family, arranged with a lad to purchase the giant turkey (the size of a boy), and sent the gift to the Cratchits. Then, with supreme joy, the reborn Ebenezer Scrooge exclaims to himself, “He shan’t know who sends it.” (“A Christmas Carol,” p. 481.) Again the word anonymous.

The sands flow through the hourglass, the clock of history moves on; yet the divine truth prevails undiminished, undiluted, unchanged.

When the magnificent ocean liner Lusitania plunged to the bottom of the Atlantic, many lives were lost with the vessel. Unknown are many deeds of valor performed by those who perished. One man who went down with the Lusitania gave his life preserver to a woman, though he could not swim a stroke. It didn’t really matter that he was Alfred Vanderbilt, the American multimillionaire. He did not give of worldly treasure; he gave his life. Said Emerson, “Rings and other jewels are not gifts, but apologies for gifts. The only gift is a portion of thyself.” (“Gifts,” in The Complete Writings of Ralph Waldo Emerson, New York: Wm. H. Wise and Co., 1929, p. 286.)

A year ago last winter, a modern jetliner faltered after takeoff and plunged into the icy Potomac River. Acts of bravery and feats of heroism were in evidence that day, the most dramatic of which was one witnessed by the pilot of a rescue helicopter. The rescue rope was lowered to a struggling survivor. Rather than grasping the lifeline to safety, the man tied the line to another, who was then lifted to safety. The rope was lowered again, and yet another was saved. Five were rescued from the icy waters. Among them was not found the anonymous hero. Unknown by name, “he left the vivid air signed with his honor.” (Stephen Spender, “I think continually of those—” in Masterpieces of Religious Verse, ed. James Dalton Morrison, New York: Harper and Brothers Publishers, p. 291.)

It is not only in dying that one can show forth the true gift. Opportunities abound in our daily lives to demonstrate our adherence to the Master’s lesson. Let me share in capsule form just three:

(1) On a winter’s morn, a father quietly awakened his two sons and whispered to them, “Boys, it snowed last night. Get dressed, and we’ll shovel the snow from our neighbors’ walks before daylight.”

The party of three, dressed warmly, and under cover of darkness, cleared the snow from the approaches to several homes. Father had given but one instruction to the boys: “Make no noise, and they will not know who helped them.” Again, the word anonymous.

(2) At a nursing home in our valley, two young men prepared the sacrament. While doing so, an elderly patient in a wheelchair spoke aloud the words, “I’m cold.” Without a moment’s hesitation, one of the young men walked over to her, removed his own jacket, placed it about the patient’s shoulders, gave her a loving pat on the arm, and then returned to the sacrament table. The sacred emblems were then blessed and passed to the assembled patients.

Following the meeting, I said to the young man, “What you did here today I shall long remember.”

He replied, “I worried that without my jacket I would not be properly dressed to bless the sacrament.”

I responded, “Never was one more properly dressed for such an occasion than were you.”

I know not his name. He remains anonymous.

(3) In far-off Europe, beyond a curtain of iron and a wall called “Berlin,” I visited, with a handful of members, a small cemetery. It was a dark night, and a cold rain had been falling throughout the entire day.

We had come to visit the grave of a missionary who many years before had died while in the service of the Lord. A hushed silence shrouded the scene as we gathered about the grave. With a flashlight illuminating the headstone, I read the inscription:

Then the light revealed that this grave was unlike any other in the cemetery. The marble headstone had been polished, weeds such as those which covered other graves had been carefully removed, and in their place was an immaculately edged bit of lawn and some beautiful flowers that told of tender and loving care. I asked, “Who has made this grave so attractive?” My query was met by silence.

At last a twelve-year-old deacon acknowledged that he wanted to render this unheralded kindness and, without prompting from parents or leaders, had done so. He said that he just wanted to do something for a missionary who gave his life while in the service of the Lord. I thanked him; and then I asked all there to safeguard his secret, that his gift might remain anonymous.

Perhaps no one in my reading has portrayed this teaching of the Master quite so memorably or so beautifully as Henry Van Dyke in his never-to-be-forgotten “The Mansion.” In this classic is featured one John Weightman, a man of means, a dispenser of political power, a successful citizen. His philosophy toward giving can be gained from his own statement: “Of course you have to be careful how you give, in order to secure the best results—no indiscriminate giving—no pennies in beggars’ hats! … Try to put your gifts where they can be identified and do good all around.” (See “The Mansion,” Unknown Quantity: A Book of Romance and Some Half-told Tales, New York: Scribner’s, 1918, pp. 337, 339.)

One evening, John Weightman sat in his comfortable chair at his library table and perused the papers before him spread. There were descriptions and pictures of the Weightman wing of the hospital and the Weightman Chair of Political Jurisprudence, as well as an account of the opening of the Weightman Grammar School. John Weightman felt satisfied.

He picked up the family Bible which lay on the table, turned to a passage and read to himself the words: “Lay not up for yourselves treasures upon earth, where moth and rust doth corrupt, and where thieves break through and steal:

“But lay up for yourselves treasures in heaven.” (Matt. 6:19–20.)

The book seemed to float away from him. He leaned forward upon the table, his head resting on his folded hands. He slipped into a deep sleep.

In his dream, John Weightman was transported to the Heavenly City. A guide met him and others whom he had known in life and advised that he would conduct them to their heavenly homes.

The group paused before a beautiful mansion and heard the guide say, “This is the home for you, Dr. McLean. Go in; there is no more sickness here, no more death, nor sorrow, nor pain; for your old enemies are all conquered. But all the good that you have done for others, all the help that you have given, all the comfort that you have brought, all the strength and love that you bestowed upon the suffering, are here; for we have built them all into this mansion for you.” (“The Mansion,” pp. 361–62.)

A devoted husband of an invalid wife was shown a lovely mansion, as were a mother, early widowed, who reared an outstanding family, and a paralyzed young woman who had lain for thirty years upon her bed—helpless but not hopeless—succeeding by a miracle of courage in her single aim: never to complain, but always to impart a bit of her joy and peace to everyone who came near her.

By this time, John Weightman was impatient to see what mansion awaited him. As he and the Keeper of the Gate walked on, the homes became smaller—then smaller. At last they stood in the middle of a dreary field and beheld a hut, hardly big enough for a shepherd’s shelter. Said the guide, “This is your mansion, John Weightman.”

In desperation, John Weightman argued, “Have you not heard that I have built a schoolhouse; a wing of a hospital; … three … churches?”

“Wait,” the guide cautioned. “… They were not ill done. But they were all marked and used as foundations for the name and mansion of John Weightman in the world. … Verily, you have had your reward for them. Would you be paid twice?”

A sadder but wiser John Weightman spoke more lowly: “What is it that counts here?”

Came the reply, “Only that which is truly given. Only that good which is done for the love of doing it. Only those plans in which the welfare of others is the master thought. Only those labors in which the sacrifice is greater than the reward. Only those gifts in which the giver forgets himself.” (“The Mansion,” pp. 364–68.)

John Weightman was awakened by the sound of the clock chiming the hour of seven. He had slept the night through. As it turned out, he yet had a life to live, love to share, and gifts to give. Oh, may we remember that—

\begin{verse}
A bell is no bell till you ring it,\\
A song is no song till you sing it,\\
And love in your heart wasn’t put there to stay,\\
Love isn’t love till you give it away.
\end{verse}

(Richard Rodgers and Oscar Hammerstein 2nd, “Sixteen Going on Seventeen.”)

May this truth guide our lives. May we look upward as we press forward in the service of our God and our fellowmen. And may we incline an ear toward Galilee, that we might hear perhaps an echo of the Savior’s teachings: “Do not your alms before men, to be seen of them.” (Matt. 6:1.) “Let not thy left hand know what thy right hand doeth.” (Matt. 6:3.) And of our good deeds: “See thou tell no man.” (Matt. 8:4.) Our hearts will then be lighter, our lives brighter, and our souls richer.

Loving service anonymously given may be unknown to man—but the gift and the giver are known to God. Of this truth I testify, in the name of Jesus Christ, amen.

\newpage
\settalk{Labels}
\setspeaker{Thomas S. Monson}
\settalkdate{October 1983}
\talkheading{Labels}{Thomas S. Monson}

The National Gallery at Trafalgar Square in London, England, is one of the truly great museums of art in all the world. The gallery proudly proclaims its Rembrandt Room and Constable Corner and urges all to take the tour of Turner’s masterpieces. Visitors come from every corner of the earth. They depart uplifted and inspired.

During a recent visit to the National Gallery, I was surprised to see displayed in a most prominent location magnificent portraits and landscapes which featured the name of no artist. Then I noticed a large placard which provided this explanation:

“This exhibition is drawn from the large number of paintings that hang in a public but somewhat neglected area of the Gallery: the lower floor. The exhibition is intended to encourage visitors to look at the paintings without being too worried about who painted them. In several instances, we do not precisely know.

“The information on labels on paintings can often affect, half-unconsciously, our estimate of them; and here labeling has been deliberately subordinate in the hope that visitors will read only after they have looked and made their own assessment of each work.”

Like the labels on paintings are the outward appearances of some men—often misleading. The Master declared to one group: “Woe unto you, scribes and Pharisees, hypocrites! for ye are like unto whited sepulchres, which indeed appear beautiful outward, but are within full of dead men’s bones, and of all uncleanness. …

“Ye … outwardly appear righteous unto men, but within ye are full of hypocrisy and iniquity.” (Matt. 23:27–28.)

Then there are those who may outwardly appear impoverished, without talent, and doomed to mediocrity. A classic label appeared beneath a picture of the boy Abraham Lincoln as he stood in front of his humble birthplace—a simple log cabin. The words read: “Ill-housed, ill-clothed, ill-fed.” Unanticipated, unspoken, and unprinted was the real label of the boy: “Destined for immortal glory.”

As the poet expressed:

\begin{verse}
Nobody knows what a boy is worth,\\
We’ll have to wait and see.\\
But every man in a noble place,\\
A boy once used to be.
\end{verse}

At another time, and in a distant place, the boy Samuel must have appeared like any lad his age as he ministered unto the Lord before Eli. As Samuel lay down to sleep and heard the voice of the Lord calling him, Samuel mistakenly thought it was aged Eli calling and responded, “Here am I.” (1 Sam. 3:4.) However, after Eli had listened to the boy’s account and told him it was of the Lord, Samuel followed Eli’s counsel and subsequently responded to the Lord’s call with the memorable reply, “Speak; for thy servant heareth.” (1 Sam. 3:10.) The record then reveals that “Samuel grew, and the Lord was with him. …

“And all Israel from Dan even to Beer-sheba knew that Samuel was established to be a prophet of the Lord.” (1 Sam. 3:19–20.)

The years rolled by, as they relentlessly do, and prophecy came to fulfillment when a lowly manger cradled a newborn child. No label could describe this event. With the birth of the babe in Bethlehem, there emerged a great endowment, a power stronger than weapons, a wealth more lasting than the coins of Caesar. This child, born in such primitive circumstances, was to be the “King of kings and Lord of lords” (1 Tim. 6:15), the promised Messiah—even Jesus Christ, the Son of God.

As a boy, Jesus was found in the temple, “sitting in the midst of the doctors, both hearing them, and asking them questions.

“And all that heard him were astonished at his understanding and answers.” And when Joseph and His mother saw Him, “they were amazed.” (See Luke 2:46–48.) To the learned doctors in the temple, the boy’s outward label may have conveyed brightness of intellect but certainly not “Son of God and future Redeemer of all mankind.”

The Messianic words of the prophet Isaiah convey a special meaning: “He hath no form nor comeliness; and when we shall see him, there is no beauty that we should desire him.” (Isa. 53:2.) Such was the foretold description of our Lord.

Matthew records the apparent necessity of that wicked multitude of sinners who would seek after the life of the Lord to conspire with the betrayer Judas, that he might point out to them who of the apostolic group was the Jesus whom they sought. These chilling verses from sacred writ torment the reader: “Now he that betrayed him gave them a sign, saying, Whomsoever I shall kiss, that same is he: hold him fast.

“And forthwith he came to Jesus, and said, Hail, master; and kissed him.

“And Jesus said unto him, Friend, wherefore art thou come? Then came they, and laid hands on Jesus, and took him.” (Matt. 26:48–50.)

The label of a traitor’s kiss had identified the Master. Judas now wore his own label of inescapable shame and revulsion.

Sometimes cities and nations bear special labels of identity. Such was a cold and very old city in eastern Canada. The missionaries called it “Stony Kingston.” There had been but one convert to the Church in six years, even though missionaries had been continuously assigned there during the entire interval. No one baptized in Kingston. Just ask any missionary who labored there. Time in Kingston was marked on the calendar like days in prison. A missionary transfer to another place—any place—would be uppermost in thoughts, even in dreams.

While I was praying about and pondering this sad dilemma, for my responsibility then as a mission president required that I pray and ponder about such things, my wife called to my attention an excerpt from the book, A Child’s Story of the Prophet Brigham Young, by Deta Petersen Neeley (Salt Lake City: Deseret News Press, 1959, p. 36). She read aloud that Brigham Young entered Kingston, Ontario, on a cold and snow-filled day. He labored there about thirty days and baptized forty-five souls. Here was the answer. If the missionary Brigham Young could accomplish this harvest, so could the missionaries of today.

Without providing an explanation, I withdrew the missionaries from Kingston, that the cycle of defeat might be broken. Then the carefully circulated word: “Soon a new city will be opened for missionary work, even the city where Brigham Young proselyted and baptized forty-five persons in thirty days.” The missionaries speculated as to the location. Their weekly letters pleaded for the assignment to this Shangri-la. More time passed. Then four carefully selected missionaries—two of them new, two of them experienced—were chosen for this high adventure. The members of the small branch pledged their support. The missionaries pledged their lives. The Lord honored both.

In the space of three months, Kingston became the most productive city of the Canadian Mission. The grey limestone buildings still stood, the city had not altered its appearance, the population remained constant. The change was one of attitude. The label of doubt yielded to the label of faith.

The branch president of the Kingston Branch of the Church wore his own identifying label. Gustav Wacker was from the old country. He spoke English with a thick accent. He never owned or drove a car. He plied the trade of a barber. The highlight of his day would be when he had the privilege of cutting the hair of a missionary. Never would there be a charge. Indeed, he would reach deep into his pockets and give the missionaries all of his tips for the day. If it were raining, as it often does in Kingston, President Wacker would call a taxi and send the missionaries to their apartment by taxi, while he himself, at day’s end, would lock the small shop and walk home—in the driving rain.

I first met Gustav Wacker when I noticed that his tithing paid was far in excess of that expected from his potential income. My efforts to explain that the Lord required no more than ten percent as tithing fell on attentive but unconvinced ears. He simply responded that he loved to pay all he could to the Lord. It amounted to about half his income. His dear wife felt exactly as he did. Their unique manner of tithing payment continued throughout their earning lives.

Gustav and Margarete Wacker established a home that was a heaven. They were not blessed with children but mothered and fathered their many Church visitors. A sophisticated and learned leader from Ottawa told me, “I like to visit President Wacker. I come away refreshed in spirit and determined to ever live close to the Lord.”

Did our Heavenly Father honor such abiding faith? The branch prospered. The membership outgrew the rented Slovakian Hall and moved into a modern and lovely chapel of their own. President and Sister Wacker had their prayers answered by serving a proselyting mission to their native Germany and later a temple mission to the beautiful temple in Washington, D.C. Then, just three months ago, his mission in mortality concluded, Gustav Wacker passed away peacefully while being held in the loving arms of his eternal companion. Only one label appears fitting for such an obedient and faithful servant: “Who honors God, God honors.” (See 1 Sam. 2:30.)

A label frequently seen and grudgingly borne is one which reads: “Handicapped.”

Years ago, President Spencer W. Kimball shared with President Gordon B. Hinckley, Elder Bruce R. McConkie, and me an experience he had in the appointment of a patriarch for the Shreveport Louisiana Stake of the Church. President Kimball described how he interviewed, how he searched, and how he prayed, that he might learn the Lord’s will concerning the selection. For some reason, none of the suggested candidates was the man for this assignment at this particular time.

The day wore on. The evening meetings began. Suddenly President Kimball turned to the stake president and asked him to identify a particular man seated perhaps two-thirds of the way back from the front of the chapel. The stake president replied that the individual was James Womack, whereupon President Kimball said, “He is the man the Lord has selected to be your stake patriarch. Please have him meet with me in the high council room following the meeting.”

Stake president Charles Cagle was startled, for James Womack did not wear the label of a typical man. He had sustained terrible injuries while in combat during World War II. He lost both hands and one arm, as well as most of his eyesight and part of his hearing. Nobody had wanted to let him in law school when he returned, yet he finished third in his class at Louisiana State University. James Womack simply refused to wear the label “Handicapped.”

That evening as President Kimball met with Brother Womack and informed him that the Lord had designated him to be the patriarch, there was a protracted silence in the room. Then Brother Womack said, “Brother Kimball, it is my understanding that a patriarch is to place his hands on the head of the person he blesses. As you can see, I have no hands to place on the head of anyone.”

Brother Kimball, in his kind and patient manner, invited Brother Womack to make his way to the back of the chair on which Brother Kimball was seated. He then said, “Now, Brother Womack, lean forward and see if the stumps of your arms will reach the top of my head.” To Brother Womack’s joy, they touched Brother Kimball, and the exclamation came forth, “I can reach you! I can reach you!”

“Of course you can reach me,” responded Brother Kimball. “And if you can reach me, you can reach any whom you bless. I will be the shortest person you will ever have seated before you.”

President Kimball reported to us that when the name of James Womack was presented to the stake conference, “the hands of the members shot heavenward in an enthusiastic vote of approval.”

The word of the Lord to the prophet Samuel at the time David was designated to be a future king of Israel provided a fitting label for the occasion. It certainly was the thought of each faithful member: “Man looketh on the outward appearance, but the Lord looketh on the heart.” (1 Sam. 16:7.)

Like a golden thread woven through the tapestry of life is the message on the label of a humble heart. It was true of the boy Samuel, it was the experience of Jesus, it was the testimony of Gustav Wacker, it marked the calling of James Womack. May it ever be the label which identifies each of us: “Lord, here am I.” In the name of Jesus Christ, amen.

\newpage
\settalk{Building Your Eternal Home}
\setspeaker{Thomas S. Monson}
\settalkdate{April 1984}
\talkheading{Building Your Eternal Home}{Thomas S. Monson}

When Jesus walked the dusty pathways of towns and villages that we now reverently call the Holy Land and taught His disciples by beautiful Galilee, He often spoke in parables, in language the people understood best. Frequently, He referred to home building in relationship to the lives of those who listened.

He declared: “Every … house divided against itself shall not stand.” (Matt. 12:25.) Later He cautioned: “Behold, mine house is a house of order, … and not a house of confusion.” (D\&C 132:8.)

In a revelation given through the Prophet Joseph Smith at Kirtland, Ohio, December 27, 1832, the Master counseled: “Organize yourselves; prepare every needful thing; and establish a house, even a house of prayer, a house of fasting, a house of faith, a house of learning, a house of glory, a house of order, a house of God.” (D\&C 88:119.)

Where could any of us locate a more suitable blueprint whereby we could wisely and properly build a house to personally occupy throughout eternity?

Such a house would meet the building code outlined in Matthew—even a house built “upon a rock” (Matt. 7:24), a house capable of withstanding the rains of adversity, the floods of opposition, and the winds of doubt everywhere present in our challenging world.

Some might question: “But that revelation was to provide guidance for the construction of a temple. Is it relevant today?”

I would respond: “Did not the Apostle Paul declare, ‘Know ye not that ye are the temple of God, and that the Spirit of God dwelleth in you?’” (1 Cor. 3:16.)

Perhaps if we consider these architectural guidelines on an individual basis, we can more readily appreciate this divine counsel from the Master Builder, the Creator of the world, our Lord and Savior, Jesus Christ.

Our inspired blueprint first cautions that our house should be a house of prayer. The Master taught:

“And when thou prayest, thou shalt not be as the hypocrites are: for they love to pray … , that they may be seen of men. …

“But thou, when thou prayest, … pray to thy Father which is in secret. …

“Use not vain repetitions. …

“After this manner … pray ye: Our Father which art in heaven, Hallowed be thy name.

“Thy kingdom come. Thy will be done in earth, as it is in heaven.

“Give us this day our daily bread.

“And forgive us our debts, as we forgive our debtors.

“And lead us not into temptation, but deliver us from evil: For thine is the kingdom, and the power, and the glory, for ever.” (Matt. 6:5–7, 9–13.)

This element of our blueprint can be taught to children when they are yet young. When our oldest son was about three, he would kneel with his mother and me in our evening prayer. I was serving as the bishop of the ward at the time, and a lovely lady in the ward, Margaret Lister, lay perilously ill with cancer. Each night we would pray for Sister Lister. One evening our tiny son offered the prayer and confused the words of the prayer with a story from a nursery book. He began: “Heavenly Father, please bless Sister Lister, Henny Penny, Chicken Licken, Turkey Lurkey, and all the little folks.” We held back the smiles that evening. Later we were humbled as Margaret Lister sustained a complete recovery. We do not demean the prayer of a child. After all, our children have more recently been with our Heavenly Father than have we.

Let our house be a house of prayer.

Our house should also be a house of fasting. This portion of the blueprint is personified in the account found in Isaiah titled the “True Fast”: “Is not this the fast that I have chosen? …

“To deal thy bread to the hungry, and that thou bring the poor that are cast out to thy house? when thou seest the naked, that thou cover him; and that thou hide not thyself from thine own flesh?”

The reward is then announced: “Then shall thy light break forth as the morning, and thine health shall spring forth speedily: and thy righteousness shall go before thee; the glory of the Lord shall be thy rereward.

“Then shalt thou call, and the Lord shall answer; thou shalt cry, and he shall say, Here I am. …

“And if thou draw out thy soul to the hungry, and satisfy the afflicted soul; then shall thy light rise in obscurity, and thy darkness be as the noonday:

“And the Lord shall guide thee continually, and satisfy thy soul in drought, … and thou shalt be like a watered garden, and like a spring of water, whose waters fail not.” (Isa. 58:6–11.)

Let our house be a house of fasting.

Our house is to be a house of faith. James recorded:

“If any of you lack wisdom, let him ask of God, that giveth to all men liberally, and upbraideth not; and it shall be given him.

“But let him ask in faith, nothing wavering. For he that wavereth is like a wave of the sea driven with the wind and tossed.” (James 1:5–6.)

A practical application of such abiding faith is found in the spirit of Nephi and his stirring declaration: “I will go and do the things which the Lord hath commanded.” (1 Ne. 3:7.) He did not waver; he believed. Is there a counterpart application even today?

Some years ago I accompanied President Hugh B. Brown on a tour of the Samoan Mission. The members and missionaries in American Samoa had advised us that a severe drought had imperiled their water supply to the point that our chapels and our school would of necessity be closed if rain did not soon fall. They asked us to unite our faith with theirs.

Signs of the drought were everywhere as we left the airport at Pago Pago and journeyed to the school at Mapasaga. The sun was shining brightly; not a cloud appeared in the azure blue sky. The members rejoiced as the meeting began. He who offered the opening prayer thanked our Heavenly Father for our safe arrival, knowing that we would somehow bring the desired rainfall. As President Brown rose to speak, the sun was soon shaded by gathering clouds. Then we heard the clap of thunder and saw the flash of lightning. The heavens opened. The rains fell. The drought ended.

Later at the airport, as we prepared for the short flight to Western Samoa, the pilot of the small plane said to the ground crew, “This is the most unusual weather pattern I have ever seen. Not a cloud is in the sky except over the Mormon school at Mapasaga. I don’t understand it!”

President Brown said to me, “Here’s your opportunity. Go help him understand.” I did so.

Our house surely is a house of faith.

Let our house be a house of learning. Said the Lord: “Seek ye out of the best books words of wisdom; seek learning, even by study and also by faith.” (D\&C 88:118.) He counseled: “Come … learn of me … and ye shall find rest unto your souls.” (Matt. 11:28–29.) No other quest for learning promises such a profound reward.

Let our house be a house of learning.

Our house is to be a house of glory. For our house to be such, we need to be square with God, fair with others, and honest with ourselves. One cannot be one person and pretend to be another. Samuel Clemens, better known as Mark Twain, had Huckleberry Finn teach us this vital lesson. Huckleberry Finn is talking:

“It made me shiver. And I about made up my mind to pray, and see if I couldn’t try to quit being the kind of a boy I was and be better. So I kneeled down. But the words wouldn’t come. Why wouldn’t they? It warn’t no use to try and hide it from Him. … I knowed very well why they wouldn’t come. … It was because I was playing double. I was letting on to give up sin, but away inside of me I was holding on to the biggest one of all. I was trying to make my mouth say I would do the right thing and the clean thing, … but deep down in me I knowed it was a lie, and He knowed it. You can’t pray a lie—I found that out.” (The Adventures of Huckleberry Finn, New York: Washington Square Press, Pocket Books, 1973, pp. 271–72; italics added.)

Someone once philosophized, “Consistency, thou art a jewel.” (The Home Book of Quotations, sel. Burton Stevenson, New York: Dodd, Mead, \& Co., 1934, p. 304.) By being consistently good, we can ensure a house of glory.

Our house is to be a house of order. “To every thing there is a season, and a time to every purpose under the heaven” (Eccl. 3:1), advised Ecclesiastes, the Preacher. Such is true in our lives. Let us provide time for family, time for work, time for study, time for service, time for recreation, time for self—but above all, time for Christ.

Then our house will be a house of order.

Finally, let our house be a house of God. Clean thoughts, noble purpose, a willing heart, and ready hands are all features of a house of God. He does not leave us to struggle alone but stands ever ready to help.

A few years ago, I was afforded the privilege to serve as a mission president and became intimately acquainted with more than four hundred missionaries. We had one young missionary who was very ill. After weeks of hospitalization, as the doctor prepared to undertake extremely serious and complicated surgery, he asked that we send for the missionary’s mother and father. He advised there was a possibility the patient would not survive the surgery.

The parents came. Late one evening, the father and I entered a hospital room in Toronto, Canada, placed our hands upon the head of the young missionary, and gave him a blessing. What happened following that blessing was a testimony to me.

The missionary was in a six-bed ward in the hospital. The other beds were occupied by five men with a variety of illnesses. The morning of his surgery, the missionary’s bed was empty. The nurse came into the room with the breakfast these men normally ate. She took a tray over to the patient in bed number one and said, “Fried eggs this morning, and I have an extra portion for you!”

The occupant of bed number one had suffered an accident with his lawnmower. Other than an injured toe, he was well physically. He said to the nurse, “I’ll not be eating this morning.”

“All right, we shall give your breakfast to your partner in bed number two.”

As she approached that patient, he said, “I think I’ll not eat this morning.”

Each of the five men declined breakfast. The young lady exclaimed, “Other mornings you eat us out of house and home, and today not one of you wants to eat! What is the reason?”

Then the man who occupied bed number six answered: “You see, bed number three is empty. Our friend is in the operating room under the surgeon’s hands. He needs all the help he can get. He is a missionary for his church, and while we have been patients in this ward, he has talked to us about the principles of his church—principles of prayer, of faith, of fasting wherein we call upon the Lord for blessings.” He continued, “We don’t know much about the Mormon church, but we have learned a great deal about our friend; and we are fasting for him today.”

The operation was a success. When I attempted to pay the doctor, he countered, “Why, it would be dishonest for me to accept a fee. I have never before performed surgery when my hands seemed to be guided by a Power which was other than my own. No,” he said, “I wouldn’t take a fee for the surgery which Someone on high literally helped me to perform.”

Such is a house of God.

This, then, is our building project. We are master builders of eternal houses, even “temples of God.” (See 1 Cor. 3:16.)

“Organize yourselves; prepare every needful thing; and establish a house, even a house of prayer, a house of fasting, a house of faith, a house of learning, a house of glory, a house of order, a house of God.” (D\&C 88:119.)

Then the Lord, even our building inspector, may say to us, as He said when He appeared to Solomon, a builder of another day: “I have hallowed this house, which thou hast built, to put my name there for ever; and mine eyes and mine heart shall be there perpetually.” (1 Kgs. 9:3.)

May we follow this divinely provided blueprint. May we be successful builders of our eternal homes, is my prayer, in the name of Jesus Christ, amen.

\newpage
\settalk{The Aaronic Priesthood Pathway}
\setspeaker{Thomas S. Monson}
\settalkdate{October 1984}
\talkheading{The Aaronic Priesthood Pathway}{Thomas S. Monson}

Every missionary in the Church is acquainted with the scriptural passage from the book of Amos: “Surely the Lord God will do nothing, but he revealeth his secret unto his servants the prophets.” (Amos 3:7.) Every member of the Church rejoices in singing the favorite hymn:

\begin{verse}
We thank thee, O God, for a prophet\\
To guide us in these latter days.
\end{verse}

(Hymns, no. 196.)

What secrets has the Lord God revealed to His prophet, our beloved leader President Spencer W. Kimball? What counsel would President Kimball provide us tonight, were he here, to guide us in these latter days? Would we listen? Would we obey? Would we be doers of the word, and not hearers only, deceiving our own selves? (See James 1:22.)

Some time ago, as the General Authorities met together on an upper floor of the temple, President Kimball stood and instructed us, saying: “Brethren, of late I have been concerned and troubled by the fact that we do not have sufficient missionaries proclaiming the message of the Restoration. I hear some parents say, ‘We’re letting our son make up his own mind regarding a mission,’ or ‘We hope our son fills a mission because it would be such a growing experience for him.’” He continued: “I have heard some young men say, ‘I think I might serve a mission if I really want to go.’” President Kimball raised his voice, stood on tiptoe—as he is prone to do when anxious to communicate with power a special thought—and said: “It doesn’t really matter whether Mother or Father thinks it might be nice for a son to serve a mission. It doesn’t really matter whether or not John, Bill, and Bob want to go—they must go!” President Kimball then proceeded to point out the missionary obligation each of us has, to repay the sacrifice and service of those missionaries who left home and family and brought the gospel to our parents or grandparents in lands near and far.

I love to read my own grandfather’s missionary journal. His first entries are classics. He wrote: “Today I married in the Salt Lake Temple the girl of my dreams.” The very next night the journal entry read: “Tonight the bishop called at our house. I have been asked to return to Scandinavia for a two- year mission. Of course I will go, and my sweet wife will remain at home and sustain me.” I am grateful for a missionary heritage.

We of the Council of the Twelve have heard President Ezra Taft Benson describe how his father was called to fill a mission. He left behind his wife, his seven children, the farm, and all that he had. Did he lose anything? President Benson tells how his mother would gather the family around the kitchen table and there, by the flickering light of an oil-fueled lamp, read the letters from her husband. Several times during the reading there would be a pause to wipe away the tears which flowed freely. The result? Each of the children later served a mission.

As we strive to respond to President Kimball’s clarion call to missionary service, perhaps we should examine the Aaronic Priesthood pathway which provides the training, quickens the desire, and leads the lad who journeys along it not only to a missionary call but also to temple marriage and, at journey’s end, even to exaltation in the celestial kingdom of God.

It is essential, even critical, that we study the Aaronic Priesthood pathway, since far too many boys falter, stumble, then fall without crossing the finish line into the quorums of the Melchizedek Priesthood. In fact, today, for the first time in the history of the Church, the prospective elders outnumber the holders of the Melchizedek Priesthood, thereby eroding the active priesthood base of the Church and curtailing the activity of loving wives and precious children.

What can we as leaders do to reverse this trend? How can we assure that every boy becomes a finisher? The place to begin is at the headwaters of the Aaronic Priesthood stream. There is an ancient Chinese proverb which purports to correctly determine the sanity of an individual. A person is shown a stream of water flowing into a stagnant pond. He is given a bucket and asked to commence to drain the pond. If he first takes steps to effectively dam the inflow to the pond, he is adjudged sane. If, on the other hand, he ignores the inflow and tries to empty the pond bucket by bucket, he is designated insane.

The best and most effective manner whereby we can solve the challenge of the growth in numbers of prospective elders is to concentrate on the Aaronic Priesthood.

The bishop, by revelation, is the president of the Aaronic Priesthood and is president of the priests in his ward. (See D\&C 107:87–88.) He cannot delegate these God-given responsibilities. However, he can place accountability with his counselors and name as quorum advisers men who can touch the lives of boys—indeed, men who are models to follow. Were I a bishop tonight, I would turn to my second counselor and say: “Brother Balmforth, you have the duty to look after the deacons in the ward. Yours is the task to ensure that every boy is worthy and is ordained a teacher when he reaches his fourteenth birthday.” Then I would address my first counselor with the thought: “Brother Hemingway, yours is the duty to make certain that every teacher is worthy and is ordained a priest when he reaches sixteen. As the bishop, I will assume the task to so labor with the young men who are priests that they are worthy and are ordained elders as they embark on their missions.”

This, then, is our assignment: to save every boy, thereby assuring a worthy husband for each of our young women, strong Melchizedek Priesthood quorums, and a missionary force trained and capable of accomplishing what the Lord expects.

A wise first step is to guide each deacon to a spiritual awareness of the sacredness of his ordained calling. In my life this was accomplished when the bishopric asked that I take the sacrament to a shut-in who lived about a mile from the chapel. That special Sunday morning, as I knocked on the door of Brother Wright and heard his feeble reply, “Come in,” I entered not only his humble cottage but also a room filled with the Spirit of the Lord. I approached his bedside and carefully placed a piece of the bread to his lips. I then held the cup of water, that he might drink. As I departed, I saw him smile as he said, “God bless you, my boy.” And God did bless me with an appreciation for the sacred emblems which continues even today.

Is every ordained teacher given the assignment to home teach? What an opportunity to prepare for a mission. What a privilege to learn the discipline of duty. A boy will automatically turn from concern for self when he is assigned to “watch over” others.

And what of the priests? These young men have the opportunity to bless the sacrament, to continue their home teaching duties, and to participate in the sacred ordinance of baptism.

I remember as a deacon watching the priests as they would officiate at the sacrament table. One priest had a lovely voice and would read the sacrament prayers with clear diction—as though he were competing in a speech contest. The older members of the ward would compliment him on his “golden voice.” I think he became a bit proud. Another priest in the ward had a hearing impediment which caused his speech to be unnatural in its sound. We deacons would twitter at times when Jack would bless the emblems. How we dared do so is beyond me: Jack had hands like a bear and could have crushed any of us. On one occasion Barry with the beautiful voice and Jack with the awkward delivery were assigned together at the sacrament table. The hymn was sung; the two priests broke the bread. Barry knelt to pray, and we closed our eyes. But nothing happened. Soon we deacons opened our eyes to see what was causing the delay. I shall ever remember Barry frantically searching the table for the little white card on which were printed the sacrament prayers. It was nowhere to be found. What to do? Barry’s face turned pink, then crimson, as the congregation began to look in his direction. Then Jack, with that bear-like hand, reached up and gently tugged Barry back to the bench. He, himself, then knelt on the little stool and began to pray: “Oh God, the Eternal Father, we ask thee in the name of thy Son, Jesus Christ, to bless and sanctify this bread to the souls of all those who partake of it. …” He continued the prayer, and the bread was then passed. Jack also blessed the water, and it was passed. What respect we deacons gained that day for Jack who, though handicapped in speech, had memorized the sacred prayers. Barry, too, had a new appreciation for Jack. A lasting bond of friendship had been established.

Beyond the influence of the bishopric and the Aaronic Priesthood quorum advisers is the impact of the home. Help of parents, when enlisted wisely, can frequently make the difference between success and failure. Our recent surveys reveal that the influence of the home surpasses all other factors as a determinant of missionary service and temple marriage.

Not to be overlooked are the strength and influence of devoted Aaronic Priesthood quorum presidencies. The revelations are crystal clear in their meaning: “Verily I say unto you, the duty of a president over the office of a deacon is to preside over twelve deacons, to sit in council with them, and to teach them their duty, edifying one another, as it is given according to the covenants.” (D\&C 107:85.) A similar charge is given to the president of the teachers quorum and to the bishop as president of the quorum of priests. (See D\&C 107:86–88.)

The stake Aaronic Priesthood committee can also provide much needed help. Stake presidents, do you ensure that your high councilors who serve on this most important committee visit the quorums of the Aaronic Priesthood on a continuing and regular basis? Do these brethren know the names of each Aaronic Priesthood boy in the stake? Generalities simply will not do. When we deal in generalities, we will never have a success; but as we deal in specifics, we will rarely have a failure.

I am reminded of the ward presided over by our own Joseph B. Wirthlin. Bishop Wirthlin had a quorum of forty-five priests. All forty-five became elders. All filled missions. The late Elder Alvin R. Dyer presided over a quorum of forty-eight priests. Forty-six of the total served full-time missions, and forty-seven married in the House of the Lord. It can indeed be done. Each boy must be saved.

When I served as a bishop, I noted one Sunday morning that one of our priests was missing from the priesthood meeting. I left the quorum in the care of the adviser and visited Richard’s home. His mother said he was working at the West Temple Garage. I drove to the garage in search of Richard and looked everywhere but could not find him. Suddenly, I had the inspiration to gaze down into the old-fashioned grease pit situated at the side of the station. From the darkness I could see two shining eyes. Then I heard Richard say: “You found me, Bishop! I’ll come up.” He never missed another priesthood meeting.

The family moved, and Richard moved with them. About a year later Bishop Arthur Spencer of the Wells Stake called and said that Richard was responding to a mission call to Mexico and asked if I would accept the family’s invitation to speak at his farewell testimonial. At the meeting, when Richard responded, he mentioned that the turning point in his determination to fill a mission came one Sunday morning—not in the chapel, but as he gazed up from the depths of a dark grease pit and found his quorum president’s outstretched hand.

John Barrie, the Scottish poet, declared: “God gave us memories, that we might have June roses in the December of our lives.” From my experience, some of the most fragrant and beautiful roses anywhere to be found bloom in profusion along the Aaronic Priesthood pathway. On this pathway there are feet to steady, hands to grasp, minds to encourage, hearts to inspire, and souls to save.

I invite each of you men to walk with me, shoulder to shoulder, together with all of the Aaronic Priesthood bearers of the Church, along this priesthood pathway which leads upward and onward toward perfection. In the name of Jesus Christ, amen.

\newpage
\settalk{“The Spirit Giveth Life”}
\setspeaker{Thomas S. Monson}
\settalkdate{April 1985}
\talkheading{“The Spirit Giveth Life”}{Thomas S. Monson}

Recently I visited the Missionary Training Center at Provo, Utah, where missionaries who have been called to serve throughout the world are devotedly learning the fundamentals of the languages spoken by the people to whom they shall teach and testify.

Vaguely familiar to me were the conversations in Spanish, French, German, and Swedish. Totally foreign to me and perhaps to most of the missionaries were the sounds of Japanese, Chinese, and Finnish. One marvels at the devotion and total concentration of these young men and women as they grapple with the unfamiliar and learn the difficult.

I am told that on occasion when a missionary in training feels that the Spanish he is called upon to master appears overwhelming or just too hard to learn, he is placed during the luncheon break next to missionaries studying the complex languages of the Orient. He listens. Suddenly Spanish becomes not too overpowering, and he eagerly returns to his study.

There is one language, however, that is understood by each missionary: the language of the Spirit. It is not learned from textbooks written by men of letters, nor is it acquired through reading and memorization. The language of the Spirit comes to him who seeks with all his heart to know God and to keep His divine commandments. Proficiency in this language permits one to breach barriers, overcome obstacles, and touch the human heart.

The Apostle Paul, in his second epistle to the Corinthians, urges that we turn from the narrow confinement of the letter of the law and seek the open vista of opportunity which the Spirit provides. I love and cherish Paul’s statement: “The letter killeth, but the spirit giveth life.” (2 Cor. 3:6.)

In a day of danger or a time of trial, such knowledge, such hope, such understanding bring comfort to the troubled mind and grieving heart. The entire message of the New Testament breathes a spirit of awakening to the human soul. Shadows of despair are dispelled by rays of hope, sorrow yields to joy, and the feeling of being lost in the crowd of life vanishes with the certain knowledge that our Heavenly Father is mindful of each of us.

The Savior provided assurance of this truth when He taught that even a sparrow shall not fall to the ground unnoticed by our Father. He then concluded the beautiful thought by saying, “Fear ye not therefore, ye are of more value than many sparrows.” (Matt. 10:29–31.)

We live in a complex world with daily challenges. There is a tendency to feel detached—even isolated—from the Giver of every good gift. We worry that we walk alone.

From the bed of pain, from the pillow wet with the tears of loneliness, we are lifted heavenward by that divine assurance and precious promise, “I will not fail thee, nor forsake thee.” (Josh. 1:5.)

Such comfort is priceless as we journey along the pathway of mortality, with its many forks and turnings. Rarely is the assurance communicated by a flashing sign or a loud voice. Rather, the language of the Spirit is gentle, quiet, uplifting to the heart and soothing to the soul.

At times, the answers to our questions and the responses to our daily prayers come to us through silent promptings of the Spirit. As William Cowper wrote:

\begin{verse}
God moves in a mysterious way\\
His wonders to perform;\\
He plants his footsteps in the sea\\
And rides upon the storm. …
\end{verse}

(Hymns, no. 48.)

We watch. We wait. We listen for that still, small voice. When it speaks, wise men and women obey. We do not postpone following promptings of the Spirit.

To address such a sacred subject, may I refer not to the writings of others, but to the actual experiences of my life. I testify to their truth, for I lived them. I share with you today three cherished examples of what President David O. McKay identified as “heart petals”—the language of the Spirit, the promptings from a heavenly source.

First, the inspiration which attends a call to serve.

Second, the gratitude of God for a life well lived.

Third, the knowledge that we do not walk alone.

Every bishop can testify to the promptings which attend calls to serve in the Church. Frequently the call seems to be not so much for the benefit of those to be taught or led as for the person who is to teach or lead.

As a bishop, I worried about any members who were inactive, not attending, not serving. Such was my thought as I drove down the street where Ben and Emily lived. They were older—even in the twilight period of life. Aches and pains of advancing years caused them to withdraw from activity to the shelter of their home—isolated, detached, shut out from the mainstream of daily life and association.

I felt the unmistakable prompting to park my car and visit Ben and Emily, even though I was on the way to a meeting. It was a sunny weekday afternoon. I approached the door to their home and knocked. Emily answered. When she recognized me, her bishop, she exclaimed, “All day long I have waited for my phone to ring. It has been silent. I hoped that the postman would deliver a letter. He brought only bills. Bishop, how did you know today was my birthday?”

I answered, “God knows, Emily, for He loves you.”

In the quiet of the living room, I said to Ben and Emily, “I don’t know why I was directed here today, but our Heavenly Father knows. Let’s kneel in prayer and ask Him why.” This we did, and the answer came. Emily was asked to sing in the choir—even to provide a solo for the forthcoming ward conference. Ben was asked to speak to the Aaronic Priesthood young men and recount a special experience in his life when his safety was assured by responding to the promptings of the Spirit. She sang. He spoke. Hearts were gladdened by the return to activity of Ben and Emily. They rarely missed a sacrament meeting from that day to the time each was called home. The language of the Spirit had been spoken. It had been heard. It had been understood. Hearts were touched and lives saved.

For my second example I turn to the release of a stake president in Star Valley, Wyoming—even the late E. Francis Winters. He had served faithfully for the lengthy term of twenty-three years. Though modest by nature and circumstance, he had been a perpetual pillar of strength to everyone in the valley. On the day of the stake conference, the building was filled to overflowing. Each heart seemed to be saying a silent thank-you to this noble leader who had given so unselfishly of his life for the benefit of others.

As I stood to speak following the reorganization of the stake presidency, I was prompted to do something I had not done before, nor have I done so since. I stated how long Francis Winters had presided in the stake; then I asked all whom he had blessed or confirmed as children to stand and remain standing. Then I asked all those persons whom President Winters had ordained, set apart, personally counseled, or blessed to please stand. The outcome was electrifying. Every person in the audience rose to his feet. Tears flowed freely—tears which communicated better than could words the gratitude of tender hearts. I turned to President and Sister Winters and said, “We are witnesses today of the prompting of the Spirit. This vast throng reflects not only individual feelings but also the gratitude of God for a life well lived.” No person who was in the congregation that day will forget how he felt when he witnessed the language of the Spirit of the Lord.

Finally, I testify that we do not walk alone.

Stan, a dear friend of mine, was taken seriously ill and rendered partially paralyzed. He had been robust in health, athletic in build, and active in many pursuits. Now he was unable to walk or to stand. His wheelchair was his home. The finest of physicians had cared for him, and the prayers of family and friends had been offered in a spirit of hope and trust. Yet Stan continued to lie in the confinement of his bed at the university hospital. He despaired.

Late one afternoon I was swimming at the Deseret Gym, gazing at the ceiling while backstroking width after width. Silently, but ever so clearly, there came to my mind the thought: “Here you swim almost effortlessly, while your friend Stan languishes in his hospital bed, unable to move.” I felt the prompting: “Get to the hospital and give him a blessing.”

I ceased my swimming, dressed, and hurried to Stan’s room at the hospital. His bed was empty. A nurse said he was in his wheelchair at the swimming pool, preparing for therapy. I hurried to the area, and there was Stan, all alone, at the edge of the deeper portion of the pool. We greeted one another and returned to his room, where a priesthood blessing was provided.

Slowly but surely, strength and movement returned to Stan’s legs. First he could stand on faltering feet. Then he learned once again to walk—step by step. Today one would not know that Stan had lain so close to death and with no hope of recovery.

Frequently Stan speaks in Church meetings and tells of the goodness of the Lord to him. To some he reveals the dark thoughts of depression which engulfed him that afternoon as he sat in his wheelchair at the edge of the pool, sentenced, it seemed, to a life of despair. He tells how he pondered the alternative. It would be so easy to propel the hated wheelchair into the silent water of the deep pool. Life would then be over. But at that precise moment he saw me, his friend. That day Stan learned literally that we do not walk alone. I, too, learned a lesson that day: Never, never, never postpone following a prompting.

As we pursue the journey of life, let us learn the language of the Spirit. May we remember and respond to the Master’s gentle invitation: “Behold, I stand at the door, and knock: if any man hear my voice, and open the door, I will come in to him.” (Rev. 3:20.) This is the language of the Spirit. He spoke it. He taught it. He lived it. May each of us do likewise, I pray in the name of Jesus Christ, amen.

\newpage
\settalk{Those Who Love Jesus}
\setspeaker{Thomas S. Monson}
\settalkdate{October 1985}
\talkheading{Those Who Love Jesus}{Thomas S. Monson}

Driving on the modern freeways during the sunshine of summer is often a pleasant experience. Frequently, one can view the grandeur of majestic mountains and the mesmerizing surf of the sea all in a single drive. However, when the traffic is heavy, the mountains and seas are set aside, and concentration is focused on the car ahead. Such was the occasion when I read with keen interest the words of a bumper sticker readily visible on the highly polished chrome bumper of a car which was weaving in and out of the traffic stream. The words were these: “Honk if you love Jesus.” No one honked. Perhaps each was disturbed by the thoughtless and rude actions of the offending driver. Then, again, would honking be an appropriate manner in which to show one’s love for the Son of God, the Savior of the world, the Redeemer of all mankind? Such was not the pattern provided by Jesus of Nazareth.

The importance of demonstrating daily a true and an abiding love was convincingly taught by the Master when the inquiring lawyer stepped forward and boldly asked him, “Master, which is the great commandment in the law?”

Matthew records that “Jesus said unto him, Thou shalt love the Lord thy God with all thy heart, and with all thy soul, and with all thy mind.

“This is the first and great commandment.

“And the second is like unto it, Thou shalt love thy neighbour as thyself.” (Matt. 22:36–39.)

Mark concludes the account with the Savior’s statement, “There is none other commandment greater than these.” (Mark 12:31.)

His answer could not be faulted. His very actions gave credence to His words. He demonstrated genuine love of God by living the perfect life, by honoring the sacred mission that was His. Never was He haughty. Never was He puffed up with pride. Never was He disloyal. Ever was He humble. Ever was He sincere. Ever was He true.

Though He was led up of the Spirit into the wilderness to be tempted by that master of deceit, even the devil; though He was physically weakened from fasting forty days and forty nights and was an hungered; yet when the evil one proffered Jesus the most alluring and tempting proposals, He gave to us a divine example of true love of God by refusing to deviate from what He knew was right. (See Matt. 4:1–11.)

Jesus, throughout His ministry, blessed the sick, restored sight to the blind, made the deaf to hear and the maimed to walk. He taught forgiveness by forgiving. He taught compassion by being compassionate. He taught devotion by giving of Himself. Jesus taught by example.

As we survey the life of our Lord, each of us could echo the words of the well-known hymn:

\begin{verse}
I stand all amazed at the love Jesus offers me,\\
Confused at the grace that so fully he proffers me.\\
I tremble to know that for me he was crucified,\\
That for me, a sinner, he suffered, he bled and died.
\end{verse}

(“I Stand All Amazed,” Hymns, 1985, no. 193.)

To demonstrate our gratitude, is it required that we, too, lay down our lives as did He? Some have.

In the beautiful city of Melbourne, Australia, there is situated in an impressive setting a historic war memorial. As one walks through the memorial’s silent corridors, one sees tablets of marble that note the deeds of valor and acts of courage of those who made the supreme sacrifice. One can almost hear the roar of the cannon, the scream of the rocket, the cry of the wounded. One can feel the exhilaration of victory and, at the same time, sense the despair of defeat.

In the center of the main hall, inscribed for all to see, is the message of the memorial. The skylight overhead permits easy reading. The words almost stand up and speak: “Greater love hath no man than this, that a man lay down his life for his friends.” (John 15:13.)

Today, the challenge which we face and must meet is not that we should go forth on the battlefield of war and lay down our lives. Rather, it is that we, on the battlefield of life, so live and serve that our lives and actions reflect a true love of God, of His Son, Jesus Christ, and of our fellowmen. This is not accomplished by clever signs printed on bumper stickers affixed to automobiles.

Jesus teaches us: “If ye love me, keep my commandments. …

“He that hath my commandments, and keepeth them, he it is that loveth me: and he that loveth me shall be loved of my Father, and I will love him, and will manifest myself to him.” (John 14:15, 21.)

Years ago we danced to a popular song, the words of which were, “It’s easy to say I love you, easy to say I’ll be true, easy to say these simple things, but prove it by the things you do.”

From our lessons learned in Primary we remember the poem entitled “Which Loved Best?”

\begin{verse}
“I love you, Mother,” said little [John];\\
Then, forgetting his work, his cap went on,\\
And he was off to the garden swing,\\
And left her the water and wood to bring.
\end{verse}

(Joy Allison, The World’s Best Loved Poems, New York: Harper and Row, 1955, pp. 243–44.)

Years pass. Childhood vanishes. Truth remains. The transition from Primary’s poems to today’s truths is not difficult. True love continues to be an outward expression of an inward conviction.

Today, on a gentle rise in the historic city of Freiberg, German Democratic Republic, there stands a beautiful, dedicated temple of God. The temple provides the ultimate—even the eternal—blessings of a loving Heavenly Father to His faithful Saints.

Ten years ago, on a Sunday morning, April 27, 1975, I stood on an outcropping of rock situated between the cities of Dresden and Meissen, high above the Elbe River. I responded to the promptings of the Holy Spirit and offered a prayer of dedication on that land and its people. That prayer noted the faith of the members. It emphasized the tender feelings of many hearts filled with an overwhelming desire to obtain temple blessings. A plea for peace was expressed. Divine help was requested. I voiced the words, “Dear Father, let this be the beginning of a new day for the members of Thy church in this land.”

Suddenly, from far below in the valley, a bell in a church steeple began to chime and the shrill crow of a rooster broke the morning silence, each heralding the commencement of a new day. Though my eyes were closed, I felt a warmth from the sun’s rays reaching my face, my hands, my arms. How could this be? An incessant rain had been falling all morning. At the conclusion of the prayer, I gazed heavenward. I noted a ray of sunshine which penetrated an opening in the heavy clouds, a ray which engulfed the spot where our small group stood. From that moment I knew divine help was at hand.

Full cooperation of government officials was forthcoming. President Spencer W. Kimball and his counselors provided enthusiastic approval. A temple was planned, a site selected, ground-breaking services held, and construction commenced. At the time of dedication, the attention of the international press was focused on this temple in its unusual setting. Words like “How?” and “Why?” were voiced frequently. This was particularly in evidence during the public open house, when 89,872 persons visited the temple. At times the waiting period stretched to three hours, occasionally in the rain. None wavered. All were shown God’s house.

During the actual dedicatory services when President Gordon B. Hinckley offered the dedicatory prayer, hymns of praise, testimonies of truth, tears of gratitude, and prayers of thanksgiving marked the historic event. To understand how, to comprehend why, it is necessary to know the faith, the devotion, the love of the members of the Church in that nation. Though fewer than five thousand in number, the activity levels exceed those found anywhere else in the world.

During the many years I have served on assignment in that area, I have noted the absence of spacious chapels with multiple teaching stations and grounds featuring the greenery of lawns and the blossoms of flowers. The meetinghouse libraries, as well as the personal libraries of our members, consist only of the standard works, a hymnbook, and one or two other volumes. These books do not remain on bookcase shelves. Their teachings are engraved on the hearts of members. They are displayed in their daily lives. Service is a privilege. A branch president, forty-two years of age, has served in his calling for twenty-one years—half his life. Never a complaint—just gratitude. In Leipzig, when the meetinghouse furnace failed one cold winter day, the meetings were not dismissed. Rather, the members met in the chill of the unheated building, sitting shoulder to shoulder, wearing their coats, singing the hymns of Zion and worshiping Him who counseled, “Be not weary in well doing,” “Follow me,” “Be thou humble; and the Lord thy God shall lead thee by the hand, and give thee answer to thy prayers.” (2 Thes. 3:13; Matt. 4:19; D\&C 112:10.)

The Apostle Paul taught the Corinthians, “If any man love God, the same is known of him.” (1 Cor. 8:3.) The love which these faithful members have for God, for His Son, Jesus Christ, and for His everlasting gospel is confirmed by their very lives. It is reminiscent of the love demonstrated by the brother of Jared as described in the Book of Mormon. The blessings of a loving, caring, and just Heavenly Father simply could not be withheld. Faith preceded the miracle. Eternal ordinances are now performed. Everlasting covenants are now made. The love of God has again blessed His people.

For those who love Jesus, these prophetic words have sublime meaning:

“Hear, O ye heavens, and give ear, O earth, and rejoice ye inhabitants thereof, for the Lord is God, and beside him there is no Savior.

“Great is his wisdom, marvelous are his ways. …

“His purposes fail not. …

“For thus saith the Lord—I, the Lord, am merciful and gracious unto those who fear me, and delight to honor those who serve me in righteousness and in truth unto the end.

“Great shall be their reward and eternal shall be their glory.” (D\&C 76:1–3, 5–6.)

Such is the blessing reserved for those who love Jesus. May each of us qualify for this great reward, this eternal glory, I pray in the name of Jesus Christ, whom I love and of whom I testify, amen.

\newpage
\settalk{A Provident Plan—A Precious Promise}
\setspeaker{Thomas S. Monson}
\settalkdate{April 1986}
\talkheading{A Provident Plan—A Precious Promise}{Thomas S. Monson}

Today, April 6, 1986, is a day of history. One hundred fifty-six years ago The Church of Jesus Christ of Latter-day Saints was organized. Numbers were few. Circumstances were modest. But the future beckoned. In solemn assembly this afternoon, President Ezra Taft Benson will be sustained by our hearts and souls, as well as by our uplifted hands, as the thirteenth President of the Church. Prayers of thanksgiving will be offered, words of wisdom provided, and songs of praise sung. Strains of “We Thank Thee, O God, for a Prophet” and “How Firm a Foundation” will emanate from this Tabernacle and reverberate throughout the lands of the earth.

It was fifty years ago this very day that the prophets of God outlined the general principles which became the “firm foundation” of the Church welfare plan. In a specially called and momentous meeting presided over by President Heber J. Grant and his counselors—J. Reuben Clark, Jr., and David O. McKay—watershed statements were presented and heaven-inspired counsel provided which have endured the passage of time, which have been rendered valid by the verdict of history, and which bear the seal of God’s approval.

On that occasion, President David O. McKay declared, “This organization is established by divine revelation, and there is nothing else in all the world that can so effectively take care of its members.” (In Henry D. Taylor, “The Church Welfare Plan,” 1984, p. 26.)

President J. Reuben Clark set the tone for the launching of this inspired effort by counseling: “[The Lord] has given us the spirituality. He has given us the actual command. … The eyes of the world are upon us. … May the Lord bless you, give us courage, give us wisdom, give us vision to carry out this great work.” (Taylor, p. 27.)

Fifty years have come and gone. Economic cycles have run their course. Societal changes have been numerous. The Church has expanded beyond the valleys of the mountains to the uttermost reaches of the earth. Membership is measured in millions. The word of God, provided on that historic day, is as an island of constancy in a sea of change.

Let us, for a moment, review the moorings, the underpinnings, even the foundation of the welfare program. Said the First Presidency in that year of announcement: “Our primary purpose was to set up, insofar as it might be possible, a system under which the curse of idleness would be done away with, the evils of a dole abolished, and independence, industry, thrift and self respect be once more established amongst our people. The aim of the Church is to help the people to help themselves.” (In Conference Report, Oct. 1936, p. 3.)

The holy scriptures leave no doubt concerning the responsibility to care for the poor, the needy, the downtrodden. The organization has been perfected, the duties defined, and the guidelines given.

I am profoundly grateful to my Heavenly Father for the privilege which has been mine to be tenderly taught and constantly counseled by the prophets of the program.

As a publisher and printer, I had the opportunity to assist President J. Reuben Clark in the preparation of his manuscript which became the monumental book Our Lord of the Gospels. What a blessing was mine to learn daily at the feet of such a master teacher and principle architect of the welfare program. Knowing that I was a newly appointed bishop presiding over a difficult ward, he emphasized the need for me to know my people, to understand their circumstances, and to minister to their needs. One day he recounted the example of the Savior as recorded in the Gospel of Luke:

“And it came to pass … that he went into a city called Nain; and many of his disciples went with him. …

“When he came nigh to the gate of the city, behold, there was a dead man carried out, the only son of his mother, and she was a widow. …

“And when the Lord saw her, he had compassion on her, and said unto her, Weep not.

“And he came and touched the bier. … And he said, Young man, I say unto thee, Arise.

“And he that was dead sat up, and began to speak. And he delivered him to his mother.” (Luke 7:11–15.)

When President Clark closed the Bible, I noticed that he was weeping. In a quiet voice, he said, “Tom, be kind to the widow and look after the poor.”

On one occasion, President Harold B. Lee, who was a stake president in the area where I was born and reared and later presided as a bishop, spoke movingly to the Aaronic Priesthood concerning how the priesthood might prepare for its role in caring for the poor. He stood at the pulpit, took the Book of Mormon in hand, and opened it to the seventeenth chapter of Alma. He then read to us concerning the sons of Mosiah:

“Now these sons of Mosiah were with Alma at the time the angel first appeared unto him; therefore Alma did rejoice exceedingly to see his brethren; and what added more to his joy, they were still his brethren in the Lord; yea, and they had waxed strong in the knowledge of the truth; for they were men of a sound understanding and they had searched the scriptures diligently, that they might know the word of God.

“But this is not all; they had given themselves to much prayer, and fasting; therefore they had the spirit of prophecy, and the spirit of revelation, and when they taught, they taught with power and authority of God.” (Alma 17:2–3.)

We had been given our pattern, provided by an inspired teacher. Reverently, he closed the covers of this sacred scripture. Like President Clark, he too had tears in his eyes.

Just a few days ago I visited with President Marion G. Romney, known throughout the Church for his ardent advocacy and knowledge of the welfare program. We spoke of the beautiful passage from Isaiah concerning the true fast:

“Is it not to deal thy bread to the hungry, and that thou bring the poor that are cast out to thy house? when thou seest the naked, that thou cover him; and that thou hide not thyself from thine own flesh?” (Isa. 58:7.)

As did President Clark, as did President Lee, President Romney wept as he spoke.

Appearing as a golden thread woven through the tapestry of the welfare program is the truth taught by the Apostle Paul: “The letter killeth, but the spirit giveth life.” (2 Cor. 3:6.)

President Ezra Taft Benson frequently counsels us: “Remember, Brethren, in this work it is the Spirit that counts.”

What has the Lord said about the spirit of this work? In a revelation given to the Prophet Joseph at Kirtland, Ohio, in June of 1831, He declared: “Remember in all things the poor and the needy, the sick and the afflicted, for he that doeth not these things, the same is not my disciple.” (D\&C 52:40.)

In that marvelous message delivered by King Benjamin, as recorded in the Book of Mormon, we read: “For the sake of retaining a remission of your sins from day to day, that ye may walk guiltless before God—I would that ye should impart of your substance to the poor, every man according to that which he hath, such as feeding the hungry, clothing the naked, visiting the sick and administering to their relief, both spiritually and temporally.” (Mosiah 4:26.)

When we depart from the Lord’s way in caring for the poor, chaos comes. Said John Goodman, president of the National Center for Political Analysis, as reported this year in a Dallas, Texas, newspaper:

“The USA’s welfare system is a disaster. It is creating poverty, not destroying it. It subsidizes divorce, unwed teenage pregnancy, the abandonment of elderly parents by their children, and the wholesale dissolution of the family. The reason? We pay people to be poor. Private charities have always been better at providing relief where it is truly needed.”

In 1982 it was my privilege to serve as a member of President Ronald Reagan’s Task Force on Private Sector Initiatives. Meeting in the White House with prominent leaders assembled from throughout the nation, President Reagan paid tribute to the welfare program of the Church. He observed: “Elder Monson is here representing The Church of Jesus Christ of Latter-day Saints. If, during the period of the Great Depression, every church had come forth with a welfare program founded on correct principles as his church did, we would not be in the difficulty in which we find ourselves today.” President Reagan praised self-sufficiency; lauded our storehouse, production, and distribution system; and emphasized family members assisting one another. He urged that in our need we turn not to government but rather to ourselves.

On another occasion in the White House, I was asked to present to a gathering of America’s religious leaders an example of our welfare program in action. I could have chosen many illustrations, but selected as typical our response to the Teton Dam disaster in Idaho. The result was dramatic. As the First Presidency stated fifty years ago, “The eyes of the world are upon us.” While this is a most important consideration, let us particularly remember that the eyes of God are similarly focused. What might He observe?

Are we generous in the payment of our fast offerings? That we should be so was taught by President Spencer W. Kimball, who urged that “instead of the amount saved by our two or more meals of fasting, perhaps much, much more—ten times more [be given] when we are in a position to do it.” (Ensign, Nov. 1977, p. 79.)

Are we prepared for the emergencies of our lives? Are our skills perfected? Do we live providently? Do we have on hand our reserve supply? Are we obedient to the commandments of God? Are we responsive to the teachings of prophets? Are we prepared to give of our substance to the poor, the needy? Are we square with the Lord?

As we look back through fifty years and reflect on the development of the welfare program, as we look forward to the years ahead, let us remember the place of the priesthood, the role of the Relief Society, and the involvement of the individual. Help from heaven will be ours.

On a cold winter’s night in 1951, there was a knock at my door. A German brother from Ogden, Utah, announced himself and said, “Are you Bishop Monson?” I answered in the affirmative. He began to weep and said, “My brother, his wife, and family are coming here from Germany. They are going to live in your ward. Will you come with us to see the apartment we have rented for them?”

On the way to the apartment, he told me he had not seen his brother for many years. Through the holocaust of World War II, his brother had been faithful to the Church, once serving as a branch president before the war took him to the Russian front.

I observed the apartment. It was cold and dreary. The paint was peeling, the wallpaper soiled, the cupboards empty. A forty-watt bulb, suspended from the living room ceiling, revealed a linoleum floor covering with a large hole in the center. I was heartsick. I thought, “What a dismal welcome for a family which has endured so much.”

My thoughts were interrupted by the brother’s statement, “It isn’t much, but it’s better than they have in Germany.” With that, the key to the apartment was left with me, along with the information that the family would arrive in Salt Lake City in three weeks—just two days before Christmas.

Sleep was slow in coming to me that night. The next morning was Sunday. In our ward welfare committee meeting, one of my counselors said, “Bishop, you look worried. Is something wrong?”

I recounted to those present my experience of the night before, revealing the details of the uninviting apartment. There were a few moments of silence. Then Brother Eardley, the group leader of the high priests, said, “Bishop, did you say that apartment was inadequately lighted and that the kitchen appliances were in need of replacement?” I answered in the affirmative. He continued, “I am an electrical contractor. Would you permit the high priests of this ward to rewire that apartment? I would also like to invite my suppliers to contribute a new stove and a new refrigerator. Do I have your permission?”

I answered with a glad “Certainly.”

Then Brother Balmforth, the seventies president, responded, “Bishop, as you know, I’m in the carpet business. I would like to invite my suppliers to contribute some carpet, and the seventies can easily lay it and eliminate that worn linoleum.”

Then Brother Bowden, the president of the elders quorum, spoke up. He was a painting contractor. He said, “I’ll furnish the paint. May the elders paint and wallpaper that apartment?”

Sister Miller, the Relief Society president, was next to speak. “We in the Relief Society cannot stand the thought of empty cupboards. May we fill them?”

The three weeks which followed are ever to be remembered. It seemed that the entire ward joined in the project. The days passed, and at the appointed time, the family arrived from Germany. Again at my door stood the brother from Ogden. With an emotion-filled voice, he introduced to me his brother, his brother’s wife, and their family. Then he asked, “Could we go visit the apartment?” As we walked up the staircase leading to the apartment, he repeated, “It isn’t much, but it’s more than they have had in Germany.” Little did he know what a transformation had taken place and that many who had participated were inside waiting for our arrival.

The door opened to reveal a newness of life. We were greeted by the aroma of freshly painted woodwork and newly papered walls. Gone was the forty-watt bulb, along with the worn linoleum it had illuminated. We stepped on carpet deep and beautiful. A walk to the kitchen presented to our view a new stove and new refrigerator. The cupboard doors were still open; however, they now revealed every shelf filled with food. As usual, the Relief Society had done its work.

In the living room, we began to sing Christmas hymns. We sang “Silent night! Holy night! All is calm, all is bright.” (Hymns, 1985, no. 204.) We sang in English; they sang in German. At the conclusion, the father, realizing that all of this was his, took me by the hand to express his thanks. His emotion was too great. He buried his head in my shoulder and repeated the words, “Mein Bruder, mein Bruder, mein Bruder.”

It was time to leave. As we walked down the stairs and out into the night air, snow was falling. Not a word was spoken. Finally, a young girl asked, “Bishop, I feel better than I have ever felt before. Can you tell me why?”

I responded with the words of the Master: “Inasmuch as ye have done it unto one of the least of these my brethren, ye have done it unto me.” (Matt. 25:40.) Suddenly there came to mind the words from “O Little Town of Bethlehem”:

\begin{verse}
How silently, how silently,\\
The wondrous gift is giv’n!\\
So God imparts to human hearts\\
The blessings of his heav’n.
\end{verse}

(Hymns, 1985, no. 208.)

Silently, wondrously, His gift had been given. Lives were blessed, needs were met, hearts were touched, and souls were saved. A provident plan had been followed. A precious promise had been fulfilled.

I testify that God lives, that Jesus is the Christ, that we are led by a prophet, that sacrifice does indeed bring forth the blessings of heaven. In the name of Jesus Christ, amen.

\newpage
\settalk{The Call of Duty}
\setspeaker{Thomas S. Monson}
\settalkdate{April 1986}
\talkheading{The Call of Duty}{Thomas S. Monson}

Whenever I have the privilege to attend this, the general priesthood meeting of the Church, I reflect on the teachings of some of the most noble of God’s leaders who have stood at this pulpit and who, from the brilliance of their minds, from the depths of their souls, and from the warmth of their hearts, have given us direction. President J. Reuben Clark, Jr., was such a man. Time and again, his fervent plea was for the priesthood of God to be united. Citing the teachings of Jesus, he inevitably admonished us, “Be one; and if ye are not one ye are not mine.” (D\&C 38:27.)

It was my great privilege to know President Clark rather well. I was his printer. On occasion, he would share with me some of his most intimate thoughts, even those scriptures around which he tailored his teachings and lived his life. Late one evening I delivered some press proofs to his office situated in his home at 80 D Street here in Salt Lake City. President Clark was reading from Ecclesiastes. He was in a quiet and reflective mood. He sat back from his large desk, which was stacked with books and papers. He held the scriptures in his hand, lifted his eyes from the printed page, and read aloud to me: “Let us hear the conclusion of the whole matter: Fear God, and keep his commandments: for this is the whole duty of man.” (Eccl. 12:13.) He exclaimed, “A treasured truth! A profound philosophy!” Through the years that conversation has remained bright in my memory. I love, I cherish the noble word duty.

The legendary General Robert E. Lee of American Civil War fame declared, “Duty is the sublimest word in our language. … You cannot do more. You should never wish to do less.”

From that same hour of history, as Abraham Lincoln left the people of Springfield to take over the nation’s presidency, he said, “Let us have faith that right makes might, and in that faith let us to the end dare to do our duty as we understand it.” (Address, Cooper Union, New York, 27 Feb. 1860.)

Time marches on. Duty keeps cadence with that march. Duty does not dim nor diminish. Catastrophic conflicts come and go, but the war waged for the souls of men continues without abatement. Like a clarion call comes the word of the Lord to you, to me, and to priesthood holders everywhere: “Wherefore, now let every man learn his duty, and to act in the office in which he is appointed, in all diligence.” (D\&C 107:99.)

The call of duty came to Adam, to Noah, to Abraham, to Moses, to Samuel, to David. It came to the Prophet Joseph Smith and to each of his successors, even to President Ezra Taft Benson. The call of duty came to the boy Nephi. Listen to his words:

“And it came to pass that I, Nephi, returned from speaking with the Lord, to the tent of my father.

“And it came to pass that he spake unto me, saying: Behold I have dreamed a dream, in the which the Lord hath commanded me that thou and thy brethren shall return to Jerusalem.

“For behold, Laban hath the record of the Jews and also a genealogy of thy forefathers, and they are engraven upon plates of brass.

“Wherefore, the Lord hath commanded me that thou and thy brothers should go unto the house of Laban, and seek the records, and bring them down hither into the wilderness.

“And now, behold thy brothers murmur, saying it is a hard thing which I have required of them; but behold I have not required it of them, but it is a commandment of the Lord.

“Therefore go, my son, and thou shalt be favored of the Lord, because thou hast not murmured.

“And it came to pass that I, Nephi, said unto my father: I will go and do the things which the Lord hath commanded, for I know that the Lord giveth no commandments unto the children of men, save he shall prepare a way for them that they may accomplish the thing which he commandeth them.” (1 Ne. 3:1–7.)

When that same call comes to you and to me, what will be our response? Will we murmur, as did Laman and Lemuel, and say, “This is a hard thing required of us”? Or will we, with Nephi, individually declare, “I will go. I will do”?

Ofttimes the wisdom of God appears as foolishness to men, but the greatest single lesson we can learn in mortality is that, when God speaks and a man obeys, man will always be right.

President John Taylor cautioned us, “If you do not magnify your calling, God will hold you responsible for those you might have saved, had you done your duty.”

The call of duty came to John E. Page when the Prophet Joseph Smith extended to him a call to serve as a missionary. John E. Page “murmured” and responded, “Brother Joseph, I can’t go on a mission to Canada. I don’t even have a coat to wear.”

The Prophet Joseph removed his own coat, handed it to Brother Page, and said, “Here, take this and the Lord will bless you.” John E. Page went on that mission to Canada and, during a two-year period, walked five thousand miles and baptized six hundred people. (See Andrew Jenson, “John E. Page,” The Historical Record, 5:57.)

A famed minister observed, “Men will work hard for money. Men will work harder for other men. But men will work hardest of all when they are dedicated to a cause. Until willingness overflows obligation, men fight as conscripts rather than following the flag as patriots. Duty is never worthily performed until it is performed by one who would gladly do more, if only he could.”

\begin{verse}
I slept and dreamt\\
That life was joy.\\
I awoke and saw\\
That life was duty.\\
I acted, and behold—\\
Duty was joy.
\end{verse}

(Rabindranath Tagore)

Robert Louis Stevenson reminded us: “I know what pleasure is, for I have done good work.”

The call of duty can come quietly as we who hold the priesthood respond to the assignments we receive. President George Albert Smith, that modest yet effective leader, declared, “It is your duty first of all to learn what the Lord wants and then by the power and strength of your holy priesthood to so magnify your calling in the presence of your fellows that the people will be glad to follow you.” (Church News, 7 Sept. 1968, p. 15.)

What does it mean to magnify a calling? It means to build it up in dignity and importance, to make it honorable and commendable in the eyes of all men, to enlarge and strengthen it, to let the light of heaven shine through it to the view of other men. And how does one magnify a calling? Simply by performing the service that pertains to it. An elder magnifies the ordained calling of an elder by learning what his duties as an elder are and then by doing them. As with an elder, so with a deacon, a teacher, a priest, a bishop, and each who holds office in the priesthood.

In 1950 the call of duty came to me as a bishop. The responsibilities were many and varied. The Doctrine and Covenants provided a sure guide. The words of the Apostle Paul to Timothy pertaining to the office of a bishop were sobering. The General Handbook was helpful. The principal areas of administration were spelled out by leaders, both stake and general: The bishop (1) is the father of the ward; (2) is the president of the Aaronic Priesthood; (3) provides for the poor, the needy; (4) is responsible for keeping proper records; and (5) is the common judge in Israel.

Then came an unusual assignment from Church headquarters. Bishops were to provide each serviceman a subscription to the Church News and the Improvement Era and were to write a personal letter to every serviceman each month. The Korean War was raging. Our ward had twenty-three members in uniform. The priesthood quorums, with effort, supplied the funds for the subscriptions to the publications. Since I had served in the Navy in World War II, I knew the importance of a letter from home. I began the task, even the duty, to write twenty-three personal letters each month. After all these years, I still have copies of many of my letters and the responses received. Tears come easily when these letters are reread. It is a joy to learn again of a soldier’s pledge to live the gospel, a sailor’s decision to keep faith with his family.

One evening I handed to a lady in the ward the stack of twenty-three letters for the current month. Her assignment was to handle the mailing and to maintain the constantly changing address file. She glanced at one envelope and, with a smile, asked, “Bishop, don’t you ever get discouraged? Here is another letter to Brother Bryson. This is the seventeenth letter you have sent to him without a reply.”

I responded, “Well, maybe this will be the month.” And it was. His reply is a keepsake, a literal treasure. It was postmarked “APO San Francisco.” He was serving far away on a distant shore, isolated, homesick, alone. He wrote: “Dear Bishop, I ain’t much at writin’ letters. [I could have told him that seventeen months earlier.] Thank you for the Church News and magazines, but most of all thank you for the personal letters. I have turned over a new leaf. I have been ordained a priest in the Aaronic Priesthood. My heart is full. I am a happy man.”

My brethren, Brother Bryson was no happier than was his bishop. I had learned the practical application of the adage, “Do your duty; that is best. Leave unto the Lord the rest.”

Years later, while attending the Salt Lake Cottonwood Stake when Elder James E. Faust served as president, I related that account in an effort to encourage attention to our servicemen. After the meeting, a fine-looking young man came forward. He took my hand in his and asked, “Bishop Monson, do you remember me?”

I replied, “Brother Bryson! How are you? What are you doing in the Church?”

With warmth and obvious pride, he responded, “I’m fine. I serve in the presidency of my elders quorum. Thank you again for your concern for me and the personal letters which you sent and which I treasure.”

\begin{verse}
Father, where shall I work today?\\
And my love flowed warm and free.\\
Then He pointed out a tiny spot\\
And said, “Tend that for me.”\\
I answered quickly, “Oh no; not that!\\
Why, no one would ever see,\\
No matter how well my work was done;\\
Not that little place for me.”\\
And the word He spoke, it was not stern;\\
He answered me tenderly:\\
“Ah, little one, search that heart of thine.\\
Art thou working for them or for me?\\
Nazareth was a little place,\\
And so was Galilee.”
\end{verse}

(Meade McGuire.)

Brethren, let us learn our duty. Let us, in the performance of our duty, follow in the footsteps of the Master. As you and I walk the pathway Jesus walked, let us listen for the sound of sandaled feet. Let us reach out for the Carpenter’s hand. Then we shall come to know Him. He may come to us as one unknown, without a name, as by the lakeside He came to those men who knew Him not. He speaks to us the same words, “Follow thou me” (John 21:22), and sets us to the task which He has to fulfill for our time. He commands, and to those who obey Him, whether they be wise or simple, He will reveal Himself in the toils, the conflicts, the sufferings that they shall pass through in His fellowship; and they shall learn by their own experience who He is.

We will discover He is more than the Babe in Bethlehem, more than the carpenter’s son, more than the greatest teacher ever to live. We will come to know Him as the Son of God, our Savior and our Redeemer. When to Him came the call of duty, He answered, “Father, thy will be done, and the glory be thine forever.” (Moses 4:2.) May we do likewise I pray, in the name of Jesus Christ, amen.

\newpage
\settalk{Courage Counts}
\setspeaker{Thomas S. Monson}
\settalkdate{October 1986}
\talkheading{Courage Counts}{Thomas S. Monson}

Tonight, those who hold the priesthood fill the Tabernacle on Temple Square, have overflowed to the adjacent Assembly Hall, and are assembled in chapels and halls ranging in size from the mammoth Marriott Center at Brigham Young University to the smallest building located many miles away. All of you have come to be uplifted, to be instructed, to be inspired. A favorite word of my nine-year-old granddaughter describes the responsibility to speak to such a vast throng: awesome.

I seek your prayers; I need your faith; I petition our Heavenly Father for that noble attribute of courage, for I know courage counts!

This truth came to me in a most vivid and dramatic manner some thirty-one years ago. I was serving as a bishop. The general session of our stake conference was being held in the Assembly Hall. Our stake presidency was to be reorganized. The Aaronic Priesthood, including members of bishoprics, were providing the music for the conference. As we concluded singing our first selection, President Joseph Fielding Smith, our conference visitor, stepped to the pulpit and read for sustaining approval the names of the new stake presidency. I am confident that the other members of the stake presidency had been made aware of their callings, but I had not. After reading my name, President Smith announced: “If Brother Monson is willing to respond to this call, we shall be pleased to hear from him now.”

As I stood at the pulpit and gazed out on that sea of faces, I remembered the song we had just sung. Its title was “Have Courage, My Boy, to Say No.” That day I selected as my acceptance theme, “Have Courage, My Boy, to Say Yes.”

Life’s journey is not traveled on a freeway devoid of obstacles, pitfalls, and snares. Rather, it is a pathway marked by forks and turnings. Decisions are constantly before us. To make them wisely, courage is needed: the courage to say no, the courage to say yes. Decisions do determine destiny.

The call for courage comes constantly to each of us. It has ever been so, and so shall it ever be. The battlefields of war witness acts of courage. Some are printed on pages of books or contained on rolls of film, while others are indelibly impressed on the human heart.

The courage of a military leader was recorded by a young infantryman wearing the gray uniform of the Confederacy during America’s Civil War. He describes the influence of General J.E.B. Stuart in these words: “At a critical point in the battle, he leaped his horse over the breastworks near my company, and when he had reached a point about the center of the brigade, while the men were loudly cheering him, he waved his hand toward the enemy and shouted, ‘Forward men. Forward! Just follow me!’

“The men were wild with enthusiasm. With courage and resolution, they poured over the breastworks after him like a raging torrent, and the objective was seized and held” (Emory M. Thomas, Bold Dragoon: The Life of J.E.B. Stuart, New York: Harper and Row, 1986).

At an earlier time, and in a land far distant, another leader issued the same plea: “Follow me” (Matt. 4:19). He was not a general of war. Rather, He was the Prince of Peace, the Son of God. Those who followed Him then, and those who follow Him now, win a far more significant victory, with consequences that are everlasting. But the need for courage is constant. Courage is ever required.

The holy scriptures portray the evidence of this truth. Joseph, son of Jacob, the same who was sold into Egypt, demonstrated the firm resolve of courage when to Potiphar’s wife, who attempted to seduce him, he declared: “How … can I do this great wickedness, and sin against God? And … he hearkened not unto her” and got out (Gen. 39:9–10).

In our day, a father applied this example of courage to the lives of his children by declaring: “If you ever find yourself where you shouldn’t ought to be, get out!”

The prophet Daniel demonstrated supreme courage by standing up for what he knew to be right and by demonstrating the courage to pray, though threatened by death were he to do so (see Dan. 6).

Courage characterized the life of Abinadi, as shown in the Book of Mormon by his willingness to offer his life rather than to deny the truth (see Mosiah 11:20; 17:20).

Who can help but be inspired by the lives of the two thousand stripling sons of Helaman who taught and demonstrated the need of courage to follow the teachings of parents, the courage to be chaste and pure? (see Alma 56).

Perhaps each of these accounts is crowned by the example of Moroni, who had the courage to persevere to the end in righteousness (see Moro. 1–10).

All were fortified by the words of Moses: “Be strong and of a good courage, fear not, nor be afraid … : for the Lord thy God, he it is that doth go with thee; he will not fail thee, nor forsake thee” (Deut. 31:6). He did not fail them. He will not fail us. He did not forsake them. He will not forsake us.

It was this knowledge that prompted the courage of Columbus—the quiet resolve to record in his ship’s log again and again: “This day we sailed on.” It was this witness that motivated the Prophet Joseph to declare, “I am going like a lamb to the slaughter; but I am calm as a summer’s morning” (D\&C 135:4).

It is this sweet assurance that can guide you and me—in our time, in our day, in our lives. Of course we will face fear, experience ridicule, and meet opposition. Let us have the courage to defy the consensus, the courage to stand for principle. Courage, not compromise, brings the smile of God’s approval. Courage becomes a living and an attractive virtue when it is regarded not only as a willingness to die manfully, but as the determination to live decently. A moral coward is one who is afraid to do what he thinks is right because others will disapprove or laugh. Remember that all men have their fears, but those who face their fears with dignity have courage as well.

From my personal chronology of courage, let me share with you two examples: one from military service, one from missionary experience.

Entering the United States Navy in the closing months of World War II was a challenging experience for me. I learned of brave deeds, acts of valor, and examples of courage. One best remembered was the quiet courage of an eighteen-year-old seaman—not of our faith—who was not too proud to pray. Of 250 men in the company, he was the only one who each night knelt down by the side of his bunk, at times amidst the jeers of the curious and the jests of unbelievers, and, with bowed head, prayed to God. He never wavered. He never faltered. He had courage.

Missionary service has ever called for courage. One who responded to this call was Randall Ellsworth. While serving in Guatemala as a missionary for The Church of Jesus Christ of Latter-day Saints, Randall Ellsworth survived a devastating earthquake that hurled a beam down on his back, paralyzing his legs and severely damaging his kidneys. He was the only American injured in the quake, which claimed the lives of some eighteen thousand persons.

After receiving emergency medical treatment, he was flown to a large hospital near his home in Rockville, Maryland. While Randall was confined there, a newscaster conducted with him an interview that I witnessed through the miracle of television. The reporter asked, “Can you walk?”

The answer: “Not yet, but I will.”

“Do you think you will be able to complete your mission?”

Came the reply: “Others think not, but I will. With the president of my church praying for me, and through the prayers of my family, my friends, and my missionary companions, I will walk, and I will return again to Guatemala. The Lord wants me to preach the gospel there for two years, and that’s what I intend to do.”

There followed a lengthy period of therapy, punctuated by heroic yet silent courage. Little by little, feeling began to return to the almost lifeless limbs. More therapy, more courage, more prayer.

At last, Randall Ellsworth walked aboard the plane that carried him back to the mission to which he had been called—back to the people whom he loved. Behind he left a trail of skeptics and a host of doubters, but also hundreds amazed at the power of God, the miracle of faith, and the example of courage.

On his return to Guatemala, Randall Ellsworth supported himself with the help of two canes. His walk was slow and deliberate. Then one day, as he stood before his mission president, Elder Ellsworth heard these almost unbelievable words spoken: “You have been the recipient of a miracle,” said the mission president. “Your faith has been rewarded. If you have the necessary confidence, if you have abiding faith, if you have supreme courage, place those two canes on my desk and walk.”

After a long pause, first one cane and then the other was placed on the desk, and a missionary walked. It was halting, it was painful—but he walked, never again to need the canes.

This spring I thought once more of the courage demonstrated by Randall Ellsworth. Years had passed since his ordeal. He was now a husband and a father. An engraved announcement arrived at my office. It read: “The President and Directors of Georgetown University announce commencement exercises of Georgetown University School of Medicine.” Randall Ellsworth received his Doctor of Medicine degree. More effort, more study, more faith, more sacrifice, more courage had been required. The price was paid, the victory won.

My brethren, let us be active participants—not mere spectators—on the stage of priesthood power. May we muster courage at the crossroads, courage for the conflicts, courage to say no, courage to say yes, for courage counts. Of this truth I testify in the name of Jesus Christ, amen.

\newpage
\settalk{Your Patriarchal Blessing: A Liahona of Light}
\setspeaker{Thomas S. Monson}
\settalkdate{October 1986}
\talkheading{Your Patriarchal Blessing: A Liahona of Light}{Thomas S. Monson}

Have you ever cleaned an attic or rummaged through an old storeroom? One discovers a bit of history and a whole lot of sentiment. A few weeks ago we emptied the attic of our mountain cabin. Seventy years of treasures, each with its own special memory, passed in review. Leading the parade was an old high chair with metal wheels. This was followed by glass milk bottles that once had pasteboard caps and by a copy of Life magazine with a story from World War II.

Featured in the magazine was an account of a once proud airplane, a mighty bomber, found rather well preserved in an isolated corner of the vast Sahara Desert. The bomber and crew had participated in the famous raid over Rumania’s Ploiesti oil fields. The craft had been struck by antiaircraft fire, which completely destroyed its communication and navigational equipment. As the stricken plane turned toward its desert landing field, a sudden sandstorm obliterated familiar points of reference. The field’s landing lights were shrouded by sand. The plane droned on, even far beyond the landing field, into the desert wastes until, with fuel exhausted, it settled on the Sahara, never to fly again. All crew members perished. Home and the safety and shelter there to be found had been denied. Victory, hopes, dreams—all had been swallowed by the silence of the desert’s dust.

Centuries earlier, a righteous and loving father by the name of Lehi took his beloved family into a desert wasteland. He journeyed in response to the voice of the Lord. But the Lord did not decree that such a “flight” be undertaken without heavenly help. The words of Nephi describe the gift provided on the morning of the historic trek:

“And it came to pass that as my father arose in the morning, and went forth to the tent door, to his great astonishment he beheld upon the ground a round ball of curious workmanship; and it was of fine brass. And within the ball were two spindles; and the one pointed the way whither we should go into the wilderness” (1 Ne. 16:10).

War and man-made means of destruction could not confuse or destroy this curious compass. Neither could the sudden desert sandstorms render useless its guiding powers. The prophet Alma explained that this “Liahona,” as it was called, was a compass prepared by the Lord. It worked for them according to their faith and pointed the way they should go (see Alma 37:38–40).

The same Lord who provided a Liahona to Lehi provides for you and for me today a rare and valuable gift to give direction to our lives, to mark the hazards to our safety, and to chart the way, even safe passage—not to a promised land, but to our heavenly home. The gift to which I refer is known as your patriarchal blessing. Every worthy member of the Church is entitled to receive such a precious and priceless personal treasure.

“Patriarchal blessings,” wrote the First Presidency in a letter to stake presidents, “contemplate an inspired declaration of the lineage of the recipient and, when so moved upon by the Spirit, an inspired and prophetic statement of the life mission of the recipient, together with such blessings, cautions and admonitions as the patriarch may be prompted to give for the accomplishment of such life’s mission, it being always made clear that the realization of all promised blessings is conditioned upon faithfulness to the gospel of our Lord, whose servant the patriarch is” (First Presidency Letter to stake presidents, 28 June 1958).

Who is this man, this patriarch, through whom such seership and priesthood power flow? How is he called? The Council of the Twelve Apostles has special responsibility pertaining to the calling of such men. From my own experience I testify that patriarchs are called of God by prophecy. How else could our Heavenly Father reveal those to whom such prophetic powers are to be given? A patriarch holds an ordained office in the Melchizedek Priesthood. The patriarchal office, however, is one of blessing—not of administration. I have never called a man to this sacred office but what I have felt the Lord’s guiding influence in the decision. May I share with you one treasured experience?

Many years back I had been assigned to name a patriarch for a stake in Logan, Utah. I found such a man, wrote his name on a slip of paper, and placed the note inside my scriptures. My further review revealed that another worthy patriarch had moved to this same area, making unnecessary the naming of a new patriarch. None was named.

Nine years later I was again assigned a stake conference in Logan. Once more a patriarch was needed for the stake I was to visit. I had been using a new set of scriptures for several years and had them in my briefcase. However, as I prepared to leave my home for the drive to Logan, I took from the bookcase shelf an older set of scriptures, leaving the new ones at home. During the conference I began my search for a patriarch: a worthy man, a blameless servant of God, one filled with faith, characterized by kindness. Pondering these requirements, I opened my scriptures and there discovered the slip of paper placed there long years before. I read the name written on the paper: Cecil B. Kenner. I asked the stake presidency if by chance Brother Kenner lived in this particular stake. I found he did. Cecil B. Kenner was that day ordained a patriarch.

Patriarchs are humble men. They are students of the scriptures. They stand before God as the means whereby the blessings of heaven can flow from that eternal source to the recipient on whose head rests the hands of the patriarch. He may not be a man of letters, a possessor of worldly wealth, or a holder of distinguished office. He, however, must be blessed with priesthood power and personal purity. To reach to heaven for divine guidance and inspiration, a patriarch is to be a man of love, a man of compassion, a man of judgment, a man of God.

A patriarchal blessing is a revelation to the recipient, even a white line down the middle of the road, to protect, inspire, and motivate activity and righteousness. A patriarchal blessing literally contains chapters from your book of eternal possibilities. I say eternal, for just as life is eternal, so is a patriarchal blessing. What may not come to fulfillment in this life may occur in the next. We do not govern God’s timetable. “For my thoughts are not your thoughts, neither are your ways my ways, saith the Lord.

“For as the heavens are higher than the earth, so are my ways higher than your ways, and my thoughts than your thoughts” (Isa. 55:8–9).

Your patriarchal blessing is yours and yours alone. It may be brief or lengthy, simple or profound. Length and language do not a patriarchal blessing make. It is the Spirit that conveys the true meaning. Your blessing is not to be folded neatly and tucked away. It is not to be framed or published. Rather, it is to be read. It is to be loved. It is to be followed. Your patriarchal blessing will see you through the darkest night. It will guide you through life’s dangers. Unlike the struggling bomber of yesteryear, lost in the desert wastes, the sands and storms of life will not destroy you on your eternal flight. Your patriarchal blessing is to you a personal Liahona to chart your course and guide your way.

In Lewis Carroll’s classic, Alice’s Adventures in Wonderland, Alice finds herself coming to a crossroads with two paths before her, each stretching onward but in opposite directions. She is confronted by the Cheshire Cat, of whom Alice asks, “Which path shall I take?”

The cat answers, “That depends where you want to go. If you do not know where you want to go, it doesn’t really matter which path you take” (see Alice’s Adventures in Wonderland, London: J.M. Dent, 1954, p. 54).

Unlike Alice, each of us knows where he or she wants to go. It does matter which way we go, for the path we follow in this life surely leads to the path we shall follow in the next.

Patience may be required as we watch, wait, and work for a promised blessing to be fulfilled.

One afternoon Percy K. Fetzer, a righteous patriarch, came to my office by appointment. He was weeping as we visited together. He explained that he had just returned from the land of Poland, where he had been privileged to give patriarchal blessings to our worthy members there. After a long pause, the patriarch revealed that he had been impressed to promise to members of a German-speaking family by the name of Konietz declarations which could not be fulfilled. He had promised missions. He had promised temple blessings. These were beyond the reach of those whom he had blessed. He whispered that he had tried to withhold the promises he knew were unattainable. It had been no use. The inspiration had come, the promises spoken, the blessings provided.

“What shall I do? What can I say?” he repeated to me.

I replied, “Brother Fetzer, these blessings have not come from you; they have been given of God. Let us kneel and pray to Him for their fulfillment.”

Within several years of that prayer, an unanticipated pact was signed between the Federal Republic of Germany and the Polish nation which provided that German nationals trapped in Poland at war’s end could now enter Germany. The Konietz family, whose members had received these special patriarchal blessings, came to live in West Germany. I had the privilege to ordain the father a bishop in the Dortmund stake of the Church. The family then made that long-awaited trek to the temple in Switzerland. They dressed in clothing of spotless white. They knelt at a sacred altar to await that ordinance which binds father, mother, brothers, and sisters not only for time, but for all eternity. He who pronounced that sacred sealing ceremony was the temple president. More than this, however, he was the same servant of the Lord, Percy K. Fetzer, who, as a patriarch years before, had provided those precious promises in the patriarchal blessings he had bestowed.

\begin{verse}
How far is Heaven?\\
It’s not very far.\\
When you live close to God,\\
It’s right where you are.
\end{verse}

Your patriarchal blessing is your passport to peace in this life. It is a Liahona of light to guide you unerringly to your heavenly home. Of these sacred truths I testify, in the name of Jesus Christ, amen.

\newpage
\settalk{Tears, Trials, Trust, Testimony}
\setspeaker{Thomas S. Monson}
\settalkdate{April 1987}
\talkheading{Tears, Trials, Trust, Testimony}{Thomas S. Monson}

Have you ever pondered the worth of a human soul? Have you ever wondered concerning the potential which lies within each of us?

Early in my service as a member of the Council of the Twelve, I was attending the conference of the Monument Park West Stake in Salt Lake City. My companion for the conference was a member of the General Church Welfare Committee, Paul C. Child. President Child was a student of the scriptures. He had been my stake president during my Aaronic Priesthood years. Now we were together as conference visitors.

When it was his opportunity to participate, President Child took the Doctrine and Covenants and left the pulpit to stand among the priesthood to whom he was directing his message. He turned to section 18 and began to read:

“Remember the worth of souls is great in the sight of God. …

“And if it so be that you should labor all your days in crying repentance unto this people, and bring, save it be one soul unto me, how great shall be your joy with him in the kingdom of my Father” (vs. 10, 15).

President Child then raised his eyes from the scriptures and asked the question of the priesthood brethren: “What is the worth of a human soul?” He avoided calling on a bishop, stake president, or high councilor for a response. Instead, he selected the president of an elders quorum—a brother who had been a bit drowsy and had missed the significance of the question.

The startled man responded: “Brother Child, could you please repeat the question?” The question was repeated: “What is the worth of a human soul?”

I knew President Child’s style. I prayed fervently for that quorum president. He remained silent for what seemed like an eternity and then declared: “Brother Child, the worth of a human soul is its capacity to become as God.”

All present pondered that reply. Brother Child returned to the stand, leaned over to me, and said: “A profound reply; a profound reply!” He proceeded with his message, but I continued to reflect on that inspired response.

To reach, to teach, to touch the precious souls whom our Father has prepared for His message is a monumental task. Success is rarely simple. Generally it is preceded by tears, trials, trust, and testimony.

Think of the magnitude of the Savior’s instruction to His Apostles:

“Go ye therefore, and teach all nations, baptizing them in the name of the Father, and of the Son, and of the Holy Ghost:

“Teaching them to observe all things whatsoever I have commanded you: and, lo, I am with you alway, even unto the end of the world” (Matt. 28:19–20).

The men to whom he gave this instruction were not owners of land, nor did they have the education of the learned. They were simple men—men of faith, men of devotion, men “called of God.”

Paul testified to the Corinthians: “Not many wise men after the flesh, not many mighty, not many noble, are called:

“But God hath chosen the foolish things of the world to confound the wise; and God hath chosen the weak things of the world to confound the things which are mighty” (1 Cor. 1:26–27).

On this, the American continent, Alma likewise counseled his son Helaman: “I say unto you, that by small and simple things are great things brought to pass” (Alma 37:6).

Then and now, servants of God take comfort from the Master’s assurance: “I am with you alway” (Matt. 28:20). This magnificent promise sustains you brethren of the Aaronic Priesthood who are called to positions of leadership in the quorums of deacons, teachers, and priests. It encourages you in your preparations to serve in the mission field. It comforts you during those moments of discouragement, which come to all. This same assurance motivates and inspires you brethren of the Melchizedek Priesthood as you lead and direct the work in the wards, the stakes, and the missions.

“Wherefore, be not weary in well-doing,” said the Lord, “for ye are laying the foundation of a great work. And out of small things proceedeth that which is great.

“Behold, the Lord requireth the heart and a willing mind” (D\&C 64:33–34).

An abiding faith, a constant trust, a fervent desire have always characterized those who serve the Lord with all their hearts.

This description typified the early beginnings of missionary work following the restoration of the gospel. As early as April of 1830, Phineas Young received a copy of the Book of Mormon from Samuel Smith, brother of the Prophet, and a few months later traveled to Upper Canada. At Kingston, he gave the first known testimony of the restored Church beyond the borders of the United States.

In 1833, the Prophet Joseph Smith, Sidney Rigdon, and Freeman Nickerson traveled to Mount Pleasant, Upper Canada. There they taught, they baptized, they organized a branch of the Church. At one time, in June of 1835, six of the Twelve held a conference in that land.

In April of 1836, Elder Heber C. Kimball and others entered the home of Parley P. Pratt and, filled with the spirit of prophecy, they placed their hands on the head of Brother Pratt and declared: “Thou shalt go to Upper Canada, even to the City of Toronto, … and there thou shalt find a people prepared for the fulness of the gospel, and they shall receive thee, and thou shalt organize the Church among them, and many shall be brought to the knowledge of the truth and shall be filled with joy; and from the things growing out of this mission, shall the fulness of the gospel spread into England, and cause a great work to be done in that land” (Autobiography of Parley P. Pratt, Salt Lake City: Deseret Book Co., 1975, pp. 130–31).

In July of this year there shall be commemorated the 150th anniversary of the beginning of the work in England. We rejoice in the tremendous accomplishments of those early missionaries and those whom the Lord prepared to play such a part in the advancement of this latter-day work.

The call to serve has ever characterized the work of the Lord. It rarely comes at a convenient time. It brings humility, it provokes prayer, it inspires commitment. The call came—to Kirtland. Revelations followed. The call came—to Missouri. Persecution prevailed. The call came—to Nauvoo. Prophets died. The call came—to the basin of the Great Salt Lake. Hardship beckoned.

That long journey, made under such difficult circumstances, was a trial of faith. But faith forged in the furnace of trials and tears is marked by trust and testimony. Only God can count the sacrifice; only God can measure the sorrow; only God can know the hearts of those who serve Him—then and now.

Lessons from the past can quicken our memories, touch our lives, and direct our actions. We are prompted to pause and remember that divinely given promise: “Wherefore, … ye are on the Lord’s errand; and whatever ye do according to the will of the Lord is the Lord’s business” (D\&C 64:29).

Such a lesson was recounted on a radio and television program many remember with fondness. The program was entitled “Death Valley Days.” The narrator, known as the Old Ranger, seemed to come right into our living rooms as he would tell the tales of the West.

On one program, the Old Ranger related how the glass was obtained for the windows of the St. George Tabernacle. The glass had been manufactured in the East. Then it had been placed on a ship in New York, which sailed forth on the long and at times perilous journey around Cape Horn and up to the West Coast of America. The precious glass, stored in cartons, was then transported to San Bernardino, California, to await the overland trek to St. George.

David Cannon and the brethren in St. George had the duty to go to San Bernardino with their teams and wagons to retrieve the glass, that the tabernacle of the Lord could be completed. One problem: They needed the then-astronomical sum of \$800 to pay for the glass. They had no money. David Cannon turned to his wife and his son and asked, “Do you think that we can raise the money, that we might obtain the glass for the tabernacle?”

His tiny boy, David, Jr., said, “Daddy, I know we can!” He then produced two cents of his own money and gave it to his father. Wilhelmina Cannon, David’s wife, went through the secret hiding places that all women have in their houses. Her search produced \$3.50 in silver. The community was scoured for money, and at length the sum of \$200 was accumulated—\$600 short of the required amount.

David Cannon sighed the sigh of despair of one who had failed although he had tried his best. The little family was really too weary to sleep and too discouraged to eat, so they prayed. Morning dawned. There gathered the teamsters with their wagons and teams, prepared to undertake the long journey to San Bernardino. But they had no \$600.

Then there came a knock at the door, and Peter Nielsen, from the nearby community of Washington, entered the house. He said to David Cannon, “Brother David, I have had a persistent dream that I should bring the money that I had saved to expand my house—bring it to you, that you would have a purpose for it.”

While all of the men gathered around the table, including little David, Jr., Peter Nielsen took out a red bandanna and dropped gold pieces, one by one, upon the table. When David Cannon counted the gold pieces, they totaled \$600—the exact amount needed to obtain the glass. Within an hour the men waved good-bye and, with their teams, set forth on their journey to San Bernardino to retrieve the glass for the tabernacle.

When that true story was told on “Death Valley Days,” young David Cannon, Jr., was then eighty-seven years of age. He listened to the story with rapt attention. I feel that in his mind, he once again heard those gold pieces, one by one, dropping upon the table as astonished men saw with their very eyes the answer to their prayers.

Tabernacles and temples are built with more than stone and mortar, wood and glass. Particularly is this true when we speak of the temple described by the Apostle Paul: “Know ye not that ye are the temple of God, and that the Spirit of God dwelleth in you?” (1 Cor. 3:16). Such temples are built with faith and fasting. They are built with service and sacrifice. They are built with trials and testimonies.

If any brethren within the sound of my voice feel unprepared, even incapable of responding to a call to serve, to sacrifice, to bless the lives of others, remember the truth: “Whom God calls, God qualifies.” He who notes the sparrow’s fall will not abandon the servant’s need.

God bless you, my brethren—you who bear the priesthood. You “are a chosen generation, a royal priesthood” (1 Pet. 2:9).

May we respond affirmatively to the Prophet Joseph, who urged: “Brethren, shall we not go on in so great a cause? Go forward and not backward. Courage, brethren; and on, on to the victory!” (D\&C 128:22). This is my earnest and humble prayer, in the name of Jesus Christ, amen.

\newpage
\settalk{The Will Within}
\setspeaker{Thomas S. Monson}
\settalkdate{April 1987}
\talkheading{The Will Within}{Thomas S. Monson}

A few months ago I stood before a capacity audience in the Marriott Center on the campus of Brigham Young University. Mine was the responsibility to speak, to lift, to motivate, to inspire. There came to my mind the realization that here were men and women of promise. They represented the hopes, the dreams, and the aspirations of parents, of family, of teachers, of God. All were participants in the passing parade of mortality. Some were gifted in the arts, others leaned toward the humanities, while some found their talents prompted a study of the natural or physical sciences. These students stood on the stage of study. Soon they would disperse to make their marks in life, to fulfill the measure of their creation, and to learn from their own lives those lessons which would prepare them for the exaltation they seek.

My thoughts turned to others striving to become master craftsmen through apprenticeship and experience. Then I reflected on that vast throng who had abandoned preparation, formed undesirable friendships, and adopted habits and practices which diverted them from that pathway which leads to perfection and enticed them along one of the many detours where sorrow, discouragement, and destruction await.

The wayward son, the willful daughter, the pouting husband, the nagging wife—all can change. There can occur a parting of the clouds, a break in the storm. Maturity comes, friendships alter, circumstances vary. “Cast in concrete” need not describe human behavior.

From the perspective of eternity, our sojourn in this life is ever so brief. Detours are costly; they must be shunned. The spiritual nature within us should not be dominated by the physical. It behooves each of us to remember who he or she is and what God expects him or her to become.

The poet Wordsworth, in his inspired Intimations of Immortality, inclined our thoughts to that heavenly home from whence each of us came:

\begin{verse}
Our birth is but a sleep and a forgetting:\\
The soul that rises with us, our life’s star,\\
Hath had elsewhere its setting,\\
And cometh from afar:\\
Not in entire forgetfulness,\\
And not in utter nakedness,\\
But trailing clouds of glory do we come\\
From God, who is our home.
\end{verse}

(“Ode: Intimations of Immortality from Recollections of Early Childhood,” lines 58–65).

In finding and trailing this spiritual contact with the infinite, we will feel the touch of inspiration and know that God will guide us as we put in Him our trust. That wise and righteous man, Job, declared the profound truth: “There is a spirit in man: and the inspiration of the Almighty giveth … understanding” (Job 32:8). It is this inspiration which we at times allow to grow dim, causing us to wander far below the level of our possibilities.

During the Great Depression, the homeless, the downtrodden, the unemployed “rode the rails” that passed not far from our home. On numerous occasions, there would be a soft knock on the back door. When I opened the door, there I would see a man, sometimes two, ill-clothed, ill-fed, ill-schooled. Generally, such a visitor held in his hand the familiar cap. His hair would be tousled, his face unshaven. The question was always the same: “Could you spare some food?” My dear mother invariably responded with a pleasant, “Come in and sit down at the table.” She would then prepare a ham sandwich, cut a piece of cake, and pour a glass of milk. Mother would ask the visitor about his home, his family, his life. She provided hope and words of encouragement. Before leaving, the visitor would pause to express a gracious thank-you. I would note that a smile of content had replaced a look of despair. Eyes that were dull now shone with new purpose. Love, that noblest attribute of the human soul, can work wonders.

In our journey on earth, we discover that life is made up of challenges—they just differ from one person to another. We are success oriented, striving to become “wonder women” and “super men.” Any intimation of failure can cause panic, even despair. Who among us cannot remember moments of failure?

One such moment came to me as a young basketball player. The game was close—hotly contested—when the coach called me from the bench to run a key play. For some reason which I shall never understand, I took the pass and dribbled the ball right through the opposing team. I jumped high toward the basket; and, as the basketball left my fingertips, I came to the abrupt realization that I was shooting for the wrong basket. I offered the shortest prayer I have ever spoken: “Dear Father, don’t let that ball go in.” My prayer was answered, but my ordeal was just beginning. I heard a loud cheer erupt from the adoring fans: “We want Monson, we want Monson, we want Monson … OUT!” The coach obliged.

Not long ago I read about an incident that occurred in the life of President Harry S. Truman after he had retired and was back in Independence, Missouri. He was at Truman Library, talking with some elementary school students and answering their questions. Finally, a question came from an owlish little boy. “Mr. President,” he said, “was you popular when you was a boy?” The President looked at the boy and answered, “Why, no. I was never popular. The popular boys were the ones who were good at games and had big, tight fists. I was never like that. Without my glasses, I was blind as a bat, and to tell the truth, I was kind of a sissy.” The little boy started to applaud, and then everyone else did, too” (Eugene W. Brice, “Good News about Failure,” Vital Speeches, 1 Feb. 1983, p. 236.)

Our responsibility is to rise from mediocrity to competence, from failure to achievement. Our task is to become our best selves. One of God’s greatest gifts to us is the joy of trying again, for no failure ever need be final. In 1902, the poetry editor of the Atlantic Monthly returned a sheaf of poems to a 28-year-old poet with this curt note: “Our magazine has no room for your vigorous verse.” The poet was Robert Frost. In 1894, the rhetoric teacher at Harrow in England wrote on a 16-year-old’s report card, “A conspicuous lack of success.” The 16-year-old was Winston Churchill.

President Theodore Roosevelt said, “It is not the critic who counts, not the man who points out how the strong man stumbles, or where the doer of deeds could have done them better. The credit belongs to the man who is actually in the arena” (Clifton Fadiman, ed., The American Treasury: 1455–1955, New York: Harper and Brothers, 1955, p. 689).

We know men and women can change—and change for the better. No more vivid example is recorded than the life of Saul of Tarsus. The sacred record reveals that Saul threatened the disciples of the Lord. Then came that light from heaven and the voice saying unto him: “Saul, Saul, why persecutest thou me?

“And he said, Who art thou, Lord? And the Lord said, I am Jesus whom thou persecutest” (Acts 9:4–5).

Saul’s answer is a model for each of us: “Lord, what wilt thou have me to do?” (v. 6). Saul the persecutor became Paul the proselyter. Night had turned to day. Darkness had yielded to light.

Simon Peter, that fisherman who left his nets and followed the Lord, had his time of struggle. He had been weak and fearful and had denied his Lord with an oath. Then there came change. Never again would he deny or desert his Lord. He found his place in the kingdom of God.

We have the example of Alma the Younger, who turned his back on sinful practices and wasteful ways. Conversion came. He became an exponent of truth. His tender words of counsel to his sons Helaman and Corianton are literary classics. To Helaman: “O, remember, my son, and learn wisdom in thy youth; yea, learn in thy youth to keep the commandments of God” (Alma 37:35). To Corianton: “Suffer not yourself to be led away by any vain or foolish thing” (Alma 39:11).

Then and now, as President David O. McKay so consistently taught, the gospel of Jesus Christ can make bad men good and good men better, can alter human nature and change human lives.

Change for the better can come to all. In December of 1985 the First Presidency proclaimed “An Invitation to Come Back.” To the inactive, the critical, the transgressor, the message declared: “‘Come back. Come back and feast at the table of the Lord, and taste again the sweet and satisfying fruits of fellowship with the Saints’” (Ensign, Mar. 1986, p. 88). Hundreds, if not thousands, have responded to this invitation. Their lives have taken on new meaning. Their families have been blessed. They have drawn closer to God.

In the private sanctuary of one’s own conscience lies that spirit, that determination to cast off the old person and to measure up to the stature of true potential. But the way is rugged, and the course is strenuous. So discovered John Helander from Goteborg, Sweden. John is twenty-six years of age and is handicapped, in that it is difficult for him to coordinate his motions.

At a youth conference in Kungsbacka, Sweden, John took part in a 1500-meter running race. He had no chance to win. Rather, his was the opportunity to be humiliated, mocked, derided, scorned. Perhaps John remembered another who lived long ago and far away. Wasn’t He mocked? Wasn’t He derided? Wasn’t He scorned? But He prevailed. He won His race. Maybe John could win his.

What a race it was! Struggling, surging, pressing, the runners bolted far beyond John. There was wonderment among the spectators. Who is this runner who lags so far behind? The participants on their second lap of this two-lap race passed John while he was but halfway through the first lap. Tension mounted as the runners pressed toward the tape. Who would win? Who would place second? Then came the final burst of speed; the tape was broken. The crowd cheered; the winner was proclaimed.

The race was over—or was it? Who is this contestant who continues to run when the race is ended? He crosses the finish line on but his first lap. Doesn’t the foolish lad know he has lost? Ever onward he struggles, the only participant now on the track. This is his race. This must be his victory. No one among the vast throng of spectators leaves. Every eye is on this valiant runner. He makes the final turn and moves toward the finish line. There is awe; there is admiration. Every spectator sees himself running his own race of life. As John approaches the finish line, the audience, as one, rises to its feet. There is a loud applause of acclaim. Stumbling, falling, exhausted but victorious, John Helander breaks the newly tightened tape. (Officials are human beings, too.) The cheering echoes for miles. And just maybe, if the ear is carefully attuned, that Great Scorekeeper—even the Lord—can be heard to say, “Well done, thou good and faithful servant” (Matt. 25:21).

Each of us is a runner in the race of life. Comforting is the fact that there are many runners. Reassuring is the knowledge that our eternal Scorekeeper is understanding. Challenging is the truth that each must run. But you and I do not run alone. That vast audience of family, friends, and leaders will cheer our courage, will applaud our determination as we rise from our stumblings and pursue our goal. The race of life is not for sprinters running on a level track. The course is marked by pitfalls and checkered with obstacles. We take confidence from the hymn:

\begin{verse}
Fear not, I am with thee; oh, be not dismayed,\\
For I am thy God and will still give thee aid.\\
I’ll strengthen thee, help thee, and cause thee to stand,\\
Upheld by my righteous, … omnipotent hand. …
\end{verse}

(“How Firm a Foundation,” Hymns, 1985, no. 85).

Let us shed any thought of failure. Let us discard any habit that may hinder. Let us seek; let us obtain the prize prepared for all, even exaltation in the celestial kingdom of God. This is my earnest prayer, in the name of Jesus Christ, amen.

\newpage
\settalk{A Doorway Called Love}
\setspeaker{Thomas S. Monson}
\settalkdate{October 1987}
\talkheading{A Doorway Called Love}{Thomas S. Monson}

Recently there moved over the wires of Associated Press a catalog of crime as the daily happenings around the world were relayed to the media and thence to the homes on every continent.

The headlines were brief. They highlighted murder, rape, robbery, molestation, fraud, deceit, and corruption. I made note of several: “Man Slays Wife and Children, Then Turns Gun on Self”; “Child Identifies Molester”; “Hundreds Lose All As Multimillion-Dollar Scam Is Exposed.” The sordid list continued. Shades of Sodom, glimpses of Gomorrah.

President Ezra Taft Benson has often stated, “We live in a wicked world.” The Apostle Paul warned, “Men shall be lovers of their own selves, covetous, boasters, proud, blasphemers, disobedient to parents, unthankful, unholy. … lovers of pleasures more than lovers of God” (2 Tim. 3:2–4).

Must we suffer the same fate as those who lived in the cities of the plain? Can we not learn the lesson taught in the time of Noah? “Is there no balm in Gilead?” (Jer. 8:22). Or is there a doorway that leads us from the morass of worldliness onward and upward to the high ground of righteousness? There echoes ever so gently to the honest mind that personal invitation of the Lord: “Behold, I stand at the door, and knock: if any man hear my voice, and open the door, I will come in to him” (Rev. 3:20). Does that doorway have a name? It surely does. I have chosen to call it “the doorway of love.”

Love is the catalyst that causes change. Love is the balm that brings healing to the soul. But love doesn’t grow like weeds or fall like rain. Love has its price. “God so loved the world, that he gave his only begotten Son, that whosoever believeth in him should not perish, but have everlasting life” (John 3:16). That Son, even the Lord Jesus Christ, gave His life that we might have eternal life, so great was His love for His Father and for us.

This same Jesus was approached by a lawyer who asked, “Master, which is the great commandment in the law?

“Jesus said unto him, Thou shalt love the Lord thy God with all thy heart, and with all thy soul, and with all thy mind.

“This is the first and great commandment.

“And the second is like unto it, Thou shalt love thy neighbour as thyself.

“On these two commandments hang all the law and the prophets” (Matt. 22:36–40).

In that tender and touching farewell, as He counseled His beloved disciples, Jesus taught: “He that hath my commandments, and keepeth them, he it is that loveth me” (John 14:21). Particularly far-reaching was the instruction, “A new commandment I give unto you, That ye love one another; as I have loved you” (John 13:34).

Little children can learn the lesson of love. Profound instruction from holy writ ofttimes is not understood by them. However, they respond readily to a favorite verse:

\begin{verse}
“I love you, Mother,” said little [John];\\
Then, forgetting his work, his cap went on,\\
And he was off to the garden swing,\\
And left her the water and wood to bring.
\end{verse}

(Joy Allison, in The World’s Best Loved Poems, comp. James Gilchrist Lawson, New York: Harper and Row, 1955, pp. 243–44)

Home should be a haven of love. Honor, courtesy, and respect symbolize love and characterize the righteous family. Fathers in such homes will not hear the denunciation of the Lord as recorded in the book of Jacob: “Ye have broken the hearts of your tender wives, and lost the confidence of your children, because of your bad examples before them; and the sobbings of their hearts ascend up to God against you” (Jacob 2:35).

In 3 Nephi the Master instructed: “There shall be no disputations among you. …

“For verily, verily I say unto you, he that hath the spirit of contention is not of me, but is of the devil, who is the father of contention, and he stirreth up the hearts of men to contend with anger, one with another.

“Behold, this is not my doctrine, to stir up the hearts of men with anger, one against another; but this is my doctrine, that such things should be done away” (3 Ne. 11:28–30).

Where love is, there is no disputation. Where love is, there is no contention. Where love is, there God will be also. Each of us has the responsibility to keep His commandments. The lessons found in scripture find fulfillment in our lives. Joseph Smith taught that “happiness is the object and design of our existence; and will be the end thereof, if we pursue the path that leads to it; and this path is virtue, uprightness, faithfulness, holiness, and keeping all the commandments of God” (Teachings of the Prophet Joseph Smith, sel. Joseph Fielding Smith, Salt Lake City: Deseret Book Co., 1938, pp. 255–56).

In the classic musical production Camelot, there is a line with words of warning for all. After the familiar triangle began to deepen regarding King Arthur, Lancelot, and Guenevere, King Arthur said, “We must not let our passions destroy our dreams.”

From that same production came another truth also spoken by Arthur as he envisioned a better world: “Violence is not strength, and compassion is not weakness.”

In this world in which we live, there is a tendency for us to describe needed change, required help, and desired relief with the familiar phrase, “They ought to do something about this.” We fail to define the word they. I love the message, “Let there be peace on earth, and let it begin with me.”

Tears came to my eyes when I read of a mere boy in one of our eastern cities who noticed a vagrant asleep on a sidewalk and who then went to his own bedroom, retrieved his own pillow, and placed it beneath the head of that one whom he knew not. Perhaps there came from the precious past the welcome words: “Inasmuch as ye have done it unto one of the least of these my brethren, ye have done it unto me” (Matt. 25:40).

I extol those who, with loving care and compassionate concern, feed the hungry, clothe the naked, and house the homeless. He who notes the sparrow’s fall will not be unmindful of such service.

\begin{verse}
A bell is no bell till you ring it,\\
A song is no song till you sing it,\\
And love in your heart wasn’t put there to stay—\\
Love isn’t love till you give it away.
\end{verse}

(“Sixteen Going on Seventeen [Reprise],” from The Sound of Music, Rodgers and Hammerstein)

From the Holy Bible we read: “And it came to pass … that [Jesus] went into a city called Nain. …

“When he came nigh to the gate of the city, behold, there was a dead man carried out, the only son of his mother, and she was a widow. …

“And when the Lord saw her, he had compassion on her, and said unto her, Weep not.

“And he came and touched the bier: and they that bare him stood still.”

In the majesty of his messianic ministry, He declared: “Young man, I say unto thee, Arise.

“And he that was dead sat up, and began to speak. And he delivered him to his mother” (Luke 7:11–15).

The desire to lift, the willingness to help, and the graciousness to give come from a heart filled with love.

The poet wrote, “Love is the most noble attribute of the human soul.” And William Shakespeare cautioned, “They do not love that do not show their love” (Two Gentlemen of Verona, act 1, scene 2, line 31).

A school teacher showed her love with her guiding philosophy: “No one fails in my class. I have the responsibility to help each student succeed.”

A priesthood quorum leader in Salt Lake City—a retired executive—said to me, “This year I have helped twelve of my brethren who were out of work to obtain permanent employment. I have never been happier in my entire life.” Short in stature, “Little Ed,” as we affectionately called him, stood tall that day as his eyes glistened and his voice quavered. He showed his love by helping those in need.

A large and tough businessman, a wholesale vendor of poultry, showed his love with a single comment made when one attempted to pay for twenty-four roasting chickens. “The chickens are going to the widows, aren’t they? There will be no charge.” As he placed them in the car trunk, he said in a faltering voice: “And there are more where these came from.”

Robert Woodruff, an executive in a former generation, traversed America with a message which he delivered to civic and business groups. The outline was simple, the message brief:

The five most important words are these: I am proud of you.

The four most important words are these: What is your opinion?

The three most important words are these: If you please.

The two most important words are these: Thank you.

To Mr. Woodruff’s list I would add, “The single most important word is love.”

A few years ago Morgan High School played Millard High for the state football championship. From his wheelchair, to which Morgan coach Jan Smith was confined, he said to his team: “This is the most important game of your lives. You lose, and you will regret it forever. You win, and you will remember it forever. Make every play as though it were all-important.”

Behind the door, his wife, whom he tenderly referred to as his chief assistant, overheard her husband say, “I love you guys. I don’t care about the ball game. I love you and want the game victory for you.” Underdog Morgan High won the football game and the state championship.

True love is a reflection of Christ’s love. In December of each year we call it the Christmas spirit. You can hear it. You can see it. You can feel it. But never alone.

One winter day, I thought back to an experience from my boyhood. I was just eleven. Our Primary president, Melissa, was an older and loving gray-haired lady. One day at Primary, Melissa asked me to stay behind and visit with her. There the two of us sat in the otherwise-empty chapel. She placed her arm about my shoulder and began to cry.

Surprised, I asked her why she was crying.

She replied, “I don’t seem to be able to encourage the Trail Builder boys to be reverent during the opening exercises of Primary. Would you be willing to help me, Tommy?”

I promised Melissa that I would. Strangely to me, but not to Melissa, that ended any problem of reverence in that Primary. She had gone to the source of the problem—me. The solution was love.

The years flew by. Marvelous Melissa, now in her nineties, lived in a nursing facility in the northwest part of Salt Lake City. Just before Christmas I determined to visit my beloved Primary president. Over the car radio, I heard the song, “Hark! the herald angels sing; Glory to the newborn King!” (Hymns, 1985, no. 209.) I reflected on the visit made by wise men those long years ago. They brought gifts of gold, of frankincense, and of myrrh. I brought only the gift of love and a desire to say thank you.

I found Melissa in the lunchroom. She was staring at her plate of food, teasing it with the fork she held in her aged hand. Not a bite did she eat. As I spoke to her, my words were met by a benign but blank stare. I took the fork in hand and began to feed Melissa, talking all the time I did so about her service to boys and girls as a Primary worker. There wasn’t so much as a glimmer of recognition, far less a spoken word. Two other residents of the nursing home gazed at me with puzzled expressions. At last they spoke, saying, “Don’t talk to her. She doesn’t know anyone—even her own family. She hasn’t said a word in all the years she’s been here.”

Luncheon ended. My one-sided conversation wound down. I stood to leave. I held her frail hand in mine, gazed into her wrinkled but beautiful countenance, and said, “God bless you, Melissa. Merry Christmas.”

Without warning, she spoke the words, “I know you. You’re Tommy Monson, my Primary boy. How I love you.” She pressed my hand to her lips and bestowed on it the kiss of love. Tears coursed down her cheeks and bathed our clasped hands. Those hands, that day, were hallowed by heaven and graced by God. The herald angels did sing. The words of the Master seemed to have a personal meaning never before fully felt: “Woman, behold thy son!” And to his disciple, “Behold thy mother!” (see John 19:26–27).

Outside the sky was blue—azure blue. The air was cool—crispy cool. The snow was white—crystal white.

From Bethlehem there seemed to echo the words:

\begin{verse}
How silently, how silently\\
The wondrous gift is giv’n!\\
So God imparts to human hearts\\
The blessings of his heav’n.
\end{verse}

(“O Little Town of Bethlehem,” Hymns, 1985, no. 208)

The wondrous gift was given, the heavenly blessing was received, the dear Christ had entered in—all through the doorway called love. I declare this solemn truth in the name of Jesus Christ, amen.

\newpage
\settalk{Missionary Memories}
\setspeaker{Thomas S. Monson}
\settalkdate{October 1987}
\talkheading{Missionary Memories}{Thomas S. Monson}

What an inspiring sight to see this historic tabernacle filled to capacity, then to realize that chapels and halls throughout the world are similarly filled by those who hold the priesthood of God. I pray for the inspiration of heaven to attend me and to direct the remarks I make.

My mind goes back in memory to a general priesthood meeting held in 1956. At that time I was serving in the stake presidency of the Temple View Stake here in Salt Lake City. Percy K. Fetzer, John R. Burt, and I, the stake presidency, had come to the Tabernacle early, that hopefully we might find a place to sit. We were among the first to enter the Tabernacle and had almost two hours to wait before the meeting would begin.

President Fetzer related to President Burt and me an experience from his missionary days in Germany. He described how one rainy night he and his companion were to present a gospel message to a group assembled in a schoolhouse. A protester had broadcast falsehoods concerning the Church, and a number of people threatened violence against the two missionaries. At a critical moment, a woman who was a widow stepped between the elders and the angry group and said, “These young men are my guests and are coming to my home now. Please make way for us to leave.”

The crowd parted, and the missionaries walked through the rainy night with their benefactress, arriving at length at her modest home. She placed their wet coats over the kitchen chairs and invited the missionaries to sit at the table while she prepared food for them. After eating, the elders presented a message to the kind lady who had befriended them. A young son of the woman was invited to come to the table, but he refused, preferring his position of solitude and warmth directly behind the kitchen stove.

President Fetzer concluded the account with the comment, “While I don’t know if that woman ever joined the Church, I’ll forever be grateful to her for her kindness that rain-drenched night thirty-three years ago.”

The brethren sitting in front of us here in the Tabernacle had been speaking to one another also. After a while, we began listening to their conversation. One asked the friend sitting next to him, “Tell me how you came to be a member of the Church.”

The brother responded, “One rainy night in Germany, my mother brought to our house two drenched missionaries whom she had rescued from a mob. Mother fed the elders, and they presented to her a message concerning the work of the Lord. They invited me to join the discussion, but I was shy and fearful, so I remained secure in my seat behind the stove. Later, when I once more heard about the Church, I remembered the courage and faith, as well as the message, of those two humble missionaries, and this led to my conversion. I suppose I’ll never meet those two missionaries here in mortality, but I’ll be forever grateful to them. I know not where they were from. I think one was named Fetzer.”

At this point, President Burt and I looked at President Fetzer and noticed the great tears which coursed down his cheeks. Without saying a word to us, President Fetzer tapped on the shoulder of the man in front of us who had just related his conversion experience. To him he then said, “I’m Bruder Fetzer. I was one of the two missionaries whom you befriended that night. I’m grateful to meet the boy who sat behind the stove—the lad who listened and who learned.”

I do not remember the messages delivered during the priesthood meeting that night, but I shall never forget the faith-filled conversation which preceded the commencement of the meeting.

The words of the Lord seemed so appropriate then. They are equally appropriate now: “And if it so be that you should labor all your days in crying repentance unto this people, and bring, save it be one soul unto me, how great shall be your joy with him in the kingdom of my Father!” (D\&C 18:15).

We are a missionary-minded people. We have a divine mandate to proclaim the message of the Restoration. You young men here this night are on the threshold of your missionary opportunity. That energetic missionary from the Book of Mormon, even Alma, provides for us a blueprint for missionary conduct: “This is my glory, that perhaps I may be an instrument in the hands of God to bring some soul to repentance; and this is my joy” (Alma 29:9).

I add my personal witness: Our missionaries are not salesmen with wares to peddle; rather, they are servants of the Most High God, with testimonies to bear, truths to teach, and souls to save.

Each missionary who goes forth in response to a sacred call becomes a servant of the Lord whose work this truly is. Do not fear, young men, for He will be with you. He never fails. He has promised: “I will go before your face. I will be on your right hand and on your left, and my Spirit shall be in your hearts, and mine angels round about you, to bear you up” (D\&C 84:88).

“And ye shall go forth in the power of my Spirit, preaching my gospel, two by two, in my name, lifting up your voices as with the sound of a trump, declaring my word like unto angels of God” (D\&C 42:6).

Fathers, bishops, quorum advisers, yours is the responsibility to prepare this generation of missionaries, to quicken in the hearts of these deacons, teachers, and priests not only an awareness of their obligation to serve, but also a vision of the opportunities and blessings which await them through a missionary call. The work is demanding, the impact everlasting. This is no time for “summer soldiers” in the army of the Lord.

The missionary recommendations that arrive daily at Church headquarters present a spectrum of preparedness. Let me share with you just one or two comments gleaned from the period I served on the Missionary Committee. One recommendation form contained this comment written by the bishop: “John is very close to his mother. She would be happy if he were assigned to a mission close to their home so she could phone weekly and visit him on occasion.” As I read this comment to President Spencer W. Kimball, who assigned the missionaries then, I wondered what his reaction would be. Would he assign the young man to California or Washington, that he might be near his Oregon home? Without raising his eyes from the assignment sheet, President Kimball said, “Please assign this young man to Johannesburg, South Africa.”

Another missionary recommendation contained the comment from the stake president: “This young man was instrumental in bringing his stepfather into the Church about a year ago. His stepfather told me it was because of Jerry getting up each Sunday morning and going to church that caused him to wonder what kind of church could have that much influence on a boy.”

In many respects, a mission is a family calling. The letters which a missionary sends to Mother and Father are packed with power—spiritual power. They are filled with faith—abiding faith. I’ve always maintained that such letters seem to pass through a heavenly post office before being delivered to home and family. Mother treasures every word. Father fills with pride. The letters are read over and over again—and are never discarded.

I trust parents will remember that their letters to a missionary son or daughter bring home and heaven close to him or to her and provide a renewal of commitment to the sacred calling of missionary. God will inspire you as you take pen in hand to express to one you love the feelings of your soul and the love of your heart.

At the funeral service for the mother of Elder Marion G. Romney, held in Provo, Utah, her son-in-law, Brother John K. Edmunds, gave the following account: “In their early married life, Brother and Sister Romney lived in Mexico. Brother Romney [like the father of President Benson] was called on a mission. There was no feasible means of support, yet he went and his wife sustained him. One day she grieved because she wanted to write her husband a letter but did not have sufficient money to buy a postage stamp. She prayed and then took a walk through the orchard that autumn day, kicking the leaves as she walked along and thinking of her husband. She noticed a shiny object on the ground and discovered it to be a coin—just the right amount for several postage stamps.”

Her letter had been written. Now, through the intervention of God, it could be mailed.

Brethren, think of the family blessings received by the Romney and Benson families, which blessings followed the commitment to missionary service.

I think of my own grandfather, Nels Monson, who waited seven years for his sweetheart to become his bride. The first entry in his missionary journal expressed eloquently his gratitude: “Today, in the Salt Lake Temple, Maria Mace became my eternal wife.” The entry written three days later was more somber: “Tonight the bishop came to our house. I have been called to serve a two-year mission to Scandinavia. My dear wife will remain at home and sustain me.” I treasure such faith. I cherish such commitment.

I commend the many couples who now go forth to serve. Leaving the comforts of home, the companionship of family, they walk hand in hand as eternal companions, but also hand in hand with God as His representatives to a faith-starved world.

To the many who contribute of their means for missionary service, I express the thanks of the Church and the sentiments of my soul. The gratitude of God may come soon. Then again, it may come as it did to Brother Fetzer—after thirty-three years. This I know: It will come. It will bless. It will comfort. It will sanctify.

Last month the Salt Lake City newspapers carried an obituary notice for Fred Sudbury. It indicated that he was survived by his wife, Pearl, and a son, Craig; that he was a member of The Church of Jesus Christ of Latter-day Saints; and that his marriage had been solemnized in the Salt Lake Temple. What the obituary notice could not adequately convey was the inspiring human drama which preceded Fred’s passing.

Some years ago, Craig Sudbury and his mother came to my office prior to Craig’s departure for the Australia Melbourne Mission. Fred Sudbury, Craig’s father, was noticeably absent. Twenty-five years earlier, Craig’s mother had married Fred, who did not share her love for the Church and, indeed, was not a member.

Craig confided to me his deep and abiding love for his parents and his hope that somehow, in some way, his father would be touched by the Spirit and open his heart to the gospel of Jesus Christ. I prayed for inspiration concerning how such a desire might be fulfilled. Such inspiration came, and I said to Craig, “Serve the Lord with all your heart. Be obedient to your sacred calling. Each week write a letter to your parents; and on occasion, write to Dad personally and let him know that you love him, and tell him why you’re grateful to be his son.” He thanked me and, with his mother, departed from the office.

I was not to see Craig’s mother for over eighteen months. She came to the office and, in sentences punctuated by tears, said to me, “It has been almost two years since Craig departed for his mission. He has never failed in writing a letter to us each week. Recently, my husband, Fred, stood for the first time in a testimony meeting and said, ‘All of you know that I am not a member of the Church, but something has happened to me since Craig left for his mission. His letters have touched my soul. May I share one with you?

“‘“Dear Dad,

“‘“Today we taught a choice family about the plan of salvation and blessings of exaltation in the celestial kingdom. For me it just wouldn’t be a celestial kingdom if you were not there. I’m grateful to be your son, Dad, and want you to know that I love you.

“‘“Your missionary son,

“‘“Craig”

“‘After twenty-six years of marriage, I have made my decision to become a member of the Church, for I know the gospel message is the word of God. My son’s mission has moved me to action. I have made arrangements for my wife and me to meet Craig when he completes his mission. I will be his final baptism as a full-time missionary of the Lord.’” He heard the message, he saw the light, he embraced the truth.

A young missionary with unwavering faith had participated with God in a modern-day miracle. His challenge to communicate with one whom he loved had been made more difficult by the barrier of the thousands of miles that lay between him and home. But the spirit of love spanned the vast expanse of the blue Pacific, and heart spoke to heart in divine dialogue.

No missionary stood so tall as did Craig Sudbury when, in far-off Australia, he helped his father into water waist-deep and, raising his right arm to the square, repeated those sacred words: “Fred Sudbury, having been commissioned of Jesus Christ, I baptize you in the name of the Father, and of the Son, and of the Holy Ghost” (see D\&C 20:73).

The prayer of a mother, the faith of a father, the service of a son brought forth the miracle of God.

“How beautiful are the feet of them that preach the gospel of peace, and bring glad tidings of good things!” (Rom. 10:15).

God bless us, my brethren, with missionary memories of stalwart service in the cause of Christ, I pray in His holy name, amen.

\newpage
\settalk{An Invitation to Exaltation}
\setspeaker{Thomas S. Monson}
\settalkdate{April 1988}
\talkheading{An Invitation to Exaltation}{Thomas S. Monson}

Everywhere, people are in a hurry. Jet-powered planes speed their precious human cargo across broad continents and vast oceans. Appointments must be kept, tourist attractions beckon, and friends and family await the arrival of a particular flight. Modern freeways with multiple lanes carry millions of automobiles, occupied by more millions of people, in a seemingly endless stream.

Does this pulsating, mobile ribbon of humanity ever come to a halt? Is the helter-skelter pace of life ever punctuated with moments of meditation—even thoughts of timeless truths?

When compared to eternal verities, the questions of daily living are really rather trivial. What shall we have for dinner? Is there a good movie playing tonight? Have you seen the television log? Where shall we go on Saturday? These questions pale in their significance when times of crisis arise, when loved ones are wounded, when pain enters the house of good health, or when life’s candle dims and darkness threatens. Then, truth and trivia are soon separated. The soul of man reaches heavenward, seeking a divine response to life’s greatest questions: Where did we come from? Why are we here? Where do we go after we leave this life?

Answers to these questions are not discovered within the covers of academia’s textbooks, by dialing Information, in tossing a coin, or through random selection of multiple-choice responses. These questions transcend mortality. They embrace eternity.

Where did we come from? This query is inevitably thought, if not spoken, by every parent or grandparent when a tiny infant utters its first cry. One marvels at the perfectly formed child. The tiny toes, the delicate fingers, the beautiful head, to say nothing of the hidden but marvelous circulatory, digestive, and nervous systems all testify to the truth of a divine Creator.

The Apostle Paul told the Athenians on Mars’ Hill that we are “the offspring of God” (Acts 17:29). Since we know that our physical bodies are the offspring of our mortal parents, we must probe for the meaning of Paul’s statement. The Lord has declared that “the spirit and the body are the soul of man” (D\&C 88:15). It is the spirit which is the offspring of God. The writer of Hebrews refers to Him as “the Father of spirits” (Heb. 12:9). The spirits of all men are literally His “begotten sons and daughters” (D\&C 76:24).

We note that inspired poets have, for our contemplation of this subject, written moving messages and recorded transcendent thoughts. William Wordsworth penned the truth:

\begin{verse}
Our birth is but a sleep and a forgetting:\\
The soul that rises with us, our life’s Star,\\
Hath had elsewhere its setting,\\
And cometh from afar:\\
Not in entire forgetfulness,\\
And not in utter nakedness,\\
But trailing clouds of glory do we come\\
From God, who is our home:\\
Heaven lies about us in our infancy!
\end{verse}

(“Ode: Intimations of Immortality from Recollections of Early Childhood,” lines 58–66)

Another writer described a newborn infant as “a sweet, new blossom of humanity, fresh fallen from God’s own home to flower here on earth.”

Parents, gazing down at a tiny infant or taking the hand of a growing child, ponder their responsibility to teach, to inspire, and to provide direction. While parents ponder, children and, particularly, youth ask the penetrating question: “Why are we here?” Usually, it is spoken silently to the soul and phrased: “Why am I here?”

How grateful we should be that a wise Creator fashioned an earth and placed us here, with a veil of forgetfulness of our previous existence, so that we might experience a time of testing, an opportunity to prove ourselves, and qualify for all that God has prepared for us to receive.

Clearly, one primary purpose of our existence upon the earth is to obtain a body of flesh and bones. In a thousand ways, we are privileged to choose for ourselves. Here we learn from the hard taskmaster of experience. We discern between good and evil. We differentiate as to the bitter and the sweet. We discover that decisions determine destiny.

While Paul taught the Philippians that man is called upon to “work out [his] own salvation with fear and trembling” (Philip. 2:12), the Master provided a guide we know as the Golden Rule: “Whatsoever ye would that men should do to you, do ye even so to them” (Matt. 7:12).

By obedience to God’s commandments, we can qualify for that “house” spoken of by Jesus when He declared: “In my Father’s house are many mansions. … I go to prepare a place for you … that where I am, there ye may be also” (John 14:2–3).

Contemplating such far-reaching matters, we reflect upon the helplessness of a newborn child. No better example can be found for total dependency. Needed is nourishment for the body and love for the soul. Mother provides both. She who, with her hand in the hand of God, descended into “the valley of the shadow of death” (Ps. 23:4), that you and I might come forth to life, is not in her maternal mission abandoned by God.

Several years ago, the Salt Lake City newspapers published an obituary notice of a close friend—a mother and wife taken by death in the prime of her life. I visited the mortuary and joined a host of persons gathered to express condolence to the distraught husband and motherless children. Suddenly the smallest child, Kelly, recognized me and took my hand in hers.

“Come with me,” she said; and she led me to the casket in which rested the body of her beloved mother. “I’m not crying, Brother Monson, and neither must you. My mommy told me many times about death and life with Heavenly Father. I belong to my mommy and my daddy. We’ll all be together again.”

Through tear-moistened eyes, I recognized a beautiful and faith-filled smile. To my young friend, whose tiny hand yet clasped mine, there would never be a hopeless dawn. Sustained by her unfailing testimony, knowing that life continues beyond the grave, she, her father, her brothers, her sisters, and indeed all who share this knowledge of divine truth, can declare to the world, “Weeping may endure for a night, but joy cometh in the morning” (Ps. 30:5).

Life moves on. Youth follows childhood, and maturity comes ever so imperceptibly.

We treasure the inspired thought:

\begin{verse}
God is a Father.\\
Man is a brother.\\
Life is a mission\\
And not a career.
\end{verse}

(In Stephen L Richards, Where Is Wisdom? Addresses of Stephen L Richards, Salt Lake City: Deseret Book Co., 1955, p. 74)

God, our Father, and Jesus Christ, our Lord, have marked the way to perfection. They beckon us to follow eternal verities and to become perfect, as they are perfect (see Matt. 5:48; 3 Ne. 12:48). We remember the inquiring lawyer who asked:

“Master, which is the great commandment in the law?

“Jesus said unto him, Thou shalt love the Lord thy God with all thy heart, and with all thy soul, and with all thy mind.

“This is the first and great commandment.

“And the second is like unto it, Thou shalt love thy neighbour as thyself” (Matt. 22:36–39).

The Apostle Paul likened life to a race with a clearly defined goal. To the Saints at Corinth he urged:

“Know ye not that they which run in a race run all, but one receiveth the prize? So run, that ye may obtain” (1 Cor. 9:24).

In our zeal, let us not overlook the sage counsel from Ecclesiastes: “The race is not to the swift, nor the battle to the strong” (Eccl. 9:11). Actually the prize belongs to him who endures to the end.

When I reflect on the race of life, I remember another type of race, even from childhood days. When I was about ten, my boyfriends and I would take pocketknives in hand and, from the soft wood of a willow tree, fashion small toy boats. With a triangular-shaped cotton sail in place, each would launch his crude craft in the race down the relatively turbulent waters of the Provo River. We would run along the river’s bank and watch the tiny vessels sometimes bobbing violently in the swift current and at other times sailing serenely as the water deepened.

During such a race, we noted that one boat led all the rest toward the appointed finish line. Suddenly, the current carried it too close to a large whirlpool, and the boat heaved to its side and capsized. Around and around it was carried, unable to make its way back into the main current. At last it came to an uneasy rest at the end of the pool, amid the flotsam and jetsam that surrounded it.

The toy boats of childhood had no keel for stability, no rudder to provide direction, and no source of power. Inevitably their destination was downstream—the path of least resistance.

Unlike toy boats, we have been provided divine attributes to guide our journey. We enter mortality not to float with the moving currents of life, but with the power to think, to reason, and to achieve.

Our Heavenly Father did not launch us on our eternal voyage without providing the means whereby we could receive from Him guidance to ensure our safe return. Yes, I speak of prayer. I speak, too, of the whisperings from that still, small voice within each of us; and I do not overlook the holy scriptures, written by mariners who successfully sailed the seas we too must cross.

At some period in our mortal mission, there appears the faltering step, the wan smile, the pain of sickness—even the fading of summer, the approach of autumn, the chill of winter, and the experience we call death.

Every thoughtful person has asked himself the question best phrased by Job of old: “If a man die, shall he live again?” (Job 14:14). Try as we may to put the question out of our thoughts, it always returns. Death comes to all mankind. It comes to the aged as they walk on faltering feet. Its summons is heard by those who have scarcely reached midway in life’s journey, and often it hushes the laughter of little children.

But what of an existence beyond death? Is death the end of all? Such a question was asked of me by a young husband and father who lay dying. I turned to the Book of Mormon and, from the book of Alma, read to him these words:

“Now, concerning the state of the soul between death and the resurrection—Behold, it has been made known unto me by an angel, that the spirits of all men, as soon as they are departed from this mortal body, yea, the spirits of all men, whether they be good or evil, are taken home to that God who gave them life.

“And then shall it come to pass, that the spirits of those who are righteous are received into a state of happiness, which is called paradise, a state of rest, a state of peace, where they shall rest from all their troubles and from all care, and sorrow” (Alma 40:11–12).

My young friend through moist eyes and with an expression of profound gratitude whispered a silent, but eloquent, “Thank you.”

After the body of Jesus had lain in the tomb for three days, the spirit again entered, and the resurrected Redeemer walked forth clothed with an immortal body of flesh and bones.

The answer to Job’s question, “If a man die, shall he live again?” came when Mary and others approached the tomb and saw two men in shining garments who spoke to them: “Why seek ye the living among the dead?

“He is not here, but is risen” (Luke 24:5–6).

Testimonies of the resurrected Lord provide comfort and understanding.

First, from the Apostle Paul:

“Christ died for our sins according to the scriptures; … he was buried, and … he rose again the third day: … he was seen of Cephas, then of the twelve: … he was seen of above five hundred brethren at once; … he was seen of James; then of all the apostles. And last of all he was seen of me also, as of one born out of due time” (1 Cor. 15:3–8).

Second, from the combined testimony of twenty-five hundred of his other sheep, as recorded in the Book of Mormon, Another Testament of Jesus Christ, the resurrected Lord “spake unto them saying:

“Arise and come forth unto me, that ye may thrust your hands into my side, and also that ye may feel the prints of the nails in my hands and in my feet, that ye may know that I am the God of Israel, and the God of the whole earth, and have been slain for the sins of the world. …

“And when they had all gone forth and had witnessed for themselves, they did cry out with one accord, saying:

“Hosanna! Blessed be the name of the Most High God! And they did fall down at the feet of Jesus, and did worship him” (3 Ne. 11:13–14, 16–17).

Third, from Joseph Smith: “After the many testimonies which have been given of him, this is the testimony, last of all, which we give of him: That He lives!

“For we saw him, even on the right hand of God; and we heard the voice bearing record that he is the Only Begotten of the Father—

“That by him, and through him, and of him, the worlds are and were created, and the inhabitants thereof are begotten sons and daughters unto God” (D\&C 76:22–24).

As the result of Christ’s victory over the grave, we shall all be resurrected. This is the redemption of the soul. Paul wrote:

“There are … celestial bodies, and bodies terrestrial: but the glory of the celestial is one, and the glory of the terrestrial is another.

“There is one glory of the sun, and another glory of the moon, and another glory of the stars: for one star differeth from another star in glory.

“So also is the resurrection of the dead” (1 Cor. 15:40–42).

It is the celestial glory which we seek. It is in the presence of God we desire to dwell. It is a forever family in which we want membership. Such blessings must be earned.

Where did we come from? Why are we here? Where do we go after this life? No longer need these universal questions remain unanswered. Our Heavenly Father rejoices for those who keep His commandments. He is concerned also for the lost child, the tardy teenager, the wayward youth, the delinquent parent. Tenderly the Master speaks to these, and indeed to all: “Come back. Come home. Come unto me.” What eternal joy awaits when we accept His divine invitation to exaltation.

I testify He is a teacher of truth—but He is more than a teacher. He is the exemplar of the perfect life—but He is more than an exemplar. He is the great physician—but He is more than a physician. He is the literal Savior of the World, the Son of God, the Prince of Peace, the Holy One of Israel, even the risen Lord, who declared: “I am Jesus Christ, whom the prophets testified shall come into the world.

“… I am the light and the life of the world” (3 Ne. 11:10–11).

“I am the first and the last; I am he who liveth, I am he who was slain; I am your advocate with the Father” (D\&C 110:4).

As his witness I testify to you that He lives, in the name of Jesus Christ, amen.

\newpage
\settalk{Happiness through Service}
\setspeaker{Thomas S. Monson}
\settalkdate{April 1988}
\talkheading{Happiness through Service}{Thomas S. Monson}

President Benson has suggested that I bear my testimony to you at this time. I am pleased to again express my witness that God lives, that Jesus is the Christ, the Son of the Living God, that this work is true, and that happiness comes through serving our Heavenly Father and serving our fellowmen.

If from this conference we can gain a new feeling of closeness to the Savior and a testimony of his divine mission, and on this Easter Sunday if we can, renewed with the spirit of the Resurrection, go forward in looking after his sheep and our family responsibilities and our church duties in a way which will be pleasing to our Heavenly Father, we ourselves will be abundantly blessed.

God bless you, my brothers and sisters, in all of your incomings and outgoings. May you have peace in your hearts, may you have tranquility in your homes, and may you have the Spirit of the Lord Jesus Christ in your souls, I ask in the name of Jesus Christ, amen.

\newpage
\settalk{You Make a Difference}
\setspeaker{Thomas S. Monson}
\settalkdate{April 1988}
\talkheading{You Make a Difference}{Thomas S. Monson}

David declares in one of his beautiful and moving psalms, “O Lord our Lord, how excellent is thy name in all the earth! …

“When I consider thy heavens, the work of thy fingers, the moon and the stars, which thou hast ordained;

“What is man, that thou art mindful of him?” (Ps. 8:1, 3–4).

Job, that righteous man of old, joined in the question when he asked, “What is man, that thou shouldest magnify him? and that thou shouldest set thine heart upon him?” (Job 7:17).

One need not grope for answers to these penetrating questions when in your presence here in the historic Tabernacle or with you in the many meeting places throughout the world where you have assembled. “Ye are a chosen generation, a royal priesthood, an holy nation” (1 Pet. 2:9). “Ye … are … a spiritual house, an holy priesthood” (1 Pet. 2:5).

As bearers of the priesthood, we have been placed on earth in troubled times. We live in a complex world with currents of conflict everywhere to be found. Political machinations ruin the stability of nations, despots grasp for power, and segments of society seem forever downtrodden, deprived of opportunity and left with a feeling of failure.

We who have been ordained to the priesthood of God can make a difference. When we qualify for the help of the Lord, we can build boys, we can mend men, we can accomplish miracles in His holy service. Our opportunities are without limit.

Though the task looms large, we are strengthened by the truth: “The greatest force in this world today is the power of God as it works through man.” If we are on the Lord’s errand, we are entitled to the Lord’s help. That divine help, however, is predicated upon our worthiness. To sail safely the seas of mortality, to perform a human rescue mission, we need the guidance of that eternal mariner—even the great Jehovah. We reach out, we reach up, to obtain heavenly help.

Are our reaching hands clean? Are our yearning hearts pure? Looking backward in time through the pages of history, we glean a lesson on worthiness from the words of the dying King Darius.

“Darius … through the proper rites had been recognized as legitimate King of Egypt; his rival Alexander had been declared the legitimate Son of Amon—he too was Pharaoh. … Alexander found the defeated Darius on the point of death in his tent, and … laid his hands upon his head to heal him, commanding him to arise and resume his kingly power, and concluding his blessing: ‘I swear unto thee, Darius, by all the gods, that I do these things truly and without faking.’ [Darius] replied with a gentle rebuke: ‘Alexander my boy … do you think you can touch heaven with those hands of yours?’” (in Hugh Nibley, Abraham in Egypt, Salt Lake City: Deseret Book Co., 1981, p. 192).

An inspiring lesson is learned from a “Viewpoint” article which appeared recently in the Church News section of the Deseret News. May I quote:

“To some it may seem strange to see ships of many nations loading and unloading cargo along the docks at Portland, Oregon. That city is 100 miles from the ocean. Getting there involves a difficult, often turbulent passage over the bar guarding the Columbia River and a long trip up the Columbia and Willamette Rivers.

“But ship captains like to tie up at Portland. They know that as their ships travel the seas, a curious salt water shellfish called a barnacle fastens itself to the hull and stays there for the rest of its life, surrounding itself with a rock-like shell. As more and more barnacles attach themselves, they increase the ship’s drag, slow its progress, decrease its efficiency.

“Periodically, the ship must go into dry dock, where with great effort the barnacles are chiseled or scraped off. It’s a difficult, expensive process that ties up the ship for days. But not if the captain can get his ship to Portland. Barnacles can’t live in fresh water. There, in the sweet, fresh waters of the Willamette or Columbia, the barnacles loosen and fall away, and the ship returns to its task lightened and renewed.

“Sins are like those barnacles. Hardly anyone goes through life without picking up some. They increase the drag, slow our progress, decrease our efficiency. Unrepented, building up one on another, they can eventually sink us.

“In His infinite love and mercy, our Lord has provided a harbor where, through repentance, our barnacles fall away and are forgotten. With our souls lightened and renewed, we can go efficiently about our work and His” (“Harbor of Forgiveness,” 30 Jan. 1988, p. 16).

A loving Heavenly Father has provided for our guidance models to follow, men who made a difference in their own times. I choose to call these noble souls “pioneers.” Webster defines a pioneer: “One who goes before, showing others the way to follow.”

With faith as their moving power, they sailed upstream against the currents of doubt which surrounded them. We cannot help but be inspired in our efforts as we remember their examples.

From Nephi: “I will go and do the things which the Lord hath commanded” (1 Ne. 3:7).

From Samuel: “To obey is better than sacrifice, and to hearken than the fat of rams” (1 Sam. 15:22).

From Paul: “For I am not ashamed of the gospel of Christ: for it is the power of God unto salvation” (Rom. 1:16).

From Job: “I know that my redeemer liveth” (Job 19:25).

From Joseph: “I am calm as a summer’s morning; I have a conscience void of offense towards God, and towards all men” (D\&C 135:4).

These noble leaders made a difference in their own times. What about today? How about me?

The world felt the quickening pace of activity when President Spencer W. Kimball declared, “We must lengthen our stride.” He stepped forward and the Church followed.

When President Ezra Taft Benson warned that we had neglected the Book of Mormon and urged every member to read and study this sacred volume, new printing presses were required to produce more and more copies of the book, as boys and girls and men and women followed the prophet in his own reading and in his inspired declaration. Every day letters arrive at the President’s office which testify to the enrichment of lives which comes from reading the Book of Mormon. They tell of families united, goals attained, and souls saved. Such is the power of a prophet.

We do not have a monopoly on goodness. There are God-fearing men and women in all nations who influence for good those with whom they associate. I think of the founder of Scouting, even Lord Baden-Powell, and those who teach and live the principles he advocated. Who can measure the far-reaching effect on human lives of the Scout oath:

“On my honor I will do my best to do my duty to God and my country and to obey the Scout Law; to help other people at all times; to keep myself physically strong, mentally awake and morally straight.”

Impossible of calculation is the result for good when men and boys observe the Scout law: trustworthy, loyal, helpful, friendly, courteous, kind, obedient, cheerful, thrifty, brave, clean and reverent.

The influence of your personal testimonies is ever so far-reaching. The Lord instructed, “The testimony which ye have borne is recorded in heaven for the angels to look upon; and they rejoice over you, and your sins are forgiven you” (D\&C 62:3).

He also cautioned us, “With some I am not well pleased, for they will not open their mouths, but they hide the talent which I have given unto them, because of the fear of man” (D\&C 60:2).

You never know when your turn will come to comply with the admonition of Peter to “be ready always to give an answer to every man that asketh you a reason of the hope that is in you” (1 Pet. 3:15).

Some years ago I had the opportunity to address a business convention in Dallas, Texas, sometimes called “the city of churches.” After the convention, I took a sightseeing bus ride about the city’s suburbs. Our driver would comment, “On the left you see the Methodist church,” or “There on the right is the Catholic cathedral.”

As we passed a beautiful red brick building situated upon a hill, the driver exclaimed, “That building is where the Mormons meet.” A lady from the rear of the bus asked, “Driver, can you tell us something more about the Mormons?” The driver steered the bus to the side of the road, turned about in his seat, and replied, “Lady, all I know about the Mormons is that they meet in that red brick building. Is there anyone on this bus who knows anything about the Mormons?”

I gazed at the expression on each person’s face for some sign of recognition, some desire to comment. I found nothing—not a sign. Then I realized the truth of the statement, “When the time for decision arrives, the time for preparation is past.” For the next fifteen minutes I had the privilege of sharing with others my testimony concerning The Church of Jesus Christ of Latter-day Saints.

The seeds of testimony frequently do not at once take root and flower. Bread cast upon the water returns, at times, only after many days.

I answered the ring of my telephone one evening to hear a voice ask, “Are you related to an Elder Monson who years ago served in the New England Mission?” I answered that such was not the case. The caller introduced himself as a Brother Leonardo Gambardella and then mentioned that an Elder Monson and an Elder Bonner called at his home long ago and bore their personal testimonies to him. He had listened but had done nothing further to apply their teachings. Subsequently he moved to California where, after thirteen years, he again found the truth and was converted and baptized. Brother Gambardella then asked if there were a way he could reach these elders who had first visited with him, that he might express to them his profound gratitude for their testimonies, which had remained with him.

I checked the records. I located the elders. Can you imagine their surprise when, now married with families of their own, I telephoned them and told them the good news—even the culmination of their early efforts. They remembered Brother Gambardella and, at my suggestion, telephoned him to extend their congratulations and welcome him into the Church.

You can make a difference. Whom the Lord calls, the Lord qualifies. This promise extends not only to missionaries, but also to home teachers, quorum leaders, presidents of branches, and bishops of wards. When we qualify ourselves by our worthiness, when we strive with faith nothing wavering to fulfill the duties appointed to us, when we seek the inspiration of the Almighty in the performance of our responsibilities, we can achieve the miraculous.

Brethren, let us hearken to the hymn, “Improve the Shining Moments”:

\begin{verse}
Time flies on wings of lightning;\\
We cannot call it back.\\
It comes, then passes forward\\
Along its onward track.
\end{verse}

(Hymns, 1985, no. 226)

As we leave this general priesthood meeting, let us all determine to shed any barnacles of sin, to prepare for our time of opportunity, and to honor the priesthood we bear through the service we render, the lives we bless, and the souls we are privileged to help save. You are “a chosen generation, a royal priesthood, an holy nation” (1 Pet. 2:9), and you can make a difference. To these truths I testify, in the name of Jesus Christ, amen.

\newpage
\settalk{Goal beyond Victory}
\setspeaker{Thomas S. Monson}
\settalkdate{October 1988}
\talkheading{Goal beyond Victory}{Thomas S. Monson}

Years ago, many of us participated as players or observers in the all-Church basketball tournaments and later in the softball tournaments. The most coveted prize was not to be adjudged first-place winner, but rather to receive the sportsmanship award. The applause of the audience was louder and longer, the smiles broader and more universal. A goal beyond victory had been won.

Lately we have received at the Office of the First Presidency letters which tell of serious arguments on the sports floor or playing field, name-calling by parents, abuse of referees, and all that characterizes poor sportsmanship. We have room for improvement, brethren, and improve we must.

In the videotape produced by the Church and entitled The Church Sports Official, there is featured this truth from the First Presidency: “Church sports activities have a unique central purpose much higher than the development of physical prowess, or even victory itself. It is to strengthen faith, build integrity, and develop in each participant the attributes of his maker.”

Brethren, it is difficult to achieve this objective if winning overshadows participation. The recreation halls in our many buildings are constructed through the tithes of the members of the Church. It is only fair that all worthy young men and young women have an opportunity to play, to learn, to develop, and to achieve.

It is not our objective to produce clones of Larry Bird or Magic Johnson—or even John Wooden or Pat Riley. When you put a player in a suit, put him in the game. Basketball begins soon. Let our teams of young men and young women be counseled appropriately. And a word or two for the spectators and coaches would not be amiss.

If I might add a personal touch, I share with you an experience that embarrassed, a game that was lost, and a lesson in not taking ourselves too seriously.

First, in a basketball game when the outcome was in doubt, the coach sent me onto the playing floor right after the second half began. I took an in-bounds pass, dribbled the ball toward the key, and let the shot fly. Just as the ball left my fingertips, I realized why the opposing guards did not attempt to stop my drive: I was shooting for the wrong basket! I offered a silent prayer: “Please, Father, don’t let that ball go in.” The ball rimmed the hoop and fell out.

From the bleachers came the call: “We want Monson, we want Monson, we want Monson—out!” The coach obliged.

I never was a basketball star. What timing—to be a freshman at the University of Utah when All-Americans Arnie Ferrin and Vern Gardner dominated the boards.

I fared much better at fast-pitch softball. My most memorable experience in softball was a thirteen-inning game I pitched in Salt Lake City on a hot Memorial Day. The game was scheduled for just seven innings, but the tied score could not be broken. In the last of the thirteenth, with two men out and a runner on third, the batter hit a high pop fly to left field. The catch was certain, I thought. And yet the ball fell through the hands of the left fielder. For thirty-eight years I have teased my friend who dropped the ball. I have promised myself I will never do so again. I’m not even going to mention his name. After all, he, too, remembers. It was only a game.

On another occasion, while pitching a game at Pioneer Park, I was absolutely stunned to see that the other team had placed a one-armed batter at the plate. Now how does a pitcher deliver the pitch to such an opponent? I tossed a gentle lob over the plate. To my amazement, the batter knocked a single, right over the second baseman’s head. My temper flared. The next batter was a returned missionary from Mexico, Homer Proctor, six foot two and about 210 pounds. I pitched him fast, high, and inside. On the first pitch, he lifted the ball right out of the park for a home run. I shall ever remember the smile of that one-armed runner, Bernell Hales, as he passed second and third and gleefully streaked for home. I felt like crying, but I broke out laughing, as did each player on both sides. We had a wonderful time.

Brethren, let’s take the necessary steps to rekindle sportsmanship, to emphasize participation, and to strive for the development of a Christlike character in each individual.

Now, there are other phases of the Lord’s work in which all members can participate, in which the growth of character is assured and the promise of life eternal bestowed. One such endeavor is referred to as the welfare program. Actually, the language of King Benjamin from the book of Mosiah provides a perfect scriptural description, even a solemn charge to each of us:

“For the sake of retaining a remission of your sins from day to day, that ye may walk guiltless before God—I would that ye should impart of your substance to the poor, every man according to that which he hath, such as feeding the hungry, clothing the naked, visiting the sick and administering to their relief, both spiritually and temporally.” (Mosiah 4:26.)

President Marion G. Romney spoke concerning the funding of caring for the needy when he said: “It has been, and now is, the desire and the objective of the Church to obtain from fast offerings the necessary funds to meet the cash needs of the welfare program. … At the present time we are not meeting this objective. We can, we ought, and we must do better. If we will double our fast offerings, we shall increase our own prosperity, both spiritually and temporally. This the Lord has promised, and this has been the record.” (“Basics of Church Welfare,” talk given to the Priesthood Board, 6 Mar. 1974, p. 10.)

Are we generous in the payment of our fast offerings? That we should be so was taught by President Joseph F. Smith. He declared that it is incumbent upon every Latter-day Saint to give to his bishop on fast day an amount equivalent to the food that he and his family would consume for the day and, if possible, a liberal donation to be so reserved and donated to the poor. (See Improvement Era, Dec. 1902, p. 148.)

President Spencer W. Kimball suggested that, in our generosity, we go beyond a minimum amount. He urged that we “give, instead of the amount we saved by our two meals of fasting, perhaps much, much more—ten times more where we are in a position to do it.” (In Conference Report, Apr. 1974, p. 184.)

The generous response of the Latter-day Saints in times of crisis is legendary. Many will remember the emergency aid provided our needy Saints in Europe following World War II. President Ezra Taft Benson directed this effort.

More recently, this generosity helped to avert starvation in Africa. Irrigation projects, producing wells, and improved agricultural methods are all part of a dream come true. Similarly, at the time of the Teton Dam disaster in Idaho, the response of the members to the call of need was overwhelming.

Today, in lands far away and right here in Salt Lake City, there are those who suffer hunger, who know want and are acquainted with poverty. Ours is the opportunity and the sacred privilege to relieve this hunger, to meet this want, to eliminate this poverty.

The Lord provided the way when He declared, “And the storehouse shall be kept by the consecrations of the church; and widows and orphans shall be provided for, as also the poor.” (D\&C 83:6.) Then the reminder, “But it must needs be done in mine own way.” (D\&C 104:16.)

In the vicinity where I once lived and served, we operated a poultry project. Most of the time it was an efficiently operated project, supplying to the storehouse thousands of dozens of fresh eggs and hundreds of pounds of dressed poultry. On a few occasions, however, the experience of being volunteer city farmers provided not only blisters on the hands, but also frustration of heart and mind.

For instance, I shall ever remember the time we gathered together the Aaronic Priesthood young men to really give the project a spring cleaning treatment. Our enthusiastic and energetic throng assembled at the project and in a speedy fashion uprooted, gathered, and burned large quantities of weeds and debris. By the light of the glowing bonfires, we ate hot dogs and congratulated ourselves on a job well done. The project was now neat and tidy. However, there was just one disastrous problem: The noise and the fires had so disturbed the fragile and temperamental population of five thousand laying hens that most of them went into a sudden molt and ceased laying. Thereafter we tolerated a few weeds, that we might produce more eggs.

No member of the Church who has canned peas, topped beets, hauled hay, or watered corn in such a cause ever forgets or regrets the experience of helping provide for those in need.

Sharing with others that which we have is not new to our generation. We need but to turn to the account found in 1 Kings to appreciate anew the principle that when we follow the counsel of the Lord, when we care for those in need, the outcome benefits all. There we read that a most severe drought had gripped the land. Famine followed. Elijah the prophet received from the Lord what to him must have been an amazing instruction: “Get thee to Zarephath … : behold, I have commanded a widow woman there to sustain thee.” (17:9.) When he had found the widow, Elijah declared: “Fetch me, I pray thee, a little water in a vessel, that I may drink.

“And as she was going to fetch it, he called to her, and said, Bring me, I pray thee, a morsel of bread in thine hand.” (Vs. 10–11.)

Her response described her pathetic situation as she explained that she was preparing a final and scanty meal for her son and for herself, and then they would die. (See v. 12.)

How implausible to her must have been Elijah’s response:

“Fear not; go and do as thou hast said: but make me thereof a little cake first, and bring it unto me, and after make for thee and for thy son.

“For thus saith the Lord God of Israel, The barrel of meal shall not waste, neither shall the cruse of oil fail, until the day that the Lord sendeth rain upon the earth.

“And she went and did according to the saying of Elijah: and she, and he, and her house, did eat many days.

“And the barrel of meal wasted not, neither did the cruse of oil fail.” (Vs. 13–16.) This is the faith that has ever motivated and inspired the welfare plan of the Lord.

Industry, thrift, self-reliance continue as guiding principles of this effort. As a people, we should avoid unreasonable debt. In a message which Elder Ezra Taft Benson delivered at a general conference more than thirty years ago, he instructed: “In the book of Kings we read about a woman who came weeping to … the prophet [of the Lord]. Her husband had died, and she owed a debt that she could not pay; and the creditor was on his way to take her two sons and sell them as slaves.

“By a miracle, [the prophet] enabled her to acquire a goodly supply of oil. And he said to her:

“‘Go, sell the oil, and pay thy debt, and live.’” (In Conference Report, Apr. 1957, p. 53.)

“Pay thy debt, and live.” (2 Kgs. 4:7.) What wise counsel for us today! Remember, the wisdom of God may appear as foolishness to men, but the greatest single lesson we can learn in mortality is that when God speaks and a man obeys, that man will always be right.

We should remember that the best storehouse system would be for every family to have a year’s supply of needed food, clothing, and, where possible, the other necessities of life. In the early Church, Paul wrote to Timothy, “If any provide not for his own, and specially for those of his own house, he hath denied the faith, and is worse than an infidel.” (1 Tim. 5:8.)

It is our sacred duty to care for our families. Often we see what might be called “parent neglect.” Too frequently the emotional, social, and, in some instances, even the material essentials of life are not provided by children to their aged parents. This is displeasing to the Lord.

The Lord’s storehouse includes the time, talents, skills, compassion, consecrated material, and financial means of faithful Church members. These resources are available to the bishop in assisting those in need. Our bishops have the responsibility to learn how to use properly these resources.

May I suggest in summary form five basic guidelines:

Let me illustrate with a sacred experience which brought these guidelines together in blessing the lives of those in need.

While serving as a bishop, one cold winter day I visited an elderly couple who lived in a two-room duplex. The modest home was heated by a small coal-burning Heatrola. As I approached the home, I met the 82-year-old husband, his aged body bent in the driving snow as he gathered a few pieces of wet coal from his exposed supply of fuel. I helped him with his burden but made a solemn resolve to do more.

I prayed and pondered, seeking a solution. Step by step the inspiration came. In the ward was an unemployed carpenter. He had no fuel for his furnace but was too proud to receive the stoker slack he needed to keep his house warm. I suggested to the carpenter a way he could work for the help he received. Would he build a coal shed for a couple in need? “Of course,” he replied.

Now where were we to obtain the materials? I approached the proprietors of a local lumberyard from whom we frequently purchased products. I remember saying to the men, “How would the two of you like to paint a bright spot on your souls this winter day?” Not knowing exactly what I meant, they agreed readily. They were invited to donate the lumber and hardware for the coal shed.

Within days the project was completed. I was invited to inspect the outcome. The coal shed was simply beautiful in its sleek covering of battleship-gray paint. The carpenter, who was a high priest, testified that he had actually felt inspired as he labored on this modest shed.

My older friend, with obvious appreciation, stroked the wall of the sturdy structure. He pointed out to me the wide door, the shiny hinges, and then opened to my view the supply of dry coal which filled the shed. In a voice filled with emotion, he said in words I shall ever treasure, “Bishop, take a look at the finest coal shed a man ever had.” Its beauty was only surpassed by the pride in the builder’s heart. And the elderly recipient labored each day at the ward chapel, dusting the benches, vacuuming the carpet runners, arranging the hymnbooks. He, too, worked for that which he had received.

Once again, the welfare plan of the Lord had blessed the lives of His children.

May our Heavenly Father guide the priesthood of this church, that we may be obedient to the revelation of the Lord to the Prophet Joseph in which we are charged to “remember in all things the poor and the needy, the sick and the afflicted, for he that doeth not these things, the same is not my disciple.” (D\&C 52:40.)

We will qualify as His disciples when we hear and heed the counsel from Isaiah describing the true fast, the spirit and the promise of the welfare effort:

“Is it not to deal thy bread to the hungry, and that thou bring the poor that are cast out to thy house? when thou seest the naked, that thou cover him; and that thou hide not thyself from thine own flesh?

“Then shall thy light break forth as the morning, and thine health shall spring forth speedily: and thy righteousness shall go before thee; the glory of the Lord shall be thy rereward.

“Then shalt thou call, and the Lord shall answer; thou shalt cry, and he shall say, Here I am. …

“And the Lord shall guide thee continually, and satisfy thy soul in drought, … and thou shalt be like a watered garden, and like a spring of water, whose waters fail not.” (Isa. 58:7–9, 11.)

May this be our blessing is my prayer, in the name of Jesus Christ, amen.

\newpage
\settalk{Hallmarks of a Happy Home}
\setspeaker{Thomas S. Monson}
\settalkdate{October 1988}
\talkheading{Hallmarks of a Happy Home}{Thomas S. Monson}

I pray for the Spirit of the Lord to be with me. “Happiness is the object and design of our existence; and will be the end thereof, if we pursue the path that leads to it; and this path is virtue, uprightness, faithfulness, holiness, and keeping all the commandments of God.” (Joseph Smith, Teachings of the Prophet Joseph Smith, sel. Joseph Fielding Smith, Salt Lake City: Deseret Book Co., 1938, pp. 255–56.)

This description of such a universal goal was provided by the Prophet Joseph Smith. It was relevant then. It is relevant now. With such a clear road map to follow, why then are there so many unhappy people? Frequently, frowns outnumber smiles and despair dampens joy. We live so far below the level of our divine possibilities. Some become confused by materialism, entangled by sin, and lost among the passing parade of humanity. Others cry out in the words of the convert of Philip of old: “How can I [find my way], except some man should guide me?” (Acts 8:31.)

Happiness does not consist of a glut of luxury, the world’s idea of a “good time.” Nor must we search for it in faraway places with strange-sounding names. Happiness is found at home.

All of us remember the home of our childhood. Interestingly, our thoughts do not dwell on whether the house was large or small, the neighborhood fashionable or downtrodden. Rather, we delight in the experiences we shared as a family. The home is the laboratory of our lives, and what we learn there largely determines what we do when we leave there.

Mrs. Margaret Thatcher, prime minister of Great Britain, expressed the profound philosophy: “The family is the building block of society. It is a nursery, a school, a hospital, a leisure centre, a place of refuge and a place of rest. It encompasses the whole of the society. It fashions our beliefs; it is the preparation for the rest of our life.” (Nicholas Wood, “Thatcher Champions the Family,” London Times, 26 May 1988, p. 24.)

“Home is where the heart is.” It does take “a heap o’ livin’” to make a house a home (Edgar A. Guest, “Home,” in The Family Book of Best-Loved Poems, ed. David L. George, Garden City, N.Y.: Doubleday, 1952, p. 151–52.) “Home, home, sweet, sweet home, Be it ever so humble, there’s no place like home.” (Hymns, 1948, no. 185.) We turn from the reverie of such pleasant recollections. We contemplate parents gone, family grown, childhood vanished. Slowly but surely we face the truth: We are responsible for the home we build. We must build wisely, for eternity is not a short voyage. There will be calm and wind, sunlight and shadows, joy and sorrow. But if we really try, our home can be a bit of heaven here on earth. The thoughts we think, the deeds we do, the lives we live influence not only the success of our earthly journey; they mark the way to our eternal goals.

Happy homes come in a variety of appearances. Some feature large families with father, mother, brothers, and sisters living together in a spirit of love. Others consist of a single parent with one or two children, while other homes have but one occupant. There are, however, identifying features which are to be found in a happy home, whatever the number or description of its family members. I refer to these as “Hallmarks of a Happy Home.” They consist of:

“Prayer is the soul’s sincere desire, Uttered or unexpressed.” (Hymns, 1985, no. 145.) So universal is its application, so beneficial its result, that prayer qualifies as the number-one hallmark of a happy home. As parents listen to the prayer of a child, they too draw close to God. These little ones, who so recently have been with their Heavenly Father, have no inhibitions in expressing to Him their feelings, their wishes, their thanks.

Family prayer is the greatest deterrent to sin, and hence the most beneficent provider of joy and happiness. The old saying is yet true: “The family that prays together stays together.”

“It is not possible for a married couple to reach happiness with eyes fixed on different stars; … they must set up a single ideal and work toward [it]. … Cease cherishing impossible fancies of impossible futures. Take the best of [your] dreams and fit them to life as it comes every day.” (Temple Bailey, “The Bride Who Makes Her Dreams Come True,” in Ladies’ Home Journal, 1912.)

On October 7, my wife, Frances, and I will have been married forty years. Our marriage took place just to the east of us in the holy temple. He who performed the ceremony, Benjamin Bowring, counseled us: “May I offer you newlyweds a formula which will ensure that any disagreement you may have will last no longer than one day? Every night kneel by the side of your bed. One night, Brother Monson, you offer the prayer, aloud, on bended knee. The next night you, Sister Monson, offer the prayer, aloud, on bended knee. I can then assure you that any misunderstanding that develops during the day will vanish as you pray. You simply can’t pray together and retain any but the best of feelings toward one another.”

When I was called to the Council of the Twelve just twenty-five years ago this weekend, President McKay asked me concerning my family. I related to him this guiding formula of prayer and bore witness to its validity. He sat back in his large leather chair and, with a smile, responded, “The same formula that has worked for you has blessed the lives of my family during all the years of our marriage.”

Prayer is the passport to spiritual power.

A second hallmark of a happy home is discovered when home is a library of learning. An essential part of our learning library will be good books.

\begin{verse}
Books are keys to wisdom’s treasure;\\
Books are gates to lands of pleasure;\\
Books are paths that upward lead;\\
Books are friends. Come, let us read.
\end{verse}

(Emilie Poulsson.)

Reading is one of the true pleasures of life. In our age of mass culture, when so much that we encounter is abridged, adapted, adulterated, shredded, and boiled down, it is mind-easing and mind-inspiring to sit down privately with a congenial book.

James A. Michener, prominent author, suggests: “A nation becomes what its young people read in their youth. Its ideals are fashioned then, its goals strongly determined.”

The Lord counseled, “Seek ye out of the best books words of wisdom; seek learning, even by study and also by faith.” (D\&C 88:118.)

The standard works offer the library of learning of which I speak. We must be careful not to underestimate the capacity of children to read and to understand the word of God.

A few months ago we took our grandchildren on an escorted tour of the Church printing facilities. There, all of us saw the missionary edition of the Book of Mormon coming off the delivery line—printed, bound, and trimmed, ready for reading. I said to a young grandson, “The operator says that you can remove one copy of the Book of Mormon to be your very own. You select the copy, and it will then be yours.”

Removing one finished copy of the book, he clutched it to his breast and said with sincerity, “I love the Book of Mormon. This is my book.”

I really don’t remember other events of that day, but none of us who was there will ever forget the honest expression from the heart of a child.

As parents, we should remember that our lives may be the book from the family library which the children most treasure. Are our examples worthy of emulation? Do we live in such a way that a son or a daughter may say, “I want to follow my dad,” or “I want to be like my mother”? Unlike the book on the library shelf, the covers of which shield the contents, our lives cannot be closed. Parents, we truly are an open book.

A third hallmark of a happy home is a legacy of love.

As a small boy, I enjoyed visiting the home of my grandmother on Bueno Avenue here in Salt Lake City. Grandmother was always so happy to see us and to draw us close to her. Seated on her lap, we listened as she read to us.

Her youngest son and his wife now occupy that same home. I visited there recently. The fireplug on the curb seemed so small compared to its size when I climbed its lofty heights those long years ago. The friendly porch was the same, the quiet, peaceful atmosphere not altered. Hanging on the kitchen wall was a framed expression which my aunt had embroidered. It carried a world of practical application: “Choose your love; love your choice.” She who prepared that message is now in frail health. Her husband, Ray, cares for her constantly and is the epitome of faithful and enduring love. She reciprocates in her own way. They live the lesson they framed.

Seemingly little lessons of love are observed by children as they silently absorb the examples of their parents. My own father, a printer, worked long and hard practically every day of his life. I’m certain that on the Sabbath he would have enjoyed just being at home. Rather, he visited elderly family members and brought cheer into their lives.

One was his uncle, who was crippled by arthritis so severe that he could not walk or care for himself. On a Sunday afternoon Dad would say to me, “Come along, Tommy; let’s take Uncle Elias for a short drive.” Boarding the old 1928 Oldsmobile, we would proceed to Eighth West, where, at the home of Uncle Elias, I would wait in the car while Dad went inside. Soon he would emerge from the house, carrying in his arms like a china doll his crippled uncle. I then would open the door and watch how tenderly and with such affection my father would place Uncle Elias in the front seat so he would have a fine view while I occupied the rear seat.

The drive was brief and the conversation limited, but oh, what a legacy of love! Father never read to me from the Bible about the good Samaritan. Rather, he took me with him and Uncle Elias in that old 1928 Oldsmobile along the road to Jericho.

When our homes carry the legacy of love, we will not receive Jacob’s chastisement as recorded in the Book of Mormon: “Ye have broken the hearts of your tender wives, and lost the confidence of your children, because of your bad examples before them; and the sobbings of their hearts ascend up to God against you.” (Jacob 2:35.)

Let us not be discouraged by the many newspaper and television accounts of discord—and sometimes cruelty—between companions and assume that virtue has vanished and love’s lamp no longer glows. Two of my dearest friends now lie in poor health and helpless. They are not alone. Their faithful companions minister to them in tender love. My friend Pres, who rarely leaves the side of his wife, said of her, “Christine is weaker but still beautiful. I love her so.” What a noble tribute to fidelity, to love, to marriage!

Another, a wife named Gertrude, makes comfortable her husband, Mark, in his room. Everything is just as he would want the room to be. She reads to him. She chats with him about the family. She once said to me during this long vigil, “I love him more than ever.”

For a beautiful example of “love at home,” we need not look beyond the family of President and Sister Benson. My wife and I were privileged to attend the Bensons’ sixty-second wedding anniversary party just three weeks ago. Children, grandchildren, and great-grandchildren rejoiced as the President and his companion held hands and led the group in singing “Keep the Home Fires Burning,” “Love’s Old Sweet Song,” and “I Love You Truly.” The entire Church can well emulate the Bensons’ example of studying the scriptures, attending the temple, and enjoying life together.

These are pictures which portray a legacy of love as a hallmark of a happy home.

A fourth hallmark of a happy home is a treasury of testimony. “The first and foremost opportunity for teaching in the Church lies in the home,” observed President David O. McKay. “A true Mormon home is one in which if Christ should chance to enter, he would be pleased to linger and to rest.” (Gospel Ideals, Salt Lake City: Improvement Era, 1953, p. 169.)

What are we doing to ensure that our homes meet this description? It isn’t enough for parents alone to have strong testimonies. Children can ride only so long on the coattails of a parent’s conviction.

President Heber J. Grant declared: “It is our duty to teach our children in their youth. I may know that the gospel is true, and so may my wife; but I want to tell you that our children will not know that the gospel is true unless they study it and gain a testimony for themselves.”

A love for the Savior, a reverence for His name, and genuine respect one for another will provide a fertile seedbed for a testimony to grow.

Learning the gospel, bearing a testimony, leading a family are rarely if ever simple processes. Life’s journey is characterized by bumps in the road, swells in the sea—even the turbulence of our times.

Some years ago, while visiting the members and missionaries in Australia, I witnessed a sublime example depicting how a treasury of testimony can bless and sanctify a home. The mission president, Horace D. Ensign, and I were traveling the long distance from Sydney to Darwin, where I was to break ground for our first chapel in that city. En route we had a scheduled stop at a mining community named Mt. Isa. As we entered the small airport at Mt. Isa, a woman and her two children approached. She said, “I am Judith Louden, a member of the Church, and these are my two children. We thought you might be on this flight, so we have come to visit with you during your brief stopover.” She explained that her husband was not a member of the Church and that she and the children were indeed the only members in the entire area. We shared lessons and bore testimony.

Time passed. As we prepared to reboard, Sister Louden looked so forlorn, so alone. She pleaded, “You can’t go yet; I have so missed the Church.” Suddenly the loudspeaker announced a thirty-minute mechanical delay of our flight. Sister Louden whispered, “My prayer has just been answered.” She then asked how she might influence her husband to show an interest in the gospel. We counseled her to include him in their home Primary lesson each week and be to him a living testimony of the gospel. I mentioned we would send to her a subscription to The Children’s Friend and additional helps for her family teaching. We urged that she never give up on her husband.

We departed Mt. Isa, a city to which I have never returned. I shall, however, always hold dear in memory that sweet mother and those precious children extending a tear-filled expression and a fond wave of gratitude and good-bye.

Several years later, while speaking at a priesthood leadership meeting in Brisbane, Australia, I emphasized the significance of gospel scholarship in the home and the importance of living the gospel and being examples of the truth. I shared with the men assembled the account of Sister Louden and the impact her faith and determination had made on me. As I concluded, I said, “I suppose I’ll never know if Sister Louden’s husband ever joined the Church, but he couldn’t have found a better model to follow.”

One of the leaders raised his hand, then stood and declared, “Brother Monson, I am Richard Louden. The woman of whom you speak is my wife. The children [his voice quavered] are our children. We are a forever family now, thanks in part to the persistence and the patience of my dear wife. She did it all.” Not a word was spoken. The silence was broken only by sniffles and muffled sobs and marked by the sight of tears streaming from every eye.

My brothers and sisters, let us determine, whatever our circumstance, to make of our houses happy homes. Let us open wide the windows of our hearts, that each family member may feel welcome and “at home.” Let us open also the doors of our very souls, that the dear Christ may enter. Remember His promise: “Behold, I stand at the door, and knock: if any man hear my voice, and open the door, I will come in to him.” (Rev. 3:20.)

How welcome He will feel, how joyful will be our lives, when the “Hallmarks of a Happy Home” greet Him, even:

That our loving Heavenly Father may bless all of us in our quest for such happy homes is my prayer, in the name of Jesus Christ, amen.

\newpage
\settalk{Go For It!}
\setspeaker{Thomas S. Monson}
\settalkdate{April 1989}
\talkheading{Go For It!}{Thomas S. Monson}

Brethren, you are an inspiring sight to behold. It is awesome to realize that in thousands of chapels throughout the world at this hour, your fellow holders of the priesthood of God are receiving this broadcast by way of satellite transmission. Your nationalities vary, and your languages are many, but a common thread binds us together. We have been entrusted to bear the priesthood and to act in the name of God. We are the recipients of a sacred trust. Much is expected of us.

Long ago, the renowned author Charles Dickens wrote of opportunities that await. In his classic volume entitled Great Expectations, Dickens described a boy by the name of Philip Pirrip, more commonly known as “Pip.” Pip was born in unusual circumstances. He was an orphan. He never met his mother or father. He never saw a picture of them. Yet he had all the normal desires of a boy. He wished with all his heart that he were a scholar. He wished that he were a gentleman. He wished that he were less ignorant. Yet all of his ambitions and all of his hopes seemed doomed to failure. Do you young men sometimes feel that way? Do those of us who are older entertain these same thoughts?

Then one day a London lawyer by the name of Jaggers approached little Pip and told him that an unknown benefactor had bequeathed to him a fortune. The lawyer put his arm around the shoulder of Pip and said to him, “My boy, you have great expectations.”

Tonight, as I look at you young men and realize who you are and what you may become, I say to you, as that lawyer said to Pip, “My boy, you have great expectations”—not as the result of an unknown benefactor, but as the result of a known Benefactor, even our Heavenly Father, and great things are expected of you.

All of us, before the period known as mortality, lived as spirit children of our Heavenly Father. In His wisdom, He has given us a record, in the book of Abraham, which tells us something of that existence:

“Now the Lord had shown unto me, Abraham, the intelligences that were organized before the world was; and among all these there were many of the noble and great ones; …

“And there stood one among them that was like unto God, and he said unto those who were with him: We will go down, for there is space there, and we will take of these materials, and we will make an earth whereon these may dwell;

“And we will prove them herewith, to see if they will do all things whatsoever the Lord their God shall command them;

“And they who keep their first estate shall be added upon; and they who keep not their first estate shall not have glory in the same kingdom with those who keep their first estate; and they who keep their second estate shall have glory added upon their heads for ever and ever.” (Abr. 3:22, 24–26.)

As we journey through mortality, let us remember from whence we came; let us be true to the trust vested in us. Let us remember who we are and what God expects us to become.

Ned Winder, a lifelong friend and formerly the executive secretary of the Missionary Department, tells of an amusing and humbling encounter which he experienced.

Two of the General Authorities, accompanied by Brother Winder, were walking down a staircase in view of a mother and her son, who were sitting on a couch facing the staircase. Seeing the brethren approach, the boy said to his mother, “Who is that first man?”

She replied, “He is Elder Marvin J. Ashton, a member of the Council of the Twelve Apostles.”

The boy continued, “Who is the man next to him?”

Mother replied, “He is Elder Loren Dunn, of the First Quorum of the Seventy.”

Then the boy concluded, “Who is the other man?”

The mother spoke more softly, yet she was still audible to Brother Winder: “Oh, he’s nobody.”

Remember, my young friends, you are somebody! You are a child of promise. You are a man of might. You are a son of God, endowed with faith, gifted with courage, and guided by prayer. Your eternal destiny is before you. The Apostle Paul speaks to you today as he spoke to Timothy long years ago: “Neglect not the gift that is in thee. … O Timothy, keep that which is committed to thy trust.” (1 Tim. 4:14; 6:20.)

As you define your goals and plan for their achievement, ponder the thought: The past is behind—learn from it; the future is ahead—prepare for it; the present is here—live in it.

At times, all of us let that enemy of achievement—even the culprit, self-defeat—dwarf our aspirations, smother our dreams, cloud our vision, and wreck our lives. The enemy’s voice whispers in our ears, “I can’t do it.” “I’m too little.” “Everyone is watching.” “I’m nobody.” This is when we need to reflect on the counsel of Maxwell Maltz, who declared:

“The most … realistic self-image of all is to conceive of yourself as ‘made in the image of God.’” You cannot sincerely hold this conviction without experiencing a profound new sense of strength and power. (Psycho-Cybernetics, Englewood Cliffs, N.J.: Prentice-Hall, 1960, p. 245.)

This is good medicine for all of us—young and old. After all, men are but boys grown older. One wife said of her husband, as he admiringly gazed at his new boat, “The bigger the boy, the bigger the toy!”

Life was never intended to consist of a glut of luxury, be an easy course, or filled only with success. There are those games which we lose, those races in which we finish last, and those promotions which never come. Such experiences provide an opportunity for us to show our determination and to rise above disappointment.

I read the other day about an athlete who is a member of LaSalle University’s wrestling team. Due to a shooting accident which occurred many years ago, he has but one leg. Does he complain? Does he curse God? Does he withdraw from the match? On the contrary, he competes with the best of them. His record this year is ten wins and eight losses. A teammate said of him, “He inspires us.”

Like some of you, I know what it is to face disappointment and youthful humiliation. As a boy, I played team softball in elementary and junior high school. Two captains were chosen, and then they, in turn, selected the players they desired on their teams. Of course, the best players were chosen first, then second and third. To be selected fourth or fifth was not too bad, but to be chosen last and relegated to a remote position in the outfield was downright awful. I know. I was there.

How I hoped that the ball would never be hit in my direction, for surely I would drop it, runners would score, and teammates would laugh.

As though it were just yesterday, I remember the moment when all that changed in my life. The game started out as I have described: I was chosen last. I made my sorrowful way to the deep pocket of right field and watched as the other team filled the bases with runners. Two batters then went down on strikes. Suddenly, the next batter hit a mighty drive. The ball was coming in my direction. Was it beyond my reach? I raced for the spot where I thought the ball would drop, uttered a silent prayer as I ran, and stretched forth my cupped hands. I surprised myself. I caught the ball! My team won the game.

This one experience bolstered my confidence, inspired my desire to practice, and led me from that last-to-be-chosen place to become a real contributor to the team.

We can experience that burst of confidence. We can feel that pride of performance. A three-word formula will help us: Never Give Up.

Opposition is ever with us. The temptation to detour from our chosen path is at times a daily confrontation. Joseph L. Townsend wrote the words of a hymn which we sing frequently:

\begin{verse}
Choose the right when a choice is placed before you.\\
In the right the Holy Spirit guides;\\
And its light is forever shining o’er you,\\
When in the right your heart confides.
\end{verse}

(“Choose the Right,” Hymns, 1985, no. 239.)

A wise father, speaking to his son, placed the question of choice in a direct setting. He counseled, “Son, if you ever find yourself in a place you shouldn’t ought to be—get out!” Good advice for a son. Good advice for a father, too.

Altogether too frequently, we are prone to place the blame on Lucifer for every temptation we encounter or every sin we commit. The words of the Apostle Paul place in perspective such thinking. To the Corinthians, Paul counseled,

“There hath no temptation taken you but such as is common to man: but God is faithful, who will not suffer you to be tempted above that ye are able; but will with the temptation also make a way to escape, that ye may be able to bear it.” (1 Cor. 10:13.)

As priesthood holders, we have a responsibility to “stand up and be counted.” Some years back, when David Kennedy was appointed as Secretary of the Treasury, a reporter attempted to entrap him with the question: “Mr. Kennedy, do you believe in prayer?”

The response was: “I do.”

Then the clever question: “Mr. Kennedy, do you pray?”

Came the firm reply: “I believe in prayer, and I pray!”

Just this past month, a mammoth 747 jetliner, while flying over the Pacific, sustained a gigantic tear in its side, ejecting nine passengers to their deaths and threatening the lives of all. When the pilot, Captain David Cronin, was interviewed, having brought the craft back safely to Honolulu, he was asked, “What did you do when the plane ripped open? How did you cope?”

Captain Cronin replied, “I prayed, then went to work.”

My brethren, this is an inspired plan for each of us to follow: Pray, and then go to work.

In the helter-skelter competitiveness of life, there is a tendency to think only of ourselves. To succumb to this philosophy narrows one’s vision and distorts a proper view of life. When concern for others replaces concern for self, our own progress is enhanced.

Tonight we have witnessed the highest honor Scouting is able to bestow, conferred upon our President, Ezra Taft Benson. This recognition is not a response to a single deed or a temporary commitment to service. Rather, it recognizes a lifetime of constant and selfless service to youth. It was said of our Lord, “He went about doing good.” President Ezra Taft Benson daily exemplifies this example of the Lord.

At the February meeting of the National Executive Board of Scouting, young men were recognized who had saved the lives of others during the past year. One of those so honored was an Aaronic Priesthood bearer—fifteen-year-old Thomas T. Nelson from Lacey, Washington. Tom had rescued two boys from a raging river which could have carried them to their deaths. I love his humble-yet-powerful response to the recognition: “I jumped in and pulled them out!”

Thousands of Scouts became heroes by blessing the lives of others during the campaign noted as “Scouting for Food.” On a given Saturday, with the campaign having been previously publicized, the homemakers of America were asked to contribute canned food to feed the hungry. Scouts became the facilitators of this objective. Hundreds of tons of food were collected, stored, and distributed. Those who gave were blessed. Those who received were fed. Those Scouts who helped to achieve the objective will never again be the same. They went about doing good.

Serving throughout the world is a great missionary force going about doing good. Missionaries teach truth. They dispel darkness. They spread joy. They bring precious souls to Christ.

Just a few weeks ago, in Guatemala City, Guatemala, I witnessed a modern miracle—even the result of God’s guidance given to His servants and the blessing of His people.

At a regional conference, almost twelve thousand members filled the Estadio del Ejercito, the local soccer stadium. The sun bathed with its rays the large gathering, while the Spirit of the Lord filled every heart. This was a day of thanksgiving, marking the forty-second anniversary of the arrival of the first missionaries to that land. John Forres O’Donnal spoke to the vast throng. He it was who, in 1946, stood alone as the only member of the Church in that nation. Personally importuning then President George Albert Smith, Brother O’Donnal facilitated the entry of the first missionaries. His wife, Carmen Galvez de O’Donnal, became the first convert and was baptized on November 13, 1948. This day of conference, as throughout the years of their marriage, she sat by her husband’s side.

While President O’Donnal spoke, my thoughts drifted back to the many missionaries who had come to this land and the hardships they endured, the sacrifices they made, and the lives they blessed. The experience of one describes the devotion of all. While I have, on a previous occasion, mentioned the experience of this missionary, following my recent visit to Guatemala I felt impressed to share it with you once again.

While serving in Guatemala as a missionary for The Church of Jesus Christ of Latter-day Saints, Randall Ellsworth survived a devastating earthquake, which hurled a beam down on his back, paralyzing his legs and severely damaging his kidneys. He was the only American injured in the quake, which claimed the lives of some eighteen thousand persons.

After receiving emergency medical treatment, Elder Ellsworth was flown to a large hospital near his home in Rockville, Maryland. While he was confined there, a newscaster conducted with him an interview that I witnessed through the miracle of television. The reporter asked, “Can you walk?”

The answer, “Not yet, but I will.”

“Do you think you will be able to complete your mission?”

Came the reply: “Others think not, but I will. With the President of my church praying for me, and through the prayers of my family, my friends, and my missionary companions, I will walk, and I will return to Guatemala. The Lord wanted me to preach the gospel there for two years, and that’s what I intend to do.”

There followed a lengthy period of therapy, punctuated by silent yet heroic courage. Little by little, the feeling began to return to the almost lifeless limbs. More therapy, more courage, more prayer.

At last Randall Ellsworth walked aboard the plane that carried him back to the mission to which he had been called, back to the people whom he loved. He left behind a trail of skeptics and a host of doubters, but also hundreds amazed at the power of God, the miracle of faith, and the reward of determination.

In Guatemala, Randall pursued his responsibilities. He walked with the use of two canes. His walk was slow and deliberate. Then one day, as he stood before his mission president, Randall Ellsworth heard him speak the almost unbelievable words, “You have been the recipient of a miracle. Your faith has been rewarded. If you have the necessary confidence, if you have abiding faith, if you have supreme courage, place those two canes on my desk—and walk.”

Slowly, Randall placed one cane and then the other on the mission president’s desk, turned toward the door and toward his future—and walked.

Today, Randall Ellsworth is a practicing physician. He is a stalwart husband and a loving father. His mission president was none other than John Forres O’Donnal—the man who helped bring to Guatemala the word of the Lord, the leader who on Sunday, March 5, 1989, addressed the throng assembled for regional conference.

Forres O’Donnal visited my office not long ago and, in his modest manner, recounted his experience with Randall Ellsworth. He then said to me, “Together we have witnessed a miracle. I have kept one of the two canes placed upon my desk that day when I challenged Elder Ellsworth to walk without them. I would like you to have the other.” With a friendly smile, he departed the office and returned home to Guatemala.

This is the cane given to me. It serves as a silent witness of our Heavenly Father’s ability to hear our prayers and to bless our lives. It is a symbol of faith. It is a reminder of courage.

Brethren of the priesthood, like the Charles Dickens character Philip Pirrip, we have great expectations. The goal of eternal life awaits. May we strive unflinchingly for its attainment. In the language of the young men assembled tonight, “Let’s go for it!” In the name of Jesus Christ, amen.

\newpage
\settalk{Thanks Be to God}
\setspeaker{Thomas S. Monson}
\settalkdate{April 1989}
\talkheading{Thanks Be to God}{Thomas S. Monson}

In the first section of the Doctrine and Covenants, we read the Lord’s promise:

“Hearken, O ye people of my church, saith the voice of him who dwells on high. … Hearken ye people from afar. …

“For verily the voice of the Lord is unto all men, and there is none to escape; and there is no eye that shall not see, neither ear that shall not hear, neither heart that shall not be penetrated. …

“And the voice of warning shall be unto all people, by the mouths of my disciples, whom I have chosen in these last days.

“And they shall go forth and none shall stay them, for I the Lord have commanded them.” (D\&C 1:1–2, 4–5.)

Exactly fifty years ago, in 1939, the heads of state in Europe solemnly returned their position papers to leather briefcases, arose from their chairs at the conference table, and returned to their respective countries. Peace had perished. Mighty armies crossed international borders. Warplanes droned overhead; giant tanks lumbered forward. World War II had begun.

Hundreds of missionaries were withdrawn from Europe and reassigned elsewhere in the world. The membership of the Church in those areas, now deprived of missionary leadership, carried on valiantly. Carnage, suffering, and death enveloped Europe.

After six terrible years, the conflict ceased and a mammoth rebuilding effort was commenced. Missionaries returned to some nations, the gospel was taught, and the Church began to grow.

In other countries, new political boundaries sprang up, borders bristled with armaments, and missionaries were denied entry. Our members there endured a period marked by patient waiting, fervent praying, and faithful living.

In October 1988, as my plane droned onward to Berlin, my thoughts were upon these nations and my heart felt concern for their people, particularly our own members who had unflinchingly borne their burdens and suffered in silence. I sat back somewhat in reverie, contemplating my lengthy assignment to the German Democratic Republic. For twenty years this had been a vital part of my ministry. My mind filled with memories. My heart overflowed with gratitude to God. I reflected on the history of the Church in the land to which I was going.

Prior to World War II, the nation we now know as the German Democratic Republic, and which some erroneously term East Germany, was the most productive area of the German-speaking world as pertained to missionary success. The city of Chemnitz, now Karl-Marx-Stadt, had as many as six large branches of members and was the greatest concentration of Latter-day Saints outside of North America. Then came the terrible destruction of World War II. After the bombs ceased and the artillery fell silent, the land was left devastated. Then, like moles from the earth came the people, bedraggled, hungry, frightened, lost. In memory one could hear the cry, “Mother, where are you? Father, where have you gone?” They were greeted by nothing but a moonscape of shell holes, jagged buildings, giant craters, and mountains of rubble. A nation lay desolate and destroyed.

About that time, the prophet of the Lord determined that one named Ezra Taft Benson would undertake a rescue mission to the struggling people. Elder Benson left his dear wife, whom he loves with all his heart, left his precious children, who were tiny at the time, and went on a mission, the length of which was uncertain. He traversed the land of German-speaking Europe—east and west. He fed the people. He clothed the people. He blessed the people. And he gave them hope. His record of service was a foundation for the progress which followed.

Another great benefactor of our German-speaking members is Walter Stover. Freely he has shared his life and generously given of his means to lift the people.

In 1968 when I made my first visit to the German Democratic Republic, tensions were high. Trust and understanding did not exist. No diplomatic relations had been established. On a cloudy and rain-filled day I journeyed to the city of Görlitz, situated deep in the German Democratic Republic near the Polish and Czech borders. I attended my first meeting with the Saints. We assembled in a small and ancient building. As the members sang the hymns of Zion, they literally filled the hall with their faith and devotion.

My heart was filled with sorrow when I realized the members had no patriarch, no wards or stakes—just branches. They could not receive temple blessings—either endowment or sealing. No official visitor had come from Church headquarters in a long time. The members could not leave their country. Yet they trusted in the Lord with all their hearts.

I stood at the pulpit, and with tear-filled eyes and a voice choked with emotion, I made a promise to the people: “If you will remain true and faithful to the commandments of God, every blessing any member of the Church enjoys in any other country will be yours.” Then I realized what I had said. That night, I dropped to my knees and pleaded with my Heavenly Father, “Father, I’m on Thy errand; this is Thy Church. I have spoken words that came not from me but from Thee and Thy Son. Wilt Thou fulfill the promise in the lives of this noble people.” Thus concluded my first visit to the German Democratic Republic.

The Lord’s promise began to unfold. A patriarch was named: Brother Percy K. Fetzer, who also was assigned as a Regional Representative for the area. Then Walter Krause, a native of that country, was ordained a patriarch. To date he has given 989 patriarchal blessings, and his wife has typed every one of them.

Time and again I paid visits to that nation. I recall leadership meetings in which the priesthood leaders eagerly ran to the front when their names were called to obtain printed instructions concerning how a quorum should operate or how a branch should function.

I remember going to a conference in the city of Annaberg. There, a sweet, older sister came forward and asked, “Are you an Apostle?”

When I answered, “Yes,” she reached in her purse and brought forth a picture of the Quorum of the Twelve Apostles. She asked, “Which one are you?”

I looked at the picture. The junior member of the Quorum of the Twelve in that picture was Elder John A. Widtsoe. She had not seen a member of the Twelve for a very long time!

Soon a member mission organization was established, the first high priest ordained, and district councils organized. In Freiberg there was created a stake of Zion and then another stake in Leipzig. Each member of the Church in the German Democratic Republic now belonged to a stake of the Church. One branch president whom I interviewed had served in this capacity for twenty-one years. He was only forty-two years of age. Half his life he had been a branch president, yet he was willing to carry on in any assignment. The members eagerly accepted their calls.

These remarkable events were preceded by a special dedication of the land.

On a Sunday morning, April 27, 1975, I stood on an outcropping of rock situated between the cities of Dresden and Meissen, high above the Elbe River, and offered a prayer on the land and its people. That prayer noted the faith of the members. It emphasized the tender feelings of many hearts filled with an overwhelming desire to obtain temple blessings. A plea for peace was expressed. Divine help was requested. I spoke the words: “Dear Father, let this be the beginning of a new day for the members of Thy Church in this land.”

Suddenly, from far below in the valley, a bell in a church steeple began to chime and the shrill crow of a rooster broke the morning silence, each heralding the commencement of a new day. Though my eyes were closed, I felt a warmth from the sun’s rays reaching my face, my hands, my arms. How could this be? An incessant rain had been falling all morning.

At the conclusion of the prayer, I gazed heavenward. I noted a ray of sunshine which streamed from an opening in the heavy clouds, a ray which engulfed the spot where our small group stood. From that moment I knew divine help was at hand.

The work moved forward. The paramount blessing needed was the privilege of our worthy members to receive their endowments and their sealings.

We explored every possibility. A trip once in a lifetime to the temple in Switzerland? Not approved by the government. Perhaps mother and father could come to Switzerland, leaving the children behind. Not right. How do you seal children to parents when they cannot kneel at an altar? It was a tragic situation. Then, through the fasting and the prayers of many members, and in a most natural manner, government leaders proposed: Rather than having your people go to Switzerland to visit a temple, why don’t you build a temple here in the German Democratic Republic? The proposal was accepted, a choice parcel of property obtained in Freiberg, and ground broken for a beautiful temple of God.

The day of dedication was an historic occasion. President Gordon B. Hinckley offered the dedicatory prayer. Heaven was close that day.

For its size, this temple is one of the busiest temples in the Church. It is the only temple where one makes an appointment to participate in an endowment session. It is the only temple I know of where stake presidents say, “What can we do? Our home teaching is somewhat down because everyone is in the temple!” When I heard that comment, I thought, “Not bad—not bad at all!”

A miracle of miracles had taken place. One more was needed. How can the Church grow without missionaries? How can our numbers increase despite an aging population? Beautiful new buildings grace the land: stake centers at Leipzig and Dresden, and chapels in Freiberg and Zwickau, with others to follow, such as a chapel under construction in the city of Plauen. A faithful brother from Plauen wrote me this poignant letter: “My parents and grandparents have served before us in this branch, but never thus far has it been possible to have our own meetinghouse. Now a long-cherished wish is being fulfilled.” After reading this touching account, the thought crossed my mind, “But what use are buildings if there are not sufficient members to occupy them?”

Such was the dilemma uppermost on my mind as my plane landed in Berlin that October afternoon. We went forward with the vital assignment to visit with the leaders of the German Democratic Republic. Our ultimate goal was to seek permission for the doorway of missionary work to open. Elder Russell M. Nelson, Elder Hans B. Ringger, and I, along with our local German Democratic Republic Church leaders, headed by President Henry Burkhardt, President Frank Apel, and President Manfred Schutze, initially met with State Secretary for Religious Affairs Kurt Löffler as he hosted a lovely luncheon in our honor. He addressed our group by saying, “We want to be helpful to you. We’ve observed you and your people for twenty years. We know you are what you profess to be: honest men and women.”

Government leaders and their wives attended the dedication of a stake center at Dresden and a chapel at Zwickau. As the Saints sang “God be with you till we meet again”—“Auf Wiedersehen, Auf Wiedersehen”—we remembered Him, the Prince of Peace, who died on the cross at Calvary. I contemplated our Lord and Savior, when He walked the path of pain, the trail of tears, even the road of righteousness. His penetrating declaration came to mind: “Peace I leave with you, my peace I give unto you: not as the world giveth, give I unto you. Let not your heart be troubled, neither let it be afraid.” (John 14:27.)

Then it was back to Berlin for the crucial meetings with the head of the nation, even Chairman Erich Honecker.

That special morning the sunlight bathed the city of Berlin. It had been raining all night, but now beauty prevailed. We were driven to the chambers of the chief representatives of the government.

Beyond the exquisite entry to the building, we were greeted by Chairman Honecker. We presented to him the statuette First Step, depicting a mother helping her child take its first step toward its father. He was highly pleased with the gift. He then escorted us into his private council room. There, around a large round table, we were seated. Others at the table included Chairman Honecker and his deputies of government.

Chairman Honecker began, “We know members of your Church believe in work; you’ve proven that. We know you believe in the family; you’ve demonstrated that. We know you are good citizens in whatever country you claim as home; we have observed that. The floor is yours. Make your desires known.”

I began, “Chairman Honecker, at the dedication and open house for the temple in Freiberg, 89,890 of your countrymen stood in line, at times up to four hours, frequently in the rain, that they might see a house of God. In the city of Leipzig, at the dedication of the stake center, 12,000 people attended the open house. In the city of Dresden there were 29,000 visitors; in the city of Zwickau, 5,300. And every week of the year 1,500 to 1,800 people visit the temple grounds in the city of Freiberg. They want to know what we believe. We would like to tell them that we believe in honoring and obeying and sustaining the law of the land. We would like to explain our desire to achieve strong family units. These are but two of our beliefs. We cannot answer questions, and we cannot convey our feelings, because we have no missionary representatives here as we do in other countries. The young men and young women whom we would like to have come to your country as missionary representatives would love your nation and your people. More particularly, they would leave an influence with your people which would be ennobling. Then we would like to see young men and young women from your nation who are members of our Church serve as missionary representatives in many nations, such as in America, in Canada, and in a host of others. They will return better prepared to assume positions of responsibility in your land.”

Chairman Honecker then spoke for perhaps thirty minutes, describing his objectives and viewpoints and detailing the progress made by his nation. At length, he smiled and addressed me and the group, saying, “We know you. We trust you. We have had experience with you. Your missionary request is approved.”

My spirit literally soared out of the room. The meeting was concluded. As we left the beautiful government chambers, Elder Russell Nelson turned to me and said, “Notice how the sunshine is penetrating this hall. It’s almost as though our Heavenly Father is saying, ‘I am pleased.’”

The black darkness of night had ended. The bright light of day had dawned. The gospel of Jesus Christ would now be carried to the millions of people in that nation. Their questions concerning the Church will be answered, and the Kingdom of God will go forth.

As I reflect on these events, my thoughts turn to the Master’s words, “In nothing doth man offend God, or against none is his wrath kindled, save those who confess not his hand in all things.” (D\&C 59:21.) I confess the hand of God in the miraculous events pertaining to the Church in the German Democratic Republic.

The faith and devotion of our members in that nation have not gone unnoticed by God. The excellent service of other General Authorities, Regional Representatives, and mission presidents has been of inestimable help. The understanding cooperation of government leaders is most appreciated. Assignments have been made to the first ten missionaries from the German Democratic Republic to serve abroad; and just three days ago, on Thursday, March 30, the first full-time missionary representatives in exactly fifty years entered the German Democratic Republic. Their mission president was there to greet them. The long period of preparation is past. The future of the Church unfolds. Thanks be to God.

From the heavens we hear anew the Lord’s declaration:

“Hear, O ye heavens, and give ear, O earth, and rejoice ye inhabitants thereof, for the Lord is God, and beside him there is no Savior.

“Great is his wisdom, marvelous are his ways, and the extent of his doings none can find out.

“His purposes fail not, neither are there any who can stay his hand. …

“For thus saith the Lord—I, the Lord, am merciful and gracious unto those who fear me, and delight to honor those who serve me in righteousness and in truth unto the end.

“Great shall be their reward and eternal shall be their glory.” (D\&C 76:1–3, 5–6.)

May this be our universal blessing, I pray in the name of Jesus Christ, amen.

\newpage
\settalk{The Service That Counts}
\setspeaker{Thomas S. Monson}
\settalkdate{October 1989}
\talkheading{The Service That Counts}{Thomas S. Monson}

While driving to the office one morning, I passed a dry-cleaning establishment which had a sign by the side of the front door. It read, “It’s the Service That Counts.” I suppose in a highly competitive field such as the dry-cleaning business and many others, the factor which distinguishes one store from another is, in actual fact, service.

The message from the small sign simply would not leave my mind. Suddenly I realized why. In actual fact it is the service that counts—the Lord’s service.

All of us admire and respect that noble king of Book of Mormon fame—even King Benjamin. How respected he must have been for the people to gather in such great numbers to hear his words and receive his counsel. I think it most interesting that the multitude “pitched their tents round about the temple, every man having his tent with the door thereof towards the temple, that thereby they might remain in their tents and hear the words which king Benjamin should speak unto them.” (Mosiah 2:6.) Even a high tower had to be erected that the people might hear his words.

In the true humility of an inspired leader, King Benjamin recounted his desire to serve his people and lead them in paths of righteousness. He then declared to them:

“Because I said unto you that I had spent my days in your service, I do not desire to boast, for I have only been in the service of God.

“And behold, I tell you these things that ye may learn wisdom; that ye may learn that when ye are in the service of your fellow beings ye are only in the service of your God.” (Mosiah 2:16–17.)

This is the service that counts, brethren—the service to which all of us have been called, the service of the Lord Jesus Christ.

As He enlists us to His cause, He invites us to draw close to Him. He speaks to you and to me:

“Come unto me, all ye that labour and are heavy laden, and I will give you rest.

“Take my yoke upon you, and learn of me; for I am meek and lowly in heart: and ye shall find rest unto your souls.

“For my yoke is easy, and my burden is light.” (Matt. 11:28–30.)

To all who go forth in His service, He provides this assurance: “I will go before your face. I will be on your right hand and on your left, and my Spirit shall be in your hearts, and mine angels round about you, to bear you up.” (D\&C 84:88.)

Many assembled tonight have responsibility to provide leadership to those holding the Aaronic Priesthood. To you I say, The finest teaching you can provide is that of a good example. Youth need fewer critics and more models to follow. All of us who are engaged in the Lord’s work have the responsibility to reach out to those who are less active and bring them to the service of the Lord. Their souls are ever so precious.

In a revelation to Joseph Smith the Prophet, Oliver Cowdery, and David Whitmer, the Lord taught:

“Remember the worth of souls is great in the sight of God;

“For, behold, the Lord your Redeemer suffered death in the flesh; wherefore he suffered the pain of all men, that all men might repent and come unto him. …

“And how great is his joy in the soul that repenteth!

“Wherefore, you are called to cry repentance unto this people.

“And if it so be that you should labor all your days in crying repentance unto this people, and bring, save it be one soul unto me, how great shall be your joy with him in the kingdom of my Father!

“And now, if your joy will be great with one soul that you have brought unto me into the kingdom of my Father, how great will be your joy if you should bring many souls unto me!” (D\&C 18:10–11, 13–16.)

Some years ago while I was attending a priesthood leadership session of the Monument Park West Stake conference, this scripture became the theme for the visitor from the Welfare Committee, my former stake president, Paul C. Child. In his accustomed style, Brother Child left the stand and began to walk down the aisle among the assembled priesthood brethren. He quoted the verse, “Remember the worth of souls is great in the sight of God.” (D\&C 18:10.) Then he asked the question, “Who can tell me the worth of a human soul?”

Every man in attendance began to think of an answer in the event Brother Child were to call on him. I had grown up under his leadership, and I knew he would never call on a high councilor or member of a bishopric; rather, he would select one who would least expect to be called. Sure enough, he called from a list he carried the name of an elders quorum president. Thunderstruck, the brother stammered as he asked, “Would you repeat the question, please?” The question was repeated, followed by an even longer pause. Suddenly the response came forth, “The worth of a human soul is its capacity to become as God.”

Brother Child closed his scripture, walked back to the pulpit, and, while passing me whispered, “A profound reply; a profound reply.”

With this perspective firmly in our minds, we are prepared to serve in the great mission of bringing souls unto Him.

Many of you hold the Aaronic Priesthood. You are preparing to serve as missionaries. Begin now to learn in your youth the joy of service in the cause of the Master.

Following Thanksgiving time a year or so ago, I received a letter from a widow whom I had known in the stake where I served in the presidency. She had just returned from a dinner sponsored by her bishopric. Her words reflect the peace she felt and the gratitude which filled her heart:

“Dear President Monson,

“I am living in Bountiful now. I miss the people of our old stake, but let me tell you of a wonderful experience I have had. In early November all the widows and older people received an invitation to come to a lovely dinner. We were told not to worry about transportation since this would be provided by the older youth in the ward.

“At the appointed hour, a very nice young man rang the bell and took me and another sister to the stake center. He stopped the car, and two other young men walked with us to the chapel where the young ladies took us to where we removed our wraps—then into the cultural hall, where we sat and visited for a few minutes. Then they took us to the tables, where we were seated on each side by either a young woman or a young man. Then we were served a lovely Thanksgiving dinner and afterward provided a choice program.

“After the program we were given our dessert—either apple or pumpkin pie. Then we left, and on the way out we were given a plastic bag with sliced turkey and two rolls. Then the young men took us home. It was such a nice, lovely evening. Most of us shed a tear or two for the love and respect we were shown.

“President Monson, when you see young people treat others like these young people did, I feel the Church is in good hands.”

I reflected on my association with this lovely widow, now grown old but ever serving the Lord. There came to mind the words from the Epistle of James: “Pure religion and undefiled before God and the Father is this, To visit the fatherless and widows in their affliction, and to keep himself unspotted from the world.” (James 1:27.)

I add my own commendation: God bless the leaders, the young men, and the young women who so unselfishly brought such joy to the lonely and such peace to their souls. Through their experience they learned the meaning of service and felt the nearness of the Lord.

One of the great missionaries of pioneer times was Joseph Millett, who served a mission to the Maritime Provinces of Canada when but eighteen years of age. His mission was marked by discouragement, yet punctuated by faith-promoting experiences—even miraculous intervention by the Lord. This lifelong servant of the Lord, who learned on his mission, and never forgot, what it is like to be in need and how to give, leaves us with this final picture of himself, taken from his personal journal and using his own words:

“One of my children came in, said that Brother Newton Hall’s folks were out of bread. Had none that day. I put … our flour in sack to send up to Brother Hall’s. Just then Brother Hall came in. Says I, ‘Brother Hall, how are you out for flour.’

“‘Brother Millett, we have none.’

“‘Well, Brother Hall, there is some in that sack. I have divided and was going to send it to you. Your children told mine that you were out.’

“Brother Hall began to cry. Said he had tried others. Could not get any. Went to the cedars and prayed to the Lord and the Lord told him to go to Joseph Millett.

“‘Well, Brother Hall, you needn’t bring this back if the Lord sent you for it. You don’t owe me for it.’”

His journal continued, “You can’t tell how good it made me feel to know that the Lord knew that there was such a person as Joseph Millett.” (In Eugene England, “Without Purse or Scrip: A 19-year-old Missionary in 1853,” New Era, July 1975, p. 28.)

Brethren, the Lord knows each of us. Do you think for a moment that He who notes the sparrow’s fall would not be mindful of our needs and our service? We simply cannot afford to attribute to the Son of God the same frailties which we find in ourselves.

A while back, my good friend G. Marion Hinckley from Utah County, my fellow trail rider, came to the office with two grandsons who were brothers, one having served an honorable mission in Japan and the other in Scotland. Brother Hinckley said, “Let me share with you a wonderful experience which came to these grandsons of mine.” His buttons were almost bursting with pride.

In faraway Japan, a commercial street photographer stopped one of the brothers, having taken a picture of him holding a small child. He offered the print for sale to the missionary and his companion. They explained that they were on a tight budget, that they were missionaries, and they directed the photographer’s attention to their nameplates. They didn’t purchase the picture.

Some months later, the brother serving in Scotland was asking two missionaries why they had arrived late for a zone meeting, when they told this story: A most persistent street photographer had attempted to sell them a picture of a missionary in Japan holding a small child. They had no interest in the picture, but to avoid arriving even later at their zone meeting, they purchased it.

“A likely story,” responded Elder Lamb, whereupon they handed him the picture. He could not believe his eyes. It was a photograph of his own brother in faraway Japan.

That day in my office they presented to my view the two pictures, and with their grandfather beaming his approval they declared, “The Lord surely is mindful of his servants the missionaries.”

As they departed my office, I thought, Yes, the Lord is mindful of his missionaries—and their fathers, their mothers, their grandparents, and all who sacrifice for their support, that precious souls may be taught and provided His gospel.

Now, many are not on the front line of missionary service in the Church callings they fill. Does God remember them also? Is He mindful of their needs and the yearnings of their hearts? What about those who have been in the limelight but have grown old with faithful service, have been released and have slipped into the anonymity of the vast congregation of Church members? To all such individuals I testify that He does remember and He does bless.

Many years ago I was assigned to divide the Modesto California Stake. The Saturday meetings had been held, the new stake presidencies selected, and preparations concluded for the announcements to be made the following morning in the Sunday session of conference.

As the Sunday session was about to begin, there went through my mind the thought that I had been in Modesto before. But when? I let my mind search back through the years for a confirmation of the thought I was thinking. Suddenly I remembered. Modesto, years before, had been a part of the San Joaquin Stake. The stake president was Clifton Rooker. I had stayed in his home during that conference. But that was many years earlier. Could my thoughts be playing tricks on my mind? I said to the stake presidency as they sat on the stand, “Is this the same stake over which Clifton Rooker once presided?”

The brethren answered, “Yes, it is. He was our former president.”

“It’s been many years since I was last here,” I said. “Is Brother Rooker with us today?”

They responded, “Oh, yes. We saw him early this morning as he came to conference.”

I asked, “Where is he seated on this day when the stake will be divided?”

“We don’t know exactly,” they replied. The response was a good one, for the building was filled to capacity.

I stepped to the pulpit and asked, “Is Clifton Rooker in the audience?” There he was—way back in the recreation hall, hardly in view of the pulpit. I felt the inspiration to say to him publicly, “Brother Rooker, we have a place for you on the stand. Would you please come forward?”

With every eye watching him, Clifton Rooker made that long walk from the rear of the building right up to the front and sat by my side. It became my opportunity to call upon him, one of the pioneers of that stake, to bear his testimony and to tell the people whom he loved that he was the actual beneficiary of the service he had rendered his Heavenly Father and which he had provided the stake members.

After the session was concluded, I said, “Brother Rooker, how would you like to step with me into the high council room and help me set apart the two new presidencies of these stakes?”

He replied, “That would be a highlight for me.”

We proceeded to the high council room. There, with his hands joining my hands and the hands of the outgoing stake presidency, we set apart to their callings the two new stake presidencies. Brother Rooker and I embraced as he said good-bye and went to his home.

Early the next morning, after I had returned to my home, I had a telephone call from the son of Clifton Rooker. “Brother Monson,” he said, “I’d like to tell you about my dad. He passed away this morning; but before he did so, he said that yesterday was the happiest day of his entire life.”

As I heard that message from Brother Rooker’s son, I paused to thank God for the inspiration which came to me to invite this good man, while he was yet alive and able to enjoy them, to come forward and receive the plaudits of the stake members whom he had served.

To all those who serve the Lord by serving their fellowmen, and to those who are the recipients of this selfless service, the Redeemer seems to be speaking to you when He declared:

“When the Son of man shall come in his glory, and all the holy angels with him, then shall he sit upon the throne of his glory:

“And before him shall be gathered all nations: and he shall separate them one from another, as a shepherd divideth his sheep from the goats:

“And he shall set the sheep on his right hand, but the goats on the left.

“Then shall the King say unto them on his right hand, Come, ye blessed of my Father, inherit the kingdom prepared for you from the foundation of the world:

“For I was an hungred, and ye gave me meat: I was thirsty, and ye gave me drink: I was a stranger, and ye took me in:

“Naked, and ye clothed me: I was sick, and ye visited me: I was in prison, and ye came unto me.

“Then shall the righteous answer him, saying, Lord, when saw we thee an hungred, and fed thee? or thirsty, and gave thee drink?

“When saw we thee a stranger, and took thee in? or naked, and clothed thee?

“Or when saw we thee sick, or in prison, and came unto thee?

“And the King shall answer and say unto them, Verily I say unto you, Inasmuch as ye have done it unto one of the least of these my brethren, ye have done it unto me.” (Matt. 25:31–40.)

That each of us may qualify for this blessing from our Lord is my prayer, in the name of Jesus Christ, amen.

\newpage
\settalk{Windows}
\setspeaker{Thomas S. Monson}
\settalkdate{October 1989}
\talkheading{Windows}{Thomas S. Monson}

While waiting my turn at an airline office in London, England, I reached forward from my chair and selected an advertising brochure from the small table which contained reading material. The publication bore the title Windows to the World. Each page contained a framed picture of a well-known and beautiful site, accompanied by a well-written description which made one desire to visit all of the locations shown. The Matterhorn in Switzerland, the Alps of New Zealand, even the Taj Mahal of India—all seemed to suggest to the reader the desirability of an immediate visit.

Windows are wonderful. They serve as a frame on which we might focus our attention. They provide a glimpse of God’s creations. The azure blue sky, the billowy, white clouds, the verdant green forest all are as framed pictures in the memory of the mind. Windows also reveal the approach of a friend, a gathering storm, a magnificent sunset—even the passing parade of life.

Windows welcome light to our lives and bring joy to our souls. The absence of windows, such as in dark prison cells, shuts out the world. When we are deprived of light, the depression of darkness encompasses us.

Windows teach lessons never to be forgotten. Ever shall I remember a visit to the home of President Hugh B. Brown. It was graduation day at Brigham Young University. He was to conduct the exercises, and I was to deliver the commencement address. I drove to President Brown’s home and escorted him to my car. Before we could drive away, however, he said to me, “Wait just a few minutes. My wife, Zina, will come to the front window.”

I glanced at the window, noted that the curtain had parted, and saw Zina Brown sitting in her wheelchair, affectionately waving a small, white handkerchief toward the gaze of her smiling husband. President Brown reached into his jacket pocket, retrieved a white handkerchief, and began to wave it gently, much to the delight of his wife. We then inched away from the curb and commenced the journey to Provo.

“What is the significance of the white-handkerchief waving?” I asked.

He replied, “Zina and I have followed that custom since we were first married. It is somewhat a symbol between us that all will be well throughout the day until we are again together at eventide.”

That day, I witnessed a window to the heart.

Some windows are sealed shut by sorrow, by pain, by neglect. The forgotten birthday, the unremembered visit, the overlooked promise—all can sow seeds of sorrow and bring to the human heart that unwelcome visitor, despair.

A national columnist one day titled her story, “What a Forgotten Birthday Can Mean,” and then quoted from a letter she had received:

“I have never written to you before, but I believe the following might interest you and your readers. I found it in an old magazine. No author’s name was mentioned—just ‘A Heavy-Hearted Observer.’

“‘Yesterday was a man’s birthday. He was ninety-one. He awakened earlier than usual, bathed, shaved, and put on his best clothes. Surely they would come today, he thought.

“‘He didn’t take his daily walk to the gas station to visit with the old-timers of the community because he wanted to be right there when they came.

“‘He sat on the front porch with a clear view of the road so he could see them coming. Surely they would come today.

“‘He decided to skip his noon nap because he wanted to be up when they came. He had six children. Two of his daughters and their married children lived within four miles. They hadn’t been to see him for such a long time. But today was his birthday. Surely they would come today.

“‘At supper time he refused to cut the cake and asked that the ice cream be left in the freezer. He wanted to wait and have dessert with them when they came.

“‘About 9 o’clock he went to his room and got ready for bed. His last words before turning out the lights were, “Promise to wake me up when they come.”

“‘It was his birthday, and he was ninety-one.’”

When I read that touching account, tears came easily. I reflected on an experience in my life, one that had a happier ending.

Each time I would visit an older widow whom I had known for many years and whose bishop I had been, my heart grieved at her utter loneliness. A favorite son of hers lived many miles away, and for years he had not visited Mother. Mattie spent long hours in a lonely vigil at her front window. Behind a frayed and frequently opened curtain, the disappointed mother would say to herself, “Dick will come; Dick will come.”

But Dick didn’t come. The years passed by one after another. Then, like a ray of sunshine, Church activity came into the life of Dick. He journeyed to Salt Lake to visit with me. He telephoned upon his arrival and, with excitement, reported the change in his life. He asked if I had time to see him if he were to come directly to my office. My response was one of gladness. However, I said, “Dick, visit your mother first, and then come to see me.” He gladly complied with my request.

Before he could get to my office, there came a phone call from Mattie, his mother. From a joyful heart came words punctuated by tears: “Tom, I knew Dick would come. I told you he would. I saw him through the window.”

Years later at Mattie’s funeral, Dick and I spoke tenderly of that experience. We had witnessed a glimpse of God’s healing power through the window of a mother’s faith in her son.

The holy scriptures are replete with sacred accounts of our Master’s love for the downtrodden and the poor of this world. Though many are forgotten by men, they are remembered by God and are ofttimes seen through the window of personal example.

Who among us can forget the timeless lesson taught by the Lord when, “in the audience of all the people he said unto his disciples,

“Beware of the scribes, which desire to walk in long robes, and love greetings in the markets, and the highest seats in the synagogues, and the chief rooms at feasts;

“Which devour widows’ houses, and for a shew make long prayers.” (Luke 20:45–47.)

“And he looked up, and saw the rich men casting their gifts into the treasury.

“And he saw also a certain poor widow casting in thither two mites.

“And he said, Of a truth I say unto you, that this poor widow hath cast in more than they all:

“For all these have of their abundance cast in unto the offerings of God: but she of her penury hath cast in all the living that she had.” (Luke 21:1–4.) What a beautiful lesson, as taught through the window of example.

At a city called Nain, the Lord opened to his disciples and to many people who followed him a window through which they might view true compassion:

“Now when he came nigh to the gate of the city, behold, there was a dead man carried out, the only son of his mother, and she was a widow: and much people of the city was with her.

“And when the Lord saw her, he had compassion on her, and said unto her, Weep not.

“And he came and touched the bier: and they that bare him stood still. And he said, Young man, I say unto thee, Arise.

“And he that was dead sat up, and began to speak. And he delivered him to his mother.” (Luke 7:12–15.)

The disciples of the Lord witnessed through the windows Jesus opened the power of God and were made partakers of this same power when, in righteousness, they ministered to the children of the Almighty.

A beautiful account, recorded in the book of Acts, tells of a disciple named Tabitha who lived at Joppa. She was described as being a woman “full of good works and almsdeeds.”

“It came to pass in those days, that she was sick, and died: whom when they had washed, they laid her in an upper chamber.

“And forasmuch as Lydda was nigh to Joppa, and the disciples had heard that Peter was there, they sent unto him two men, desiring him that he would not delay to come to them.

“Then Peter arose and went with them. When he was come, they brought him into the upper chamber: and all the widows stood by him weeping, and shewing the coats and garments which [Tabitha] made, while she was with them.” (Could we not say this was a window through which Peter glimpsed the industry of Tabitha’s life?) “Peter put them all forth, and kneeled down, and prayed; and turning him to the body said, Tabitha, arise. And she opened her eyes: and when she saw Peter, she sat up.

“And he gave her his hand, and lifted her up, and when he had called the saints and widows, presented her alive.

“And it was known throughout all Joppa; and many believed in the Lord.” (Acts 9:36–42.)

Would it not be ever so sad if such a window to priesthood power, to faith, to healing, were to be restricted to Joppa alone? Are these sacred and moving accounts recorded only for our uplift and enlightenment? Can we not apply such mighty lessons to our daily lives?

When we catch the vision regarding the worth of human souls, when we realize the truth of the adage, “God’s sweetest blessings always flow through hands that serve Him here below,” then we have quickened within our souls the desire to do good, the willingness to serve, and the yearning to lift to a higher plane the children of God.

Such was the experience of William Norris, formerly the chairman of a large computer manufacturing firm and a friend of many years. Mr. Norris determined to build a plant in an area of extreme poverty. The neighborhood was predominantly composed of a minority race—unmarried women with children, uneducated, uncared-for, but needing help. These women became the work force in the production of high-tech computers.

I had the privilege to be hosted by Mr. Norris and to be given a tour of his new facility. I was impressed with the employment provided—but more impressed with the company nursery, which occupied a wing of the building. Here, while their mothers worked, children received schooling, including proficiency with computers. Since most of the children did not have fathers and grandfathers who cared, retired grandfathers in the community were invited to have lunch with them. The children were benefited, and the grandfathers had a special blessing brought into their lives.

As a result of Mr. Norris’s dream, the chain of poverty was broken. Children learned to earn. It was as though William Norris had personally blessed the life of each worker. Through the window provided by Mr. Norris—even love in action—I saw demonstrated the philosophical and practical truth: The bottom line of living is giving.

As we go about our daily lives, we discover countless opportunities to follow the example of the Savior. When our hearts are in tune with His teachings, we discover the unmistakable nearness of His divine help. It is almost as though we are on the Lord’s errand; and we then discover that, when we are on the Lord’s errand, we are entitled to the Lord’s help.

Through the years, the offices I have occupied have been decorated with lovely paintings of peaceful and pastoral scenes. However, there is one picture that always hangs on the wall which I face when seated behind my desk. It is a constant reminder of Him whom I serve, for it is a picture of our Lord and Savior, Jesus Christ. When confronted with a vexing problem or difficult decision, I always gaze at that picture of the Master and silently ask myself the question, “What would He have me do?” No longer does doubt linger, nor does indecision prevail. The way to go is clear, and the pathway before me beckons.

Some months back I sat in my office chair reading the daily mail. I opened a letter from Martha Sharp of Wellsville, Utah, and read her entreaty seeking a blessing for her grown son, Steven, who was a patient at University Hospital in Salt Lake City. She described Steven’s spiritual and physical needs and the likelihood that he would suffer the amputation of his foot. Her tears were felt in each word, and her feelings of love marked every sentence. Hers was a request which the Spirit simply did not allow me to delegate.

When I entered Steven’s hospital room that night, I saw a man who just seemed built to ride a horse. Sensing this, I began to chat with him about a Western adventure film I had seen recently. I described the beautiful horses ridden by the principal characters. A warm smile came over Steven’s face. Not until that moment did I note on his nightstand a book he had been reading. It was the book from which the film we had been discussing was made. Our conversation was warm and free from that point forward.

In describing his condition, Steven commented, “I hope they leave enough of my foot so that I can get it into a stirrup.” I assured him we would remember his name when the First Presidency and Council of the Twelve met in the holy temple and that my wife and I would personally remember him in our prayers. I told him that he had a wonderful mother, who loved him and remembered him in his need, and a Heavenly Father who also loved and remembered him. Steven began to weep. A special spirit filled the room. A blessing was given, a heart cleansed, a memory of home and family rekindled, and a mother comforted.

As I departed the hospital, situated high on the east bench of Salt Lake City, I gazed at the panoramic view of the valley before me. The miles collapsed; the stars drew near. I could almost see through the window of mortality the expanse of eternity. One star shone especially bright. It seemed to light the way and mark the path to Wellsville. I remembered the poem from Primary days:

\begin{verse}
Star light, star bright,\\
The first star I see tonight,\\
I wish I may, I wish I might,\\
Have the wish I wish tonight.
\end{verse}

What was my wish? That Martha Sharp might receive the welcome message, “Your son loves you.”

From sacred soil far away, and from a timeless truth taught long ago, came the message, “With God all things are possible.” (Matt. 19:26.)

Once more a gentle but unseen hand had opened a window to the soul, that precious lives might receive blessings heaven-sent.

He beckons to each of us and extends the warm invitation not only to gaze at the beauty seen through the windows He opens, but also to pass through them to the priceless opportunities He provides to bless the lives of others.

That each may experience this privilege is my humble prayer, in the name of Jesus Christ, amen.

\newpage
\settalk{“A Little Child Shall Lead Them”}
\setspeaker{Thomas S. Monson}
\settalkdate{April 1990}
\talkheading{“A Little Child Shall Lead Them”}{Thomas S. Monson}

During the Galilean ministry of our Lord and Savior, the disciples came unto Him, saying: “Who is the greatest in the kingdom of heaven?

“And Jesus called a little child unto him, and set him in the midst of them,

“And said, Verily I say unto you, Except ye be converted, and become as little children, ye shall not enter into the kingdom of heaven.

“Whosoever therefore shall humble himself as this little child, the same is greatest in the kingdom of heaven.

“And whoso shall receive one such little child in my name receiveth me.

“But whoso shall offend one of these little ones which believe in me, it were better for him that a millstone were hanged about his neck, and that he were drowned in the depth of the sea” (Matt. 18:1–6).

Recently, as I read the daily newspaper, my thoughts turned to this passage and the firm candor of the Savior’s declaration. In one column of the newspaper I read of a custody battle between the mother and father of a child. Accusations were made, threats hurled, and anger displayed as parents moved here and there on the international scene with the child spirited away from one continent to another.

A second story told of a twelve-year-old lad who was beaten and set on fire because he refused a neighborhood bully’s order to take drugs. Hospitalized, his condition remains critical.

Still a third report told of a father’s sexual molestation of his small child.

These are reported cases of child abuse. There are many more never reported but equally as serious. A physician revealed to me the large number of children who are brought to the emergency rooms of local hospitals in your city and mine. In many cases guilty parents provide fanciful accounts of the child falling from his high chair or stumbling over a toy and striking his head. Altogether too frequently it is discovered that the parent was the abuser and the innocent child the victim. Shame on the perpetrators of such vile deeds. God will hold such strictly accountable for their actions.

President Ezra Taft Benson is one who exemplifies a true love for these little ones. To see the tiny tots gather at his side, extend a small hand to be held in his or to kiss his cheek, is to see the love adults should have for these children. No one in the presence of President Benson refers to a child as a “kid.” His correction for such a remark is sure and to the point. A visiting ambassador from another nation errantly made this slip. He was corrected with love.

When we realize just how precious children are, we will not find it difficult to follow the pattern of the Master in our association with them. Not long ago, a sweet scene took place at the Salt Lake Temple. Children, who had been ever so tenderly cared for by faithful workers in the temple nursery, were now leaving in the arms of their mothers and fathers. One child turned to the lovely women who had been so kind to the children and, with a wave of her arm, spoke the feelings of her heart as she exclaimed, “Goodnight, angels.”

The poet described a child so recently with its Heavenly Father as “a sweet new blossom of humanity, fresh fallen from God’s own home to flower on earth.”

Who among us has not praised God and marveled at His powers when an infant is held in one’s arms? That tiny hand, so small yet so perfect, instantly becomes the topic of conversation. No one can resist placing his little finger in the clutching hand of an infant. A smile comes to the lips, a certain glow to the eyes, and one appreciates the tender feelings which prompted the poet to pen the lines:

\begin{verse}
Our birth is but a sleep and a forgetting:\\
The soul that rises with us, our life’s star,\\
Hath had elsewhere its setting,\\
And cometh from afar;\\
Not in entire forgetfulness,\\
And not in utter nakedness,\\
But trailing clouds of glory do we come\\
From God, who is our home.
\end{verse}

[William Wordsworth, “Ode: Intimations of Immortality from Recollections of Early Childhood”]

When the disciples of Jesus attempted to restrain the children from approaching the Lord, He declared:

“Suffer the little children to come unto me, and forbid them not: for of such is the kingdom of God.

“Verily I say unto you, Whosoever shall not receive the kingdom of God as a little child, he shall not enter therein.

“And he took them up in his arms, put his hands upon them, and blessed them” (Mark 10:14–16).

What a magnificent pattern for us to follow.

My heart burned warmly within me when the First Presidency approved the allocation of a substantial sum from your special fast-offering contributions to join with those funds from Rotary International, that polio vaccine might be provided and the children living in Kenya immunized against this vicious crippler and killer of children.

I thank God for the work of our doctors who leave for a time their own private practices and journey to distant lands to minister to children. Cleft palates and other deformities which would leave a child impaired physically and damaged psychologically are skillfully repaired. Despair yields to hope. Gratitude replaces grief. These children can now look in the mirror and marvel at a miracle in their own lives.

In a recent meeting, I told of a dentist in my ward who each year visits the Philippine Islands to work his skills without compensation to provide corrective dentistry for children. Smiles are restored, spirits lifted, and futures enhanced. I did not know the daughter of this dentist was in the congregation to which I was speaking. At the conclusion of my remarks, she came forward and, with a broad smile of proper pride, said, “You have been speaking of my father. How I love him and what he is doing for children!”

In the faraway islands of the Pacific, hundreds who were near-blind now see because a missionary said to his physician brother-in-law, “Leave your wealthy clientele and the comforts of your palatial home and come to these special children of God who need your skills and need them now.” The ophthalmologist responded without a backward glance. Today he comments quietly that this visit was the best service he ever rendered and the peace which came to his heart the greatest blessing of his life.

Tears come easily to me when I read of a father who has donated one of his own kidneys in the hope that his son might have a more abundant life. I drop to my knees at night and add my prayer of faith in behalf of a mother in our community who journeyed to Chicago that she might provide part of her liver to her daughter in a delicate and potentially life-threatening surgery. She, who already had gone down into the valley of the shadow of death to bring forth this child into mortality, again put her hand in the hand of God and placed her own life in jeopardy for her child. Never a complaint, but ever a willing heart and a prayer of faith.

Elder Russell M. Nelson, upon returning from Romania, shared with us the pitiable plight of orphan children in that land—perhaps thirty thousand in the city of Bucharest alone. He visited one such orphanage and arranged that the Church might provide vaccine, medical dressings, and other urgently needed supplies. Certain couples will be identified and called to fill special missions to these children. I can think of no more Christlike service than to hold a motherless child in one’s arms or to take a fatherless boy by the hand.

We need not be called to missionary service, however, in order to bless the lives of children. Our opportunities are limitless. They are everywhere to be found—sometimes very close to home.

Last summer I received a letter from a woman who has emerged from a long period of Church inactivity. She is ever so anxious for her husband, who as yet is not a member of the Church, to share the joy she now feels.

She wrote of a trip which she, her husband, and their three sons made from the family home to Grandmother’s home in Idaho. While driving through Salt Lake City, they were attracted by the message which appeared on a billboard. The message invited them to visit Temple Square. Bob, the nonmember husband, made the suggestion that a visit would be pleasant. The family entered the visitors’ center, and Father took two sons up a ramp that one called “the ramp to heaven.” Mother and three-year-old Tyler were a bit behind the others, they having paused to appreciate the beautiful paintings which adorned the walls. As they walked toward the magnificent sculpture of Thorvaldsen’s Christus, tiny Tyler bolted from his mother and ran to the base of the Christus, while exclaiming, “It’s Jesus! It’s Jesus!” As Mother attempted to restrain her son, Tyler looked back toward her and his father and said, “Don’t worry. He likes children.”

After departing the center and again making their way along the freeway toward Grandmother’s, Tyler moved to the front seat next to his father. Dad asked him what he liked best about their adventure on Temple Square. Tyler smiled up at him and said, “Jesus.”

“How do you know that Jesus likes you, Tyler?”

Tyler, with a most serious expression on his face, looked up at his father’s eyes and answered, “Dad, didn’t you see his face?” Nothing else needed to be said.

As I read this account, I thought of the statement from the book of Isaiah: “And a little child shall lead them” (Isa. 11:6).

The words of a Primary hymn express the feelings of a child’s heart:

\begin{verse}
Tell me the stories of Jesus I love to hear,\\
Things I would ask him to tell me if he were here.\\
Scenes by the wayside, tales of the sea,\\
Stories of Jesus, tell them to me.\\
Oh, let me hear how the children stood round his knee.\\
I shall imagine his blessings resting on me;\\
Words full of kindness, deeds full of grace,\\
All in the lovelight of Jesus’ face.
\end{verse}

[“Tell Me the Stories of Jesus,”Children’s Songbook (1989), p. 57]

I know of no more touching passage in scripture than the account of the Savior blessing the children, as recorded in 3 Nephi. The Master spoke movingly to the vast multitude of men, women, and children. Then, responding to their faith and the desire that He tarry longer, He invited them to bring to Him their lame, their blind, and their sick, that He might heal them. With joy they accepted His invitation. The record reveals that “he did heal them every one” (3 Ne. 17:9). There followed His mighty prayer to His Father. The multitude bore record: “The eye hath never seen, neither hath the ear heard, before, so great and marvelous things as we saw and heard Jesus speak unto the Father” (v. 16).

Concluding this magnificent event, Jesus “wept, … and he took their little children, one by one, and blessed them, and prayed unto the Father for them. …

“And he spake unto the multitude, and said unto them: Behold your little ones.

“And as they looked to behold they cast their eyes towards heaven, and they saw the heavens open, and they saw angels descending out of heaven … ; and they came down and encircled those little ones … ; and the angels did minister unto them” (vs. 21, 23–24).

Over and over in my mind I pondered the phrase, “Whosoever shall not receive the kingdom of God as a little child, he shall not enter therein” (Mark 10:15).

One who fulfilled in his life this admonition of the Savior was a missionary, Thomas Michael Wilson. He is the son of Willie and Julia Wilson, Route 2, Box 12, Lafayette, Alabama. Elder Wilson completed his earthly mission on January 13, 1990. When he was but a teenager, and he and his family were not yet members of the Church, he was stricken with cancer, followed by painful radiation therapy, and then blessed remission. This illness caused his family to realize that not only is life precious but that it can also be short. The family began to look to religion to help them through this time of tribulation. Subsequently they were introduced to the Church and baptized. After accepting the gospel, young Brother Wilson yearned for the opportunity of being a missionary. A mission call came for him to serve in the Utah Salt Lake City Mission. What a privilege to represent the family and the Lord as a missionary!

Elder Wilson’s missionary companions described his faith as like that of a child—unquestioning, undeviating, unyielding. He was an example to all. After eleven months, illness returned. Bone cancer now required the amputation of his arm and shoulder. Yet he persisted in his missionary labors.

Elder Wilson’s courage and consuming desire to remain on his mission so touched his nonmember father that he investigated the teachings of the Church and also became a member.

An anonymous caller brought to my attention Elder Wilson’s plight. She said she didn’t want to leave her name and indicated she’d never before called a General Authority. However, she said, “You don’t often meet someone of the caliber of Elder Wilson.”

I learned that an investigator whom Elder Wilson had taught was baptized at the baptistry on Temple Square but then wanted to be confirmed by Elder Wilson, whom she respected so much. She, with a few others, journeyed to Elder Wilson’s bedside in the hospital. There, with his remaining hand resting upon her head, Elder Wilson confirmed her a member of The Church of Jesus Christ of Latter-day Saints.

Elder Wilson continued month after month his precious but painful service as a missionary. Blessings were given; prayers were offered. The spirit of his fellow missionaries soared. Their hearts were full. They lived closer to God.

Elder Wilson’s physical condition deteriorated. The end drew near. He was to return home. He asked to serve but one additional month. What a month this was! Like a child trusting implicitly its parents, Elder Wilson put his trust in God. He whom Thomas Michael Wilson silently trusted opened the windows of heaven and abundantly blessed him. His parents, Willie and Julia Wilson, and his brother Tony came to Salt Lake City to help their son and brother home to Alabama. However, there was yet a prayed-for, a yearned-for, blessing to be bestowed. The family invited me to come with them to the Jordan River Temple, where those sacred ordinances which bind families for eternity, as well as for time, were performed.

I said good-bye to the Wilson family. I can see Elder Wilson yet as he thanked me for being with him and his loved ones. He said, “It doesn’t matter what happens to us in this life as long as we have the gospel of Jesus Christ and live it.” What courage. What confidence. What love. The Wilson family made the long trek home to Lafayette, where Elder Thomas Michael Wilson slipped from here to eternity.

President Kevin K. Meadows, Elder Wilson’s branch president, presided at the funeral services. The words of his subsequent letter to me I share with you today: “On the day of the funeral, I took the family aside and expressed to them, President Monson, the sentiments you sent to me. I reminded them of what Elder Wilson had told you that day in the temple, that it did not matter whether he taught the gospel on this or the other side of the veil, so long as he could teach the gospel. I gave to them the inspiration you provided from the writings of President Joseph F. Smith—that Elder Wilson had completed his earthly mission and that he, as all ‘faithful elders of this dispensation, when they depart from mortal life, continue their labors in the preaching of the gospel of repentance and redemption, through the sacrifice of the Only Begotten Son of God, among those who are in darkness and under the bondage of sin in the great world of the spirits of the dead’ [D\&C 138:57]. The spirit bore record that this was the case. Elder Thomas Michael Wilson was buried with his missionary name tag in place.”

When Elder Wilson’s mother and his father visit that rural cemetery and place flowers of remembrance on the grave of their son, I feel certain they will remember the day he was born, the pride they felt, and the genuine joy that was theirs. This tiny child they will remember became the mighty man who later brought to them the opportunity to achieve celestial glory. Perhaps on these pilgrimages, when emotions are close to the surface and tears cannot be restrained, they will again thank God for their missionary son, who never lost the faith of a child, and then ponder deep within their hearts the Master’s words, “And a little child shall lead them” (Isa. 11:6).

Peace will then be their blessing. It will be our blessing, also, as we remember and follow the Prince of Peace. That we may do so is my sincere prayer. In the name of Jesus Christ, amen.

\newpage
\settalk{Conference Is Here}
\setspeaker{Thomas S. Monson}
\settalkdate{April 1990}
\talkheading{Conference Is Here}{Thomas S. Monson}

President Benson has suggested that I commence this conference with a brief message given in his behalf and that I convey to listeners and viewers far and near his greeting, his love, and his blessing.

The spirit of spring is very much in evidence here on historic Temple Square. The manicured lawns have discarded their drab winter color and now appear as a carpet of green accentuating the elegant flower beds with their brilliant blooms. It is a period of renewal, a time of gratitude, and a season for reflection.

The world has experienced sweeping changes since last we met. A wall in Berlin has crumbled. Families now may join together on either side and experience the joy they have long been deprived of. In Poland, Hungary, Czechoslovakia, and the German Democratic Republic, the bells of freedom have sounded, heralding a new day for our time.

All of us remember, President Benson, that dark period following World War II when our members were near starvation and bordering on despair. Then you undertook your dramatic assignment to supervise the distribution of food, clothing, and medical supplies from the storehouse of the Church to the war-devastated families in Europe.

Your words, President, echo loud and clear: “We must ‘dedicate our strength to serving the needs, rather than the fears, of the world.’ … I believe errands of mercy, such as the distribution of food, housing, and clothing to those in need, are rendered most effectively when handled by private individuals and organizations such as the Church” (The Teachings of Ezra Taft Benson [Salt Lake City: Bookcraft, 1988], p. 261).

In the spirit of President Benson’s counsel, we have a responsibility to extend help as well as hope to the hungry, to the homeless, and to the downtrodden both at home and abroad. Such assistance is being provided for the blessing of all. In a host of cities, where need has outdistanced help, lives have been lifted, hearts have been touched, and the frown of despair has been transformed to the smile of confidence, thanks to the generosity of the Church membership in the payment of their fast offerings as the Lord has commanded.

To the youth of the Church—President Benson has long been your champion and advocate. On a previous occasion he summed up the feelings of all your leaders when he declared, “Beloved youth, you will have your trials and temptations through which you must pass, but there are great moments of eternity which lie ahead. You have our love and our confidence. We pray that you will be prepared for the reins of leadership. We say to you, ‘Arise and shine forth’ (D\&C 115:5) and be a light unto the world, a standard to others” (“To ‘the Rising Generation,’” New Era, June 1986, p. 8).

My young brothers and sisters, from the days President Benson was a Scoutmaster to the present period of presiding over the entire Church, he has not forgotten you. He rejoices in your achievements, he admires your strengths. He is your friend and your advocate.

To the parents of the Church—President Benson has long urged that a good example is the best teacher. I have heard him offer sublime prayers to our Heavenly Father. Simple supplication, generous gratitude mark these petitions. Children joining parents in prayer will tend to be united with their families and found following the teachings of the Lord.

How President Benson and his beloved wife, Flora, enjoy attending the temple each week! His feeling for the temple is found in his statement: “I love the temples of God. This is the closest place to heaven on earth—the house of the Lord” (The Teachings of Ezra Taft Benson [Salt Lake City: Bookcraft, 1988], p. 253).

As this conference commences, I join President Benson, and know that I reflect the feelings of President Hinckley and all other General Authorities as well, in declaring our love for Heavenly Father’s children everywhere. Perhaps never in history has the need for cooperation, understanding, and goodwill among all people—nations and individuals alike—been so urgent as today. It is not only fitting, it is imperative that we emphasize the ideal of brotherhood and the responsibility true brotherhood confers upon us all.

As Edwin Markham observed:

\begin{verse}
There is a destiny that makes us brothers;\\
None goes his way alone:\\
All that we send into the lives of others\\
Comes back into our own.
\end{verse}

Let us live the commandments of God. Let us follow in the footsteps of His Son and our Savior, even Jesus Christ the Lord. As we sincerely and fervently seek Him, we shall indeed find Him.

He may come to us as one unknown, without a name—as of old, by the lakeside, He came to those men who knew Him not. He speaks to us the same words, “Follow thou me” (John 21:22), and sets us to the task which He has to fulfill for our time. He commands, and to those who obey Him, whether they be wise or simple, He will reveal Himself in the toils, the conflicts, the sufferings which they shall pass through in His fellowship; and they shall learn in their own experience who He is.

I bear to you my witness that God does live, that Jesus is the Christ, our Redeemer, and that we are led today by God’s prophet, even President Ezra Taft Benson. In the name of Jesus Christ, amen.

\newpage
\settalk{My Brother’s Keeper}
\setspeaker{Thomas S. Monson}
\settalkdate{April 1990}
\talkheading{My Brother’s Keeper}{Thomas S. Monson}

The Holy Bible is an inspiration to me. This sacred book has inspired the minds of men and has motivated readers to live the commandments of God and to love one another. It is printed in greater quantities, is translated into more languages, and has touched more human hearts than any other volume.

Particularly do I enjoy reading from the book of Genesis the account describing the creation of the world. Ponder the power of that culminating declaration: “God created man in his own image, in the image of God created he him; male and female created he them. And God blessed them” (Gen. 1:27–28).

Joy turns to sadness as we learn of Abel’s tragic death at the hands of his brother Cain. Chapters of counsel, lessons for living, guidance from God are found in one brief verse: “And the Lord said unto Cain, Where is Abel thy brother? And he said, I know not: Am I my brother’s keeper?” (Gen. 4:9).

These two significant questions are asked, then answered, in themes taught throughout the scriptures. One such example is found in the life of Joseph and his brothers. We will recall that Joseph was especially loved by his father, Jacob, which occasioned bitterness and jealousy on the part of his brothers. There followed the plot to slay him, which eventually placed Joseph in a pit without food and without water to sustain life. Upon the arrival of a passing caravan of merchants, Joseph’s brothers determined to sell him rather than to leave him to die. Twenty pieces of silver extricated Joseph from the pit and placed him eventually in the house of Potiphar in the land of Egypt. There Joseph prospered, for “the Lord was with Joseph” (Gen. 39:2).

After the years of plenty, there followed the years of famine. In the midst of this latter period, when the brothers of Joseph came to Egypt to buy corn, they were blessed by this favored man in Egypt—even their own brother. Joseph could have dealt harshly with his brethren for the callous and cruel treatment he had earlier received from them. However, he was kind and gracious to his brethren and won their favor and support with these words and actions:

“Now therefore be not grieved, nor angry with yourselves, that ye sold me hither: for God did send me before you to preserve life. …

“And God sent me before you to preserve you a posterity in the earth, and to save your lives by a great deliverance.

“So now it was not you that sent me hither, but God. …

“Moreover [Joseph] kissed all his brethren, and wept upon them: and after that his brethren talked with him” (Gen. 45:5, 7–8, 15).

They had found their brother. Joseph in very deed was his brothers’ keeper.

In the touching account of the good Samaritan, Jesus teaches vividly the interpretation of the lesson, “Thou shalt love thy neighbour as thyself” (Matt. 19:19). Answered effectively is the haunting question, “Am I my brother’s keeper?”

An entire vista of opportunity is unfolded to our view when we contemplate the magnitude of King Benjamin’s admonition, recorded in the Book of Mormon: “When ye are in the service of your fellow beings ye are only in the service of your God” (Mosiah 2:17).

Just last week the First Presidency and the Council of the Twelve were provided the opportunity to view the new Church history exhibit situated in the museum just west of Temple Square. I loved the replica of the entry to the Fourth Ward—one of the original wards in the valley. I noted with keen interest the lighted map which plotted the pioneer trek from Nauvoo. However, my heart was truly touched when I gazed at an actual handcart displayed in a place of honor. The handcart communicated to me a silent yet eloquent account of its long and momentous journey.

Let us for a moment join Captain Edward Martin and the handcart company he led. While we will not feel the pangs of hunger which they felt or experience the bitter cold that penetrated their weary bodies, we will emerge from our visit with a better appreciation of hardship borne, courage demonstrated, and faith fulfilled. We will witness with tear-filled eyes a dramatic answer to the question “Am I my brother’s keeper?”

“The handcarts moved on November 3 and reached the river, filled with floating ice. To cross would require more courage and fortitude, it seemed, than human nature could muster. Women shrank back and men wept. Some pushed through, but others were unequal to the ordeal.

“Three eighteen-year-old boys belonging to the relief party came to the rescue; and to the astonishment of all who saw, carried nearly every member of that ill-fated handcart company across the snow-bound stream. The strain was so terrible, the exposure so great, that in later years all the boys died from the effects of it. When President Brigham Young heard of this heroic act, he wept like a child, and later declared publicly, ‘That act alone will ensure C. Allen Huntington, George W. Grant, and David P. Kimball an everlasting salvation in the Celestial Kingdom of God, worlds without end’” (LeRoy R. Hafen and Ann W. Hafen, Handcarts to Zion [Glendale, California: The Arthur H. Clark Company, 1960], pp. 132–33).

Our service to others may not be so dramatic, but we can bolster human spirits, clothe cold bodies, feed hungry people, comfort grieving hearts, and lift to new heights precious souls.

Junius Burt of Salt Lake City, a longtime worker in the Streets Department, related a touching and inspirational experience. He declared that on a cold winter morning, the street cleaning-crew of which he was a member was removing large chunks of ice from the street gutters. The regular crew was assisted by temporary laborers who desperately needed the work. One such wore only a lightweight sweater and was suffering from the cold. A slender man with a well-groomed beard stopped by the crew and asked the worker, “You need more than that sweater on a morning like this. Where is your coat?” The man replied that he had no coat to wear. The visitor then removed his own overcoat, handed it to the man and said, “This coat is yours. It is heavy wool and will keep you warm. I just work across the street.” The street was South Temple. The good Samaritan who walked into the Church Administration Building to his daily work and without his coat was President George Albert Smith of The Church of Jesus Christ of Latter-day Saints. His selfless act of generosity revealed his tender heart. Surely he was his brother’s keeper.

In December of 1989, the beautiful and long-awaited Las Vegas temple was dedicated in inspiring sessions, which continued for three days. The messages and music in the dedicatory sessions lifted each heart heavenward and prompted the listener to keep the commandments of God and to emulate the example of righteous living taught by Jesus of Nazareth. Thoughts of self yielded to consideration for others. One sermon stressed the injunction of the Lord as recorded in Matthew:

“Lay not up for yourselves treasures upon earth, where moth and rust doth corrupt, and where thieves break through and steal:

“But lay up for yourselves treasures in heaven, where neither moth nor rust doth corrupt, and where thieves do not break through nor steal:

“For where your treasure is, there will your heart be also” (Matt. 6:19–21).

After the session during which this passage of scripture had been presented, a handwritten letter, carefully tucked away in a sealed envelope, was handed to me by an usher. May I share with you the contents of this touching letter:

“Dear President Monson:

“My husband and I feel the completion and dedication of this beautiful Las Vegas Nevada Temple is the finest gift we could receive during this sacred season. Temples are such a sweet gift to all the world; and as you spoke of righteous Saints who are worthy to obtain the blessings of the Lord’s house but lack the financial means to attend a temple, our hearts were so touched.

“President Monson, there must be a family somewhere who needs to attend the temple, because as my dear companion and I spoke of our great joy during this special Christmas season, we both commented as to how any store-bought gift would pale in comparison to what we have received in these dedicatory services. Instead of spending our budgeted Christmas funds for some gift from a local store, we would like to give you this \$500 to help some family waiting to be endowed and sealed for all eternity. We appreciate your assisting us in our gifts to each other this year.”

The letter was unsigned. The givers remain anonymous. Perhaps today this brother may be viewing this session of general conference. If so, he may be pleased to learn that this gift has made it possible for a worthy family from the Villa Real District of the Portugal Porto Mission to journey to the temple and receive their precious temple blessings. To the unknown givers of this priceless gift I extend my thanks for being your brother’s keeper. I have the inner feeling that your Christmas season was marked by joy and filled with peace.

We have no way of knowing when our privilege to extend a helping hand will unfold before us. The road to Jericho each of us travels bears no name, and the weary traveler who needs our help may be one unknown. Altogether too frequently the recipient of kindness shown fails to express his feelings, and we are deprived of a glimpse of greatness and a touch of tenderness that motivates us to go and do likewise. Genuine gratitude was expressed by the writer of a letter received recently at Church headquarters. No return address was shown, but the postmark was from Portland, Oregon:

“To the Office of the First Presidency:

“Salt Lake City showed me Christian hospitality once during my wandering years.

“On a cross-country journey by bus to California, I stepped down in the terminal in Salt Lake City, sick and trembling from aggravated loss of sleep caused by a lack of necessary medication. In my headlong flight from a bad situation in Boston, I had completely forgotten my supply.

“In the Temple Square Hotel restaurant, I sat dejectedly, cheekbones propped on fists, staring at a cup of coffee I really didn’t want. Out of the corner of my eye I saw a couple approach my table. ‘Are you all right, young man?’ the woman asked. I raised up, crying and a bit shaken, and related my story and the predicament I was in then. They listened carefully and patiently to my nearly incoherent ramblings, and then they took charge. They must have been prominent citizens. They spoke with the restaurant manager, then told me I could have all I wanted to eat there for five days. They took me next door to the hotel desk and got me a room for five days. Then they drove me to a clinic and saw that I was provided with the medications I needed—truly my basic lifeline to sanity and comfort.

“While I was recuperating and building my strength, I made it a point to attend the daily Tabernacle organ recitals. The celestial voicing of that instrument from the faintest intonation to the mighty full organ is the most sublime sonority of my acquaintance. I have acquired albums and tapes of the Tabernacle organ and the choir which I can rely upon any time to soothe and buttress a sagging spirit.

“On my last day at the hotel, before I resumed my journey, I turned in my key; and there was a message for me from that couple: ‘Repay us by showing gentle kindness to some other troubled soul along your road.’ That was my habit, but I determined to be more keenly on the lookout for someone who needed a lift in life.

“I wish you well. I don’t know if these are indeed the ‘latter days’ spoken of in the scriptures, but I do know that two members of your church were saints to me in my desperate hours of need. I just thought you might like to know.”

What a touching account. There comes to mind the experience of Jesus, when ten lepers were cleansed.

“And one of them, when he saw that he was healed, turned back, and with a loud voice glorified God,

“And fell down on his face at his feet. …

“And [Jesus] said unto him, Arise, go thy way: thy faith hath made thee whole” (Luke 17:15–16, 19).

The desire to help another, the quest for the lost sheep, may not always yield success at once. On occasion progress is slow—even indiscernible. Such was the experience of my longtime friend Gil Warner. He was serving as a newly called bishop when “Douglas,” a member of his ward, transgressed and was deprived of his Church membership. Father was saddened; Mother was totally devastated. Douglas soon thereafter moved from the state. The years hurried by, but Bishop Warner, now a member of a high council, never ceased to wonder what had become of Douglas.

In 1975, I attended the stake conference of the Parleys stake and held a priesthood leadership meeting early on the Sunday morning. I spoke of the Church discipline system and the need to labor earnestly and lovingly to rescue any who had strayed. Gil Warner asked to speak and then outlined the story of Douglas. He concluded with the question, “Who has the responsibility to work with Douglas and bring him back to Church membership?” Gil advised me later that my response to his question was direct and given without hesitation: “It is your responsibility, Gil, for you were his bishop, and he knew you cared.”

Unbeknownst to Gil Warner, Douglas’s mother had, the previous week, fasted and prayed that a man would be raised up to help save her son. Gil discovered this when he felt prompted to call her to report his determination to be of help.

Gil began his odyssey of redemption. Douglas was contacted by him. Old times, happy times, were remembered. Testimony was expressed, love was conveyed, and confidence instilled. The pace was excruciatingly slow. Discouragement frequently entered the scene; but, step by step, Douglas made headway. At long last prayers were answered, efforts rewarded, and victory attained. Douglas was approved for baptism.

The baptismal date was set, family members gathered, and former bishop Gil Warner flew to Seattle for the occasion. Can we appreciate the supreme joy felt by Bishop Warner as he, dressed in white, stood with Douglas in water waist deep and, raising his right arm to the square, repeated those sacred words, “Having been commissioned of Jesus Christ, I baptize you in the name of the Father, and of the Son, and of the Holy Ghost” (D\&C 20:73).

He that was lost was found. A 26-year mission, marked by love and pursued with determination, had been successfully completed. Gil Warner said to me, “This was one of the greatest days of my life. I know the joy promised by the Lord when He declared, ‘And if it so be that you should labor all your days … and bring, save it be one soul unto me, how great shall be your joy with him in the kingdom of my Father!’” (D\&C 18:15).

Were the Lord to say to Gil Warner today, as He said to Adam’s son long years ago, “Where is Douglas, thy brother?” Bishop Warner could reply, “I am my brother’s keeper, Lord. Behold Douglas, Thy son.”

May all of us who hold the priesthood of God demonstrate by our lives that we are our brothers’ keepers, I pray, in the name of Jesus Christ, amen.

\newpage
\settalk{The Lord’s Way}
\setspeaker{Thomas S. Monson}
\settalkdate{April 1990}
\talkheading{The Lord’s Way}{Thomas S. Monson}

As a foundation for my remarks this evening, I turn to the scriptures, that we might concentrate our thoughts on a passage familiar to most and applicable to all.

In the book of Malachi, the Lord instructs: “Even from the days of your fathers ye are gone away from mine ordinances, and have not kept them. Return unto me, and I will return unto you, saith the Lord of hosts. But ye said, Wherein shall we return?

“Will a man rob God? Yet ye have robbed me. But ye say, Wherein have we robbed thee? In tithes and offerings. …

“Bring ye all the tithes into the storehouse, that there may be meat in mine house, and prove me now herewith, saith the Lord of hosts, if I will not open you the windows of heaven, and pour you out a blessing, that there shall not be room enough to receive it.

“And I will rebuke the devourer for your sakes, and he shall not destroy the fruits of your ground; neither shall your vine cast her fruit before the time in the field, saith the Lord of hosts. …

“For ye shall be a delightsome land, saith the Lord of hosts.” (Mal. 3:7–8, 10–12.)

Elder Packer has mentioned briefly the matter of tithing. President Hinckley will amplify in greater detail this commandment of the Lord which pertains to the financing of His church here on earth. I set forth the passage that we might recognize that this is the Lord’s way—tithes and offerings—that His work may go forward so that His people may be blessed.

The newly announced local unit budget allowance program is but one of several carefully studied and prayerfully implemented steps taken by the Church to relieve the membership of financial burdens which some simply could not carry.

First to be introduced was the consolidated meeting schedule, that the time of Church members could be conserved and the cost of attending meetings reduced.

Second, there followed the introduction of increased Church participation in the construction of meetinghouses. Many of you will recall the time when meetinghouses were constructed under a participation ratio of 50:50, where the Church contributed half of the cost, with the other half coming from the members of the units who would occupy the building. This moved gradually to a 60:40 ratio, then to a 70:30 ratio, then to a 96:4 ratio, and finally to the welcome announcement that the total cost of building sites and the construction of buildings would be lifted from the local units altogether and provided in full through the tithes of the Church.

Third, there was eliminated the per-capita welfare assessment utilized to provide commodities to be distributed to the needy through the welfare program. Generous fast offerings would supplant the commodity budget.

Finally, the new local unit budget allowance program will replace local ward and stake budgets, with many costs heretofore borne by individual Church members now being covered through their tithes.

These steps were preceded by lengthy discussion and fervent prayer on the part of those sustained as prophets, seers, and revelators. In the words of the Apostle Paul, “This thing was not done in a corner.” (Acts 26:26.)

The details relating to the new budget allowance program have been communicated to bishops and stake presidents, that the membership of the Church in the United States and Canada might be advised. I shall not recount that detail this evening. However, a few comments on the objectives and the scope of the program might prove helpful.

Through the faith of Church members in the payment of tithing, it is now possible for the Church to provide in full to the wards and stakes in the United States and Canada the total costs incident to chapel site purchases; the construction of approved meetinghouses; provision of all utility and maintenance costs, including repairs and renovations; as well as the majority of custodial care of our buildings. This announcement was received with rejoicing, for the amounts of dollars involved were most substantial. Thus, a heavy burden was lifted from the backs of the members—made possible by increased faithfulness in the payment of tithing.

Every missionary who labors with diligence and love will benefit as he brings to our meetings investigators who will be able to worship in a dedicated chapel and learn of the saving principles of the gospel in the comfort of well-lighted and adequately heated buildings. Gone will be that ever-so-lengthy period when schools, lodges, and clubs had, of necessity, to be rented for worship services while a struggling branch or ward had to obtain from its membership—many of whom were recent converts—a substantial portion of the costs incident to site acquisition, meetinghouse construction and maintenance, the provision of utilities, as well as expecting them to pay an honest tithing, a generous fast offering, and a host of other related contributions.

Well do I remember the comment of a family in the mission field who were investigating Church membership. The missionaries brought them to the basement of the local Moose Lodge, where the branch met, and said to them, “This is where you will find the Spirit of the Lord—here in His true church.” Hesitatingly, but with curiosity, the parents turned to the moose-head on the wall and asked, “What is the significance of the animal head as pertains to your religious beliefs?” When the missionaries explained that these were temporary meeting facilities, the next question was, “Is your church here in Sudbury on a temporary basis?” The new program will help to eliminate this problem.

How grateful I am for this giant step forward in funding all such costs through tithing—even the Lord’s way.

Not so well understood, and perhaps less appreciated, is the announcement pertaining to local unit budgets. It will be helpful if we keep in mind the principles that govern the budget allowance program:

One of our objectives has been to insure that all budget costs be funded either through the 100-percent reimbursed items or the per-person budget allowance and that there be no separate assessments or fund-raising activities to support the programs of the Church. An exception would be that relating to our affiliation with the Scouting program, which has as a basic tenet that a boy earns his own way. Permitted under the budget allowance program is the financing of prescribed Cub, Scout, Varsity, and Explorer activities. This same exception is made for Young Women for camping activities outlined in the Young Women Handbook. It is the desire that restraint be used in programming youth activities and that consistency between Young Women and Young Men programs be achieved.

The primary responsibility for building testimonies and providing faith-building experiences in our members, including our youth, resides in the home. The Church should continue to support the determination of the family to do this. Priesthood leaders will wish to increase their efforts to build strong, gospel-centered homes. Families vary in size and composition. All are to receive our devoted attention. The building of testimonies is not related to financial costs. It is not necessary to buy the activity of our youth. Our youth activities depart from the pattern of the world.

To measure the goodness of life by its delights and pleasures is to apply a false standard. The abundant life does not consist of a glut of luxury. It does not make itself content with commercially produced pleasure, mistaking it for joy and happiness.

To find real happiness, we must seek for it in a focus outside ourselves. No one has learned the meaning of living until he has surrendered his ego to the service of his fellowmen. Service to others is akin to duty, the fulfillment of which brings true joy.

In some respects, many of our youth activities in recent years have supplanted the home and family. There has been a tendency to trend in our thinking to the position that an activity must be exotic to be successful. Faraway places with strange sounding names beckon as a Pied Piper for our youth to follow. Featured in our Church publications at times are glowing accounts of excursions to Hawaii, the Sacred Grove, historical sites, and other tempting locations. The word spreads, the cost escalates, and yearning increases, while objectives dim and time commitments of leaders and youth border on the burdensome. Errantly, we have used the term “super-activity” to encourage the exotic rather than the practical.

Many units are now planning major youth conferences on a two-year or three-year basis rather than each year. Some have discovered that through careful scheduling, there are sites and facilities very close to home available for productive youth activities. One stake reported holding its youth conference at the stake center, utilizing the parking lot and grounds for some of the functions and the recreational hall and chapel for others. The report: “One of the finest youth conferences we have ever held!”

When we turn our attention to outdoor encampments, let us remember that the same moon, the same stars shine forth from the heavens from hilltops close to home as the ones which shine over the Himalayas. The campfire glow, the sharing experience, lessons from leaders, and that inner feeling of closeness to God do not depend on distance. They are available to all.

In every location there are places of historical significance which can provide a focal point for a successful activity. You can identify such treasures. Even the local cemetery is a backdrop for effective teaching.

Ever shall I remember a bus trek from Salt Lake City to the Clarkston Cemetery near Logan, Utah, which involved all the Aaronic Priesthood in the stake. There, in the quiet of Clarkston, we gathered the youth around the grave of Martin Harris, one of the Three Witnesses of the Book of Mormon.

With rapt attention the young men listened to a brief retelling of the experiences and testimony of Martin Harris. The value of the Book of Mormon in the mind of each youth soared. Then we walked a few paces to a pioneer grave with a very small marker. It bore the name of John P. Malmberg and contained the verse:

\begin{verse}
A light from our household is gone;\\
A voice we loved is stilled.\\
A place is vacant in our hearts\\
That never can be filled.
\end{verse}

We talked with the boys about sacrifice, about dedication to the truth. Duty, honor, service, and love all were taught by that tombstone. I can, in the bright memory provided by God, see each boy reach for a handkerchief to wipe away a tear. Heard yet are the sniffles which testified that hearts were touched and commitments made. Each youth had determined to be a pioneer—even one who goes before, showing others the way to follow.

We retired as a group to a local park, where all enjoyed a feast of food. Before returning homeward, we stopped at the grounds of the Logan Temple. There we stretched out on the spacious lawn and gazed at a sky of blue, marked by white, billowy clouds hurried along on their journey by a steady breeze. We admired the beauty of the temple. We talked of sacred ordinances and eternal covenants. Silent pledges were made. Lessons were learned. Hearts were touched. Parents, grandparents, great-grandparents were remembered. Thoughts turned to the Master. His presence was close. His gentle invitation, “Follow me,” was somehow heard and felt.

Such experiences are available to all youth and their leaders. Their financial cost is minimal. The eternal dividends they yield are enormous. Souls indeed are precious in the sight of God. Ours is the power to help save these souls.

A letter from a member in the eastern part of the United States, received just last week, touched my heart. May I share this letter with you:

“Dear President Monson,

“I apologize for taking from your busy schedule, but felt I would be a most ungrateful father if I didn’t take a few minutes to express my sincere and personal thanks to you and the Brethren for the recent announcement on ward and stake budgets.

“Yes, we too are grateful to the Lord for this blessed and inspired day—not so much for the financial relief, but more for the hopeful reduction in ward and stake activities that will allow families to return home.

“It seems that over the years, we have become so activity-conscious that, unfortunately, though well-intended, our focus has been redirected from basic gospel principles to social events and concerns. I rejoice in the thought that the Lord is indeed causing local leaders to return back to basics, that our meetings and activities will focus on the Master and His blessed life. I sense this recent announcement is a big step in that direction.

“Please know of our great love and respect for you and the Brethren. We rejoice in the continual revelation the Lord is sending to His saints. Our prayers continue to be with the prophet and his associates in this great latter-day work.”

My brothers and sisters, these are momentous times in the history of the Church. The Lord has opened the windows of heaven and showered us with His benevolent blessings. Let us be faithful in our tithing obligation to the Lord. Let us be generous in our fast offerings, that the poor and needy may be blessed. Let us help in proclaiming His glorious message to all the world. Then we can be the beneficiaries of the Lord’s beautiful promise found in Isaiah. He gave the quiet assurance, Thou shalt “call, and the Lord shall answer; thou shalt cry, and he shall say, Here I am …

“Then shall thy light rise in obscurity, and thy darkness be as the noonday:

“And the Lord shall guide thee continually, and satisfy thy soul in drought, … and thou shalt be like a watered garden, and like a spring of water, whose waters fail not.” (Isa. 58:9–11.)

That we may, as a people, merit the blessings of God, I pray in the name of Jesus Christ, amen.

\newpage
\settalk{Days Never to Be Forgotten}
\setspeaker{Thomas S. Monson}
\settalkdate{October 1990}
\talkheading{Days Never to Be Forgotten}{Thomas S. Monson}

As the year 1990 moves inexorably toward its close, all members of The Church of Jesus Christ of Latter-day Saints can pause and reflect on the momentous happenings that have occurred in our time, in our day, and in our lives.

In the month of May, my wife and I were in the historic city of Berlin. We boarded a taxi and asked that the driver take us to the Berlin Wall. When the driver failed to respond to the direction provided, again the desired destination was given. Still no movement. Then he turned toward us and, in halting English, explained, “I can’t. The wall is kaput—gone!” We drove to the Brandenburg Gate. We viewed its restoration. We gazed from West Berlin to East Berlin, now one Berlin, and reflected on the events which followed the wall’s demise: a new mission of the Church established in Poland, another in Hungary, yet another in Greece, and a mission reestablished in Czechoslovakia. And now, official recognition of our Leningrad Branch in the Soviet Union. Who, except the Lord Himself, could have foreseen these historic events? It was He who declared, “This gospel of the kingdom shall be preached in all the world for a witness unto all nations” (Matt. 24:14). Surely the purposes of the Lord continue to unfold to our view if we but have eyes that truly see and hearts that know and feel.

Another transcendent blessing came the last weekend of August when a magnificent temple of the Lord was dedicated in Toronto, Ontario, Canada. In its gleaming glory, the temple seems to beckon to each who views its splendor, “Come! Come to the house of the Lord. Here is found ‘rest for the weary and peace for the soul.’”

And how the people did come! First they thronged to the public open house, where reverently and quietly they viewed the interior of the temple and learned the purpose for its erection and of the blessings which the temple can provide. One visitor described the temple’s beauty with the words, “This is a center of serenity.”

As she was about to leave the temple, a young Asian girl said, “Mommy, this is beautiful here. I don’t want to go.”

One woman surprised an usher with her request: “I have been so impressed with what I have seen. How do I join your church?”

Then came the faithful membership of the Church to the dedicatory sessions. From Ontario and Quebec they came. Others traveled from those cities in the United States which are a part of the temple district. Some journeyed to Toronto from the distant Maritime Provinces of Canada. None who came returned home disappointed.

A boy of tender years viewing the cornerstone-laying ceremony was, by the spirit of inspiration, called to take trowel in hand and assist in the sealing of the cornerstone.

Dora Valencia, who had lain four years in the Ajax Ontario Hospital, mustered her courage and fulfilled the desire to attend. From her hospital bed, which was wheeled into the celestial room, she not only basked in the spirit found there, but she also helped to provide that spirit. As I walked past her, upon leaving the room, and gazed at her expression of profound gratitude to the Lord, I bent low and took her hand in mine. Heaven was very near.

Angelic choirs lifted spirits heavenward as they sang the beautiful “Hosanna Anthem.” When the congregation joined with the choir to sing “The Spirit of God like a fire is burning,” no eye remained dry and no heart untouched.

Speakers recounted the history of the Church in the Toronto area, and the beautiful dedicatory prayer given at each session whispered peace. The words of Oliver Cowdery, spoken of another time, seemed to capture the spirit of the dedication: “These were days never to be forgotten” (JS—H 1:71, note).

As we recount the history of the Church in eastern Canada, we come to appreciate the tender feelings of the members of the Church on having a temple in their midst.

As early as April 1830, Phineas Young received a copy of the Book of Mormon from Samuel Smith, brother of the Prophet, and a few months later traveled to upper Canada. At Kingston, he gave the first known testimony of the restored Church beyond the borders of the United States.

The Prophet Joseph Smith, with Sidney Rigdon and Freeman Nickerson, was in Brantford and Mt. Pleasant, Ontario, in 1833. Joseph and Sidney had long been absent from their families and felt great concern about them. In the revelation we now know as the 100th section of the Doctrine and Covenants, the Lord counseled: “Verily, thus saith the Lord unto you, my friends Sidney and Joseph, your families are well; they are in mine hands, …

“Therefore, follow me, and listen to the counsel which I shall give unto you.

“Behold, … I have much people in this place, in the regions round about; and an effectual door shall be opened … in this eastern land” (D\&C 100:1–3).

Toward the people, the Prophet evidenced the same kind feelings that the Lord had shown to him and Sidney Rigdon. Of them he makes entries in his journal, such as, “The people were very tender and inquiring,” and again, “O God, seal our testimony to their hearts” (History of the Church, 1:422–23).

In 1836 Parley P. Pratt went to Canada following a great prophecy uttered by Heber C. Kimball in which Brother Pratt was instructed to go to Toronto. He was told that he would there find people waiting for him who would receive the gospel, and that from there the gospel would spread into England, where a great work would be done. In Toronto he found President John Taylor, the Fieldings, and many others.

In August of the next year, 1837, the Prophet Joseph Smith, with Sidney Rigdon and Thomas B. Marsh, then President of the Twelve Apostles, visited Toronto. Riding in a carriage and holding evening meetings by candlelight, they visited the churches. Elder Taylor accompanied them. He said: “This was as great a treat to me as I ever enjoyed. I had daily opportunities of conversing with them, of listening to their instructions, and in participating in the rich stores of intelligence that flowed continually from the Prophet Joseph.”

Recounting this history brings to my mind the experience of John E. Page as the Prophet Joseph Smith called him to serve a mission in Canada. “But I can’t go on a mission to Canada, Brother Joseph,” protested John E. Page. “I don’t even have a coat to wear.”

“Here,” said Joseph Smith, removing his own coat, “take this, and the Lord will bless you.”

John E. Page left Kirtland, Ohio, May 31, 1836, on his first mission as an elder of the Church. He labored in Canada for two years. During that time, he traveled over five thousand miles, mostly on foot, and baptized some six hundred people.

One of the great families to join the Church in Canada was that of Archibald Gardner. From his journal, we learn of the family’s experience in Canada during the year 1843.

Robert Gardner describes the day of their baptism: “We went about a mile and a half into the woods to find a suitable stream. We cut a hole through ice eighteen inches thick. My brother William baptized me. … I was confirmed while sitting on a log beside the stream. …

“I cannot describe my feelings at the time and for a long time afterwards. I felt like a little child and was very careful of what I thought or said or did lest I might offend my Father in Heaven. Reading the Scriptures and secret prayer occupied my leisure time. I kept a pocket Testament constantly with me. When something on a page impressed me supporting Mormonism, I turned down a corner. Soon I could hardly find a desired passage. I had nearly all the pages turned down. I had no trouble believing the Book of Mormon. Everytime I took the book to read I had a burning testimony in my bosom of its truthfulness.”

Archibald Gardner added: “[My] mother … [accepted] the Gospel at once and whole heartedly, after hearing it. … Not long after contacting the new faith she became desperately ill, so ill that her life was despaired of. She insisted on being baptized. The neighbors said that if we put her in the water they would have us tried for murder as she would surely die. Nevertheless, well bundled up, and tucked into a sleigh, we drove her two miles to the place appointed. Here a hole was cut in the ice and she was baptized in the presence of a crowd of doubters who had come to witness her demise. She was taken home. Her bed was prepared but she said, ‘No, I do not need to go to bed. I am quite well.’ And she was” (in Delilah G. Hughes, The Life of Archibald Gardner: Utah Pioneer of 1847 [Draper, Utah: Review and Preview Publishers, 1970], pp. 25–27).

Down through the years, this same spirit of faith and confidence in the Lord has continued. During the period 1959 to 1962, my family and I lived in Toronto, where I served as the mission president. We are witnesses to the love God has for the Saints in that area. May I describe some of these “never to be forgotten” events?

One situation featured the Donald Mabey family. Brother Mabey had moved his family from Salt Lake City to North Bay, Ontario, because of a business transfer by his company. Don was an elder in the Church but had been less than fully active in priesthood callings. He was about thirty-five years of age at the time and had a lovely family. The North Bay Branch was a struggling unit desperately in need of priesthood leadership. When I attended that branch and recognized this fact, I held an interview with Brother Mabey and said to him, “I am calling you to serve in the presidency of the North Bay Branch.”

He replied, “I can’t do it.”

I asked, “Why?”

He answered, “I have never done it before.”

“That’s no hindrance,” I responded. I took fresh hope from Don’s name, Mabey, and the words of a once-popular ballad, “Please don’t say no—say maybe.”

Brother Mabey said yes. Today he is a high priest living here in the West. All of his family members have entered temple doors and have received temple blessings.

Another evidence of faith took place when I first visited the St. Thomas Branch of the mission, situated about 120 miles from Toronto. My wife and I had been invited to attend the branch sacrament meeting and to speak to the members there. As we drove along a fashionable street, we saw many church buildings and wondered which one was ours. None was. We located the address which had been provided and discovered it to be a decrepit lodge hall. Our branch met in the basement of the lodge hall and was composed of perhaps twenty-five members, twelve of whom were in attendance. The same individuals conducted the meeting, blessed and passed the sacrament, offered the prayers, and sang the songs.

At the conclusion of the services, the branch president, Irving Wilson, asked if he could meet with me. At this meeting, he handed to me a copy of the Improvement Era, forerunner of today’s Ensign. Pointing to a picture of one of our new chapels in Australia, President Wilson declared, “This is the building we need here in St. Thomas.”

I smiled and responded, “When we have enough members here to justify and to pay for such a building, I am sure we will have one.” At that time, the local members were required to raise 30 percent of the cost of the site and the building, in addition to the payment of tithing and other offerings.

He countered, “Our children are growing to maturity. We need that building, and we need it now!”

I provided encouragement for them to grow in numbers by their personal efforts to fellowship and teach. The outcome is a classic example of faith, coupled with effort and crowned with testimony.

President Wilson requested six additional missionaries to be assigned to St. Thomas. When this was accomplished, he called the missionaries to a meeting in the back room of his small jewelry store, where they knelt in prayer. He then asked one elder to hand to him the yellow-page telephone directory, which was on a nearby table. President Wilson took the book in hand and observed, “If we are ever to have our dream building in St. Thomas, we will need a Latter-day Saint to design it. Since we do not have a member who is an architect, we will simply have to convert one.” With his finger moving down the column of listed architects, he paused at one name and said, “This is the one we will invite to my home to hear the message of the Restoration.”

President Wilson followed the same procedure with regard to plumbers, electricians, and craftsmen of every description. Nor did he neglect other professions, feeling a desire for a well-balanced branch. The individuals were invited to his home to meet the missionaries, the truth was taught, testimonies were borne and conversion resulted. Those newly baptized then repeated the procedure themselves, inviting others to listen, week after week and month after month.

The St. Thomas Branch experienced marvelous growth. Within two and one-half years, a site was obtained, a beautiful building was constructed, and an inspired dream became a living reality. That branch is now a thriving ward in a stake of Zion.

When I reflect on the town of St. Thomas, I dwell not on the ward’s hundreds of members and many dozens of families; rather, in memory I return to that sparse sacrament meeting in the lodge-hall basement and the Lord’s promise, “Where two or three are gathered together in my name, there am I in the midst of them” (Matt. 18:20).

Temples like the Toronto Temple are built with stone, glass, wood, and metal. But they are also a product of faith and an example of sacrifice. The funds to build temples come from all tithe payers and consist of the widow’s mite, children’s pennies, and workmen’s dollars—all sanctified by faith.

Whenever I attend a temple dedication, I think of Brother and Sister Gustav and Margarete Wacker of Kingston, Ontario. He was once the branch president of the Kingston Branch. He was from the old country. He spoke English with a thick accent. He never owned or drove a car. He plied the trade of a barber. He made but little money cutting hair near an army base at Kingston. How he loved the missionaries! The highlight of his day would be when he had the privilege to cut the hair of a missionary. Never would there be a charge. When they would make a feeble attempt to pay him, he would say, “Oh, no; it is a joy to cut the hair of a servant of the Lord.” Indeed, he would reach deep into his pockets and give the missionaries all of his tips for the day. If it were raining, as it often does in Kingston, President Wacker would call a taxi and send the missionaries to their apartment by cab, while he, himself, at day’s end would lock the small shop and walk home—alone in the driving rain.

I first met Gustav Wacker when I noticed that his tithing was far in excess of that expected from his potential income. My efforts to explain to him that the Lord required no more than a tenth fell on attentive but unconvinced ears. He simply responded that he loved to pay all he could to the Lord. It amounted to about a third of his income. His dear wife felt exactly as he did. Their unique manner of tithing payment continued.

Gustav and Margarete Wacker established a home that was a heaven. They were not blessed with children but mothered and fathered their many Church visitors. A sophisticated and learned Church leader from Ottawa told me, “I like to visit the Wacker home. I come away refreshed in spirit and determined to ever live close to the Lord.”

Did our Heavenly Father honor such abiding faith? The branch prospered. The membership outgrew the rented Slovakian Hall where they met and moved into a modern and lovely chapel of their own to which the branch members had contributed their share and more, that it might grace the city of Kingston. President and Sister Wacker had their prayers answered by serving a proselyting mission to their native Germany and later a temple mission to that beautiful temple in Washington, D.C. Then, in 1983, his mission in mortality concluded, Gustav Wacker peacefully passed away while being held in the loving arms of his eternal companion, dressed in his white temple suit, there in the Washington Temple.

All of this history and much more crowded my mind during the dedication services of the Toronto Temple. I reflected on the many nationalities represented by our members there. English, Scottish, German, French, and Italian predominated, but there were also members from Greece, Hungary, Finland, Holland, Estonia, and Poland. Surely, Toronto is an example of the promise of the Lord found in Jeremiah: “I will take you one of a city, and two of a family, and I will bring you to Zion” (Jer. 3:14). This He has done; and from this Zion called Toronto, the word now goes forth in these native tongues to the home nations of those He has gathered.

When I prepared to leave Toronto following the concluding dedicatory session, I gazed upward toward heaven, that I might offer a silent prayer of gratitude to God for His watchful care, His bounteous blessings, and for “days never to be forgotten.” High above the gleaming white temple, which personifies purity and reflects righteousness, is the gold-leafed statue of the Angel Moroni. I remembered being told that from that height of 105 feet, on a clear day one can see all the way to Cumorah. I noted that in Moroni’s hand was his familiar trumpet. He was gazing homeward—homeward to Cumorah. The beautiful Toronto Temple prepares all who enter to return homeward—homeward to heaven, homeward to family, homeward to God.

That all of us may travel safely to our eternal home is my humble prayer. In the name of Jesus Christ, amen.

\newpage
\settalk{“God Be with You Till We Meet Again”}
\setspeaker{Thomas S. Monson}
\settalkdate{October 1990}
\talkheading{“God Be with You Till We Meet Again”}{Thomas S. Monson}

There is a loneliness in the empty chair between President Hinckley and me, and we feel it in our hearts. I wish I could take you with President Hinckley and me to President Benson’s hospital room, which we visited a few days ago. I think the picture of tranquillity and love which was there would be very beneficial for all members of the Church to see. President Benson lay on his hospital bed, his left hand held by a noble son and his right hand by a beautiful daughter as she read to him from the Book of Mormon. In the background, a recording of Tabernacle Choir music played softly. It was just a little bit of heaven.

As we come to the close of another conference, our spirits have been lifted, our minds inspired, and our souls filled.

The messages delivered at this pulpit have provided words of counsel and guidance for our journey through mortality. The prayers have been offered with humility, and their petitions reflect the feelings of our hearts. The angelic music provided by the choirs at each session has confirmed the Lord’s words that “the song of the righteous is a prayer unto me, and it shall be answered with a blessing upon their heads” (D\&C 25:12).

We sincerely regret that President Ezra Taft Benson has been unable to be with us here in the Tabernacle. Nonetheless, we have felt his spirit throughout the proceedings. His love of the Lord, for the membership of the Church, and for God’s children everywhere is legendary. His many acts of kindness have blessed the lives of those with whom he has met everywhere he has gone.

One Friday, he and Sister Benson followed their usual practice of attending a session at the Jordan River Temple. While there, President Benson was approached by a young man who greeted him with joy in his heart and announced that he had been called to fill a full-time mission. President Benson took the newly called missionary by the hand and, with a smile on his lips, declared, “Take me with you! Take me with you!” That missionary testified that, in a way, he took President Benson with him on his mission, since this greeting demonstrated President Benson’s abiding love, his devotion to missionary work, and his desire to ever be found in the service of the Lord.

With the rapidly developing changes on the face of Europe, we remember President Benson’s great service to the hungry and to the homeless on that continent at the close of World War II. In attendance today is one who was the recipient of such service. She recently wrote to President Benson: “This is the first time in my life that I am here in Salt Lake City to attend general conference. I hope you will remember our first acquaintance in the autumn of 1946 in Langen, Germany. You and I will never forget the remarkable days following the Second World War. We will never forget your help for the refugees in those sad days. Now, forty-four years have gone, and we have both grown older. I wish you happiness and the blessings of the Lord all the days of your life and send you all my love.”

If President Benson were here at the pulpit at this, the conclusion of the final session of this glorious conference, he would extend to you his love, his admonitions, and his blessing. May I, with President Benson’s own words, provide you his counsel:

“Let us be valiant in our testimony of Jesus all the days of our lives” (Come unto Christ [Salt Lake City: Deseret Book Co., 1983], p. 16).

“His word is one of the most valuable gifts He has given us. I urge you to recommit yourselves to a study of the scriptures. Immerse yourselves in them daily so you will have the power of the Spirit to attend you. … Read them in your families and teach your children to love and treasure them” (“The Power of the Word,” Ensign, May 1986, p. 82).

“It is soul-satisfying to know that God is mindful of us and ready to respond when we place our trust in Him and do that which is right. There is no place for fear among men and women who place their trust in the Almighty and who do not hesitate to humble themselves in seeking divine guidance through prayer. Though persecutions arise, though reverses come, in prayer we can find reassurance, for God will speak peace to the soul. That peace, that spirit of serenity, is life’s greatest blessing” (“Pray Always,” Ensign, Feb. 1990, p. 5).

He continues: “I am getting older and less vigorous and am so grateful for your prayers and for the support of my younger Brethren. I thank the Lord for renewing my body from time to time so that I can still help build His kingdom. … God willing, I intend to spend all my remaining days in that glorious effort” (in Conference Report, Oct. 1988, p. 5; or Ensign, Nov. 1988, p. 6).

President Benson is a man of love, and this love he would have me extend to you in his behalf. He has a beautiful voice and has often sung the melodic strains of a favorite hymn:

\begin{verse}
God be with you till we meet again;\\
By his counsels guide, uphold you;\\
With his sheep securely fold you.\\
God be with you till we meet again.
\end{verse}

[Hymns, 1985, no. 152]

To the membership of the Church and to God’s children everywhere, our prophet, President Ezra Taft Benson, conveys to you the tender feelings of his heart, his gratitude for your prayers, and his abiding love. God be with you, brothers and sisters, till we meet again, in the name of Jesus Christ, amen.

\newpage
\settalk{That We May Touch Heaven}
\setspeaker{Thomas S. Monson}
\settalkdate{October 1990}
\talkheading{That We May Touch Heaven}{Thomas S. Monson}

One of my most vivid memories is attending priesthood meeting as a newly ordained deacon and singing the opening hymn, “Come, All Ye Sons of God Who Have Received the Priesthood.” Tonight, to this capacity audience assembled in the Tabernacle and in chapels worldwide, I echo the spirit of that special hymn and say to you, Come, all ye sons of God who have received the priesthood; let us consider our callings, let us reflect on our responsibilities, let us determine our duty, and let us follow Jesus Christ, our Lord.

While we may differ in age, in custom, or in nationality, we hold membership in the same church and are united as one in our priesthood callings.

Two weeks ago I attended a sacrament meeting where the children responded to the theme, I Belong to The Church of Jesus Christ of Latter-day Saints. These boys and girls demonstrated they were in training for service to the Lord and to others. The music was beautiful, the recitations skillfully rendered, and the spirit heaven-sent. My eleven-year-old grandson had spoken of the First Vision as he presented his part on the program. Afterward, as he came to his parents and grandparents, I said to him, “Tommy, I think you are almost ready to be a missionary.”

He replied, “Not yet; there is much I have to learn.”

To help him and all youth prepare for their service to God, a new booklet, entitled For the Strength of Youth, has been published under the direction of the First Presidency and the Quorum of the Twelve. The booklet features standards from the writings and teachings of Church leaders and from scriptures, adherence to which will bring the blessings of our Heavenly Father and the guidance of His Son to each of us.

May I share with you, as I shared with the sisters in the women’s meeting held last week, portions of the introduction to this new guide to your mortal journey, this new road map to help you chart an undeviating course toward eternal life. The statement by the First Presidency begins:

“Our beloved young men and women,

“We want you to know that we love you. We have great confidence in you. …

“We desire everything in this world for you that is right and good. … You are choice spirits who have been held in reserve to come forth in this day when the temptations, responsibilities, and opportunities are the very greatest. You are at a critical time in your lives. …

“We counsel you to [be] morally clean. …

“You cannot do wrong and feel right. It is impossible! Years of happiness can be lost in the foolish gratification of a momentary desire for pleasure. …

“You can avoid the burden of guilt and sin and all of the attending heartaches … as you keep the standards outlined in the scriptures and emphasized in this pamphlet.

“We pray that you—the young and rising generation—will keep your bodies and minds clean, free from the contaminations of the world, that you will be fit and pure vessels to bear triumphantly the responsibilities of the kingdom of God in preparation for the second coming of our Savior” (For the Strength of Youth [pamphlet, 1990], pp. 3–5).

May I review with you, the young men of the Church, these special standards referred to in the introduction just read? There are twelve items, followed by a conclusion. I shall treat briefly each standard.

\intalkheader{1. Dating}

Begin to prepare for a temple marriage. Proper dating is a part of that preparation. In cultures where dating is appropriate, do not date until you are sixteen years old. Not all teenagers need to date or even want to. When you begin dating, go in groups or on double dates. Make sure your parents meet and become acquainted with those you date. Because dating is a preparation for marriage, date only those who have high standards.

Be careful to go to places where there is a good environment, where you won’t be faced with temptation.

A wise father said to his son, “If you ever find yourself in a place where you shouldn’t ought to be, get out!” Good advice for all of us.

\intalkheader{2. Dress and Appearance}

Servants of the Lord have always counseled us to dress modestly to show respect for our Heavenly Father and for ourselves. The way you dress sends messages about yourself to others and often influences the way you and others act. Dress in such a way as to bring out the best in yourself and those around you. Avoid extremes in clothing and appearance.

\intalkheader{3. Friendshipping}

Everyone needs good friends. Your circle of friends will greatly influence your thinking and behavior, just as you will theirs. When you share common values with your friends, you can strengthen and encourage each other. Treat everyone with kindness and dignity. Many nonmembers have come into the Church through friends who have involved them in Church activities.

\intalkheader{4. Honesty}

The oft-repeated adage is ever true: “Honesty’s the best policy” (Miguel de Cervantes, in Familiar Quotations, 14th ed., comp. John Bartlett [Boston: Little, Brown and Co., 1968], p. 197). A Latter-day Saint young man lives as he teaches and as he believes. He is honest with others. He is honest with himself. He is honest with God. He is honest by habit and as a matter of course. When a difficult decision must be made, he never asks himself, What will others think? but rather, What will I think of myself?

For some, there will come the temptation to dishonor a personal standard of honesty. In a business law class at the university I attended, I remember that one particular classmate never prepared for the class discussions. I thought to myself, How is he going to pass the final examination?

I discovered the answer when he came to the classroom for the final examination, on a winter’s day, wearing on his bare feet only a pair of sandals. I was surprised and watched him as the class began. All of his books had been placed upon the floor. He slipped the sandals from his feet; and then, with toes that he had trained and had prepared with glycerine, he skillfully turned the pages of one of the books which he had placed on the floor, thereby viewing the answers to the examination questions.

He received one of the highest grades in that course on business law. But the day of reckoning came. Later, as he prepared to take his comprehensive examination, for the first time the dean of his particular discipline said, “This year I shall depart from tradition and shall conduct an oral, rather than a written, test.” Our favorite, trained-toe expert found that he had his foot in his mouth on that occasion and failed the examination.

\intalkheader{5. Language}

How you speak and the words you use tell much about the image you choose to portray. Use language to build and uplift those around you. Profane, vulgar, or crude language and inappropriate or off-color jokes are offensive to the Lord. Never misuse the name of God or Jesus Christ. The Lord said, “Thou shalt not take the name of the Lord thy God in vain” (Ex. 20:7).

\intalkheader{6. Media: Movies, Television, Radio, Videocassettes, Books and Magazines}

Our Heavenly Father has counseled us to seek after “anything virtuous, lovely, or of good report or praiseworthy” (A of F 1:13). Whatever you read, listen to, or watch makes an impression on you.

Pornography is especially dangerous and addictive. Curious exploration of pornography can become a controlling habit, leading to coarser material and to sexual transgression.

Don’t be afraid to walk out of a movie, turn off a television set, or change a radio station if what’s being presented does not meet your Heavenly Father’s standards. In short, if you have any question about whether a particular movie, book, or other form of entertainment is appropriate, don’t see it, don’t read it, don’t participate.

Recently there appeared in the newspaper an observation by comedian Steve Allen. It describes one of the greater problems of our time:

“Steve Allen doesn’t find anything funny about television’s trend toward stronger language and adult-oriented themes. The veteran comedian lashed out at current television trends in an opinion piece published in the Los Angeles Times.

“The ‘flow is carrying all of us right into the sewer,’ he wrote. ‘The very sort of language parents forbid their children to use is now being encouraged not only by anything-goes cable entrepreneurs but the once high-minded networks,’ said Allen. ‘Shows that depict children and others using vulgar language only point up the collapse of the American family,’ he said.”

Perhaps Mr. Allen was referring to a review in a recent issue of Newsweek magazine entitled “A Season on the Brink.” “Desperate to outrun [competition], the Big Three [networks] launch lineups that are rocking, ribald, real … and risky,” reads the subheadline. A summary statement declares, “The networks … are suddenly turning the airwaves blue” (Harry F. Waters, 3 Sept. 1990, pp. 70–71).

\intalkheader{7. Mental and Physical Health}

The Apostle Paul declared, “Know ye not that ye are the temple of God, and that the Spirit of God dwelleth in you? … The temple of God is holy, which temple ye are” (1 Cor. 3:16–17). Nutritious meals, regular exercise, and appropriate sleep are necessary for a strong body, just as consistent scripture study and prayer strengthen the mind and spirit.

Hard drugs, wrongful use of prescription drugs, alcohol, coffee, tea, and tobacco products destroy your physical, mental, and spiritual well-being. Any form of alcohol, including beer, is harmful to your spirit and your body. Tobacco can enslave you, weaken your lungs, and shorten your life.

\intalkheader{8. Music and Dancing}

Music can help you draw closer to your Heavenly Father. It can be used to educate, edify, inspire, and unite. However, music can, by its tempo, beat, intensity, and lyrics, dull your spiritual sensitivity. You cannot afford to fill your minds with unworthy music. Dancing can be enjoyable and provide an opportunity to meet new people and strengthen friendships. Plan and attend dances where dress, grooming, lighting, dancing styles, lyrics, and music contribute to an atmosphere in which the Spirit of the Lord may be present.

\intalkheader{9. Sexual Purity}

Because sexual intimacy is so sacred, the Lord requires self-control and purity before marriage, as well as full fidelity after marriage. In dating, treat your date with respect, and expect your date to show that same respect for you. Tears inevitably follow transgression. Men, take care not to make women weep, for God counts their tears.

President David O. McKay advised, “I implore you to think clean thoughts.” He then made this significant declaration of truth: “Every action is preceded by a thought. If we want to control our actions, we must control our thinking.” Brethren, fill your minds with good thoughts, and your actions will be proper. May each one of you be able to echo in truth the line from Tennyson spoken by Sir Galahad: “My strength is as the strength of ten, because my heart is pure” (in Familiar Quotations, p. 647).

From ancient times comes an example which emphasizes this truth. Darius, through the proper rites, had been recognized as legitimate king of Egypt. His rival, Alexander, had been declared legitimate son of Ammon; he, too, was Pharaoh. Alexander found the defeated Darius on the point of death and laid his hands upon his head to heal him, commanding him to arise and resume his kingly power, concluding, “I swear unto thee, Darius, by all the gods, that I do these things truly and without fakery,” to which Darius replied with a gentle rebuke, “Alexander, my boy … do you think you can touch heaven with those hands of yours?”

Brethren, are we prepared to touch heaven as we fill our priesthood callings?

Recently, the author of a paper on teenage sexuality summed up his research by saying that he doesn’t see any major reduction ahead in the sexual activity of teenagers, in part because society sends teens a mixed message: advertisements and the mass media convey “very heavy messages that sexual activity is acceptable and expected,” inducements that sometimes drown out the warnings of experts and the pleas of parents. The Lord cuts through all the media messages with clear and precise language when He declares to us, “Be … clean” (3 Ne. 20:41).

Whenever temptation comes, remember the wise counsel of the Apostle Paul, who declared, “There hath no temptation taken you but such as is common to man: but God is faithful, who will not suffer you to be tempted above that ye are able; but will with the temptation also make a way to escape, that ye may be able to bear it” (1 Cor. 10:13).

\intalkheader{10. Sunday Behavior}

The Lord has given the Sabbath day for your benefit and has commanded you to keep it holy. Many activities are appropriate for the Sabbath. Bear in mind, however, that Sunday is not a holiday. Sunday is a holy day.

\intalkheader{11. Spiritual Help}

When you were confirmed a member of the Church, you received the right to the companionship of the Holy Ghost. He can help you make good choices. When challenged or tempted, you do not need to feel alone. Remember that prayer is the passport to spiritual power.

\intalkheader{12. Repentance}

If any has stumbled in his journey, there is a way back. The process is called repentance. Our Savior died to provide you and me that blessed gift. Though the path is difficult, the promise is real: “Though your sins be as scarlet, they shall be as white as snow” (Isa. 1:18).

Don’t put your eternal life at risk. Keep the commandments of God. If you have sinned, the sooner you begin to make your way back, the sooner you will find the sweet peace and joy that come with the miracle of forgiveness.

These, then, are the standards found in For the Strength of Youth. Joy and happiness come from living the way the Lord wants you to live and from service to God and others.

Our beloved President Ezra Taft Benson sends to you his greetings. He loves you. He trusts you. And how might you return that love, that trust?

You have a heritage: Honor it.

You will meet sin: Shun it.

You have the truth: Live it.

You have a testimony: Share it.

Spiritual strength frequently comes through selfless service. Some years ago, I visited the California Mission, where I interviewed a young missionary from Georgia. I recall saying to him, “Do you send a letter home to your parents every week?”

He replied, “Yes, Brother Monson.”

Then I asked, “Do you enjoy receiving letters from home?”

He didn’t answer. At length, I inquired, “When was the last time you had a letter from home?”

With a quavering voice, he responded, “I’ve never had a letter from home. Father’s just a deacon, and Mother’s not a member of the Church. They pleaded with me not to come. They said that if I left on a mission they would not be writing to me. What shall I do?”

I offered a silent prayer to my Heavenly Father: “What shall I tell this young servant of Thine, who has sacrificed everything to serve Thee?” And the inspiration came. I said, “Elder, you send a letter home to your mother and father every week of your mission. Tell them what you are doing. Tell them how much you love them, and then bear your testimony to them.”

He asked, “Will they then write to me?”

I responded, “Then they will write to you.”

We parted, and I went on my way. Months later I was attending a stake conference in Southern California when a young man came up to me and said, “Brother Monson, do you remember me? I’m the young missionary who had not received a letter from my mother or my father during my first nine months in the mission field. I’m the one to whom you said, ‘Send a letter home every week, Elder, and your parents will write to you.’” Then he asked, “Do you remember that promise, Elder Monson?”

I remembered. I inquired, “Have you heard from your parents?”

He reached into his pocket and took out a sheaf of letters with an elastic band around them, took a letter from the top of the stack and said, “Have I heard from my parents! Listen to this letter from my mother: ‘Son, we so much enjoy your letters. We’re proud of you, our missionary. Guess what? Dad has been ordained a priest. He’s preparing to baptize me. I’m meeting with the missionaries; and one year from now we want to come to California as you complete your mission, for we, with you, would like to become a forever family by entering the temple of the Lord.’” Then the young man put his hand in mine and asked, “Brother Monson, does Heavenly Father always answer prayers and fulfill Apostles’ promises?”

I replied, “When one has faith as you have demonstrated, our Heavenly Father hears such prayers and answers in His own way.”

Clean hands, a pure heart, and a willing mind had touched heaven. A blessing, heaven-sent, had answered the fervent prayer of a missionary’s humble heart.

Brethren, it is my prayer that we may so live that we, too, may touch heaven and be similarly blessed. In the name of Jesus Christ, amen.

\newpage
\settalk{The Lighthouse of the Lord}
\setspeaker{Thomas S. Monson}
\settalkdate{October 1990}
\talkheading{The Lighthouse of the Lord}{Thomas S. Monson}

My dear sisters, the spirit which permeates this meeting here in the historic Tabernacle and in hundreds of chapels and stake centers in many parts of the world is a reflection of your strength, your devotion, your goodness. To quote the words of the Lord: “Ye are the salt of the earth. … Ye are the light of the world. … Let your light so shine before men, that they may see your good works, and glorify your Father which is in heaven.” (Matt. 5:13–14, 16.)

Some of you are just approaching young womanhood, soon to leave the comfort of Primary and enter the exciting and challenging years as young women of The Church of Jesus Christ of Latter-day Saints. Others here include those not yet married, many of whom are your teachers. There are also mothers, grandmothers, even great-grandmothers, who, with an occasional tear, think back to the summertime of their youth and ponder the words of Longfellow:

\begin{verse}
How beautiful is youth! how bright it gleams\\
With its illusions, aspirations, dreams!\\
Book of Beginnings, Story without End;\\
Each maid a heroine, and each man a friend!
\end{verse}

(Henry Wadsworth Longfellow, Morituri Salutamus, in The Complete Poetical Works of Longfellow, Cambridge, Mass.: The Riverside Press, 1922, p. 311.)

All of you are sisters to one another and daughters of our Heavenly Father. It is with a humble and prayerful heart that I stand before you. I have always loved the words frequently quoted by President David O. McKay as he described you: “Woman was taken out of man—not out of his feet to be trampled underfoot, but out of his side to be equal to him, under his arm to be protected, and near his heart to be loved.”

But the thought that never fails to stir my soul is the simple and sage advice: “Men should take care not to make women weep, for God counts their tears.”

Do we in attendance tonight know who we are and what God expects us to become? Remember that the recognition of a power higher than oneself does not in any sense debase; rather, it exalts. If we will but realize that we have been created in the image of God, we will not find Him difficult to approach, for God did create “man in his own image, in the image of God created he him; male and female created he them.” (Gen. 1:27.) This knowledge, acquired through faith, will bring inner calm and profound peace.

Just twenty years ago many of you had not yet commenced your journey through mortality. Your abode was a heavenly home. We know relatively little concerning the details of our existence there—only that we were among those who loved us and were concerned for our eternal well-being. Then there arrived the period where earth life became necessary to our progress. Farewells were no doubt spoken, expressions of confidence given, and graduation to mortality achieved.

What a commencement service awaited each of us! Loving parents joyously welcomed us to our earthly home. Tender care and affectionate embrace awaited our every whim. Someone described a newborn child as “a sweet, new blossom of Humanity, Fresh fallen from God’s own home, to flower on earth.” (Gerald Massey, “Wooed and Won,” in The Home Book of Quotations, sel. Burton Stevenson, New York: Dodd, Mead and Co., 1956, p. 121.)

Those first years were precious, special years. Satan had no power to tempt us. We had not yet become accountable but were innocent before God. They were learning years.

Soon we entered that period some have labeled “the terrible teens.” I prefer “the terrific teens.” What a time of opportunity, a season of growth, a semester of development, marked by the acquisition of knowledge and the quest for truth.

No one has described these years as being easy. Indeed, they have become increasingly more difficult. The world seems to have slipped from the moorings of safety and drifted from the harbor of peace.

Permissiveness, immorality, pornography, and the power of peer pressure cause many to be tossed about on a sea of sin and crushed on the jagged reefs of lost opportunities, forfeited blessings, and shattered dreams.

Anxiously you ask, “Is there a way to safety? Can someone guide me? Is there an escape from threatened destruction? The answer is a resounding yes! I counsel you: Look to the lighthouse of the Lord. There is no fog so dense, no night so dark, no gale so strong, no mariner so lost but what its beacon light can rescue. It beckons through the storms of life. It calls, “This way to safety; this way to home.”

The lighthouse of the Lord sends forth signals readily recognized and never failing. These words of warning, these safety standards, are printed in a small booklet soon to be distributed and entitled For the Strength of Youth.

May I share with you the introduction to the booklet, prepared by the First Presidency of the Church:

“Our beloved young men and women,

“We want you to know that we love you. We have great confidence in you. …

“We desire everything in this world for you that is right and good. You are not just ordinary young men and women. You are choice spirits who have been held in reserve to come forth in this day when the temptations, responsibilities, and opportunities are the very greatest. You are at a critical time in your lives. This is a time for you not only to live righteously but also to set an example for your peers. …

“God loves you. … His desire … is to have you return to Him pure and undefiled, having proven yourselves worthy of an eternity of joy in His presence. …

“We counsel you to choose to live a morally clean life. …

“You cannot do wrong and feel right. It is impossible! Years of happiness can be lost in the foolish gratification of a momentary desire for pleasure. …

“You can avoid the burden of guilt and sin and all of the attending heartaches … as you keep the standards outlined in the scriptures and emphasized in this pamphlet. …

“We pray that you—the young and rising generation—will keep your bodies and minds clean, free from the contaminations of the world, that you will be fit and pure vessels to bear triumphantly the responsibilities of the kingdom of God in preparation for the second coming of our Savior.” (For the Strength of Youth, 1990, p. 1.)

May I review with you, the women of the Church, these special standards referred to in the introduction just read? There are twelve items, followed by a conclusion. I shall treat briefly each standard.

These, then, are the standards found in For the Strength of Youth. Joy and happiness come from living the way the Lord wants you to live and from service to God and others.

Our beloved President Ezra Taft Benson sends to you his love. He is your advocate for all that is good and clean and wholesome. He loves you. He trusts you. And how might you return that love, that trust?

You have a heritage: Honor it.

You will meet sin: Shun it.

You have the truth: Live it.

You have a testimony: Share it.

Spiritual strength frequently comes through selfless service. An actual experience involving young women, their teacher, and a widow illustrates this truth.

As the Christmas season approached, a teacher of Laurels arranged a visit to bring joy to the heart of a lonely widow, Jane. The girls busied themselves preparing delicious cookies, special refreshments—even a Christmas tree with ornaments to be placed thereon. They did not forget a beautiful corsage which they knew would brighten the soul of the special woman they planned to visit.

With their parcels tucked tightly under each arm, the girls and their teacher made their way up the never-ending flights of stairs which led to Jane’s apartment. There was an interminable delay as Jane’s aged feet made their way to the door. As she opened it, she greeted each of the beautiful young ladies and made them welcome in her humble apartment. Their smiles reflected the goodness of their hearts as they erected the Christmas tree and placed upon it the decorations they had so carefully carried. Then the packaged gifts were placed beneath its outstretched branches. I was there. I had never seen a more beautiful tree, for no tree had ever before been decorated with such love, such Christlike care and concern. The teacher slipped into the kitchen; and, aided by three of her girls, the refreshments were prepared and a feast enjoyed.

Then the dear widow gathered her guests around her to share with them her life’s precious memories. She told how, as a young girl in far-off Scotland, she had heard the missionaries, embraced the truth which they taught—even suffered the gibes and comments which adherence to a then-unpopular faith inevitably provoked. She told them how the entire Sabbath day was taken just to travel and to attend the meetings of her newfound faith. Instinctively, the girls compared the ease with which they made their way each Sunday to their chapel.

When Jane told them of the voyage to America, described the storm-tossed Atlantic and the warm feeling which came to the heart when the Statue of Liberty was first glimpsed, I noted that the girls were visibly touched. Tears brimmed in their eyes and pledges were made within their hearts—pledges to do that which is correct, to be honorable, to live true to the faith, and to abide by their standards.

As the evening came to an end, there were kisses and embraces; and then each young lady filed silently from the doorway and made her way down the staircases to the street outside. They left behind a mother filled with the goodness of the world, with love rekindled, with faith again inspired. I’m certain this was one of the happiest days of her life. That night the corsage was carefully and tenderly placed in safekeeping. It had become a symbol of all that is good and clean and wholesome.

Outside the snow was falling, and the girls could hear the crunch of their own footsteps on the snow-covered pavement. Words didn’t come easily, and then one Laurel girl asked, “Why is it I feel better than I’ve ever felt before?” Others nodded the same curiosity. I answered them, “Remember the words of the Master: ‘Inasmuch as ye have done it unto one of the least of these … , ye have done it unto me.’” (Matt. 25:40.) The words from the hymn “O Little Town of Bethlehem” seemed fitting. They were repeated:

\begin{verse}
How silently, how silently\\
The wondrous gift is giv’n!\\
So God imparts to human hearts\\
The blessings of his heav’n.\\
No ear may hear his coming;\\
But in this world of sin,\\
Where meek souls will receive him, still\\
The dear Christ enters in.
\end{verse}

(Hymns, 1985, no. 208.)

The dear Christ had indeed entered in—entered a humble home, entered a widow’s heart, and entered, forever to remain, in the soul of each girl. The lighthouse of the Lord had shown the way.

That this same spirit, even the Christ spirit, may ever be ours is my humble prayer. In the name of Jesus Christ, amen.

\newpage
\settalk{A Royal Priesthood}
\setspeaker{Thomas S. Monson}
\settalkdate{April 1991}
\talkheading{A Royal Priesthood}{Thomas S. Monson}

It will now be my opportunity to speak to you, brethren. As I assume this responsibility I must confess I do so in all humility and with a prayer in my heart that our Heavenly Father will bless me.

My brethren of the priesthood worldwide, you are an inspiring audience. To use a word favored by the youth of today, it is an “awesome” responsibility to speak to you.

There is a look of determination about you. You know who you are and what God expects you to become. As I consider the number of young men of the Aaronic Priesthood assembled tonight, I see a great future for you.

When I was about nine years old and attending elementary school here in Salt Lake City, all of the youth in the city’s schools were asked to fill out a form indicating what we wanted to be when we grew up. The lists were then to be placed in a waterproof metal box and buried beneath a new flagpole which graced the entrance to the City and County Building grounds. Years later, the box was to be opened and its contents made available.

As I sat with pencil in hand, I thought of the question, “What do I want to be when I grow up?” Almost without hesitation, I wrote the word cowboy. At lunch that day I reported to my mother my response. I can almost see Mother now as she admonished me, “You get right back to school and change that to banker or lawyer!” I obeyed Mother, and all dreams of being a cowboy vanished forever.

One of greater childhood determination was Steve Alford, who plays for the Dallas Mavericks team in the National Basketball Association. He remembers telling his eighth grade counselor, as she completed a career path form for him, that he was going to be an NBA player. She responded, “I can’t put that answer down.” Steve Alford replied, “Then leave it blank, ‘cause that’s what I’m going to do!” And he did.

One of the great leaders of our time, President Harold B. Lee, in a devotional address at BYU, spoke of a Latter-day Saint young man who, during World War II, was in England. He had gone to an officers’ club where they were holding a riotous kind of celebration. He noticed, off to the side, a young British officer who didn’t seem to appreciate the party at all. So he walked over to him and said, “You don’t seem to be enjoying this kind of party.” And this young British officer straightened himself a few inches taller than he was before and replied, “No, sir; I can’t engage in this kind of party because, you see, I belong to the royal household of England.” As our Latter-day Saint young man walked away, he said to himself, “Neither can I, because I belong to the royal household of the kingdom of God” (“Be Loyal to the Royal within You,” in Speeches of the Year, 1973 [Provo: Brigham Young University Press, 1973], p. 100).

Perhaps the young man remembered the bold declaration of the Apostle Peter: “Ye are a chosen generation, a royal priesthood, an holy nation, a peculiar people; that ye should shew forth the praises of him who hath called you out of darkness into his marvellous light” (1 Pet. 2:9). Brethren, be loyal to the royal within you.

My thoughts of late have focused upon the words of the Savior during the week of the atoning sacrifice, when He said: “Come, ye blessed of my Father, inherit the kingdom prepared for you from the foundation of the world:

“For I was an hungred, and ye gave me meat: I was thirsty, and ye gave me drink: I was a stranger, and ye took me in:

“Naked, and ye clothed me: I was sick, and ye visited me: I was in prison, and ye came unto me.

“Then shall the righteous answer him, saying, Lord, when saw we thee an hungred, and fed thee? or thirsty, and gave thee drink?

“When saw we thee a stranger, and took thee in? or naked, and clothed thee?

“Or when saw we thee sick, or in prison, and came unto thee?

“And the King shall answer and say unto them, Verily I say unto you, Inasmuch as ye have done it unto one of the least of these my brethren, ye have done it unto me” (Matt. 25:34–40).

Brethren, we wish to commend you for your faith in living the law of the fast and your generosity in the contribution of your fast offerings. We also compliment you deacons and teachers who assist in collecting the fast offerings in many parts of the world. The welfare program is divinely inspired, and those in need are being assisted by bishops who follow the inspiration of the Spirit and the principles of welfare in responding to those needs.

Beyond the ongoing assistance provided through the use of your regular fast offering contributions—and this assistance is most substantial—I felt tonight you would appreciate being informed of the current status of the special fasts and the donations affiliated with them. The proceeds from the two special fast days in 1985 and donations to the special relief of the suffering since that time have totaled \$13,145,527. The contributions have been utilized in the following locations: Africa, \$8,662,765, with the balance of the expenditures being distributed in the United States, Latin America, Asia, Europe, and the Middle East, with the total expended to date \$11,460,760, and a balance of \$1,684,767.

Let me share with you a little more detail concerning some of the projects and the people who have been blessed through your generosity.

In the fertile lowlands of eastern Guatemala, near the city of San Esteban, the Church and the Ezra Taft Benson Agriculture and Food Institute are helping poor rural farm families to increase agricultural production. By teaching techniques for improved soil preparation, fertilization, and irrigation, small farms achieve balanced cropping that provides better nutrition for families and additional feed for livestock.

At the outset, 160 families benefited from this instruction and assistance. Within a short time, the number of families will reach 400. As knowledge and skills are imparted among neighbors, many thousands stand to benefit.

Released from the confinement of poverty and want, they will then be better able to receive the spiritual gifts He holds in store for them. We, by our efforts to assist them, will better understand His words, “I was in prison, and ye came unto me” (Matt. 25:36).

The children in African nations are receiving immunizations in an effort to eradicate common communicable diseases by the end of the century. A specific project involves a cooperative effort with Rotary International’s Polio Plus endeavor. The Church has purchased sufficient polio serum to immunize 300,000 children. Gas and electric refrigerators have been placed in rural health outposts to keep vaccines viable until they are administered to the children. You, my brethren, and your families helped to bring this dream to reality.

Closer to this tabernacle, caring dentists joined together to provide free dental care to residents of an urban homeless shelter. These dentists, hygienists, and other professionals volunteer their time and skills. The Church has helped to provide the needed dental supplies.

These efforts not only relieve discomfort and pain, they also brighten the smiles, lift the spirits, and gladden the hearts of homeless patients. The words of the Master bring peace to the souls of all who participate in such endeavors: “I was a stranger, and ye took me in” (Matt. 25:35).

In the Philippines, the Church provides assistance to the Mabuhay Deseret Foundation, which aids hundreds of children to receive operations to repair deformed palates and lips and to correct untreated fractures or burns. Children once shunned now live normal lives. The spring of their step and the sound of their joy seem to echo, “I was sick, and ye visited me” (Matt. 25:36).

Generous contributions of wearing apparel to Deseret Industries are being used to clothe men, women, and children around the world. Clothing is sorted, sized, and shipped to locations as far distant as Romania, Peru, Zimbabwe, and Sierra Leone, as well as to cities in North America. This clothing has warmed and comforted those exiled in refugee centers and orphanages. The bright patterns and sound fabrics considered surplus by the donors are now new and wonderful attire to the aged and impoverished. Meaning is given to the words, “I was … naked and ye clothed me” (Matt. 25:35–36).

The Church’s humanitarian efforts are reaching the hungry and homeless of many American cities. Throughout the state of Utah, among the border towns of Texas, Arizona, and California, and into the communities of Appalachia, food and clothing are donated through private voluntary organizations or directly to children’s homes, food banks, and soup kitchens. Much of this food starts its long journey on production projects managed by local agent stakes. Food is processed and packaged in Church canneries and distributed through storehouses, where Church welfare recipients and volunteers labor to assist their poor and needy neighbors within and outside the Church. Many could say with feeling, “I was an hungred, and ye gave me meat” (Matt. 25:35).

Far away in the foothills on the western slopes of Mt. Kenya, along the fringe of the colossal Rift Valley, pure water is coming to the thirsty people. A potable water project has changed the lives of 1,100 families. In cooperation with TechnoServe, a private voluntary organization, the Church is assisting in a project that will pipe drinkable water through twenty-five miles of pipes to waiting homes in a fifteen-village area. The simple blessing of safe drinking water recalls the words of the Savior, “I was thirsty, and ye gave me drink” (Matt. 25:35).

In behalf of the hundreds of thousands who have benefited by your generous fast offering contributions—children who now walk, who smile, who are fed and clothed; and parents who now may live normal lives with their children—I extend to you, the priesthood of this Church, the heartfelt expression of so many: “Thank you, and may God bless you.”

Two thousand years ago, Jesus of Nazareth sat by a well in Samaria and talked to a woman about living waters: “Jesus … said unto her, Whosoever drinketh of this water shall thirst again:

“But whosoever drinketh of the water that I shall give him shall never thirst; but the water that I shall give him shall be in him a well of water springing up into everlasting life” (John 4:13–14).

The gospel of the Lord Jesus Christ provides all of us this cherished blessing. King Benjamin, in his memorable message, declared, “When ye are in the service of your fellow beings ye are only in the service of your God” (Mosiah 2:17).

Brethren of the priesthood, each of us is so employed. Ours is the responsibility to teach, to lift, to build, and to inspire our fellowmen, for “the worth of souls is great in the sight of God” (D\&C 18:10).

There are examples all around us of those who have recognized in others the need, even the thirst, for these “living waters,” who have through their own lives and service quenched this thirst and blessed these lives.

An example of true love and inspired teaching was found in the life of the late James Collier, who had, through his personal efforts, reactivated a large number of brethren in the Bountiful, Utah, area. I was invited by Brother Collier to address those who had now been ordained elders and who, with their wives and families, had been to the Salt Lake Temple to receive those eternal covenants and blessings for which they had so earnestly strived.

At the banquet honoring this achievement, I could see and I could feel the love that Jim had for those whom he had taught and rescued. Unfortunately, Jim Collier at that time was afflicted with a terminal illness and had to persuade the doctors to allow him to leave the hospital to attend this final night of recognition. As he stood at the pulpit, a large smile came over his face. With tear-filled eyes, he expressed his love to the group. There wasn’t a dry eye to be found. Brother Collier quipped, “Everyone wants to go to the celestial kingdom, but no one wants to die to get there.” Lowering his voice, he continued, “I’m prepared to go, but I will be there waiting on the other side to greet each of you, my beloved friends.” He returned to the hospital. His funeral service was held just a few weeks later.

May I conclude with two experiences from my own life: one from boyhood, one from manhood.

When I was a deacon, I loved baseball; in fact, I still do. I had a fielder’s glove inscribed with the name “Mel Ott.” He was the Darryl Strawberry of my day. My friends and I would play ball in a small alleyway behind the houses where we lived. The quarters were cramped but all right, provided you hit straight away to center field. However, if you hit the ball to the right of center, disaster was at the door. Here lived a lady who would watch us play, and as soon as the ball rolled to her porch her English setter would retrieve the ball and present it to Mrs. Shinas as she opened the door. Into her house Mrs. Shinas would return and add the ball to the many she had previously confiscated. She was our nemesis, the destroyer of our fun—even the bane of our existence. None of us had a good word for Mrs. Shinas, but we had plenty of bad words for her. The windows of her house received more special soap treatment on Halloween than did any other. None of us would speak to Mrs. Shinas, and she never spoke to us. She was hampered by a stiff leg which impaired her walking and must have caused her great pain. She and her husband had no children, lived secluded lives, and rarely came out of their house.

This private war continued for some time—perhaps two years—and then an inspired thaw melted the ice of winter and brought a springtime of good feelings to the stalemate. One night as I performed my daily task of hand-watering our front lawn, holding the nozzle of the hose in hand as was the style at that time, I noticed that Mrs. Shinas’s lawn was dry and turning brown. I honestly don’t know what came over me, but I took a few more minutes and, with our hose, watered her lawn. This I did each night, and then when autumn came, I hosed her lawn free of leaves as I did ours and stacked the leaves in piles at the street’s edge to be burned or gathered. During the entire summer I had not seen Mrs. Shinas. We had long since given up playing ball in the alley. We had run out of baseballs and had no money to buy more.

Then early one evening, her front door opened, and Mrs. Shinas beckoned for me to jump the small fence and come to her front porch. This I did, and as I approached her, Mrs. Shinas invited me into her living room, where I was asked to sit in a comfortable chair. She went to the kitchen and returned with a large box filled with baseballs and softballs, representing several seasons of her confiscation efforts. The filled box was presented to me; however, the treasure was not to be found in the gift, but rather in her voice. I saw for the first time a smile come across the face of Mrs. Shinas, and she said, “Tommy, I want you to have these baseballs, and I want to thank you for being kind to me.” I expressed my own gratitude to her and walked from her home a better boy than when I entered. No longer were we enemies. Now we were friends. The Golden Rule had again succeeded.

Brethren, at times those who most need our help appear to be least anxious to receive it. As I departed for the mission field to preside in Toronto, Canada, if anyone had asked me who of all the people I knew I would consider least likely to join the Church, I would have included the name of Shelley, a man I had known for many years. His sweet wife had tried in vain to interest him in the Church. A lovely daughter and precious son had both put forth their best efforts, with no perceptible change. Perhaps Shelley just couldn’t express his inner feelings or demonstrate positive emotions. In the ward, every effort had been expended, but to no avail. Shelley remained on the outside.

Perhaps it was the loss of his son to cancer which made the difference, or maybe the friendly conversation of a school crossing guard with whom Shelley visited sometimes in the morning and sometimes in the afternoon. Then again, faithful home teachers in the ward to which Shelley and his family had moved had helped to bring about the quiet miracle.

After an absence of three years, my family and I returned to our home in Salt Lake City. Time passed, and the next conversation I had with my friend Shelley was after I was called to the Twelve. One evening I received a telephone call from him. In his characteristic direct way, he asked if I would perform the ordinance in the temple which would seal his family for all eternity. I responded, “Shelley, that would be a privilege for me, but first you must become a member of the Church.” Can you imagine my surprise when he replied, “I have joined the Church. I now hold the Melchizedek Priesthood and am very active.”

What a special blessing, to welcome Shelley, his wife, Eugenia, his daughter, Utahna, and, by proxy, his son, Robert, to a beautiful sealing room in the Salt Lake Temple. The blessings of eternity were bestowed. Just three years later, I spoke at Shelley’s funeral services. He had progressed from doubt to faith and now had looked upward and gone forward, bidding farewell to mortality and receiving a welcome to paradise. Today he is with his beloved Eugenia, and they are with Robert and wait one day to welcome Utahna. When I reflect on the life of Shelley, I feel a debt of gratitude to that humble crossing guard, to those faithful home teachers, to that patient wife and daughter, and to all who made a difference in the unfolding of eternal blessings for Shelley and his family.

Our Lord and Savior said, “Come, follow me” (Luke 18:22). When we accept His invitation and walk in His footsteps, He will direct our paths. His gentle voice guides us in life’s journey and reminds us of our duty: “Lay not up for yourselves treasures upon earth, where moth and rust doth corrupt, and where thieves break through and steal:

“But lay up for yourselves treasures in heaven, where neither moth nor rust doth corrupt, and where thieves do not break through nor steal:

“For where your treasure is, there will your heart be also” (Matt. 6:19–21).

May we hear His voice. May we follow His example. May we live His teachings. Then we will be as the Apostle Peter declared, even a “royal priesthood.” May each of us earn the tribute spoken of our Lord: He “went about doing good, … for God was with him” (Acts 10:38). This is my prayer in the name of Jesus Christ, amen.

\newpage
\settalk{Never Alone}
\setspeaker{Thomas S. Monson}
\settalkdate{April 1991}
\talkheading{Never Alone}{Thomas S. Monson}

This Sabbath day has been designated as a day of thanksgiving, a day of gratitude—even a day of prayer. We pause, we ponder, we reflect on the blessings an all-wise Heavenly Father has bestowed upon us, His children, by bringing peace to the battlefield of war and comfort to the hearts of so many in this wonderful world where we live and which we call home.

Today knees will bow, bells will peal, hearts will swell, and voices will proclaim the glorious message “Thanks be to God.” In the United States of America, a grateful nation and a thankful president will give utterance to the tender feelings felt by all in a world that welcomed peace.

Who among us will ever forget the touching and vivid pictures of husbands and fathers bidding good-bye to weeping wives and wondering children as fond farewells dominated every newscast and printed story. The children cried but did not know why. Wives wept because they did know the danger, the loneliness, the fear that awaited.

With the wave of a hand and a somewhat forced smile, the men and the women of the military went off to war. Their farewell expressions even now ring the conviction of their hearts: “I love my country.” “I’m proud to serve.” “I’ll be home soon.” “Try not to worry.”

But worry they did. Constant bombardment not only by bombs and missiles but by the press and over the television provoked the haunting questions, “Was the downed pilot my husband?” “Was the navigator taken captive my son?”

In her classic poem “The Gate of the Year,” the poetess M. Louise Haskins summed up the feelings of all touched by the conflict and concerned for the safety of loved ones. She penned the comforting lines:

\begin{verse}
And I said to the man who stood at the gate of the year:\\
“Give me a light, that I may tread safely into the unknown!”\\
And he replied:\\
“Go out into the darkness and put your hand into the Hand of God.\\
“That shall be to you better than light and safer than a known way.”
\end{verse}

[In Masterpieces of Religious Verse, ed. James Dalton Morrison (New York: Harper and Brothers, 1948), p. 92]

At last the guns fell silent. Aircraft remained grounded. Mobile patrols halted. A quiet calm settled over the battlefield. The din of war succumbed to the silence of peace.

A scene on the cruel desert sands—and a sentence uttered from the heart—spoke volumes. An American soldier looked down at his vanquished enemy prisoner, touched the man’s shoulder, and reassured him with the words, “It’s all right; it’s all right.”

Every man and woman embroiled in that conflict thought of home, of family, and of friends. The embers of longing for loved ones glowed brightly and were found on every face. Love replaced hate, warmth filled every heart, and compassion overflowed every soul.

The words of King Arthur, from Lerner and Loewe’s long-running musical Camelot, left the stage and found deep meaning on a far distant desert: “Violence is not strength, and compassion is not weakness.”

The account of a homecoming as related by successful prison warden Kenyon J. Scudder brings to the surface tender feelings held in the heart:

A friend of his happened to be sitting in a railroad coach next to a young man who was obviously depressed. Finally the young man revealed that he was a paroled convict returning from a distant prison. His imprisonment had brought shame to his family, and they had neither visited him nor written often. He hoped, however, that this was only because they were too poor to travel and too uneducated to write. He hoped, despite the evidence, that they had forgiven him.

To make it easy for them, however, he had written to them asking that they put up a signal for him when the train passed their little farm on the outskirts of town. If his family had forgiven him, they were to put up a white ribbon in the big apple tree which stood near the tracks. If they didn’t want him to return, they were to do nothing, and he would remain on the train as it traveled onward.

As the train neared his hometown, the suspense became so great that he couldn’t bear to look out of his window. He exclaimed, “In just five minutes the engineer will sound the whistle indicating our approach to the long bend which opens into the valley I know as home. Will you watch for the apple tree at the side of the track?” His companion said he would; they exchanged places. The minutes seemed like hours, but then there came the shrill sound of the train whistle. The young man asked, “Can you see the tree? Is there a white ribbon?”

Came the reply, “I see the tree. I see not one white ribbon, but many. There is a white ribbon on every branch. Son, someone surely does love you.”

In that instant, all the bitterness that had poisoned a life was dispelled. “I felt as if I had witnessed a miracle,” the other man said. Indeed, he had witnessed a miracle. (See John Kord Lagemann, “Forgiveness: The Saving Grace,” The Reader’s Digest, Mar. 1961, pp. 41–42.)

Today a yellow ribbon has replaced one that is white. However, the message is the same: “Welcome home!” Men, women, and children everywhere are tying yellow ribbons around everything. Not only are they being tied around trees, but also around lampposts, street signs, and mailboxes—even about the necks of pets. So overwhelming is the demand for yellow ribbon material that busy suppliers working around the clock cannot meet the need. A classic yellow bow was one which completely girdled a large plane bringing soldiers safely home. I have surmised that each one who tenderly tied a yellow bow was singing, humming, or at least thinking of the words of the song “Tie a Yellow Ribbon ‘Round the Old Oak Tree.”

In the warm and poignant airport scene of a family awaiting a returning father and husband, smiles and tears of gratitude were everywhere to be found. My eye caught the expression of a small boy holding aloft a stick around which was tied a yellow ribbon. No words could describe the unspoken feeling. It is the welcome home of the heart that brings tears to every eye and peace to every soul.

Children have the capacity for compassion. They have no fear to express their genuine feelings. In the popular movie entitled Home Alone, a scene near the end grips the viewer’s emotions and causes that familiar lump to fill the throat. The scene takes place in a chapel; the time is Christmas; the two lonely characters are seated next to one another on a church bench. The older man, who lives by himself, is estranged from family and bereft of friends. His next-door neighbor, played by McCaulay Culkin, is the lad left “home alone” by his family, which had departed for a European vacation, inadvertently forgetting this one small family member.

The boy asks the lonely man if he has any family. The gentleman explains quietly that he and his son and his son’s family have parted ways and no longer communicate. In the innocence of youth, the boy blurts out the plea, “Why don’t you just call your son and tell him you are sorry and invite him home for Christmas!”

The old man sighs and responds, “I’m too afraid he would say no.” The fear of failure had blocked the ability to express love and to voice an apology.

The viewer is left to wonder concerning the outcome of the conversation, but not for too long. Christmas comes; the boy’s family returns. He is pictured at an upstairs bedroom window looking in the direction of the old man’s sidewalk. Suddenly he views a tender scene as the neighbor welcomes his returning son, his daughter-in-law, and their children. Son embraces father, and the old man buries his head against the shoulder of his precious son. As they turn to walk on, the old neighbor looks upward to the bedroom window of the house next door and sees his small friend observing the private miracle of forgiveness. Their eyes meet, their hands express a gentle greeting of gratitude. “Welcome home” replaces “Home alone.”

One emerges from the theater with moist eyes. As the brightness of day envelops the silent throng, perhaps there are those whose thoughts turn to that man of miracles, that teacher of truth—even the Lord of lords, Jesus Christ. I know my thoughts did.

I reflected on the Savior’s capacity for compassion. In Galilee “there came a leper to him, beseeching him, and kneeling down to him, and saying unto him, If thou wilt, thou canst make me clean.

“And Jesus, moved with compassion, put forth his hand, and touched him, and saith unto him, I will; be thou clean.

“And as soon as he had spoken, immediately the leprosy departed from him, and he was cleansed” (Mark 1:40–42).

On this, the American continent, Jesus appeared to a multitude and said:

“Have ye any that are sick among you? Bring them hither. Have ye any that are lame, or blind, or halt, or maimed, or leprous, or that are withered, or that are deaf, or that are afflicted in any manner? Bring them hither and I will heal them, for I have compassion upon you. …

“… And he did heal them every one. …

“And they did all, both they who had been healed and they who were whole, bow down at his feet, and did worship him; and as many as could come for the multitude did kiss his feet, insomuch that they did bathe his feet with their tears” (3 Ne. 17:7, 9–10).

Few accounts of the Master’s ministry touch me more than His example of compassion shown to the grieving widow at Nain:

“And it came to pass the day after, that he went into a city called Nain; and many of his disciples went with him, and much people.

“Now when he came nigh to the gate of the city, behold, there was a dead man carried out, the only son of his mother, and she was a widow: and much people of the city was with her.

“And when the Lord saw her, he had compassion on her, and said unto her, Weep not.

“And he came and touched the bier: and they that bare him stood still. And he said, Young man, I say unto thee, Arise.

“And he that was dead sat up, and began to speak. And he delivered him to his mother” (Luke 7:11–15).

What power, what tenderness, what compassion did our Master and Exemplar thus demonstrate! We, too, can bless if we will but follow His noble example. Opportunities are everywhere. Needed are eyes to see the pitiable plight and ears to hear the silent pleadings of a broken heart. Yes, and a soul filled with compassion, that we might communicate not only eye to eye or voice to ear but, in the majestic style of the Savior, even heart to heart.

Within walking distance from this tabernacle is a shelter for the homeless, a dental clinic, a soup kitchen. The compassion of this community is in evidence there each day. The Church and its members join with others not of our particular faith to bless the lives of those in need. A few streets beyond stands the regional bishops’ storehouse, stocked with commodities representing your generosity. No one leaves there without food or clothing or without gratitude to God.

Another place of refuge located nearby is Neighborhood House, a nondenominational care center where generous women share their time and their means to teach preschool children whose single mothers work to provide for their own. This organization also brings joy to the elderly who assemble there to exchange views and to listen to presentations and entertainment. These noble women bring the light of hope to the lives of the depressed, the downtrodden of society, and to children who will be the parents of tomorrow.

Without exception, those compassionate souls who feed the hungry, clothe the weary, and relieve the suffering of fellow beings exclaim, “I have never before felt more blessed, more rewarded, or so at peace.” A writer expressed the feeling:

\begin{verse}
I have wept in the night\\
For the shortness of sight\\
That to somebody’s need made me blind;\\
But I never have yet\\
Felt a tinge of regret\\
For being a little too kind.
\end{verse}

[Anonymous, quoted by Richard L. Evans, Improvement Era, May 1960, p. 340]

Similar projects are to be found in every community. The need beckons. We as a people need but to respond.

Recently two envelopes arrived at my office, sent by persons who preferred to remain anonymous. Each contained a number of one-hundred-dollar bills and a brief message expressing gratitude to God for His kind blessings and a desire that the money enclosed enable needy persons to receive their temple blessings. If these couples are viewing the conference, I am pleased to report that families in Bolivia and in Portugal will now be able to travel to temples in Lima, Peru, and Frankfurt, Germany, to fulfill this wish and achieve eternal blessings.

Perhaps these compassionate, anonymous donors would appreciate the thoughts of Henry Burton, who wrote the lines:

\begin{verse}
Have you had a kindness shown?\\
Pass it on.\\
’Twas not given for thee alone,\\
Pass it on.\\
Let it travel down the years,\\
Let it wipe another’s tears,\\
’Till in heav’n the deed appears—\\
Pass it on.
\end{verse}

[In Masterpieces of Religious Verse, pp. 389–90]

One Sunday morning in a nursing home in the valley, I witnessed the presentation of a beautiful gift as a young girl shared her musical talent with those lonely and elderly men and women who yearned not for food or for clothing but for someone who cared, someone who shared, and someone who provided a “hyacinth” for the soul.

A hush fell over the wheelchair-confined audience as the girl took bow in hand and played on her violin a beautiful melody. At the conclusion, one patient audibly declared, “My dear, that was lovely.” Then she began to clap her hands to express approval. A second patient joined in clapping, then a third, a fourth, and soon everyone applauded.

Together the young girl and I walked out of the nursing home. She said to me, “I have never played better. I have never felt better.” She had been guided by God and led by the Lord. Aches, pains, despair, and sadness had been conquered. Compassion had gained the victory.

Today, and in the tomorrows which lie ahead, we shall rejoice in the return to their homes and families of all who served in Desert Storm. They heard the call of duty. They fought the fight of the brave. They return victorious. To those who lost loved ones in Desert Storm or, for that matter, in any storm of deprivation, our heartfelt compassion goes out to you.

A story that moved across the wires revealed that a Methodist Sunday School teacher was the first U.S. soldier killed. One of the last was a soldier whose dad called her “Angel.” Of the 182 soldiers who died, there were those with cut-short honeymoons. Some left behind expectant wives. Some had put dreams on hold.

Now there is a widow in Virginia who has buried her only son, a young man in western Pennsylvania whose wedding plans have been permanently tucked away, a wife in Alaska soon due to deliver a baby her husband will never hold.

There is no satisfactory answer to the unspoken question, “Of the thousands and thousands of soldiers, why is mine among those not coming back?” Expressed is the lament, “A light from our household is gone; a voice we loved is stilled. A place is vacant in our hearts that never can be filled.” Lamenting the terrible sacrifice of any armed conflict, one writer penned the lines, “War leaves nothing but dead ends on the roads to all our fondest hopes and our brightest dreams” (Dennis Smith, “When Uncle Louis Came Home from Belgium,” Deseret News, 11 Jan. 1991, p. C1).

The Holy Bible furnishes a formula which eases the pain and heals the hearts of those who grieve:

“Trust in the Lord with all thine heart; and lean not unto thine own understanding.

“In all thy ways acknowledge him, and he shall direct thy paths” (Prov. 3:5–6).

To all who have loved and lost on either side of this tragic conflict, your grief can be assuaged. There is balm in Gilead. There awaits the promise of a new day. There echoes from a land not far from where your loved ones fell even a promise of peace, spoken by our Lord, the Prince of Peace:

“Peace I leave with you, my peace I give unto you: not as the world giveth, give I unto you. Let not your heart be troubled, neither let it be afraid” (John 14:27).

“In my Father’s house are many mansions: if it were not so, I would have told you. I go to prepare a place for you … that where I am, there ye may be also” (John 14:2–3).

His love, His promise, His presence is as a yellow ribbon, tied with care and marked with compassion. To your loved ones He has beckoned, “Welcome home.” To you He speaks the heavenly and divine assurance: “I am with you; you are never alone.”

“Weeping may endure for a night, but joy cometh in the morning” (Ps. 30:5).

To these words I add my witness: God lives, and His Son, Jesus Christ, is our Savior and Redeemer. Tonight my wife and I shall join millions of you as we kneel in solemn prayer and supplication. We shall acknowledge His holy hand in our lives. And from our hearts will come our expression of gratitude, “Thanks be to God.” In the name of Jesus Christ, amen.

\newpage
\settalk{The Power of Prayer}
\setspeaker{Thomas S. Monson}
\settalkdate{April 1991}
\talkheading{The Power of Prayer}{Thomas S. Monson}

My beloved brothers and sisters, it has been customary for the President of the Church to welcome you and deliver a message at the commencement of conference. With all his noble heart, President Benson would desire to stand at this pulpit and bear to you his witness concerning the truth of this work, the gratitude he feels for your prayers, and his fervent hope that all may so live as to merit and receive the abundant blessings a loving Heavenly Father desires to bestow.

Tears come easily to the eyes of our prophet when he receives letters written by children in which they express their greetings and send to him their love. President Benson is a family man who loves his children, grandchildren, and great-grandchildren and, indeed, children everywhere in this wonderful world in which we live.

President Benson has suggested that I begin this conference with a brief message in his behalf. He is pleased that the president of the United States has proclaimed that yesterday, today, and tomorrow be designated as days of national prayer and that sincere expressions of gratitude ascend to heaven for the end of the war in the Middle East. The First Presidency has commented: “We are thankful for the resolution of the war, and it is our fervent hope and prayer that all nations involved will work in concert for a lasting peace. The collective prayers of the nation and the world should focus not only on a lasting peace but also on the needs of the many on both sides who lost loved ones and endured suffering in the conflict.”

President Benson has stated: “The price of peace is righteousness. Men and nations may loudly proclaim, ‘Peace, peace,’ but there shall be no peace until individuals nurture in their souls those principles of personal purity, integrity, and character which foster the development of peace. Peace cannot be imposed. It must come from the lives and hearts of men. There is no other way” (The Teachings of Ezra Taft Benson [Salt Lake City: Bookcraft, 1988], p. 703).

President Benson has urged: “If we would advance in holiness—increase in favor with God—nothing can take the place of prayer. … Give prayer—daily prayer, secret prayer—a foremost place in your lives. Let no day pass without it. Communion with the Almighty has been a source of strength, inspiration, and enlightenment through the world’s history to men and women who have shaped the destinies of individuals and nations for good” (God, Family, Country: Our Three Great Loyalties [Salt Lake City: Deseret Book Co., 1974], p. 8).

In speaking to a large audience in São Paulo, Brazil, some time ago, President Benson testified:

“All through my life the counsel to depend on prayer has been prized above almost any other advice I have ever received. It has become an integral part of me, an anchor, a constant source of strength and the basis of my knowledge of things divine.

“Our Heavenly Father is always near. … Thank God we can reach out and tap that unseen power, without which no man can do his best” (address delivered to temple workers and Church employees at São Paulo, Brazil, 20 Nov. 1982).

President Benson has frequently quoted the words of a favorite hymn:

\begin{verse}
Prayer is the soul’s sincere desire,\\
Uttered or unexpressed,\\
The motion of a hidden fire\\
That trembles in the breast.
\end{verse}

[Hymns, 1985, no. 145]

The words of testimony spoken by President Benson are particularly appropriate on this special day of prayer and thanksgiving. He said: “I testify that there is a God in heaven who hears and answers prayers. I know this to be true, for He has answered mine. I would humbly urge all within the sound of my voice to keep in close touch with our Father in Heaven, through prayer” (São Paulo, Brazil, 20 Nov. 1982).

I heartily endorse this plea of our prophet and President, in the name of Jesus Christ, amen.

\newpage
\settalk{“Called to Serve”}
\setspeaker{Thomas S. Monson}
\settalkdate{October 1991}
\talkheading{“Called to Serve”}{Thomas S. Monson}

One cannot gaze into the faces of this vast congregation of men assembled in the historic Tabernacle and envision the unseen audiences meeting in other locations throughout the world without feeling your strength, recognizing your faith, and knowing of your spiritual power, even the power of the priesthood.

All of us are familiar with the beautiful account found in Matthew: “And Jesus, walking by the sea of Galilee, saw two brethren, Simon called Peter, and Andrew his brother, casting a net into the sea: for they were fishers.

“And he saith unto them, Follow me, and I will make you fishers of men.

“And they straightway left their nets, and followed him” (Matt. 4:18–20).

Brethren, we too have been called to follow Him as fishers of men, laborers in the vineyard to build boys and mend men and bring all unto Christ. We are stirred in our souls as we repeat the words of the well-known refrain:

\begin{verse}
Called to serve Him, heav’nly King of glory,\\
Chosen e’er to witness for his name,\\
Far and wide we tell the Father’s story,\\
Far and wide his love proclaim.
\end{verse}

[Hymns, 1985, no. 249]

It is no small thing to extend to another a call to serve; neither is it insignificant to receive such a call. President Spencer W. Kimball often taught, “Let there be no ‘ditch bank’ appointments in this Church.” Calls to serve are to be preceded by careful thought and earnest prayer. As the Lord declared, “Remember the worth of souls is great in the sight of God” (D\&C 18:10).

Some of you are called to serve the young men who hold the Aaronic Priesthood. These precious young men come in all sizes and with varying dispositions and different backgrounds. Yours is the privilege to know them individually and to motivate and lead each youth in his quest to qualify for the Melchizedek Priesthood, a successful mission, a temple marriage, a life of service, and a testimony of truth. Let us remember that a boy is the only known substance from which a man can be made.

\begin{verse}
Nobody knows what a boy is worth;\\
We’ll have to wait and see.\\
But every man in a noble place\\
A boy once used to be.
\end{verse}

A proper perspective of our young men is absolutely essential for those called to serve them. They are young, pliable, eager, and filled with unlimited energy. Sometimes they make mistakes. I remember a meeting where we of the First Presidency and the Twelve were reviewing a youthful mistake made by a missionary. The tone was serious and rather critical when Elder LeGrand Richards said: “Now, brethren, if the good Lord wanted to put a forty-year-old head on a nineteen-year-old body, He would have done so. But He didn’t. He placed a nineteen-year-old head on a nineteen-year-old body, and we should be a bit more understanding.” The mood of the group changed, the problem was solved, and we moved on with the meeting.

The years in the Aaronic Priesthood are growing years. They are years of maturing, learning, developing. They are years of emotional highs and lows, a period when wise counseling and proper example by an inspired leader can work wonders and lift lives.

The quorum meetings of the Aaronic Priesthood provide you advisers and members of bishoprics with ideal opportunities to teach and train these young men in gospel scholarship and in dedicated service. Be examples worthy of emulation. Youth need fewer critics and more models to follow. “Teach ye diligently,” said the Lord, “and my grace shall attend you” (D\&C 88:78).

These young men of the Aaronic Priesthood, many of whom are assembled here tonight, have a vital interest in athletics. The Church recognizes this fact and provides through its activities and athletic programs an opportunity for participation and growth. The enormous financial investment in physical facilities made by the Church, with the anticipation that all may benefit, can provide fellowship and brotherhood as well as the development of athletic skills. These goals, however, are defeated if winning the game overshadows participation in the game. Young men come to play—not to sit on the bench. Ours is the privilege to provide this opportunity.

I remember in my youth a basketball team from the Twenty-fifth Ward of the Pioneer Stake that had ten young men participating. A wise leader decided not to play just the five best, with the other five substituting here and there. Rather, he formed two teams with balanced ability and age. One team of five played the first and third periods, while the remaining team of five played the second and fourth periods. It was not a contest between bench warmers and active players, but a situation where morale was high, playing time was equal, and games were played and won in the right spirit. No participant in Church-sponsored athletic contests should warm the bench for the entire game.

Scouting is another area of vital interest to young men. Much has been said in the media of late regarding Scouting. Let me affirm that The Church of Jesus Christ of Latter-day Saints has not diminished in any way its support of the Scouting movement. President Spencer W. Kimball declared that the Church endorses Scouting “and will seek to provide leadership which will help boys keep close to their families and close to the Church as they develop the qualities of citizenship and character and fitness which Scouting represents. …

“We’ve remained strong and firm in our support of this great movement for boys and the Oath and the Law which are at its center” (in Conference Report, Apr. 1977, pp. 50–51; or Ensign, May 1977, p. 36).

President Ezra Taft Benson described Scouting as “a noble program,” saying, “It is a builder of character, not only in the boys, but also in the men who provide the leadership” (… So Shall Ye Reap [Salt Lake City: Deseret Book Co., 1960], p. 138).

Brethren, if ever there were a time when the principles of Scouting were vitally needed—that time is now. If ever there were a generation who would benefit by keeping physically strong, mentally awake, and morally straight—that generation is the present generation.

A few years ago a Scouting skill saved a life—in my own family. My nephew’s son, eleven-year-old Craig Dearden, successfully completed his requirements for Scouting’s swimming award. His father beamed his approval, while mother tenderly placed an affectionate kiss. Little did those attending the court of honor realize the life-or-death impact of that award. Later that very afternoon, it was Craig who spotted a dark object at the deep end of the swimming pool. It was Craig who, without fear, plunged into the pool to investigate and brought to the surface his own little brother. Tiny Scott was so still, so blue, so lifeless. Recalling the life-saving procedures he had learned and practiced, Craig and others responded in the true tradition of Scouting. Suddenly there was a cry, breathing, movement, life. Is Scouting relevant? Ask a mother, a father, a family who know a Scouting skill saved a son and brother.

Many of you are serving as members of bishoprics, of high councils, and as officers in priesthood quorums. At times, your tasks may seem overwhelming, and discouragement may creep into your lives. Our Heavenly Father has inspired your call and desires that you succeed. Through His beloved Son, our Savior, we learn: “Therefore, O ye that embark in the service of God, see that ye serve him with all your heart, might, mind and strength, that ye may stand blameless before God at the last day.

“Therefore, if ye have desires to serve God ye are called to the work” (D\&C 4:2–3).

In a revelation to the Prophet Joseph Smith, the Lord counseled: “Wherefore, be not weary in well-doing, for ye are laying the foundation of a great work. And out of small things proceedeth that which is great.

“Behold, the Lord requireth the heart and a willing mind” (D\&C 64:33–34).

Through humble prayer, diligent preparation, and faithful service, we can succeed in our sacred callings. Some priesthood bearers are gifted with the ability to reach out to the less active and renew the faith and rekindle the desire to once again return to the fold. Give such specially endowed brethren an assignment which will utilize this talent. Other brethren have the ability to work with youth, to win their respect, prompt their determination to overcome temptation, and lead with love these choice young spirits as they travel along that pathway which, when followed, provides eternal life. The Lord will hear your prayers and guide your decisions, for this is His work in which we are engaged.

I have frequently said that there is no feeling to surpass that feeling which engulfs us when we recognize that we have been on the Lord’s errand and He has allowed us to help fulfill His purposes.

Every bishop can testify to the promptings which attend calls to serve in the Church. Frequently the call seems to be for the benefit not so much of those to be taught or led as for the person who is to teach or lead.

As a bishop, I worried about any members who were inactive, not attending, not serving. Such was my thought one day as I drove down the street where Ben and Emily lived. They were older—even in the twilight period of life. Aches and pains of advancing years caused them to withdraw from activity to the shelter of their home—isolated, detached, shut out from the mainstream of daily life and association.

That day I felt the unmistakable prompting to park my car and visit Ben and Emily, even though I was on my way to a meeting. It was a sunny weekday afternoon. I approached the door to their home and knocked. Emily answered. When she recognized me, her bishop, she exclaimed: “All day long I have waited for my phone to ring. It has been silent. I hoped that the postman would deliver a letter. He brought only bills. Bishop, how did you know today is my birthday?”

I answered, “God knows, Emily, for He loves you.”

In the quiet of their living room, I said to Ben and Emily: “I don’t know why I was directed here today, but our Heavenly Father knows. Let’s kneel in prayer and ask Him why.” This we did, and the answer came. Emily, who had a beautiful voice, was asked to sing in the choir—even to provide a solo for the forthcoming ward conference. Ben was asked to speak to the Aaronic Priesthood young men and recount a special experience in his life when his safety was assured by responding to the promptings of the Spirit.

She sang. He spoke. Hearts were gladdened by the return to activity of Ben and Emily. They rarely missed a sacrament meeting from that day to the time each was called home. The language of the Spirit had been spoken. It had been heard. It had been understood. Hearts were touched and lives saved.

As priesthood leaders, we soon discover that some of our work, though not recorded on any written report, is of vital significance. The visits to the homes of quorum members, blessing the sick, helping a member with a project, or comforting grieving hearts when a loved one passes on are all sacred privileges of priesthood service. True, they may not be recorded on a written report, but more important, they find lodgment in the soul and bring joy to the heart. They are also known of the Lord.

Should our load seem heavy or the results of our efforts discouraging, we may well recall the words of President Kimball to some who noted his undeviating devotion to his calling even in his advancing years: “My life is like my shoes—to be worn out in service” (“He Is at Peace,” Ensign, Dec. 1985, p. 41).

I trust that all young men here tonight are preparing now to serve a full-time mission in the service of the Lord. ElRay L. Christiansen often said, “Your mission is the mold in which your life will be cast.” Prepare to serve worthily, with an eye single to the glory of God and His purposes. You will never know the full influence of your testimony and your service, but you will return with gladness for having had the privilege of responding to a sacred call to serve the Master. You will be forever loved by those to whom you bring the light of truth. Your teachings will be found in their service. Your examples will be guides to follow. Your faith will prompt courage to meet life’s challenges.

Let me provide an example. When I first visited Czechoslovakia, accompanied by Hans B. Ringger, long before the freedom bell sounded, I was met by Jiri Snederfler, our leader through this dark period, and Sister Olga Snederfler, his wife. I went to their home in Prague where the branch met. Displayed on the walls of the room in which we assembled were picture after picture of the Salt Lake Temple. I said to Sister Snederfler, “Your husband must truly love the temple.”

She responded, “I, too; I, too.”

We sat down for some soup Sister Snederfler had prepared, after which she brought out a treasure trove: an album containing individual pictures of the missionaries who were serving there in 1950 when the government edict came for the mission to be closed. As she slowly leafed through the pictures of different missionaries, she would say, “Wonderful boy, wonderful boy.”

Brother Snederfler has been a courageous Church leader in Czechoslovakia and has been willing to put everything on the line for the gospel. When the opportunity came that we would seek recognition for the Church in that country, the government leaders, then Communist, said: “Don’t send an American. Don’t send a German. Don’t send a Swiss. Send a citizen of Czechoslovakia.”

There were ominous implications in that particular statement because to have admitted that you were a church leader during this period of the prohibition of religion was tantamount to possible imprisonment. And yet, this call came to Brother Snederfler to be the designated person to go before the government and to forthrightly state that he was the leader of The Church of Jesus Christ of Latter-day Saints for all of Czechoslovakia and that he was seeking recognition for his church. He later told me that he had been somewhat frightened and had asked for the prayers of his brothers and sisters in the Prague Branch. He went to his sweet wife, Olga, and said to her: “I love you. I don’t know when, or if, I’ll be back, but I love the gospel, and I must follow my Savior.”

With that spirit of faith and devotion, Brother Snederfler went before the government officials and acknowledged that he was the leader of the Church and that he was there to seek a restoration of the recognition the Church had enjoyed long years before. In the meantime Elder Russell M. Nelson had been working tirelessly to bring about the desired decision. Later, Brother Snederfler heard the good news: “Your church is again recognized in Czechoslovakia.” How eager Brother Snederfler was to tell his dear wife and the other stalwart members of the Church the wonderful news that once again missionaries could come to Czechoslovakia and the Church could provide a haven for freedom of worship in that nation. It was a happy day for Czechoslovakia.

Where are Jiri and Olga Snederfler today? Just last month they responded to their calls to serve as temple president and matron of the Freiberg Germany Temple, where faithful members of the Church in Germany, Czechoslovakia, and surrounding nations attend. These two saintly souls find themselves each day in the Lord’s house they so dearly love.

And what of Richard Winder, one of the former missionaries of whom Olga Snederfler exclaimed, “Wonderful boy, wonderful boy”? He is now the president of the Czechoslovakia Prague Mission, responding to the call to him and his wife, Barbara, to reopen the work in that country.

To the Snederflers, the Winders, and to all who willingly respond to the sacred call of service comes the commendation of the Lord: “I, the Lord, am merciful and gracious unto those who fear me, and delight to honor those who serve me in righteousness and in truth unto the end.

“Great shall be their reward and eternal shall be their glory” (D\&C 76:5–6).

\begin{verse}
Called to know the richness of his blessing—\\
Sons and daughters, children of a King—\\
Glad of heart, his holy name confessing,\\
Praises unto him we bring.
\end{verse}

[Hymns, 1985, no. 249]

May we ever be found serving faithfully, I pray humbly in the name of Jesus Christ, amen.

\newpage
\settalk{Precious Children—A Gift from God}
\setspeaker{Thomas S. Monson}
\settalkdate{October 1991}
\talkheading{Precious Children—A Gift from God}{Thomas S. Monson}

From the book of Matthew we learn that after Jesus and His disciples descended from the Mount of Transfiguration, they paused at Galilee, then came to Capernaum. The disciples said unto Jesus, “Who is the greatest in the kingdom of heaven?

“And Jesus called a little child unto him, and set him in the midst of them,

“And said, Verily I say unto you, Except ye be converted, and become as little children, ye shall not enter into the kingdom of heaven.

“Whosoever therefore shall humble himself as this little child, the same is greatest in the kingdom of heaven.

“And whoso shall receive one such little child in my name receiveth me.

“But whoso shall offend one of these little ones which believe in me, it were better for him that a millstone were hanged about his neck, and that he were drowned in the depth of the sea” (Matt. 18:1–6).

I think it significant that Jesus loved these little ones who so recently had left the preexistence to come to earth. Children then and children now bless our lives, kindle our love, and prompt our good deeds.

Is it any wonder that the poet Wordsworth speaks thus of our birth: “Trailing clouds of glory do we come / From God, who is our home” (William Wordsworth, “Ode: Intimations of Immortality from Recollections of Early Childhood”).

Most of these little ones come to parents who eagerly await their arrival, mothers and fathers who rejoice to be a part of that miracle we call birth. No sacrifice is too great, no pain too severe, no waiting too long.

No wonder we are shocked when a wire story originating from a city in America informs that “a newborn girl who was wrapped in a paper bag and dumped in a garbage can is under close observation at a hospital. The child is doing well. ‘She’s a real beautiful, healthy baby,’ a hospital spokesman said Wednesday. Police said the infant was discovered after trash men emptied the garbage can into the back end of their dump truck and saw something move in the debris. Authorities are looking for the mother.”

It is our solemn duty, our precious privilege, even our sacred opportunity to welcome to our homes and to our hearts the children who grace our lives.

Our children have three classrooms of learning which are quite distinct one from another. I speak of the classroom at school, the classroom in church, and the classroom called home.

The Church has always had a vital interest in public education and encourages its members to participate in parent-teacher activities and other events designed to improve the education of our youth.

There is no more important aspect of public education than the teacher who has the opportunity to love, to teach, and to inspire eager boys and girls and young men and young women. President David O. McKay said: “Teaching is the noblest profession in the world. Upon the proper education of youth depend the permanency and purity of home, the safety and perpetuity of the nation. The parent gives the child an opportunity to live; the teacher enables the child to live well” (Gospel Ideals [Salt Lake City: Improvement Era, 1953], p. 436). I trust we shall recognize their importance and their vital mission by providing adequate facilities, the finest of books, and salaries which show our gratitude and our trust.

Each of us remembers with affection the teachers of our youth. I think it amusing that my elementary school music teacher was a Miss Sharp. She had the capacity to infuse within her pupils a love for music and taught us to identify musical instruments and their sounds. I well recall the influence of a Miss Ruth Crow who taught the subject of health. Though these were depression times, she ensured that each sixth-grade student had a dental health chart. She personally checked each pupil for dental health and made certain that through public or private resources, no child went without proper dental care. As Miss Burkhaus, who taught geography, rolled down the maps of the world and with her pointer marked the capital cities of nations and the distinctive features of each country, language, and culture, little did I anticipate or dream that one day I would visit these lands and peoples.

Oh, the importance in the lives of our children of teachers who lift their spirits, sharpen their intellects, and motivate their very lives!

The classroom at church adds a vital dimension to the education of every child and youth. In this setting each teacher can provide an upward reach to those who listen to her lessons and feel the influence of her testimony. In Primary, Sunday School, Young Women meetings and those of the Aaronic Priesthood, well-prepared teachers, called under the inspiration of the Lord, can touch each child, each youth, and prompt all to “seek … out of the best books words of wisdom; seek learning, even by study and also by faith” (D\&C 88:118). A word of encouragement here and a spiritual thought there can affect a precious life and leave an indelible imprint upon an immortal soul.

Many years ago, at a Church magazines awards banquet, we sat with President and Sister Harold B. Lee. President Lee said to our teenage daughter Ann: “The Lord has blessed you with a beautiful face and body. Keep the inside just as beautiful as the outside, and you will be blessed with true happiness.” This master teacher left with Ann an inspired guide to the celestial kingdom of our Heavenly Father.

The humble and inspired teacher in the church classroom can instill in her pupils a love for the scriptures. Why, the teacher can bring the Apostles of old and the Savior of the world not only into the classroom but also into the hearts, the minds, the souls of our children.

Perhaps most significant of all classrooms is the classroom of the home. It is in the home that we form our attitudes, our deeply held beliefs. It is in the home that hope is fostered or destroyed. Our homes are the laboratories of our lives. What we do there determines the course of our lives when we leave home. Dr. Stuart E. Rosenberg wrote in his book The Road to Confidence, “Despite all new inventions and modern designs, fads and fetishes, no one has yet invented, or will ever invent, a satisfying substitute for one’s own family.”

A happy home is but an earlier heaven. President George Albert Smith asked, “[Do] we want our homes to be happy[?] If we do, let them be the abiding place of prayer, thanksgiving and gratitude” (in Conference Report, Apr. 1944, p. 32).

There are those situations where children come to mortality with a physical or mental handicap. Try as we will, it is not possible to know why or how such events occur. I salute those parents who without complaint take such a child into their arms and into their lives and provide that added measure of sacrifice and love to one of Heavenly Father’s children.

This past summer at Aspen Grove Family Camp, I observed a mother patiently feeding a teenage daughter injured at birth and totally dependent upon Mother. Mother administered each spoonful of food, each swallow of water while holding steady the head and neck of her daughter. Silently I thought to myself, “For seventeen years, Mother has provided this service and all others to her daughter, never thinking of her own comfort, her own pleasure, her own food.” May God bless such mothers, such fathers, such children. And He will.

Parents everywhere realize that the most powerful combination of emotions in the world is not called out by any grand cosmic event, nor is it found in novels or history books, but merely by a parent gazing down upon a sleeping child.

When doing so, the truth of the words of Charles M. Dickinson come to mind:

\begin{verse}
They are idols of hearts and of households;\\
They are angels of God in disguise;\\
The sunlight still sleeps in their tresses,\\
His glory still gleams in their eyes;\\
These truants from home and from Heaven,\\
They have made me more manly and mild;\\
And I know now how Jesus could liken\\
The kingdom of God to a child.
\end{verse}

[“The Children”]

In our daily experiences with children, we discover they are most perceptive and often utter profound truths. Charles Dickens, the author of the classic A Christmas Carol, illustrated this fact when he described the humble Bob Cratchit family assembling for a rather meager but long-anticipated Christmas dinner. Bob, the father, was returning home with his frail son Tiny Tim upon his shoulder. Tiny Tim “bore a little crutch, and had his limbs supported by an iron frame.” Bob’s wife asked of him, “And how did little Tim behave?”

“‘As good as gold,’ said Bob, ‘and better. Somehow he gets thoughtful, sitting by himself so much, and thinks the strangest things you ever heard. He told me, coming home, that he hoped the people saw him in the church, because he was a cripple, and it might be pleasant to them to remember upon Christmas Day who made lame beggars walk, and blind men see’” (Christmas Carol and Cricket on the Hearth [New York: Grosset and Dunlop, n.d.], pp. 50–51).

Charles Dickens himself said, “I love these little people, and it is not a slight thing when they who are so fresh from God love us.”

Children express their love in original and innovative ways. On my birthday a few weeks ago, a precious little girl presented me with her handwritten birthday card and enclosed in the envelope a tiny toy padlock which she liked and thought I would enjoy receiving as a gift.

“Of all the dear sights in the world, nothing is so beautiful as a child when it is giving something. Any small thing it gives. A child gives the world to you. It opens the world to you as if it were a book you’d never been able to read. But when a gift must be found, it is always some absurd little thing, pasted on crooked, … an angel looking like a clown. A child has so little that it can give, because it never knows it has given you everything” (Margaret Lee Runbeck, Bits \& Pieces, 20 Sept. 1990).

Such was Jenny’s gift to me.

Children seem to be endowed with abiding faith in their Heavenly Father and His capacity and desire to answer their sweet prayers. It has been my personal experience that when a child prays, God listens.

Let me share with you the experience of Barry Bonnell and Dale Murphy, well-known professional baseball players formerly with the Atlanta Braves baseball club. Each is a convert to the Church, Dale Murphy having been baptized by Barry Bonnell.

“An experience occurred during the 1978 season that Barry described as ‘life changing.’ He was struggling terribly, batting about .200. Because of his poor performance, he was down on himself and felt miserable. He really didn’t want to go when Dale Murphy asked him to ‘come along to the hospital,’ but he went anyway. There he met Ricky Little, a stalwart [Atlanta] Braves’ supporter, but a youngster afflicted with leukemia. It was readily apparent that Ricky was near death. Barry felt a deep desire to think of something comforting to say but nothing seemed adequate. Finally, he asked if there was anything they could do. The youngster hesitated, and then asked if they would each hit a home run for him during the next game. Barry said [later], ‘That request wasn’t such a hard thing for Dale, who in fact hit two homers that night, but I was struggling at the plate and hadn’t hit a homer all year. Then I felt a warm feeling come over me and I told Ricky to count on it.’” That night, Barry hit his only home run of the season. (Jim Ison, Mormons in the Major Leagues [Cincinnati: Action Sports, 1991], p. 21.) A child’s prayer had been answered; a child’s wish had been fulfilled.

If only all children had loving parents, safe homes, and caring friends, what a wonderful world would be theirs. Unfortunately, not all children are so bounteously blessed. Some children witness their fathers savagely beating their mothers, while others are on the receiving end of such abuse. What cowardice, what depravity, what shame!

Local hospitals everywhere receive these little ones, bruised and battered, accompanied by bald-faced lies that the child “ran into the door” or “fell down the stairs.” Liars, bullies who abuse children, they will one day reap the whirlwind of their foul deeds. The quiet, the hurt, the offended child victim of abuse and at times incest must receive help.

A district judge, in a letter to me, declared: “Sexual abuse of children is one of the most depraved, destructive, and demoralizing crimes in civilized society. There is an alarming increase of reported physical, psychological, and sexual abuse of children. Our courts are becoming inundated with this repulsive behavior.”

The Church does not condone such heinous and vile conduct. Rather, we condemn in the harshest of terms such treatment of God’s precious children. Let the child be rescued, nurtured, loved, and healed. Let the offender be brought to justice, to accountability for his actions and receive professional treatment to curtail such wicked and devilish conduct. When you and I know of such conduct and fail to take action to eradicate it, we become part of the problem. We share part of the guilt. We experience part of the punishment.

I trust I have not spoken too harshly, but I love these little ones and know that the Lord loves them too. No more touching account of this love can be found than the experience of Jesus blessing the children as described in 3 Nephi. It tells of Jesus healing the sick, teaching the people, and praying to Heavenly Father for them. But then let me quote the precious words:

“[Jesus] took their little children, one by one, and blessed them, and prayed unto the Father for them.

“And when he had done this he wept again;

“And he spake unto the multitude, and said unto them: Behold your little ones.

“And as they looked to behold they cast their eyes towards heaven, and they saw the heavens open, and they saw angels descending out of the heaven as it were in the midst of fire; … and the angels did minister unto them” (3 Ne. 17:21–24).

You may ask, “Do such things occur even today?” Let me share with you the beautiful account of a grandmother and a grandfather now serving a mission and the manner in which their little grandson was blessed. The missionary grandfather wrote:

“My wife, Deanna, and I are now serving a mission in Jackson, Ohio. One of our big concerns as we accepted a mission call was our family. We would not be there when they had problems.

“Just before we went on our mission, our grandson R. J., who was two-and-a-half years old, had to have surgery to correct a crossed eye. His mother asked me to go with them because R. J. and I are real buddies. The operation went well, but R. J. did cry before and after the surgery because none of the family could go into the operating room, and he was afraid.

“About six months later, while we were still on our mission, R. J. needed the other eye corrected. His mother phoned and expressed her desire for me to be there to go with them for the second operation. Of course, distance and the mission prevented me from being with him. Deanna and I fasted and prayed for the Lord to comfort our grandson during his operation.

“We called shortly after the surgery was over and found that R. J. had remembered the previous experience and did not want to leave his parents. But as soon as he entered the operating room, he quieted down. He lay down on the operating table, took off his glasses for them, and went through the operation with a calm spirit. We were very thankful; our prayers had been answered.

“A couple of days later, we called our daughter and asked about R. J. He was doing fine, and she related this incident to us: In the afternoon after the operation, R. J. awakened and told his mother that Grandpa was there during the operation. He said, ‘Grandpa was there and made it all right.’ You see, the Lord made the anesthesiologist appear to that little boy as though he were his grandpa, but his grandpa and grandma were on a mission 1,800 miles away.”

Grandpa may not have been by your bedside, R. J., but you were in his prayers and in his thoughts. You were cradled in the hand of the Lord and blessed by the Father of us all.

My dear brothers and sisters, may the laughter of children gladden our hearts. May the faith of children soothe our souls. May the love of children prompt our deeds. “Children are an heritage of the Lord” (Ps. 127:3). May our Heavenly Father ever bless these sweet souls, these special friends of the Master, is my humble and earnest prayer. In the name of Jesus Christ, amen.

\newpage
\settalk{“The Lord Bless You”}
\setspeaker{Thomas S. Monson}
\settalkdate{October 1991}
\talkheading{“The Lord Bless You”}{Thomas S. Monson}

Traditionally the President of the Church provides closing remarks at the conclusion of conference. How we would enjoy hearing from President Benson! We are grateful that now, in his ninety-third year, he is free from pain, able to move about, meet on occasion with the First Presidency and Council of the Twelve, and be the recipient of your faith, your prayers, and your expressions of love for him. All of us were gratified that President Benson was able to attend a portion of the opening session on Saturday and on Sunday morning.

Since he is unable to address the conference prior to its adjournment, I have been asked to respond in his behalf. I seek the inspiration of the Lord as I convey to you the prophet’s love and counsel.

This has been a glorious conference. The Brethren have been inspired in their utterances, the prayers offered have been spoken from the heart, and the music truly has been “the song of the righteous” and “a prayer unto [the Lord]” (D\&C 25:12). Our sincere appreciation and gratitude go to each who has in any way made the conference most memorable.

We miss the association of Elder Derek A. Cuthbert, who was called to his heavenly home on April 7, 1991. We recall his testimony of truth concerning this work and marvel at all he was able to accomplish, even with limitations of health. In our prayers we remember his dear wife, Muriel, and each member of his family.

The conference speakers have emphasized the troubles of our times and the necessity to make certain our lives are lived in conformity with the principles of the gospel, that all of us may merit the companionship of the Lord to guide us on our earthly journey and qualify through our obedience for the blessings He desires to bestow upon us.

President Benson has frequently emphasized the importance of the family. He declared: “Remember, the family is one of God’s greatest fortresses against the evils of our day. Help keep your family strong and close and worthy of our Father in Heaven’s blessings. As you do, you will receive faith and strength which will bless your lives forever” (in Conference Report, Apr. 1986, p. 56; or Ensign, May 1986, p. 43).

“[Our] homes … need also the blessings which come from daily communion with God. … The differences and irritations of the day melt away as families approach the throne of heaven together. Unity increases. The ties of love and affection are re-enforced and the peace of heaven enters” (… So Shall Ye Reap [Salt Lake City: Deseret Book Co., 1960], p. 107).

President Benson has always emphasized the strength of youth and our responsibility to youth. He counseled: “One great thing the Lord requires of each of us is to provide a home where a happy, positive influence for good exists. In future years the costliness of home furnishings or the number of bathrooms will not matter much, but what will matter significantly is whether our children felt love and acceptance in the home. It will greatly matter whether there was happiness and laughter, or bickering and contention” (in Conference Report, Apr. 1981, p. 46; or Ensign, May 1981, p. 34).

How President Benson loves meeting and shaking hands with children and youth! He has traveled throughout the Church and has always taken great pleasure in singing to the children the song “A Mormon Boy.”

President Benson receives many letters from children. Sometimes they are humorous, other times tender. When President Benson was hospitalized and the doctors provided a pacemaker to help regulate his heart, one little girl wrote in and said, “Dear President Benson, I know you will be all right because the Bible says, ‘Blessed are the pacemakers.’”

He wept when I shared with him a letter I received from a child’s father. The letter began: “This past April, my wife and I were watching the Sunday afternoon session of conference. Our three-year-old son, Christopher, was standing on a chair at the kitchen counter playing with Play-Doh, listening to conference on the radio. As we entered the kitchen at the end of President Benson’s comments to the children, Christopher reported excitedly, ‘That man on the radio said that even when we make mistakes, our Heavenly Father still loves us.’ That simple statement has left a lasting and meaningful impression on our young son. I can still ask him today what President Benson said and receive the same enthusiastic reply. It is a comfort to him to know that he has a kind and loving Father in Heaven.”

This touching account is representative of the personal influence for good President Benson has ever been. He is gentle. He is kind. He is loving. He is your friend and my friend, and he knows the Lord, our Savior. I am certain I speak for him and for all the Brethren as this conference concludes: “The Lord bless [you], and keep [you]:

“The Lord make his face shine upon [you], and be gracious unto [you]:

“The Lord lift up his countenance upon [you], and give [you] peace” (Num. 6:24–26).

In the name of Jesus Christ, amen.

\newpage
\settalk{An Attitude of Gratitude}
\setspeaker{Thomas S. Monson}
\settalkdate{April 1992}
\talkheading{An Attitude of Gratitude}{Thomas S. Monson}

On this Sabbath day our thoughts turn to Him who atoned for our sins, who showed us the way to live and how to pray, and who demonstrated by His own actions the blessings of service. Born in a stable, cradled in a manger, this Son of God, even Jesus Christ the Lord, yet beckons to each of us to follow Him.

In the book of Luke, chapter 17, we read:

“And it came to pass, as he went to Jerusalem, that he passed through the midst of Samaria and Galilee.

“And as he entered into a certain village, there met him ten men that were lepers, which stood afar off:

“And they lifted up their voices, and said, Jesus, Master, have mercy on us.

“And when he saw them, he said unto them, Go shew yourselves unto the priests. And it came to pass, that, as they went, they were cleansed.

“And one of them, when he saw that he was healed, turned back, and with a loud voice glorified God,

“And fell down on his face at his feet, giving him thanks: and he was a Samaritan.

“And Jesus answering said, Were there not ten cleansed? but where are the nine?

“There are not found that returned to give glory to God, save this stranger.

“And he said unto him, Arise, go thy way: thy faith hath made thee whole” (Luke 17:11–19).

Through divine intervention, those who were lepers were spared from a cruel, lingering death and given a new lease on life. The expressed gratitude by one merited the Master’s blessing, the ingratitude shown by the nine, His disappointment.

Like the leprosy of yesteryear are the plagues of today. They linger; they debilitate; they destroy. They are to be found everywhere. Their pervasiveness knows no boundaries. We know them as selfishness, greed, indulgence, cruelty, and crime, to identify but a few. Surfeited with their poison, we tend to criticize, to complain, to blame, and, slowly but surely, to abandon the positives and adopt the negatives of life.

A popular refrain from the 1940s captured the thought:

\begin{verse}
Accentuate the positive;\\
Eliminate the negative.\\
Latch on to the affirmative;\\
Don’t mess with Mr. In-between.
\end{verse}

Good advice then. Good advice now.

This is a wonderful time to be living here on earth. Our opportunities are limitless. While there are some things wrong in the world today, there are many things right, such as teachers who teach, ministers who minister, marriages that make it, parents who sacrifice, and friends who help.

We can lift ourselves, and others as well, when we refuse to remain in the realm of negative thought and cultivate within our hearts an attitude of gratitude. If ingratitude be numbered among the serious sins, then gratitude takes its place among the noblest of virtues.

A favorite hymn always lifts our spirits, kindles our faith, and inspires our thoughts:

\begin{verse}
When upon life’s billows you are tempest tossed,\\
When you are discouraged, thinking all is lost,\\
Count your many blessings; name them one by one,\\
And it will surprise you what the Lord has done. …
\end{verse}

[Hymns, 1985, no. 241]

Well could we reflect upon our lives as individuals. We will soon discover much to prompt our personal gratitude.

First, there is gratitude for our mothers.

Mother, who willingly made that personal journey into the valley of the shadow of death to take us by the hand and introduce us to birth—even to mortal life—deserves our undying gratitude. One writer summed up our love for mother when he declared, “God could not be everywhere, and so He gave us mothers.”

While on the cruel cross of Calvary, suffering intense pain and anguish, Jesus “saw his mother, and the disciple standing by, whom he loved, he saith unto his mother, Woman, behold thy son! Then saith he to the disciple, Behold thy mother!” (John 19:26–27). What a divine example of gratitude and love!

My own mother may not have read to me from the scriptures; rather, she taught me by her life and actions what the “Good Book” contains. Care for the poor, the sick, the needy were everyday dramas never to be forgotten.

Second, let us reflect gratitude for our fathers.

Father, like Mother, is ever willing to sacrifice his own comfort for that of his children. Daily he toils to provide the necessities of life, never complaining, ever concerned for the well-being of his family. This love for children, this desire to see them well and happy, is a constant in a time of change.

On occasion I have observed parents shopping to clothe a son about to enter missionary service. The new suits are fitted, the new shoes are laced, and shirts, socks, and ties are bought in quantity. I met one father who said to me, “Brother Monson, I want you to meet my son.” Pride popped his buttons; the cost of the clothing emptied his wallet; love filled his heart. Tears filled my eyes when I noticed that his suit was old, his shoes well worn; but he felt no deprivation. The glow on his face was a memory to cherish.

As I reflect on my own father, I remember he yielded his minuscule discretionary time to caring for a crippled uncle, aged aunts, and his family. He served in the ward Sunday School presidency, always preferring to work with the children. He, like the Master, loved children. I never heard from his lips one word of criticism of another. He personified in his life the work ethic. I join you in an expression of gratitude for our fathers.

Third, all of us remember with gratitude our teachers.

The teacher not only shapes the expectations and ambitions of pupils; the teacher also influences their attitudes toward their future and themselves. If the teacher loves the students and has high expectations of them, their self-confidence will grow, their capabilities will develop, and their future will be assured. A citation to such a teacher could well read: “She created in her room an atmosphere where warmth and acceptance weave their magic spell; where growth and learning, the soaring of the imagination, and the spirit of the young are assured.”

May I express public gratitude for three of my own teachers. I thank G. Homer Durham, my history professor. He taught the truth, “The past is behind; learn from it.” He loved his subject; he loved his students. The love in his classroom opened the windows of my mind, that learning might enter.

O. Preston Robinson, my professor of marketing, instilled in his students that the future is ahead and we are to prepare for it. When he entered the classroom, his presence was like a welcome breath of fresh air. He instilled a spirit of “You can do it.” His life reflected his teaching—that of friendly persuasion. He taught truth. He inspired effort. He prompted love.

Then there was a Sunday School teacher—never to be forgotten, ever to be remembered. We met for the first time on a Sunday morning. She accompanied the Sunday School president into the classroom and was presented to us as a teacher who actually requested the opportunity to teach us. We learned that she had been a missionary and loved young people. Her name was Lucy Gertsch. She was beautiful, soft-spoken, and interested in us. She asked each class member to introduce himself or herself, and then she asked questions that gave her an understanding and an insight into the background of each boy, each girl. She told us of her childhood in Midway, Utah; and as she described that beautiful valley, she made its beauty live, and we desired to visit the green fields she loved so much. She never raised her voice. Somehow rudeness and boisterousness were incompatible with the beauty of her lessons. She taught us that the present is here and that we must live in it. She made the scriptures actually come to life. We became personally acquainted with Samuel, David, Jacob, Nephi, and the Lord Jesus Christ. Our gospel scholarship grew. Our deportment improved. Our love for Lucy Gertsch knew no bounds.

We undertook a project to save nickels and dimes for what was to be a gigantic party. Sister Gertsch kept a careful record of our progress. As boys and girls with typical appetites, we converted in our minds the monetary totals to cakes, cookies, pies, and ice cream. This was to be a glorious occasion—the biggest party ever. Never before had any of our teachers even suggested a social event like this one was going to be.

The summer months faded into autumn; autumn turned to winter. Our party goal had been achieved. The class had grown. A good spirit prevailed.

None of us will forget that gray morning in January when our beloved teacher announced to us that the mother of one of our classmates had passed away. We thought of our own mothers and how much they meant to us. We felt sorrow for Billy Devenport in his great loss.

The lesson that Sunday was from the book of Acts, chapter 20, verse 35: “Remember the words of the Lord Jesus, how he said, It is more blessed to give than to receive.” At the conclusion of the presentation of a well-prepared lesson, Lucy Gertsch commented on the economic situation of Billy’s family. These were depression times; money was scarce. With a twinkle in her eyes, she asked, “How would you like to follow this teaching of the Lord? How would you feel about taking your party fund and, as a class, giving it to the Devenports as an expression of our love?” The decision was unanimous. We counted very carefully each penny and placed the total sum in a large envelope.

Ever shall I remember the tiny band walking those three city blocks, entering Billy’s home, greeting him, his brother, sisters, and father. Noticeably absent was his mother. Always I shall treasure the tears which glistened in the eyes of each one present as the white envelope containing our precious party fund passed from the delicate hand of our teacher to the needy hand of a grief-stricken father. We fairly skipped our way back to the chapel. Our hearts were lighter than they had ever been, our joy more full, our understanding more profound. This simple act of kindness welded us together as one. We learned through our own experience that indeed it is more blessed to give than to receive.

The years have flown. The old chapel is gone, a victim of industrialization. The boys and girls who learned, who laughed, who grew under the direction of that inspired teacher of truth have never forgotten her love or her lessons.

Even today when we sing that old favorite—

\begin{verse}
Thanks for the Sabbath School. Hail to the day\\
When evil and error are fleeing away.\\
Thanks for our teachers who labor with care\\
That we in the light of the gospel may share.
\end{verse}

[Hymns, 1985, no. 278]

—we think of Lucy Gertsch, our Sunday School teacher, for we loved Lucy, and Lucy loved us.

Let us ever have an attitude of gratitude for our teachers.

Fourth, let us have gratitude for our friends. Our most cherished friend is our partner in marriage. This old world would be so much better off today if kindness and deference were daily a reflection of our gratitude for wife, for husband.

The Lord spoke the word friend almost with a reverence. He said, “Ye are my friends, if ye do whatsoever I command you” (John 15:14).

True friends put up with our idiosyncrasies. They have a profound influence in our lives.

Oscar Benson, a Scouter of renown, had a hobby of interviewing men on death row in various prisons throughout the country. He once reported that 125 of these men had said they had never known a decent man.

In the depths of World War II, I experienced an expression of true friendship. Jack Hepworth and I were teenagers. We had grown up in the same neighborhood. One afternoon I saw Jack running down the sidewalk toward me. When we met, I saw that there were tears in his eyes. In a voice choked with emotion, he blurted out the words, “Tom, my brother Joe, who is in the Navy Air Corps, has been killed in a fiery plane crash!” We embraced. We wept. We sorrowed. I felt highly complimented that instinctively Jack, my friend, felt the urgency to share with me his grief. We can all be grateful for such friends.

Fifth, may we acknowledge gratitude for our country—the land of our birth.

When we ponder that vast throng who have died honorably defending home and hearth, we contemplate those immortal words, “Greater love hath no man than this, that a man lay down his life for his friends” (John 15:13). The feelings of heartfelt gratitude for the supreme sacrifice made by so many cannot be confined to a Memorial Day, a military parade, or a decorated grave.

At the famed Theatre Royal, situated on Drury Lane in London, England, is a beautifully framed plaque containing words which touch my very soul and prompt feelings of deep gratitude:

Sixth and finally—even supremely—let us reflect gratitude for our Lord and Savior, Jesus Christ. His glorious gospel provides answers to life’s greatest questions: Where did we come from? Why are we here? Where does my spirit go when I die? His called missionaries bring to those who live in darkness the light of divine truth:

\begin{verse}
Go, ye messengers of glory;\\
Run, ye legates of the skies.\\
Go and tell the pleasing story\\
That a glorious angel flies,\\
Great and mighty, great and mighty,\\
With a message from the skies.
\end{verse}

[Hymns, 1985, no. 262]

He taught us how to pray. He taught us how to live. He taught us how to die. His life is a legacy of love. The sick He healed; the downtrodden He lifted; the sinner He saved.

Only He stood alone. Some Apostles doubted; one betrayed Him. The Roman soldiers pierced His side. The angry mob took His life. There yet rings from Golgotha’s hill His compassionate words, “Father, forgive them; for they know not what they do” (Luke 23:34).

Earlier, perhaps perceiving the culmination of His earthly mission, He spoke the lament, “Foxes have holes, and birds of the air have nests; but the Son of man hath not where to lay his head” (Luke 9:58). “No room in the inn” was not a singular expression of rejection—just the first. Yet He invites you and me to host Him. “Behold, I stand at the door, and knock: if any man hear my voice, and open the door, I will come in to him, and will sup with him, and he with me” (Rev. 3:20).

Who was this Man of sorrows, acquainted with grief? Who is this King of glory, this Lord of hosts? He is our Master. He is our Savior. He is the Son of God. He is the author of our salvation. He beckons, “Follow me” (Matt. 4:19). He instructs, “Go, and do thou likewise” (Luke 10:37). He pleads, “Keep my commandments” (John 14:15).

Let us follow Him. Let us emulate His example. Let us obey His word. By so doing, we give to Him the divine gift of gratitude.

My sincere prayer is that we may, in our individual lives, reflect that marvelous virtue: an attitude of gratitude. In the name of Jesus Christ, amen.

\newpage
\settalk{Memories of Yesterday, Counsel for Today}
\setspeaker{Thomas S. Monson}
\settalkdate{April 1992}
\talkheading{Memories of Yesterday, Counsel for Today}{Thomas S. Monson}

How our beloved prophet and President, Ezra Taft Benson, would enjoy standing at this pulpit to open a glorious conference of the Church. President Benson, we love you; we pray for you; we are anxious to follow your inspired direction.

This morning I pray for heavenly help as I respond to President Benson’s assignment to speak in his behalf. I shall attempt to express his thoughts and counsel, largely in his own words.

This year we commemorate the 150th anniversary of the founding of the Relief Society. Women of the Church rejoice as they reflect on past achievements of their organization and, with foresight coupled with faith, meet today’s challenges and plan for future accomplishments.

President Benson has singled out two members of the Relief Society for his personal tribute. He said: “I pay grateful tribute to two elect women who have influenced my life—my mother, and my own sweetheart and eternal companion. I thank God that they have used their womanly attributes of compassion and charity to bless my life and the lives of all their posterity.”

Reminiscing of boyhood days, President Benson recalls:

“Mother was Relief Society president in the ward, a small but solid country ward. I remember how important Father considered her work in that assignment.

“Father gave to me, as the oldest [child], the responsibility of harnessing the horse and getting the buggy ready for Mother’s … weekly Relief Society meetings. …

“At that time I was not tall enough to buckle the collar or put the bridle on the horse without getting on the fence or [standing on] a box.

“In addition, I was to take half a bushel of wheat from our granary and put it in the back of the buggy. In those days the Relief Society sisters were building up a storage of wheat against a time of need. …

“When Mother was called to visit the sick in the ward or to help mothers with new babies, it was always by horse and buggy. As the buggy rolled down the dirt road, the circling wheels left a track that stayed even after the buggy had disappeared. Mother’s influence has also stayed—in my life and in the countless lives she blessed through compassionate service and example.”

I find it interesting that Ezra Benson, the boy who assisted his mother and the Relief Society gather and store wheat for a future day of hunger, was Ezra Benson the Apostle, who years later directed a massive distribution of wheat and other essentials to the famished of Europe following World War II.

Of his companion, Flora, President Benson has said: “I honor and acknowledge my precious wife. … Her loving devotion, inspiration, faith and loyal support have contributed to whatever success may be ours.”

Thinking of the example of his own mother and that of his beloved and faithful wife, Flora, President Benson has offered ten specific suggestions for mothers as they guide their precious children:

Though President Benson has addressed these suggestions primarily to mothers, I am confident he would expect those of us who are men and fathers bearing the holy priesthood to do our part, along with each son and daughter, to implement them and bring to fruition their divine objectives.

President Benson leaves us this counsel:

Brothers and sisters, “make it a family objective to all be together in the celestial kingdom. Strive to make your home a little bit of heaven on earth so that after this life is over, you may be able to say:

\begin{verse}
We are all here!\\
Father, mother, sister, brother,\\
All who hold each other dear.\\
Each chair is filled—\\
We’re all at home. …\\
We’re all—all here.
\end{verse}

God bless you, President Benson, in the name of Jesus Christ, amen.

\newpage
\settalk{The Spirit of Relief Society}
\setspeaker{Thomas S. Monson}
\settalkdate{April 1992}
\talkheading{The Spirit of Relief Society}{Thomas S. Monson}

Today our souls have reached toward heaven. We have been blessed with beautiful music and inspired messages. The Spirit of the Lord is here.

I bring to you noble sisters of the Relief Society the greetings of President Ezra Taft Benson, who by special wire is viewing these proceedings at his apartment; President Gordon B. Hinckley, who is on assignment abroad; and all the General Authorities of the Church. We commend you. We pray for you. We are proud of you.

President Elaine Jack, Chieko Okazaki, Aileen Clyde—thank heaven for your dear mothers, your teachers, your youth leaders who recognized in you your potential.

To paraphrase a thought:

\begin{verse}
You never know what a girl is worth,\\
You’ll just have to wait and see;\\
But every woman in a noble place,\\
A girl once used to be.
\end{verse}

Years ago I saw a photograph of a Sunday School class in the Sixth Ward of the Pioneer Stake in Salt Lake City. The photograph was taken in 1905. A sweet girl, her hair in pigtails, was shown on the front row. Her name was Belle Smith. Later, as Belle Smith Spafford, general president of the Relief Society, she wrote: “Never have women had greater influence than in today’s world. Never have the doors of opportunity opened wider for them. This is an inviting, exciting, challenging and demanding period of time for women. It is a time rich in rewards if we keep our balance, learn the true values of life, and wisely determine priorities.”

The Apostle Paul gave us this caution: “The letter [of the law] killeth, but the spirit giveth life.” (2 Cor. 3:6.) The spirit of Relief Society is being made manifest today, in our time. We see stirrings of strength, we hear the rustling of a resurrection, we observe the dawning of a new day.

In the Church News, Sister Irene Maximova, Relief Society president in the St. Petersburg (Russia) Ward, reported some changes she sees in the lives of women after they join the Church: “They have more compassion for other people. I see increased consideration and respect. They are more occupied with scriptures and spiritual matters. … As Church members in Russia, we must always remember the Lord’s commandments to love God and to love our neighbors. … For 70 years our society lost those good qualities.”

In that same issue of the Church News was the dramatic announcement that three new missions would soon be opened in what was the Soviet Union. This has now been accomplished. Branches of the Church will be organized, the waters of baptism will welcome those who are prepared, Relief Society membership will soar, and souls will be saved.

In this, your sesquicentennial year, I compliment you on your carefully chosen theme to eliminate illiteracy. Those of us who can read and write do not appreciate the deprivation of those who cannot read, who cannot write. They are shrouded by a dark cloud which stifles their progress, dulls their intellect, and dims their hopes. Sisters of the Relief Society, you can lift this cloud of despair and welcome heaven’s divine light as it shines upon your sisters.

Several months ago I was in Monroe, Louisiana, attending a regional conference. It was a beautiful occasion. At the airport on my way home, I was approached by a lovely black member of the Church who said, smiling broadly, “President Monson, before I joined the Church and became a member of Relief Society, I could not read nor write. None of my family could. You see, we were all poor sharecroppers. President, my white Relief Society sisters—they taught me to read. They taught me to write. Now I help teach my white sisters how to read and how to write.” I reflected on the supreme joy she must have felt when she opened her Bible and read for the first time the words of the Lord:

“Come unto me, all ye that labour and are heavy laden, and I will give you rest.

“Take my yoke upon you, and learn of me; for I am meek and lowly in heart: and ye shall find rest unto your souls.

“For my yoke is easy, and my burden is light.” (Matt. 11:28–30.)

That day in Monroe, Louisiana, I received a confirmation by the Spirit of your exalted objective.

In planning the women’s curriculum, these guidelines have been followed with resolute care:

In keeping with this statement, may I today issue to you sisters of the Relief Society four challenges for our times:

First: Share your talents.

Second: Sustain your husband.

Third: Strengthen your home.

Fourth: Serve your God.

Share your talents. Each of you, single or married, regardless of age, has the opportunity to learn and grow. Expand your knowledge, both intellectual and spiritual, to the full stature of your divine potential. There is no limit to your influence for good. Share your talents, for that which we willingly share, we keep. But that which we selfishly keep, we lose.

Sustain your husband. Both husband and wife should appreciate that “woman was taken out of man; not out of his feet to be trampled underfoot, but out of his side to be equal to him, under his arm to be protected, and near his heart to be loved” [see Matthew Henry’s Commentary on the Whole Bible (1991), 1:16]. Be patient, be tender, be loving, be considerate, be understanding. Be your best self as you sustain your husband, remembering that children often outgrow their need for affection, but husbands never do.

Many members of Relief Society do not have husbands. Death, divorce, and indeed lack of opportunity to marry have, in many instances, made it necessary for a woman to stand alone. In reality, she need not stand alone, for a loving Heavenly Father will be by her side to give direction to her life and provide peace and assurance in those quiet moments when loneliness is found and when compassion is needed.

Strengthen your home. Home, that marvelous place, was meant to be a haven called heaven where the Spirit of the Lord might dwell.

Too frequently, women underestimate their influence for good. Well could you follow the formula given by the Lord: “Establish a house, even a house of prayer, a house of fasting, a house of faith, a house of learning, a house of glory, a house of order, a house of God.” (D\&C 88:119.)

In such a house will be found happy, smiling children who have been taught, by precept and example, the truth. In a Latter-day Saint home, children are not simply tolerated, but welcomed; not commanded, but encouraged; not driven, but guided; not neglected, but loved.

Serve your God. You cannot serve your neighbor without demonstrating your love for God. Service is a product of love. So long as we love, we serve. As James Russell Lowell stated so beautifully in his classic poem, The Vision of Sir Launfal, “Not what we give, but what we share. For the gift without the giver is bare.” “All the beautiful sentiments in the world weigh less than a single lovely action.”

\begin{verse}
Go gladden the lonely, the dreary;\\
Go comfort the weeping, the weary;\\
Go scatter kind deeds on your way.\\
Oh, make the world brighter today.
\end{verse}

The heart of compassionate service, one of the hallmark creeds of Relief Society, is the gift of oneself. Emerson explained, “Rings and … jewels are not gifts, but apologies for gifts. The only [true] gift is a portion of thyself.”

Sisters, will you accept these four challenges to (1) share your talents; (2) sustain your husband; (3) strengthen your home; and (4) serve your God. As you do, the blessings of heaven will attend.

Perhaps I could illustrate. A number of years ago I received a rather unique and frightening assignment. Folkman D. Brown, then our Director of Mormon Relationships for the Boy Scouts of America, came to my office, having learned that I was about to depart for a lengthy assignment visiting the missions of New Zealand. He told me of his sister, Belva Jones, who had been stricken with terminal cancer and who knew not how to “break the sad news” to her only son—a missionary in far-off New Zealand. Her wish, even her plea, was that he remain in the mission field and serve faithfully. She worried about his reaction, for the missionary, Elder Ryan Jones, had lost his father just a year earlier to the same dread disease.

I accepted the responsibility to inform Elder Jones of his mother’s illness and to convey to him her wish that he remain in New Zealand until his service there was completed. After a missionary meeting held adjacent to the majestically beautiful New Zealand Temple, I met privately with Elder Jones and, as gently as I could, explained the situation of his mother. Naturally, there were tears—not all his—but then the handclasp of assurance and the pledge: “Tell my mother I shall serve, I shall pray, and I shall see her again.”

I returned to Salt Lake City just in time to attend a conference of the Lost River Stake in Idaho. As I sat on the stand with the stake president, Burns Beal, my attention was drawn to the east side of the chapel, where the morning sunlight seemed to bathe an occupant of a front bench. President Beal introduced the woman as Belva Jones and said, “She has a missionary son in New Zealand. She is very ill and has requested a blessing.”

Prior to that moment, I had not known where Belva Jones lived. My assignment that weekend could have been to any of many stakes. Yet the Lord, in His own way, had answered the prayer of faith of a devoted Relief Society member. Following the meeting, we had a most delightful visit together. I reported, word for word, the reaction and resolve of her son Ryan. A blessing was provided. A prayer was offered. A witness was received that Belva Jones would live to see Ryan again.

This privilege she enjoyed. Just one month prior to her passing, Ryan returned, having successfully completed his mission.

I never think of the Lost River Stake but what I see again in my memory that modest sister made beautiful by her faith. Our Father had used the radiance of His sunlight to make known His purpose. I shall not forget Belva Jones. Here was one who shared her talents freely. Here was one who sustained her husband—and then her son—in their priesthood callings. Here was one who strengthened her home, even in the absence of a husband and father. Here was one who continued to serve her God and all others. Here was one who exemplified the spirit of Relief Society.

Dear sisters of Relief Society, move with vision, fueled by faith, into your next 150 years. To all of you I repeat that old, but ever welcome wish: Happy 150th birthday!

May “the Lord bless [you], and keep [you]: The Lord make his face shine upon [you], and be gracious unto [you]: The Lord lift up his countenance upon [you], and give [you] peace.” (Num. 6:24–26.)

In the name of the Prince of Peace, Jesus Christ the Lord, amen.

\newpage
\settalk{To Learn, To Do, To Be}
\setspeaker{Thomas S. Monson}
\settalkdate{April 1992}
\talkheading{To Learn, To Do, To Be}{Thomas S. Monson}

Truly a royal priesthood has assembled tonight. The Tabernacle on Temple Square is filled to overflowing, and the Assembly Hall is occupied, as are chapels throughout many countries in the world. In all likelihood this is the largest assemblage of priesthood holders ever to come together. Your devotion to your sacred callings is inspiring. Your desire to learn your duty is evident. The purity of your souls brings heaven closer to you and your families.

These are difficult economic times. Cutbacks in industry, layoffs on a substantial scale, and the resultant dislocation of families become a serious challenge. We must make certain that those for whom we share responsibilities do not go hungry or unclothed or unsheltered. When the priesthood of this Church work together as one in meeting these vexing conditions, near miracles take place.

We urge all Latter-day Saints to be prudent in their planning, to be conservative in their living, and to avoid excessive or unnecessary debt. The financial affairs of the Church are being managed in this manner, for we are aware that your tithing and other contributions have not come without sacrifice and are sacred funds.

Let us make of our homes sanctuaries of righteousness, places of prayer, and abodes of love, that we might merit the blessings that can come only from our Heavenly Father. We need His guidance in our daily lives.

In this vast throng is priesthood power and the capacity to reach out and share the glorious gospel with others. We have the hands to lift others from complacency and inactivity. We have the hearts to serve faithfully in our priesthood callings and thereby inspire others to walk on higher ground and to avoid the swamps of sin which threaten to engulf so many. The worth of souls is indeed great in the sight of God. Ours is the precious privilege, armed with this knowledge, to make a difference in the lives of others. The words found in Ezekiel could well pertain to all of us who follow the Savior in this sacred work:

“A new heart also will I give you, and a new spirit will I put within you. …

“And I will put my spirit within you, and cause you to walk in my statutes, and ye shall keep my judgments, and do them.

“And ye shall dwell in the land that I gave to your fathers; and ye shall be my people, and I will be your God” (Ezek. 36:26–28).

How might we merit this promise? What will qualify us to receive this blessing? Is there a guide to follow? May I suggest three imperatives for our consideration? They apply to the deacon as well as the high priest. They are within our reach. A kind Heavenly Father will help us in our quest.

First: Learn what we should learn!

Second: Do what we should do!

Third: Be what we should be!

Let us in some detail discuss these objectives, that we might be profitable servants in the sight of our Lord.

Learn what we should learn. Do what we should do. Be what we should be. By so doing, the blessings of heaven will attend. We will know that we do not serve alone. He who notes the sparrow’s fall will, in His own way, acknowledge our service.

Let me share with you, brethren, a touching experience that illustrates this assurance.

Brother Edwin Q. Cannon, Jr., was a missionary to Germany in 1938, where he loved the people and served faithfully. At the conclusion of his mission, he returned home to Salt Lake City. He married and commenced his own business.

Forty years passed by. One day Brother Cannon came to my office and said he had been pruning his missionary slides. Among those slides he had kept since his mission were several which he could not specifically identify. Every time he had planned to discard those few slides, he had been impressed to keep them, although he was at a loss as to why. They were photographs taken by Brother Cannon during his mission when he served in Stettin, Germany, and were of a family—a mother, a father, a small girl, a small boy. Brother Cannon knew their surname was Berndt but could remember nothing more about them. He indicated that he understood there was a Berndt who was a Regional Representative in Germany, and he thought, although the possibility was remote, that this Berndt might have some connection with the Berndts who had lived in Stettin and who were depicted in the photographs. Before disposing of the slides, he thought he would check with me.

I told Brother Cannon I was leaving shortly for Berlin, where I anticipated that I would see Dieter Berndt, the Regional Representative, and that I would show the slides to him to see if there were any relationship and if he wanted them. There was a possibility I would also see Brother Berndt’s sister, who was married to Dietmar Matern, a stake president in Hamburg.

The Lord didn’t even let me get to Berlin before His purposes were accomplished. I was in Zurich, Switzerland, boarding the flight to Berlin, when who should also board the plane but Dieter Berndt. He sat next to me, and I told him I had some old slides of people named Berndt from Stettin. I handed them to him and asked if he could identify those shown in the photographs. As he looked at them carefully he began to weep. He said, “Our family lived in Stettin during the war. My father was killed when an Allied bomb struck the plant where he worked. Not long afterward, the Russians invaded Poland and the area of Stettin. My mother took my sister and me and fled from the advancing enemy. Everything had to be left behind, including any photographs we had. Brother Monson, I am the little boy pictured in these slides, and my sister is the little girl. The man and the woman are our dear parents. Until today, I have had no photographs of our childhood in Stettin or of my father.”

Wiping away my own tears, I told Brother Berndt the slides were his. He placed them carefully and lovingly in his briefcase.

At the next general conference, when Dieter Berndt, Regional Representative, visited Salt Lake City, he paid a visit to Brother and Sister Edwin Cannon, Jr., that he might express in person his own gratitude for the inspiration that came to Brother Cannon to retain these precious slides and that he followed that inspiration in keeping them for forty years.

William Cowper penned the lines:

\begin{verse}
God moves in a mysterious way\\
His wonders to perform;\\
He plants his footsteps in the sea\\
And rides upon the storm. …
\end{verse}

[Hymns (1948), no. 48]

I leave with you my testimony that this work in which we are engaged is true. The Lord is at the helm. May we ever follow Him is my sincere prayer, in the name of Jesus Christ, amen.

\newpage
\settalk{“An Example of the Believers”}
\setspeaker{Thomas S. Monson}
\settalkdate{October 1992}
\talkheading{“An Example of the Believers”}{Thomas S. Monson}

This has been a beautiful and rewarding meeting. I endorse the counsel provided by President Howard W. Hunter and that given by each of the sisters who has addressed us. As I contemplate the vast audience assembled tonight, I ponder the words of President Heber J. Grant, who declared: “I have often felt that a photograph of our dear sisters, with the intelligent, Godlike faces they possess, would be a testimony to all the world of the integrity of our people.”

We would certainly need the widest wide-angle lens to include all of you in one photograph. Such is not available to us, but with God, all things are possible. In His infinite vision, He literally can view all of us and bless all of us. All we need do is to so live that we merit the blessings ever predicated on our faithfulness to His commandments.

Said President George Albert Smith: “I desire to impress on you daughters of God … that if this world is to endure, you must keep the faith. If this world is to be happy, you will have to set the pace for that happiness. … If we are to maintain our physical strength and mental power and spiritual joy, it will have to be on the Lord’s terms.” (Relief Society Magazine, Dec. 1945, p. 719.) Perhaps a young lady had this thought in mind when she spoke the feelings of her yearning heart: “What we really and truly need is less criticism and more models to follow.”

Frequently we are too quick to criticize, too prone to judge, and too ready to abandon an opportunity to help, to lift, and, yes, even to save. Some point the accusing finger at the wayward or unfortunate and in derision say, “Oh, she will never change. She has always been a bad one.” A few see beyond the outward appearance and recognize the true worth of a human soul. When they do, miracles occur. The downtrodden, the discouraged, the helpless become “no more strangers and foreigners, but fellowcitizens with the saints, and of the household of God.” (Eph. 2:19.) True love can alter human lives and change human nature.

This truth was portrayed so beautifully on the stage in My Fair Lady. Eliza Doolittle, the flower girl, spoke to one for whom she cared: “You see, really and truly, apart from the things anyone can pick up—the dressing and the proper way of speaking—the difference between a lady and a flower girl is not how she behaves, but how she’s treated. I shall always be a flower girl to Professor Higgins because he always treats me as a flower girl and always will; but I know I can be a lady to you because you always treat me as a lady and always will.”

The Apostle Paul wrote an epistle to his beloved companion Timothy in which he provided inspired counsel equally as applicable to you and me today as it was to Timothy. Listen carefully to his words: “Neglect not the gift that is in thee,” “but be thou an example of the believers, in word, in conversation, in charity, in spirit, in faith, in purity.” (1 Tim. 4:14, 12.)

We need not wait for a cataclysmic event, a dramatic occurrence in the world in which we live, or a special invitation to be an example—even a model to follow. Our opportunities lie before us here and now. But they are perishable. Likely they will be found in our own homes and in the everyday actions of our lives. Our Lord and Master marked the way: “[He] went about doing good.” (Acts 10:38.) He in very deed was a model to follow—even an example of the believers.

Are we?

Happiness abounds when there is genuine respect one for another. Wives draw closer to their husbands, and husbands are more appreciative of their wives, and children are happy, as children are meant to be. Where there is respect in the home, children do not find themselves in that dreaded “never never land”—never the object of concern, never the recipient of proper parental guidance.

To those who are not yet married, I counsel: People who marry in the hope of forming a permanent partnership require certain skills and attitudes of mind. They must be skillful in adapting to each other. They need capacity to work out mutual problems. They need willingness to give and take in the search for harmony. They need unselfishness of the highest sort, with thought for one’s partner taking the place of desire for oneself.

Many years ago I had the opportunity to deliver a commencement address to a graduating class. I had gone to the home of President Hugh B. Brown that we might drive together to the university where he was to conduct the exercises and I was to speak. As President Brown entered my car, he said, “Wait a moment.” He looked toward the large bay window of his lovely home, and then I realized what he was looking for. The curtain parted, and I saw Sister Zina Brown, his beloved companion of well over fifty years, at the window, propped up in a wheelchair, waving a little white handkerchief. President Brown took from his inside coat pocket a white handkerchief, which he waved to her in return. Then, with a smile, he said to me, “Let’s go.”

As we drove, I asked President Brown to tell me about the sign of the white handkerchiefs. He related to me the following incident: “The first day after Sister Brown and I were married, as I went to work I heard a tap at the window, and there was Zina, waving a white handkerchief. I found mine and waved in reply. From that day until this I have never left my home without that little exchange between my wife and me. It is a symbol of our love one for another. It is an indication to one another that all will be well until we are joined together at eventide.” Yes, a model to follow, “an example of the believers.”

To you young women in attendance tonight, you, too, can be a model—even an example. We are all aware that we live in a time when there are those who mock virtue, who peddle pornography under the guise of art or culture, who turn a blind eye, a deaf ear, and a calloused heart to the teachings of Jesus and a code of decency. Many of our young people are tugged in the wrong direction and enticed to partake of the sins of the world. Yearningly such individuals seek for the strength of those who have the ability to stand firm for truth. Through righteous living and by extending the helping hand and the understanding heart, you can rescue, you can save. How great will then be your joy. How eternal will be the blessing you will have conferred.

Some women face illness and incapacity, even to the point of being bedfast. Even so, there is the privilege to rise above affliction and to be a true example of faith, of love, and of service. Such was the partnership of Virginia and her husband, Eugene Jelesnik. They for many years worked together in bringing the gift of song and the joy of music to thousands of servicemen and women and to audiences from stages worldwide. Then illness and advancing age forced Virginia to remain within four walls—bedfast. But her spirit could not be held hostage by an impaired body. She continued to encourage her husband and to be his inspiration and constant support. All who are the beneficiaries of Eugene’s community concerts and his civic service marvel at his energy, his enthusiasm, and his kindness. In his many responsibilities, Virginia was ever a source of his strength.

While the Apostle Paul urged that we be examples of the believers, he didn’t restrict the boundaries of our service or limit the extent of our influence.

In July of this year, my wife and I attended an honor banquet where individuals were recognized for their quiet service, their selfless sacrifice, their untold devotion to lifting others to a higher plane of living with no thought of aggrandizement or personal reward. One Native American lady had literally given much of her life to teaching boys and girls of her native race how to live, how to love, and how to serve. Her response when recognized for her accomplishments bespoke her humility. Quietly and sincerely, she said just two special words: “Thank you.”

Another beautiful woman was honored for her caring, her serving, and her leadership. As a nurse she comforted the wounded in World War II. As a wife and partner with her husband, she built a worldwide business which blessed the lives of many. And today she, as a widow, continues daily service to her state and community. She seems to always be smiling. Perhaps this is because she has found the key to happiness. She has always been a missionary. She has ever been there when needed.

Yet another, we learned, had quietly yet effectively labored with love to ensure that the rights of abused children should not be neglected or abandoned.

There were others. All qualified for the definition of a pioneer—namely, “one who goes before, showing others the way to follow.”

During the banquet and program, I sat next to a well-known personality, Flip Harmon, and his wife, Lois. Flip has been involved with the direction of the Days of ’47 celebration for forty-three years, this being an annual July 24th activity in Salt Lake City. Since Flip was up and around the room fulfilling his official duties, I had the privilege of talking with Lois. She mentioned that she and family members were in attendance at every presentation of the famous rodeo which is one of the highlights of the Days of ’47 celebration. Now, a rodeo is nice once in a while—but every night? I asked Lois how she endured the schedule. Her response was from the heart. “This is Flip’s life, and I want to be part of it. He counts on me.” The night I had attended the rodeo with Sister Monson, my Aunt Blanche (age ninety-five), and our grandchildren, Lois was surrounded by children and precious grandchildren. She was the epitome of happiness. Now, during our luncheon conversation, Lois volunteered to me a few details about her husband. She said Flip had an angel mother who prayed fervently for her sons as they served their country during wartime. When Flip returned home, he and Lois were married. A busy life and welcome children followed. Each year as their wedding anniversary approached, Flip would say to Lois, “What gift do you want for our anniversary?” Each year the answer was the same, “A temple sealing.” The gift was not given.

Then one year, as the perennial question was asked, “What do you want, Lois, for our anniversary?” and the usual response was given, “To go to the temple of God together,” Flip’s reply was unexpected: “Fine. I’ll prepare for such an event.” They were sealed for time and eternity in the holy house of God on their twenty-ninth anniversary. Later, Flip served as a bishop. Each remains faithful to the other and loyal to the Lord.

As Lois continued, I noticed tears brimming in her eyes. She said, “You know Flip always wears cowboy boots. At the end of each day he would sit in the chair before the fireplace, where he would take off his boots and then read the paper. He would never put away the boots, no matter how many times I mentioned the subject. Years ago that would bother me. But not anymore. Today I just love those boots. Tender are my feelings and full is my heart as I willingly and lovingly pick them up and put them away each evening.”

Now tears were brimming in my eyes. Unexpectedly, Lois Harmon was asked to come to the podium, where she was given signal honor for her silent service. A beautiful bouquet of red roses was presented to her. Flip was asked to respond. His expression was from his heart. It was as though the two of them were alone in the large hotel dining area. “Lois is the light of my life. She’s my eternal partner.” (The word partner seemed to fit with the cowboy boots.) “We’ll be together forever.” Patience was rewarded. Love was expressed. Heaven was near.

My dear sisters, young and those just a bit older, though your circumstances may differ and your opportunities may vary, you can be models to follow, even “examples of the believers.”

In the Holy Temple just east of the Tabernacle on Temple Square in Salt Lake City, a beautiful tribute was paid to two of our sisters serving in the nursery. They, of course, were dressed in white, as were the children who had, that evening, been sealed to their parents. As the sisters bade their farewell to the children, one tiny girl, from a faith-filled heart, said to them, “Good night, angels.” May I borrow her words and say to you sisters worldwide, “Good night, angels.” In the name of Jesus Christ, amen.

\newpage
\settalk{At Parting}
\setspeaker{Thomas S. Monson}
\settalkdate{October 1992}
\talkheading{At Parting}{Thomas S. Monson}

Traditionally, the President of the Church, the Lord’s prophet, seer, and revelator, provides the concluding expressions of a general conference and gives his blessing to all. Humbly and respectfully I respond to the assignment to represent him at this time.

This has been a glorious conference. The prayers have been sincere and from the heart; the music and singing have lifted us heavenward and given us an upward reach we thought perhaps was beyond our grasp. The Brethren who have spoken, and Sister Jepsen, have declared the word of God and touched our hearts with their inspired messages. We are all better for having been a part of the conference.

President Benson’s chair has remained unoccupied during the conference sessions, which brings some sadness to our hearts. His ready smile, the wave of his hand, the declarations of truth that have marked his influence have been missed. However, President Benson, we are pleased and grateful that you have been a part of the conference through television. Our hearts go out to you in the passing of your beloved eternal companion, Flora. How thankful we are for the sacred covenant that binds you two sweethearts together for all eternity! The entire Church joins in a mighty prayer to our Heavenly Father that you may be cradled in the palm of His hand and blessed according to your need and His divine purposes. We sustain you. We follow you. We love you—our prophet.

President Benson revered President David O. McKay, who supervised his missionary labors in Great Britain those long years ago. President McKay closed a conference with these words: “As we come to this parting hour, I hope that the teachings and life of the Master seem to you all to be more beautiful, more necessary, and more applicable to human happiness than ever before. … Accepting him as my Redeemer, Savior, and Lord, I accept his gospel as the plan of salvation, as the one perfect way to human happiness and peace.”

President Joseph Fielding Smith, for whom President Benson had such great love, said as he concluded a conference: “Now I pray that our Father in heaven will bless his people—bless them abundantly and in full measure. I pray that the Saints shall stand firm against the pressures and enticements of the world; that they shall put first in their lives the things of God’s kingdom; that they shall be true to every trust and keep every covenant.”

President Harold B. Lee, boyhood friend and companion, and later esteemed associate of President Benson in the Lord’s work, declared: “I can’t leave this conference without saying to you that I have a conviction that the Master hasn’t been absent from us on these occasions. This is his church. … He isn’t an absentee master; he is concerned about us. He wants us to follow where he leads.”

President Spencer W. Kimball, who was sustained as an Apostle and member of the Council of the Twelve at the same time as President Benson, closed a general conference by saying: “As each one of these wonderful sermons has been rendered I’ve listened with [rapt] attention, and I have made up my mind that I shall go home and be a [better] man than I have ever been before.”

President Benson, these have been declarations from four of your associates who have been an ongoing influence in your life. You, yourself, have said in a similar close of a conference: “May we all go to our homes rededicated to the sacred mission of the Church as so beautifully set forth in these conference sessions—to ‘invite all to come unto Christ’ (D\&C 20:59), ‘yea, come unto Christ, and be perfected in him’ (Moro. 10:32).”

My brothers and sisters, I know the love President Benson has for you, for the Lord, and for His work. He would urge us to keep the commandments, sanctify our homes, and perfect our lives. May we, in unity, as members of The Church of Jesus Christ of Latter-day Saints, achieve these three objectives. Doing so will bring joy to our souls, peace to our prophet’s heart, and the smile of God’s approval on our efforts.

\begin{verse}
Sing we now at parting\\
One more strain of praise.\\
To our Heavenly Father\\
Sweetest songs we’ll raise.\\
For his loving kindness,\\
For his tender care,\\
Let our songs of gladness\\
Fill this Sabbath air.
\end{verse}

The work is true. Jesus is the Christ. Ezra Taft Benson is a prophet of God. I so testify and pray that heaven’s blessings may attend all of us in the name of Jesus Christ, amen.

\newpage
\settalk{Miracles—Then and Now}
\setspeaker{Thomas S. Monson}
\settalkdate{October 1992}
\talkheading{Miracles—Then and Now}{Thomas S. Monson}

Almost forty years ago I received an invitation to meet with President J. Reuben Clark, Jr., a Counselor in the First Presidency of the Church, a statesman of towering stature, and a scholar of international renown. My profession then was in the field of printing and publishing. President Clark made me welcome in his office and then produced from his old rolltop desk a large sheaf of handwritten notes, many of them made when he was a law student long years before. He proceeded to outline for me his goal of producing a harmony of the Gospels. This goal was achieved with his monumental work Our Lord of the Gospels.

Recently I took down from my library shelf a personally inscribed, leather-bound copy of this classic treatment of the life of Jesus of Nazareth. As I perused the many pages, I paused at the section entitled “The Miracles of Jesus.” I remembered as though it were yesterday President Clark asking me to read to him several of these accounts while he sat back in his large leather chair and listened. This was a day in my life never to be forgotten.

President Clark asked me to read aloud the account found in Luke concerning the man filled with leprosy. I proceeded to read:

“And it came to pass, when he was in a certain city, behold a man full of leprosy: who seeing Jesus fell on his face, and besought him, saying, Lord, if thou wilt, thou canst make me clean.

“And he put forth his hand, and touched him, saying, I will: be thou clean. And immediately the leprosy departed from him” (Luke 5:12–13).

He asked that I continue reading from Luke concerning the man afflicted with palsy and the enterprising manner in which he was presented for the attention of the Lord:

“And, behold, men brought in a bed a man which was taken with a palsy: and they sought means to bring him in, and to lay him before him.

“And when they could not find by what way they might bring him in because of the multitude, they went upon the housetop, and let him down through the tiling with his couch into the midst before Jesus.

“And when he saw their faith, he said unto him, Man, thy sins are forgiven thee” (Luke 5:18–20).

There followed snide comments from the Pharisees concerning who had the right to forgive sins. Jesus silenced their bickering by saying: “Whether [it] is easier, to say, Thy sins be forgiven thee; or to say, Rise up and walk?

“But that ye may know that the Son of man hath power upon earth to forgive sins, (he said unto the sick of the palsy,) I say unto thee, Arise, and take up thy couch, and go into thine house.

“And immediately he rose up before them, and took up that whereon he lay, and departed to his own house, glorifying God” (Luke 5:23–25).

President Clark removed from his pocket a handkerchief and wiped the tears from his eyes. He commented, “As we grow older, tears come more frequently.” After a few words of goodbye, I departed from his office, leaving him alone with his thoughts and his tears.

As I reflect on this experience, my heart fills with gratitude to the Lord for His divine intervention to relieve the suffering, heal the sick, and raise the dead. I grieve, however, for the many, similarly afflicted, who knew not how to find the Master, to learn of His teachings, and to become the beneficiaries of His power. I remember that President Clark himself suffered heartache and pain in the tragic death at Pearl Harbor of his son-in-law, Mervyn S. Bennion, captain of the battleship West Virginia. That day there had been no ram in the thicket, no steel to stop the shrapnel, no miracle to heal the wounds of war. But faith never wavered, and answered prayers provided the courage to carry on.

So it is today. In our lives, sickness comes to loved ones, accidents leave their cruel marks of remembrance, and tiny legs that once ran are imprisoned in a wheelchair.

Mothers and fathers who anxiously await the arrival of a precious child sometimes learn that all is not well with this tiny infant. A missing limb, sightless eyes, a damaged brain, or the term “Down’s syndrome” greets the parents, leaving them baffled, filled with sorrow, and reaching out for hope.

There follows the inevitable blaming of oneself, the condemnation of a careless action, and the perennial questions: Why such a tragedy in our family? Why didn’t I keep her home? If only he hadn’t gone to that party. How did this happen? Where was God? Where was a protecting angel? If, why, where, how—those recurring words do not bring back the lost son, the perfect body, the plans of parents, or the dreams of youth. Self-pity, personal withdrawal, or deep despair will not bring the peace, the assurance, or help which are needed. Rather, we must go forward, look upward, move onward, and rise heavenward.

It is imperative that we recognize that whatever has happened to us has happened to others. They have coped and so must we. We are not alone. Heavenly Father’s help is near.

Perhaps no other has been so afflicted as the man Job, who was described as “perfect and upright, and one that feared God, and eschewed evil” (Job 1:1). He prospered by every measurement. In other words, he had it all made. Then came the loss of literally everything: his wealth, his family, his health. At one time the suggestion was made that he “curse God, and die” (Job 2:9). Job’s summation of his faith, after ordeals demanded of few others, is a testimony of truth, a proclamation of courage, and a declaration of trust:

“Oh that my words were now written! oh that they were printed in a book!

“That they were graven with an iron pen and lead in the rock for ever!

“For I know that my redeemer liveth, and that he shall stand at the latter day upon the earth:

“And though after my skin worms destroy this body, yet in my flesh shall I see God:

“Whom I shall see for myself, and mine eyes shall behold, and not another” (Job 19:23–27).

Let me share with you a brief look into the lives of others, to learn that after the tears of a day of despair, a night of sorrow, “joy cometh in the morning” (Ps. 30:5).

Just two years ago, Eve Gail McDaniel and her parents, Bishop and Sister Jerry Lee McDaniel of the Reedsport Oregon Ward, came to my office and presented as a contribution to the Church Historical Department a copy of the Book of Mormon which Eve had written, by hand, and placed in three large binders. Eve, then 28, was born September 18, 1962. A case of meningitis when she was a baby resulted in brain damage. She cannot read, but she copied the entire Book of Mormon, letter by letter, over a period of about eighteen months. In doing so, she learned to recognize certain words and phrases, such as commandments and nevertheless. Her favorite—and she glowed as she repeated the phrase—was “And it came to pass.” Eve reflected the joy of accomplishment, even the smile of success. Her parents rejoiced in her gladness of heart and buoyancy of spirit. Heaven was very near.

On another occasion, near the Christmas season, I had the opportunity to meet in the Church Office Building with a group of handicapped children. There were about sixty in the group. My heart literally melted as I met with them. They sang for me “I Am a Child of God,” “Rudolph, the Red-nosed Reindeer,” and “As I Have Loved You, Love One Another.” There was such an angelic expression on their faces and such a simple trust expressed in their comments that I felt I was on sacred ground. They presented to me a beautiful booklet where each one had prepared a special page illustrating those blessings for which he or she was most thankful at Christmastime. I commend the many teachers and families who work behind the scenes in bringing a measure of comfort, purpose, and joy to these special children. They brightened my entire day.

Several years ago, Brigham Young University honored with a presidential citation Sarah Bagley Shumway, a truly remarkable woman of our time. The citation contained the words: “It is often within our homes and among our own family members that the eternally significant—but usually unheralded—dramas of daily living occur. The people in these plain but important places bring stability to the present and promise to the future. Their lives are filled with struggle and deep feeling as they face circumstances that rarely fit neatly within the formulae of plays, films and newscasts. But their victories, however slight, strengthen the boundaries through which the history of future generations must pass.”

Sarah married H. Smith Shumway, then her “friend and sweetheart of nine years,” in 1948. The courtship was longer than most because Smith, an infantry officer in World War II, was blinded and severely wounded by a land-mine explosion in the advance on Paris, France. During his long rehabilitation, Sarah learned braille so that she could correspond with him in privacy. She couldn’t tolerate the idea of others reading her letters aloud to the man she loved.

Something of the spirit of this young couple comes to us in the simple candor of Smith Shumway’s proposal of marriage. Finally home in Wyoming after the war, he told Sarah, “If you will drive the car and sort the socks and read the mail, I will do the rest.” She accepted the offer.

Years of study led to a successful career, eight accomplished children, a host of grandchildren, and lives of service. The Shumways, along life’s pathway, have faced problems of a child with severe deafness, a missionary son developing cancer, and a twin granddaughter injured at birth.

My family and I had the privilege to meet the entire Shumway clan at Aspen Grove a year ago. It was our joy to be with them. Each wore an identifying T-shirt on which was a map depicting the location of each child and family, along with the names of all. Brother Shumway, with justifiable pride, pointed to the location on his shirt of his precious ones and beamed the smile of gladness. Only then did I ponder that he had never seen any of his children or grandchildren. Or had he? While his eyes had never beheld them, in his heart he knew them and he loved them.

At an evening of entertainment, the Shumway family was on the stage at Aspen Grove. The children were asked, “What was it like growing up in a household with a sightless father?” One daughter smiled and said, “When we were little, occasionally we felt Daddy should not have too much dessert at dinner, so without telling him, we would trade our smaller helping with his larger one. Maybe he knew, but he never complained.”

One child touched our hearts when she recounted, “When I was about five years old, I remember my father holding my hand and walking me around the neighborhood, and I never realized he was blind because he talked about the birds and other things. I always thought he held my hand because he loved me more than other fathers loved their children.”

Today Brother Shumway is a patriarch. Who would you guess learned typing skills so as to be able to type the many blessings he gives? You’re correct: his beloved wife, Sarah.

Smith and Sarah Shumway and their family are examples of rising above adversity and sorrow, overcoming the tragedy of war-inflicted impairment, and walking bravely the higher roadway of life.

Ella Wheeler Wilcox, the poetess, wrote:

\begin{verse}
It is easy enough to be pleasant,\\
When life flows by like a song,\\
But the man worth while is one who will smile,\\
When everything goes dead wrong.\\
For the test of the heart is trouble,\\
And it always comes with the years,\\
And the smile that is worth the praises of earth\\
Is the smile that shines through tears.
\end{verse}

[“Worth While]

May I conclude with the inspiring example of Melissa Engle of West Valley, Utah. Melissa is featured in the August 1992 issue of the New Era. She tells her own story:

“When I was born I only had a thumb on my right hand because the umbilical cord got wrapped around my fingers and [severed them]. My dad wanted to find something I could do to strengthen my hand and make it useful. Playing the violin seemed like a natural because I wouldn’t have to finger with both hands, like you would with a flute. …

“I’ve been playing for about eight years now. I take private lessons, and I have to work at things like a paper route to help pay for them. I get to [my violin] lessons by riding a bus across town. …

“A highlight [of my life] was Interlochen, located on a lake in Michigan, one of the best music camps in the world for [youth]. I sent in my application for the eight weeks of intensive music training and couldn’t believe I [was] accepted.

“The only problem was money. It costs thousands of dollars, and there was no way [I could] make that much before the deadline. So I prayed and prayed, and about a week before I had to send in the money, I was called into the office of a man who had a grant for someone with a handicap who was pursuing the arts. That, to me, was a miracle. … I’m really grateful for it” (“Something You Really Love,” New Era, Aug. 1992, pp. 30–31).

Melissa, when she received the grant, turned to her mother, who had been anxious not to see her daughter disappointed and had thus attempted to curb her enthusiasm and hope, and said, “Mother, I told you Heavenly Father answers prayers, for look how He has answered mine.”

He that notes a sparrow’s fall had fulfilled a child’s dream, answered a child’s prayer.

To all who have suffered silently from sickness, to you who have cared for those with physical or mental impairment, who have borne a heavy burden day by day, year by year, and to you noble mothers and dedicated fathers—I salute you and pray God’s blessings to ever attend you. To the children, particularly they who cannot run and play and frolic, come the reassuring words: “Dearest children, God is near you, Watching o’er you day and night” (Hymns, 1985, no. 96).

There will surely come that day, even the fulfillment of the precious promise from the Book of Mormon:

“The soul shall be restored to the body, and the body to the soul; yea, and every limb and joint shall be restored to its body; yea, even a hair of the head shall not be lost; but all things shall be restored to their proper and perfect frame. …

“And then shall the righteous shine forth in the kingdom of God” (Alma 40:23, 25).

From the Psalm echoes the assurance: “My help cometh from the Lord, which made heaven and earth. …

“He that keepeth thee will not slumber.

“Behold, he that keepeth Israel shall neither slumber nor sleep” (Ps. 121:2–4).

Through the years the Latter-day Saints have taken comfort from the favorite hymn remembered from our youth:

\begin{verse}
When upon life’s billows you are tempest-tossed,\\
When you are discouraged, thinking all is lost,\\
Count your many blessings; name them one by one,\\
And it will surprise you what the Lord has done. …
\end{verse}

[Hymns, 1985, no. 241]

To any who from anguish of heart and sadness of soul have silently asked, “Heavenly Father, are you really there? … Do you hear and answer every … prayer?” (Children’s Songbook, p. 12), I bear to you my witness that He is there. He does hear and answer every prayer. His Son, the Christ, burst the bands of our earthly prisons. Heaven’s blessings await you. In the name of Jesus Christ, amen.

\newpage
\settalk{The Priesthood in Action}
\setspeaker{Thomas S. Monson}
\settalkdate{October 1992}
\talkheading{The Priesthood in Action}{Thomas S. Monson}

What a glorious sight is before me tonight! Here in the Tabernacle on Temple Square, in the Assembly Hall, at the BYU Marriott Center, and gathered together in chapels scattered throughout the world is a might army of men—even the royal army of the Lord. We have been entrusted with the priesthood. We have been prepared for duty. We have been called to serve.

The experience of the boy Samuel, as he responded to the Lord’s call, has ever been an inspiration to me, as it has no doubt been to each holder of the priesthood. We remember that the child Samuel ministered unto the Lord before Eli. One evening as the boy slept, the Lord called him by name: “Samuel.” And he answered, “Here am I.” Thinking that Eli had called him, Samuel ran to him and repeated the declaration, “Here am I.” He was advised to return to his sleep.

Three times the voice of the Lord came to him, with the same response. Then the Lord called a fourth time, repeating the boy’s name twice: “Samuel, Samuel.”

The lad’s answer, as before, is a classic example for you and me. He responded: “Speak; for thy servant heareth.

“And the Lord said to Samuel, Behold, I will do a thing in Israel, at which both the ears of every one that heareth it shall tingle” (see 1 Sam. 3:1–11).

Most of you young men will one day receive a call to serve a mission. How I pray that your response will be as was Samuel’s: “Here am I. … Speak; for thy servant heareth.” Then will heavenly help be yours. Every missionary strives to be the missionary his mother thinks he is, the missionary his father hopes he is—even the missionary the Lord knows he can become.

I remember a missionary recommendation for one young man on which the bishop had written: “This candidate is the finest I have ever recommended. He has served as an officer in the deacons, teachers, and priests quorums of which he has been a member. He excelled scholastically and athletically in high school. I know of no finer young man. P.S. I am proud to be his father.” President Spencer W. Kimball, then chairman of the Missionary Committee, mused, “I hope his parents will be content with his assigned mission. I know of no opening for him this morning in the celestial kingdom.”

Yes, sometimes expectations of those who love us are a bit beyond our capacity. Years ago, before a temple was completed in South Africa, the Saints planning to visit a temple had to travel the long and costly route to London, England, or, later, to São Paulo, Brazil. When I visited South Africa, they, with all the strength of their hearts and souls, petitioned me to importune President Kimball to seek the heavenly inspiration to erect a temple in their country. I assured them this was a matter for the Lord and His prophet. They responded, “We have faith in you, Brother Monson. Please help us.”

Upon returning to Salt Lake City, I discovered that a proposed temple for South Africa had already been approved and was to be announced immediately. When this occurred, I received a telegram from our members in South Africa. It read, “Thank you, Elder Monson. We knew you could do it!” You know, I believe I never did convince them that though I approved of the proposal, I did not bring it about.

Every call to serve is a human drama in the life of the recipient. I am certain that such has been the case with each of the Brethren who earlier today were sustained as new General Authorities. Let me share with you some marvelous lessons from the life of one of these Brethren, Jay E. Jensen, as recently reported in the Church News (“Spiritual Foundation Set Early in Life,” 8 Aug. 1992, pp. 6, 14).

Elder Jensen speaks of turning points in his life. His spiritual awakening began when he was a small boy growing up in Mapleton, Utah. His parents held family night long before it became a Church program. He recalled that his father read to him lessons from the Book of Mormon. His mother’s deep love for books also had a favorable impact on her son. However, it was when he read for himself Joseph Smith’s account of the First Vision that the witness of its truth became a reality.

Upon graduation from high school, young Jay and his sweetheart, Lona, decided to get married and not wait for a call to serve a mission. “It nearly broke my father’s heart,” Elder Jensen related. “Mother told me that Dad just wept.”

Two weeks later, and before wedding plans were finalized, Jay and Lona attended a sacrament meeting where a returned missionary reported his mission. The Spirit touched their hearts. They concluded to postpone marriage. Jay arose, went to the bishop’s office, and reported for missionary service. The rest is history. Jay served in the Spanish-American Mission.

Lona moved to California for employment and served a stake mission. Upon the completion of Jay’s mission, they were married in the Manti Temple. Elder Jensen’s father lived long enough to see his son serve an honorable mission and marry in the temple. Sister Jensen has often said that sending her husband-to-be on a mission was the hardest thing she ever did, but that it was the most rewarding. “I would never do it differently. We could never have been as happy otherwise.”

Today, Jay and Lona serve in Guatemala. He is a member of the Central America Area Presidency.

Reflecting on these turning points in the lives of Jay and Lona Jensen, we recall the observation, “The door of history turns on small hinges,” and so do people’s lives.

Fathers, grandfathers, are we reading to our sons and grandsons the word of the Lord? Returned missionaries, do your messages and your lives inspire others to stand up and serve? Brethren, are we sufficiently in tune with the Spirit that when the Lord calls, we can hear, as did Samuel, and declare, “Here am I”? Do we have the fortitude and the faith, whatever our callings, to serve with unflinching courage and unshakable resolve? When we do, the Lord can work His mighty miracles through us.

One such miracle is taking place in the southern part of the United States in the area once referred to as the Confederacy. It pertains to family history and temple work. During the period between 1860 and 1865, this region literally became saturated with the blood of America’s youth as soldiers by the hundreds of thousands perished. Even today, the earth here and there reveals a timeworn uniform button, a belt buckle, a spent bullet. But what of the men who fell while in the flower of their youth? Many had never married. Who was to do their temple work? Were they forever to be denied the blessings of eternal ordinances?

William D. Taylor, a Canadian with no ties to either side of the conflict that raged so long ago, found himself, together with wife and family, living in the old South and suddenly filled with a compelling interest in those who died while so young in years. An urgency came upon Brother Taylor to do something personally, a call to silent service.

In a letter to me dated July 20, 1992, Brother Taylor wrote: “It’s been approximately one year since I last gave you an update on the extraction work that is being done for the Confederate soldiers (approximately four years since this project was started). The extraction has been progressing at a steady pace. As of this writing, we have sent for temple work just over 101,000 names. I am thankful for being allowed to do this work. It brings me joy unparalleled to anything I have ever known. It’s hard to put my feelings into words. I exult when another regiment is prepared and ready to be sent to the temple, and my soul is pained when the information in the regimental history is insufficient for a soldier’s work to be submitted.”

A poet’s words expressed Brother Taylor’s feelings:

\begin{verse}
There I see them marching down the lane,\\
One in blue and one in gray,\\
Now arm and arm again,\\
And there I see them rising toward the Son,\\
Proud Rebels and proud Yankees,\\
Silent journey just begun.
\end{verse}

[David Matthews, “Road to Gettysburg”]

Brethren, let me share with you a description of priesthood service pertaining to this work, as described by a priesthood leader. He wrote: “On Saturday afternoon our Aaronic Priesthood young men and their leaders assembled at the temple to perform the baptismal work for the fallen soldiers. What a marvelous sight it was to see these young Aaronic Priesthood brethren being baptized by their own priesthood leaders. In almost every case, when the young brother had finished his 14 or 15 names, he would turn and embrace his leader and shed a few tears of joy. What an example of true priesthood love and service! I had the experience of being a witness at the font and gained firsthand knowledge of this and, in a few cases, the undeniable witness of the Spirit that those young soldiers who had died had accepted the baptisms that were being performed in their behalf by our Aaronic Priesthood brethren.

“We wrote down the name of each soldier who was baptized that glorious day so that the young men could have a brief history of the soldiers for whom they were baptized. I have no doubt that this experience will have a lifelong effect for good for all those who participated.”

The statement of President Joseph F. Smith, in speaking of the redemption of the dead, provides a touching explanation of the joy felt by all who participate in this and other similar endeavors: “Through our efforts in their behalf their chains of bondage will fall from them, and the darkness surrounding them will clear away, that light may shine upon them and they shall hear in the spirit world of the work that has been done for them by their children here, and will rejoice with you in your performance of these duties.”

Brother William Taylor, I salute you for your leadership in bringing eternal blessings to your “troops,” who must indeed call your name blessed.

When one holds the priesthood of God, he never knows when his moment of service may come. The challenge is to be ready to serve. On August 24, Hurricane Andrew slammed into the Florida coast south of Miami. Wind gusts exceeded two hundred miles per hour. It became the most costly disaster in United States history. Eighty-seven thousand homes were destroyed, leaving 150,000 homeless. Damages are estimated at 30 billion dollars. One hundred seventy-eight member homes were damaged, with forty-six of them destroyed.

A spearhead unit was deployed from the Church welfare facility in Atlanta before the storm hit, and it arrived at its appointed location just as the winds abated. The truck carried food, water, bedding, tools, and medical supplies—the first relief shipment to arrive in the disaster area.

Local priesthood and Relief Society leaders organized rapidly to assess injuries and damage and to assist in the cleanup effort. Three large waves of member volunteers, numbering over five thousand, labored shoulder to shoulder with disaster-stricken residents, helping to repair three thousand homes, a Jewish synagogue, a Pentecostal church, and two schools. Forty-six missionaries from the Florida Fort Lauderdale Mission worked full time for more than two weeks unloading supply trucks, serving as interpreters, providing security and traffic control, and assisting with repairs.

Time will permit but a glance at several heartwarming accounts pertaining to this tremendous example of the priesthood in action.

God bless Elder Morrison, his counselors, and all fellow priesthood leaders; missionaries, both elders and sisters; and all the many thousands who have served so magnificently and unstintingly. Truly these responded as did Samuel: “Here am I.”

The cleanup following Hurricane Andrew continues, as does the work of repair pertaining to the devastation wrought by Hurricane Iniki, which struck the island of Kauai in the Hawaiian Islands.

In these cataclysmic events and in the quiet challenges of individual lives, the priesthood is truly in action. Let us never despair, for this is the work of the Lord in which we are engaged. It has been said, “The Lord shapes the back to bear the burden placed upon it.” The Master’s counsel to all of us assembled tonight, to whom priesthood authority has been given and of whom priesthood service is expected, brings peace to the heart and comfort to the soul:

“Come unto me, all ye that labour and are heavy laden, and I will give you rest.

“Take my yoke upon you, and learn of me; for I am meek and lowly in heart: and ye shall find rest unto your souls.

“For my yoke is easy, and my burden is light” (Matt. 11:28–30).

To this divine truth I testify as I leave my witness with you that this work is true, that the priesthood does combine and present to our Heavenly Father a mighty army of righteousness, in the name of Jesus Christ, amen.

\newpage
\settalk{Gifts}
\setspeaker{Thomas S. Monson}
\settalkdate{April 1993}
\talkheading{Gifts}{Thomas S. Monson}

President David O. McKay would frequently suggest the need for us to turn from the hectic day-to-day schedule filled with letters to answer, calls to be made, people to see, meetings to attend, and take time to meditate, to ponder, and to reflect on the eternal truths and the sources of the joy and happiness which comprise each person’s quest.

When we do, the mundane, the mechanical, the repetitious patterns of our lives yield to the spiritual qualities, and we acquire a much-needed dimension which inspires our daily living. When I follow this counsel, thoughts of family, experiences with friends, and treasured memories of special days and quiet nights course through my mind and bring a sweet repose to my being.

The Christmas season, with its special meaning, inevitably prompts a tear, inspires a renewed commitment to God, and provides, borrowing the words from the lovely song “Calvary,” “rest to the weary and peace to the soul.”

I reflect on the contrasts of Christmas. The extravagant gifts, expensively packaged and professionally wrapped, reach their zenith in the famed commercial catalogs carrying the headline “For the person who has everything.” In one such reading I observed a four-thousand-square-foot home wrapped with a gigantic ribbon and comparable greeting card which said, “Merry Christmas.” Other items included diamond-studded clubs for the golfer, a Caribbean cruise for the traveler, and a luxury trip to the Swiss Alps for the adventurer. Such seemed to fit the theme of a Christmas cartoon which showed the Three Wise Men traveling to Bethlehem with gift boxes on their camels. One says, “Mark my words, Balthazar, we’re starting something with these gifts that’s going to get way out of hand!”

Then there is the remembered Christmas tale of O. Henry about a young husband and wife who lived in abject poverty yet who wanted to give one another a special gift. But they had nothing to give. Then the husband had a ray of inspiration: “I shall provide my dear wife a beautiful ornamental comb to adorn her magnificent long black hair.” The wife also received an idea: “I shall obtain a lovely chain for my husband’s prized watch which he values so highly.”

Christmas day came; the treasured gifts were exchanged. Then the surprise ending, so typical of O. Henry’s short stories: The wife had shorn her long hair and sold it to obtain funds to purchase the watch chain, only to discover that her husband had sold his watch, that he might purchase the comb to adorn her beautiful long hair, which now she did not have.

At home in a hidden-away corner, I have a small black walking stick with an imitation silver handle. It once belonged to a distant relative. Why do I keep it for a period now spanning sixty years? There is a special reason. You see, as a very small boy I participated in a Christmas pageant in our ward. I was privileged to be one of the Three Wise Men. With a bandanna about my head, Mother’s Chickering piano bench cover draped over my shoulder, and the black cane in my hand, I spoke my assigned lines: “Where is he that is born King of the Jews? for we have seen his star in the east, and are come to worship him.” I don’t recall all of the words in that pageant, but I vividly remember the feelings of my heart as the three of us “wise men” looked upward and saw a star, journeyed across the stage, found Mary with the young child Jesus, then fell down and worshiped him and opened our treasures and presented gifts: gold, frankincense, and myrrh.

I especially liked the fact that we did not return to the evil Herod to betray the baby Jesus, but obeyed God and departed another way.

The years have flown by, the events of a busy life take their proper places in the hallowed halls of memory, but the Christmas cane continues to occupy a special place in my home; and in my heart is a commitment to Christ.

For a few moments, may we set aside the catalogs of Christmas, with their gifts of exotic description. Let’s even turn from the flowers for Mother, the special tie for Father, the cute doll, the train that whistles, the long-awaited bicycle—even the “Star Trek” books and videos—and direct our thoughts to those God-given gifts that endure. I have chosen from a long list just four:

First, the gift of birth. It has been universally bestowed on each of us. Ours was the divine privilege to depart our heavenly home to tabernacle in the flesh and to demonstrate by our lives our worthiness and qualifications to one day return to Him, precious loved ones, and a kingdom called celestial. Our mothers and our fathers bestowed this marvelous gift on us. Ours is the responsibility to show our gratitude by the actions of our lives.

My own father, a printer, gave to me a copy of a piece he had printed. It was entitled “A Letter from a Father” and concluded with this thought: “Perhaps my greatest hope as a parent is to have such a relationship with you that when the day comes and you look down into the face of your first child, you will feel deep within you the desire to be to your child the kind of parent your dad has tried to be to you. What greater compliment could any man ask? Love, Dad.”

Our gratitude to Mother for the gift of birth is equal or beyond that owed to Father. She who looked upon us as “a sweet new blossom of humanity, fresh fallen from God’s own home, to flower on earth” and cared for our every need, comforted our every cry, and later rejoiced in any of our accomplishments and wept over our failures and disappointments, occupies a singular place of honor in our hearts.

A passage from 3 John sets forth the formula whereby we might express to our parents our gratitude for the gift of birth: “I have no greater joy than to hear that my children walk in truth.” Let us so walk. Let us so honor the givers of this priceless gift of birth.

Second, the gift of peace. In the raucous world in which we live, the din of traffic, the blaring commercials of the media, and the sheer demands placed on our time—to say nothing of the problems of the world—cause headache, inflict pain, and sap our strength to cope. The burden of sickness or the grief of mourning a loved one departed brings us to our knees seeking heavenly help. With the ancients we may wonder, “Is there no balm in Gilead?” There is a certain sadness, even hopelessness, in the verse:

\begin{verse}
There is never a life without sadness,\\
There is never a heart free from pain;\\
If one seeks in this world for true solace,\\
He seeks it forever in vain.
\end{verse}

He who was burdened with sorrow and acquainted with grief speaks to every troubled heart and bestows the gift of peace: “Peace I leave with you, my peace I give unto you: not as the world giveth, give I unto you. Let not your heart be troubled, neither let it be afraid.”

He sends forth His word through the missionaries serving far and wide, proclaiming His gospel of good tidings and salutation of peace. Vexing questions such as “Whence did I come?” “What is the purpose of my being?” and “Whither go I after death?” are answered by His special servants. Frustration flees, doubt disappears, and wonder wanes when truth is taught in boldness, yet in a spirit of humility, by those who have been called to serve the Prince of Peace—even the Lord Jesus Christ. His gift is bestowed individually: “Behold, I stand at the door, and knock: if any man hear my voice, and open the door, I will come in to him.”

The passport to peace is the practice of prayer. The feelings of the heart, humbly expressed rather than a mere recitation of words, provide the peace we seek.

In Shakespeare’s Hamlet, the wicked King Claudius kneels and tries to pray, but he rises and says: “My words fly up, my thoughts remain below: Words without thoughts never to heaven go.”

One who received and welcomed the gift of peace was Joseph Millett, an early missionary to the Maritime Provinces of Canada, who learned, while there and in his later experiences in life, of the need to rely on heavenly help. An experience which he recalled in his journal is a beautiful illustration of simple yet profound faith:

“One of my children came in, said that Brother Newton Hall’s folks were out of bread. Had none that day. I put … our flour in [a] sack to send up to Brother Hall’s. Just then Brother Hall came in. Says I, ‘Brother Hall, how are you [fixed] for flour.’ ‘Brother Millett, we have none.’ ‘Well, Brother Hall, there is some in that sack. I have divided [it] and was going to send it to you. Your children told mine that you were out.’ Brother Hall began to cry. Said he had tried others. Could not get any. Went to the cedars and prayed to the Lord and the Lord told him to go to Joseph Millett. ‘Well, Brother Hall, you needn’t bring this back if the Lord sent you for it. You don’t owe me for it.’ You can’t tell how good it made me feel to know that the Lord knew that there was such a person as Joseph Millett.”

Prayer brought the gift of peace to Newton Hall and to Joseph Millett.

Third, the gift of love. “Master, which is the great commandment in the law?” queried the lawyer who spoke to Jesus. Came the prompt reply: “Thou shalt love the Lord thy God with all thy heart, and with all thy soul, and with all thy mind.

“This is the first and great commandment.

“And the second is like unto it, Thou shalt love thy neighbour as thyself.”

On another occasion, the Lord taught, “He that hath my commandments, and keepeth them, he it is that loveth me.” The scriptures are filled with the importance of love and its relevance in our lives. The Book of Mormon teaches that charity is the pure love of Christ. The Master Himself provided an ideal pattern for us to follow. Of Him it was said that He “went about doing good, … for God was with him.”

A few lines from the favorite musical The Sound of Music suggest a course of action all might well follow:

\begin{verse}
A bell is no bell till you ring it,\\
A song is no song till you sing it,\\
And love in your heart wasn’t put there to stay—\\
Love isn’t love till you give it away.
\end{verse}

A segment of our society desperately yearning for an expression of true love is found among those growing older, and particularly when they suffer from pangs of loneliness. The chill wind of dying hopes and vanished dreams whistles through the ranks of the elderly and those who approach the declining side of the summit of life.

“What they need in the loneliness of their older years is, in part at least, what we needed in the uncertain years of our youth: a sense of belonging, an assurance of being wanted, and the kindly ministrations of loving hearts and hands; not merely dutiful formality, not merely a room in a building, but room in someone’s heart and life. …

“We cannot bring them back the morning hours of youth. But we can help them live in the warm glow of a sunset made more beautiful by our thoughtfulness, by our provision, and by our active and unfeigned love.” So wrote Elder Richard L. Evans some years ago.

At times an awareness of the elderly is brought into focus by a reminder from one ever so young. May I share with you a Pakistani folktale which illustrates this truth:

An ancient grandmother lived with her daughter and grandson. As she grew frail and feeble, instead of being a help around the house, she became a constant trial. She broke plates and cups, lost knives, spilled water. One day, exasperated because the old woman had broken another precious plate, the daughter sent the grandson to buy his grandmother a wooden plate. The boy hesitated because he knew a wooden plate would humiliate his grandmother. But his mother insisted, so off he went. He returned bringing not one, but two wooden plates.

“I only asked you to buy one,” his mother said. “Didn’t you hear me?”

“Yes,” said the boy. “But I bought the second one so there would be one for you when you get old.”

Frequently we are inclined to wait a lifetime to express love for the kindness or help given by another even long years before. Perhaps just such an experience prompted George Herbert to say, “Thou that hast given so much to [me], give one thing more … a grateful heart.”

The story is told of a group of men who were talking about people who had influenced their lives and to whom they were grateful. One man thought of a high-school teacher who had introduced him to Tennyson. He decided to write and thank her.

In time, written in a feeble scrawl, came this letter:

“My Dear Willie:

“I can’t tell you how much your note meant to me. I am in my eighties, living alone in a small room, cooking my own meals, lonely, and like the last leaf lingering behind. You will be interested to know that I taught school for fifty years, and yours is the first note of appreciation I have ever received. It came on a blue, cold morning, and it cheered me as nothing has for years.”

As I read this account, I thought of the treasured line, “The Lord has two homes: heaven and a grateful heart.”

Much more could be said pertaining to the gift of love. However, a favorite verse sums up rather well this precious gift:

\begin{verse}
I have wept in the night\\
For the shortness of sight\\
That to somebody’s need made me blind;\\
But I never have yet\\
Felt a tinge of regret\\
For being a little too kind.
\end{verse}

Fourth, the gift of life—even immortality. Our Heavenly Father’s plan contains the ultimate expressions of true love. All that we hold dear—even our families, our friends, our joy, our knowledge, our testimonies—would vanish were it not for our Father and His Son, the Lord Jesus Christ. Among the most cherished thoughts and writings in this world is the divine statement of truth: “For God so loved the world, that he gave his only begotten Son, that whosoever believeth in him should not perish, but have everlasting life.”

This precious Son, our Lord and Savior, atoned for our sins and the sins of all. That memorable night in Gethsemane His suffering was so great, His anguish so consuming that He pleaded, “Father, if it be possible, let this cup pass from me: nevertheless not as I will, but as thou wilt.” Later, on the cruel cross, He died that we might live, and live everlastingly. Resurrection morning was preceded by pain, by suffering in accordance with the divine plan of God. Before Easter there had to be a cross. The world has witnessed no greater gift, nor has it known more lasting love.

Nephi gives to us our charge: “Ye must press forward with a steadfastness in Christ, having a perfect brightness of hope, and a love of God and of all men. … If ye shall press forward, feasting upon the word of Christ, and endure to the end, behold, thus saith the Father: Ye shall have eternal life.

“And now, behold, … this is the way; and there is none other way nor name given under heaven whereby man can be saved in the kingdom of God.”

I close with the words of a revered prophet, even President Harold B. Lee: “Life is God’s gift to man. What we do with our life is our gift to God.”

May we give generously to Him, as He has so abundantly given to us, by living and loving as He and His Son have so patiently taught, is my earnest prayer. In the name of Jesus Christ, amen.

\newpage
\settalk{Search and Rescue}
\setspeaker{Thomas S. Monson}
\settalkdate{April 1993}
\talkheading{Search and Rescue}{Thomas S. Monson}

During the seemingly never-ending years of the Vietnam conflict, we frequently heard through the media’s blaring voice the term search and destroy. This phrase helped explain to the public the peculiar nature of combat in that area of dense jungle, oppressive heat, and debilitating disease.

This war was not marked by large-scale battles on open terrain. Rather, the enemy was often not visible—but nonetheless highly dangerous—thereby leading to the concept of search and destroy. Casualties were high, suffering rampant, and destruction everywhere to be found. We will never know how many cried out their own expression of the biblical question, “Is there no balm in Gilead?” The world sighed profound relief when conflict ceased and peace prevailed.

I was thinking of the term search and destroy this past winter as I visited with a neighbor and friend in beautiful Heber Valley east of Salt Lake City. Some snowmobile adventurers had been lost for a several-day period in the backcountry of high winds, penetrating cold, and eerie silence. My friend Johnny told me of the desperate plight of the lost and referred to the anxiety of their families. He mentioned that he was a member of the county search and rescue force, whose members left their businesses and farms and went in search of the lost and missing.

The searchers had prayed for a break in the winter weather, knowing the critical element of time in such a rescue. Their prayers were answered; the weather cleared. Surveying each grid of the vast area through high-powered field glasses as the helicopter flew back and forth through the mountains and the ravines, the search party finally spotted the lost party. Then came the difficult task of reaching and retrieving the courageous group. All was well. The lost were found. Lives were spared. Worry and fear yielded to joy and jubilation.

Johnny, with heartfelt emotion, said to me, “I love to search and rescue. Just to look into the faces of those who could have died and feel, as well as see, their profound gratitude fills my body and soul with compassion and thanksgiving. I’ve never before experienced anything quite like it.”

Perhaps he was witnessing the personal understanding of the Lord’s pronouncement, “Remember the worth of souls is great in the sight of God.” Or maybe Johnny was feeling the penetrating declaration of the Prophet Joseph Smith, who said, “It is better to save the life of a man than to raise one from the dead.”

My thoughts turned to that favorite song from Sunday School, the one that always brings tears to my eyes and compassion to my heart:

\begin{verse}
Dear to the heart of the Shepherd,\\
Dear are the “ninety and nine”;\\
Dear are the sheep that have wandered\\
Out in the desert to pine.\\
Hark! he is earnestly calling,\\
Tenderly pleading today:\\
“Will you not seek for my lost ones,\\
Off from my shelter astray?”
\end{verse}

The next verse reflects our response to the Shepherd’s plea:

\begin{verse}
Green are the pastures inviting;\\
Sweet are the waters and still.\\
Lord, we will answer thee gladly,\\
“Yes, blessed Master, we will!\\
Make us thy true undershepherds;\\
Give us a love that is deep.\\
Send us out into the desert,\\
Seeking thy wandering sheep.”
\end{verse}

Tonight I express the gratitude of the First Presidency and Council of the Twelve Apostles and all the General Authorities of the Church to members worldwide for your generosity and sacrifice in contributing your time, talents, and means through fast offerings and other service to alleviate suffering and to bless lives.

In the past twelve months, for example, the LDS Church participated in more that 350 hunger relief, community development, and in-kind projects in Asia, Eastern Europe, Africa, Latin America, the Caribbean, and the United States and Canada.

Included in the 1992 projects were such diverse activities as shipping more than 7.6 million pounds of sorted, used clothing—more than 190 container loads—to overseas and domestic destinations for distribution to refugees, displaced families, and other needy. Special attention was given to needs in Africa, where clothing, blankets, and other supplies and more than a million pounds of food were authorized for famine relief and community development. Another half-million pounds of food were contributed to food banks and feeding programs for the homeless and other needs in the United States and abroad.

Couples are now serving on full-time humanitarian service missions in Europe, Africa, Asia, Mongolia, and Latin America. Individual doctors, nurses, educators, and others have served on short-term consulting assignments with government ministries, hospitals, schools, and other institutions in many countries. Some projects have attacked the causes of poverty and suffering by supporting community development efforts of the local people.

Though the Church sometimes implements programs directly, efforts to carry out relief and development projects are often handled through agencies which have established reputations for honest, effective service, including the American Red Cross, International Red Cross and Red Crescent Societies, Salvation Army, Catholic Relief Services, Catholic Community Services, and other religious and civic organizations. All of this is in addition to the vast help extended by bishops of wards, presidents of branches, and leaders of missions to members of the Church throughout the world. The words of a Western Hemisphere prophet, uttered centuries ago, are still heard and followed today. King Benjamin reminded his people that “when ye are in the service of your fellow beings ye are only in the service of your God.”

From that same sacred record we contemplate the words spoken of the people during the reign of Alma, son of Alma: “They did not send away any who were naked, or that were hungry, or that were athirst, or that were sick, or that had not been nourished; and they did not set their hearts upon riches; therefore they were liberal to all, both old and young, both bond and free, both male and female, whether out of the church or in the church, having no respect to persons as to those who stood in need.”

The book of Luke, in one chapter, provides us two related parables which prompt our thinking and guide our footsteps in following the Master. First is the parable of the lost sheep, and second the parable of the prodigal son. The Lord began:

“What man of you, having an hundred sheep, if he lose one of them, doth not leave the ninety and nine in the wilderness, and go after that which is lost, until he find it?

“And when he hath found it, he layeth it on his shoulders, rejoicing.

“And when he cometh home, he calleth together his friends and neighbours, saying unto them, Rejoice with me; for I have found my sheep which was lost.

“I say unto you, that likewise joy shall be in heaven over one sinner that repenteth, more than over ninety and nine just persons, which need no repentance.”

In the parable of the prodigal son, we remember that one son wasted his substance and was reduced to near starvation. I ponder the line “and no man gave unto him.” Finally, when he came to himself he returned to the land of his father, expecting nothing but a rebuke and reprimand.

“And he arose, and came to his father. But when he was yet a great way off, his father saw him, and had compassion, and ran, and fell on his neck, and kissed him.

“And the son said unto him, Father, I have sinned against heaven, and in thy sight, and am no more worthy to be called thy son.

“But the father said to his servants, Bring forth the best robe, and put it on him; and put a ring on his hand, and shoes on his feet:

“And bring hither the fatted calf, and kill it; and let us eat, and be merry:

“For this my son was dead, and is alive again; he was lost, and is found.”

To the faithful son who was a bit critical of his father’s actions toward his brother came the same response: “Thy brother was dead, and is alive again; and was lost, and is found.”

Could I leave that long-distant time and faraway place to share with you examples of the guiding influence of the Master Shepherd and how we, in the fulfillment of our assignments, whatever they may be in His service, can see the evidence of His divine help and feel the touch of His gentle hand.

I served as a bishop during the period of the Korean War. We had received from Church headquarters a letter indicating that bishops should send a personal letter to each serviceman every month, along with a copy of the Church magazine at that time, the Improvement Era, and a subscription to the Church News. That took a little doing. In our large ward we had about eighteen servicemen. We did not have much money. The priesthood quorums, with effort, supplied funds for the subscriptions to the publications, and I took care of the letter writing. From my experience in the navy at the end of a previous war, I knew the importance of receiving word from home.

One day the sister who took the shorthand for those individually dictated letters said to me, “Bishop Monson, don’t you ever get discouraged?”

I said, “No, I don’t. Why?”

“Do you realize,” she explained, “that this is the seventeenth consecutive monthly letter you have sent to Lawrence Bryson, and you have never received a reply?”

I said, “Well, send number seventeen. It might do the job.” And it did. I received a reply from an APO number, San Francisco. Brother Bryson, far away in the Pacific, had written a short letter which began: “Dear Bishop, I ain’t much at writing letters [I could have told him that seventeen months sooner], but today has been a special day. I have been ordained a teacher in the Aaronic Priesthood. My group leader has stayed close to me, and I am grateful to him.” Then he said, “By the way, thanks for the Church News. Thanks for the magazine. But a special thanks for your letter which comes each month.”

Years later at a stake conference in the Cottonwood Stake, when Elder James E. Faust was stake president, I mentioned that experience in a stake priesthood meeting. A man came up after the meeting and said, “Do you remember me?”

I looked at him. It had probably been twenty-two years since I’d seen him. I said, “Lawrence Bryson!”

He said, “That’s me. Thanks for the letters. That’s why I’m here today.”

Where is Lawrence Bryson now? He and his wife are currently serving full-time missions. Their lives demonstrate full activity in the Church. They are searching for sheep that are lost. I think they will know where to find them. I know they will save them.

I still have that wonderful letter written to me from Lawrence Bryson and dated “Christmas Day, December 25, 1953.” It was one of the most treasured Christmas gifts ever received by me. Sure, you sometimes wonder after seventeen letters have been sent why no reply has come, but I remembered a line of truth: “The wisdom of God may appear as foolishness to men. But the greatest single lesson we can learn in mortality is that when God speaks and a man obeys, that man will always be right.” The leaders of the Church had spoken. We as bishops needed only to obey. The blessing was sure to follow.

Brethren, in our priesthood callings I am confident we at times wonder if we are affecting the lives of others in a favorable manner. The instructor in the quorum who prepares so diligently, the home teachers who put aside their own convenience and carry a message to the families upon whom they call, and the quorum officers who reach out to rescue perhaps will never fully know the far-reaching influence of their service. This is particularly true of the faithful missionaries who day after day carry on in the service of the Master. Never complaining, ever serving, always sacrificing for the benefit of others, these noble servants deserve our undying gratitude and our fervent prayers.

The simple words from Ecclesiastes, or the Preacher, carry an assurance that brings comfort and inspires effort: “Cast thy bread upon the waters: for thou shalt find it after many days.”

Such was my experience as pertains to President George H. Watson, who today serves as first counselor in the Naperville Illinois Stake presidency.

Brother Watson wrote a letter to me, never mailed, dated 3 October 1978, which tells of his conversion to the Church and of his baptism, which took place in the summer of 1959 in eastern Canada, where I served as the mission president at that time. I did not receive this letter until this past year, when it was carried to me by Elder John E. Fowler, who discovered its existence while visiting with the Watson family following a stake conference in Naperville. Both Brother Watson and I have some modest reluctance in sharing with you his private letter, but feeling the impression that the account would help to encourage many of you brethren participating in this worldwide priesthood meeting this evening, we shall do so.

I will conclude by reading President Watson’s own words. He wrote:

“Dear Elder Monson:

“This is a letter out of the blue. Its purpose is to thank you for the letters you wrote some twenty years ago—one to me and the other about me—and to let you know the effect they had on my life.

“My name is George Watson. In 1957, at the age of twenty-one, I emigrated from Ireland, where I had grown up, to Canada. The main purpose of going to Canada was to put together sufficient money to do postgraduate work at London University.

“The firm for which I worked was in Niagara Falls, and I found a room at the ridiculously inexpensive cost of \$6.00 per week. The only drawback was that I had to drive the landlady—age seventy-three—to church each Sunday in St. Catharines, Ontario.

“I soon found this chore to be very annoying, as she used the twenty-five-minute drive to try to get me to see the missionaries from her church. I resisted this very effectively for better than a year, until one day she told me that there were two young ladies coming to supper, and would I care to join them. It is very difficult to be rude to lady missionaries!

“I did a great deal of thinking over the next few months and decided that although what eleven sets of missionaries were telling me felt right, I would have to give up too much, besides which I was fed up running my landlady to church. In order to stop her asking for the ride, I decided to take her half an hour late on the next Sunday and to go in and sit with her in an open-neck shirt, sneakers, and sports slacks. I thought this would embarrass her and she would not ask me again.

“My plan worked perfectly, except that she was not annoyed at being late, and I made as much impact as a damp squid. We arrived just as the Sunday School was splitting for class. I would not go into class and spent my time talking to a very fine man who was crippled and who ‘understood’ me. As I was to return to Ireland eight days later (July 1959), he suggested that I should join the Church on the Saturday before I left. He was to call and confirm this during the week, but I effectively countered this by not answering the phone all week. On Sunday, after a sleepless night, I phoned him to apologize and was baptized in Hamilton virtually on the way to the airport—knowing that I would never meet any Mormons in Ireland and that the Church would lose track of me.

“I have no idea, President Monson, where you found my address in Ireland, but on the Friday after I returned, I had a letter from you welcoming me into the Church, and on Sunday at 9:00 a.m. there was a knock on the door and a President Lynn stood on the doorstep saying he had had a letter from President Monson in Toronto asking him to watch over me.

“The next few months or years were traumatic. Three meetings on a Sunday were entirely unreasonable; no way would I speak in front of that group; they can’t expect more than 10 percent. Even more traumatic, my girlfriend set out to show me how ridiculous I was. She ended up being baptized.

“We now live in Illinois with three wonderful children. I often sit and ponder why the Lord has blessed us so greatly. We have all had reason to feel His sustaining hand in difficult times.

“Although it is unlikely that we will ever meet, I would like to very sincerely thank you for taking the trouble to write those two letters. They have completely changed the course of our lives. I am grateful for the knowledge of the Savior’s purpose in coming to earth, my relationship to Him, and what He expects of me. The courage and steadfastness of Joseph Smith, the Prophet, and the knowledge that he imparted to us will always be a source of inspiration to me. I am thrilled at the opportunity of serving in the Lord’s Church.

“May the Lord continue to bless you in His work, and thank you for the effect you have had on my life.”

“[signed] George Watson”

This past Christmas, when George Watson and his beloved Chloe came to Salt Lake City to visit two of their children and a son-in-law, they came to my office, that we might formally meet. They expressed their testimonies and again conveyed their thanks for all who had participated in this human drama, this miracle in our time. Tears flowed, prayers were offered, and gratitude conveyed.

It was an appropriate season of the year for our visit together, when all Christendom pauses for a brief moment and remembers Him—even Jesus Christ—who died that we might have eternal life. He who notes the fall of the sparrow surely orchestrated the search-and-rescue mission that brought the Watson family to His fold. May we ever be found in His service and on His errand is my humble prayer, in the name of Jesus Christ, amen.

\newpage
\settalk{The Temple of the Lord}
\setspeaker{Thomas S. Monson}
\settalkdate{April 1993}
\talkheading{The Temple of the Lord}{Thomas S. Monson}

My beloved brothers and sisters, it is customary for the President of the Church to open each conference, to greet the Saints worldwide, and to set the tone of all that follows. Since President Benson is unable to be with us in person, I respond to his invitation to speak in his behalf. For the most part, I will present his actual words.

Last Wednesday, President Hinckley and I had a most delightful visit with President Benson. He greeted us warmly, flashed that friendly smile all of us love, and made us feel most welcome. When President Hinckley outlined the plans for conference and asked the President if it was his wish that we go forward with the arrangements and extend his love to all, he responded with a resounding, “Yes!” We understand his concerns. We share his love, and we bring to you his blessings. This giant of the Lord merits our constant prayers and our abiding faith.

On Friday, March 26, Sister Monson and I attended and participated in the ribbon-cutting ceremony formally opening a truly magnificent exhibit in the museum west of Temple Square. It is entitled “The Mountain of the Lord’s House” and depicts the fascinating forty-year saga required for the construction of the Salt Lake Temple. Where possible, I urge all of you to see the exhibit and feel the spirit it conveys. Tuesday, April 6, the Salt Lake Temple will have a birthday. One hundred years will have passed since that glorious day when it was dedicated.

While visiting the exhibit, a reporter asked me the question, “Would President Benson like this exhibit?”

I answered, “He would love it!”

President Benson has always loved temples and temple work. When he felt better, each Friday he and Sister Benson would enter the temple to participate in a session. We knew our First Presidency meeting that morning must accommodate this commitment. One morning I commented that I had to get busy and do some of my own family names that were prepared. With a smile and a twinkle in his eye, the President said, “Brother Monson, if you’re too busy, why not let Sister Benson and me do your names for you.” Needless to say, we found time to do the work ourselves.

President Benson’s own expressions indicate this love for temples. He reflected:

“I remember so well, as a little boy, coming in from the field and approaching the old farm house. … I could hear my mother singing, ‘Have I Done Any Good in the World Today?’ … I can still see her in my mind’s eye bending over the ironing board … with beads of perspiration on her forehead.” She was ironing long strips of white cloth, with newspapers on the floor to keep them clean. “When I asked her what she was doing, she said, ‘These are temple robes, my son. Your father and I are going to the temple at Logan.’

“Then she put the old flatiron on the stove, drew a chair close to mine, and told me about temple work—how important it is to be able to go to the temple and participate in the sacred ordinances performed there. She also expressed her fervent hope that some day her children and grandchildren and great-grandchildren would have the opportunity to enjoy those priceless blessings.” He continued, “I am happy to say that her fondest hopes in large measure have been realized.”

On another occasion, President Benson instructed us: “Sometimes in the peace of lovely temples, the serious problems of life find their solutions. [At times] pure knowledge flows to us there under the influence of the Spirit.” Said he: “I am grateful to the Lord for temples. The blessings of the House of the Lord are eternal. They are of the highest importance to us because it is in the temples that we obtain God’s greatest blessings pertaining to eternal life. Temples really are the gateways to heaven.”

He said: “May we remember always, as we [visit and work in these glorious temples], that the veil may become very thin between this world and the spirit world. I know this is true.” He declared, “It is well also that we keep in mind that it is all one great program on both sides of the veil and it is not too important whether we serve here or over there, as long as we serve with all our heart, might, mind, and strength.”

President Benson, your words are welcomed. We have heard them. We shall follow them. They, like the temples you so much love, are as a refuge from life’s storms—even a never-failing beacon guiding us to safety.

I echo the feelings of one and all, President Benson, in saying we love you and ever pray for you. In the name of Jesus Christ, amen.

\newpage
\settalk{Meeting Life’s Challenges}
\setspeaker{Thomas S. Monson}
\settalkdate{October 1993}
\talkheading{Meeting Life’s Challenges}{Thomas S. Monson}

Thirty years ago this conference, I was called and sustained as a member of the Council of the Twelve Apostles. On that occasion I asked earnestly for your faith and your prayers. And today as my opportunity to speak to you has come, I renew that request, that I may have your faith and prayers.

Just a month ago, while celebrating a national holiday, Elder Russell M. Nelson and I found ourselves with our children and grandchildren in a swimming pool filled with warm water and with a breathtaking view of an azure blue sky overhead. Mostly we were keeping a watchful eye on the little ones, much like a mother hen tracks the movement of her chicks. I said to Elder Nelson, “Isn’t it interesting that even though parents are watching their children, we assume the need to give overall supervision of our respective flock of grandchildren.” We had a wonderful time watching children at play and listening to their expressions of delight.

Then I noticed among those in the pool a father holding his severely handicapped son, moving the boy’s shrunken, tiny body back and forth in the pool. Other family members helped, and the lad obviously enjoyed the fun. He, however, was totally dependent. No sound of exuberant joy came forth from his lips, no splash of playful movement emanated from his almost lifeless limbs. Stricken as an infant with severe illness, he was left speechless, brain-damaged, and potentially a burden to loved ones. The boy’s grandfather said to me, “He is my grandson. All in our family love him. We enjoy his company; we respond to his needs. He is a blessing in our lives.”

Soon the crowd began to leave the pool. Laughter and play ceased. A silence shrouded the scene as the afternoon sun began its descent and the chill air reminded me it was time to go. But this tender view of love and devotion remained with me.

My thoughts turned to a place far distant and to a time long ago—even to another pool called Bethesda. The book of John describes what occurred there:

“Now there is at Jerusalem by the sheep market a pool, which is called in the Hebrew tongue Bethesda, having five porches.

“In these lay a great multitude of impotent folk, of blind, halt, withered, waiting for the moving of the water.

“For an angel went down at a certain season into the pool, and troubled the water: whosoever then first after the troubling of the water stepped in was made whole of whatsoever disease he had.

“And a certain man was there, which had an infirmity thirty and eight years.

“When Jesus saw him lie, and knew that he had been now a long time in that case, he saith unto him, Wilt thou be made whole?

“The impotent man answered him, Sir, I have no man, when the water is troubled, to put me into the pool: but while I am coming, another steppeth down before me.

“Jesus saith unto him, Rise, take up thy bed, and walk.

“And immediately the man was made whole, and took up his bed, and walked.”

Another scene of suffering and sorrow is found in the famous Tate Gallery in London, England. There adorns the wall of a much-traveled corridor a masterpiece entitled Sickness and Health. The painting portrays an organ-grinder with his monkey and a group of happy, healthy children frolicking and showing their amusement at the monkey’s antics. In the background is a small, pale-faced boy confined to a wheelchair, unable to play, unable to join in the fun of the other children. Feelings of empathy and silent tears of sadness overcome those who gaze upon the scene and sense the unspoken feelings of the sick boy’s heart.

Who can count the boys and girls, the men and women, where sickness has left its mark, rendering strong limbs lifeless and causing loved ones to shed tears of sorrow and offer prayers of faith for them?

Illness is not the only culprit that intrudes and alters our lives. In our hectic and fast-moving world, accidents can in an instant inflict pain, destroy happiness, and curtail our future. Such was the experience of young Robert Hendricks. Healthy and carefree three years ago, a sudden, three-car accident left him with brain damage, limited use of his limbs, and impaired speech. Summoned to his side by his mother, who pleaded her despair, I gazed at his almost-lifeless form as he lay on the white hospital bed in the critical care unit. Life supports functioning, his head swathed in bandages, his future was not only in doubt, but death appeared certain.

The hoped-for miracle, however, did occur. Heavenly help was forthcoming. Robert lived. His recovery has been labored and slow—but steady. A devoted friend, who was bishop at the time of the accident, has cared for Robert each week, getting him ready and driving him to his Sunday Church meetings—always patient, ever faithful.

One day Robert’s former bishop brought him to my office, since Robert wanted to meet with me, not having remembered that I saw him that night of crisis in the hospital. He and the dedicated bishop sat down, and Robert “talked” with me through a small electronic machine on which he spelled out his thoughts and they were then printed on strips of paper. He spelled out on the machine the love he has for his mother, his thanks for helping hands and willing hearts which have aided him, and his gratitude to a kind and caring Heavenly Father who has sustained him through his prayers. Here are some of his less private and personal messages: “I’m coming along pretty good, considering what I’ve been through.” Another: “I know that I will be able to help people and make some difference in people’s lives, and that’s great.” Another: “I don’t really know just how fortunate I am, but in my prayers I am told to just keep pushing on.”

At the conclusion of our visit, the bishop said, “Robert would like to surprise you.” Robert stood and, with considerable effort, said aloud, “Thank you.” A broad smile crossed his face. He was on the way back. “Thanks be to God” were the only words I could utter. Later I prayed aloud, “Thanks be also for loving bishops, kind teachers, and skilled specialists.”

Today, Robert, through the help of his former bishop, his current bishop, and others, has been to the temple. He has learned the computer. He is enrolled in computer study at college. He was also aided along the way by Deseret Industries helpers who provided encouragement and taught him essential skills. Now, with the support of a cane, Robert walks. He has learned to talk, though in halting phrases and with great effort. His progress has been phenomenal.

At times illness and accident take the lives of those whom they strike. Place and station, age and whereabouts make no difference. Death comes to the aged as they walk on faltering feet. Its summons is heard by those who have scarcely reached midway in life’s journey, and often it hushes the laughter of little children.

Throughout the world there is enacted daily the sorrowful scene of loved ones mourning as they bid farewell to a son, a daughter, a brother, a sister, a mother, a father, or a cherished friend.

Let us look in on one such scene which took place just last month in the Sunset Gardens Cemetery. Gathered were friends and family of Roger S. Olson, whose casket, bedecked with flowers, contained his earthly body. Claudia, his wife, six precious children, and family, friends, and associates stood in silence.

Just a few days earlier, Roger had left for his work, where he was a talented and recognized authority in his field of specialized photography. An accident resulted in the helicopter crash which took his life—all in the twinkling of an eye and without advance warning. Filled with grief but comforted by faith, those who had loved and lived together had bid but a temporary farewell to husband and father. They are sustained by the knowledge the skeptic rejects. They treasure the account recorded in Luke which describes that most significant event following the crucifixion and burial of our Lord and Savior Jesus Christ:

“Now upon the first day of the week, very early in the morning, [Mary Magdalene and the other Mary] came unto the sepulchre.” To their astonishment, the body of their Lord was gone. Luke records that two men in shining garments stood by them and said: “Why seek ye the living among the dead? He is not here, but is risen.”

Against the philosophy rampant in today’s world—a doubting of the authenticity of the Sermon on the Mount, an abandonment of Christ’s teaching, a denial of God, and a rejection of His laws—the Olsons and true believers everywhere treasure the testimonies of eyewitnesses to His resurrection. Stephen, doomed to the cruel death of a martyr, looked up to heaven and cried, “Behold, I see the heavens opened, and the Son of man standing on the right hand of God.”

Saul, on the road to Damascus, had a vision of the risen, exalted Christ. Peter and John also testified of the risen Christ. And in our dispensation, the Prophet Joseph Smith bore eloquent testimony of the Son of God, for he saw Him and heard the Father introduce him: “This is My Beloved Son. Hear Him!”

As we ponder the events that can befall all of us—even sickness, accident, death, and a host of lesser challenges—we can say, with Job of old, “Man is born unto trouble.” Needless to add, that reference to man in the King James Version of the book of Job encompasses women as well. It may be safely assumed that no person has ever lived entirely free of suffering and tribulation. Nor has there ever been a period in human history that did not have its full share of turmoil, ruin, and misery.

When the pathway of life takes a cruel turn, there is the temptation to think or speak the phrase, “Why me?” Self-incrimination is a common practice, even when we may have had no control over our difficulty. Socrates is quoted as saying, “If we were all to bring our misfortunes into a common store, so that each person should receive an equal share in the distribution, the majority would be glad to take up their own and depart.”

However, at times there appears to be no light at the tunnel’s end—no dawn to break the night’s darkness. We feel surrounded by the pain of broken hearts, the disappointment of shattered dreams, and the despair of vanished hopes. We join in uttering the biblical plea, “Is there no balm in Gilead?” We are inclined to view our own personal misfortunes through the distorted prism of pessimism. We feel abandoned, heartbroken, alone.

To all who so despair, may I offer the assurance of the Psalmist’s words: “Weeping may endure for a night, but joy cometh in the morning.”

Whenever we are inclined to feel burdened down with the blows of life’s fight, let us remember that others have passed the same way, have endured, and then have overcome.

Job was a perfect and an upright man who “feared God, and eschewed evil.” Pious in his conduct, prosperous in his fortune, Job was to face a test which would tempt any man. Shorn of his possessions, scorned by his friends, afflicted by his suffering, even tempted by his wife, Job was to declare from the depths of his noble soul, “Behold, my witness is in heaven, and my record is on high.” “I know that my redeemer liveth.”

Turning to our own time, let me share with you an example of faith, of courage, of compassion, of victory. It illustrates how it is possible to meet life’s challenges—head-on. It exemplifies the ability to suffer physical impairment, endure pain and suffering, and yet never complain. Such are Wendy Bennion of Sandy, Utah, and Jami Palmer of Park Valley, Utah. Both are teenagers; both have borne similar afflictions. Their situations run almost parallel. Since Wendy’s battle has been of a longer duration, I shall speak today of her.

Stricken with cancer at a tender age, subjected to long periods of chemotherapy, Wendy persevered valiantly. Teachers cooperated, parents and family helped—but the mainstay in her affliction has been her indomitable spirit. Wendy has brought cheer to others similarly afflicted. She has prayed for them; she has sustained them with her own example and faith.

After Wendy completed eighteen months of chemotherapy, a balloon-launching party was held in her honor. The public media covered the event. One of the many balloons launched that day was found miles away by Jayne Johnson. It had landed in her backyard, and she discovered it just as she was starting her own chemotherapy treatments. She wrote to Wendy, indicating she had been feeling sad and frightened but that finding the balloon and the note inside—which told about Wendy, her cancer, and the completion of her treatments—had given her the strength and that Wendy was a real inspiration to her. Wendy said, “I think she was supposed to find that balloon so that she would know that it’s not the end of the world and that people do get better.”

Though Wendy’s cancer recurred and a second round of therapy was needed, this choice young lady has not wavered, nor has she shrunk from her course. Rarely have I witnessed one with such courage, such determination, such faith. The same can be said of Jami Palmer. They personify the words of the poetess Ella Wheeler Wilcox, who wrote:

\begin{verse}
It is easy enough to be pleasant,\\
When life flows by like a song,\\
But the man worth while is one who will smile,\\
When everything goes dead wrong.\\
For the test of the heart is trouble,\\
And it always comes with the years,\\
And the smile that is worth the praises of earth\\
Is the smile that shines through tears.
\end{verse}

There is one life that sustains those who are troubled or beset with sorrow and grief—even the Lord Jesus Christ. Foretelling His coming, the prophet Isaiah records: “He hath no form nor comeliness; and when we shall see him, there is no beauty that we should desire him.

“He is despised and rejected of men; a man of sorrows, and acquainted with grief: and we hid as it were our faces from him; he was despised, and we esteemed him not.

“Surely he hath borne our griefs, and carried our sorrows: yet we did esteem him stricken, smitten of God, and afflicted.

“But he was wounded for our transgressions, he was bruised for our iniquities: the chastisement of our peace was upon him; and with his stripes we are healed.”

Yes, our Lord and Savior, Jesus Christ, is our Exemplar and our strength. He is the light that shineth in darkness. He is the Good Shepherd. Though engaged in His majestic ministry, He embraced the opportunity to lift burdens, provide hope, mend bodies, and restore life.

Few accounts of the Master’s ministry touch me more than His example of compassion shown to the grieving widow at Nain: “And it came to pass … that he went into a city called Nain; and many of his disciples went with him, and much people.

“Now when he came nigh to the gate of the city, behold, there was a dead man carried out, the only son of his mother, and she was a widow: and much people of the city was with her.

“And when the Lord saw her, he had compassion on her, and said unto her, Weep not.

“And he came and touched the bier: and they that bare him stood still. And he said, Young man, I say unto thee, Arise.

“And he that was dead sat up, and began to speak. And he delivered him to his mother.”

What power, what tenderness, what compassion did our Master thus demonstrate! We, too, can bless if we will but follow His noble example. Opportunities are everywhere. Needed are eyes to see the pitiable plight and ears to hear the silent pleadings of a broken heart—yes, and a soul filled with compassion, that we might communicate not only eye to eye or voice to ear but, in the majestic style of the Savior, even heart to heart.

His words become our guide: “In the world ye shall have tribulation: but be of good cheer; I have overcome the world.”

He lives. He will sustain each of us. May we keep His commandments. May we ever follow Him and merit His companionship, that we may successfully meet and overcome life’s challenges, I pray humbly, in His holy name, the name of Jesus Christ, amen.

\newpage
\settalk{The Upward Reach}
\setspeaker{Thomas S. Monson}
\settalkdate{October 1993}
\talkheading{The Upward Reach}{Thomas S. Monson}

My dear friends and fellow Scouters, Jere Ratcliffe, Bud Reid, and Mike Hoover, you honor me tonight by your attendance and with your remarks. I am humbled by the presentation of the Bronze Wolf Award. I know that as you bestow this honor, you are also expressing gratitude to the Church and to leaders past and present who have permitted me to serve on the National Executive Board these past twenty-four years and to follow in the footsteps of President Ezra Taft Benson and President George Albert Smith, who preceded me in this appointment. As a member of the International Committee of the board, I have had the privilege to travel to many lands and to witness the favorable influence of Scouting in the lives of young men of many languages, races, and cultures.

As a church, we do rather well in carrying the Scouting program in the United States and Canada. With the help of Jacques Moreillon, secretary general of the World Organization of the Scout Movement, we are taking steps to expand the influence of Scouting to our young men worldwide.

I love the inspired words of President Spencer W. Kimball as he spoke to Church members throughout the world: “The Church of Jesus Christ of Latter-day Saints affirms the continued support of Scouting and will seek to provide leadership which will help boys keep close to their families and close to the Church as they develop the qualities of citizenship and character and fitness which Scouting represents. … We have remained strong and firm in our support of this great movement for boys and of the Oath and the Law which are at its center.” Tonight we renew that endorsement.

Would you permit me to relate just one personal experience. When I was fourteen years old, our troop went to Big Cottonwood Canyon on a Scout outing. After setting up camp, our leader said to me, “Monson, you like to fish. I’m giving you two fishing flies—a black gnat and a white miller. Now you catch enough fish to feed this troop for the next three days, and I’ll pick all of you up on Saturday.” He departed. I never questioned his charge. I knew if I did my part I’d catch the fish and feed the troop. And I did. I was a man before I realized it just isn’t proper for the Scoutmaster to bail out on the boys. But what a learning experience it was for us.

The paintings of Norman Rockwell on the cover of The Saturday Evening Post magazine or in Boy’s Life always brought tender feelings to me. Of the two paintings I most admire, one is of a Scoutmaster sitting by the dying embers of the bonfire and observing his boys—fast asleep in their small tents. The sky is filled with stars, the tousled heads of the boys illumined by the fire’s glow. The Scoutmaster’s countenance reflects his love, his faith, his devotion. The scene recalls the thought, “The greatest gift a man can give a boy is his willingness to share a part of his life with him.”

The other painting is of a small lad, clad in the oversized Scout uniform of his older brother. He is looking at himself in a mirror which adorns the wall, his tiny arm raised in the Scout salute. It could well be entitled “Following in the Footsteps of Scouting.”

In this world where some misguided men and women strive to tear down and destroy great movements such as Scouting, I am pleased to stand firm for an organization that teaches duty to God and country, that embraces the Scout Law; yes, an organization whose motto is “Be prepared” and whose slogan is “Do a good turn daily.”

The Aaronic Priesthood prepares boys for manhood and the weightier duties of the Melchizedek Priesthood. Scouting helps our boys to walk uprightly the priesthood path to exaltation. Along that path there will be turns and detours, requiring decisions of utmost importance. Heavenly inspiration will provide a road map that will ensure the accuracy of our choices. There comes a time in the life of every young man for serious contemplation and wise evaluation concerning his future—for decisions determine destiny.

Tonight in this vast priesthood audience are those who have successfully navigated the pathways of their youth. Such men of experience and faith are needed as examples for those who look to them for guidance and safety. Brethren, are we prepared for our leadership opportunity—even our life-saving privilege? The need for our help is here and now.

In cities across the land and in nations throughout the world, there has occurred a deterioration of the home and family. Abandoned in many instances is the safety net of personal and family prayer. A macho-inspired attitude of “I can go it alone” or “I don’t need the help of anyone” dominates the daily philosophy of many. Frequently there is a rebellion against long-established traditions of decency and order, and the temptation to run with the crowd is overwhelming. Such a destructive philosophy, this formula for failure, can lead to ruin unless men of faith, filled with love, step forward to show a faltering boy the right way to go. Remember the verse:

\begin{verse}
He stood at the crossroads all alone,\\
The sunlight in his face.\\
He had no thought for the world unknown—\\
He was set for a manly race.\\
But the roads stretched east, and the roads stretched west,\\
And the lad knew not which road was best;\\
So he chose the road that led him down,\\
And he lost the race and victor’s crown.\\
He was caught at last in an angry snare\\
Because no one stood at the crossroads there\\
To show him the better road.
\end{verse}

Those who hold the Melchizedek Priesthood are not the only resource with the strength to lift, the wisdom to guide, and the ability to save. Many of you young men comprise the presidencies of quorums of deacons, quorums of teachers, and hold leadership positions assisting the bishops in guiding quorums of priests. As you magnify your callings with respect to aiding those over whom you preside, heavenly help will be forthcoming. Remember that throughout the ages of time, our Heavenly Father has shown His confidence in those of tender years.

The boy Samuel must have appeared like any boy his age as he ministered unto the Lord before Eli. As Samuel lay down to sleep and heard the voice of the Lord calling him, Samuel mistakenly thought it was aged Eli and responded, “Here am I.” However, after Eli listened to the boy’s account and told him it was of the Lord, Samuel followed Eli’s counsel and subsequently responded to the Lord’s call with the memorable reply, “Speak; for thy servant heareth.” The record then reveals that “Samuel grew, and the Lord was with him.”

Contemplate for a moment the far-reaching effect of the prayer of a boy, born in the year of our Lord one thousand eight hundred and five in Sharon, Windsor County, state of Vermont—even Joseph Smith, the first prophet of this dispensation. The Father and the Son appeared to him, and divine guidance was provided—all for the purpose of exalting the children of God.

We remember with gratitude that night of nights which marked the fulfillment of prophecy when a lowly manger cradled a newborn child. With the birth of the babe in Bethlehem, there emerged a great endowment, a power stronger than weapons, a wealth more lasting than the coins of Caesar. This child, born in such primitive circumstances, was to be the King of Kings and the Lord of Lords, the promised Messiah—even Jesus Christ, the Son of God.

As a boy, Jesus was found “in the temple, sitting in the midst of the doctors, both hearing them, and asking them questions.

“And all that heard him were astonished at his understanding and answers.

“And when [Joseph and His mother] saw him, they were amazed. …

“And Jesus increased in wisdom and stature, and in favour with God and man.”

He “went about doing good, … for God was with him.”

I mention these powerful examples so that every young man within the sound of my voice may know for himself his own strength when God is with him.

As each realizes his own potential and what our Heavenly Father expects of him—a determination to live proper standards, to be true to one’s best self, and to act always in accordance with a high sense of true values—there will follow incomparable joy and lasting peace.

A four-point guide will help focus our attention on such a goal:

First, be where we ought to be. A wise father counseled his son, “If you ever find yourself where you shouldn’t be, then get out!” Choose your friends carefully, for you will tend to be like them and be found where they choose to go.

Second, say what we ought to say. What we say and how we say it tend to reflect what we are. In the life of the Apostle Peter, when he attempted to distance himself from Jesus and pretended to be other than what he was, his tormenters detected his true identity with the penetrating statement, “Thy speech bewrayeth thee.” The words we utter will reflect the feelings of our hearts, the strength of our character, and the depth of our testimonies.

Third, do what we ought to do. Pierre, one of the central characters in Tolstoy’s War and Peace, torn by spiritual agonies, cries out to God, “Why is it that I know what is right and I do what is wrong?” Pierre needed a mind-set, a resolve—even a stiffening of his backbone. One clever with words put it this way as he paraphrased the familiar counsel “Never put off ’til tomorrow what you should do today,” by adding, “Why do we not put off ’til tomorrow what we shouldn’t do today!”

Then there is the excuse of the weak: “The devil made me do it.” It is only when we take charge of our own actions that we direct them in the proper course.

Fourth, be what we ought to be. The Apostle Paul counseled his beloved young friend Timothy, “Let no man despise thy youth; but be thou an example of the believers, in word, in conversation, in charity, in spirit, in faith, in purity.” Peter asked the question, “What manner of persons ought ye to be in all holy conversation and godliness?” Then Peter’s life answered convincingly his own question. The Master’s own voice queried: “What manner of men ought ye to be? Verily I say unto you, even as I am.”

On occasion, when I have met with young men, I have been asked the question, “Brother Monson, is there one thing I can do to help me pattern my life and live up to my full potential?” As I have searched memory’s corridors for an answer to such a question, I have recalled an experience of a few years ago. A group of friends were trail riding on strong Morgan horses when we came to a clearing which opened on a lush grass meadow with a small, clear stream meandering through it. No mule deer could wish for a better home. However, there was a danger lurking. The wily deer can detect the slightest movement in the surrounding bush; he can hear the crack of a twig and discern the scent of man. He is vulnerable from but one direction—overhead. In a mature tree, hunters had erected a platform high above the enticing spot. Though in many places this is illegal, the hunter can take his prey as it comes to eat and to drink. No twig would break, no movement disturb, no scent reveal the hunter’s whereabouts. Why? The magnificent buck deer, with its highly developed senses to warn of impending danger, does not have the capacity to look directly upward and thus detect his enemy. Man is not so restricted. His greatest safety is found in his ability and his desire to “look to God and live.”

Wrote the poet:

\begin{verse}
But chief of all thy wondrous works,\\
Supreme of all thy plan,\\
Thou hast put an upward reach\\
Into the heart of man.
\end{verse}

May I conclude with a heart-tugging account of one small boy, a Cub Scout whose love of Scouting brought him and those who knew him and loved him closer to God as he reached upward and stepped over the limits of mortality and entered the broad expanse of eternity, clad in the uniform he loved and wearing the honor he had won—in Scouting.

In October 1992, nine-year-old Jared Barney passed away as a result of brain cancer. He had, in his short life, endured multiple surgeries, along with radiation and chemotherapy treatments. His last surgery was August 9, 1992. A month after that, an MRI picked up six new tumors, two of which were already quite large.

The radiation and chemotherapy made Jared very ill. The surgeries were difficult, but he always bounced back very quickly. Although he suffered much pain, the Lord blessed and sustained him.

Jared had a special spirit that drew others to him. He never complained about how he felt or about having to be sick or about the treatments he had to have. When asked how he was doing, he always said, “Good,” no matter how he felt. He was ever known for his contagious smile. The Light of Christ was in his eyes.

May I quote from Jared’s mother, Olivia, who wrote concerning his last days: “Our many prayers were answered in behalf of our little son. We prayed that he would be able to walk, talk, and see until the end, and then that the Lord would take him quickly. He was able to do all of these things, and we are so thankful to the Lord for answering our prayers. Jared loved life so much, and we wanted him to be able to enjoy it fully until the end.

“Jared had earned some Cub Scout awards three weeks prior to his passing. He had earned his Bear badge, his Faith in God, a Gold Arrow Point, and two Silver Arrow Points. We know that he loved to get those awards. He was failing quickly, and he wouldn’t even let himself sleep until he could attend the pack meeting held on October 14, 1992, to achieve his awards. At the pack meeting, he raised his hand three times and told everyone how long he had waited for these awards and how happy he was to get them. When we returned home, he asked me to sew his badges on that very night. I did. Then he prayed that Heavenly Father would let him sleep because he was so tired. He said that three times. He went to sleep and never moved all night. From then on he slept most of the time until his passing.

“We buried him in his Cub Scout shirt with those long-awaited emblems sewn and pinned on the front. He had a beautiful service. Many were present, for he had made so many friends in the community through his example of courage and faith.”

Such was the influence of an inspired program in the life of a tiny boy and his family.

To all the Aaronic Priesthood assembled tonight with your fathers and your leaders, the priesthood program of the Church, with its accompanying activities, including Scouting, will help and not hinder you as you journey through life. May each one of us resolve to follow the example of our Lord and Savior, Jesus Christ, keep His commandments, and live His teachings, that we may inherit the greatest of all gifts, eternal life with God. In the name of Jesus Christ, amen.

\newpage
\settalk{The Path to Peace}
\setspeaker{Thomas S. Monson}
\settalkdate{April 1994}
\talkheading{The Path to Peace}{Thomas S. Monson}

On this beautiful Easter morning, prayers of gratitude for the life and mission of our Lord and Savior, Jesus Christ, fill the Sabbath air while strains of inspiring music comfort our hearts and whisper to our souls the ageless salutation, “Peace be unto you.”

In a world where peace is such a universal quest, we sometimes wonder why violence walks our streets, accounts of murder and senseless killings fill the columns of our newspapers, and family quarrels and disputes mar the sanctity of the home and smother the tranquillity of so many lives.

Perhaps we stray from the path which leads to peace and find it necessary to pause, to ponder, and to reflect on the teachings of the Prince of Peace and determine to incorporate them in our thoughts and actions and to live a higher law, walk a more elevated road, and be better disciples of Christ.

The ravages of hunger in Somalia, the brutality of hate in Bosnia, and the ethnic struggles across the globe remind us that the peace we seek will not come without effort and determination. Anger, hatred, and contention are foes not easily subdued. These enemies inevitably leave in their destructive wake tears of sorrow, the pain of conflict, and the shattered hopes of what could have been. Their sphere of influence is not restricted to the battlefields of war but can be observed altogether too frequently in the home, around the hearth, and within the heart. So soon do many forget and so late do they remember the counsel of the Lord: “There shall be no disputations among you, …

“For verily, verily I say unto you, he that hath the spirit of contention is not of me, but is of the devil, who is the father of contention, and he stirreth up the hearts of men to contend with anger, one with another.

“Behold, this is not my doctrine, to stir up the hearts of men with anger, one against another; but this is my doctrine, that such things should be done away.”

As we turn backward the clock of time, we recall that some fifty-five years ago a desperately arranged peace, a conference of peace, convened in the Bavarian city of Munich. Leaders of the European powers assembled even as the world tottered on the brink of war. Their purpose, openly stated, was to pursue a course which they felt would avert war and maintain peace. Mistrust, intrigue, a quest for power doomed to failure that conference. The outcome was not “peace in our time” but rather war and destruction to a degree not previously experienced. Overlooked, or at least set aside, was the hauntingly touching appeal of one who had fallen in an earlier war. He seemed to be writing in behalf of millions of comrades—friend and foe alike:

\begin{verse}
In Flanders fields the poppies blow\\
Between the crosses, row on row,\\
That mark our place; and in the sky\\
The larks, still bravely singing, fly\\
Scarce heard amid the guns below.
\end{verse}

Are we doomed to repeat the mistakes of the past? After such a brief interval of peace following World War I came the cataclysm of World War II. In fact, this June will mark the fiftieth anniversary of the famed landings of Allied forces on the beaches of Normandy. Tens of thousands of dignitaries and veterans will flock to the scene as the landings are reenacted. One writer observed:

“Lower Normandy has more than its share of [hallowed dead. Their bodies] lie in graves from Falaise to Cherbourg: 13,796 Americans, 17,958 British, 8,658 Canadian, 650 Polish, and around 65,000 Germans, more than 106,000 dead in all, and that is just the military, all killed in the space of a summer holiday.” Similar accounts could be written describing the terrible losses in other theaters of combat in that same conflict.

The famed statesman William Gladstone described the formula for peace when he declared: “We look forward to the time when the power of love will replace the love of power. Then will our world know the blessings of peace.”

World peace, though a lofty goal, is but an outgrowth of the personal peace each individual seeks to attain. I speak not of the peace promoted by man, but peace as promised of God. I speak of peace in our homes, peace in our hearts, even peace in our lives. Peace after the way of man is perishable. Peace after the manner of God will prevail.

We are reminded that “anger doesn’t solve anything. It builds nothing, but it can destroy everything.” The consequences of conflict are so devastating that we yearn for guidance—even a way to ensure our success as we seek the path to peace. What is the way to obtain such a universal blessing? Are there prerequisites? Let us remember that to obtain God’s blessings, one must do God’s bidding. May I suggest three ideas to prompt our thinking and guide our footsteps:

First: Search inward. Self-evaluation is always a difficult procedure. We are so frequently tempted to gloss over areas which demand correction and dwell endlessly on our individual strengths. President Ezra Taft Benson counsels us:

“The price of peace is righteousness. Men and nations may loudly proclaim, ‘Peace, peace,’ but there shall be no peace until individuals nurture in their souls those principles of personal purity, integrity, and character which foster the development of peace. Peace cannot be imposed. It must come from the lives and hearts of men. There is no other way.”

Elder Richard L. Evans observed: “To find peace—the peace within, the peace that passeth understanding—men must live in honesty, honoring each other, honoring obligations, working willingly, loving and cherishing loved ones, serving and considering others, with patience, with virtue, with faith and forbearance, with the assurance that life is for learning, for serving, for repenting, and improving. And God be thanked for the blessed principle of repenting and improving, which is a way that is open to us all.”

The place of parents in the home and family is of vital importance as we examine our personal responsibilities in this regard. Recently a distinguished group met in conference to examine the increase of violence in the lives of individuals, particularly the young. Some observations from their deliberations are helpful to us as we examine our priorities:

“A society that views graphic violence as entertainment … should not be surprised when senseless violence shatters the dreams of its youngest and brightest. …

“Unemployment and despair can lead to desperation. But most people will not commit desperate acts if they have been taught that dignity, honesty and integrity are more important than revenge or rage; if they understand that respect and kindness ultimately give one a better chance at success. …

“The women of the anti-violence summit have hit on the solution—the only one that can reverse a downward spiral of destructive behavior and senseless pain. A return to old-fashioned family values will work wonders.”

So frequently we mistakenly believe that our children need more things, when in reality their silent pleadings are simply for more of our time. The accumulation of wealth or the multiplication of assets belies the Master’s teaching:

“Lay not up for yourselves treasures upon earth, where moth and rust doth corrupt, and where thieves break through and steal:

“But lay up for yourselves treasures in heaven, where neither moth nor rust doth corrupt, and where thieves do not break through nor steal:

“For where your treasure is, there will your heart be also.”

The other evening I saw large masses of parents and children crossing an intersection in Salt Lake City en route to the Delta Center to see the Disney on Ice production of Beauty and the Beast. I actually pulled my car over to the curb to watch the gleeful throng. Fathers, who I am certain were cajoled into going to the event, held tightly in their hands the small and clutching hands of their precious children. Here was love in action. Here was an unspoken sermon of caring. Here was a rearranging of time as a God-given priority.

Truly peace will reign triumphant when we improve ourselves after the pattern taught by the Lord. Then we will appreciate the deep spirituality hidden behind the simple words of a familiar song: “There is beauty all around when there’s love at home.”

Second: Reach outward. Though exaltation is a personal matter, and while individuals are saved not as a group but indeed as individuals, yet one cannot live in a vacuum. Membership in the Church calls forth a determination to serve. A position of responsibility may not be of recognized importance, nor may the reward be broadly known. Service, to be acceptable to the Savior, must come from willing minds, ready hands, and pledged hearts.

Occasionally discouragement may darken our pathway; frustration may be a constant companion. In our ears there may sound the sophistry of Satan as he whispers, “You cannot save the world; your small efforts are meaningless. You haven’t time to be concerned for others.” Trusting in the Lord, let us turn our heads from such falsehoods and make certain our feet are firmly planted in the path of service and our hearts and souls dedicated to follow the example of the Lord. In moments when the light of resolution dims and when the heart grows faint, we can take comfort from His promise: “Be not weary in well-doing. … Out of small things proceedeth that which is great.

“Behold, the Lord requireth the heart and a willing mind.”

During the past year, the Primary organization has conducted an effort to have the children become better acquainted with the holy temples of God. Frequently this has entailed a visit to the temple grounds. The laughter of small children, the joy of unfettered youth, and the exuberance of energy displayed by them gladdened the heart of this observer. As a loving teacher guided a boy or girl to the large door of the Salt Lake Temple and the little one reached out and up to touch the temple, I could almost see the Master welcoming the little children to His side and could almost hear His comforting words: “Suffer the little children to come unto me, and forbid them not: for of such is the kingdom of God.”

Number three: Look heavenward. As we do, we find it comforting and satisfying to communicate with our Heavenly Father through prayer, that path to spiritual power—even a passport to peace. We are reminded of His beloved Son, the Prince of Peace, that pioneer who literally showed the way for others to follow. His divine plan can save us from the Babylons of sin, complacency, and error. His example points the way. When faced with temptation, He shunned it. When offered the world, He declined it. When asked for His life, He gave it.

On one significant occasion, Jesus took a text from Isaiah: “The Spirit of the Lord God is upon me; because the Lord hath anointed me to preach good tidings unto the meek; he hath sent me to bind up the brokenhearted, to proclaim liberty to the captives, and the opening of the prison to them that are bound”—a clear pronouncement of the peace that passeth all understanding.

Frequently, death comes as an intruder. It is an enemy that suddenly appears in the midst of life’s feast, putting out its lights and its gaiety. Death lays its heavy hand upon those dear to us and, at times, leaves us baffled and wondering. In certain situations, as in great suffering and illness, death comes as an angel of mercy. But to those bereaved, the Master’s promise of peace is the comforting balm which heals: “Peace I leave with you, my peace I give unto you: not as the world giveth, give I unto you. Let not your heart be troubled, neither let it be afraid.” “I go to prepare a place for you … ; that where I am, there ye may be also.”

How I pray that all who have loved then lost might know the reality of the Resurrection and have the unshakable knowledge that families can be forever. One such was a Major Sullivan Ballou, who, during the time of the American Civil War, wrote a touching letter to his wife—just one week before he was killed in the Battle of Bull Run. With me, feel the love of his soul, his trust in God, his courage, his faith.

“July 14, 1861

“Camp Clark, Washington

“My very dear Sarah:

“The indications are very strong that we shall move in a few days—perhaps tomorrow. Lest I should not be able to write again, I feel impelled to write a few lines that may fall under your eye when I shall be no more.

“I have no misgivings about, or lack of confidence in, the cause in which I am engaged, and my courage does not halt or falter. … I am … perfectly willing … to lay down all my joys in this life, to help maintain this Government. …

“Sarah, my love for you is deathless; it seems to bind me with mighty cables that nothing but Omnipotence could break; and yet my love of Country comes over me like a strong wind and bears me unresistibly on with all these chains to the battle field.

“The memories of the blissful moments I have spent with you come creeping over me, and I feel most gratified to God and to you that I have enjoyed them so long. And hard it is for me to give them up and burn to ashes the hopes of future years, when, God willing, we might still have lived and loved together, and seen our sons grown up to honorable manhood around us. I have, I know, but few and small claims upon Divine Providence, but something whispers to me—perhaps it is the wafted prayer of my little Edgar, that I shall return to my loved ones unharmed. If I do not, my dear Sarah, never forget how much I love you, and when my last breath escapes me on the battle field, it will whisper your name. Forgive [me] my … faults, and the many pains I have caused you. How thoughtless and foolish I have oftentimes been! How gladly would I wash out with my tears every little spot upon your happiness. …

“But, O Sarah! If the dead can come back to this earth and the unseen around those they loved, I shall always be near you; in the gladdest days and in the darkest nights … always, always, and if there be a soft breeze upon your cheek, it shall be my breath, as the cool air fans your throbbing temple, it shall be my spirit passing by. Sarah, do not mourn me dead; think I am gone and wait for thee, for we shall meet again.”

The darkness of death can ever be dispelled by the light of revealed truth. “I am the resurrection, and the life,” spoke the Master. “He that believeth in me, though he were dead, yet shall he live: And whosoever liveth and believeth in me shall never die.”

Added to His own words are those of the angels, spoken to the weeping Mary Magdalene and the other Mary as they approached the tomb to care for the body of their Lord: “Why seek ye the living among the dead? He is not here, but is risen.”

Such is the message of Easter morn. He lives! And because He lives all shall indeed live again. This knowledge provides the peace for loved ones of those whose graves are marked by the crosses of Normandy, those hallowed resting places in Flanders fields where the poppies blow in springtime, and for those who rest in countless other locations, including the depths of the sea. “Oh, sweet the joy this sentence gives: ‘I know that my Redeemer lives!’” In the name of Jesus Christ, amen.

\newpage
\settalk{The Priesthood—A Sacred Trust}
\setspeaker{Thomas S. Monson}
\settalkdate{April 1994}
\talkheading{The Priesthood—A Sacred Trust}{Thomas S. Monson}

What a solemn thought, to contemplate the vast priesthood audience assembled here in the Tabernacle on Temple Square and gathered in hundreds of buildings throughout the world! I sincerely pray for the Spirit of the Lord to guide my remarks this evening.

The presence of those who hold the Aaronic Priesthood brings to mind my own experiences as I graduated from Primary, having memorized the Articles of Faith, and then received the Aaronic Priesthood and the office and calling of a deacon. To pass the sacrament was a privilege, and to gather fast offerings a sacred trust. I was set apart as the secretary of the deacons quorum and, at that moment, felt that boyhood had passed and young manhood had begun.

Can you young men realize the shock I felt, while attending an officers’ meeting of our ward conference, when a member of the stake presidency, after calling upon the priesthood and auxiliary leaders to speak, without warning read my name and office, inviting me to give an account of my stewardship and to express my feelings regarding my calling as secretary of the deacons quorum and thus a ward officer. I don’t recall what I said, but a sense of responsibility engulfed me, never to depart thereafter.

I sincerely hope that each deacon, teacher, and priest is aware of the significance of his priesthood ordination and the privilege which is his to fulfill a vital role in the life of every member through his participation in administration and passing of the sacrament each Sunday.

At the time I held the Aaronic Priesthood, it seemed we always sang the same hymns during the opening exercises of priesthood meeting. They were “Come, All Ye Sons of God,” “Come, All Ye Sons of Zion,” “How Firm a Foundation,” “Israel, Israel, God Is Calling,” and a few others. Our voices were not the best, nor was volume adequate, but we learned the words and remembered the message of each.

I smile when I reflect on an account I heard concerning Brother Thales Smith and his service in a bishopric with Bishop Israel Heaton. Sister Heaton called Brother Smith one Sunday morning and mentioned that her husband was ill and unable to attend priesthood meeting. Brother Smith reported this to the brethren assembled that morning and asked the brother who was to offer the invocation to remember Bishop Israel Heaton in the prayer. Then he announced that the opening hymn would be “Israel, Israel, God Is Calling.” I suppose the smiles outnumbered any frowns. By the way, Bishop Heaton recovered.

The opening exercises of priesthood meeting may be brief but should be held in each ward without fail. It brings to the hearts and souls of all assembled a spirit of unity, the brotherhood of priesthood, and a beautiful reminder of our sacred duties.

All who hold the priesthood have opportunities for service to our Heavenly Father and to His children here on earth. It is contrary to the spirit of service to live selfishly within ourselves and disregard the needs of others. The Lord will guide us and make us equal to the challenges before us. Remember His promise and counsel: “The power and authority of the higher, or Melchizedek Priesthood, is to hold the keys of all the spiritual blessings of the church—

“To have the privilege of receiving the mysteries of the kingdom of heaven, to have the heavens opened unto them, to commune with the general assembly and church of the Firstborn, and to enjoy the communion and presence of God the Father, and Jesus the mediator of the new covenant.”

To merit this blessing, it is necessary for each of us to recall who is the Giver of every gift and the Provider of every blessing. “The worth of souls is great in the sight of God” is not an idle phrase but a heaven-sent declaration for our enlightenment and guidance. We must ever remember who we are and what God expects us to become. This pearl of philosophy is hidden away in the delightful musical Fiddler on the Roof as the peasant father Tevye counsels his growing daughters. Other contemporary plays carry thoughts worthy of our consideration as we prepare for service.

From the production Camelot comes the observation, “Violence is not strength, and compassion is not weakness.” From Shenandoah, “If we don’t try, we don’t do; and if we don’t do, then why are we here?” Eliza Doolittle, the pupil of Professor Henry Higgins in My Fair Lady, observes of Colonel Pickering her philosophy: “The difference between a lady and a flower girl is not how she behaves but how she is treated. I shall always be a flower girl to Professor Higgins because he always treats me as a flower girl and always will. But I know that I shall always be a lady to Colonel Pickering because he always treats me as a lady and always will.” Again from Camelot, King Arthur pleaded with Guinevere, “We must not let our passions destroy our dreams.” The list continues. In reality, each magnificent observation is but a paraphrase of the teachings of our Lord, Jesus Christ. He is our exemplar and our guide. It is in His footsteps we must walk to be successful in our priesthood callings.

May I share with you tonight words of wisdom from fellow servants who labored in the ranks but who have now gone to their eternal reward.

First, from a wise stake president to a young bishop: “The work is all-consuming, but ever realize three guidelines to be a successful bishop: feed the poor, have no favorites, and tolerate no iniquity.” Commenting on this last guideline, President Spencer W. Kimball declared, “When dealing with transgression, apply a bandage large enough to cover the wound—no larger, no smaller.”

Second, prior to the creation of the Toronto Ontario Stake in 1960, Elder ElRay L. Christiansen, then an Assistant to the Council of the Twelve, recounted for the benefit of priesthood leaders a lesson from his own life when he was called to preside over the East Cache Stake in Logan, Utah. He mentioned that he and his counselors met to discuss what the stake members most needed and which principles of the gospel the stake presidency should stress. Their opinions varied from sacrament meeting attendance to observance of the Sabbath day, with a lot of territory in between. At length they agreed that the principle most needed was spirituality. They appreciated the truth found in the observation: When one deals in generalities, he will rarely have a success; but when he deals in specifics, he will rarely have a failure.

The four-year plan of President Christiansen and his counselors was refined in a splendid fashion. Year one: We shall increase the spirituality of the membership of the East Cache Stake by every family having family prayer. Year two: We shall increase the spirituality of the membership of the East Cache Stake by every member attending sacrament meeting weekly. Year three: We shall increase the spirituality of the membership of the East Cache Stake by each member paying an honest tithing. Year four: We shall increase the spirituality of the membership of the East Cache Stake by each member honoring the Sabbath day and keeping it holy. Each was the theme for the entire year; emphasis was given constantly.

After the four-year program was concluded, all four of the specific objectives had been attained, but of even greater significance, the spirituality of the membership of the East Cache Stake had shown marked improvement.

Spirituality is not bestowed simply by wishing; rather, it comes quietly and imperceptibly by serving. The Lord counseled, “Therefore, if ye have desires to serve God ye are called to the work.” Many years ago, while attending a district conference in Ottawa, Canada, I called two men from the small Cornwall Branch to serve in responsible positions in the Lord’s service. I jotted down their heartfelt response and share with you tonight the words of yesteryear. From John Brady: “I have covenanted; I will faithfully serve.” From Walter Danic: “The gospel is the most important thing in my life; I will serve.”

President John Taylor provided rather direct counsel to those of us who hold the priesthood: “If you do not magnify your callings, God will hold you responsible for those whom you might have saved had you done your duty.”

Somehow I feel that if we will always remember who it is we serve, and on whose errand we are, we will draw closer to the source of the inspiration we seek, even our Master and Savior.

President Harold B. Lee had a marked influence on Sister Monson and me and our three children. On rather brief occasions he commented to each of our children in a tone which reflected deep spirituality, genuine interest, and inspired counsel.

Our youngest son, Clark, was about to turn twelve when we chanced to meet Brother Lee in the parking lot of the Church Office Building. He asked Clark how old he was. Clark answered, “Soon to be twelve.”

Came the question: “What happens to you when you turn twelve?”

The response: “I’ll receive the Aaronic Priesthood and be ordained a deacon.”

With a warm smile and the clasp of his hand, Brother Lee said, “Bless you, my boy.”

Our daughter, Ann, as a young teenager was with her mother and me when we encountered Brother Lee, and proper introductions were made. Brother Lee took our daughter’s hand in his and, with a lovely smile, said to her, “You, my dear one, are beautiful inside as well as outside. What a choice young lady you are.”

In a more solemn setting, Brother Lee met me one evening on the steps of the LDS Hospital in Salt Lake City. By appointment we were to give a blessing to my eldest son, Tom, who was then in his later teens. Surgery awaited which could be of a most serious nature. Brother Lee took my hand before we ascended the stairs and, looking me straight in the eye, said, “Tom, there is no place I would rather be at this moment than by your side to participate with you in providing a sacred priesthood blessing to your son.”

We then went to the room, where he said to Tom, “We are about to give you a blessing, even to provide a priesthood ordinance. We approach this privilege in humility, for we remember the counsel of the Prophet Joseph Smith, who said that when those who hold the priesthood place their hands on the head of a person in this sacred ordinance, it is as though the hands of the Lord are placed thereon.” The blessing was given; the surgery turned out to be minor. But lessons were learned, spirituality of a great leader was observed, and a model to follow was provided.

Brethren, there are tens of thousands of priesthood holders scattered among you who, through indifference, hurt feelings, shyness, or weakness, cannot bless to the fullest extent their wives and children—without considering the lives of others they could lift and bless. Ours is the solemn duty to bring about a change, to take such an individual by the hand and help him arise and be well spiritually. As we do so, sweet wives will call our names blessed, and grateful children will marvel at the change in Daddy as lives are altered and souls are saved.

When I visited stake conferences as a member of the Twelve, I always took note of those stakes which had excelled in bringing to activity those brethren whose talents and potential leadership had lain dormant. Inevitably I would ask, “How were you able to achieve success? What did you do and how did you do it?” One such stake was the North Carbon Stake when President Cecil Broadbent presided. Eighty-seven men had been reactivated and, with their wives and children, went to the Manti Temple in the space of one year. President Broadbent, upon hearing my questions, turned to his counselor, President Stanley Judd, a large and good-natured coal miner, and said, “This is President Judd’s responsibility in the stake presidency. He will answer.”

As I restated my questions to President Judd, I concluded with the plea, “Will you tell me how you did it?”

With a smile, he replied, “No.” I was stunned! Then he said, “If I tell you how we did it, then you will tell others, and they will surpass our record.” I was still stunned. Then, with a twinkle in his eye, this wonderful man added, “However, Brother Monson, if you will give me two tickets to general conference, I’ll tell you how we did it.”

The tickets were provided; the success pattern was revealed. However, President Judd felt the contract was open-ended and asked for and received from me two tickets for each conference until he was eventually ordained a patriarch.

The formula was the same, generally speaking, in each successful stake with regard to this phase of the work. It consisted of four ingredients: one, put forth your efforts at the ward level; two, involve the ward bishop; three, provide inspired teaching; and four, do not attempt to concentrate on all the brethren at once; rather, work with a few husbands and their wives at a given time and then have them help you as you work with others.

High-powered sales techniques are not the answer in priesthood leadership; rather, devotion to duty, continuous effort, abundant love, and personal spirituality combine to touch the heart, prompt the change, and bring to the table of the Lord His hungry children who have wandered in the wilderness of the world but who now have returned “home.”

Long years ago I reorganized the Star Valley Wyoming Stake at the time the legendary leader President E. Francis Winters was released. He had served faithfully and with distinction for many years.

The Sabbath day dawned; the members came from far and wide and crowded into the Afton, Wyoming, chapel. Every available space was taken. As the reorganization of the stake presidency was concluded, I did something I had not done before. I felt impressed to conduct a modest exercise, and I asked publicly, “Would all of you who have been given a name or been baptized or confirmed by Francis Winters please stand and remain standing.” Many stood. Then I continued, “Now will all of you who have been ordained or set apart by Francis Winters please stand and remain standing.” Another large number swelled the ranks of those standing. “Finally, will all of you who have received a blessing under the hands of Francis Winters please stand and remain standing.” All the remainder stood.

I turned to President Winters and, with tears coursing down my cheeks, said to him, “President Winters, you see before you the result of your ministry as stake president. The Lord is pleased.” Silence prevailed. Heads nodded their approval as sobs were then heard and handkerchiefs retrieved from every purse and pocket. It was one of the most spiritually rewarding experiences of my life. No one in that vast throng will ever forget how he or she felt at that hour.

After the work of the conference had been concluded, good-byes were said, and I began the drive home. I found myself singing the favorite hymn from the Sunday School days of my youth:

\begin{verse}
Thanks for the Sabbath School. Hail to the day\\
When evil and error are fleeing away.\\
Thanks for our teachers who labor with care\\
That we in the light of the gospel may share. …
\end{verse}

And then I literally boomed the chorus:

\begin{verse}
Join in the jubilee; mingle in song.\\
Join in the joy of the Sabbath School throng.\\
Great be the glory of those who do right,\\
Who overcome evil, in good take delight.
\end{verse}

I was all alone in the car—or was I? The miles hurried by. In silent reverie, I reflected on the events of the conference. Francis Winters, a bookkeeper at the community cheese factory, a man of modest means and humble home, had walked the path that Jesus walked, and like the Master he “went about doing good.” He qualified for the Savior’s description of Nathanael as he approached Him from afar: “Behold an Israelite indeed, in whom is no guile!”

Brethren, my prayer tonight is that all of us, in whatever capacities we serve in the Church, may merit the gentle touch on our shoulder of the Master’s hand and qualify for that same salutation received by Nathanael. That we, at the conclusion of life’s journey, may hear those divinely spoken words, “Well done, thou good and faithful servant” is my prayer in the name of Jesus Christ, amen.

\newpage
\settalk{What He Would Have Us Do}
\setspeaker{Thomas S. Monson}
\settalkdate{April 1994}
\talkheading{What He Would Have Us Do}{Thomas S. Monson}

We have missed Elder Marvin J. Ashton and another familiar soul at our conferences, D. Arthur Haycock, each of whom has passed away since our last conference. Our hearts and our prayers go out to Sister Ashton, Sister Haycock, and all who have loved and lost someone during this period.

President Benson’s chair, situated between President Hinckley and me, has been unoccupied at this conference, although he has viewed the conference proceedings at his apartment by television. Our hearts are full of love for the prophet of God, and his teachings ring in our ears. If he were standing before us at this moment at the conclusion of the conference, I believe he would say, “Lord, it has been good for us to be here.”

We have sustained with our uplifted hands and also with our hearts those called to new positions of responsibility.

As we leave for our homes, may we travel in peace and safety. May we be obedient to the commandments of God. As we reflect upon the messages of conference, we see woven like a beautiful golden thread in a fine tapestry the mission of the Lord Jesus Christ, the sanctity of the home, and the importance of obedience to divine teachings.

I like the thought, “Before Easter, there must be a cross.” And many have heavy crosses to bear. With the birth of the Babe in Bethlehem, there emerged a great endowment—a power stronger than weapons, a wealth more lasting than the coins of Caesar. He may come to us as one unknown, without a name, as by the lakeside He came to those men who knew Him not. He speaks to us the same words, “Follow thou me,” and sets us to the task which He has to fulfill for our time. He commands, and to those who obey Him, whether they be wise or simple, He will reveal Himself in the toils, the conflicts, the sufferings that they shall pass through in His fellowship; and they shall learn in their own experience who He is.

May we praise His name, follow His example, and incorporate His truths into our lives, and then this conference will have been successful. That such may be our experience, I pray in His worthy name—even Jesus Christ—amen.

\newpage
\settalk{My Brother’s Keeper}
\setspeaker{Thomas S. Monson}
\settalkdate{October 1994}
\talkheading{My Brother’s Keeper}{Thomas S. Monson}

My dear brethren, I am confident that you, as I, have seen the newscasts on television and have heard them on radio, have read feature articles published by weekly and monthly magazines, and have observed the glaring headlines in daily newspapers. They all describe the fighting in Bosnia, tribal conflicts in Africa, and extensive flooding in Georgia and Florida. The parade of devastation, loss of homes, damage to farms, ruination of businesses, and, above all, frightful human suffering and death continues almost without interruption.

After expressions of sorrow, the shaking of one’s head in incredible disbelief, and, yes, even the wringing of the hands in frustration, the question is asked, “When are they going to do something about this terrible suffering?”

Long years ago a similar question was posed and preserved in holy writ, even the Bible: “And Cain talked with Abel his brother: and it came to pass, when they were in the field, that Cain rose up against Abel his brother, and slew him.

“And the Lord said unto Cain, Where is Abel thy brother? And he said, I know not: Am I my brother’s keeper?”

This evening I felt to present to you a response to this question which represents a collective reply from Church members everywhere and from the Church itself. But first a brief background.

In March of 1967, early in my service as a member of the Council of the Twelve, I was attending a conference of the Monument Park West Stake in Salt Lake City. My companion for the conference was a member of the General Church Welfare Committee, Paul C. Child. President Child was a student of the scriptures. He had been my stake president during my Aaronic Priesthood years. Now we were together as conference visitors.

When it was his opportunity to participate, President Child took in hand the Doctrine and Covenants and left the pulpit to stand among the priesthood brethren to whom he was directing his message. He turned to section 18 and began to read:

“Remember the worth of souls is great in the sight of God. …

“And if it so be that you should labor all your days in crying repentance unto this people, and bring, save it be one soul unto me, how great shall be your joy with him in the kingdom of my Father!”

President Child then raised his eyes from the scriptures and asked the brethren, “What is the worth of a human soul?” He avoided calling on a bishop, a stake president, or a high councilor for a response. Instead he selected the president of an elders quorum, a brother who had been a bit drowsy and had missed the significance of the question.

The startled man responded, “Brother Child, could you please repeat the question?”

The question was repeated: “What is the worth of a human soul?”

I knew President Child’s style. I prayed fervently for that quorum president. He remained silent for what seemed like an eternity and then declared, “Brother Child, the worth of a human soul is its capacity to become as God.”

All present pondered that reply. Brother Child returned to the stand, leaned over to me, and said, “A profound reply; a profound reply!” He proceeded with his message, but I continued to reflect on that inspired response.

Another pioneer in Church welfare, Walter Stover, who died some months ago at the same age as President Ezra Taft Benson, was one who understood the worth of a human soul. At his funeral service this tribute was paid to Brother Stover: “He had the ability to see Christ in every face he encountered, and he acted accordingly. Legendary are his acts of compassionate help and his talent to lift heavenward every person whom he met. His guiding light was the Master’s voice speaking, ‘Inasmuch as ye have done it unto one of the least of these … , ye have done it unto me.’”

The publication Times and Seasons, in its March 1842 issue, proclaimed the following: “Respecting how much a man … shall give … we have no special instructions … ; he is to feed the hungry, to clothe the naked, to provide for the widow, to dry up the tear of the orphan, to comfort the afflicted, whether in this church, or in any other [church], or in no church at all, wherever he finds them.”

Since the two special fast days in 1985, called for by the First Presidency, humanitarian efforts by members of The Church of Jesus Christ of Latter-day Saints have reached into every corner of the globe. Millions of the earth’s needy have been blessed as members of the Church have consecrated their means to provide life-sustaining food and clothing, establish immunization and infant feeding programs, teach basic literacy, dig freshwater wells, foster village banks, create new jobs, sustain hospitals and orphanages, teach basic self-reliance, and act in many other ways to help Heavenly Father’s children improve their lives both spiritually and temporally.

The scope of humanitarian aid given is dramatic:

All of the foregoing is in addition to the conventional welfare program of the Church, fundamentally financed through regular fast-offering contributions.

The examples of humanitarian aid and on-the-scene testimonials are inspiring and heartwarming.

Following its colonial period, a series of tribal conflicts has decimated the population of Rwanda in Africa. In the spring of this year, open hostilities resumed, resulting in the deaths of more than half a million people. Refugees huddle in squalid and unhealthy camps within the borders of neighboring Zaire, Uganda, Tanzania, and Burundi.

Joining with the efforts of other agencies in the international community, this church has committed \$1.2 million in goods and cash for refugee relief. Most of the promised assistance has already been consigned or shipped through four helping agencies—even Catholic Relief Services, the International Committee of the Red Cross, C.A.R.E., and the United Nations High Commissioner for Refugees. Continuing efforts are planned to help stem the tide of pain among these children of our Heavenly Father.

In Yugoslavia, following the demise of the former government, the country disintegrated into ethnic factions. The resulting civil conflict has claimed thousands of lives and inflicted hardship, heartache, and suffering upon millions.

Working with seven different humanitarian agencies, the Church has provided, since 1991, food, clothing, blankets, hygiene kits, and medical supplies valued at \$850,000. This is in addition to personal contributions by our members in other European nations.

In May 1993 Danijela Curcic of Zagreb, Croatia, wrote this letter addressed to Church headquarters, expressing her gratitude for food shared by the Saints.

“Dear Charitable Persons,

“I would like to thank you for every good thing that you’ve done for the people in my country. This horrible civil war is a crime which doesn’t spare anything and anybody. Uncounted numbers of refugees, thousands of dead children are about us everywhere. I respect with all my heart you dear friends because you showed you care. It’s easier and doesn’t hurt as much when you’re aware of the fact that there are nice people who are willing to help you.”

Closer to home, but serviced by conventional welfare procedures, are the victims of the devastating south Georgia flood of 1994. Thirty-five thousand families were evacuated, five thousand people found temporary refuge in two of our chapels, and nine eighteen-wheel truckloads of food and supplies were provided by the Church, primarily to other than members of our church.

Our own Church spearhead unit, carrying emergency welfare supplies, was on site with everything requested just five hours after being activated by the Area President.

On the first weekend of the flood, 500 member volunteers assisted in the cleanup of 1,569 damaged houses. The next weekend, more than 5,500 volunteers arrived and helped—all from units of the Church from a wide area well beyond the stricken region.

Priesthood volunteers from the Jacksonville Florida West Stake worked all weekend cleaning up a house which had been nearly submerged by the flood. The owner, a retired nonmember named Davis, was overwhelmed by the help provided. When the work was completed, the brethren asked Mr. Davis if they could bless his house. They gathered together, and the bishop pronounced a blessing on the home and on the family. Tears ran down Mr. Davis’s cheeks, and the Spirit was very strong. Each of the volunteers hugged him and told him how glad they were to have been of help. He said they had done more than they could ever know and that he didn’t know how to thank them enough.

The response of the membership of the Church, and particularly the priesthood performance in such situations, touches the heart and is a marvel to behold. Thus it has ever been.

From an earlier period, following the carnage of World War II, Elder Ezra Taft Benson led Church response in providing food, medicine, and clothing—totaling two million in 1940s dollars and requiring 133 boxcars to transport it—to the cold and starving members in Europe. This desperately needed aid saved lives, rescued the dispirited, and brought a newness of hope and quickened prayers of thanksgiving and expressions of profound gratitude from one and all. “Charity never faileth.”

During a drive to amass warm clothing to ship to suffering Saints, Elder Harold B. Lee and Elder Marion G. Romney took President George Albert Smith to Welfare Square in Salt Lake City. They were impressed by the generous response of the membership of the Church to the clothing drive and the preparations for sending the goods overseas. They watched President Smith observing the workers as they packaged this great volume of donated clothing and shoes. They saw tears running down his face. After a few moments, President Smith removed a new overcoat that he had on and said, “Please ship this also.”

The Brethren said to him, “No, President, no; don’t send that; it’s cold and you need your coat.”

But President Smith would not take it back.

The Apostle Paul’s admonition surely was fulfilled that day: “Be thou an example of the believers, in word, in conversation, in charity, in spirit, in faith, in purity.”

Two weeks ago Elder Dallin H. Oaks, Elder Robert K. Dellenbach, and I attended a regional conference in Holland. While meeting with the Saints, I recalled the miracle of the potatoes which took place in that nation in November of 1947.

In the first week of November 1947, ten huge trucks moved across Holland. They headed east and contained a costly cargo—seventy-five tons of potatoes, a gift from the Dutch Church members to the Saints in Germany.

Many months earlier, in the spring of 1947, the members within the Netherlands Mission were asked to begin a welfare project of their own, now that they had received much needed welfare supplies from the members in America. The proposal was welcomed with enthusiasm. The priesthood went to work, and within a short time every quorum had found a suitable piece of land for the project. The recommended crop: potatoes. At the various branches of the Church there was singing, speaking, and praying, at the end of which the potatoes were entrusted to the soil. Soon there came news of good prospects for the harvest, and cautious estimates were made as to how large the yield would be.

During the time the potatoes were growing, Walter Stover, president of the East German Mission, visited the Netherlands Mission in Holland. During his visit, with tears in his eyes, he told of the hunger of the Church members in Germany. They were in worse condition than the Saints in the Netherlands. Supplies had not yet reached the Saints in Germany as quickly as they had the Saints in Holland.

When Cornelius Zappey, the Netherlands Mission president, heard the condition of the German Saints, he couldn’t help but have compassion toward them, knowing how they had suffered. The thought came; the action followed: “Let’s give our potatoes to the members of the Church in Germany.” I’m sure he worried, for the German armies and the Dutch armies had been in conflict with each other. The Dutch had been starving. Would they respond? A Dutch widow who had received a sack of the potatoes heard that the bulk of the potatoes was to be given to the members in Germany, and she stepped forward and said, “My potatoes must be with them.” And this hungry widow returned her sack of potatoes.

What are the words of the Lord pertaining to such an act? “Verily I say unto you, That this poor widow hath cast more in, than all they which have cast into the treasury. … She of her want did cast in all that she had.”

It was President J. Reuben Clark, Jr., who in 1936 declared: “The real long term objective of the Welfare Plan is the building of character in the members of the Church, givers and receivers, rescuing all that is finest down deep inside of them, and bringing to flower and fruitage the latent richness of the spirit, which after all is the mission and purpose and reason for being of this Church.”

“Am I my brother’s keeper?” This ageless question has been answered! From the psalm of David comes the precious promise:

“Blessed is he that considereth the poor: the Lord will deliver him in time of trouble.

“The Lord will preserve him, and keep him alive; and he shall be blessed upon the earth: and thou wilt not deliver him unto the will of his enemies.

“The Lord will strengthen him.”

Brethren, may the Lord strengthen each of us who holds the priesthood, that each may learn his duty as his brother’s keeper and be found on the Lord’s errand, I humbly pray in the name of Jesus Christ, amen.

\newpage
\settalk{The Fatherless and the Widows—Beloved of God}
\setspeaker{Thomas S. Monson}
\settalkdate{October 1994}
\talkheading{The Fatherless and the Widows—Beloved of God}{Thomas S. Monson}

Many years ago I attended a large gathering of Church members in the city of Berlin, Germany. A spirit of quiet reverence permeated the gathering as an organ prelude of hymns was played. I gazed at those who sat before me. There were mothers and fathers and relatively few children. The majority of those who sat on crowded benches were women about middle age—and alone. Suddenly it dawned on me that perhaps these were widows, having lost their husbands during World War II. My curiosity demanded an answer to my unexpressed thought, so I asked the conducting officer to take a sort of standing roll call. When he asked all those who were widows to please arise, it seemed that half the vast throng stood. Their faces reflected the grim effect of war’s cruelty. Their hopes had been shattered, their lives altered, and their future had in a way been taken from them. Behind each countenance was a personal travail of tears. I addressed my remarks to them and to all who have loved, then lost, those most dear.

Frederick W. Babbel, who accompanied Elder Ezra Taft Benson on his postwar visit to Europe to assist the struggling Saints, recounts in his book On Wings of Faith one heartrending account. A woman, even the mother of four small children, had been newly widowed. Her husband, young and handsome, whom she loved more than life itself, had been killed during the final days of the frightful battles in their homeland of East Prussia. She and her children were forced to flee to West Germany, a distance of a thousand miles. The weather was mild as they began their long and difficult trek on foot. Constantly being faced with dangers from panicky refugees and marauding troops was difficult enough, but then came the cold of winter, with its accompanying snow and ice. Her resources were meager; now they were gone. All she had was her strong faith in God and in the gospel as revealed to the latter-day prophet Joseph Smith.

And then one morning the unthinkable happened. She awakened with a chill in her heart. The tiny form of her three-year-old daughter was cold and still, and she realized that death had claimed her. With great effort the mother prepared a shallow grave and buried her precious child.

Death, however, was to be her companion again and again on the journey. Her seven-year-old perished, and then her five-year-old. Her despair was all-consuming. Finally, as she was reaching the end of her travel, the baby died in her arms. She had lost her husband and all her children. She had given up all her earthly goods, her home, and even her homeland.

From the depths of her despair, she knelt and prayed more fervently than she had ever prayed in her life: “Dear Heavenly Father, I do not know how I can go on. I have nothing left—except my faith in thee. I feel amidst the desolation of my soul an overwhelming gratitude for the atoning sacrifice of thy Son, Jesus Christ. I know that because he suffered and died, I shall live again with my family; that because he broke the chains of death, I shall see my children again in the flesh and will have the joy of raising them. Though I do not at this moment wish to live, I will do so, that we may be reunited as a family and return, together, to thee.” This prayer, this testimony sustained her until finally she reached Karlsruhe, her destination.

Though perhaps not so cruel and dramatic, yet equally poignant, are the lives described in the obituaries of our day and time when the uninvited enemy called death enters the stage of our mortal existence and snatches from our grasp a loving husband or precious wife and frequently, in the young exuberance of life, our children and grandchildren. Death shows no mercy. Death is no respecter of persons, but in its insidious way it visits all. At times it is after long-suffering and is a blessing, while in other instances those in the prime of life are taken by its grasp.

As of old, the heartbroken frequently and silently repeat the ancient question: “Is there no balm in Gilead?” “Why me; why now?” The words of a beautiful hymn provide a partial answer:

\begin{verse}
Where can I turn for peace? Where is my solace\\
When other sources cease to make me whole?\\
When with a wounded heart, anger, or malice,\\
I draw myself apart, Searching my soul? …
\end{verse}

The plight of the widow is a recurring theme through holy writ. Our hearts go out to the widow at Zarephath. Gone was her husband. Consumed was her scant supply of food. Starvation and death awaited. But then came God’s prophet with the seemingly brazen command that the widow woman should feed him. Her response is particularly touching: “As the Lord thy God liveth, I have not a cake, but an handful of meal in a barrel, and a little oil in a cruse: and, behold, I am gathering two sticks, that I may go in and dress it for me and my son, that we may eat it, and die.”

The reassuring words of Elijah penetrated her very being:

“Fear not; go and do as thou hast said: but make me thereof a little cake first, and bring it unto me, and after make for thee and for thy son.

“For thus saith the Lord God of Israel, The barrel of meal shall not waste, neither shall the cruse of oil fail. …

“And she went and did according to the saying of Elijah. …

“And the barrel of meal wasted not, neither did the cruse of oil fail.”

Like the widow at Zarephath was the widow of Nain. The New Testament of our Lord records a moving and soul-stirring account of the Master’s tender regard for the grieving widow:

“And it came to pass … that he went into a city called Nain; and many of his disciples went with him, and much people.

“Now when he came nigh to the gate of the city, behold, there was a dead man carried out, the only son of his mother, and she was a widow: and much people of the city was with her.

“And when the Lord saw her, he had compassion on her, and said unto her, Weep not.

“And he came and touched the bier: and they that bare him stood still. And he said, Young man, I say unto thee, Arise.

“And he that was dead sat up, and began to speak. And he delivered him to his mother.”

What power, what tenderness, what compassion did our Master and Exemplar demonstrate. We too can bless if we will but follow his noble example. Opportunities are everywhere. Needed are eyes to see the pitiable plight, ears to hear the silent pleadings of a broken heart, yes, and a soul filled with compassion, that we might communicate not only eye to eye or voice to ear but in the majestic style of the Savior, even heart to heart.

The word widow appears to have had a most significant meaning to our Lord. He cautioned his disciples to beware the example of the scribes, who feigned righteousness by their long apparel and their lengthy prayers, but who devoured the houses of widows.

To the Nephites came the direct warning, “I will come near to you to judgment; and I will be a swift witness against … those that oppress … the widow.”

And to the Prophet Joseph Smith he directed, “The storehouse shall be kept by the consecrations of the church; and widows and orphans shall be provided for, as also the poor.”

The widow’s home is generally not large or ornate. Frequently it is a modest one in size and humble in appearance. Often it is tucked away at the top of the stairs or the back of the hallway and consists of but one room. To such homes he sends you and me.

There may exist an actual need for food, clothing, even shelter. Such can be supplied. Almost always there remains the hope for that special hyacinth to feed the soul.

\begin{verse}
Go, gladden the lonely, the dreary;\\
Go, comfort the weeping, the weary;\\
Go, scatter kind deeds on your way;\\
Oh, make the world brighter today!
\end{verse}

Let us remember that after the funeral flowers fade, the well-wishes of friends become memories, and the prayers offered and words spoken dim in the corridors of the mind. Those who grieve frequently find themselves alone. Missed is the laughter of children, the commotion of teenagers, and the tender, loving concern of a departed companion. The clock ticks more loudly, time passes more slowly, and four walls do indeed a prison make.

Hopefully all of us may again hear the echo of words spoken by the Master, inspiring us to good deeds: “Inasmuch as ye have done it unto one of the least of these … , ye have done it unto me.”

The late Elder Richard L. Evans left for our contemplation and action this admonition:

“We who are [young] should never become so blindly absorbed in our own pursuits as to forget that there are still with us those who will live in loneliness unless we let them share our lives as once they let us share theirs. …

“We cannot bring them back the morning hours of youth. But we can help them live in the warm glow of a sunset made more beautiful by our thoughtfulness, by our provision, and by our active and unfeigned love. Life in its fulness is a loving ministry of service from generation to generation. God grant that those who belong to us may never be left in loneliness.”

Long years ago a severe drought struck the Salt Lake Valley. The commodities at the storehouse on Welfare Square had not been of their usual quality, nor were they found in abundance. Many products were missing, especially fresh fruit. As I was a young bishop, worrying about the needs of the many widows in my ward, my prayer one evening is especially sacred to me. I pleaded for these widows, who were among the finest women I knew in mortality and whose needs were simple and conservative, because they had no resources on which they might rely.

The next morning I received a call from a ward member, a proprietor of a produce business situated in our ward. “Bishop,” he said, “I would like to send a semitrailer filled with oranges, grapefruit, and bananas to the bishops’ storehouse to be given to those in need. Could you make arrangements?” Could I make arrangements! The storehouse was alerted, and then each bishop was telephoned and the entire shipment distributed. Bishop Jesse M. Drury, that beloved welfare pioneer and storekeeper, said he had never witnessed a day like it before. He described the occasion with one word: “Wonderful!”

The wife of that generous businessman is today a widow. I know the decision her husband and she made has brought her sweet memories and comforting peace to her soul.

I express my sincere appreciation to one and all who are mindful of the widow. To the thoughtful neighbors who invite a widow to dinner and to that royal army of noble women, the visiting teachers of the Relief Society, I add, may God bless you for your kindness and your love unfeigned toward her who reaches out and touches vanished hands and listens to voices forever stilled. The words of the Prophet Joseph Smith describe their mission: “I attended by request, the Female Relief Society, whose object is the relief of the poor, the destitute, the widow and the orphan, and for the exercise of all benevolent purposes.”

Thank you to thoughtful and caring bishops who ensure that no widow’s cupboard is empty, no house unwarmed, no life unblessed. I admire the ward leaders who invite the widows to all social activities, often providing a young Aaronic Priesthood lad to be a special escort for the occasion.

Frequently the need of the widow is not one of food or shelter but of feeling a part of ongoing events. President Bryan Richards of Salt Lake City, now serving as a mission president, brought to my office a sweet widow whose husband had passed away during a full-time mission they were serving. President Richards explained that her financial resources were adequate and that she desired to contribute to the Church’s General Missionary Fund the proceeds of two insurance policies on the life of her departed husband. I could not restrain my tears when she meekly advised me, “This is what I wish to do. It is what my missionary-minded husband would like.”

The gift was received and entered as a most substantial donation to missionary service. I saw the receipt made in her name, but I believe in my heart it was also recorded in heaven. I invited her and President Richards to follow me to the unoccupied First Presidency Council room in the Church Administration Building. The room is beautiful and peaceful. I asked this sweet widow to sit in the chair usually occupied by our church President. I felt he would not mind, for I knew his heart. As she sat ever so humbly in the large leather chair, she gripped each armrest with a hand and declared, “This is one of the happiest days of my life.” It was also such for President Richards and for me.

I never travel to work along busy Seventh East in Salt Lake City but what I see in my mind’s eye a thoughtful daughter, afflicted with arthritis and carrying in her hand a plate of warm food to her aged mother, who lived across the busy thoroughfare. She has now gone home to that mother who preceded her in passing. But her lesson was not lost on her daughters, who delight their widowed father by cleaning his house each week, inviting him to dinners in their homes, and sharing with him the laughter of good times together, leaving in that widower’s heart a prayer of gratitude for his daughters, the light of his life. Fathers experience loneliness as well as mothers.

One evening at Christmastime, my wife and I visited a nursing home in Salt Lake City. We looked in vain for a 95-year-old widow whose memory had become clouded and who could not speak a word. An attendant led us in our search, and we found Nell in the dining room. She had eaten her meal; she was sitting silently, staring into space. She did not show us any sign of recognition. As I reached to take her hand, she withdrew it. I noticed that she held firmly to a Christmas greeting card. The attendant smiled and said, “I don’t know who sent that card, but she will not lay it aside. She doesn’t speak, but pats the card and holds it to her lips and kisses it.” I recognized the card. It was one my wife, Frances, had sent to Nell the week before. We left Maytime Manor more filled with the Christmas spirit than when we entered. We kept to ourselves the mystery of that special card and the life it had gladdened and the heart it had touched. Heaven was nearby.

We need not wait for Christmas; we need not postpone till Thanksgiving Day our response to the Savior’s tender admonition, “Go, and do thou likewise.”

As we follow in his footsteps, as we ponder his thoughts and his deeds, as we keep his commandments, we will be blessed. The grieving widow, the fatherless child, and the lonely of heart everywhere will be gladdened, comforted, and sustained through our service, and we will experience a deeper understanding of the words recorded in the Epistle of James: “Pure religion and undefiled before God and the Father is this, To visit the fatherless and widows in their affliction, and to keep himself unspotted from the world.”

May the peace promised by the Savior be the gift of one and all this Sabbath day and always is my fervent and humble prayer, in the name of Jesus Christ, amen.

\newpage
\settalk{A Time to Choose}
\setspeaker{Thomas S. Monson}
\settalkdate{April 1995}
\talkheading{A Time to Choose}{Thomas S. Monson}

What an imposing and inspiring throng are you young women, your mothers, and your leaders as you meet together this evening on this sacred occasion. All of us have been uplifted by the beauty—even the majesty—portrayed in what we have heard, what we have seen, and what we have felt during this conference. I seek the help of our Heavenly Father as I respond briefly to the privilege to speak to you.

I know it is important for me to keep in mind your perspective. This truth I learned from a granddaughter. I was speaking to her family about the importance of having sufficient numbers of young men and young women in a ward to maximize social opportunities and to learn together the principles of the gospel. I commented, “Why, do you know that when Elder Joseph Wirthlin was a bishop here in Salt Lake City, he had a full quorum of forty-eight boys who were priests.”

My granddaughter, who had been listening but saying little, suddenly exclaimed, “Oh, that would be wonderful!”

I came to appreciate the importance of having the right perspective. It has been said that the young want to change the world—and the old want to change the young!

Oh, the dreams of our youth! How beautiful and how perishable they are. Today, however, some youth are drifting on a sea of chance, with waves of temptation threatening to engulf them. A nationally prominent journalist described situations in our country by saying, “They are indicative of the days in which we are living: … days of compromise and diluting of principles; days when sin is labeled as error, when morality is relative and when materialism emphasizes the value of expedience and the shirking of responsibility.”

You young women ask silently, “What can I do to insure my eternal joy? Can you help me?” I offer four suggestions:

First, study diligently. All that has been said this evening points to the holy scriptures as an unfailing guide in our lives. Become acquainted with the lessons the scriptures teach. Learn the background and setting of the Master’s parables and the prophets’ admonitions. Study them as though each were speaking to you, for such is the truth.

For example, let us hearken to the gentle yet persuasive appeal of the Apostle Paul as he counsels his young friend Timothy: “Let no man despise thy youth; but be thou an example of the believers, in word, in conversation, in charity, in spirit, in faith, in purity.”

Crash courses on scripture study are not nearly so effective as the day-to-day reading and application of the scriptures in our lives. Also, there are lessons to be learned when we study good literature. One of the most popular musicals of our time is Fiddler on the Roof, by Joseph Stein.

The gaiety of the dance, the rhythm of the music, the excellence of the acting all fade in significance when Tevye, the father, speaks what to me becomes the message of the musical. He gathers his lovely daughters to his side and, in the simplicity of his peasant surroundings, counsels them as they prepare for their future. “Remember, in Anatevka each one of you knows who she is and what God expects her to become.”

Contemplating our earthly life, could not we well consider Tevye’s statement and respond, “Here, each one of you knows who she is and what God expects her to become.” Study diligently.

Second, choose carefully. All of you commenced an awesome and vital undertaking when you left the spirit world and entered the stage of mortality. Loving parents made you welcome. Inspired teachers taught you truth. True friends provided counsel. Yet life’s choices remain for each one to make. No choice is insignificant, for we become what we think about. Our choices determine our destiny.

Several years ago I held in my hand such a guide to choice. It was a volume of scripture we commonly call the Triple Combination, containing the Book of Mormon, the Doctrine and Covenants, and the Pearl of Great Price. The book was a gift from a loving father to a beautiful, blossoming daughter who followed carefully his advice. On the flyleaf page of the book, her father had written in his own hand these inspired words:

To My Dear Maurine,

That you may have a constant measure by which to judge between truth and the errors of man’s philosophies, and thus grow in spirituality as you increase in knowledge, I give you this sacred book to read frequently and cherish throughout your life.

Lovingly your father,

/s/ Harold B. Lee

Young women, choose carefully your friends, for they help to determine your future. Choose to honor your father and your mother, as Heavenly Father would have you do. They love you and would not knowingly lead you astray.

In Lewis Carroll’s delightful classic Alice’s Adventures in Wonderland, Alice finds herself coming to a crossroads with two paths before her, each stretching onward but in opposite directions. She is confronted by the Cheshire cat, of whom Alice asks, “Which path should I take?”

The cat answers, “That depends where you want to go. If you do not know where you want to go, it doesn’t really matter which path you take.”

Unlike Alice, each of you knows where you want to go. It does matter which way we go, for the path we follow in this life surely leads to the path we shall follow in the next. Choose carefully.

Third, pray fervently. Each of you is a daughter of God, created in His image. Yours is a celestial journey. Heavenly Father wants you to check in with Him through sincere and fervent prayer. Remember, you are never alone. Never forget that you are loved. Never doubt that someone surely cares for you.

Your challenges are real, your concerns important, your need for answers vital. Youth should become acquainted with the ninth section of the Doctrine and Covenants. It has a lesson for each of you. When you contemplate a decision, go to your Heavenly Father in the manner in which the Prophet Joseph indicated that the Lord advised him. The Lord said to the Prophet Joseph: “Behold, I say unto you, that you must study it out in your mind; then you must ask me if it be right, and if it is right I will cause that your bosom shall burn within you; therefore, you shall feel that it is right.” The Lord continues, “But if it be not right you shall have no such feelings, but you shall have a stupor of thought that shall cause you to forget the thing which is wrong.”

That counsel will guide you. It has guided me. Pray fervently.

Fourth and finally, act wisely. Take the Lord as your guide. Do not lend a listening ear to the persuasive voice of that evil one who would entice you to depart from your standards, your home-inspired teachings, and your philosophy of life. Rather, remember that gentle and ever genuine invitation from the Redeemer, “Come, follow me.” Follow Him, and you will be acting wisely and will be blessed eternally.

Along your pathway of life you will observe that you are not the only traveler. There are others who need your help. There are feet to steady, hands to grasp, minds to encourage, hearts to inspire, and souls to save.

Recently I saw a young teen-aged friend, Jami Palmer, whom I have known for several years. She is recovering from cancer. She has endured the diagnosis. She has undergone surgery and painful chemotherapy. Today she is a bright, beautiful young lady and looking to the future with confidence and with faith. I learned that in her darkest hour, when any future appeared somewhat grim, her leg where the cancer was situated would require multiple surgeries. A long-planned hike with her Young Women class up to Timpanogos Cave was out of the question—she thought. Jami told her friends they would have to undertake the hike without her. I’m confident there was a catch in her voice and disappointment in her heart. But then the other young women responded emphatically, “No, Jami, you are going with us!”

“But I can’t walk,” came the anguished reply.

“Then, Jami, we’ll carry you to the top!” And they did.

Today, the hike is a memory, but in reality it is much more. James Barrie, the Scottish poet, declared: “God gave us memories, that we might have June roses in the December of our lives.” None of those precious young women will ever forget that memorable day when, I am confident, a loving Heavenly Father looked down with a smile of approval and was well pleased.

That our Heavenly Father may ever bless you precious young women, that He may inspire your dear mothers and guide your teachers and watch over you always, is my sincere prayer. In the name of Jesus Christ, amen.

\newpage
\settalk{Mercy—The Divine Gift}
\setspeaker{Thomas S. Monson}
\settalkdate{April 1995}
\talkheading{Mercy—The Divine Gift}{Thomas S. Monson}

Not long ago, I read a lengthy report concerning the violence and bloodshed that continue to stalk the land of what was once Yugoslavia. The killing and maiming seem to go on despite the efforts put forth to bring peace. The account of a sniper taking deadly aim and snuffing out the life of a small child brought sorrow to my soul. I silently asked, Where to be found is that divine attribute of mercy?

The cruelty of war seems to bring forth hatred toward others and disregard for human life. It has ever been so. Yet in such degradation at times there shines forth the inextinguishable light of mercy.

During the television documentaries shown throughout the fiftieth anniversary of the D-Day invasion of Normandy, the terrible toll in human life was graphically illustrated, and gripping firsthand experiences of soldiers who were there were shared. I particularly remember the comments of an American infantryman who told that, after a day of ferocious fighting, he glanced up from his shallow foxhole to see an enemy soldier with his gun barrel leveled at the American’s heart. Said the infantryman: “I felt I was soon to cross over that bridge of death which leads to eternity. Incredibly my enemy, in broken English, said to me, ‘Soldier, for you this war is over!’ He took me prisoner and thus saved my life. Such mercy I shall remember forever.”

At an earlier time and in a different conflict—namely the American Civil War—a historically documented account illustrates courage, coupled with mercy.

From December 11 to 13, 1862, the Union forces attacked Marye’s Heights, a large hill overlooking the town of Fredericksburg, Virginia, where six thousand Rebels awaited them. The Southern troops were in secure defensive positions behind a stone wall which meandered along the foot of the hill. In addition, they stood four deep on a sunken road behind the wall, out of sight of Union forces.

The Union troops—over forty thousand strong—launched a series of suicidal attacks across open ground. They were mowed down by a scythe of shot; none got closer than forty yards from the stone wall.

Soon the ground in front of the Confederate positions was littered with hundreds, then thousands, of fallen Union soldiers in their blue uniforms—over twelve thousand before sunset. Crying for help, the wounded lay in the bitter cold throughout that terrible night.

The next day, a Sunday, dawned cold and foggy. As the morning fog lifted, the agonized cries of the wounded could still be heard. Finally, a young Confederate soldier, a nineteen-year-old sergeant, had had all he could take. The young man’s name was Richard Rowland Kirkland. To his commanding officer, Kirkland exclaimed, “All night and all day I have heard those poor people crying for water, and I can stand it no longer. I … ask permission to go and give them water.” His request was initially denied on the grounds that it was too dangerous. Finally, however, permission was granted, and soon thousands of amazed men on both sides saw the young soldier, with several canteens draped around his neck, climb over the wall and walk to the nearest wounded Union soldier. He raised the stricken man’s head, gently gave him a drink, and covered him with his own overcoat. Then he moved to the next of the wounded—and the next and the next. As Kirkland’s purpose became clear, fresh cries of “Water, water, for God’s sake, water!” arose all over the field.

The Union soldiers were at first too surprised to shoot. Soon they began to cheer the young Southerner as they saw what he was doing. For more than an hour and a half, Sergeant Kirkland continued his work of mercy.

Tragically, Richard Kirkland was himself killed a few months later at the battle of Chicamauga. His last words to his companions were, “Save yourselves, and tell my pa I died right.”

Kirkland’s Christlike compassion made his name synonymous with mercy for a post–Civil War generation, both North and South. He became known by soldiers on both sides of the conflict as “the angel of Marye’s Heights.” His loving errand of mercy is commemorated by a bronze monument which stands today in front of the stone wall at Fredericksburg. It depicts Sergeant Kirkland lifting the head of a wounded Union soldier to give him a drink of refreshing water. A tablet to Kirkland’s honor hangs in the Episcopal church in Gettysburg, Pennsylvania. With simple eloquence, it captures the essence of the young soldier’s mission of mercy. It reads: “A hero of benevolence, at the risk of his own life, he gave his enemy drink at Fredericksburg.”

The words of William Shakespeare describe Kirkland’s deed:

\begin{verse}
The quality of mercy is not strain’d;\\
It droppeth as the gentle rain from heaven\\
Upon the place beneath. It is twice blest:\\
It blesseth him that gives and him that takes: …\\
It is an attribute to God himself.
\end{verse}

Two brilliant and faith-filled counselors to President David O. McKay spoke to us everlasting counsel concerning the greatest act of mercy ever known to man. President Stephen L Richards said, “The Savior himself declared that he came to fulfill the law, not to do away with it, but with the law he brought the principle of mercy to temper its enforcement, and to bring hope and encouragement to [the] offenders for forgiveness through [mercy and] repentance.”

President J. Reuben Clark, Jr., testified: “You know, I believe that the Lord will help us. I believe if we go to him, he will give us wisdom, if we are living righteously. I believe he will answer our prayers. I believe that our Heavenly Father wants to save every one of his children. I do not think he intends to shut any of us off because of some slight transgression, some slight failure to observe some rule or regulation. There are the great elementals that we must observe, but he is not going to be captious about the lesser things.

“I believe that his juridical concept of his dealings with his children could be expressed in this way: I believe that in his justice and mercy he will give us the maximum reward for our acts, give us all that he can give, and in the reverse, I believe that he will impose upon us the minimum penalty which it is possible for him to impose.”

“I often think that one of the most beautiful things in the Christ’s life was his words on the cross, when, suffering under the agony of a death that is said to have been the most painful that the ancients could devise, death on the cross, after he had been unjustly, illegally, contrary to all the rules of mercy, condemned and then crucified, when he had been nailed to the cross and was about to give up his life, he said to his Father in heaven, as those who were within hearing testify: ‘… Father, forgive them; for they know not what they do.’ (Luke 23:34.)”

In the Book of Mormon, Alma describes beautifully the foregoing with his words: “The plan of mercy could not be brought about except an atonement should be made; therefore God himself atoneth for the sins of the world, to bring about the plan of mercy, to appease the demands of justice, that God might be a perfect, just God, and a merciful God also.”

From the springboard of such knowledge we ask ourselves, Why, then, do we see on every side those instances where people decline to forgive one another and show forth the cleansing act of mercy and forgiveness? What blocks the way for such healing balm to cleanse human wounds? Is it stubbornness? Could it be pride? Maybe hatred has yet to melt and disappear. “Blame keeps wounds open. … Only forgiveness heals!”

Recently I read where an elderly man disclosed at the funeral of his brother, with whom he had shared, from early manhood, a small, one-room cabin near Canisteo, New York, that following a quarrel, they had divided the room in half with a chalk line and neither had crossed the line nor spoken a word to the other since that day—sixty-two years before! What a human tragedy—all for the want of mercy and forgiveness.

At times the need for mercy can be found close to home and in simple settings. We have a four-year-old grandson named Jeffrey. One day his fifteen-year-old brother, Alan, had just completed, on the family computer, a most difficult and rather ingenious design of an entire city. When Alan slipped out of the room for just a moment, little Jeffrey approached the computer and accidentally erased the program. Upon his return, Alan was furious when he observed what his brother had done. Sensing that his doom was at hand, Jeffrey raised his finger and, pointing it toward Alan, declared from his heart and soul, “Remember, Alan, Jesus said, ‘Don’t hurt little boys.’” Alan began to laugh; anger subsided; mercy prevailed.

There are those among us who torture themselves through their inability to show mercy and to forgive others some supposed offense or slight, however small it may be. At times the statement is made, “I never can forgive [this person or that person].” Such an attitude is destructive to an individual’s well-being. It can canker the soul and ruin one’s life. In other instances, an individual can forgive another but cannot forgive himself. Such a situation is even more destructive.

Early in my ministry as a member of the Council of the Twelve, I took to President Hugh B. Brown the experience of a fine person who could not serve in a ward position because he could not show mercy to himself. He could forgive others but not himself; mercy was seemingly beyond his grasp. President Brown suggested that I visit with that individual and counsel him along these lines: “I, the Lord, will forgive whom I will forgive, but of you it is required to forgive all men.” Then from Isaiah and the Doctrine and Covenants: “Though your sins be as scarlet, they shall be as white as snow; though they be red like crimson, they shall be as wool.” “Behold, he who has repented of his sins, the same is forgiven, and I, the Lord, remember them no more.”

With a pensive expression on his face, President Brown added, “Tell that man that he should not persist in remembering that which the Lord has said He is willing to forget.” Such counsel will help to cleanse the soul and renew the spirit of any who applies it.

The Prophet Joseph urged: “Be merciful and you shall find mercy. Seek to help save souls, not to destroy them: for verily you know, that ‘there is more joy in heaven, over one sinner that repents, than there is over ninety and nine just persons [who] need no repentance.’”

At times a small mistake can fester and bring distress and heartache to him or her who harbors and dwells on the matter, leaving it uncorrected. All of us are subject to such an experience. Let me share with you an example with a beautiful ending. I recently received a note, with a key enclosed, which read:

“Dear President Monson, Thirteen years ago this summer my husband and I stayed at the Hotel Utah. As a memento of our vacation, I took this hotel key and have felt bad about it ever since. I know that the Church owns the former Hotel Utah, and so I am returning this key to you—to the Church—in an effort to set this right. I am so sorry for having taken the key. Please, please, forgive me.”

I thought to myself, What honesty; what a sweet spirit the writer must possess. I replied as follows:

“Dear Sister, Thank you for your thoughtful note and for the Hotel Utah key which you returned. My heart was touched by your sincerity. Though the key itself weighed very little, apparently this has been a heavy burden for you to carry for such a long time. Though the key was of very little worth, its return is of far greater value. I am honored to accept the key and know that you are certainly forgiven. Please accept the enclosed gift with my warmest wishes.”

The key was returned to her, mounted on an attractive plaque.

Should you or I have erred or spoken harshly to another, it is good to take steps to straighten out the matter and to move onward with our lives. “He [who] cannot forgive others breaks the bridge over which he himself must pass if he would ever reach heaven; for every one has need to be forgiven.”

One of the most touching examples of mercy and forgiveness is the well-remembered experience in the life of Jesus when he “went unto the mount of Olives.

“And early in the morning he came again into the temple, and all the people came unto him; and he sat down, and taught them.

“And the scribes and Pharisees brought unto him a woman taken in adultery; and when they had set her in the midst,

“They say unto him, Master, this woman was taken in adultery, in the very act.

“Now Moses in the law commanded us, that such should be stoned: but what sayest thou?

“This they said, tempting him, that they might have to accuse him. But Jesus stooped down, and with his finger wrote on the ground, as though he heard them not.

“So when they continued asking him, he lifted up himself, and said unto them, He that is without sin among you, let him first cast a stone at her.

“And again he stooped down, and wrote on the ground.

“And they which heard it, being convicted by their own conscience, went out one by one, beginning at the eldest, even unto the last: and Jesus was left alone, and the woman standing in the midst.

“When Jesus had lifted up himself, and saw none but the woman, he said unto her, Woman, where are … thine accusers? hath no man condemned thee?

“She said, No man, Lord. And Jesus said unto her, Neither do I condemn thee: go, and sin no more.”

The sands of time quickly erased what the Savior had written, but forever will be remembered the mercy He showed.

\begin{verse}
I stand all amazed at the love Jesus offers me,\\
Confused at the grace that so fully he proffers me.\\
I tremble to know that for me he was crucified,\\
That for me, a sinner, he suffered, he bled and died. …\\
I think of his hands pierced and bleeding to pay the debt!\\
Such mercy, such love, and devotion can I forget?\\
No, no, I will praise and adore at the mercy seat,\\
Until at the glorified throne I kneel at his feet.
\end{verse}

This same Jesus, “seeing the multitudes, he went up into a mountain: and when he was set, his disciples came unto him:

“And he opened his mouth, and taught them, saying, …

“Blessed are the merciful: for they shall obtain mercy.”

My sincere and humble prayer this Sabbath day is that each of us may be the provider and the recipient of mercy—the divine gift. In the name of Jesus Christ, amen.

\newpage
\settalk{That All May Hear}
\setspeaker{Thomas S. Monson}
\settalkdate{April 1995}
\talkheading{That All May Hear}{Thomas S. Monson}

Brethren, you are an inspiring sight to behold. It is gratifying to realize that in thousands of chapels throughout the world, holders of the priesthood of God are receiving this broadcast by way of satellite transmission. Your nationalities vary and your languages are many, but a common thread binds us together. We have been entrusted to bear the priesthood and to act in the name of God. We are the recipients of a sacred trust. Much is expected of us.

With moist eyes and tender hearts, we have said farewell to that gentle giant of a man, even a prophet of God, President Howard W. Hunter. We have sustained this day President Gordon B. Hinckley as the President of the Church and the prophet, seer, and revelator of God. I know that President Hinckley has been called of our Heavenly Father as the prophet and that he will lead us along those pathways the Savior has outlined. The work will go forward and the people will be blessed. It is an honor and distinct privilege to serve with President Gordon B. Hinckley and with President James E. Faust in the First Presidency of the Church.

Long years ago a divine command was given by our Lord and Savior, Jesus Christ, as He said to His beloved eleven disciples: “Go ye therefore, and teach all nations, baptizing them in the name of the Father, and of the Son, and of the Holy Ghost: Teaching them to observe all things whatsoever I have commanded you: and, lo, I am with you alway, even unto the end of the world.” Mark records that “they went forth, and preached every where, the Lord working with them.”

This sacred charge has not been rescinded. Rather, it has been reemphasized. The Prophet Joseph Smith set forth the purpose of the Church when he declared, “It is the bringing of men and women to a knowledge of the eternal truth that Jesus is the Christ, the Redeemer and Savior of the world, and that only through belief in Him, and faith which manifests itself in good works, can men and nations enjoy peace.”

Does the world in which we live stand in need of the teachings of the gospel of Jesus Christ? Almost everywhere one looks there appears an erosion not only of the environment but, even more seriously, an erosion of spirituality and of compliance with eternal commandments. One sees a blatant disregard for the precious souls of mankind.

It is almost as though the faces of many have been turned away from Him—even the Lord—who solemnly declared, “The worth of souls is great in the sight of God.” The gentle words “Come, follow me” fall on many with stopped ears and closed hearts. Such seem to be attuned to another voice.

Do you, with me, remember the story from childhood days of that persuasive musician, the Pied Piper of Hamelin? You will recall that he entered Hamelin and offered, for a specified sum of money, to rid the town of the vermin with which it was plagued. When the contract was agreed upon, he played his pipe and the rats came swarming from the buildings and followed him to the river, where they drowned. When the town leaders refused to pay him for his services, he returned to play his pipe and led the precious children away from the safety of their families and their homes, never to return.

Are there Pied Pipers even today? Are they playing alluring music to lead, to their own destruction, those who listen and follow? These “pipers” pipe the tunes of pride and pleasure, of selfishness and greed and leave in their wake confused minds, troubled hearts, empty lives, and destroyed dreams.

The deep yearning of countless numbers is expressed in the plea of one who spoke to Philip of old: “How can I [find my way], except some man should guide me?”

Brethren of the priesthood, the world is in need of your help. There are feet to steady, hands to grasp, minds to encourage, hearts to inspire, and souls to save. The harvest truly is great. Let there be no mistake about it; the missionary opportunity of a lifetime is yours. The blessings of eternity await you. Yours is the privilege to be not spectators but participants on the stage of priesthood service.

To those of you who hold the Aaronic Priesthood, I say, prepare for your full-time missions. You will become a part of that valiant missionary army of the Lord which now numbers fifty thousand strong.

How might you best respond? May I suggest a formula that will insure your success as missionaries:

First, prepare with purpose. Remember the qualifying statement of the Master: “Behold, the Lord requireth the heart and a willing mind.” Missionary work is difficult. It taxes one’s energies, it strains one’s capacity, it demands one’s best effort—frequently a second effort. No other labor requires longer hours or greater devotion or such sacrifice and fervent prayer.

President John Taylor summed up the requirements: “The kind of men we want as bearers of this Gospel message are men who have faith in God; men who have faith in their religion; men who honor their priesthood; men in whom the people who know them have faith and in whom God has confidence. … We want men full of the Holy Ghost and the power of God. … Men who bear the words of life among the nations, ought to be men of honor, integrity, virtue and purity; and this being the command of God to us, we shall try and carry it out.”

Second, teach with testimony. Peter and John, those converted fishermen who became Apostles, were warned not to preach Jesus Christ and Him crucified. Their answer was firm: “Whether it be right in the sight of God to hearken unto you more than unto God, judge ye. For we cannot but speak the things which we have seen and heard.”

Paul the Apostle, that sterling testifier of truth, was speaking to all of us—members and missionaries alike—when he counseled his beloved friend Timothy, “Be thou an example of the believers, in word, in conversation, in charity, in spirit, in faith, in purity.”

Elder Delbert L. Stapley, who served as a member of the Council of the Twelve a number of years ago, quoted Paul in his epistle to the Romans: “For I am not ashamed of the gospel of Christ: for it is the power of God unto salvation.” Then Elder Stapley added: “If we are not ashamed of the gospel of Christ, then we should not be ashamed to live it. And if we are not ashamed to live it, then we should not be ashamed to share it.”

Third, labor with love. There is no substitute for love. Often this love is kindled in youth by a mother, expanded by a father, and kept vibrant through service to God. Remember the Lord’s counsel: “And faith, hope, charity and love, with an eye single to the glory of God, qualify him for the work.” Well might each of us ask himself: Today have I increased in faith, in hope, in charity, in love? When our lives comply with God’s standard and we labor with love to bring souls unto Him, those within our sphere of influence will never speak the lament, “The harvest is past, the summer is ended, and we are not saved.”

Young missionaries always have an idea as to where they would love to serve. Usually it’s a faraway place with a strange-sounding name.

One day I was in the men’s suit department of a large store when I encountered two missionaries with their mothers. It isn’t difficult to spot missionaries or their mothers. The two elders were conversing, and one said to the other, “Where are you going on your mission?”

Came the reply, “I’m going to Austria.”

The first missionary responded, “You lucky dog, going to Austria! Those beautiful Austrian Alps, that wonderful music, those delightful people! I wish I were going there.”

“Where are you going?” said the missionary assigned to Austria.

“California,” came the answer. “You know, less than two hours away by plane. We go there every year for a vacation.”

I could see by the expression on the mothers’ faces and the near tears of one of the missionaries that it was time for me to intervene. “Did you say California?” I asked. “Why, I once supervised that area. You have an inspired call, Elder. Do you realize what you will have in California to help you? You’ll have chapels and stake centers that dot the land, and they’ll be filled with Latter-day Saints who can be inspired to be fellow missionaries with you in sharing the gospel. You are a very fortunate missionary to be going there.” I glanced at the other mother, who said, “Brother Monson, say something about Austria, quick!” I did so.

Young men, wherever you are called will be right for you, and you will learn to love your mission.

Brethren, all of us can participate, as may our wives and children, in bringing souls to Christ through cooperative effort with the stake and full-time missionaries. One highly successful and rewarding way is through the conducting of open house events in our buildings. You bishops of wards and presidents of stakes have had a video presentation provided you, featuring Elder Jeffrey R. Holland. It is an excellent tool to be used in a missionary open house. Use it. The membership of the Church will grow and the Spirit of the Lord will be among us as we do so.

Prepare with purpose. Teach with testimony. Labor with love. I testify to the truth of this formula and, indeed, this divine work of the Lord.

Many years ago I boarded a plane in San Francisco en route to Los Angeles. As I sat down, the seat next to mine was empty. Soon, however, there occupied that seat a most lovely young lady. As the plane became airborne, I noticed that she was reading a book. As one is wont to do, I glanced at the title: A Marvelous Work and a Wonder. I mustered up my courage and said to her, “Excuse me. You must be a Mormon.”

She replied, “Oh, no. Why do you ask?”

I said, “Well, you’re reading a book written by LeGrand Richards, a very prominent leader of The Church of Jesus Christ of Latter-day Saints.”

She responded, “Is that right? A friend gave this book to me, but I don’t know much about it. However, it has aroused my curiosity.”

I wondered silently, Should I be forward and say more about the Church? The words of the Apostle Peter crossed my mind: “Be ready always to give an answer to every [one] that asketh you a reason of the hope that is in you.” I decided that now was the time for me to share my testimony with her. I told her that it had been my privilege years before to assist Elder Richards in printing that book. I mentioned the great missionary spirit of this man and told her of the many thousands of people who had embraced the truth after reading that which he had prepared. Then it was my privilege, during the remainder of the flight, to answer her questions relative to the Church—intelligent questions which came from the heart, which I perceived was a heart seeking truth. I asked if I might have an opportunity to have the missionaries call upon her. I asked if she would like to attend one of our wards in San Francisco, where she lived. Her answers were affirmative. She gave me her name—Yvonne Ramirez—and indicated that she was a flight attendant on her way to an assignment.

Upon returning home I wrote to the mission president and the stake president, advising them of my conversation and that I had written to her and sent along some suggested reading. Incidentally, young men, I recommended that rather than sending two elders to this pretty off-duty flight attendant and her pretty roommate, two lady missionaries be assigned to call.

Several months passed by. Then I received a telephone call from the stake president, who asked, “Brother Monson, do you remember sitting next to a flight attendant on a trip from San Francisco to Los Angeles earlier this fall?” I answered affirmatively. He continued, “I thought you would like to know that Yvonne Ramirez has just become the most recently baptized and confirmed member of the Church. She would like to speak with you.”

A sweet voice came on the line: “Brother Monson, thank you for sharing with me your testimony. I am the happiest person in all the world.”

As tears filled my eyes and gratitude to God enlarged my soul, I thanked her and commended her on her search for truth and, having found it, her decision to enter those waters which cleanse and purify and provide entrance to eternal life.

I sat silently for a few minutes after replacing the telephone receiver. The words of our Savior coursed through my mind: “And whoso receiveth you, there I will be also, for I will go before your face. I will be on your right hand and on your left, and my Spirit shall be in your hearts, and mine angels round about you, to bear you up.”

Such is the promise to all of us when we pursue our missionary opportunities and follow the counsel and obey the commandments of Jesus of Nazareth, our Savior and our King. He lives—I so testify. In the name of the Lord Jesus Christ, amen.

\newpage
\settalk{Patience—A Heavenly Virtue}
\setspeaker{Thomas S. Monson}
\settalkdate{October 1995}
\talkheading{Patience—A Heavenly Virtue}{Thomas S. Monson}

Recently I met an old friend I had not seen for some time. He greeted me with the salutation, “How is the world treating you?” I don’t recall the specifics of my reply, but his provocative question caused me to reflect on my many blessings and my gratitude for life itself and the privilege and opportunity to serve.

At times the response to this same question brings an unanticipated answer. Some years ago I attended a stake conference in Texas. I was met at the airport by the stake president, and while we were driving to the stake center, I said, “President, how is everything going for you?”

He responded, “I wish you had asked me that question a week earlier, for this week has been rather eventful. On Friday I was terminated from my employment, this morning my wife came down with bronchitis, and this afternoon the family dog was struck and killed by a passing car. Other than these things, I guess everything is all right.”

Life is full of difficulties, some minor and others of a more serious nature. There seems to be an unending supply of challenges for one and all. Our problem is that we often expect instantaneous solutions to such challenges, forgetting that frequently the heavenly virtue of patience is required.

The counsel heard in our youth is still applicable today and should be heeded. “Hold your horses,” “Keep your shirt on,” “Slow down,” “Don’t be in such a hurry,” “Follow the rules,” “Be careful” are more than trite expressions. They describe sincere counsel and speak the wisdom of experience.

The mindless and reckless speeding of a youth-filled car down a winding and hazardous canyon road can bring a sudden loss of control, the careening of the car with its precious cargo over the precipice, and the downward plunge that ofttimes brings permanent incapacity, perhaps premature death, and grieving hearts of loved ones. The glee-filled moment can turn in an instant to a lifetime of regret.

Oh, precious youth, please give life a chance. Apply the virtue of patience.

In sickness, with its attendant pain, patience is required. If the only perfect man who ever lived—even Jesus of Nazareth—was called upon to endure great suffering, how can we, who are less than perfect, expect to be free of all such challenges?

Who can count the vast throngs of the lonely, the aged, the helpless—those who feel abandoned by the caravan of life as it moves relentlessly onward and then disappears beyond the sight of those who ponder, who wonder, and who sometimes question as they are left alone with their thoughts. Patience can be a helpful companion during such stressful times.

Occasionally I visit nursing homes, where long-suffering is found. While attending Sunday services at one facility, I noticed a young girl who was to play her violin for the comfort of those assembled. She told me she was nervous and hoped she could do her best. As she played, one called out, “Oh, you are so pretty, and you play so beautifully.” The strains of the moving bow across the taut strings and the elegant movement of the young girl’s fingers seemed inspired by the impromptu comment. She played magnificently.

Afterward I congratulated her and her gifted accompanist. They responded, “We came to cheer the frail, the sick, and the elderly. Our fears vanished as we played. We forgot our own cares and concerns. We may have cheered them, but they truly did inspire us.”

Sometimes the tables are reversed. A dear and cherished young friend, Wendy Bennion of Salt Lake City, was such an example. Just the day before yesterday, she quietly departed mortality and returned “to that God who gave [her] life” (Alma 40:11). She had struggled for over five long years in her battle with cancer. Ever cheerful, always reaching out to help others, never losing faith, she had a contagious smile that attracted others to her as a magnet attracts metal shavings. While Wendy was ill and in pain, a friend of hers, feeling downcast with her own situation, visited her. Nancy, Wendy’s mother, knowing Wendy was in extreme pain, felt that perhaps the friend had stayed too long. She asked Wendy, after the friend had left, why she had allowed her to stay so long when she herself was in so much pain. Wendy’s response: “What I was doing for my friend was a lot more important than the pain I was having. If I can help her, then the pain is worth it.” Wendy’s attitude was reminiscent of Him who bore the sorrows of the world, who patiently suffered excruciating pain and disappointment, but who, with silent step of His sandaled feet, passed by a man who was blind from birth, restoring his sight. He approached the grieving widow of Nain and raised her son from the dead. He trudged up Calvary’s steep slope, carrying His own cruel cross, undistracted by the constant jeers and taunting that accompanied His every step. For He had an appointment with divine destiny. In a very real way He visits us, each one, with His teachings. He brings cheer and inspires goodness. He gave His precious life that the grave would be deprived of its victory, that death would lose its sting, that life eternal would be our gift.

Taken from the cross, buried in a borrowed tomb, this man of sorrows, acquainted with grief, arose on the morning of the third day. His resurrection was discovered by Mary and the other Mary when they approached the tomb. The great stone blocking the entrance had been rolled away. Came the query of two angels who stood by in shining garments: “Why seek ye the living among the dead? He is not here, but is risen.”

Paul declared to the Hebrews, “Wherefore seeing we also are compassed about with so great a cloud of witnesses, let us lay aside every weight, and the sin which doth so easily beset us, and let us run with patience the race that is set before us.”

Perhaps there has never occurred such a demonstration of patience as that exemplified by Job, who was described in the Holy Bible as being perfect and upright and one who feared God and eschewed evil. He was blessed with great wealth and riches in abundance. Satan obtained leave from the Lord to try to tempt Job. How great was Job’s misery, how terrible his loss, how tortured his life. Urged by his wife to curse God and die, his reply bespoke his faith: “I know that my redeemer liveth, and that he shall stand at the latter day upon the earth: And though after my skin worms destroy this body, yet in my flesh shall I see God.” What faith, what courage, what trust. Job lost possessions—all of them. Job lost his health—all of it. Job honored the trust given him. Job personified patience.

Another who portrayed the virtue of patience was the Prophet Joseph Smith. After his supernal experience in the grove called Sacred, where the Father and the Son appeared to him, he was called upon to wait. At length, after Joseph suffered through over three years of derision for his beliefs, the angel Moroni appeared to him. And then more waiting and patience were required. Let us remember the counsel found in Isaiah: “My thoughts are not your thoughts, neither are your ways my ways, saith the Lord. For as the heavens are higher than the earth, so are my ways higher than your ways, and my thoughts than your thoughts.”

Today in our hurried and hectic lives, we could well go back to an earlier time for the lesson taught us regarding crossing dangerous streets. “Stop, look, and listen” were the watchwords. Could we not apply them now? Stop from a reckless road to ruin. Look upward for heavenly help. Listen for His invitation: “Come unto me, all ye that labour and are heavy laden, and I will give you rest.”

He will teach us the truth of the beautiful verse:

\begin{verse}
Life is real! Life is earnest!\\
And the grave is not its goal;\\
Dust thou art, to dust returnest,\\
Was not spoken of the soul.
\end{verse}

We will learn that each of us is precious to our Elder Brother, even the Lord Jesus Christ. He truly loves us.

His life is the flawless example of one afflicted with sorrows and disappointments, who nonetheless provided the example of forgetting self and serving others. The remembered verse of childhood echoes afresh:

\begin{verse}
Yes, Jesus loves me!\\
Yes, Jesus loves me!\\
Yes, Jesus loves me!\\
The Bible tells me so!
\end{verse}

And so does the Book of Mormon, so does the Doctrine and Covenants, and so does the Pearl of Great Price. Let the scriptures be your guide, and you will never find yourself traveling the road to nowhere.

Today some are out of work, out of money, out of self-confidence. Hunger haunts their lives, and discouragement dogs their paths. But help is here—even food for the hungry, clothing for the naked, and shelter for the homeless.

Thousands of tons move outward from our Church storehouses weekly—even food, clothing, medical equipment and supplies to the far corners of the earth and to empty cupboards and needy people closer to home.

I witness the motivation which prompts busy and talented dentists and doctors on a regular basis to leave their practices and donate their skills to those who need such help. They travel to faraway places to repair cleft palates, correct malformed bones, and restore crippled bodies—all for the love of God’s children. The afflicted who have patiently waited for corrective help are blessed by these “angels in disguise.”

In the words of a well-known song, I wish you could “come fly with me” to eastern Germany, where I visited last month. As we traveled along the autobahns, I reflected on a time twenty-seven years before when I saw on the same autobahns just trucks carrying armed soldiers and policemen. Barking dogs everywhere strained on their leashes, and informers walked the streets. Back then, the flame of freedom had flickered and burned low. A wall of shame sprang up, and a curtain of iron came down. Hope was all but snuffed out. Life, precious life, continued on in faith, nothing wavering. Patient waiting was required. An abiding trust in God marked the life of each Latter-day Saint.

When I made my initial visit beyond the wall, it was a time of fear on the part of our members as they struggled in the performance of their duties. I found the dullness of despair on the faces of many passersby but a bright and beautiful expression of love emanating from our members. In Görlitz the building in which we met was shell-pocked from the war, but the interior reflected the tender care of our leaders in bringing brightness and cleanliness to an otherwise shabby and grimy structure. The Church had survived both the war and the Cold War which followed. The singing of the Saints brightened every soul. They sang the old Sunday School favorite:

\begin{verse}
If the way be full of trial; Weary not!\\
If it’s one of sore denial, Weary not!\\
If it now be one of weeping,\\
There will come a joyous greeting,\\
When the harvest we are reaping—Weary not!
\end{verse}

I was touched by their sincerity. I was humbled by their poverty. They had so little. My heart filled with sorrow because they had no patriarch. They had no wards or stakes—just branches. They could not receive temple blessings—neither endowment nor sealing. No official visitor had come from Church headquarters in a long time. The members were forbidden to leave the country. Yet they trusted in the Lord with all their hearts, and they leaned not to their own understanding. In all their ways they acknowledged Him, and He directed their paths. I stood at the pulpit, and with tear-filled eyes and a voice choked with emotion, I made a promise to the people: “If you will remain true and faithful to the commandments of God, every blessing any member of the Church enjoys in any other country will be yours.”

That night as I realized what I had promised, I dropped to my knees and prayed, “Heavenly Father, I’m on Thy errand; this is Thy church. I have spoken words that came not from me, but from Thee and Thy Son. Wilt Thou, therefore, fulfill the promise in the lives of this noble people.” There coursed through my mind the words from the psalm: “Be still, and know that I am God.” The heavenly virtue of patience was required.

Little by little the promise was fulfilled. First, patriarchs were ordained, then lesson manuals produced. Wards were formed and stakes created. Chapels and stake centers were begun, completed, and dedicated. Then, miracle of miracles, a holy temple of God was permitted, designed, constructed, and dedicated. Finally, after an absence of fifty years, approval was granted for full-time missionaries to enter the nation and for local youth to serve elsewhere in the world. Then, like the wall of Jericho, the Berlin Wall crumbled, and freedom, with its attendant responsibilities, returned.

All of the parts of the precious promise of twenty-seven years earlier were fulfilled, save one. Tiny Görlitz, where the promise had been given, still had no chapel of its own. Now even that dream became a reality. The building was approved and completed. Dedication day dawned. Just a month ago, Sister Monson and I, along with Elder and Sister Dieter Uchtdorf, held a meeting of dedication in Görlitz. The same songs were sung as were rendered twenty-seven years earlier. The members knew the significance of the occasion, marking the total fulfillment of the promise. They wept as they sang. The song of the righteous was indeed a prayer unto the Lord and had been answered with a blessing upon their heads.

At the conclusion of the meeting we were reluctant to leave. As we did so, seen were the waving hands of all, heard were the words, “Auf Wiedersehen, auf Wiedersehen; God be with you till we meet again.”

Patience, that heavenly virtue, had brought to humble Saints its heaven-sent reward. The words of Rudyard Kipling’s “Recessional” seemed so fitting:

\begin{verse}
The tumult and the shouting dies;\\
The captains and the kings depart.\\
Still stands thine ancient sacrifice,\\
An humble and a contrite heart.\\
Lord God of Hosts, be with us yet,\\
Lest we forget, lest we forget.
\end{verse}

In the name of Jesus Christ, amen.

\newpage
\settalk{Who Honors God, God Honors}
\setspeaker{Thomas S. Monson}
\settalkdate{October 1995}
\talkheading{Who Honors God, God Honors}{Thomas S. Monson}

It is no small undertaking to stand before you this evening. I am impressed by your faith, in awe of your potential, and inspired by your devotion to duty in the cause of the Master.

A dear personal friend and associate in the work of the Lord, Elder Bruce R. McConkie, had a favorite hymn which he enjoyed hearing sung. He said the words of the hymn prompted him to do his best. Listen to just two verses:

\begin{verse}
Ye who are called to labor and minister for God,\\
Blest with the royal priesthood, appointed by his word\\
To preach among the nations the news of gospel grace,
\end{verse}

What a mighty promise these precious words proclaim. They apply to you young men who bear the Aaronic Priesthood and to your fathers and other brethren who have received the Melchizedek Priesthood.

It seems like yesterday that I was secretary of the deacons quorum in my ward. We were tutored by wise and patient men who taught us from the holy scriptures, even men who knew us well. These men who took time to listen and to laugh, to build and to inspire, emphasized that we, like the Lord, could increase in wisdom and stature and in favor with God and man. They were examples to us. Their lives were a reflection of their testimonies.

Youth is a time for growth. Our minds during these formative years are receptive to truth, but they are also receptive to error. The responsibility to choose rests with each deacon, teacher, and priest. As the years go by, the choices become increasingly complex, and at times we may be tempted to waver. The need for a personal code of honor is demanded not only on a daily basis but frequently many times in a given day.

The counsel found in one of the hymns sung frequently in our meetings provides an inspired guide:

\begin{verse}
Choose the right when a choice is placed before you.\\
In the right the Holy Spirit guides;\\
And its light is forever shining o’er you,\\
When in the right your heart confides.
\end{verse}

A spirit of determination to do the right thing can come in earliest boyhood. At the cemetery, following a lovely funeral I attended, there stood near the open grave a small lad. His face was one of innocence, and his shining eyes showed the promise of a bright future. I said to him, “You, my boy, are going to make a great missionary. How old are you?”

He answered, “Ten.”

“In nine years we’re going to be looking for you to serve a mission,” I countered.

His response was immediate and told me something about him. He said, “Brother Monson, you won’t have to look for me ’cause I’ll be looking for you.” Young men, some lessons in life are learned from your parents, while others you learn in school or in church. There are, however, certain moments when you know our Heavenly Father is doing the teaching and you are His student. Let me share with you tonight such a lesson, effectively taught and everlastingly learned. The lesson pertains to swimming but goes far beyond that skill.

I learned to swim in the swift-running currents of the Provo River in beautiful Provo Canyon. The “old swimming hole” was in a deep portion of the river, formed by a large rock which had fallen into the river, I assume, when the workmen constructing the railroad were blasting through the canyon. The pool was dangerous, what with its depth of sixteen feet, its current, which moved swiftly against the large rock, and the sucking action of the whirlpools below the rock. It was not a place for a novice or the inexperienced swimmer.

One warm summer afternoon when I was about twelve or thirteen, I took a large, inflated inner tube from a tractor tire, slung it over my shoulder, and walked barefoot up the railroad track which followed the course of the river. I entered the water about a mile above the swimming hole, sat comfortably in the tube, and enjoyed a leisurely float down the river. The river held no fear for me, for I knew its secrets.

That day the Greek-speaking people in Utah held a reunion at Vivian Park in Provo Canyon, as they did every year. Native food, games, and dances were the order of the day. But some left the party to try swimming in the river. When they arrived at the swimming hole, it was deserted, for afternoon shadows were beginning to envelop it.

As my inflated tube bobbed up and down, I was about to enter the swiftest portion of the river just at the head of the swimming hole when I heard frantic cries, “Save her! Save her!” A young lady swimmer, accustomed to the still waters of a gymnasium swimming pool, had fallen from the rock into the treacherous whirlpools. None of the party could swim to save her. Suddenly I appeared on the potentially tragic scene. I saw the top of her head disappearing under the water for the third time, there to descend to a watery grave. I stretched forth my hand, grasped her hair, and lifted her over the side of the tube and into my arms. At the pool’s lower end, the water was slower as I paddled the tube, with my precious cargo, to her waiting relatives and friends. They threw their arms around the water-soaked girl and kissed her, crying, “Thank God! Thank God you are safe!” Then they hugged and kissed me. I was embarrassed and quickly returned to the tube and continued my float down to the Vivian Park bridge. The water was frigid, but I was not cold, for I was filled with a warm feeling. I realized that I had participated in the saving of a life. Heavenly Father had heard the cries, “Save her! Save her,” and permitted me, a deacon, to float by at precisely the time I was needed. That day I learned that the sweetest feeling in mortality is to realize that God, our Heavenly Father, knows each one of us and generously permits us to see and to share His divine power to save.

Pray always in the performance of your priesthood responsibilities, and you will never be in the position of Alice in Wonderland. As Lewis Carroll tells us, Alice was following a path through a forest in Wonderland when it divided into two directions. Standing irresolute, she inquired of the Cheshire Cat, which had suddenly appeared in a nearby tree, which path she should take. “Where do you want to go?” asked the cat.

“I don’t know,” said Alice.

“Then,” said the cat, “it really doesn’t matter, does it?”

We who hold the priesthood know where it is we wish to go. Our objective is the celestial kingdom of our Heavenly Father. Ours is the sacred duty to follow the well-defined path that leads to it.

Soon you will be ready to serve a mission. It’s wonderful that you are willing and prepared to serve wherever the Spirit of the Lord directs. This alone is a modern miracle, considering the times in which we live.

Missionary work is hard work. Missionary service is demanding and requires long hours of study and preparation, that the missionary himself might match the divine message he proclaims. It is a labor of love but also of sacrifice and devotion to duty.

An anxious mother of a prospective missionary once asked me what I would recommend her son learn before the arrival of his missionary call. I am certain she anticipated a profound response which would contain the more familiar requirements for service of which we are all aware. However, I said, “Teach your son how to cook, but more particularly, teach him how to get along with others. He will be happier and more productive if he learns these two vital skills.”

Young men, you are preparing for your missions when you learn your duties as deacons, teachers, and priests and then perform those duties with determination and love, knowing you are on the Lord’s errand.

Sometimes the lessons will come quietly. A few weeks ago I was visiting a sacrament meeting at a care facility in Salt Lake City. The priests at the sacrament table were sitting quietly prior to performing their duties when the opening hymn was announced. A patient near the front of the large room had difficulty opening his hymnbook. Without so much as a question, one of the young men slipped to the patient’s side and, gently turning the pages to the correct hymn, placed the disabled man’s finger at the beginning of the first verse of the hymn. They exchanged an understanding smile, and the priest returned to his seat. This modest gesture of love and helpfulness impressed me. I congratulated him and said, “You are going to be an effective missionary.”

Some missionaries are gifted with the power of expression, while others have a superior knowledge of the gospel. Some, however, are late bloomers who day by day become more proficient and successful. Avoid the temptation of ladder climbing in the mission leadership ranks. It matters little whether you are a district or zone leader or assistant to the president. The important thing is that each one does his very best in the work to which he has been called. Why, I had some missionaries who were so adept at training new missionaries that I couldn’t spare them for other leadership assignments.

Entering the mission field can sometimes be an overpowering and frightening experience. President Harold B. Lee was talking to me one day concerning those who feel inadequate and are worried when they receive an assignment in the Church. He counseled, “Remember, whom the Lord calls, the Lord qualifies.”

When I served as president of the Canadian Mission, headquartered in Toronto, one missionary came to our mission without some of the talents of others, yet he devotedly plunged into his missionary labors. The work was difficult for him; however, he valiantly struggled to be his best self.

At a zone conference, with a General Authority attending, the missionaries had not done too well in a scripture quiz conducted by the visitor. The visitor, with a little sarcasm, commented, “Why, I don’t believe this group knows even the names of the basic missionary pamphlets and their authors.”

Well, that was the proverbial “straw” that broke the camel’s back. I spoke up: “I think they do know them.”

“Well, we will see,” he said, and then he had the missionaries stand. In making a selection of a missionary to prove the point, none of the bright-appearing, experienced, polished missionaries was selected, but rather my new missionary, who had a hard time gaining knowledge of such things, was singled out. My heart literally sank. I looked at the pleading expression on the elder’s face; I knew that he was paralyzed with fear. How I prayed—oh, how I prayed: “Heavenly Father, come to his rescue.” And He did. After a long pause, the visitor said, “Who authored the pamphlet The Plan of Salvation?”

After what seemed like an eternity, the trembling missionary responded, “John Morgan.”

“Who wrote Which Church Is Right?”

Again the pause, and then the reply, “Mark E. Petersen.”

“How about The Lord’s Tenth?”

“James E. Talmage wrote that one,” came the response.

And so it went through the list of missionary pamphlets we used. Finally came the question, “Is there another pamphlet?”

“Yes. It’s called After Baptism, What?”

“Who wrote it?”

Without hesitation the missionary answered, “The name of the author isn’t shown in the pamphlet, but my mission president told me it was written by Elder Mark E. Petersen by assignment from President David O. McKay.”

But what about the missionary? He completed an honorable mission and returned to his home in the West. Later he was called to serve as the bishop of his ward. Every year I receive a Christmas card from him and his wife and family. He always signs his name and then adds this comment, “From your best missionary.”

Each year when that Christmas card arrives, I think of that experience, and the lesson from First Samuel in the Holy Bible penetrates my soul. You will recall that the prophet Samuel was directed by the Lord to go to Bethlehem, even to Jesse, with the revelation that a king would be found among the sons of Jesse. Samuel did as the Lord had commanded him. Each of Jesse’s sons was introduced to Samuel—even seven of them. Though they were fair and qualified in appearance, Samuel was told by the Lord that none was to be chosen. “And Samuel said unto Jesse, Are here all thy children? And he said, There remaineth yet the youngest, and, behold, he keepeth the sheep. And Samuel said unto Jesse, Send and fetch him. … And he sent, and brought him in. … And the Lord said, Arise, anoint him: for this is he.”

The lesson for us to learn is found in the sixteenth chapter of First Samuel, verse seven: “Man looketh on the outward appearance, but the Lord looketh on the heart.”

As bearers of the priesthood, all of us united as one can qualify for the guiding influence of our Heavenly Father as we pursue our respective callings. We are engaged in the work of the Lord Jesus Christ. We, like those of olden times, have answered His call. We are on His errand. We shall succeed in the solemn charge given by Mormon to declare the Lord’s word among His people. He wrote: “Behold, I am a disciple of Jesus Christ, the Son of God. I have been called of him to declare his word among his people, that they might have everlasting life.”

May we ever remember the truth, “Who honors God, God honors.” In the name of Jesus Christ, amen.

\newpage
\settalk{Duty Calls}
\setspeaker{Thomas S. Monson}
\settalkdate{April 1996}
\talkheading{Duty Calls}{Thomas S. Monson}

What a vast audience is attending this general priesthood meeting this evening. The Apostle Peter aptly described you: “Ye are a chosen generation, a royal priesthood, an holy nation, a peculiar people; that ye should shew forth the praises of him who hath called you out of darkness into his marvellous light.”

In the Sunday School of our youth, we often sang the hymn:

\begin{verse}
We are all enlisted till the conflict is o’er;\\
Happy are we! Happy are we!\\
Soldiers in the army, there’s a bright crown in store;\\
We shall win and wear it by and by. …
\end{verse}

Brethren, when we contemplate the wonderful world in which we live and then realize the tumultuous times which beset us, joyful are we to know that Jesus, our leader, ever is near. We live in a world of waste. Too often our natural resources are squandered. We live in a world of want. Some enjoy the lap of luxury, yet others stare starvation in the face. Food, shelter, clothing, and love are not found by all. Unrelieved suffering, unnecessary illness, and premature death stalk too many. We live in a world of wars. Some are political in nature, while others are economic by definition. The greatest battle of all, however, is for the souls of mankind.

Our Captain, even the Lord Jesus Christ, declared:

“Remember the worth of souls is great in the sight of God. …

“And if it so be that you should labor all your days in crying repentance unto this people, and bring, save it be one soul unto me, how great shall be your joy with him in the kingdom of my Father!

“And now, if your joy will be great with one soul that you have brought unto me into the kingdom of my Father, how great will be your joy if you should bring many souls unto me!”

He called fishermen at Galilee to leave their nets and follow Him, declaring, “I will make you fishers of men.” And so He did. He sent His beloved Apostles into all the world to proclaim His glorious gospel. And He issues a call to each of us to “come join the ranks.” He provides our battle plan, with the admonition, “Wherefore, now let every man learn his duty, and to act in the office in which he is appointed, in all diligence.” I love and cherish the noble word duty.

President John Taylor cautioned us, “If you do not magnify your callings, God will hold you responsible for those whom you might have saved had you done your duty.”

Another President, even George Albert Smith, said, “It is your duty first of all to learn what the Lord wants and then by the power and strength of His holy Priesthood to [so] magnify your calling in the presence of your fellows in such a way that the people will be glad to follow you.”

How does one magnify a calling? Simply by performing the service that pertains to it.

We have accepted the call; we have been ordained; we bear the priesthood.

President Stephen L Richards spoke often to holders of the priesthood and emphasized his philosophy pertaining to it. He declared, “The Priesthood is usually simply defined as ‘the power of God delegated to man.’” He continues: “This definition, I think, is accurate. But for practical purposes I like to define the Priesthood in terms of service and I frequently call it ‘the perfect plan of service.’”

You may well ask, “Where does the path of duty lie?” Brethren, I believe with all my heart that two markers define the path: the Duty to Prepare and the Duty to Serve. Let us elaborate on these two markers.

First is the Duty to Prepare. The Lord counseled us, “Seek ye out of the best books words of wisdom; seek learning, even by study and also by faith.”

Preparation for life’s opportunities and responsibilities has never been more vital. We live in a changing society. Intense competition is a part of life. The role of husband, father, grandfather, provider, and protector is vastly different from what it was a generation ago. Preparation is not a matter of perhaps or maybe. It is a mandate. The old phrase “Ignorance is bliss” is forever gone. Preparation precedes performance.

All of us who hold the priesthood are now, or surely will be, teachers of truth. The Lord advised:

“Teach ye diligently and my grace shall attend you, that you may be instructed more perfectly in theory, in principle, in doctrine, in the law of the gospel, in all things that pertain unto the kingdom of God, that are expedient for you to understand, …

“That ye may be prepared in all things when I shall send you again to magnify the calling whereunto I have called you, and the mission with which I have commissioned you.”

Second is the Duty to Serve.

The First Presidency, comprised of Joseph F. Smith, Anthon H. Lund, and Charles W. Penrose, in February 1914 declared, “Priesthood is not given for the honor or aggrandizement of man, but for the ministry of service among those for whom the bearers of that sacred commission are called to labor.”

Now some of you may be shy by nature or consider yourselves inadequate to respond affirmatively to a calling. Remember that this work is not yours and mine alone. It is the Lord’s work, and when we are on the Lord’s errand, we are entitled to the Lord’s help. Remember that whom the Lord calls, the Lord qualifies.

At times, the Lord needs a little help to assist some as to the validity of this truth. I recall when I served as chairman of the Church Missionary Committee that I received a telephone call from a member of the presidency of the Missionary Training Center at Provo, Utah. He advised that a particular young man called to a Spanish-speaking mission was having difficulty applying himself to his language study and had declared, “I never can learn Spanish!” The leader asked, “What do you recommend we do?”

I thought for a moment, then suggested, “Place him tomorrow as an observer in a class of missionaries struggling to learn Japanese, and then advise me of his reaction.”

My caller responded within 24 hours with the report, “The missionary was only in the Japanese language class one-half day when he called me and excitedly said, ‘Place me back in the Spanish class! I know I can learn that language.’” And he did.

While the formal classroom may be intimidating at times, some of the most effective teaching and learning takes place other than in the chapel or the classroom.

Many of you hold the Aaronic Priesthood. You are preparing to become missionaries. Begin now to learn in your youth the joy of service in the cause of the Master. Could I share with you an example of such service.

Following Thanksgiving time a few years ago, I received a letter from a widow whom I had known in the stake where I served in the presidency. She had just returned from a dinner sponsored by her bishopric. Her words reflect the peace she felt and the gratitude which filled her heart:

“Dear President Monson,

“I am living in Bountiful now. I miss the people of our old stake, but let me tell you of a wonderful experience I have had. In early November, all the widows and older people received an invitation to come to a lovely dinner. We were told not to worry about transportation, since this would be provided by the older youth in the ward.

“At the appointed hour, a very nice young man rang the bell and took me and another sister to the stake center. He stopped the car, and two other young men walked with us to the building. Inside, they escorted us to the tables, where seated on each side of us was either a young woman or a young man. We were served a lovely Thanksgiving dinner and afterward provided a choice program.

“Then the young men took us home. It was such a nice evening. Most of us shed a tear or two for the love and respect we were shown.

“President Monson, when you see young people treat others like these young people did, I feel the Church is in good hands.”

There came to mind the words from the Epistle of James: “Pure religion and undefiled before God and the Father is this, To visit the fatherless and widows in their affliction, and to keep himself unspotted from the world.”

I add my own commendation: God bless the leaders, the young men, and the young women who so unselfishly brought such joy to the lonely and such peace to their souls. Through their own experience, they learned the meaning of service and felt the nearness of the Lord.

In 1962, having returned home from presiding over the Canadian Mission of the Church, I received a telephone call from Elder Marion G. Romney. He advised me that the First Presidency had named me as a member of the Adult Correlation Committee of the Church, which committee had the specific assignment to work on the preparation of a new concept—even home teaching. Thus began a most interesting and rewarding experience for me. Each phase of our work, when completed, was reviewed by the First Presidency and the Council of the Twelve. In the spring of 1963, our work was done and a number of us were called to serve on a new committee—the Priesthood Home Teaching Committee—and assigned to go among the stakes of the Church, teaching and encouraging its implementation.

President David O. McKay met with all of the General Authorities of the Church and with the representatives of the committee. He counseled those assembled: “Home teaching is one of our most urgent and most rewarding opportunities to nurture and inspire, to counsel and direct our Father’s children. … It is a divine service, a divine call. It is our duty as Home Teachers to carry the divine spirit into every home and heart.”

In 1987 President Ezra Taft Benson counseled the brethren attending the general priesthood meeting: “Home teaching is not to be undertaken casually. A home teaching call is to be accepted as if extended to you personally by the Lord Jesus Christ.” He quoted the familiar passage from section 20 of the Doctrine and Covenants, wherein the Lord declared to the priesthood:

“Watch over the church always, and be with and strengthen them;

“And see that there is no iniquity in the church … ;

“And see that the church meet together often, and also see that all the members do their duty.”

“And visit the house of each member, exhorting them to pray vocally and in secret and attend to all family duties.”

Recently our grandchildren received their report cards. They displayed them with satisfaction to their parents and to us. Tonight I would like all of the priesthood to mark their own grade on the report card of home teaching. Are you ready? Yes or No answers are sufficient.

The quiz could continue, but I sense the questions have been adequate to prompt a mental review and to motivate improved performance.

I am aware that we at headquarters authorized some modifications in the home teaching effort where priesthood numbers were very few—even to permitting a wife to accompany her husband where another companion from the priesthood was not available. But these exceptions were to be just that: exceptions—not the rule. We urge that an active bearer of the Melchizedek Priesthood have assigned with him a teacher or a priest or a prospective elder, complying with the scripture, “And if any man among you be strong in the Spirit, let him take with him him that is weak, that he may be edified in all meekness, that he may become strong also.” This is priesthood home teaching as it generally is meant to function.

Should we feel the assignment too arduous or time-consuming, let me share with you the experience of a faithful home teacher and his companion in what was then East Germany.

Brother Johann Denndorfer had been converted to the Church in Germany, and following World War II, he found himself virtually a prisoner in his own land—the land of Hungary in the city of Debrecen. How he wanted to visit the temple! How he desired to receive his spiritual blessings! Request after request to journey to the temple in Switzerland had been denied, and he almost despaired. Then his home teacher visited. Brother Walter Krause went from the northeastern portion of Germany all the way to Hungary. He had said to his home teaching companion, “Would you like to go home teaching with me this week?”

His companion said, “When will we leave?”

“Tomorrow,” replied Brother Krause.

“When will we come back?” asked the companion.

“Oh, in about a week—if we get back then!”

And away they went to visit Brother Denndorfer. He had not had home teachers since before the war. Now, when he saw the servants of the Lord, he was overwhelmed. He did not shake hands with them; rather, he went to his bedroom and took from a secret hiding place his tithing that he had saved from the day he became a member of the Church and returned to Hungary. He presented the tithing to his home teachers and said: “Now I am current with the Lord. Now I feel worthy to shake the hands of servants of the Lord!”

Brother Krause asked him about his desire to attend the temple in Switzerland. Brother Denndorfer said: “It’s no use. I have tried and tried. The government has even confiscated my Church books, my greatest treasure.”

Brother Krause, a patriarch, provided Brother Denndorfer with a patriarchal blessing. At the conclusion of the blessing, he said to Brother Denndorfer, “Approach the government again about going to Switzerland.” And Brother Denndorfer submitted the request once again to the authorities. This time approval came, and with joy Brother Denndorfer went to the Swiss Temple and stayed a month. He received his own endowment, his deceased wife was sealed to him, and he was able to accomplish the work for hundreds of his ancestors. He returned to his home renewed in body and in spirit.

And what about the home teachers who undertook this historic and inspired visit to their brother, Johann Denndorfer?

Knowing personally each member of this human drama, I wouldn’t be a bit surprised to learn that on the way from Debrecen, Hungary, to their home in East Germany, they sang aloud: “Dangers may gather—why should we fear? Jesus, our Leader, ever is near. He will protect us, comfort, and cheer. We’re joyfully, joyfully marching to our home.”

Brethren of the priesthood, may all of us remember our duty to prepare and our duty to serve, that we may merit the Lord’s approbation, “Well done, thou good and faithful servant.” In the name of Jesus Christ, amen.

\newpage
\settalk{The Way of the Master}
\setspeaker{Thomas S. Monson}
\settalkdate{April 1996}
\talkheading{The Way of the Master}{Thomas S. Monson}

During the later Judean ministry of the Lord, “a certain lawyer stood up, and tempted him, saying, Master, what shall I do to inherit eternal life?

“He said unto him, What is written in the law? how readest thou?

“And he answering said, Thou shalt love the Lord thy God with all thy heart, and with all thy soul, and with all thy strength, and with all thy mind; and thy neighbour as thyself.

“And he said unto him, Thou hast answered right: this do, and thou shalt live.

“But he, willing to justify himself, said unto Jesus, And who is my neighbour?

“And Jesus answering said, A certain man went down from Jerusalem to Jericho, and fell among thieves, which stripped him of his raiment, and wounded him, and departed, leaving him half dead.

“And by chance there came down a certain priest that way: and when he saw him, he passed by on the other side.

“And likewise a Levite, when he was at the place, came and looked on him, and passed by on the other side.

“But a certain Samaritan, as he journeyed, came where he was: and when he saw him, he had compassion on him,

“And went to him, and bound up his wounds, pouring in oil and wine, and set him on his own beast, and brought him to an inn, and took care of him.

“And on the morrow when he departed, he took out two pence, and gave them to the host, and said unto him, Take care of him; and whatsoever thou spendest more, when I come again, I will repay thee.

“Which now of these three, thinkest thou, was neighbour unto him that fell among the thieves?

“And he said, He that shewed mercy on him. Then said Jesus unto him, Go, and do thou likewise.”

Times change, the years roll by, circumstances vary—but the Master’s counsel to the lawyer applies to you and to me just as surely as though we heard His voice speaking directly to us this Easter morn.

How might we fulfill today the first part of the divine commandment to love the Lord our God?

The Lord declared: “He that hath my commandments, and keepeth them, he it is that loveth me”; “Come, follow me”; “I have set an example for you”; “I am the light which ye shall hold up—that which ye have seen me do.” What, indeed, did He do?

Born in a stable, cradled in a manger, He brought to fulfillment the prophecies of the ages. Shepherds came with haste to adore Him. Wise men from the East came bearing for Him precious gifts; the meridian of time had dawned.

With the birth of the babe in Bethlehem, there emerged a great endowment, a power stronger than weapons, a wealth more lasting than the coins of Caesar. This child was to be the King of Kings and Lord of Lords, the promised Messiah—even Jesus Christ, the Son of God.

The holy scriptures inform us that “Jesus increased in wisdom and stature, and in favour with God and man.” He was baptized by John. He “went about doing good.” At Nain He raised from death to life the widow’s son and presented him to her. At Bethesda He took compassion on the crippled man who had no hope to get to the pool of promise. He extended His hand; He lifted him up. He healed him from his infirmity.

Then came the Garden of Gethsemane with its exceeding anguish. He wrought the great Atonement as He took upon Himself the sins of all. He did for us what we could not do for ourselves.

Then came the cruel cross of Golgotha. In His final hours of mortality, He brought comfort to the malefactor, saying, “To day shalt thou be with me in paradise.” He remembered His mother in that eloquent sermon of love personified: “When Jesus therefore saw his mother, and the disciple standing by, whom he loved, he saith unto his mother, Woman, behold thy son! Then saith he to the disciple, Behold thy mother! And from that hour that disciple took her unto his own home.” He died—the Great Redeemer died!

Two questions, spoken at an earlier time, roll as thunder to the ears of each of us: “What think ye of Christ?” and “What shall [we] do … with Jesus?” I proffer these three suggestions:

By learning of Him, by believing in Him, by following Him, there is the capacity to become like Him. The countenance can change, the heart can be softened, the step can be quickened, the outlook enhanced. Life becomes what it should become. Change is at times imperceptible, but it does take place.

The Savior’s entire ministry exemplified love of neighbor, the second part of that lesson given to the inquiring lawyer—spoken of as the “royal law.”

A blind man healed, the daughter of Jairus raised, and the lepers cleansed—all were neighbors of Jesus. Neighbor also was the woman at the well. He, the perfect man, standing before a confessed sinner, extended a hand. She was the traveler; He was the good Samaritan. And so the caravan of His kindness continued.

What about our time and place? Do neighbors await our love, our kindness, our help?

A few years ago I read a Reuters news service account of an Alaska Airlines nonstop flight from Anchorage to Seattle, carrying 150 passengers, which was diverted to a remote town on a mercy mission to rescue a badly injured boy. Two-year-old Elton Williams III had severed an artery in his arm when he fell on a piece of glass while playing near his home in Yakutat, 450 miles south of Anchorage. Medics at the scene asked the airline to evacuate the boy. As a result, the Anchorage-to-Seattle flight was diverted to Yakutat.

The medics said the boy was bleeding badly and probably would not live through the flight to Seattle, so the plane flew 200 miles to Juneau, the nearest city with a hospital. The flight then went on to Seattle, with the passengers arriving two hours late, most missing their connections. But none complained. In fact, they dug into their pocketbooks and took up a collection for the boy and his family.

Later, as the flight was about to land in Seattle, the passengers broke into a cheer when the pilot said he had received word by radio that Elton was going to be all right. Surely love of neighbor was in evidence.

A man was asked one day, “Who is your next-door neighbor?”

He said, “I don’t know his name, but his children run across my lawn and his dog keeps me awake at night!”

Another man, in a different mood, wrote silently one night in his journal: “I thought the house across the street was empty until yesterday. Black crepe on the door made me aware that someone had been living there.”

A poet set to verse the sorrow of opportunities forever lost:

\begin{verse}
Around the corner I have a friend,\\
In this great city that has no end;\\
Yet days go by, and weeks rush on,\\
And before I know it, a year is gone,\\
And I never see my old friend’s face,\\
For Life is a swift and terrible race.\\
He knows I like him just as well\\
As in the days when I rang his bell\\
And he rang mine.\\
We were younger then,\\
And now we are busy, tired men:\\
Tired with playing a foolish game,\\
Tired with trying to make a name.\\
“To-morrow,” I say, “I will call on Jim,\\
Just to show that I’m thinking of him.”\\
But to-morrow comes—and to-morrow goes,\\
And the distance between us grows and grows.
\end{verse}

Long years ago I was touched by a story which illustrated love of neighbor between a small boy named Paul and a telephone operator he had never met. These were the days many will remember with nostalgia but which a new generation will never experience.

Paul related the story: “When I was quite young, my father had one of the first telephones in our neighborhood. I remember that the shiny receiver hung on the side of the box. I was too little to reach the telephone, but I used to listen with fascination when Mother would talk to it. Then I discovered that somewhere inside the wonderful device lived an amazing person. Her name was ‘Information, Please,’ and there was nothing she did not know. ‘Information, Please’ could supply anybody’s number and the correct time.

“I learned that if I stood on a stool, I could reach the telephone. I called ‘Information, Please’ for all sorts of things. I asked her for help with my geography, and she told me where Philadelphia was. She helped me with my arithmetic, too.

“Then there was the time that Petey, our pet canary, died. I called ‘Information, Please’ and told her the sad story. She listened and then said the usual things grown-ups say to soothe a child. But I was unconsoled. ‘Why is it that birds should sing so beautifully and bring joy to all families, only to end up as a heap of feathers, feet up, on the bottom of the cage?’ I asked.

“She must have sensed my deep concern, for she said quietly, ‘Paul, always remember that there are other worlds in which to sing.’ Somehow I felt better.

“All this took place in a small town near Seattle. Then we moved across the country to Boston. I missed my friend very much. ‘Information, Please’ belonged to that old wooden box back home, and I somehow never thought of trying to call her. The memories of those childhood conversations never really left me; often in moments of doubt and perplexity I would recall the serene sense of security I had then. I appreciated now how patient, understanding, and kind she was to have spent her time on a little boy.

“Later, when I went west to college, my plane made a stop in Seattle,” Paul continued. “I called ‘Information, Please,’ and when, miraculously, I heard that familiar voice, I said to her, ‘I wonder if you have any idea how much you meant to me during that time?’

“‘I wonder,’ she said, ‘if you know how much your calls meant to me. I never had any children, and I used to look forward to your calls.’ I told her how often I had thought of her over the years, and I asked if I could call her again when I came back west.

“‘Please do,’ she said. ‘Just ask for Sally.’

“Only three months later I was back in Seattle. A different voice answered, ‘Information,’ and I asked for Sally. ‘Are you a friend?’ the woman asked.

“‘Yes, a very old friend,’ I replied.

“‘Then I’m sorry to have to tell you. Sally has only been working part-time the last few years because she was ill. She died five weeks ago.’ But before I could hang up, she said, ‘Wait a minute. Did you say your name was Paul?’

“‘Yes,’ I responded.

“‘Well, Sally left a message for you. She wrote it down. Here it is—I’ll read it. Tell him I still say there are other worlds in which to sing. He’ll know what I mean.

“I thanked her and hung up,” said Paul. “I did know what Sally meant.”

Sally, the telephone operator, and Paul, the boy—the man—were in reality good Samaritans to each other.

There are indeed other worlds in which to sing. Our Lord and Savior brought to each of us the reality of this truth. To the grieving Martha He comforted: “I am the resurrection, and the life: he that believeth in me, though he were dead, yet shall he live:

“And whosoever liveth and believeth in me shall never die.”

If we truly seek our Lord and Savior, we shall surely find Him. He may come to us as one unknown, without a name, as of old by the lakeside He came to those men who knew Him not. He speaks to us the same words, “Follow thou me,” and sets us to the tasks which He has to fulfill for our time. He commands, and to those who obey Him, whether they be wise or simple, He will reveal Himself in the toils, the conflicts, the sufferings which they shall pass through in His fellowship, and they shall learn in their own experience who He is.

On this Easter Sabbath, we remember loved ones who have gone from our midst. Cherished memories of happy days, followed by lonely nights, long years, and pensive thoughts, turn our hearts to Him who promised: “Peace I leave with you, my peace I give unto you: not as the world giveth, give I unto you. Let not your heart be troubled, neither let it be afraid.” “In my Father’s house are many mansions: if it were not so, I would have told you. I go to prepare a place for you … ; that where I am, there ye may be also.”

He who taught us to love the Lord our God with all our hearts, and with all our souls, and with all our strength, and with all our minds, and our neighbors as ourselves, is a Teacher of truth—but He is more than a teacher. He is the Exemplar of the perfect life—but He is more than an exemplar. He is the Great Physician—but He is more than a physician. He is the literal Savior of the world, the Son of God, the Prince of Peace, the Holy One of Israel, even the risen Lord, who declared: “Behold, I am Jesus Christ, whom the prophets testified shall come into the world. … I am the light and the life of the world.” “I am the first and the last; I am he who liveth, I am he who was slain; I am your advocate with the Father.”

This Easter morning, as His witness, I testify to you that He lives and that through Him, we too shall live. In the name of Jesus Christ, amen.

\newpage
\settalk{“Be Thou an Example”}
\setspeaker{Thomas S. Monson}
\settalkdate{October 1996}
\talkheading{“Be Thou an Example”}{Thomas S. Monson}

As I contemplate the vast audience assembled for this general priesthood meeting of the Church, I seek the help of our Heavenly Father. I approach my responsibility to speak to you with the deepest humility.

Of late I have been studying the teachings of the early Apostles, including their calls, their ministries, and their very lives. It is a fascinating experience and brings one closer to the Lord Jesus Christ.

Tonight may I share with you a profound plea given by the Apostle Paul to his beloved Timothy. Paul’s words are applicable to each of us: “Be thou an example of the believers, in word, in conversation, in charity, in spirit, in faith, in purity. … Neglect not the gift that is in thee. … Meditate upon these things.”

Brethren, ours is the opportunity to learn, the privilege to obey, and the duty to serve. In our time there are feet to steady, hands to grasp, minds to encourage, hearts to inspire, and souls to save.

For example, consider the law of tithing. The honest payment of tithing provides a person the inner strength and commitment to comply with the other commandments.

President Gordon B. Hinckley has declared: “There has been laid upon the Church a tremendous responsibility. Tithing is the source of income for the Church to carry forward its mandated activities. The need is always greater than the availability. God help us to be faithful in observing this great principle which comes from him with his marvelous promise.”

From Malachi we read: “Will a man rob God? Yet ye have robbed me. But ye say, Wherein have we robbed thee? In tithes and offerings. … Bring ye all the tithes into the storehouse, that there may be meat in mine house, and prove me now herewith, saith the Lord of hosts, if I will not open you the windows of heaven, and pour you out a blessing, that there shall not be room enough to receive it.”

All of us can afford to pay tithing. In reality, none of us can afford not to pay tithing. The Lord will strengthen our resolve. He will open a way to comply.

May I share with you a letter I received some months ago which provides such an example? The letter begins:

“We live on the edge of a small town, and our neighbor uses our pasture for his cattle and as payment provides us with all the beef we want. Each time we get new meat, we have some of the present supply left over; and since we live in a student ward, we take meat to some students we feel might have use for some good beef.

“During the time my wife was serving in a Relief Society presidency, her secretary was a student’s wife—the mother of eight children. Her husband, Jack, had recently been called as ward clerk.

“My wife had always prayed to know which students might need our help with our excess meat. When she told me she felt we should give some meat to Jack and his family, I was very concerned that we might offend them. So was she. We both were worried because they were a very independent family.

“A few days later, my wife said she still felt we should take the meat to them, and I reluctantly agreed to go along. When we delivered the meat, my wife’s hands were actually shaking, and I was very nervous. The children opened the door, and when they heard why we were there, they began dancing around. The parents were reserved but pleasant. When we drove away, my wife and I both were so relieved and happy that they had accepted our gift.

“A few months later our friend Jack got up in testimony meeting and related the following. He said that all his life he had had a hard time paying tithing. With such a large family, they used all the money he made just to get by. When he became ward clerk, he saw all the other people paying tithing and felt he needed to also. He did so for a couple of months, and all was well. Then one month he had a problem. In his job, he completed work and was paid a few months later. He could see that the family was going to be far short of money. He and his wife decided to share the problem with their children. If they paid their tithing, they would run out of food on about the 20th of the month. If they didn’t pay their tithing, they could buy enough food to last until the next paycheck. Jack said he wanted to buy [the] food, but the children said they wanted to pay tithing—so Jack paid the tithing, and they all prayed.

“A few days after paying their tithing, we had shown up with our package of meat for them. With the meat, added to what they had, there was no problem having enough food until the next paycheck.

“There are so many lessons here for me—for instance, always listen to my wife—but for me the most important is that the prayers of people are almost always answered by the actions of others.”

I recognize that there are thousands of missionaries attending this priesthood meeting tonight. I wish to share a word or two especially with you. During the time I served as a mission president, and then in thousands of missionary interviews as a member of the Twelve, I said to the missionaries I interviewed, “When you return to your home, I ask that you make three commitments.” Eagerly, without knowing what the commitments were, they would nod their approval. I then shared with them this counsel:

I could say much more concerning tithing, but tonight I would also wish to address the other part of Malachi’s message—namely, offerings.

The concept of fast offerings appears as early as the time of Isaiah when, speaking of the true fast, he encouraged people to fast and “to deal thy bread to the hungry, and … bring the poor that are cast out to thy house.” The Prophet Joseph instituted the practice of collecting fast offerings for the poor in Kirtland, Ohio. And later at Nauvoo, Illinois, the Quorum of the Twelve Apostles sent a general letter to the Church defining “the principle of fasts,” stating: “Let this be an ensample to all saints, and there will never be any lack for bread: When the poor are starving, let those who have, fast one day and give what they otherwise would have eaten to the bishops for the poor, and every one will abound for a long time; and this is one great and important principle of fasts approved of the Lord. And so long as the saints will all live to this principle with glad hearts and cheerful countenances they will always have an abundance.”

The prophets of our day and time have been equally specific. Harold B. Lee counseled: “When you think about it, there is so much promised in the gospel for so little required on our part. For example, the ordinance of baptism is given us for the remission of sins, for entrance into the kingdom—a new birth; the gift of the Holy Ghost gives us the right to companionship with one of the Godhead; administration to the sick qualifies the individual with faith for a special blessing; by paying our tithing, the windows of heaven may be opened unto us; by fasting and by paying our fast offerings, we are told that then we might call on the Lord and He will hear our cry and our call.”

President Lee’s successor in the Presidency of the Church, President Spencer W. Kimball, said: “We wish to remind all the Saints of the blessings that come from observing the regular fast and contributing as generous a fast offering as we can, and as we are in a position to give. Wherever we can, we should give many times the value of the meals from which we abstained.” President Kimball added: “Collecting fast offerings is an important [duty]. I thought it was a great honor to be a deacon. My father was always considerate … and … permitted me to take the buggy and horse to gather fast offerings. My responsibility included that part of the town in which I lived, but it was quite a long walk to the homes, and a sack of flour or a bottle of fruit or vegetables or bread became quite heavy as it accumulated. So the buggy was very comfortable and functional. We have changed to cash in later days, but it was commodities in my day. It was a very great honor to do this service for my Heavenly Father: and though times have changed, when money is given generally instead of commodities, it is still a great honor to perform this service.”

I imagine you young deacons today also wouldn’t mind taking a horse and buggy to gather fast offerings!

I remember when, as a young deacon, I would cover a portion of the ward on fast Sunday morning, giving the small envelope to each family, waiting while a contribution was placed in it, and then returning it to the bishop. On one such occasion, an elderly member, Brother Wright, welcomed me at the door and, with aged hands, fumbled at the tie of the envelope and placed within it a quarter. His eyes fairly twinkled as he made his contribution. He told me of a time years before when the Relief Society president, Sister Balmforth, with food collected from those who had given, carried to his home in a small red wagon food for his cupboard and provided gratitude for his soul. He described her as “an angel sent from heaven.” I have not forgotten Eddie Wright.

Deacons and others of the Aaronic Priesthood who perform today this sacred service, please know this to be a sacred duty. I recall that as a bishop, one morning the boys in the ward over which I presided had assembled—sleepy-eyed, a bit disheveled, and mildly complaining about arising so early to fulfill their assignment. Not a word of reproof was spoken, but during the following week we escorted the boys to Welfare Square in Salt Lake City for a guided tour. They saw firsthand a lame sister operating the telephone switchboard, an older man stocking shelves, women arranging clothing to be distributed—even a blind person placing labels on cans of food. Here were individuals earning their sustenance through their contributed labors. A penetrating silence came over the boys as they witnessed how their effort each month helped to collect the sacred fast offering funds which aided the needy and provided employment for those who otherwise would be idle.

From that hallowed day forward, we no longer had to urge our deacons with regard to collecting fast offerings. On fast Sunday mornings, they were present at 7:00 a.m., dressed in their Sunday best, anxious to do their duty as holders of the Aaronic Priesthood. No longer were they simply distributing and collecting envelopes. They were helping to provide food for the hungry and shelter for the homeless—all after the way of the Lord. Their smiles were more frequent, their pace more eager, their very souls more subdued. Perhaps now they were marching to the beat of a different drummer; perhaps now they better understood the classic passage, “Inasmuch as ye have done it unto one of the least of these my brethren, ye have done it unto me.”

In the vicinity where I lived and served, we operated a poultry project. Most of the time it was an efficiently operated welfare project, supplying to the storehouse thousands of dozens of fresh eggs and hundreds of pounds of dressed poultry. On a few occasions, however, the experience of being volunteer city farmers provided not only blisters on the hands, but frustration of heart and mind. For instance, I shall ever remember the time we gathered together the teenaged Aaronic Priesthood young men to really give the poultry project a spring cleaning. Our enthusiastic and energetic throng gathered at the project and in a speedy fashion uprooted, gathered, and burned large quantities of weeds and debris. By the light of the glowing bonfires we ate hot dogs and congratulated ourselves on a job well done. The project was now neat and tidy. However, there was just one disastrous problem. The noise and the fires had so disturbed the fragile and temperamental population of several thousand laying hens that most of them went into a sudden molt and ceased laying. Thereafter we tolerated a few weeds, that we might produce more eggs.

No member of The Church of Jesus Christ of Latter-day Saints who has canned peas, topped beets, hauled hay, or shoveled coal in such a cause ever forgets or regrets the experience of helping provide for those in need. Devoted men and women help to operate this vast and inspired welfare program. In reality, the plan would never succeed on effort alone, for this program operates through faith after the way of the Lord.

Brethren, you and your families are to be commended for the manner in which you also contribute generously to the humanitarian efforts of the Church throughout the world. We provide essential help to the needy in times of natural disasters, starvation, sickness, and events that can strike anywhere. Emergency food supplies, clothing, shelter, and medical equipment bring succor to the suffering and peace to the recipient and to the giver—even the peace promised of the Lord. Projects provided by your generosity bring health and happiness through the drilling of wells to provide uncontaminated water to those who have never had such. Children walk who once would have been crippled by polio, thanks to your contributions which provided the vaccine to prevent such tragedies.

Should you be in Salt Lake City, visit the Sort Center, where millions of pounds of contributed clothing are received, sorted, packed, and shipped to the needy throughout the world as well as to pockets of poverty situated closer to home. One is reminded of the statement made by the Prophet Joseph: “A man filled with the love of God is not content with blessing his family alone, but ranges through the whole world, anxious to bless the whole human race.”

Most of you are home teachers. You are the eyes and ears of the bishops in seeking out the poor and the afflicted. While doing their duty, vigilant home teachers have observed unemployed fathers anxious to obtain work; distraught mothers seeing their tiny broods suffer; children crying from hunger, inadequately clothed to protect them from the cold of winter. In one instance, all of the family members were sleeping on the floor because they had no beds. Without delay, needed help was provided.

Remember the counsel from King Benjamin described in Mosiah: “Ye yourselves will succor those that stand in need of your succor; ye will administer of your substance unto him that standeth in need; and ye will not suffer that the beggar putteth up his petition to you in vain, and turn him out to perish.”

Fortunately and commendably, the Church is doing more than it has ever done to relieve suffering, to satisfy hunger, to prevent and cure illness, and to bless those in need. There is more to do.

Brethren, my prayer is that we be “example[s] of the believers, in word, in conversation, in charity, in spirit, in faith, in purity.” Then shall we be recipients of the Lord’s promise:

“I, the Lord, am merciful and gracious unto those who fear me, and delight to honor those who serve me in righteousness and in truth unto the end.

“Great shall be their reward and eternal shall be their glory.”

In the name of Jesus Christ, amen.

\newpage
\settalk{Christ at Bethesda’s Pool}
\setspeaker{Thomas S. Monson}
\settalkdate{October 1996}
\talkheading{Christ at Bethesda’s Pool}{Thomas S. Monson}

One of the most famous art galleries in the world is the National Gallery of Art, situated adjacent to Trafalgar Square in the city of London, England. The gallery has on display many priceless masterpieces.

Just a few weeks ago my wife, Frances, and I visited the National Gallery and admired the display of inspired genius which met our gaze and touched our hearts. A large painting occupied most of the wall of one room. It was an incomparable piece by the renowned Bartolomé Esteban Murillo, completed in the year 1670 and titled Christ Healing the Paralytic at the Pool of Bethesda. The centuries have not dimmed its beauty, dulled its appeal, nor diminished its impact.

I could not avert my eyes, nor could I transfer my thoughts. I was carried back through time as I saw the crippled man lying on his crude crutch with his arms extended and his hands upturned as he appealed to the Savior of the world. The words and thoughts expressed in the book of John coursed through my mind. I share them with you this morning:

“Now there is at Jerusalem by the sheep market a pool, which is called in the Hebrew tongue Bethesda, having five porches.

“In these lay a great multitude of impotent folk, of blind, halt, withered, waiting for the moving of the water.

“For an angel went down at a certain season into the pool, and troubled the water: whosoever then first after the troubling of the water stepped in was made whole of whatsoever disease he had.

“And a certain man was there, which had an infirmity thirty and eight years.

“When Jesus saw him lie, and knew that he had been now a long time in that case, he saith unto him, Wilt thou be made whole?

“The impotent man answered him, Sir, I have no man, when the water is troubled, to put me into the pool: but while I am coming, another steppeth down before me.

“Jesus saith unto him, Rise, take up thy bed, and walk.

“And immediately the man was made whole, and took up his bed, and walked.”

At length, after pondering this scripture, I left the reverie of the room; however, the impact of that masterpiece was indelibly impressed on my soul.

I have thought since of the majesty of the Master’s command, the tenderness of His heart, and the incredible joy His act had brought to the afflicted man.

\begin{verse}
Jesus, the very thought of thee\\
With sweetness fills my breast;\\
But sweeter far thy face to see\\
And in thy presence rest.
\end{verse}

Do we remember the question posed by one Pontius Pilate as he spoke to those who would shed the blood of Jesus and thus end His mortal life? “What shall I do then with Jesus which is called Christ? They all say unto him, Let him be crucified.” And so He was.

The question each of us must answer is the same: What shall I do with Jesus? He Himself has provided us the answer: “Follow me, and do the things which ye have seen me do.”

The mortal mission of our Lord was foretold by the holy prophets, as was His birth. For generations, enlightened mankind in the Old and the New World anxiously sought the fulfillment of prophecies uttered by righteous men inspired of Almighty God.

Then came that heavenly pronouncement to the “shepherds abiding in the field, keeping watch over their flock by night. … For unto you is born this day in the city of David a Saviour, which is Christ the Lord.” Born in a stable, cradled in a manger, He came forth from heaven to live on earth as mortal man and to establish the kingdom of God. His glorious gospel reshaped the thinking of the world. He blessed the sick; He caused the lame to walk, the blind to see, the deaf to hear. He even raised the dead to life. He provided for you and for me the greatest gift we shall ever receive: the Atonement and all that it conveys. He willingly died that we might forever live.

From time to time the question has been posed, “If Jesus appeared to you today, what questions would you ask of Him?”

My answer has always been, “I would not utter a word. I would listen to Him.”

Down through the generations of time, the message from Jesus has been the same. To Peter by the shores of beautiful Galilee, He said, “Follow me.” To Philip of old came the call, “Follow me.” To Levi who sat at receipt of customs came the instruction, “Follow me.” And to you and to me, if we but listen, shall come that same beckoning invitation, “Follow me.”

“Jesus increased in wisdom and stature, and in favour with God and man.” Have we? Of Him it was said that He “went about doing good.” Do we?

His beloved Apostles noted well His example. He lived “not to be ministered unto, but to minister”; not to receive, but to give; not to save His life, but to pour it out for others. It has been said, “If they would see the star that should at once direct their feet and influence their destiny, they must look for it, not in the changing skies [of] outward circumstance, but each in the depth of his own heart and after the pattern provided by the Master.”

Reflect for a moment on the experience of Peter at the Gate Beautiful of the temple. One sympathizes with the plight of the man lame from birth who each day was carried to the temple gate that he might ask alms of all who entered. That he asked alms of Peter and John as they approached him indicates he regarded them no differently from others who must have passed him each day. I love Peter’s simple and direct instruction: “Look on us.” The lame man gave heed to them.

“Then Peter said, Silver and gold have I none; but such as I have give I thee: In the name of Jesus Christ of Nazareth rise up and walk.

“And he took him by the right hand, and lifted him up. …

“… He … stood, and walked, and entered with them into the temple.”

Not all who approached the Master abided by His divine direction:

“And when he was gone forth into the way, there came one running, and kneeled to him, and asked him, Good Master, what shall I do that I may inherit eternal life?

“And Jesus said unto him, …

“Thou knowest the commandments, Do not commit adultery, Do not kill, Do not steal, Do not bear false witness, Defraud not, Honour thy father and mother.

“And he answered and said unto him, Master, all these have I observed from my youth.

“Then Jesus beholding him loved him, and said unto him, One thing thou lackest: go thy way, sell whatsoever thou hast, and give to the poor, and thou shalt have treasure in heaven: and come, take up the cross, and follow me.

“And he was sad at that saying, and went away grieved: for he had great possessions.”

Some time ago I received a touching letter from Randy Spaulding, who lived in northern Utah. The letter explained the composition of his family and then the gradual onset of an illness that took his father from a healthy, strong individual to a weak and crippled middle-aged man. The father’s physical condition deteriorated until he could not work, could not walk, became confined to a wheelchair, and was almost helpless.

Randy told how the family and ward members have taken over the care of the farm and have provided much help to the family. Father can no longer speak; Mother is his constant provider of care—yet neither of them has uttered or written those words, “Why us?”

Let me return to Randy Spaulding’s actual words. He wrote: “One morning as I was thinking about the mundane things of life and hurrying out the door to begin the day, I happened to notice my father sitting in the corner of the room reading his scriptures. I stopped and went over to speak to him. I noticed the difficult circumstances he was under. With his right hand, he was trying to hold up his head enough to see me and read the Book of Mormon. I learned that at one of the most trying times, he still had enough faith to read about a God of love, a God of miracles who heals and makes us whole, and a God of life—eternal life. My father still believes. Oh, how I long to take him back in time to the Pool of Bethesda and to ask our Master if He would please have mercy on us, so that my father, also, could take up his bed and walk.”

His letter continued: “That day I returned to my bedroom and thanked my Heavenly Father for a father and mother second to none.”

Let us remember that it was not the waters of Bethesda’s pool which healed the impotent man. Rather, his blessing came through the touch of the Master’s hand. From the beautiful Psalm we learn: “Lord, thou hast heard the desire of the humble: thou wilt prepare their heart, thou wilt cause thine ear to hear.”

He has heard, and He indeed has blessed you and yours. An angel wife and mother who, without stint, sacrifices her own comfort for the blessing of her eternal companion; neighbors with hands that help, hearts that feel, and whose feet and talents all come quickly to rescue—are manifested blessings of the Lord’s promises. Though Bethesda beckons, the Lord has heard. Said He: “Verily, verily, I say unto you, even as you desire of me so it shall be unto you.”

President Harold B. Lee comforted us with these words: “[Those] who have been denied blessings … in this life—who say in their heart, if I could have done, I would have done, or I would give if I had, but I cannot for I have not—the Lord will bless you as though you had done, and the world to come will compensate for those who desire in their hearts the righteous blessings that they were not able to have because of no fault of their own.”

On every side there are those who suffer pain, who endure debilitating illness, who battle the demon of depression. Our hearts go out to all. Our prayers ascend in their behalf. Hands that help are extended.

I love the sentiment contained in the words of the poem entitled “Living What We Pray For”:

\begin{verse}
I knelt to pray when day was done\\
And prayed, “O Lord, bless everyone;\\
Lift from each saddened heart the pain,\\
And let the sick be well again.”
\end{verse}

When I read the phrase from this poem “I hid my face and cried,” the hallowed halls of memory prompt me to share a tender, personal account with you.

Long years ago, when I served as a bishop, I received notification that Mary Watson, a member of my ward, was a patient in the county hospital. When I went to visit her, I discovered her in a large room with so many beds that it was difficult to single her out. As I identified her bed and approached her, I said, “Hello, Mary.”

She replied, “Hello, Bishop.”

I noticed that a patient in the bed next to Mary Watson covered her face with the bedsheet.

I gave Mary Watson a blessing, shook her hand, and said, “Good-bye,” but I could not leave her side. It was as though an unseen hand were resting on my shoulder, and I felt within my soul that I was hearing these words: “Go over to the next bed where the little lady covered her face when you came in.” I did so. I have learned in my life never to postpone a prompting.

I approached the bedside of the other patient, gently tapped her shoulder, and carefully pulled back the sheet which had covered her face. Lo and behold! She too was a member of my ward. I had not known she was a patient in the hospital. Her name was Kathleen McKee. When her eyes met mine, she exclaimed through her tears, “Oh, Bishop, when you entered that door, I felt you had come to see me and bless me in response to my prayers. I was rejoicing inside to think that you would know I was here, but when you stopped at the other bed, my heart sank, and I knew that you had not come to see me.”

I said to Kathleen McKee: “It does not matter that I didn’t know you were here. It is important, however, that our Heavenly Father knew and that you had prayed silently for a priesthood blessing. It was He who prompted me to intrude on your privacy.”

A blessing was given; a prayer was answered. I bestowed a kiss on her forehead and left the hospital with gratitude in my heart for the promptings of the Spirit. It would be the last time I was to see Kathleen McKee in mortality—but not the last time I heard from her.

Upon her death, the hospital called with this message: “Bishop Monson, Kathleen McKee died tonight. She made arrangements that we were to notify you, should she pass away. She left for you a key to her basement apartment.”

Kathleen McKee had no immediate family. With my sweet wife accompanying me, I visited her humble apartment. I turned the key in the door, opened it, and switched on the light. There in her immaculate two-room apartment, I saw a small table with a note resting beneath an Alka-Seltzer bottle. The note, written in her own hand, said: “Bishop, my tithing is in this envelope, and the Alka-Seltzer bottle contains coins covering my fast offering. I am square with the Lord.” The receipts were written.

The sweetness of the night has not been forgotten. Tears of gratitude to God filled my very soul.

A message in a birthday card which I received a few weeks ago, from parents who last year lost a beautiful daughter to cancer, expresses this profound thought:

“‘And what is as important as knowledge?’ asked the mind.

“‘Caring and seeing with the heart,’ answered the soul.”

This expression describes Bethesda’s blessing. Of this divine truth I testify. In the name of Jesus Christ, amen.

\newpage
\settalk{Faith in Every Footstep: The Epic Pioneer Journey [Video Presentation]}
\setspeaker{Unknown}
\settalkdate{April 1997}
\talkheading{Faith in Every Footstep: The Epic Pioneer Journey [Video Presentation]}{Unknown}

\intalkheader{Narrator: President Gordon B. Hinckley}

The epic pioneer journey of the Latter-day Saints began on the banks of the Mississippi River. Here at Nauvoo they had transformed a swamp into a thriving community of commerce and fellowship. But Nauvoo was not to be a final home, merely a brief rest for a season. The severe persecution that had driven the Saints from Missouri again threatened their lives and their city. The Prophet Joseph Smith and Hyrum were martyred at Carthage Jail on June 27, 1844. Life in Nauvoo was drawing to a close.

Sunday, February 1, 1846, the Saints worshiped together in the City of Joseph. The next day, Brigham Young directed families to be ready to leave with only four hours’ notice.

In the bone-chilling cold of that bitter winter, the exodus began. Many of the Saints gathered their belongings and closed the doors of their dwellings for the last time as they turned to what lay across the river—and west.

Nauvoo is peaceful now. Homes and shops have been lovingly restored. This is a place that speaks of industry and commitment. I see their courage and craftsmanship as they built a city to God.

How the Saints must have felt, leaving so much behind—the fields they had cultivated, the trees they had planted, the temple they had built. The men, women, and children walked out of their beautiful homes, climbed aboard their wagons, drove down to the river, there to cross and move slowly over the soil of Iowa, looking back now and again at what they were leaving and would never see again.

Leaving Nauvoo was a remarkable act of faith. There was much of hardship ahead for these pioneers, but they had faith in their leaders and faith in the Lord and His goodness—faith that He would once again lead His people to the promised land, faith that they would not falter or fall. So they walked out into a wilderness, their journey marked by faith in every footstep.

\intalkheader{Narrator: President Thomas S. Monson, First Counselor in the First Presidency}

The way west was slow. Many were ill prepared for the grueling trek. Freezing temperatures, incessant rain, and mud up to the knees tried even the hardiest emigrant. They struggled for 131 days just to cross Iowa.

Like the army of Israel of old, they had their cloud by day and pillar of fire by night. Out of the travail of Iowa came the hymn that echoes down the generations: “Come, come, ye Saints, no toil nor labor fear; / But with joy wend your way” (Hymns, no. 30).

Stopping at Garden Grove and Mount Pisgah to set up stations for those who would follow, these faithful pioneers pressed on to the banks of the Missouri and temporary respite for the winter.

Here at Winter Quarters was Zion in the wilderness. President Brigham Young organized the people and pooled their meager resources. Yet despite all they could do, sickness and death stalked the camps.

This monument is placed directly over the graves of an unknown child and seven other pioneers. My heart is deeply touched as I realize just how high a price those noble Saints paid in responding to the call of the prophet to leave their homes and journey west.

So many struggled and lost so much. Truly, these noble pioneers walked a path of pain and a trail of tears. Their journey was over, but their names live on as testaments of their love of truth and faith in the Lord.

\intalkheader{Narrator: President James E. Faust, Second Counselor in the First Presidency}

When spring came that April of 1847, the Quorum of the Twelve, under the direction of Brigham Young, handpicked a vanguard company and left Winter Quarters with 143 men, 3 women, and 2 young boys, 72 wagons, 93 horses, 66 oxen, 52 mules, 19 cows, 17 dogs, and some chickens.

Between that refuge and the promise of Zion stood a vast plain and the fertile Platte River, their lifeline as they pushed farther into the American West. Moving across Nebraska, they marked the rolling miles and journeyed past Chimney Rock, a solitary formation jutting out of the prairie.

Into this land speckled with sage and air swirling with dust, tired oxen lumbered, wagon wheels creaked, brave men and women toiled, and occasionally wolves howled. Even today, signs of their crossing are carved into the landscape.

The pioneers left the North Platte and now followed the Sweetwater, a stream they would ford many times. Camping at the round outcropping called Independence Rock, a few of these 19th-century travelers left their names on the granite stone.

Past Independence Rock, the wagons skirted the side of Devil’s Gate, a deep gash in the hillside and often mentioned in their journals. The trail soon turned upward and increasingly rocky.

Here at Rocky Ridge is holy ground. This very spot is one of the highest points on the trail west. The pioneers who came over this ridge faced discouragement, some even death, as they inched their way up this sharp slope. I hold in my hand a square nail and a piece of metal jolted loose from a wagon or a handcart. Imagine facing this ridge in a wagon. Then imagine pulling a handcart.

For some, the punishing climb of Rocky Ridge would be fatal. The Martin and Willie Handcart Companies of 1856 were caught in early blizzards near this summit. Rescue came from Salt Lake but too late to save close to 200 souls who perished in the cold and deep snow.

Martin’s Cove sheltered many during that agonizing and poignant time. A memorial at Rock Creek honors those buried here for their faith in the face of enormous adversity.

In the heroic effort of the handcart pioneers, we learn a great truth. All must pass through a refiner’s fire, and the insignificant and unimportant in our lives can melt away like dross and make our faith bright, intact, and strong. There seems to be a full measure of anguish, sorrow, and often heartbreak for everyone, including those who earnestly seek to do right and be faithful. Yet this is part of the purging to become acquainted with God.

\intalkheader{Narrator: President Monson}

With the Wind River peaks to the north, the pioneer trail crossed South Pass—the only major break between mountain ranges and the most direct route to the Great Basin. Entering northeastern Utah, they worked their way slowly through Echo Canyon, a narrow passageway flanked by red, overhanging cliffs.

This final stretch would try what little strength was left. Ahead loomed a broken succession of hills piled on hills, and mountains in every direction. Hearts full of enthusiasm to be so near their journey’s end often sank as they knew there was only one way to go: up and over.

On this high summit they named Big Mountain, the pioneers gazed for the first time on their new home: a glistening mountain valley on the far horizon. What joy they must have felt! The countless sacrifices and struggles along the way were nearly over. The Salt Lake Valley was in sight. Although much hardship still lay ahead, they had endured. With feet worn and weary with fatigue, they had kept step with their faith.

Big Mountain holds a special place in my heart. A pioneer ancestor, Gibson Condie, came over this summit on his way to help rescue the stranded handcart pioneers. At the call of the prophet, he journeyed to this very spot in the bitter winter of 1856. The snow was 16 feet deep on the road. How grateful I am for this pioneer ancestor, who, leaving the comfort of home and family, risked his own safety to help those in such desperate need.

\intalkheader{Narrator: President Faust}

President Young arrived in the valley on a Saturday, July 24th. These pioneers had come so far and given so much, and they paused on the Sabbath to worship and give thanks for their safe arrival.

They came “one of a city, and two of a family” (Jer. 3:14) across a continent to a new life in the desert. What else but a divine restoration would prompt such an endeavor and require such a sacrifice? They had walked with faith, knowing that God lives and He knew where those steps would take them.

Now in this valley home, they took fresh courage for the tasks ahead. There were shelters to erect, land to cultivate, crops to plant, and the temple to build.

\intalkheader{Narrator: President Hinckley}

Rising above the Salt Lake Valley is a dome-shaped peak. Brigham Young saw it in a vision before the Saints left Nauvoo. He saw an ensign descend upon the hill and heard the voice of Joseph Smith say, “Build under [that] point … and you will prosper and have peace” (quoted by George A. Smith, in Deseret News [semiweekly], 29 June 1869, 3).

When Brigham Young first arrived in the valley, he immediately recognized the peak. On the morning of July 26, 1847, the men who would eventually comprise the new First Presidency, along with several members of the Twelve, climbed its slopes.

This small group of priesthood leaders gazed out upon the valley below. “This is whereon we will plant the soles of our feet,” President Young said, “and where the Lord will place his name among his people” (quoted by Erastus Snow, in Deseret News, 22 Oct. 1873, 5).

As I now stand at Ensign Peak and see the valley below, I marvel at the foresight of that little group. These prophets, dressed in old, travel-worn clothes, standing in boots they had worn for more than a thousand miles, spoke of a millennial vision. It was both bold and audacious. It was almost unbelievable.

Here they were, almost a thousand miles from the nearest settlement to the east and almost eight hundred miles from the Pacific coast. They were in an untried climate. They had never raised a crop here. They had not built a structure of any kind.

They were exiles, driven from their fair city on the Mississippi into this desert region of the West. But they were possessed of a vision drawn from the scriptures and words of revelation: “And he shall set up an ensign for the nations, and shall assemble the outcasts of Israel, and gather together the dispersed of Judah from the four corners of the earth” (Isa. 11:12).

This great pioneering movement of more than a century ago goes forward with latter-day pioneers. Today pioneer blood flows in our veins just as it did with those who walked west. It’s the essence of our courage to face modern-day mountains and our commitment to carry on. The faith of those early pioneers burns still, and nations are being blessed by latter-day pioneers who possess a clear vision of this work of the Lord.

The footsteps that made such a deep impression over the heartland of America make similar impressions in countries across the world—from Belgium to Brazil and France to the Philippines. Step by faithful step, we walk together toward a glorious destiny, building the kingdom of God on earth and preparing the minds and hearts of people everywhere to come unto Christ, the Redeemer and Savior of the world.

\newpage
\settalk{Pioneers All}
\setspeaker{Thomas S. Monson}
\settalkdate{April 1997}
\talkheading{Pioneers All}{Thomas S. Monson}

What a glorious sight you are. I realize that beyond this pioneer tabernacle many thousands are assembled in chapels and in other settings throughout much of the world. I pray for heavenly help as I respond to the opportunity to address you.

Your leaders have done so well tonight, but then we men realize that this is typical of the sisters. My congratulations go out to each of you who has had a part in the preparations for this conference and to those who participated on the program.

In his classic poem, Henry Wadsworth Longfellow described youth and the future. He wrote:

\begin{verse}
How beautiful is youth! how bright it gleams\\
With its illusions, aspirations, dreams!\\
Book of Beginnings, Story without End,\\
Each maid a heroine, and each man a friend!
\end{verse}

The First Presidency declared, on April 6, 1942: “How glorious and near to the angels is youth that is clean. This youth will have joy unspeakable here and eternal happiness hereafter.”

We’ve heard much about the pioneers of 1847 and their trek across the plains and entrance into the Salt Lake Valley. We shall hear more as this sesquicentennial year moves along.

Not surprisingly, as the pioneer theme is presented, each goes back in memory to his or her own family line. There are usually examples to identify and which fit the definition of a pioneer: “one who goes before, showing others the way to follow.” Some, if not all, made great sacrifices to leave behind comfort and ease and respond to that clarion call of their newly found faith.

Two of my own great-grandparents fit the mold of many. Gibson and Cecelia Sharp Condie lived in Clackmannan, Scotland. Their families were engaged in coal mining—at peace with the world, surrounded by relatives and friends, and housed in fairly comfortable quarters in a land they loved. They listened to the message of the missionaries from The Church of Jesus Christ of Latter-day Saints and were converted to the depths of their very souls. They heard the call to journey to Zion and knew they must answer that call.

They sold their possessions and prepared for a hazardous voyage across the mighty Atlantic Ocean. With five children, they boarded a sailing vessel, all their worldly possessions in a tiny trunk. They traveled 3,000 miles across the waters, eight long, weary weeks on a treacherous sea—night and day nothing but water—eight weeks of watching and waiting, with poor food, poor water, and no help beyond the length and breadth of that small sailing vessel.

In the midst of this soul-trying situation, their son, Nathaniel, sickened and died. My great-grandparents loved that son just as much as your parents love you; and when his eyes were closed in death, their hearts were torn asunder. To add to their grief, the law of the sea must be obeyed. Wrapped in a canvas weighed down with iron, his body was consigned to a watery grave. As they sailed away, only those parents knew the crushing blow dealt to wounded hearts. Gibson Condie and his good wife were comforted by the words “Not my will, but Thy will, O Father.”

That first trek of 1847, organized and led by Brigham Young, is described by historians as one of the great epics of United States history. Mormon pioneers by the hundreds suffered and died from disease, exposure, or starvation. There were some who, lacking wagons and teams, literally walked the 1,300 miles across the plains and through the mountains, pushing and pulling handcarts.

As the long, painful struggle approached its welcome end, a jubilant spirit filled each heart. Tired feet and weary bodies somehow found new strength.

Time-marked pages of a dusty pioneer journal speak movingly to us: “We bowed ourselves down in humble prayer to Almighty God with hearts full of thanksgiving to Him, and dedicated this land unto Him for the dwelling place of His people.”

We honor those who endured incredible hardships. We praise their names and reflect on their sacrifices.

What about our time? Are there pioneering experiences for us? Will future generations reflect with gratitude on our efforts, our examples? You young women, wherever you are this night, can indeed be pioneers in courage, in faith, in charity, in determination.

You can strengthen one another; you have the capacity to notice the unnoticed. When you have eyes to see, ears to hear, and hearts to feel, you can reach out and rescue others of your age.

From Proverbs comes the counsel “Ponder the path of thy feet.”

I hope that you young people recognize the strength and the power of your testimonies. Several years ago I was in the nation of Czechoslovakia. There, in an inspiring meeting held in Prague under dangerous circumstances and when freedom was curtailed, I met a young woman whose name is Olga. She was about 25 years of age at the time and had, in the previous two years, brought to membership in the Church 16 young men and young women her own age. As I met with them, I knew they were truly converted to the gospel. I felt they would be the foundation of the Church in Czechoslovakia. They learned the truth of the gospel and felt the strength of testimony—all from Olga. When I complimented Olga and thanked her for having a testimony she is willing to share, she said, “Oh, Brother Monson, I have 14 others with whom I am working!” Later I learned that almost all of those 14 became members of the Church. The light of Christ shone in Olga’s eyes as she encouraged others to “come unto him.”

My young sisters, we really don’t know how much good we can do until we put forth the effort. Our testimonies can penetrate the hearts of others and can bring to them the blessings which will prevail in this troubled world and which will guide them to exaltation.

Recently I heard from a teenaged friend, Jami Palmer, whom I have known for a number of years. When she was 12, she was diagnosed with cancer. She underwent grueling and painful treatments for many months. Today she is bright, beautiful, and looking to the future with confidence and with faith.

In one of her darkest hours, when any future appeared somewhat grim, she learned that she must undergo months of chemotherapy, followed by an 11-hour surgery to save her leg. A long-planned hike with her Young Women class up to Timpanogos Cave was out of the question—she thought. Jami told her friends they would have to undertake the hike without her. Surely there was a catch in her voice and disappointment in her heart. But then the other young women responded emphatically, “No, Jami. You are going with us!”

“But I can’t walk,” came the anguished reply.

“Then, Jami, we’ll carry you to the top!” And they did.

The hike is now a memory, but in reality it is much more. James Barrie, the Scottish poet, declared, “God gave us memories, that we might have June roses in the December of our lives.” None of those precious young women will ever forget that memorable day when, I am confident, a loving Heavenly Father looked down with a smile of approval and was well pleased.

Today Jami is an accomplished pianist, vocalist, and athlete. She is an officer and spokesperson for the Make-A-Wish Foundation.

In preparing to speak to you tonight, I turned to the scriptures for inspiration. The word come, I discovered, was frequently used. The Lord said, “Come unto me.” He invited, “Come, … learn of me,” and then, “Come, follow me.” My plea is that we would come to the Lord.

I counsel you to honor your father and your mother. May I share with you an example of honoring one’s mother. Some years ago Ruth Fawson, mother of six, underwent life-threatening surgery. Her devoted husband and her three sons and three daughters were all at the hospital. The physicians and nurses explained to the family that they could return to their homes and that the staff was prepared to care adequately for Sister Fawson. The family expressed their thanks to the hospital staff but indicated a determination for at least one of its number to be present at all times. A daughter expressed the feelings of all: “We wanted to be there when Mother awakened and stretched forth her hand, so that it would be our hands she would grasp, it would be our smiles she would see, it would be our words she would hear, it would be our love she would feel.” “Honour thy father and thy mother.”

In the Clarkston, Utah, cemetery, Martin Harris, one of the Three Witnesses of the Book of Mormon, is buried. Behind his imposing and beautiful monument are the graves of others. One contains the tender inscription: “A light from our household is gone; a voice we loved is stilled. A place is vacant in our hearts that never can be filled.”

My dear young sisters, don’t wait until that light from your household is gone; don’t wait until that voice you love is stilled before you say, “I love you, Mother; I love you, Father.” Now is the time to think and the time to thank. I trust you do both.

Essential to your success and happiness is the advice “Choose your friends with caution.” In a survey made in selected wards and stakes of the Church, we learned a most significant fact: those persons whose friends married in the temple usually married in the temple, while those persons whose friends did not marry in the temple usually did not marry in the temple. The influence of one’s friends appeared to be a highly dominant factor—even more so than parental urging, classroom instruction, or proximity to a temple.

I am pleased that many of your leaders from Young Women are here or are viewing and listening in so many locations. I paraphrase a well-known poem originally written to leaders of boys. I feel this poem is worthy for you and your young women:

\begin{verse}
She stood at the crossroads all alone,\\
The sunlight in her face.\\
She had no thought for the world unknown—\\
She was set for a noble race.\\
But the roads stretched east and the roads stretched west,\\
And the girl knew not which road was best,\\
So she chose the road that led her down,\\
And she lost the race and the victor’s crown.\\
She was caught at last in an angry snare\\
Because no one stood at the crossroads there\\
To show her the better road.
\end{verse}

Noble leaders of young women, you stand at the crossroads in the lives of those whom you teach. Inscribed on the wall of Stanford University Memorial Hall is this truth: “We must teach our youth that all that is not eternal is too short, and all that is not infinite is too small.”

President Hinckley emphasized our responsibilities when he declared: “In this work there must be commitment. There must be devotion. We are engaged in a great eternal struggle that concerns the very souls of the sons and daughters of God. We are not losing. We are winning. We will continue to win if we will be faithful and true. … There is nothing the Lord has asked of us that in faith we cannot accomplish.”

A human drama illustrating the bond between the teacher and the young women in her class has been an inspiration to me, as I know it will be to you. It is the account of a first-year Beehive in Young Women. I share it with you, using her own words:

“One day, a few months before my 12th birthday, I noticed a note card on the dresser of the room I shared with my older sister. It read, ‘I’m happy to be your teacher and hope that we have a great year in Mutual.’ It was signed ‘Baur Dee.’

“I soon learned that all of the girls loved Baur Dee. They visited her at home, wanted to sit with her in church, and stayed after Mutual each Wednesday to talk with her.

“Looking back so many years, I am amazed that I still have such a vivid memory of my earliest real meeting with Baur Dee. That first night, as I walked in the front door of our ward building to attend Mutual, she stood waiting to greet me. I noticed for the first time the smile which always transformed her appearance from average to beautiful. ‘Welcome,’ she said to me. ‘I’m so glad you’re in my class. We’re going to have a great time!’ There was no adjustment period for me from Primary to Mutual. I felt right at home from that moment.

“Over the next few weeks, I joined the other girls as one of Baur Dee’s fans. At the time, I didn’t try to figure out her popularity. So many years later, though, I believe I understand. She really, truly cared about each one of us, and we knew it.

“Baur Dee suffered from a disease called nephritis—a disease which not too many years later would be treated with dialysis and often cured with a kidney transplant. But for Baur Dee there was no cure, no miracle. She passed away peacefully. She was 27 years old.

“After the funeral services, as we girls stood somberly around the open grave at the cemetery, we made a vow that we would visit Baur Dee’s final resting place together every Memorial Day throughout our lives and that we would never, ever allow her memory to die.”

Forty years have gone by since Baur Dee, this teacher of girls, passed away—yet the pledge lives on. One of her girls has said: “Wherever I go, whatever I do, something of Baur Dee goes with me and with each of her ‘girls.’ She lives on in us and in those with whom we have shared her lessons.” As Henry Brooks Adams observed, “A teacher affects eternity; [she] can never tell where [her] influence stops.”

Tonight, may all who hear my voice know that this work is of our Heavenly Father. He loves you. He hears your prayers. He knows your thoughts and actions. I testify that Christ is our Redeemer. I know that President Gordon B. Hinckley is God’s prophet.

I close with a scriptural passage, from Alma in the Book of Mormon, which expresses my love for you: “I perceive that ye are in the paths of righteousness; I perceive that ye are in the path which leads to the kingdom of God.”

To all of you noble pioneers who go before, showing others the way to follow, I urge, “Carry on.” In the name of Jesus Christ, amen.

\newpage
\settalk{They Showed the Way}
\setspeaker{Thomas S. Monson}
\settalkdate{April 1997}
\talkheading{They Showed the Way}{Thomas S. Monson}

This year, 1997, commemorates the 150th anniversary since the pioneers, under the inspired leadership of Brigham Young, entered the valley of the Great Salt Lake and proclaimed: “This is the right place. Drive on.” Much will be said at this conference concerning that epochal event, and thanks will be given to God for His watchful care and guidance.

On this beautiful Sabbath morning I wish to make a few remarks concerning “other pioneers” who preceded that trek. In doing so, I pause and ponder the dictionary definition of the word pioneer: “one who goes before, [showing] others [the way] to follow.”

Let us turn back the clock of time and journey to other places, that we might review several who I feel meet the high standard of the word pioneer.

Such a one was Moses. Raised in Pharaoh’s court and learned in all the wisdom of the Egyptians, he became mighty in words and deeds. One cannot separate Moses, the great lawgiver, from the tablets of stone provided him by God and on which were written the Ten Commandments. They were binding then; they are binding now.

Moses endured constant frustration as some of his trusted followers returned to their previous ways. Though he was disappointed in their actions, yet he loved them and led them, even the children of Israel, from their Egyptian bondage. Certainly Moses qualifies as a pioneer.

Another who qualifies is Ruth, who forsook her people, her kindred, and her country in order to accompany her mother-in-law, Naomi—worshiping Jehovah in His land and adopting the ways of His people. How very important was Ruth’s obedience to Naomi and the resulting marriage to Boaz by which Ruth—the foreigner and a Moabite convert—became a great-grandmother of David and therefore an ancestress of Jesus Christ. The book of the Holy Bible that bears her name contains language poetic in style, reflective of her spirit of determination and courage. “And Ruth said, Intreat me not to leave thee, or to return from following after thee: for whither thou goest, I will go; and where thou lodgest, I will lodge: thy people shall be my people, and thy God my God: Where thou diest, will I die, and there will I be buried: the Lord do so to me, and more also, if ought but death part thee and me.”

Yes, Ruth, precious Ruth, was a pioneer.

Other faithful women also qualify, such as Mary, the mother of Jesus; Mary Magdalene; Esther; and Elisabeth. Let us not overlook Abraham, Isaac, and Jacob, nor fail to include Isaiah, Jeremiah, Ezekiel, and some from a later period.

We remember John the Baptist. His clothing was simple, his life spartan, his message brief: faith, repentance, baptism by immersion, and the bestowal of the Holy Ghost by an authority greater than that possessed by himself. He declared, “I am not the Christ, but … I am sent before him.” “I indeed baptize you with water; but one mightier than I cometh … : he shall baptize you with the Holy Ghost and with fire.”

The River Jordan marked the historic meeting place when Jesus came down from Galilee to be baptized of John. At first John pleaded with the Master: “I have need to be baptized of thee, and comest thou to me?” Came the response: “It becometh us to fulfil all righteousness. … And Jesus, when he was baptized, went up straightway out of the water: and, lo, the heavens were opened unto him, and he saw the Spirit of God descending like a dove, and lighting upon him: “And lo a voice from heaven, saying, This is my beloved Son, in whom I am well pleased.”

John freely declared and taught, “Behold the Lamb of God, which taketh away the sin of the world.”

Of John the Lord declared, “Among them that are born of women there hath not risen a greater than John the Baptist.” Like so many other pioneers through the annals of history, John wore the martyr’s crown.

Many who were pioneers in spirit and action were called by Jesus to be His Apostles. Much could be told of each.

Peter was among the first of Jesus’ disciples. Peter the fisherman, in response to a divine call, laid aside his nets and hearkened to the Master’s declaration: Come “follow me, and I will make you [a fisher] of men.” I never think of Peter without admiring his testimony of the Lord: “Thou art the Christ, the Son of the living God.”

John the Beloved is the only one of the Twelve recorded as being at the Crucifixion of Christ. From the cruel cross, Jesus uttered the magnificent charge to John, referring to His mother, Mary: “Behold thy mother,” and to Mary, “Behold thy son.”

The Apostles went before, showing others the way to follow. They were pioneers.

History records, however, that most men did not come unto Christ, nor did they follow the way He taught. Crucified was the Lord, slain were most of the Apostles, rejected was the truth. The bright sunlight of enlightenment slipped away, and the lengthening shadows of a black night enshrouded the earth.

Generations before, Isaiah had prophesied, “Darkness shall cover the earth, and gross darkness the people.” Amos had foretold of a famine in the land, “not a famine of bread, nor a thirst for water, but of hearing the words of the Lord.” The dark ages of history seemed never to end. Would no heavenly messengers make their appearance?

In due time honest men with yearning hearts, at the peril of their very lives, attempted to establish points of reference, that they might find the true way. The day of the Reformation was dawning, but the path ahead was difficult. Persecutions would be severe, personal sacrifice overwhelming, and the cost beyond calculation. The reformers were pioneers, blazing wilderness trails in a desperate search for those lost points of reference which, they felt, when found would lead mankind back to the truth Jesus taught.

John Wycliffe, Martin Luther, Jan Hus, Zwingli, Knox, Calvin, and Tyndale all pioneered the period of the Reformation. Significant was the declaration of Tyndale to his critics: “I will cause a boy that driveth the plough shall know more of the scripture than thou doest.”

Such were the teachings and lives of the great reformers. Their deeds were heroic, their contributions many, their sacrifices great—but they did not restore the gospel of Jesus Christ.

Of the reformers, one could ask: “Was their sacrifice in vain? Was their struggle futile?” I answer with a reasoned “no.” The Holy Bible was now within the grasp of the people. Each person could better find his or her way. Oh, if only all could read and all could understand! But some could read, and others could hear, and all had access to God through prayer.

The long-awaited day of restoration did indeed come. But let us review that significant event in the history of the world by recalling the testimony of the plowboy who became a prophet, the witness who was there—even Joseph Smith.

Describing his experience, Joseph said: “I was one day reading the Epistle of James, first chapter and fifth verse, … If any of you lack wisdom, let him ask of God, that giveth to all men liberally, and upbraideth not; and it shall be given him.”

“At length I came to the conclusion that I must either remain in darkness and confusion, or else I must do as James directs, that is, ask of God. …

“… I retired to the woods to make the attempt. It was on the morning of a beautiful, clear day, early in the spring of eighteen hundred and twenty. …

“… I kneeled down and began to offer up the desires of my heart to God. …

“… I saw a pillar of light exactly over my head, above the brightness of the sun, which descended gradually until it fell upon me.

“… When the light rested upon me I saw two Personages, whose brightness and glory defy all description, standing above me in the air. One of them spake unto me, calling me by name and said, pointing to the other—This is My Beloved Son. Hear Him!”

The Father and the Son, Jesus Christ, had appeared to Joseph Smith. The morning of the dispensation of the fulness of times had come, dispelling the darkness of the long generations of spiritual night.

Volumes have been written concerning the life and accomplishments of Joseph Smith, but for our purposes here today perhaps a highlight or two will suffice: He was visited by the angel Moroni. He translated, from the precious plates to which he was directed, the Book of Mormon, with its new witness of Christ to all the world. He was the instrument in the hands of the Lord through whom came mighty revelations pertaining to the establishment of The Church of Jesus Christ of Latter-day Saints. In the course of his ministry he was visited by John the Baptist, Moses, Elijah, Peter, James, and John, that the restoration of all things might be accomplished. He endured persecution; he suffered grievously, as did his followers. He trusted in God. He was true to his prophetic calling. He commenced a marvelous missionary effort to the entire world, which today brings light and truth to the souls of mankind. At length, Joseph Smith died the martyr’s death, as did his brother Hyrum.

Joseph Smith was a pioneer indeed.

Turning the pages of scriptural history from beginning to end, we learn of the ultimate pioneer—even Jesus Christ. His birth was foretold by the prophets of old; His entry upon the stage of life was announced by an angel. His life and His ministry have transformed the world.

With the birth of the babe in Bethlehem, there emerged a great endowment, a power stronger than weapons, a wealth more lasting than the coins of Caesar. This child was to be the King of Kings and Lord of Lords, the Promised Messiah, even Jesus Christ, the Son of God. Born in a stable, cradled in a manger, He came forth from heaven to live on earth as mortal man and to establish the kingdom of God. During His earthly ministry, He taught men the higher law. His glorious gospel reshaped the thinking of the world. He blessed the sick. He caused the lame to walk, the blind to see, the deaf to hear. He even raised the dead to life.

One sentence from the book of Acts speaks volumes: Jesus “went about doing good, … for God was with him.”

He taught us to pray: “Our Father which art in heaven, Hallowed be thy name. Thy kingdom come. Thy will be done in earth, as it is in heaven.”

In the garden known as Gethsemane, where His suffering was so great that blood came from His pores, He pleaded as He prayed, “Father, if thou be willing, remove this cup from me: nevertheless not my will, but thine, be done.”

He taught us to serve: “Inasmuch as ye have done it unto one of the least of these my brethren, ye have done it unto me.”

He taught us to forgive: “I, the Lord, will forgive whom I will forgive, but of you it is required to forgive all men.”

He taught us to love: “Thou shalt love the Lord thy God with all thy heart, and with all thy soul, and with all thy mind. This is the first and great commandment. And the second is like unto it, Thou shalt love thy neighbour as thyself.” Like the true pioneer He was, He invited, “Come, follow me.”

Let us turn to Capernaum. There Jairus, a ruler of the synagogue, came to the Master, saying, “My little daughter lieth at the point of death: I pray thee, come and lay thy hands on her, that she may be healed; and she shall live.” Then came the news from the ruler’s house: “Thy daughter is dead.”

Christ responded, “Be not afraid, only believe.” He came to the house, passed by the mourners, and said to them: “Why make ye this ado, and weep? the damsel is not dead, but sleepeth. And they laughed him to scorn,” knowing that she was dead. “He … put them all out. … And he took [her] by the hand, and said unto her, … Damsel, I say unto thee, arise. … And they were astonished.”

It is emotionally draining for me to recount the events leading up to the Crucifixion of the Master. I cringe when I read the words of Pilate responding to cries of the throng: “Crucify him. … Crucify him.” Pilate “took water, and washed his hands before the multitude, saying, I am innocent of the blood of this just person: see ye to it.” Jesus was mocked. He was spit upon and a crown of thorns placed upon His head. He was given vinegar to drink. They crucified Him.

His body was placed in a borrowed tomb, but no tomb could hold the body of the Lord. On the morning of the third day came the welcome message to Mary Magdalene, to Mary the mother of James, and to other women who were with them as they came to the tomb, saw the large entrance stone rolled away, and noted the tomb was empty. Two angels said to the weeping women: “Why seek ye the living among the dead? He is not here, but is risen.”

Yes, the Lord had indeed risen. He appeared to Mary; He was seen by Cephas, or Peter, then by His brethren of the Twelve. He was seen by Joseph Smith, who declared: “This is the testimony, last of all, which we give of him: That he lives! For we saw him, even on the right hand of God.”

Our Mediator, our Redeemer, our Brother, our Advocate with the Father died for our sins and the sins of all mankind. The Atonement of Jesus Christ is the foreordained but voluntary act of the Only Begotten Son of God. He offered His life as a redeeming ransom for us all.

His mission, His ministry among men, His teachings of truth, His acts of mercy, His unwavering love for us prompt our gratitude and warm our hearts. Jesus Christ, Savior of the world—even the Son of God—was and is the ultimate pioneer, for He has gone before, showing all others the way to follow. May we ever follow Him, I pray, in the name of Jesus Christ, amen.

\newpage
\settalk{They Will Come}
\setspeaker{Thomas S. Monson}
\settalkdate{April 1997}
\talkheading{They Will Come}{Thomas S. Monson}

Several years ago an unusual motion picture swept the theaters in this and in other lands. It was entitled Field of Dreams and was the story of a young man who revered the baseball players of his youth and, from this foundation, carved out a large section from his cornfield and located there a full-blown baseball diamond. People mocked his foolishness and ridiculed his lack of common sense. The film goes on to show the many challenges he faced in completing the project and readying the baseball diamond for view. His was not an easy task. During the period of doubt as to the future success of his dream, he was driven by the reassuring words, “If you build it, [they] will come.” And come they did. Travelers by the thousands visited this unique place, which was filled with baseball’s many memories.

Lately, I have reflected on the importance of building a bridge to the heart of a person. I think of the nearly 55,000 full-time missionaries from our faith who are assigned over much of the world with the divine commission to teach, to testify, and to baptize. Theirs is a bridge-building task, awesome to behold and somewhat overwhelming to contemplate. With God’s mandate ringing in their ears, with the Lord’s instruction penetrating their hearts, they move forward in their lofty callings. They ponder the Lord’s words: “Remember the worth of souls is great in the sight of God.”

“Go ye therefore, and teach all nations, baptizing them in the name of the Father, and of the Son, and of the Holy Ghost:

“Teaching them to observe all things whatsoever I have commanded you: and, lo, I am with you alway, even unto the end of the world.”

Last year was the centennial of Utah statehood, and many ambassadors from other countries made a visit to our state capitol and also to the Church Administration Building. Many also toured the Missionary Training Center at Provo, Utah. They visited the classes of learning; they heard the testimonies of those going to their respective fields of labor. They marveled at the language proficiency, faith, and love exhibited by the missionaries. One ambassador stated, “I observed a sense of purpose, a commitment to prepare and to serve, and a joyful heart in each missionary.”

These missionaries go forward with faith. They know their duty. They understand that they are a vital link for the persons they will meet as missionaries and in the teaching and testifying they will experience as they bring others to the truth of the gospel.

They yearn for more persons to teach. They pray for the essential help each member can give to the conversion process.

The decision to change one’s life and come unto Christ is, perhaps, the most important decision of mortality. Such a dramatic change is taking place daily throughout the world.

Alma chapter 5, verse 13, describes this personal miracle: “And behold, … a mighty change was … wrought in their hearts, and they humbled themselves and put their trust in the true and living God.”

The covenant of baptism spoken of by Alma causes all of us to probe the depths of our souls:

“Now, as ye are desirous to come into the fold of God, and to be called his people, and are willing to bear one another’s burdens, that they may be light;

“Yea, and are willing to mourn with those that mourn; yea, and comfort those that stand in need of comfort, and to stand as witnesses of God at all times and in all things, and in all places …

“Now I say unto you, if this be the desire of your hearts, what have you against being baptized in the name of the Lord, as a witness before him that ye have entered into a covenant with him, that ye will serve him and keep his commandments, that he may pour out his Spirit more abundantly upon you?”

Our studies reveal that most of those who embrace the message of the missionaries have had other exposures to The Church of Jesus Christ of Latter-day Saints—perhaps hearing the magnificent Tabernacle Choir perform, maybe reading and viewing press reports of our well-traveled President Gordon B. Hinckley and his skillful participation in broad-ranging interviews, or just knowing another person who is a member and for whom respect exists. We, as members, should be at our best. Our lives should reflect the teachings of the gospel, and our hearts and voices ever be ready to share the truth.

Fellowshipping of the investigator should begin well before baptism. The teachings of the missionaries often need the second witness of a new convert to the Church. It has been my experience that such a witness, borne from the heart of one who has undergone this mighty change himself, brings resolve and commitment. When I served as mission president in eastern Canada, we found that in Toronto, as well as in most of the cities of Ontario and Quebec, there was no dearth of willing helpers to accompany the missionaries and to fellowship the investigators, welcome them to meetings, and introduce them to the ward or branch officers and members. Fellowshipping, friendshipping, and reactivating are ongoing in the daily life of a Latter-day Saint.

Each new convert should be provided a calling in the Church. Such brings interest, stability, and growth. The task may be somewhat simple, such as that given to Jacob de Jager when he and his family became members in Toronto. He held lofty posts in business, but his first calling in the Church was to put the hymnbooks in place along the pews. He took his assignment seriously. In recollecting this first calling, he said, “I had to be present each week, or the hymnbooks would remain undistributed.” As you know, Elder de Jager later served many years as a member of the First Quorum of the Seventy. Though he had many demanding responsibilities as a General Authority, he never forgot his first calling in the Church.

The unseen hand of the Lord guides the efforts of those who strive to learn and live the truth of the gospel. As a mission president, I received a weekly letter from each missionary. One that pleased me greatly came from a young elder serving in Hamilton. He and his companion were working with a lovely family, a young couple with two children. The couple felt that the message was true, and they could not deny their desire to be baptized. The wife, however, worried about her mother and father in faraway western Canada, fearing she and her husband would be disowned by her parents for joining the Church. She took pen in hand and jotted a note to her parents in Vancouver. The note read something like this:

“Dear Mother and Father,

“I want to thank you with all of my heart for your kindness and for your understanding and for the teachings which you gave me in my youth. John and I have come across a great truth, The Church of Jesus Christ of Latter-day Saints. We have studied the discussions, and our baptism will take place next Saturday night. We hope you will understand. In fact, we hope that you will welcome the missionaries in your home as we welcomed them in ours.”

The letter was sealed with a tear, a stamp was affixed, and it was mailed to Vancouver. On the very day it was received in Vancouver, the couple in Hamilton received a letter from the wife’s mother and father. They wrote:

“We are far away from you, or we would surely talk to you in person. We want you to know that missionaries from The Church of Jesus Christ of Latter-day Saints have called at our home, and we cannot deny the validity of their message. We have set a date for our baptism to take place next week. We hope you will understand and not be unduly critical of our decision. This gospel means so much to us and has brought such happiness into our lives that we pray someday you might also agree to learn more about it.”

Can you imagine what happened when the couple in Hamilton received that letter from the wife’s parents? They phoned Mother and Dad, and there were many tears of joy shed. I am sure there was a long-distance embrace, for both families became members of the Church.

You see, our Heavenly Father knows who we are, His sons and His daughters. He wants to bring into our lives the blessings for which we qualify, and He can do it. He can accomplish anything.

A visible and tender act of fellowshipping was witnessed in the ancient city of Rome. Some years ago, Sister Monson and I met with over 500 members there in a district conference. The presiding officer at that time was Leopoldo Larcher, a wonderful Italian. His brother had been working as a guest employee in the auto plants in Germany when two missionaries taught him the gospel. He went back to Italy and taught the gospel to his brother. Leopoldo accepted and sometime later became the president of the Italy Rome Mission and then the Italy Catania Mission.

During that meeting, I noticed that in the throng were many who were wearing a white carnation. I said to Leopoldo, “What is the significance of the white carnation?”

He said, “Those are new members. We provide a white carnation to every member who has been baptized since our last district conference. Then all the members and the missionaries know that these people are especially to be fellowshipped.”

I watched those new members being embraced, being greeted, being spoken to. They were no more strangers nor foreigners; they were “fellowcitizens with the saints, and of the household of God.”

Beyond the new convert to the Church are some who have drifted from that pathway which upward leads and, for one reason or another, have become less active for months, even years. Perhaps they were not fellowshipped; maybe friends departed from their lives. Whatever the reason, the fact remains: We need them, and they need us. Missionaries can effectively visit the homes where these individuals reside. When they approach, those within the shelter of home may come to remember the glorious feelings which came over them when they first heard the principles of the gospel taught to them. The missionaries can teach such individuals and witness the changes which come into their lives as they return to activity.

They need friends with testimonies. They need to know that we truly care for the one.

Aaronic Priesthood quorum advisers and Young Women teachers are on the line of battle, and miracles are within their grasp. Who is the teacher you best remember from your youth? I would guess that in all probability it was the one who knew your name, who welcomed you to class, who was interested in you as a person, and who truly cared. When a leader walks the pathway of mortality with a precious youth alongside, there develops a bond of commitment between the two that shields the youth from the temptations of sin and keeps him or her walking steadfastly on the path that leads onward, upward, and unswervingly to eternal life. Build a bridge to each youth.

All of us here and abroad this evening must answer the call of our prophet, President Gordon B. Hinckley, to spare no effort in fellowshipping and reactivating those who need our help, our labors, and our testimonies.

Let me share with you visits to two stake conferences where I evidenced the miracle which can take place when we take to heart the words of the pioneer hymn “Put Your Shoulder to the Wheel.”

One visit was to the Millcreek stake in Salt Lake City some years ago. Just over 100 brethren who were prospective elders had been ordained elders during the preceding year. I asked President James Clegg the secret of his success. He was too modest to take the credit. His counselor revealed that President Clegg, recognizing the challenge, had undertaken to personally call and arrange a private appointment between him and each prospective elder. During the appointment, President Clegg would mention the temple of the Lord, the saving ordinances and covenants emphasized there, and would conclude with this question: “Wouldn’t you desire to take your sweet wife and your children to the house of the Lord, that you might be a forever family throughout all eternity?” An acknowledgment followed, the reactivation process was pursued, and the goal was obtained.

The other visit was to the North Carbon stake in Price, Utah, also many years ago. I noted during my visit that they had rescued 86 men from the prospective elders in one year and had taken them and their wives to the Manti Temple. I said to Cecil Broadbent, the president, “How did you do it, President?”

He said, “I didn’t. My counselor, President Judd, did.”

President Judd was a large, ruddy-faced Welsh coal miner. I said to him, “President Judd, will you tell me how you were able to rescue 86 brethren in one year?”

I sat anticipating his answer, and he said, “No!”

I was stunned. I’d never had anyone say no so directly in my life. I asked, “Why not?”

He said, “Then you’ll tell the other stake presidents you visit, and we won’t lead the Church in reactivation.” He was smiling, though, so I knew it was half in jest. He said, “I’ll make a deal with you, Brother Monson. I’ll tell you how we rescued 86 men in one year if you’ll get me two tickets to general conference.”

I said, “You’re on!” And so he told me. What he didn’t tell me is that he intended to collect interest every conference for the next 10 years. He came faithfully every six months for his two tickets.

In both the Millcreek and the North Carbon stakes, as well as in others which have been successful in this phase of the work, four principles have prevailed:

In building a bridge to the investigator, the new convert, or the less-active member, when we do our part the Lord does His. I testify concerning this truth.

When I served as a bishop, I noted one Sunday morning that one of our priests was missing from priesthood meeting. I left the quorum in the care of the adviser and visited Richard’s home. His mother said he was working at the West Temple Garage.

I drove to the garage in search of Richard and looked everywhere, but I could not find him. Suddenly I had the inspiration to gaze down into the old-fashioned grease pit situated at the side of the station. From the darkness I could see two shining eyes. Then I heard Richard say, “You found me, Bishop! I’ll come up.” After that he rarely missed a priesthood meeting.

The family moved to a nearby stake. Time passed, and I received a phone call informing me that Richard had been called to serve a mission in Mexico, and I was invited by the family to speak at his farewell testimonial. At the meeting, when Richard responded, he mentioned that the turning point in his determination to fill a mission came one Sunday morning—not in the chapel, but as he gazed up from the depths of a dark grease pit and found his quorum president’s outstretched hand.

Through the years, Richard has stayed in touch with me, telling of his testimony, his family, and his faithful service in the Church, including his calling as a bishop.

My beloved brethren, let us, with faith unwavering and with love unstinting, be bridge builders to the hearts of those with whom we labor. As in the movie Field of Dreams, if we build it, they will come. Of this truth I testify, in the name of Jesus Christ, amen.

\newpage
\settalk{Home Teaching—a Divine Service}
\setspeaker{Thomas S. Monson}
\settalkdate{October 1997}
\talkheading{Home Teaching—a Divine Service}{Thomas S. Monson}

This has been a conference session marked by spirituality, and I know that you, as I, have been edified. The statement has been made, “Where the President is, there is strength, and to know that he is with us and is presiding is a strength to the entire Church.”

President Hinckley programmed an energy-consuming regimen during the past year and has borne his witness to vast thousands of members and others throughout the nations of the world. For many, the experience was one never before enjoyed by devoted members in faraway places with strange-sounding names. He appreciates our prayers in his behalf.

In addition to so many other responsibilities, the President of the Church receives a great deal of correspondence each day. I am reminded of one such letter and share it with you. I have changed the name of the young man who wrote the letter. It begins:

“Dear President,

“Hi. My name is David Smith. I live in an area where the starlings are very bad, and they are making nests in my step-grandpa’s boat and in my dad’s barn and all over the place. My step-grandpa and my dad both think I should shoot them, but my mom doesn’t. I know the law says it is okay, but I am not asking your opinion as a hunter. I am asking your opinion as a Church leader.

“Sincerely, David Smith

“P.S. A starling is a black bird that eats other bird’s eggs and other bad things.”

Each letter which comes in is answered. A response to this particular letter was sent by the secretary to the First Presidency, F. Michael Watson.

“Dear David:

“I have been asked to acknowledge your letter of April 30 addressed to the President of the Church about the problems you have been having with starlings.

“The Church does not have an official policy on this matter. The Brethren feel it should be left up to your parents to give you appropriate guidance.

“I hope this information is helpful to you.

“Sincerely yours, F. Michael Watson”

President Hinckley cannot possibly answer every letter personally, nor can he be everywhere. Neither can those of us who assist him reach each member in every nation. However, the wisdom of the Lord provided us guidelines whereby we who hold the priesthood of God can serve, can teach, can testify to the families of the Church. Yes, I speak of home teaching.

Let us review the counsel of the Lord and His prophets concerning this vital endeavor.

The bishop of each ward in The Church of Jesus Christ of Latter-day Saints assigns priesthood holders as home teachers to visit the homes of members every month. They go in pairs; often a youth holding the Aaronic Priesthood accompanies an adult holding the Melchizedek Priesthood.

The home teaching program is a response to modern revelation commissioning those ordained to the priesthood to “teach, expound, exhort, baptize, and watch over the church; … and visit the house of each member, and exhort them to pray vocally and in secret and attend to all family duties; … to watch over the church always, and be with and strengthen them; and see that there is no iniquity in the church, neither hardness with each other, neither lying, backbiting, nor evil speaking.”

President David O. McKay admonished: “Home teaching is one of our most urgent and most rewarding opportunities to nurture and inspire, to counsel and direct our Father’s children. … It is a divine service, a divine call. It is our duty as home teachers to carry the divine spirit into every home and heart. To love the work and do our best will bring unbounded peace, joy, and satisfaction to a noble, dedicated teacher of God’s children.”

From the Book of Mormon, Alma “consecrated all their priests and all their teachers; and none were consecrated except they were just men. Therefore they did watch over their people, and did nourish them with things pertaining to righteousness.”

In performing our home teaching responsibilities, we are wise if we learn and understand the challenges of the members of each family. A home teaching visit is also more likely to be successful if an appointment is made in advance.

The late John R. Burt, with whom I served for many years in ward and stake positions, told of an experience when as a boy he went home teaching with a devout and outspoken high priest—without warning—to a less-active family. They had come at a bad time. A poker game was under way in a smoke-filled living room. As the home teachers viewed the room, the high priest senior companion turned to young Brother Burt and exclaimed, “This congregation needs repentance! Please lead us in singing a hymn.”

Instead, the junior companion said, “I think we had best leave and come back another night.”

Some years ago, when the Missionary Executive Committee was comprised of Spencer W. Kimball, Gordon B. Hinckley, and Thomas S. Monson, Brother and Sister Hinckley hosted a dinner for the committee members and our wives. We had just finished a lovely dinner in the beautiful home—which Brother Hinckley constructed and on which he did most of the actual work—when suddenly there was a knock at the door. President Hinckley opened the door and noted his home teacher standing there. The home teacher said, “I don’t have with me my companion, but I felt I should come tonight. I didn’t know you would be entertaining company.”

President Hinckley graciously invited the home teacher to come in and sit down and instruct three Apostles and their wives concerning our duty as members. With a bit of trepidation, the home teacher did his best. President Hinckley thanked him for coming, after which the home teacher made a prompt retreat.

Abraham Lincoln offered this wise counsel, which surely applies to home teachers: “If you would win a man to your cause, first convince him that you are his sincere friend.” President Ezra Taft Benson urged, “Above all, be a genuine friend to the individuals and families you teach.”

As the Savior declared to us, “I will call you friends, for you are my friends.” A friend makes more than a dutiful visit each month. A friend is more concerned about helping people than getting credit. A friend cares. A friend loves. A friend listens. And a friend reaches out.

Some who are here will remember the story President Romney used to tell about the so-called home teacher who once went to the Romney home on a cold night. He kept his hat in his hand and shifted nervously when invited to sit down and give his message. “Well, I’ll tell you, Brother Romney,” he responded. “It’s cold outside, and I left my car engine running so it wouldn’t stop. I just came by so I could tell the bishop I made my calls.”

Brother Romney, after relating this experience in a meeting of priesthood holders, then said, “We can do better than that, brethren—much better!”

Home teaching answers many prayers and permits us to see the occurrence of living miracles. Let me illustrate using situations with which I have been intimately acquainted in years past, as well as in the present period of time.

The proprietor of Dick’s Cafe in St. George, Utah, is such an example. Dick Hammer came to Utah during the Depression years with the Civilian Conservation Corps. During that period, he met and married a Latter-day Saint young woman. He opened his cafe, which became a popular meeting spot. Home teacher to the Hammer family was Willard Milne. Since I knew Dick Hammer and had printed his menus, I would ask my friend Brother Milne when I visited St. George, “How is our friend Dick Hammer coming?”

The reply would generally be, “Slowly.”

The years passed by, and just a year or two ago Willard said to me: “Brother Monson, Dick Hammer is converted and is going to be baptized. He is in his 90th year, and we have been friends all our adult lives. His decision warms my heart. I’ve been his home teacher for many years—perhaps 15 years.”

Brother Hammer was indeed baptized and a year later entered that beautiful St. George Temple and there received his endowment and sealing blessings.

I asked Willard, “Did you ever become discouraged teaching for such a long time?”

He replied, “No, it was worth the effort. I am a happy man.”

Some years ago, before my leaving to become president of the Canadian Mission headquartered in Toronto, Ontario, I had developed a friendship with a man by the name of Shelley who lived in the ward but did not embrace the gospel, irrespective of the fact that his wife and children had done so. As I served as a mission president, had I been asked to name anyone I knew not likely to become a member of the Church, I believe I would have thought of Shelley.

After I was called to the Twelve, I received a telephone call from Shelley. He said, “Bishop, will you seal my wife, my family, and me in the Salt Lake Temple?”

I answered hesitantly, “But, Shelley, you first must be a baptized member of the Church.”

He laughed and responded, “Oh, I took care of that while you were in Canada. My home teacher was a school crossing guard, and every weekday as he and I would visit at the crossing, we would discuss the gospel.”

I had the privilege to see this miracle with my own eyes and feel the joy with my heart and soul. The sealings were performed; a family was united. Shelley died not long after this period, but not before he publicly thanked his home teachers for their faithful service.

Elder Mark E. Petersen, when discussing activation of members, would frequently declare, “The challenge is one of lack of conversion.” We, the priesthood of the Church, cannot afford to leave families in a cocoon, isolated from the body of the Church.

Long years ago, Joseph Lyon of Salt Lake City shared with me the lesson of a lecture which a minister from another faith observed as he spoke to the Associated Credit Men of Salt Lake. The minister boldly proclaimed: “Mormonism is the greatest philosophy in the world today. The biggest test for the Church will come with the advent of television and radio, which tend to keep people away from the Church.” He then proceeded to relate what I’ve called the “hot coals” story. He described a warm fireplace where the pieces of wood had burned brightly, with the embers still glowing and giving off heat. He then observed that by taking in hand brass tongs, he could remove one of the hot embers. That ember would then slowly pale in light and turn black. No longer would it glow. No longer would it warm. He then pointed out that by returning the black, cold ember to the bed of living coals, the dark ember would begin to glow and brighten and warm. He concluded: “People are somewhat like the coals of a fire. Should they absent themselves from the warmth and spirit of the active church membership, they will not contribute to the whole, but in their isolation will be changed. As with the embers removed from the heat of the fire, as they distance themselves from the intensity of the spirit generated by the active membership, they will lose that warmth and spirit.”

The reverend closed his comments by observing, “People are more important than the embers of a fire.”

As years come and then go and life’s challenges become more difficult, the visits of home teachers to those who have absented themselves from Church activity can be the key which will eventually open the doors to their return.

With this thought in mind, can we brethren not reach out to those for whom we are responsible and bring them to the table of the Lord to feast on His word and to enjoy the companionship of His Spirit, and be “no more strangers and foreigners, but fellowcitizens with the saints, and of the household of God”?

President Ezra Taft Benson said that home teaching is “priesthood compassionate service.” Not long ago I received a touching letter from Sister Mori Farmer. It tells of two home teachers and the loving service they provided the Farmer family during a time when the family was experiencing some difficult financial circumstances. At the time the service was provided, the Farmer family was out of town attending a family reunion.

I share with you first a letter written to the Farmer family by their home teachers, which the family found taped to their garage door when they returned home. It begins: “We hope you had a great family reunion. While you were gone, we and about fifty of our friends had a great party at your house. We want to thank you from the bottom of our hearts for the years of unselfish service you both have given to us. You have been Christlike examples of untiring service to others. We can never repay you for that—but just thought we’d like to say thanks. Signed, your home teachers.”

I quote now from Sister Mori Farmer’s letter to me:

“[After reading the note from our home teachers] we entered the house with great anticipation. What we found shocked us so much we were at a loss for words. I stayed up all night crying over the generosity of the people in our ward.

“Our home teachers had decided that they would fix our carpet while we were away. They had moved the furniture out into the front yard so the carpet could get stretched and finished. One man in the ward stopped and asked what was going on. He returned later with several hundred dollars’ worth of paint and said, ‘We might as well paint the house while everything is out.’ Others saw the cars out front and stopped to see what was going on, and by week’s end 50 people were busy repairing, painting, cleaning, and sewing.

“Our friends and fellow ward members had fixed our poorly laid carpet, painted the entire house, repaired holes in the drywall, oiled and varnished our kitchen cabinets, put curtains on three windows in the kitchen and family room, done all the laundry, cleaned every room in the house, had the carpets cleaned, fixed broken door latches, and on and on. In trying to make a list of all the wonderful things they did for us, we filled three pages. All of this had been accomplished between Wednesday and our return on Sunday.

“Almost everyone we talked to told us, with tears in their eyes, what a spiritual experience it had been to participate. We have been truly humbled by this experience. As we look around our home, we are reminded of their kindness and of the great sacrifice of time, talents, and money they made for our family. Our home teachers have truly been angels in our lives, and we will never forget them and the wonderful things they have done for us.”

Other instances could well be cited. However, I turn to one example to describe the type of home teachers we should be. There is one Teacher whose life overshadows all others. He taught of life and death, of duty and destiny. He lived not to be served, but to serve; not to receive, but to give; not to save His life, but to sacrifice it for others. He described a love more beautiful than lust, a poverty richer than treasure. It was said of this Teacher that He taught with authority and not as did the scribes. In today’s world, when many men are greedy for gold and for glory and are dominated by the philosophies of men, remember that this Teacher never wrote—once only He wrote on the sand, and wind destroyed forever His handwriting. His laws were not inscribed upon stone, but upon human hearts. I speak of the Master Teacher, even Jesus Christ, the Son of God, the Savior and Redeemer of all mankind. The biblical account says of Him, He “went about doing good.” With Him as our unfailing Guide and Exemplar, we shall qualify for His divine help in our home teaching. Lives will be blessed. Hearts will be comforted. Souls will be saved.

In the name of Jesus Christ, amen.

\newpage
\settalk{Teach the Children}
\setspeaker{Thomas S. Monson}
\settalkdate{October 1997}
\talkheading{Teach the Children}{Thomas S. Monson}

In Salt Lake City, a touch of autumn is in the air. Daylight hours grow fewer and the weather turns cooler, reminding one and all that winter is just around the corner. The Christmas season will soon be upon us.

Inevitably, the spirit of Christmas inspires kind deeds, touches human hearts, and prompts one’s mind to reach back to that humble stable in faraway Bethlehem, to a time when the prophecies of the prophets, both in that area and here on the American continent, became a living reality. Christ the Lord was born.

Precious little is written concerning the childhood of Jesus. One might suppose that His birth was so revolutionary in its magnitude as to dominate accounts of His boyhood. We marvel at the mature wisdom of the boy who, leaving Joseph and Mary, was found in the temple, “sitting in the midst of the doctors,” teaching them the gospel. When Mary and Joseph expressed their concern about His absence, He asked of them the penetrating question: “Wist ye not that I must be about my Father’s business?”

The sacred record declares of Him: “Jesus increased in wisdom and stature, and in favour with God and man.” An obscure passage describes the transition from child to man: He “went about doing good.”

Because of Jesus Christ the world has changed—the divine Atonement has been made, the price of sin has been paid, and the fearful spectacle of death yields to the light of truth and the assurance of resurrection. Though the years roll by, His birth, His ministry, His legacy continue to guide the destiny of all who follow Him as He so invitingly urged.

Children are born each day—even each hour—to mothers who have, with their hand in God’s hand, entered the valley of the shadow of death, that they might bring forth a son, a daughter, to grace a family, a home, and in a way a portion of the earth.

Those precious days of infancy bond mother and father to son or daughter. Every smile is noted, every fear comforted, every hunger abated. Step by step the child grows. The poet wrote that each child is “a sweet new blossom of Humanity, / Fresh fallen from God’s own home to flower on earth.”

The child grows in wisdom and also in stature. Learning and doing become priorities to be addressed.

There are those who dismiss these responsibilities, feeling they can be deferred until the child grows up. Not so, the evidence reveals. Prime time for teaching is fleeting. Opportunities are perishable. The parent who procrastinates the pursuit of his responsibility as a teacher may, in years to come, gain bitter insight into Whittier’s expression: “Of all sad words of tongue or pen, / The saddest are these: ‘It might have been!’ ”

Dr. Glenn Doman, a prominent author and renowned scientist, reported a lifetime of research in the statement: “The newborn child is almost an exact duplicate of an empty … computer, although superior to such a computer in almost every way. … What is placed in the child’s brain during the first eight years of his life is probably there to stay." “If you put misinformation into his brain during [this period], it is extremely difficult to erase it.”

This evidence should provoke a renewal of commitment in every parent: “I must be about my Father’s business.” Children learn through gentle direction and persuasive teaching. They search for models to imitate, knowledge to acquire, things to do, and teachers to please.

Parents and grandparents fill the role of teacher. So do siblings of the growing child. In this regard, I offer four simple suggestions for your consideration:

First, teach prayer. “Prayer is the simplest form of speech / That infant lips can try; / Prayer, the sublimest strains that reach / The Majesty on high.”

We learn to pray by praying. One can devote countless hours to examining the experiences of others, but nothing penetrates the human heart as does a personal, fervent prayer and its heaven-sent response.

Such was the example of the boy Samuel. Such was the experience of young Nephi. Such was the far-reaching prayer of the youth Joseph Smith. Such can be the blessing of one who prays. Teach prayer.

Next, inspire faith. This sesquicentennial year of the epic pioneer trek to the valley of the Great Salt Lake has inspired more music, more drama, more involvement by youth and adults than perhaps any other occasion in our history. We as families have learned more of Church history—the glory and the suffering, the hardship and sorrow, then victory upon arrival in the valley—than can be estimated. Some years ago, Bryant S. Hinckley, the father of our President, prepared a book entitled The Faith of Our Pioneer Fathers. Accounts which the volume contains are so well written and set forth. This past year they were retold by the score. Countless members looked back on their own pioneer heritage. Hundreds of youth—even thousands throughout the world—pulled and pushed handcarts and walked their own pioneer trail.

I think that there isn’t a member of this Church today who has not been touched by the year now drawing to its close. Those who did so much for the good of all surely had as their objective to inspire faith. They met the goal in a magnificent manner.

Third, live truth. At times the most effective lesson in living truth is found close to the home and dear to the heart.

At the funeral service of a noble General Authority, H. Verlan Andersen, a tribute was expressed by a son. It has application wherever we are and whatever we are doing. It is the example of personal experience.

The son of Elder Andersen related that years earlier, he had a special school date on a Saturday night. He borrowed from his father the family car. As he obtained the car keys and headed for the door, his father said, “The car will need more gas before tomorrow. Be sure to fill the tank before coming home.”

Elder Andersen’s son then related that the evening activity was wonderful. Friends met, refreshments were served, and all had a good time. In his exuberance, however, he failed to follow his father’s instruction and add fuel to the car’s tank before returning home.

Sunday morning dawned. Elder Andersen discovered the gas gauge showed empty. The son saw his father put the car keys on the table. In the Andersen family the Sabbath day was a day for worship and thanksgiving, and not for purchases.

As the funeral message continued, Elder Andersen’s son declared, “I saw my father put on his coat, bid us good-bye, and walk the long distance to the chapel, that he might attend an early meeting.” Duty called. Truth was not held slave to expedience.

In concluding his funeral message, Elder Andersen’s son said, “No son ever was taught more effectively by his father than I was on that occasion. My father not only knew the truth, but he also lived it.” Live truth.

Finally, honor God. No one can surpass the Lord Jesus Christ in setting an example of living this goal. The fervency of His prayer at Gethsemane says it all: “Father, if thou be willing, remove this cup from me: nevertheless not my will, but thine, be done.” His example on the cruel cross of Golgotha speaks volumes: “Father, forgive them; for they know not what they do.”

The Master taught so everlastingly to all who would listen a simple yet profound truth as recorded in Matthew. We learn that after Jesus and His disciples descended from the Mount of Transfiguration, they paused at Galilee and then went to Capernaum. The disciples said unto Jesus, “Who is the greatest in the kingdom of heaven?

“And Jesus called a little child unto him, and set him in the midst of them,

“And said, Verily I say unto you, Except ye be converted, and become as little children, ye shall not enter into the kingdom of heaven.

“Whosoever therefore shall humble himself as this little child, the same is greatest in the kingdom of heaven.

“And whoso shall receive one such little child in my name receiveth me.”

I think it significant that Jesus so loved these little ones who recently had left the preexistence to come to earth. Children then and children now bless our lives, kindle our love, and prompt good deeds.

Is it any wonder that the poet Wordsworth speaks thus of our birth: “Trailing clouds of glory do we come / From God, who is our home: / Heaven lies about us in our infancy!”

It is in the home that we form our attitudes, our deeply held beliefs. It is in the home that hope is fostered or destroyed. Wrote Dr. Stuart E. Rosenberg in his book The Road to Confidence: “Despite all new inventions and modern designs, fads and fetishes, no one has yet invented, or will ever invent, a satisfying substitute for one’s own family.”

We ourselves can learn from our children and grandchildren. They have no fear. They have no doubt concerning our Heavenly Father’s love for them. They love Jesus and want to be like Him.

Our grandson, six-year-old Jeffrey Monson Dibb, accompanied by his six-year-old girlfriend, paused at an end table in his house on which there was a picture of Elder Jeffrey R. Holland. The young girl pointed to the picture and asked, “Who is that man?”

Jeff replied, “Oh, that’s Elder Jeffrey Holland of the Quorum of Twelve Apostles. He’s named after me!”

This same namesake of Elder Holland’s, along with his girlfriend, went for a walk one day. They marched up the front steps of a home, not knowing who lived there or what affiliation they might have with the Church. They knocked on the front door, and a woman answered. Without the slightest hesitation, Jeff Dibb said to her, “We are the visiting home teachers. May we come in?” They were ushered into the living room and were asked to be seated. With total faith the children addressed the woman, “Do you have a treat for us?” What could she do? She produced a treat, and they had a nice conversation. The impromptu teachers departed, uttering a sincere “Thank you.”

“Come back again,” they heard the woman say, with a smile on her face.

“We will,” came the reply.

The parents of the two youngsters heard of the incident. I am certain they were restrained in counseling the little ones. Perhaps they remembered the words from the scriptures: “And a little child shall lead them.”

The sound of laughing children, joyfully playing together, can give the impression that childhood is free from trouble and sorrow. Not so. Children’s hearts are tender. They long for the companionship of other children. In the famous Victoria and Albert Museum in London hangs a masterpiece on canvas. Its title is simply Sickness and Health. Depicted is a small girl in a wheelchair. Her face is pale; her countenance reflects sadness. She watches an organ-grinder perform while two little girls, carefree and happy, frolic and dance.

Sadness and sorrow at times come to all, including children. But children are resilient. They bear up beautifully to shoulder the burden they may be called upon to endure. Perhaps the lovely psalm describes this virtue: “Weeping may endure for a night, but joy cometh in the morning.”

May I now paint a picture of such a situation. In faraway Bucharest, Romania, Dr. Lynn Oborn, volunteering at an orphanage, was attempting to teach little Raymond, who had never walked, how to use his legs. Raymond had been born with severe clubfeet and was completely blind. Recent orthopedic surgery performed by Dr. Oborn had corrected the clubfeet, but Raymond was still unable to use his legs. Dr. Oborn knew that a child-size walker would enable Raymond to get on his feet, but such a walker was not available anywhere in Romania. I’m sure fervent prayers were offered by this doctor who had done all he could without a walking aid for the boy. Blindness can hamper a child, but inability to walk, to run, to play can injure his precious spirit.

Let us turn now to Provo, Utah. The Richard Headlee family, learning of the suffering and pitiful conditions in Romania, joined with others to assemble a 40-foot container filled with 40,000 pounds of needed supplies, including food, clothing, medicine, blankets, and toys. The project deadline arrived, and the container had to be shipped that day.

No one involved with the project knew of the particular need for a child-size walker. However, at the last possible moment, a family brought forth a child’s walker and placed it in the container.

When the anxiously awaited container arrived at the orphanage in Bucharest, Dr. Oborn was present as it was opened. Every item it contained would be put to immediate use at the orphanage. As the Headlee family introduced themselves to Dr. Oborn, he said, “Oh, I hope you brought me a child’s walker for Raymond!”

One of the Headlee family members responded, “I can vaguely remember something like a walker, but I don’t know its size.” Another family member was dispatched back into the container, crawling among all the bales of clothes and boxes of food, searching for the walker. When he found it, he lifted it up and cried out, “It’s a little one!” Cheers erupted—which quickly turned to tears, for they all knew they had been part of a modern-day miracle.

There may be some who say, “We don’t have miracles today.” But the doctor whose prayers were answered would respond, “Oh, yes we do, and Raymond is walking!” She who was inspired to give the walker was a willing vessel and surely would agree.

Who was the angel of mercy touched by the Lord to play a vital role in this human drama? Her name is Kristin. She is the daughter of Kurt and Melodie Bestor. Kristin was born with spina bifida, as was her younger sister, Erika. The two children have spent long days and worrisome nights in the hospital. Modern medicine, lovingly practiced, along with help from our Heavenly Father have brought a measure of mobility to each. Neither is downhearted. Both inspire others to carry on. Last month Kristin and Erika entertained guests celebrating the 75th anniversary of Primary Children’s Medical Center. They sang with their father and mother, and then the girls movingly sang a duet. Each person in the audience had red-rimmed eyes; handkerchiefs were everywhere displayed. These girls, this family, had overcome sorrow and brought joy to the lives of others.

Kristin’s father said to me that evening, “President Monson, meet Kristin. She is the one who felt impressed to send her walker to Romania, hoping that some child there would be benefited.”

I spoke to Kristin as she sat in her wheelchair. “Thank you for listening to the Spirit of the Lord. You have been the instrument in the Lord’s hands to answer a doctor’s prayer, a child’s wish.”

Later, as I walked out of that celebration held for the benefit of children, I looked upward toward the heavens and offered my own “Thank you” to God for children, for families, for miracles in our time.

Let us earnestly follow His direction: “Suffer the little children to come unto me, and forbid them not: for of such is the kingdom of God.”

A popular song includes the words, “There are angels among us.” These angels are quite frequently the precious little ones God, our Father, has entrusted to our earthly care. To them may we teach prayer, inspire faith, live truth, and honor God. Then we shall have heavenly homes and forever families. For what higher gift could we wish? For what greater blessing could we pray? None!

In the name of Jesus Christ, amen.

\newpage
\settalk{The Mighty Strength of the Relief Society}
\setspeaker{Thomas S. Monson}
\settalkdate{October 1997}
\talkheading{The Mighty Strength of the Relief Society}{Thomas S. Monson}

Sisters, you are assembled tonight in one of the largest gatherings of Relief Society sisters ever held. Your conference has been uplifting and inspirational.

This evening marks the first general meeting conducted by your new presidency: President Mary Ellen Smoot and her counselors, Virginia Jensen and Sheri Dew. Predecessor presidencies have also served with distinction. We are honored tonight by their presence and their service.

A thought has gone through my mind as I’ve prepared for this opportunity. I’ve expressed it this way: Remember the past; learn from it. Contemplate the future; prepare for it. Live in the present; serve in it. Therein is the mighty strength of the Relief Society of this Church.

From the early days of the Restoration, the prophets of God have stressed the importance of your organization. President Brigham Young counseled: “Now, Bishops, you have smart women for wives. … Let them organize Female Relief Societies in the various wards. We have many talented women among us, and we wish their help in this matter. Some may think this is a trifling thing, but it is not; and you will find that the sisters will be the mainspring of the movement.”

President Lorenzo Snow taught that the Relief Society exemplifies pure religion. “The Apostle James said that ‘pure religion and undefiled before God … is this, to visit the fatherless and widows in their affliction, and to keep himself unspotted from the world.’[] … The members of the Relief Society have most surely exemplified in their lives pure and undefiled religion; for they have ministered to those in affliction, they have thrown their arms of love around the fatherless and the widows, and they have kept themselves unspotted from the world. I can testify that there are no purer and more God-fearing women in the world than are to be found within the ranks of the Relief Society.”

I can attest to the truth of President Snow’s statement. Relief Society has always been comprised of those who put others first and self last. I remember that when I was a small boy during the Depression, my mother was the secretary-treasurer of the ward Relief Society. Back then dues were paid to assist those in need. Mother was not really a bookkeeper, so Dad would help her. The individual contributions were never so much as a dollar, but rather would be a quarter, a dime, a nickel, a few pennies.

I learned many lessons from my mother. I must have been a very active boy, for Mother was always telling me, “Slow down, Tommy, slow down. You’re on the verge of Saint Vitus’ dance!” You know, I never did know what Saint Vitus’ dance was. All I knew was that Mother said I was on the verge of it—and the way she spoke the words, I assumed it was a drastic ailment.

Since we lived just a block or two from the railroad tracks, frequently men, unemployed, without funds for food, would leave the train and come to our house for something to eat. Such men were always polite. They offered to do some work for the food. Indelibly imprinted on my mind is the picture of a gaunt and hungry man standing at our kitchen door, hat in hand, pleading for food. Mother would welcome such a visitor and would direct him to the kitchen sink to wash up while she prepared food for him to eat. She never skimped on quality or quantity; the visitor ate exactly the same lunch as did my father. As he wolfed down the food, Mother took the opportunity to counsel him to return to his home and his family. When he left the table, he had been nourished physically and spiritually. These men never failed to say thank you. Tears in their eyes revealed ever so silently the gratitude of their hearts.

But what of today? Are there not hungry souls to feed? Are there not greetings to share? Are there not visits to be made? As I contemplate the Relief Society of today, humbled by my privilege to speak to you, I turn to our Heavenly Father for His divine guidance.

In this spirit, I have felt to provide each member of the Relief Society throughout the world three goals to meet:

Let us consider these three goals. First, gain knowledge through study. In a vital revelation which has universal impact, the Lord declared: “And as all have not faith, seek ye diligently and teach one another words of wisdom; yea, seek ye out of the best books words of wisdom; seek learning, even by study and also by faith.”

Elder Adam S. Bennion, a member of the Council of the Twelve several decades ago and a well-educated teacher and leader, urged: “God help us to appreciate the value of true education. If we were in this building and had only one window, we would see only one corner of the universe. A man who has not been trained looks out on life through the little window of narrow experience. It is the business of education to fill the building of life with windows so that we can look out on the universe at every angle.

“Having left this world, we shall enter heaven, the bigger school, and I hope we shall all be passed by into the bigger field with a development, which shall be inscribed, ‘This man or this woman, in the school of life, achieved all the opportunities that were his [or hers].’”

An example of a narrow window of vision being replaced by vision unlimited took place at the Monroe, Louisiana, airport several years ago. I was on my way home from a regional meeting and met a lovely African-American sister who approached me and said joyfully, “President Monson, before I joined the Church and became a member of the Relief Society, I could not read. I could not write. None of my family could. You see, we were all poor sharecroppers. President, my white Relief Society sisters—they taught me to read. They taught me to write. Now I help teach other white sisters how to read and how to write.” I reflected on the supreme happiness she must have felt when she opened her Bible and read for the first time the words of the Lord: “Come unto me, all ye that labour and are heavy laden, and I will give you rest. Take my yoke upon you, and learn of me; for I am meek and lowly in heart: and ye shall find rest unto your souls. For my yoke is easy, and my burden is light.”

That day in Monroe, Louisiana, I received a confirmation by the Spirit of your exalted objective of improving literacy among your sisters.

Each of you, single or married, regardless of age, has the opportunity to learn and to grow. Expand your knowledge, both intellectual and spiritual, to the full stature of your divine potential.

The Holy Ghost will be your constant guide when difficult decisions are to be made. “For they that are wise … have taken the Holy Spirit for their guide.”

Be true to your ideals, for “ideals are like stars; you will not succeed in touching them with your hands. But … following them you will reach your destiny.” Spiritual help is but a prayer away.

During the next two years, as Sister Smoot has declared, members of the Relief Society and holders of the Melchizedek Priesthood will each be studying the teachings of the prophet Brigham Young. The manual has been painstakingly prepared; it is beautifully printed and bound, with highly relevant discussion items featured. The lesson material will be taught during the Relief Society period on two Sundays of the month and likewise to the brethren of the Melchizedek Priesthood for two Sundays. On the remaining Sundays, the conventional matters of Relief Society and priesthood quorum work will go forward.

Years ago I saw a photograph of a Sunday School class in the Sixth Ward of the Pioneer Stake in Salt Lake City, the ward where my family lived. The photograph was taken in 1905. A sweet girl, her hair in pigtails, was shown on the front row. Her name was Belle Smith. Later, as Belle Smith Spafford, general president of the Relief Society, she wrote: “Never have women had greater influence than in today’s world. Never have the doors of opportunity opened wider for them. This is an inviting, exciting, challenging, and demanding period of time for women. It is a time rich in rewards if we keep our balance, learn the true values of life, and wisely determine priorities.”

My dear sisters, this is your day, this is your time. The holy scriptures adorn our bookshelves. Make certain they provide nourishment to our minds and guidance for our lives. Our goal, gain knowledge through study.

Second, make home a heaven.

Speaking to a general session of conference in 1945, a day or so after Relief Society conference, President George Albert Smith said: “Yesterday this house was filled with the daughters of Zion, and I say without hesitation that you could find no more beautiful picture of womankind in all the world than was here yesterday afternoon [in the conference of the women’s Relief Societies of the Church]. These faithful wives, these faithful daughters, assume their portion of the burden and carry it on. They make their homes a heaven.”

My dear sisters, home—that marvelous place—was meant to be a haven called heaven where the Spirit of the Lord might dwell.

Too frequently, women underestimate their influence for good. Well could you follow the formula given by the Lord: “Establish a house, even a house of prayer, a house of fasting, a house of faith, a house of learning, a house of glory, a house of order, a house of God.”

In such a house will be found happy, smiling children who have been taught, by precept and example, the truth. In a Latter-day Saint home children are not simply tolerated, but welcomed; not commanded, but encouraged; not driven, but guided; not neglected, but loved.

Years pass, and children become more independent. They move away from mother’s protective care, but they are ever influenced by mother’s teachings, mother’s example, and mother’s love. Some appear by their actions to have forgotten this influence. However far from the hearth of home the wanderer travels, the word mother mentally and emotionally brings him or her homeward once again. And mother, as always, stands with welcoming arms.

President Stephen L Richards declared: “The various organizations of the Church … however much of good they may accomplish, can in no sense take the place of the home. They cannot be proxy for parents. …

“I believe in the home as the foundation of society, as the cornerstone of the nation, and as the primary institution of the Church. I cannot conceive of a great people without great, good homes. I believe that the first calling of man and woman is to form a good home.”

Speaking in a Relief Society conference in 1953, Sister Spafford declared, “Mothers, you must feel your testimony before you can influence or give of that testimony to your children.”

There are many women in the Relief Society who are not married. Death, divorce, and indeed lack of opportunity to marry have in many instances made it necessary for a woman to stand alone. In reality, she need not stand alone, for a loving Heavenly Father will be by her side to give direction to her life and provide peace and assurance in those quiet moments where loneliness is found and where compassion is needed.

President Joseph Fielding Smith, in speaking to the single sisters who have never had the opportunity to marry, promised: “If in your hearts you feel that the gospel is true, and would under proper conditions receive these ordinances and sealing blessings in the temple of the Lord; and that is your faith and your hope and your desire, and that does not come to you now; the Lord will make it up, and you shall be blessed—for no blessing shall be withheld.”

Let us make home a heaven.

Goal number three, find joy in service.

The Prophet Joseph Smith recorded that on March 24, 1842, he responded to an invitation to attend Relief Society, “whose object is the relief of the poor, the destitute, the widow and the orphan, and for the exercise of all benevolent purposes. … [The Relief Society sisters] will fly to the relief of the stranger; … they will dry up the tears of the orphan and make the widow’s heart to rejoice.”

At times the call to service extended to a member of the Relief Society is a bit unusual. Such an assignment I share with you in closing.

When I was bishop of the Sixth-Seventh Ward in Salt Lake City, back when we had a Relief Society Magazine, I noted that our record for subscriptions to that publication was low. Prayerfully my counselors and I analyzed the names of the individuals whom we could call to be magazine representative, and the inspiration dictated that Elizabeth Keachie should be given the assignment. She responded affirmatively to the call. She and her sister-in-law Helen Ivory, also a member of the ward, commenced to canvass the entire ward, house by house, street by street, and block by block. The result was phenomenal. We had more subscriptions to the Relief Society Magazine than had been recorded by all the other units of our stake combined.

I congratulated Elizabeth Keachie one Sunday evening and said to her, “Your task is done.”

She replied, “Not yet, Bishop. There are two blocks we have not yet covered.”

When she told me which blocks they were, I said, “Sister Keachie, no one lives on those blocks. That area is all industrial.”

“Just the same,” she said, “I’ll feel better if I can go and check them myself.”

Sister Keachie and Sister Ivory, on a rainy day, covered those final two blocks but discovered no homes. As they were about to discontinue their search, they noted a driveway which was strewn with mud puddles from a recent storm. It was next to a foundry. Sister Keachie gazed down the driveway perhaps 60 feet and could just make out a garage with a curtain at the window.

Deciding to investigate, the two sweet sisters walked through the mud to a point where the entire garage could be seen. Now they noticed a door, not visible from the street, which had been cut into the side of the garage. They noticed a chimney with smoke rising from it.

They knocked at the door. A man of about 65 years of age, William Ringwood, answered. They presented their story concerning the need of every home having the Relief Society Magazine. William Ringwood replied, “You’d better ask my father.” Ninety-three-year-old Charles W. Ringwood then came to the door and also listened to the message. He subscribed.

Elizabeth Keachie reported to me the presence of these two men in our ward. When I requested their membership certificates from the Membership Department at the Presiding Bishopric’s Office, I was told that the certificates had remained in the lost file of the Presiding Bishopric’s Office for many years.

On Sunday morning Elizabeth Keachie brought to our priesthood meeting Charles and William Ringwood—the first time they had been inside a chapel for a long while. Charles Ringwood, 93, was the oldest deacon I had ever met, and his son was the oldest male member holding no priesthood I had ever met.

The elder Brother Ringwood was ordained a priest and then an elder. I shall never forget his interview with respect to seeking a temple recommend. He handed me a silver dollar which he took from an old worn leather coin purse and said, “This is my fast offering.”

I replied, “Oh, Brother Ringwood, you owe no fast offering. You need it yourself.”

“I want to receive the blessings, not keep the money,” he responded.

It was my opportunity to take Charles Ringwood to the Salt Lake Temple and to attend with him the endowment session. That same evening Elizabeth Keachie served as proxy for the deceased Sister Ringwood.

At the conclusion of the ceremony, Charles Ringwood said to me, “Bishop, I told my wife just before she died 16 years ago that I would not delay in getting this work done. I am happy this has been accomplished.”

Within two months, Charles W. Ringwood passed away. At his funeral service, I noticed his family sitting on the front row of the mortuary chapel, but I also noticed two sweet ladies sitting near the rear—Elizabeth Keachie and Helen Ivory. As I gazed upon those two sweet women, I thought of the 76th section of the Doctrine and Covenants: “I, the Lord, am merciful and gracious unto those who fear me, and delight to honor those who serve me in righteousness and in truth unto the end. Great shall be their reward and eternal shall be their glory.” I testify that we can find joy in service.

Sisters, may we gain knowledge through study. May we make home a heaven. May we find joy in service. By so doing, we shall experience the fulfillment of the Lord’s promise: “I, the Lord, am well pleased.” In the name of Jesus Christ, amen.

\newpage
\settalk{In Harm’s Way}
\setspeaker{Thomas S. Monson}
\settalkdate{April 1998}
\talkheading{In Harm’s Way}{Thomas S. Monson}

On July 16, 1945, the USS Indianapolis departed the Mare Island Naval Shipyard in California on a secret cargo mission to Tinian Island in the Marianas. The cargo included highly sophisticated equipment which could well bring an end to the Second World War, with all its suffering, remorse, and death. The ship delivered its cargo on July 26 and was heading, unescorted, toward Leyte in the Philippines.

Because they were traveling through hostile waters in the Philippine Sea, the captain had discretionary orders to follow a zigzag course of travel to prevent detection by and attack from the enemy. He failed to do so. Just before midnight on Sunday, July 29, 1945, as the Indianapolis continued toward Leyte Gulf, the heavy cruiser was discovered by an enemy submarine. Easily avoiding detection while submerging to periscope depth, the submarine fired a fanwise salvo of six torpedoes from 1,500 yards. As the torpedoes struck the target, explosions of ammunition and aviation fuel ripped away the cruiser’s bow and destroyed its power center. Without power, the radio officer was unable to send a distress signal. The order to abandon ship, when it came, had to be passed by word of mouth because all communications were down. Just 12 minutes after being hit, the stern rose up a hundred feet straight into the air, and the ship plunged into the depths of the sea.

Of the nearly 1,200-man crew, approximately 400 were killed instantly or went down with the ship. About 800 survived the sinking and went into the water.

Four days later, on August 2, 1945, the pilot of a Lockheed Ventura, flying on patrol, noticed an unusual oil slick on the water’s surface and followed it for 15 miles. Then the plane’s occupants spotted those men who had managed to survive since the Indianapolis had gone down.

A major rescue effort began. Ships hurried to the area, and planes were dispatched to drop food, water, and survival gear to the men. Of the approximately 800 who had gone into the water, only 316 remained alive. The rest had been claimed by the perilous, shark-infested sea.

Two weeks later World War II was over. The sinking of the Indianapolis, called “the final great naval tragedy of World War II,” is now legend.

Are there lessons for our lives in the horrific experience of those men aboard the Indianapolis? They were in harm’s way. Danger lurked; the enemy stalked. The vessel sailed on, disregarding the command to zigzag, and thus it became an easy target. Catastrophe was the result.

On the day the Indianapolis sailed toward Leyte, I enlisted in the United States Navy. At the Naval Training Station near San Diego, California, I endured the extreme discipline of boot camp and the intense training for combat.

At last our first liberty came, and we were advised that all those who could swim could now take the navy bus to San Diego, while those sailors who could not were to remain for swimming training. How pleased I was that I could swim and had done so for many years. Then came an unexpected order. We who answered that we could swim were marched away—not to the waiting bus, but rather to the base swimming pool. We assembled at the pool’s deep end, were told to undress, and then were commanded to jump in one at a time and swim the length of the pool. Most accomplished the feat with little effort and anticipated eagerly the bus ride to San Diego. But there were men who had been untruthful, who had answered they could swim when in reality they could not. For them, the petty officers waited until they were about to go under the water for the second or third time before proffering a bamboo pole to tow them to safety. The lesson learned? Tell the truth. It could ultimately save your life if you were in harm’s way.

Our journey through mortality will at times place us in harm’s way. Is there a road map to safety? Are there those to whom we can look for help?

May I offer to you tonight six road signs which, when observed and followed, will guide you to safety. They are:

There is a battle of significant consequence taking place in the lives of young men today. In simple terms, it is the struggle between doing right or doing wrong.

At an earlier time, Moroni offered the advice: “For behold, the Spirit of Christ is given to every man, that he may know good from evil; wherefore, I show unto you the way to judge; for every thing which inviteth to do good, and to persuade to believe in Christ, is sent forth by the power and gift of Christ; wherefore ye may know with a perfect knowledge it is of God.

“But whatsoever thing persuadeth men to do evil, and believe not in Christ, and deny him, and serve not God, then ye may know with a perfect knowledge it is of the devil.”

May I share a thought or two concerning each of the six road signs previously mentioned to keep you from harm’s way.

Brethren, are we prepared for the voyage of life? The sea of life can at times become turbulent. Crashing waves of emotional conflict may break all around us. Chart your course, be cautious, and follow the safety measures outlined.

In so doing, we will sail safely the seas of life and arrive at home port—even the celestial kingdom of God. Then, as mariners of mortality, may we hear the plaudit, “Well done, thou good and faithful servant: … enter thou into the joy of thy lord.”

For this blessing I fervently pray, in the name of Jesus Christ, amen.

\newpage
\settalk{Look to God and Live}
\setspeaker{Thomas S. Monson}
\settalkdate{April 1998}
\talkheading{Look to God and Live}{Thomas S. Monson}

I commence my message this morning with a question: Have you ever taken a vacation with your entire family? If not, you are in for some surprises when you do. My wife and I a few years ago joined our children, their companions, and the grandchildren at Disneyland in southern California. Beyond the entrance to the famous theme park, the group rushed to what was then the newest feature—Star Tours. You enter a simulated rocket, take your seat, and fasten your seat belt. All of a sudden the entire vehicle begins to vibrate violently. I think the mechanical voice which comes over the loudspeaker calls it “heavy turbulence.” (I have never returned to this featured ride. I get all the real turbulence I can handle just flying from place to place fulfilling my responsibilities.)

After recuperating for a few minutes, we journeyed to the feature with the longest line. It is called Splash Mountain. The crowd filed round and round in a serpentine pattern. The music, which was piped through the loudspeakers to the waiting throng, contained the words of the song:

\begin{verse}
Zip-a-dee-doo-dah, zip-a-dee-ay,\\
My, oh my, what a wonderful day!\\
Plenty of sunshine, headin’ my way,\\
Zip-a-dee-doo-dah, zip-a-dee-ay!
\end{verse}

By now we were ready to board the boat which would carry us in a vertical dive that evoked screams from the passengers in the boat ahead as it roared down the waterfall and glided to a stop in the water below. Just before taking the plunge, however, I noticed on one wall a small sign declaring a profound truth: “You can’t run away from trouble; there’s no place that far!”

These few words have remained with me. They pertain not only to the theme of Splash Mountain but also to our sojourn in mortality.

Life is a school of experience, a time of probation. We learn as we bear our afflictions and live through our heartaches.

As we ponder the events that can befall all of us—even sickness, accident, death, and a host of other challenges—we can say with Job of old: “Man is born unto trouble.” Job was a “perfect and upright” man who “feared God, and eschewed evil.” Pious in his conduct, prosperous in his fortune, Job was to face a test which could have destroyed anyone. Shorn of his possessions, scorned by his friends, afflicted by his suffering, shattered by the loss of his family, he was urged to “curse God, and die.” He resisted this temptation and declared from the depths of his noble soul: “Behold, my witness is in heaven, and my record is on high.”

“I know that my redeemer liveth.” Job kept the faith.

It may safely be assumed that no person has ever lived entirely free of suffering and tribulation, nor has there ever been a period in human history that did not have its full share of turmoil, ruin, and misery.

When the pathway of life takes a cruel turn, there is the temptation to ask the question “Why me?” Self-incrimination is a common practice, even when we may have had no control over our difficulty. At times there appears to be no light at the tunnel’s end, no dawn to break the night’s darkness. We feel surrounded by the pain of broken hearts, the disappointment of shattered dreams, and the despair of vanished hopes. We join in uttering the biblical plea, “Is there no balm in Gilead … ?” We feel abandoned, heartbroken, alone.

To all who so despair, may I offer the assurance found in the psalm, “Weeping may endure for a night, but joy cometh in the morning.”

Whenever we are inclined to feel burdened down with the blows of life, let us remember that others have passed the same way, have endured, and then have overcome.

There seems to be an unending supply of trouble for one and all. Our problem is that we often expect instantaneous solutions, forgetting that frequently the heavenly virtue of patience is required.

Do any of the following challenges sound familiar to you?

The list is endless. In the world of today there is at times a tendency to feel detached—even isolated—from the Giver of every good gift. We worry that we walk alone. You ask, “How can we cope?” What brings to us ultimate comfort is the gospel.

From the bed of pain, from the pillow wet with tears, we are lifted heavenward by that divine assurance and precious promise, “I will not fail thee, nor forsake thee.”

Such comfort is priceless as we journey along the pathway of mortality, with its many forks and turnings. Rarely is the assurance communicated by a flashing sign or a loud voice. Rather, the language of the Spirit is gentle, quiet, uplifting to the heart, and soothing to the soul.

Lest we question the Lord concerning our troubles, let us remember that the wisdom of God may appear as foolishness to men; but the greatest single lesson we can learn in mortality is that when God speaks and a man obeys, that man will always be right.

The experience of Elijah the Tishbite is illustrative of this truth. In the midst of a terrible famine, drought, and the despair of hunger, suffering, and perhaps even death, “the word of the Lord came unto him, saying, Arise, get thee to Zarephath, … and dwell there: behold, I have commanded a widow woman there to sustain thee.”

Elijah didn’t question the Lord. “He arose and went to Zarephath. And when he came to the gate of the city, behold, the widow woman was there gathering of sticks: and he called to her, and said, Fetch me, I pray thee, a little water in a vessel, that I may drink.

“And as she was going to fetch it, he called to her, and said, Bring me, I pray thee, a morsel of bread in thine hand.

“And she said, As the Lord thy God liveth, I have not a cake, but an handful of meal in a barrel, and a little oil in a cruse: and, behold, I am gathering two sticks, that I may go in and dress it for me and my son, that we may eat it, and die.

“And Elijah said unto her, Fear not; go and do as thou hast said: but make me thereof a little cake first, and bring it unto me, and after make for thee and for thy son.

“For thus saith the Lord God of Israel, The barrel of meal shall not waste, neither shall the cruse of oil fail, until the day that the Lord sendeth rain upon the earth.”

She did not question the improbable promise. “She went and did according to the saying of Elijah: and she, and he, and her house, did eat many days.

“And the barrel of meal wasted not, neither did the cruse of oil fail, according to the word of the Lord, which he spake by Elijah.”

Let us now fast-forward the pages of history to that special night when shepherds were abiding with their flocks and heard the holy pronouncement: “Fear not: for, behold, I bring you good tidings of great joy, which shall be to all people.

“For unto you is born this day in the city of David a Saviour, which is Christ the Lord.”

With the birth of the babe in Bethlehem, there emerged a great endowment—a power stronger than weapons, a wealth more lasting than the coins of Caesar. The long-foretold promise was fulfilled; the Christ child was born.

The sacred record reveals that the boy Jesus “increased in wisdom and stature, and in favour with God and man.” At a later time, a quiet entry records that He “went about doing good.”

Out of Nazareth and down through the generations of time come His excellent example, His welcome words, His divine deeds. They inspire patience to endure affliction, strength to bear grief, courage to face death, and confidence to meet life. In this world of chaos, of trial, of uncertainty, never has our need for such divine guidance been more desperate.

Lessons from Nazareth, Capernaum, Jerusalem, and Galilee transcend the barriers of distance, the passage of time, the limits of understanding as they bring to troubled hearts a light and a way. Ahead lay Gethsemane’s garden and Golgotha’s hill.

The biblical account reveals: “Then cometh Jesus with them unto a place called Gethsemane, and saith unto the disciples, Sit ye here, while I go and pray yonder.

“And he took with him Peter, [James, and John] and began to be sorrowful and very heavy.

“Then saith he unto them, My soul is exceeding sorrowful, even unto death: tarry ye here, and watch with me.

“And he went a little further, … and prayed, saying,”

“Father, if thou be willing, remove this cup from me: nevertheless not my will, but thine, be done.

“And there appeared an angel unto him from heaven, strengthening him.

“And being in an agony he prayed more earnestly: and his sweat was as it were great drops of blood falling down to the ground.”

What suffering, what sacrifice, what anguish did He endure to atone for the sins of the world!

For our benefit, the poet wrote:

\begin{verse}
In golden youth when seems the earth\\
A summer-land of singing mirth,\\
When souls are glad and hearts are light,\\
And not a shadow lurks in sight,\\
We do not know it, but there lies\\
Somewhere veiled ‘neath evening skies\\
A garden which we all must see—\\
The garden of Gethsemane. …
\end{verse}

The mortal mission of the Savior of the world drew rapidly to its close. Ahead lay Calvary’s cross, the acts of depravity committed by those who thirsted for the blood of the Son of God. His divine response is a simple but profoundly significant prayer: “Father, forgive them; for they know not what they do.”

The conclusion came: “Father, into thy hands I commend my spirit: and having said thus,” the Great Redeemer died. He was buried in a tomb. He rose on the morning of the third day. He was seen by His disciples. Words that linger from that epochal event course through the annals of time and bring to our souls even today the comfort, the assurance, the balm, the certainty, “He is not here: … he is risen.” Resurrection became a reality for all.

Last week I received a faith-filled letter from Laurence M. Hilton. May I share with you the account of surviving personal tragedy with faith, nothing wavering.

In 1892, Thomas and Sarah Hilton, Laurence’s grandparents, went to Samoa, where Thomas was set apart as mission president after their arrival. They brought with them a baby daughter; two sons were born to them while they served there. Tragically, all three died in Samoa, and in 1895 the Hiltons returned from their mission childless.

David O. McKay was a friend of the family and was deeply touched by their loss. In 1921, as part of a world tour of visits to the members of the Church in many nations, Elder McKay stopped in Samoa, accompanied by Elder Hugh J. Cannon. Before leaving on his tour, he had promised the now-widowed Sister Hilton that he would personally visit the graves of her three children. I share with you the letter David O. McKay wrote to her from Samoa:

“Dear Sister Hilton:

“Just as the descending rays of the late afternoon sun touched the tops of the tall coconut trees, Wednesday, May 18th, 1921, a party of five stood with bowed heads in front of the little Fagali’i Cemetery. … We were there, as you will remember, in response to a promise I made you before I left home.

“The graves and headstones are in a good state of preservation. … I reproduce here a copy I made as I stood … outside the stone wall surrounding the spot.

\begin{verse}
Janette Hilton\\
Bn: Sept. 10, 1891\\
Died: June 4, 1892\\
“Rest, darling Jennie”
\end{verse}

“As I looked at those three little graves, I tried to imagine the scenes through which you passed during your young motherhood here in old Samoa. As I did so, the little headstones became monuments not only to the little babes sleeping beneath them, but also to a mother’s faith and devotion to the eternal principles of truth and life. Your three little ones, Sister Hilton, in silence most eloquent and effective, have continued to carry on your noble missionary work begun nearly 30 years ago, and they will continue as long as there are gentle hands to care for their last earthly resting place.

\begin{verse}
By loving hands their dying eyes were closed;\\
By loving hands their little limbs composed;\\
By foreign hands their humble graves adorned;\\
By strangers honored, and by strangers mourned.\\
“Tofa Soifua,\\
“David O. McKay”
\end{verse}

This touching account conveys to the grieving heart “the peace … which passeth all understanding.” Our Heavenly Father lives. Jesus Christ the Lord is our Savior and Redeemer. He guided the Prophet Joseph. He guides His prophet today, even President Gordon B. Hinckley. Of a truth I bear this personal witness.

That we may shoulder our sorrows, bear our burdens, and face our fears—as did our Savior—is my prayer. I know that He lives. In the name of Jesus Christ, amen.

\newpage
\settalk{Think to Thank}
\setspeaker{Thomas S. Monson}
\settalkdate{October 1998}
\talkheading{Think to Thank}{Thomas S. Monson}

In a land far away, and at a time long ago, Jesus journeyed to Jerusalem. “He passed through the midst of Samaria and Galilee.

“And as he entered into a certain village, there met him ten men that were lepers, which stood afar off:

“And they lifted up their voices, and said, Jesus, Master, have mercy on us.

“And when he saw them, he said unto them, Go shew yourselves unto the priests. And it came to pass, that, as they went, they were cleansed.

“And one of them, when he saw that he was healed, turned back, and with a loud voice glorified God,

“And fell down on his face at his feet, giving him thanks: and he was a Samaritan.

“And Jesus answering said, Were there not ten cleansed? but where are the nine?

“There are not found that returned to give glory to God, save this stranger.

“And he said unto him, Arise, go thy way: thy faith hath made thee whole.”

From the 30th Psalm, David pledges, “O Lord my God, I will give thanks unto thee for ever.”

The Apostle Paul, in his epistle to the Corinthians, proclaimed, “Thanks be unto God for his unspeakable gift.” And to the Thessalonians, “In every thing give thanks: for this is the will of God.”

My brothers and sisters, do we give thanks to God “for his unspeakable gift” and His rich blessings so abundantly bestowed upon us?

Do we pause and ponder Ammon’s words? “Now my brethren, we see that God is mindful of every people, whatsoever land they may be in; yea, he numbereth his people … over all the earth. Now this is my joy, and my great thanksgiving; yea, and I will give thanks unto my God forever.”

Robert W. Woodruff, a prominent business leader of a former time, toured the United States giving a lecture which he entitled “A Capsule Course in Human Relations.” In his message, he said that the two most important words in the English language are these: “Thank you.”

Gracias, danke, merci—whatever language is spoken, “thank you” frequently expressed will cheer your spirit, broaden your friendships, and lift your lives to a higher pathway as you journey toward perfection. There is a simplicity—even a sincerity—when “thank you” is spoken.

The beauty and eloquence of an expression of gratitude is reflected in a newspaper story of some years ago:

The District of Columbia police auctioned off about 100 unclaimed bicycles Friday. “One dollar,” said an 11-year-old boy as the bidding opened on the first bike. The bidding, however, went much higher. “One dollar,” the boy repeated hopefully each time another bike came up.

The auctioneer, who had been auctioning stolen or lost bikes for 43 years, noticed that the boy’s hopes seemed to soar higher whenever a racer-type bicycle was put up.

Then there was just one racer left. The bidding went to eight dollars. “Sold to that boy over there for nine dollars!” said the auctioneer. He took eight dollars from his own pocket and asked the boy for his dollar. The youngster turned it over in pennies, nickels, dimes, and quarters—took his bike, and started to leave. But he went only a few feet. Carefully parking his new possession, he went back, gratefully threw his arms around the auctioneer’s neck, and cried.

When was the last time we felt gratitude as deeply as did this boy? The deeds others perform in our behalf might not be as poignant, but certainly there are kind acts that warrant our expressions of gratitude.

The song frequently sung in the Sunday School of our youth placed the spirit of thanksgiving into the depths of our souls:

\begin{verse}
When upon life’s billows you are tempest-tossed,\\
When you are discouraged, thinking all is lost,\\
Count your many blessings; name them one by one,\\
And it will surprise you what the Lord has done.
\end{verse}

Astronaut Gordon Cooper, while orbiting the earth over 30 years ago, offered this sweet and simple prayer of thanks: “Father, thank You, especially for letting me fly this flight. Thank You for the privilege of being able to be in this position: to be up in this wondrous place, seeing all these many startling, wonderful things that You have created.”

We are thankful for blessings we cannot measure, for gifts we cannot appraise, “for books, music, art, and for the great inventions which make these blessings available; … for the laughter of little children; … for the … means for relieving human suffering … and increasing … the enjoyment of life; … for everything good and uplifting.”

The prophet Alma urged, “Counsel with the Lord in all thy doings, and he will direct thee for good; yea, when thou liest down at night lie down unto the Lord, that he may watch over you in your sleep; and when thou risest in the morning let thy heart be full of thanks unto God; and if ye do these things, ye shall be lifted up at the last day.”

I would like to mention three instances where I believe a sincere “thank you” could lift a heavy heart, inspire a good deed, and bring heaven’s blessings closer to the challenges of our day.

First, may I ask that we express thanks to our parents for life, for caring, for sacrificing, for laboring to provide a knowledge of our Heavenly Father’s plan for happiness.

From Sinai the words thunder to our conscience, “Honour thy father and thy mother: that thy days may be long upon the land which the Lord thy God giveth thee.”

I know of no sweeter expression toward a parent than that spoken by our Savior upon the cross: “When Jesus therefore saw his mother, and the disciple standing by, whom he loved, he saith unto his mother, Woman, behold thy son!

“Then saith he to the disciple, Behold thy mother! And from that hour that disciple took her unto his own home.”

Next, have we thought on occasion of a certain teacher at school or at church who seemed to quicken our desire to learn, who instilled in us a commitment to live with honor?

The story is told of a group of men who were talking about people who had influenced their lives and for whom they were grateful. One man thought of a high school teacher who had introduced him to Tennyson. He decided to write and thank her. In time, written in a feeble scrawl, came the teacher’s reply:

“My dear Willie:

“I can’t tell you how much your note meant to me. I am in my 80s, living alone in a small room, cooking my own meals, lonely and like the last leaf lingering behind. You will be interested to know that I taught school for 50 years, and yours is the first note of appreciation I have ever received. It came on a blue, cold morning, and it cheered me as nothing has for years.”

We owe an eternal debt of gratitude to all of those, past and present, who have given so much of themselves, that we might have so much ourselves.

Third, I mention an expression of “thank you” to one’s peers. The teenage years can be difficult for the teens themselves as well as for their parents. These are trying times in the life of a boy or a girl. Each boy wants to make the football team; each girl wants to be the beauty queen. “Many are called, but few are chosen” could have an application here.

Let me share with you a modern-day miracle which occurred a year or so ago at Murray High School near Salt Lake City, where every person was a winner, and not a loser was to be found.

A newspaper article highlighted the event. It was entitled “Tears, Cheers and True Spirit: Students Elect 2 Disabled Girls to Murray Royalty.” The article began:

“Ted and Ruth Eyre did what any parents would do.

“When their daughter, Shellie, became a finalist for Murray High School homecoming queen, they counseled her to be a good sport in case she didn’t win. They explained only one girl among the 10 … would be selected queen. …

“As student body officers crowned the school’s homecoming [royalty] in the school gym Thursday night, Shellie Eyre experienced, instead, inclusion. The 17-year-old senior, born with Down syndrome, was selected by fellow students as homecoming queen. …

“… As Ted Eyre escorted his daughter onto the gym floor as the candidates were introduced, the gym erupted into deafening cheers and applause. They were greeted with a standing ovation.”

Similar standing ovations were extended to Shellie’s attendants, one of whom, April Perschon, has physical and mental disabilities resulting from a brain hemorrhage suffered when she was just 10 years old.

When the ovations had ceased, the school vice principal Glo Merrill said, “‘Tonight … the students voted on inner beauty.’ …

“Obviously moved, parents, school administrators and students wept openly.” Said one student, “‘I’m so happy. I cried when they came out. I think Murray High is so awesome to do this.’”

I extend a heartfelt “thank you” to one and all who made this night one ever to be remembered. To paraphrase the Scottish poet James Barrie, “God gave us memories, that we might have June roses in the December of our lives.”

In August of this year, there occurred a tragedy in Salt Lake County. It was reported in the local and national press. Five beautiful little girls—so young, so vibrant, so loving—hiding away, as children often do in their games of hide-and-seek, entered the trunk of a parent’s car. The trunk lid was pulled shut, they were unable to escape, and all perished from heat exhaustion.

The entire community was so kind, so thoughtful, so caring in the passing of Alisha, Ashley, McKell, Audrey, and Jaesha. Flowers, food, calls, visits, and prayers were shared.

On the Sunday after the devastating event occurred, long lines of automobiles filled with grieving occupants drove ever so slowly past the Smith home, the scene of the accident. Sister Monson and I wished to be among those who expressed condolences in this way. As we drove by, we felt we were on holy ground. We literally crept along at a snail’s pace along the street. It was as though we could visualize a traffic sign reading, “Please drive slowly; children at play.” Tears filled our eyes and compassion flowed from our hearts.

At the funeral, as well as the evening prior, thousands passed by the caskets and expressed support for the grieving parents and grandparents. In two of the three families, the deceased children were all the children they had.

Frequently death comes as an intruder. It is an enemy that suddenly appears in the midst of life’s feast, putting out its lights and gaiety. It visits the aged as they walk on faltering feet. Its summons is heard by those who have scarcely reached midway in life’s journey, and often it hushes the laughter of little children.

At the funeral services for the five little angels, I counseled: “There is one phrase which should be erased from your thinking and from the words you speak aloud. It is the phrase ‘If only.’ It is counterproductive and is not conducive to the spirit of healing and of peace. Rather, recall the words of Proverbs: ‘Trust in the Lord with all thine heart; and lean not unto thine own understanding. In all thy ways acknowledge him, and he shall direct thy paths.’”

Before the closing of the caskets, I noted that each child held a favorite toy, a soft gift to cuddle. I reflected on the words of the poet Eugene Field:

\begin{verse}
The little toy dog is covered with dust,\\
But sturdy and staunch he stands;\\
And the little toy soldier is red with rust,\\
And his musket moulds in his hands.\\
Time was when the little toy dog was new,\\
And the soldier was passing fair,\\
And that was the time when our Little Boy Blue\\
Kissed them and put them there.
\end{verse}

The little toy dog and the soldier fair may wonder, but God in His infinite mercy has not left grieving loved ones to wonder. He has provided truth. He will inspire an upward reach, and His outstretched arms will embrace you. Jesus promises to one and all who grieve, “I will not leave you comfortless: I will come to you.”

There is only one source of true peace. I am certain that the Lord, who notes the fall of a sparrow, looks with compassion upon those who have been called upon to part—even temporarily—from their precious children. The gifts of healing and of peace are desperately needed, and Jesus, through His Atonement, has provided them for one and all.

The Prophet Joseph Smith spoke inspired words of revelation and comfort:

“All children who die before they arrive at the years of accountability are saved in the celestial kingdom of heaven.”

“The mother [and father] who laid down [their] little child[ren], being deprived of the privilege, the joy, and the satisfaction of bringing [them] up to manhood or womanhood in this world, would, after the resurrection, have all the joy, satisfaction and pleasure, and even more than it would have been possible to have had in mortality, in seeing [their] child[ren] grow to the full measure of the stature of [their] spirit[s].” This is as the balm of Gilead to those who grieve, to those who have loved and lost precious children.

The Psalmist provided this assurance: “Weeping may endure for a night, but joy cometh in the morning.”

Said the Lord: “Peace I leave with you, my peace I give unto you: not as the world giveth, give I unto you. Let not your heart be troubled, neither let it be afraid.”

“In my Father’s house are many mansions: if it were not so, I would have told you. I go to prepare a place for you … that where I am, there ye may be also.”

I express my profound thanks to a loving Heavenly Father, who gives to you, to me, and to all who sincerely seek, the knowledge that death is not the end, that His Son—even our Savior, Jesus Christ—died that we might live. Temples of the Lord dot the lands of many countries. Sacred covenants are made. Celestial glory awaits the obedient. Families can be together—forever.

The Master invites one and all:

“Come unto me, all ye that labour and are heavy laden, and I will give you rest.

“Take my yoke upon you, and learn of me; for I am meek and lowly in heart: and ye shall find rest unto your souls.”

That all may do so is my humble prayer of thanks, in the name of Jesus Christ, amen.

\newpage
\settalk{Today Determines Tomorrow}
\setspeaker{Thomas S. Monson}
\settalkdate{October 1998}
\talkheading{Today Determines Tomorrow}{Thomas S. Monson}

It is a joy and a privilege for me to stand before you, such a vast audience of priesthood holders both seen and unseen. General Church priesthood meetings have always been a treat—from Aaronic Priesthood days until the present. To “come, listen to a prophet’s voice, and hear the word of God,” as a song from our hymnbook states, is a cherished blessing.

We sustain Gordon B. Hinckley as the President of The Church of Jesus Christ of Latter-day Saints and as the prophet, seer, and revelator of the Church in our time. A letter which I received from a proud father tells of an experience with his then five-year-old son and the boy’s love for the President of the Church and desire to emulate the President’s example. The father wrote:

“When Christopher was five years old, he would get ready for church on Sundays mostly by himself. On one particular Sunday, he decided that he wanted to wear a suit and tie, which to that point he had never done. He scoured the closet on his own for a hand-me-down tie and produced a rather used clip-on one that he didn’t need to create a knot for. He attached the tie to his white shirt, then capped it off with the small navy jacket that had hung for years in the boys’ closet.

“On his own, he went into the bathroom and painstakingly combed his blonde hair to perfection. About that time, I came into the bathroom to finish getting ready myself. I found Christopher beaming at himself in the mirror. Without taking his eyes off his reflection, he proclaimed proudly, ‘Look, Papa—Christopher B. Hinckley!’” And Father realized that a boy had been watching the prophet of the Lord.

Our children are watching. They are absorbing eternal lessons. They are shaping their futures. What is the example we are presenting to them?

Years ago when our youngest son, Clark, was attending a religion class at Brigham Young University, the instructor, during a lecture, asked Clark, “What is an example of life with your father that you best remember?”

The instructor later wrote to me and told me of the reply which Clark had given to the class. Said Clark: “When I was a deacon in the Aaronic Priesthood, my dad and I went pheasant hunting near Malad, Idaho. The day was Monday—the last day of the season. We walked through countless fields in search of pheasants but only saw a few, and these we missed. Dad then said to me, ‘Clark, let’s unload our guns, and we’ll place them in this ditch. Then we’ll kneel down to pray.’ I thought Dad would pray for more pheasants, but I was wrong. He explained to me that Elder Richard L. Evans was gravely ill and that at 12 noon on that particular Monday the members of the Quorum of the Twelve—wherever they may be at the time—were to kneel and, in a way, together unite in a fervent prayer of faith for Elder Evans. Removing our caps, we knelt, we prayed.”

I well remember the occasion, but I never dreamed a son was watching, was learning, was building his own testimony.

In analyzing the statistical performance of those who hold the Aaronic Priesthood as deacons, teachers, and priests, we become concerned when significant numbers of deacons slip into inactivity and fail to be ordained teachers at the proper time. The same is true with some who are teachers but not ordained priests—and particularly priests who never receive the Melchizedek Priesthood. Brethren, this should never be. We have an awesome responsibility to guide and inspire these young men on the priesthood trail, that no avalanche of sin or error will deter their progress or sweep them away from their eternal goals.

Bishops and bishops’ counselors, will you undertake a study of the activity levels of each Aaronic Priesthood young man and outline your own plan to ensure the progress and activity of each one?

One newly called bishop, in his first meeting with his counselors, declared, “The Aaronic Priesthood is a prime responsibility of ours.” To the second counselor, he directed, “I ask you to be personally responsible to ensure that every deacon, at the appropriate age, be worthy and be ordained a teacher.” To the other counselor, he said, “Will you please do the same as pertains to the teachers, that they may, on schedule, be worthy and be ordained priests.” Then the bishop continued, “I will take the same responsibility for the priests to receive the Melchizedek Priesthood and be ordained elders. Together, and with God’s help, we can do it.” And they did.

Our youth need less criticism and more models to follow. You advisers to the Aaronic Priesthood quorums are teachers and models for the young men. Do you know the gospel? Have you prepared the lesson? Do you know each boy and prayerfully determine how you might reach his mind, his heart, and help fashion his future?

Remember, it isn’t sufficient to assume that when you teach, the boy is listening to what you say. Let me illustrate:

In what we call the west boardroom of the Church Administration Building, there hangs a lovely painting rendered by the artist Harry Anderson. The painting depicts Jesus sitting on a small stone wall with numerous children gathered around, knowing they are the object of His love. Each time I gaze at that painting, I think of the passage of scripture, “Suffer the little children to come unto me, and forbid them not: for of such is the kingdom of God.”

On one occasion I had given a priesthood blessing in that room to a small lad who was soon to undergo major surgery. I directed his attention and that of his parents to the painting of Jesus and the children. I then made a few remarks concerning the Savior and His never-failing love. I asked the boy if he had any questions. “Yes,” he replied seriously. “Brother Monson, how does a boy go about getting a little goat and a leash for it like that one in the painting?”

For a moment I was stunned by the unanticipated question, a little deflated concerning my teaching ability, but then I responded: “Jesus gives to you and me gifts far more important than a goat on a leash. He provides a road map to heaven. His teachings, His example, His love are far greater gifts than that offered by the world.”

“Come, follow me,” He invited. And we are wise when we follow Him!

Let all young men who bear the Aaronic Priesthood learn and live the Savior’s teachings and prepare to receive the Melchizedek Priesthood.

May I share with you brethren my personal experience as a teachers quorum president? The member of the bishopric who had responsibility for us invited the new presidency and secretary to come to his home for leadership training. He wanted our ideas concerning how we should go about our newly given duties. We obliged—on condition that he would invite his wife, Nettie, to serve us some of the meat pies for which she was famous. This he agreed to do. Brethren, isn’t it remarkable how we men will obligate our wives to do things—often without notice? The resulting meeting was one of the best I have ever attended. We were taught to the level of our understanding and inspired to look after our quorum members.

After a delicious meat pie smothered with gravy, we asked the bishop’s counselor and his wife to join in a game of Monopoly®. I am certain they had other things to do, but they willingly complied with our request.

I don’t remember who won the Monopoly game, but I have never forgotten the lessons learned that night in Church government and in the administration of a priesthood quorum.

During the fervor of the early years of World War II, one of our teachers quorum members, Fritz, wanted to defend our country but didn’t want to wait until he reached the minimum age required to serve. He falsified his age and enlisted in the United States Navy. Soon he found himself far away in the Pacific sea battles. The vessel on which he served was sent to the bottom, with many hands lost. Fritz survived and later appeared in our quorum meeting in full uniform, with battle ribbons affixed. I remember asking Fritz, “Fritz, do you have any advice for us?” We were all on the very doorstep of mandatory military service.

Fritz thought for a moment and then said, “Never lie about your age or about anything else!” That one-sentence declaration is remembered yet.

Young men between the ages of 12 and 17 are in a time of preparation and personal spiritual growth. Accordingly, the purposes of the Aaronic Priesthood are to help each person who is ordained:

Serving throughout the world is a great missionary force, going about doing good as did the Savior. Missionaries teach truth. They dispel darkness. They spread joy. They bring precious souls to Christ.

On that special day when a mission call is received, parents, brothers and sisters, and grandparents gather around the prospective missionary and note his nervousness as he carefully opens the letter of call. There is a pause, and then he announces where the prophet of the Lord has assigned him to serve. Feelings are very near to the surface. Tears come easily, and the family rejoices in the bond of love and the goodness of God.

The full-time missionaries and all others engaged in the work of the Lord have answered His call. We are on His errand. We shall succeed in the solemn charge given by Mormon to declare the Lord’s word among the people. Wrote Mormon: “Behold, I am a disciple of Jesus Christ, the Son of God. I have been called of him to declare his word among his people, that they might have everlasting life.”

In 1926 President Fred Tadje, president of the German-Austrian Mission, called a mission conference to be held at Dresden, Germany, in August. The missionaries were to walk to this conference from their fields of labor basically “without purse or scrip,” although they had to carry a small amount of money or they could be arrested as vagabonds.

Elder Alfred Lippold and his companion, Elder Parker Thomas, took the north route. Somewhere along the way, the two called at a home where they met a woman and her eight children. She told the elders that her husband had left her and the children and that they were now without money. After she had let them in, the woman said: “If you travel without purse or scrip, then you must be hungry. Sit down.” She gave each of them a big slice of bread with plum jam on it. The missionaries blessed the breakfast and in the blessing on the food asked the Lord to give the woman what she needed.

The missionaries then departed. After they had walked about a mile, Elder Thomas said, “I must go back,” which he did without explanation.

On his return, Elder Lippold asked, “Why did you go back?”

Elder Thomas explained: “In our prayer we asked that the woman be given what she needed. I had what was needed—a \$20 bill. It was in my pocket, and I went back to give it to her. It would have burned a hole in my pocket.”

Thirty years ago I had responsibility for much of the work in the South Pacific. A Brother J. Vernon Monson was called, together with his wife, to journey to faraway Rarotonga in the Cook Islands, there to serve as district president.

Later, in a letter to me, he reported: “We are most grateful for the progress being made, and I would especially like to mention the goodwill and wonderful relations that have developed with the representatives of government and the business community toward us and the Church.

“One thing climaxed the development of this public acceptance,” he wrote. “It was in having our nephew and niece, Dr. and Mrs. Odeen Manning, render an outstanding service here in the Cook Islands. Dr. Manning is an ophthalmologist, and I wrote to him outlining a proposal whereby he might render service to the people of Rarotonga. My proposal included the following: (1) no remuneration; (2) he must pay his own expenses; (3) that he turn his practice over to the other doctors to handle for the three months he would be away; (4) we would furnish them free board and room while in Rarotonga; and (5) that he bring his own surgical instruments, as none would be available in Rarotonga.”

Brother Vernon Monson’s letter to me continued: “The Mannings airmailed their reply in two words: ‘Offer accepted.’ As preparations began, the government of the Cook Islands assigned competent doctors to assist Dr. Manning and to learn from him. In all, 284 patients were examined, with most being fitted for glasses. Fifty-three patients had serious eye operations, such as cataract surgery.

“The entire three-month program was wonderful and most heartwarming. Truly we were blessed. It has buoyed up the Saints, who gained new pride in being members of a faith which would bring medical service to these islands.” The letter ended.

Years later, my wife and I were guests on a BYU-sponsored cruise to the Holy Land. One evening as we were seated on the ship’s deck, the man sitting next to us turned to me and said, “Elder Monson, my name is Odeen Manning from Woodland Hills, California. I am an ophthalmologist by profession and served a brief medical mission to Rarotonga when my uncle and aunt were serving there.”

I acknowledged that I was aware of his sacrifice and his service. I asked Dr. Manning, “As you reflect on this experience, would you wish to share with me your feelings concerning it?”

He responded with emotion, saying, “It was the most spiritually rewarding experience of my life.”

I believe it was more than coincidence that my wife and I would be on the cruise vessel at that particular time and in that particular area of the deck, sitting next to a man we never before had met. Heaven was close as Dr. Manning and I embraced, and thanks were expressed for his service—not only to those who were blind and now could see, but also to our Lord and Savior. As Jacob declared, “Great are the promises of the Lord unto them who are upon the isles of the sea.”

Of Him who delivered each of us from endless death, even Jesus Christ, I testify that He is a teacher of truth—but He is more than a teacher. He is the Exemplar of the perfect life—but He is more than an exemplar. He is the Great Physician—but He is more than a physician. He who rescued the “lost battalion” of mankind is the literal Savior of the world, the Son of God, the Prince of Peace, the Holy One of Israel—even the risen Lord—who declared, “I am the first and the last; I am he who liveth, I am he who was slain; I am your advocate with the Father.”

My dear brethren, let each of us:

By so doing we can become like Him. Of this truth I solemnly bear witness in the name of Jesus Christ, amen.

\newpage
\settalk{For I Was Blind, but Now I See}
\setspeaker{Thomas S. Monson}
\settalkdate{April 1999}
\talkheading{For I Was Blind, but Now I See}{Thomas S. Monson}

When Jesus walked and taught among men, He spoke in language easily understood. Whether He was journeying along the dusty way from Perea to Jerusalem, addressing the multitudes on the shore of the Sea of Galilee, or pausing beside Jacob’s well in Samaria, He taught in parables. Jesus spoke frequently of having hearts that could know and feel, ears that were capable of hearing, and eyes that could truly see.

One not so blessed with the gift of sight was the blind man who, in an effort to sustain himself, sat day in and day out at his usual place on the edge of a busy sidewalk in one of our large cities. In one hand he held an old felt hat filled with pencils. With his other hand he held out a tin cup. His simple appeal to the passerby was brief and to the point. It had a certain finality to it, almost a tone of despair. The message was contained on the small placard held about his neck by a string. It read, “I am blind.”

Most did not stop to buy his pencils or to place a coin in the tin cup. They were too busy, too occupied by their own problems. That tin cup had never been filled or even half-filled. Then one beautiful spring day a man paused and, with a marking pen, added several new words to the shabby sign. No longer did it read, “I am blind.” Now the message read, “It is springtime and I am blind.” The cup was soon filled to overflowing. Perhaps the busy people were touched by Charles L. O’Donnell’s exclamation, “I have never been able to school my eyes against young April’s blue surprise.” To each, however, the coins were a poor substitute for the desired ability to actually restore sight.

Each of us knows those who do not have sight. We also know many others who have their eyesight but who walk in darkness at noonday. These in this latter group may never carry the usual white cane and carefully make their way to the sound of the familiar “tap, tap, tap.” They may not have a faithful Seeing Eye dog by their side nor carry a sign about their neck which reads, “I am blind,” but blind they surely are. Some have been blinded by anger, others by indifference, by revenge, by hate, by prejudice, by ignorance, by neglect of precious opportunities. Of such the Lord said, “Their ears are dull of hearing, and their eyes they have closed; lest at any time they should see with their eyes, and hear with their ears, and should understand with their heart, and should be converted, and I should heal them.”

Well might each lament, “It is springtime, the gospel of Jesus Christ has been restored, and yet I am blind.” Some, like the friend of Philip of old, call out, “How can I [find my way], except some man should guide me?”

Many years ago, while attending a stake conference, I noticed that a counselor in the stake presidency was blind. He functioned beautifully, performing his duties as though he had sight. It was a stormy night as we met in the stake office situated on the second floor of the building. Suddenly there was a loud clap of thunder. The lights in the building almost immediately went out. Instinctively I reached out for our sightless leader, and I said, “Here, take my arm and I will help you down the stairway.”

I’m certain he must have had a smile on his face as he responded, “No, Brother Monson, give me your arm, that I might help you.” And he added, “You are now in my territory.”

The storm abated, the lights returned, but I shall never forget the trek down those stairs, guided by the man who was sightless yet filled with light.

Long ago and at a place far distant, as Jesus passed by He saw a man who was blind from birth. His disciples questioned the Master as to why this person was blind. Had he sinned or had his parents sinned, causing him to have this affliction?

“Jesus answered, Neither hath this man sinned, nor his parents: but that the works of God should be made manifest in him. …

“As long as I am in the world, I am the light of the world.

“When he had thus spoken, he spat on the ground, and made clay of the spittle, and he anointed the eyes of the blind man with the clay,

“And said unto him, Go, wash in the pool of Siloam. … He went his way therefore, and washed, and came seeing.”

A great dispute ensued among the Pharisees concerning this miracle:

“Then again called they the man that was blind, and said unto him, Give God the praise: we know that this man [Jesus] is a sinner.

“He answered and said, Whether he be a sinner or no, I know not: one thing I know, that, whereas I was blind, now I see.”

One thinks of the fisherman called Simon, better known to you and to me as Peter, chief among the Apostles. Doubting, disbelieving, impetuous Peter, in fulfillment of the Master’s prophecy, indeed did deny Him thrice. Amidst the pushing, the jeers, and the blows, “the Lord in the agony of His humiliation, in the majesty of His silence, … ‘turned and looked upon Peter.’” As one chronologist described the change: “It was enough. … [Peter] ‘knew no more danger, he feared no more death.’ … [He] rushed forth into the night … ‘to meet the morning dawn.’ … This broken-hearted penitent [stood] before the tribunal of his own conscience, and there his old life, his old shame, his old weakness, his old self was doomed to that death of godly sorrow which was to issue in a new and a [nobler] birth.”

The Apostle Paul had a similar experience to that of Peter. From the day of his conversion until the day of his death, Paul urged men to “put off … the old man” and to “put on the new man, which after God is created in righteousness and true holiness.”

Simon the fisherman had become Peter the Apostle. Saul the persecutor had become Paul the proselyter.

The passage of time has not altered the capacity of the Redeemer to change men’s lives. As He said to the dead Lazarus, so He says to you and to me: “Come forth.”

Said President Harold B. Lee: “Every soul who walks the earth, wherever he lives, in whatever nation he may have been born, no matter whether he be in riches or in poverty, had at birth an endowment of that first light which is called the Light of Christ, the Spirit of Truth, or the Spirit of God—that universal light of intelligence with which every soul is blessed. Moroni spoke of that Spirit when he said:

“‘For behold, the Spirit of Christ is given to every man, that he may know good from evil; wherefore, I show unto you the way to judge; for every thing which inviteth to do good, and to persuade to believe in Christ, is sent forth by the power and gift of Christ; wherefore ye may know with a perfect knowledge it is of God.’”

You and I know those who qualify for the Savior’s blessing in accordance with this definition.

Such was Walter Stover of Salt Lake City. Born in Germany, Walter embraced the gospel message and came to America. He established his own business. He gave freely of his time and of his means.

Following World War II Walter Stover was called to return to his native land. He directed the Church in that nation and blessed the lives of all whom he met and with whom he served. With his own funds he constructed two chapels in Berlin—a beautiful city that had been so devastated by the conflict. He planned a gathering in Dresden for all the members of the Church from that nation and then chartered a train to bring them from all around the land so they could meet, partake of the sacrament, and bear witness of the goodness of God to them.

At the funeral services for Walter Stover, his son-in-law Thomas C. LeDuc said of him, “He had the ability to see Christ in every face he encountered, and he acted accordingly.”

The poet wrote:

\begin{verse}
I met a stranger in the night, whose lamp had ceased to shine;\\
I paused and let him light his lamp from mine.\\
A tempest sprang up later on, and shook the world about,\\
And when the wind was gone, my lamp was out.\\
But back came to me the stranger—his lamp was glowing fine;\\
He held the precious flame and lighted mine.
\end{verse}

Perhaps the moral of this poem is simply that if you want to give a light to others, you have to glow yourself.

When the Prophet Joseph Smith went into a grove of trees made sacred by what occurred there, he described the event:

“It was on the morning of a beautiful, clear day, early in the spring of eighteen hundred and twenty. It was the first time in my life that I had made such an attempt, for amidst all my anxieties I had never as yet made the attempt to pray vocally.”

After enduring a harrowing experience from an unseen power, Joseph continued:

“I saw a pillar of light exactly over my head, above the brightness of the sun, which descended gradually until it fell upon me. …

“… When the light rested upon me I saw two Personages, whose brightness and glory defy all description, standing above me in the air. One of them spake unto me, calling me by name and said, pointing to the other—This is My Beloved Son. Hear Him!”

Joseph listened. Joseph learned.

On occasion I will be asked, “Brother Monson, if the Savior appeared to you, what questions would you ask of Him?”

My reply is always the same: “I would ask no question of Him. Rather, I would listen!”

Late one evening on a Pacific isle, a small boat slipped silently to its berth at the crude pier. Two Polynesian women helped Meli Mulipola from the boat and guided him to the well-worn pathway leading to the village road. The women marveled at the bright stars which twinkled in the midnight sky. The friendly moonlight guided them along their way. However, Meli Mulipola could not appreciate these delights of nature—the moon, the stars, the sky—for he was blind.

His vision had been normal until that fateful day when, while working on a pineapple plantation, light turned suddenly to darkness and day became perpetual night. He had learned of the restoration of the gospel and the teachings of The Church of Jesus Christ of Latter-day Saints. His life had been brought into compliance with these teachings.

He and his loved ones had made this long voyage, having learned that one who held the priesthood of God was visiting among the islands. He sought a blessing under the hands of those who held the sacred priesthood. His wish was granted. Tears streamed from his sightless eyes and coursed down his brown cheeks, tumbling finally upon his native dress. He dropped to his knees and prayed: “Oh, God, thou knowest I am blind. Thy servants have blessed me that if it be thy will, my sight may return. Whether in thy wisdom I see light or whether I see darkness all the days of my life, I will be eternally grateful for the truth of thy gospel which I now see and which provides me the light of life.”

He arose to his feet, thanked us for providing the blessing, and disappeared into the dark of the night. Silently he came; silently he departed. But his presence I shall never forget. I reflected upon the message of the Master: “I am the light of the world: he that followeth me shall not walk in darkness, but shall have the light of life.”

Today is a day of temple building. Never before have so many temples been erected and dedicated. President Gordon B. Hinckley, God’s prophet on this earth, has a vision of the vital ordinances performed in such houses of the Lord. Temples will bless all who attend them and who sacrifice for their completion. The light of Christ will shine on all—even those who have gone beyond. President Joseph F. Smith, speaking of work for the dead, declared, “Through our efforts in their behalf their chains of bondage will fall from them, and the darkness surrounding them will clear away, that light may shine upon them and they shall hear in the spirit world of the work that has been done for them by their children here, and will rejoice with you in your performance of these duties.”

The Apostle Paul urged, “Be thou an example of the believers.” And from James: “Be ye doers of the word, and not hearers only, deceiving your own selves.”

I close with the words of the poet Minnie Louise Haskins, who wrote:

\begin{verse}
And I said to the man who stood at the gate of the year:\\
“Give me a light, that I may tread safely into the unknown!”\\
And he replied:\\
“Go out into the darkness and put your hand into the Hand of God.\\
That shall be to you better than light and safer than a known way.”\\
So, I went forth, and finding the Hand of God, trod gladly into the night.\\
And He led me toward the hills and the breaking of day in the lone East.
\end{verse}

On this Easter morning and always, may our light so shine that we glorify our Heavenly Father and His Son, Jesus Christ, whose name is the only name under heaven whereby we might be saved.

That we may ever walk in the footsteps of Jesus Christ is my humble prayer, in His holy name, amen.

\newpage
\settalk{The Priesthood—Mighty Army of the Lord}
\setspeaker{Thomas S. Monson}
\settalkdate{April 1999}
\talkheading{The Priesthood—Mighty Army of the Lord}{Thomas S. Monson}

I am honored tonight to be with the vast army of priesthood bearers who daily respond to calls to serve, who teach diligently as the Lord has commanded, and who labor mightily to bring a correction course to a specific challenge which the Church must meet—namely, to live in the world without being of the world.

In this day in which we live, the floodwaters of immorality, irresponsibility, and dishonesty lap at the very moorings of our individual lives. If we do not safeguard those moorings, if we do not have deeply entrenched foundations to withstand such eroding influences, we are going to experience difficulty.

One of the greatest safeguards we have in the Church is a strong, firm, committed, dedicated, and testifying Melchizedek Priesthood base.

In my office I have two small earthen containers. One is filled with water I retrieved from the Dead Sea. The other contains water from the Sea of Galilee. Occasionally I will shake one of the bottles to ensure that the water has not diminished. When I follow this practice, my mind turns to these two different bodies of water. The Dead Sea is void of life. The Sea of Galilee is filled with life and with memories of the mission of the Lord Jesus Christ.

There is another body of water found throughout the Church today. I speak of the pool of prospective elders in each ward and each stake. Picture in your mind a river of water gushing into the pool. Then consider a trickle of water emerging from that stagnated pool—a trickle which represents those going forward into the Melchizedek Priesthood. The pool of prospective elders is becoming larger and wider and deeper more rapidly than any of us can fully appreciate.

It is essential, even critical, that we study the Aaronic Priesthood pathway, since far too many boys falter, stumble, then fall without advancing into the quorums of the Melchizedek Priesthood, thereby eroding the active priesthood base of the Church and curtailing the activity of loving wives and precious children.

What can we as leaders do to reverse this trend? The place to begin is at the headwaters of the Aaronic Priesthood stream. There is an ancient proverb which purports to correctly determine the sanity of an individual. A person is shown a stream of water flowing into a stagnant pond. He is given a bucket and asked to commence to drain the pond. If he first takes steps to effectively dam the inflow to the pond, he is adjudged sane. If on the other hand he ignores the inflow and tries to empty the pond bucket by bucket, he is designated as insane.

The bishop, by revelation, is the president of the Aaronic Priesthood and is president of the priests quorum in his ward. He cannot delegate these God-given responsibilities. However, he can place accountability with those called as quorum advisers, men who can touch the lives of boys.

The bishop’s counselors, other ward officers and teachers, and particularly the fathers and the mothers of our young men can be of immeasurable help. Also very effective can be the service rendered by Aaronic Priesthood quorum presidencies.

This, then, is our goal: to save every young man, thereby assuring a worthy husband for each of our young women, strong Melchizedek Priesthood quorums, and a missionary force trained and capable of accomplishing what the Lord expects.

A wise first step is to guide each deacon to a spiritual awareness of the sacredness of his ordained calling. In one ward, this lesson was effectively taught pertaining to the collection of fast offerings.

On fast day the ward members were visited by deacons and teachers so that each family could make a contribution. The deacons were a bit disgruntled, having to arise earlier than usual to fulfill this assignment.

The inspiration came for the bishopric to take a busload of the deacons and teachers to Welfare Square here in Salt Lake City. Here they saw needy children receiving new shoes and other items of clothing. Here they witnessed empty baskets being filled with groceries. There was no money exchanged. One brief comment was made: “Young men, this is what the money you collect on fast day provides—even food, clothing, and shelter.” The Aaronic Priesthood young men smiled more, stepped higher, and served with a willing mind in the filling of their assignments.

A question: Is every ordained teacher given the assignment to home teach? What an opportunity to prepare for a mission. What a privilege to learn the discipline of duty. A boy will automatically turn from concern for self when he is assigned to “watch over” others.

And what of the priests? These young men have the opportunity to bless the sacrament, to continue their home teaching duties, and to participate in the sacred ordinance of baptism.

I remember as a deacon watching the priests as they would officiate at the sacrament table. One priest by the name of Barry had a lovely voice and would read the sacrament prayers with clear diction—as though he were competing in a speech contest. The other members of the ward, particularly the older sisters, would compliment him on his “golden voice.” I think he became a bit proud. Jack, another priest in the ward, was hearing impaired, which caused his speech to be unnatural in its sound. We deacons would twitter at times when Jack would bless the emblems. How we dared to do so is beyond me, for Jack had hands like a bear and could have crushed any one of us.

On one occasion Barry, with the beautiful voice, and Jack, with the awkward delivery, were assigned together at the sacrament table. The hymn was sung; the two priests broke the bread. Barry knelt to pray, and we closed our eyes. But nothing happened. Soon we deacons opened our eyes to see what was causing the delay. I shall ever remember the picture of Barry frantically searching the table for the little white card on which were printed the sacrament prayers. It was nowhere to be found. What to do? Barry’s face turned pink and then crimson as the congregation began to look in his direction.

Then Jack, with that bearlike hand, reached up and gently tugged Barry back onto the bench. He himself then knelt on the little footstool and began to pray: “O God, the Eternal Father, we ask thee in the name of thy Son, Jesus Christ, to bless and sanctify this bread to the souls of all those who partake of it.” He continued the prayer, and the bread was passed. Jack also blessed the water, and it was passed. What respect we deacons gained that day for Jack, who though handicapped in speech, had memorized the sacred prayers! Barry, too, had a new appreciation for Jack. A lasting bond of friendship had been established.

Beyond the influence of the bishopric and the Aaronic Priesthood quorum advisers is the impact of the home. Help of parents, when enlisted wisely, can frequently make the difference between success and failure. A survey we conducted recently reveals that the influence of the home is a dominant factor in determining missionary service and temple marriage.

I know in my experience of only three wards with a full complement of 48 priests. These wards were presided over by Joseph B. Wirthlin, Alfred B. Smith, and Alvin R. Dyer. Almost without exception, each young man filled a mission and married in the temple. One of the keys to their success was to call to service as Aaronic Priesthood advisers men who were models for the young men to follow. An ideal model is a returned missionary, fresh from his mission and filled with testimony, where a young Aaronic Priesthood holder can say, “That’s the man I want to follow.”

As we dam off that inflow of Aaronic Priesthood streaming into the pool of prospective elders, we will solve more problems than we realize. We will ensure that every young man will more likely than not go on a mission and will marry in the temple. Then there will not be that disproportionate number of worthy young women with few worthy young men to select as an eternal companion. We are not talking about a boy; we’re talking about husbands, fathers, grandfathers, patriarchs to their own families. Let’s put a solid foundation beneath our Aaronic Priesthood young men.

Let us not overlook the adult converts to the Church who receive the Aaronic Priesthood but who are not ordained to the office of elder in a timely fashion. They then join the brethren who remain in that stagnant pool of inactivity. There are those wards and stakes which have rescued vast numbers of fine men who had felt trapped by no outlet in the pond. In traveling the Church, I kept records of those units which had caught the vision of this rescue effort. All of them had similar experiences. They learned that the rescue work is best done one-on-one and at the ward level. The bishop has to be involved, for isn’t he the president of the Aaronic Priesthood as well as the presiding high priest of his ward?

Worthy and well-prepared instructors must be called to help in such a critical effort. Brethren, prayerfully analyze your situation and then call to the colors those whom the Lord has prepared to go forth to serve and to save. “Remember the worth of souls is great in the sight of God.” Ponder the joy that comes to a wife and children when Daddy sees the light, mends his ways, and follows in the footsteps of Jesus Christ our Lord.

An example of true love and inspired teaching was found in the life of the late James Collier, who had through his personal efforts reactivated a large number of brethren in Bountiful, Utah. I was invited by Brother Collier to address those who had now been ordained elders and who, with their wives and families, had been to the Salt Lake Temple to receive those eternal covenants and blessings for which they had so earnestly strived.

At the banquet honoring this achievement, I could see and I could feel the love that Jim had for those whom he had taught and rescued and the love they had for him. Unfortunately, Jim Collier at that time was afflicted with a terminal illness and had to persuade the doctors to allow him to leave the hospital to attend this final night of recognition.

As Jim stood at the pulpit, a large smile came over his face. With emotion he expressed his love to the group. There wasn’t a dry eye to be found. Brother Collier quipped, “Everyone wants to go to the celestial kingdom, but no one wants to die to get there.” Then, lowering his voice, Jim continued, “I’m prepared to go, and I will be there waiting on the other side to greet each of you, my beloved friends.”

Jim returned to the hospital. His funeral service was held just a short time later.

In fulfilling our responsibility to those who bear the Aaronic Priesthood, both the youth and the prospective elders, I urge that we remember that there is no need for us to walk alone. We can look up and reach out for divine help. “The recognition of [a] power higher than man … does not in any sense debase him. If in his faith he ascribes beneficence and high purpose to the power which is superior to himself, he envisions a higher destiny and nobler attributes for his kind and is stimulated and encouraged in the struggle of existence. … He must seek[,] believing, praying, and hoping that he will find. No such sincere, prayerful effort will go unrequited—that is the very constitution of the philosophy of faith.” So taught President Stephen L Richards.

A line from the delightful play The King and I gives us encouragement in our labors. The King of Siam lay dying. With him is Anna, his English tutor, whose son asks her the question, “Was he as good … as he could have been?” Anna answers wistfully, “I don’t think any man has ever been as good … as he could have been—but this one [really] tried.”

The Prophet Joseph declared, “Happiness is the object and design of our existence; and will be the end thereof, if we pursue the path that leads to it; and this path is virtue, uprightness, faithfulness, holiness, and keeping all the commandments of God.”

Let us walk these clearly defined paths. To help us do so we can follow the shortest sermon in the world. It can be found on a common traffic sign. It reads, “Keep Right.”

This advice was found and followed by Joe, who had been asked to get up at six in the morning and drive a crippled child 50 miles to a hospital. He didn’t want to do it, but he didn’t know how to say no. A woman carried the child out to the car and set him next to the driver’s seat, mumbling thanks through her tears. Joe said everything would be all right and drove off quickly.

After a mile or so, the child inquired shyly, “You’re God, aren’t you?”

“I’m afraid not, little fellow,” replied Joe.

“I thought you must be God,” said the child. “I heard Mother praying next to my bed and asking God to help me get to the hospital, so I could get well and play with the other boys. Do you work for God?”

“Sometimes, I guess,” said Joe, “but not regularly. I think I’m going to work for Him a lot more from now on.”

My brethren, will you? Will I? Will we? I pray humbly, yet earnestly, that we will.

In the name of the Lord Jesus Christ, amen.

\newpage
\settalk{Your Celestial Journey}
\setspeaker{Thomas S. Monson}
\settalkdate{April 1999}
\talkheading{Your Celestial Journey}{Thomas S. Monson}

My dear sisters, what a blessing is mine to stand before you this evening and to contemplate that in addition to all assembled here in the Tabernacle, there are many thousands watching and listening to the proceedings by way of satellite transmission. I pray for the help of the Lord.

Henry Wadsworth Longfellow, in a classic poem, described you and your future. He wrote:

\begin{verse}
How beautiful is youth! how bright it gleams\\
With its illusions, aspirations, dreams!\\
Book of Beginnings, Story without End,\\
Each maid a heroine, and each man a friend!
\end{verse}

Precious young women, your mothers, your teachers, and Young Women leaders, may I leave with you some thoughts and suggestions to guide your footsteps through mortality and to the celestial kingdom of our Heavenly Father.

I have carefully chosen four action-packed objectives for your guidance and eternal joy. They are:

First, let us discuss the plea, gaze upward.

Our Heavenly Father has placed an upward reach in every one of us. The words of scripture speak loud and clear: “Look to God and live.” No problem is too small for His attention nor so large that He cannot answer the prayer of faith. Prayer surely is the passport to spiritual power. You can pray with purpose when you realize who you are and what Heavenly Father wants you to become.

You will not find it difficult to approach Him with your sincere prayer as you remember the words of the Apostle Paul, “Know ye not that ye are the temple of God, and that the Spirit of God dwelleth in you?”

If you want to please our Heavenly Father, honor your father and your mother, as He has commanded. They love you dearly. Your joy is their joy, and your sorrow is their sorrow. They want for you the heavenly guidance the Lord provides.

I’ve heard some frustrated parents speaking of a daughter or son as being in the “terrible teens.” I much prefer to describe you as the “terrific teens.”

Life was never meant to be all smiles and happiness. There will come those teaching times to each of you when you will witness the love of your mother, the strength of your father, and the inspiration of God.

I sought permission from Elder Russell M. Nelson to share with you a lesson of sorrow, tempered by knowledge of our Heavenly Father’s plan.

Elder and Sister Nelson have been blessed with nine daughters, followed by one son. They are a happy family, a close-knit family. When the children were younger, they gathered around Mother and Father one evening, and Father proceeded to teach them. He said, “Many couples are being called to serve as missionaries and, in the case of mission presidents, to take their children with them to the areas of their assignment.” Then Dad posed the critical question: “If your mother and I were called to such an assignment, would you be willing to go with us?”

He awaited their responses. One daughter said, “Daddy, they wouldn’t call you, since I’m a cheerleader at high school!”

An older child added, “I couldn’t go. I’m a student at the university.”

The teenage responses continued, until little Emily, with the purity of her soul, answered, “Daddy, if you were called, I would go with you.”

Actually, each of the children would be willing to go, but Emily brought tender tears with her profound yet simple reply.

The years moved along hurriedly. The children married. Grandchildren arrived. Then dreaded cancer struck Emily, and after a valiant and courageous battle she was called home.

Elder Nelson spoke at the funeral services. I’ve never heard a finer or more tender message. He spoke of the plan of salvation and described the promises of God pertaining to the eternal nature of the family. Quietly he said, “Emily has just graduated a little early from mortality.” What a teaching moment!

As the large family walked behind the casket, Elder Nelson carried in his arms two of Emily’s small children. All in attendance became part of truth taught and lessons learned. We were inspired to gaze heavenward.

Second, look inward.

May each ask herself this question, Do I know where I want to go, what I want to be, what I want to do?

The Lord has answered such questions: “Seek ye out of the best books words of wisdom; seek learning, even by study and also by faith.”

The holy scriptures, the guidance of your parents, and the diligent teaching you receive in Primary, Young Women, Sunday School, sacrament meeting, and seminary will fortify you in your determination to be your best self.

Study with purpose, both in church and in school. Write down your goals and what you plan to do to achieve them. Aim high, for you are capable of eternal blessings.

It must not be expected that the road of life spreads itself in an unobstructed view before the person starting her journey. You must anticipate coming upon forks and turnings in the road. But you cannot hope to reach your desired journey’s end if you think aimlessly about whether to go east or west. You must make your decisions purposefully.

As Lewis Carroll tells us in his well-known Alice’s Adventures in Wonderland, Alice was following a path through a forest in Wonderland when it divided in two directions. Standing irresolute, she inquired of the Cheshire cat, which had suddenly appeared in a nearby tree, which path she should take. “Where do you want to go?” asked the cat.

“I don’t know,” said Alice.

“Then,” said the cat, “it really doesn’t matter, does it?”

We know where we want to go. Do we have the resolution—even the faith—to get there?

“Come, … learn of me,” said the Lord. “Come, follow me,” He urged. By responding affirmatively to His gentle invitation, each of you will be ready to move to our next objective and reach out to serve.

The Apostle Paul provided you this wise counsel: “Let no man despise thy youth; but be thou an example of the believers, in word, in conversation, in charity, in spirit, in faith, in purity.”

Young sisters, your opportunities to reach outward and bless the lives of others are limitless. Think, for example, of the privilege you have to attend the holy temple, there to reach out to others who have passed beyond by serving as proxies to provide them the blessings of baptism.

One morning as I walked to the temple, I saw a group of young women who, early that morning, had participated in baptisms for those who had passed beyond. Their hair was wet. Their smiles were radiant. Their hearts were filled with joy. One girl turned back to face the temple and expressed her feelings. “This has been the happiest day of my life,” she said.

There are other opportunities to serve the living. You can do so and can bring untold joy to them. Extended care facilities become a home for those who are ill or aged and require such care. They yearn for the days of their youth. They long for the company of their families and the comforts of their homes.

At a church service I attended in a care center, after the wheelchair-bound residents received the sacrament, a young woman your age played a solo on her violin. The elderly sisters were so appreciative. They declared aloud their gratitude with comments such as “Beautiful,” “Wonderful,” “I love you.” Such distractions did not deter the violinist; rather, they enabled her to reach new heights in her performance.

That day she said to me: “I have never played better in my life. Something seemed to lift me beyond myself and my own abilities. I felt the inspiration of my Heavenly Father’s love.”

I reminded her, “When you are in the service of your fellow beings you are only in the service of your God.”

She nodded her acknowledgement, carefully placed her violin in its case, and, with tears of joy coursing down her cheeks, returned to her seat.

May we remember to reach outward.

Finally, press forward. Avoid the tendency to postpone a prompting or an opportunity to grow and to serve. Procrastination is truly a thief of time. Meet the daily challenges of your lives. How long has it been since you looked into the eyes of your mother and, holding nothing back, spoke those welcome words, “Mother, I truly love you”? How about Father, who daily toils to provide for you? Fathers appreciate hearing those same precious words from the lips of a child, “I love you.”

It is too easy to take parents for granted and to fail to realize just how much they mean to you and you to them. An illustration of this occurred in a classroom. One of the questions, after a study of magnets at Olympus Junior High, was, “What begins with M and picks things up?” More than a third of the students answered, “Mother.”

Move against temporary trials or stoppages which impede your progress.

A blessing you can qualify to receive is your patriarchal blessing. Your parents and your bishop will know when the time is right for you to receive it. A patriarchal blessing contains chapters from your life’s book of possibilities. To you it will be as a lighthouse on a hill, warning of dangers, and directing you to the tranquility of safe harbors. It is a prophetic utterance from the lips of one called and ordained to provide you such a blessing.

May I take the opportunity to express, from each of you young women, a heartfelt “thank you” to your parents, to your teachers, to your leaders. They are role models for you. They know there will be disappointments, days that are downers, and personal frustrations in your lives. They will show you the way to rise above such experiences and continue on that high road of life which leads upward and onward to celestial glory. Remember that once you have experienced excellence, you will never again be content with mediocrity.

Some years ago a lovely young woman, Jami Palmer, then 12 years of age, was wheeled into my office by her parents. She had been diagnosed with cancer. Surgery would be required. The treatments would be many and the time of recovery long. It was a solemn moment as we visited. Father requested me to join him in blessing his crestfallen daughter who had just had her dreams, her hopes, her plans placed on hold. All of us were weeping. The priesthood blessing was provided.

I have maintained contact with Jami and her family. The years have flown by. She has rendered unlimited service to others through being a spokesperson for the Make-A-Wish Foundation, which blesses youth afflicted with life-threatening diseases. Jami has grown into a beautiful young woman. She is now a student at Brigham Young University. She is healthy. She has been through the refiner’s fire and has had her life prolonged. She gives thanks to all who aided her through these difficult years and especially to her Heavenly Father for her very life today.

A turning point in Jami’s life came early in her treatment for cancer. She and the youth in her ward had planned a hike to Timpanogos Cave. You who have made that hike know the way is steep, and it seems to take forever to reach the cave. Sadly Jami said to her friends, “I won’t be able to make the hike with you.”

“Why not?” they asked.

Jami replied, “I can’t walk.”

There was a silent moment, and then one replied, “Jami, if you can’t walk, then we’ll carry you.” And they did—up and back!

Young women, will you gaze upward, look inward, reach outward, and press forward? As you do, great shall be your reward and eternal shall be your glory.

I bear to you, my beloved sisters, my witness that Heavenly Father lives, that Jesus is the Christ, and that we are led today by a prophet for our time, even President Gordon B. Hinckley. In the name of Jesus Christ, amen.

\newpage
\settalk{Becoming Our Best Selves}
\setspeaker{Thomas S. Monson}
\settalkdate{October 1999}
\talkheading{Becoming Our Best Selves}{Thomas S. Monson}

During a time long past, and in a place far away, our Lord and Savior, Jesus Christ, taught the multitudes and His disciples “the way, the truth, and the life.” He provided counsel with His holy words. He lived an example for us with His exemplary life. On occasion the Lord would ask another this question: “What manner of persons ought ye to be?”

During His ministry on the American continent, He added significant words when He answered the same question: “What manner of men ought ye to be? Verily I say unto you, even as I am.”

In His earthly ministry, the Master outlined how we should live, how we should teach, how we should serve, and what we should do so that we could become our best selves.

One such lesson comes from the book of John in the Holy Bible: “Philip findeth Nathanael, and saith unto him, We have found him, of whom Moses in the law, and the prophets, did write, Jesus of Nazareth, the son of Joseph.

“And Nathanael said unto him, Can there any good thing come out of Nazareth? Philip saith unto him, Come and see.

“Jesus saw Nathanael coming to him, and saith of him, Behold an Israelite indeed, in whom is no guile!”

In our mortal journey, the advice of the Apostle Paul provides heavenly guidance: “Whatsoever things are true, whatsoever things are honest, whatsoever things are just, whatsoever things are pure, whatsoever things are lovely, whatsoever things are of good report; if there be any virtue, and if there be any praise, think on these things.” Then came the concluding charge: “Those things, which ye have both learned, and received, and heard, and seen in me, do: and the God of peace shall be with you.”

In the search for our best selves, several questions will guide our thinking: “Am I what I want to be? Am I closer to the Savior today than I was yesterday? Will I be closer yet tomorrow? Do I have the courage to change for the better?”

It is time to choose an oft-forgotten path, the path we might call “The Family Way,” so that our children and grandchildren might indeed grow to their full potential. There is a national—even an international—tide running. It carries the unspoken message, “Return to your roots, to your families, to lessons learned, to lives lived, to examples shown, even family values.” Often it is just a matter of coming home—coming home to attics not recently examined, to diaries seldom read, to photo albums almost forgotten.

The Scottish poet James Barrie wrote, “God gave us memories, that we might have June roses in the December of our lives.” What memories do we have of Mother? Father? Grandparents? Family? Friends?

What lessons have we learned from our fathers? Years ago a father asked Elder ElRay L. Christiansen what name he could suggest for his newly acquired boat. Brother Christiansen suggested, “Why not call it The Sabbath Breaker?” I’m confident the would-be sailor pondered whether his pride and joy would be a Sabbath breaker or a Sabbath keeper. Whatever his decision, it no doubt left a lasting impression upon his children.

Yet another father taught a son a never-to-be-forgotten lesson in obedience and, by example, to honor the Sabbath day. I learned of this at the funeral service of a noble General Authority, H. Verlan Andersen. A tribute was paid to him by one of his sons. It has application wherever we are and whatever we are doing. It is the example of personal experience.

The son of Elder Andersen related that years earlier he had a special school date on a Saturday night. He borrowed from his father the family car. As he obtained the car keys and was heading for the door, his father said, “The car will need more gasoline before tomorrow. Be sure to fill the tank before coming home.”

Elder Andersen’s son related that the evening activity was wonderful. Friends met, refreshments were served, and all had a good time. In his exuberance, however, he failed to follow his father’s instruction to add fuel to the car’s tank before returning home.

Sunday morning dawned. Elder Andersen discovered the gas gauge showed empty. The son saw his father walk back into the house and put the car keys on the table. In the Andersen home, the Sabbath day was a day for worship and thanksgiving, and not for purchases.

As the funeral message continued, Elder Andersen’s son declared, “I saw my father put on his coat, bid us good-bye, and then walk the long distance to the chapel, that he might attend an early meeting.” Duty called. Truth was not held slave to expedience.

In concluding his funeral message, he said: “No son was ever taught more effectively by his father than I was on that occasion. My father not only knew the truth—he lived it.”

It is in the home that we form our attitudes, our deeply held beliefs. It is in the home that hope is fostered or destroyed.

Our homes are to be more than sanctuaries; they should also be places where God’s Spirit can dwell, where the storm stops at the door, where love reigns and peace dwells.

Not long ago a young mother wrote to me: “Sometimes I wonder if I make a difference in my children’s lives. Especially as a single mother working two jobs to make ends meet, I sometimes come home to confusion, but I never give up hope.

“My children and I were watching a television broadcast of general conference, and you were speaking about prayer. My son made the statement, ‘Mother, you’ve already taught us that.’ I said, ‘What do you mean?’ And he replied: ‘Well, you’ve taught us to pray and showed us how, but the other night I came to your room to ask something and found you on your knees praying to Heavenly Father. If He’s important to you, He’ll be important to me.’” The letter concluded, “I guess you never know what kind of influence you’ll be until a child observes you doing yourself what you have tried to teach him to do.” What a magnificent lesson a child learned from his mother.

As a boy I made a startling discovery in Sunday School one Mother’s Day which has remained with me all through the years. Melvin, a sightless brother in the ward, a talented vocalist, would stand and face the congregation as though he were seeing one and all. He would then sing “That Wonderful Mother of Mine.” The bright, glowing embers of memory penetrated human hearts. Men reached for their handkerchiefs; women’s eyes brimmed with tears.

We deacons would go among the congregation carrying a small geranium in a clay pot for presentation to each mother. Some of the mothers were young, some were middle-aged, some were barely hanging on to life in their old age. I became aware that the eyes of each mother were kind eyes. The words of each mother were “Thank you.” I felt the spirit of the statement “When someone gives another person a flower, the fragrance of the flower lingers on the hands of the giver.” I have not forgotten the lesson learned, nor shall I ever forget it.

Some mothers, some fathers, some children, some families are called upon to bear a heavy burden here in mortality. Such a family was the Borgstrom family in northern Utah. The time was World War II. Fierce battles raged in various parts of the world.

Tragically, the Borgstroms lost four of their five sons who were serving in the armed forces. Within a six-month period, all four sons gave their lives—each in a different part of the world.

Following the war, the bodies of the four Borgstrom brothers were brought home to Tremonton, and an appropriate service was conducted, filling the Garland Utah Tabernacle. General Mark Clark attended the service. He later spoke with tenderness these words: “I flew to Garland the morning of June 26. Met with the family, including among others the mother, father, and two remaining sons, … one a lad in his teens. I had never met a more stoic family group.

“As the four flag-draped coffins were lined up in front of us in the church, and as I sat by these brave parents, I was deeply impressed by their understanding, by their faith, and their pride in these magnificent sons who had made the supreme sacrifice for principles which had been instilled in them by noble parents since childhood.

“During the luncheon period, Mrs. Borgstrom turned to me and said in a low voice, ‘Are you going to take my young one?’ I answered in a whisper that as long as I remained in command of the army on the West Coast, if her boy were called I would do my best to have him assigned to duty at home.

“In the middle of this whispered conversation with the mother, the father suddenly leaned forward and said to Mrs. Borgstrom: ‘Mother, I have overheard your conversation with the general about our youngest. We know that if and when his country needs him, he will go.’

“I could hardly contain my emotions. Here were parents with four sons lying dead from wounds received in battle and yet were ready to make the last sacrifice if their country required it.”

It is the gospel of the Lord Jesus Christ that touched home and heart that ever-to-be-remembered day.

The years have come and the years have gone, but the need for a testimony of the gospel continues paramount. As we move toward the future, we must not neglect the lessons of the past. Our Heavenly Father gave His Son. The Son of God gave His life. We are asked by Them to give our lives, as it were, in Their divine service. Will you? Will I? Will we? There are lessons to be taught, there are kind deeds to be done, there are souls to be saved.

Let us remember the counsel of King Benjamin: “When ye are in the service of your fellow beings ye are only in the service of your God.” Reach out to rescue those who need your help. Lift such to the higher road and the better way. As we sing in Primary, “Lead me, guide me, walk beside me, help me find the way. Teach me all that I must do to live with him someday.”

Real faith is not restricted to childhood, but rather applies to all. We learn from the Proverbs: “Trust in the Lord with all thine heart; and lean not unto thine own understanding.

“In all thy ways acknowledge him, and he shall direct thy paths.” When we do, we will come to realize that we have been on His holy errand, that His divine purposes have been fulfilled, and that we have shared in that fulfillment.

May I illustrate this truth with a personal experience. Many years ago, while serving as a bishop, I felt impressed to call upon Augusta Schneider, a widow from the Alsace-Lorraine area of Europe who spoke very little English, although she was fluent in French and German. For years after that first impression I would visit with her at Christmastime. On one occasion Augusta said, “Bishop, I have something of great value to me which I would like to present to you.” She then went to a special place in her modest apartment and retrieved the gift. It was a beautiful piece of felt, perhaps six by eight inches in size, to which she had pinned the medals her husband had been presented for his service as a member of the French forces in World War I. She said, “I would like you to have this personal treasure which is so close to my heart.” I protested politely and suggested there must be some member of her extended family to whom the gift should be given. “No,” she replied firmly, “the gift is yours, for you have the soul of a Frenchman.”

Shortly after presenting this special gift to me, Augusta departed mortality and went home to that God who gave her life. Occasionally I would wonder concerning her declaration that I had “the soul of a Frenchman.” I didn’t have the slightest idea what that meant. I still don’t.

Many years later I had the privilege to accompany President Ezra Taft Benson to the dedication of the Frankfurt Germany Temple, which temple would serve German-, French-, and Dutch-speaking members. In packing for the trip, I felt impressed to take along the gift of medals, without any thought concerning what I would do with them. I’d had them a number of years.

In a French-speaking dedication session, the temple was filled. The singing and messages presented were beautiful. Gratitude for God’s blessings penetrated each heart. I saw from my conducting notes that the session included members from the Alsace-Lorraine area.

During my remarks I observed that the organist had the name of Schneider. I therefore related the account of my association with Augusta Schneider, then stepped to the organ and presented the organist with the medals, along with the charge that since his name was Schneider, he had a responsibility to pursue the Schneider name in his genealogical activities. The Spirit of the Lord confirmed in our hearts that this was a special session. Brother Schneider had a difficult time preparing to play the closing number of the dedicatory service, so moved was he by the Spirit which we felt there in the temple.

I knew that the treasured gift—even the widow’s mite, for it was all Augusta Schneider had—was placed in the hand of one who would ensure that many with the souls of Frenchmen would now receive the blessings the holy temples provide, both to the living and for those who have passed beyond mortality.

I testify that with God, all things are possible. He is our Heavenly Father; His Son is our Redeemer. As we strive to learn His truths and then to live them, our lives and the lives of others will be abundantly blessed.

I declare in all soberness that Gordon B. Hinckley is a true prophet for our time and is guided in the great work going forward under his leadership.

May we ever remember that obedience to God’s commandments brings forth the blessings promised.

May each of us qualify to receive them, I pray, in the name of Jesus Christ, amen.

\newpage
\settalk{Priesthood Power}
\setspeaker{Thomas S. Monson}
\settalkdate{October 1999}
\talkheading{Priesthood Power}{Thomas S. Monson}

Brethren of the priesthood, assembled here and worldwide, I am humbled by the responsibility which is mine to address a few remarks to you. I pray for the Spirit of the Lord to attend me as I do so.

Some of you are deacons; others are teachers or priests—all offices in the Aaronic Priesthood. Many of you are elders, seventies, or high priests. Much is expected of each of us.

In a proclamation of the First Presidency and the Council of the Twelve Apostles issued on April 6, 1980, this declaration of testimony and truth was set forth:

“We solemnly affirm that The Church of Jesus Christ of Latter-day Saints is in fact a restoration of the Church established by the Son of God, when in mortality he organized his work upon the earth; that it carries his sacred name, even the name of Jesus Christ; that it is built upon a foundation of Apostles and prophets, he being the chief cornerstone; that its priesthood, in both the Aaronic and Melchizedek orders, was restored under the hands of those who held it anciently: John the Baptist, in the case of the Aaronic; and Peter, James, and John in the case of the Melchizedek.”

On October 6, 1889, President George Q. Cannon expressed this plea:

“I want to see the power of the Priesthood strengthened. … I want to see this strength and power diffused through the entire body of the Priesthood, reaching from the head down to the least and most humble deacon in the Church. Every man should seek for and enjoy the revelations of God, the light of heaven shining in his soul and giving unto him knowledge concerning his duties, concerning that portion of the work of God that devolves upon him in his Priesthood.”

The Lord Himself summed up our responsibility when He, in the revelation on the priesthood, urged, “Wherefore, now let every man learn his duty, and to act in the office in which he is appointed, in all diligence.”

Brethren of the Aaronic Priesthood, whether deacon, teacher, or priest, learn your duty. Brethren of the Melchizedek Priesthood, learn your duty.

Some years ago, as our youngest son, Clark, was approaching his 12th birthday, he and I were leaving the Church Administration Building when President Harold B. Lee greeted us. I mentioned to President Lee that Clark would soon be 12, whereupon President Lee asked him, “What happens to you, Clark, when you turn 12?” This was one of those times when a father prays that a son will be inspired to give a proper response. Without hesitation Clark said to President Lee, “I will be ordained a deacon.”

The answer was the one President Lee had sought. He then counseled our son, “Remember, it is a great blessing to hold the priesthood.”

I hope with all my heart and soul that every young man who receives the priesthood will honor that priesthood and be true to the trust which is conveyed when it is conferred.

Forty-four years ago I heard William J. Critchlow Jr., then president of the South Ogden Stake, speak to the brethren in the general priesthood session of conference, and retell a story concerning trust, honor, and duty. May I share the story with you. Its simple lesson applies to us today, as it did then.

“Rupert stood by the side of the road watching an unusual number of people hurry past. At length he recognized a friend. ‘Where are all of you going in such a hurry?’ he asked.

“The friend paused. ‘Haven’t you heard?’ he said.

“‘I’ve heard nothing,’ Rupert answered.

“‘Well,’ continued [the] friend, ‘the King has lost his royal emerald. Yesterday he attended a wedding of the nobility and wore the emerald on the slender golden chain around his neck. In some way the emerald became loosened from the chain. Everyone is searching, for the King has offered a reward … to the one who finds it. Come, we must hurry.’

“‘But I cannot go without asking Grandmother,’ faltered Rupert.

“‘Then I cannot wait. I want to find the emerald,’ replied his friend.

“Rupert hurried back to the cabin at the edge of the woods to seek his grandmother’s permission. ‘If I could find it, we could leave this hut with its dampness and buy a piece of land up on the hillside,’ he pleaded with Grandmother.

“But his grandmother shook her head. ‘What would the sheep do?’ she asked. ‘Already they are restless in the pen, waiting to be taken to the pasture—and please do not forget to take them to water when the sun shines high in the heavens.’

“Sorrowfully, Rupert took the sheep to the pasture, and at noon he led them to the brook in the woods. There he sat on a large stone by the stream. ‘If I could only have had a chance to look for the King’s emerald,’ he thought. Turning his head to gaze down at the sandy bottom of the brook, suddenly he stared into the water. What was it? It could not be! He leaped into the water, and his gripping fingers held something that was green, with a slender bit of gold chain. ‘The King’s emerald!’ he shouted. ‘It must have been flung from the chain when the King [astride his horse, galloped across the bridge spanning the stream, and the current carried] it here.’

“With shining eyes Rupert ran to his grandmother’s hut to tell her of his great find. ‘Bless you, my boy,’ she said, ‘but you never would have found it if you had not been doing your duty, herding the sheep.’ And Rupert knew that this was the truth.”

The lesson to be learned from this story is found in the familiar couplet: “Do your duty; that is best. Leave unto the Lord the rest.”

Should there be anyone who feels he is too weak to change the onward and downward course of his life, or should there be those who fail to resolve to do better because of that greatest of fears—the fear of failure—there is no more comforting assurance to be had than these words of the Lord: “My grace is sufficient for all men that humble themselves before me; for if they humble themselves before me, and have faith in me, then will I make weak things become strong unto them.”

Miracles are everywhere to be found when priesthood callings are magnified. When faith replaces doubt, when selfless service eliminates selfish striving, the power of God brings to pass His purposes.

The priesthood is not really so much a gift as it is a commission to serve, a privilege to lift, and an opportunity to bless the lives of others.

Brethren, let us who have responsibility with the Aaronic Priesthood young men not only provide them opportunities to learn but also set before them examples worthy of emulation.

For those of us who hold the Melchizedek Priesthood, our privilege to magnify our callings is ever present. We are shepherds watching over Israel. The hungry sheep look up, ready to be fed the bread of life. Are we prepared to feed the flock of God? It is imperative that we recognize the worth of a human soul, that we never give up on one of His precious sons.

May I share with you a letter from a young man which reflects the spirit of love and which helped to make firm a testimony of the gospel:

“Dear President Monson:

“Thank you for speaking to us at the National Scouting Jamboree held at Fort A. P. Hill, Virginia. On the tour that we took we saw a lot of famous places like Niagara Falls, the Statue of Liberty, the Liberty Bell, and many other places. The one I enjoyed the most was the Sacred Grove. Our parents had written us all letters to read by ourselves while in the grove. After I had finished the letter my parents had written to me, I knelt in prayer. I asked if the Church was really true and if Joseph Smith really did see a vision and is a true prophet of God, and also if President Hinckley is a true prophet of God. Right after I was done praying I felt this feeling of the Spirit that these things were indeed true. I had prayed before about the same things but never received such a powerful answer. There was no way that I could deny that this Church is true or that President Hinckley is a prophet of God.

“I feel so blessed to be a member of this Church. Thanks again for attending the Jamboree.

“Sincerely,

“Chad D. Olson

“P.S. We gave our tour guide and our bus driver a copy of the Book of Mormon with our testimonies in it. They are the greatest! I want to be a missionary.”

Like Joseph Smith, this young man had retired to a sacred grove and prayed for answers to questions phrased by his inquiring mind. Once more a prayer was answered and a confirmation of the truth was gained.

There are many less-active members who wander in the wilderness of wonder or who struggle in the swamp of sin. One such member wrote to me:

“I’m afraid to be alone. The gospel has never left my heart, even though it has left my life. I ask for your prayers. I would be happy to even eat the crumbs that fall from the table of the lowliest member of the Church, because he has more than I have now. I used to think position and responsibility were important in the Church, but now I know I was wrong all the time. It was membership, priesthood power, fatherhood, and service. I know where the Church is, but sometimes I think I need someone else to show me the way, encourage me, take away my fear, and bear testimony to me. I thought the Church was lost, when really it was only me.”

The call of duty can come quietly as we who hold the priesthood respond to the assignments we receive. President George Albert Smith, that modest yet effective leader, declared, “It is your duty first of all to learn what the Lord wants and then by the power and strength of [your] holy Priesthood to magnify your calling in the presence of your fellows in such a way that the people will be glad to follow you.”

What does it mean to magnify a calling? It means to build it up in dignity and importance, to make it honorable and commendable in the eyes of all men, to enlarge and strengthen it to let the light of heaven shine through it to the view of other men. And how does one magnify a calling? Simply by performing the service that pertains to it. An elder magnifies the ordained calling of an elder by learning what his duties as an elder are and then by doing them. As with an elder, so with a deacon, a teacher, a priest, a bishop, and each who holds office in the priesthood.

As we remember, Paul, who was known as Saul, was on his way to Damascus to persecute the Christians there. As he journeyed close to the city of Damascus, a bright light shone round about him, and he fell to the earth, stunned, and he heard a voice saying, “Saul, Saul, why persecutest thou me?” And Saul asked, “Who art thou, Lord?” And the voice said, “I am Jesus.”

A penitent Saul asked, “Lord, what wilt thou have me to do?” With the Lord’s answer, Saul the persecutor became Paul the proselytizer and began his great missionary endeavors.

Brethren, it is in doing—not just dreaming—that lives are blessed, others are guided, and souls are saved. “Be ye doers of the word, and not hearers only, deceiving your own selves,” added James.

May all of us assembled tonight in this priesthood meeting make a renewed effort to qualify for the Lord’s guidance in our lives. There are many out there who plead and pray for help. There are those who are discouraged, those who are beset by poor health and challenges of life which leave them in despair.

I’ve always believed in the truth of the words, “God’s sweetest blessings always go by hands that serve him here below.” Let us have ready hands, clean hands, and willing hands, that we may participate in providing what our Heavenly Father would have others receive from Him.

I conclude with an example in my own life. Once I had a treasured friend who seemed to experience more of life’s troubles and frustrations than he could bear. Finally he lay in the hospital, terminally ill. I knew not that he was there.

Sister Monson and I had gone to that same hospital to visit another person who was very ill. As we exited the hospital and proceeded to where our car was parked, I felt the distinct impression to return and make inquiry concerning whether Hyrum Adams might be a patient there. Long years before, I had learned never, never, to postpone a prompting from the Lord. It was late, but a check with the desk clerk confirmed that indeed Hyrum was a patient.

We proceeded to his room, knocked on the door, and opened it. We were not prepared for the sight that awaited us. Balloon bouquets were everywhere. Prominently displayed on the wall was a poster with the words “Happy Birthday” written on it. Hyrum was sitting up in his hospital bed, his family members by his side. When he saw us, he said, “Why, Brother Monson, how in the world did you know that this is my birthday?” I smiled but I left the question unanswered.

Those in the room who held the Melchizedek Priesthood surrounded this, their father and my friend, and a priesthood blessing was given.

After tears were shed, smiles of gratitude exchanged, and tender hugs received and given, I leaned over to Hyrum and spoke softly to him: “Hyrum, remember the words of the Lord, for they will sustain you. He promised, ‘I will not leave you comfortless: I will come to you.’”

May each of us ever be on the Lord’s errand and thereby be entitled to the Lord’s help, I pray humbly. In the name of Jesus Christ, amen.

\newpage
\settalk{Your Eternal Home}
\setspeaker{Thomas S. Monson}
\settalkdate{April 2000}
\talkheading{Your Eternal Home}{Thomas S. Monson}

One day during the personal ministry of our Savior, He took Peter, James, and John “up into an high mountain …

“And was transfigured before them: and his face did shine as the sun, and his raiment was white as the light.

“And, behold, there appeared unto them Moses and Elias talking with him.

“Then answered Peter, and said unto Jesus, Lord, it is good for us to be here.”

Today, on this historic occasion, we assemble in this magnificent Conference Center and in the overflow facilities on Temple Square and throughout the world.

Tears moisten our eyes and gratitude fills our hearts as we echo the title of a beautiful hymn, “Thanks Be to God!” The erection of this edifice has long been in the planning stage. We have needed a much larger building to accommodate those who attend conference and other activities throughout the year. Workmen with finely honed skills have labored with their hearts and hands to provide a structure worthy of His divine approbation, “Well done, thou good and faithful servant.”

When Jesus ministered among men at a time long ago and a place far away, He often spoke in parables, in language the people best understood. Oftentimes He referred to home building in relationship to the lives of those who listened. Wasn’t He frequently known as “the carpenter’s son”? He declared, “Every … house divided against itself shall not stand.” Later He cautioned, “Behold, mine house is a house of order, saith the Lord God, and not a house of confusion.”

In a revelation given through the Prophet Joseph Smith at Kirtland, Ohio, December 27, 1832, the Master counseled, “Organize yourselves; prepare every needful thing; and establish a house, even a house of prayer, a house of fasting, a house of faith, a house of learning, a house of glory, a house of order, a house of God.”

Where could any of us locate a more suitable blueprint whereby he or she could wisely and properly build a house to personally occupy throughout eternity?

In a very real sense, we are builders of eternal houses. We are apprentices to the trade—not skilled craftsmen. We need divine help if we are to build successfully. The words of instruction provided by the Apostle Paul give the assurance we need: “Know ye not that ye are the temple of God, and that the Spirit of God dwelleth in you?”

When we remember that each of us is literally a spirit son or daughter of God, we will not find it difficult to approach our Heavenly Father in prayer. He appreciates the value of this raw material which we call life. “Remember the worth of souls is great in the sight of God.” His pronouncement finds lodgment in our souls and inspires purpose in our lives.

There is a Teacher who will guide our efforts if we will but place our faith in Him—even the Lord Jesus Christ. He invites us: “Come unto me, all ye that labour and are heavy laden, and I will give you rest.

“Take my yoke upon you, and learn of me; for I am meek and lowly in heart: and ye shall find rest unto your souls.

“For my yoke is easy, and my burden is light.”

It was said of Jesus that He “increased in wisdom and stature, and in favour with God and man.” Do we have the determination to do likewise? One line of holy writ contains a tribute to our Lord and Savior, of whom it was said, “[He] went about doing good.”

Paul, in his epistle to his beloved Timothy, outlined a way whereby we could become our better selves and, at the same time, provide assistance to others who silently ponder and then audibly ask the question, “How can I [find my way], except some man should guide me?”

The answer, given by Paul to Timothy, provides an inspired charge to each of us. Let us take heed of his wise counsel: “Be thou an example of the believers, in word, in conversation, in charity, in spirit, in faith, in purity.”

Let us examine this solemn instruction which, in a very real sense, is given to us.

First, be an example in word. “Let your words tend to edifying one another,” said the Lord.

Do we remember the counsel of a favorite Sunday School hymn?

\begin{verse}
Oh, the kind words we give shall in memory live\\
And sunshine forever impart.\\
Let us oft speak kind words to each other;\\
Kind words are sweet tones of the heart.
\end{verse}

Consider the observation of Mary Boyson Wall, who celebrated her 105th birthday a few years ago. She married Don Harvey Wall in the Salt Lake Temple in 1913. They celebrated their 81st wedding anniversary shortly before Don died at age 103. In a Church News article “she attributed longevity in life and in their marriage to speaking kind words. She said, ‘I think that helped us through because we tried to help each other and not say unkind words to each other.’”

Second, be an example in conversation. In a general conference in October 1987, President Gordon B. Hinckley declared: “Foul talk defiles the man who speaks it. If you have that habit, how do you break it? You begin by making a decision to change. The next time you are prone to use words you know to be wrong, simply stop. Keep quiet or say what you have to say in a different way.”

François de la Rochefoucauld observed, “One of the reasons why so few people are to be found who seem sensible and pleasant in conversation is that almost everybody is thinking about what he wants to say himself, rather than about answering clearly what is said to him.”

Third, be an example in charity.

From Corinthians comes the beautiful truth, “Charity never faileth.”

Satisfying to the soul is the ready response the Church has made to disasters of nature, such as in Mozambique, Madagascar, Venezuela, and many other locations. Frequently we have arrived first on the scene following such disasters, and with the most help. There are other organizations which likewise respond in a generous fashion.

What is charity? Moroni, in writing a few of the words of his father, Mormon, recorded, “Charity is the pure love of Christ, and it endureth forever.”

One who exemplified charity in his life was President George Albert Smith. Immediately following World War II, the Church had a drive to amass warm clothing to ship to suffering Saints in Europe. Elder Harold B. Lee and Elder Marion G. Romney took President George Albert Smith to Welfare Square in Salt Lake City to view the results. They were impressed by the generous response of the membership of the Church. They watched President Smith observing the workers as they packaged this great volume of donated clothing and shoes. They saw tears running down his face. After a few moments, President George Albert Smith removed his own new overcoat and said, “Please ship this also.”

The Brethren said to him, “No, President, no; don’t send that; it’s cold and you need your coat.”

But President Smith would not take it back; and so his coat, with all the others, was sent to Europe, where the nights were long and dark and food and clothing were scarce. Then the shipments arrived. Joy and thanksgiving were expressed aloud, as well as in secret prayer.

Fourth, be an example in spirit. The Psalmist wrote, “Create in me a clean heart, O God; and renew a right spirit within me.”

As a 17-year-old, I enlisted in the United States Navy and attended boot camp in San Diego, California. For the first three weeks, one felt as though the navy were trying to kill rather than train him on how to stay alive.

I shall ever remember the first Sunday at San Diego. The chief petty officer said to us, “Today everybody goes to church.” We then lined up in formation on the drill ground. The petty officer shouted, “All of you who are Catholics—you meet in Camp Decatur. Forward, march! And don’t come back until three!” A large number marched out. He then said, “All of you who are of the Jewish faith—you meet in Camp Henry. Forward, march! And don’t come back until three!” A smaller contingent moved out. Then he said, “The rest of you Protestants meet in the theaters in Camp Farragut. Forward, march! And don’t come back until three o’clock!”

There flashed through my mind the thought, “Monson, you’re not Catholic. You’re not Jewish. You’re not a Protestant.” I elected to stand fast. It seemed as though hundreds of men marched by me. Then I heard the sweetest words which the petty officer ever uttered in my presence. He said, “And what do you men call yourselves?” He used the plural—men. This was the first time I knew that anyone else was standing behind me on that drill ground. In unison we said, “We’re Mormons.” He scratched his head, an expression of puzzlement on his face, and said, “Well, go and find somewhere to meet—and don’t come back until three o’clock.” We marched away. One could almost count cadence to the rhyme learned in Primary:

\begin{verse}
Dare to be a Mormon;\\
Dare to stand alone.\\
Dare to have a purpose firm,\\
And dare to make it known.
\end{verse}

Fifth, be an example in faith.

President Stephen L Richards, speaking of faith, declared: “The recognition of power higher than man himself does not in any sense debase him. If in his faith he ascribes beneficence and high purpose to the power which is superior to himself, he envisions a higher destiny and nobler attributes for his kind and is stimulated and encouraged in the struggle of existence. … He must seek believing, praying, and hoping that he will find. No such sincere, prayerful effort will go unrequited—that is the very constitution of the philosophy of faith.” Divine favor will attend those who humbly seek it.

Minnie Louise Haskins set forth this principle in a lovely poem: “I said to the man who stood at the gate of the year: ‘Give me a light that I may tread safely into the unknown.’ And he replied: ‘Go out into the darkness and put your hand into the hand of God. That shall be to you better than light and safer than a known way.’”

Finally, be an example in purity.

“Who shall ascend into the hill of the Lord? or who shall stand in his holy place?

“He that hath clean hands, and a pure heart; who hath not lifted up his soul unto vanity, nor sworn deceitfully.

“He shall receive the blessing from the Lord, and righteousness from the God of his salvation.”

As President David O. McKay observed: “The safety of our nation depends upon the purity and strength of the home; and I thank God for the teachings of the … Church in relation to home building, and the impression that kind parents have made, that the home must be the most sacred place in the world. Our people are home-builders, and they are taught everywhere, from childhood to old age, that the home should be kept pure and safe from the evils of the world.”

Many years ago I attended a stake conference in Star Valley, Wyoming, where the stake presidency was reorganized. The stake president who was being released, E. Francis Winters, had served faithfully for the lengthy term of 23 years. Though modest by nature and circumstance, he had been a perpetual pillar of strength to everyone in the valley. On the day of the stake conference, the building was filled to overflowing. Each heart seemed to be saying a silent thank-you to this noble leader who had given so unselfishly of his life for the benefit of others.

As I stood to speak, I was prompted to do something I had not done before, nor have I done so since. I stated how long Francis Winters had presided in the stake; then I asked all whom he had blessed or confirmed as children to stand and remain standing. Then I asked all those persons whom President Winters had ordained, set apart, personally counseled, or blessed to please stand. The outcome was electrifying. Every person in the audience rose to his or her feet. Tears flowed freely—tears which communicated better than could words the gratitude of tender hearts. I turned to President and Sister Winters and said, “We are witnesses today of the prompting of the Spirit. This vast throng reflects not only individual feelings but also the gratitude of God for a life well lived.” No person who was in the congregation that day will forget how he or she felt when we witnessed the language of the Spirit of the Lord.

Here, in Francis Winters, was “an example of the believers, in word, in conversation, in charity, in spirit, in faith, in purity.”

\begin{verse}
True to the faith that our parents have cherished,\\
True to the truth for which martyrs have perished,\\
To God’s command,\\
Soul, heart, and hand,\\
Faithful and true we will ever stand.
\end{verse}

That each of us may do so is my humble prayer, in the name of Jesus Christ, amen.

\newpage
\settalk{Your Eternal Voyage}
\setspeaker{Thomas S. Monson}
\settalkdate{April 2000}
\talkheading{Your Eternal Voyage}{Thomas S. Monson}

One of my most vivid memories was attending priesthood meeting as a newly ordained deacon and singing the opening hymn, “Come, all ye sons of God who have received the priesthood.” Tonight, to the capacity audience assembled in this magnificent Conference Center and in chapels worldwide, I echo the spirit of that special hymn and say to you, “Come, all ye sons of God who have received the priesthood, let us consider our callings, let us reflect on our responsibilities, let us determine our duty, and let us follow Jesus Christ our Lord.”

While we may differ in age, in custom, or in nationality, we are united as one in our priesthood callings.

As bearers of the priesthood, we have been placed on earth in troubled times. We live in a complex world, with currents of conflict everywhere to be found. Political machinations ruin the stability of nations, despots grasp for power, and segments of our society seem forever downtrodden, deprived of opportunity, and left with a feeling of failure.

We who have been ordained to the priesthood of God can make a difference. When we qualify for the help of the Lord, we can build boys. We can mend men. We can accomplish miracles in His holy service. Our opportunities are without limit.

Though the task seems large, we are strengthened by this truth: “The greatest force in this world today is the power of God as it works through man.” If we are on the Lord’s errand, we are entitled to the Lord’s help. That divine help, however, is predicated upon our worthiness. To sail safely the seas of mortality, to perform a human rescue mission, we need the guidance of that eternal mariner—even the great Jehovah. We reach out, we reach up, to obtain heavenly help.

Are our reaching hands clean? Are our yearning hearts pure? Looking backward in time through the pages of history, a lesson on worthiness is gleaned from the words of the dying King Darius. “Darius, … through the proper rites had been recognized as legitimate King of Egypt; his rival Alexander [the Great] had been declared … legitimate Son of Amon—he too was Pharaoh. … Alexander[, finding] the defeated Darius on the point of death … , laid his hands upon his head to heal him, commanding him to arise and resume his kingly power, … concluding … : ‘I swear unto thee, Darius, by all the gods that I do these things truly and without faking. …’ [Darius] replied with a gentle rebuke: ‘Alexander my boy … do you think you can touch heaven with those hands of yours?’”

An inspiring lesson is learned from a “Viewpoint” article which appeared some time ago in the Church News. May I quote:

“To some it may seem strange to see ships of many nations loading and unloading cargo along the docks at Portland, Ore. That city is 100 miles from the ocean. Getting there involves a difficult, often turbulent passage over the bar guarding the Columbia River and a long trip up the Columbia and Willamette Rivers.

“But ship captains like to tie up at Portland. They know that as their ships travel the seas, a curious saltwater shellfish called a barnacle fastens itself to the hull and stays there for the rest of its life, surrounding itself with a rocklike shell. As more and more [of these] barnacles attach themselves, they increase the ship’s drag, slow its progress, decrease its efficiency.

“Periodically, the ship must go into dry dock, where with great effort the barnacles are chiseled or scraped off. It’s a difficult, expensive process that ties up the ship for days.

“But not if the captain can get his ship to Portland. Barnacles can’t live in fresh water. There, in the sweet, fresh waters of the Willamette or Columbia, the barnacles [die and some] fall away, [while those that remain are easily removed. Thus,] the ship returns to its task lightened and renewed.

“Sins are like those barnacles. Hardly anyone goes through life without picking up some. They increase the drag, slow our progress, decrease our efficiency. Unrepented, building up one on another, they can eventually sink us.

“In His infinite love and mercy, our Lord has provided a harbor where, through repentance, our barnacles fall away and are forgotten. With our souls lightened and renewed, we can go efficiently about our work and His.”

The priesthood represents a mighty army of righteousness—even a royal army. We are led by a prophet of God, even President Gordon B. Hinckley. In supreme command is our Lord and Savior, Jesus Christ. Our marching orders are clear. They are concise. Matthew describes our challenge in these words from the Master:

“Go ye therefore, and teach all nations, baptizing them in the name of the Father, and of the Son, and of the Holy Ghost:

“Teaching them to observe all things whatsoever I have commanded you: and, lo, I am with you alway, even unto the end of the world.”

“And they went forth, and preached every where, the Lord working with them.”

The call to serve has ever characterized the work of the Lord. It rarely comes at a convenient time. It prompts humility; it invites prayer; it inspires commitment. The call came—to Kirtland. Revelations followed. The call came—to Missouri. Persecution prevailed. The call came—to Nauvoo. Prophets died. The call came—to the basin of the Great Salt Lake. Hardship beckoned.

That long journey, made under such difficult circumstances, was a trial of faith. But faith forged in the furnace of trials and tears is marked by trust and testimony. Only God can count the sacrifice; only He can measure the sorrow; only He can know the hearts of those who serve Him—then and now.

Lessons from the past can quicken our memories, touch our lives, and direct our actions. We are prompted to pause and remember that divinely given promise: “Wherefore, … ye are on the Lord’s errand; and whatever ye do according to the will of the Lord is the Lord’s business.”

Many in this vast audience of priesthood bearers are holders of the Aaronic Priesthood—even deacons, teachers, and priests. Young men, some lessons in life are learned from your parents, while others you learn in school or in church. There are, however, certain moments when you know our Heavenly Father is doing the teaching and you are His student. The thoughts we think, the feelings we feel—even the deeds we do in boyhood—can affect our lives forever.

When I was a deacon, I loved baseball. In fact, I still do. I had a fielder’s glove inscribed with the name Mel Ott. He was the premier player of my day. My friends and I would play ball in a small alleyway behind the houses where we lived. Our playing field was cramped, but all right, provided you hit straightaway to center field. However, if you hit the ball to the right of center, disaster was at the door. Here lived Mrs. Shinas, who, from her kitchen window, would watch us play; and as soon as the ball rolled to her porch, her large dog would retrieve the ball and present it to her as she opened the door. Into her house Mrs. Shinas would return and add the ball to the many she had previously confiscated. She was our nemesis, the destroyer of our fun—even the bane of our existence. None of us had a good word for Mrs. Shinas, but we had plenty of bad words for her. None of us would speak to her, and she never spoke to us. She was hampered by a stiff leg which impaired her walking and must have caused her great pain. She and her husband had no children, lived secluded lives, and rarely came out of their house.

This private war continued for some time—perhaps two years—and then an inspired thaw melted the ice of winter and brought a springtime of good feelings to the stalemate.

One evening as I performed my daily task of watering our front lawn, holding the nozzle of the hose in the hand as was the style at that time, I noticed that Mrs. Shinas’s lawn was dry and beginning to turn brown. I honestly don’t know, brethren, what came over me, but I took a few more minutes and, with our hose, watered her lawn. I continued to do this throughout the summer, and then when autumn came I hosed her lawn free of leaves as I did ours and stacked the leaves in piles at the street’s edge to be gathered. During the entire summer I had not seen Mrs. Shinas. We boys had long since given up playing ball in the alleyway. We had run out of baseballs and had no money to buy more.

Early one evening, Mrs. Shinas’s front door opened, and she beckoned for me to jump the small fence and come to her front porch. This I did. As I approached her, she invited me into her living room, where I was asked to sit in a comfortable chair. She treated me to cookies and milk. Then she went to the kitchen and returned with a large box filled with baseballs and softballs, representing several seasons of her confiscation efforts. The filled box was presented to me. The treasure, however, was not to be found in the gift but rather in her words. I saw for the first time a smile come across the face of Mrs. Shinas, and she said, “Tommy, I want you to have these baseballs, and I want to thank you for being kind to me.” I expressed my own gratitude to her and walked from her home a better boy than when I entered. No longer were we enemies. Now we were friends. The Golden Rule had again succeeded.

Fathers, bishops, quorum advisers—yours is the responsibility to prepare this generation of missionaries, to quicken in the hearts of these deacons, teachers, and priests not only an awareness of their obligation to serve but also a vision of the opportunities and blessings which await them through a mission call. The work is demanding, the impact everlasting. This is no time for “summer soldiers” in the army of the Lord.

Each missionary who goes forth in response to a sacred call becomes a servant of the Lord, whose work this truly is. Do not fear, young men, for He will be with you. He never fails. He has promised: “I will go before your face. I will be on your right hand and on your left, and my Spirit shall be in your hearts, and mine angels round about you, to bear you up.”

Brethren, we have no way of knowing when our privilege to extend a helping hand will unfold before us. The road to Jericho each of us travels bears no name, and the weary traveler who needs our help may be one unknown. Altogether too frequently the recipient of kindness shown fails to express his feelings, and we are deprived of a glimpse of greatness and a touch of tenderness that motivates us to go and do likewise.

Two thousand years ago, Jesus of Nazareth sat by a well in Samaria and talked there to a woman:

“Jesus … said unto her, Whosoever drinketh of this water shall thirst again:

“But whosoever drinketh of the water that I shall give him shall never thirst; but the water that I shall give him shall be in him a well of water springing up into everlasting life.”

Should there be anyone who feels he is too weak to change the onward and downward course of his life, or should there be those who fail to resolve to do better because of that greatest of fears, the fear of failure, there is no more comforting assurance to be had than the words of the Lord: “My grace,” said He, “is sufficient for all men that humble themselves before me; for if they humble themselves before me, and have faith in me, then will I make weak things become strong unto them.”

Through humble prayer, diligent preparation, and faithful service, we can succeed in our sacred callings.

Remember how the captains of oceangoing vessels burdened by the weight of barnacles set a course to the fresh waters of the Columbia and Willamette Rivers to rid themselves of these impediments of progress? Let us, in our own lives and in our service in the Lord’s work, shed the barnacles of doubt, laziness, fear, and sin by plying the living waters of the gospel of Jesus Christ. We know their names: faith, prayer, charity, obedience, and love—to identify but a few. The lighthouse of the Lord Jesus Christ marks the way. His beacon light will guide our course to celestial glory.

May we be wise mariners as we go forth on such a voyage. Let us be pure vessels before the Lord. Let us recognize and respond to the needs of the widow; the cry of the child; the plight of the unemployed; the burden of the sick, the confined, the aged, the poor, the hungry, the lame, and the forgotten. They are remembered by our Heavenly Father and His Beloved Son, Jesus Christ. May you and I follow Their divine examples. Heavenly peace will then be our blessing, in the name of Jesus Christ, amen.

\newpage
\settalk{Dedication Day}
\setspeaker{Thomas S. Monson}
\settalkdate{October 2000}
\talkheading{Dedication Day}{Thomas S. Monson}

A favorite hymn describes the tender feelings of my heart and soul this beautiful day of dedication. I think the words will describe your feelings also:

\begin{verse}
On this day of joy and gladness,\\
Lord, we praise thy holy name;\\
In this sacred place of worship,\\
We thy glories loud proclaim!
\end{verse}

Charles C. Rich, on April 7, 1863, spoke of the need for a tabernacle in which to meet. He declared: “What shall I say in regard to the Tabernacle? We can see at once that we … need the blessings of such a house at the present time. If we put it off, when will it be built? When that house is built we can then enjoy the benefits and blessings which it will afford. The same principle may be applied to everything we take in hand, and with which we have to do, whether it be to build a Temple, a Tabernacle, to send teams to the frontiers to gather the poor, or to do any other work that is required of us. Nothing that is required will be performed until we go to work and do something ourselves. We have no other people to lean upon, and therefore it remains for us to go to work and perform well our part.”

They went to work!

Thanks be to God for our noble prophet, President Gordon B. Hinckley, who, with the foresight of a seer, recognized the need for this magnificent facility and, with the help of many others, “went to work.” The result is before us today and will be dedicated this morning.

As a symbol of our gratitude, as an expression of our love for the Lord, could we not rededicate our lives and our homes in a like manner?

The Apostle Paul, in his Epistle to the Corinthians, added an apostolic dimension to our building commitment when he declared, “Know ye not that ye are the temple of God, and that the Spirit of God dwelleth in you?”

The need for personal dedication and recommitment is essential in today’s society. Just a hurried glance at several newspaper stories describes our plight.

From the Associated Press came the following: “In the name of free speech, the Supreme Court struck down a federal law that shielded children from sex-oriented cable TV channels.”

From the San Jose Mercury News came this story: “Germany may be the economic engine of Europe, but on Sundays it stops. … But global market forces are beginning to disturb Germany’s traditional day of rest. With … American-style [7-day-a-week] shopping [already being offered], and the Internet providing 24-hour access to the world’s goods, such rigid store regulations ‘are like a castle from the old century.’ … To vie with other world-class cities, Berlin must be more aggressive. … ‘We want to make more money.’”

As we view the disillusionment that engulfs countless thousands today, we are learning the hard way what an ancient prophet wrote out for us 3,000 years ago: “He that loveth silver shall not be satisfied with silver; nor he that loveth abundance with increase.”

The revered Abraham Lincoln accurately described our plight: “We have been the recipients of the choicest bounties of Heaven; we have been preserved these many years in peace and prosperity; we have grown in numbers, wealth, and power. … But we have forgotten God. We have forgotten the gracious hand which preserved us in peace and multiplied and enriched and strengthened us, and we have vainly imagined, in the deceitfulness of our hearts, that all these blessings were produced by some superior wisdom and virtue of our own. Intoxicated with unbroken success, we have become too self-sufficient to feel the necessity of redeeming and preserving grace, too proud to pray to the God that made us.”

When the seas of life are stormy, a wise mariner seeks a port of peace. The family, as we have traditionally known it, is such a refuge of safety. “The home is the basis of a righteous life, and no other instrumentality can take its place or fulfill its essential functions.” Actually, a home is much more than a house. A house is built of lumber, brick, and stone. A home is made of love, sacrifice, and respect. A house can be a home, and a home can be a heaven when it shelters a family. When true values and basic virtues undergird the families of society, hope will conquer despair, and faith will triumph over doubt.

Such values, when learned and lived in our families, will be as welcome rain to parched soil. Love will be engendered; loyalty to one’s best self will be enhanced; and those virtues of character, integrity, and goodness will be fostered. The family must hold its preeminent place in our way of life because it’s the only possible base upon which a society of responsible human beings has ever found it practicable to build for the future and maintain the values they cherish in the present.

Happy homes come in a variety of appearances. Some feature families with father, mother, brothers, and sisters living together in a spirit of love. Others consist of a single parent with one or two children, while other homes have but one occupant. There are, however, identifying features which are to be found in a happy home, whatever the number or description of its family members. These identifying features are:

On this, the American continent, Jacob, the brother of Nephi, declared, “Look unto God with firmness of mind, and pray unto him with exceeding faith.”

A prominent judge was asked what we, as citizens of the countries of the world, could do to reduce crime and disobedience to law and to bring peace and contentment into our lives and into our nations. He thoughtfully replied, “I would suggest a return to the old-fashioned practice of family prayer.”

Concerning making our personal lives and our homes libraries of learning, the Lord counseled, “Seek ye out of the best books words of wisdom; seek learning, even by study and also by faith.”

The standard works offer the library of learning of which I speak. We must be careful not to underestimate the capacity of children to read and to understand the word of God.

As parents, we should remember that our lives may be the book from the family library which the children most treasure. Are our examples worthy of emulation? Do we live in such a way that a son or a daughter may say, “I want to follow my dad,” or “I want to be like my mother”? Unlike the book on the library shelf, the covers of which shield its contents, our lives cannot be closed. Parents, we truly are an open book in the library of learning of our homes.

Next, do we exemplify the legacy of love? Do our homes? Bernadine Healy, in a commencement address, gave this counsel: “As a physician, who has been deeply privileged to share the most profound moments of people’s lives including their final moments, let me tell you a secret. People facing death don’t think about what degrees they have earned, what positions they have held, or how much wealth they have accumulated. At the end, what really matters is who you loved and who loved you. That circle of love is everything, and is a great measure of a past life. It is the gift of greatest worth.”

Our Lord and Savior’s message was one of love. It can be as a light to our personal pathway to exaltation.

Near the end of his life, one father looked back on how he had spent his time on earth. An acclaimed, respected author of numerous scholarly works, he said, “I wish I had written one less book and taken my children fishing more often.”

Time passes quickly. Many parents say that it seems like yesterday that their children were born. Now those children are grown, perhaps with children of their own. “Where did the years go?” they ask. We cannot call back time that is past, we cannot stop time that now is, and we cannot experience the future in our present state. Time is a gift, a treasure not to be put aside for the future but to be used wisely in the present.

Have we cultivated a spirit of love in our homes? Observed President David O. McKay, “A true Mormon home is one in which if Christ should chance to enter, he would be pleased to linger and to rest.”

What are we doing to ensure that our homes meet this description? Do we ourselves meet it?

On the journey along the pathway of life, there are casualties. Some depart from the road markers which lead to life eternal, only to discover that the detour chosen ultimately leads to a dead end. Indifference, carelessness, selfishness, and sin all take their costly toll in human lives. There are those who, for unexplained reasons, march to the sound of a different drummer, later to learn they have followed the Pied Piper of sorrow and suffering.

Today there goes forth from this pulpit an invitation to people throughout the world: Come from your wandering way, weary traveler. Come to the gospel of Jesus Christ. Come to that heavenly haven called home. Here you will discover the truth. Here you will learn the reality of the Godhead, the comfort of the plan of salvation, the sanctity of the marriage covenant, the power of personal prayer. Come home.

From our youth, many of us may remember the story of a very young boy who was abducted from his parents and his home and taken to a village situated far away. Under these conditions, the small boy grew to young manhood without a knowledge of his actual parents or earthly home.

But where was home to be found? Where were his mother and father to be discovered? Oh, if only he could remember even their names, his task would be less hopeless. Desperately he sought to recall even a glimpse of his childhood.

Like a flash of inspiration, he remembered the sound of a bell which from the tower atop the village church pealed its welcome each Sabbath morning. From village to village the young man wandered, ever listening for that familiar bell to chime. Some bells were similar, others far different from the sound he remembered.

At length the weary young man stood one Sunday morning before a church of a typical town. He listened carefully as the bell began to peal. The sound was familiar. It was unlike any other he had heard, save that bell which pealed in the memory of his childhood days. Yes, it was the same bell. Its ring was true. His eyes filled with tears. His heart rejoiced in gladness. His soul overflowed with gratitude. The young man dropped to his knees, looked upward beyond the bell tower—even toward heaven—and in a prayer of gratitude whispered, “Thanks be to God. I’m home.”

Like the peal of a remembered bell will be the truth of the gospel of Jesus Christ to the soul of him who earnestly seeks. Many of you have traveled long in a personal quest for that which rings true. The Church of Jesus Christ of Latter-day Saints sends forth to you an earnest appeal: Open your doors to the missionaries. Open your minds to the word of God. Open your hearts—even your very souls—to the sound of that still, small voice which testifies of truth. As the prophet Isaiah promised, “Thine ears shall hear a word … , saying, This is the way, walk ye in it.” Then, like the boy of whom I’ve spoken, you too will, on bended knee, say to your God and mine, “I’m home.”

May such be the blessing of all, I pray in the name of Jesus Christ, amen.

\newpage
\settalk{The Call to Serve}
\setspeaker{Thomas S. Monson}
\settalkdate{October 2000}
\talkheading{The Call to Serve}{Thomas S. Monson}

What a privilege is mine to stand before you tonight in this magnificent Conference Center and in assemblies throughout the world. What a mighty body of priesthood!

For a text, I turn to the words spoken through the Prophet Joseph Smith and found in the 107th section of the Doctrine and Covenants. They apply to all of us, whether bearers of the Aaronic Priesthood or the Melchizedek Priesthood: “Wherefore, now let every man learn his duty, and to act in the office in which he is appointed, in all diligence.”

President Wilford Woodruff declared: “All the organizations of the priesthood have power. The deacon has power, through the priesthood which he holds. So has the teacher. They have power to go before the Lord and have their prayers heard and answered, as well as the prophet, the seer, or the revelator has. … It is by this priesthood that men have ordinances conferred upon them, that their sins are forgiven, and that they are redeemed. For this purpose has it been revealed and sealed upon our heads.”

Those who bear the Aaronic Priesthood should be given opportunities to magnify their callings in that priesthood.

For example, when I was ordained a deacon, our bishopric stressed the sacred responsibility which was ours to pass the sacrament. Emphasized was proper dress, a dignified bearing, and the importance of being clean inside and out.

As we were taught the procedure in passing the sacrament, we were told that we were assisting every member in a renewal of the covenant of baptism, with its responsibilities and blessings. We were also told how we should assist a particular brother—Louis—who had a palsied condition, that he might have the opportunity to partake of the sacred emblems.

How I remember being assigned to pass the sacrament to the row where Louis sat. I was hesitant as I approached this wonderful brother, and then I saw his smile and the eager expression of gratitude that showed his desire to partake. Holding the tray in my left hand, I took a piece of bread and pressed it to his open lips. The water was later served in the same way. I felt I was on holy ground. And indeed I was. The privilege to pass the sacrament to Louis made better deacons of us all.

Noble leaders of young men, you stand at the crossroads in the lives of those whom you teach. Inscribed on the wall of Stanford University Memorial Church is this truth, that we must teach our youth that all that is not eternal is too short, and all that is not infinite is too small.

President Gordon B. Hinckley emphasized our responsibilities when he declared: “In this work there must be commitment. There must be devotion. We are engaged in a great eternal struggle that concerns the very souls of the sons and daughters of God. We are not losing. We are winning. We will continue to win if we will be faithful and true. … There is nothing the Lord has asked of us that in faith we cannot accomplish.”

Brethren, is every ordained teacher given the assignment to home teach? What an opportunity to prepare for a mission. What a privilege to learn the discipline of duty. A boy will automatically turn from concern for self when he is assigned to “watch over” others.

And what of the priests? These young men have the opportunity to bless the sacrament, to continue their home teaching duties, and to participate in the sacred ordinance of baptism.

We can strengthen one another; we have the capacity to notice the unnoticed. When we have eyes that see, ears that hear, and hearts that know and feel, we can reach out and rescue those for whom we have responsibility.

From Proverbs comes the counsel, and I love it, “Ponder the path of thy feet.”

I revere the priesthood of Almighty God. I have witnessed its power. I have seen its strength. I have marveled at the miracles it has wrought.

Fifty years ago, I knew a young man—even a priest—who held the authority of the Aaronic Priesthood. As the bishop, I was his quorum president. Robert stuttered and stammered, void of control. He was self-conscious, shy, fearful of himself and all others, and this impediment was devastating to him. Never did he fulfill an assignment; never would he look another in the eye; always he would gaze downward. Then one day, through a set of unusual circumstances, he accepted an assignment to perform the priestly responsibility to baptize another.

I sat next to Robert in the baptistry of the Salt Lake Tabernacle. He was dressed in immaculate white, prepared for the ordinance he was to perform. I leaned over and asked him how he felt. He gazed at the floor and stuttered almost uncontrollably that he felt terrible, terrible.

We both prayed fervently that he would be made equal to his task. Suddenly the clerk said, “Nancy Ann McArthur will now be baptized by Robert Williams, a priest.”

Robert left my side, stepped into the font, took little Nancy by the hand, and helped her into that water which cleanses human lives and provides a spiritual rebirth. He spoke the words, “Nancy Ann McArthur, having been commissioned of Jesus Christ, I baptize you in the name of the Father, and of the Son, and of the Holy Ghost. Amen.” Not once did he stutter! Not once did he falter! A modern miracle had been witnessed. Robert then performed the baptismal ordinance for two or three other children in the same fashion.

In the dressing room, as I congratulated Robert, I expected to hear this same uninterrupted flow of speech. I was wrong. He gazed downward and stammered his reply of gratitude.

To each of you brethren this evening, I testify that when Robert acted in the authority of the Aaronic Priesthood, he spoke with power, with conviction, and with heavenly help.

We must provide for our young men of the Aaronic Priesthood faith-building experiences. They seek to have the opportunity we have had to feel the Spirit of the Lord helping them.

I remember when I was assigned to give my first talk in church. I was given the liberty to choose my subject. I’ve always liked birds, so I thought of the Seagull Monument. In preparation, I went to Temple Square and looked at the monument. First I was attracted to all the coins in the water surrounding the monument. I wondered how they would be retrieved and who would retrieve them. I shall not confess any thought of taking them.

Then I looked upward at the seagulls atop that monument. I tried in my boyish mind to imagine what it would be like to be a pioneer watching the first year’s growth of precious grain being devoured by crickets and then seeing those seagulls, with their lofty wings, descending upon the fields and eating the crickets. I loved the account. I sat down with a pencil in hand and wrote out a two-and-one-half-minute talk. I’ve never forgotten the seagulls. I’ve never forgotten the crickets. I’ve never forgotten my knees knocking together as I gave that talk. I’ve never forgotten the experience of letting some of my innermost feelings be expressed verbally at the pulpit. I would urge that we give the Aaronic Priesthood an opportunity to think, to reason, and to serve.

President David O. McKay remarked: “God help us all to be true to the ideals of the priesthood—Aaronic and Melchizedek. May he help us to magnify our callings and to inspire men by our actions—not only members of the Church, but all men everywhere—to live higher and better lives, to help them all to be better husbands, better neighbors, better leaders, under all conditions.”

The world seems to have slipped from the moorings of safety and drifted from the harbor of peace. Permissiveness, immorality, pornography, and the power of peer pressure cause many to be tossed about on a sea of sin and crushed on the jagged reefs of lost opportunities, forfeited blessings, and shattered dreams.

Anxiously some may ask, “Is there a way to safety?” “Can someone guide me?” “Is there an escape from threatened destruction?” The answer, brethren, is a resounding “Yes!” Look to the lighthouse of the Lord. There is no fog so dense, no night so dark, no gale so strong, no mariner so lost but what its beacon light can rescue. It beckons through the storms of life. The lighthouse of the Lord sends forth signals readily recognized and never failing.

There are many such signals. I name but three. Note them carefully; exaltation may depend upon them—yours and mine:

First: Prayer provides peace.

Second: Faith precedes the miracle.

And third: Honesty is the best policy.

First, concerning prayer—Adam prayed; Jesus prayed; Joseph prayed. We know the outcome of their prayers. He who notes the fall of a sparrow surely hears the pleadings of our hearts. Remember the promise: “If any of you lack wisdom, let him ask of God, that giveth to all men liberally, and upbraideth not; and it shall be given him.”

Next, faith precedes the miracle. It has ever been so and shall ever be. It was not raining when Noah was commanded to build an ark. There was no visible ram in the thicket when Abraham prepared to sacrifice his son Isaac. Two heavenly personages were not yet seen when Joseph knelt and prayed. First came the test of faith—and then the miracle.

Remember that faith and doubt cannot exist in the same mind at the same time, for one will dispel the other. Cast out doubt. Cultivate faith.

Finally, honesty is the best policy. I learned this truth in a dramatic manner during boot camp when I served in the navy 55 years ago. After those first three weeks of isolated training, the good news came that we would have our first liberty and could visit the city of San Diego. All of the men were most eager for this change of pace. As we prepared to board the buses to town, the petty officer commanded, “Now all of you men who know how to swim, you stand over here. You will go into San Diego for liberty. Those of you who don’t know how to swim, you line up over there. You will go to the swimming pool and have a lesson on how to swim. Only when you learn to swim will you be permitted liberty.”

I had been a swimmer most of my life, so I prepared to get on the bus to town; but then that petty officer said to our group, “One more thing before we board the buses. Follow me. Forward, march!” He marched us right to the swimming pool and had us take our clothing off and stand at the edge of the deep end of the pool. Then he directed, “Jump in and swim the length of the pool.” In that group, all of whom could supposedly swim (at least they had so declared), were about 10 who had thought they could fool somebody. They did not really know how to swim. In the water they went, voluntarily or otherwise. Catastrophe was at the door. The petty officers let them go under once or twice before they extended the bamboo pole to pull them to safety. With a few choice words, which I shall not repeat, they then said, “That will teach you to tell the truth!”

How grateful I was that I had told the truth, that I knew how to swim and made it easily to the other end of the pool. Such lessons teach us to be true—true to the faith, true to the Lord, true to our companions, true to all that is sacred and dear to us. That lesson has never left me.

The lighthouse of the Lord beckons us to safety and eternal joy as we are guided by its never-failing signals:

I testify to you this night that Jesus is indeed the Christ, our beloved Redeemer and Savior. We are led by a prophet of Almighty God—even President Gordon B. Hinckley. I know you share this same conviction.

I close by reading a simple yet profound letter that reflects our love for our prophet and his leadership:

“Dear President Monson,

“Five years ago, President Hinckley was sustained as prophet, seer, and revelator. For me that was an extraordinary occasion which had to do with your calling for the sustaining vote of the Church.

“On that particular morning, I needed to haul hay for my livestock. I was enjoying conference on my truck radio. I had picked up the hay, backed into the barn, and was throwing down hay bales from the back of the truck. When you called for the brethren of the priesthood, ‘wherever you are,’ to prepare to sustain the prophet, I wondered if you meant me. I wondered if the Lord would be offended because I was sweaty and covered with dust. But I took you at your word and climbed down from the truck.

“I shall never forget standing alone in the barn, hat in hand, with sweat running down my face, with arm to the square to sustain President Hinckley. Tears mixed with sweat as I sat for several minutes contemplating this sacred occasion.”

He continued:

“In our lives, we place ourselves at particular places when events of large consequence occur. That has happened to me, but none more spiritual or tender or memorable than that morning in the barn with only cows and a roan horse looking on.

“Sincerely,

“Clark Cederlof”

President Hinckley, we the priesthood brethren of the Church do love and sustain you. I so testify, in the name of Jesus Christ, amen.

\newpage
\settalk{Compassion}
\setspeaker{Thomas S. Monson}
\settalkdate{April 2001}
\talkheading{Compassion}{Thomas S. Monson}

Oklahoma City, Oklahoma, is a most interesting place. In company with Elders Richard G. Scott, Rex D. Pinegar, and Larry W. Gibbons, I presided at a regional conference there just a short time ago. The facility in which we met was packed with members of the Church and other interested persons. The singing by the choir was heavenly, the spoken word inspiring, and the sweet spirit which prevailed during the conference will long be remembered.

I reflected on my previous visits to this location, the beauty of the state song—“Oklahoma,” from the musical production of Rodgers and Hammerstein—and the wonderful hospitality of the people there.

This community’s spirit of compassionate help was tested in the extreme, however, on April 19, 1995, when a terrorist-planted bomb destroyed the Alfred P. Murrah Federal Building in downtown Oklahoma City, taking 168 persons to their deaths and injuring countless others.

Following the regional conference in Oklahoma City, I was driven to the entrance of a beautiful and symbolic memorial which graces the area where the Murrah building once stood. It was a dreary, rainy day, which tended to underscore the pain and suffering which had occurred there. The memorial features a 400-foot reflecting pool. On one side of the pool are 168 empty glass and granite chairs in honor of each of the people killed. These are placed, as far as can be determined, where the fallen bodies were found.

On the opposite side of the pool there stands, on a gentle rise of ground, a mature American elm tree—the only nearby tree to survive the destruction. It is appropriately and affectionately named “The Survivor Tree.” In regal splendor it honors those who survived the horrific blast.

My host directed my attention to the inscription above the gate of the memorial:

He then, with tears in his eyes and with a faltering voice, declared, “This community, and all the churches and citizens in it, have been galvanized together. In our grief we have become strong. In our spirit we have become united.”

We concluded that the best word to describe what had taken place was compassion.

My thoughts turned to the musical play Camelot. King Arthur, in his dream of a better world, an ideal relationship one with another, said, as he envisioned the purpose of the Round Table, “Violence is not strength, and compassion is not weakness.”

A stirring account which illustrates this statement is found in the Old Testament of the Holy Bible. Joseph was especially loved by his father, Jacob, which occasioned bitterness and jealousy on the part of his brothers. There followed the plot to slay Joseph, which eventually placed Joseph in a deep pit without food or water to sustain life. Upon the arrival of a passing caravan of merchants, Joseph’s brothers determined to sell Joseph rather than leaving him to die. Twenty pieces of silver extricated Joseph from the pit and placed him eventually in the house of Potiphar in the land of Egypt. There Joseph prospered, for “the Lord was with Joseph.”

After the years of plenty, there followed the years of famine. In the midst of this latter period, when the brothers of Joseph came to Egypt to buy corn, they were blessed by this favored man in Egypt—even their own brother. Joseph could have dealt harshly with his brothers for the callous and cruel treatment he had earlier received from them. However, he was kind and gracious to them and won their favor and support with these words and actions:

“Now therefore be not grieved, nor angry with yourselves, that ye sold me hither: for God did send me before you to preserve life. …

“And God sent me before you to preserve you a posterity in the earth, and to save your lives by a great deliverance.” Joseph exemplified the magnificent virtue of compassion.

During the meridian of time, when Jesus walked the dusty pathways of the Holy Land, He often spoke in parables.

Said He: “A certain man went down from Jerusalem to Jericho, and fell among thieves, which stripped him of his raiment, and wounded him, and departed, leaving him half dead.

“And by chance there came down a certain priest that way: and when he saw him, he passed by on the other side.

“And likewise a Levite, when he was at the place, came and looked on him, and passed by on the other side.

“But a certain Samaritan, as he journeyed, came where he was: and when he saw him, he had compassion on him,

“And went to him, and bound up his wounds, pouring in oil and wine, and set him on his own beast, and brought him to an inn, and took care of him.

“And on the morrow when he departed, he took out two pence, and gave them to the host, and said unto him, Take care of him; and whatsoever thou spendest more, when I come again, I will repay thee.”

Well could the Savior say to us, “Which now of these three, thinkest thou, was neighbour unto him that fell among the thieves?”

No doubt our response would be, “He that shewed mercy on him.”

Now, as then, Jesus would say to us, “Go, and do thou likewise.”

Jesus provided us many examples of compassionate concern. The crippled man at the pool of Bethesda; the woman taken in adultery; the woman at Jacob’s well; the daughter of Jairus; Lazarus, brother of Mary and Martha—each represented a casualty on the Jericho road. Each needed help.

To the cripple at Bethesda, Jesus said, “Rise, take up thy bed, and walk.” To the sinful woman came the counsel, “Go, and sin no more.” To help her who came to draw water, He provided a well of water “springing up into everlasting life.” To the dead daughter of Jairus came the command, “Damsel, I say unto thee, arise.” To the entombed Lazarus, “Come forth.”

The Savior has always shown unlimited capacity for compassion.

On this, the American continent, Jesus appeared to a multitude and said:

“Have ye any that are sick among you? Bring them hither. Have ye any that are lame, or blind, or halt, or maimed, or leprous, or that are withered, or that are deaf, or that are afflicted in any manner? Bring them hither and I will heal them, for I have compassion upon you. …

“… And he did heal them every one.”

One may well ask the penetrating question: These accounts pertain to the Redeemer of the world. Can there actually occur in my own life, on my own Jericho road, such a treasured experience?

I phrase my answer in the words of the Master, “Come and see.”

We have no way of knowing when our privilege to extend a helping hand will unfold before us. The road to Jericho each of us travels bears no name, and the weary traveler who needs our help may be one unknown.

Genuine gratitude was expressed by the writer of a letter received some time ago at Church headquarters. No return address was shown, no name, but the postmark was from Portland, Oregon:

“To the Office of the First Presidency:

“Salt Lake City showed me Christian hospitality once during my wandering years.

“On a cross-country journey by bus to California, I stepped down in the terminal in Salt Lake City, sick and trembling from aggravated loss of sleep caused by a lack of necessary medication. In my headlong flight from a bad situation in Boston, I had completely forgotten my supply.

“In the Temple Square Hotel restaurant, I sat dejectedly. Out of the corner of my eye I saw a couple approach my table. ‘Are you all right, young man?’ the woman asked. I raised up, crying and a bit shaken, related my story and the predicament I was in then. They listened carefully and patiently to my nearly incoherent ramblings, and then they took charge. They spoke with the restaurant manager, then told me I could have all I wanted to eat there for five days. They took me next door to the hotel desk and got me a room for five days. Then they drove me to a clinic and saw that I was provided with the medications I needed—truly my basic lifeline to sanity and comfort.

“While I was recuperating and building my strength, I made it a point to attend the daily Tabernacle organ recitals. The celestial voicing of that instrument from the faintest intonation to the mighty full organ is the most sublime sonority of my acquaintance. I have acquired albums and tapes of the Tabernacle organ and the choir which I can rely upon any time to soothe and buttress a sagging spirit.

“On my last day at the hotel, before I resumed my journey, I turned in my key; and there was a message for me from that couple: ‘Repay us by showing gentle kindness to some other troubled soul along your road.’ That was my habit, but I determined to be more keenly on the lookout for someone who needed a lift in life.

“I wish you well. I don’t know if these are indeed the ‘latter days’ spoken of in the scriptures, but I do know that two members of your church were saints to me in my desperate hours of need. I just thought you might like to know.”

What an example of caring compassion.

At one privately owned and operated care facility, compassion reigned supreme. The proprietress was Edna Hewlett. There was a waiting list of patients who desired to live out their remaining days under her tender care, for she was an angelic person. She would wash and style the hair of every patient. She cleansed elderly bodies and dressed them with bright and clean clothing.

Through the years, in visiting the widows of the ward over which I once presided, I would generally start my visits at Edna’s facility. She would welcome me with a cheery smile and take me to the living room where a number of the patients were seated. I always had to begin with Jeannie Burt, who was the oldest—102 when she died. She had known me and my family from the time I was born.

On one occasion, Jeannie asked with her thick Scottish brogue, “Tommy, have you been to Edinburgh lately?”

I replied, “Yes, not too long ago I was there.”

“Isn’t it beautiful!” she responded.

Jeannie closed her aged eyes in an expression of silent reverie. Then she became serious. “I’ve paid in advance for my funeral—in cash. You are to speak at my funeral and you are to recite ‘Crossing the Bar’ by Tennyson. Now let’s hear it!”

It seemed every eye was upon me, and surely this was the case. I took a deep breath and began:

\begin{verse}
Sunset and evening star,\\
And one clear call for me!\\
And may there be no moaning of the bar,\\
When I put out to sea.
\end{verse}

Jeannie’s smile was benign and heavenly—then she declared, “Oh, Tommy, that was nice. But see that you practice a wee bit before my funeral!” This I did.

At some period in our mortal mission, there appears the faltering step, the wan smile, the pain of sickness—even the fading of summer, the approach of autumn, the chill of winter, and the experience we call death, which comes to all mankind. It comes to the aged as they walk on faltering feet. Its summons is heard by those who have scarcely reached midway in life’s journey. Often it hushes the laughter of little children.

Throughout the world there is enacted daily the sorrowful scene of loved ones mourning as they bid farewell to a son, a daughter, a brother, a sister, a mother, a father, or a cherished friend.

From the cruel cross, the Savior’s tender words of farewell to His mother are particularly poignant:

“When Jesus therefore saw his mother, and the disciple standing by, whom he loved, he saith unto his mother, Woman, behold thy son!

“Then saith he to the disciple, Behold thy mother! And from that hour that disciple took her unto his own home.”

Let us remember that after the funeral flowers fade, the well-wishes of friends become memories and the prayers offered and words spoken dim in the corridors of the mind. Those who grieve frequently find themselves alone. Missed is the laughter of children, the commotion of teenagers, and the tender, loving concern of a departed companion. The clock ticks more loudly, time passes more slowly, and four walls can indeed a prison make.

I extol those who, with loving care and compassionate concern, feed the hungry, clothe the naked, and house the homeless. He who notes the sparrow’s fall will not be unmindful of such service.

In our Father’s compassion and according to His divine plan, holy temples bring to His children the peace which surpasses understanding.

Today, under the leadership of President Gordon B. Hinckley, the number of new temples constructed and under construction staggers the mind to contemplate. Heavenly Father’s compassionate concern for His children here on earth and for those who have gone beyond mortality merits our gratitude.

Thanks be to our Lord and Savior Jesus Christ for His life, for His gospel, for His example, and for His blessed Atonement.

I return in my thoughts to Oklahoma City. To me, it is beyond mere coincidence that now a temple of the Lord, in all its beauty, stands in that city as a heaven-sent beacon to mark the way to joy here on earth and eternal joy hereafter. Let us remember the words from the Psalms, “Weeping may endure for a night, but joy cometh in the morning.”

In a very real way, the Master speaks to us: “Behold, I stand at the door, and knock: if any man hear my voice, and open the door, I will come in to him.”

Let us listen for His knock. Let us open the door of our hearts, that He—the living example of true compassion—may enter, I sincerely pray, in the name of Jesus Christ, amen.

\newpage
\settalk{To the Rescue}
\setspeaker{Thomas S. Monson}
\settalkdate{April 2001}
\talkheading{To the Rescue}{Thomas S. Monson}

Mine is the overwhelming and humbling responsibility tonight to address you, my dear brethren who hold the priesthood of God and who have assembled here in the Conference Center and throughout the world.

Some of you are deacons, perhaps newly ordained; others of you are high priests who have served long and faithfully in sacred callings. All have assembled that we might better learn our duty.

Brethren, the world is in need of your help. There are feet to steady, hands to grasp, minds to encourage, hearts to inspire, and souls to save. The blessings of eternity await you. Yours is the privilege to be not spectators but participants on the stage of priesthood service.

President Wilford Woodruff declared: “All the organizations of the Priesthood have power. The Deacon has power, through the Priesthood which he holds. So has the Teacher. They have power to go before the Lord and have their prayers heard and answered, as well as the Prophet. … It is by this Priesthood that men have ordinances conferred upon them, that their sins are forgiven, and that they are redeemed. For this purpose it has been revealed and sealed upon our heads.”

Once I heard from a newly ordained deacon soon after he had received the Aaronic Priesthood. He said, “Today is my first day to pass the sacrament. I can’t wait. I know it is a very holy ordinance, so I’ll treat it with care. I have a true testimony of the Church, and I hope to go on a mission soon.”

May I share with you tonight, brethren, a letter which I received some time ago, written by a husband who strayed far from the priesthood path of service and duty. It typifies the plea of too many of our brethren. He wrote:

“Dear President Monson:

“I had so much and now have so little. I am unhappy and feel as though I am failing in everything. The gospel has never left my heart, even though it has left my life. I ask for your prayers.

“Please don’t forget those of us who are out here—the lost Latter-day Saints. I know where the Church is, but sometimes I think I need someone else to show me the way, encourage me, take away my fear, and bear testimony to me.”

While reading this letter, I returned in my thoughts to a visit to one of the great art galleries of the world—even the famed Victoria and Albert Museum in London, England. There, exquisitely framed, was a masterpiece painted in 1831 by Joseph Mallord William Turner. The painting features heavy-laden black clouds and the fury of a turbulent sea portending danger and death. A light from a stranded vessel gleams far off. In the foreground, tossed high by incoming waves of foaming water, is a large lifeboat. The men pull mightily on the oars as the lifeboat plunges into the tempest. On the shore there stand a wife and two children, wet with rain and whipped by wind. They gaze anxiously seaward. In my mind I abbreviated the name of the painting. To me, it became To the Rescue.

Amidst the storms of life, danger lurks; and men, like boats, find themselves stranded and facing destruction. Who will man the lifeboats, leaving behind the comforts of home and family, and go to the rescue?

President John Taylor cautioned us, “If you do not magnify your callings, God will hold you responsible for those whom you might have saved had you done your duty.”

Brethren, our task is not insurmountable. We are on the Lord’s errand, and therefore we are entitled to the Lord’s help. But we must try. From the stage play Shenandoah comes the spoken line which inspires: “If we don’t try, then we don’t do; and if we don’t do, then why are we here?”

When the Master ministered among men, He called fishermen at Galilee to leave their nets and follow Him, declaring, “I will make you fishers of men.” And so He did. Tonight He issues a call to each of us to “come join the ranks.” He provides our battle plan with His admonition, “Wherefore, now let every man learn his duty, and to act in the office in which he is appointed, in all diligence.”

I love and cherish the noble word duty. Let us hearken to the stirring reminder found in the epistle of James: “Be ye doers of the word, and not hearers only, deceiving your own selves.”

There is an old song of my vintage. It’s entitled “Wishing Will Make It So.” It’s not true. Wishing will not make it so. The Lord expects our thinking. He expects our action. He expects our labors. He expects our testimonies. He expects our devotion. Unfortunately, there are those who have departed from the track of priesthood activity. Let us help them back to that path that leads to life eternal. Let us build that strong Melchizedek Priesthood base which will be the foundation of Church activity and growth. It will be the underpinning to strengthen every family, every home, every quorum in every land.

Brethren, we can reach out to those for whom we are responsible and bring them to the table of the Lord, there to feast on His word and to enjoy the companionship of His Spirit and be “no more strangers and foreigners, but fellowcitizens with the saints, and of the household of God.”

The passage of time has not altered the capacity of the Redeemer to change men’s lives—our lives and the lives of those with whom we labor. As He said to the dead Lazarus, so He says today: “Come forth.” Come forth from the despair of doubt. Come forth from the sorrow of sin. Come forth from the death of disbelief. Come forth to a newness of life. Come forth.

We will discover that those whom we serve, who have felt through our labors the touch of the Master’s hand, somehow cannot explain the change which comes into their lives. There is a desire to serve faithfully, to walk humbly, and to live more like the Savior. Having received their spiritual eyesight and glimpsed the promises of eternity, they echo the words of the blind man to whom Jesus restored sight, who said, “One thing I know, that, whereas I was blind, now I see.”

How can we account for these miracles? Why the upsurge of activity in men long dormant? The poet, speaking of death, wrote, “God … touched him, and he slept.” I say, speaking of this new birth, “God touched them, and they awakened.”

Two fundamental reasons largely account for these changes of attitudes, of habits, of actions. First, men have been shown their eternal possibilities and have made the decision to achieve them. Men cannot really long rest content with mediocrity once they see excellence is within their reach.

Second, other men have followed the admonition of the Savior and have loved their neighbors as themselves and helped to bring their neighbors’ dreams to fulfillment and their ambitions to realization.

The catalyst in this process has been—and will continue to be—the principle of love.

Another principle of truth which will guide us in our determination is that boys and men can change. I’m reminded of the words of a prison warden who taught this fact. A critic who knew of Warden Duffy’s efforts to rehabilitate men said, “Don’t you know that leopards can’t change their spots?”

Warden Duffy responded, “You should know I don’t work with leopards. I work with men, and men change every day.”

Many years ago, before leaving to become president of the Canadian Mission, headquartered in Toronto, Ontario, I had developed a friendship with a man by the name of Shelley, who lived in my ward but did not embrace the gospel, irrespective of the fact that his wife and children had done so. Shelley had been known as the toughest man in town when he was young. He was quite a pugilist. His fights were rarely in the ring but rather elsewhere. Try as I might, I could not bring about a change in Shelley’s attitude. The task appeared hopeless. In time, Shelley and his family moved from our ward.

After I had returned from Canada and was called to the Twelve, I received a telephone call from Shelley. He said, “Will you seal my wife and me and our family in the Salt Lake Temple?”

I answered hesitatingly, “Shelley, you first must be a baptized member of the Church.”

He laughed and responded, “Oh, I took care of that while you were in Canada. My home teacher was a school crossing guard, and every weekday as he and I would visit at the crossing, we would discuss the gospel.”

The sealings were performed; a family was united; joy followed.

Abraham Lincoln offered this wise counsel, which surely applies to home teachers: “If you would win a man to your cause, first convince him that you are his sincere friend.”

A friend makes more than a dutiful visit each month. A friend is more concerned about helping people than getting credit. A friend cares. A friend loves. A friend listens. And a friend reaches out.

There are brethren in every ward who seem to have a special skill and aptitude to penetrate the outer shell and reach the heart. Such was Raymond L. Egan, who served as my counselor in the bishopric. He loved to befriend and activate in the Church the father of a family and thereby bring into the fold a dear wife and precious children as well. This wonderful phenomenon occurred many times right up until Brother Egan departed mortality.

There are other ways as well by which one might lift and serve. On one occasion, I was speaking with a retired executive I had known for a long time. I asked him, “Ed, what are you doing in the Church?” He replied, “I have the best assignment in the ward. My responsibility is to help men who are unemployed find permanent employment. This year I have helped 12 of my brethren who were out of work to obtain good jobs. I have never been happier in my entire life.” Short in stature, “Little Ed,” as we affectionately called him, stood tall that evening as his eyes glistened and his voice quavered. He showed his love by helping those in need. He restored human dignity. He opened doors for those who knew not how to do so themselves.

I truly believe that those who have the ability to reach out and to lift up have found the formula descriptive of Brother Walter Stover—a man who spent his entire life in service to others. At Brother Stover’s funeral, his son-in-law paid tribute to him in these words: “Walter Stover had the ability to see Christ in every face he encountered, and he treated each person accordingly.” Legendary are his acts of compassionate help and his talent to lift heavenward every person whom he met. His guiding light was the Master’s voice speaking, “Inasmuch as ye have done it unto one of the least of these … , ye have done it unto me.”

Brethren, acquire the language of the Spirit. It is not learned from textbooks written by men of letters, nor is it acquired through reading and memorization. The language of the Spirit comes to him who seeks with all his heart to know God and keep His divine commandments. Proficiency in this “language” permits one to breach barriers, overcome obstacles, and touch the human heart.

In a day of danger or a time of trial, such knowledge, such hope, such understanding bring comfort to a troubled soul and a grieving heart. Shadows of despair are dispelled by rays of hope; sorrow yields to joy; and the feeling of being lost in the crowd of life vanishes with the certain knowledge that our Heavenly Father is mindful of each of us.

In closing, I return to the painting by Turner. In a very real sense, those persons stranded on the vessel which had run aground in the storm-tossed sea are like many young men—and older men as well—who await rescue by those of us who have the priesthood responsibility to man the lifeboats. Their hearts yearn for help. Mothers and fathers pray for their sons. Wives and children plead to heaven that Daddy and others may be reached.

Tonight I pray that all of us who hold the priesthood may sense our responsibilities and, as one, follow our Leader—even the Lord Jesus Christ, and His prophet, President Gordon B. Hinckley—to the rescue.

In the name of Jesus Christ, amen.

\newpage
\settalk{“Be Thou an Example”}
\setspeaker{Thomas S. Monson}
\settalkdate{October 2001}
\talkheading{“Be Thou an Example”}{Thomas S. Monson}

Tonight we have been inspired by the stirring messages of the general presidency of the Relief Society of the Church. Their plea that all of us be steadfast and immovable is wise counsel, that we might meet the turbulence of our times and indeed be citadels of constancy midst a sea of change.

Let us review words of wisdom written by the Apostle Paul to his beloved Timothy:

“Now the Spirit speaketh expressly, that in the latter times some shall depart from the faith, giving heed to seducing spirits, and doctrines of devils;

“Speaking lies in hypocrisy; having their conscience seared with a hot iron.”

Then came Paul’s rallying call to Timothy—equally applicable to each one of us: “Be thou an example of the believers, in word, in conversation, in charity, in spirit, in faith, in purity.”

With you dear sisters assembled here in the Conference Center and in congregations throughout the world, I share a three-part formula to serve as an unfailing guide to meet this challenge issued by the Apostle Paul:

First, fill your mind with truth. We do not find truth groveling through error. Truth is found by searching, studying, and living the revealed word of God. We adopt error when we mingle with error. We learn truth when we associate with truth.

The Savior of the world instructed, “Seek ye out of the best books words of wisdom; seek learning, even by study and also by faith.” He added, “Search the scriptures; for in them ye think ye have eternal life: and they are they which testify of me.”

He invites each of us, “Learn of me, and listen to my words; walk in the meekness of my Spirit, and you shall have peace in me.”

One from pioneer times who exemplified the charge heard this evening to be steadfast and immovable and who filled her mind, heart, and soul with truth was Catherine Curtis Spencer. Her husband, Orson Spencer, was a sensitive, well-educated man. She had been reared in Boston and was cultured and refined. She had six children. Her delicate health declined from exposure and from the hardships encountered after leaving Nauvoo. Elder Spencer wrote to her parents and asked if she could return to live with them while he established a home for her in the West. Their reply: “Let her renounce her degrading faith, and she can come back—but never until she does.”

Sister Spencer would not renounce her faith. When her parents’ letter was read to her, she asked her husband to get his Bible and read to her from the book of Ruth as follows: “Intreat me not to leave thee, or to return from following after thee: for whither thou goest, I will go; and where thou lodgest, I will lodge: thy people shall be my people, and thy God my God.”

Outside the storm raged, the wagon covers leaked, and friends held milk pans over Sister Spencer’s head to keep her dry. In these conditions and without a word of complaint, she closed her eyes for the last time.

Though we may not necessarily be called upon to forfeit our lives, let us remember that He hears our silent prayers. He who observes our unheralded acts will reward us openly when the need comes.

We live in turbulent times. Often the future is unknown; therefore, it behooves us to prepare for uncertainties. Statistics reveal that at some time, because of the illness or death of your husband or because of economic necessity, you may find yourself in the role of financial provider. I urge you to pursue your education and learn marketable skills so that, should an emergency arise, you are prepared to provide.

Your talents will expand as you study and learn. You will be able to better assist your children in their learning, and you will have peace of mind in knowing that you have prepared yourself for the eventualities that you may encounter in life.

To illustrate the second part of our formula—namely, fill your heart with love—I turn to a beautiful account recorded in the book of Acts which tells of a disciple named Tabitha, or Dorcas, who lived at Joppa. She was described as being a woman “full of good works and almsdeeds.”

“It came to pass in those days, that she was sick, and died: whom when they had washed, they laid her in an upper chamber.

“And forasmuch as … the disciples had heard that Peter was there, they sent unto him two men, desiring him that he would not delay to come to them.

“Then Peter arose and went with them. When he was come, they brought him into the upper chamber: and all the widows stood by him weeping, and shewing the coats and garments which [Tabitha] made, while she was with them.

“But Peter put them all forth, and kneeled down, and prayed; and turning him to the body said, Tabitha, arise. And she opened her eyes: and when she saw Peter, she sat up.

“And he gave her his hand, and lifted her up, and when he had called the saints and widows, presented her alive.

“And it was known throughout all Joppa; and many believed in the Lord.”

To me the scriptural reference to Tabitha, which describes her as a woman “full of good works and almsdeeds,” defines some of the fundamental responsibilities of Relief Society; namely, the relief of suffering, the caring for the poor, and all which that implies. Women of Relief Society, you truly are angels of mercy. This is demonstrated on a grand scale through the humanitarian outreach to the cold, the hungry, and to suffering wherever it is found. Your labors are also very much in evidence in our wards and in our stakes and missions. Every bishop in the Church could testify of this truth.

I remember when, as a young deacon, I would cover a portion of the ward on fast Sunday morning, giving the small envelope to each family, waiting while a contribution was placed in the envelope and then returning it to the bishop. On one such occasion, an elderly member, Brother Wright, who lived alone, welcomed me at the door and, with aged hands, fumbled at the tie of the envelope and placed within it a small sum. His eyes fairly glistened as he made his contribution. He invited me to sit down and then told me of a time many years before when his cupboard had been empty of food. In his hunger, he had prayed to Heavenly Father for food to eat. Not long thereafter, he gazed out his front window and beheld someone approaching his door, pulling behind her a red-colored wagon. It was Sister Balmforth, the Relief Society president, who had pulled that wagon almost half a mile over the railroad tracks and to his door. The wagon overflowed with food collected from the sisters of the ward Relief Society, with which Sister Balmforth filled the empty shelves in Brother Wright’s kitchen. He described her to me as “an angel sent from heaven.”

Sisters, you are the epitome of love. You brighten your homes, you lead with kindness your children; and while your husbands may be head of the home, you surely are the heart of the home. Together, through respect for each other and sharing of responsibilities, you make an unbeatable team.

To me it is significant that when children need care and loving attention, they turn to you—their mothers. Even the wayward son or neglectful daughter, when he or she recognizes the need to return to the embrace of family, almost inevitably comes to Mother, who has never given up on her child.

Mother’s love brings out the best in a child. You become the model for your children to follow.

The first word a child learns and utters is usually the dear expression “Mama.” To me it is significant that on the battlefields of war or in peace, frequently when death is about to overtake a son, his final word is usually “Mother.” Sisters, what a noble role is yours. I testify that your hearts are filled with love.

To the third part of our formula—namely, fill your life with service—I mention two separate examples. One features a teacher and the profound influence she has had in the lives of those whom she taught, while the other pertains to a missionary couple whose service helped to bring the light of the gospel to those who had lived in spiritual darkness.

Many years ago there was a young woman, Baur Dee Sheffield, who taught in Mutual. She had no children of her own, though she and her husband dearly longed for children. Her love was expressed through devotion to her special young women as each week she taught them eternal truths and lessons of life. Then came illness, followed by death. She was but 27.

Each year, on Memorial Day, her Mutual girls made a pilgrimage of prayer to the graveside of their teacher, always leaving flowers and a little card signed “To Baur Dee, from your girls.” First there were 10 girls who went, then five, then two, and eventually just one, who continues to visit each Memorial Day, always placing on the grave a bouquet of flowers and a card, inscribed as always, “To Baur Dee, from your girls.”

One year, nearly 25 years after Baur Dee’s death, the only one of “her girls” who continued to visit the grave realized she would be away on Memorial Day and decided to visit her teacher’s grave a few days early. She had gathered flowers, tied them with a ribbon, attached a card, and was putting on her jacket to leave when her doorbell rang. She opened the door and was greeted by one of her visiting teachers, Colleen Fuller, who said she had experienced difficulty getting together with her visiting teaching partner and so had decided to come alone and unannounced in an effort to complete her visiting teaching before the end of the month. As Colleen was invited in, she noticed the jacket and flowers and apologized for obviously interrupting whatever had been planned.

“Oh, no problem,” came the response. “I’m just on my way to the cemetery to put flowers on the grave of the woman who was my Mutual teacher, who had a profound influence on me and the other girls she taught. Originally about 10 of us visited her grave each year to express our love and thanks to her, but now I represent the group.”

Colleen asked, “Could your teacher’s name have been Baur Dee?”

“Why, yes,” came the answer. “How did you know?”

With a catch in her voice, Colleen said, “Baur Dee was my aunt—my mother’s sister. Every Memorial Day since she died, my family has found on her grave a bouquet of flowers and a card inscribed from Baur Dee’s girls. They’ve always wanted to know who these girls were so they could thank them for remembering Baur Dee. Now I can let them know.”

Said American author Thornton Wilder, “The highest tribute to the dead is not grief but gratitude.”

The second example of lives filled with service, with which I shall conclude, is the missionary experience of Juliusz and Dorothy Fussek, who were called to fill an 18-month mission in Poland. Brother Fussek was born in Poland. He spoke the language. He loved the people. Sister Fussek was born in England and knew little of Poland and nothing of its people.

Trusting in the Lord, they embarked on their assignment. The living conditions were primitive, the work lonely, their task immense. A mission had not at that time been fully established in Poland. The assignment given the Fusseks was to prepare the way so that the mission could be expanded and gain permanence, that other missionaries be called to serve, people taught, converts baptized, branches established, and chapels erected.

Did Elder and Sister Fussek despair because of the enormity of their assignment? Not for a moment. They knew their calling was from God, they prayed for His divine help, and they devoted themselves wholeheartedly to their work. They remained in Poland not 18 months, but rather served for five years. All of the foregoing objectives were realized. Such came about following an earlier meeting where Elders Russell M. Nelson, Hans B. Ringger, and I, accompanied by Elder Fussek, met with Minister Adam Wopatka of the Polish government, and we heard him say, “Your church is welcome here. You may build your buildings, you may send your missionaries. You are welcome in Poland. This man,” pointing to Juliusz Fussek, “has served your church well, as has his wife. You can be grateful for their example and their work.”

Like the Fusseks, let us do what we should do in the work of the Lord. Then we can, with Juliusz and Dorothy Fussek, echo the Psalm: “My help cometh from the Lord.”

Dear sisters, you indeed are “examples of the believers.” May our Heavenly Father bless each of you, married or single, in your homes, in your families, in your very lives—that you may merit the glorious salutation of the Savior of the world: “Well done, thou good and faithful servant.” For this I pray, as I bless you, in the name of Jesus Christ, amen.

\newpage
\settalk{Duty Calls}
\setspeaker{Thomas S. Monson}
\settalkdate{October 2001}
\talkheading{Duty Calls}{Thomas S. Monson}

My dear brethren, it is an awesome responsibility, and yet a precious privilege, to respond to the assignment to speak to you this evening. The excitement, the anticipation of general conference, including the general priesthood meeting—whether in person or by satellite or television—brings joy to our hearts.

The Lord has made it clear what our responsibilities are and has given to us in the 107th section of the Doctrine and Covenants a solemn charge: “Wherefore, now let every man learn his duty, and to act in the office in which he is appointed, in all diligence.”

At times the performance of duty, the response to a divine call, or the reaction to a spiritual prompting is not overwhelming. On occasion, however, the duty to respond is downright overpowering. I experienced such a situation prior to the general conference of April 1966. That’s 35 years ago, but I remember it vividly.

I had received my assignment to speak at one of the conference sessions and had prepared and committed to memory a message entitled “Meeting Your Goliath.” This was based on the account of the famous battle waged by David and Goliath of olden times.

Then I received a telephone call from President David O. McKay. The conversation went about like this: “Brother Monson, this is President McKay calling. How are you?”

I took a deep breath and answered, “Oh, I’m fine, President, and looking forward to conference.”

“That’s why I’m calling, Brother Monson. The Saturday morning session will be rebroadcast on Sunday as our Easter message to the world. I will be speaking to an Easter theme and would like you to join me and speak during that important session to that type of theme.”

“Of course, President. I will be happy to do so.”

That’s when the extent of this brief conversation really dawned on me. All of a sudden “Meeting Your Goliath” didn’t quite measure up to an Easter message. I knew I must begin to prepare all over again. There was so little time. Indeed, my “Goliath” stood before me.

That night I cleared the kitchen table and placed my typewriter on the tabletop along with a ream of bond paper, with a trusty wastepaper basket by my side to hold all the false starts that accompany such an assignment of preparation. I began at about 7:00 p.m. and had not written a satisfactory line by 1:00 a.m. The wastebasket was filled, but my mind certainly was not. What was I to do? The clock was running—indeed, it was racing. I paused to pray.

Soon thereafter there came to my mind the sadness of my neighbors Mark and Wilma Shumway in the recent loss of their youngest child. I thought to myself, “Perhaps I could speak directly to them and peripherally to all others, for who hasn’t lost a dear one and had occasion to grieve?” My fingers raced over the typewriter keyboard but could barely keep up with my thoughts.

As the first dim light of morning peered through our kitchen window, I had finished the message. The task remained to learn it and then deliver it to the world. Rarely have I struggled so hard to fill a prophetic assignment. However, Heavenly Father had heard my prayer. I shall never forget the experience.

Two landmark passages of scripture filled my soul as the conference session concluded. They are both familiar to you, brethren. They have no expiration date attached to them. First, from Nephi of old: “I will go and do the things which the Lord hath commanded, for I know that the Lord giveth no commandments unto the children of men, save he shall prepare a way for them that they may accomplish the thing which he commandeth them.”

Second is the promise of the Lord Himself to you and to me from the Doctrine and Covenants: “I will go before your face. I will be on your right hand and on your left, and my Spirit shall be in your hearts, and mine angels round about you, to bear you up.”

Many of us assembled tonight hold the Melchizedek Priesthood, while others bear the Aaronic Priesthood. All of us have a solemn duty to honor the priesthood and labor to bring many precious souls unto the Lord. We remember His declaration: “The worth of souls is great in the sight of God.” Are we doing all we should? Do we remember the words of President John Taylor: “If you do not magnify your callings, God will hold you responsible for those whom you might have saved had you done your duty.”

The desire to help another, the quest for the lost sheep, may not always yield success at once. On occasion, progress is slow—even indiscernible. Such was the experience of my longtime friend Gill Warner. He was a newly called bishop when Douglas, a member of his ward serving in the mission field, transgressed and was deprived of his Church membership. Father was saddened; Mother was totally devastated. Douglas soon thereafter moved from the state. The years hurried by, but Bishop Warner, now a member of a high council, never ceased to wonder what had become of Douglas.

In 1975 I attended the stake conference of Brother Warner’s stake and held a priesthood leadership meeting early on Sunday morning. I spoke of the Church discipline system and the need to labor earnestly and lovingly to rescue any who had strayed. Gill Warner raised his hand and outlined the story of Douglas. He concluded by posing a question to me: “Do I have any responsibility to work with Douglas to help bring him back to Church membership?”

Gill reminded me later that my response to his question was direct and given without hesitation. I said, “As his former bishop and one who knew and loved him, I would think you would wish to do all you could to bring him back.”

Unbeknownst to Gill Warner, Douglas’s mother had, the previous week, fasted and prayed that a man would be raised up to help save her son. Gill discovered this when, following the meeting, he felt prompted to call her to report his determination to be of help.

Gill began his odyssey of redemption. Douglas was contacted by him. Old times, happy times, were remembered. Testimony was expressed, love was conveyed, and confidence instilled. The pace was excruciatingly slow. Discouragement frequently entered the scene; but step-by-step Douglas made headway. At long last prayers were answered, efforts rewarded, and victory attained. Douglas was approved for baptism.

The baptismal date was set, family members gathered, and former bishop Gill Warner flew to the city where Douglas lived and performed the ordinance.

Bishop Warner, through the love of his heart and with a sense of responsibility to a former priest in the Aaronic Priesthood—even the quorum over which he presided—went to the rescue, that not one would be lost.

There may be others, but I have personally known three bishops who, when they presided over their wards, there were in the priests quorum 48 or more young men, or in other words, a full scripturally defined quorum of priests. These three bishops were Alvin R. Dyer, Joseph B. Wirthlin, and Alfred B. Smith. Were they overwhelmed by their task? Not at all. Through their diligent efforts and with the help of caring parents and the blessings of the Lord, these bishops guided each member of their respective priests quorum—almost without exception—to the ordination of elder in the Melchizedek Priesthood, service in the mission field, and marriage in the temple of the Lord. While Brother Dyer and Brother Smith have gone to their eternal reward, Elder Joseph B. Wirthlin, a member of the Quorum of the Twelve Apostles, is present with us tonight. Elder Wirthlin, your service and leadership with these young men, now grown older, will never be forgotten.

As a 12-year-old boy, I had the privilege to serve as the secretary of my deacons quorum. I recall with joy the many assignments we members of that quorum had the opportunity to fill. Passing the sacred sacrament, collecting the monthly fast offerings, and looking after one another come readily to mind. The most frightening one, however, happened at the leadership session of our ward conference. The member of our stake presidency presiding was William F. Perschon. He called on a number of the ward officers to speak. Then, without the slightest warning, President Perschon stood and said, “We will now hear from Thomas S. Monson, secretary of the deacons quorum, to give us an accounting of his service and bear his testimony.” I don’t recall a thing I said, but I have never forgotten the experience.

Brethren, remember the Apostle Peter’s admonition, “Be ready always to give an answer to every man that asketh you a reason of the hope that is in you.”

During World War II, as a teenager I was privileged to serve as president of the teachers quorum. I was asked to learn and then apply the counsel from the Doctrine and Covenants, section 107, verse 86: “The duty of the president over the office of the teachers is to preside over … the teachers, and to sit in council with them, teaching them the duties of their office, as given in the covenants.” I tried my best to live up to this defined duty.

In that quorum was a young man, Fritz Hoerold. He was short in stature but tall on courage. Soon after Fritz turned 17, he enlisted in the United States Navy and was off to training. He then found himself on a large battleship in a number of those bloody engagements in the Pacific. His ship was severely damaged, and many sailors were killed or wounded.

Fritz returned home on leave following such an engagement and came back to our teachers quorum. The quorum adviser invited him to speak to us. My, he looked resplendent in his Navy blues with appropriate war ribbons attached. I remember that I asked Fritz to tell us any thoughts he had for our benefit. After all, we were about the same age. With a wry smile, he responded, “Don’t volunteer for anything!”

I didn’t see Fritz again from that time when we were 17 until, a few years ago, I read a magazine article concerning those same battles at sea. I wondered if Fritz Hoerold were even still alive and, if so, if he lived somewhere in Salt Lake City. Through a telephone call I found him and sent the magazine to him. He and his wife expressed to me their thanks. Having learned that Fritz had not as yet been ordained an elder and hence had never been to the temple, I wrote a letter encouraging him to qualify for the blessings of the temple. On two occasions we happened to see one another at restaurants. His dear wife, Joyce, always urged me, “Keep working with this man of mine.” His daughters voiced their approval of their mother’s urging. I kept up my encouragement.

Just a few weeks ago, I saw in the newspaper obituary columns that Joyce, Fritz’s wife, had passed away. How I wished that I had been more successful with my private project to get Fritz to the temple. I noted the time and place of Sister Hoerold’s funeral service, rescheduled other appointments, and went to the service. Immediately upon seeing me, Fritz made a beeline to my side. We both shed a few tears. He asked me to be the final speaker.

When I arose to speak, I looked at Fritz and his family and said, “Fritz, I am here today as the president of the teachers quorum of which you and I were once members.” I proffered how he and his family could become a “forever family” through temple ordinances—ordinances at which I pledged to officiate when that time came.

I concluded my remarks, choking back the tears of emotion, by saying to Fritz in the hearing and view of his family and all in attendance, “Fritz, my dear friend and fellow sailor, you have courage, you have determination. You put your life on the line for your country in a time of peril. Now, Fritz, you must heed the call of the bos’n’s whistle: ‘All aboard—anchors aweigh’—for your journey to exaltation. Joyce is there waiting for you. I know your dear children and grandchildren are praying for you. Fritz, as your teachers quorum president of long ago, I will strive with all my heart and soul to make certain you don’t miss the ship that will carry you and your loved ones to celestial glory.”

I gave to him a Navy salute. Fritz stood and returned the salute.

Brethren, may each of us be obedient to the easily remembered couplet: “Do your duty, that is best. Leave unto the Lord the rest.” For this I pray in the name of Jesus Christ, amen.

\newpage
\settalk{Now Is the Time}
\setspeaker{Thomas S. Monson}
\settalkdate{October 2001}
\talkheading{Now Is the Time}{Thomas S. Monson}

As I stand before you this morning, my thoughts return to the time of my youth when in Sunday School we often sang the lovely hymn:

\begin{verse}
Welcome, welcome, Sabbath morning;\\
Now we rest from ev’ry care.\\
Welcome, welcome is thy dawning,\\
Holy Sabbath, day of prayer.
\end{verse}

This Sabbath day I pray for an interest in your faith and prayers as I respond to the invitation to address you.

All of us have been dramatically affected by the tragic events of that fateful day, September 11, 2001. Suddenly, without warning, devastating destruction left death in its wake and snuffed out the lives of enormous numbers of men, women, and children. Evaporated were well-laid plans for pleasant futures. Substituted therefor were tears of sorrow and cries of pain from wounded souls.

Countless are the reports we have heard during the past three and a half weeks of those who were touched in some way—either directly or indirectly—by the events of that day. I should like to share with you the comments of a Church member, Rebecca Sindar, who was on a flight from Salt Lake City to Dallas on the morning of Tuesday, September 11. The flight was interrupted, as were all flights in the air at the time of the tragedies, and the plane grounded in Amarillo, Texas. Sister Sindar reports: “We all left the plane and found televisions in the airport, where we crowded around to see the broadcast of what had happened. People were lined up to call loved ones to assure them we were safely on the ground. I shall always remember the 12 or so missionaries who were on their way to the mission field on our flight. They made phone calls, and then we saw them huddled in a circle in a corner of the airport, kneeling in prayer together. How I wish I could have captured that moment to share with the mothers and fathers of those sweet young men as they saw the need for prayer right away.”

My brothers and sisters, death eventually comes to all mankind. It comes to the aged as they walk on faltering feet. Its summons is heard by those who have scarcely reached midway in life’s journey, and often it hushes the laughter of little children. Death is one fact that no one can escape or deny.

Frequently death comes as an intruder. It is an enemy that suddenly appears in the midst of life’s feast, putting out its lights and gaiety. Death lays its heavy hand upon those dear to us and at times leaves us baffled and wondering. In certain situations, as in great suffering and illness, death comes as an angel of mercy. But for the most part, we think of it as the enemy of human happiness.

The darkness of death can ever be dispelled by the light of revealed truth. “I am the resurrection, and the life,” spoke the Master. “He that believeth in me, though he were dead, yet shall he live: And whosoever liveth and believeth in me shall never die.”

This reassurance—yes, even holy confirmation—of life beyond the grave could well provide the peace promised by the Savior when He assured His disciples: “Peace I leave with you, my peace I give unto you: not as the world giveth, give I unto you. Let not your heart be troubled, neither let it be afraid.”

Out of the darkness and the horror of Calvary came the voice of the Lamb, saying, “Father, into thy hands I commend my spirit.” And the dark was no longer dark, for He was with His Father. He had come from God, and to Him He had returned. So also those who walk with God in this earthly pilgrimage know from blessed experience that He will not abandon His children who trust in Him. In the night of death, His presence will be “better than [a] light and safer than a known way.”

Saul, on the road to Damascus, had a vision of the risen, exalted Christ. Later, as Paul, defender of truth and fearless missionary in the service of the Master, he bore witness of the risen Lord as he declared to the Saints at Corinth: “Christ died for our sins according to the scriptures;

“… He was buried, and … he rose again the third day according to the scriptures:

“… He was seen of Cephas, then of the twelve:

“After that, he was seen of above five hundred brethren at once. …

“After that, he was seen of James; then of all the apostles.

“And last of all he was seen of me.”

In our dispensation this same testimony was spoken boldly by the Prophet Joseph Smith, as he and Sidney Rigdon testified:

“And now, after the many testimonies which have been given of him, this is the testimony, last of all, which we give of him: That he lives!

“For we saw him, even on the right hand of God; and we heard the voice bearing record that he is the Only Begotten of the Father—

“That by him, and through him, and of him, the worlds are and were created, and the inhabitants thereof are begotten sons and daughters unto God.”

This is the knowledge that sustains. This is the truth that comforts. This is the assurance that guides those bowed down with grief out of the shadows and into the light. It is available to all.

How fragile life; how certain death. We do not know when we will be required to leave this mortal existence. And so I ask, “What are we doing with today?” If we live only for tomorrow, we’ll have a lot of empty yesterdays today. Have we been guilty of declaring, “I’ve been thinking about making some course corrections in my life. I plan to take the first step—tomorrow”? With such thinking, tomorrow is forever. Such tomorrows rarely come unless we do something about them today. As the familiar hymn teaches:

\begin{verse}
There are chances for work all around just now,\\
Opportunities right in our way.\\
Do not let them pass by, saying, “Sometime I’ll try,”\\
But go and do something today.
\end{verse}

Let us ask ourselves the questions: “Have I done any good in the world today? Have I helped anyone in need?” What a formula for happiness! What a prescription for contentment, for inner peace—to have inspired gratitude in another human being.

Our opportunities to give of ourselves are indeed limitless, but they are also perishable. There are hearts to gladden. There are kind words to say. There are gifts to be given. There are deeds to be done. There are souls to be saved.

As we remember that “when ye are in the service of your fellow beings ye are only in the service of your God,” we will not find ourselves in the unenviable position of Jacob Marley’s ghost, who spoke to Ebenezer Scrooge in Dickens’ immortal A Christmas Carol. Marley spoke sadly of opportunities lost. Said he: “Not to know that any Christian spirit working kindly in its little sphere, whatever it may be, will find its mortal life too short for its vast means of usefulness. Not to know that no space of regret can make amends for one life’s opportunities misused! Yet such was I. Oh! such was I!”

Marley added: “Why did I walk through crowds of fellow-beings with my eyes turned down, and never raise them to that blessed Star which led the Wise Men to a poor abode? Were there no poor homes to which its light would have conducted me!”

Fortunately, as we know, Ebenezer Scrooge changed his life for the better. I love his line, “I am not the man I was.”

Why is the story A Christmas Carol so popular? Why is it ever new? I personally feel it is inspired of God. It brings out the best within human nature. It gives hope. It motivates change. We can turn from the paths which would lead us down and, with a song in our hearts, follow a star and walk toward the light. We can quicken our step, bolster our courage, and bask in the sunlight of truth. We can hear more clearly the laughter of little children. We can dry the tear of the weeping. We can comfort the dying by sharing the promise of eternal life. If we lift one weary hand which hangs down, if we bring peace to one struggling soul, if we give as did the Master, we can—by showing the way—become a guiding star for some lost mariner.

Because life is fragile and death inevitable, we must make the most of each day.

There are many ways in which we can misuse our opportunities. Some time ago I read a tender story written by Louise Dickinson Rich which vividly illustrates this truth. She wrote:

“My grandmother had an enemy named Mrs. Wilcox. Grandma and Mrs. Wilcox moved, as brides, into next-door houses on … Main Street of the tiny town in which they were to live out their lives. I don’t know what started the war [between them]—and I don’t think that by the time I came along, over 30 years later, they themselves remembered what started it. …

“… This was no polite sparring match. This was … total war. Nothing in town escaped repercussion. The 300-year-old church, which had lived through the Revolution, the Civil War, and the Spanish-American War, almost went down when Grandma and Mrs. Wilcox fought the Battle of the Ladies’ Aid. Grandma won that engagement, but it was a hollow victory. Mrs. Wilcox, since she couldn’t be president, resigned … in a huff. … What’s the fun of running a thing if you can’t force your … enemy to ‘eat crow’?

“Mrs. Wilcox won the Battle of the Public Library, getting her niece, Gertrude, appointed librarian instead of my Aunt Phyllis. The day Gertrude took over was the day Grandma stopped reading library books. [They became] ‘filthy germy things’ … overnight.

“The Battle of the High School was a draw. The principal got a better job and left before Mrs. Wilcox succeeded in having him ousted, or Grandma in having him given life tenure of office.

“… When, as children, we visited my grandmother, part of the fun was making faces at Mrs. Wilcox’s … grandchildren. … One banner day, we put a snake into the Wilcox rain barrel. My grandmother made token protests, but we sensed tacit sympathy. …

“Don’t think for a minute that this was a one-sided campaign. Mrs. Wilcox had grandchildren, too. … Grandma didn’t get off scot free. … Never a windy washday went by [that the clothesline didn’t mysteriously break, with the clothes falling in the dirt].

“I don’t know how Grandma could have borne her troubles so long if it hadn’t been for the household page of her daily Boston newspaper. This household page was a wonderful institution. Besides the usual cooking hints and cleaning advice, it had a department composed of letters from readers to each other. The idea was that if you had a problem—or even only some steam to blow off—you wrote a letter to the paper, signing some fancy name like Arbutus. That was Grandma’s pen name. Then some of the other ladies who had the same problem wrote back and told you what they had done about it, signing themselves One Who Knows or Xanthippe or whatever. Very often, the problem disposed of, you kept on for years writing to each other through the columns of the paper, telling each other about your children and your canning and your new dining room suite. That’s what happened to Grandma. She and a woman called Sea Gull corresponded for a quarter of a century. Sea Gull was Grandma’s true … friend.

“When I was about 16, Mrs. Wilcox died. In a small town, no matter how much you have hated your next-door neighbor, it is only common decency to run over and see what practical service you can do the bereaved.

“Grandma, neat in a percale apron to show that she meant what she said about being put to work, crossed the two lawns to the Wilcox house, where the Wilcox daughters set her to cleaning the already immaculate front parlor for the funeral. And there on the parlor table in the place of honor was a huge scrapbook, and in the scrapbook, pasted neatly in parallel columns were [Grandma’s] letters to Sea Gull over the years and Sea Gull’s letters to her. [Though neither woman had known it,] Grandma’s worst enemy had been her best friend.

“That was the only time I remembered seeing my grandmother cry. I didn’t know then exactly what she was crying about, but I do now. She was crying for all the wasted years that could never be salvaged.”

My brothers and sisters, may we resolve from this day forward to fill our hearts with love. May we go the extra mile to include in our lives any who are lonely or downhearted or who are suffering in any way. May we “[cheer] up the sad and [make] someone feel glad.” May we live so that when that final summons is heard, we may have no serious regrets, no unfinished business, but will be able to say with the Apostle Paul, “I have fought a good fight, I have finished my course, I have kept the faith.” In the name of Jesus Christ, amen.

\newpage
\settalk{Hidden Wedges}
\setspeaker{Thomas S. Monson}
\settalkdate{April 2002}
\talkheading{Hidden Wedges}{Thomas S. Monson}

In April 1966, at the Church’s annual general conference, Elder Spencer W. Kimball gave a memorable address. He quoted an account written by Samuel T. Whitman entitled “Forgotten Wedges.” Today I too have chosen to quote from Samuel T. Whitman, followed by examples from my own life.

Whitman wrote: “The ice storm [that winter] wasn’t generally destructive. True, a few wires came down, and there was a sudden jump in accidents along the highway. … Normally, the big walnut tree could easily have borne the weight that formed on its spreading limbs. It was the iron wedge in its heart that caused the damage.

“The story of the iron wedge began years ago when the white-haired farmer [who now inhabited the property on which it stood] was a lad on his father’s homestead. The sawmill had then only recently been moved from the valley, and the settlers were still finding tools and odd pieces of equipment scattered about. …

“On this particular day, it was a faller’s wedge—wide, flat, and heavy, a foot or more long, and splayed from mighty poundings,” which the lad found in the south pasture. A faller’s wedge, used to help fell a tree, is inserted in a cut made by a saw and then struck with a sledge hammer to widen the cut.

“Because he was already late for dinner, the lad laid the wedge … between the limbs of the young walnut tree his father had planted near the front gate. He would take the wedge to the shed right after dinner, or sometime when he was going that way.

“He truly meant to, but he never did. [The wedge] was there between the limbs, a little tight, when he attained his manhood. It was there, now firmly gripped, when he married and took over his father’s farm. It was half grown over on the day the threshing crew ate dinner under the tree. … Grown in and healed over, the wedge was still in the tree the winter the ice storm came.

“In the chill silence of that wintry night, … one of the three major limbs split away from the trunk and crashed to the ground. This so unbalanced the remainder of the top that it, too, split apart and went down. When the storm was over, not a twig of the once-proud tree remained.

“Early the next morning, the farmer went out to mourn his loss. …

“Then, his eyes caught sight of something in the splintered ruin. ‘The wedge,’ he muttered reproachfully. ‘The wedge I found in the south pasture.’ A glance told him why the tree had fallen. Growing, edge-up in the trunk, the wedge had prevented the limb fibers from knitting together as they should.”

My dear brothers and sisters, there are hidden wedges in the lives of many whom we know—yes, perhaps in our own families.

Let me share with you the account of a lifelong friend, now departed from mortality. His name was Leonard. He was not a member of the Church, although his wife and children were. His wife served as a Primary president; his son served an honorable mission. His daughter and his son married companions in solemn ceremonies and had families of their own.

Everyone who knew Leonard liked him, as did I. He supported his wife and children in their Church assignments. He attended many Church-sponsored events with them. He lived a good and a clean life, even a life of service and kindness. His family, and indeed many others, wondered why Leonard had gone through mortality without the blessings the gospel brings to its members.

In Leonard’s advanced years, his health declined. Eventually he was hospitalized, and life was ebbing away. In what turned out to be my last conversation with Leonard, he said, “Tom, I’ve known you since you were a boy. I feel persuaded to explain to you why I have never joined the Church.” He then related an experience of his parents which took place many, many years before. Reluctantly, the family had reached a point where they felt it was necessary to sell their farm, and an offer had been received. Then a neighboring farmer asked that the farm be sold to him instead—although at a lesser price—adding, “We’ve been such close friends. This way, if I own the property, I’ll be able to watch over it.” At length Leonard’s parents agreed, and the farm was sold. The buyer—even the neighbor—held a responsible position in the Church, and the trust this implied helped to persuade the family to sell to him, even though they did not realize as much money from the sale as they would have if they had sold to the first interested buyer. Not long after the sale was made, the neighbor sold both his own farm and the farm acquired from Leonard’s family in a combined parcel which maximized the value and hence the selling price. The long-asked question of why Leonard had never joined the Church had been answered. He always felt that his family had been deceived by the neighbor.

He confided to me following our conversation that he felt a great burden had at last been lifted as he prepared to meet his Maker. The tragedy is that a hidden wedge had kept Leonard from soaring to greater heights.

I am acquainted with a family which came to America from Germany. The English language was difficult for them. They had but little by way of means, but each was blessed with the will to work and with a love of God.

Their third child was born, lived but two months, and then died. Father was a cabinetmaker and fashioned a beautiful casket for the body of his precious child. The day of the funeral was gloomy, thus reflecting the sadness they felt in their loss. As the family walked to the chapel, with Father carrying the tiny casket, a small number of friends had gathered. However, the chapel door was locked. The busy bishop had forgotten the funeral. Attempts to reach him were futile. Not knowing what to do, the father placed the casket under his arm and, with his family beside him, carried it home, walking in a drenching rain.

If the family were of a lesser character, they could have blamed the bishop and harbored ill feelings. When the bishop discovered the tragedy, he visited the family and apologized. With the hurt still evident in his expression, but with tears in his eyes, the father accepted the apology, and the two embraced in a spirit of understanding. No hidden wedge was left to cause further feelings of anger. Love and acceptance prevailed.

The spirit must be freed from tethers so strong and feelings never put to rest, so that the lift of life may give buoyancy to the soul. In many families, there are hurt feelings and a reluctance to forgive. It doesn’t really matter what the issue was. It cannot and should not be left to injure. Blame keeps wounds open. Only forgiveness heals. George Herbert, an early 17th-century poet, wrote these lines: “He that cannot forgive others breaks the bridge over which he himself must pass if he would ever reach heaven, for every one has need of forgiveness.”

Beautiful are the words of the Savior as He was about to die upon the cruel cross. Said He, “Father, forgive them; for they know not what they do.”

There are some who have difficulty forgiving themselves and who dwell on all of their perceived shortcomings. I quite like the account of a religious leader who went to the side of a woman who lay dying, attempting to comfort her—but to no avail. “I am lost,” she said. “I’ve ruined my life and every life around me. There is no hope for me.”

The man noticed a framed picture of a lovely girl on the dresser. “Who is this?” he asked.

The woman brightened. “She is my daughter, the one beautiful thing in my life.”

“And would you help her if she were in trouble or had made a mistake? Would you forgive her? Would you still love her?”

“Of course I would!” cried the woman. “I would do anything for her. Why do you ask such a question?”

“Because I want you to know,” said the man, “that figuratively speaking, Heavenly Father has a picture of you on His dresser. He loves you and will help you. Call upon Him.”

A hidden wedge to her happiness had been removed.

In a day of danger or a time of trial, such knowledge, such hope, such understanding will bring comfort to the troubled mind and grieving heart. The entire message of the New Testament breathes a spirit of awakening to the human soul. Shadows of despair are dispelled by rays of hope, sorrow yields to joy, and the feeling of being lost in the crowd of life vanishes with the certain knowledge that our Heavenly Father is mindful of each of us.

The Savior provided assurance of this truth when He taught that even a sparrow shall not fall to the ground unnoticed by our Father. The Savior then concluded the beautiful thought by saying, “Fear ye not therefore, ye are of more value than many sparrows.”

Some time ago I read the following Associated Press dispatch, which appeared in the newspaper. An elderly man disclosed at the funeral of his brother, with whom he had shared, from early manhood, a small, one-room cabin near Canisteo, New York, that following a quarrel, they had divided the room in half with a chalk line and neither had crossed the line or spoken a word to the other since that day—62 years before. What a powerful and destructive hidden wedge.

As Alexander Pope wrote, “To err is human; to forgive, divine.”

Sometimes we can take offense so easily. On other occasions we are too stubborn to accept a sincere apology. Who will subordinate ego, pride, and hurt—then step forward with, “I am truly sorry! Let’s be as we once were: friends. Let’s not pass to future generations the grievances, the anger of our time.” Let’s remove any hidden wedges that can do nothing but destroy.

Where do hidden wedges originate? Some come from unresolved disputes, which lead to ill feelings, followed by remorse and regret. Others find their beginnings in disappointments, jealousies, arguments, and imagined hurts. We must solve them—lay them to rest and not leave them to canker, fester, and ultimately destroy.

A lovely lady of more than 90 years visited with me one day and unexpectedly recounted several regrets. She mentioned that many years earlier a neighboring farmer, with whom she and her husband had occasionally disagreed, asked if he could take a shortcut across her property to reach his own acreage. She paused in her narrative and, with a tremor in her voice, said, “Tommy, I didn’t let him cross our property but required him to take the long way around—even on foot—to reach his property. I was wrong and I regret it. He’s gone now, but oh, I wish I could say to him, ‘I’m so sorry.’ How I wish I had a second chance.”

As I listened to her, the words written by John Greenleaf Whittier came into my mind: “Of all sad words of tongue or pen, / The saddest are these: ‘It might have been!’”

From 3 Nephi in the Book of Mormon comes this inspired counsel: “There shall be no disputations among you. … For verily, verily I say unto you, he that hath the spirit of contention is not of me, but is of the devil, who is the father of contention, and he stirreth up the hearts of men to contend with anger, one with another. Behold, this is not my doctrine, to stir up the hearts of men with anger, one against another; but this is my doctrine, that such things should be done away.”

Let me conclude with an account of two men who are heroes to me. Their acts of courage were not performed on a national scale, but rather in a peaceful valley known as Midway, Utah.

Long years ago, Roy Kohler and Grant Remund served together in Church capacities. They were the best of friends. They were tillers of the soil and dairymen. Then a misunderstanding arose which became somewhat of a rift between them.

Later, when Roy Kohler became grievously ill with cancer and had but a limited time to live, my wife, Frances, and I visited Roy and his wife, and I gave him a blessing. As we talked afterward, Brother Kohler said, “Let me tell you about one of the sweetest experiences I have had during my life.” He then recounted to me his misunderstanding with Grant Remund and the ensuing estrangement. His comment was, “We were sort of on the outs with each other.”

“Then,” continued Roy, “I had just put up our hay for the winter to come, when one night, as a result of spontaneous combustion, the hay caught fire, burning the hay, the barn, and everything in it right to the ground. I was devastated,” said Roy. “I didn’t know what in the world I would do. The night was dark, except for the dying embers of the fire. Then I saw coming toward me from the road, in the direction of Grant Remund’s place, the lights of tractors and heavy equipment. As the ‘rescue party’ turned in our drive and met me amidst my tears, Grant said, ‘Roy, you’ve got quite a mess to clean up. My boys and I are here. Let’s get to it.’” Together they plunged to the task at hand. Gone forever was the hidden wedge which had separated them for a short time. They worked throughout the night and into the next day, with many others in the community joining in.

Roy Kohler has passed away, and Grant Remund is getting older. Their sons have served together in the same ward bishopric. I truly treasure the friendship of these two wonderful families.

May we ever be exemplary in our homes and faithful in keeping all of the commandments, that we may harbor no hidden wedges but rather remember the Savior’s admonition: “By this shall all men know that ye are my disciples, if ye have love one to another.”

This is my plea and my prayer, in the name of Jesus Christ, amen.

\newpage
\settalk{Pathways to Perfection}
\setspeaker{Thomas S. Monson}
\settalkdate{April 2002}
\talkheading{Pathways to Perfection}{Thomas S. Monson}

Our Young Women presidency have done so well, haven’t they? I sustain and endorse all that you have heard from these splendid women today. They are truly servants of our Heavenly Father and have presented His holy word.

“Happiness,” the Prophet Joseph Smith wrote, “is the object and design of our existence; and will be the end thereof, if we pursue the path that leads to it; and this path is virtue, uprightness, faithfulness, holiness, and keeping all the commandments of God.”

But how does one find that pathway, and what’s more, how does one stay on that pathway which leads to perfection?

In Lewis Carroll’s classic Alice’s Adventures in Wonderland, Alice finds herself coming to a crossroads with two paths before her, each stretching onward but in opposite directions. She is confronted by the Cheshire Cat, of whom she asks, “Which path shall I take?”

The cat answers: “That depends where you want to go. If you do not know where you want to go, it doesn’t really matter which path you take!”

Unlike Alice, each of you knows where you want to go. It does matter which way you go, for the path you follow in this life leads to the path you will follow in the next.

A lilting ballad, popular many years ago, contains the provocative line, “If wishing can make it so, then keep on wishing and cares will go.” Another formula for failure comes from the more recent song, “Don’t worry; be happy!”

Our theme for this evening, “Stand Ye in Holy Places,” is more appropriate. I also appreciate the words which follow: “Stand ye in holy places, and be not moved.”

President George Albert Smith, eighth President of the Church, urged: “Let us plant our feet upon the highway that leads to happiness and the celestial kingdom, not just occasionally, but every day, and every hour, because if we will stay on the Lord’s side of the line, if we will remain under the influence of our Heavenly Father, the adversary cannot even tempt us. But if we go into the devil’s territory … we will be unhappy and that unhappiness will increase as the years go by, unless we repent of our sins and turn to the Lord.”

In speaking to young men of the Aaronic Priesthood, I have frequently quoted the advice of a father to a precious son: “If you ever find yourself where you hadn’t ought to be—then get out!” The same truth is applicable to you young women here in the Conference Center and to you assembled in meetinghouses throughout the world.

I have always felt that if we speak in generalities, we rarely have success; but if we speak in specifics, we will rarely have a failure. Therefore, I urge that you exemplify in your lives four tested, specific virtues. They are:

First, an attitude of gratitude. In the book of Luke, chapter 17, we read the account of the 10 lepers. The Savior, in traveling toward Jerusalem, passed through Galilee and Samaria and entered a certain village where He was met on the outskirts by 10 lepers who were forced, because of their condition, to live apart from others. They stood “afar off” and cried, “Jesus, Master, have mercy on us.”

The Savior, full of sympathy and love for them, said, “Go shew yourselves unto the priests,” and as they went they discovered that they were healed. The scriptures tell us, “One of them, when he saw that he was healed, turned back, and with a loud voice glorified God, and fell down on his face at [the Master’s] feet, giving him thanks: and he was a Samaritan.”

The Savior responded, “Were there not ten cleansed? but where are the nine? There are not found that returned to give glory to God, save this stranger. And he said unto him, Arise, go thy way: thy faith hath made thee whole.”

Through divine intervention, those who were lepers were spared from a cruel, lingering death and given a new lease on life. The gratitude expressed by one merited the Master’s blessing, the ingratitude by the nine His disappointment.

Like the leprosy of yesteryear are the plagues of today. They linger; they debilitate; they destroy. They are to be found everywhere. Their pervasiveness knows no boundaries. We know them as selfishness, greed, indulgence, cruelty, and crime—to identify but a few.

At a regional conference, President Gordon B. Hinckley declared: “We live in a world of so much filth. It is everywhere. It is on the streets. It is on television. It is in books and magazines. … It is like a great flood, ugly and dirty and mean, engulfing the world. We have got to stand above it. … The world is slipping in its moral standards. That can only bring misery. The way to happiness lies in a return to strong family life and the observance of moral standards, the value of which has been proven through centuries of time.”

By following President Hinckley’s counsel, we can make this a wonderful time to be living here on earth. Our opportunities are limitless. There are so many things right—such as teachers who teach, friends who help, marriages that make it, and parents who sacrifice.

Be grateful for your mother, for your father, for your family, and for your friends. Express gratitude for your Young Women teachers. They love you; they pray for you; they serve you. You are precious in their sight and in the sight of your Heavenly Father. He hears your prayers. He extends to you His peace and His love. Stay close to Him and to His Son, and you will not walk alone.

Second, a longing for learning.

The Apostle Paul said to Timothy, “Let no man despise thy youth; but be thou an example of the believers.”

President Stephen L Richards, who was a counselor in the First Presidency many years ago, was a profound thinker. He said, “Faith and doubt cannot exist in the same mind at the same time, for one will dispel the other.” My advice is to seek faith and dispel doubt.

The Lord counseled, “Seek ye out of the best books words of wisdom; seek learning, even by study and also by faith.”

We can find truth in the scriptures, the teachings of the prophets, the instructions from our parents, and the inspiration that comes to us as we bend our knees and seek the help of God.

We must be true to our ideals, for ideals are like the stars: you can’t touch them with your hands, but by following them you reach your destination.

Many of your teachers are assembled with you this evening. I trust that each teacher would fit the description written of one: “She created in her classroom an atmosphere where warmth and acceptance weave their magic spell; where growth and learning, the soaring of the imagination, and the spirit of the young are assured.”

Third, may we discuss a devotion to discipline.

Our Heavenly Father has given to each of us the power to think and reason and decide. With such power, self-discipline becomes a necessity.

Each of us has the responsibility to choose. You may ask, “Are decisions really that important?” I say to you, decisions determine destiny. You can’t make eternal decisions without eternal consequences.

May I provide a simple formula by which you can measure the choices which confront you. It’s easy to remember: “You can’t be right by doing wrong; you can’t be wrong by doing right.” Your personal conscience always warns you as a friend before it punishes you as a judge.

The Lord, in a revelation given through Joseph Smith the Prophet, counseled: “That which doth not edify is not of God, and is darkness. That which is of God is light.”

Some foolish persons turn their backs on the wisdom of God and follow the allurement of fickle fashion, the attraction of false popularity, and the thrill of the moment. Courage is required to think right, choose right, and do right, for such a course will rarely, if ever, be the easiest to follow.

The battle for self-discipline may leave you a bit bruised and battered but always a better person. Self-discipline is a rigorous process at best; too many of us want it to be effortless and painless. Should temporary setbacks afflict us, a very significant part of our struggle for self-discipline is the determination and the courage to try again.

My dear young sisters, I know of no truer description of you than that expressed by the First Presidency on April 6, 1942: “How glorious and near to the angels is youth that is clean; this youth has joy unspeakable here and eternal happiness hereafter.”

Eternal life in the kingdom of our Father is your goal, and self-discipline will surely be required if you are to achieve it.

Finally, let each of us cultivate a willingness to work. President J. Reuben Clark, many years ago a counselor in the First Presidency, said: “I believe that we are here to work, and I believe there is no escape from it. I think that we cannot get that thought into our souls and into our beings too soon. Work we must, if we shall succeed or if we shall advance. There is no other way.”

“Put your shoulder to the wheel, push along” is more than a line from a favorite hymn; it is a summons to work.

Perhaps an example would be helpful. Procrastination is truly a thief of time—especially when it comes to downright hard work. I speak of the need to study diligently as you prepare for the tests of school and, indeed, the tests of life.

I know of a university student who was so busy with the joys of student life that preparation for an exam was postponed. The night before, she realized the hour was late and the preparation was not done. She rationalized, “Now what is more important—my health, which requires that I must sleep, or the drudgery of study?” Well, you can probably guess the outcome. Sleep won, study failed, and the test was a personal disaster. Work we must.

This, then, is the suggested formula:

There will come into every life moments of despair and the need for direction from a divine source—even an unspoken plea for help. With all my heart and soul I testify to you that our Heavenly Father loves you, is mindful of you, and will not abandon you.

Let me illustrate with a personal and treasured experience. For many years my assignments took me into that part of Germany which was behind what was called the Iron Curtain. Under Communist control, those who lived in that area of Germany had lost nearly all of their freedoms. Activities of youth were restricted; all actions were monitored.

Shortly after I assumed my responsibilities for that area, I attended a most uplifting conference held in that part of Germany. Following the inspirational songs and the spoken word, I felt the impression to meet briefly outside of the old building with the precious teenage youth. They were relatively few in number but listened to every word I spoke. They had hungered for the word and encouragement of an Apostle of the Lord.

Prior to attending the conference, before leaving the United States, I felt the prompting to buy three cartons of chewing gum. I purchased three flavors: Doublemint, Spearmint, and Juicy Fruit. Now, as the gathering of the youth was concluded, I distributed carefully to each youth two sticks of gum—something they had never before tasted. They received the gift with joy.

The years went by. I returned to Dresden—the site of our earlier conference. Now we had chapels; now the people had freedom. They had a temple. Germany was no longer separated by political boundaries but had become one nation. The youth were now adults with children of their own.

Following a large and inspirational conference, a mother and her daughter sought me out to speak to me. The daughter, who was about your age and who spoke some English, said to me, “President Monson, do you remember long ago holding a brief gathering of youth following a district conference, where you gave to each boy and each girl two sticks of chewing gum?”

I responded, “Oh, yes, I surely do remember.”

She continued, “My mother was one to whom you gave that gift. She told me that she rationed in little pieces one stick of gum. She mentioned how sweet to the taste it was and so precious to her.” Then, under the approving smile of her dear mother, she handed to me a small box. As I opened the lid of the box, there I beheld the other stick of gum, still with its wrapper after nearly 20 years. And then she said, “My mother and I want you to have this,” she said.

The tears flowed; embraces followed.

The mother then spoke to me: “Before you came to our conference so many years ago, I had prayed to my Heavenly Father to know that He indeed cared about me. I saved that gift so that I might remember and teach my daughter that Heavenly Father does hear our prayers.”

I hold before you tonight that gift—even a symbol of faith and assurance of the heavenly help our Father and His Son, Jesus Christ, will provide you.

On this Easter eve, may our thoughts turn to Him who atoned for our sins, who showed us the way to live, how to pray, and who demonstrated by His own actions how we might do so. Born in a stable, cradled in a manger, this Son of God—even Jesus Christ the Lord—beckons to each of us to follow Him. “Oh, sweet the joy this sentence gives: ‘I know that my Redeemer lives!’” In the name of Jesus Christ, amen.

\newpage
\settalk{They Pray and They Go}
\setspeaker{Thomas S. Monson}
\settalkdate{April 2002}
\talkheading{They Pray and They Go}{Thomas S. Monson}

My brethren, I am honored by the privilege to speak to you this evening. What a joy to see this magnificent Conference Center filled to its capacity with those young and old who hold the priesthood of God. To realize that similar throngs are assembled throughout the world brings to me an overwhelming sense of responsibility. I pray that the inspiration of the Lord will guide my thoughts and inspire my words.

Many years ago, on an assignment to Tahiti, I was talking to our mission president, President Raymond Baudin, about the Tahitian people. They are known as some of the greatest seafaring people in all the world. Brother Baudin, who speaks French and Tahitian but little English, was trying to describe to me the secret of the success of the Tahitian sea captains. He said, “They are amazing. The weather may be terrible, the vessels may be leaky, there may be no navigational aids except their inner feelings and the stars in the heavens, but they pray and they go.” He repeated that phrase three times. There is a lesson in that statement. We need to pray, and then we need to act. Both are important.

The promise from the book of Proverbs gives us courage:

“Trust in the Lord with all thine heart; and lean not unto thine own understanding. In all thy ways acknowledge him, and he shall direct thy paths.”

We need but to turn to the account found in 1 Kings to appreciate anew the principle that when we follow the counsel of the Lord, when we pray and then go, the outcome benefits all. There we read that a most severe drought had gripped the land. Famine followed. Elijah the prophet received from the Lord what to him must have been an amazing instruction: “Get thee to Zarephath … : behold, I have commanded a widow woman there to sustain thee.” When he had found the widow, Elijah declared, “Fetch me, I pray thee, a little water in a vessel, that I may drink.

“And as she was going to fetch it, he called to her, and said, Bring me, I pray thee, a morsel of bread in thine hand.”

Her response described her desperate situation as she explained that she was preparing a final and scanty meal for her son and for herself, and then they would die.

How implausible to her must have been Elijah’s response: “Fear not; go and do as thou hast said: but make me thereof a little cake first, and bring it unto me, and after make for thee and for thy son.

“For thus saith the Lord God of Israel, The barrel of meal shall not waste, neither shall the cruse of oil fail, until the day that the Lord sendeth rain upon the earth.

“And she went and did according to the saying of Elijah: and she, and he, and her house, did eat many days.

“And the barrel of meal wasted not, neither did the cruse of oil fail.”

If I were to ask you which of all the passages in the Book of Mormon is the most widely read, I venture it would be the account found in 1 Nephi concerning Nephi, his brothers, his father, and the command to obtain from Laban the plates of brass. Perhaps this is because most of us, from time to time, pledge to again read the Book of Mormon. Usually we begin with 1 Nephi. In reality, the passages found therein portray beautifully the need to pray and then to go and do. Said Nephi, “I will go and do the things which the Lord hath commanded, for I know that the Lord giveth no commandments unto the children of men, save he shall prepare a way for them that they may accomplish the thing which he commandeth them.”

We remember the commandment. We remember the response. We remember the outcome.

In our day and our time, there are many examples concerning the experiences of those who pray and then go and do. I share with you a touching account of a fine family that lived in the beautiful city of Perth, Australia. In 1957, four months before the dedication of the New Zealand temple, Donald Cummings, the father, was the president of the member district in Perth. He and his wife and family were determined to attend the dedication of the temple, although they were of very modest financial means. They began to pray, to work, and to save. They sold their only car and gathered together every penny they could, but a week before their scheduled departure, they were still 200 pounds short. Through two unexpected gifts of 100 pounds each, they met their goal just in time. Because Brother Cummings couldn’t get time off work for the trip, he decided to quit his job.

They traveled by train across the vast Australian continent, arriving at Sydney, where they joined other members also traveling to New Zealand. Brother Cummings and his family were among the first Australians to be baptized for the dead in the New Zealand temple. They were among the first ones to be endowed in the New Zealand temple from far-off Perth, Australia. They prayed, they prepared, and then they went.

When the Cummings family returned to Perth, Brother Cummings obtained a new and better job. He was still serving as district president nine years later when it was my privilege to call him as the first president of the Perth Australia Stake. I think it significant that he is now the first president of the Perth Australia Temple.

From the movie Shenandoah come the spoken words which inspire: “If we don’t try, we don’t do; and if we don’t do, then why are we here?”

There are now more than 60,000 full-time missionaries serving the Lord throughout the world. Many of this vast throng are listening in tonight and viewing this priesthood session of general conference. They pray and then they go, trusting in the Lord concerning where they are sent and trusting in their mission president as to where they serve within their missions. Among the many revelations concerning their sacred callings are two passages which are favorites of mine. Both are from the Doctrine and Covenants.

The first is from section 100. You will remember that Joseph Smith and Sidney Rigdon had been absent from their families for some time, and they were concerned about them. The Lord revealed unto them this assurance, which is comforting to missionaries throughout the Church: “Verily, thus saith the Lord unto you, my friends … , your families are well; they are in mine hands, and I will do with them as seemeth me good; for in me there is all power.”

The second is from the 84th section of the Doctrine and Covenants: “Whoso receiveth you, there I will be also, for I will go before your face. I will be on your right hand and on your left, and my Spirit shall be in your hearts, and mine angels round about you, to bear you up.”

Inspiring is the missionary service rendered by Walter Krause, who lives in Prenzlau, Germany. Brother Krause, whose dedication to the Lord is legendary, is now 92 years of age. As a patriarch, he has given more than a thousand patriarchal blessings to members living throughout many parts of Europe.

Homeless following World War II, like so many others at that time, Brother Krause and his family lived in a refugee camp in Cottbus and began to attend church there. He was immediately called to lead the Cottbus branch. Four months later, in November of 1945, the country still in ruins, district president Richard Ranglack came to Brother Krause and asked him what he would think about going on a mission. Brother Krause’s answer reflects his commitment to the Church. Said he: “I don’t have to think about it at all. If the Lord needs me, I’ll go.”

He set out on December 1, 1945, with 20 German marks in his pocket and a piece of dry bread. One of the branch members had given him a winter coat left over from a son who had fallen in the war. Another member, who was a shoemaker, gave him a pair of shoes. With these and with two shirts, two handkerchiefs, and two pairs of stockings, he left on his mission.

Once, in the middle of winter, he walked from Prenzlau to Kammin, a little village in Mecklenburg, where 46 attended the meetings which were held. He arrived long after dark that night after a six-hour march over roads, paths, and finally across plowed fields. Just before he reached the village, he came to a large, white, flat area which made for easy walking, and he soon arrived at a member’s home to stay the night.

The next morning the game warden knocked on the door of the member’s house, asking, “Do you have a guest?”

“Yes,” came the reply.

The game warden continued, “Then come and take a look at his tracks.” The large, flat area on which Brother Krause had walked was actually a frozen lake, and some time earlier the warden had chopped a large hole in the middle of the lake for fishing. The wind had driven snow over the hole and covered it so that Brother Krause could not have seen his danger. His tracks went right next to the edge of the hole and straight to the house of the member, without his knowing anything about it. Weighed down by his backpack and his rubber boots, he would certainly have drowned had he gone one step further toward the hole he couldn’t see. He commented later that this event caused quite a stir in the village at the time.

Brother Krause’s entire life has been to pray and then to go.

Should any of us feel inadequate or tend to doubt the ability to respond to a priesthood call to serve the Lord, let this divine truth be remembered: “With God all things are possible.”

Not long ago I learned of the passing of James Womack, the patriarch of the Shreveport Louisiana Stake. He had served long and had blessed ever so many lives. Years before, President Spencer W. Kimball shared with President Gordon B. Hinckley, Elder Bruce R. McConkie, and me an experience he had in the appointment of a patriarch for the Shreveport Louisiana Stake of the Church. President Kimball described how he interviewed, how he searched, and how he prayed that he might learn the Lord’s will concerning the selection. For some reason, none of the suggested candidates was the man for this assignment at this particular time.

The day wore on; the evening meetings began. Suddenly President Kimball turned to the stake president and asked him to identify a particular man seated perhaps two-thirds of the way back from the front of the chapel. The stake president replied that the individual was James Womack, whereupon President Kimball said: “He is the man the Lord has selected to be your stake patriarch. Please have him meet with me in the high council room following the meeting.”

Stake president Charles Cagle was startled, for James Womack did not wear the label of a typical man. He had sustained terrible injuries while in combat during World War II. He lost both hands and part of an arm, as well as most of his eyesight and part of his hearing. Nobody had wanted to let him into law school when he returned, yet he finished third in his class at Louisiana State University.

That evening as President Kimball met with Brother Womack and informed him that the Lord had designated him to be the patriarch, there was a protracted silence in the room. Then Brother Womack said: “Brother Kimball, it is my understanding that a patriarch is to place his hands on the head of the person he blesses. As you can see, I have no hands to place on the head of anyone.”

Brother Kimball, in his kind and patient manner, invited Brother Womack to stand behind the chair on which Brother Kimball was seated. He then said, “Now, Brother Womack, lean forward and see if the stumps of your arms will reach the top of my head.” To Brother Womack’s joy, they touched Brother Kimball’s head, and the exclamation came forth, “I can reach you! I can reach you!”

“Of course you can reach me,” responded Brother Kimball. “And if you can reach me, you can reach any whom you bless. I will probably be the shortest person you will ever have seated before you.”

President Kimball reported to us that when the name of James Womack was presented to the stake conference, “the hands of the members shot heavenward in an enthusiastic vote of approval.”

Remembered were the words of the Lord to the prophet Samuel at the time David was designated to be a future king of Israel: “Man looketh on the outward appearance, but the Lord looketh on the heart.”

Brethren, whatever our calling, regardless of our fears or anxieties, let us pray and then go and do, remembering the words of the Master, even the Lord Jesus Christ, who promised, “I am with you alway, even unto the end of the world.”

In the epistle of James we are counseled, “Be ye doers of the word, and not hearers only, deceiving your own selves.”

Let us, as a mighty body of priesthood, be doers of the word, and not hearers only. Let us pray; then let us go and do.

In the name of Jesus Christ, amen.

\newpage
\settalk{Models to Follow}
\setspeaker{Thomas S. Monson}
\settalkdate{October 2002}
\talkheading{Models to Follow}{Thomas S. Monson}

Many years ago I marveled at the cover of one of our Church publications which featured a magnificent reproduction of a Carl Bloch painting. The scene which the artist captured in his mind and then—with a touch of the Master’s hand—transferred to canvas depicted Elisabeth, wife of Zacharias, welcoming Mary, the mother of Jesus. Both were to bear sons—miracle births.

The son born of Elisabeth became known as John the Baptist. As with Jesus, son of Mary, so with John—precious little is recorded of their years of youth. A single sentence tells us all that we know of John’s life from his birth to his public ministry: “And the child grew, and waxed strong in spirit, and was in the deserts till the day of his shewing unto Israel.”

John’s message was brief. He preached faith, repentance, baptism by immersion, and the bestowal of the Holy Ghost by an authority greater than that possessed by himself. “I am not the Christ,” he told his faithful disciples, “but … I am sent before him.” “I indeed baptize you with water; but one mightier than I cometh … : he shall baptize you with the Holy Ghost and with fire.”

Then occurred the baptism of Christ by John the Baptist. Later Jesus testified, “Among them that are born of women there hath not risen a greater than John the Baptist.”

All of us living in the world today need points of reference, even models to follow. John the Baptist provides for us a flawless example of unfeigned humility, as he deferred always to the One who was to follow—the Savior of mankind.

As we learn of others who trusted God and followed His teachings, the Spirit whispers to our souls, “Be still, and know that I am God.” As they resolutely kept His commandments and trusted in Him, they were blessed. When we follow their examples, we too will be similarly blessed in our day and in our time. Each one becomes a model to follow.

All of us love the beautiful account from the Holy Bible of Abraham and Isaac. How exceedingly difficult it must have been for Abraham, in obedience to God’s command, to take his beloved Isaac into the land of Moriah, there to present him as a burnt offering. Can you imagine the heaviness of his heart as he gathered the wood for the fire and journeyed to the appointed place? Surely pain must have racked his body and tortured his mind as he “bound Isaac … and laid him on the altar upon the wood [and] stretched forth his hand, and took the knife to slay his son.” How glorious was the pronouncement, and with what wondered welcome did it come: “Lay not thine hand upon the lad, neither do thou any thing unto him: for now I know that thou fearest God, seeing thou hast not withheld thy son, thine only son from me.”

Abraham qualifies as a model of unquestioning obedience.

If any of us feels his challenges are beyond his capacity to meet them, let him or her read of Job. By so doing, there comes the feeling, “If Job could endure and overcome, so will I.”

Job was a “perfect and upright” man who “feared God, and eschewed evil.” Pious in his conduct, prosperous in his fortune, Job was to face a test which could have destroyed anyone. Shorn of his possessions, scorned by his friends, afflicted by his suffering, shattered by the loss of his family, he was urged to “curse God, and die.” He resisted this temptation and declared from the depths of his noble soul, “Behold, my witness is in heaven, and my record is on high.” “I know that my redeemer liveth.”

Job became a model of unlimited patience. To this day we refer to those who are long-suffering as having the patience of Job. He provides an example for us to follow.

“A just man and perfect in his generations,” one who “walked with God,” was the prophet Noah. Ordained to the priesthood at an early age, “he became a preacher of righteousness and declared the gospel of Jesus Christ, … teaching faith, repentance, baptism, and the reception of the Holy Ghost.” He warned that failure to heed his message would bring floods upon those who heard his voice, and yet they hearkened not to his words.

Noah heeded God’s command to build an ark, that he and his family might be spared destruction. He followed God’s instructions to gather into the ark two of every living creature, that they also might be saved from the floodwaters.

Said President Spencer W. Kimball: “As yet there was no evidence of rain and flood. … [Noah’s] warnings were considered irrational. … How foolish to build an ark on dry ground with the sun shining and life moving forward as usual! But time ran out. … The floods came. The disobedient … were drowned. The miracle of the ark followed the faith manifested in its building.”

Noah had the unwavering faith to follow God’s commandments. May we ever do likewise. May we remember that the wisdom of God ofttimes appears as foolishness to men, but the greatest lesson we can learn in mortality is that when God speaks and we obey, we will always be right.

A model of ideal womanhood is Ruth. Sensing the grief-stricken heart of her mother-in-law Naomi—who suffered the loss of each of her two fine sons—feeling perhaps the pangs of despair and loneliness that plagued the very soul of Naomi, Ruth uttered what has become that classic statement of loyalty: “Intreat me not to leave thee, or to return from following after thee: for whither thou goest, I will go; and where thou lodgest, I will lodge: thy people shall be my people, and thy God my God.” Ruth’s actions demonstrated the sincerity of her words.

Through Ruth’s undeviating loyalty to Naomi, she was to marry Boaz, by which she—the foreigner and Moabite convert—became a great-grandmother of David and, therefore, an ancestor of our Savior, Jesus Christ.

I now turn to a mighty Book of Mormon prophet—even Nephi, son of Lehi and Sariah. He was faithful and obedient to God, courageous and bold. When given the difficult assignment to obtain the plates of brass from Laban, he did not murmur but declared, “I will go and do the things which the Lord hath commanded, for I know that the Lord giveth no commandments unto the children of men, save he shall prepare a way for them that they may accomplish the thing which he commandeth them.” Perhaps this act of courage prompted a verse of counsel for us found in the hymn “The Iron Rod”:

\begin{verse}
To Nephi, seer of olden time,\\
A vision came from God. …\\
Hold to the rod, the iron rod;\\
’Tis strong, and bright, and true.\\
The iron rod is the word of God;\\
’Twill safely guide us through.
\end{verse}

Nephi personified unflagging determination.

No description of models for us to follow would be complete without including Joseph Smith, the first prophet of this dispensation. When but 14 years of age, this courageous young man entered a grove of trees, which later would be called sacred, and received an answer to his sincere prayer.

There followed for Joseph unrelenting persecution as he related to others the account of the glorious vision he received in that grove. Yet, although he was ridiculed and scorned, he stood firm. Said he, “I had seen a vision; I knew it, and I knew that God knew it, and I could not deny it, neither dared I do it.”

Step by step, facing opposition at nearly every turn and yet always guided by the hand of the Lord, Joseph organized The Church of Jesus Christ of Latter-day Saints. He proved courageous in all that he did.

Toward the end of his life, as he was led away with his brother Hyrum to Carthage Jail, he bravely faced what he undoubtedly knew lay ahead for him, and he sealed his testimony with his blood.

As we face life’s tests, may we ever emulate that undaunted courage epitomized by the Prophet Joseph Smith.

There stands before us today another prophet of God—even our beloved President Gordon B. Hinckley. He has presided over the largest expansion of the Church—both numerically and geographically—in our history. He has traversed frontiers not heretofore crossed and has visited with government leaders and with members the world over. His love for the people transcends the barriers of language and culture.

With prophetic vision, he has instituted the Perpetual Education Fund, which breaks the cycle of poverty for our members in many areas of the world and provides skills and training which qualify young men and young women for gainful employment. This inspired plan has kindled the light of hope in the eyes of those who felt doomed to mediocrity but who now have an opportunity for a brighter future.

President Hinckley has labored unceasingly to bring sacred blessings to members of the Church worldwide by providing temples that are within the reach of all. He has the capacity to lift to a higher plane those from all walks of life, regardless of the faith to which they ascribe. He is a model of unfailing optimism, and we revere him as prophet, seer, and revelator.

The unique qualities possessed by these men and women whom I have mentioned can be of invaluable assistance to us as we face the problems and the trials which lie ahead. May I illustrate by mentioning the experience of the Jerome Kenneth Pollard family of Oakland, California.

This past May, as Elder Taavili Joseph Samuel Pollard was traveling to the mission office on the last day of his mission in Zimbabwe, the mission car he was driving somehow spun out of control and hit a tree. A passerby was able to rescue Elder Pollard’s companion, but Elder Pollard, who was unconscious, was trapped in the car, which burst into flames. Elder Pollard perished. His mother had passed away eight years earlier; hence, his father was rearing the family alone. A brother was serving in the West Indies Mission.

When the news of Elder Pollard’s death reached his father, this humble man—who had already lost his wife—called the son serving in the West Indies Mission to let him know of his brother’s death. Over that long-distance telephone line, Brother Pollard and his son, no doubt grief stricken and heartsick, sang together “I Am a Child of God.” Before concluding the call, the father offered a prayer to Heavenly Father, thanking Him for His blessings and seeking His divine comfort.

Brother Pollard later commented that he knew his family would be all right, for they have strong testimonies of the gospel and of the plan of salvation.

My brothers and sisters, in this marvelous dispensation of the fulness of times, as we journey through mortality and face the trials and challenges of the future, may we remember the examples of these models to follow which I have referred to this morning. May we have the unfeigned humility of John the Baptist, the unquestioning obedience of Abraham, the unlimited patience of Job, the unwavering faith of Noah, the undeviating loyalty of Ruth, the unflagging determination of Nephi, the undaunted courage of the Prophet Joseph Smith, and the unfailing optimism of President Hinckley. Such will be as a fortress of strength to us throughout our lives.

May we ever be guided by the supreme Exemplar, even the son of Mary, the Savior Jesus Christ—whose very life provided a perfect model for us to follow.

Born in a stable, cradled in a manger, He came forth from heaven to live on earth as a mortal man and to establish the kingdom of God. During His earthly ministry, He taught men the higher law. His glorious gospel reshaped the thinking of the world. He blessed the sick; He caused the lame to walk, the blind to see, the deaf to hear. He even raised the dead to life.

What was the reaction to His message of mercy, His words of wisdom, His lessons of life? There were a precious few who appreciated Him. They bathed His feet. They learned His word. They followed His example.

Then there were those who denied Him. When asked by Pilate, “What shall I do then with Jesus which is called Christ?” they cried, “Crucify him.” They mocked Him. They gave Him vinegar to drink. They reviled Him. They smote Him with a reed. They did spit upon Him. They crucified Him.

Down through the generations of time, the message from Jesus has been the same. To Peter and Andrew by the shores of the beautiful Sea of Galilee, He said, “Follow me.” To Philip of old came the call, “Follow me.” To the Levite who sat at receipt of customs came the instruction, “Follow me.” And to you and to me, if we but listen, will come that same beckoning invitation, “Follow me.”

My prayer today is that we shall do so. In the sacred name of Jesus Christ, amen.

\newpage
\settalk{Peace, Be Still}
\setspeaker{Thomas S. Monson}
\settalkdate{October 2002}
\talkheading{Peace, Be Still}{Thomas S. Monson}

The singing of the men’s choir this evening has lighted memory’s fire and brought to my mind the songs I sang when I was a boy. With fervor we would render:

\begin{verse}
Put your shoulder to the wheel; push along.\\
Do your duty with a heart full of song.\\
We all have work; let no one shirk.\\
Put your shoulder to the wheel.
\end{verse}

We had a chorister who taught us boys how to sing. We had to sing. Sister Stella Waters would wave the baton within inches of our noses and beat time with a heavy foot that made the floor creak.

If we responded properly, Sister Waters let us choose a favorite hymn to sing. Inevitably, the selection was:

\begin{verse}
Master, the tempest is raging!\\
The billows are tossing high!\\
The sky is o’ershadowed with blackness.\\
No shelter or help is nigh.\\
Carest thou not that we perish?\\
How canst thou lie asleep\\
When each moment so madly is threat’ning\\
A grave in the angry deep?
\end{verse}

And then the assuring chorus:

\begin{verse}
The winds and the waves shall obey thy will:\\
Peace, be still, peace, be still.\\
Whether the wrath of the storm-tossed sea\\
Or demons or men or whatever it be,\\
No waters can swallow the ship where lies\\
The Master of ocean and earth and skies.\\
They all shall sweetly obey thy will:\\
Peace, be still; peace, be still.\\
They all shall sweetly obey thy will:\\
Peace, peace, be still.
\end{verse}

As a boy, I could fathom somewhat the danger of a storm-tossed sea. However, I had but little understanding of other demons which can stalk our lives, destroy our dreams, smother our joys, and detour our journey toward the celestial kingdom of God.

A list of destructive demons is lengthy; and each man, young or old, knows the ones with which he must contend. I’ll name but a few: the Demon of Greed; the Demon of Dishonesty; the Demon of Debt; the Demon of Doubt; the Demon of Drugs; and those twin Demons of Immodesty and Immorality. Each of these demons can wreak havoc with our lives. A combination of them can spell utter destruction.

Concerning greed, the counsel from Ecclesiastes speaks caution: “He that loveth silver shall not be satisfied with silver; nor he that loveth abundance with increase.”

Jesus counseled, “Take heed, and beware of covetousness: for a man’s life consisteth not in the abundance of the things which he possesseth.”

We must learn to separate need from greed.

When we speak of the demon of dishonesty, we can find it in a variety of locations. One such place is in school. Let us avoid cheating, falsifying, taking advantage of others, or anything like unto it. Let integrity be our standard.

In decision making, ask not, “What will others think?” but rather, “What will I think of myself?”

Enticements to embrace the demon of debt are thrust upon us many times each day. I quote the counsel from President Gordon B. Hinckley:

“I am troubled by the huge consumer installment debt which hangs over the people of the nation, including our own people. …

“We are beguiled by seductive advertising. Television carries the enticing invitation to borrow up to 125 percent of the value of one’s home. But no mention is made of interest. …

“I recognize that it may be necessary to borrow to get a home, of course. But let us buy a home that we can afford and thus ease the payments which will constantly hang over our heads without mercy or respite for as long as 30 years.”

I would add: We must not allow our yearnings to exceed our earnings.

In discussing the demon of drugs, I include, of course, alcohol. Drugs impair our ability to think, to reason, and to make prudent and wise choices. Often they result in violence and in child and wife abuse, and they can provoke conduct which brings pain and suffering to those who are innocent. “Just say no to drugs” is an effective statement of one’s determination. And this can be buttressed by the scripture:

“Know ye not that ye are the temple of God, and that the Spirit of God dwelleth in you?

“If any man defile the temple of God, him shall God destroy; for the temple of God is holy, which temple ye are.”

When I consider the demons who are twins—even immodesty and immorality—I should make them triplets and include pornography. They all three go together.

In the interpretation of Lehi’s dream, we find a rather apt description of the destructiveness of pornography: “And the mists of darkness are the temptations of the devil, which blindeth the eyes, and hardeneth the hearts of the children of men, and leadeth them away into broad roads, that they perish and are lost.”

A modern-day Apostle, Hugh B. Brown, has declared, “Any immodesty inducing impure thoughts is a desecration of the body—that temple in which the Holy Spirit may dwell.”

I commend to you tonight a jewel from the Improvement Era. It was published in 1917 but is equally applicable here and now: “The current and common custom of indecency in dress, the flood of immoral fiction in printed literature, in the drama, and notably in [motion] picture[s] … , the toleration of immodesty in every-day conversation and demeanor, are doing deadly work in the fostering of soul-destroying vice.”

Alexander Pope, in his inspired Essay on Man, declared:

\begin{verse}
Vice is a monster of so frightful mien,\\
As to be hated needs but to be seen;\\
Yet seen too oft, familiar with her face,\\
We first endure, then pity, then embrace.
\end{verse}

Perhaps a fitting summation pertaining to this demon can be found in the First Epistle of Paul to the Corinthians: “There hath no temptation taken you but such as is common to man: but God is faithful, who will not suffer you to be tempted above that ye are able; but will with the temptation also make a way to escape, that ye may be able to bear it.”

For each of us it is infinitely better to hear and heed the call of conscience, for conscience always warns us as a friend before punishing us as a judge.

The Lord Himself gives us the final word: “Be ye clean that bear the vessels of the Lord.”

Brethren, there is one responsibility that no man can evade. That is the effect of personal influence.

Our influence is surely felt in our respective families. Sometimes we fathers forget that once we too were boys, and boys at times can be vexing to parents.

I recall how much, as a youngster, I liked dogs. One day I took my wagon and placed a wooden orange crate in it and went looking for dogs. At that time, dogs were everywhere to be found: at school, walking along the sidewalks, or exploring vacant lots, of which there were many. As I would find a dog and capture it, I placed it in the crate, took it home, locked it in the coal shed, and turned the latch on the door. That day I think I brought home six dogs of varying sizes and made them my prisoners after this fashion. I had no idea what I would do with all those dogs, so I didn’t reveal my deed to anyone.

Dad came home from work and, as was his custom, took the coal bucket and went to the coal shed to fill it. Can you imagine his shock and utter consternation as he opened the door and immediately faced six dogs, all attempting to escape at once? As I recall, Dad flushed a little bit, and then he calmed down and quietly told me, “Tommy, coal sheds are for coal. Other people’s dogs rightfully belong to them.” By observing him, I learned a lesson in patience and calmness.

It is a good thing I did, for a similar event occurred in my life with our youngest son, Clark.

Clark has always liked animals, birds, reptiles—anything that is alive. Sometimes that resulted in a little chaos in our home. One day in his boyhood he came home from Provo Canyon with a water snake, which he named Herman.

Right off the bat Herman got lost. Sister Monson found him in the silverware drawer. Water snakes have a way of being where you least expect them. Well, Clark moved Herman to the bathtub, put a plug in the drain, put a little water in, and had a sign taped to the back of the tub which read, “Don’t use this tub. It belongs to Herman.” So we had to use the other bathroom while Herman occupied that sequestered place.

But then one day, to our amazement, Herman disappeared. His name should have been Houdini. He was gone! So the next day Sister Monson cleaned up the tub and prepared it for normal use. Several days went by.

One evening I decided it was time to take a leisurely bath, so I filled the tub with a lot of warm water, and then I peacefully lay down in the tub for a few moments of relaxation. I was lying there just pondering, when the soapy water reached the level of the overflow drain and began to flow through it. Can you imagine my surprise when, with my eyes focused on that drain, Herman came swimming out, right for my face? I yelled out to my wife, “Frances! Here comes Herman!”

Well, Herman was captured again, put in a foolproof box, and we made a little excursion to Vivian Park in Provo Canyon and there released Herman into the beautiful waters of the South Fork Creek. Herman was never again to be seen by us.

There appears in the Doctrine and Covenants, section 107, verse 99, a brief but direct admonition to each priesthood bearer: “Wherefore, now let every man learn his duty, and to act in the office in which he is appointed, in all diligence.” I have always taken this charge seriously and have attempted to live up to its direction.

In the recesses of my mind, I hear over and over again the guiding direction which President John Taylor gave to the brethren of the priesthood: “If you do not magnify your callings, God will hold you responsible for those you might have saved, had you done your duty.”

In the performance of our responsibilities, I have learned that when we heed a silent prompting and act upon it without delay, our Heavenly Father will guide our footsteps and bless our lives and the lives of others. I know of no experience more sweet or feeling more precious than to heed a prompting only to discover that the Lord has answered another person’s prayer through you.

Perhaps just one example will suffice. One day just over a year ago, after taking care of matters at the office, I felt a strong impression to visit an aged widow who was a patient at St. Joseph Villa here in Salt Lake City. I drove there directly.

When I went to her room, I found it empty. I asked an attendant concerning her whereabouts and was directed to a lounge area. There I found this sweet widow visiting with her sister and another friend. We had a pleasant conversation together.

As we were talking, a man came to the door of the room to obtain a can of soda water from the vending machine. He glanced at me and said, “Why, you are Tom Monson.”

“Yes,” I replied. “And you look like a Hemingway.” He acknowledged that he was Stephen Hemingway, the son of Alfred Eugene Hemingway, who had served as my counselor when I was a bishop many years ago and whom I called Gene. Stephen told me that his father was there in the same facility and was near death. He had been calling my name, and the family had wanted to contact me but had been unable to find a telephone number for me.

I excused myself immediately and went with Stephen up to the room of my former counselor, where others of his children were also gathered, his wife having passed away some years previous. The family members regarded my meeting Stephen in the lounge area as a response by our Heavenly Father to their great desire that I would see their father before he died and answer his call. I too felt that this was the case, for if Stephen had not entered the room in which I was visiting at precisely the time he did, I would not have known that Gene was even in that facility.

We gave a blessing to him. A spirit of peace prevailed. We had a lovely visit, after which I left.

The following morning a phone call revealed that Gene Hemingway had passed away—just 20 minutes after he had received the blessing from his son and me.

I expressed a silent prayer of thanks to Heavenly Father for His guiding influence which prompted my visit to St. Joseph Villa and led me to my dear friend Alfred Eugene Hemingway.

I like to think that Gene Hemingway’s thoughts that evening, as we basked in the Spirit’s glow, participated in humble prayer, and pronounced a priesthood blessing, echoed the words mentioned in the hymn “Master, the Tempest Is Raging,” which I cited at the beginning of my message:

\begin{verse}
Linger, O blessed Redeemer!\\
Leave me alone no more,\\
And with joy I shall make the blest harbor\\
And rest on the blissful shore.
\end{verse}

I still love that hymn and testify to you tonight as to the comfort it offers:

\begin{verse}
Whether the wrath of the storm-tossed sea\\
Or demons or men or whatever it be,\\
No waters can swallow the ship where lies\\
The Master of ocean and earth and skies.\\
They all shall sweetly obey thy will:\\
Peace, be still.
\end{verse}

His words in holy writ are sufficient: “Be still, and know that I am God.” I testify to this truth in the name of Jesus Christ, amen.

\newpage
\settalk{In Search of Treasure}
\setspeaker{Thomas S. Monson}
\settalkdate{April 2003}
\talkheading{In Search of Treasure}{Thomas S. Monson}

When I was a boy I enjoyed reading Treasure Island by Robert Louis Stevenson. I also saw adventure movies where several individuals had separate pieces of a well-worn map which led the way to buried treasure if only the pieces could be found and put together.

I recall listening to a 15-minute radio program each weekday afternoon. The program of which I speak was Jack Armstrong, the All-American Boy. It began with the jingle, “Have you tried Wheaties, the best breakfast food in the land?” Then, in a voice filled with mystery, there emanated from the radio the message, “We now join Jack and Betty as they approach the fabulous secret entry to the elephants’ burial ground, where a treasure is concealed. But wait; danger lurks on the path ahead.”

Nothing could tear me away from this program. It was as though I were leading the search for the hidden treasure of precious ivory.

At another time and in a different setting, the Savior of the world spoke of treasure. In His Sermon on the Mount He declared:

“Lay not up for yourselves treasures upon earth, where moth and rust doth corrupt, and where thieves break through and steal:

“But lay up for yourselves treasures in heaven, where neither moth nor rust doth corrupt, and where thieves do not break through nor steal:

“For where your treasure is, there will your heart be also.”

The promised reward was not a treasure of ivory, gold, or silver. Neither did it consist of acres of land or a portfolio of stocks and bonds. The Master spoke of riches within the grasp of all—even joy unspeakable here and eternal happiness hereafter.

Today I have chosen to provide the three pieces of your treasure map to guide you to your eternal happiness. They are:

Let us consider each segment of the map.

First, learn from the past. Each of us has a heritage—whether from pioneer forebears, later converts, or others who helped to shape our lives. This heritage provides a foundation built of sacrifice and faith. Ours is the privilege and responsibility to build on such firm and stable footings.

A story written by Karen Nolen, which appeared in the New Era in 1974, tells of a Benjamin Landart who, in 1888, was 15 years old and an accomplished violinist. Living on a farm in northern Utah with his mother and seven brothers and sisters was sometimes a challenge to Benjamin, as he had less time than he would have liked to play his violin. Occasionally his mother would lock up the violin until he had his farm chores done, so great was the temptation for Benjamin to play it.

In late 1892 Benjamin was asked to travel to Salt Lake to audition for a place with the territorial orchestra. For him, this was a dream come true. After several weeks of practicing and prayers, he went to Salt Lake in March of 1893 for the much anticipated audition. When he heard Benjamin play, the conductor, a Mr. Dean, told Benjamin he was the most accomplished violinist he had heard west of Denver. He was told to report to Denver for rehearsals in the fall and learned that he would be earning enough to keep himself, with some left over to send home.

A week after Benjamin received the good news, however, his bishop called him into his office and asked if he couldn’t put off playing with the orchestra for a couple of years. He told Benjamin that before he started earning money there was something he owed the Lord. He then asked Benjamin to accept a mission call.

Benjamin felt that giving up his chance to play in the territorial orchestra would be almost more than he could bear, but he also knew what his decision should be. He promised the bishop that if there were any way to raise the money for him to serve, he would accept the call.

When Benjamin told his mother about the call, she was overjoyed. She told him that his father had always wanted to serve a mission but had been killed before that opportunity had come to him. However, when they discussed the financing of the mission, her face clouded over. Benjamin told her he would not allow her to sell any more of their land. She studied his face for a moment and then said, “Ben, there is a way we can raise the money. This family [has] one thing that is of great enough value to send you on your mission. You will have to sell your violin.”

Ten days later, on March 23, 1893, Benjamin wrote in his journal: “I awoke this morning and took my violin from its case. All day long I played the music I love. In the evening when the light grew dim and I could see to play no longer, I placed the instrument in its case. It will be enough. Tomorrow I leave [for my mission].”

Forty-five years later, on June 23, 1938, Benjamin wrote in his journal: “The greatest decision I ever made in my life was to give up something I dearly loved to the God I loved even more. He has never forgotten me for it.”

Learn from the past.

Second, prepare for the future. We live in a changing world. Technology has altered nearly every aspect of our lives. We must cope with these advances—even these cataclysmic changes—in a world of which our forebears never dreamed.

Remember the promise of the Lord: “If ye are prepared ye shall not fear.” Fear is a deadly enemy of progress.

It is necessary to prepare and to plan so that we don’t fritter away our lives. Without a goal, there can be no real success. One of the best definitions of success I have ever heard goes something like this: Success is the progressive realization of a worthy ideal. Someone has said the trouble with not having a goal is that you can spend your life running up and down the field and never crossing the goal line.

Years ago there was a romantic and fanciful ballad that contained the words, “Wishing will make it so / Just keep on wishing / And care will go.” I want to state here and now that wishing will not replace thorough preparation to meet the trials of life. Preparation is hard work but absolutely essential for our progress.

Our journey into the future will not be a smooth highway which stretches from here to eternity. Rather, there will be forks and turnings in the road, to say nothing of the unanticipated bumps. We must pray daily to a loving Heavenly Father, who wants each of us to succeed in life.

Prepare for the future.

Third, live in the present. Sometimes we let our thoughts of tomorrow take up too much of today. Daydreaming of the past and longing for the future may provide comfort but will not take the place of living in the present. This is the day of our opportunity, and we must grasp it.

Professor Harold Hill, in Meredith Willson’s The Music Man, cautioned, “You pile up enough tomorrows, and you’ll find you’ve collected a lot of empty yesterdays.”

There is no tomorrow to remember if we don’t do something today, and to live most fully today, we must do that which is of greatest importance. Let us not procrastinate those things which matter most.

I recently read the account of a man who, just after the passing of his wife, opened her dresser drawer and found there an item of clothing she had purchased when they visited the eastern part of the United States nine years earlier. She had not worn it but was saving it for a special occasion. Now, of course, that occasion would never come.

In relating the experience to a friend, the husband said, “Don’t save something only for a special occasion. Every day in your life is a special occasion.”

That friend later said those words changed her life. They helped her to cease putting off the things most important to her. Said she: “Now I spend more time with my family. I use crystal glasses every day. I’ll wear new clothes to go to the supermarket if I feel like it. The words ‘someday’ and ‘one day’ are fading from my vocabulary. Now I take the time to call my relatives and closest friends. I’ve called old friends to make peace over past quarrels. I tell my family members how much I love them. I try not to delay or postpone anything that could bring laughter and joy into our lives. And each morning, I say to myself that this could be a special day. Each day, each hour, each minute is special.”

A wonderful example of this philosophy was shared by Arthur Gordon many years ago in a national magazine. He wrote:

“When I was around thirteen and my brother ten, Father had promised to take us to the circus. But at lunchtime there was a phone call; some urgent business required his attention downtown. We braced ourselves for disappointment. Then we heard him say [into the phone], ‘No, I won’t be down. It’ll have to wait.’

“When he came back to the table, Mother smiled. ‘The circus keeps coming back, you know,’ [she said].

“‘I know,’ said Father. ‘But childhood doesn’t.’”

Elder Monte J. Brough of the First Quorum of the Seventy tells of a summer at his childhood home in Randolph, Utah, when he and his younger brother, Max, decided to build a tree house in a large tree in the backyard. They made plans for the most wonderful creation of their lives. They gathered building materials from all over the neighborhood and carried them up to a part of the tree where two branches provided an ideal location for the house. It was difficult, and they were anxious to complete their work. The vision of the finished tree house provided tremendous motivation for them to complete the project.

They worked all summer, and finally in the fall just before school began for the new year, their house was completed. Elder Brough said he will never forget the feelings of joy and satisfaction which were theirs when they finally were able to enjoy the fruit of their work. They sat in the tree house, looked around for a few minutes, climbed down from the tree—and never returned. The completed project, as wonderful as it was, could not hold their interest for even one day. In other words, the process of planning, gathering, building, and working—not the completed project—provided the enduring satisfaction and pleasure they had experienced.

Let us relish life as we live it and, as did Elder Brough and his brother, Max, find joy in the journey.

The old adage “Never put off until tomorrow what you can do today” is doubly important when it comes to expressing our love and affection—in word and in deed—to family members and friends. Said author Harriet Beecher Stowe, “The bitterest tears shed over graves are for words left unsaid and deeds left undone.”

A poet set to verse the sorrow of opportunities forever lost. I quote a portion:

\begin{verse}
Around the corner I have a friend,\\
In this great city that has no end;\\
Yet days go by, and weeks rush on,\\
And before I know it, a year is gone,\\
And I never see my old friend’s face,\\
For Life is a swift and terrible race. …\\
But to-morrow comes—and to-morrow goes,\\
And the distance between us grows and grows.
\end{verse}

Just a little over a year ago, I determined that I would not put off any longer a visit with a dear friend whom I hadn’t seen for many years. I had been meaning to visit him in California but just had not gotten around to it.

Bob Biggers and I met when we were both in the Classification Division at the United States Naval Training Center in San Diego, California, toward the close of World War II. We were good friends from the beginning. He visited in Salt Lake once before he married, and we remained friends through correspondence from the time I was discharged in 1946. My wife, Frances, and I have exchanged Christmas cards every year with Bob and his wife, Grace.

Finally, at the beginning of January 2002, I scheduled a stake conference visit to Whittier, California, where the Biggers live. I telephoned my friend Bob, now 80 years old, and arranged for Frances and me to meet him and Grace, that we might reminisce concerning former days.

We had a delightful visit. I took with me a number of photographs which had been taken when we were in the Navy together over 55 years earlier. We identified the men we knew and provided each other an update on their whereabouts as best we could. Although not a member of our Church, Bob remembered going to a sacrament meeting with me those long years before when we were stationed in San Diego.

As Frances and I said our good-byes to Bob and Grace, I felt an overwhelming sense of peace and joy at having finally made the effort to see once again a friend who had been cherished from afar throughout the years.

One day each of us will run out of tomorrows. Let us not put off what is most important.

Live in the present.

Your treasure map is now in place: Learn from the past, prepare for the future, live in the present.

I conclude where I began. From our Lord and Savior:

“Lay not up for yourselves treasures upon earth, where moth and rust doth corrupt, and where thieves break through and steal:

“But lay up for yourselves treasures in heaven, where neither moth nor rust doth corrupt, and where thieves do not break through nor steal:

“For where your treasure is, there will your heart be also.”

My brothers and sisters, from the depths of my soul I bear you my personal witness: God is our Father; His Son is our Savior and Redeemer; we are led by a prophet for our time, even President Gordon B. Hinckley.

In the name of Jesus Christ, amen.

\newpage
\settalk{Stand in Your Appointed Place}
\setspeaker{Thomas S. Monson}
\settalkdate{April 2003}
\talkheading{Stand in Your Appointed Place}{Thomas S. Monson}

We are assembled this evening as a mighty body of the priesthood, both here in the Conference Center and in locations throughout the world. Some hold the Aaronic Priesthood, while others bear the Melchizedek Priesthood.

President Stephen L. Richards, who served as a counselor to President David O. McKay, declared, “The Priesthood is usually simply defined as ‘the power of God delegated to man.’” He continues: “This definition, I think, is accurate. But for practical purposes I like to define the Priesthood in terms of service and I frequently call it ‘the perfect plan of service.’ … It is an instrument of service. … And the man who fails to use it is apt to lose it, for we are plainly told by revelation that he who neglects it ‘shall not be counted worthy to stand.’”

In the Pioneer Stake, located in Salt Lake City and where I received both the Aaronic and Melchizedek Priesthood, we were taught to become familiar with the scriptures, including sections 20, 84, and 107 of the Doctrine and Covenants. In these sections we learn about priesthood and Church government.

Tonight I wish to emphasize one verse from section 107: “Wherefore, now let every man learn his duty, and to act in the office in which he is appointed, in all diligence.”

President Harold B. Lee frequently taught: “When one becomes a holder of the priesthood, he becomes an agent of the Lord. He should think of his calling as though he were on the Lord’s errand.”

We also learn from these sections the duties of quorum presidencies and the fact that we are responsible for others besides ourselves.

I firmly believe that the Church today is stronger than it has ever been. Activity levels of our youth testify that this is a generation of faith and devotion to truth. Yet there are some who drop by the wayside, who find other interests that persuade them to neglect their Church duties. We must not lose such precious souls.

There are growing numbers among the prospective elders who are not found in Church meetings nor filling Church assignments. This situation can and must be remedied. The task is ours. Responsibility needs to be assigned and effort put forth without delay.

The presidencies of the Aaronic Priesthood quorums, under the leadership of the bishopric and quorum advisers, can be empowered to reach out and rescue.

Said the Lord, “Remember the worth of souls is great in the sight of God; … and how great is his joy in the soul that repenteth!”

Sometimes the task appears overwhelming. We can take fresh courage from the experience of Gideon of old, who, with his modest force, was to do battle with the Midianites and the Amalekites. You will remember how Gideon and his army faced an overwhelming strength of forces vastly superior in equipment and in number. The book of Judges in the Old Testament records that the united enemy, the Midianites and the Amalekites, “lay along in the valley like grasshoppers for multitude; and their camels were without number, as the sand by the sea side for multitude.” Gideon went to Almighty God for his strength.

To his surprise, Gideon was advised by the Lord that his forces were too many in number for the Lord to deliver the enemy into their hands, lest they say, “Mine own hand hath saved me.” Gideon was instructed to proclaim to his people: “Whosoever is fearful and afraid, let him return and depart … from mount Gilead. And there returned of the people twenty and two thousand; and there remained ten thousand.”

Then the Lord said, “The people are yet too many.” He instructed Gideon to take the men to water to observe the manner in which they should drink of the water. Those who lapped the water were placed in one group, and those who bowed down upon their knees to drink were placed in another. The Lord said unto Gideon, “By the three hundred men that lapped will I save you, and deliver the Midianites into thine hand: and let all the other people go every man unto his place.”

Gideon returned to his forces and said to them, “Arise; for the Lord hath delivered into your hand the host of Midian.” Then he divided the 300 men into three companies, and he put a trumpet in every man’s hand, with empty pitchers and lamps within the pitchers. And he said unto them:

“Look on me, and do likewise: and, behold, when I come to the outside of the camp, it shall be that, as I do, so shall ye do.

“When I blow with a trumpet, I and all that are with me, then blow ye the trumpets also on every side … and say, The sword of the Lord, and of Gideon.” He then said in effect, “Follow me.” His exact words were, “As I do, so shall ye do.”

At the leader’s signal, the host of Gideon did blow on the trumpets and did break the pitchers and did shout, “The sword of the Lord, and of Gideon.” The scripture records the outcome of this decisive battle: “And they stood every man in his place,” and the victory was won.

Home teaching is part of today’s plan to rescue. When it was introduced by President David O. McKay to all of the General Authorities, he counseled: “Home teaching is one of our most urgent and most rewarding opportunities to nurture and inspire, to counsel and direct our Father’s children. … [It] is a divine service, a divine call. It is our duty as Home Teachers to carry the divine spirit into every home and heart.”

In certain areas where adequate Melchizedek Priesthood strength is missing, stake presidents and bishops, coordinating with the mission president, may use full-time missionaries to visit less-active and part-member families. Not only does this rekindle the missionary spirit in the home, but it also provides an ideal opportunity for quality referrals to be obtained.

Over the years as I have visited many stakes throughout the world, there have been those stakes where ward and stake leaders, out of necessity or in response to duty, stopped wringing their hands, rolled up their sleeves, and, with the Lord’s help, went to work and brought precious men to qualify for the Melchizedek Priesthood and, with their wives and children, to enter the holy temple for their endowments and sealings.

In brief form I will mention several examples.

On a visit to the Millcreek Stake in Salt Lake City some years ago, I learned that just over 100 brethren who were prospective elders had been ordained elders during the preceding year. I asked President James Clegg the secret of his success. Although he was too modest to take the credit, one of his counselors revealed that President Clegg, recognizing the challenge, had undertaken to personally call and arrange a private appointment between him and each prospective elder. During the appointment, President Clegg would mention the temple of the Lord, the saving ordinances and covenants emphasized there, and would conclude with this question: “Wouldn’t you desire to take your sweet wife and your precious children to the house of the Lord, that you might be a forever family throughout the eternities?” An acknowledgment followed, the reactivation process was pursued, and the goal was achieved.

In 1952 the majority of the families in the Rose Park Third Ward were members whose fathers or husbands held only the Aaronic Priesthood, rather than the Melchizedek Priesthood. Brother L. Brent Goates was called to serve as the bishop. He invited a less-active brother in the ward, Ernest Skinner, to assist in activating the 29 adult brethren in the ward who held the office of teacher in the Aaronic Priesthood and to help these men and their families get to the temple. As a less-active member himself, Brother Skinner was reluctant at first but finally indicated that he would do what he could. He began personally visiting with the less-active adult teachers, trying to help them see their role as priesthood leaders in their homes and as husbands and fathers to their families. He soon enlisted some of the less-active brethren to assist him in his assignment. One by one they became fully active again and took their families to the temple.

One day the ward clerk came out of a grocery checking line to greet the last of the group to go to the temple. Commenting on his position as the last, the man said: “I stood by and watched as all of that group became active in our ward and went to the temple. If only I had been able to imagine how beautiful it was in the temple, and how it would change my life forever, I never would have been the last of 29 to be sealed in the temple.”

In each of these accounts, there were four elements which led them to success:

Brethren, let us remember the counsel of King Benjamin: “When ye are in the service of your fellow beings ye are only in the service of your God.”

Let us reach out to rescue those who so need our help and lift them to the higher road and the better way. Let us focus our thinking on the needs of priesthood holders and their wives and children who have slipped from the path of activity. May we listen to the unspoken message from their hearts:

\begin{verse}
Lead me, guide me, walk beside me,\\
Help me find the way.\\
Teach me all that I must do\\
To live with him someday.
\end{verse}

The work of reactivation is no task for the idler or dreamer. Children grow, parents age, and time waits for no man. Don’t postpone a prompting; rather, act on it, and the Lord will open the way.

Frequently the heavenly virtue of patience is required. As a bishop I felt prompted one day to call on a man whose wife was somewhat active, as were the children. This man, however, had never responded. It was a hot summer’s day when I knocked on the screen door of Harold G. Gallacher. I could see Brother Gallacher sitting in his chair reading the newspaper. “Who is it?” he queried without looking up.

“Your bishop,” I replied. “I’ve come to get acquainted and to urge your attendance with your family at our meetings.”

“No, I’m too busy,” came the disdainful response. He never looked up. I thanked him for listening and departed the doorstep.

The Gallacher family moved to California shortly thereafter. The years went by. Then, as a member of the Quorum of the Twelve, I was working in my office one day when my secretary called, saying: “A Brother Gallacher who once lived in your ward would like to talk to you. He’s here in my office.”

I responded, “Ask him if his name is Harold G. Gallacher, who, with his family, lived at Vissing Place on West Temple and Fifth South.”

She said, “He is the man.”

I asked her to send him in. We had a pleasant conversation together concerning his family. He told me, “I’ve come to apologize for not getting out of my chair and letting you in the door that summer day long years ago.” I asked him if he was active in the Church. With a wry smile, he replied: “I’m now second counselor in my ward bishopric. Your invitation to come out to church, and my negative response, so haunted me that I determined to do something about it.”

Harold and I visited together on numerous occasions before he passed away. The Gallachers and their children filled many callings in the Church. One of the youngest grandchildren is now serving a full-time mission.

To the many missionaries who may be listening this evening, I share the observation that the seeds of testimony frequently do not immediately take root and flower. Bread cast upon the water returns, at times, only after many days. But it does return.

I answered the ring of my telephone one evening to hear a voice ask, “Are you related to an Elder Monson who years ago served in the New England Mission?”

I answered that such was not the case. The caller introduced himself as a Brother Leonardo Gambardella and then mentioned that an Elder Monson and an Elder Bonner called at his home long ago and bore their testimonies to him and his wife. They had listened but had done nothing further to apply their teachings. Subsequently they moved to California, where, some 13 years later, they again found the truth and were converted and baptized. Brother Gambardella then asked if there were any way he could reach the elders who first had visited with them, that he might express his profound gratitude for their testimonies, which had remained with him and his wife.

I checked the records. I located the elders, now married with families of their own. Can you imagine their surprise when I telephoned them and told them the good news—even the culmination of their early efforts? They instantly remembered the Gambardellas. I arranged a conference telephone call so they could personally extend their congratulations and welcome them into the Church. They did. There were tears, but they were tears of joy.

Edwin Markham penned these lines:

\begin{verse}
There is a destiny that makes us brothers;\\
None goes his way alone:\\
All that we send into the lives of others\\
Comes back into our own.
\end{verse}

Tonight I pray that all of us who hold the priesthood may sense our responsibilities, that we, like Gideon of old, may stand every man in his appointed place and, as one, follow our Leader—even the Lord Jesus Christ—and His prophet, President Gordon B. Hinckley. May we reach out and rescue those who have fallen by the wayside, that not one precious soul will be lost.

In the name of Jesus Christ, amen.

\newpage
\settalk{Bring Him Home}
\setspeaker{Thomas S. Monson}
\settalkdate{October 2003}
\talkheading{Bring Him Home}{Thomas S. Monson}

My dear brethren, it is a humbling experience to stand before you this evening and to realize that beyond the imposing audience in this, the Conference Center, many hundreds of thousands of priesthood bearers are similarly assembled throughout the world.

While contemplating the responsibility to speak to you, I recalled a definition of priesthood authority declared by President Stephen L Richards. Said he: “The Priesthood is usually simply defined as ‘the power of God delegated to man.’ This definition, I think, is accurate. But for practical purposes I like to define the Priesthood in terms of service and I frequently call it ‘the perfect plan of service.’”

Whether we hold the office of a deacon in the Aaronic Priesthood or that of an elder in the Melchizedek Priesthood, we are duty bound by the Lord’s revelation found in the 107th section of the Doctrine and Covenants, verse 99: “Wherefore, now let every man learn his duty, and to act in the office in which he is appointed, in all diligence.”

As our youngest son, Clark, was approaching his 12th birthday, he and I were leaving the Church Administration Building when President Harold B. Lee approached and greeted us. I mentioned that Clark would soon be 12, whereupon President Lee turned to him and asked, “What happens to you when you turn 12?”

This was one of those times when a father prays that a son will be inspired to give a proper response. Clark, without hesitation, said to President Lee, “I will be ordained a deacon!”

The answer was the one President Lee had sought. He then counseled our son, “Remember, it is a great blessing to hold the priesthood.”

When I was a boy, I looked forward to passing the sacrament to the ward members. We deacons were trained as to our duties. One of the men in our ward, Louis, suffered from palsy. His head and hands shook so violently that he could not, by himself, partake of the sacrament. Each deacon knew that his duty in serving Louis was to hold the bread to his lips so that he might partake and to similarly place the cup of water to his mouth with one hand, while steadying his head with the other, the tray being held by another deacon while doing so. Always Louis would say, “Thank you.”

It was 40 years ago this conference time when President David O. McKay called me to serve as a member of the Quorum of the Twelve Apostles. At the first meeting of the Presidency and Twelve which I attended where the sacrament was served, President McKay announced, “Before we partake of the sacrament, I would like to ask our newest member of this body, Brother Monson, if he would instruct the First Presidency and Twelve on the atoning sacrifice of our Lord and Savior Jesus Christ.” It was then that I gained a true understanding of the old adage: “When the time for decision arrives, the time for preparation is past.” It was also the time to remember the counsel found in 1 Peter: “Be ready always to give an answer to every man that asketh you a reason of the hope that is in you.”

I began my remarks by referring to a letter which I had received from one of the servicemen from our ward who was serving on the front lines in Korea during that sometimes forgotten war. The writer told how, amidst the shelling on Sunday morning, several in his platoon partook of the bread and then the water, both served from a helmet. Each remembered the significance of the blessing pronounced on the sacred emblems and his individual responsibility to keep the commandments of the Lord and to follow the Lord’s example of service to others.

The memory of that particular experience with the First Presidency and Quorum of the Twelve has not dimmed in the intervening 40 years.

To those who have been absent from home and family, whether in the military, on missions, or for other purposes, the holiday season brings forth a yearning—even a longing—to be together with loved ones. To hear the laughter of children, to witness the expression of love by parents, and to feel the embrace of brothers and sisters provide a preview of heaven and the eternal joy to be found there.

One December evening, while waiting to board a plane en route to the United States, Sister Monson and I were standing in the stifling heat and humidity of Singapore, when over the airport loudspeaker system came a familiar, lilting melody, with Bing Crosby singing the words:

\begin{verse}
I’ll be home for Christmas;\\
You can plan on me.\\
Please have snow and mistletoe\\
And presents on the tree.\\
Christmas Eve will find me\\
Where the love-light gleams.\\
I’ll be home for Christmas\\
If only in my dreams.
\end{verse}

The First Presidency has long emphasized the statement, “The home is the basis of a righteous life, and no other instrumentality can take its place or fulfill its essential functions.”

There are those families comprised of mothers and fathers, sons and daughters who have, through thoughtless comment, isolated themselves from one another. An account of how such a tragedy was narrowly averted occurred many years ago in the life of a young man who, for purposes of privacy, I shall call Jack.

Throughout Jack’s life, he and his father had many serious arguments. One day, when he was 17, they had a particularly violent one. Jack said to his father, “This is the straw that breaks the camel’s back. I’m leaving home, and I shall never return.” So saying, he went to the house and packed his bag. His mother begged him to stay; he was too angry to listen. He left her crying at the doorway.

Leaving the yard, he was about to pass through the gate when he heard his father call to him, “Jack, I know that a large share of the blame for your leaving rests with me. For this I am truly sorry. I want you to know that if you should ever wish to return home, you’ll always be welcome. And I’ll try to be a better father to you. I want you to know that I’ll always love you.”

Jack said nothing but went to the bus station and bought a ticket to a distant point. As he sat on the bus, watching the miles go by, he commenced to think about the words of his father. He began to realize how much love it had required for him to do what he had done. Dad had apologized. He had invited him back and left the words ringing in the summer air: “I love you.”

It was then that Jack realized that the next move was up to him. He knew the only way he could ever find peace with himself was to demonstrate to his father the same kind of maturity, goodness, and love that Dad had shown toward him. Jack got off the bus. He bought a return ticket and went back.

He arrived shortly after midnight, entered the house, turned on the light. There in the rocking chair sat his father, his head in his hands. As he looked up and saw Jack, he arose from the chair and they rushed into each other’s arms. Jack often said, “Those last years that I was home were among the happiest of my life.”

We could say that here was a boy who overnight became a man. Here was a father who, suppressing passion and bridling pride, rescued his son before he became one of that vast “lost battalion” resulting from fractured families and shattered homes. Love was the binding band, the healing balm. Love so often felt, so seldom expressed.

From Mount Sinai there thunders in our ears, “Honour thy father and thy mother.” And later from the Lord the injunction, “Live together in love.”

Brethren, ours is the responsibility, yes, even the solemn duty, to reach out to those who have slipped into inactivity or strayed from the family circle.

Recall with me the beautiful words of the Lord’s revelation from section 18 of the Doctrine and Covenants: “Remember the worth of souls is great in the sight of God. …

“And if it so be that you should labor all your days in crying repentance unto this people, and bring, save it be one soul unto me, how great shall be your joy with him in the kingdom of my Father!

“And now, if your joy will be great with one soul that you have brought unto me into the kingdom of my Father, how great will be your joy if you should bring many souls unto me!”

As presidencies of Aaronic Priesthood quorums, as advisers to these quorums, we can, with the Lord’s help, reach out and rescue those for whom we have responsibility. Young men, with a smile on your face and determination in your heart, you can take, arm in arm, a less-active boy and together come to priesthood meeting and learn of the Lord and what He has prepared for you to do. You are entitled to His divine help, for He has promised you: “I will go before your face. I will be on your right hand and on your left, and my Spirit shall be in your hearts, and mine angels round about you, to bear you up.”

Brethren of the Melchizedek Priesthood, you have the same sacred charge and obligation as pertains to your duties to other men and to their families. And you have the same promise of the Lord to attend your efforts.

As you succeed, you will be answering a mother’s prayer, the tender though unexpressed feelings of children’s hearts; and your names will forever be honored by those whom you reach out and help.

Let me share with you a rather private but joyful example from my own experience.

As a bishop, I worried about any members who were inactive, not attending, not serving. Such was my thought one day as I drove down the street where Ben and Emily Fullmer lived. Aches and pains of advancing years caused them to withdraw from activity to the shelter of their home—isolated, detached, shut out from the mainstream of daily life and association. Ben and Emily had not been in our sacrament meeting for many years. Ben, a former bishop, would sit constantly in his front room reading and memorizing the New Testament.

I was en route from my uptown sales office to our plant on Industrial Road. For some reason I had driven down First West, a street which I never had traveled before to reach the destination of our plant. Then I felt the unmistakable prompting to park my car and visit Ben and Emily, even though I was on my way to a meeting. I did not heed the impression at first but drove on for two more blocks; however, when the impression came again, I returned to their home.

It was a sunny weekday afternoon. I approached the door to their home and knocked. I heard the tiny fox terrier dog bark at my approach. Emily welcomed me in. Upon seeing me, she exclaimed, “All day long I have waited for my phone to ring. It has been silent. I hoped the postman would deliver a letter. He brought only bills. Bishop, how did you know today is my birthday?”

I answered, “God knows, Emily, for He loves you.”

In the quiet of their living room, I said to Ben and Emily, “I really don’t know why I was directed here today, but I was. Our Heavenly Father knows. Let’s kneel in prayer and ask Him why.” This we did, and the answer came. As we arose from our knees, I said to Brother Fullmer, “Ben, would you come to priesthood meeting when we meet with all the priesthood and relate to our Aaronic Priesthood boys the story you once told me when I was a boy, how you and a group of boys were en route to the Jordan River to swim one Sunday, but you felt the Spirit direct you to attend Sunday School. And you did. One of the boys who failed to respond to that Spirit drowned that Sunday. Our boys would like to hear your testimony.”

“I’ll do it,” he responded.

I then said to Sister Fullmer, “Emily, I know you have a beautiful voice. My mother has told me so. Our ward conference is a few weeks away, and our choir will sing. Would you join the choir and attend our ward conference and perhaps sing a solo?”

“What will the number be?” she inquired.

“I don’t know,” I said, “but I’d like you to sing it.”

She sang. He spoke to the Aaronic Priesthood. Hearts were gladdened by the return to activity of Ben and Emily. They rarely missed a sacrament meeting from that day forward. The language of the Spirit had been spoken. It had been heard. It had been understood. Hearts were touched and souls saved. Ben and Emily Fullmer had come home.

One of the longest-running musicals in history is Les Misérables. The story is set in the period of the French Revolution. The principal character in the musical is Jean Valjean. In his heartfelt concern for the young man Marius, who is going off to battle, he expresses in song a sincere prayer:

\begin{verse}
God on high,\\
Hear my prayer;\\
In my need\\
You have always been there.
\end{verse}

Brethren, as we go forward as bearers of the priesthood of God, learning our duty and then reaching out to our brethren who stand in need of our help, let us look upward to our Heavenly Father, who is the Father of us all. We may not hear His voice, but we will remember His salutation, “Well done, thou good and faithful servant.”

And within our hearts we will recognize His unspoken plea: Bring him home. In the name of Jesus Christ, amen.

\newpage
\settalk{The Bridge Builder}
\setspeaker{Thomas S. Monson}
\settalkdate{October 2003}
\talkheading{The Bridge Builder}{Thomas S. Monson}

Many years ago I read a book entitled The Way to the Western Sea, by David S. Lavender. It provides a fascinating account of the epic journey of Meriwether Lewis and William Clark as they led their famed expedition across the North American continent to discover an overland route to the Pacific Ocean.

Their trek was a nightmare of backbreaking toil, deep gorges which had to be crossed, and extensive travel by foot, carrying with them their supply-laden boats to find the next stream on which to make their way.

As I read of their experiences, I frequently mused, “If only there were modern bridges to span the gorges of the raging waters.” There came to my mind thoughts of magnificent bridges of our time which accomplish this task with ease: the beautiful Golden Gate Bridge of San Francisco fame; the sturdy Sydney, Australia, Harbour Bridge; and others in many lands.

In reality, we are all travelers—even explorers of mortality. We do not have the benefit of previous personal experience. We must pass over steep precipices and turbulent waters in our own journey here on earth.

Perhaps such a somber thought inspired the poet Will Allen Dromgoole’s classic poem entitled “The Bridge Builder.”

\begin{verse}
An old man, going a lone highway,\\
Came at the evening, cold and gray,\\
To a chasm, vast and deep and wide,\\
Through which was flowing a sullen tide.\\
The old man crossed in the twilight dim;\\
The sullen stream had no fears for him;\\
But he turned when safe on the other side\\
And built a bridge to span the tide.
\end{verse}

The message of the poem has prompted my thinking and comforted my soul, for our Lord and Savior, Jesus Christ, was the supreme architect and builder of bridges for you, for me, for all mankind. He has built the bridges over which we must cross if we are to reach our heavenly home.

The Savior’s mission was foretold. Matthew recorded, “And she shall bring forth a son, and thou shalt call his name Jesus: for he shall save his people from their sins.”

There followed the miracle of His birth and the gathering of the shepherds who came with haste to that stable, to that mother, to that child. Even the Wise Men, journeying from the East, followed that star and bestowed their precious gifts upon the young child.

The scripture records that Jesus “grew, and waxed strong in spirit, filled with wisdom: and the grace of God was upon him” and that He “went about doing good.”

What personal bridges did He build and cross here in mortality, showing us the way to follow? He knew mortality would be filled with dangers and difficulties. He declared: “Come unto me, all ye that labour and are heavy laden, and I will give you rest. Take my yoke upon you, and learn of me; for I am meek and lowly in heart: and ye shall find rest unto your souls. For my yoke is easy, and my burden is light.”

Jesus provided the Bridge of Obedience. He was an unfailing example of personal obedience as He kept the commandments of His Father.

When He was led of the Spirit into the wilderness to be tempted of Satan, He was weak from fasting. Satan was at his seductive best in the offerings he proffered. His first was to satisfy the Savior’s physical needs, including His hunger. To this the Savior replied, “It is written, Man shall not live by bread alone, but by every word that proceedeth out of the mouth of God.”

Next Satan offered power. Responded the Savior, “It is written again, Thou shalt not tempt the Lord thy God.”

Finally the Savior was offered wealth and earthly glory. His response: “Get thee hence, Satan: for it is written, Thou shalt worship the Lord thy God, and him only shalt thou serve.”

The Apostle Paul was inspired of the Lord to declare for our time, as well as for his: “There hath no temptation taken you but such as is common to man: but God is faithful, who will not suffer you to be tempted above that ye are able; but will with the temptation also make a way to escape, that ye may be able to bear it.”

Lest we equivocate, I mention a comment from ABC Nightline’s Ted Koppel: “What Moses brought down from Mt. Sinai were not the Ten Suggestions [but the Ten] Commandments!”

A bit of subtle humor is found in an account of a conversation between Mark Twain and a friend. Said the wealthy friend to Twain, “Before I die, I mean to make a pilgrimage to the Holy Land. I will climb to the top of Mount Sinai and read the Ten Commandments aloud.”

Replied Twain, “Why don’t you stay home and keep them!”

The second bridge provided by the Master for us to cross is the Bridge of Service. We look to the Savior as our example of service. Although He came to earth as the Son of God, He humbly served those around Him. He came forth from heaven to live on earth as mortal man and to establish the kingdom of God. His glorious gospel reshaped the thinking of the world. He blessed the sick; He caused the lame to walk, the blind to see, the deaf to hear. He even raised the dead to life.

In the 25th chapter of the book of Matthew, the Savior tells us this concerning the faithful who will be on His right hand at His triumphal return:

“Then shall the King say unto them … , Come, ye blessed of my Father, inherit the kingdom prepared for you from the foundation of the world:

“For I was an hungred, and ye gave me meat: I was thirsty, and ye gave me drink: I was a stranger, and ye took me in:

“Naked, and ye clothed me: I was sick, and ye visited me: I was in prison, and ye came unto me.

“Then shall the righteous answer him, saying, Lord, when saw we thee an hungred, and fed thee? or thirsty, and gave thee drink?

“When saw we thee a stranger, and took thee in? or naked, and clothed thee?

“Or when saw we thee sick, or in prison, and came unto thee?

“And the King shall answer and say unto them, Verily I say unto you, Inasmuch as ye have done it unto one of the least of these my brethren, ye have done it unto me.”

Elder Richard L. Evans once counseled, “We can’t do everything for everyone everywhere, but we can do something for someone somewhere.”

May I share with you an account of an opportunity of service which came to me unexpectedly and in an unusual manner. I received a telephone call from a granddaughter of an old friend. She asked, “Do you remember Francis Brems, who was your Sunday School teacher?” I told her that I did. She continued, “He is now 105 years of age. He lives in a small care center but meets with the entire family each Sunday, where he delivers a Sunday School lesson. Last Sunday, Grandpa announced to us, ‘My dears, I am going to die this week. Will you please call Tommy Monson and tell him this. He’ll know what to do.’”

I visited Brother Brems the very next evening. I could not speak to him, for he was deaf. I could not write a message for him to read, for he was blind. What was I to do? I was told that his family communicated with him by taking the finger of his right hand and then tracing on the palm of his left hand the name of the person visiting and then any message. I followed the procedure and took his finger and spelled on the palm of his hand T-O-M-M-Y M-O-N-S-O-N. Brother Brems became excited and, taking my hands, placed them on his head. I knew his desire was to receive a priesthood blessing. The driver who had taken me to the care center joined me as we placed our hands on the head of Brother Brems and provided the desired blessing. Afterward, tears streamed from his sightless eyes. He grasped our hands, and we read the movement of his lips. The message: “Thank you so much.”

Within that very week, just as Brother Brems had predicted, he passed away. I received the telephone call and then met with the family as funeral arrangements were made. How thankful I am that a response to render service was not delayed.

The bridge of service invites us to cross over it frequently.

Finally, the Lord provided us the Bridge of Prayer. He directed, “Pray always, and I will pour out my Spirit upon you, and great shall be your blessing.”

I share with you an account described in a mother’s letter to me relating to prayer. She wrote:

“Sometimes I wonder if I make a difference in my children’s lives. Especially as a single mother working two jobs to make ends meet, I sometimes come home to confusion, but I never give up hope.

“My children and I were watching a television broadcast of general conference, and you were speaking about prayer. My son made the statement, ‘Mother, you’ve already taught us that.’ I said, ‘What do you mean?’ And he replied, ‘Well, you’ve taught us to pray and showed us how, but the other night I came to your room to ask something and found you on your knees praying to Heavenly Father. If He’s important to you, He’ll be important to me.’”

The letter concluded, “I guess you never know what kind of influence you’ll be until a child observes you doing yourself what you have tried to teach him to do.”

No relating of a prayer touches me so deeply as the prayer offered by Jesus in the Garden of Gethsemane. I believe Luke describes it best:

“He … went … to the mount of Olives; and his disciples also followed him.

“And when he was at the place, he said unto them, Pray that ye enter not into temptation.

“And he was withdrawn from them about a stone’s cast, and kneeled down, and prayed,

“Saying, Father, if thou be willing, remove this cup from me: nevertheless not my will, but thine, be done.

“And there appeared an angel unto him from heaven, strengthening him.

“And being in an agony he prayed more earnestly: and his sweat was as it were great drops of blood falling down to the ground.”

In due time came the trek to the cross. What suffering He endured as He made His burdensome way, carrying His own cross. Heard were the words He uttered upon the cross: “Father, forgive them; for they know not what they do.”

At length Jesus declared, “It is finished: and he bowed his head, and gave up the ghost.”

These events, coupled with His glorious Resurrection, completed the final bridge of our trilogy: The Bridge of Obedience, the Bridge of Service, the Bridge of Prayer.

Jesus, the Bridge Builder, spanned that vast chasm we call death. “For as in Adam all die, even so in Christ shall all be made alive.” He did for us what we could not do for ourselves; hence, mankind can cross the bridges He built—into life eternal.

I close by paraphrasing the poem “The Bridge Builder”:

\begin{verse}
“You have crossed the chasm, deep and wide—\\
Why build you the bridge at the eventide?”
\end{verse}

That we may have the wisdom and determination to cross the bridges the Savior built for each of us is my sincere prayer, in the name of Jesus Christ, amen.

\newpage
\settalk{The Call for Courage}
\setspeaker{Thomas S. Monson}
\settalkdate{April 2004}
\talkheading{The Call for Courage}{Thomas S. Monson}

Brethren, you are an inspiring sight to behold. It is awesome to realize that in thousands of chapels throughout the world at this hour, your fellow holders of the priesthood of God are receiving this broadcast by way of satellite transmission. Your nationalities vary, and your languages are many, but a common thread binds us together. We have been entrusted to bear the priesthood and to act in the name of God. We are the recipients of a sacred trust. Much is expected of us.

Long ago the renowned author Charles Dickens wrote of opportunities that await. In his classic volume entitled Great Expectations, Dickens described a boy by the name of Philip Pirrip, more commonly known as Pip. Pip was born in unusual circumstances. He was an orphan. He wished with all his heart that he were a scholar and a gentleman. Yet all of his ambitions and all of his hopes seemed doomed to failure. Do you young men sometimes feel that way? Do those of us who are older entertain these same thoughts?

Then one day a London lawyer by the name of Jaggers approached little Pip and told him that an unknown benefactor had bequeathed to him a fortune. The lawyer put his arm around the shoulder of Pip and said to him, “My boy, you have great expectations.”

Tonight, as I look at you young men and realize who you are and what you may become, I declare, “You have great expectations”—not as the result of an unknown benefactor but as the result of a known benefactor, even our Heavenly Father, and great things are expected of you.

Life’s journey is not traveled on a freeway devoid of obstacles, pitfalls, and snares. Rather, it is a pathway marked by forks and turnings. Decisions are constantly before us. To make them wisely, courage is needed: the courage to say no, the courage to say yes. Decisions do determine destiny.

The call for courage comes constantly to each of us. It has ever been so, and so shall it ever be.

The courage of a military leader was recorded by a young infantryman wearing the gray uniform of the Confederacy during America’s Civil War. He describes the influence of General J. E. B. Stuart in these words:

“[At a critical point in the battle,] he waved his hand toward the enemy and shouted, ‘Forward men! Forward! Just follow me!’ …

“… With courage and resolution [they followed] after him like a wide raging torrent,” and the objective was seized and held.

At an earlier time, in a land far distant, another leader issued the same plea: “Follow me.” He was not a general of war. Rather, He was the Prince of Peace, the Son of God. Those who followed Him then and those who follow Him now win a far more significant victory, with consequences that are everlasting. The need for courage is constant.

The holy scriptures portray the evidence of this truth. Joseph, son of Jacob, the same who was sold into Egypt, demonstrated the firm resolve of courage when to Potiphar’s wife, who attempted to seduce him, he declared: “How … can I do this great wickedness, and sin against God? And … he hearkened not unto her … and got … out.”

In our day a father applied this example of courage to the lives of his children by declaring, “If you ever find yourself where you shouldn’t be, get out!”

Who can help but be inspired by the lives of the 2,000 stripling sons of Helaman, who taught and demonstrated the need of courage to follow the teachings of parents, the courage to be chaste and pure?

Perhaps each of these accounts is crowned by the example of Moroni, who had the courage to persevere to the end in righteousness.

All were fortified by the words of Moses: “Be strong and of a good courage, fear not, nor be afraid … : for the Lord thy God, he it is that doth go with thee; he will not fail thee, nor forsake thee.” He did not fail them. He will not fail us. He did not forsake them. He will not forsake us.

It is this sweet assurance that can guide you and me—in our time, in our day, in our lives. Of course we will face fear, experience ridicule, and meet opposition. Let us have the courage to defy the consensus, the courage to stand for principle. Courage, not compromise, brings the smile of God’s approval. Courage becomes a living and an attractive virtue when it is regarded not only as a willingness to die manfully, but also as a determination to live decently. A moral coward is one who is afraid to do what he thinks is right because others will disapprove or laugh. Remember that all men have their fears, but those who face their fears with dignity have courage as well.

From my personal chronology of courage, let me share with you an example from military service.

Entering the United States Navy in the closing months of World War II was a challenging experience for me. I learned of brave deeds, acts of valor, and examples of courage. One best remembered was the quiet courage of an 18-year-old seaman—not of our faith—who was not too proud to pray. Of 250 men in the company, he was the only one who each night knelt down by the side of his bunk—at times amidst the jeers of the curious, the jests of unbelievers—and, with bowed head, prayed to God. He never wavered. He never faltered. He had courage.

I love these words from the poet Ella Wheeler Wilcox:

\begin{verse}
It is easy enough to be pleasant,\\
When life flows by like a song,\\
But the man worth while is one who will smile,\\
When everything goes dead wrong.
\end{verse}

Such a man was Paul Tingey. Just a month ago I attended his funeral services here in Salt Lake City. Paul grew up in a fine Latter-day Saint home and served an honorable mission for the Lord in Germany. A companion of his in the mission field was Elder Bruce D. Porter of the First Quorum of the Seventy. Elder Porter described Elder Tingey as one of the most dedicated and successful missionaries he ever knew.

At the conclusion of his mission, Elder Tingey returned home, completed his studies at the university, married his sweetheart, and together with her reared their family. He served as a bishop and was successful in his vocation.

Then, without much warning, the symptoms of a dreaded disease struck his nervous system—even multiple sclerosis. Held captive by this malady, Paul Tingey struggled valiantly but then was confined to a care facility for the remainder of his life. There he cheered up the sad and made everyone feel glad. Whenever I attended Church meetings there, Paul lifted my spirits, as he did all others.

When the World Olympics came to Salt Lake City in 2002, Paul was selected to carry the Olympic torch for a specified distance. When this was announced at the care facility, a cheer erupted from those patients assembled and a hearty round of applause echoed through the halls. As I congratulated Paul, he said with his limited diction, “I hope I don’t drop the torch!”

Brethren, Paul Tingey didn’t drop the Olympic torch. What’s more, he carried bravely the torch he was handed in life and did so to the day of his passing.

Spirituality, faith, determination, courage—Paul Tingey had them all.

Someone has said that courage is not the absence of fear but the mastery of it. At times, courage is needed to rise from failure, to strive again.

As a young teenager I participated in a Church basketball game. When the outcome was in doubt, the coach sent me onto the playing floor right after the second half began. I took an inbounds pass, dribbled the ball toward the key, and let the shot fly. Just as the ball left my fingertips, I realized why the opposing guards did not attempt to stop my drive: I was shooting for the wrong basket! I offered a silent prayer: “Please, Father, don’t let that ball go in.” The ball rimmed the hoop and fell out.

From the bleachers came the call, “We want Monson, we want Monson, we want Monson—out!” The coach obliged.

Many years later, as a member of the Council of the Twelve, I joined other General Authorities in visiting a newly completed chapel where, as an experiment, we were trying out a tightly woven carpet on the gymnasium floor.

While several of us were examining the floor, Bishop J. Richard Clarke, who was then in the Presiding Bishopric, suddenly threw the basketball to me with a challenge: “I don’t believe you can hit the basket, standing where you are!”

I was some distance behind what is now the professional three-point line. I had never made such a basket in my entire life. Elder Mark E. Petersen of the Twelve called out to the others, “I think he can!”

My thoughts returned to my embarrassment of years before, shooting toward the wrong basket. Nevertheless, I aimed and let that ball fly. Through the net it went!

Throwing the ball in my direction, Bishop Clarke once more issued the challenge: “I know you can’t do that again!”

Elder Petersen spoke up, “Of course he can!”

The words of the poet echoed in my heart:

\begin{verse}
Lead us, O lead us,\\
Great Molder of men,\\
Out of the shadow\\
To strive once again.
\end{verse}

I shot the ball. It soared toward the basket and went right through.

That ended the inspection visit.

At lunchtime, Elder Petersen said to me, “You know, you could have been a starter in the NBA.”

Winning or losing in basketball fades from our thoughts when we contemplate our duties as bearers of the priesthood of God—both the Aaronic and Melchizedek Priesthood. We have a solemn duty to prepare ourselves through compliance with the commandments of the Lord and in responding to the calls we receive to serve Him.

We who have been ordained to the priesthood of God can make a difference. When we qualify for the help of the Lord, we can build boys, we can mend men, we can accomplish miracles in His holy service. Our opportunities are without limit.

Though the task seems large, we are strengthened by the truth “The greatest force in this world today is the power of God as it works through man.” If we are on the Lord’s errand, we are entitled to the Lord’s help. That divine help, however, is predicated upon our worthiness. To sail safely the seas of mortality, to perform a human rescue mission, we need the guidance of that eternal mariner, even the great Jehovah. We look up, we reach out to obtain heavenly help.

Are our reaching hands clean? Are our yearning hearts pure? Looking backward in time through the pages of history, we find a lesson on worthiness gleaned from the words of the dying King Darius. Through proper rites, Darius had been recognized as legitimate king of Egypt. His rival, Alexander the Great, had been declared legitimate son of Amon. He too was Pharaoh. Alexander, finding the defeated Darius on the point of death, laid his hands upon his head to heal him, commanding him to arise and resume his kingly power, concluding, “I swear unto thee, Darius, by all the gods that I do these things truly and without faking.” Darius replied with a gentle rebuke: “Alexander my boy … do you think you can touch heaven with those hands of yours?”

Brethren, as we learn our duty and magnify the callings which have come to us, the Lord will guide our efforts and touch the hearts of those whom we serve.

Many years ago I would visit an older widow named Mattie, whom I had known for many years and whose bishop I had been. My heart grieved at her utter loneliness. A precious son of hers lived many miles away, and for years he had not visited his mother. Mattie spent long hours in a lonely vigil at her front window. Behind a frayed and frequently opened curtain, the disappointed mother would say to herself, “Dick will come; Dick will come.”

But Dick didn’t come. The years passed by one after another. Then, like a ray of sunshine, Church activity came into the life of Dick, one of my former Aaronic Priesthood boys, who now lived in Houston, Texas, far away from his mother. He journeyed to Salt Lake to visit with me. He telephoned upon his arrival and, with excitement, reported the change in his life. He asked if I had time to see him if he were to come directly to my office. My response was one of gladness. However, I said, “Dick, first visit your mother and then come to see me.” He gladly complied with my request.

Before he could get to my office, there came a phone call from Mattie, his mother. From a joyful heart came words punctuated by tears: “Bishop, I knew Dick would come. I told you he would. I saw him coming through the window.”

Not many years later at Mattie’s funeral, Dick and I spoke tenderly of that experience. We had witnessed a glimpse of God’s healing power through the window of a mother’s faith in her son.

Time marches on. Duty keeps cadence with that march. Duty does not dim nor diminish. Catastrophic conflicts come and go, but the war waged for the souls of men continues without abatement. Like a clarion call comes the word of the Lord to you, to me, and to priesthood holders everywhere: “Wherefore, now let every man learn his duty, and to act in the office in which he is appointed, in all diligence.”

May we each have the courage to do so, I pray, in the name of Jesus Christ, amen.

\newpage
\settalk{Your Personal Influence}
\setspeaker{Thomas S. Monson}
\settalkdate{April 2004}
\talkheading{Your Personal Influence}{Thomas S. Monson}

My dear brothers and sisters, both within my view and assembled throughout the world, I seek an interest in your prayers and your faith as I respond to the assignment and privilege to address you.

More than 40 years ago, when President David O. McKay extended to me a call to the Quorum of the Twelve Apostles, he warmly welcomed me with a heartfelt smile and a tender embrace. Among the sacred counsel he extended was the declaration, “There is one responsibility that no one can evade. That is the effect of one’s personal influence.”

The calling of the early Apostles reflected the influence of the Lord. When He sought a man of faith, He did not select him from the throng of the self-righteous who were found regularly in the synagogue. Rather, He called him from among the fishermen of Capernaum. Peter, Andrew, James, and John heard the call, “Follow me, and I will make you fishers of men.” They followed. Simon, man of doubt, became Peter, Apostle of faith.

When the Savior was to choose a missionary of zeal and power, He found him not among His advocates but amidst His adversaries. Saul of Tarsus—the persecutor—became Paul the proselyter. The Redeemer chose imperfect men to teach the way to perfection. He did so then; He does so now.

He calls you and me to serve Him here below and sets us to the task He would have us fulfill. The commitment is total. There is no conflict of conscience.

As we follow that Man of Galilee—even the Lord Jesus Christ—our personal influence will be felt for good wherever we are, whatever our callings.

Our appointed task may appear insignificant, unnecessary, unnoticed. Some may be tempted to question:

\begin{verse}
“Father, where shall I work today?”\\
And my love flowed warm and free.\\
Then he pointed out a tiny spot\\
And said, “Tend that for me.”\\
I answered quickly, “Oh no, not that!\\
Why, no one would ever see,\\
No matter how well my work was done.\\
Not that little place for me.”\\
And the word he spoke, it was not stern; …\\
“Art thou working for them or for me?\\
Nazareth was a little place,\\
And so was Galilee.”
\end{verse}

The family is the ideal place for teaching. It is also a laboratory for learning. Family home evening can bring spiritual growth to each member.

“The home is the basis of a righteous life, and no other instrumentality can take its place or fulfill its essential functions.” Such truth has been taught by many Presidents of the Church.

It is in the home where fathers and mothers can teach provident living to their children. Sharing of tasks and helping one another set a pattern for future families as children grow, marry, and leave home. The lessons learned in the home are those that last the longest. President Gordon B. Hinckley continues to stress the avoidance of unnecessary debt, the fallacy of living beyond our means, and the temptation to let our wants become our necessities.

The Apostle Paul’s exhortation to his beloved Timothy provides the counsel that will enable our personal influence to find lodgment in the hearts of those with whom we associate: “Be thou an example of the believers, in word, in conversation, in charity, in spirit, in faith, in purity.”

When I was a boy, our family lived in the Sixth-Seventh Ward of the Pioneer Stake. The ward population was rather transient, which resulted in an accelerated rate of turnover with respect to the teachers in the Sunday School. As boys and girls we would just become acquainted with a particular teacher and grow to appreciate him or her when the Sunday School superintendent would visit the class and introduce a new teacher. Disappointment filled each heart, and a breakdown of discipline resulted.

Prospective teachers, hearing of the unsavory reputation of our particular class, would graciously decline to serve or suggest the possibility of teaching a different class where the students were more manageable. We took delight in our newly found status and determined to live up to the fears of the faculty.

One Sunday morning a lovely young lady accompanied the superintendent into the classroom and was presented to us as a teacher who requested the opportunity to teach us. We learned that she had been a missionary and loved young people. Her name was Lucy Gertsch. She was beautiful, soft-spoken, and interested in us. She asked each class member to introduce himself, and then she asked questions which gave her an understanding and insight into the background of each. She told us of her girlhood in Midway, Utah, and as she described that beautiful valley she made its beauty live within us and we desired to visit the green fields she loved so much.

When Lucy taught, she made the scriptures actually live. We became personally acquainted with Samuel, David, Jacob, Nephi, Joseph Smith, and the Lord Jesus Christ. Our gospel scholarship grew. Our deportment improved. Our love for Lucy Gertsch knew no bounds.

We undertook a project to save nickels and dimes for what was to be a gigantic Christmas party. Sister Gertsch kept a careful record of our progress. As boys with typical appetites, we converted in our minds the monetary totals to cakes, cookies, pies, and ice cream. This was to be a glorious event. Never before had any of our teachers even suggested a social event like this was to be.

The summer months faded into autumn. Autumn turned to winter. Our party goal had been achieved. The class had grown. A good spirit prevailed.

None of us will forget that gray morning when our beloved teacher announced to us that the mother of one of our classmates had passed away. We thought of our own mothers and how much they meant to us. We felt sincere sorrow for Billy Devenport in his great loss.

The lesson this Sunday was from the book of Acts, chapter 20, verse 35: “Remember the words of the Lord Jesus, how he said, It is more blessed to give than to receive.” At the conclusion of the presentation of a well-prepared lesson, Lucy Gertsch commented on the economic situation of Billy’s family. These were Depression times, and money was scarce. With a twinkle in her eyes, she asked: “How would you like to follow this teaching of our Lord? How would you feel about taking our party fund and, as a class, giving it to the Devenports as an expression of our love?” The decision was unanimous. We counted so carefully each penny and placed the total sum in a large envelope. A beautiful card was purchased and inscribed with our names.

This simple act of kindness welded us together as one. We learned through our own experience that it is indeed more blessed to give than to receive.

The years have flown. The old chapel is gone, a victim of industrialization. The boys and girls who learned, who laughed, who grew under the direction of that inspired teacher of truth have never forgotten her love or her lessons. Her personal influence for good was contagious.

A General Authority whose personal influence was felt far and wide was the late President Spencer W. Kimball. He really made a difference in the lives of countless individuals.

When I was a bishop, the telephone rang one day, and the caller identified himself as Elder Spencer W. Kimball. He said, “Bishop Monson, in your ward is a trailer court, and in a little trailer in that court—the smallest trailer of all—is a sweet Navajo widow, Margaret Bird. Would you have your Relief Society president visit her and invite her to come to Relief Society and to participate with the sisters?” We did. Margaret Bird came and found a warm welcome.

Elder Kimball called on another occasion. “Bishop Monson,” he said, “I have learned that there are two Samoan boys living in a downtown hotel. They’re going to get in trouble. Will you make them members of your ward?”

I found these two boys at midnight sitting on the steps of the hotel playing ukuleles and singing. They became members of our ward. Eventually each of them married in the temple and served valiantly. Their influence for good was widespread.

When I was first called as a bishop, I discovered that our record for subscriptions to the Relief Society Magazine in the Sixth-Seventh Ward had been at a low ebb. Prayerfully we analyzed the names of individuals whom we could call to be magazine representative. The inspiration dictated that Elizabeth Keachie should be given the assignment. As her bishop, I approached her with the task. She responded, “Bishop Monson, I’ll do it.”

Elizabeth Keachie was of Scottish descent, and when she replied, “I’ll do it,” one knew she indeed would. She and her sister-in-law, Helen Ivory—neither more than five feet tall—commenced to walk the ward, house by house, street by street, and block by block. The result was phenomenal. We had more subscriptions to the Relief Society Magazine than had been recorded by all the other units of the stake combined.

I congratulated Elizabeth Keachie one Sunday evening and said to her, “Your task is done.”

She replied, “Not yet, Bishop. There are two square blocks we have not yet covered.”

When she told me which blocks they were, I said, “Oh, Sister Keachie, no one lives on those blocks. They are totally industrial.”

“Just the same,” she said, “I’ll feel better if Nell and I go and check them ourselves.”

On a rainy day she and Nell covered those final two blocks. On the first one she found no home, nor did she on the second. She and Sister Ivory paused, however, at a driveway which was muddy from a recent storm. Sister Keachie gazed about 100 feet (30 m) down the driveway, which was adjacent to a machine shop, and there noticed a garage. This was not a normal garage, however, in that there was a curtain at the window.

She turned to her companion and said, “Nell, shall we go and investigate?”

The two sweet sisters then walked down the muddy driveway 40 feet (12 m) to a point where the entire view of the garage could be seen. Now they noticed a door which had been cut into the side of the garage, which door was unseen from the street. They also noticed that there was a chimney with smoke rising from it.

Elizabeth Keachie knocked at the door. A man 68 years of age, William Ringwood, answered. They then presented their story concerning the need of every home having the Relief Society Magazine. William Ringwood replied, “You’d better ask my father.”

Ninety-four-year-old Charles W. Ringwood then came to the door and also listened to the message. He subscribed.

Elizabeth Keachie reported to me the presence of these two men in our ward. When I requested their membership certificates from Church headquarters, I received a call from the Membership Department at the Presiding Bishopric’s Office. The clerk said, “Are you sure you have living in your ward Charles W. Ringwood?”

I replied that I did, whereupon she reported that the membership certificate for him had remained in the “lost and unknown” file of the Presiding Bishopric’s Office for the previous 16 years.

On Sunday morning Elizabeth Keachie and Nell Ivory brought to our priesthood meeting Charles and William Ringwood. This was the first time they had been inside a chapel for many years. Charles Ringwood was the oldest deacon I had ever met. His son was the oldest male member holding no priesthood I had ever met.

It became my opportunity to ordain Brother Charles Ringwood a teacher and then a priest and finally an elder. I shall never forget his interview with respect to seeking a temple recommend. He handed me a silver dollar, which he took from an old, worn leather coin purse, and said, “This is my fast offering.”

I said, “Brother Ringwood, you owe no fast offering. You need it yourself.”

“I want to receive the blessings, not retain the money,” he responded.

It was my opportunity to take Charles Ringwood to the Salt Lake Temple and to attend with him the endowment session.

Within a few months, Charles W. Ringwood passed away. At his funeral service I noticed his family sitting on the front rows in the mortuary chapel, but I noticed also two sweet women sitting near the rear of the chapel, Elizabeth Keachie and Helen Ivory.

As I gazed upon those two faithful and dedicated women and contemplated their personal influence for good, the promise of the Lord filled my very soul: “I, the Lord, am merciful and gracious unto those who fear me, and delight to honor those who serve me in righteousness and in truth unto the end. Great shall be their reward and eternal shall be their glory.”

There is one, above all others, whose personal influence covers the continents, spans the oceans, and penetrates the hearts of true believers. He atoned for the sins of mankind.

I testify that He is a teacher of truth—but He is more than a teacher. He is the Exemplar of the perfect life—but He is more than an exemplar. He is the Great Physician—but He is more than a physician. He is the literal Savior of the world, the Son of God, the Prince of Peace, the Holy One of Israel, even the risen Lord, who declared:

“I am Jesus Christ, whom the prophets testified shall come into the world. … I am the light and the life of the world.”

“I am the first and the last; I am he who liveth, I am he who was slain; I am your advocate with the Father.”

As His witness, I testify to you that He lives! In His holy name—even Jesus Christ, the Savior—amen.

\newpage
\settalk{Anxiously Engaged}
\setspeaker{Thomas S. Monson}
\settalkdate{October 2004}
\talkheading{Anxiously Engaged}{Thomas S. Monson}

My dear brethren, it is a solemn and somewhat humbling experience to stand before you this evening and respond to the invitation to teach and to testify concerning the sacred privilege which is ours to bear the priesthood of God. I pray for your faith and your prayers in my behalf.

Beyond those who hold the Aaronic and Melchizedek Priesthood who are in attendance this evening here in this beautiful Conference Center or seated in locations worldwide, there are vast numbers of priesthood bearers who, for whatever reason, have drifted from their duties and have chosen to pursue other pathways.

The Lord speaks rather plainly to us to reach out and rescue such individuals and bring them and theirs to the table of the Lord. We well could pay heed to the Lord’s divine instructions when He declared, “Wherefore, now let every man learn his duty, and to act in the office in which he is appointed, in all diligence.” He added:

“For behold, it is not meet that I should command in all things; for he that is compelled in all things, the same is a slothful and not a wise servant; wherefore he receiveth no reward.

“Verily I say, men should be anxiously engaged in a good cause, and do many things of their own free will, and bring to pass much righteousness;

“For the power is in them, wherein they are agents unto themselves. And inasmuch as men do good they shall in nowise lose their reward.”

The sacred scriptures provide for you and me a model to follow when they declare, “Jesus increased in wisdom and stature, and in favour with God and man.” And He “went about doing good, … for God was with him.”

I have observed in studying the life of the Master that His lasting lessons and His marvelous miracles usually occurred when He was doing His Father’s work. On the way to Emmaus He appeared with a body of flesh and bones. He partook of food and testified of His divinity. All of this took place after He had exited the tomb.

At an earlier time, it was while He was on the road to Jericho that He restored sight to one who was blind.

The Savior was ever up and about—teaching, testifying, and saving others. Such is our individual duty as members of priesthood quorums today.

In a proclamation of the First Presidency and the Quorum of the Twelve Apostles issued on April 6, 1980, this declaration of testimony and truth was set forth:

“We solemnly affirm that The Church of Jesus Christ of Latter-day Saints is in fact a restoration of the Church established by the Son of God, when in mortality he organized his work upon the earth; that it carries his sacred name, even the name of Jesus Christ; that it is built upon a foundation of Apostles and prophets, he being the chief cornerstone; that its priesthood, in both the Aaronic and Melchizedek orders, was restored under the hands of those who held it anciently: John the Baptist, in the case of the Aaronic; and Peter, James, and John in the case of the Melchizedek.”

On October 6, 1889, President George Q. Cannon expressed this plea:

“I want to see the power of the Priesthood strengthened. … I want to see this strength and power diffused through the entire body of the Priesthood, reaching from the head down to the least and most humble deacon in the Church. Every man should seek for and enjoy the revelations of God, the light of heaven shining in his soul and giving unto him knowledge concerning his duties, concerning that portion of the work of God that devolves upon him in his Priesthood.”

I share with you tonight two experiences from my life—one which took place when I was a boy and the other pertaining to a friend of mine who was a husband and father of children.

Not long after my ordination as a teacher in the Aaronic Priesthood, I was called to serve as president of the quorum. Our adviser, Harold, was interested in us, and we knew it. One day he said to me, “Tom, you enjoy raising pigeons, don’t you?”

I responded with a warm, “Yes.”

Then he proffered, “How would you like me to give you a pair of purebred Birmingham Roller pigeons?”

This time I answered, “Yes, Sir!” You see, the pigeons I had were just the common variety, trapped on the roof of the Grant Elementary School.

He invited me to come to his home the next evening. The following day was one of the longest in my young life. I was awaiting my adviser’s return from work an hour before he arrived home. He took me to his pigeon loft, which was in the upper area of a small barn located at the rear of his yard. As I looked at the most beautiful pigeons I had yet seen, he said, “Select any male, and I will give you a female which is different from any other pigeon in the world.” I made my selection. He then placed in my hand a tiny hen pigeon. I asked what made her so different. He responded, “Look carefully, and you’ll notice that she has but one eye.” Sure enough, one eye was missing, a cat having done the damage. “Take them home to your loft,” he counseled. “Keep them in for about 10 days, and then turn them out to see if they will remain at your place.”

I followed Harold’s instructions. Upon his release, the male pigeon strutted about the roof of the loft, then returned inside to eat. But the one-eyed female was gone in an instant. I called Harold and asked, “Did that one-eyed pigeon return to your loft?”

“Come on over,” he said, “and we’ll have a look.”

As we walked from his kitchen door to the loft, my adviser commented, “Tom, you are the president of the teachers quorum.” This, of course, I already knew. Then he added, “What are you going to do to activate Bob, who is a member of your quorum?”

I answered, “I’ll have him at quorum meeting this week.”

Then he reached up to a special nest and handed me the one-eyed pigeon. “Keep her in a few more days and try again.” This I did, and once more she disappeared. Again the experience: “Come on over, and we’ll see if she returned home.” Came the comment as we walked to the loft, “Congratulations on getting Bob to priesthood meeting. Now what are you and Bob going to do to activate Bill?”

“We’ll have him there next week,” I volunteered.

This experience was repeated over and over again. I was a grown man before I fully realized that indeed Harold, my adviser, had given me a special pigeon, the only pigeon in his loft he knew would return every time she was released. It was his inspired way of having an ideal personal priesthood interview with the president of the teachers quorum every two weeks. I owe a lot to that one-eyed pigeon. I owe more to that quorum adviser. He had the patience and the skill to help me prepare for the responsibilities which lay ahead.

Fathers, grandfathers, we have an even greater responsibility to guide our precious sons and grandsons. They need our help; they need our encouragement; they need our example. It has been wisely said that our youth need fewer critics and more models to follow.

Now for the illustration pertaining to those men whose habits and lives include but little Church attendance or Church activity of any kind. The ranks of these prospective elders have grown larger. This is because of those younger boys of the Aaronic Priesthood quorums who are lost along the Aaronic Priesthood pathway and also those grown men who are baptized but do not persevere in activity and faith so that they might be ordained elders.

I not only reflect on the hearts and souls of such individual men, but also sorrow for their sweet wives and growing children. These men await a helping hand, an encouraging word, and a personal testimony of truth expressed from a heart filled with love and a desire to lift and to build.

Shelley, my friend, was such a person. His wife and children were fine members, but all efforts to motivate him toward baptism and then priesthood blessings had miserably failed.

But then Shelley’s mother died. Shelley was so sorrowful that he retired to a special room at the mortuary where the funeral was being held. We had wired the proceedings to this room so that he might mourn alone and where no one could see him weep with sorrow. As I comforted him in that room before going to the pulpit, he gave me a hug, and I knew a tender chord had been touched.

Time passed. Shelley and his family moved to another part of the city. I was called to preside over the Canadian Mission and, together with my family, moved to Toronto, Canada, for a three-year period.

When I returned and after I was called to the Twelve, Shelley telephoned me. He said, “Bishop, will you seal my wife, my family, and me in the Salt Lake Temple?”

I answered hesitantly, “But Shelley, you must first be baptized a member of the Church.”

He laughed and responded, “Oh, I took care of that while you were in Canada. I sort of snuck up on you. There was this home teacher who called on us regularly and taught me the truths of the Church. He was a school crossing guard and helped the small children across the street each morning when they went to school and each afternoon when they went home. He asked me to help him. During the intervals when there was no child crossing, he gave me additional instruction pertaining to the Church.”

I had the privilege to see this miracle with my own eyes and feel the joy with my heart and soul. The sealings were performed; a family was united. Shelley died not too long after this period. I had the privilege of speaking at his funeral services. I shall ever see, in memory’s eye, the body of my friend Shelley lying in his casket, dressed in his temple clothing. I readily admit the presence of tears—tears of gratitude, for the lost had been found.

Those who have felt the touch of the Master’s hand somehow cannot explain the change which comes into their lives. There is a desire to live better, to serve faithfully, to walk humbly, and to be more like the Savior. Having received their spiritual eyesight and glimpsed the promises of eternity, they echo the words of the blind man to whom Jesus restored sight: “One thing I know, that, whereas I was blind, now I see.”

How can we account for these miracles? Why the upsurge of activity in men long dormant? The poet, speaking of death, wrote, “God … touch’d him, and he slept.” I say, speaking of this new birth, “God touched them, and they awakened.”

Two fundamental reasons largely account for these changes of attitudes, of habits, of actions.

First, men have been shown their eternal possibilities and have made the decision to achieve them. They cannot really long rest content with mediocrity once excellence is within their reach.

Second, other men and women and, yes, young people have followed the admonition of the Savior and have loved their neighbors as themselves and helped to bring their neighbors’ dreams to fulfillment and their ambitions to realization.

The catalyst in this process has been the principle of love.

The passage of time has not altered the capacity of the Redeemer to change men’s lives. As He said to the dead Lazarus, so He says to you and to me, “Come forth.” I add: Come forth from the despair of doubt. Come forth from the sorrow of sin. Come forth from the death of disbelief. Come forth to a newness of life.

As we do and direct our footsteps along the paths which Jesus walked, let us remember the testimony Jesus gave: “Behold, I am Jesus Christ, whom the prophets testified shall come into the world. … I am the light and … life of the world.” “I am the first and the last; I am he who liveth, I am he who was slain; I am your advocate with the Father.”

There are quorum members and those who should be our quorum members who require our help. John Milton wrote in his poem “Lycidas,” “The hungry sheep look up, and are not fed.” The Lord Himself said to Ezekiel the prophet, “Woe be to the shepherds of Israel that … feed not the flock.”

My brethren of the priesthood, the task is ours. Let us remember and never forget, however, that such an undertaking is not insurmountable. Miracles are everywhere to be seen when priesthood callings are magnified. When faith replaces doubt, when selfless service eliminates selfish striving, the power of God brings to pass His purposes. We are on the Lord’s errand. We are entitled to the Lord’s help. But we must try. From the play Shenandoah comes the spoken line which inspires: “If we don’t try, then we don’t do; and if we don’t do, then why are we here?”

Let us, one and all, be doers of the word and not hearers only. Let us follow the example of our President, Gordon B. Hinckley, the Lord’s prophet.

May we, as did the Savior’s followers of old, respond to the invitation, “Follow me, and I will make you fishers of men.” That we may do so is my prayer, in the name of Jesus Christ, amen.

\newpage
\settalk{Choose You This Day}
\setspeaker{Thomas S. Monson}
\settalkdate{October 2004}
\talkheading{Choose You This Day}{Thomas S. Monson}

My dear brothers and sisters, both within my view and assembled throughout the world, I seek an interest in your faith and prayers as I respond to the assignment and privilege to address you. First, however, I should like to extend a personal welcome to Elders Dieter Uchtdorf and David Bednar, our new members of the Quorum of the Twelve Apostles.

I have been thinking recently about choices and their consequences. It has been said that the gate of history turns on small hinges, and so do people’s lives. The choices we make determine our destiny.

Joshua of old declared, “Choose you this day whom ye will serve; … but as for me and my house, we will serve the Lord.”

All of us commenced an awesome and vital journey when we left the spirit world and entered this often challenging stage called mortality. We brought with us that great gift from God—our agency. Said the prophet Wilford Woodruff: “God has given unto all of His children … individual agency. … [We] possessed it in the heaven of heavens before the world was, and the Lord maintained and defended it there against the aggression of Lucifer. … By virtue of this agency you and I and all mankind are made responsible beings, responsible for the course we pursue, the lives we live, the deeds we do.”

Brigham Young said, “All must use [this agency] in order to gain exaltation in [God’s] kingdom; inasmuch as [we] have the power of choice [we] must exercise that power.”

The scriptures tell us that we are free to act for ourselves, “to choose the way of everlasting death or the way of eternal life.”

A familiar hymn provides inspiration in the choices we make:

\begin{verse}
Choose the right when a choice is placed before you.\\
In the right the Holy Spirit guides;\\
And its light is forever shining o’er you,\\
When in the right your heart confides. …
\end{verse}

Do we have a guide to help us choose the right and avoid dangerous detours? Positioned on the wall of my office, directly opposite my desk, is a lovely print of the Savior, painted by Heinrich Hofmann. I love the painting, which I have had since I was a 22-year-old bishop and which I have taken with me wherever I have been assigned to labor. I have tried to pattern my life after the Master. Whenever I have a difficult decision to make, I have looked at that picture and asked myself, “What would He do?” Then I try to do it. We can never go wrong when we choose to follow the Savior.

Some choices may seem more important than others, but no choice is insignificant.

Some years ago I held in my hand a guide which, if followed, will never fail in helping us to make correct choices. It was a volume of scripture we commonly call the triple combination, containing the Book of Mormon, Doctrine and Covenants, and Pearl of Great Price. This book was a gift from a loving father to a precious daughter who followed carefully his advice. On the flyleaf page, her father had written in his own hand these inspired words:

“To my dear Maurine,

“That you may have a constant measure by which to judge between truth and the errors of man’s philosophies, and thus grow in spirituality as you increase in knowledge, I give you this sacred book to read frequently and cherish throughout your life.

“Lovingly your father,

“Harold B. Lee”

As members of The Church of Jesus Christ of Latter-day Saints, our goal is to obtain celestial glory.

Let us not find ourselves as indecisive as is Alice in Lewis Carroll’s classic Alice’s Adventures in Wonderland. You will remember that she comes to a crossroads with two paths before her, each stretching onward but in opposite directions. She is confronted by the Cheshire cat, of whom Alice asks, “Which path shall I follow?”

The cat answers: “That depends where you want to go. If you do not know where you want to go, it doesn’t matter which path you take.”

Unlike Alice, we all know where we want to go, and it does matter which way we go, for the path we follow in this life surely leads to the path we will follow in the next.

Each of us should remember that he or she is a son or daughter of God, endowed with faith, gifted with courage, and guided by prayer. Our eternal destiny is before us. The Apostle Paul speaks to us today as he spoke to Timothy long years ago: “Neglect not the gift that is in thee.” “O Timothy, keep that which is committed to thy trust.”

At times many of us let that enemy of achievement—even the culprit “self-defeat”—dwarf our aspirations, smother our dreams, cloud our vision, and impair our lives. The enemy’s voice whispers in our ears, “You can’t do it.” “You’re too young.” “You’re too old.” “You’re nobody.” This is when we remember that we are created in the image of God. Reflection on this truth provides a profound sense of strength and power.

Mine was the privilege to know rather intimately President J. Reuben Clark Jr., who served for so many years as a member of the First Presidency. As I assisted him in the preparation for printing his monumental books, priceless lessons were learned. One day while in a somber, reflective mood, President Clark asked if I could arrange for the printing of a picture suitable for framing. The picture was to feature the lions of Persepolis guarding the ruins of a crumbled glory. President Clark wished to have printed with the picture—between the decaying arches of a civilization that was no more—a number of his favorite scriptures, chosen from his vast knowledge of holy writ. I felt you would wish to know his selections. There were three—two from Ecclesiastes and one from the Gospel of John.

First, from Ecclesiastes: “Fear God, and keep his commandments: for this is the whole duty of man.”

Second, “Vanity of vanities, saith the Preacher, vanity of vanities; all is vanity.”

Third, from John: “And this is life eternal, that they might know thee the only true God, and Jesus Christ, whom thou hast sent.”

An earlier prophet, even Moroni, writing in what is now the Book of Mormon, counseled, “And now, I would commend you to seek this Jesus of whom the prophets and apostles have written, that the grace of God the Father, and also the Lord Jesus Christ, and the Holy Ghost, which beareth record of them, may be and abide in you forever.”

President David O. McKay counseled: “‘The greatest battle of life is fought within the silent chambers of your own soul.’ … It is a good thing to sit down and commune with yourself, to come to an understanding with yourself and decide in that silent moment what your duty is to your family, to your Church, to your country, and … to your fellowmen.”

The boy prophet Joseph Smith sought heavenly help by entering a grove which then became sacred. Do we need similar strength? Does each need to seek his or her own “Sacred Grove”? A place where communication between God and man can go forth unimpeded, uninterrupted, and undisturbed is such a grove.

In the New Testament we learn that it is impossible to take a right attitude toward Christ without taking an unselfish attitude toward men. In the book of Matthew, Jesus taught, “Inasmuch as ye have done it unto one of the least of these my brethren, ye have done it unto me.”

When the Savior sought a man of faith, He did not select him from the throng of self-righteous who were found regularly in the synagogue. Rather, He called him from among the fishermen of Capernaum. While teaching on the seashore, He saw two ships standing by the lake. He entered one and asked its owner to put it out a little from the land so He might not be pressed upon by the crowd. After teaching further, He said to Simon, “Launch out into the deep, and let down your nets.”

Simon answered, “Master, we have toiled all the night, and have taken nothing: nevertheless at thy word I will let down the net.

“And when they had this done, they inclosed a great multitude of fishes. …

“When Simon Peter saw it, he fell down at Jesus’ knees, saying, Depart from me; for I am a sinful man, O Lord.”

Came the reply, “Follow me, and I will make you fishers of men.”

Simon the fisherman had received his call. Doubting, disbelieving, unschooled, untrained, impetuous Simon did not find the way of the Lord a highway of ease nor a path free from pain. He was to hear the rebuke, “O thou of little faith.” Yet when the Master asked him, “Whom say ye that I am?” Peter answered, “Thou art the Christ, the Son of the living God.”

Simon, man of doubt, had become Peter, Apostle of faith. Peter made his choice.

When the Savior was to choose a missionary of zeal and power, He found him not among His advocates but amidst His adversaries. The experience of Damascus’s way changed Saul. Of him the Lord declared, “He is a chosen vessel unto me, to bear my name before the Gentiles, and kings, and the children of Israel.”

Saul the persecutor became Paul the proselyter. Paul made his choice.

Acts of selfless service are performed daily by countless members of the Church. There are many which are freely given, with no fanfare or boasting, but rather through quiet love and tender care. Let me share with you the example of one who made such a simple yet profound choice to serve.

A few years ago, Sister Monson and I were in the city of Toronto, where we once lived when I was the mission president. Olive Davies, the wife of the first stake president in Toronto, was gravely ill and preparing to pass from this life. Her illness required her to leave her cherished home and enter a hospital which could provide the care she needed. Her only child lived with her own family far away in the West.

I attempted to comfort Sister Davies, but she had present with her the comfort she longed to have. A stalwart grandson sat silently next to his grandmother. I learned he had spent most of the summer away from his university studies, that he might serve his grandmother’s needs. I said to him, “Shawn, you will never regret your decision. Your grandmother feels you are heaven-sent, an answer to her prayers.”

He replied, “I chose to come because I love her and know this is what my Heavenly Father would have me do.”

Tears were near the surface. Grandmother told us how she enjoyed being helped by her grandson and introducing him to each employee and every patient in the hospital. Hand in hand, they walked the halls, and during the night he was close by.

Olive Davies has passed on to her reward, there to meet her faithful husband and together continue an eternal journey. In a grandson’s heart there will ever remain those words, “Choose the right when a choice is placed before you. In the right the Holy Spirit guides.”

Such are foundation stones in building one’s personal temple. As the Apostle Paul counseled, “Know ye not that ye are the temple of God, and that the Spirit of God dwelleth in you?”

May I leave with you today a simple yet far-reaching formula to guide you in the choices of life:

Fill your minds with truth.

Fill your hearts with love.

Fill your lives with service.

By doing so, may we one day hear the plaudit from our Lord and Savior, “Well done, good and faithful servant; thou hast been faithful over a few things, I will make thee ruler over many things: enter thou into the joy of thy lord.”

In the name of Jesus Christ, amen.

\newpage
\settalk{If Ye Are Prepared Ye Shall Not Fear}
\setspeaker{Thomas S. Monson}
\settalkdate{October 2004}
\talkheading{If Ye Are Prepared Ye Shall Not Fear}{Thomas S. Monson}

It is a privilege to stand before you at this general Relief Society conference. I recognize that beyond you who are gathered in this Conference Center, there are many thousands watching and listening to the proceedings by way of satellite transmission.

As I speak to you tonight, I realize that as a man I am in the minority and must be cautious in my comments. I feel much the same as the shy country cousin who came to visit his relative in a large cosmopolitan city. He had not sought his kinsman for some years and was startled when a young boy answered the ringing of the doorbell. The lad asked him in; and after they were comfortably seated, he inquired, “Who are you, anyway?”

The visitor answered, “I’m a cousin on your father’s side,” whereupon the boy replied, “Mister, in this house, that puts you on the wrong side!”

I trust that tonight, in this house, I might be found on the right side, even the Lord’s side.

Years ago I saw a photograph of a Sunday School class in the Sixth Ward of the Pioneer Stake in Salt Lake City. The photograph was taken in 1905. A sweet girl, her hair in pigtails, was shown on the front row. Her name was Belle Smith. Later, as Belle Smith Spafford, general president of the Relief Society, she wrote: “Never have women had greater influence than in today’s world. Never have the doors of opportunity opened wider for them. This is an inviting, exciting, challenging, and demanding period of time for women. It is a time rich in rewards if we keep our balance, learn the true values of life, and wisely determine priorities.”

The Relief Society organization has had a goal to help eliminate illiteracy. Those of us who can read and write do not appreciate the deprivation of those who cannot read, who cannot write. They are shrouded by a dark cloud which stifles their progress, dulls their intellect, and dims their hopes. Sisters of the Relief Society, you can lift this cloud of despair and welcome heaven’s divine light as it shines upon your sisters.

Some years ago I was in Monroe, Louisiana, attending a regional conference. It was a beautiful occasion. At the airport on my way home, I was approached by a lovely African-American woman—a member of the Church—who said, smiling broadly, “President Monson, before I joined the Church and became a member of Relief Society, I could not read nor write. None of my family could. You see, we were all poor sharecroppers. President, my white Relief Society sisters—they taught me to read. They taught me to write. Now I help teach my white sisters how to read and how to write.” I reflected on the supreme joy she must have felt when she opened her Bible and read for the first time the words of the Lord:

“Come unto me, all ye that labour and are heavy laden, and I will give you rest.

“Take my yoke upon you, and learn of me; for I am meek and lowly in heart: and ye shall find rest unto your souls.

“For my yoke is easy, and my burden is light.”

That day in Monroe, Louisiana, I received a confirmation by the Spirit of the exalted objective of the Relief Society to help eliminate illiteracy.

The poet wrote:

\begin{verse}
You may have tangible wealth untold;\\
Caskets of jewels and coffers of gold.\\
Richer than I you can never be—\\
I had a Mother who read to me.
\end{verse}

Another added this poignant verse:

\begin{verse}
But think of the fate of a different child,\\
Whose manner is meek, whose temper is mild,\\
While yet instilled with that same special need,\\
Was born to a mother who could not read.
\end{verse}

Parents everywhere have a concern for their children and for their eternal happiness. This is depicted in the musical Fiddler on the Roof, one of the longest running musicals in the history of the stage.

One laughs as he observes the old-fashioned father of a Jewish family in Russia as he attempts to cope with the changing times brought forcibly home to him by his beautiful teenage daughters.

The gaiety of the dance, the rhythm of the music, the excellence of the acting all fade in their significance when old Tevye speaks what to me becomes the message of the musical. He gathers his lovely daughters to his side and, in the simplicity of his peasant surroundings, counsels them as they ponder their future. Remember, cautions Tevye, “in Anatevka … everyone knows who he is and what God expects him to do.”

You, my beloved sisters, know who you are and what God expects you to become. Your challenge is to bring all for whom you are responsible to a knowledge of this truth. The Relief Society of this, the Lord’s Church, can be the means to achieve such a goal.

“The first and foremost opportunity for teaching in the Church lies in the home,” observed President David O. McKay. “A true Mormon home is one in which if Christ should chance to enter, he would be pleased to linger and to rest.”

What are we doing to ensure that our homes meet this description? It isn’t enough for parents alone to have strong testimonies. Children can ride only so long on the coattails of a parent’s conviction.

President Heber J. Grant declared: “It is our duty to teach our children in their youth. … I may know that the gospel is true, and so may my wife; but I want to tell you that our children will not know that the gospel is true, unless they study it and gain a testimony for themselves.”

A love for the Savior, a reverence for His name, and genuine respect one for another will provide a fertile seedbed for a testimony to grow.

Learning the gospel, bearing a testimony, leading a family are rarely if ever simple processes. Life’s journey is characterized by bumps in the road, swells in the sea—even the turbulence of our times.

Some years ago, while visiting the members and missionaries in Australia, I witnessed a sublime example depicting how a treasury of testimony can bless and sanctify a home. The mission president, Horace D. Ensign, and I were traveling by plane the long distance from Sydney to Darwin, where I was to break ground for our first chapel in that city. En route we had a scheduled fueling stop at a remote mining community named Mt. Isaiah As we entered the small airport, a woman and her two young children approached. She said, “I am Judith Louden, a member of the Church, and these are my children. We thought you might be on this flight, so we have come to visit with you during your brief stopover.” She explained that her husband was not a member of the Church and that she and the children were indeed the only members in the entire area. We shared experiences and bore testimony.

Time passed. As we prepared to reboard, Sister Louden looked so forlorn, so alone. She pleaded, “You can’t go yet; I have so missed the Church.” Suddenly, over the loudspeaker there was announced a 30-minute mechanical delay of our flight. Sister Louden whispered, “My prayer has been answered.” She then asked how she might influence her husband to show an interest in the gospel. We counseled her to include him in their home Primary lesson each week and be to him a living testimony of the gospel. I mentioned we would send to her a subscription to the Children’s Friend and additional helps for her family teaching. We urged that she never give up on her husband.

We departed Mt. Isa, a city to which I have never returned. I shall, however, always hold dear in memory that sweet mother and those precious children extending a tear-filled expression and a fond wave of gratitude and good-bye.

Several years later, while speaking at a priesthood leadership meeting in Brisbane, Australia, I emphasized the significance of gospel scholarship in the home and the importance of living the gospel and being examples of the truth. I shared with the men assembled the account of Sister Louden and the impact her faith and determination had had on me. As I concluded, I said, “I suppose I’ll never know if Sister Louden’s husband ever joined the Church, but he couldn’t have found a better model to follow than his wife.”

One of the leaders raised his hand, then stood and declared, “Brother Monson, I am Richard Louden. The woman of whom you speak is my wife. The children [his voice quavered] are our children. We are a forever family now, thanks in part to the persistence and the patience of my dear wife. She did it all.” Not a word was spoken. The silence was broken only by sniffles and marked by the sight of tears.

We do live in turbulent times. Often the future is unknown; therefore, it behooves us to prepare for uncertainties. Statistics reveal that at some time, for a variety of reasons, you may find yourself in the role of financial provider. I urge you to pursue your education and learn marketable skills so that, should such a situation arise, you are prepared to provide.

The role of women is unique. The renowned American essayist, novelist, and historian, Washington Irving, stated: “There is one in the world who feels for him who is sad a keener pang than he feels for himself; there is one to whom reflected joy is better than that which comes direct; there is one who rejoices in another’s honor more than in any which is one’s own; there is one on whom transcendent excellence sheds no beam but that of delight; there is one who hides another’s infirmities more faithfully than one’s own; there is one who loses all sense of self in the sentiment of kindness, tenderness, and devotion to another. That one is woman.”

Said President Gordon B. Hinckley: “God planted within women something divine that expresses itself in quiet strength, in refinement, in peace, in goodness, in virtue, in truth, in love. And all of these remarkable qualities find their truest and most satisfying expression in motherhood.”

Being a mother has never been an easy role. Some of the oldest writings in the world admonish us not to forsake the law of our mother, instruct us that a foolish son is the heaviness of his mother, and warn us not to ignore our mother when she is old.

The scriptures also remind us that what we learn from our mothers comprises our very core values, as with the 2,000 stripling sons and warriors of Helaman, who “had been taught by their mothers, that if they did not doubt, God would deliver them.” And He did!

Many members of Relief Society do not have husbands. Death, divorce, or lack of opportunity to marry have, in many instances, made it necessary for a woman to stand alone. Additionally, there are those who have just come from the Young Women program. In reality, no one need stand alone, for a loving Heavenly Father will be by her side to give direction to her life and provide peace and assurance in those quiet moments where loneliness is found and where compassion is needed. Also significant is the fact that the women of Relief Society stand side by side as sisters. May you ever be there to care for each other, to recognize one another’s needs. May you be sensitive to the circumstances of each, realizing that some women are facing particular challenges, but that every woman is a valued daughter of our Heavenly Father.

As I conclude my remarks, may I share with you an experience of several years ago which depicted the strength of you dear sisters in Relief Society.

During 1980, the sesquicentennial year of the organization of the Church, each member of the Relief Society general board was asked to write a personal letter to the sisters of the Church in the year 2030—50 years hence. The following is an excerpt from the letter written by Sister Helen Lee Goates:

“Our world of 1980 is filled with uncertainty, but I am determined to live each day with faith and not fear, to trust the Lord and to follow the counsel of our prophet today. I know that God lives, and I love Him with all my soul. I am so grateful that the gospel was restored to the earth 150 years ago and that I can enjoy the blessings of membership in this great Church. I am grateful for the priesthood of God, having felt its power throughout my life.

“I am at peace in my world and pray that you may be sustained in yours by firm testimonies and unwavering convictions of the gospel of Jesus Christ.”

Helen Lee Goates passed away in April of the year 2000. Shortly before her impending death from cancer, Sister Monson and I visited with her and her husband and family. She appeared calm and at peace. She told us she was prepared to go and looked forward to seeing once again her parents and other loved ones who had preceded her. In her life Sister Goates exemplified the nobility of Latter-day Saint women. In her passing she personified your theme: “If ye are prepared ye shall not fear.”

I bear to you, my beloved sisters, my witness that Heavenly Father lives, that Jesus is the Christ, and that we are led today by a prophet for our time—even President Gordon B. Hinckley. Safe journey to you as you travel along life’s pathway, I pray, in the name of Jesus Christ, amen.

\newpage
\settalk{Be Thou an Example}
\setspeaker{Thomas S. Monson}
\settalkdate{April 2005}
\talkheading{Be Thou an Example}{Thomas S. Monson}

My dear sisters, both those of you assembled in the magnificent Conference Center and those receiving the proceedings by satellite throughout the world, I pray for an interest in your prayers, that I may rise to the responsibility which is mine to address you.

We have been edified and inspired by the messages of the Young Women presidency, the beautiful music rendered, and the very spirit of this meeting. We have received a renewed appreciation for the Prophet Joseph Smith, for his life, and for the restored gospel of Jesus Christ.

The First Presidency of the Church loves you and has confidence in you and in your leaders. You are an example of righteousness in a world which desperately needs your influence and your strength.

Perhaps your battle cry might well be the charge given by the Apostle Paul to his beloved Timothy: “Be thou an example of the believers, in word, in conversation, in charity, in spirit, in faith, in purity.”

Today, permissiveness, immorality, pornography, and the power of peer pressure cause many to be tossed on a sea of sin and crushed on the jagged reefs of lost opportunities, forfeited blessings, and shattered dreams.

Precious young women, and you mothers, Young Women leaders, and advisers, may I leave with you a code of conduct to guide your footsteps safely through mortality and to the celestial kingdom of our Heavenly Father. I have divided my code of conduct into four parts:

First, you have a heritage; honor it. There come thundering to our ears the words from Mount Sinai: “Honour thy father and thy mother.”

My, how your parents love you, how they pray for you. Honor them.

How do you honor your parents? I like the words of William Shakespeare: “They do not love that do not show their love.” There are countless ways in which you can show true love to your mothers and your fathers. You can obey them and follow their teachings, for they will never lead you astray. You can treat them with respect. They have sacrificed much and continue to sacrifice in your behalf.

Be honest with your mother and your father. One reflection of such honesty with parents is to communicate with them. Avoid the silent treatment. The clock ticks more loudly, its hands move more slowly when the night is dark, the hour is late, and a precious daughter has not yet come home. If you are detained, make a telephone call: “Mom, Dad, we’re OK. Just stopped for something to eat. Don’t worry; we’re fine. Be home soon.”

A number of years ago, while attending a youth gathering at the Clarkston, Utah, cemetery, where each of the group viewed the memorial which marks the grave of Martin Harris, one of the Three Witnesses to the Book of Mormon, I noticed another marker—a small stone in which was inscribed a name and this poignant verse: “A light from our household is gone; a voice we loved is stilled. A place is vacant in our hearts that never can be filled.”

Don’t wait until that light from your household is gone; don’t wait until that voice you know is stilled before you say, “I love you, Mother; I love you, Father.” Now is the time to think and the time to thank. I trust you will do both. You have a heritage; honor it.

Next in our code of conduct: You will meet temptation; withstand it.

The Prophet Joseph Smith faced temptation. Can you imagine the ridicule, the scorn, the mocking that must have been heaped upon him as he declared that he had seen a vision? I suppose it became almost unbearable for the boy. He no doubt knew that it would be easier to retract his statements concerning the vision and just get on with a normal life. He did not, however, give in. These are his words: “I had actually seen a light, and in the midst of that light I saw two Personages, and they did in reality speak to me; and though I was hated and persecuted for saying that I had seen a vision, yet it was true. … I had seen a vision; I knew it, and I knew that God knew it, and I could not deny it.” Joseph Smith taught courage by example. He faced temptation and withstood it.

Many of you are familiar with the play Camelot. I’d like to share with you one of my favorite lines from this production. As the difficulties among King Arthur, Sir Lancelot, and Queen Guinevere deepen, King Arthur cautions, “We must not let our passions destroy our dreams.” This plea I would leave with you tonight. Do not let your passions destroy your dreams. Withstand temptation.

Remember the words from the Book of Mormon: “Wickedness never was happiness.”

Essential to your success and happiness is the advice “Choose your friends with caution.” We tend to become like those whom we admire, and they are usually our friends. We should associate with those who, like us, are planning not for temporary convenience, shallow goals, or narrow ambition—but rather with those who value the things that matter most, even eternal objectives.

Maintain an eternal perspective. Let there be a temple marriage in your future. There is no scene so sweet, no time so sacred as that very special day of your marriage. Then and there you glimpse celestial joy. Be alert; do not permit temptation to rob you of this blessing.

Make every decision you contemplate pass this test: What does it do to me? What does it do for me? And let your code of conduct emphasize not, “What will others think?” but rather, “What will I think of myself?” Be influenced by that still, small voice. Remember that one with authority placed his hands on your head at the time of your confirmation and said, “Receive the Holy Ghost.” Open your hearts, even your very souls, to the sound of that special voice which testifies of truth. As the prophet Isaiah promised, “Thine ears shall hear a word … saying, This is the way, walk ye in it.”

The tenor of our times is permissiveness. All around us we see the idols of the movie screen, the heroes of the athletic field—those whom many young people long to emulate—as disregarding the laws of God and rationalizing away sinful practices, seemingly with no ill effect. Don’t you believe it! There is a time of reckoning—even a balancing of the ledger. Every Cinderella has her midnight—it’s called Judgment Day, even the Big Exam of Life. Are you prepared? Are you pleased with your own performance?

Help can come to you from many sources. One is your patriarchal blessing. Such a blessing contains chapters from your book of eternal possibilities. Read your blessing frequently. Study it carefully. Be guided by its cautions. Live to merit its promises.

Now, if any has stumbled in her journey, there is a way back. The process is called repentance. Our Savior died to provide you and me that blessed gift. Though the path is difficult, the promise is real: “Though your sins be as scarlet, they shall be as white as snow.” “And I will remember [them] no more.” You will meet temptation; it is my prayer that you will withstand it.

Next in our code of conduct: You know the truth; live it.

After Joseph Smith’s vision in the Sacred Grove, he received no additional communication for three years. Can you imagine how you would feel if you had seen God the Father and Jesus Christ, His Son, if Christ had spoken to you, and then you had no additional word or communication for three years? Would you begin to doubt? Would you wonder or question why? The Prophet Joseph Smith did not wonder; he did not question; he did not doubt the Lord. He had received the truth, and he lived it.

My dear young friends, you have been reserved to come forth at this particular time when the gospel of Jesus Christ has been restored to the earth. Speaking of the gospel and of testimony, President Gordon B. Hinckley said: “[The] thing which we call testimony … is as real and powerful as any force on the earth. … It is found in young and old. … It brings with it the assurance that life is purposeful, that some things are of far greater importance than others, that we are on an eternal journey, that we are answerable unto God.”

You have been taught the truths of the gospel by your parents and by your teachers in the Church. You will continue to find truth in the scriptures, in the teachings of the prophets, and through the inspiration which comes to you as you bend your knees and seek the help of God.

Remember, faith and doubt cannot exist in the mind at the same time, for one will dispel the other. Cast out doubt. Cultivate faith. Strive always to retain that childlike faith which can move mountains and bring heaven closer to heart and home.

When firmly planted, your testimony of the gospel, of the Savior, and of our Heavenly Father will influence all that you do throughout your life. It will help to determine how you spend your time and with whom you choose to associate. It will affect the way you treat your family, how you interact with others. It will bring love, peace, and joy into your life. It should help you determine to be modest in your dress and in your speech. In the past year or so we have noticed a dramatic change in the way some of our young women are dressing. Styles in clothing change; fads come and go; but if the dress styles are immodest, it is important that our young women avoid them. When you dress modestly, you show respect for your Heavenly Father and for yourself. At this time, when dress fashions are styled after the skimpy clothing some of the current movie and music idols are wearing, it may be difficult to find modest apparel in clothing stores. However, it is possible, and it is important. The Apostle Paul declared: “Know ye not that ye are the temple of God, and that the Spirit of God dwelleth in you? … The temple of God is holy, which temple ye are.” You know the truth; live it.

Finally, you possess a testimony; share it. Never underestimate the far-reaching influence of your testimony. You can strengthen one another; you have the capacity to notice the unnoticed. When you have eyes to see, ears to hear, and hearts to feel, you can reach out and rescue others of your age.

To illustrate, may I share with you an experience which took place several years ago when Sister Monson had been hospitalized because of a fall. She asked me to go to the supermarket and purchase a few items. This was something I had not done before. I had a shopping list which included potatoes. I promptly found a grocery cart and placed a number of potatoes in it. I knew nothing of the plastic bags in which purchases are normally placed. As I moved the cart along, the potatoes fell out and onto the floor, exiting through two rather small openings in the back of the cart. A dutiful clerk hurried to my aid and called out, “Let me help you!” I tried to explain to her that my cart was defective. It was only then that I was told that all the carts had those two holes in the back and that they were meant for the legs of children.

Next the clerk took my list and helped me find each item. Then she said, “You are Bishop Monson, aren’t you?”

I answered that many years earlier I had indeed been a bishop. She continued: “At that time I lived on Gale Street in your ward and was not a member of the Church. You made certain the girls who were members contacted me each week and took me with them to Mutual and other activities. They were fine young women whose friendship and kindness touched my heart. I want to let you know that the fellowshipping you arranged for me led to my being baptized and confirmed a member of the Church. What a blessing this has been in my life,” she said, “and I thank you for your kindness.”

You can share your testimony in many ways—by the words you speak, by the example you set, by the manner in which you live your life.

May each of us emulate the Prophet Joseph’s great example. He taught the truth; he lived the truth; he shared the truth. You possess a testimony; share it.

My dear sisters, may God bless you. We love you; we pray for you. Remember that you do not walk alone. The Lord has promised you: “I will go before your face. I will be on your right hand and on your left, and my Spirit shall be in your hearts, and mine angels round about you, to bear you up.”

Tomorrow is Easter. On this Easter eve, may our thoughts turn to Him who atoned for our sins, who showed us the way to live, how to pray, and who demonstrated by His own actions how we might do so. Born in a stable, cradled in a manger, the Son of God beckons to each of us to follow Him. “Oh, sweet the joy this sentence gives: ‘I know that my Redeemer lives.’” May His Spirit be with you always, I pray, in His holy name—even Jesus Christ the Lord—amen.

\newpage
\settalk{Constant Truths for Changing Times}
\setspeaker{Thomas S. Monson}
\settalkdate{April 2005}
\talkheading{Constant Truths for Changing Times}{Thomas S. Monson}

My dear brothers and sisters, both within my view and assembled throughout the world, I seek an interest in your prayers and your faith as I respond to the assignment and privilege to speak to you.

I begin by expressing commendation to all of you. In this challenging world, the youth of the Church are the very best ever. The faith, the service, and the actions of our members are praiseworthy. We are a prayerful and faith-filled people, ever striving to be decent and honest. We take care of each other. We try to show love to our neighbors.

However, lest we become complacent, may I quote from 2 Nephi in the Book of Mormon:

“At that day shall [the devil] … lull them away into carnal security, that they will say: All is well in Zion; yea, Zion prospereth, all is well—and thus the devil cheateth their souls.”

Someone has said that our complacency tree has many branches, and each spring more buds come into bloom.

We cannot afford to be complacent. We live in perilous times; the signs are all around us. We are acutely aware of the negative influences in our society that stalk traditional families. At times television and movies portray worldly and immoral heroes and heroines and attempt to hold up as role models some actors and actresses whose lives are anything but exemplary. Why should we follow a blind guide? Radios blare forth much denigrating music with blatant lyrics, dangerous invitations, and descriptions of almost every type of evil imaginable.

We, as members of The Church of Jesus Christ of Latter-day Saints, must stand up to the dangers which surround us and our families. To aid us in this determination, I offer several suggestions, as well as some examples from my own life.

I begin with family home evening. We cannot afford to neglect this heaven-inspired program. It can bring spiritual growth to each member of the family, helping him or her to withstand the temptations which are everywhere. The lessons learned in the home are those that last the longest. As President Gordon B. Hinckley and his predecessors have stated, “The home is the basis of a righteous life, and no other instrumentality can take its place or fulfill its essential functions.”

Dr. Glenn J. Doman, noted author and medical authority, wrote: “The newborn child is almost an exact duplicate of an empty … computer, although superior to such a computer in almost every way. … What is placed in the child’s [mind] during the first eight years of life is probably there to stay.” “If you put misinformation into his [mind] during [this period], it is extremely difficult to erase it.” Dr. Doman added that the most receptive age in human life is that of two or three years.

I like this thought: “Your mind is a cupboard, and you stock the shelves.” Let us make certain that our cupboard shelves, and those of our family members, are stocked with the things which will provide safety to our souls and enable us to return to our Father in Heaven. Such shelves could well be stocked with gospel scholarship, faith, prayer, love, service, obedience, example, and kindness.

Next, I address the subject of debt. This is a day of borrowing, a day when multiple credit card offers arrive in our mailboxes each week. They generally offer a very low rate of interest which may apply for a short period of time; but what one usually doesn’t realize is that after that period has expired, the rates increase dramatically. I share with you a statement made by President J. Reuben Clark Jr., who many years ago was a member of the First Presidency. Its truth is timeless. Said he:

“It is a rule of our financial and economic life in all the world that interest is to be paid on borrowed money. …

“Interest never sleeps nor sickens nor dies; it never goes to the hospital; it works on Sundays and holidays; it never takes a vacation; it never visits nor travels; it takes no pleasure; it is never laid off work nor discharged from employment; it never works on reduced hours. … Once in debt, interest is your companion every minute of the day and night; you cannot shun it or slip away from it; you cannot dismiss it; it yields neither to entreaties, demands, or orders; and whenever you get in its way or cross its course or fail to meet its demands, it crushes you.”

My brothers and sisters, I’m appalled at some of the advertising I see and hear advocating home equity loans. Simply put, they are second mortgages on homes. The promotion for such loans is designed to tempt us to borrow more in order to have more. What is never mentioned is the fact that, should one be unable to make this “second” house payment, one is in danger of losing his house.

Avoid the philosophy and excuse that yesterday’s luxuries have become today’s necessities. They aren’t necessities unless we ourselves make them such. Many of our young couples today want to begin with multiple cars and the type of home Mother and Dad worked a lifetime to obtain. Consequently, they enter into long-term debt on the basis of two salaries. Perhaps too late they find that changes do come—women have children, sickness stalks some families, jobs are lost, natural disasters and other situations occur—and no longer can the mortgage payment, based on the income from two salaries, be made.

It is essential for us to live within our means.

Next, I have felt impressed to speak to mothers, to fathers, to sons, and to daughters.

I would say to each mother, each father—be a good listener. Communication is so vital today in our fast-paced world. Take time to listen. And to you children, talk to your mother and to your father. It may be difficult to realize, but your parents have lived through many of the same challenges which you face today. Often they see the big picture more clearly than you can. They pray for you each day and are entitled to the inspiration of our Heavenly Father in providing you counsel and advice.

Mothers, share household duties. It is often easier to do everything yourself than to persuade your children to help, but it is so essential for them to learn the importance of doing their share.

Fathers, I would counsel you to demonstrate love and kindness to your wife. Be patient with your children. Don’t indulge them to excess, for they must learn to make their own way in the world.

I would encourage you to be available to your children. I have heard it said that no man, as death approaches, has ever declared that he wished he had spent more time at the office.

I love the following example, taken from an article entitled “A Day at the Beach” by Arthur Gordon. Said he:

“When I was around thirteen and my brother ten, Father had promised to take us to the circus. But at lunchtime there was a phone call; some urgent business required his attention downtown. We braced ourselves for disappointment. Then we heard him say, ‘No, I won’t be down. It’ll have to wait.’

“When he came back to the table, Mother smiled [and said,] ‘The circus keeps coming back, you know.’

“‘I know,’ said Father. ‘But childhood doesn’t.’”

My brothers and sisters, time with your children is fleeting. Do not put off being with them now. Someone put it another way: Live only for tomorrow, and you will have a lot of empty yesterdays today.

Parents, help your children set goals concerning school and careers. Help your sons learn manners and respect for women and children.

Said President Hinckley: “As we train a new generation, so will the world be in a few years. If you are worried about the future, then look to the upbringing of your children.”

The Apostle Paul’s statement to his beloved Timothy could well apply: “Be thou an example of the believers, in word, in conversation, in charity, in spirit, in faith, in purity.”

Parents, live your lives in such a way that your children will find you an example worthy of emulation.

I admonish all families: search out your heritage. It is important to know, as far as possible, those who came before us. We discover something about ourselves when we learn about our ancestors.

I recall as a boy hearing of the experiences of my Miller ancestors. In the spring of 1848, my great-great-grandparents Charles Stewart Miller and Mary McGowan Miller, joined the Church in their native Scotland; left their home in Rutherglen, Scotland; and journeyed across the Atlantic Ocean. They reached the port of New Orleans and traveled up the Mississippi River to St. Louis, Missouri, with a group of Saints, arriving there in 1849. One of their 11 children, Margaret, would become my great-grandmother.

When the family arrived in St. Louis, planning to earn enough money to make their way to the Salt Lake Valley, a plague of cholera struck the area. The Miller family was hard-hit: in the space of two weeks, mother, father, and two of their sons died. My great-grandmother Margaret Miller was 13 years old at the time.

Because of all the deaths in the area, there were no caskets available—at any price. The older surviving boys dismantled the family’s oxen pens in order to make crude caskets for the family members who had passed away.

The nine remaining orphaned Miller children and the husband of one of the older daughters left St. Louis in the spring of 1850 with four oxen and one wagon, arriving finally in the Salt Lake Valley that same year.

I owe such a debt of gratitude to these and other noble forebears who loved the gospel and who loved the Lord so deeply that they were willing to sacrifice all they had, including their own lives, for The Church of Jesus Christ of Latter-day Saints. How grateful I am for the temple ordinances which bind us together for all eternity.

I emphasize how essential is the work we do in the temples of the Lord for our kindred dead.

Just two months ago today, members of my family gathered together in the Salt Lake Temple to perform sealings for some of our deceased ancestors. This was one of the most spiritual experiences our family has had together and enhanced the love we have for one another and the obligation which is ours to live worthy of our heritage.

Years ago when our youngest son, Clark, was attending a religion class at Brigham Young University, the instructor, during a lecture, asked him, “Clark, what is an example of life with your father that you best remember?”

The instructor later wrote to me and told me of the reply which Clark had given to the class. Said Clark: “When I was a deacon in the Aaronic Priesthood, my father and I went pheasant hunting near Malad, Idaho. The day was Monday—the last day of the pheasant hunting season. We walked through numerous fields in search of pheasants but saw only a few, and those we missed. Dad then said to me, ‘Clark,’ he looked at his watch, ‘let’s unload our guns, and we’ll place them in this ditch. Then we’ll kneel down to pray.’ I thought Dad would pray for more pheasants, but I was wrong. He explained to me that Elder Richard L. Evans of the Quorum of the Twelve was gravely ill and that at 12:00 noon on that particular Monday the members of the Quorum of the Twelve—wherever they may be—were to kneel and, in a way, together unite in a fervent prayer of faith for Elder Evans. Removing our caps, we knelt, we prayed.”

I well remember the occasion, but I never dreamed a son was watching, was learning, was building his own testimony.

Several years ago we had a young paperboy who didn’t always deliver the paper in the manner intended. Instead of getting the paper on the porch, he sometimes accidentally threw it into the bushes or even close to the street. Some on his paper route decided to start a petition of complaint. One day a delegation came to our home and asked my wife, Frances, to sign the petition. She declined, saying, “Why, he’s just a little boy, and the papers are so heavy for him. I would never be critical of him, for he tries his best.” The petition, however, was signed by many of the others on the paper route and sent to the boy’s supervisors.

Not many days afterward, I came home from work and found Frances in tears. When she was finally able to talk, she told me that she had just learned that the body of the little paperboy had been found in his garage, where he had taken his own life. Apparently the criticism heaped upon him had been too much for him to bear. How grateful we were that we had not joined in that criticism. What a vivid lesson this has always been regarding the importance of being nonjudgmental and treating everyone with kindness.

The Savior should be our example. As is recorded of Him, He “increased in wisdom and stature, and in favour with God and man.” He “went about doing good, … for God was with him.”

Remember that ofttimes the wisdom of God appears as foolishness to men, but the greatest single lesson we can learn in mortality is that when God speaks and a man obeys, that man will always be right.

May we ever follow the Prince of Peace, who literally showed the way for us to follow, for by doing so, we will survive these turbulent times. His divine plan can save us from the dangers which surround us on every side. His example points the way. When faced with temptation, He shunned it. When offered the world, He declined it. When asked for His life, He gave it.

Now is the time. This is the place. May we follow Him, I pray, in the name of Jesus Christ, amen.

\newpage
\settalk{The Sacred Call of Service}
\setspeaker{Thomas S. Monson}
\settalkdate{April 2005}
\talkheading{The Sacred Call of Service}{Thomas S. Monson}

I too wish to express my welcome to those who have been called to new assignments at this conference and my hearty congratulations to those who have received honorable releases from their service. The work moves forward. We love each of you.

My dear brethren, I am honored by the privilege to speak to you this evening. What a joy to see this magnificent Conference Center filled to capacity with those both young and old who hold the priesthood of God. To realize that similar throngs are assembled throughout the world brings to me an overwhelming sense of responsibility. I pray that the inspiration of the Lord will guide my thoughts and inspire my words.

President Joseph F. Smith made the following statement concerning the priesthood. Said he: “The Holy Priesthood is that authority which God has delegated to man, by which he may speak the will of God. … It is sacred, and it must be held sacred by the people. It should be honored and respected by them, in whomsoever it is held.”

The oath and covenant of the priesthood pertains to all of us. To those who hold the Melchizedek Priesthood, it is a declaration of our requirement to be faithful and obedient to the laws of God and to magnify the callings which come to us. To those who hold the Aaronic Priesthood, it is a pronouncement concerning future duty and responsibility, that they may prepare themselves here and now.

Said President Marion G. Romney, a former member of the First Presidency: “Every bearer of the Melchizedek Priesthood should give diligent and solemn heed to the implications of this oath and covenant which he has received. Failure to observe the obligations imposed by it is sure to bring disappointment, sorrow, and suffering.”

Added President Spencer W. Kimball: “One breaks [his] priesthood covenant by transgressing commandments—but also by leaving undone his duties. Accordingly, to break this covenant one needs only to do nothing.”

A famed minister observed: “Men will work hard for money. [Men] will work harder for other men. But men will work hardest of all when they are dedicated to a cause. … Duty is never worthily performed until it is performed by one who would gladly do more if only he could.”

The performance of one’s duty brings a sense of happiness and peace. Wrote the poet:

\begin{verse}
I slept and dreamt that life was joy.\\
I awoke and saw that life was duty.\\
I acted, and behold—\\
Duty was joy.
\end{verse}

The call of duty can come quietly as we who hold the priesthood respond to the assignments we receive. President George Albert Smith, that modest yet effective leader, declared, “It is your duty first of all to learn what the Lord wants and then by the power and strength of [your] holy Priesthood to [so] magnify your calling in the presence of your fellows … that the people will be glad to follow you.”

What does it mean to magnify a calling? It means to build it up in dignity and importance, to make it honorable and commendable in the eyes of all men, to enlarge and strengthen it, to let the light of heaven shine through it to the view of other men.

And how does one magnify a calling? Simply by performing the service that pertains to it. An elder magnifies the ordained calling of an elder by learning what his duties as an elder are and then by doing them. As with an elder, so with a deacon, a teacher, a priest, a bishop, and each who holds office in the priesthood.

Poet and author Robert Louis Stevenson reminded us, “I know what pleasure is, for I have done good work.”

Brethren, let us remember the counsel of King Benjamin: “When ye are in the service of your fellow beings ye are only in the service of your God.”

Let us reach out to rescue those who need our help and lift them to the higher road and the better way. Let us focus our thinking on the needs of priesthood holders and their wives and children who have slipped from the path of activity. May we listen to the unspoken message from their hearts. You will find it to be familiar:

\begin{verse}
Lead me, guide me, walk beside me,\\
Help me find the way.\\
Teach me all that I must do\\
To live with him someday.”
\end{verse}

The work of reactivation is no task for the idler or daydreamer. Children grow, parents age, and time waits for no man. Do not postpone a prompting; rather, act on it, and the Lord will open the way.

Frequently the heavenly virtue of patience is required. As a bishop I felt prompted one day to call on a man whose wife was somewhat active, as were the children. This man, however, had never responded. It was a hot summer’s day when I knocked on the screen door of Harold G. Gallacher. I could see Brother Gallacher sitting in his chair reading the newspaper. “Who is it?” he queried, without looking up.

“Your bishop,” I replied. “I’ve come to get acquainted and to urge your attendance with your family at our meetings.”

“No, I’m too busy,” came the disdainful response. He never looked up. I thanked him for listening and departed the doorstep.

The Gallacher family moved to California shortly thereafter. Many years went by. Then, as a member of the Quorum of the Twelve, I was working in my office one day when my secretary called, saying, “A Brother Gallacher who once lived in your ward is here in the office and would like to talk to you.”

I responded, “Ask him if his name is Harold G. Gallacher who, with his family, once lived at Vissing Place on West Temple and Fifth South.”

She said, “He is the man.”

I asked her to send him in. We had a pleasant conversation together concerning his family. He told me, “I’ve come to apologize for not getting out of my chair and letting you in the door that summer day long years ago.” I asked him if he was active in the Church. With a smile, he replied: “I’m a counselor in my ward bishopric. Your invitation to come out to church, and my negative response, so haunted me that I determined to do something about it.”

Harold and I visited together on numerous occasions before he passed away. The Gallachers and their children filled many callings in the Church.

President Stephen L Richards, who served as a counselor to President David O. McKay, declared, “The Priesthood is usually simply defined as ‘the power of God delegated to man.’” He continues: “This definition, I think, is accurate. But for practical purposes I like to define the Priesthood in terms of service, and I frequently call it ‘the perfect plan of service.’ … It is an instrument of service … and the man who fails to use it is apt to lose it, for we are plainly told by revelation that he who neglects it ‘shall not be counted worthy to stand.’”

This past January I had the privilege of witnessing a profound act of service in the life of a woman who had lived in my ward when I served as bishop many years ago. Her name is Adele, and she and her two grown daughters—one of whom is handicapped—have lived for many years in the Rose Park area of the Salt Lake Valley. Adele, who is a widow, has struggled financially, and her life has often been difficult.

I had received a telephone call from an individual involved with the Gingerbread House Project inviting me to the unveiling of Adele’s home, the renovation of which had been undertaken during a period of just over three days and nights by many kind and generous individuals, all working voluntarily with materials donated by numerous local businesses. During the time the makeover of her home had been accomplished, Adele and her two daughters had been hosted in a city a number of miles away, where they themselves had received some pampering.

I was present when the limousine bearing Adele and her daughters arrived on the scene. The group which had been waiting for them included not only family and friends but also many of the craftsmen who had worked night and day on the project. It was obvious they were pleased with the result and were anxious to see the reaction of Adele and her daughters.

The women stepped from the car, blindfolds in place. What a thrilling moment it was when the blindfolds were removed and Adele and her daughters turned around and saw their new home. They were absolutely stunned by the magnificent project which had been completed, including a redesign of the front, an extension of the home itself, and a new roof. The outside looked new and immaculate. They could not help but cry.

I accompanied Adele and others as we entered the home and were amazed at what had been accomplished to beautify and enhance the surroundings. The walls had been painted, the floor coverings changed. There were new furnishings, new curtains, new drapes. The cupboards in the kitchen had been replaced; there were new countertops and new appliances. The entire house had been done over from top to bottom, each room spotless and beautiful. Adele and her daughters were literally overcome. However, just as poignant and touching were the expressions on the faces of those who had worked feverishly to make the house new. Tears welled in their eyes as they witnessed the joy they had brought into the lives of Adele and her daughters. Not only had a widow’s burden been made lighter, but countless other lives were touched in the process. All were better people for having participated in this effort.

President Harold B. Lee, one of the great teachers in the Church, gave us this easy-to-understand counsel regarding the priesthood. Said he: “You see, when one becomes a holder of the priesthood, he becomes an agent of the Lord. He should think of his calling as though he were on the Lord’s errand.”

Now, some of you may be shy by nature, perhaps feeling yourselves inadequate to respond affirmatively to a calling. Remember that this work is not yours and mine alone. It is the Lord’s work, and when we are on the Lord’s errand, brethren, we are entitled to the Lord’s help. Remember that the Lord will shape the back to bear the burden placed upon it.

While the formal classroom may be intimidating at times, some of the most effective teaching takes place other than in the chapel or the classroom. Well do I remember that during the spring season some years ago, members of my ward and an adjoining ward took all the Aaronic Priesthood, who eagerly looked forward to an annual outing commemorating the restoration of the Aaronic Priesthood. On this particular occasion we journeyed by bus 90 miles north to the Clarkston, Utah, cemetery. There, in the quiet of that beautiful setting, we gathered the youth around the grave of Martin Harris, one of the Three Witnesses of the Book of Mormon. While we surrounded the beautiful granite shaft which marks his grave, Elder Glen L. Rudd, then the bishop of the other ward, presented the background of the life of Martin Harris and read from the Book of Mormon his testimony and that of Oliver Cowdery and David Whitmer. The young men listened with rapt attention, realizing they were standing at the grave site of one who had seen an angel and had actually beheld the plates with his own eyes. They reverently touched the granite marker designating the grave and pondered the words they had heard and the feelings they had felt.

Then we walked a short distance to a pioneer grave. The marker bore the name of John P. Malmberg and contained the verse:

\begin{verse}
A light from our household is gone.\\
A voice we loved is stilled.\\
A place is vacant in our hearts\\
That never can be filled.
\end{verse}

We talked with the boys about sacrifice, about dedication to truth. Duty, honor, service, and love—all were taught by that tombstone. In memory’s eye I can see the boys reach for their handkerchiefs to wipe away a tear. Heard yet are the sniffles which testified that hearts were touched and commitments made. I believe each youth had determined to be a pioneer—one who goes before, showing others the way to follow.

We then retired as a group to a local park, where all enjoyed a picnic lunch. Before turning homeward, we stopped at the grounds of the beautiful Logan temple. It was a warm day. I invited the boys to stretch out on the spacious lawn and with me gaze at a sky of blue, marked by white, billowy clouds hurried along on their journey by a steady breeze. We admired the beauty of this magnificent pioneer temple. We talked of sacred ordinances and eternal covenants. Lessons were learned. Hearts were touched. Covenants and promises became much more than words. The desire to be worthy to enter temple doors lodged in those youthful hearts. Thoughts turned to the Master; His presence was close. His gentle invitation “Follow me” was somehow heard and felt.

To all who willingly respond to the sacred call of service comes the promise: “I, the Lord, am merciful and gracious unto those who fear me, and delight to honor those who serve me in righteousness and in truth unto the end.

“Great shall be their reward and eternal shall be their glory.”

My sincere prayer is that all of us may qualify for this divine promise, in the name of Jesus Christ, our Savior, amen.

\newpage
\settalk{Do Your Duty—That Is Best}
\setspeaker{Thomas S. Monson}
\settalkdate{October 2005}
\talkheading{Do Your Duty—That Is Best}{Thomas S. Monson}

Brethren of the priesthood, assembled here in the Conference Center and worldwide, I am humbled by the responsibility which is mine to address a few remarks to you. I pray for the Spirit of the Lord to attend me as I do so.

I am aware that our audience this evening ranges from the most recently ordained deacon to the eldest high priest. To each, the restoration of the Aaronic Priesthood to Joseph Smith and Oliver Cowdery by John the Baptist and the Melchizedek Priesthood to Joseph and Oliver by Peter, James, and John are sacred and treasured events.

To you deacons, may I say that I recall the time when I was ordained a deacon. Our bishopric stressed the sacred responsibility which was ours to pass the sacrament. Emphasized were proper dress, a dignified bearing, and the importance of being clean inside and out. As we were taught the procedure in passing the sacrament, we were told how we should assist Louis McDonald, a particular brother in our ward who was afflicted with a palsied condition, that he might have the opportunity to partake of the sacred emblems.

How I remember being assigned to pass the sacrament to the row where Brother McDonald sat. I was fearful and hesitant as I approached this wonderful brother, and then I saw his smile and the eager expression of gratitude that showed his desire to partake. Holding the tray in my left hand, I took a small piece of bread and pressed it to his lips. The water was later served in the same way. I felt I was on holy ground. And indeed I was. The privilege to pass the sacrament to Brother McDonald made better deacons of us all.

Just two months ago, on Sunday, July 31, I was at Fort A. P. Hill, Virginia, attending an LDS sacrament meeting held during the National Scout Jamboree. My purpose in being there was to speak to the 5,000 Latter-day Saint young men and their leaders who had spent the previous week participating in the activities of the jamboree. They sat reverently in a natural amphitheater as an impressive 400-voice Aaronic Priesthood chorus sang:

\begin{verse}
A Mormon boy, a Mormon boy,\\
I am a Mormon boy.\\
I might be envied by a king,\\
For I am a Mormon boy.
\end{verse}

The sacrament was blessed, with 65 priests officiating at the many large sacrament tables which had been placed throughout the assembled group. Approximately 180 deacons then passed the sacrament. Within the time it would take to handle the passing of the sacrament in a crowded ward chapel, this large gathering was served. What an awe-inspiring sight I witnessed that morning as these Aaronic Priesthood young men participated in this holy ordinance.

It is important for each deacon to be guided to a spiritual awareness of the sacredness of his ordained calling. In one ward the lesson was effectively taught pertaining to the collection of fast offerings.

On fast day the ward members were visited by deacons and teachers so that each family could make a contribution. The deacons were a bit disgruntled, having to arise earlier than usual to fulfill this assignment.

The inspiration came for the bishopric to take a busload of the deacons and teachers to Welfare Square in Salt Lake City. Here they saw needy children receiving new shoes and other items of clothing. Here they witnessed empty baskets being filled with groceries. There was no money exchanged. One brief comment was made: “Young men, this is what the money you collect on fast day provides—even food, clothing, and shelter for those who are in need.” The Aaronic Priesthood young men smiled more, stepped higher, and served more willingly in filling their assignments.

Now, pertaining to the teachers and priests, every one of you should be given the assignment to home teach with a companion who holds the Melchizedek Priesthood. What an opportunity to prepare for a mission. What a privilege to learn the discipline of duty. A young man will automatically turn from concern for self when he is assigned to “watch over” others.

President David O. McKay counseled: “Home teaching is one of our most urgent and most rewarding opportunities to nurture and inspire, to counsel and direct our Father’s children. … [It] is a divine service, a divine call. It is our duty as Home Teachers to carry the divine spirit into every home and heart.”

Home teaching answers many prayers and permits us to see the occurrence of living miracles.

As I think of home teaching, I am reminded of a man by the name of Johann Denndorfer from Debrecen, Hungary. He had been converted to the Church years before in Germany, and now, following World War II, he found himself virtually a prisoner in his own land of Hungary. How he longed for contact with the Church. Then his home teachers visited. Brother Walter Krause and his companion went from the northeastern portion of Germany all the way to Hungary to fulfill their home teaching assignment. Before they left from their homes in Germany, Brother Krause had said to his companion, “Would you like to go home teaching with me this week?”

His companion asked, “When will we leave?”

Brother Krause’s response: “Tomorrow.”

Then came the question, “When will we come back?”

Brother Krause did not hesitate; he said, “Oh, in about a week.”

And away they went to visit Brother Denndorfer and others. Brother Denndorfer had not had home teachers since before the war. Now, when he saw the servants of the Lord, he was overwhelmed. He did not shake hands with them; rather, he went to his bedroom and took from a secret hiding place his tithing that he had saved for years. This tithing he gave to his home teachers, and then he said, “Now I can shake your hands.”

Now a word for the priests in the Aaronic Priesthood. You young men have the opportunity to bless the sacrament, to continue your home teaching duties, and to participate in the sacred ordinance of baptism.

Fifty-five years ago I knew a young man, Robert Williams, who held the office of priest in the Aaronic Priesthood. As the bishop, I was his quorum president. When he spoke, Robert stuttered and stammered, void of control. He was self-conscious, shy, fearful of himself and everybody else; this impediment was devastating to him. Rarely did he accept an assignment; never would he look another person in the eye; always would he gaze downward. Then one day, through a set of unusual circumstances, he accepted an assignment to perform the responsibility to baptize another.

I sat next to Robert in the baptistry of the Salt Lake Tabernacle. I knew he needed all the help he could get. He was dressed in immaculate white, prepared for the ordinance he was to perform. I asked him how he felt. He gazed at the floor and stuttered almost uncontrollably that he felt terrible.

We both prayed fervently that he would be made equal to his task. The clerk then said, “Nancy Ann McArthur will now be baptized by Robert Williams, a priest.”

Robert left my side, stepped into the font, took little Nancy by the hand, and helped her into that water which cleanses human lives and provides a spiritual rebirth. He spoke the words, “Nancy Ann McArthur, having been commissioned of Jesus Christ, I baptize you in the name of the Father, and of the Son, and of the Holy Ghost. Amen.”

And he baptized her. Not once did he stutter! Not once did he falter! A modern miracle had been witnessed. Robert then performed the baptismal ordinance for two or three other children in the same fashion.

In the dressing room, I hurried to congratulate Robert. I expected to hear this same uninterrupted flow of speech. I was wrong. He gazed downward and stammered his reply of gratitude.

I testify to you that when Robert acted in the authority of the Aaronic Priesthood, he spoke with power, with conviction, and with heavenly help.

Just over two years ago it was my privilege to speak at the funeral services for Robert Williams and to pay tribute to this faithful priesthood holder who tried his best throughout his life to honor his priesthood.

Some of you young men here tonight may be shy by nature or might consider yourselves inadequate to respond to a calling. Remember that this work is not yours and mine alone. We can look up and reach out for divine help.

Like some of you, I know what it is to face disappointment and youthful humiliation. As a boy, I played team softball in elementary and junior high school. Two captains were chosen, and then they, in turn, selected the players they desired on their respective teams. Of course, the best players were chosen first, then second, and third. To be selected fourth or fifth was not too bad, but to be chosen last and relegated to a remote position in the outfield was downright awful. I know; I was there.

How I hoped the ball would never be hit in my direction, for surely I would drop it, runners would score, and teammates would laugh.

As though it were just yesterday, I remember the very moment when all that changed in my life. The game started out as I have described: I was chosen last. I made my sorrowful way to the deep pocket of right field and watched as the other team filled the bases with runners. Two batters then went down on strikes. Suddenly, the next batter hit a mighty drive. I even heard him say, “This will be a home run.” That was humiliating, since the ball was coming in my direction. Was it beyond my reach? I raced for the spot where I thought the ball would drop, uttered a prayer while running, and stretched forth my cupped hands. I surprised myself. I caught the ball! My team won the game.

This one experience bolstered my confidence, inspired my desire to practice, and led me from that last-to-be-chosen place to become a real contributor to the team.

We can experience that burst of confidence. We can feel that pride of performance. A three-word formula will help us: Never give up.

From the play Shenandoah comes the spoken line which inspires: “If we don’t try, then we don’t do; and if we don’t do, then why are we here?”

Miracles are everywhere to be found when priesthood callings are magnified. When faith replaces doubt, when selfless service eliminates selfish striving, the power of God brings to pass His purposes. The priesthood is not really so much a gift as it is a commission to serve, a privilege to lift, and an opportunity to bless the lives of others.

The call of duty can come quietly as we who hold the priesthood respond to the assignments we receive. President George Albert Smith, that modest yet effective leader, declared, “It is your duty first of all to learn what the Lord wants and then by the power and strength of His holy Priesthood to magnify your calling in the presence of your fellows in such a way that the people will be glad to follow you.”

And how does one magnify a calling? Simply by performing the service that pertains to it. An elder magnifies the ordained calling of an elder by learning what his duties as an elder are and then by doing them. As with an elder, so with a deacon, a teacher, a priest, a bishop, and each who holds office in the priesthood.

Brethren, it is in doing, not just dreaming, that lives are blessed, others are guided, and souls are saved. “Be ye doers of the word, and not hearers only, deceiving your own selves,” counseled James.

May all within the sound of my voice make a renewed effort to qualify for the Lord’s guidance in our lives. There are many who plead and pray for help. There are those who are discouraged and in need of a helping hand.

Many years ago when I served as a bishop, I presided over a large ward with over 1,000 members, including 87 widows. On one occasion I was visiting, along with one of my counselors, a widow and her mature handicapped daughter. As we left their apartment, a lady from the apartment across the hall was standing outside her door and stopped us. She spoke with a foreign accent and asked if I were a bishop; I replied that I was. She told me that she noticed I often visited with others. Then she said, “No one visits me or my bedfast husband. Do you have time to come in and visit with us, even though we are not members of your church?”

As we entered her apartment, we noticed that she and her husband were listening to the Tabernacle Choir on the radio. We talked with the couple for a while, then provided a blessing to the husband.

Following that initial visit I stopped by as often as I could. The couple eventually met with the missionaries, and the wife, Angela Anastor, was baptized. Sometime later her husband passed away, and I had the privilege of conducting and speaking at his funeral services. Sister Anastor, with her knowledge of the Greek language, later was to translate the widely used pamphlet Joseph Smith Tells His Own Story into the Greek language.

I love the motto “Do [your] duty; that is best; Leave unto [the] Lord the rest!”

Active service in the Aaronic Priesthood will prepare you young men to receive the Melchizedek Priesthood, to serve missions, and to marry in the holy temple.

You will ever remember your Aaronic Priesthood quorum advisers and your fellow quorum members, thereby experiencing the truth “God gave us memories, that we might have June roses in the December of our lives.”

Young men of the Aaronic Priesthood, your future beckons; prepare for it. May Heavenly Father ever guide you as you do so. May He guide all of us as we strive to honor the priesthood which we hold and to magnify our callings, I pray humbly, in the name of Jesus Christ, amen.

\newpage
\settalk{The Prophet Joseph Smith: Teacher by Example}
\setspeaker{Thomas S. Monson}
\settalkdate{October 2005}
\talkheading{The Prophet Joseph Smith: Teacher by Example}{Thomas S. Monson}

My brothers and sisters, in this bicentennial year of his birth, I should like to speak of our beloved Prophet Joseph Smith.

On December 23, 1805, Joseph Smith Jr. was born in Sharon, Vermont, to Joseph Smith Sr. and Lucy Mack Smith. On the day of his birth, as the proud parents looked down upon this tiny baby, they could not have known what a profound impact he would have upon the world. A choice spirit had come to dwell in its earthly tabernacle; he has affected our lives and has taught us—through his own example—essential lessons. Today I should like to share a few of those lessons with you.

When Joseph was about six or seven years old, he and his brothers and sisters were stricken with typhus fever. Although the others recovered readily, Joseph was left with a painful sore on his leg. The doctors, using the best medicine they had, treated him, and yet the sore persisted. In order to save Joseph’s life, they said, he would have to lose his leg. Thankfully, however, soon after that diagnosis, the doctors returned to the Smith home and reported that there was a new procedure which might save Joseph’s leg. They wanted to operate immediately and had brought some cord with which to tie little Joseph to the bed so that he wouldn’t thrash about, since they had nothing with which to dull the pain. Young Joseph, however, told them, “You won’t need to tie me.”

The doctors suggested he take some brandy or wine so that the pain might not be so severe. “No,” young Joseph replied. “If my father will sit on the bed and hold me in his arms, I will do whatever is necessary.” Joseph Smith Sr. held in his arms his small child, and the doctors removed the diseased piece of bone. Although young Joseph was lame for some time afterward, he was healed. At such a young age and countless other times throughout his life, Joseph Smith taught us courage—by example.

Before Joseph’s 15th year, his family moved to Manchester, New York. He later described the great religious revival which seemed everywhere present at this time and of prime concern to nearly everyone. Joseph himself longed to know which church he should join. He writes in his history:

“I often said to myself: … Who of all these parties are right; or, are they all wrong together? If any one of them be right, which is it, and how shall I know it?

“While I was laboring under the extreme difficulties caused by … these parties of religionists, I was one day reading the Epistle of James, first chapter and fifth verse … : If any of you lack wisdom, let him ask of God, that giveth to all men liberally, and upbraideth not; and it shall be given him.”

Joseph reported that he knew he must either put the Lord to the test and ask Him or perhaps choose to remain in darkness forever. Early one morning he stepped into a grove, now called sacred, and knelt and prayed, having faith that God would give him the enlightenment which he so earnestly sought. Two personages appeared to Joseph—the Father and the Son—and he was told, in answer to his question, that he was to join none of the churches, for none of them was true. The Prophet Joseph Smith taught us the principle of faith—by example. His simple prayer of faith on that spring morning in 1820 brought about this marvelous work which continues today throughout the world.

A few days after his prayer in the Sacred Grove, Joseph Smith gave an account of his vision to a preacher with whom he was acquainted. To his surprise, his communication was treated with “contempt” and “was the cause of great persecution, which continued to increase.” Joseph, however, did not waver. He later wrote, “I had actually seen a light, and in the midst of that light I saw two Personages, and they did in reality speak to me; and though I was hated and persecuted for saying that I had seen a vision, yet it was true. … For I had seen a vision; I knew it, and I knew that God knew it, and I could not deny it.” Despite the physical and mental punishment at the hands of his opponents which the Prophet Joseph Smith endured throughout the remainder of his life, he did not falter. He taught honesty—by example.

After that great First Vision, the Prophet Joseph received no additional communication for three years. However, he did not wonder; he did not question; he did not doubt the Lord. He waited patiently. He taught us the heavenly virtue of patience—by example.

Following the visits of the angel Moroni to young Joseph and his acquisition of the plates, Joseph commenced the difficult assignment of translation. One can but imagine the dedication, the devotion, and the labor required to translate in fewer than 90 days this record of over 500 pages covering a period of 2,600 years. I love the words Oliver Cowdery used to describe the time he spent assisting Joseph with the translation of the Book of Mormon: “These were days never to be forgotten—to sit under the sound of a voice dictated by the inspiration of heaven, awakened the utmost gratitude of this bosom!” The Prophet Joseph Smith taught us diligence—by example.

As we know, the Prophet Joseph sent forth missionaries to preach the restored gospel. He himself served a mission in Upper New York and in Canada with Sidney Rigdon. He not only inspired others to volunteer for missions, but he also taught the importance of missionary work—by example.

I think one of the sweetest lessons taught by the Prophet Joseph, and yet one of the saddest, occurred close to the time of his death. He had seen in vision the Saints leaving Nauvoo and going to the Rocky Mountains. He was anxious that his people be led away from their tormentors and into this promised land which the Lord had shown him. He no doubt longed to be with them. However, he had been issued an arrest warrant on trumped up charges. Despite many appeals to Governor Ford, the charges were not dismissed. Joseph left his home, his wife, his family, and his people and gave himself up to the civil authorities, knowing he would probably never return.

These are the words he spoke as he journeyed to Carthage: “I am going like a lamb to the slaughter; but I am calm as a summer’s morning; I have a conscience void of offense towards God, and towards all men.”

In Carthage Jail he was incarcerated with his brother Hyrum and others. On June 27, 1844, Joseph, Hyrum, John Taylor, and Willard Richards were together there when an angry mob stormed the jail, ran up the stairway, and began firing through the door of the room they occupied. Hyrum was killed, and John Taylor was wounded. Joseph Smith’s last great act here upon the earth was one of selflessness. He crossed the room, most likely “thinking that it would save the lives of his brethren in the room if he could get out, … and sprang into the window when two balls pierced him from the door, and one entered his right breast from without.” He gave his life; Willard Richards and John Taylor were spared. “Greater love hath no man than this, that a man lay down his life for his friends.” The Prophet Joseph Smith taught us love—by example.

In retrospect, over 160 years later, although the events of June 27, 1844, were tragic, we are provided comfort as we realize that Joseph Smith’s martyrdom was not the last chapter in this account. Although those who sought to take his life felt that the Church would collapse without him, his powerful testimony of truth, the teachings he translated, and his declaration of the Savior’s message go on today in the hearts of over 12 million members throughout the world, who proclaim him a prophet of God.

The testimony of the Prophet Joseph continues to change lives. Some years ago I served as the president of the Canadian Mission. In Ontario, Canada, two of our missionaries were proselyting door-to-door on a cold, snowy afternoon. They had not had any measure of success. One elder was experienced; one was new.

The two called at the home of Mr. Elmer Pollard, and he, feeling sympathy for the almost frozen missionaries, invited them in. They presented their message and asked if he would join in prayer. He agreed, on the provision that he could offer the prayer.

The prayer he offered astonished the missionaries. He said, “Heavenly Father, bless these two unfortunate, misguided missionaries that they may return to their homes and not waste their time telling the people of Canada about a message which is so fantastic and about which they know so little.”

As they arose from their knees, Mr. Pollard asked the missionaries never to return to his home. As they left, he said mockingly to them, “You can’t tell me you really believe that Joseph Smith was a prophet of God anyway!” and he slammed the door.

The missionaries had walked but a short distance when the junior companion said timidly, “Elder, we didn’t answer Mr. Pollard.”

The senior companion responded: “We’ve been rejected. Let’s move on.”

The young missionary persisted, however, and the two returned to Mr. Pollard’s door. Mr. Pollard answered the knock and angrily said, “I thought I told you young men never to return!”

The junior companion then said, with all the courage he could muster, “Mr. Pollard, when we left your door, you said that we didn’t really believe Joseph Smith was a prophet of God. I want to testify to you, Mr. Pollard, that I know Joseph Smith was a prophet of God, that by inspiration he translated the sacred record known as the Book of Mormon, that he did see God the Father and Jesus the Son.” The missionaries then departed the doorstep.

I heard this same Mr. Pollard in a testimony meeting state the experiences of that memorable day. He said, “That evening, sleep would not come. I tossed and turned. Over and over in my mind I heard the words, ‘Joseph Smith was a prophet of God. I know it. … I know it. … I know it.’ I could scarcely wait for morning to come. I telephoned the missionaries, using their number which was printed on the small card containing the Articles of Faith. They returned, and this time my wife, my family, and I joined in the discussion as earnest seekers of truth. As a result, we have all embraced the gospel of Jesus Christ. We shall ever be grateful to the testimony of truth brought to us by those two courageous, humble missionaries.”

In the 135th section of the Doctrine and Covenants we read the words of John Taylor concerning the Prophet Joseph: “Joseph Smith, the Prophet and Seer of the Lord, has done more, save Jesus only, for the salvation of men in this world, than any other man that ever lived in it.”

I love the words of President Brigham Young, who said, “I feel like shouting Hallelujah, all the time, when I think that I ever knew Joseph Smith, the Prophet whom the Lord raised up and ordained, and to whom he gave keys and power to build up the Kingdom of God on earth.”

To this fitting tribute to our beloved Joseph, I add my own testimony that I know he was God’s prophet, chosen to restore the gospel of Jesus Christ in these latter days. I pray that as we celebrate the 200th anniversary of his birth, we may learn from his life. May we incorporate into our own lives the divine principles which he so beautifully taught—by example—that we, ourselves, might live more completely the gospel of Jesus Christ. May our lives reflect the knowledge we have that God lives, that Jesus Christ is His Son, that Joseph Smith was a prophet, and that we are led today by another prophet of God—even President Gordon B. Hinckley.

This conference marks 42 years since I was called to the Quorum of the Twelve Apostles. In my first meeting with the First Presidency and Quorum of the Twelve in the temple, the hymn which we sang, honoring Joseph Smith the Prophet, was and is a favorite of mine. I close with a verse from that hymn:

\begin{verse}
Praise to the man who communed with Jehovah!\\
Jesus anointed that Prophet and Seer.\\
Blessed to open the last dispensation,\\
Kings shall extol him, and nations revere.
\end{verse}

I testify of this solemn truth, in the name of Jesus Christ, amen.

\newpage
\settalk{Our Sacred Priesthood Trust}
\setspeaker{Thomas S. Monson}
\settalkdate{April 2006}
\talkheading{Our Sacred Priesthood Trust}{Thomas S. Monson}

Some years ago as our youngest son, Clark, was approaching his 12th birthday, he and I were leaving the Church Administration Building when President Harold B. Lee approached and greeted us. I mentioned to President Lee that Clark would soon be 12, whereupon President Lee turned to him and asked, “What happens to you when you turn 12?”

This was one of those times when a father prays that a son will be inspired to give a proper response. Clark, without hesitation, said to President Lee, “I will be ordained a deacon!”

The answer was the one for which I had prayed and which President Lee had sought. He then counseled our son, “Remember, it is a great blessing to hold the priesthood.”

I hope with all my heart and soul that every young man who receives the priesthood will honor that priesthood and be true to the trust which is conveyed when it is conferred. May each of us who holds the priesthood of God know what he believes. As the Apostle Peter admonished, may we ever be ready “to give an answer to every man that asketh you a reason of the hope that is in you.” There will be occasions in each of our lives when we will be called upon to explain or to defend our beliefs. When the time for performance arrives, the time for preparation is past.

Most of you young men will have the opportunity to share your testimonies when you serve as missionaries throughout the world. Prepare now for that wonderful privilege.

I have experienced many opportunities. One occurred 21 years ago, prior to the time when the German Democratic Republic—or East Germany, as it was more commonly known—was freed from Communist rule. I was visiting with the East German state secretary, Minister Gysi. At that time our temple at Freiberg, in East Germany, was under construction, along with two or three meetinghouses. Minister Gysi and I visited on a number of subjects, including our worldwide building program. He then asked, “Why is your church so wealthy that you can afford to build buildings in our country and throughout the world? How do you get your money?”

I answered that the Church is not wealthy but that we follow the ancient biblical principle of tithing, which principle is reemphasized in our modern scripture. I explained also that our Church has no paid ministry and indicated that these were two reasons why we were able to build the buildings then under way, including the beautiful temple at Freiberg.

Minister Gysi was most impressed with the information I presented, and I was very grateful I was able to answer his questions.

The opportunity to declare a truth may come when we least expect it. Let us be prepared.

On one occasion, President David O. McKay was asked by a woman not a member of the Church what specific belief set apart the teachings of the Church from those of any other faith. In speaking of this later, President McKay indicated that he had felt impressed to answer, “That which differentiates the beliefs of my church from those of others is divine authority by direct revelation.”

Where could we find a more significant example of divine authority by direct revelation than in the events which occurred that “beautiful, clear day, early in the spring of eighteen hundred and twenty,” when the lad Joseph Smith retired to the woods to pray. His words describing that moment in history are overpowering: “I saw two Personages, whose brightness and glory defy all description, standing above me in the air. One of them spake unto me, calling me by name and said, pointing to the other—This is My Beloved Son. Hear Him!”

Our thoughts turn to the visit of that heavenly messenger John the Baptist on May 15, 1829. There on the bank of the Susquehanna River, near Harmony, Pennsylvania, John laid his hands upon Joseph Smith and Oliver Cowdery and ordained them, saying, “Upon you my fellow servants, in the name of Messiah I confer the Priesthood of Aaron, which holds the keys of the ministering of angels, and of the gospel of repentance, and of baptism by immersion for the remission of sins.” The messenger announced that he acted under the direction of Peter, James, and John, who held the keys of the Melchizedek Priesthood. Ordination and baptism followed. This is yet another example of divine authority by direct revelation.

In due time, Peter, James, and John were sent to bestow the blessings of the Melchizedek Priesthood. These Apostles sent by the Lord ordained and confirmed Joseph and Oliver to be Apostles and special witnesses of His name. Divine authority by direct revelation characterized this sacred visitation.

As a result of these experiences, all of us carry the requirement—even the blessed opportunity and solemn duty—to be true to the trust we have received.

President Brigham Young declared, “The Priesthood of the Son of God is … the law by which the worlds are, were, and will continue for ever and ever.” President Joseph F. Smith, expanding on this theme, advised, “It is nothing more nor less than the power of God delegated to man by which man can act in the earth for the salvation of the human family, in the name of the Father and the Son and the Holy Ghost, and act legitimately; not assuming that authority, nor borrowing it from generations that are dead and gone, but authority that has been given in this day in which we live by ministering angels and spirits from above, direct from the presence of Almighty God.”

As I approached my 18th birthday and prepared to enter military service in World War II, I was recommended to receive the Melchizedek Priesthood. Mine was the task to telephone President Paul C. Child, my stake president, for an interview. He was one who loved and understood the holy scriptures, and it was his intent that all others should similarly love and understand them. As I knew from others of his rather detailed and searching interviews, our telephone conversation went something like this:

“Hello, President Child. This is Brother Monson. I have been asked by the bishop to visit with you relative to being ordained an elder.”

“Fine, Brother Monson. When can you see me?”

Knowing that his sacrament meeting time was 4:00 and desiring minimum exposure of my scriptural knowledge to his review, I suggested, “How would 3:00 be?”

His response: “Oh, Brother Monson, that would not provide us sufficient time to peruse the scriptures. Could you please come at 2:00 and bring with you your personally marked set of scriptures?”

Sunday finally arrived, and I visited President Child’s home. I was greeted warmly, and then the interview began. He said, “Brother Monson, you hold the Aaronic Priesthood.” Of course, I knew that. He continued, “Have you ever had an angel minister to you?”

My reply, “I’m not sure.”

“Do you know,” said he, “that you are entitled to such?”

Came my response: “No.”

Then he instructed, “Brother Monson, repeat from memory the 13th section of the Doctrine and Covenants.”

I began, “Upon you my fellow servants, in the name of Messiah I confer the Priesthood of Aaron, which holds the keys of the ministering of angels …”

“Stop,” President Child directed. Then in a calm, kindly tone, he counseled, “Brother Monson, never forget that as a holder of the Aaronic Priesthood you are entitled to the ministering of angels. Now continue the passage.”

I recited from memory the remainder of the section. President Child said, “Splendid.” He then discussed with me several other sections of the Doctrine and Covenants pertaining to the priesthood. It was a long interview, but I have never forgotten it. At the conclusion, President Child put his arm around my shoulder and said, “You are now ready to receive the Melchizedek Priesthood. Remember that the Lord blesses the person who serves Him.”

Many years later, Paul C. Child, then of the Priesthood Welfare Committee, and I attended a stake conference together. At the priesthood leadership session, when it was his turn to speak, he took his scriptures in hand and walked from the stand into the congregation. Knowing President Child as I did, I knew what he was going to do. He quoted from the Doctrine and Covenants, including section 18 concerning the worth of a soul, indicating that we should labor all our days to bring souls unto the Lord. He then turned to one elders quorum president and asked, “What is the worth of a soul?”

The stunned quorum president hesitated as he formulated his reply. I had a prayer in my heart that he would be able to answer the question. He finally responded, “The worth of a soul is its capacity to become as God.”

Brother Child closed his scriptures, walked solemnly and quietly up the aisle and back to the stand. As he passed by me, he said, “A most profound reply.”

We need to know the oath and covenant of the priesthood because it pertains to all of us. To those who hold the Melchizedek Priesthood, it is a declaration of our requirement to be faithful and obedient to the laws of God and to magnify the callings which come to us. To those who hold the Aaronic Priesthood, it is a pronouncement concerning future duty and responsibility, that they may prepare themselves here and now.

This oath and covenant is set forth by the Lord in these words:

“For whoso is faithful unto the obtaining these two priesthoods of which I have spoken, and the magnifying their calling, are sanctified by the Spirit unto the renewing of their bodies.

“They become the sons of Moses and of Aaron and the seed of Abraham, and the church and kingdom, and the elect of God.

“And also all they who receive this priesthood receive me, saith the Lord;

“For he that receiveth my servants receiveth me;

“And he that receiveth me receiveth my Father;

“And he that receiveth my Father receiveth my Father’s kingdom; therefore all that my Father hath shall be given unto him.”

The late Elder Delbert L. Stapley of the Quorum of the Twelve once observed: “There are two main requirements of this oath and covenant. First is faithfulness, which denotes obedience to the laws of God and connotes true observance of all gospel standards. …

“The second requirement … is to magnify one’s calling. To magnify is to honor, to exalt and glorify, and cause to be held in greater esteem or respect. It also means to increase the importance of, to enlarge and make greater.”

The Prophet Joseph Smith was once asked, “Brother Joseph, you frequently urge that we magnify our callings. What does this mean?” He is said to have replied, “To magnify a calling is to hold it up in dignity and importance, that the light of heaven may shine through one’s performance to the gaze of other men. An elder magnifies his calling when he learns what his duties as an elder are and then performs them.”

Those who bear the Aaronic Priesthood should be given opportunity to magnify their callings in that priesthood.

One Sunday two years ago I was attending sacrament meeting in my ward. That’s a rarity. There were three priests at the sacrament table, with the young man in the center being somewhat handicapped in movement but particularly so in speech. He tried twice to bless the bread but stumbled badly each time, no doubt embarrassed by his inability to give the prayer perfectly. One of the other priests then took over and gave the blessing on the bread.

During the passing of the bread, I thought to myself, “I just can’t let that young man experience failure at the sacrament table.” I had a strong feeling that if I didn’t doubt, he would be able to bless the water effectively. Inasmuch as I was on the stand near the sacrament table, I leaned over and said to the priest closest to me, pointing to the young man who had experienced the difficulty, “Let him bless the water; it’s a shorter prayer.” And then I prayed. I didn’t want a double failure. I love that passage of scripture which tells us that we should not doubt but believe.

When it was time to bless the water, that young man knelt again and gave the prayer, perhaps somewhat haltingly but without missing a word. I rejoiced silently. While the deacons were passing the trays, I looked over at the boy and gave him a thumbs-up. He gave me a broad smile. When the young men were excused to sit with their families, he sat on the row between his mother and father. What a joy it was to see his mother give him a big smile and a warm hug, while his father congratulated him and put his arm around his shoulder. All three of them looked in my direction, and I gave them all a thumbs-up. I could see the mother and father wiping tears from their eyes. I felt impressed that this young man would do just fine in the future.

The priesthood is not really so much a gift as it is a commission to serve, a privilege to lift, and an opportunity to bless the lives of others.

Not long ago I received a letter concerning a choice young deacon, Isaac Reiter, and the deacons, teachers, and priests who served, lifted, and blessed his life and their own lives.

Isaac fought cancer from the time he was seven months old until his death at age 13. When he and his family moved to a home near a hospital so that Isaac could receive proper medical attention, the Aaronic Priesthood members in the nearby ward were asked to provide the sacrament to them each Sunday. This weekly ordinance became a favorite of the Aaronic Priesthood holders who participated. Along with their leaders and Isaac’s family, they would gather around Isaac’s hospital bed, sing hymns, and share testimonies. Then the sacrament would be blessed. Isaac always insisted that, as a deacon, he pass the sacrament to his family and to those who had brought it. As he lay in his bed, he gathered the strength to hold a plate of either the blessed bread or water. All present would come to Isaac and partake of the sacrament from the plate. Nurses and other medical staff soon began to participate in the meeting as they realized that Isaac was close to his Heavenly Father and always honored Him. Though weak and in pain, Isaac always held himself with the honor of someone holding a royal priesthood.

Isaac was a great example to the young men in the ward. They saw his desire to fulfill his duties, even on his deathbed, and they realized that those duties were really privileges. They began showing up earlier in order to prepare the sacrament and to be in their seats on time. There was more reverence.

Isaac Reiter became a living sermon concerning honoring the priesthood. At his funeral, it was said that throughout his life he had one foot in heaven. No doubt he continues to magnify his duties and assist in the work beyond the veil.

For those of us who hold the Melchizedek Priesthood, our privilege to magnify our callings is ever present. We are shepherds watching over Israel. The hungry sheep do look up, ready to be fed the bread of life. Are we prepared, brethren, to feed the flock of God? It is imperative that we recognize the worth of a human soul, that we never give up on one of His precious sons.

Should there be anyone who feels he is too weak to do better because of that greatest of fears, the fear of failure, there is no more comforting assurance to be had than the words of the Lord: “My grace is sufficient for all men that humble themselves before me; for if they humble themselves before me, and have faith in me, then will I make weak things become strong unto them.”

Miracles are everywhere to be found when priesthood callings are magnified. When faith replaces doubt, when selfless service eliminates selfish striving, the power of God brings to pass His purposes. Whom God calls, God qualifies.

May our Heavenly Father ever bless, ever inspire, and ever lead all who hold His precious priesthood is my sincere prayer, and I offer it in the name of the Lord Jesus Christ, amen.

\newpage
\settalk{True to the Faith}
\setspeaker{Thomas S. Monson}
\settalkdate{April 2006}
\talkheading{True to the Faith}{Thomas S. Monson}

Many years ago, on an assignment to the beautiful islands of Tonga, I was privileged to visit our Church school, the Liahona High School, where our youth are taught by teachers with a common bond of faith—providing training for the mind and preparation for life. On that occasion, entering one classroom, I noticed the rapt attention the children gave their native instructor. His textbook and theirs lay closed upon the desks. In his hand he held a strange-appearing fishing lure fashioned from a round stone and large seashells. This, I learned, was a maka-feke, an octopus lure. In Tonga, octopus meat is a delicacy.

The teacher explained that Tongan fishermen glide over a reef, paddling their outrigger canoes with one hand and dangling the maka-feke over the side with the other. An octopus dashes out from its rocky lair and seizes the lure, mistaking it for a much-desired meal. So tenacious is the grasp of the octopus and so firm is its instinct not to relinquish the precious prize that fishermen can flip it right into the canoe.

It was an easy transition for the teacher to point out to the eager and wide-eyed youth that the evil one—even Satan—has fashioned so-called maka-fekes with which to ensnare unsuspecting persons and take possession of their destinies.

Today we are surrounded by the maka-fekes which the evil one dangles before us and with which he attempts to entice us and then to ensnare us. Once grasped, such maka-fekes are ever so difficult—and sometimes nearly impossible—to relinquish. To be safe, we must recognize them for what they are and then be unwavering in our determination to avoid them.

Constantly before us is the maka-feke of immorality. Almost everywhere we turn, there are those who would have us believe that what was once considered immoral is now acceptable. I think of the scripture, “Wo unto them that call evil good, and good evil, that put darkness for light, and light for darkness.” Such is the maka-feke of immorality. We are reminded in the Book of Mormon that chastity and virtue are precious above all things.

When temptation comes, remember the wise counsel of the Apostle Paul, who declared, “There hath no temptation taken you but such as is common to man: but God is faithful, who will not suffer you to be tempted above that ye are able; but will with the temptation also make a way to escape, that ye may be able to bear it.”

Next, the evil one also dangles before us the maka-feke of pornography. He would have us believe that the viewing of pornography really hurts no one. How applicable is Alexander Pope’s classic An Essay on Man:

\begin{verse}
Vice is a monster of so frightful mien,\\
As to be hated needs but to be seen;\\
Yet seen too oft, familiar with her face,\\
We first endure, then pity, then embrace.
\end{verse}

Some publishers and printers prostitute their presses by printing millions of pieces of pornography each day. No expense is spared to produce a product certain to be viewed, then viewed again. One of the most accessible sources of pornography today is the Internet, where one can turn on a computer and instantly have at his fingertips countless sites featuring pornography. President Gordon B. Hinckley has said: “I fear this may be going on in some of your homes. It is vicious. It is lewd and filthy. It is enticing and habit-forming. It will take [you] down to destruction as surely as anything in this world. It is foul sleaze that makes its exploiters wealthy, its victims impoverished.”

Tainted as well is the movie producer, the television programmer, or the entertainer who promotes pornography. Long gone are the restraints of yesteryear. So-called realism is the quest, with the result that today we are surrounded by this filth.

Avoid any semblance of pornography. It will desensitize the spirit and erode the conscience. We are told in the Doctrine and Covenants, “That which doth not edify is not of God, and is darkness.” Such is pornography.

I mention next the maka-feke of drugs, including alcohol. Once grasped, this maka-feke is particularly difficult to abandon. Drugs and alcohol cloud thinking, remove inhibitions, fracture families, shatter dreams, and shorten life. They are everywhere to be found and are placed purposely in the pathway of vulnerable youth.

Each one of us has a body that has been entrusted to us by a loving Heavenly Father. We have been commanded to care for it. Can we deliberately abuse or injure our bodies without being held accountable? We cannot! The Apostle Paul declared: “Know ye not that ye are the temple of God, and that the Spirit of God dwelleth in you? …

“The temple of God is holy, which temple ye are.” May we keep our bodies—our temples—fit and clean, free from harmful substances which destroy our physical, mental, and spiritual well-being.

The final maka-feke I wish to mention today is one which can crush our self-esteem, ruin relationships, and leave us in desperate circumstances. It is the maka-feke of excessive debt. It is a human tendency to want the things which will give us prominence and prestige. We live in a time when borrowing is easy. We can purchase almost anything we could ever want just by using a credit card or obtaining a loan. Extremely popular are home equity loans, where one can borrow an amount of money equal to the equity he has in his home. What we may not realize is that a home equity loan is equivalent to a second mortgage. The day of reckoning will come if we have continually lived beyond our means.

My brothers and sisters, avoid the philosophy that yesterday’s luxuries have become today’s necessities. They aren’t necessities unless we make them so. Many enter into long-term debt, only to find that changes occur: people become ill or incapacitated, companies fail or downsize, jobs are lost, natural disasters befall us. For many reasons, payments on large amounts of debt can no longer be made. Our debt becomes as a Damocles sword hanging over our heads and threatening to destroy us.

I urge you to live within your means. One cannot spend more than one earns and remain solvent. I promise you that you will then be happier than you would be if you were constantly worrying about how to make the next payment on nonessential debt. In the Doctrine and Covenants we read: “Pay the debt thou hast contracted. … Release thyself from bondage.”

There are, of course, countless other maka-fekes which the evil one dangles before us to lead us from the path of righteousness. However, our Heavenly Father has given us life and with it the capacity to think, to reason, and to love. We have the power to resist any temptation and the ability to determine the path we will take, the direction we will travel. Our goal is the celestial kingdom of God. Our purpose is to steer an undeviating course in that direction.

To all who walk the pathway of life, our Heavenly Father cautions: beware the detours, the pitfalls, the traps. Cunningly positioned are those cleverly disguised maka-fekes beckoning us to grasp them and to lose that which we most desire. Do not be deceived. Pause to pray. Listen to that still, small voice which speaks to the depths of our souls the Master’s gentle invitation, “Come, follow me.” By doing so, we turn from destruction, from death, and find happiness and life everlasting.

Yet there are those who do not hear, who will not obey, who listen to the enticings of the evil one, who grasp those maka-fekes until they cannot let go, until all is lost. I think of that person of power, that cardinal of the cloth, even Cardinal Wolsey. The prolific pen of William Shakespeare described the majestic heights, the pinnacle of power to which Cardinal Wolsey ascended. That same pen told how principle was eroded by vain ambition, by expediency, by a clamor for prominence and prestige. Then came the tragic descent, the painful lament of one who had gained everything, then lost it all.

To Cromwell, his faithful servant, Cardinal Wolsey speaks:

\begin{verse}
O Cromwell, Cromwell!\\
Had I but served my God with half the zeal\\
I served my king, He would not in mine age\\
Have left me naked to mine enemies.
\end{verse}

That inspired mandate which would have led Cardinal Wolsey to safety was ruined by the pursuit of power and prominence, the quest for wealth and position. Like others before him and many more yet to follow, Cardinal Wolsey fell.

In an earlier time and by a wicked king, a servant of God was tested. Aided by the inspiration of heaven, Daniel interpreted to King Belshazzar the writing on the wall. Concerning the proffered rewards—even a royal robe and a necklace of gold—Daniel said, “Let thy gifts be to thyself, and give thy rewards to another.”

Darius, a later king, also honored Daniel, elevating him to the highest position of prominence. There followed the envy of the crowd, the jealousy of princes, and the scheming of ambitious men.

Through trickery and flattery, King Darius signed a proclamation providing that anyone who made a request of any god or man, except the king, should be thrown into the lions’ den. Prayer was forbidden. In such matters, Daniel took direction not from an earthly king but from the King of heaven and earth, his God. Overtaken in his daily prayers, Daniel was brought before the king. Reluctantly, the penalty was pronounced. Daniel was to be thrown into the lions’ den.

I love the biblical account which follows:

“The king arose very early in the morning, and went in haste unto the den of lions.

“And when he came to the den, he cried with a lamentable voice … , O Daniel, … is thy God, whom thou servest continually, able to deliver thee from the lions?

“Then said Daniel unto the king, …

“My God hath sent his angel, and hath shut the lions’ mouths, that they have not hurt me. …

“Then was the king exceeding glad. … Daniel was taken up out of the den, and no manner of hurt was found upon him, because he believed in his God.”

In a time of critical need, Daniel’s determination to remain true and faithful provided divine protection and a sanctuary of safety.

The clock of history, like the sands of the hourglass, marks the passage of time. A new cast occupies the stage of life. The problems of our day loom ominously before us. Surrounded by the challenges of modern living, we look heavenward for that unfailing sense of direction, that we might chart and follow a wise and proper course. Our Heavenly Father will not deny our petition.

When I think of righteous individuals, the names of Gustav and Margarete Wacker come readily to mind. Let me describe them. I first met the Wackers when I was called to preside over the Canadian Mission in 1959. They had immigrated to Kingston, Ontario, Canada, from their native Germany.

Brother Wacker earned his living as a barber. His means were limited, but he and Sister Wacker always paid more than a tenth as tithing. As branch president, Brother Wacker started a missionary fund, and for months at a time he was the only contributor. When there were missionaries in the city, the Wackers fed and cared for them, and the missionaries never left the Wacker home without some tangible donation to their work and welfare.

Gustav and Margarete Wacker’s home was a heaven. They were not blessed with children, but they mothered and fathered their many Church visitors. Men and women of learning and sophistication sought out these humble, unlettered servants of God and counted themselves fortunate if they could spend an hour in their presence. The Wackers’ appearance was ordinary, their English halting and somewhat difficult to understand, their home unpretentious. They didn’t own a car or a television, nor did they do any of the things to which the world usually pays attention. Yet the faithful beat a path to their door in order to partake of the spirit that was there.

In March of 1982, Brother and Sister Wacker were called to serve as full-time ordinance workers in the Washington D.C. Temple. On June 29, 1983, while Brother and Sister Wacker were still serving in this temple assignment, Brother Wacker, with his beloved wife at his side, peacefully passed from mortality to his eternal reward. Fitting are the words, “Who honors God, God honors.”

My brothers and sisters, let us resolve here and now to follow that straight path which leads home to the Father of us all so that the gift of eternal life—life in the presence of our Heavenly Father—may be ours. Should there be those things which need to be changed or corrected in order to do so, I encourage you to take care of them now.

In the words of a familiar hymn, may we ever be

\begin{verse}
True to the faith that our parents have cherished,\\
True to the truth for which martyrs have perished,\\
To God’s command,\\
Soul, heart, and hand,\\
Faithful and true we will ever stand.
\end{verse}

That each of us may do so is my humble prayer, in the name of Jesus Christ, amen.

\newpage
\settalk{How Firm a Foundation}
\setspeaker{Thomas S. Monson}
\settalkdate{October 2006}
\talkheading{How Firm a Foundation}{Thomas S. Monson}

My dear brothers and sisters, both within my view and assembled throughout the world, I seek an interest in your faith and prayers as I respond to the assignment and privilege to address you.

In 1959, not long after I began my service as president of the Canadian Mission, headquartered in Toronto, Ontario, Canada, I met N. Eldon Tanner, a prominent Canadian who just months later would be called as an Assistant to the Quorum of the Twelve Apostles, then to the Quorum of the Twelve, and then as a counselor to four Presidents of the Church.

At the time I met him, President Tanner was president of the vast Trans-Canada Pipelines, Ltd., and president of the Canada Calgary Stake. He was known as “Mr. Integrity” in Canada. During that first meeting, we discussed, among other subjects, the cold Canadian winters, where storms rage, temperatures can linger well below freezing for weeks at a time, and where icy winds lower those temperatures even further. I asked President Tanner why the roads and highways in western Canada basically remained intact during such winters, showing little or no signs of cracking or breaking, while the road surfaces in many areas where winters are less cold and less severe developed cracks and breaks and potholes.

Said he, “The answer is in the depth of the base of the paving materials. In order for them to remain strong and unbroken, it is necessary to go very deep with the foundation layers. When the foundations are not deep enough, the surfaces cannot withstand the extremes of weather.”

Over the years I have thought often of this conversation and of President Tanner’s explanation, for I recognize in his words a profound application for our lives. Stated simply, if we do not have a deep foundation of faith and a solid testimony of truth, we may have difficulty withstanding the harsh storms and icy winds of adversity which inevitably come to each of us.

Mortality is a period of testing, a time to prove ourselves worthy to return to the presence of our Heavenly Father. In order for us to be tested, we must face challenges and difficulties. These can break us, and the surface of our souls may crack and crumble—that is, if our foundations of faith, our testimonies of truth are not deeply embedded within us.

We can rely on the faith and testimony of others only so long. Eventually we must have our own strong and deeply placed foundation, or we will be unable to withstand the storms of life, which will come. Such storms come in a variety of forms. We may be faced with the sorrow and heartbreak of a wayward child who chooses to turn from the pathway leading to eternal truth and rather travel the slippery slopes of error and disillusionment. Sickness may strike us or a loved one, bringing suffering and sometimes death. Accidents may leave their cruel marks of remembrance or may snuff out life. Death comes to the aged as they walk on faltering feet. Its summons is heard by those who have scarcely reached midway in life’s journey, and often it hushes the laughter of little children.

At times there appears to be no light at the tunnel’s end, no dawn to break the night’s darkness. We feel surrounded by the pain of broken hearts, the disappointment of shattered dreams, and the despair of vanished hopes. We join in uttering the biblical plea, “Is there no balm in Gilead?” (Jeremiah 8:22). We are inclined to view our own personal misfortunes through the distorted prism of pessimism. We feel abandoned, heartbroken, alone.

How can we build a foundation strong enough to withstand such vicissitudes of life? How can we maintain the faith and testimony which will be required, that we might experience the joy promised to the faithful? Constant, steady effort is necessary. Most of us have experienced inspiration so strong that it brings tears to our eyes and a determination to ever remain faithful. I have heard the statement, “If I could just keep these feelings with me always, I would never have trouble doing what I should.” Such feelings, however, can be fleeting. The inspiration we feel during these conference sessions may diminish and fade as Monday comes and we face the routines of work, of school, of managing our homes and families. Such can easily take our minds from the holy to the mundane, from that which uplifts to that which, if we allow it, will chip away at our testimonies, our strong foundations.

Of course we do not live in a world where we experience nothing but the spiritual, but we can fortify our foundations of faith, our testimonies of truth, so that we will not falter, we will not fail. How, you may ask, can we most effectively gain and maintain the foundation needed to survive spiritually in the world in which we live?

May I offer three guidelines to help us in our quest.

First, fortify your foundation through prayer. “Prayer is the soul’s sincere desire, uttered or unexpressed” (“Prayer Is the Soul’s Sincere Desire,” Hymns, no. 145).

As we pray, let us really communicate with our Father in Heaven. It is easy to let our prayers become repetitious, expressing words with little or no thought behind them. When we remember that each of us is literally a spirit son or daughter of God, we will not find it difficult to approach Him in prayer. He knows us; He loves us; He wants what is best for us. Let us pray with sincerity and meaning, offering our thanks and asking for those things we feel we need. Let us listen for His answers, that we may recognize them when they come. As we do, we will be strengthened and blessed. We will come to know Him and His desires for our lives. By knowing Him, by trusting His will, our foundations of faith will be strengthened. If any one of us has been slow to hearken to the counsel to pray always, there is no finer hour to begin than now. William Cowper declared, “Satan trembles when he sees the weakest saint upon his knees” (in William Neil, comp., Concise Dictionary of Religious Quotations [1974], 144).

Let us not neglect our family prayers. Such is an effective deterrent to sin, and thence a most beneficent provider of joy and happiness. That old saying is yet true: “The family that prays together stays together.” By providing an example of prayer to our children, we will also be helping them to begin their own deep foundations of faith and testimonies, which they will need throughout their lives.

My second guideline: Let us study the scriptures and “meditate therein day and night,” as counseled by the Lord in the book of Joshua (1:8).

In 2005, hundreds of thousands of Latter-day Saints accepted President Gordon B. Hinckley’s challenge to read the Book of Mormon by the end of the year. I do believe December of 2005 would set an all-time record for hours devoted to meeting the challenge on time. We were blessed when we completed the task; our testimonies were strengthened, our knowledge increased. I would encourage all of us to continue to read and study the scriptures, that we might understand them and apply in our lives the lessons we find there. I paraphrase the poet James Phinney Baxter:

\begin{verse}
Who learns and learns but never knows\\
Is like the one who plows and plows but never sows.
\end{verse}

(“The Baxter Collection,” Baxter Memorial Library, Gorham, Maine)

Spending time each day in scripture study will, without doubt, strengthen our foundations of faith and our testimonies of truth.

Recall with me the joy Alma experienced as he was journeying from the land of Gideon southward to the land of Manti and met the sons of Mosiah. Alma had not seen them for some time, and he was overjoyed to discover that they were “still his brethren in the Lord; yea, and they had waxed strong in the knowledge of the truth; for they were men of a sound understanding and they had searched the scriptures diligently, that they might know the word of God” (see Alma 17:1–2).

May we also know the word of God and conduct our lives accordingly.

My third guideline for building a strong foundation of faith and testimony involves service.

While driving to the office one morning, I passed a dry-cleaning establishment which had a sign in the window. It read, “It’s the Service That Counts.” The sign’s message simply would not leave my mind. Suddenly I realized why. In actual fact it is the service that counts—the Lord’s service.

In the Book of Mormon we read of noble King Benjamin. In the true humility of an inspired leader, he recounted his desire to serve his people and lead them in paths of righteousness. He then declared to them:

“Because I said unto you that I had spent my days in your service, I do not desire to boast, for I have only been in the service of God.

“And behold, I tell you these things that ye may learn wisdom; that ye may learn that when ye are in the service of your fellow beings ye are only in the service of your God” (Mosiah 2:16–17).

This is the service that counts, the service to which all of us have been called: the service of the Lord Jesus Christ.

Along your pathway of life you will observe that you are not the only traveler. There are others who need your help. There are feet to steady, hands to grasp, minds to encourage, hearts to inspire, and souls to save.

Thirteen years ago it was my privilege to provide a blessing to a beautiful 12-year-old young lady, Jami Palmer. She had just been diagnosed with cancer and was frightened and bewildered. She subsequently underwent surgery and painful chemotherapy. Today she is cancer-free and is a bright, beautiful 26-year-old who has accomplished much in her life. Some time ago, I learned that in her darkest hour, when any future appeared somewhat grim, she learned that her leg where the cancer was situated would require multiple surgeries. A long-planned hike with her Young Women class up a rugged trail to Timpanogos Cave—located in the Wasatch Mountains about 40 miles south of Salt Lake City, Utah—was out of the question, she thought. Jami told her friends they would have to undertake the hike without her. I’m confident there was a catch in her voice and disappointment in her heart. But then the other young women responded emphatically, “No, Jami, you are going with us!”

“But I can’t walk,” came the anguished reply.

“Then, Jami, we’ll carry you to the top!” And they did.

Today, the hike is a memory, but in reality it is much more. James Barrie, the Scottish poet, declared, “God gave us memories, that we might have June roses in the December of our lives” (paraphrasing James Barrie, in Laurence J. Peter, comp., Peter’s Quotations: Ideas for Our Time [1977], 335). None of those precious young women will ever forget that memorable day when a loving Heavenly Father looked down with a smile of approval and was well pleased.

As He enlists us to His cause, He invites us to draw close to Him, and we feel His Spirit in our lives.

As we establish a firm foundation for our lives, let us each one remember His precious promise:

\begin{verse}
Fear not, I am with thee; oh, be not dismayed,\\
For I am thy God and will still give thee aid.\\
I’ll strengthen thee, help thee, and cause thee to stand,\\
Upheld by my righteous, omnipotent hand.
\end{verse}

(“How Firm a Foundation,” Hymns, no. 85)

May each of us qualify for this blessing, I humbly pray, in the name of Jesus Christ, our Savior, amen.

\newpage
\settalk{True to Our Priesthood Trust}
\setspeaker{Thomas S. Monson}
\settalkdate{October 2006}
\talkheading{True to Our Priesthood Trust}{Thomas S. Monson}

A few weeks ago at a fast and testimony meeting at our ward, I watched a little boy on the back row mustering up courage to bear his testimony. He made three or four false starts and then sat down. Finally it was his turn. He squared his little shoulders, walked bravely up the aisle to the stand, took the two steps up to the level of the pulpit, stepped over and put his hands on the pulpit, gazed into the congregation, smiled—and then turned around, went back off those two steps and down the same aisle to his mother and father. I looked at you tonight in this vast Conference Center and thought of those listening in and could appreciate more fully the actions of that little boy.

My brethren, I am honored by the privilege to speak to you this evening. I have contemplated what I might say to you. There has come to my mind a favorite scripture from Ecclesiastes: “Fear God, and keep his commandments: for this is the whole duty of man” (Ecclesiastes 12:13). I love, I cherish the noble word duty.

The legendary General Robert E. Lee of American Civil War fame declared: “Duty is the sublimest word in our language. … You cannot do more. You should never wish to do less” (in John Bartlett, Familiar Quotations [1968], 620).

Each of us has duties associated with the sacred priesthood which we bear. Whether we bear the Aaronic or the Melchizedek Priesthood, much is expected of each of us. The Lord Himself summed up our responsibility when He, in the revelation on the priesthood, urged, “Wherefore, now let every man learn his duty, and to act in the office in which he is appointed, in all diligence” (D\&C 107:99).

I hope with all my heart and soul that every young man who receives the priesthood will honor that priesthood and be true to the trust which is conveyed when it is conferred.

Fifty-one years ago I heard William J. Critchlow Jr., then president of the South Ogden Stake who would later become an Assistant to the Quorum of the Twelve, speak to the brethren of the general priesthood session of conference and retell a story concerning trust, honor, and duty. May I share the story with you. Its simple lesson applies to us today, as it did then.

“[Young] Rupert stood by the side of the road watching an unusual number of people hurry past. At length he recognized a friend. ‘Where are all of you going in such a hurry?’ he asked.

“The friend paused. ‘Haven’t you heard?’ he said.

“‘I’ve heard nothing,’ Rupert answered.

“‘Well,’ continued [the] friend, ‘the King has lost his royal emerald! Yesterday he attended a wedding of the nobility and wore the emerald on the slender golden chain around his neck. In some way the emerald became loosened from the chain. Everyone is searching, for the King has offered a reward … to the one who finds it. Come, we must hurry.’

“‘But I cannot go without asking Grandmother,’ faltered Rupert.

“‘Then I cannot wait. I want to find the emerald,’ replied his friend.

“Rupert hurried back to the cabin at the edge of the woods to seek his grandmother’s permission. ‘If I could find it we could leave this hut with its dampness and buy a piece of land up on the hillside,’ he pleaded with Grandmother.

“But his grandmother shook her head. ‘What would the sheep do?’ she asked. ‘Already they are restless in the pen, waiting to be taken to the pasture, and please do not forget to take them to water when the sun shines high in the heavens.’

“Sorrowfully, Rupert took the sheep to the pasture, and at noon he led them to the brook in the woods. There he sat on a large stone by the stream. ‘If I could only have had a chance to look for the King’s emerald!’ he thought. Turning his head to gaze down at the sandy bottom of the brook, suddenly he stared into the water. What was it? It could not be! He leaped into the water, and his gripping fingers held something that was green with a slender bit of gold chain [that had been broken]. ‘The King’s emerald!’ he shouted. ‘It must have been flung from the chain when the King [astride his horse galloped across the bridge spanning the stream, and the current carried] it here.’

“With shining eyes Rupert ran to his grandmother’s hut to tell her of his great find. ‘Bless you, my boy,’ she said, ‘but you never would have found it if you had not been doing your duty, herding the sheep.’ And Rupert knew that this was the truth” (in Conference Report, Oct. 1955, 86; paragraphing, capitalization, and punctuation altered).

The lesson to be learned from this story is found in the familiar couplet: “Do [your] duty; that is best; Leave unto [the] Lord the rest!” (Henry Wadsworth Longfellow, “The Legend Beautiful,” in The Complete Poetical Works of Longfellow [1893], 258).

To you who are or have been presidents of your quorums, may I suggest that your duty does not end when your term of office concludes. That relationship with your quorum members, your duty to them, continues throughout your life.

During the time I was a teacher in the Aaronic Priesthood, I was called to be president of the quorum. With the urging and assistance of a dedicated and inspired quorum adviser, I worked diligently to ensure that each of the young men attended our meetings regularly. Two of them were a particular challenge, but with our perseverance and love and a little persuasion, they began to attend meetings and participate in quorum activities. However, as time passed and they left the ward to pursue education and employment, each of them drifted back into inactivity.

Over the years I have seen each of these two dear friends at various functions. Whenever I do, I place a hand on their shoulder and remind them, “I’m still your quorum president, and I won’t let go. You mean so much to me, and I want you to enjoy the blessings which come with activity in the Church.” They know I love them and that I’ll never ever give up on them.

For those of us who hold the Melchizedek Priesthood, our privilege to magnify our callings is ever present. We are shepherds watching over Israel. The hungry sheep look up, ready to be fed the bread of life.

Many years ago, on a Halloween night, it was my privilege to be of assistance to one who had temporarily lost his way and needed a helping hand to return. I was driving home from the office rather late. I had been stalling on Halloween, letting my wife handle the trick-or-treat visitors. As I passed St. Mark’s Hospital in Salt Lake City, I remembered that a dear friend, Max, lay ill in that very hospital. As he and I had become acquainted years before, we discovered that we had grown up in the same ward, although at different times. By the time I was born, Max and his parents had moved from the ward.

That Halloween night, I drove into the parking lot and entered the hospital. As I stopped at the desk to inquire as to his room number, I was informed that when Max had registered at the hospital, he had listed as his religious preference not LDS but rather another church.

I entered Max’s room and greeted him. I told him how proud I was to be his friend and how much I cared about him. I talked about his career in banking and as an orchestra leader on the side. I discovered that he had been offended by a comment or two from others and so had decided to attend another church. I said to him, “Max, you hold the Melchizedek Priesthood. I would like to give you a blessing tonight.” He agreed, and the blessing was provided. He then informed me that his wife, Bernice, was also very ill and was, in fact, in an adjoining room. At my invitation, Max joined me in giving a blessing to her. He asked me to help him. I coached him. He anointed his wife. There were tears and embraces all around as I sealed the anointing with Max, his hands on his wife’s head with mine, making that Halloween evening one ever to be remembered.

As I left the hospital that night, I stopped at the desk and told the receptionist that with the permission of Max and his wife the record should be changed to reflect their membership in The Church of Jesus Christ of Latter-day Saints. I waited and I watched until it was changed.

My friends Max and Bernice are now both on the other side of the veil, but they spent the last period of their lives active and happy and receiving the blessings which come with testimonies of the gospel and attendance at church.

Brethren, our task is to reach out to those who, for whatever reason, are in need of our help. Our challenge is not insurmountable. We are on the Lord’s errand, and therefore we are entitled to the Lord’s help. But we must try. From the play Shenandoah comes the spoken line which inspires: “If we don’t try, then we don’t do; and if we don’t do, then why are we here?”

Ours is the responsibility to so conduct our lives that when the call comes to provide a priesthood blessing or to assist in any way, we are worthy to do so. We have been told that truly we cannot escape the effect of our personal influence. We must be certain that our influence is positive and uplifting.

Are our hands clean? Are our hearts pure? Looking backward in time through the pages of history, we find a lesson on worthiness gleaned from the words of the dying King Darius. Through the proper rites, Darius had been recognized as legitimate king of Egypt. His rival, Alexander the Great, had been declared legitimate son of Amon. He too was Pharaoh. Alexander, finding the defeated Darius on the point of death, laid his hands upon his head to heal him, commanding him to arise and resume his kingly power, concluding, “I swear unto thee, Darius, by all the gods, that I do these things truly and without fakery.”

Darius replied with a gentle rebuke: “Alexander, my boy, … do you think you can touch heaven with those hands of yours?” (Adapted from Hugh Nibley, Abraham in Egypt [1981], 192.)

The call of duty can come quietly as we who hold the priesthood respond to the assignments we receive. President George Albert Smith, that modest yet effective leader and eighth President of the Church, declared, “It is your duty first of all to learn what the Lord wants and then by the power and strength of His holy Priesthood to magnify your calling in the presence of your fellows in such a way that the people will be glad to follow you” (in Conference Report, Apr. 1942, 14).

And how does one magnify a calling? Simply by performing the service that pertains to it.

Brethren, it is in doing—not just dreaming—that lives are blessed, others are guided, and souls are saved. “Be ye doers of the word, and not hearers only, deceiving your own selves,” declared James (James 1:22).

May all of us assembled tonight in this priesthood meeting make a renewed effort to qualify for the Lord’s guidance in our lives. There are so many out there who plead and pray for help. There are those who are discouraged, those who long to return but who don’t know how to begin.

I’ve always believed in the truth of the words “God’s sweetest blessings always go by hands that serve him here below” (Whitney Montgomery, “Revelation,” in Best-Loved Poems of the LDS People, ed. Jack M. Lyon and others [1996], 283). Let us have ready hands, clean hands, and willing hearts, that we may participate in providing what our Heavenly Father would have others receive from Him.

I conclude with an example from my own life. I once had a treasured friend who seemed to experience more of life’s troubles and frustrations than he could bear. Finally he lay in the hospital terminally ill. I knew not that he was there.

Sister Monson and I had gone to that same hospital to visit another person who was very ill. As we exited the hospital and proceeded to where our car was parked, I felt the distinct impression to return and make inquiry concerning whether my friend Hyrum might still be a patient there. A check with the clerk at the desk confirmed that Hyrum was indeed a patient there after many weeks.

We proceeded to his room, knocked on the door, and opened it. We were not prepared for the sight that awaited us. Balloon bouquets were everywhere. Prominently displayed on the wall was a poster with the words “Happy Birthday, Daddy” written on it. Hyrum was sitting up in his hospital bed, his family members by his side. When he saw us, he said, “Brother Monson, how in the world did you know that today is my birthday?” I smiled, but I left the question unanswered.

Those in the room who held the Melchizedek Priesthood surrounded this, their father and grandfather and my friend, and a priesthood blessing was given.

After tears were shed, smiles of gratitude exchanged, and tender hugs received and given, I leaned over to Hyrum and spoke softly to him: “Remember the words of the Lord, for they will sustain you. He promised you, ‘I will not leave you comfortless: I will come to you’ (John 14:18).”

Time marches on. Duty keeps cadence with that march. Duty does not dim nor diminish. Catastrophic conflicts come and go, but the war waged for the souls of men continues without abatement. Like a clarion call comes the word of the Lord to you and to me, and to priesthood holders everywhere. I reiterate that word: “Wherefore, now let every man learn his duty, and to act in the office in which he is appointed, in all diligence” (D\&C 107:99).

Brethren, let us learn our duties. Let us ever be worthy to perform those duties and, in so doing, follow in the footsteps of the Master. When to Him came the call of duty, He answered, “Father, thy will be done, and the glory be thine forever” (Moses 4:2). May we do likewise, I pray humbly, in the name of Jesus Christ the Lord, amen.

\newpage
\settalk{I Know That My Redeemer Lives!}
\setspeaker{Thomas S. Monson}
\settalkdate{April 2007}
\talkheading{I Know That My Redeemer Lives!}{Thomas S. Monson}

Recently I was looking through some family photo albums. Cherished memories flooded my mind as I came across image after image of loved ones gathered at family outings, birthdays, reunions, anniversaries. Since these photographs were taken, some of those beloved family members have departed this life. I thought of the words of the Lord, “Thou shalt live together in love, insomuch that thou shalt weep for the loss of them that die.” I miss each one who has left our family circle.

Though difficult and painful, death is an essential part of our mortal experience. We began our sojourn here by leaving our premortal existence and coming to this earth. The poet Wordsworth captured that journey in his inspired ode to immortality. He wrote:

\begin{verse}
Our birth is but a sleep and a forgetting:\\
The Soul that rises with us, our life’s Star,\\
Hath had elsewhere its setting,\\
And cometh from afar:\\
Not in entire forgetfulness,\\
And not in utter nakedness,\\
But trailing clouds of glory do we come\\
From God, who is our home:\\
Heaven lies about us in our infancy!
\end{verse}

Life moves on. Youth follows childhood, and maturity comes ever so imperceptibly. As we search and ponder the purpose and the problems of life, all of us sooner or later face the question of the length of life and of a personal, everlasting life. These questions most insistently assert themselves when loved ones leave us or when we face leaving those we love.

At such times, we ponder the universal question, best phrased by Job of old, who centuries ago asked, “If a man die, shall he live again?”

Today, as always, the skeptic’s voice challenges the word of God, and each must choose to whom he will listen. Clarence Darrow, the famous lawyer and agnostic, declared, “No life is of much value, and … every death is [but a] little loss.” Schopenhauer, the German philosopher and pessimist, wrote, “To desire immortality is to desire the eternal perpetuation of a great mistake.” And to their words are added those of new generations, as foolish men crucify the Christ anew—for they modify His miracles, doubt His divinity, and reject His Resurrection.

Robert Blatchford, in his book God and My Neighbor, attacked with vigor accepted Christian beliefs, such as God, Christ, prayer, and immortality. He boldly asserted, “I claim to have proved everything I set out to prove so fully and decisively that no Christian, however great or able he may be, can answer my arguments or shake my case.” He surrounded himself with a wall of skepticism. Then a surprising thing happened. His wall suddenly crumbled to dust. He was left exposed and undefended. Slowly he began to feel his way back to the faith he had scorned and ridiculed. What had caused this profound change in his outlook? His wife died. With a broken heart, he went into the room where lay all that was mortal of her. He looked again at the face he loved so well. Coming out, he said to a friend: “It is she, and yet it is not she. Everything is changed. Something that was there before is taken away. She is not the same. What can be gone if it be not the soul?”

Later he wrote: “Death is not what some people imagine. It is only like going into another room. In that other room we shall find … the dear women and men and the sweet children we have loved and lost.”

Against the doubting in today’s world concerning Christ’s divinity, we seek a point of reference, an unimpeachable source, even a testimony of eyewitnesses. Stephen, from biblical times, doomed to the cruel death of a martyr, looked up to heaven and cried, “I see the heavens opened, and the Son of man standing on the right hand of God.”

Who can help but be convinced by the stirring testimony of Paul to the Corinthians? He declared “that Christ died for our sins according to the scriptures; and that he was buried, and that he rose again the third day according to the scriptures: and … was seen of Cephas, then of the twelve: … And,” said Paul, “last of all he was seen of me.”

In our dispensation, this same testimony was spoken boldly by the Prophet Joseph Smith, as he and Sidney Rigdon testified, “And now, after the many testimonies which have been given of him, this is the testimony, last of all, which we give of him: That he lives!”

This is the knowledge that sustains. This is the truth that comforts. This is the assurance that guides those who are bowed down with grief—out of the shadows and into the light.

On Christmas Eve, 1997, I met a remarkable family. Each member of the family had an unshakable testimony of the truth and of the reality of the Resurrection. The family consisted of a mother and father and four children. Each of the children—three sons and a daughter—had been born with a rare form of muscular dystrophy, and each was handicapped. Mark, who was then 16 years old, had undergone spinal surgery in an effort to help him move about more freely. The other two boys—Christopher, age 13, and Jason, age 10—were to leave for California in a few days to undergo similar surgery. The only daughter, Shanna, was then five years old—a beautiful child. All of the children were intelligent and faith-filled, and it was obvious that their parents, Bill and Sherry, were proud of each one. We visited for a while, and the special spirit of that family filled my office and my heart. The father and I gave blessings to the two boys who were facing surgery, and then the parents asked if little Shanna could sing for me. Her father mentioned that she had diminished lung capacity and that it might be difficult for her, but that she wanted to try. To the accompaniment of a recorded cassette, and in a beautiful, clear voice—never missing a note—she sang of a brighter future:

\begin{verse}
On a beautiful day that I dream about\\
In a world I would love to see\\
Is a beautiful place where the sun comes out\\
And it shines in the sky for me.\\
On this beautiful winter’s morning,\\
If my wish could come true somehow,\\
Then the beautiful day that I dream about\\
Would be here and now.
\end{verse}

The emotions of all of us were very near the surface as she finished. The spirituality of this visit set the tone for my Christmas that year.

I kept in touch with the family, and when the oldest son, Mark, turned 19, arrangements were made for him to serve a special mission at Church headquarters. Eventually, the other two brothers also had an opportunity to serve such missions.

Nearly a year ago, Christopher, who was then 22 years old, succumbed to the disease with which each of the children has been afflicted. And then, last September, I received word that little Shanna, now 14 years old, had passed away. At the funeral services, Shanna was honored by beautiful tributes. Leaning on the pulpit for support, each of her surviving brothers, Mark and Jason, shared poignant family experiences. Shanna’s mother sang a lovely musical number as part of a duet. Her father and grandfather gave touching sermons. Though their hearts were broken, each bore powerful and deep-felt testimony of the reality of the Resurrection and of the actuality that Shanna lives still, as does her brother Christopher, each awaiting a glorious reunion with their beloved family.

When it was my time to speak, I recounted that visit the family made to my office nearly nine years earlier and spoke of the lovely song Shanna sang on that occasion. I concluded with the thought: “Because our Savior died at Calvary, death has no hold upon any one of us. Shanna lives, whole and well, and for her that beautiful day she sang about on a special Christmas Eve in 1997, the day she dreamed about, is here and now.”

My brothers and sisters, we laugh, we cry, we work, we play, we love, we live. And then we die. Death is our universal heritage. All must pass its portals. Death claims the aged, the weary and worn. It visits the youth in the bloom of hope and the glory of expectation. Nor are little children kept beyond its grasp. In the words of the Apostle Paul, “It is appointed unto men once to die.”

And dead we would remain but for one Man and His mission, even Jesus of Nazareth. Born in a stable, cradled in a manger, His birth fulfilled the inspired pronouncements of many prophets. He was taught from on high. He provided the life, the light, and the way. Multitudes followed Him. Children adored Him. The haughty rejected Him. He spoke in parables. He taught by example. He lived a perfect life.

Though the King of kings and Lord of lords had come, He was accorded by some the greeting given to an enemy, a traitor. There followed a mockery which some called a trial. Cries of “Crucify him, crucify him” filled the air. Then commenced the climb to Calvary’s hill.

He was ridiculed, reviled, mocked, jeered, and nailed to a cross amidst shouts of “Let Christ the King of Israel descend now from the cross, that we may see and believe.” “He saved others; himself he cannot save.” His response: “Father, forgive them; for they know not what they do.” “Into thy hands I commend my spirit: and having said thus, he gave up the ghost.” His body was placed by loving hands in a sepulchre hewn of stone.

On the first day of the week, very early in the morning, Mary Magdalene and Mary the mother of James, along with others, came to the sepulchre. To their astonishment, the body of their Lord was gone. Luke records that two men in shining garments stood by them and said: “Why seek ye the living among the dead? He is not here, but is risen.”

Next week the Christian world will celebrate the most significant event in recorded history. The simple pronouncement, “He is not here, but is risen,” was the first confirmation of the literal Resurrection of our Lord and Savior, Jesus Christ. The empty tomb that first Easter morning brought comforting assurance, an affirmative answer to Job’s question, “If a man die, shall he live again?”

To all who have lost loved ones, we would turn Job’s question to an answer: If a man die, he shall live again. We know, for we have the light of revealed truth. “I am the resurrection, and the life,” spoke the Master. “He that believeth in me, though he were dead, yet shall he live: and whosoever liveth and believeth in me shall never die.”

Through tears and trials, through fears and sorrows, through the heartache and loneliness of losing loved ones, there is assurance that life is everlasting. Our Lord and Savior is the living witness that such is so.

With all my heart and the fervency of my soul, I lift up my voice in testimony as a special witness and declare that God does live. Jesus is His Son, the Only Begotten of the Father in the flesh. He is our Redeemer; He is our Mediator with the Father. He it was who died on the cross to atone for our sins. He became the firstfruits of the Resurrection. Because He died, all shall live again. “Oh, sweet the joy this sentence gives: ‘I know that my Redeemer lives!’” May the whole world know it and live by that knowledge, I humbly pray, in the name of Jesus Christ, the Lord and Savior, amen.

\newpage
\settalk{Tabernacle Memories}
\setspeaker{Thomas S. Monson}
\settalkdate{April 2007}
\talkheading{Tabernacle Memories}{Thomas S. Monson}

My brothers and sisters, both here in the Tabernacle and listening by various means throughout the world, it is a joy for me to stand before you once again in this magnificent building. In this setting one cannot help but feel the spirit of the early Saints who constructed this beautiful house of worship, as well as all those who over the years have labored to preserve and beautify it.

I have been thinking recently of the many significant events in my life which are associated with the Salt Lake Tabernacle. Although there are far too many for me to mention today, I would like to share just a few.

I recall the time I approached baptism, when I was eight years of age. My mother talked with me about repentance and about the meaning of baptism; and then, on a Saturday in September of 1935, she took me on a streetcar to the Tabernacle baptistry which, until recently, was here in this building. At the time it was not as customary as it is now for fathers to baptize their children, since the ordinance was generally performed on a Saturday morning or afternoon, and many fathers were working at their daily professions or trades. I dressed in white and was baptized. I remember that day as though it were yesterday and the happiness I felt at having had this ordinance performed.

Over the years and particularly during the time I served as a bishop, I witnessed many other baptisms in the Tabernacle font. Each was a special and inspiring occasion, and each served to remind me of my own baptism.

In April of 1950, my wife, Frances, and I were in attendance at the Sunday afternoon session of general conference, held in this building. President George Albert Smith was the President of the Church, and in closing the conference, he delivered an inspiring and powerful message concerning the Resurrection of our Lord and Savior Jesus Christ. Before he concluded his remarks, however, he sounded a prophetic warning. Said he: “It will not be long until calamities will overtake the human family unless there is speedy repentance. It will not be long before those who are scattered over the face of the earth by millions will die … because of what will come” (in Conference Report, Apr. 1950, 169). These were alarming words, for they came from a prophet of God.

Two and a half months after that general conference, on June 25, 1950, war broke out in Korea—a war which would eventually claim an estimated 2.5 million lives. This event prompted me to reflect on the statement President Smith made as we sat in this building that spring day.

I attended many general conference sessions in the Tabernacle, always being edified and inspired by the words of the Brethren. Then, in October of 1963, President David O. McKay invited me to his office and extended to me a call to serve as a member of the Quorum of the Twelve Apostles. He asked that I keep this sacred call confidential, revealing it to no one except my wife, and that I be present for general conference in the Tabernacle the next day, when my name would be read aloud.

The following morning I came into the Tabernacle not knowing exactly where to sit. Being a member of the Priesthood Home Teaching Committee, I determined that I would be seated among the members of that committee. I noticed a friend of mine by the name of Hugh Smith, who was also a member of the Priesthood Home Teaching Committee. He motioned for me to sit by him. I couldn’t say a thing to him about my call, but I sat down.

During the session, the members of the Quorum of the Twelve Apostles were sustained and, of course, my name was read. I believe the walk from the audience to the stand was the longest walk of my life.

It has been nearly 44 years since that conference. Until the year 2000, when the Conference Center was dedicated, it was my privilege to deliver 101 general conference messages from the pulpit in this building, not including those given at general auxiliary conferences and other meetings held here. My remarks today bring the total to 102. I have had many spiritual experiences over the years as I have stood here.

During the message I delivered at general conference in October 1975, I felt prompted to direct my remarks to a little girl with long blonde hair who was seated in the balcony of this building. I called the attention of the audience to her and felt a freedom of expression which testified to me that this small girl needed the message I had in mind concerning the faith of another young lady.

At the conclusion of the session, I returned to my office and found waiting for me a young child by the name of Misti White, together with her grandparents and an aunt. As I greeted them, I recognized Misti as the one in the balcony to whom I had directed my remarks. I learned that as her eighth birthday approached, she was in a quandary concerning whether or not to be baptized. She felt she would like to be baptized, and her grandparents, with whom she lived, wanted her to be baptized, but her less-active mother suggested she wait until she was 18 years of age to make the decision. Misti had told her grandparents, “If we go to conference in Salt Lake City, maybe Heavenly Father will let me know what I should do.”

Misti and her grandparents and her aunt had traveled from California to Salt Lake City for conference and were able to obtain seats in the Tabernacle for the Saturday afternoon session. This was where they were seated when my attention was drawn to Misti and my decision made to speak to her.

As we continued our visit after the session, Misti’s grandmother said to me, “I think Misti has something she would like to tell you.” This sweet young girl said, “Brother Monson, while you were speaking in conference, you answered my question. I want to be baptized!”

The family returned to California, and Misti was baptized and confirmed a member of The Church of Jesus Christ of Latter-day Saints. Through all the years since, Misti has remained true and faithful to the gospel of Jesus Christ. Fourteen years ago, it was my privilege to perform her temple marriage to a fine young man, and together they are rearing five beautiful children, with another one on the way.

My brothers and sisters, I feel privileged to be standing once again at the Tabernacle pulpit in this building which holds for me such wonderful memories. The Tabernacle is a part of my life—a part which I cherish.

I have been honored and pleased during my lifetime to raise my arm to the square in sustaining nine Church Presidents as their names have been read. This morning I joined you in sustaining once again our beloved prophet, President Gordon B. Hinckley. It is a joy and a privilege to serve by his side, along with President Faust.

As this building is rededicated today, may we pledge to rededicate our lives to the work of our Lord and Savior Jesus Christ, who so willingly died that we might live. May we follow in His footsteps each day, I pray humbly in the name of Jesus Christ, amen.

\newpage
\settalk{The Priesthood—a Sacred Gift}
\setspeaker{Thomas S. Monson}
\settalkdate{April 2007}
\talkheading{The Priesthood—a Sacred Gift}{Thomas S. Monson}

Brethren, we are assembled this evening as a mighty body of the priesthood, both here in the Conference Center and in locations throughout the world. I am honored by the privilege to speak to you. I pray that the inspiration of the Lord will guide my thoughts and inspire my words.

During the past several weeks, as I have contemplated what I might say to you tonight, I have thought repeatedly of the blessing which is ours to be bearers of the sacred priesthood of God. When we look at the world as a whole, with a population of over six and a half billion people, we realize that we comprise a very small, select group. We who hold the priesthood are, in the words of the Apostle Peter, “a chosen generation, a royal priesthood.”

President Joseph F. Smith defined the priesthood as “the power of God delegated to man by which man can act in the earth for the salvation of the human family, … by which [men] may speak the will of God as if the angels were here to speak it themselves; by which men are empowered to bind on earth and it shall be bound in heaven, and to loose on earth and it shall be loosed in heaven.” President Smith added, “[The priesthood] is sacred, and it must be held sacred by the people.”

My brethren, the priesthood is a gift which brings with it not only special blessings but also solemn responsibilities. It is our responsibility to conduct our lives so that we are ever worthy of the priesthood we bear. We live in a time when we are surrounded by much that is intended to entice us into paths which may lead to our destruction. To avoid such paths requires determination and courage.

Courage counts. This truth came to me in a most vivid and dramatic manner many years ago. I was serving as a bishop at the time. The general session of our stake conference was being held in the Assembly Hall on Temple Square in Salt Lake City. Our stake presidency was to be reorganized. The Aaronic Priesthood, including members of bishoprics, were providing the music for the conference. As we concluded singing our first selection, President Joseph Fielding Smith, our conference visitor, stepped to the pulpit and read for sustaining approval the names of the new stake presidency. He then mentioned that Percy Fetzer, who became our new stake president, and John Burt, who became the first counselor—each of whom had been counselors in the previous presidency—had been made aware of their new callings before the conference began. However, he indicated that I, who had been called to be second counselor in the new presidency, had no previous knowledge of the calling and was hearing of it for the first time as my name was read for sustaining vote. He then announced, “If Brother Monson is willing to respond to this call, we will be pleased to hear from him now.”

As I stood at the pulpit and gazed out on that sea of faces, I remembered the song we had just sung. It pertained to the Word of Wisdom and was titled “Have Courage, My Boy, to Say No.” That day I selected as my acceptance theme “Have Courage, My Boy, to Say Yes.” The call for courage comes constantly to each of us—the courage to stand firm for our convictions, the courage to fulfill our responsibilities, the courage to honor our priesthood.

Wherever we go, our priesthood goes with us. Are we standing in “holy places”? Said President J. Reuben Clark Jr., who served for many years as a counselor in the First Presidency: “The Priesthood is not like a suit of clothes that you can lay off and take back on. … Depending upon ourselves [it is] an everlasting endowment.” He continued: “If we really had that … conviction … that we could not lay [the priesthood] aside, and that God would hold us responsible if we [demeaned] it, it would save us from doing a good many things, save us from going a good many places. If, every time we started a little detour away from the straight and narrow, we would remember, ‘I am carrying my Priesthood here. Should I?’ it would not take us long to work back into the straight and narrow.”

President Spencer W. Kimball said: “There is no limit to the power of the priesthood which you hold. The limit comes in you if you do not live in harmony with the Spirit of the Lord and you limit yourselves in the power you exert.”

My brethren of the priesthood—from the youngest to the oldest—are you living your life in accordance with that which the Lord requires? Are you worthy to bear the priesthood of God? If you are not, make the decision here and now, muster the courage it will take, and institute whatever changes are necessary so that your life is what it should be. To sail safely the seas of mortality, we need the guidance of that eternal mariner—even the great Jehovah. If we are on the Lord’s errand, we are entitled to the Lord’s help.

His help has come to me on countless occasions throughout my life. During the final phases of World War II, I turned 18 and was ordained an elder—one week before I departed for active duty with the navy. A member of my ward bishopric was at the train station to bid me farewell. Just before train time, he placed in my hand a book which I hold before you tonight. Its title: The Missionary’s Hand Book. I laughed and commented, “I’ll be in the navy—not on a mission.” He answered, “Take it anyway. It may come in handy.”

It did. During basic training our company commander instructed us concerning how we might best pack our clothing in a large seabag. He then advised, “If you have a hard, rectangular object you can place in the bottom of the bag, your clothes will stay more firm.” I thought, “Where am I going to find a hard, rectangular object?” Suddenly I remembered just the right rectangular object—The Missionary’s Hand Book. And thus it served for 12 weeks at the bottom of that seabag.

The night preceding our Christmas leave, our thoughts were, as always, on home. The barracks were quiet. Suddenly I became aware that my buddy in the adjoining bunk—a member of the Church, Leland Merrill—was moaning in pain. I asked, “What’s the matter, Merrill?”

He replied, “I’m sick. I’m really sick.”

I advised him to go to the base dispensary, but he answered knowingly that such a course would prevent him from being home for Christmas. I then suggested he be quiet so that we didn’t awaken the entire barracks.

The hours lengthened; his groans grew louder. Then, in desperation, he whispered, “Monson, aren’t you an elder?” I acknowledged this to be so, whereupon he pleaded, “Give me a blessing.”

I became very much aware that I had never given a blessing. I had never received such a blessing; I had never witnessed a blessing being given. My prayer to God was a plea for help. The answer came: “Look in the bottom of the seabag.” Thus, at 2:00 a.m. I emptied on the deck the contents of the bag. I then took to the night-light that hard, rectangular object, The Missionary’s Hand Book, and read how one blesses the sick. With about 120 curious sailors looking on, I proceeded with the blessing. Before I could stow my gear, Leland Merrill was sleeping like a child.

The next morning, Merrill smilingly turned to me and said, “Monson, I’m glad you hold the priesthood!” His gladness was only surpassed by my gratitude—gratitude not only for the priesthood but for being worthy to receive the help I required in a time of desperate need and to exercise the power of the priesthood.

Brethren, our Lord and Savior said, “Come, follow me.” When we accept His invitation and walk in His footsteps, He will direct our paths.

In April of 2000, I felt such direction. I had received a phone call from Rosa Salas Gifford, whom I did not know. She explained that her parents had been visiting from Costa Rica for a few months and that just a week prior to her call, her father, Bernardo Agusto Salas, had been diagnosed with liver cancer. She indicated that the doctors had informed the family that her father would live just a few more days. Her father’s great desire, she explained, was to meet me before he died. She left her address and asked if I could come to her home in Salt Lake City to visit with her father.

Because of meetings and obligations, it was rather late when I left my office. Instead of going straight home, however, I felt impressed that I should drive further south and visit Brother Salas that very evening. With the address in hand, I attempted to locate the residence. In rather heavy traffic and with dimming light, I drove past the location where the road to the house should have been. I could see nothing. However, I don’t give up easily. I drove around the block and came back. Still nothing. One more time I tried and still no sign of the road. I began to feel that I would be justified in turning toward home. I had made a gallant effort but had been unsuccessful in finding the address. Instead, I offered a silent prayer for help. The inspiration came that I should approach the area from the opposite direction. I drove a distance and turned the car around so that I was now on the other side of the road. Going in this direction, the traffic was much lighter. As I neared the location once again, I could see, through the faint light, a street sign that had been knocked down—it was lying on its side at the edge of the road—and a nearly invisible, weed-covered track leading to a small apartment building and a single, tiny residence some distance from the main road. As I drove toward the buildings, a small girl in a white dress waved to me, and I knew that I had found the family.

I was ushered into the home and then to the room where Brother Salas lay. Surrounding the bed were three daughters and a son-in-law, as well as Sister Salas. All but the son-in-law were from Costa Rica. Brother Salas’s appearance reflected the gravity of his condition. A damp rag with frayed edges—not a towel or a washcloth but a damp rag with frayed edges—rested upon his forehead, emphasizing the humble economic circumstances of the family.

With some prompting, Brother Salas opened his eyes, and a wan smile graced his lips as I took him by the hand. I spoke the words, “I have come to meet you.” Tears welled up in his eyes and in mine.

I asked if a blessing would be desired, and the unanimous answer from the family members was affirmative. Since the son-in-law did not hold the priesthood, I proceeded by myself to provide a priesthood blessing. The words seemed to flow freely under the direction of the Spirit of the Lord. I included the Savior’s words found in the Doctrine and Covenants, section 84, verse 88: “I will go before your face. I will be on your right hand and on your left, and my Spirit shall be in your hearts, and mine angels round about you, to bear you up.” Following the blessing, I offered a few words of comfort to the grieving family members. I spoke carefully so they could understand my English. And then, with my limited Spanish language ability, I let them know that I loved them and that our Heavenly Father would bless them.

I asked for the family Bible and directed their attention to 3 John, verse 4: “I have no greater joy than to hear that my children walk in truth.” I said to them, “This is what your husband and father would have you remember as he prepares to depart this earthly existence.”

With tears streaming down her face, Brother Salas’s sweet wife then asked if I would write down the references for the two scriptures I had shared with them so that the family might read them again. Not having anything handy on which I could write, Sister Salas reached into her purse and drew from it a slip of paper. As I took it from her, I noticed it was a tithing receipt. My heart was touched as I realized that, despite the extremely humble circumstances in which the family lived, they were faithful in paying their tithes.

After a tender farewell, I was escorted to my car. As I drove homeward, I reflected on the special spirit we had felt. I experienced, as well, as I have many times before, a sense of gratitude that my Heavenly Father had answered another person’s prayer through me.

My brethren, let us ever remember that the priesthood of God which we bear is a sacred gift which brings to us and to those we serve the blessings of heaven. May we, in whatever place we may be, honor and protect that priesthood. May we ever be on the Lord’s errand, that we might ever be entitled to the Lord’s help.

There is a war being waged for men’s souls—yours and mine. It continues without abatement. Like a clarion call comes the word of the Lord to you and to me and to priesthood holders everywhere: “Wherefore, now let every man learn his duty, and to act in the office in which he is appointed, in all diligence.”

May we each have the courage to do so, I pray in the name of Jesus Christ, amen.

\newpage
\settalk{A Royal Priesthood}
\setspeaker{Thomas S. Monson}
\settalkdate{October 2007}
\talkheading{A Royal Priesthood}{Thomas S. Monson}

Brethren, as I gaze from one end to the other of this majestic building, I can only say, you are an inspiring sight to behold. It is amazing to realize that in thousands of chapels throughout the world, others of you—holders of the priesthood of God—are receiving this broadcast by way of satellite transmission. Nationalities vary and languages are many, but a common thread binds us together. We have been entrusted to bear the priesthood and to act in the name of God. We are the recipients of a sacred trust. Much is expected of us.

We who hold the priesthood of God and honor it are among those who have been reserved for this special period in history. The Apostle Peter described us in the second chapter of 1 Peter, the ninth verse: “Ye are a chosen generation, a royal priesthood, an holy nation, a peculiar people; that ye should shew forth the praises of him who hath called you out of darkness into his marvellous light.”

How might you and I qualify ourselves to be worthy of that designation, “a royal priesthood”? What are the characteristics of a true son of the living God? Tonight I would like us to consider just some of those very characteristics.

Times may change, circumstances may alter, but the marks of a true holder of the priesthood of God remain constant.

May I suggest that first of all every one of us develop the mark of vision. One writer said that the door of history turns on small hinges, and so do people’s lives. If we were to apply that maxim to our lives, we could say that we are the result of many small decisions. In effect, we are the product of our choices. We must develop the capacity to recall the past, to evaluate the present, and to look into the future in order to accomplish in our lives what the Lord would have us do.

You young men holding the Aaronic Priesthood should have the ability to envision the day when you will hold the Melchizedek Priesthood and then prepare yourselves as deacons, as teachers, as priests to receive the holy Melchizedek Priesthood of God. You have the responsibility to be ready, when you receive the Melchizedek Priesthood, to respond to a call to serve as a missionary by accepting it and then fulfilling it. How I pray that every boy and every man will have the mark of vision.

The second principle I should like to emphasize as a characteristic of a true priesthood holder of God is the mark of effort. It is not enough to want to make the effort and to say we’ll make the effort. We must actually make the effort. It’s in the doing, not just the thinking, that we accomplish our goals. If we constantly put our goals off, we will never see them fulfilled. Someone put it this way: Live only for tomorrow, and you will have a lot of empty yesterdays today.

In July of 1976, runner Garry Bjorklund was determined to qualify for the U.S. Olympic team’s 10,000-meter race which would be run at the Montreal Olympics. Halfway through the grinding qualifying race, however, he lost his left shoe. What would you and I do if that were our experience? I suppose he could have given up and stopped. He could have blamed his bad luck and lost the opportunity of participating in the greatest race of his life, but this champion athlete did not do that. He ran on without his shoe. He knew that he would have to run faster than he had ever run in his life. He knew that his competitors now had an advantage that they did not have at the beginning of the race. Over that cinder track he ran, with one shoe on and one shoe off, finishing third and qualifying for the opportunity to participate in the race for the gold medal. His own running time was the best he had ever recorded. He put forth the effort necessary to achieve his goal.

As priesthood holders we may find that there are times in our lives when we falter, when we become weary or fatigued, or when we suffer a disappointment or a heartache. When that happens, I would hope that we will persevere with even greater effort toward our goal.

At one time or another each of us will be called to fill a position in the Church, whether as a deacons quorum president, a teachers quorum secretary, a priesthood adviser, a class teacher, a bishop. I could name more, but you get the picture. I was just 22 years of age when I was called to be the bishop of the Sixth-Seventh Ward in Salt Lake City. With 1,080 members in the ward, a great deal of effort was required to make certain that every matter which needed to be handled was taken care of and every member of the ward felt included and watched over. Although the assignment was monumental in scope, I did not let it overwhelm me. I went to work, as did others, and did all I could to serve. Each of us can do the same, regardless of the calling or assignment.

Just last year I decided to see how many residential dwellings were still standing from the period between 1950 and 1955 when I served as bishop of that same area. I drove slowly around each of the blocks that once comprised the ward. I was surprised to observe in my search that of all the houses and apartment buildings where our 1,080 members had lived, only three dwellings were still standing. At one of those houses, the grass was overgrown, the trees unpruned, and I found no one was living there. Of the other two houses remaining, one was boarded up and unoccupied, and the other housed some sort of a modest business office.

I parked my car, turned off the ignition, and just sat there for a long while. I could picture in my mind each house, each apartment building, each member who lived there. While the homes and the buildings were gone, the memories were still very vivid concerning the families who resided in each dwelling. I thought of the words of the author James Barrie, who wrote that God gave us memories that we might have June roses in the December of our lives. How grateful I was for the opportunity to serve in that assignment. Such can be the blessing of each of us if we put forth in our assignments our very best efforts.

The mark of effort is required of every priesthood holder.

The third principle I would like to emphasize is the mark of faith. We must have faith in ourselves, faith in the ability of our Heavenly Father to bless us and to guide us in our endeavors. Many years ago the writer of a psalm wrote a beautiful truth: “It is better to trust in the Lord than to put confidence in man. It is better to trust in the Lord than to put confidence in princes.” In other words, let us put our confidence in the ability of the Lord to guide us. Friendships, we know, may alter and change, but the Lord is constant.

Shakespeare, in his play King Henry the Eighth, taught this truth through Cardinal Wolsey—a man who enjoyed great prestige and pride because of his friendship with the king. When the friendship ended, Cardinal Wolsey was stripped of his authority, resulting in a loss of prominence and prestige. He was one who had gained everything and then lost all. In the sorrow of his heart, he spoke a real truth to his servant, Cromwell. He said:

\begin{verse}
O Cromwell, Cromwell!\\
Had I but served my God with half the zeal\\
I served my king, He would not in mine age\\
Have left me naked to mine enemies.
\end{verse}

I trust we shall have the mark of faith in every heart represented here tonight.

I add to my list the mark of virtue. The Lord indicated that we should let virtue garnish our thoughts unceasingly.

I recall a priesthood meeting held in the Tabernacle in Salt Lake City when I was a holder of the Aaronic Priesthood. The President of the Church was speaking to the priesthood, and he made a statement I have never forgotten. He said, in essence, that men who commit sexual sin or other sins do not do so in the twinkling of an eye. He emphasized that our actions are preceded by our thoughts, and when we commit sin, it is because we have first thought of committing that particular sin. Then the President declared that the way to avoid sin is to keep our thinking pure. The scripture tells us that as we think in our hearts, so are we. We must have the mark of virtue.

If we are to be missionaries in the kingdom of our Heavenly Father, we must be entitled to the companionship of His Holy Spirit, and we have been told precisely that His Spirit will not dwell in impure or unholy tabernacles.

Finally, may I add the mark of prayer. The desire to communicate with one’s Heavenly Father is a mark of a true priesthood holder of God.

As we offer unto the Lord our family and our personal prayers, let us do so with faith and trust in Him. Let us remember the injunction of the Apostle Paul to the Hebrews: “For he that cometh to God must believe that he is, and that he is a rewarder of them that diligently seek him.” If any of us has been slow to hearken to the counsel to pray always, there is no finer hour to begin than now. William Cowper declared, “Satan trembles when he sees the weakest saint upon his knees.” Those who feel that prayer might denote a physical weakness should consider that a man never stands taller than when he is upon his knees.

May we ever remember:

\begin{verse}
Prayer is the soul’s sincere desire,\\
Uttered or unexpressed,\\
The motion of a hidden fire\\
That trembles in the breast. …
\end{verse}

As we cultivate the mark of prayer, we will receive the blessings our Heavenly Father has for us.

In conclusion, may we have vision. May we put forth effort. May we exemplify faith and virtue and ever make prayer a part of our lives. Then we shall indeed be a royal priesthood. This would be my prayer, my personal prayer this evening, and I offer it from my heart in the name of Jesus Christ, amen.

\newpage
\settalk{Mrs. Patton—the Story Continues}
\setspeaker{Thomas S. Monson}
\settalkdate{October 2007}
\talkheading{Mrs. Patton—the Story Continues}{Thomas S. Monson}

I miss my colleague James E. Faust today and express my love to his dear wife and family and am assured he is serving the Lord elsewhere. I welcome the newly sustained General Authorities, President Eyring, Elder Cook, and Elder González, and assure them they have my full support.

Thirty-eight years ago, at a general conference held in the Tabernacle on Temple Square, I spoke of one of my childhood friends, Arthur Patton, who died at a young age. The talk was titled “Mrs. Patton, Arthur Lives.” I addressed my remarks to Arthur’s mother, Mrs. Patton, who was not a member of the Church. Although I had little hope that Mrs. Patton would actually hear my talk, I wanted to share with all who were within the sound of my voice the glorious gospel message of hope and love. Recently I have felt impressed to refer once again to Arthur and to relate to you what transpired following my original message.

First may I tell you about Arthur. He had blond, curly hair and a smile as big as all outdoors. He stood taller than any boy in the class. I suppose this is how, in 1940, as the great conflict which became World War II was overtaking much of Europe, Arthur was able to fool the recruiting officers and enlist in the navy at the tender age of 15. To Arthur and most of the boys, the war was a great adventure. I remember how striking he appeared in his navy uniform. How we wished we were older or at least taller so we too could enlist.

Youth is a very special time of life. As Longfellow wrote:

\begin{verse}
How beautiful is youth! how bright it gleams\\
With its illusions, aspirations, dreams!\\
Book of Beginnings, Story without End,\\
Each maid a heroine, and each man a friend!
\end{verse}

Arthur’s mother was so proud of the blue star which graced her living room window. It represented to every passerby that her son wore the uniform of his country and was actively serving. When I would pass the house, she often opened the door and invited me in to read the latest letter from Arthur. Her eyes would fill with tears; I would then be asked to read aloud. Arthur meant everything to his widowed mother.

I can still picture Mrs. Patton’s coarse hands as she would carefully replace the letter in its envelope. These were hardworking hands; Mrs. Patton was a cleaning woman for a downtown office building. Each day of her life except Sundays she could be seen walking along the sidewalk, pail and brush in hand, her gray hair pulled back into a tight bob, her shoulders weary from work and stooped with age.

In March 1944, with the war now raging, Arthur was transferred from the USS Dorsey, a destroyer, to the USS White Plains, an aircraft carrier. While at Saipan in the South Pacific, the ship was attacked. Arthur was one of those on board who was lost at sea.

The blue star was taken from its hallowed spot in the front window of the Patton home. It was replaced by one of gold, indicating that he whom the blue star represented had been killed in battle. A light went out in the life of Mrs. Patton. She groped in utter darkness and deep despair.

With a prayer in my heart, I approached the familiar walkway to the Patton home, wondering what words of comfort could come from the lips of a mere boy.

The door opened, and Mrs. Patton embraced me as she would her own son. Home became a chapel as a grief-stricken mother and a less-than-adequate boy knelt in prayer.

As we arose from our knees, Mrs. Patton gazed into my eyes and spoke: “Tommy, I belong to no church, but you do. Tell me, will Arthur live again?” To the best of my ability, I testified to her that Arthur would indeed live again.

In general conference those long years ago, as I related this account, I mentioned that I had lost track of Mrs. Patton but that I wanted to once more answer her question “Will Arthur live again?”

I referred to the Savior of the world, who walked the dusty paths of villages we now reverently call the Holy Land; who caused the blind to see, the deaf to hear, the lame to walk, and the dead to live; to Him who tenderly and lovingly assured us, “I am the way, the truth, and the life.”

I explained that the plan of life and an explanation of its eternal course come to us from the Master of heaven and earth, even Jesus Christ the Lord. To understand the meaning of death, we must appreciate the purpose of life.

I indicated that in this dispensation the Lord declared: “And now, verily I say unto you, I was in the beginning with the Father, and am the Firstborn.” “Man was also in the beginning with God.”

Jeremiah the prophet recorded:

“The word of the Lord came unto me, saying,

“Before I formed thee … I knew thee; and before thou camest forth … I sanctified thee, and I ordained thee a prophet unto the nations.”

From that majestic world of spirits we enter the grand stage of life to prove ourselves obedient to all things commanded of God. During mortality we grow from helpless infancy to inquiring childhood and then to reflective maturity. We experience joy and sorrow, fulfillment and disappointment, success and failure. We taste the sweet, yet sample the bitter. This is mortality.

Then to each life comes the experience known as death. None is exempt. All must pass its portals.

To most, there is something sinister and mysterious about this unwelcome visitor called death. Perhaps it is a fear of the unknown which causes many to dread its coming.

Arthur Patton died quickly. Others linger. We know, through the revealed word of God, that “the spirits of all men, as soon as they are departed from this mortal body, … are taken home to that God who gave them life.”

I assured Mrs. Patton and all others who were listening that God would never forsake them—that He sent His Only Begotten Son into the world to teach us by example the life we should live. His Son died upon the cross to redeem all mankind. His words to the grieving Martha and to His disciples today bring comfort to us:

“I am the resurrection, and the life: he that believeth in me, though he were dead, yet shall he live:

“And whosoever liveth and believeth in me shall never die.”

“In my Father’s house are many mansions: if it were not so, I would have told you. I go to prepare a place for you.

“… I will come again, and receive you unto myself; that where I am, there ye may be also.”

I reiterated the testimonies of John the Revelator and Paul the Apostle. John recorded:

“I saw the dead, small and great, stand before God; …

“And the sea gave up the dead which were in it.”

Paul declared, “As in Adam all die, even so in Christ shall all be made alive.”

I explained that until the glorious Resurrection morning, we walk by faith. “For now we see through a glass, darkly; but then face to face.”

I reassured Mrs. Patton that Jesus invited her and all others:

“Come unto me, all ye that labour and are heavy laden, and I will give you rest.

“Take my yoke upon you, and learn of me; for I am meek and lowly in heart: and ye shall find rest unto your souls.”

As part of my message, I explained to Mrs. Patton that such knowledge would sustain her in her heartache—that she would never be in the tragic situation of the disbeliever who, having lost a son, was heard to say as she watched the casket lowered into mother earth, “Good-bye, my boy. Good-bye forever.” Rather, with head erect, courage undaunted, and faith unwavering, she could lift her eyes as she looked beyond the gently breaking waves of the blue Pacific and whisper, “Good-bye, Arthur, my precious son. Good-bye—until we meet again.”

I quoted the words of Tennyson, as though spoken to her by Arthur:

\begin{verse}
Sunset and evening star,\\
And one clear call for me!\\
And may there be no moaning of the bar,\\
When I put out to sea, …
\end{verse}

As I concluded my message those long years ago, I expressed to Mrs. Patton my personal testimony as a special witness, telling her that God our Father was mindful of her—that through sincere prayer she could communicate with Him; that He too had a Son who died, even Jesus Christ the Lord; that Christ is our advocate with the Father, the Prince of Peace, our Savior and divine Redeemer; and that one day we would see Him face-to-face.

I hoped that my message to Mrs. Patton would reach and touch others who had lost a loved one.

And now, my brothers and sisters, I share with you the rest of this account. I delivered my message on April 6, 1969. Again, I had little or no hope that Mrs. Patton would actually hear the talk. I had no reason to think she would listen to general conference. As I have mentioned, she was not a member of the Church. And then I learned that something akin to a miracle had taken place. Having no idea whatsoever who would be speaking at conference or what subjects they might speak about, Latter-day Saint neighbors of Mrs. Terese Patton in California, where she had moved, invited her to their home to listen to a session of conference with them. She accepted their invitation and thus was listening to the very session where I directed my remarks to her personally.

During the first week of May 1969, to my astonishment and joy, I received a letter postmarked Pomona, California, and dated April 29, 1969. It was from Mrs. Terese Patton. I share with you a part of that letter:

“Dear Tommy,

“I hope you don’t mind my calling you Tommy, as I always think of you that way. I don’t know how to thank you for the comforting talk you gave.

“Arthur was 15 years old when he enlisted in the navy. He was killed one month before his 19th birthday on July 5, 1944.

“It was wonderful of you to think of us. I don’t know how to thank you for your comforting words, both when Arthur died and again in your talk. I have had many questions over the years, and you have answered them. I am now at peace concerning Arthur. … God bless and keep you always.

“Love,

“Terese Patton”

My brothers and sisters, I do not believe it was a coincidence that I was impressed to give that particular message at the April 1969 general conference. Nor do I believe it was a coincidence that Mrs. Terese Patton was invited by neighbors to join them in their home for that particular session of conference. I am certain our Heavenly Father was mindful of her needs and wanted her to hear the comforting truths of the gospel.

Although Mrs. Patton has long since left mortality, I have felt a strong impression to share with you the manner in which our Heavenly Father blessed and provided for her, a widow, in her need. With all the strength of my soul I testify that our Heavenly Father loves each one of us. He hears the prayers of humble hearts; He hears our cries for help, as He heard Mrs. Patton. His Son, our Savior and Redeemer, speaks to each of us today: “Behold, I stand at the door, and knock: if any man hear my voice, and open the door, I will come in to him.”

Will we listen for that knock? Will we hear that voice? Will we open that door to the Lord, that we may receive the help He is so ready to provide? I pray that we will, in the sacred name of Jesus Christ, amen.

\newpage
\settalk{Three Goals to Guide You}
\setspeaker{Thomas S. Monson}
\settalkdate{October 2007}
\talkheading{Three Goals to Guide You}{Thomas S. Monson}

This evening our souls have reached toward heaven. We have been blessed with beautiful music and inspired messages. The Spirit of the Lord is here.

Sisters Julie Beck, Silvia Allred, Barbara Thompson—thank heaven for your dear mothers and fathers, your teachers, your youth leaders, and others who recognized in you your potential.

To paraphrase a thought:

\begin{verse}
You never know what a girl is worth,\\
You’ll have to wait and see;\\
But every woman in a noble place,\\
A girl once used to be.
\end{verse}

It is a great privilege for me to be with you. I recognize that beyond you who are gathered in the Conference Center, there are many thousands watching and listening to the proceedings by way of satellite transmission.

As I speak to you, I realize that as a man I am in the minority and must be cautious in my comments. I’m reminded of the man who walked into a bookstore and asked the clerk—a woman—for help: “Have you got a book titled Man, the Master of Women?” The clerk looked him straight in the eye and said sarcastically, “Try the fiction section!”

I assure you tonight that I honor you, the women of the Church, and am well aware, to quote William R. Wallace, that “the hand that rocks the cradle is the hand that rules the world.”

In 1901 President Lorenzo Snow said: “The members of the Relief Society have … ministered to those in affliction, they have thrown their arms of love around the fatherless and the widows, and they have kept themselves unspotted from the world. I can testify that there are no purer and more God-fearing women in the world than are to be found within the ranks of the Relief Society.”

As in President Snow’s time, there are, here and now, visits to be made, greetings to be shared, and hungry souls to be fed. As I contemplate the Relief Society of today, humbled by my privilege to speak to you, I turn to our Heavenly Father for His divine guidance.

In this spirit, I have felt to provide each member of the Relief Society throughout the world three goals to meet:

Let us consider each of these goals. First, study diligently. The Savior of the world instructed: “Seek ye out of the best books words of wisdom; seek learning, even by study and also by faith.” He added: “Search the scriptures; for in them ye think ye have eternal life: and they are they which testify of me.”

A study of the scriptures will help our testimonies and the testimonies of our family members. Our children today are growing up surrounded by voices urging them to abandon that which is right and to pursue, instead, the pleasures of the world. Unless they have a firm foundation in the gospel of Jesus Christ, a testimony of the truth, and a determination to live righteously, they are susceptible to these influences. It is our responsibility to fortify and protect them.

To an alarming extent, our children today are being educated by the media, including the Internet. In the United States, it is reported that the average child watches approximately four hours of television daily, much of the programming being filled with violence, alcohol and drug use, and sexual content. Watching movies and playing video games is in addition to the four hours. And the statistics are much the same for other developed countries. The messages portrayed on television, in movies, and in other media are very often in direct opposition to that which we want our children to embrace and hold dear. It is our responsibility not only to teach them to be sound in spirit and doctrine but also to help them stay that way, regardless of the outside forces they may encounter. This will require much time and effort on our part—and in order to help others, we ourselves need the spiritual and moral courage to withstand the evil we see on every side.

We live in the time spoken of in 2 Nephi, chapter 9:

“O the vainness, and the frailties, and the foolishness of men! When they are learned they think they are wise, and they hearken not unto the counsel of God, for they set it aside, supposing they know of themselves, wherefore, their wisdom is foolishness and it profiteth them not. And they shall perish.

“But to be learned is good if they hearken unto the counsels of God.”

Required is the courage to hold fast to our standards despite the derision of the world. Said President J. Reuben Clark Jr., for many years a member of the First Presidency: “Not unknown are cases where [those] of presumed faith … have felt that, since by affirming their full faith they might call down upon themselves the ridicule of their unbelieving colleagues, they must either modify or explain away their faith or destructively dilute it, or even pretend to cast it away. Such are hypocrites.”

There comes to mind the powerful verses found in 2 Timothy, in the New Testament, chapter 1, verses 7 and 8:

“For God hath not given us the spirit of fear; but of power, and of love, and of a sound mind.

“Be not thou therefore ashamed of the testimony of our Lord.”

Beyond our study of spiritual matters, secular learning is also essential. Often the future is unknown; therefore, it behooves us to prepare for uncertainties. Statistics reveal that at some time, because of the illness or death of a husband or because of economic necessity, you may find yourself in the role of financial provider. Some of you already occupy that role. I urge you to pursue your education—if you are not already doing so or have not done so—that you might be prepared to provide if circumstances necessitate such.

Your talents will expand as you study and learn. You will be able to better assist your families in their learning, and you will have peace of mind in knowing that you have prepared yourself for the eventualities that you may encounter in life.

I reiterate: Study diligently.

The second goal I wish to mention: Pray earnestly. The Lord directed, “Pray always, and I will pour out my Spirit upon you, and great shall be your blessing.”

Perhaps there has never been a time when we had greater need to pray and to teach our family members to pray. Prayer is a defense against temptation. It is through earnest and heartfelt prayer that we can receive the needed blessings and the support required to make our way in this sometimes difficult and challenging journey we call mortality.

We can teach the importance of prayer to our children and grandchildren both by word and by example. I share with you a lesson in teaching by example as described in a mother’s letter to me relating to prayer. “Dear President Monson: Sometimes I wonder if I make a difference in my children’s lives. Especially as a single mother working two jobs to make ends meet, I sometimes come home to confusion, but I never give up hope.”

Her letter continues as she describes how she and her children were watching general conference, where I was speaking about prayer. Her son made the comment, “Mother, you’ve already taught us that.” She asked, “What do you mean?” Her son replied, “Well, you’ve taught us to pray and showed us how, but the other night I came to your room to ask something and found you on your knees praying to Heavenly Father. If He’s important to you, He’ll be important to me.” The letter concluded, “I guess you never know what kind of influence you’ll be until a child observes you doing yourself what you have tried to teach him to do.”

Some years ago, just before leaving Salt Lake to attend the annual meetings of Boy Scouts of America in Atlanta, Georgia, I decided to take with me enough copies of the New Era so that I might share with Scouting officials this excellent publication. When I arrived at the hotel in Atlanta, I opened the package of magazines. I found that my secretary, for no accountable reason, had put in the package two extra copies of the June issue, an issue that featured temple marriage. I left the two copies in the hotel room and, as planned, distributed the other copies.

On the final day of meetings, I had no desire to attend the scheduled luncheon but felt compelled to return to my room. The telephone was ringing as I entered. The caller was a member of the Church who had heard I was in Atlanta. She introduced herself and asked if I could provide a blessing for her 10-year-old daughter. I agreed readily, and she indicated that she and her husband, their daughter, and their son would come immediately to my hotel room. As I waited, I prayed for help. The applause of the convention was replaced by the feelings of peace which accompanied prayer.

Then came the knock at the door and the privilege which was mine to meet a choice family. The 10-year-old daughter walked with the aid of crutches. Cancer had required the amputation of her left leg; however, her countenance was radiant, her trust in God unwavering. A blessing was provided. Mother and son knelt by the side of the bed while the father and I placed our hands on the tiny daughter. We were directed by the Spirit of God. We were humbled by its power.

I felt the tears course down my cheeks and tumble upon my hands as they rested on the head of that beautiful child of God. I spoke of eternal ordinances and family exaltation. The Lord prompted me to urge this family to enter the holy temple of God. At the conclusion of the blessing, I learned that such a temple visit was planned. Questions pertaining to the temple were asked. I heard no heavenly voice, nor did I see a vision. Yet there came clearly into my mind the words, “Refer to the New Era.” I looked toward the dresser, and there were the two extra copies of the temple issue of the New Era. One copy was given to the daughter and the other to her parents. We reviewed them together.

The family said farewell, and once again the room was still. A prayer of gratitude came easily and, once more, the resolve to ever provide a place for prayer.

My dear sisters, do not pray for tasks equal to your abilities, but pray for abilities equal to your tasks. Then the performance of your tasks will be no miracle, but you will be the miracle.

Pray earnestly.

Finally, serve willingly. You are a mighty force for good, one of the most powerful in the entire world. Your influence ranges far beyond yourself and your home and touches others all around the globe. You have reached out to your brothers and sisters across streets, across cities, across nations, across continents, across oceans. You personify the Relief Society motto: “Charity never faileth.”

You are, of course, surrounded by opportunities for service. No doubt at times you recognize so many such opportunities that you may feel somewhat overwhelmed. Where do you begin? How can you do it all? How do you choose, from all the needs you observe, where and how to serve?

Often small acts of service are all that is required to lift and bless another: a question concerning a person’s family, quick words of encouragement, a sincere compliment, a small note of thanks, a brief telephone call. If we are observant and aware, and if we act on the promptings which come to us, we can accomplish much good. Sometimes, of course, more is needed.

I learned recently of loving service given to a mother when her children were very young. Frequently she would be up in the middle of the night tending to the needs of her little ones, as mothers do. Often her friend and neighbor across the street would come over the next day and say, “I saw your lights on in the middle of the night and know you were up with the children. I’m going to take them to my house for a couple of hours while you take a nap.” Said this grateful mother: “I was so thankful for her welcome offer, it wasn’t until this had happened many times that I realized if she had seen my lights on in the middle of the night, she was up with one of her children as well and needed a nap just as much as I did. She taught me a great lesson, and I’ve since tried to be as observant as she was in looking for opportunities to serve others.”

Countless are the acts of service provided by the vast army of Relief Society visiting teachers. A few years ago I heard of two of them who aided a grieving widow, Angela, the granddaughter of a cousin of mine. Angela’s husband and a friend of his had gone snowmobiling and had become victims of suffocation through a snowslide. Each of them left a pregnant wife—in Angela’s case, their first child, and in the other case, a wife not only expecting a child but also the mother of a toddler. In the funeral held for Angela’s husband, the bishop reported that upon hearing of the tragic accident, he had gone immediately to Angela’s home. Almost as soon as he arrived, the doorbell sounded. The door was opened, and there stood Angela’s two visiting teachers. The bishop said he watched as they so sincerely expressed to Angela their love and compassion. The three women cried together, and it was apparent that these two fine visiting teachers cared deeply about Angela. As perhaps only women can, they gently indicated—without being asked—exactly what help they would be providing. That they would be close by as long as Angela needed them was obvious. The bishop expressed his deep gratitude in knowing they would be a real source of comfort to her in the days ahead.

Such acts of love and compassion are repeated again and again by the wonderful visiting teachers of this Church—not always in such dramatic situations but just as genuinely, nevertheless.

I extol you who, with loving care and compassionate concern, feed the hungry, clothe the naked, and house the homeless. He who notes the sparrow’s fall will not be unmindful of such service. The desire to lift, the willingness to help, and the graciousness to give come from a heart filled with love. Serve willingly.

Our beloved prophet, even President Gordon B. Hinckley, said of you, “God planted within women something divine that expresses itself in quiet strength, in refinement, in peace, in goodness, in virtue, in truth, in love.”

My dear sisters, may our Heavenly Father bless each of you, married or single, in your homes, in your families, in your very lives—that you may merit the glorious salutation of the Savior of the World: “Well done, thou good and faithful servant” I pray, as I bless you and also the dear wife of James E. Faust, his beloved Ruth, who is here tonight on the front row, and their family, in the name of Jesus Christ, amen.

\newpage
\settalk{Abundantly Blessed}
\setspeaker{Thomas S. Monson}
\settalkdate{April 2008}
\talkheading{Abundantly Blessed}{Thomas S. Monson}

I’ve been attending conference for a long time, but I think I’ve never felt quite as richly blessed as during this session. We’ve had rapid-fire messages from a lot of speakers, but every one touched on a very important subject. We’ve had a smorgasbord today of faith, of love, and of counsel. Let’s incorporate these things in our lives.

Brother Ballard, several years ago my dear wife went to the hospital. She left a note behind for the children: “Dear children, do not let Daddy touch the microwave”—followed by a comma, “or the stove, or the dishwasher, or the dryer.” I’m embarrassed to add any more to that list.

I think it was Brother Uchtdorf who said, “You told the audience today about your heritage on your mother’s side. What about your father’s side?” So I conclude with just a word or two about my father’s side.

My father’s father came from Sweden, and his wife from England. They met on the ship coming over. He waited for her to grow up, and then he proposed marriage. They were married in the Salt Lake Temple, and he wrote in his journal, “Today is the happiest day of my life. My sweetheart and I were married for time and eternity in the holy temple.”

Three days later, on April 23, 1898, he wrote, “Took the train at the Rio Grande Western Depot enroute eventually to Scandinavia, where I have been called as a missionary.” Off he went to Sweden, leaving his bride of three days.

His journal, written in pencil, came to me from an uncle who somehow chose me to receive his father’s journal. The most frequent entry in the journal was, “My feet are wet.” But the most beautiful entry said: “Today we went to the Jansson home. We met Sister Jansson. She had a lovely dinner for us. She is a good cook.” And then he said, “The children all sang or played a harmonica or did a little dance, and then she paid her tithing. Five krona for the Lord and one for my companion, Elder Ipson, and one for me.” And then there were listed the names of the children.

When I read that in the journal, there was my wife’s father’s name as one who was in that household, one who probably sang a song, one who became the father of only one daughter, the girl whom I married.

The first day I saw Frances, I knew I’d found the right one. The Lord brought us together later, and I asked her to go out with me. I went to her home to call on her. She introduced me, and her father said, “‘Monson’—that’s a Swedish name, isn’t it?”

I said, “Yes.”

He said, “Good.”

Then he went into another room and brought out a picture of two missionaries with their top hats and their copies of the Book of Mormon.

“Are you related to this Monson,” he said, “Elias Monson?”

I said, “Yes, he’s my grandfather’s brother. He too was a missionary in Sweden.”

Her father wept. He wept easily. He said, “He and his companion were the missionaries who taught the gospel to my mother and my father and all of my brothers and sisters and to me.” He kissed me on the cheek. And then her mother cried, and she kissed me on the other cheek. And then I looked around for Frances. She said, “I’ll go get my coat.”

My sweet Frances had a terrible fall a few years ago. She went to the hospital. She lay in a coma for about 18 days. I sat by her side. She never moved a muscle. The children cried, the grandchildren cried, and I wept. Not a movement.

And then one day she opened her eyes. I set a speed record in getting to her side. I gave her a kiss and a hug, and I said, “You’re back. I love you.” And she said, “I love you too, Tom, but we’re in serious trouble.” I thought, What do you know about trouble, Frances? She said, “I forgot to mail in our fourth-quarter income tax payment.”

I said to her, “Frances, if you had said that before you extended a kiss to me and told me you love me, I might have left you here.”

Brethren, let’s treat our wives with dignity and with respect. They’re our eternal companions. Sisters, honor your husbands. They need to hear a good word. They need a friendly smile. They need a warm expression of true love.

Leaving my own family for a moment, my brothers and sisters, this has been a wonderful conference. We have been edified by wise and inspired messages. Our testimonies have been strengthened. I believe we are all the more determined to live the principles of the gospel of Jesus Christ.

Not only have we been blessed by the fine talks which have been given; we have also been uplifted by the beautiful music which has been provided. We are abundantly blessed in the Church by those who share their musical talents with us. Every choir and chorus has performed so well during the past two days.

I express my great love for all those who have participated and to all of you who have listened. I have felt your prayers in my behalf and have been sustained and blessed during the two months since our beloved President Hinckley left us. Once again, I appreciate your sustaining vote.

I cannot adequately express my gratitude for the Restoration of the gospel in these latter days and for what that has meant in my life. Each of us has been influenced and shaped as we have followed the Savior and have adhered to the principles of His gospel.

To you who are parents, I say, show love to your children. You know you love them, but make certain they know it as well. They are so precious. Let them know. Call upon our Heavenly Father for help as you care for their needs each day and as you deal with the challenges which inevitably come with parenthood. You need more than your own wisdom in rearing them.

We commend our wonderful young people who stand up to the iniquity in the world and who live the commandments to the best of their ability.

To you who are able to attend the temple, I would counsel you to go often. Doing so will help to strengthen marriages and families.

Let us be kind to one another, be aware of each other’s needs, and try to help in that regard.

My dear brothers and sisters, I love you, and I pray for you. Please pray for me. And together we will reap the blessings our Heavenly Father has in store for each one of us. This is my prayer, my plea as I add my testimony. This work is true. In the name of Jesus Christ, amen.

\newpage
\settalk{Examples of Righteousness}
\setspeaker{Thomas S. Monson}
\settalkdate{April 2008}
\talkheading{Examples of Righteousness}{Thomas S. Monson}

Tonight I am aware that you, my brethren, both here in the Conference Center and in thousands of other locations, represent the largest gathering of the priesthood ever to assemble. We are a part of the greatest brotherhood in all the world. How fortunate and blessed we are to be holders of the priesthood of God.

We have been instructed and uplifted as we have listened to inspired messages. I pray that I might have an interest in your faith and prayers as I share with you those thoughts and feelings that have been in my mind lately as I have prepared to address you.

As bearers of the priesthood, we have been placed on earth in troubled times. We live in a complex world with currents of conflict everywhere to be found. Political machinations ruin the stability of nations, despots grasp for power, and segments of society seem forever downtrodden, deprived of opportunity, and left with a feeling of failure.

We who have been ordained to the priesthood of God can make a difference. When we qualify for the help of the Lord, we can build boys, we can mend men, we can accomplish miracles in His holy service. Our opportunities are without limit.

Ours is the task to be fitting examples. We are strengthened by the truth that the greatest force in the world today is the power of God as it works through man. If we are on the Lord’s errand, brethren, we are entitled to the Lord’s help. Never forget that truth. That divine help, of course, is predicated upon our worthiness. Each must ask: Are my hands clean? Is my heart pure? Am I a worthy servant of the Lord?

We are surrounded by so much that is designed to divert our attention from those things which are virtuous and good and to tempt us with that which would cause us to be unworthy to exercise the priesthood we bear. I speak not just to the young men of the Aaronic Priesthood but to those of all ages. Temptations come in various forms throughout our lives.

Brethren, are we qualified at all times to perform the sacred duties associated with the priesthood we bear? Young men—you who are priests—are you clean in body and spirit as you sit at the sacrament table on Sunday and bless the emblems of the sacrament? Young men who are teachers, are you worthy to prepare the sacrament? Deacons, as you pass the sacrament to the members of the Church, do you do so knowing that you are spiritually qualified to do so? Does each of you fully understand the importance of all the sacred duties you perform?

My young friends, be strong. The philosophies of men surround us. The face of sin today often wears the mask of tolerance. Do not be deceived; behind that facade is heartache, unhappiness, and pain. You know what is right and what is wrong, and no disguise, however appealing, can change that. The character of transgression remains the same. If your so-called friends urge you to do anything you know to be wrong, you be the one to make a stand for right, even if you stand alone. Have the moral courage to be a light for others to follow. There is no friendship more valuable than your own clear conscience, your own moral cleanliness—and what a glorious feeling it is to know that you stand in your appointed place clean and with the confidence that you are worthy to do so.

Brethren of the Melchizedek Priesthood, do you strive diligently each day to live as you should? Are you kind and loving to your wife and your children? Are you honest in your dealings with those around you—at all times and in all circumstances?

If any of you has slipped along the way, there are those who will help you to once again become clean and worthy. Your bishop or branch president is anxious and willing to help and will, with understanding and compassion, do all within his power to assist you in the repentance process, that you may once again stand in righteousness before the Lord.

Many of you will remember President N. Eldon Tanner, who served as a counselor to four Presidents of the Church. He provided an undeviating example of righteousness throughout a career in industry, during service in the government in Canada, and consistently in his private life. He gave us this inspired counsel:

“Nothing will bring greater joy and success than to live according to the teachings of the gospel. Be an example; be an influence for good. …

“Every one of us has been foreordained for some work as [God’s] chosen servant on whom he has seen fit to confer the priesthood and power to act in his name. Always remember that people are looking to you for leadership and you are influencing the lives of individuals either for good or for bad, which influence will be felt for generations to come.”

My brethren, I reiterate that, as holders of the priesthood of God, it is our duty to live our lives in such a way that we may be examples of righteousness for others to follow. As I have pondered how we might best provide such examples, I have thought of an experience I had some years ago while attending a stake conference. During the general session, I observed a young boy sitting with his family on the front row of the stake center. I was seated on the stand. As the meeting progressed, I began to notice that if I crossed one leg over the other, the young boy would do the same thing. If I reversed the motion and crossed the other leg, he would follow suit. I would put my hands in my lap, and he would do the same. I rested my chin in my hand, and he also did so. Whatever I did, he would imitate my actions. This continued until the time approached for me to address the congregation. I decided to put him to the test. I looked squarely at him, certain I had his attention, and then I wiggled my ears. He made a vain attempt to do the same, but I had him! He just couldn’t quite get his ears to wiggle. He turned to his father, who was sitting next to him, and whispered something to him. He pointed to his ears and then to me. As his father looked in my direction, obviously to see my ears wiggle, I sat solemnly with my arms folded, not moving a muscle. The father glanced back skeptically at his son, who looked slightly defeated. He finally gave me a sheepish grin and shrugged his shoulders.

I have thought about that experience over the years as I’ve contemplated how, particularly when we’re young, we tend to imitate the example of our parents, our leaders, our peers. The prophet Brigham Young said: “We should never permit ourselves to do anything that we are not willing to see our children do. We should set them an example that we wish them to imitate.”

To you who are fathers of boys or who are leaders of boys, I say, strive to be the kind of example the boys need. The father, of course, should be the prime example, and the boy who is blessed with a worthy father is fortunate indeed. Even an exemplary family, however, with diligent and faithful father and mother, can use all the supportive help they can get from good men who genuinely care. There is also the boy who has no father or whose father is not currently providing the type of example needed. For that boy, the Lord has provided a network of helpers within the Church—bishops, advisers, teachers, Scoutmasters, home teachers. When the Lord’s program is in effect and properly working, no young man in the Church should be without the influence of good men in his life.

The effectiveness of an inspired bishop, adviser, or teacher has very little to do with the outward trappings of power or an abundance of this world’s goods. The leaders who have the most influence are usually those who set hearts afire with devotion to the truth, who make obedience to duty seem the essence of manhood, who transform some ordinary routine occurrence so that it becomes a vista where we see the person we aspire to be.

Not to be overlooked—and in fact our primary example—is our Savior, Jesus Christ. His birth was foretold by prophets; angels heralded the announcement of His earthly ministry. He “grew, and waxed strong in spirit, filled with wisdom: and the grace of God was upon him.”

Baptized of John in the river known as Jordan, He commenced His official ministry to men. To the sophistry of Satan, Jesus turned His back. To the duty designated by His Father, He turned His face, pledged His heart, and gave His life. And what a sinless, selfless, noble, and divine life it was. Jesus labored. Jesus loved. Jesus served. Jesus testified. What finer example could we strive to emulate? Let us begin now, this very night, to do so. Cast off forever will be the old self and with it defeat, despair, doubt, and disbelief. To a newness of life we come—a life of faith, hope, courage, and joy. No task looms too large; no responsibility weighs too heavily; no duty is a burden. All things become possible.

Many years ago I spoke of one who took his example from the Savior, one who stood firm and true, strong and worthy through the storms of life. He courageously magnified his priesthood callings. He provides an example to each of us. His name was Thomas Michael Wilson, the son of Willie and Julia Wilson of Lafayette, Alabama.

When he was but a teenager and he and his family were not yet members of the Church, he was stricken with cancer, followed by painful radiation therapy and then blessed remission. This illness caused his family to realize that not only is life precious but that it can also be short. They began to look to religion to help them through this time of tribulation. Subsequently, they were introduced to the Church, and eventually all but the father were baptized. After accepting the gospel, young Brother Wilson yearned for the opportunity of being a missionary, even though he was older than most young men when they begin their missionary service. At the age of 23, he received a mission call to serve in the Utah Salt Lake City Mission.

Elder Wilson’s missionary companions described his faith as unquestioning, undeviating, and unyielding. He was an example to all. However, after 11 months of missionary service, illness returned. Bone cancer now required the amputation of his arm and shoulder. Yet he persisted in his missionary labors.

Elder Wilson’s courage and consuming desire to remain on his mission so touched his nonmember father that he investigated the teachings of the Church and also became a member.

I learned that an investigator whom Elder Wilson had taught was baptized but then wanted to be confirmed by Elder Wilson, whom she respected so much. She, with a few others, journeyed to Elder Wilson’s bedside in the hospital. There, with his remaining hand resting upon her head, Elder Wilson confirmed her a member of The Church of Jesus Christ of Latter-day Saints.

Elder Wilson continued month after month his precious but painful service as a missionary. Blessings were given; prayers were offered. Because of his example of dedication, his fellow missionaries lived closer to God.

Elder Wilson’s physical condition deteriorated. The end drew near, and he was to return home. He asked to serve but one additional month, and his request was granted. He put his faith in God, and He whom Thomas Michael Wilson silently trusted opened the windows of heaven and abundantly blessed him. His parents, Willie and Julia Wilson, and his brother Tony came to Salt Lake City to help their son and brother home to Alabama. However, there was yet a prayed-for, a yearned-for blessing to be bestowed. The family invited me to come with them to the Jordan River Temple, where those sacred ordinances which bind families for eternity, as well as for time, were performed.

I said good-bye to the Wilson family. I can see Elder Wilson yet as he thanked me for being with him and his loved ones. He said, “It doesn’t matter what happens to us in this life as long as we have the gospel of Jesus Christ and live it. It doesn’t matter whether I teach the gospel on this or the other side of the veil, so long as I can teach it.” What courage. What confidence. What love. The Wilson family made the long trek home to Lafayette, where Elder Thomas Michael Wilson slipped from here to eternity. He was buried there with his missionary tag in place.

My brethren, as we now leave this general priesthood meeting, let us all determine to prepare for our time of opportunity and to honor the priesthood we bear through the service we render, the lives we bless, and the souls we are privileged to help save. You “are a chosen generation, a royal priesthood, an holy nation,” and you can make a difference. To these truths I testify in the name of Jesus Christ, our Savior, amen.

\newpage
\settalk{Looking Back and Moving Forward}
\setspeaker{Thomas S. Monson}
\settalkdate{April 2008}
\talkheading{Looking Back and Moving Forward}{Thomas S. Monson}

I think this has been a remarkable session. The messages have been inspiring; the music has been beautiful, the testimonies sincere. I think anyone who has attended this session will never forget it—for the Spirit we’ve felt.

My beloved brothers and sisters, over 44 years ago, in October of 1963, I stood at the pulpit in the Tabernacle, having just been sustained as a member of the Quorum of the Twelve Apostles. On that occasion I mentioned a small sign I had seen on another pulpit. The words on the sign were these: “Who stands at this pulpit, let him be humble.” I assure you that I was humbled by my call to the Twelve at that time. However, as I stand at this pulpit today, I address you from the absolute depths of humility. I feel very keenly my dependence upon the Lord. I humbly seek the guidance of the Spirit as I share with you the feelings of my heart.

Just two months ago we said farewell to our dear friend and leader Gordon B. Hinckley, the 15th President of The Church of Jesus Christ of Latter-day Saints, an outstanding ambassador of truth to the entire world and beloved of all. We miss him. More than 53,000 men, women, and children journeyed to the beautiful Hall of the Prophets in this very building to pay their last respects to this giant of the Lord, who now belongs to the ages.

With the passing of President Hinckley, the First Presidency was dissolved. President Eyring and I, who served as counselors to President Hinckley, returned to our places in the Quorum of the Twelve Apostles, and that quorum became the presiding authority of the Church.

On Saturday, February 2, 2008, funeral services for President Hinckley were held in this magnificent Conference Center—a building which will ever stand as a monument to his foresight and vision. During the funeral, beautiful and loving tributes were paid to this man of God.

The following day, all 14 ordained Apostles living on the earth assembled in an upper room of the Salt Lake Temple. We met in a spirit of fasting and prayer. During that solemn and sacred gathering, the Presidency of the Church was reorganized in accordance with well-established precedent, after the pattern which the Lord Himself put in place.

Members of the Church around the world convened yesterday in a solemn assembly. You raised your hands in a sustaining vote to approve the action which was taken in that meeting in the temple to which I have just referred. As your hands were raised toward heaven, my heart was touched. I felt your love and support, as well as your commitment to the Lord.

I know without question, my brothers and sisters, that God lives. I testify to you that this is His work. I testify as well that our Savior Jesus Christ is at the head of this Church, which bears His name. I know that the sweetest experience in all this life is to feel His promptings as He directs us in the furtherance of His work. I felt those promptings as a young bishop, guided to the homes where there was spiritual—or perhaps temporal—want. I felt them again as a mission president in Toronto, Canada, working with wonderful missionaries who were a living witness and testimony to the world that this work is divine and that we are led by a prophet. I have felt them throughout my service in the Twelve and in the First Presidency and now as President of the Church. I testify that each one of us can feel the Lord’s inspiration as we live worthily and strive to serve Him.

I am keenly aware of the 15 men who preceded me as President of the Church. Many of them I have known personally. I have had the blessing and privilege of serving as a counselor to three of them. I am grateful for the abiding legacy left by each one of those 15 men. I have the sure knowledge, as I am confident they had, that God directs His prophet. My earnest prayer is that I might continue to be a worthy instrument in His hands to carry on this great work and to fulfill the tremendous responsibilities which come with the office of President.

I thank the Lord for wonderful counselors. President Henry B. Eyring and President Dieter F. Uchtdorf are men of great ability and sound understanding. They are counselors in the true sense of the word. I value their judgment. I believe they have been prepared by the Lord for the positions they now occupy. I love the members of the Quorum of the Twelve Apostles and treasure my association with them. They too are dedicated to the work of the Lord and are spending their lives in His service. I look forward to serving with Elder Christofferson, who has now been called to that quorum and who has received your sustaining vote. He too has been prepared for the position to which he has been called. It has also been a joy to serve with the members of the Quorums of the Seventy and with the Presiding Bishopric. New members of the Seventy have been called and were sustained yesterday, and I look forward to associating with them in the work of the Master.

A sweet spirit of unity exists among the General Authorities. The Lord has declared, “If ye are not one ye are not mine.” We will continue to be united in one purpose—namely, the furtherance of the work of the Lord.

I feel to express thanks to my Heavenly Father for His countless blessings to me. I can say, as did Nephi of old, that I was born of goodly parents, whose own parents and grandparents were gathered out of the lands of Sweden and Scotland and England by dedicated missionaries. As those missionaries bore humble testimonies, they touched the hearts and the spirits of my forebears. After joining the Church, these noble men, women, and children made their way to the valley of the Great Salt Lake. Many were the trials and heartaches they encountered along the way.

In the spring of 1848, my great-great-grandparents, Charles Stewart Miller and Mary McGowan Miller, who had joined the Church in their native Scotland, left their home in Rutherglen, Scotland, and journeyed to St. Louis, Missouri, with a group of Saints, arriving there in 1849. One of their 11 children, Margaret, would become my great-grandmother.

While the family was in St. Louis working to earn enough money to complete their journey to the Salt Lake Valley, a plague of cholera swept through the area, leaving death and heartache in its wake. The Miller family was hard hit. In the space of two weeks, four of the family members succumbed. The first, on June 22, 1849, was 18-year-old William. Five days later Mary McGowan Miller, my great-great-grandmother and the mother of the family, died. Two days afterward, 15-year-old Archibald passed away, and five days after his death, my great-great-grandfather, Charles Stewart Miller, father of the family, succumbed. The children who survived were left orphans, including my great-grandmother Margaret, who was 13 years old at the time.

Because of so many deaths in the area, there were no caskets available, at any price, in which to bury the deceased family members. The older surviving boys dismantled the family’s oxen pens in order to make caskets for the family members who had passed away.

Little is recorded of the heartache and struggles of the nine remaining Miller children as they continued to work and save for that journey their parents and brothers would never make. We know that they left St. Louis in the spring of 1850 with four oxen and one wagon, arriving finally in the Salt Lake Valley that same year.

Others of my ancestors faced similar hardships. Through it all, however, their testimonies remained steadfast and firm. From all of them I received a legacy of total dedication to the gospel of Jesus Christ. Because of these faithful souls, I stand before you today.

I thank my Father in Heaven for my sweet companion, Frances. This October she and I will celebrate 60 wonderful years of marriage. Although my Church service began at an early age, she has never once complained when I’ve left home to attend meetings or to fulfill an assignment. For many years my assignments as a member of the Twelve took me away from Salt Lake City often—sometimes for five weeks at a time—leaving her alone to care for our small children and our home. Beginning when I was called as a bishop at the age of 22, we have seldom had the luxury of sitting together during a Church service. I could not have asked for a more loyal, loving, and understanding companion.

I express gratitude to my Heavenly Father for our three children and their companions, for eight wonderful grandchildren, and for four beautiful great-grandchildren.

It’s difficult for me to find the words to convey to you, my brothers and sisters, my heartfelt appreciation for the lives you live, for the good you do, for the testimonies you bear. You serve one another willingly. You are dedicated to the gospel of Jesus Christ.

During more than 44 years as a General Authority, I have had the opportunity to travel the world over. One of my greatest joys has been to meet with you, the members, wherever you may be—to feel of your spirit and your love. I look forward to many more such opportunities.

Throughout the journey along the pathway of life, there are casualties. Some depart from the road markers which point toward life eternal, only to discover the detour chosen ultimately leads to a dead end. Indifference, carelessness, selfishness, and sin all take their costly toll in human lives.

Change for the better can come to all. Over the years we have issued appeals to the less active, the offended, the critical, the transgressor—to come back. “Come back and feast at the table of the Lord, and taste again the sweet and satisfying fruits of fellowship with the Saints.”

In the private sanctuary of one’s own conscience lies that spirit, that determination to cast off the old person and to measure up to the stature of true potential. In this spirit, we again issue that heartfelt invitation: Come back. We reach out to you in the pure love of Christ and express our desire to assist you and to welcome you into full fellowship. To those who are wounded in spirit or who are struggling and fearful, we say, Let us lift you and cheer you and calm your fears. Take literally the Lord’s invitation, “Come unto me, all ye that labour and are heavy laden, and I will give you rest. Take my yoke upon you, and learn of me; for I am meek and lowly in heart: and ye shall find rest unto your souls. For my yoke is easy, and my burden is light.”

It was said of the Savior that He “went about doing good … for God was with him.” May we follow that perfect example. In this sometimes precarious journey through mortality, may we also follow that advice from the Apostle Paul which will help to keep us safe and on course: “Whatsoever things are true, whatsoever things are honest, whatsoever things are just, whatsoever things are pure, whatsoever things are lovely, whatsoever things are of good report; if there be any virtue, and if there be any praise, think on these things.”

I would encourage members of the Church wherever they may be to show kindness and respect for all people everywhere. The world in which we live is filled with diversity. We can and should demonstrate respect toward those whose beliefs differ from ours.

May we also demonstrate kindness and love within our own families. Our homes are to be more than sanctuaries; they should also be places where God’s Spirit can dwell, where the storm stops at the door, where love reigns and peace dwells.

The world can at times be a frightening place in which to live. The moral fabric of society seems to be unraveling at an alarming speed. None—whether young or old or in between—is exempt from exposure to those things which have the potential to drag us down and destroy us. Our youth, our precious youth, in particular, face temptations we can scarcely comprehend. The adversary and his hosts seem to be working nonstop to cause our downfall.

We are waging a war with sin, my brothers and sisters, but we need not despair. It is a war we can and will win. Our Father in Heaven has given us the tools we need in order to do so. He is at the helm. We have nothing to fear. He is the God of light. He is the God of hope. I testify that He loves us—each one.

Mortality is a period of testing, a time to prove ourselves worthy to return to the presence of our Heavenly Father. In order to be tested, we must sometimes face challenges and difficulties. At times there appears to be no light at the tunnel’s end—no dawn to break the night’s darkness. We feel surrounded by the pain of broken hearts, the disappointment of shattered dreams, and the despair of vanished hopes. We join in uttering the biblical plea, “Is there no balm in Gilead?” We are inclined to view our own personal misfortunes through the distorted prism of pessimism. We feel abandoned, heartbroken, alone. If you find yourself in such a situation, I plead with you to turn to our Heavenly Father in faith. He will lift you and guide you. He will not always take your afflictions from you, but He will comfort and lead you with love through whatever storm you face.

With all my heart and the fervency of my soul, I lift my voice in testimony today as a special witness and declare that God does live. Jesus is His Son, the Only Begotten of the Father in the flesh. He is our Redeemer; He is our Mediator with the Father. He loves us with a love we cannot fully comprehend, and because He loves us, He gave His life for us. My gratitude to Him is beyond expression.

I invoke His blessings upon you, my beloved brothers and sisters, in your homes, in your work, in your service to one another and to the Lord Himself. Together we shall move forward doing His work.

I pledge my life, my strength—all that I have to offer—in serving Him and in directing the affairs of His Church in accordance with His will and by His inspiration, and I do so in His holy name, even the Lord Jesus Christ, amen.

\newpage
\settalk{Finding Joy in the Journey}
\setspeaker{Thomas S. Monson}
\settalkdate{October 2008}
\talkheading{Finding Joy in the Journey}{Thomas S. Monson}

My dear brothers and sisters, I am humbled as I stand before you this morning. I ask for your faith and prayers in my behalf as I speak about those things which have been on my mind and which I have felt impressed to share with you.

I begin by mentioning one of the most inevitable aspects of our lives here upon the earth, and that is change. At one time or another we’ve all heard some form of the familiar adage: “Nothing is as constant as change.”

Throughout our lives, we must deal with change. Some changes are welcome; some are not. There are changes in our lives which are sudden, such as the unexpected passing of a loved one, an unforeseen illness, the loss of a possession we treasure. But most of the changes take place subtly and slowly.

This conference marks 45 years since I was called to the Quorum of the Twelve Apostles. As the junior member of the Twelve then, I looked up to 14 exceptional men, who were senior to me in the Twelve and the First Presidency. One by one, each of these men has returned home. When President Hinckley passed away eight months ago, I realized that I had become the senior Apostle. The changes over a period of 45 years that were incremental now seem monumental.

This coming week Sister Monson and I will celebrate our 60th wedding anniversary. As I look back to our beginnings, I realize just how much our lives have changed since then. Our beloved parents, who stood beside us as we commenced our journey together, have passed on. Our three children, who filled our lives so completely for many years, are grown and have families of their own. Most of our grandchildren are grown, and we now have four great-grandchildren.

Day by day, minute by minute, second by second we went from where we were to where we are now. The lives of all of us, of course, go through similar alterations and changes. The difference between the changes in my life and the changes in yours is only in the details. Time never stands still; it must steadily march on, and with the marching come the changes.

This is our one and only chance at mortal life—here and now. The longer we live, the greater is our realization that it is brief. Opportunities come, and then they are gone. I believe that among the greatest lessons we are to learn in this short sojourn upon the earth are lessons that help us distinguish between what is important and what is not. I plead with you not to let those most important things pass you by as you plan for that illusive and nonexistent future when you will have time to do all that you want to do. Instead, find joy in the journey—now.

I am what my wife, Frances, calls a “show-a-holic.” I thoroughly enjoy many musicals, and one of my favorites was written by the American composer Meredith Willson and is entitled The Music Man. Professor Harold Hill, one of the principal characters in the show, voices a caution that I share with you. Says he, “You pile up enough tomorrows, and you’ll find you’ve collected a lot of empty yesterdays.”

My brothers and sisters, there is no tomorrow to remember if we don’t do something today.

I’ve shared with you previously an example of this philosophy. I believe it bears repeating. Many years ago, Arthur Gordon wrote in a national magazine, and I quote:

“When I was around thirteen and my brother ten, Father had promised to take us to the circus. But at lunchtime there was a phone call; some urgent business required his attention downtown. We braced ourselves for disappointment. Then we heard him say [into the phone], ‘No, I won’t be down. It’ll have to wait.’

“When he came back to the table, Mother smiled. ‘The circus keeps coming back, you know,’ [she said.]

“‘I know,’ said Father. ‘But childhood doesn’t.’”

If you have children who are grown and gone, in all likelihood you have occasionally felt pangs of loss and the recognition that you didn’t appreciate that time of life as much as you should have. Of course, there is no going back, but only forward. Rather than dwelling on the past, we should make the most of today, of the here and now, doing all we can to provide pleasant memories for the future.

If you are still in the process of raising children, be aware that the tiny fingerprints that show up on almost every newly cleaned surface, the toys scattered about the house, the piles and piles of laundry to be tackled will disappear all too soon and that you will—to your surprise—miss them profoundly.

Stresses in our lives come regardless of our circumstances. We must deal with them the best we can. But we should not let them get in the way of what is most important—and what is most important almost always involves the people around us. Often we assume that they must know how much we love them. But we should never assume; we should let them know. Wrote William Shakespeare, “They do not love that do not show their love.” We will never regret the kind words spoken or the affection shown. Rather, our regrets will come if such things are omitted from our relationships with those who mean the most to us.

Send that note to the friend you’ve been neglecting; give your child a hug; give your parents a hug; say “I love you” more; always express your thanks. Never let a problem to be solved become more important than a person to be loved. Friends move away, children grow up, and loved ones pass on. It’s so easy to take others for granted—until that day when they’re gone from our lives and we are left with feelings of “what if” and “if only.” Said author Harriet Beecher Stowe, “The bitterest tears shed over graves are for words left unsaid and deeds left undone.”

In the 1960s during the Vietnam War, Church member Jay Hess, an airman, was shot down over North Vietnam. For two years his family had no idea whether he was dead or alive. His captors in Hanoi eventually allowed him to write home but limited his message to less than 25 words. What would you and I say to our families if we were in the same situation—not having seen them for over two years and not knowing if we would ever see them again? Wanting to provide something his family could recognize as having come from him and also wanting to give them valuable counsel, Brother Hess wrote—and I quote: “These things are important: temple marriage, mission, college. Press on, set goals, write history, take pictures twice a year.”

Let us relish life as we live it, find joy in the journey, and share our love with friends and family. One day each of us will run out of tomorrows.

In the book of John in the New Testament, chapter 13, verse 34, the Savior admonishes us, “As I have loved you, … love one another.”

Some of you may be familiar with Thornton Wilder’s classic drama Our Town. If you are, you will remember the town of Grover’s Corners, where the story takes place. In the play, Emily Webb dies in childbirth, and we read of the lonely grief of her young husband, George, left with their four-year-old son. Emily does not wish to rest in peace; she wants to experience again the joys of her life. She is granted the privilege of returning to earth and reliving her 12th birthday. At first it is exciting to be young again, but the excitement wears off quickly. The day holds no joy now that Emily knows what is in store for the future. It is unbearably painful to realize how unaware she had been of the meaning and wonder of life while she was alive. Before returning to her resting place, Emily laments, “Do … human beings ever realize life while they live it—every, every minute?”

Our realization of what is most important in life goes hand in hand with gratitude for our blessings.

Said one well-known author: “Both abundance and lack [of abundance] exist simultaneously in our lives, as parallel realities. It is always our conscious choice which secret garden we will tend … when we choose not to focus on what is missing from our lives but are grateful for the abundance that’s present—love, health, family, friends, work, the joys of nature, and personal pursuits that bring us [happiness]—the wasteland of illusion falls away and we experience heaven on earth.”

In the Doctrine and Covenants, section 88, verse 33, we are told: “For what doth it profit a man if a gift is bestowed upon him, and he receive not the gift? Behold, he rejoices not in that which is given unto him, neither rejoices in him who is the giver of the gift.”

The ancient Roman philosopher Horace admonished, “Whatever hour God has blessed you with, take it with grateful hand, nor postpone your joys from year to year, so that in whatever place you have been, you may say that you have lived happily.”

Many years ago I was touched by the story of Borghild Dahl. She was born in Minnesota in 1890 of Norwegian parents and from her early years suffered severely impaired vision. She had a tremendous desire to participate in everyday life despite her handicap and, through sheer determination, succeeded in nearly everything she undertook. Against the advice of educators, who felt her handicap was too great, she attended college, receiving her bachelor of arts degree from the University of Minnesota. She later studied at Columbia University and the University of Oslo. She eventually became the principal of eight schools in western Minnesota and North Dakota.

She wrote the following in one of the 17 books she authored: “I had only one eye, and it was so covered with dense scars that I had to do all my seeing through one small opening in the left of the eye. I could see a book only by holding it up close to my face and by straining my one eye as hard as I could to the left.”

Miraculously, in 1943—when she was over 50 years old—a revolutionary procedure was developed which finally restored to her much of the sight she had been without for so long. A new and exciting world opened up before her. She took great pleasure in the small things most of us take for granted, such as watching a bird in flight, noticing the light reflected in the bubbles of her dishwater, or observing the phases of the moon each night. She closed one of her books with these words: “Dear … Father in heaven, I thank Thee. I thank Thee.”

Borghild Dahl, both before and after her sight was restored, was filled with gratitude for her blessings.

In 1982, two years before she died, at the age of 92 her last book was published. Its title: Happy All My Life. Her attitude of thankfulness enabled her to appreciate her blessings and to live a full and rich life despite her challenges.

In 1 Thessalonians in the New Testament, chapter 5, verse 18, we are told by the Apostle Paul, “In every thing give thanks: for this is the will of God.”

Recall with me the account of the 10 lepers:

“And as [Jesus] entered into a certain village, there met him ten men that were lepers, which stood afar off:

“And they lifted up their voices, and said, Jesus, Master, have mercy on us.

“And when he saw them, he said unto them, Go shew yourselves unto the priests. And it came to pass, that, as they went, they were cleansed.

“And one of them, when he saw that he was healed, turned back, and with a loud voice glorified God,

“And fell down on his face at his feet, giving him thanks: and he was a Samaritan.

“And Jesus answering said, Were there not ten cleansed? but where are the nine?

“There are not found that returned to give glory to God, save this stranger.”

Said the Lord in a revelation given through the Prophet Joseph Smith, “In nothing doth man offend God, or against none is his wrath kindled, save those who confess not his hand in all things.” May we be found among those who give our thanks to our Heavenly Father. If ingratitude be numbered among the serious sins, then gratitude takes its place among the noblest of virtues.

Despite the changes which come into our lives and with gratitude in our hearts, may we fill our days—as much as we can—with those things which matter most. May we cherish those we hold dear and express our love to them in word and in deed.

In closing, I pray that all of us will reflect gratitude for our Lord and Savior, Jesus Christ. His glorious gospel provides answers to life’s greatest questions: Where did we come from? Why are we here? Where does my spirit go when I die?

He taught us how to pray. He taught us how to serve. He taught us how to live. His life is a legacy of love. The sick He healed; the downtrodden He lifted; the sinner He saved.

The time came when He stood alone. Some Apostles doubted; one betrayed Him. The Roman soldiers pierced His side. The angry mob took His life. There yet rings from Golgotha’s hill His compassionate words, “Father, forgive them; for they know not what they do.”

Earlier, perhaps perceiving the culmination of His earthly mission, He spoke the lament, “Foxes have holes, and the birds of the air have nests; but the Son of man hath not where to lay his head.” “No room in the inn” was not a singular expression of rejection—just the first. Yet He invites you and me to receive Him. “Behold, I stand at the door, and knock: if any man hear my voice, and open the door, I will come in to him, and will sup with him, and he with me.”

Who was this Man of sorrows, acquainted with grief? Who is the King of glory, this Lord of hosts? He is our Master. He is our Savior. He is the Son of God. He is the Author of our Salvation. He beckons, “Follow me.” He instructs, “Go, and do thou likewise.” He pleads, “Keep my commandments.”

Let us follow Him. Let us emulate His example. Let us obey His word. By so doing, we give to Him the divine gift of gratitude.

Brothers and sisters, my sincere prayer is that we may adapt to the changes in our lives, that we may realize what is most important, that we may express our gratitude always and thus find joy in the journey. In the name of Jesus Christ, amen.

\newpage
\settalk{To Learn, to Do, to Be}
\setspeaker{Thomas S. Monson}
\settalkdate{October 2008}
\talkheading{To Learn, to Do, to Be}{Thomas S. Monson}

You’ve seen a witness tonight of the strength of the two counselors in this First Presidency. I stand before you and declare this First Presidency is united as one under the direction of the Lord Jesus Christ.

I want to especially thank this missionary choir. I had an experience I think they may be interested in, and you may find it interesting also. Many years ago I had a desperate call from the head of the missionary training center. He said, “President Monson, I have a missionary who is going home. Nothing can prevent him from quitting.”

I replied, “Well, that’s not singular. It’s happened before. What’s his problem?”

He said, “He’s been called to a Spanish-speaking mission, and he’s absolutely certain he cannot learn Spanish.”

I said, “I have a suggestion for you. Tomorrow morning have him attend a class learning Japanese. And then have him report to you at 12:00 noon.”

The next morning he phoned at 10:00! He said, “The young man is here with me now, and he wants me to know he’s absolutely certain he can learn Spanish.”

When there’s a will, there’s a way.

Now, as I speak to you tonight, truly you are a royal priesthood, assembled in many places but in unity. In all likelihood this is the largest assemblage of priesthood holders ever to come together. Your devotion to your sacred callings is inspiring. Your desire to learn your duty is evident. The purity of your souls brings heaven closer to you and your families.

Many areas of the world have experienced difficult economic times. Businesses have failed, jobs have been lost, and investments have been jeopardized. We must make certain that those for whom we share responsibility do not go hungry or unclothed or unsheltered. When the priesthood of this Church works together as one in meeting these vexing conditions, near miracles take place.

We urge all Latter-day Saints to be prudent in their planning, to be conservative in their living, and to avoid excessive or unnecessary debt. The financial affairs of the Church are being managed in this manner, for we are aware that your tithing and other contributions have not come without sacrifice and are sacred funds.

Let us make our homes sanctuaries of righteousness, places of prayer, and abodes of love that we might merit the blessings that can come only from our Heavenly Father. We need His guidance in our daily lives.

In this vast throng are priesthood power and the capacity to reach out and share the glorious gospel with others. As has been mentioned, we have the hands to lift others from complacency and inactivity. We have the hearts to serve faithfully in our priesthood callings and thereby inspire others to walk on higher ground and to avoid the swamps of sin which threaten to engulf so many. The worth of souls is indeed great in the sight of God. Ours is the precious privilege, armed with this knowledge, to make a difference in the lives of others. The words found in Ezekiel could well pertain to all of us who follow the Savior in this sacred work:

“A new heart … will I give you, and a new spirit will I put within you. …

“And I will put my spirit within you, and cause you to walk in my statutes, and ye shall keep my judgments, and do them.

“And ye shall dwell in the land that I gave to your fathers; and ye shall be my people, and I will be your God.”

How might we merit this promise? What will qualify us to receive this blessing? Is there a guide to follow?

May I suggest three imperatives for our consideration. They apply to the deacon as well as to the high priest. They are within our reach. A kind Heavenly Father will help us in our quest.

First, learn what we should learn.

Second, do what we should do.

And third, be what we should be.

Let us discuss these objectives, that we might be profitable servants in the sight of our Lord.

First, learn what we should learn. The Apostle Paul placed an urgency on our efforts to learn. He said to the Philippians, “One thing I do, forgetting those things which are behind, and reaching forth unto those things which are before, I press toward the mark for the prize of the high calling of God in Christ Jesus.” And to the Hebrews he urged, “Lay aside … sin[;] … let us run with patience the race … set before us, looking [for an example] unto Jesus the author and finisher of our faith.”

President Stephen L Richards, who served for many years in the Quorum of the Twelve Apostles and then in the First Presidency, spoke often to holders of the priesthood and emphasized his philosophy pertaining to it. He declared: “The Priesthood is usually simply defined as ‘the power of God delegated to man.’ This definition, I think, is accurate.”

He continued: “But for practical purposes I like to define the Priesthood in terms of service and I frequently call it ‘the perfect plan of service.’ I do so because it seems to me that it is only through the utilization of the divine power conferred on men that they may ever hope to realize the full import and vitality of this endowment. It is an instrument of service … and the man who fails to use it is apt to lose it, for we are plainly told by revelation that he who neglects it ‘shall not be counted worthy to stand.’”

President Harold B. Lee, 11th President of the Church and one of the great teachers in the Church, put his counsel in easy-to-understand terms. Said he: “When one becomes a holder of the priesthood, he becomes an agent of the Lord. He should think of his calling as though he were on the Lord’s errand.”

Now, some of you may be shy by nature or consider yourselves inadequate to respond affirmatively to a calling. Remember that this work is not yours and mine alone. It is the Lord’s work, and when we are on the Lord’s errand, we are entitled to the Lord’s help. Remember that the Lord will shape the back to bear the burden placed upon it.

While the formal classroom may be intimidating at times, some of the most effective teaching takes place other than in the chapel or the classroom. Well do I remember that some years ago, members holding the Aaronic Priesthood would eagerly look forward to an annual outing commemorating the restoration of the Aaronic Priesthood. By the busload the young men of our stake journeyed 90 miles (145 km) north to the Clarkston Cemetery, where we viewed the grave of Martin Harris, one of the Three Witnesses of the Book of Mormon. While surrounding the beautiful granite shaft which marks the grave, a high councilor would present background concerning the life of Martin Harris, read from the Book of Mormon his testimony, and then bear his own witness to the truth. The young men listened with rapt attention, touched the granite marker, and pondered the words they had heard and the feelings they had felt.

At a park in Logan, lunch was enjoyed. The group of young men would then lie down on the lawn at the Logan Temple and gaze upward at its lofty spires. Often beautiful white clouds would hurry past the spires, moved along by a gentle breeze. The purpose of temples was taught. Covenants and promises became much more than words. The desire to be worthy to enter those temple doors entered those youthful hearts. Heaven was very close. Learning what we should learn was assured.

Number two, do what we should do. In a revelation on priesthood, given through Joseph Smith the Prophet, recorded as the 107th section of the Doctrine and Covenants, “learning” moves to “doing” as we read, “Wherefore, now let every man learn his duty, and to act in the office in which he is appointed, in all diligence.”

Each priesthood holder attending this session tonight has a calling to serve, to put forth his best efforts in the work assigned to him. No assignment is menial in the work of the Lord, for each has eternal consequences. President John Taylor warned us, “If you do not magnify your callings, God will hold you responsible for those whom you might have saved had you done your duty.” And who of us can afford to be responsible for the delay of eternal life of a human soul? If great joy is the reward of saving one soul, then how terrible must be the remorse of those whose timid efforts have allowed a child of God to go unwarned or unaided so that he has to wait till a dependable servant of God comes along.

The old adage is ever true: “Do your duty, that is best; leave unto the Lord the rest.”

Most service given by priesthood holders is accomplished quietly, without fanfare. A friendly smile, a warm handclasp, a sincere testimony of truth can literally lift lives, change human nature, and save precious souls.

An example of such service was the missionary experience of Juliusz and Dorothy Fussek, who were called to fill a two-year mission in Poland. Brother Fussek was born in Poland. He spoke the language. He loved the people. Sister Fussek was English and knew little of Poland and its people.

Trusting in the Lord, they embarked on their assignment. The living conditions were primitive, the work lonely, their task immense. A mission had not at that time been established in Poland. The assignment given the Fusseks was to prepare the way, that a mission could be established so that other missionaries could be called to serve, people could be taught, converts could be baptized, branches could be established, and chapels could be erected.

Did Elder and Sister Fussek despair because of the enormity of their assignment? Not for a moment. They knew their calling was from God. They prayed for His divine help, and they devoted themselves wholeheartedly to their work. They remained in Poland not two years but five years. All of the foregoing objectives were realized.

Elders Russell M. Nelson, Hans B. Ringger, and I, accompanied by Elder Fussek, met with Minister Adam Wopatka of the Polish government, and we heard him say, “Your church is welcome here. You may build your buildings; you may send your missionaries. You are welcome in Poland. This man,” pointing to Juliusz Fussek, “has served your church well. You can be grateful for his example and his work.”

Like the Fusseks, let us do what we should do in the work of the Lord. Then we can, with Juliusz and Dorothy Fussek, echo the Psalm: “My help cometh from the Lord, which made heaven and earth … : he that keepeth thee will not slumber. Behold, he that keepeth Israel shall neither slumber nor sleep.”

Third, be what we should be. Paul counseled his beloved friend and associate Timothy, “Be thou an example of the believers, in word, in conversation, in charity, in spirit, in faith, in purity.”

I would urge all of us to pray concerning our assignments and to seek divine help, that we might be successful in accomplishing that which we are called to do. Someone has said that “the recognition of power higher than man himself does not in any sense debase him.” He must seek, believe in, pray, and hope that he will find. No such sincere, prayerful effort will go unanswered: that is the very constitution of the philosophy of faith. Divine favor will attend those who humbly seek it.

From the Book of Mormon comes counsel that says it all. The Lord speaks: “Therefore, what manner of men ought ye to be? Verily I say unto you, even as I am.”

And what manner of man was He? What example did He set in His service? From John chapter 10 we learn:

“I am the good shepherd: the good shepherd giveth his life for the sheep.

“But he that is an hireling, and not the shepherd, whose own the sheep are not, seeth the wolf coming, and leaveth the sheep, and fleeth: and the wolf catcheth them, and scattereth the sheep.

“The hireling fleeth, because he is an hireling, and careth not for the sheep.”

Said the Lord: “I am the good shepherd, and know my sheep, and am known of mine.

“As the Father knoweth me, even so know I the Father: and I lay down my life for the sheep.”

Brethren, may we learn what we should learn, do what we should do, and be what we should be. By so doing, the blessings of heaven will attend. We will know that we are not alone. He who notes the sparrow’s fall will, in His own way, acknowledge us.

Several years ago I received a letter from a longtime friend. He bore his testimony in that letter. I would like to share part of it with you tonight, since it illustrates the strength of the priesthood in one who learned what he should learn, who did what he should do, and who always tried to be what he should be. I shall read excerpts of that letter from my friend Theron W. Borup, who passed away three years ago at the age of 90:

“At the age of eight, when I was baptized and received the Holy Ghost, I was much impressed about being good and able to have the Holy Ghost to be a help throughout my life. I was told that the Holy Ghost associated only in good company and that when evil entered our lives, he would leave. Not knowing when I would need his promptings and guidance, I tried to so live that I would not lose this gift. On one occasion it saved my life.

“During World War II, I was an engineer-gunner in a B-24 bomber fighting in the South Pacific. … One day there was an announcement that the longest bombing flight ever made would be attempted to knock out an oil refinery. The promptings of the Spirit told me I would be assigned on this flight but that I would not lose my life. At the time I was the president of the LDS group.

“The combat was ferocious as we flew over Borneo. Our plane was hit by attacking planes and soon burst into flames, and the pilot told us to prepare to jump. I went out last. We were shot at by enemy pilots as we floated down. I had trouble inflating my life raft. Bobbing up and down in the water, I began to drown and passed out. I came to momentarily and cried, ‘God save me!’ … Again I tried inflating the life raft and this time was successful. With just enough air in it to keep me afloat, I rolled over on top of it, too exhausted to move.

“For three days we floated about in enemy territory with ships all about us and planes overhead. Why they couldn’t see a yellow group of rafts on blue water is a mystery,” he wrote. “A storm came up, and waves thirty feet high almost tore our rafts apart. Three days went by with no food or water. The others asked me if I prayed. I answered that I did pray and we would indeed be rescued. That evening we saw our submarine that was there to rescue us, but it passed by. The next morning it did [the same. We knew] this was the last day [it would] be in the area. Then came the promptings of the Holy Ghost. ‘You have the priesthood. Command the sub to pick you up.’ Silently I prayed, ‘In the name of Jesus Christ, and by the power of the priesthood, turn about and pick us up.’ In a few minutes, they were alongside of us. When on deck, the captain … said, ‘I don’t know how we ever found you, for we were not even looking for you.’ I knew.”

I leave with you my testimony that this work in which we are engaged is true. The Lord is at the helm. That we may ever follow Him is my sincere prayer, and I ask it in the name of Jesus Christ, amen.

\newpage
\settalk{Until We Meet Again}
\setspeaker{Thomas S. Monson}
\settalkdate{October 2008}
\talkheading{Until We Meet Again}{Thomas S. Monson}

Brothers and sisters, I know you will agree with me that this has been a most inspiring conference. We have felt the Spirit of the Lord in rich abundance these past two days as our hearts have been touched and our testimonies of this divine work have been strengthened. I am certain I represent the membership of the Church everywhere in expressing appreciation to the brethren and sisters who have addressed us. I am reminded of the words of Moroni found in the Book of Mormon: “And their meetings were conducted by the church after the manner of the workings of the Spirit, and by the power of the Holy Ghost; for as the power of the Holy Ghost led them whether to preach, or to exhort, or to pray, or to supplicate, or to sing, even so it was done.”

May we long remember what we have heard during this general conference. The messages which have been given will be printed in next month’s Ensign and Liahona magazines. I urge you to study them and to ponder their teachings.

To you brethren who have been released at this conference, we express our deep appreciation for your many years of dedicated service. The entire membership of the Church has benefited from your countless contributions.

I assure you that our Heavenly Father is mindful of the challenges we face in the world today. He loves each of us and will bless us as we strive to keep His commandments and seek Him through prayer.

We are a global church, brothers and sisters. Our membership is found throughout the world. May we be good citizens of the nations in which we live and good neighbors in our communities, reaching out to those of other faiths, as well as to our own. May we be men and women of honesty and integrity in everything we do.

There are those throughout the world who are hungry; there are those who are destitute. Working together, we can alleviate suffering and provide for those in need. In addition to the service you give as you care for one another, your contributions to the funds of the Church enable us to respond almost immediately when disasters occur anywhere in the world. We are nearly always among the first on the scene to provide whatever assistance we can. We thank you for your generosity.

There are other difficulties in the lives of some. Particularly among the young, there are those who are tragically involved in such things as drugs, immorality, gangs, and all the serious problems that go with them. In addition, there are those who are lonely, including widows and widowers, who long for the company and the concern of others. May we ever be mindful of the needs of those around us and be ready to extend a helping hand and a loving heart.

Brothers and sisters, how blessed we are that the heavens are indeed open, that the restored Church of Jesus Christ is upon the earth today, and that the Church is founded upon the rock of revelation. We know that continuous revelation is the very lifeblood of the gospel of Jesus Christ.

May each of us go safely to our homes. May we live together in peace and harmony and love. May we strive each day to follow the Savior’s example.

God bless you, my brothers and sisters. I thank you for your prayers in my behalf and in behalf of all the General Authorities. We are deeply grateful for you.

In one of Christopher Marlowe’s plays, The Tragical History of Dr. Faustus, there is portrayed an individual, Dr. Faustus, who chose to ignore God and follow the pathway of Satan. At the end of his wicked life, and facing the frustration of opportunities lost and punishment certain to come, he lamented, “[There is] more searing anguish than [flaming] fire—eternal exile from God.”

My brothers and sisters, just as eternal exile from God may be the most searing anguish, so eternal life in the presence of God is our most precious goal.

With all my heart and soul I pray that we might continue to persevere in the pursuit of this most precious prize.

I bear witness to you that this work is true, that our Savior lives, and that He guides and directs His Church here upon the earth. I bid you farewell, my dear brothers and sisters, until we meet again in six months. In the name of Jesus of Nazareth, our Savior, our Redeemer whom we serve, amen.

\newpage
\settalk{Welcome to Conference}
\setspeaker{Thomas S. Monson}
\settalkdate{October 2008}
\talkheading{Welcome to Conference}{Thomas S. Monson}

My dear brothers and sisters, the past six months since last we met seem to have flown by. Much has transpired as the work of the Lord has moved forward uninterrupted.

It has been my privilege, accompanied by my counselors and by other General Authorities, to dedicate three new temples: in Curitiba, Brazil; in Panama City, Panama; and in Twin Falls, Idaho—bringing to 128 the number of temples in operation throughout the world.

The evening before each of the temple dedications took place, magnificent cultural events were held. In Curitiba, Brazil, 4,330 members from the temple district, supported by a choir of 1,700 voices, presented a most inspirational program through song, dance, and video. The enormous soccer stadium where the event took place was filled with spectators. The wind had been blowing, and rain threatened. I offered a silent prayer asking Heavenly Father to look with mercy upon those who had prepared so diligently for our entertainment and whose costumes and presentations would be damaged if a heavy rain or wind enveloped them. He honored that prayer, and it wasn’t until the end of the show and later on that evening that rain fell in abundance.

A history of the Church in Brazil was portrayed in song and dance. A particularly moving scene was the portrayal of Elders James E. Faust and William Grant Bangerter, who served as missionaries in Curitiba in 1940. As their photos were displayed on large screens, a tremendous cheer went up from the audience. All in all, it was a glorious event.

In Panama City, Panama, the evening before the dedication of the temple there, we watched some 900 of our youth, gathered from across Panama. They were dressed in colorful folkloric costumes as they danced and presented messages of family, fellowship, and faith. We learned that they had been practicing for a year. They came from points as distant as the San Blas Islands and the Changuinola region in northeast Panama. The trip to the capital city for the San Blas youth exacted three days of travel over land and sea. The event was magnificent and inspiring.

In preparation for our most recent temple dedication, in Twin Falls, Idaho, local Church members constructed a huge stage at the Filer, Idaho, fairgrounds and filled the dirt arena with sod and other decorations, including a large waterfall to represent Shoshone Falls, a popular landmark located two miles (3 km) from the new temple. The evening of the performance, 3,200 young men and young women entered the arena waving blue and white ribbons, turning the arena into a representation of a great river of flowing water. Titled “Living Water,” from John 4:10, 14, the celebration brought together youth from 14 stakes in the new temple district. They depicted through song and dance both their dependence for their spiritual lives on the living water from the Savior and their dependence for their physical lives on the mountain streams and rivers in their area. Those of us privileged to witness this event were uplifted and edified.

I am an advocate for such events. They enable our youth to participate in something they truly find unforgettable. The friendships they form and the memories they make will be theirs forever.

Next month the Mexico City Mexico Temple will be rededicated following extensive renovations. In the coming months, the construction of other temples will be completed, and open houses and dedications will take place.

This morning I am pleased to announce five new temples for which sites have been acquired and which, in coming months and years, will be built in the following locations: Calgary, Alberta, Canada; Córdoba, Argentina; the greater Kansas City area; Philadelphia, Pennsylvania; and Rome, Italy.

Brothers and sisters, our missionary force, serving throughout the world, continues to seek out those who are searching for the truths which are found in the gospel of Jesus Christ. The Church is steadily growing; it has since its organization over 178 years ago.

It has been my privilege during the past six months to meet with leaders of countries and with representatives of governments. Those with whom I’ve met feel kindly toward the Church and our members, and they have been cooperative and accommodating. There remain, however, areas of the world where our influence is limited and where we are not allowed to share the gospel freely. As did President Spencer W. Kimball over 32 years ago, I urge you to pray for the opening of those areas, that we might share with them the joy of the gospel. As we prayed then in response to President Kimball’s pleadings, we saw miracles unfold as country after country, formerly closed to the Church, was opened. Such will transpire again as we pray with faith.

Now, my brothers and sisters, we have come here to be instructed and inspired. Some of you are new in the Church. We welcome you. Some of you are struggling with problems, with challenges, with disappointments, with losses. We love you and pray for you. Many messages will be shared during the next two days. I can assure you that those men and women who will speak to you have prayed about what they should say. They have been inspired and impressed as they have sought heaven’s help and direction.

Our Heavenly Father is mindful of each one of us and our needs. May we be filled with His Spirit as we partake of the proceedings of this, the 178th Semiannual General Conference of the Church. This is my sincere prayer, and I offer it in the name of Jesus Christ, amen.

\newpage
\settalk{Be of Good Cheer}
\setspeaker{Thomas S. Monson}
\settalkdate{April 2009}
\talkheading{Be of Good Cheer}{Thomas S. Monson}

My dear brothers and sisters, I express my love to you. I am humbled by the responsibility to address you, and yet I am grateful for the opportunity to do so.

Since last we met together in a general conference six months ago, there have been continuing signs that circumstances in the world aren’t necessarily as we would wish. The global economy, which six months ago appeared to be sagging, seems to have taken a nosedive, and for many weeks now the financial outlook has been somewhat grim. In addition, the moral footings of society continue to slip, while those who attempt to safeguard those footings are often ridiculed and, at times, picketed and persecuted. Wars, natural disasters, and personal misfortunes continue to occur.

It would be easy to become discouraged and cynical about the future—or even fearful of what might come—if we allowed ourselves to dwell only on that which is wrong in the world and in our lives. Today, however, I’d like us to turn our thoughts and our attitudes away from the troubles around us and to focus instead on our blessings as members of the Church. The Apostle Paul declared, “God hath not given us the spirit of fear; but of power, and of love, and of a sound mind.”

None of us makes it through this life without problems and challenges—and sometimes tragedies and misfortunes. After all, in large part we are here to learn and grow from such events in our lives. We know that there are times when we will suffer, when we will grieve, and when we will be saddened. However, we are told, “Adam fell that men might be; and men are, that they might have joy.”

How might we have joy in our lives, despite all that we may face? Again from the scriptures: “Wherefore, be of good cheer, and do not fear, for I the Lord am with you, and will stand by you.”

The history of the Church in this, the dispensation of the fulness of times, is replete with the experiences of those who have struggled and yet who have remained steadfast and of good cheer as they have made the gospel of Jesus Christ the center of their lives. This attitude is what will pull us through whatever comes our way. It will not remove our troubles from us but rather will enable us to face our challenges, to meet them head on, and to emerge victorious.

Too numerous to mention are the examples of all the individuals who have faced difficult circumstances and yet who have persevered and prevailed because their faith in the gospel and in the Savior has given them the strength they have needed. This morning, however, I’d like to share with you three such examples.

First, from my own family, I mention a touching experience that has always been an inspiration to me.

My maternal great-grandparents Gibson and Cecelia Sharp Condie lived in Clackmannan, Scotland. Their families were engaged in coal mining. They were at peace with the world, surrounded by relatives and friends, and were housed in fairly comfortable quarters in a land they loved. Then they listened to the message of the missionaries from The Church of Jesus Christ of Latter-day Saints and, to the depths of their very souls, were converted. They heard the call to gather to Zion and knew they must answer that call.

Sometime around 1848, they sold their possessions and prepared for the hazardous voyage across the mighty Atlantic Ocean. With five small children, they boarded a sailing vessel, all their worldly possessions in one tiny trunk. They traveled 3,000 miles (4,800 km) across the waters—eight long, weary weeks on a treacherous sea, watching and waiting, with poor food, poor water, and no help beyond the length and breadth of that small ship.

In the midst of this soul-trying situation, one of their young sons became ill. There were no doctors, no stores at which they might purchase medicine to ease his suffering. They watched, they prayed, they waited, and they wept as day by day his condition deteriorated. When his eyes were at last closed in death, their hearts were torn asunder. To add to their grief, the laws of the sea must be obeyed. Wrapped in a canvas weighed down with iron, the little body was consigned to a watery grave. As they sailed away, only those parents knew the crushing blow dealt to wounded hearts. However, with a faith born of their deep conviction of the truth and their love of the Lord, Gibson and Cecelia held on. They were comforted by the words of the Lord: “In the world ye shall have tribulation: but be of good cheer; I have overcome the world.”

How grateful I am for ancestors who had the faith to leave hearth and home and to journey to Zion, who made sacrifices I can scarcely imagine. I thank my Heavenly Father for the example of faith, of courage, and of determination Gibson and Cecelia Sharp Condie provide for me and for all their posterity.

I introduce next a gentle, faith-filled man who epitomized the peace and joy which the gospel of Jesus Christ can bring into one’s life.

Late one evening on a Pacific isle, a small boat slipped silently to its berth at the crude pier. Two Polynesian women helped Meli Mulipola from the boat and guided him to the well-worn pathway leading to the village road. The women marveled at the bright stars, which twinkled in the midnight sky. The moonlight guided them along their way. However, Meli Mulipola could not appreciate these delights of nature—the moon, the stars, the sky—for he was blind.

Brother Mulipola’s vision had been normal until a fateful day when, while working on a pineapple plantation, light turned suddenly to darkness and day became perpetual night. He was depressed and despondent until he learned the good news of the gospel of Jesus Christ. His life was brought into compliance with the teachings of the Church, and he once again felt hope and joy.

Brother Mulipola and his loved ones had made a long voyage, having learned that one who held the priesthood of God was visiting among the islands of the Pacific. He sought a blessing, and it was my privilege, along with another who held the Melchizedek Priesthood, to provide that blessing to him. As we finished, I noted that tears were streaming from his sightless eyes, coursing down his brown cheeks, and tumbling finally upon his native dress. He dropped to his knees and prayed: “O God, Thou knowest I am blind. Thy servants have blessed me that my sight might return. Whether in Thy wisdom I see light or whether I see darkness all the days of my life, I will be eternally grateful for the truth of Thy gospel, which I now see and which provides the light of my life.”

He rose to his feet and, smiling, thanked us for providing the blessing. He then disappeared into the still of the night. Silently he came; silently he departed. But his presence I shall never forget. I reflected upon the message of the Master: “I am the light of the world: he that followeth me shall not walk in darkness, but shall have the light of life.”

My brothers and sisters, each of us has that light in his or her life. We are not left to walk alone, no matter how dark our pathway.

I love the words penned by M. Louise Haskins:

\begin{verse}
And I said to the man who stood at the gate of the year:\\
“Give me a light, that I may tread safely into the unknown!”\\
And he replied:\\
“Go out into the darkness and put your hand into the Hand of God.\\
That shall be to you better than [a] light and safer than a known way.”
\end{verse}

The setting for my final example of one who persevered and ultimately prevailed, despite overwhelmingly difficult circumstances, begins in East Prussia following World War II.

In about March 1946, less than a year after the end of the war, Ezra Taft Benson, then a member of the Quorum of the Twelve, accompanied by Frederick W. Babbel, was assigned a special postwar tour of Europe for the express purpose of meeting with the Saints, assessing their needs, and providing assistance to them. Elder Benson and Brother Babbel later recounted, from a testimony they heard, the experience of a Church member who found herself in an area no longer controlled by the government under which she had resided.

She and her husband had lived an idyllic life in East Prussia. Then had come the second great world war within their lifetimes. Her beloved young husband was killed during the final days of the frightful battles in their homeland, leaving her alone to care for their four children.

The occupying forces determined that the Germans in East Prussia must go to Western Germany to seek a new home. The woman was German, and so it was necessary for her to go. The journey was over a thousand miles (1,600 km), and she had no way to accomplish it but on foot. She was allowed to take only such bare necessities as she could load into her small wooden-wheeled wagon. Besides her children and these meager possessions, she took with her a strong faith in God and in the gospel as revealed to the latter-day prophet Joseph Smith.

She and the children began the journey in late summer. Having neither food nor money among her few possessions, she was forced to gather a daily subsistence from the fields and forests along the way. She was constantly faced with dangers from panic-stricken refugees and plundering troops.

As the days turned into weeks and the weeks to months, the temperatures dropped below freezing. Each day, she stumbled over the frozen ground, her smallest child—a baby—in her arms. Her three other children struggled along behind her, with the oldest—seven years old—pulling the tiny wooden wagon containing their belongings. Ragged and torn burlap was wrapped around their feet, providing the only protection for them, since their shoes had long since disintegrated. Their thin, tattered jackets covered their thin, tattered clothing, providing their only protection against the cold.

Soon the snows came, and the days and nights became a nightmare. In the evenings she and the children would try to find some kind of shelter—a barn or a shed—and would huddle together for warmth, with a few thin blankets from the wagon on top of them.

She constantly struggled to force from her mind overwhelming fears that they would perish before reaching their destination.

And then one morning the unthinkable happened. As she awakened, she felt a chill in her heart. The tiny form of her three-year-old daughter was cold and still, and she realized that death had claimed the child. Though overwhelmed with grief, she knew that she must take the other children and travel on. First, however, she used the only implement she had—a tablespoon—to dig a grave in the frozen ground for her tiny, precious child.

Death, however, was to be her companion again and again on the journey. Her seven-year-old son died, either from starvation or from freezing or both. Again her only shovel was the tablespoon, and again she dug hour after hour to lay his mortal remains gently into the earth. Next, her five-year-old son died, and again she used her tablespoon as a shovel.

Her despair was all-consuming. She had only her tiny baby daughter left, and the poor thing was failing. Finally, as she was reaching the end of her journey, the baby died in her arms. The spoon was gone now, so hour after hour she dug a grave in the frozen earth with her bare fingers. Her grief became unbearable. How could she possibly be kneeling in the snow at the graveside of her last child? She had lost her husband and all her children. She had given up her earthly goods, her home, and even her homeland.

In this moment of overwhelming sorrow and complete bewilderment, she felt her heart would literally break. In despair she contemplated how she might end her own life, as so many of her fellow countrymen were doing. How easy it would be to jump off a nearby bridge, she thought, or to throw herself in front of an oncoming train.

And then, as these thoughts assailed her, something within her said, “Get down on your knees and pray.” She ignored the prompting until she could resist it no longer. She knelt and prayed more fervently than she had in her entire life:

“Dear Heavenly Father, I do not know how I can go on. I have nothing left—except my faith in Thee. I feel, Father, amidst the desolation of my soul, an overwhelming gratitude for the atoning sacrifice of Thy Son, Jesus Christ. I cannot express adequately my love for Him. I know that because He suffered and died, I shall live again with my family; that because He broke the chains of death, I shall see my children again and will have the joy of raising them. Though I do not at this moment wish to live, I will do so, that we may be reunited as a family and return—together—to Thee.”

When she finally reached her destination of Karlsruhe, Germany, she was emaciated. Brother Babbel said that her face was a purple-gray, her eyes red and swollen, her joints protruding. She was literally in the advanced stages of starvation. In a Church meeting shortly thereafter, she bore a glorious testimony, stating that of all the ailing people in her saddened land, she was one of the happiest because she knew that God lived, that Jesus is the Christ, and that He died and was resurrected so that we might live again. She testified that she knew if she continued faithful and true to the end, she would be reunited with those she had lost and would be saved in the celestial kingdom of God.

From the holy scriptures we read, “Behold, the righteous, the saints of the Holy One of Israel, they who have believed in [Him], they who have endured the crosses of the world, … they shall inherit the kingdom of God, … and their joy shall be full forever.”

I testify to you that our promised blessings are beyond measure. Though the storm clouds may gather, though the rains may pour down upon us, our knowledge of the gospel and our love of our Heavenly Father and of our Savior will comfort and sustain us and bring joy to our hearts as we walk uprightly and keep the commandments. There will be nothing in this world that can defeat us.

My beloved brothers and sisters, fear not. Be of good cheer. The future is as bright as your faith.

I declare that God lives and that He hears and answers our prayers. His Son, Jesus Christ, is our Savior and our Redeemer. Heaven’s blessings await us. In the name of Jesus Christ, amen.

\newpage
\settalk{Be Your Best Self}
\setspeaker{Thomas S. Monson}
\settalkdate{April 2009}
\talkheading{Be Your Best Self}{Thomas S. Monson}

My beloved brethren of the priesthood assembled here in this full Conference Center and in locations throughout the world, I am humbled by the responsibility which is mine to address you. I endorse those messages which have already been presented and express to each of you my sincere love, as well as my appreciation for your faith and your devotion.

Brethren, our responsibilities as bearers of the priesthood are most significant, as outlined in the Doctrine and Covenants: “The power and authority of the higher, or Melchizedek Priesthood, is to hold the keys of all the spiritual blessings of the church.” And further, “The power and authority of the lesser, or Aaronic Priesthood, is to hold the keys of the ministering of angels, and to administer in outward ordinances, the letter of the gospel, the baptism of repentance for the remission of sins, agreeable to the covenants and commandments.”

In 1958 Elder Harold B. Lee, who later served as the 11th President of the Church, described the priesthood as “the Lord’s … troops against the forces of evil.”

President John Taylor stated that “the power manifested by the priesthood is simply the power of God.”

These stirring declarations from prophets of God help us to understand that each man and each boy who holds the priesthood of God must be worthy of that great privilege and responsibility. Each must strive to learn his duty and then to do it to the best of his ability. As we do so, we provide the means by which our Heavenly Father and His Son, Jesus Christ, can accomplish Their work here upon the earth. It is we who are Their representatives here.

In the world today we face difficulties and challenges, some of which can seem truly daunting. However, with God on our side we cannot fail. As we bear His holy priesthood worthily, we will be victorious.

Now to you who hold the Aaronic Priesthood, may I say that I sincerely hope each of you is aware of the significance of your priesthood ordination. Yours is a vital role in the life of every member of your ward as you participate in the administration and passing of the sacrament each Sunday.

I had the privilege to serve as the secretary of my deacons quorum. I recall the many assignments we members of that quorum had the opportunity to fill. Passing the sacred sacrament, collecting the monthly fast offerings, and looking after one another come readily to mind. The most frightening one, however, happened at the leadership session of our ward conference. The member of our stake presidency who was presiding called on a number of the ward officers to speak. They did so. Then, without the slightest warning, he stood and said, “We will now call on one of our younger ward officers, Thomas S. Monson, secretary of the deacons quorum, to give us an accounting of his service and to bear his testimony.” I don’t remember a single thing I said, but I have never forgotten the experience or the lesson that it taught me. It was the Apostle Peter who said, “Be ready always to give an answer to every man that asketh you a reason of the hope that is in you.”

In an earlier generation, the Lord gave this promise to holders of the priesthood: “I will go before your face. I will be on your right hand and on your left, and my Spirit shall be in your hearts, and mine angels round about you, to bear you up.”

This is not a time for fear, brethren, but rather a time for faith—a time for each of us who holds the priesthood to be his best self.

Although our journey through mortality will at times place us in harm’s way, may I offer you tonight three suggestions which, when observed and followed, will lead us to safety. They are:

1. Study diligently.

2. Pray fervently.

3. Live righteously.

These suggestions are not new; they have been taught and repeated again and again. If we incorporate them into our lives, however, we will have the strength to withstand the adversary. Should we ignore them, we will be opening the door for Satan to have influence and power over us.

First, study diligently. Every holder of the priesthood should participate in daily scripture study. Crash courses are not nearly so effective as the day-to-day reading and application of the scriptures in our lives. Become acquainted with the lessons the scriptures teach. Learn the background and setting of the Master’s parables and the prophets’ admonitions. Study them as though they were speaking to you, for such is the truth.

The prophet Lehi and his son Nephi were each shown in vision the importance of obtaining and then holding fast to the word of God. Concerning the rod of iron shown him, Nephi said this to his disbelieving brothers, Laman and Lemuel:

“And I said unto them that [the rod] was the word of God; and whoso would hearken unto the word of God, and would hold fast unto it, they would never perish; neither could the temptations and the fiery darts of the adversary overpower them unto blindness, to lead them away to destruction.

“Wherefore, I, Nephi, did exhort them to give heed unto the word of the Lord; yea, I did exhort them with all the energies of my soul, and with all the faculty which I possessed, that they would give heed to the word of God and remember to keep his commandments always in all things.”

I promise you, whether you hold the Aaronic or the Melchizedek Priesthood, that if you will study the scriptures diligently, your power to avoid temptation and to receive direction of the Holy Ghost in all you do will be increased.

Second, pray fervently. With God, all things are possible. Men of the Aaronic Priesthood, men of the Melchizedek Priesthood, remember the prayer of the Prophet Joseph, offered in that grove called sacred. Look around you and see the result of that answered prayer.

Adam prayed; Jesus prayed. We know the outcome of their prayers. He who notes the fall of a sparrow surely hears the pleadings of our hearts. Remember the promise: “If any of you lack wisdom, let him ask of God, that giveth to all men liberally, and upbraideth not; and it shall be given him.”

To those within the sound of my voice who are struggling with challenges and difficulties large and small, prayer is the provider of spiritual strength; it is the passport to peace. Prayer is the means by which we approach our Father in Heaven, who loves us. Speak to Him in prayer and then listen for the answer. Miracles are wrought through prayer.

Sister Daisy Ogando lives in New York City, home to more than eight million people. Some years ago Sister Ogando met with the missionaries and was taught the gospel. Gradually, she and the missionaries lost contact. Time passed. Then, in 2007, the principles of the gospel she had been taught by the missionaries stirred within her heart.

One day while getting into a taxi, Daisy saw the missionaries at a distance, but she was unable to make contact with them before they disappeared from view. She prayed fervently to our Heavenly Father and promised Him that if He would somehow direct the missionaries to her once again, she would open her door to them. She returned home that day with faith in her heart that God would hear and answer her prayer.

In the meantime, two young missionaries who had been sincerely praying and working to find people to teach were one day examining the tracting records of missionaries who had previously served in their area. As they did so, they came across the name of Daisy Ogando. When they approached her apartment the very afternoon that Sister Ogando offered that simple but fervent prayer, she opened the door and said those words that are music to every missionary who has ever heard them: “Elders, come in. I’ve been waiting for you!”

Two fervent prayers were answered, contact was reestablished, missionary lessons were taught, and arrangements were made for Daisy and her son Eddy to be baptized.

Remember to pray fervently.

My final suggestion, my brethren: live righteously. Isaiah, that great prophet of the Old Testament, gave this stirring charge to holders of the priesthood: “Touch no unclean thing; … be ye clean, that bear the vessels of the Lord.” That’s about as straight as it could be given.

Holders of the priesthood may not necessarily be eloquent in their speech. They may not hold advanced degrees in difficult fields of study. They may very well be men of humble means. But God is no respecter of persons, and He will sustain His servants in righteousness as they avoid the evils of our day and live lives of virtue and purity. May I illustrate.

Some 900 miles (1,400 km) north of Salt Lake City is the beautiful city of Calgary, Alberta, Canada, home of the famous Calgary Stampede, one of Canada’s largest annual events and the world’s largest outdoor rodeo. The 10-day event features a rodeo competition, exhibits, agricultural competitions, and chuck wagon races. The Stampede Parade, which occurs on opening day, is one of the festival’s oldest and largest traditions. The parade follows a nearly three-mile (5-km) route in downtown Calgary, with attendance reaching 350,000 spectators, many dressed in western attire.

Several years ago, a marching band from a large high school in Utah had auditioned for and had received one of the coveted entries to march in the Calgary Stampede Parade. Months of fund-raising, early-morning practices up and down the streets, and other preparations were undertaken in order for the band to travel to Calgary and participate in the parade, where one band would be selected to receive the first-place honor.

Finally the day for departure arrived, with the eager students and their leaders boarding the buses and heading north for the long journey to Calgary.

While en route, the caravan stopped in Cardston, Alberta, Canada, where the group remained for an overnight stay. The local Relief Society sisters there prepared sack lunches for the band members to enjoy before departing again. Brad, one of the band members, who was a priest in the Aaronic Priesthood, was not hungry and decided to keep his lunch until later.

Brad liked to sit in the back of the bus. As he took his usual seat there in preparation for the remainder of the journey to Calgary, he tossed his sack lunch on the shelf behind the last row of seats. There the lunch sat by the rear window as the July afternoon sun shone through. Unfortunately, the sack lunch contained an egg salad sandwich. For those of you who don’t understand the significance of this, may I just say that egg salad must be refrigerated. If it is not, and if it is subjected to high heat such as that which would be produced by the sun beating through a bus window on a sunny day, it becomes a rather efficient incubator for various strains of bacteria that can result in what may commonly be referred to as food poisoning.

Sometime before arriving in Calgary, Brad grew hungry. Remembering the sack lunch, he gulped down the egg salad sandwich. As the buses arrived in Calgary and drove around the city, the members of the band grew excited—all except for Brad. Unfortunately, all that grew within him were severe stomach pains and other discomforts associated with food poisoning. You know what they are.

Upon arriving at their destination, the band members exited the bus. Brad, however, did not. Although he knew his fellow band members were counting on him to play his drum in the parade the following morning, Brad was doubled over in pain and was too sick to leave the bus. Providentially for him, two of his friends, Steve and Mike, who had recently graduated from high school and who had also recently been ordained to the office of elder in the Melchizedek Priesthood, found that Brad was missing and decided to look for him.

Finding Brad in the rear of the bus and learning what the problem was, Steve and Mike felt helpless. Finally it occurred to them that they were elders and held the power of the Melchizedek Priesthood to bless the sick. Despite their total lack of experience in giving a priesthood blessing, these two new elders had faith in the power they held. They laid their hands on Brad’s head and, invoking the authority of the Melchizedek Priesthood, in the name of Jesus Christ uttered the simple words to bless Brad to be made well.

From that moment, Brad’s symptoms were completely gone. The next morning he took his place with the rest of the band members and proudly marched down the streets of Calgary. The band received first-place honors and the coveted blue ribbon. Far more important, however, was that two young, inexperienced but worthy priesthood holders had answered the call to represent the Lord in serving their fellow man. When it was necessary for them to exercise their priesthood in behalf of one who was desperately in need of their help, they were able to respond because they lived their lives righteously.

Brethren, are we prepared for our journey through life? The pathway can at times be difficult. Chart your course, be cautious, and determine to study diligently, pray fervently, and live righteously.

Let us never despair, for the work in which we are engaged is the work of the Lord. It has been said, “The Lord shapes the back to bear the burden placed upon it.”

The strength which we earnestly seek in order to meet the challenges of a complex and changing world can be ours when, with fortitude and resolute courage, we stand and declare with Joshua, “As for me and my house, we will serve the Lord.” To this divine truth I testify and do so in the name of Jesus Christ, our Lord, amen.

\newpage
\settalk{May You Have Courage}
\setspeaker{Thomas S. Monson}
\settalkdate{April 2009}
\talkheading{May You Have Courage}{Thomas S. Monson}

My dear young sisters, what a glorious sight you are. I realize that beyond this magnificent Conference Center many thousands are assembled in chapels and in other settings throughout much of the world. I pray for heavenly help as I respond to the opportunity to address you.

We have heard timely, inspiring messages from your general Young Women leaders. These are choice women, called and set apart to guide and teach you. They love you, as do I.

You have come to this earth at a glorious time. The opportunities before you are nearly limitless. Almost all of you live in comfortable homes, with loving families, adequate food, and sufficient clothing. In addition, most of you have access to amazing technological advances. You communicate through cell phones, text messaging, instant messaging, e-mailing, blogging, Facebook, and other such means. You listen to music on your iPods and MP3 players. This list, of course, represents but a few of the technologies which are available to you.

All of this is a little daunting to someone such as I, who grew up when radios were generally large floor models and when there were no televisions to speak of, let alone computers or cell phones. In fact, when I was your age, telephone lines were mostly shared. In our family, if we wanted to make a telephone call, we would have to pick up the phone and listen first to make certain no other family was using the line, for several families shared one line.

I could go on all night talking about the differences between my generation and yours. Suffice it to say that much has changed between the time I was your age and the present.

Although this is a remarkable period when opportunities abound, you also face challenges which are unique to this time. For instance, the very technological tools I have mentioned provide opportunities for the adversary to tempt you and to ensnare you in his web of deceit, thereby hoping to take possession of your destiny.

As I contemplate all that you face in the world today, one word comes to my mind. It describes an attribute needed by all of us but one which you—at this time of your life and in this world—will need particularly. That attribute is courage.

Tonight I’d like to talk with you about the courage you will need in three aspects of your lives:

May I speak first about the courage to refrain from judging others. Oh, you may ask, “Does this really take courage?” And I would reply that I believe there are many times when refraining from judgment—or gossip or criticism, which are certainly akin to judgment—takes an act of courage.

Unfortunately, there are those who feel it necessary to criticize and to belittle others. You have, no doubt, been with such people, as you will be in the future. My dear young friends, we are not left to wonder what our behavior should be in such situations. In the Sermon on the Mount, the Savior declared, “Judge not.” At a later time He admonished, “Cease to find fault one with another.” It will take real courage when you are surrounded by your peers and feeling the pressure to participate in such criticisms and judgments to refrain from joining in.

I would venture to say that there are young women around you who, because of your unkind comments and criticism, are often left out. It seems to be the pattern, particularly at this time in your lives, to avoid or to be unkind to those who might be judged different, those who don’t fit the mold of what we or others think they should be.

The Savior said:

“A new commandment I give unto you, That ye love one another. …

“By this shall all men know that ye are my disciples, if ye have love one to another.”

Mother Teresa, a Catholic nun who worked among the poor in India most of her life, spoke this truth: “If you judge people, you have no time to love them.”

A friend told me of an experience she had many years ago when she was a teenager. In her ward was a young woman named Sandra who had suffered an injury at birth, resulting in her being somewhat mentally handicapped. Sandra longed to be included with the other girls, but she looked handicapped. She acted handicapped. Her clothing was always ill fitting. She sometimes made inappropriate comments. Although Sandra attended their Mutual activities, it was always the responsibility of the teacher to keep her company and to try to make her feel welcome and valued, since the girls did not.

Then something happened: a new girl of the same age moved into the ward. Nancy was a cute, redheaded, self-confident, popular girl who fit in easily. All the girls wanted to be her friend, but Nancy didn’t limit her friendships. In fact, she went out of her way to befriend Sandra and to make certain she always felt included in everything. Nancy seemed to genuinely like Sandra.

Of course the other girls took note and began wondering why they hadn’t ever befriended Sandra. It now seemed not only acceptable but desirable. Eventually they began to realize what Nancy, by her example, was teaching them: that Sandra was a valuable daughter of our Heavenly Father, that she had a contribution to make, and that she deserved to be treated with love and kindness and positive attention.

By the time Nancy and her family moved from the neighborhood a year or so later, Sandra was a permanent part of the group of young women. My friend said that from then on she and the other girls made certain no one was ever left out, regardless of what might make her different. A valuable, eternal lesson had been learned.

True love can alter human lives and change human nature.

My precious young sisters, I plead with you to have the courage to refrain from judging and criticizing those around you, as well as the courage to make certain everyone is included and feels loved and valued.

I turn next to the courage you will need to be chaste and virtuous. You live in a world where moral values have, in great measure, been tossed aside, where sin is flagrantly on display, and where temptations to stray from the strait and narrow path surround you. Many are the voices telling you that you are far too provincial or that there is something wrong with you if you still believe there is such a thing as immoral behavior.

Isaiah declared, “Woe unto them that call evil good, and good evil; that put darkness for light, and light for darkness.”

Great courage will be required as you remain chaste and virtuous amid the accepted thinking of the times.

In the world’s view today there is little thought that young men and young women will remain morally clean and pure before marriage. Does this make immoral behavior acceptable? Absolutely not!

The commandments of our Heavenly Father are not negotiable!

Powerful is this quote from news commentator Ted Koppel, host of ABC’s Nightline program for many years. Said he:

“We have actually convinced ourselves that slogans will save us. ‘Shoot up if you must; but use a clean needle.’ ‘Enjoy sex whenever with whomever you wish; but [protect yourself].’

“No. The answer is no. Not no because it isn’t cool or smart or because you might end up in jail or dying in an AIDS ward—but no, because it’s wrong. …

“What Moses brought down from Mt. Sinai were not the Ten Suggestions, they are Commandments. Are, not were.”

My sweet young sisters, maintain an eternal perspective. Be alert to anything that would rob you of the blessings of eternity.

Help in maintaining the proper perspective in these permissive times can come to you from many sources. One valuable resource is your patriarchal blessing. Read it frequently. Study it carefully. Be guided by its cautions. Live to merit its promises. If you have not yet received your patriarchal blessing, plan for the time when you will receive it, and then cherish it.

If any has stumbled in her journey, there is a way back. The process is called repentance. Our Savior died to provide you and me that blessed gift. The path may be difficult, but the promise is real: “Though your sins be as scarlet, they shall be as white as snow.” “And I will remember [them] no more.”

Some years ago another First Presidency made this statement, and your First Presidency today echoes the appeal. I quote: “To the youth … , we plead with you to live clean [lives], for the unclean life leads only to suffering, misery, and woe physically,—and spiritually it is the path to destruction. How glorious and near to the angels is youth that is clean; this youth has joy unspeakable here and eternal happiness hereafter. Sexual purity is youth’s most precious possession; it is the foundation of all righteousness.”

May you have the courage to be chaste and virtuous.

My final plea tonight is that you have the courage to stand firm for truth and righteousness. Because the trend in society today is away from the values and principles the Lord has given us, you will almost certainly be called upon to defend that which you believe. Unless the roots of your testimony are firmly planted, it will be difficult for you to withstand the ridicule of those who challenge your faith. When firmly planted, your testimony of the gospel, of the Savior, and of our Heavenly Father will influence all that you do throughout your life. The adversary would like nothing better than for you to allow derisive comments and criticism of the Church to cause you to question and doubt. Your testimony, when constantly nourished, will keep you safe.

Recall with me Lehi’s vision of the tree of life. He saw that many who had held to the iron rod and had made their way through the mists of darkness, arriving at last at the tree of life and partaking of the fruit of the tree, did then “cast their eyes about as if they were ashamed.” Lehi wondered as to the cause of their embarrassment. As he looked about, he “beheld, on the other side of the river of water, a great and spacious building. …

“And it was filled with people, both old and young, both male and female; and their manner of dress was exceedingly fine; and they were in the attitude of mocking and pointing their fingers towards those who … were partaking of the fruit.”

The great and spacious building in Lehi’s vision represents those in the world who mock God’s word and who ridicule those who embrace it and who love the Savior and live the commandments. What happens to those who are ashamed when the mocking occurs? Lehi tells us, “And after they had tasted of the fruit they were ashamed, because of those that were scoffing at them; and they fell away into forbidden paths and were lost.”

My beloved young sisters, with the courage of your convictions, may you declare with the Apostle Paul, “I am not ashamed of the gospel of Christ: for it is the power of God unto salvation.”

Lest you feel inadequate for the tasks which lie ahead, I remind you of another of the Apostle Paul’s stirring statements from which we might draw courage: “For God hath not given us the spirit of fear; but of power, and of love, and of a sound mind.”

In closing may I share with you the account of a brave young woman whose experience has stood through the ages as an example of the courage to stand for truth and righteousness.

Most of you are familiar with the Old Testament account of Esther. It is a very interesting and inspiring record of a beautiful young Jewish girl whose parents had died, leaving her to be raised by an older cousin, Mordecai, and his wife.

Mordecai worked for the king of Persia, and when the king was looking for a queen, Mordecai took Esther to the palace and presented her as a candidate, advising her not to reveal that she was Jewish. The king was pleased with Esther above all the others and made Esther his queen.

Haman, the chief prince in the king’s court, became increasingly angry with Mordecai because Mordecai would not bow down and pay homage to him. In retribution, Haman convinced the king—in a rather devious manner—that there were “certain people” in all 127 provinces of the kingdom whose laws were different from others’ and that they would not obey the king’s laws and should be destroyed. Without naming these people to the king, Haman was, of course, referring to the Jews, including Mordecai.

With the king’s permission to handle the matter, Haman sent letters to the governors of all of the provinces, instructing them “to destroy, to kill, and to cause to perish, all Jews, both young and old, little children and women, … [on] the thirteenth day of the twelfth month.”

Through a servant, Mordecai sent word to Esther concerning the decree against the Jews, requesting that she go in to the king to plead for her people. Esther was at first reluctant, reminding Mordecai that it was against the law for anyone to go unbidden into the inner court of the king. Punishment by death would be the result—unless the king were to hold out his golden scepter, allowing the person to live.

Mordecai’s response to Esther’s hesitation was to the point. He replied to her thus:

“Think not … that thou shalt escape in the king’s house, more than all the Jews.

“For if thou altogether holdest thy peace at this time, … thou and thy father’s house shall be destroyed.”

And then he added this searching question: “Who knoweth whether thou art come to the kingdom for such a time as this?”

In response, Esther asked Mordecai to gather all the Jews he could and to ask them to fast three days for her and said that she and her handmaids would do the same. She declared, “I [will] go in unto the king, which is not according to the law: and if I perish, I perish.” Esther had gathered her courage and would stand firm and immovable for that which was right.

Physically, emotionally, and spiritually prepared, Esther stood in the inner court of the king’s house. When the king saw her, he held out his golden scepter, telling her that he would grant whatever request she had. She invited the king to a feast she had arranged, and during the feast she revealed that she was a Jew. She also exposed Haman’s underhanded plot to exterminate all of the Jews in the kingdom. Esther’s plea to save herself and her people was granted.

Esther, through fasting, faith, and courage, had saved a nation.

You will probably not be called upon to put your life on the line, as did Esther, for that which you believe. You will, however, most likely find yourself in situations where great courage will be required as you stand firm for truth and righteousness.

Again, my dear young sisters, although there have always been challenges in the world, many of those which you face are unique to this time. But you are some of our Heavenly Father’s strongest children, and He has saved you to come to the earth “for such a time as this.” With His help, you will have the courage to face whatever comes. Though the world may at times appear dark, you have the light of the gospel, which will be as a beacon to guide your way.

My earnest prayer is that you will have the courage required to refrain from judging others, the courage to be chaste and virtuous, and the courage to stand firm for truth and righteousness. As you do so, you will be “an example of the believers,” and your life will be filled with love and peace and joy. May this be so, my beloved young sisters, I ask in the name of Jesus Christ, our Savior, amen.

\newpage
\settalk{Until We Meet Again}
\setspeaker{Thomas S. Monson}
\settalkdate{April 2009}
\talkheading{Until We Meet Again}{Thomas S. Monson}

My beloved brothers and sisters, my heart is full and my feelings tender as we conclude this great general conference.

We have been richly blessed as we have listened to the counsel and testimonies of those who have spoken to us. I believe we are all more determined to live the principles of the gospel of Jesus Christ.

I express my sincere thanks to each one who participated in the conference, including those Brethren who offered prayers.

The music has been magnificent. How grateful I am for those blessed with musical talents who are willing to share their talents with others. I am reminded of the scripture found in the Doctrine and Covenants: “For my soul delighteth in the song of the heart; yea, the song of the righteous is a prayer unto me, and it shall be answered with a blessing upon their heads.”

May we long remember that which we have heard during this conference. I remind you that the messages will be printed in next month’s Ensign and Liahona magazines. I urge you to study the messages and to ponder their teachings and then to apply them in your life.

I want you to know how much I love and appreciate my devoted counselors, President Henry B. Eyring and President Dieter F. Uchtdorf. They are men of wisdom and understanding. Their service is invaluable. I love and support my Brethren of the Quorum of the Twelve Apostles. During this conference we sustained a new member of that Quorum. He is completely dedicated to the work of the Lord, and I testify that he is the man our Heavenly Father wants to fill this position at this time.

I express my love to the members of the Seventy and the Presiding Bishopric. They serve selflessly and so effectively. Similarly, I pay tribute to the general auxiliary officers. In accordance with our policy of rotation, we have sustained new general presidencies of the Young Men and of the Sunday School. We look forward to working with them. We thank those who were released from these positions at this conference and who served so faithfully in these capacities.

My brothers and sisters, may we strive to live closer to the Lord. May we remember to “pray always lest [we] enter into temptation.”

To you parents, express your love to your children. Pray for them that they may be able to withstand the evils of the world. Pray that they may grow in faith and testimony. Pray that they may pursue lives of goodness and of service to others.

Children, let your parents know you love them. Let them know how much you appreciate all they have done and continue to do for you.

Now, a word of caution to all—both young and old, both male and female. We live at a time when the adversary is using every means possible to ensnare us in his web of deceit, trying desperately to take us down with him. There are many pathways along which he entices us to go—pathways that can lead to our destruction. Advances in many areas that can be used for good can also be used to speed us along those heinous pathways.

I feel to mention one in particular, and that is the Internet. On one hand, it provides nearly limitless opportunities for acquiring useful and important information. Through it we can communicate with others around the world. The Church itself has a wonderful Web site, filled with valuable and uplifting information and priceless resources.

On the other hand, however—and extremely alarming—are the reports of the number of individuals who are utilizing the Internet for evil and degrading purposes, the viewing of pornography being the most prevalent of these purposes. My brothers and sisters, involvement in such will literally destroy the spirit. Be strong. Be clean. Avoid such degrading and destructive types of content at all costs—wherever they may be! I sound this warning to everyone, everywhere. I add—particularly to the young people—that this includes pornographic images transmitted via cell phones.

My beloved friends, under no circumstances allow yourselves to become trapped in the viewing of pornography, one of the most effective of Satan’s enticements. And if you have allowed yourself to become involved in this behavior, cease now. Seek the help you need to overcome and to change the direction of your life. Take the steps necessary to get back on the strait and narrow, and then stay there.

May we say, with Joshua of old, “Choose you this day whom ye will serve; … but as for me and my house, we will serve the Lord.”

Now, my brothers and sisters, we have built temples throughout the world and will continue to do so. To you who are worthy and able to attend the temple, I would admonish you to go often. The temple is a place where we can find peace. There we receive a renewed dedication to the gospel and a strengthened resolve to keep the commandments.

What a privilege it is to be able to go to the temple, where we may experience the sanctifying influence of the Spirit of the Lord. Great service is given when we perform vicarious ordinances for those who have gone beyond the veil. In many cases we do not know those for whom we perform the work. We expect no thanks, nor do we have the assurance that they will accept that which we offer. However, we serve, and in that process we attain that which comes of no other effort: we literally become saviors on Mount Zion. As our Savior gave His life as a vicarious sacrifice for us, so we, in some small measure, do the same when we perform proxy work in the temple for those who have no means of moving forward unless something is done for them by those of us here on the earth.

I am deeply grateful that as a church we continue to extend humanitarian aid where there is great need. We have done much in this regard and have blessed the lives of thousands upon thousands of our Father’s children who are not of our faith as well as those who are. We intend to continue to help wherever such is needed. We express gratitude to you for your contributions in this regard.

How grateful I am, my brothers and sisters, for the Restoration of the gospel in this dispensation and for all the blessings that have come into my life and into your lives as a result. We are a blessed people, for we have the sure knowledge that God lives and that Jesus is the Christ.

May heaven’s blessings be with you. May your homes be filled with harmony and love. May you constantly nourish your testimonies that they might be a protection to you against the adversary.

As your humble servant, I desire with all my heart to do God’s will and to serve Him and to serve you.

Now, my brothers and sisters, conference is over. As we return to our homes, may we do so safely.

I love you. I pray for you. I would ask that you would remember me and all the General Authorities in your prayers. Until we meet again in six months’ time, I ask the Lord’s blessings to be upon all of us, and I do it in the name of Jesus Christ the Lord, our Savior, amen.

\newpage
\settalk{Welcome to Conference}
\setspeaker{Thomas S. Monson}
\settalkdate{April 2009}
\talkheading{Welcome to Conference}{Thomas S. Monson}

My dear brothers and sisters, as we open this, the 179th Annual General Conference, we note with sadness the absence of Elder Joseph B. Wirthlin of the Quorum of the Twelve Apostles. We mourn his passing. We miss him. We extend our love to his family. I have no doubt that he is carrying on this great work on the other side of the veil.

Because of the passing of Elder Wirthlin, there exists a vacancy in the Quorum of the Twelve Apostles. After much fasting and prayer, we have called Elder Neil Linden Andersen to fill this vacancy. We present his name to you this morning for your sustaining vote. All those of you who feel you can sustain him in this sacred calling will please signify by the uplifted hand. Any who may be opposed may signify by the same sign.

We thank you for your sustaining vote. Elder Andersen’s name will be included when the officers of the Church are read this afternoon.

Elder Andersen, we invite you now to take your place on the stand with the members of the Twelve. We look forward to hearing from you in the Sunday morning session of conference.

Since we met six months ago, my brothers and sisters, I have traveled to Mexico City, Mexico, with President and Sister Henry B. Eyring, to rededicate the temple there. For many months it had been undergoing extensive renovations.

The evening before the rededication, a magnificent cultural event was held in the Aztec Stadium. Approximately 87,000 spectators squeezed into the open-air stadium, and a cast of more than 8,000 young people participated in the program, which featured an 80-minute display of music, dance, and Mexican history.

President Eyring and I were each presented a serape and a sombrero. Outfitted in this native costume, I couldn’t resist serenading the group with an impromptu version of “El Rancho Grande,” which I had originally learned in my ninth-grade Spanish class. I shall not do that today.

Each of the two dedicatory sessions the following day was filled with the Spirit of the Lord.

Just two weeks ago, in 12 sessions we dedicated the Draper Utah Temple, a magnificent structure nestled in the foothills of the mountains in the south portion of the Salt Lake Valley. There were approximately 685,000 people—members and nonmembers alike—who attended the open house. Over 365,000 members were present at the dedicatory sessions, including the sessions broadcast by satellite to various stake centers. The Spirit of the Lord was present in rich abundance as the temple was dedicated.

In the near future, we will be dedicating the Oquirrh Mountain Utah Temple, and then in the coming months and years there will be many more dedications. We look forward to these occasions. There is something about a temple dedication which prompts a reevaluation of one’s own performance and a sincere desire to do even better.

Now, my brothers and sisters, I am pleased to report that the Church is doing very well. The work of the Lord continues to move forward uninterrupted.

We now have approximately 53,000 missionaries serving in 348 missions throughout the world. We take most seriously the Savior’s mandate when He said, “Go ye therefore, and teach all nations, baptizing them in the name of the Father, and of the Son, and of the Holy Ghost.” We are deeply grateful for the labors of our missionaries and for the sacrifices which they and their families make in order for them to serve.

We also have countless volunteers and missionaries in nonproselyting activities. These are generally mature individuals who donate their time and talents in order to further the work of the Lord and to bless our Heavenly Father’s children. How thankful we are for the valuable services these individuals are providing.

The Perpetual Education Fund, established in 2001, continues to move forward. Since its inception, 35,600 young men and young women have been enrolled in the program and have trained to improve their skills and their employment opportunities. Thus far, 18,900 have finished that training. On average, with the 2.7 years of education they are now receiving, they are increasing their income by three to four times. What a blessing this is in their lives! This is indeed an inspired program.

My brothers and sisters, I thank you for your faith and devotion to the gospel. I thank you for the love and care you show to one another. I thank you for the service you provide in your wards and branches and in your stakes and districts. It is such service that enables the Lord to accomplish His purposes here upon the earth.

I express my thanks to you for your kindnesses to me wherever I go. I thank you for your prayers in my behalf. I have felt those prayers and am most grateful for them.

Now, my brothers and sisters, we are anxious to listen to the messages which will be presented to us during the next two days, that we might be taught and inspired and have a renewed determination to live the gospel and to serve the Lord. Those who will address us have sought heaven’s help and direction as they have prepared their messages. They have been impressed concerning that which they will share with us.

To those of you who are new in the Church, we welcome you. To those of you who are struggling with challenges or with disappointments or with losses, we pray for you. Our Heavenly Father loves each of us and is mindful of our needs. May we be filled with His Spirit as we listen to that which will be presented. Such is my prayer this morning as we open this great conference. I also add a fond remembrance of President Gordon B. Hinckley, who preceded me as President of the Church. I’m sure he’s serving well on the other side. In the name of our Lord and Savior, Jesus Christ, amen.

\newpage
\settalk{Closing Remarks}
\setspeaker{Thomas S. Monson}
\settalkdate{October 2009}
\talkheading{Closing Remarks}{Thomas S. Monson}

My heart is full as we come to the close of this conference. We have been richly taught and spiritually edified as we have listened to the messages which have been presented and the testimonies which have been borne. We express thanks to each one who has participated, including those Brethren offering prayers.

Once again the music has been wonderful. I express my personal gratitude for those willing to share with us their talents, touching and inspiring us in the process. The beautiful music they produce enhances and enriches each session of conference.

We remind you that the messages we have heard during this conference will be printed in the November issues of the Ensign and Liahona magazines. As we read and study them, we will be additionally taught and inspired. May we incorporate into our daily lives the truths found therein.

We express to those Brethren who have been released during this conference our deep appreciation. They have served well and have made significant contributions to the work of the Lord. Their dedication has been complete. We thank them from the bottom of our hearts.

We live at a time when many in the world have slipped from the moorings of safety found in compliance with the commandments. It is a time of permissiveness, with society in general routinely disregarding and breaking the laws of God. We often find ourselves swimming against the current, and sometimes it seems as though the current could carry us away.

I am reminded of the words of the Lord found in the book of Ether in the Book of Mormon. Said the Lord, “Ye cannot cross this great deep save I prepare you against the waves of the sea, and the winds which have gone forth, and the floods which shall come.” My brothers and sisters, He has prepared us. If we heed His words and live the commandments, we will survive this time of permissiveness and wickedness—a time which can be compared with the waves and the winds and the floods that can destroy. He is ever mindful of us. He loves us and will bless us as we do what is right.

How grateful we are that the heavens are indeed open, that the gospel of Jesus Christ has been restored, and that the Church is founded on the rock of revelation. We are a blessed people, with apostles and prophets upon the earth today.

Now, as we leave this conference, I invoke the blessings of heaven upon each of you. May all of you return safely to your homes. As you ponder the things you have heard during this conference, may you say, with the people of King Benjamin who all cried with one voice, “We believe all the words which thou hast spoken unto us; and also, we know of their surety and truth, because of the Spirit of the Lord Omnipotent, which has wrought a mighty change in us … that we have no more disposition to do evil, but to do good continually.” May every man and woman, boy and girl leave this conference a better person than he or she was when it began two days ago.

I love you, my brothers and sisters. I pray for you. I would ask once again that you would remember me and all the General Authorities in your prayers. We are one with you in moving forward this marvelous work. I testify to you that we are all in this together and that every man, woman, and child has a part to play. May God give us the strength and the ability and the determination to play our part well.

I bear my testimony to you that this work is true, that our Savior lives, and that He guides and directs His Church here upon the earth. I leave with you my witness and my testimony that God our Eternal Father lives and loves us. He is indeed our Father, and He is personal and real.

May God bless you. May His promised peace be with you now and always.

I bid you farewell until we meet again in six months’ time, and do so in the name of Jesus Christ, our Savior and Redeemer and our Advocate with the Father, amen.

\newpage
\settalk{School Thy Feelings, O My Brother}
\setspeaker{Thomas S. Monson}
\settalkdate{October 2009}
\talkheading{School Thy Feelings, O My Brother}{Thomas S. Monson}

Brethren, we are assembled as a mighty body of the priesthood, both here in the Conference Center and in locations throughout the world. We have heard inspired messages this evening, and I express my appreciation to those Brethren who have addressed us. I am honored, yet humbled, by the privilege to speak to you, and I pray that the inspiration of the Lord may attend me.

Recently as I watched the news on television, I realized that many of the lead stories were similar in nature in that the tragedies reported all basically traced back to one emotion: anger. The father of an infant had been arrested for physical abuse of the baby. It was alleged that the baby’s crying had so infuriated him that he had broken one of the child’s limbs and several ribs. Alarming was the report of growing gang violence, with the number of gang-related killings having risen sharply. Another story that night involved the shooting of a woman by her estranged husband, who was reportedly in a jealous rage after finding her with another man. Then, of course, there was the usual coverage of wars and conflicts throughout the world.

I thought of the words of the Psalmist: “Cease from anger, and forsake wrath.”

Many years ago, a young couple called my office and asked if they could come in for counseling. They indicated they had suffered a tragedy in their lives and that their marriage was in serious jeopardy. An appointment was arranged.

The tension between this husband and wife was apparent as they entered my office. Their story unfolded slowly at first as the husband spoke haltingly and the wife cried quietly and participated very little in the conversation.

The young man had returned from serving a mission and was accepted to a prestigious university in the eastern part of the United States. It was there, in a university ward, that he had met his future wife. She was also a student at the university. After a year of dating, they journeyed to Utah and were married in the Salt Lake Temple, returning east shortly afterward to finish their schooling.

By the time they graduated and returned to their home state, they were expecting their first child and the husband had employment in his chosen field. The wife gave birth to a baby boy. Life was good.

When their son was about 18 months old, they decided to take a short vacation to visit family members who lived a few hundred miles away. This was at a time when car seats for children and seat belts for adults were scarcely heard of, let alone used. The three members of the family all rode in the front seat with the toddler in the middle.

Sometime during the trip, the husband and wife had a disagreement. After all these years, I cannot recall what caused it. But I do remember that their argument escalated and became so heated that they were eventually yelling at one another. Understandably, this caused their young son to begin crying, which the husband said only added to his anger. Losing total control of his temper, he picked up a toy the child had dropped on the seat and flung it in the direction of his wife.

He missed hitting his wife. Instead, the toy struck their son, with the result that he was brain damaged and would be handicapped for the rest of his life.

This was one of the most tragic situations I had ever encountered. I counseled and encouraged them. We talked of commitment and responsibility, of acceptance and forgiveness. We spoke of the affection and respect which needed to return to their family. We read words of comfort from the scriptures. We prayed together. Though I have not heard from them since that day so long ago, they were smiling through their tears as they left my office. All these years I’ve hoped they made the decision to remain together, comforted and blessed by the gospel of Jesus Christ.

I think of them whenever I read the words: “Anger doesn’t solve anything. It builds nothing, but it can destroy everything.”

We’ve all felt anger. It can come when things don’t turn out the way we want. It might be a reaction to something which is said of us or to us. We may experience it when people don’t behave the way we want them to behave. Perhaps it comes when we have to wait for something longer than we expected. We might feel angry when others can’t see things from our perspective. There seem to be countless possible reasons for anger.

There are times when we can become upset at imagined hurts or perceived injustices. President Heber J. Grant, seventh President of the Church, told of a time as a young adult when he did some work for a man who then sent him a check for \$500 with a letter apologizing for not being able to pay him more. Then President Grant did some work for another man—work which he said was 10 times more difficult, involving 10 times more labor and a great deal more time. This second man sent him a check for \$150. Young Heber felt he had been treated most unfairly. He was at first insulted and then incensed.

He recounted the experience to an older friend, who asked, “Did that man intend to insult you?”

President Grant replied, “No. He told my friends he had rewarded me handsomely.”

To this the older friend replied, “A man’s a fool who takes an insult that isn’t intended.”

The Apostle Paul asks in Ephesians, chapter 4, verse 26 of the Joseph Smith Translation: “Can ye be angry, and not sin? let not the sun go down upon your wrath.” I ask, is it possible to feel the Spirit of our Heavenly Father when we are angry? I know of no instance where such would be the case.

From 3 Nephi in the Book of Mormon, we read:

“There shall be no disputations among you. …

“For verily, verily I say unto you, he that hath the spirit of contention is not of me, but is of the devil, who is the father of contention, and he stirreth up the hearts of men to contend with anger, one with another.

“Behold, this is not my doctrine, to stir up the hearts of men with anger, one against another; but this is my doctrine, that such things should be done away.”

To be angry is to yield to the influence of Satan. No one can make us angry. It is our choice. If we desire to have a proper spirit with us at all times, we must choose to refrain from becoming angry. I testify that such is possible.

Anger, Satan’s tool, is destructive in so many ways.

I believe most of us are familiar with the sad account of Thomas B. Marsh and his wife, Elizabeth. Brother Marsh was one of the first modern-day Apostles called after the Church was restored to the earth. He eventually became President of the Quorum of the Twelve Apostles.

While the Saints were in Far West, Missouri, Elizabeth Marsh, Thomas’s wife, and her friend Sister Harris decided they would exchange milk in order to make more cheese than they otherwise could. To be certain all was done fairly, they agreed that they should not save what were called the strippings, but that the milk and strippings should all go together. Strippings came at the end of the milking and were richer in cream.

Sister Harris was faithful to the agreement, but Sister Marsh, desiring to make some especially delicious cheese, saved a pint of strippings from each cow and sent Sister Harris the milk without the strippings. This caused the two women to quarrel. When they could not settle their differences, the matter was referred to the home teachers to settle. They found Elizabeth Marsh guilty of failure to keep her agreement. She and her husband were upset with the decision, and the matter was then referred to the bishop for a Church trial. The bishop’s court decided that the strippings were wrongfully saved and that Sister Marsh had violated her covenant with Sister Harris.

Thomas Marsh appealed to the high council, and the men comprising this council confirmed the bishop’s decision. He then appealed to the First Presidency of the Church. Joseph Smith and his counselors considered the case and upheld the decision of the high council.

Elder Thomas B. Marsh, who sided with his wife through all of this, became angrier with each successive decision—so angry, in fact, that he went before a magistrate and swore that the Mormons were hostile toward the state of Missouri. His affidavit led to—or at least was a factor in—Governor Lilburn Boggs’s cruel extermination order, which resulted in over 15,000 Saints being driven from their homes, with all the terrible suffering and consequent death that followed. All of this occurred because of a disagreement over the exchange of milk and cream.

After 19 years of rancor and loss, Thomas B. Marsh made his way to the Salt Lake Valley and asked President Brigham Young for forgiveness. Brother Marsh also wrote to Heber C. Kimball, First Counselor in the First Presidency, of the lesson he had learned. Said Brother Marsh: “The Lord could get along very well without me and He … lost nothing by my falling out of the ranks; But O what have I lost?! Riches, greater riches than all this world or many planets like this could afford.”

Apropos are the words of the poet John Greenleaf Whittier: “Of all sad words of tongue or pen, the saddest are these: ‘It might have been!’”

My brethren, we are all susceptible to those feelings which, if left unchecked, can lead to anger. We experience displeasure or irritation or antagonism, and if we so choose, we lose our temper and become angry with others. Ironically, those others are often members of our own families—the people we really love the most.

Many years ago I read the following Associated Press dispatch which appeared in the newspaper: An elderly man disclosed at the funeral of his brother, with whom he had shared, from early manhood, a small, one-room cabin near Canisteo, New York, that following a quarrel, they had divided the room in half with a chalk line, and neither had crossed the line or spoken a word to the other since that day—62 years before. Just think of the consequence of that anger. What a tragedy!

May we make a conscious decision, each time such a decision must be made, to refrain from anger and to leave unsaid the harsh and hurtful things we may be tempted to say.

I love the words of the hymn written by Elder Charles W. Penrose, who served in the Quorum of the Twelve and in the First Presidency during the early years of the 20th century:

\begin{verse}
School thy feelings, O my brother;\\
Train thy warm, impulsive soul.\\
Do not its emotions smother,\\
But let wisdom’s voice control.\\
School thy feelings; there is power\\
In the cool, collected mind.\\
Passion shatters reason’s tower,\\
Makes the clearest vision blind.
\end{verse}

Each of us is a holder of the priesthood of God. The oath and covenant of the priesthood pertains to all of us. To those who hold the Melchizedek Priesthood, it is a declaration of our requirement to be faithful and obedient to the laws of God and to magnify the callings which come to us. To those who hold the Aaronic Priesthood, it is a pronouncement concerning future duty and responsibility, that you may prepare yourselves here and now.

This oath and covenant is set forth by the Lord in these words:

“For whoso is faithful unto the obtaining these two priesthoods of which I have spoken, and the magnifying their calling, are sanctified by the Spirit unto the renewing of their bodies.

“They become the sons of Moses and of Aaron and the seed of Abraham, and the church and kingdom, and the elect of God.

“And also all they who receive this priesthood receive me, saith the Lord;

“For he that receiveth my servants receiveth me;

“And he that receiveth me receiveth my Father;

“And he that receiveth my Father receiveth my Father’s kingdom; therefore all that my Father hath shall be given unto him.”

Brethren, great promises await us if we are true and faithful to the oath and covenant of this precious priesthood which we hold. May we be worthy sons of our Heavenly Father. May we ever be exemplary in our homes and faithful in keeping all of the commandments, that we may harbor no animosity toward any man but rather be peacemakers, ever remembering the Savior’s admonition, “By this shall all men know that ye are my disciples, if ye have love one to another.” This is my plea tonight at the conclusion of this great priesthood meeting, and it’s also my humble and sincere prayer, for I love you, brethren, with all my heart and soul. And I pray our Heavenly Father’s blessing to attend each of you in your life, in your home, in your heart, in your soul, in the name of Jesus Christ, amen.

\newpage
\settalk{Welcome to Conference}
\setspeaker{Thomas S. Monson}
\settalkdate{October 2009}
\talkheading{Welcome to Conference}{Thomas S. Monson}

My beloved brothers and sisters, I extend my greetings to all of you as we commence this, the 179th Semiannual General Conference of The Church of Jesus Christ of Latter-day Saints.

How grateful I am for the age in which we live—an age of such advanced technology that we are able to address you across the world. As the General Authorities and auxiliary leaders stand here in the Conference Center in Salt Lake City, our voices will be reaching you by various means, including radio, television, satellite transmission, and the Internet. Although we will be speaking to you in English, you will be hearing us in some 92 languages.

Since last we met in April of this year, we have dedicated the beautiful Oquirrh Mountain Utah Temple in South Jordan, Utah. Sandwiched between the Draper Utah Temple dedication in March and this most recent dedication of the Oquirrh Mountain Utah Temple in August, a spectacular two-night cultural event was held, featuring youth from both temple districts. The productions retraced the rich legacy of Utah through song and dance. All told, approximately 14,000 youth participated over the two nights.

We continue to build temples. We desire that as many members as possible have an opportunity to attend the temple without having to travel inordinate distances. Worldwide, 83 percent of our members live within 200 miles (320 km) of a temple. That percentage will continue to increase as we construct new temples around the world. Currently there are 130 temples in operation, with 16 announced or under construction. This morning I am pleased to announce 5 additional temples for which sites are being acquired and which, in coming months and years, will be built in the following locations: Brigham City, Utah; Concepción, Chile; Fortaleza, Brazil; Fort Lauderdale, Florida; and Sapporo, Japan.

Millions of ordinances are performed in the temples each year in behalf of our deceased loved ones. May we continue to be faithful in performing such ordinances for those who are unable to do so for themselves. I love the words of President Joseph F. Smith as he spoke of temple service and of the spirit world beyond mortality. Said he, “Through our efforts in their behalf their chains of bondage will fall from them, and the darkness surrounding them will clear away, that light may shine upon them and they shall hear in the spirit world of the work that has been done for them by their [people] here, and will rejoice with you in your performance of these duties.”

Brothers and sisters, the Church continues to grow, as it has since being organized over 179 years ago. It is changing the lives of more and more people every year and is spreading far and wide over the earth as our missionary force seeks out those who are looking for the truths which are found in the gospel of Jesus Christ. We call upon all members of the Church to befriend the new converts, to reach out to them, to surround them with love, and to help them feel at home.

I would ask that your faith and prayers continue to be offered in behalf of those areas where our influence is limited and where we are not allowed to share the gospel freely at this time. Miracles can occur as we do so.

Now, my brothers and sisters, we are anxious to listen to the messages which will be presented to us during the next two days. Those who will address us have sought heaven’s help and direction as they have prepared their messages. They have been impressed concerning that which they will share with us. That we may be filled with the Spirit of the Lord as we listen and learn is my prayer in the name of Jesus Christ, amen.

\newpage
\settalk{What Have I Done for Someone Today?}
\setspeaker{Thomas S. Monson}
\settalkdate{October 2009}
\talkheading{What Have I Done for Someone Today?}{Thomas S. Monson}

My beloved brothers and sisters, I greet you this morning with love in my heart for the gospel of Jesus Christ and for each of you. I am grateful for the privilege to stand before you, and I pray that I might effectively communicate to you that which I have felt prompted to say.

A few years ago I read an article written by Jack McConnell, MD. He grew up in the hills of southwest Virginia in the United States as one of seven children of a Methodist minister and a stay-at-home mother. Their circumstances were very humble. He recounted that during his childhood, every day as the family sat around the dinner table, his father would ask each one in turn, “And what did you do for someone today?” The children were determined to do a good turn every day so they could report to their father that they had helped someone. Dr. McConnell calls this exercise his father’s most valuable legacy, for that expectation and those words inspired him and his siblings to help others throughout their lives. As they grew and matured, their motivation for providing service changed to an inner desire to help others.

Besides Dr. McConnell’s distinguished medical career—where he directed the development of the tuberculosis tine test, participated in the early development of the polio vaccine, supervised the development of Tylenol, and was instrumental in developing the magnetic resonance imaging procedure, or MRI—he created an organization he calls Volunteers in Medicine, which gives retired medical personnel a chance to volunteer at free clinics serving the working uninsured. Dr. McConnell said his leisure time since he retired has “evaporated into 60-hour weeks of unpaid work, but [his] energy level has increased and there is a satisfaction in [his] life that wasn’t there before.” He made this statement: “In one of those paradoxes of life, I have benefited more from Volunteers in Medicine than my patients have.” There are now over 70 such clinics across the United States.

Of course, we can’t all be Dr. McConnells, establishing medical clinics to help the poor; however, the needs of others are ever present, and each of us can do something to help someone.

The Apostle Paul admonished, “By love serve one another.” Recall with me the familiar words of King Benjamin in the Book of Mormon: “When ye are in the service of your fellow beings ye are only in the service of your God.”

The Savior taught His disciples, “For whosoever will save his life shall lose it: but whosoever will lose his life for my sake, the same shall save it.”

I believe the Savior is telling us that unless we lose ourselves in service to others, there is little purpose to our own lives. Those who live only for themselves eventually shrivel up and figuratively lose their lives, while those who lose themselves in service to others grow and flourish—and in effect save their lives.

In the October 1963 general conference—the conference at which I was sustained as a member of the Quorum of the Twelve Apostles—President David O. McKay made this statement: “Man’s greatest happiness comes from losing himself for the good of others.”

Often we live side by side but do not communicate heart to heart. There are those within the sphere of our own influence who, with outstretched hands, cry out, “Is there no balm in Gilead?”

I am confident it is the intention of each member of the Church to serve and to help those in need. At baptism we covenanted to “bear one another’s burdens, that they may be light.” How many times has your heart been touched as you have witnessed the need of another? How often have you intended to be the one to help? And yet how often has day-to-day living interfered and you’ve left it for others to help, feeling that “oh, surely someone will take care of that need”?

We become so caught up in the busyness of our lives. Were we to step back, however, and take a good look at what we’re doing, we may find that we have immersed ourselves in the “thick of thin things.” In other words, too often we spend most of our time taking care of the things which do not really matter much at all in the grand scheme of things, neglecting those more important causes.

Many years ago I heard a poem which has stayed with me, by which I have tried to guide my life. It’s one of my favorites:

\begin{verse}
I have wept in the night\\
For the shortness of sight\\
That to somebody’s need made me blind;\\
But I never have yet\\
Felt a tinge of regret\\
For being a little too kind.
\end{verse}

My brothers and sisters, we are surrounded by those in need of our attention, our encouragement, our support, our comfort, our kindness—be they family members, friends, acquaintances, or strangers. We are the Lord’s hands here upon the earth, with the mandate to serve and to lift His children. He is dependent upon each of us.

You may lament: I can barely make it through each day, doing all that I need to do. How can I provide service for others? What can I possibly do?

Just over a year ago, I was interviewed by the Church News prior to my birthday. At the conclusion of the interview, the reporter asked what I would consider the ideal gift that members worldwide could give to me. I replied, “Find someone who is having a hard time or is ill or lonely, and do something for him or her.”

I was overwhelmed when this year for my birthday I received hundreds of cards and letters from members of the Church around the world telling me how they had fulfilled that birthday wish. The acts of service ranged from assembling humanitarian kits to doing yard work.

Dozens and dozens of Primaries challenged the children to provide service, and then those acts of service were recorded and sent to me. I must say that the methods for recording them were creative. Many came in the form of pages put together into various shapes and sizes of books. Some contained cards or pictures drawn or colored by the children. One very creative Primary sent a large jar containing hundreds of what they called “warm fuzzies,” each one representing an act of service performed during the year by one of the children in the Primary. I can only imagine the happiness these children experienced as they told of their service and then placed a “warm fuzzy” in the jar.

I share with you just a few of the countless notes contained in the many gifts I received. One small child wrote, “My grandpa had a stroke, and I held his hand.” From an 8-year-old girl: “My sister and I served my mom and family by organizing and cleaning the toy closet. It took us a few hours and we had fun. The best part was that we surprised my mom and made her happy because she didn’t even ask us to do it.” An 11-year-old girl wrote: “There was a family in my ward that did not have a lot of money. They have three little girls. The mom and dad had to go somewhere, so I offered to watch the three girls. The dad was just about to hand me a \$5 bill. I said, ‘I can’t take [it].’ My service was that I watched the girls for free.” A Primary child in Mongolia wrote that he had brought in water from the well so his mother would not have to do so. From a 4-year-old boy, no doubt written by a Primary teacher: “My dad is gone for army training for a few weeks. My special job is to give my mom hugs and kisses.” Wrote a 9-year-old girl: “I picked strawberries for my great-grandma. I felt good inside!” And another: “I played with a lonely kid.”

From an 11-year-old boy: “I went to a lady’s house and asked her questions and sang her a song. It felt good to visit her. She was happy because she never gets visitors.” Reading this particular note reminded me of words penned long ago by Elder Richard L. Evans of the Quorum of the Twelve. Said he: “It is difficult for those who are young to understand the loneliness that comes when life changes from a time of preparation and performance to a time of putting things away. … To be so long the center of a home, so much sought after, and then, almost suddenly to be on the sidelines watching the procession pass by—this is living into loneliness. … We have to live a long time to learn how empty a room can be that is filled only with furniture. It takes someone … beyond mere hired service, beyond institutional care or professional duty, to thaw out the memories of the past and keep them warmly living in the present. … We cannot bring them back the morning hours of youth. But we can help them live in the warm glow of a sunset made more beautiful by our thoughtfulness … and unfeigned love.”

My birthday cards and notes came also from teenagers in Young Men and Young Women classes who made blankets for hospitals, served in food pantries, were baptized for the dead, and performed numerous other acts of service.

Relief Societies, where help can always be found, provided service above and beyond that which they would normally have given. Priesthood groups did the same.

My brothers and sisters, my heart has seldom been as touched and grateful as it was when Sister Monson and I literally spent hours reading of these gifts. My heart is full now as I speak of the experience and contemplate the lives which have been blessed as a result, for both the giver and the receiver.

The words from the 25th chapter of Matthew come to mind:

“Come, ye blessed of my Father, inherit the kingdom prepared for you from the foundation of the world:

“For I was an hungred, and ye gave me meat: I was thirsty, and ye gave me drink: I was a stranger, and ye took me in:

“Naked, and ye clothed me: I was sick, and ye visited me: I was in prison, and ye came unto me.

“Then shall the righteous answer him, saying, Lord, when saw we thee an hungred, and fed thee? or thirsty, and gave thee drink?

“When saw we thee a stranger, and took thee in? or naked, and clothed thee?

“Or when saw we thee sick, or in prison, and came unto thee?

“And the King shall answer and say unto them, Verily I say unto you, Inasmuch as ye have done it unto one of the least of these my brethren, ye have done it unto me.”

My brothers and sisters, may we ask ourselves the question which greeted Dr. Jack McConnell and his brothers and sisters each evening at dinnertime: “What have I done for someone today?” May the words of a familiar hymn penetrate our very souls and find lodgment in our hearts:

\begin{verse}
Have I done any good in the world today?\\
Have I helped anyone in need?\\
Have I cheered up the sad and made someone feel glad?\\
If not, I have failed indeed.\\
Has anyone’s burden been lighter today\\
Because I was willing to share?\\
Have the sick and the weary been helped on their way?\\
When they needed my help was I there?
\end{verse}

That service to which all of us have been called is the service of the Lord Jesus Christ.

As He enlists us to His cause, He invites us to draw close to Him. He speaks to you and to me:

“Come unto me, all ye that labour and are heavy laden, and I will give you rest.

“Take my yoke upon you, and learn of me; for I am meek and lowly in heart: and ye shall find rest unto your souls.

“For my yoke is easy, and my burden is light.”

If we truly listen, we may hear that voice from far away say to us, as it spoke to another, “Well done, thou good and faithful servant.” That each may qualify for this blessing from our Lord is my prayer, and I offer it in His name, even Jesus Christ, our Savior, amen.

\newpage
\settalk{A Word at Closing}
\setspeaker{Thomas S. Monson}
\settalkdate{April 2010}
\talkheading{A Word at Closing}{Thomas S. Monson}

This has been a wonderful closing session. I’ve seldom heard such fine sermons taught in so few words as we’ve experienced today. We’re all here because we love the Lord. We want to serve Him. Our Heavenly Father is mindful of us. Of that I testify. I acknowledge His hand in all things.

One brief scripture:

“Trust in the Lord with all thine heart; and lean not unto thine own understanding.

“In all thy ways acknowledge him, and he shall direct thy paths.”

That has been the story of my life.

My dear brothers and sisters, we come now to the conclusion of a most uplifting and inspiring conference. After listening to the counsel and testimonies of those who have spoken to us, I believe we have been richly blessed and are all more determined to live the principles of the gospel of Jesus Christ. It has been good for us to be here. We express our gratitude to each one who has spoken to us, as well as to those who have offered prayers.

The music has been magnificent. I am reminded of the scripture found in the Doctrine and Covenants: “For my soul delighteth in the song of the heart; yea, the song of the righteous is a prayer unto me, and it shall be answered with a blessing upon their heads.”

Remember that the messages we have heard during this conference will be printed in the May issues of the Ensign and Liahona magazines. I urge you to study the messages, to ponder their teachings, and then to apply them in your life.

I know you join with me in expressing gratitude to those brethren and sisters who have been released during this conference. They have served well and have made significant contributions to the work of the Lord. Their dedication has been complete. We thank them from the bottom of our hearts.

Now, we have also sustained, by uplifted hands, brethren and sisters who have been called to new positions during this conference. We want them to know that we look forward to working with them in the cause of the Master.

My brothers and sisters, today, as we look at the world around us, we are faced with problems which are serious and of great concern to us. The world seems to have slipped from the moorings of safety and drifted from the harbor of peace.

Permissiveness, immorality, pornography, dishonesty, and a host of other ills cause many to be tossed about on a sea of sin and crushed on the jagged reefs of lost opportunities, forfeited blessings, and shattered dreams.

My counsel for all of us is to look to the lighthouse of the Lord. There is no fog so dense, no night so dark, no gale so strong, no mariner so lost but what its beacon light can rescue. It beckons through the storms of life. The lighthouse of the Lord sends forth signals readily recognized and never failing.

I love the words found in Psalms: “The Lord is my rock, and my fortress, and my deliverer; my God, my strength, in whom I will trust; … I will call upon the Lord … so [I shall] be saved from mine enemies.”

The Lord loves us, my brothers and sisters, and will bless us as we call upon Him.

How grateful we are for the restored gospel of Jesus Christ and for all the good it brings into our lives. The Lord has poured out His blessings upon us as a people. I bear my testimony to you that this work is true, that our Savior lives, and that He guides and directs His Church here upon the earth.

Now, as we come to the final moments of this conference, my heart is full and my feelings tender. I express my love and gratitude to you. Thank you for your prayers in my behalf and in behalf of all of the General Authorities of the Church. The Lord hears your prayers and blesses us and directs us in the affairs of His kingdom here upon the earth. For this we are deeply grateful.

As we leave this conference, I invoke the blessings of heaven upon each of you. As you return to your homes around the world, I pray our Heavenly Father will bless you and your families. May the messages and spirit of this conference find expression in all that you do—in your homes, in your work, in your meetings, and in all your comings and goings.

I love you. I pray for you. May God bless you. May His promised peace be with you now and always, in the name of Jesus Christ, amen.

\newpage
\settalk{He Is Risen!}
\setspeaker{Thomas S. Monson}
\settalkdate{April 2010}
\talkheading{He Is Risen!}{Thomas S. Monson}

This has been a remarkable session. In behalf of all who participated thus far in word or music, as the President of the Church, I have chosen simply to say to you at this moment just two words, known as the two most important words in the English language. To Sister Cheryl Lant and her counselors, the choir, the musicians, the speakers, those words are “Thank you.”

Many years ago, while in London, England, I visited the famed Tate art gallery. Works by Gainsborough, Rembrandt, Constable, and other renowned artists were displayed in room after room. I admired their beauty and recognized the skill which had been required to create these masterpieces. Tucked away in a quiet corner of the third floor, however, was a painting which not only caught my attention but also captured my heart. The artist, Frank Bramley, had painted a humble cottage facing a windswept sea. Two women, the mother and the wife of an absent fisherman, had watched and waited the night through for his return. Now the night had passed, and the realization had set in that he had been lost at sea and would not return. Kneeling at the side of her mother-in-law, her head buried in the lap of the older woman, the young wife wept in despair. The spent candle on the window ledge told of the fruitless vigil.

I sensed the young woman’s heartache; I felt her grief. The hauntingly vivid inscription which the artist gave to his work told the tragic story. It read, A Hopeless Dawn.

Oh, how the young woman longed for the comfort, even the reality, of Robert Louis Stevenson’s “Requiem”:

\begin{verse}
Home is the sailor, home from sea,\\
And the hunter home from the hill.
\end{verse}

Among all the facts of mortality, none is so certain as its end. Death comes to all; it is our “universal heritage; it may claim its victim[s] in infancy or youth, [it may visit] in the period of life’s prime, or its summons may be deferred until the snows of age have gathered upon the … head; it may befall as the result of accident or disease, … or … through natural causes; but come it must.” It inevitably represents a painful loss of association and, particularly in the young, a crushing blow to dreams unrealized, ambitions unfulfilled, and hopes vanquished.

What mortal being, faced with the loss of a loved one or, indeed, standing himself or herself on the threshold of infinity, has not pondered what lies beyond the veil which separates the seen from the unseen?

Centuries ago the man Job—so long blessed with every material gift, only to find himself sorely afflicted by all that can befall a human being—sat with his companions and uttered the timeless, ageless question, “If a man die, shall he live again?” Job spoke what every other living man or woman has pondered.

This glorious Easter morning I’d like to consider Job’s question—“If a man die, shall he live again?”—and provide the answer which comes not only from thoughtful consideration but also from the revealed word of God. I begin with the essentials.

If there is a design in this world in which we live, there must be a Designer. Who can behold the many wonders of the universe without believing that there is a design for all mankind? Who can doubt that there is a Designer?

In the book of Genesis we learn that the Grand Designer created the heaven and the earth: “And the earth was without form, and void; and darkness was upon the face of the deep.”

“Let there be light,” said the Grand Designer, “and there was light.” He created a firmament. He separated the land from the waters and said, “Let the earth bring forth grass, … the fruit tree yielding fruit after his kind, whose seed is in itself.”

Two lights He created—the sun and the moon. Came the stars by His design. He called for living creatures in the water and fowls to fly above the earth. And it was so. He made cattle, beasts, and creeping things. The design was nearly complete.

Last of all, He created man in His own image—male and female—with dominion over all other living things.

Man alone received intelligence—a brain, a mind, and a soul. Man alone, with these attributes, had the capacity for faith and hope, for inspiration and ambition.

Who could persuasively argue that man—the noblest work of the Great Designer, with dominion over all living things, with a brain and a will, with a mind and a soul, with intelligence and divinity—should come to an end when the spirit forsakes its earthly temple?

To understand the meaning of death, we must appreciate the purpose of life. The dim light of belief must yield to the noonday sun of revelation, by which we know that we lived before our birth into mortality. In our premortal state, we were doubtless among the sons and daughters of God who shouted for joy because of the opportunity to come to this challenging yet necessary mortal existence. We knew that our purpose was to gain a physical body, to overcome trials, and to prove that we would keep the commandments of God. Our Father knew that because of the nature of mortality, we would be tempted, would sin, and would fall short. So that we might have every chance of success, He provided a Savior, who would suffer and die for us. Not only would He atone for our sins, but as a part of that Atonement, He would also overcome the physical death to which we would be subject because of the Fall of Adam.

Thus, more than 2,000 years ago, Christ, our Savior, was born to mortal life in a stable in Bethlehem. The long-foretold Messiah had come.

There was very little written of the boyhood of Jesus. I love the passage from Luke: “And Jesus increased in wisdom and stature, and in favour with God and man.” And from the book of Acts, there is a short phrase concerning the Savior which has a world of meaning: “[He] went about doing good.”

He was baptized by John in the river Jordan. He called the Twelve Apostles. He blessed the sick. He caused the lame to walk, the blind to see, the deaf to hear. He even raised the dead to life. He taught, He testified, and He provided a perfect example for us to follow.

And then the mortal mission of the Savior of the world drew to its close. A last supper with His Apostles took place in an upper room. Ahead lay Gethsemane and Calvary’s cross.

No mere mortal can conceive the full import of what Christ did for us in Gethsemane. He Himself later described the experience: “[The] suffering caused myself, even God, the greatest of all, to tremble because of pain, and to bleed at every pore, and to suffer both body and spirit.”

Following the agony of Gethsemane, now drained of strength, He was seized by rough, crude hands and taken before Annas, Caiaphas, Pilate, and Herod. He was accused and cursed. Vicious blows further weakened His pain-racked body. Blood ran down His face as a cruel crown fashioned of sharp thorns was forced onto His head, piercing His brow. And then once again He was taken to Pilate, who gave in to the cries of the angry mob: “Crucify him, crucify him.”

He was scourged with a whip into whose multiple leather strands sharp metals and bones were woven. Rising from the cruelty of the scourge, with stumbling steps He carried His own cross until He could go no farther and another shouldered the burden for Him.

Finally, on a hill called Calvary, while helpless followers looked on, His wounded body was nailed to a cross. Mercilessly He was mocked and cursed and derided. And yet He cried out, “Father, forgive them; for they know not what they do.”

The agonizing hours passed as His life ebbed. From His parched lips came the words, “Father, into thy hands I commend my spirit: and having said thus, he gave up the ghost.”

As the serenity and solace of a merciful death freed Him from the sorrows of mortality, He returned to the presence of His Father.

At the last moment, the Master could have turned back. But He did not. He passed beneath all things that He might save all things. His lifeless body was hurriedly but gently placed in a borrowed tomb.

No words in Christendom mean more to me than those spoken by the angel to the weeping Mary Magdalene and the other Mary when, on the first day of the week, they approached the tomb to care for the body of their Lord. Spoke the angel:

“Why seek ye the living among the dead?

“He is not here, but is risen.”

Our Savior lived again. The most glorious, comforting, and reassuring of all events of human history had taken place—the victory over death. The pain and agony of Gethsemane and Calvary had been wiped away. The salvation of mankind had been secured. The Fall of Adam had been reclaimed.

The empty tomb that first Easter morning was the answer to Job’s question, “If a man die, shall he live again?” To all within the sound of my voice, I declare, If a man die, he shall live again. We know, for we have the light of revealed truth.

“For since by man came death, by man came also the resurrection of the dead.

“For as in Adam all die, even so in Christ shall all be made alive.”

I have read—and I believe—the testimonies of those who experienced the grief of Christ’s Crucifixion and the joy of His Resurrection. I have read—and I believe—the testimonies of those in the New World who were visited by the same risen Lord.

I believe the testimony of one who, in this dispensation, spoke with the Father and the Son in a grove now called sacred and who gave his life, sealing that testimony with his blood. Declared he:

“And now, after the many testimonies which have been given of him, this is the testimony, last of all, which we give of him: That he lives!

“For we saw him, even on the right hand of God; and we heard the voice bearing record that he is the Only Begotten of the Father.”

The darkness of death can always be dispelled by the light of revealed truth. “I am the resurrection, and the life,” spoke the Master. “Peace I leave with you, my peace I give unto you.”

Over the years I have heard and read testimonies too numerous to count, shared with me by individuals who testify of the reality of the Resurrection and who have received, in their hours of greatest need, the peace and comfort promised by the Savior.

I will mention just part of one such account. Two weeks ago I received a touching letter from a father of seven who wrote about his family and, in particular, his son Jason, who had become ill when 11 years of age. Over the next few years, Jason’s illness recurred several times. This father told of Jason’s positive attitude and sunny disposition, despite his health challenges. Jason received the Aaronic Priesthood at age 12 and “always willingly magnified his responsibilities with excellence, whether he felt well or not.” He received his Eagle Scout Award when he was 14 years old.

Last summer, not long after Jason’s 15th birthday, he was once again admitted to the hospital. On one of his visits to see Jason, his father found him with his eyes closed. Not knowing whether Jason was asleep or awake, he began talking softly to him. “Jason,” he said, “I know you have been through a lot in your short life and that your current condition is difficult. Even though you have a giant battle ahead, I don’t ever want you to lose your faith in Jesus Christ.” He said he was startled as Jason immediately opened his eyes and said, “Never!” in a clear, resolute voice. Jason then closed his eyes and said no more.

His father wrote: “In this simple declaration, Jason expressed one of the most powerful, pure testimonies of Jesus Christ that I have ever heard. … As his declaration of ‘Never!’ became imprinted on my soul that day, my heart filled with joy that my Heavenly Father had blessed me to be the father of such a tremendous and noble boy. … [It] was the last time I heard him declare his testimony of Christ.”

Although his family was expecting this to be just another routine hospitalization, Jason passed away less than two weeks later. An older brother and sister were serving missions at the time. Another brother, Kyle, had just received his mission call. In fact, the call had come earlier than expected, and on August 5, just a week before Jason’s passing, the family gathered in his hospital room so that Kyle’s mission call could be opened there and shared with the entire family.

In his letter to me, this father included a photograph of Jason in his hospital bed, with his big brother Kyle standing beside the bed, holding his mission call. This caption was written beneath the photograph: “Called to serve their missions together—on both sides of the veil.”

Jason’s brother and sister already serving missions sent beautiful, comforting letters home to be shared at Jason’s funeral. His sister, serving in the Argentina Buenos Aires West Mission, as part of her letter, wrote: “I know that Jesus Christ lives, and because He lives, all of us, including our beloved Jason, will live again too. … We can take comfort in the sure knowledge we have that we have been sealed together as an eternal family. … If we do our very best to obey and do better in this life, we will see [him again].” She continued: “[A] scripture that I have long loved now takes on new significance and importance at this time. … [From] Revelation chapter 21, verse 4: ‘And God shall wipe away all tears from their eyes; and there shall be no more death, neither sorrow, nor crying, neither shall there be any more pain: for the former things are passed away.’”

My beloved brothers and sisters, in our hour of deepest sorrow, we can receive profound peace from the words of the angel that first Easter morning: “He is not here: for he is risen.”

\begin{verse}
He is risen! He is risen!\\
Tell it out with joyful voice.\\
He has burst his three days’ prison;\\
Let the whole wide earth rejoice.\\
Death is conquered; man is free.\\
Christ has won the victory!
\end{verse}

As one of His special witnesses on earth today, this glorious Easter Sunday, I declare that this is true, in His sacred name—even the name of Jesus Christ, our Savior—amen.

\newpage
\settalk{Preparation Brings Blessings}
\setspeaker{Thomas S. Monson}
\settalkdate{April 2010}
\talkheading{Preparation Brings Blessings}{Thomas S. Monson}

Brethren, you who are here in the Conference Center in Salt Lake City are an inspiring sight to behold. It is amazing to realize that in thousands of chapels throughout the world, others of you—fellow holders of the priesthood of God—are receiving this broadcast by way of satellite transmission. Your nationalities vary, and your languages are many, but a common thread binds us together. We have been entrusted to bear the priesthood and to act in the name of God. We are the recipients of a sacred trust. Much is expected of us.

One of my most vivid memories is attending priesthood meeting as a newly ordained deacon and singing the opening hymn “Come, All Ye Sons of God.” Tonight I echo the spirit of that special hymn and say to you, “Come, all ye sons of God who have received the priesthood.” Let us consider our callings, let us reflect on our responsibilities, and let us follow Jesus Christ, our Lord.

Twenty years ago I attended a sacrament meeting where the children responded to the theme “I Belong to The Church of Jesus Christ of Latter-day Saints.” These boys and girls demonstrated they were in training for service to the Lord and to others. The music was beautiful, the recitations skillfully rendered, and the spirit heaven-sent. One of my grandsons, who was 11 years old at that time, had spoken of the First Vision as he presented his part on the program. Afterward, as he came to his parents and grandparents, I said to him, “Tommy, I think you are almost ready to be a missionary.”

He replied, “Not yet. I still have a lot to learn.”

Through the years that followed, Tommy did learn, thanks to his parents and to teachers and advisers at church, who were dedicated and conscientious. When he was old enough, he was called to serve a mission. He did so in a most honorable fashion.

Young men, I admonish you to prepare for service as a missionary. There are many tools to help you learn the lessons which will be beneficial to you as well as helping you to live the life you will need to have lived to be worthy. One such tool is the booklet entitled For the Strength of Youth, published under the direction of the First Presidency and Quorum of the Twelve Apostles. It features standards from the writings and teachings of Church leaders and from scripture, adherence to which will bring the blessings of our Heavenly Father and the guidance of His Son to each of us. In addition, there are lesson manuals, carefully prepared after prayerful consideration. Families have family home evenings, where gospel principles are taught. Almost all of you have the opportunity to attend seminary classes taught by dedicated teachers who have much to share.

Begin to prepare for a temple marriage as well as for a mission. Proper dating is a part of that preparation. In cultures where dating is appropriate, do not date until you are 16 years old. “Not all teenagers need to date or even want to. … When you begin dating, go in groups or on double dates. … Make sure your parents meet [and become acquainted with] those you date.” Because dating is a preparation for marriage, “date only those who have high standards.”

Be careful to go to places where there is a good environment, where you won’t be faced with temptation.

A wise father said to his son, “If you ever find yourself in a place where you shouldn’t ought to be, get out!” Good advice for all of us.

Servants of the Lord have always counseled us to dress appropriately to show respect for our Heavenly Father and for ourselves. The way you dress sends messages about yourself to others and often influences the way you and others act. Dress in such a way as to bring out the best in yourself and those around you. Avoid extremes in clothing and appearance, including tattoos and piercings.

Everyone needs good friends. Your circle of friends will greatly influence your thinking and behavior, just as you will theirs. When you share common values with your friends, you can strengthen and encourage each other. Treat everyone with kindness and dignity. Many nonmembers have come into the Church through friends who have involved them in Church activities.

The oft-repeated adage is ever true: “Honesty [is] the best policy.” A Latter-day Saint young man lives as he teaches and as he believes. He is honest with others. He is honest with himself. He is honest with God. He is honest by habit and as a matter of course. When a difficult decision must be made, he never asks himself, “What will others think?” but rather, “What will I think of myself?”

For some, there will come the temptation to dishonor a personal standard of honesty. In a business law class at the university I attended, I remember that one particular classmate never prepared for the class discussions. I thought to myself, “How is he going to pass the final examination?”

I discovered the answer when he came to the classroom for the final exam on a winter’s day wearing on his bare feet only a pair of sandals. I was surprised and watched him as the class began. All of our books had been placed upon the floor, as per the instruction. He slipped the sandals from his feet; and then, with toes that he had trained and had prepared with glycerin, he skillfully turned the pages of one of the books which he had placed on the floor, thereby viewing the answers to the examination questions.

He received one of the highest grades in that course on business law. But the day of reckoning came. Later, as he prepared to take his comprehensive exam, for the first time the dean of his particular discipline said, “This year I will depart from tradition and will conduct an oral, rather than a written, test.” Our favorite trained-toe expert found that he had his foot in his mouth on that occasion and failed the exam.

How you speak and the words you use tell much about the image you choose to portray. Use language to build and uplift those around you. Profane, vulgar, or crude language and inappropriate or off-color jokes are offensive to the Lord. Never misuse the name of God or Jesus Christ. The Lord said, “Thou shalt not take the name of the Lord thy God in vain.”

Our Heavenly Father has counseled us to seek after “anything virtuous, lovely, or of good report or praiseworthy.” Whatever you read, listen to, or watch makes an impression on you.

Pornography is especially dangerous and addictive. Curious exploration of pornography can become a controlling habit, leading to coarser material and to sexual transgression. Avoid pornography at all costs.

Don’t be afraid to walk out of a movie, turn off a television set, or change a radio station if what’s being presented does not meet your Heavenly Father’s standards. In short, if you have any question about whether a particular movie, book, or other form of entertainment is appropriate, don’t see it, don’t read it, don’t participate.

The Apostle Paul declared: “Know ye not that ye are the temple of God, and that the Spirit of God dwelleth in you? … The temple of God is holy, which temple ye are.” Brethren, it is our responsibility to keep our temples clean and pure.

Hard drugs, wrongful use of prescription drugs, alcohol, coffee, tea, and tobacco products destroy your physical, mental, and spiritual well-being. Any form of alcohol is harmful to your spirit and your body. Tobacco can enslave you, weaken your lungs, and shorten your life.

Music can help you draw closer to your Heavenly Father. It can be used to educate, edify, inspire, and unite. However, music can, by its tempo, beat, intensity, and lyrics, dull your spiritual sensitivity. You cannot afford to fill your minds with unworthy music.

Because sexual intimacy is so sacred, the Lord requires self-control and purity before marriage as well as full fidelity after marriage. In dating, treat your date with respect and expect your date to show that same respect for you. Tears inevitably follow transgression.

President David O. McKay, ninth President of the Church, advised, “I implore you to think clean thoughts.” He then made this significant declaration of truth: “Every action is preceded by a thought. If we want to control our actions, we must control our thinking.” Brethren, fill your minds with good thoughts, and your actions will be proper. May each of you be able to echo in truth the line from Tennyson spoken by Sir Galahad: “My strength is as the strength of ten, because my heart is pure.”

Not long ago the author of a paper on teenage sexuality summed up his research by saying that society sends teens a mixed message: advertisements and the mass media convey “very heavy messages that sexual activity is acceptable and expected,” inducements that sometimes drown out the warnings of experts and the pleas of parents. The Lord cuts through all the media messages with clear and precise language when He declares to us, “Be ye clean.”

Whenever temptation comes, remember the wise counsel of the Apostle Paul, who declared, “There hath no temptation taken you but such as is common to man: but God is faithful, who will not suffer you to be tempted above that ye are able; but will with the temptation also make a way to escape, that ye may be able to bear it.”

When you were confirmed a member of the Church, you received the right to the companionship of the Holy Ghost. He can help you make good choices. When challenged or tempted, you do not need to feel alone. Remember that prayer is the passport to spiritual power.

If any has stumbled in his journey, there is a way back. The process is called repentance. Our Savior died to provide you and me that blessed gift. Though the path is difficult, the promise is real: “Though your sins be as scarlet, they shall be as white as snow.”

Don’t put your eternal life at risk. Keep the commandments of God. If you have sinned, the sooner you begin to make your way back, the sooner you will find the sweet peace and joy that come with the miracle of forgiveness. Happiness comes from living the way the Lord wants you to live and from service to God and others.

Spiritual strength frequently comes through selfless service. Some years ago I visited what was then called the California Mission, where I interviewed a young missionary from Georgia. I recall saying to him, “Do you send a letter home to your parents every week?”

He replied, “Yes, Brother Monson.”

Then I asked, “Do you enjoy receiving letters from home?”

He didn’t answer. At length I inquired, “When was the last time you had a letter from home?”

With a quavering voice, he responded, “I’ve never had a letter from home. Father’s just a deacon, and Mother’s not a member of the Church. They pleaded with me not to come. They said that if I left on a mission, they would not be writing to me. What shall I do, Brother Monson?”

I offered a silent prayer to my Heavenly Father: “What should I tell this young servant of Thine, who has sacrificed everything to serve Thee?” And the inspiration came. I said, “Elder, you send a letter home to your mother and father every week of your mission. Tell them what you are doing. Tell them how much you love them and then bear your testimony to them.”

He asked, “Will they then write to me?”

I responded, “Then they will write to you.”

We parted and I went on my way. Months later I was attending a stake conference in Southern California, when a young missionary came up to me and said, “Brother Monson, do you remember me? I’m the missionary who had not received a letter from my mother or my father during my first nine months in the mission field. You told me, ‘Send a letter home every week, Elder, and your parents will write to you.’” Then he asked, “Do you remember that promise, Elder Monson?”

I remembered. I inquired, “Have you heard from your parents?”

He reached into his pocket and took out a sheaf of letters with an elastic band around them, took a letter from the top of the stack, and said, “Have I heard from my parents! Listen to this letter from my mother: ‘Son, we so much enjoy your letters. We’re proud of you, our missionary. Guess what? Dad has been ordained a priest. He’s preparing to baptize me. I’m meeting with the missionaries; and one year from now we want to come to California as you complete your mission, for we, with you, would like to become a forever family by entering the temple of the Lord.’” This young missionary asked, “Brother Monson, does Heavenly Father always answer prayers and fulfill Apostles’ promises?”

I replied, “When one has faith as you have demonstrated, our Heavenly Father hears such prayers and answers in His own way.”

Clean hands, a pure heart, and a willing mind had touched heaven. A blessing, heaven-sent, had answered the fervent prayer of a missionary’s humble heart.

Brethren, it is my prayer that we may so live that we too may touch heaven and be similarly blessed, each and every one, in the name of the Giver of all blessings, even Jesus Christ, amen.

\newpage
\settalk{Welcome to Conference}
\setspeaker{Thomas S. Monson}
\settalkdate{April 2010}
\talkheading{Welcome to Conference}{Thomas S. Monson}

How good it is, my beloved brothers and sisters, to meet together once again. This conference marks 180 years since the Church was organized. How grateful we are for the Prophet Joseph Smith, who sought for the truth, who found it, and who, under the direction of the Lord, restored the gospel and organized the Church.

The Church has grown steadily since that day in 1830. It continues to change the lives of more and more people every year and to spread across the earth as our missionary force seeks out those who are searching for the truth. Once again we call upon the members of the Church to reach out to the new converts or to those making their way back into the Church, to surround them with love and to help them feel at home.

Thank you, my brothers and sisters, for your faith and devotion to the gospel of Jesus Christ. Thank you for all that you do in your wards and branches, in your stakes and districts. You serve willingly and well and accomplish great good. May the Lord bless you as you strive to follow Him and to obey His commandments.

Since last we met, the Church has continued to provide much-needed humanitarian assistance in various locations around the world. In the past three months alone, humanitarian assistance has been provided in French Polynesia, Mongolia, Bolivia, Peru, Arizona, Mexico, Portugal, and Uganda, among other areas. Most recently we have assisted in Haiti and Chile following devastating earthquakes and tsunamis in those areas. We express our love to our Church members who have suffered in these disasters. You are in our prayers. We express profound gratitude to all of you for your willingness to assist with our humanitarian efforts by sharing your resources and, in many cases, your time, your talents, and your expertise.

This year marks 25 years since our humanitarian program became part of our welfare effort. The number of individuals assisted by this program could never adequately be measured. We will always strive to be among the first on the scene of disasters, wherever they may occur.

The Church continues to grow and to move forward. The building of temples is an indication of such growth. Recently we announced a new temple which will be built in Payson, Utah. We also announced major renovations which will be made to the Ogden Utah Temple. Within the next three months we will dedicate new temples in Vancouver, British Columbia; in the Gila Valley, Arizona; and in Cebu City in the Philippines. Later in the year other temples will be dedicated or rededicated. We will continue to build temples throughout the world as our membership grows. Each year millions of ordinances are performed in the temples for our deceased loved ones. May we continue to be faithful in performing such ordinances for those who are unable to do so for themselves.

Many of you are aware that a short time after October conference, my dear wife, Frances, suffered a fall, which left her with a broken hip and a broken shoulder. After two successful surgeries and several weeks of hospitalization, she was able to return home. She is doing well and continues to make progress toward a full recovery. She was able to attend the general Young Women meeting last Saturday and plans to attend a session or two this weekend. In fact, at the last minute she said, “I’m going today!” And she’s here! She joins me in expressing our deep gratitude to our Heavenly Father and to all of you for your prayers and your well wishes in her behalf.

Now, brothers and sisters, we have come here to be instructed and inspired. We welcome those of you who are new in the Church. Others of you are struggling with problems, with challenges, with disappointments, with losses. We love you and we pray for you. Many messages, covering a variety of gospel topics, will be given during the next two days. Those men and women who will speak to you have sought heaven’s help concerning the messages they will give.

It is my prayer that we may be filled with His Spirit as we listen and learn. That this may be so, I pray in the name of Jesus Christ, our Lord and Savior, amen.

\newpage
\settalk{As We Meet Together Again}
\setspeaker{Thomas S. Monson}
\settalkdate{October 2010}
\talkheading{As We Meet Together Again}{Thomas S. Monson}

My beloved brothers and sisters, we welcome you to general conference, which is being heard and seen by various means throughout the world. We express thanks to all who have to do with the complicated logistics of this great undertaking.

Since April when we last met, the work of the Church has moved forward unhindered. It has been my privilege to dedicate four new temples. Accompanied by my counselors and other General Authorities, I have traveled to Gila Valley, Arizona; to Vancouver, British Columbia; to Cebu City in the Philippines; and to Kyiv, Ukraine. The temple in each of these locations is magnificently beautiful. Each one is blessing the lives of our members and is an influence for good upon those not of our faith.

The evening prior to each temple dedication, we were privileged to view a cultural celebration, participated in by our young people and some of our not-so-young people. These events were generally held in large stadiums, although in Kyiv we met in a beautiful palace. The dancing, singing, musical performances, and displays were excellent. I express my commendation and love to all who were involved.

Each temple dedication was a spiritual feast. We felt the Spirit of the Lord at all of them.

Next month we will rededicate the Laie Hawaii Temple, one of our oldest temples, which has undergone extensive renovations during many months. We look forward to that sacred occasion.

We continue to build temples. This morning I am pleased to announce five additional temples for which sites are being acquired and which, in coming months and years, will be built in the following locations: Lisbon, Portugal; Indianapolis, Indiana; Urdaneta, Philippines; Hartford, Connecticut; and Tijuana, Mexico.

The ordinances performed in our temples are vital to our salvation and to the salvation of our deceased loved ones. May we continue faithful in attending the temples, which are being built closer and closer to our members.

Now, before we hear from our speakers this morning, may I mention a matter close to my heart and which deserves our serious attention. I speak of missionary work.

First, to young men of the Aaronic Priesthood and to you young men who are becoming elders: I repeat what prophets have long taught—that every worthy, able young man should prepare to serve a mission. Missionary service is a priesthood duty—an obligation the Lord expects of us who have been given so very much. Young men, I admonish you to prepare for service as a missionary. Keep yourselves clean and pure and worthy to represent the Lord. Maintain your health and strength. Study the scriptures. Where such is available, participate in seminary or institute. Familiarize yourself with the missionary handbook Preach My Gospel.

A word to you young sisters: while you do not have the same priesthood responsibility as do the young men to serve as full-time missionaries, you also make a valuable contribution as missionaries, and we welcome your service.

And now to you mature brothers and sisters: we need many, many more senior couples. To the faithful couples now serving or who have served in the past, we thank you for your faith and devotion to the gospel of Jesus Christ. You serve willingly and well and accomplish great good.

To those of you who are not yet to the season of life when you might serve a couples mission, I urge you to prepare now for the day when you and your spouse might do so. As your circumstances allow, as you are eligible for retirement, and as your health permits, make yourselves available to leave home and give full-time missionary service. There are few times in your lives when you will enjoy the sweet spirit and satisfaction that come from giving full-time service together in the work of the Master.

Now, my brothers and sisters, may you be attuned to the Spirit of the Lord as we hear from His servants during the next two days. That this may be the blessing of each, I pray humbly in the name of Jesus Christ, amen.

\newpage
\settalk{Charity Never Faileth}
\setspeaker{Thomas S. Monson}
\settalkdate{October 2010}
\talkheading{Charity Never Faileth}{Thomas S. Monson}

Our souls have rejoiced tonight and reached toward heaven. We have been blessed with beautiful music and inspired messages. The Spirit of the Lord is here. I pray for His inspiration to be with me now as I share with you some of my thoughts and feelings.

I begin with a short anecdote which illustrates a point I should like to make.

A young couple, Lisa and John, moved into a new neighborhood. One morning while they were eating breakfast, Lisa looked out the window and watched her next-door neighbor hanging out her wash.

“That laundry’s not clean!” Lisa exclaimed. “Our neighbor doesn’t know how to get clothes clean!”

John looked on but remained silent.

Every time her neighbor would hang her wash to dry, Lisa would make the same comments.

A few weeks later Lisa was surprised to glance out her window and see a nice, clean wash hanging in her neighbor’s yard. She said to her husband, “Look, John—she’s finally learned how to wash correctly! I wonder how she did it.”

John replied, “Well, dear, I have the answer for you. You’ll be interested to know that I got up early this morning and washed our windows!”

Tonight I’d like to share with you a few thoughts concerning how we view each other. Are we looking through a window which needs cleaning? Are we making judgments when we don’t have all the facts? What do we see when we look at others? What judgments do we make about them?

Said the Savior, “Judge not.” He continued, “Why beholdest thou the mote that is in thy brother’s eye, but considerest not the beam that is in thine own eye?” Or, to paraphrase, why beholdest thou what you think is dirty laundry at your neighbor’s house but considerest not the soiled window in your own house?

None of us is perfect. I know of no one who would profess to be so. And yet for some reason, despite our own imperfections, we have a tendency to point out those of others. We make judgments concerning their actions or inactions.

There is really no way we can know the heart, the intentions, or the circumstances of someone who might say or do something we find reason to criticize. Thus the commandment: “Judge not.”

Forty-seven years ago this general conference, I was called to the Quorum of the Twelve Apostles. At the time, I had been serving on one of the general priesthood committees of the Church, and so before my name was presented, I sat with my fellow members of that priesthood committee, as was expected of me. My wife, however, had no idea where to go and no one with whom she could sit and, in fact, was unable to find a seat anywhere in the Tabernacle. A dear friend of ours, who was a member of one of the general auxiliary boards and who was sitting in the area designated for the board members, asked Sister Monson to sit with her. This woman knew nothing of my call—which would be announced shortly—but she spotted Sister Monson, recognized her consternation, and graciously offered her a seat. My dear wife was relieved and grateful for this kind gesture. Sitting down, however, she heard loud whispering behind her as one of the board members expressed her annoyance to those around her that one of her fellow board members would have the audacity to invite an “outsider” to sit in this area reserved only for them. There was no excuse for her unkind behavior, regardless of who might have been invited to sit there. However, I can only imagine how that woman felt when she learned that the “intruder” was the wife of the newest Apostle.

Not only are we inclined to judge the actions and words of others, but many of us judge appearances: clothing, hairstyles, size. The list could go on and on.

A classic account of judging by appearance was printed in a national magazine many years ago. It is a true account—one which you may have heard but which bears repeating.

A woman by the name of Mary Bartels had a home directly across the street from the entrance to a hospital clinic. Her family lived on the main floor and rented the upstairs rooms to outpatients at the clinic.

One evening a truly awful-looking old man came to the door asking if there was room for him to stay the night. He was stooped and shriveled, and his face was lopsided from swelling—red and raw. He said he’d been hunting for a room since noon but with no success. “I guess it’s my face,” he said. “I know it looks terrible, but my doctor says it could possibly improve after more treatments.” He indicated he’d be happy to sleep in the rocking chair on the porch. As she talked with him, Mary realized this little old man had an oversized heart crowded into that tiny body. Although her rooms were filled, she told him to wait in the chair and she’d find him a place to sleep.

At bedtime Mary’s husband set up a camp cot for the man. When she checked in the morning, the bed linens were neatly folded and he was out on the porch. He refused breakfast, but just before he left for his bus, he asked if he could return the next time he had a treatment. “I won’t put you out a bit,” he promised. “I can sleep fine in a chair.” Mary assured him he was welcome to come again.

In the several years he went for treatments and stayed in Mary’s home, the old man, who was a fisherman by trade, always had gifts of seafood or vegetables from his garden. Other times he sent packages in the mail.

When Mary received these thoughtful gifts, she often thought of a comment her next-door neighbor made after the disfigured, stooped old man had left Mary’s home that first morning. “Did you keep that awful-looking man last night? I turned him away. You can lose customers by putting up such people.”

Mary knew that maybe they had lost customers once or twice, but she thought, “Oh, if only they could have known him, perhaps their illnesses would have been easier to bear.”

After the man passed away, Mary was visiting with a friend who had a greenhouse. As she looked at her friend’s flowers, she noticed a beautiful golden chrysanthemum but was puzzled that it was growing in a dented, old, rusty bucket. Her friend explained, “I ran short of pots, and knowing how beautiful this one would be, I thought it wouldn’t mind starting in this old pail. It’s just for a little while, until I can put it out in the garden.”

Mary smiled as she imagined just such a scene in heaven. “Here’s an especially beautiful one,” God might have said when He came to the soul of the little old man. “He won’t mind starting in this small, misshapen body.” But that was long ago, and in God’s garden how tall this lovely soul must stand!

Appearances can be so deceiving, such a poor measure of a person. Admonished the Savior, “Judge not according to the appearance.”

A member of a women’s organization once complained when a certain woman was selected to represent the organization. She had never met the woman, but she had seen a photograph of her and didn’t like what she saw, considering her to be overweight. She commented, “Of the thousands of women in this organization, surely a better representative could have been chosen.”

True, the woman who was chosen was not “model slim.” But those who knew her and knew her qualities saw in her far more than was reflected in the photograph. The photograph did show that she had a friendly smile and a look of confidence. What the photograph didn’t show was that she was a loyal and compassionate friend, a woman of intelligence who loved the Lord and who loved and served His children. It didn’t show that she volunteered in the community and was a considerate and concerned neighbor. In short, the photograph did not reflect who she really was.

I ask: if attitudes, deeds, and spiritual inclinations were reflected in physical features, would the countenance of the woman who complained be as lovely as that of the woman she criticized?

My dear sisters, each of you is unique. You are different from each other in many ways. There are those of you who are married. Some of you stay at home with your children, while others of you work outside your homes. Some of you are empty nesters. There are those of you who are married but do not have children. There are those who are divorced, those who are widowed. Many of you are single women. Some of you have college degrees; some of you do not. There are those who can afford the latest fashions and those who are lucky to have one appropriate Sunday outfit. Such differences are almost endless. Do these differences tempt us to judge one another?

Mother Teresa, a Catholic nun who worked among the poor in India most of her life, spoke this profound truth: “If you judge people, you have no time to love them.” The Savior has admonished, “This is my commandment, That ye love one another, as I have loved you.” I ask: can we love one another, as the Savior has commanded, if we judge each other? And I answer—with Mother Teresa: no, we cannot.

The Apostle James taught, “If any … among you seem to be religious, and bridleth not his tongue, but deceiveth his own heart, this man’s [or woman’s] religion is vain.”

I have always loved your Relief Society motto: “Charity never faileth.” What is charity? The prophet Mormon teaches us that “charity is the pure love of Christ.” In his farewell message to the Lamanites, Moroni declared, “Except ye have charity ye can in nowise be saved in the kingdom of God.”

I consider charity—or “the pure love of Christ”—to be the opposite of criticism and judging. In speaking of charity, I do not at this moment have in mind the relief of the suffering through the giving of our substance. That, of course, is necessary and proper. Tonight, however, I have in mind the charity that manifests itself when we are tolerant of others and lenient toward their actions, the kind of charity that forgives, the kind of charity that is patient.

I have in mind the charity that impels us to be sympathetic, compassionate, and merciful, not only in times of sickness and affliction and distress but also in times of weakness or error on the part of others.

There is a serious need for the charity that gives attention to those who are unnoticed, hope to those who are discouraged, aid to those who are afflicted. True charity is love in action. The need for charity is everywhere.

Needed is the charity which refuses to find satisfaction in hearing or in repeating the reports of misfortunes that come to others, unless by so doing, the unfortunate one may be benefited. The American educator and politician Horace Mann once said, “To pity distress is but human; to relieve it is godlike.”

Charity is having patience with someone who has let us down. It is resisting the impulse to become offended easily. It is accepting weaknesses and shortcomings. It is accepting people as they truly are. It is looking beyond physical appearances to attributes that will not dim through time. It is resisting the impulse to categorize others.

Charity, that pure love of Christ, is manifest when a group of young women from a singles ward travels hundreds of miles to attend the funeral services for the mother of one of their Relief Society sisters. Charity is shown when devoted visiting teachers return month after month, year after year to the same uninterested, somewhat critical sister. It is evident when an elderly widow is remembered and taken to ward functions and to Relief Society activities. It is felt when the sister sitting alone in Relief Society receives the invitation, “Come—sit by us.”

In a hundred small ways, all of you wear the mantle of charity. Life is perfect for none of us. Rather than being judgmental and critical of each other, may we have the pure love of Christ for our fellow travelers in this journey through life. May we recognize that each one is doing her best to deal with the challenges which come her way, and may we strive to do our best to help out.

Charity has been defined as “the highest, noblest, strongest kind of love,” the “pure love of Christ … ; and whoso is found possessed of it at the last day, it shall be well with [her].”

“Charity never faileth.” May this long-enduring Relief Society motto, this timeless truth, guide you in everything you do. May it permeate your very souls and find expression in all your thoughts and actions.

I express my love to you, my sisters, and pray that heaven’s blessings may ever be yours. In the name of Jesus Christ, amen.

\newpage
\settalk{The Divine Gift of Gratitude}
\setspeaker{Thomas S. Monson}
\settalkdate{October 2010}
\talkheading{The Divine Gift of Gratitude}{Thomas S. Monson}

This has been a marvelous session. When I was appointed President of the Church, I said, “I’ll take one assignment for myself. I’ll be the adviser for the Tabernacle Choir.” I’m very proud of my choir!

My mother once said of me, “Tommy, I’m very proud of all that you’ve done. But I have one comment to make to you. You should have stayed with the piano.”

So I went to the piano and played a number for her: “Here we go, [here we go] to a birthday party.” Then I gave her a kiss on the forehead, and she embraced me.

I think of her. I think of my father. I think of all those General Authorities who’ve influenced me, and others, including the widows whom I visited—85 of them—with a chicken for the oven, sometimes a little money for their pocket.

I visited one late one night. It was midnight, and I went to the nursing home, and the receptionist said, “I’m sure she’s asleep, but she told me to be sure to awaken her, for she said, ‘I know he’ll come.’”

I held her hand; she called my name. She was wide awake. She pressed my hand to her lips and said, “I knew you’d come.” How could I not have come?

Beautiful music touches me that way.

My beloved brothers and sisters, we have heard inspired messages of truth, of hope, and of love. Our thoughts have turned to Him who atoned for our sins, who showed us the way to live and how to pray, and who demonstrated by His own actions the blessings of service—even our Lord and Savior, Jesus Christ.

In the book of Luke, chapter 17, we read of Him:

“And it came to pass, as he went to Jerusalem, that he passed through the midst of Samaria and Galilee.

“And as he entered into a certain village, there [he met] ten men that were lepers, which stood afar off:

“And they lifted up their voices, and said, Jesus, Master, have mercy on us.

“And when he saw them, he said unto them, Go shew yourselves unto the priests. And it came to pass, that, as they went, they were cleansed.

“And one of them, when he saw that he was healed, turned back, and with a loud voice glorified God,

“And fell down on his face at his feet, giving him thanks: and he was a Samaritan.

“And Jesus answering said, Were there not ten cleansed? but where are the nine?

“There are not found that returned to give glory to God, save this stranger.

“And he said unto him, Arise, go thy way: thy faith hath made thee whole.”

Through divine intervention, those who were lepers were spared from a cruel, lingering death and given a new lease on life. The expressed gratitude by one merited the Master’s blessing; the ingratitude shown by the nine, His disappointment.

My brothers and sisters, do we remember to give thanks for the blessings we receive? Sincerely giving thanks not only helps us recognize our blessings, but it also unlocks the doors of heaven and helps us feel God’s love.

My beloved friend President Gordon B. Hinckley said, “When you walk with gratitude, you do not walk with arrogance and conceit and egotism, you walk with a spirit of thanksgiving that is becoming to you and will bless your lives.”

In the book of Matthew in the Bible, we have another account of gratitude, this time as an expression from the Savior. As He traveled in the wilderness for three days, more than 4,000 people followed and traveled with Him. He took compassion on them, for they may not have eaten during the entire three days. His disciples, however, questioned, “Whence should we have so much bread in the wilderness, as to fill so great a multitude?” Like many of us, the disciples saw only what was lacking.

“And Jesus saith unto them, How many loaves have ye? And [the disciples] said, Seven, and a few little fishes.

“And [Jesus] commanded the multitude to sit down on the ground.

“And he took the seven loaves and the fishes, and gave thanks, and brake them, and gave to his disciples, and the disciples to the multitude.”

Notice that the Savior gave thanks for what they had—and a miracle followed: “And they did all eat, and were filled: and they took up of the broken meat that was left seven baskets full.”

We have all experienced times when our focus is on what we lack rather than on our blessings. Said the Greek philosopher Epictetus, “He is a wise man who does not grieve for the things which he has not, but rejoices for those which he has.”

Gratitude is a divine principle. The Lord declared through a revelation given to the Prophet Joseph Smith:

“Thou shalt thank the Lord thy God in all things. …

“And in nothing doth man offend God, or against none is his wrath kindled, save those who confess not his hand in all things.”

In the Book of Mormon we are told to “live in thanksgiving daily, for the many mercies and blessings which [God] doth bestow upon you.”

Regardless of our circumstances, each of us has much for which to be grateful if we will but pause and contemplate our blessings.

This is a wonderful time to be on earth. While there is much that is wrong in the world today, there are many things that are right and good. There are marriages that make it, parents who love their children and sacrifice for them, friends who care about us and help us, teachers who teach. Our lives are blessed in countless ways.

We can lift ourselves and others as well when we refuse to remain in the realm of negative thought and cultivate within our hearts an attitude of gratitude. If ingratitude be numbered among the serious sins, then gratitude takes its place among the noblest of virtues. Someone has said that “gratitude is not only the greatest of virtues, but the parent of all the others.”

How can we cultivate within our hearts an attitude of gratitude? President Joseph F. Smith, sixth President of the Church, provided an answer. Said he: “The grateful man sees so much in the world to be thankful for, and with him the good outweighs the evil. Love overpowers jealousy, and light drives darkness out of his life.” He continued: “Pride destroys our gratitude and sets up selfishness in its place. How much happier we are in the presence of a grateful and loving soul, and how careful we should be to cultivate, through the medium of a prayerful life, a thankful attitude toward God and man!”

President Smith is telling us that a prayerful life is the key to possessing gratitude.

Do material possessions make us happy and grateful? Perhaps momentarily. However, those things which provide deep and lasting happiness and gratitude are the things which money cannot buy: our families, the gospel, good friends, our health, our abilities, the love we receive from those around us. Unfortunately, these are some of the things we allow ourselves to take for granted.

The English author Aldous Huxley wrote, “Most human beings have an almost infinite capacity for taking things for granted.”

We often take for granted the very people who most deserve our gratitude. Let us not wait until it is too late for us to express that gratitude. Speaking of loved ones he had lost, one man declared his regret this way: “I remember those happy days, and often wish I could speak into the ears of the dead the gratitude which was due to them in life, and so ill returned.”

The loss of loved ones almost inevitably brings some regrets to our hearts. Let’s minimize such feelings as much as humanly possible by frequently expressing our love and gratitude to them. We never know how soon it will be too late.

A grateful heart, then, comes through expressing gratitude to our Heavenly Father for His blessings and to those around us for all that they bring into our lives. This requires conscious effort—at least until we have truly learned and cultivated an attitude of gratitude. Often we feel grateful and intend to express our thanks but forget to do so or just don’t get around to it. Someone has said that “feeling gratitude and not expressing it is like wrapping a present and not giving it.”

When we encounter challenges and problems in our lives, it is often difficult for us to focus on our blessings. However, if we reach deep enough and look hard enough, we will be able to feel and recognize just how much we have been given.

I share with you an account of one family which was able to find blessings in the midst of serious challenges. This is an account I read many years ago and have kept because of the message it conveys. It was written by Gordon Green and appeared in an American magazine over 50 years ago.

Gordon tells how he grew up on a farm in Canada, where he and his siblings had to hurry home from school while the other children played ball and went swimming. Their father, however, had the capacity to help them understand that their work amounted to something. This was especially true after harvesttime when the family celebrated Thanksgiving, for on that day their father gave them a great gift. He took an inventory of everything they had.

On Thanksgiving morning he would take them to the cellar with its barrels of apples, bins of beets, carrots packed in sand, and mountains of sacked potatoes as well as peas, corn, string beans, jellies, strawberries, and other preserves which filled their shelves. He had the children count everything carefully. Then they went out to the barn and figured how many tons of hay there were and how many bushels of grain in the granary. They counted the cows, pigs, chickens, turkeys, and geese. Their father said he wanted to see how they stood, but they knew he really wanted them to realize on that feast day how richly God had blessed them and had smiled upon all their hours of work. Finally, when they sat down to the feast their mother had prepared, the blessings were something they felt.

Gordon indicated, however, that the Thanksgiving he remembered most thankfully was the year they seemed to have nothing for which to be grateful.

The year started off well: they had leftover hay, lots of seed, four litters of pigs, and their father had a little money set aside so that someday he could afford to buy a hay loader—a wonderful machine most farmers just dreamed of owning. It was also the year that electricity came to their town—although not to them because they couldn’t afford it.

One night when Gordon’s mother was doing her big wash, his father stepped in and took his turn over the washboard and asked his wife to rest and do her knitting. He said, “You spend more time doing the wash than sleeping. Do you think we should break down and get electricity?” Although elated at the prospect, she shed a tear or two as she thought of the hay loader that wouldn’t be bought.

So the electrical line went up their lane that year. Although it was nothing fancy, they acquired a washing machine that worked all day by itself and brilliant lightbulbs that dangled from each ceiling. There were no more lamps to fill with oil, no more wicks to cut, no more sooty chimneys to wash. The lamps went quietly off to the attic.

The coming of electricity to their farm was almost the last good thing that happened to them that year. Just as their crops were starting to come through the ground, the rains started. When the water finally receded, there wasn’t a plant left anywhere. They planted again, but more rains beat the crops into the earth. Their potatoes rotted in the mud. They sold a couple of cows and all the pigs and other livestock they had intended to keep, getting very low prices for them because everybody else had to do the same thing. All they harvested that year was a patch of turnips which had somehow weathered the storms.

Then it was Thanksgiving again. Their mother said, “Maybe we’d better forget it this year. We haven’t even got a goose left.”

On Thanksgiving morning, however, Gordon’s father showed up with a jackrabbit and asked his wife to cook it. Grudgingly she started the job, indicating it would take a long time to cook that tough old thing. When it was finally on the table with some of the turnips that had survived, the children refused to eat. Gordon’s mother cried, and then his father did a strange thing. He went up to the attic, got an oil lamp, took it back to the table, and lighted it. He told the children to turn out the electric lights. When there was only the lamp again, they could hardly believe that it had been that dark before. They wondered how they had ever seen anything without the bright lights made possible by electricity.

The food was blessed, and everyone ate. When dinner was over, they all sat quietly. Wrote Gordon:

“In the humble dimness of the old lamp we were beginning to see clearly again. …

“It [was] a lovely meal. The jack rabbit tasted like turkey and the turnips were the mildest we could recall. …

“… [Our] home … , for all its want, was so rich [to] us.”

My brothers and sisters, to express gratitude is gracious and honorable, to enact gratitude is generous and noble, but to live with gratitude ever in our hearts is to touch heaven.

As I close this morning, it is my prayer that in addition to all else for which we are grateful, we may ever reflect our gratitude for our Lord and Savior, Jesus Christ. His glorious gospel provides answers to life’s greatest questions: Where did we come from? Why are we here? Where do our spirits go when we die? That gospel brings to those who live in darkness the light of divine truth.

He taught us how to pray. He taught us how to live. He taught us how to die. His life is a legacy of love. The sick He healed; the downtrodden He lifted; the sinner He saved.

Ultimately, He stood alone. Some Apostles doubted; one betrayed Him. The Roman soldiers pierced His side. The angry mob took His life. There yet rings from Golgotha’s hill His compassionate words: “Father, forgive them; for they know not what they do.”

Who was this “man of sorrows, … acquainted with grief”? “Who is this King of glory,” this Lord of lords? He is our Master. He is our Savior. He is the Son of God. He is the Author of Our Salvation. He beckons, “Follow me.” He instructs, “Go, and do thou likewise.” He pleads, “Keep my commandments.”

Let us follow Him. Let us emulate His example. Let us obey His words. By so doing, we give to Him the divine gift of gratitude.

My sincere, heartfelt prayer is that we may in our individual lives reflect that marvelous virtue of gratitude. May it permeate our very souls, now and evermore. In the sacred name of Jesus Christ, our Savior, amen.

\newpage
\settalk{The Three Rs of Choice}
\setspeaker{Thomas S. Monson}
\settalkdate{October 2010}
\talkheading{The Three Rs of Choice}{Thomas S. Monson}

My beloved brethren of the priesthood, my earnest prayer tonight is that I might enjoy the help of our Heavenly Father in giving utterance to those things which I feel impressed to share with you.

I have been thinking recently about choices and their consequences. Scarcely an hour of the day goes by but what we are called upon to make choices of one sort or another. Some are trivial, some more far-reaching. Some will make no difference in the eternal scheme of things, and others will make all the difference.

As I’ve contemplated the various aspects of choice, I’ve put them into three categories: first, the right of choice; second, the responsibility of choice; and third, the results of choice. I call these the three Rs of choice.

I mention first the right of choice. I am so grateful to a loving Heavenly Father for His gift of agency, or the right to choose. President David O. McKay, ninth President of the Church, said, “Next to the bestowal of life itself, the right to direct that life is God’s greatest gift to man.”

We know that we had our agency before this world was and that Lucifer attempted to take it from us. He had no confidence in the principle of agency or in us and argued for imposed salvation. He insisted that with his plan none would be lost, but he seemed not to recognize—or perhaps not to care—that in addition, none would be any wiser, any stronger, any more compassionate, or any more grateful if his plan were followed.

We who chose the Savior’s plan knew that we would be embarking on a precarious, difficult journey, for we walk the ways of the world and sin and stumble, cutting us off from our Father. But the Firstborn in the Spirit offered Himself as a sacrifice to atone for the sins of all. Through unspeakable suffering He became the great Redeemer, the Savior of all mankind, thus making possible our successful return to our Father.

The prophet Lehi tells us: “Wherefore, men are free according to the flesh; and all things are given them which are expedient unto man. And they are free to choose liberty and eternal life, through the great Mediator of all men, or to choose captivity and death, according to the captivity and power of the devil; for he seeketh that all men might be miserable like unto himself.”

Brethren, within the confines of whatever circumstances we find ourselves, we will always have the right to choose.

Next, with the right of choice comes the responsibility to choose. We cannot be neutral; there is no middle ground. The Lord knows this; Lucifer knows this. As long as we live upon this earth, Lucifer and his hosts will never abandon the hope of claiming our souls.

Our Heavenly Father did not launch us on our eternal journey without providing the means whereby we could receive from Him God-given guidance to assist in our safe return at the end of mortal life. I speak of prayer. I speak too of the whisperings from that still, small voice within each of us, and I do not overlook the holy scriptures, written by mariners who successfully sailed the seas we too must cross.

Each of us has come to this earth with all the tools necessary to make correct choices. The prophet Mormon tells us, “The Spirit of Christ is given to every man, that he may know good from evil.”

We are surrounded—even at times bombarded—by the messages of the adversary. Listen to some of them; they are no doubt familiar to you: “Just this once won’t matter.” “Don’t worry; no one will know.” “You can stop smoking or drinking or taking drugs any time you want.” “Everybody’s doing it, so it can’t be that bad.” The lies are endless.

Although in our journey we will encounter forks and turnings in the road, we simply cannot afford the luxury of a detour from which we may never return. Lucifer, that clever pied piper, plays his lilting melody and attracts the unsuspecting away from the safety of their chosen pathway, away from the counsel of loving parents, away from the security of God’s teachings. He seeks not just the so-called refuse of humanity; he seeks all of us, including the very elect of God. King David listened, wavered, and then followed and fell. So did Cain in an earlier era and Judas Iscariot in a later one. Lucifer’s methods are cunning; his victims, numerous.

We read of him in 2 Nephi: “Others will he pacify, and lull them away into carnal security.” “Others he flattereth away, and telleth them there is no hell … until he grasps them with his awful chains.” “And thus the devil cheateth their souls, and leadeth them away carefully down to hell.”

When faced with significant choices, how do we decide? Do we succumb to the promise of momentary pleasure? To our urges and passions? To the pressure of our peers?

Let us not find ourselves as indecisive as is Alice in Lewis Carroll’s classic Alice’s Adventures in Wonderland. You will remember that she comes to a crossroads with two paths before her, each stretching onward but in opposite directions. She is confronted by the Cheshire cat, of whom Alice asks, “Which path shall I follow?”

The cat answers, “That depends where you want to go. If you do not know where you want to go, it doesn’t matter which path you take.”

Unlike Alice, we all know where we want to go, and it does matter which way we go, for by choosing our path, we choose our destination.

Decisions are constantly before us. To make them wisely, courage is needed—the courage to say no, the courage to say yes. Decisions do determine destiny.

I plead with you to make a determination right here, right now, not to deviate from the path which will lead to our goal: eternal life with our Father in Heaven. Along that straight and true path there are other goals: missionary service, temple marriage, Church activity, scripture study, prayer, temple work. There are countless worthy goals to reach as we travel through life. Needed is our commitment to reach them.

Finally, brethren, I speak of the results of choice. All of our choices have consequences, some of which have little or nothing to do with our eternal salvation and others of which have everything to do with it.

Whether you wear a green T-shirt or a blue one makes no difference in the long run. However, whether you decide to push a key on your computer which will take you to pornography can make all the difference in your life. You will have just taken a step off the straight, safe path. If a friend pressures you to drink alcohol or to try drugs and you succumb to the pressure, you are taking a detour from which you may not return. Brethren, whether we are 12-year-old deacons or mature high priests, we are susceptible. May we keep our eyes, our hearts, and our determination focused on that goal which is eternal and worth any price we will have to pay, regardless of the sacrifice we must make to reach it.

No temptation, no pressure, no enticing can overcome us unless we allow such. If we make the wrong choice, we have no one to blame but ourselves. President Brigham Young once expressed this truth by relating it to himself. Said he: “If Brother Brigham shall take a wrong track, and be shut out of the Kingdom of heaven, no person will be to blame but Brother Brigham. I am the only being in heaven, earth, or hell, that can be blamed.” He continued: “This will equally apply to every Latter-day Saint. Salvation is an individual operation.”

The Apostle Paul has assured us, “There hath no temptation taken you but such as is common to man: but God is faithful, who will not suffer you to be tempted above that ye are able; but will with the temptation also make a way to escape, that ye may be able to bear it.”

We have all made incorrect choices. If we have not already corrected such choices, I assure you that there is a way to do so. The process is called repentance. I plead with you to correct your mistakes. Our Savior died to provide you and me that blessed gift. Although the path is not easy, the promise is real: “Though your sins be as scarlet, they shall be as white as snow.” “And I, the Lord, remember them no more.” Don’t put your eternal life at risk. If you have sinned, the sooner you begin to make your way back, the sooner you will find the sweet peace and joy that come with the miracle of forgiveness.

Brethren, you are of a noble birthright. Eternal life in the kingdom of our Father is your goal. Such a goal is not achieved in one glorious attempt but rather is the result of a lifetime of righteousness, an accumulation of wise choices, even a constancy of purpose. As with anything really worthwhile, the reward of eternal life requires effort.

The scriptures are clear:

“Ye shall observe to do … as the Lord your God hath commanded you: ye shall not turn aside to the right hand or to the left.

“Ye shall walk in all the ways which the Lord your God hath commanded you.”

In closing may I share with you an example of one who determined early in life what his goals would be. I speak of Brother Clayton M. Christensen, a member of the Church who is a professor of business administration in the business school at Harvard University.

When he was 16 years old, Brother Christensen decided, among other things, that he would not play sports on Sunday. Years later, when he attended Oxford University in England, he played center on the basketball team. That year they had an undefeated season and went through to the British equivalent of what in the United States would be the NCAA basketball tournament.

They won their games fairly easily in the tournament, making it to the final four. It was then that Brother Christensen looked at the schedule and, to his absolute horror, saw that the final basketball game was scheduled to be played on a Sunday. He and the team had worked so hard to get where they were, and he was the starting center. He went to his coach with his dilemma. His coach was unsympathetic and told Brother Christensen he expected him to play in the game.

Prior to the final game, however, there was a semifinal game. Unfortunately, the backup center dislocated his shoulder, which increased the pressure on Brother Christensen to play in the final game. He went to his hotel room. He knelt down. He asked his Heavenly Father if it would be all right, just this once, if he played that game on Sunday. He said that before he had finished praying, he received the answer: “Clayton, what are you even asking me for? You know the answer.”

He went to his coach, telling him how sorry he was that he wouldn’t be playing in the final game. Then he went to the Sunday meetings in the local ward while his team played without him. He prayed mightily for their success. They did win.

That fateful, difficult decision was made more than 30 years ago. Brother Christensen has said that as time has passed, he considers it one of the most important decisions he ever made. It would have been very easy to have said, “You know, in general, keeping the Sabbath day holy is the right commandment, but in my particular extenuating circumstance, it’s okay, just this once, if I don’t do it.” However, he says his entire life has turned out to be an unending stream of extenuating circumstances, and had he crossed the line just that once, then the next time something came up that was so demanding and critical, it would have been so much easier to cross the line again. The lesson he learned is that it is easier to keep the commandments 100 percent of the time than it is 98 percent of the time.

My beloved brethren, may we be filled with gratitude for the right of choice, accept the responsibility of choice, and ever be conscious of the results of choice. As bearers of the priesthood, all of us united as one can qualify for the guiding influence of our Heavenly Father as we choose carefully and correctly. We are engaged in the work of the Lord Jesus Christ. We, like those of olden times, have answered His call. We are on His errand. We shall succeed in the solemn charge, “Be ye clean, that bear the vessels of the Lord.” That this may be so is my solemn and humble prayer, in the name of Jesus Christ, our Master, amen.

\newpage
\settalk{Till We Meet Again}
\setspeaker{Thomas S. Monson}
\settalkdate{October 2010}
\talkheading{Till We Meet Again}{Thomas S. Monson}

My brothers and sisters, my heart is full as we bring to a close this wonderful general conference of the Church. We have been spiritually fed as we have listened to the counsel and testimonies of those who have participated in each session. I am certain I speak for all members everywhere when I express deep appreciation for the truths we have been taught. We could echo the words, found in the Book of Mormon, of those who heard the sermon of the great King Benjamin and “cried with one voice, saying: Yea, we believe all the words which thou hast spoken unto us; and also, we know of their surety and truth, because of the Spirit of the Lord Omnipotent.”

I hope that we will take the time to read the conference talks, which will be reprinted in the November issue of the Ensign and Liahona magazines, for they are deserving of our careful study.

What a blessing it is that we have been able to meet here, in this magnificent Conference Center, in peace and comfort and safety. We have had unprecedented coverage of the conference, reaching across the continents and the oceans to people everywhere. Though we are far removed from many of you, we feel of your spirit and send our love and appreciation to you.

To our Brethren who have been released at this conference, may I express the heartfelt gratitude of all of us for your many years of devoted service. Countless are those who have been blessed by your contributions to the work of the Lord.

The Tabernacle Choir and other choirs which participated in the sessions have provided truly heavenly music that has enhanced and beautified all else which has taken place. I thank you for sharing with us your musical talents and abilities.

I love and appreciate my faithful counselors, President Henry B. Eyring and President Dieter F. Uchtdorf. They are truly men of wisdom and understanding, and their service is invaluable. I could not do all that I am called upon to do without their support and assistance. I love and admire my Brethren of the Quorum of the Twelve Apostles and all in the Quorums of the Seventy and in the Presiding Bishopric. They serve selflessly and effectively. I similarly express my appreciation for the women and men who serve as general auxiliary officers.

How blessed we are to have the restored gospel of Jesus Christ. It provides answers to questions concerning where we came from, why we are here, and where we will go when we pass from this life. It provides meaning and purpose and hope to our lives.

We live in a troubled world, a world of many challenges. We are here on this earth to deal with our individual challenges to the best of our ability, to learn from them, and to overcome them. Endure to the end we must, for our goal is eternal life in the presence of our Father in Heaven. He loves us and wants nothing more than for us to succeed in this goal. He will help us and bless us as we call upon Him in our prayers, as we study His words, and as we obey His commandments. Therein is found safety; therein is found peace.

May God bless you, my brothers and sisters. I thank you for your prayers in my behalf and in behalf of all of the General Authorities. We are deeply grateful for you and for all that you do to further the kingdom of God on earth.

May heaven’s blessings be with you. May your homes be filled with love and courtesy and with the Spirit of the Lord. May you constantly nourish your testimonies of the gospel, that they will be a protection to you against the buffetings of Satan.

Conference is now over. As we return to our homes, may we do so safely. May the spirit we have felt here be and abide with us as we go about those things which occupy us each day. May we show increased kindness toward one another; may we ever be found doing the work of the Lord.

I love you; I pray for you. I bid you farewell till we meet again in six months’ time. In the name of our Lord and Savior, even Jesus Christ, amen.

\newpage
\settalk{At Parting}
\setspeaker{Thomas S. Monson}
\settalkdate{April 2011}
\talkheading{At Parting}{Thomas S. Monson}

My brothers and sisters, my heart is full as we come to the close of this conference. We have felt the Spirit of the Lord in rich abundance. I express my appreciation and that of members of the Church everywhere to each one who has participated, including those who have offered prayers. May we long remember the messages we have heard. As we receive the issues of the Ensign and Liahona magazines which will contain these messages in written form, may we read and study them.

Once again the music in all of the sessions has been wonderful. I express my personal gratitude for those willing to share with us their talents, touching and inspiring us in the process.

We have sustained, by uplifted hand, Brethren who have been called to new positions during this conference. We want them to know that we look forward to working with them in the cause of the Master.

I express my love and appreciation for my devoted counselors, President Henry B. Eyring and President Dieter F. Uchtdorf. They are men of wisdom and understanding. Their service is invaluable. I love and support my Brethren of the Quorum of the Twelve Apostles. They serve most effectively, and they are completely dedicated to the work. I also express my love to the members of the Seventy and the Presiding Bishopric.

We face many challenges in the world today, but I assure you that our Heavenly Father is mindful of us. He loves each of us and will bless us as we seek Him through prayer and strive to keep His commandments.

We are a global church. Our membership is found throughout the world. May we be good citizens of the nations in which we live and good neighbors in our communities, reaching out to those of other faiths as well as to those of our own. May we be examples of honesty and integrity wherever we go and in whatever we do.

Thank you for your prayers in my behalf, brothers and sisters, and in behalf of all of the General Authorities of the Church. We are deeply grateful for you and for all that you do to further the work of the Lord.

As you return to your homes, may you do so safely. May the blessings of heaven be upon you.

Now, before we leave today, may I share with you my love for the Savior and for His great atoning sacrifice for us. In three weeks’ time the entire Christian world will be celebrating Easter. I believe that none of us can conceive the full import of what Christ did for us in Gethsemane, but I am grateful every day of my life for His atoning sacrifice in our behalf.

At the last moment, He could have turned back. But He did not. He passed beneath all things that He might save all things. In doing so, He gave us life beyond this mortal existence. He reclaimed us from the Fall of Adam.

To the depths of my very soul, I am grateful to Him. He taught us how to live. He taught us how to die. He secured our salvation.

As I close, may I share with you touching words written by Emily Harris which describe so well my feelings as Easter comes:

\begin{verse}
The linen which once held Him is empty.\\
It lies there,\\
Fresh and white and clean.\\
The door stands opened.\\
The stone is rolled away,\\
And I can almost hear the angels singing His praises.\\
Linen cannot hold Him.\\
Stone cannot hold Him.\\
The words echo through the empty limestone chamber,\\
“He is not here.”\\
The linen which once held Him is now empty.\\
It lies there,\\
Fresh and white and clean\\
And oh, hallelujah, it is empty.
\end{verse}

Blessings to you, my brothers and sisters. In the name of Jesus Christ, our Savior, amen.

\newpage
\settalk{It’s Conference Once Again}
\setspeaker{Thomas S. Monson}
\settalkdate{April 2011}
\talkheading{It’s Conference Once Again}{Thomas S. Monson}

When this building was planned, we thought we’d never fill it. Just look at it now.

My beloved brothers and sisters, how good it is to be together once again as we begin the 181st Annual General Conference of The Church of Jesus Christ of Latter-day Saints.

The past six months seem to have passed rapidly as I’ve been busy with many responsibilities. One of the great blessings during this time was to rededicate the beautiful Laie Hawaii Temple, which had been undergoing extensive renovations for nearly two years. I was accompanied by President and Sister Henry B. Eyring, Elder and Sister Quentin L. Cook, and Elder and Sister William R. Walker. During the evening prior to the rededication, which took place during November, we watched 2,000 young people from the temple district as they filled the Cannon Activities Center on the BYU–Hawaii campus and performed for us. Their production was titled “The Gathering Place” and creatively and masterfully recounted significant events in local Church history and the history of the temple. What a wonderful evening it was!

The following day was a spiritual feast as the temple was rededicated in three sessions. The Spirit of the Lord was with us in rich abundance.

We continue to build temples. It is my privilege this morning to announce three additional temples for which sites are being acquired and which, in coming months and years, will be built in the following locations: Fort Collins, Colorado; Meridian, Idaho; and Winnipeg, Manitoba, Canada. They will certainly be a blessing to our members in those areas.

Each year millions of ordinances are performed in the temples. May we continue to be faithful in performing such ordinances, not only for ourselves but also for our deceased loved ones who are unable to do so for themselves.

The Church continues to provide humanitarian aid in times of disaster. Most recently our hearts and our help have gone out to Japan following the devastating earthquake and tsunami and the resultant nuclear challenges. We have distributed over 70 tons of supplies, including food, water, blankets, bedding, hygiene items, clothing, and fuel. Our young single adults have volunteered their time to locate missing members using the Internet, social media, and other modern means of communication. Members are delivering aid via scooters provided by the Church to areas that are difficult to reach by car. Service projects to assemble hygiene kits and cleaning kits are being organized in multiple stakes and wards in Tokyo, Nagoya, and Osaka. Thus far, over 40,000 hours of service have been donated by more than 4,000 volunteers. Our help will be ongoing in Japan and in any other areas where there is need.

My brothers and sisters, I thank you for your faith and devotion to the gospel, for the love and care you show to one another, and for the service you provide in your wards and branches and stakes and districts. Thank you, as well, for your faithfulness in paying your tithes and offerings and for your generosity in contributing to the other funds of the Church.

As of the end of the year 2010, there were 52,225 missionaries serving in 340 missions throughout the world. Missionary work is the lifeblood of the kingdom. May I suggest that if you are able, you might consider making a contribution to the General Missionary Fund of the Church.

Now, brothers and sisters, we are anxious to listen to the messages which will be presented to us today and tomorrow. Those who will address us have sought heaven’s help and direction as they have prepared their messages. That we may be filled with the Spirit of the Lord and be uplifted and inspired as we listen and learn is my prayer. In the name of Jesus Christ, amen.

\newpage
\settalk{Priesthood Power}
\setspeaker{Thomas S. Monson}
\settalkdate{April 2011}
\talkheading{Priesthood Power}{Thomas S. Monson}

I prayed and studied long about what I might say tonight. I wish not to offend anyone. I thought, “What are the challenges we have? What do I deal with every day that causes me to weep sometimes late into the night?” I thought that I would try to address a few of those challenges tonight. Some will apply to the young men. Some will apply to those who are middle-aged. Some will apply to those who are a little bit above middle age. We don’t talk about old age.

And so I simply want to begin by declaring, it has been good for us to be together this evening. We’ve heard wonderful and timely messages concerning the priesthood of God. I, with you, have been uplifted and inspired.

Tonight I wish to address matters which have been much on my mind of late and which I have felt impressed to share with you. In one way or another, they all relate to the personal worthiness required to receive and exercise the sacred power of the priesthood which we hold.

May I begin by reciting to you from section 121 of the Doctrine and Covenants:

“The rights of the priesthood are inseparably connected with the powers of heaven, and … the powers of heaven cannot be controlled nor handled only upon the principles of righteousness.

“That they may be conferred upon us, it is true; but when we undertake to cover our sins, or to gratify our pride, our vain ambition, or to exercise control or dominion or compulsion upon the souls of the children of men, in any degree of unrighteousness, behold, the heavens withdraw themselves; the Spirit of the Lord is grieved; and when it is withdrawn, Amen to the priesthood or the authority of that man.”

Brethren, that is the definitive word of the Lord concerning His divine authority. We cannot be in doubt as to the obligation this places upon each of us who bear the priesthood of God.

We have come to the earth in troubled times. The moral compass of the masses has gradually shifted to an “almost anything goes” position.

I’ve lived long enough to have witnessed much of the metamorphosis of society’s morals. Where once the standards of the Church and the standards of society were mostly compatible, now there is a wide chasm between us, and it’s growing ever wider.

Many movies and television shows portray behavior which is in direct opposition to the laws of God. Do not subject yourself to the innuendo and outright filth which are so often found there. The lyrics in much of today’s music fall in the same category. The profanity so prevalent around us today would never have been tolerated in the not-too-distant past. Sadly, the Lord’s name is taken in vain over and over again. Recall with me the commandment—one of the ten—which the Lord revealed to Moses on Mount Sinai: “Thou shalt not take the name of the Lord thy God in vain; for the Lord will not hold him guiltless that taketh his name in vain.” I am sorry that any of us is subjected to profane language, and I plead with you not to use it. I implore you not to say or to do anything of which you cannot be proud.

Stay completely away from pornography. Do not allow yourself to view it, ever. It has proven to be an addiction which is more than difficult to overcome. Avoid alcohol and tobacco or any other drugs, also addictions which you would be hard-pressed to conquer.

What will protect you from the sin and evil around you? I maintain that a strong testimony of our Savior and of His gospel will help see you through to safety. If you have not read the Book of Mormon, read it. I will not ask for a show of hands. If you do so prayerfully and with a sincere desire to know the truth, the Holy Ghost will manifest its truth to you. If it is true—and it is—then Joseph Smith was a prophet who saw God the Father and His Son, Jesus Christ. The Church is true. If you do not already have a testimony of these things, do that which is necessary to obtain one. It is essential for you to have your own testimony, for the testimonies of others will carry you only so far. Once obtained, a testimony needs to be kept vital and alive through obedience to the commandments of God and through regular prayer and scripture study. Attend church. You young men, attend seminary or institute if such is available to you.

Should there be anything amiss in your life, there is open to you a way out. Cease any unrighteousness. Talk with your bishop. Whatever the problem, it can be worked out through proper repentance. You can become clean once again. Said the Lord, speaking of those who repent, “Though your sins be as scarlet, they shall be as white as snow,” “and I, the Lord, remember them no more.”

The Savior of mankind described Himself as being in the world but not of the world. We also can be in the world but not of the world as we reject false concepts and false teachings and remain true to that which God has commanded.

Now, I have thought a lot lately about you young men who are of an age to marry but who have not yet felt to do so. I see lovely young ladies who desire to be married and to raise families, and yet their opportunities are limited because so many young men are postponing marriage.

This is not a new situation. Much has been said concerning this matter by past Presidents of the Church. I share with you just one or two examples of their counsel.

Said President Harold B. Lee, “We are not doing our duty as holders of the priesthood when we go beyond the marriageable age and withhold ourselves from an honorable marriage to these lovely women.”

President Gordon B. Hinckley said this: “My heart reaches out to … our single sisters, who long for marriage and cannot seem to find it. … I have far less sympathy for the young men, who under the customs of our society have the prerogative to take the initiative in these matters but in so many cases fail to do so.”

I realize there are many reasons why you may be hesitating to take that step of getting married. If you are concerned about providing financially for a wife and family, may I assure you that there is no shame in a couple having to scrimp and save. It is generally during these challenging times that you will grow closer together as you learn to sacrifice and to make difficult decisions. Perhaps you are afraid of making the wrong choice. To this I say that you need to exercise faith. Find someone with whom you can be compatible. Realize that you will not be able to anticipate every challenge which may arise, but be assured that almost anything can be worked out if you are resourceful and if you are committed to making your marriage work.

Perhaps you are having a little too much fun being single, taking extravagant vacations, buying expensive cars and toys, and just generally enjoying the carefree life with your friends. I’ve encountered groups of you running around together, and I admit that I’ve wondered why you aren’t out with the young ladies.

Brethren, there is a point at which it’s time to think seriously about marriage and to seek a companion with whom you want to spend eternity. If you choose wisely and if you are committed to the success of your marriage, there is nothing in this life which will bring you greater happiness.

When you marry, brethren, you will wish to marry in the house of the Lord. For you who hold the priesthood, there should be no other option. Be careful lest you destroy your eligibility to be so married. You can keep your courtship within proper bounds while still having a wonderful time.

Now, brethren, I turn to another subject about which I feel impressed to address you. In the three years since I was sustained as President of the Church, I believe the saddest and most discouraging responsibility I have each week is the handling of cancellations of sealings. Each one was preceded by a joyous marriage in the house of the Lord, where a loving couple was beginning a new life together and looking forward to spending the rest of eternity with each other. And then months and years go by, and for one reason or another, love dies. It may be the result of financial problems, lack of communication, uncontrolled tempers, interference from in-laws, entanglement in sin. There are any number of reasons. In most cases divorce does not have to be the outcome.

The vast majority of requests for cancellations of sealings come from women who tried desperately to make a go of the marriage but who, in the final analysis, could not overcome the problems.

Choose a companion carefully and prayerfully; and when you are married, be fiercely loyal one to another. Priceless advice comes from a small framed plaque I once saw in the home of an uncle and aunt. It read, “Choose your love; love your choice.” There is great wisdom in those few words. Commitment in marriage is absolutely essential.

Your wife is your equal. In marriage neither partner is superior nor inferior to the other. You walk side by side as a son and a daughter of God. She is not to be demeaned or insulted but should be respected and loved. Said President Gordon B. Hinckley: “Any man in this Church who … exercises unrighteous dominion over [his wife] is unworthy to hold the priesthood. Though he may have been ordained, the heavens will withdraw, the Spirit of the Lord will be grieved, and it will be amen to the authority of the priesthood of that man.”

President Howard W. Hunter said this about marriage: “Being happily and successfully married is generally not so much a matter of marrying the right person as it is being the right person.” I like that. “The conscious effort to do one’s part fully is the greatest element contributing to success.”

Many years ago in the ward over which I presided as the bishop, there lived a couple who often had very serious, heated disagreements. I mean real disagreements. Each of the two was certain of his or her position. Neither one would yield to the other. When they weren’t arguing, they maintained what I would call an uneasy truce.

One morning at 2:00 a.m. I had a telephone call from the couple. They wanted to talk to me, and they wanted to talk right then. I dragged myself from bed, dressed, and went to their home. They sat on opposite sides of the room, not speaking to each other. The wife communicated with her husband by talking to me. He replied to her by talking to me. I thought, “How in the world are we going to get this couple together?”

I prayed for inspiration, and the thought came to me to ask them a question. I said, “How long has it been since you have been to the temple and witnessed a temple sealing?” They admitted it had been a very long time. They were otherwise worthy people who held temple recommends and who went to the temple and did ordinance work for others.

I said to them, “Will you come with me to the temple on Wednesday morning at 8:00? We will witness a sealing ceremony there.”

In unison they asked, “Whose ceremony?”

I responded, “I don’t know. It will be for whoever is getting married that morning.”

On the following Wednesday at the appointed hour, we met at the Salt Lake Temple. The three of us went into one of the beautiful sealing rooms, not knowing a soul in the room except Elder ElRay L. Christiansen, then an Assistant to the Quorum of the Twelve, a General Authority position which existed at that time. Elder Christiansen was scheduled to perform a sealing ceremony for a bride and groom in that very room that morning. I am confident the bride and her family thought, “These must be friends of the groom” and that the groom’s family thought, “These must be friends of the bride.” My couple were seated on a little bench with about a full two feet (0.6 m) of space between them.

Elder Christiansen began by providing counsel to the couple who were being married, and he did so in a beautiful fashion. He mentioned how a husband should love his wife, how he should treat her with respect and courtesy, honoring her as the heart of the home. Then he talked to the bride about how she should honor her husband as the head of the home and be of support to him in every way.

I noticed that as Elder Christiansen spoke to the bride and the groom, my couple moved a little closer together. Soon they were seated right next to one another. What pleased me is that they had both moved at about the same rate. By the end of the ceremony, my couple were sitting as close to each other as though they were the newlyweds. Each was smiling.

We left the temple that day, and no one ever knew who we were or why we had come, but my friends were holding hands as they walked out the front door. Their differences had been set aside. I had not had to say one word. You see, they remembered their own wedding day and the covenants they had made in the house of God. They were committed to beginning again and trying harder this time around.

If any of you are having difficulty in your marriage, I urge you to do all that you can to make whatever repairs are necessary, that you might be as happy as you were when your marriage started out. We who are married in the house of the Lord do so for time and for all eternity, and then we must put forth the necessary effort to make it so. I realize that there are situations where marriages cannot be saved, but I feel strongly that for the most part they can be and should be. Do not let your marriage get to the point where it is in jeopardy.

President Hinckley taught that it is up to each of us who hold the priesthood of God to discipline ourselves so that we stand above the ways of the world. It is essential that we be honorable and decent men. Our actions must be above reproach.

The words we speak, the way we treat others, and the way we live our lives all impact our effectiveness as men and boys holding the priesthood.

The gift of the priesthood is priceless. It carries with it the authority to act as God’s servants, to administer to the sick, to bless our families, and to bless others as well. Its authority can reach beyond the veil of death, on into the eternities. There is nothing else to compare with it in all this world. Safeguard it, treasure it, live worthy of it.

My beloved brethren, may righteousness guide our every step as we journey through life. Today and always, may we be worthy recipients of the divine power of the priesthood we bear. May it bless our lives and may we use it to bless the lives of others, as did He who lived and died for us—even Jesus Christ, our Lord and Savior. This is my prayer in His sacred name, His holy name, amen.

\newpage
\settalk{The Holy Temple—a Beacon to the World}
\setspeaker{Thomas S. Monson}
\settalkdate{April 2011}
\talkheading{The Holy Temple—a Beacon to the World}{Thomas S. Monson}

My beloved brothers and sisters, I extend my love and greetings to each of you and pray that our Heavenly Father will guide my thoughts and inspire my words as I speak to you today.

May I begin by making a comment or two concerning the fine messages we have heard this morning from Sister Allred and Bishop Burton and others pertaining to the Church’s welfare program. As indicated, this year marks the 75th anniversary of this inspired program, which has blessed the lives of so many. It was my privilege to know personally some of those who pioneered this great endeavor—men of compassion and foresight.

As both Bishop Burton and Sister Allred and others mentioned, the bishop of the ward is given the responsibility to care for those in need who reside within the boundaries of his ward. Such was my privilege when I presided as a very young bishop in Salt Lake City over a ward of 1,080 members, including 84 widows. There were many who needed assistance. How grateful I was for the welfare program of the Church and for the help of the Relief Society and the priesthood quorums.

I declare that the welfare program of The Church of Jesus Christ of Latter-day Saints is inspired of Almighty God.

Now, my brothers and sisters, this conference marks three years since I was sustained as President of the Church. Of course they have been busy years, filled with many challenges but also with countless blessings. The opportunity I have had to dedicate and rededicate temples has been among the most enjoyable and sacred of these blessings, and it is concerning the temple that I wish to speak to you today.

During the October general conference in 1902, Church President Joseph F. Smith expressed in his opening address the hope that one day we would “have temples built in the various parts of the [world] where they are needed for the convenience of the people.”

During the first 150 years following the organization of the Church, from 1830 to 1980, 21 temples were built, including the temples in Kirtland, Ohio, and Nauvoo, Illinois. Contrast that with the 30 years since 1980, during which 115 temples were built and dedicated. With the announcement yesterday of 3 new temples, there are additionally 26 temples either under construction or in preconstruction stages. These numbers will continue to grow.

The goal President Joseph F. Smith hoped for in 1902 is becoming a reality. Our desire is to make the temple as accessible as possible to our members.

One of the temples currently under construction is in Manaus, Brazil. Many years ago I read of a group of over a hundred members who left Manaus, located in the heart of the Amazon rain forest, to travel to what was then the closest temple, located in São Paulo, Brazil—nearly 2,500 miles (4,000 km) from Manaus. Those faithful Saints journeyed by boat for four days on the Amazon River and its tributaries. After completing this journey by water, they boarded buses for another three days of travel—over bumpy roads, with very little to eat, and with nowhere comfortable to sleep. After seven days and nights, they arrived at the temple in São Paulo, where ordinances eternal in nature were performed. Of course their return journey was just as difficult. However, they had received the ordinances and blessings of the temple, and although their purses were empty, they themselves were filled with the spirit of the temple and with gratitude for the blessings they had received. Now, many years later, our members in Manaus are rejoicing as they watch their own temple take shape on the banks of the Rio Negro. Temples bring joy to our faithful members wherever they are built.

Reports of the sacrifices made in order to receive the blessings found only in temples of God never fail to touch my heart and bring to me a renewed sense of thankfulness for temples.

May I share with you the account of Tihi and Tararaina Mou Tham and their 10 children. The entire family, except for one daughter, joined the Church in the early 1960s, when missionaries came to their island, located about 100 miles (160 km) south of Tahiti. Soon they began to desire the blessings of an eternal family sealing in the temple.

At that time the nearest temple to the Mou Tham family was the Hamilton New Zealand Temple, more than 2,500 miles (4,000 km) to the southwest, accessible only by expensive airplane travel. The large Mou Tham family, which eked out a meager living on a small plantation, had no money for airplane fare, nor was there any opportunity for employment on their Pacific island. So Brother Mou Tham and his son Gérard made the difficult decision to travel 3,000 miles (4,800 km) to work in New Caledonia, where another son was already employed.

The three Mou Tham men labored for four years. Brother Mou Tham alone returned home only once during that time, for the marriage of a daughter.

After four years, Brother Mou Tham and his sons had saved enough money to take the family to the New Zealand Temple. All who were members went except for one daughter, who was expecting a baby. They were sealed for time and eternity, an indescribable and joyful experience.

Brother Mou Tham returned from the temple directly to New Caledonia, where he worked for two more years to pay for the passage of the one daughter who had not been at the temple with them—a married daughter and her child and husband.

In their later years Brother and Sister Mou Tham desired to serve in the temple. By that time the Papeete Tahiti Temple had been constructed and dedicated, and they served four missions there.

My brothers and sisters, temples are more than stone and mortar. They are filled with faith and fasting. They are built of trials and testimonies. They are sanctified by sacrifice and service.

The first temple to be built in this dispensation was the temple at Kirtland, Ohio. The Saints at the time were impoverished, and yet the Lord had commanded that a temple be built, so build it they did. Wrote Elder Heber C. Kimball of the experience, “The Lord only knows the scenes of poverty, tribulation and distress which we passed through to accomplish it.” And then, after all that had been painstakingly completed, the Saints were forced to leave Ohio and their beloved temple. They eventually found refuge—although it would be temporary—on the banks of the Mississippi River in the state of Illinois. They named their settlement Nauvoo, and willing to give their all once again and with their faith intact, they erected another temple to their God. Persecutions raged, however, and with the Nauvoo Temple barely completed, they were driven from their homes once again, seeking refuge in a desert.

The struggle and the sacrifice began once again as they labored for 40 years to erect the Salt Lake Temple, which stands majestically on the block just south of those of us who are here today in the Conference Center.

Some degree of sacrifice has ever been associated with temple building and with temple attendance. Countless are those who have labored and struggled in order to obtain for themselves and for their families the blessings which are found in the temples of God.

Why are so many willing to give so much in order to receive the blessings of the temple? Those who understand the eternal blessings which come from the temple know that no sacrifice is too great, no price too heavy, no struggle too difficult in order to receive those blessings. There are never too many miles to travel, too many obstacles to overcome, or too much discomfort to endure. They understand that the saving ordinances received in the temple that permit us to someday return to our Heavenly Father in an eternal family relationship and to be endowed with blessings and power from on high are worth every sacrifice and every effort.

Today most of us do not have to suffer great hardships in order to attend the temple. Eighty-five percent of the membership of the Church now live within 200 miles (320 km) of a temple, and for a great many of us, that distance is much shorter.

If you have been to the temple for yourselves and if you live within relatively close proximity to a temple, your sacrifice could be setting aside the time in your busy lives to visit the temple regularly. There is much to be done in our temples in behalf of those who wait beyond the veil. As we do the work for them, we will know that we have accomplished what they cannot do for themselves. President Joseph F. Smith, in a mighty declaration, stated, “Through our efforts in their behalf their chains of bondage will fall from them, and the darkness surrounding them will clear away, that light may shine upon them and they shall hear in the spirit world of the work that has been done for them by their children here, and will rejoice with you in your performance of these duties.” My brothers and sisters, the work is ours to do.

In my own family, some of our most sacred and treasured experiences have occurred when we have joined together in the temple to perform sealing ordinances for our deceased ancestors.

If you have not yet been to the temple or if you have been but currently do not qualify for a recommend, there is no more important goal for you to work toward than being worthy to go to the temple. Your sacrifice may be bringing your life into compliance with what is required to receive a recommend, perhaps by forsaking long-held habits which disqualify you. It may be having the faith and the discipline to pay your tithing. Whatever it is, qualify to enter the temple of God. Secure a temple recommend and regard it as a precious possession, for such it is.

Until you have entered the house of the Lord and have received all the blessings which await you there, you have not obtained everything the Church has to offer. The all-important and crowning blessings of membership in the Church are those blessings which we receive in the temples of God.

Now, my young friends who are in your teenage years, always have the temple in your sights. Do nothing which will keep you from entering its doors and partaking of the sacred and eternal blessings there. I commend those of you who already go to the temple regularly to perform baptisms for the dead, arising in the very early hours of the morning so you can participate in such baptisms before school begins. I can think of no better way to start a day.

To you parents of young children, may I share with you some sage advice from President Spencer W. Kimball. Said he: “It would be a fine thing if … parents would have in every bedroom in their house a picture of the temple so [their children] from the time [they are] infant[s] could look at the picture every day [until] it becomes a part of [their lives]. When [they reach] the age that [they need] to make [the] very important decision [concerning going to the temple], it will have already been made.”

Our children sing in Primary:

\begin{verse}
I love to see the temple.\\
I’ll go inside someday.\\
I’ll cov’nant with my Father;\\
I’ll promise to obey.
\end{verse}

I plead with you to teach your children of the temple’s importance.

The world can be a challenging and difficult place in which to live. We are often surrounded by that which would drag us down. As you and I go to the holy houses of God, as we remember the covenants we make within, we will be more able to bear every trial and to overcome each temptation. In this sacred sanctuary we will find peace; we will be renewed and fortified.

Now, my brothers and sisters, may I mention one more temple before I close. In the not-too-distant future as new temples take shape around the world, one will rise in a city which came into being over 2,500 years ago. I speak of the temple which is now being built in Rome, Italy.

Every temple is a house of God, filling the same functions and with identical blessings and ordinances. The Rome Italy Temple, uniquely, is being built in one of the most historic locations in the world, a city where the ancient Apostles Peter and Paul preached the gospel of Christ and where each was martyred.

Last October, as we gathered on a lovely pastoral site in the northeast corner of Rome, it was my opportunity to offer a prayer of dedication as we prepared to break the ground. I felt impressed to call upon Italian senator Lucio Malan and Rome’s vice-mayor Giuseppe Ciardi to be among the first to turn a shovelful of earth. Each had been a part of the decision to allow us to build a temple in their city.

The day was overcast but warm, and although rain threatened, not more than a drop or two fell. As the magnificent choir sang in Italian the beautiful strains of “The Spirit of God,” one felt as though heaven and earth were joined in a glorious hymn of praise and gratitude to Almighty God. Tears could not be restrained.

In a coming day, the faithful in this, the Eternal City, will receive ordinances eternal in nature in a holy house of God.

I express my undying gratitude to my Heavenly Father for the temple now being built in Rome and for all of our temples, wherever they are. Each one stands as a beacon to the world, an expression of our testimony that God, our Eternal Father, lives, that He desires to bless us and, indeed, to bless His sons and daughters of all generations. Each of our temples is an expression of our testimony that life beyond the grave is as real and as certain as is our life here on earth. I so testify.

My beloved brothers and sisters, may we make whatever sacrifices are necessary to attend the temple and to have the spirit of the temple in our hearts and in our homes. May we follow in the footsteps of our Lord and Savior, Jesus Christ, who made the ultimate sacrifice for us, that we might have eternal life and exaltation in our Heavenly Father’s kingdom. This is my sincere prayer, and I offer it in the name of our Savior, Jesus Christ the Lord, amen.

\newpage
\settalk{As We Meet Again}
\setspeaker{Thomas S. Monson}
\settalkdate{October 2011}
\talkheading{As We Meet Again}{Thomas S. Monson}

It is good, brothers and sisters, to welcome you to the 181st Semiannual General Conference of The Church of Jesus Christ of Latter-day Saints.

This conference marks 48 years—think of it, 48 years—since I was called to the Quorum of the Twelve Apostles by President David O. McKay. That was in October of 1963. It seems impossible that so many years have come and gone since then.

When we’re busy, time seems to pass far too quickly, and the past six months have been no exception for me. One of the highlights during that period was the opportunity I had to rededicate the Atlanta Georgia Temple on May 1. I was accompanied by Elder and Sister M. Russell Ballard, Elder and Sister Walter F. González, and Elder and Sister William R. Walker.

During the cultural celebration entitled “Southern Lights,” held the evening prior to the rededication, 2,700 young men and young women from throughout the temple district performed. It was one of the most outstanding programs I have seen and had the audience on its feet several times for standing ovations.

The following day the temple was rededicated in two sessions, where the Spirit of the Lord was with us in rich abundance.

During the latter part of August, President Henry B. Eyring dedicated the San Salvador El Salvador Temple. He was accompanied by Sister Eyring and by Elder and Sister D. Todd Christofferson, Elder and Sister William R. Walker, and Sister Silvia H. Allred of the Relief Society general presidency and her husband, Jeffry. President Eyring reported that it was a most spiritual event.

In the latter part of this year, President Dieter F. Uchtdorf and Sister Uchtdorf will travel with other General Authorities to Quetzaltenango, Guatemala, where he will dedicate our temple there.

The building of temples continues uninterrupted, brothers and sisters. Today it is my privilege to announce several new temples.

First, may I mention that no Church-built facility is more important than a temple. Temples are places where relationships are sealed together to last through the eternities. We are grateful for all the many temples across the world and for the blessing they are in the lives of our members.

Late last year the Provo Tabernacle in Utah County was seriously damaged by a terrible fire. This wonderful building, much beloved by generations of Latter-day Saints, was left with only the exterior walls standing. After careful study, we have decided to rebuild it, with full preservation and restoration of the exterior, to become the second temple of the Church in the city of Provo. The existing Provo Temple is one of the busiest in the Church, and a second temple there will accommodate the increasing numbers of faithful Church members who are attending the temple from Provo and the surrounding communities.

I am also pleased to announce new temples in the following locations: Barranquilla, Colombia; Durban, South Africa; Kinshasa in the Democratic Republic of the Congo; and Star Valley, Wyoming. In addition, we are moving forward on our plans for a temple to be built in Paris, France.

Details of these temples will be provided in the future as site and other necessary approvals are obtained.

I have mentioned in previous conferences the progress we are making in placing temples closer to our members. Although they are readily available to many members in the Church, there are still areas of the world where temples are so distant from our members that they cannot afford the travel required to get to them. They are thus unable to partake of the sacred and eternal blessings temples provide. To help in this regard, we have available what is called the General Temple Patron Assistance Fund. This fund provides a one-time visit to the temple for those who otherwise would not be able to go to the temple and yet who long desperately for that opportunity. Any who might wish to contribute to this fund can simply write in the information on the normal contribution slip which is given to the bishop each month.

Now, brothers and sisters, it is my prayer that we may be filled with the Spirit of the Lord as we listen to the messages today and tomorrow and learn those things the Lord would have us know. This I pray for in the name of Jesus Christ, amen.

\newpage
\settalk{Dare to Stand Alone}
\setspeaker{Thomas S. Monson}
\settalkdate{October 2011}
\talkheading{Dare to Stand Alone}{Thomas S. Monson}

My beloved brethren, it is a tremendous privilege to be with you tonight. We who hold the priesthood of God form a great bond and brotherhood.

We read in the Doctrine and Covenants, section 121, verse 36, “that the rights of the priesthood are inseparably connected with the powers of heaven.” What a wonderful gift we have been given—to hold the priesthood, which is “inseparably connected with the powers of heaven.” This precious gift, however, brings with it not only special blessings but also solemn responsibilities. We must conduct our lives so that we are ever worthy of the priesthood we bear. We live in a time when we are surrounded by much that is intended to entice us into paths which may lead to our destruction. To avoid such paths requires determination and courage.

I recall a time—and some of you here tonight will also—when the standards of most people were very similar to our standards. No longer is this true. I recently read an article in the New York Times concerning a study which took place during the summer of 2008. A distinguished Notre Dame sociologist led a research team in conducting in-depth interviews with 230 young adults across America. I believe we can safely assume that the results would be similar in most parts of the world.

I share with you just a portion of this very telling article:

“The interviewers asked open-ended questions about right and wrong, moral dilemmas and the meaning of life. In the rambling answers, … you see the young people groping to say anything sensible on these matters. But they just don’t have the categories or vocabulary to do so.

“When asked to describe a moral dilemma they had faced, two-thirds of the young people either couldn’t answer the question or described problems that are not moral at all, like whether they could afford to rent a certain apartment or whether they had enough quarters to feed the meter at a parking spot.”

The article continues:

“The default position, which most of them came back to again and again, is that moral choices are just a matter of individual taste. ‘It’s personal,’ the respondents typically said. ‘It’s up to the individual. Who am I to say?’

“Rejecting blind deference to authority, many of the young people have gone off to the other extreme [saying]: ‘I would do what I thought made me happy or how I felt. I have no other way of knowing what to do but how I internally feel.’”

Those who conducted the interviews emphasized that the majority of the young people with whom they spoke had “not been given the resources—by schools, institutions [or] families—to cultivate their moral intuitions.”

Brethren, none within the sound of my voice should be in any doubt concerning what is moral and what is not, nor should any be in doubt about what is expected of us as holders of the priesthood of God. We have been and continue to be taught God’s laws. Despite what you may see or hear elsewhere, these laws are unchanging.

As we go about living from day to day, it is almost inevitable that our faith will be challenged. We may at times find ourselves surrounded by others and yet standing in the minority or even standing alone concerning what is acceptable and what is not. Do we have the moral courage to stand firm for our beliefs, even if by so doing we must stand alone? As holders of the priesthood of God, it is essential that we are able to face—with courage—whatever challenges come our way. Remember the words of Tennyson: “My strength is as the strength of ten, because my heart is pure.”

Increasingly, some celebrities and others who—for one reason or another—are in the public eye have a tendency to ridicule religion in general and, at times, the Church in particular. If our testimonies are not firmly enough rooted, such criticisms can cause us to doubt our own beliefs or to waver in our resolves.

In Lehi’s vision of the tree of life, found in 1 Nephi 8, Lehi sees, among others, those who hold to the iron rod until they come forth and partake of the fruit of the tree of life, which we know is a representation of the love of God. And then, sadly, after they partake of the fruit, some are ashamed because of those in the “great and spacious building,” who represent the pride of the children of men, who are pointing fingers at them and scoffing at them; and they fall away into forbidden paths and are lost. What a powerful tool of the adversary is ridicule and mockery! Again, brethren, do we have the courage to stand strong and firm in the face of such difficult opposition?

I believe my first experience in having the courage of my convictions took place when I served in the United States Navy near the end of World War II.

Navy boot camp was not an easy experience for me, nor for anyone who endured it. For the first three weeks I was convinced my life was in jeopardy. The navy wasn’t trying to train me; it was trying to kill me.

I shall ever remember when Sunday rolled around after the first week. We received welcome news from the chief petty officer. Standing at attention on the drill ground in a brisk California breeze, we heard his command: “Today everybody goes to church—everybody, that is, except for me. I am going to relax!” Then he shouted, “All of you Catholics, you meet in Camp Decatur—and don’t come back until three o’clock. Forward, march!” A rather sizeable contingent moved out. Then he barked out his next command: “Those of you who are Jewish, you meet in Camp Henry—and don’t come back until three o’clock. Forward, march!” A somewhat smaller contingent marched out. Then he said, “The rest of you Protestants, you meet in the theaters at Camp Farragut—and don’t come back until three o’clock. Forward, march!”

Instantly there flashed through my mind the thought, “Monson, you are not a Catholic; you are not a Jew; you are not a Protestant. You are a Mormon, so you just stand here!” I can assure you that I felt completely alone. Courageous and determined, yes—but alone.

And then I heard the sweetest words I ever heard that chief petty officer utter. He looked in my direction and asked, “And just what do you guys call yourselves?” Until that very moment I had not realized that anyone was standing beside me or behind me on the drill ground. Almost in unison, each of us replied, “Mormons!” It is difficult to describe the joy that filled my heart as I turned around and saw a handful of other sailors.

The chief petty officer scratched his head in an expression of puzzlement but finally said, “Well, you guys go find somewhere to meet. And don’t come back until three o’clock. Forward, march!”

As we marched away, I thought of the words of a rhyme I had learned in Primary years before:

\begin{verse}
Dare to be a Mormon;\\
Dare to stand alone.\\
Dare to have a purpose firm;\\
Dare to make it known.
\end{verse}

Although the experience turned out differently from what I had expected, I had been willing to stand alone, had such been necessary.

Since that day, there have been times when there was no one standing behind me and so I did stand alone. How grateful I am that I made the decision long ago to remain strong and true, always prepared and ready to defend my religion, should the need arise.

Lest we at any time feel inadequate for the tasks ahead for us, brethren, may I share with you a statement made in 1987 by then Church President Ezra Taft Benson as he addressed a large group of members in California. Said President Benson:

“In all ages prophets have looked down through the corridors of time to our day. Billions of the deceased and those yet to be born have their eyes on us. Make no mistake about it—you are a marked generation.”

“For nearly six thousand years, God has held you in reserve to make your appearance in the final days before the second coming of the Lord. Some individuals will fall away, but the kingdom of God will remain intact to welcome the return of its Head—even Jesus Christ.

“While [this] generation will be comparable in wickedness to the days of Noah, when the Lord cleansed the earth by flood, there is a major difference this time: [it is that] God has saved for the final inning some of His strongest … children, who will help bear off the kingdom triumphantly.”

Yes, brethren, we represent some of His strongest children. Ours is the responsibility to be worthy of all the glorious blessings our Father in Heaven has in store for us. Wherever we go, our priesthood goes with us. Are we standing in holy places? Please, before you put yourself and your priesthood in jeopardy by venturing into places or participating in activities which are not worthy of you or of that priesthood, pause to consider the consequences. Each of us has had conferred upon him the Aaronic Priesthood. In the process, each received the power which holds the keys to the ministering of angels. Said President Gordon B. Hinckley:

“You cannot afford to do anything that would place a curtain between you and the ministering of angels in your behalf.

“You cannot be immoral in any sense. You cannot be dishonest. You cannot cheat or lie. You cannot take the name of God in vain or use filthy language and still have the right to the ministering of angels.”

If any of you has stumbled in your journey, I want you to understand without any question whatsoever that there is a way back. The process is called repentance. Our Savior gave His life to provide you and me that blessed gift. Despite the fact that the repentance path is not easy, the promises are real. We have been told: “Though your sins be as scarlet, they shall be as white as snow.” “And I will remember [them] no more.” What a statement. What a blessing. What a promise.

There may be those of you who are thinking to yourselves, “Well, I’m not living all the commandments, and I’m not doing everything I should, and yet my life is going along just fine. I think I can have my cake and eat it too.” Brethren, I promise you that this will not work in the long run.

Not too many months ago I received a letter from a man who once thought he could have it both ways. He has now repented and has brought his life into compliance with gospel principles and commandments. I want to share with you a paragraph from his letter, for it represents the reality of flawed thinking: “I have had to learn for myself (the hard way) that the Savior was absolutely correct when He said, ‘No man can serve two masters: for either he will hate the one, and love the other; or else he will hold to the one, and despise the other. Ye cannot serve God and mammon.’ I tried, about as hard as anyone ever has, to do both. In the end,” said he, “I had all of the emptiness, darkness, and loneliness that Satan provides to those who believe his deceptions, illusions, and lies.”

In order for us to be strong and to withstand all the forces pulling us in the wrong direction or all the voices encouraging us to take the wrong path, we must have our own testimony. Whether you are 12 or 112—or anywhere in between—you can know for yourself that the gospel of Jesus Christ is true. Read the Book of Mormon. Ponder its teachings. Ask Heavenly Father if it is true. We have the promise that “if ye shall ask with a sincere heart, with real intent, having faith in Christ, he will manifest the truth of it unto you, by the power of the Holy Ghost.”

When we know the Book of Mormon is true, then it follows that Joseph Smith was indeed a prophet and that he saw God the Eternal Father and His Son, Jesus Christ. It also follows that the gospel was restored in these latter days through Joseph Smith—including the restoration of both the Aaronic and Melchizedek Priesthoods.

Once we have a testimony, it is incumbent upon us to share that testimony with others. Many of you brethren have served as missionaries throughout the world. Many of you young men will yet serve. Prepare yourselves now for that opportunity. Make certain you are worthy to serve.

If we are prepared to share the gospel, we are ready to respond to the counsel of the Apostle Peter, who urged, “Be ready always to give an answer to every man that asketh you a reason of the hope that is in you.”

We will have opportunities throughout our lives to share our beliefs, although we don’t always know when we will be called upon to do so. Such an opportunity came to me in 1957, when I worked in the publishing business and was asked to go to Dallas, Texas, sometimes called “the city of churches,” to address a business convention. Following the conclusion of the convention, I took a sightseeing bus ride through the city’s suburbs. As we passed the various churches, our driver would comment, “On the left you see the Methodist church” or “There on the right is the Catholic cathedral.”

As we passed a beautiful red brick building situated upon a hill, the driver exclaimed, “That building is where the Mormons meet.” A lady in the rear of the bus called out, “Driver, can you tell us something more about the Mormons?”

The driver pulled the bus over to the side of the road, turned around in his seat, and replied, “Lady, all I know about the Mormons is that they meet in that red brick building. Is there anyone on this bus who knows anything more about the Mormons?”

I waited for someone to respond. I gazed at the expression on each person’s face for some sign of recognition, some desire to comment. Nothing. I realized it was up to me to do as the Apostle Peter suggested, to “be ready always to give an answer to every man that asketh you a reason of the hope that is in you.” I also realized the truth of the adage “When the time for decision arrives, the time for preparation is past.”

For the next 15 or so minutes, I had the privilege of sharing with those on the bus my testimony concerning the Church and our beliefs. I was grateful for my testimony and grateful that I was prepared to share it.

With all my heart and soul, I pray that every man who holds the priesthood will honor that priesthood and be true to the trust which was conveyed when it was conferred. May each of us who holds the priesthood of God know what he believes. May we ever be courageous and prepared to stand for what we believe, and if we must stand alone in the process, may we do so courageously, strengthened by the knowledge that in reality we are never alone when we stand with our Father in Heaven.

As we contemplate the great gift we have been given—“the rights of the priesthood … inseparably connected with the powers of heaven”—may our determination ever be to guard and defend it and to be worthy of its great promises. Brethren, may we follow the Savior’s instruction to us found in the book of 3 Nephi: “Hold up your light that it may shine unto the world. Behold I am the light which ye shall hold up—that which ye have seen me do.”

That we may ever follow that light and hold it up for all the world to see is my prayer and my blessing upon all who hear my voice, in the name of Jesus Christ, amen.

\newpage
\settalk{Stand in Holy Places}
\setspeaker{Thomas S. Monson}
\settalkdate{October 2011}
\talkheading{Stand in Holy Places}{Thomas S. Monson}

My beloved brothers and sisters, we have heard fine messages this morning, and I commend each who has participated. We’re particularly delighted to have Elder Robert D. Hales with us once again and feeling improved. We love you, Bob.

As I pondered what I would like to say to you this morning, I have felt impressed to share certain thoughts and feelings which I consider to be pertinent and timely. I pray that I may be guided in my remarks.

I have lived on this earth for 84 years now. To give you a little perspective, I was born the same year Charles Lindbergh flew the first solo nonstop flight from New York to Paris in a single-engine, single-seat monoplane. Much has changed during the 84 years since then. Man has long since been to the moon and back. In fact, yesterday’s science fiction has become today’s reality. And that reality, thanks to the technology of our times, is changing so fast we can barely keep up with it—if we do at all. For those of us who remember dial telephones and manual typewriters, today’s technology is more than merely amazing.

Also evolving at a rapid rate has been the moral compass of society. Behaviors which once were considered inappropriate and immoral are now not only tolerated but also viewed by ever so many as acceptable.

I recently read in the Wall Street Journal an article by Jonathan Sacks, Britain’s chief rabbi. Among other things, he writes: “In virtually every Western society in the 1960s there was a moral revolution, an abandonment of its entire traditional ethic of self-restraint. All you need, sang the Beatles, is love. The Judeo-Christian moral code was jettisoned. In its place came [the adage]: [Do] whatever works for you. The Ten Commandments were rewritten as the Ten Creative Suggestions.”

Rabbi Sacks goes on to lament:

“We have been spending our moral capital with the same reckless abandon that we have been spending our financial capital. …

“There are large parts of [the world] where religion is a thing of the past and there is no counter-voice to the culture of buy it, spend it, wear it, flaunt it, because you’re worth it. The message is that morality is passé, conscience is for wimps, and the single overriding command is ‘Thou shalt not be found out.’”

My brothers and sisters, this—unfortunately—describes much of the world around us. Do we wring our hands in despair and wonder how we’ll ever survive in such a world? No. Indeed, we have in our lives the gospel of Jesus Christ, and we know that morality is not passé, that our conscience is there to guide us, and that we are responsible for our actions.

Although the world has changed, the laws of God remain constant. They have not changed; they will not change. The Ten Commandments are just that—commandments. They are not suggestions. They are every bit as requisite today as they were when God gave them to the children of Israel. If we but listen, we hear the echo of God’s voice, speaking to us here and now:

“Thou shalt have no other gods before me.

“Thou shalt not make unto thee any graven image. …

“Thou shalt not take the name of the Lord thy God in vain. …

“Remember the sabbath day, to keep it holy. …

“Honour thy father and thy mother. …

“Thou shalt not kill.

“Thou shalt not commit adultery.

“Thou shalt not steal.

“Thou shalt not bear false witness. …

“Thou shalt not covet.”

Our code of conduct is definitive; it is not negotiable. It is found not only in the Ten Commandments but also in the Sermon on the Mount, given to us by the Savior when He walked upon the earth. It is found throughout His teachings. It is found in the words of modern revelation.

Our Father in Heaven is the same yesterday, today, and forever. The prophet Mormon tells us that God is “unchangeable from all eternity to all eternity.” In this world where nearly everything seems to be changing, His constancy is something on which we can rely, an anchor to which we can hold fast and be safe, lest we be swept away into uncharted waters.

It may appear to you at times that those out in the world are having much more fun than you are. Some of you may feel restricted by the code of conduct to which we in the Church adhere. My brothers and sisters, I declare to you, however, that there is nothing which can bring more joy into our lives or more peace to our souls than the Spirit which can come to us as we follow the Savior and keep the commandments. That Spirit cannot be present at the kinds of activities in which so much of the world participates. The Apostle Paul declared the truth: “The natural man receiveth not the things of the Spirit of God: for they are foolishness unto him: neither can he know them, because they are spiritually discerned.” The term natural man can refer to any of us if we allow ourselves to be so.

We must be vigilant in a world which has moved so far from that which is spiritual. It is essential that we reject anything that does not conform to our standards, refusing in the process to surrender that which we desire most: eternal life in the kingdom of God. The storms will still beat at our doors from time to time, for they are an inescapable part of our existence in mortality. We, however, will be far better equipped to deal with them, to learn from them, and to overcome them if we have the gospel at our core and the love of the Savior in our hearts. The prophet Isaiah declared, “The work of righteousness shall be peace; and the effect of righteousness quietness and assurance for ever.”

As a means of being in the world but not being of the world, it is necessary that we communicate with our Heavenly Father through prayer. He wants us to do so; He’ll answer our prayers. The Savior admonished us, as recorded in 3 Nephi 18, to “watch and pray always lest ye enter into temptation; for Satan desireth to have you. …

“Therefore ye must always pray unto the Father in my name;

“And whatsoever ye shall ask the Father in my name, which is right, believing that ye shall receive, behold it shall be given unto you.”

I gained my testimony of the power of prayer when I was about 12 years old. I had worked hard to earn some money and had managed to save five dollars. This was during the Great Depression, when five dollars was a substantial sum of money—especially for a boy of 12. I gave all my coins, which totaled five dollars, to my father, and he gave me in return a five-dollar bill. I know there was something specific I planned to purchase with the five dollars, although all these years later I can’t recall what it was. I just remember how important that money was to me.

At the time, we did not own a washing machine, so my mother would send to the laundry each week our clothes which needed to be washed. After a couple of days, a load of what we called “wet wash” would be returned to us, and Mother would hang the items on our clothesline out back to dry.

I had tucked my five-dollar bill in the pocket of my jeans. As you can probably guess, my jeans were sent to the laundry with the money still in the pocket. When I realized what had happened, I was sick with worry. I knew that pockets were routinely checked at the laundry prior to washing. If my money was not discovered and taken during that process, I knew it was almost certain the money would be dislodged during washing and would be claimed by a laundry worker who would have no idea to whom the money should be returned, even if he had the inclination to do so. The chances of getting back my five dollars were extremely remote—a fact which my dear mother confirmed when I told her I had left the money in my pocket.

I wanted that money; I needed that money; I had worked very hard to earn that money. I realized there was only one thing I could do. In my extremity I turned to my Father in Heaven and pleaded with Him to keep my money safe in that pocket somehow until our wet wash came back.

Two very long days later, when I knew it was about time for the delivery truck to bring our wash, I sat by the window, waiting. As the truck pulled up to the curb, my heart was pounding. As soon as the wet clothes were in the house, I grabbed my jeans and ran to my bedroom. I reached into the pocket with trembling hands. When I didn’t find anything immediately, I thought all was lost. And then my fingers touched that wet five-dollar bill. As I pulled it from the pocket, relief flooded over me. I offered a heartfelt prayer of gratitude to my Father in Heaven, for I knew that He had answered my prayer.

Since that time of long ago, I have had countless prayers answered. Not a day has gone by that I have not communicated with my Father in Heaven through prayer. It is a relationship I cherish—one I would literally be lost without. If you do not now have such a relationship with your Father in Heaven, I urge you to work toward that goal. As you do so, you will be entitled to His inspiration and guidance in your life—necessities for each of us if we are to survive spiritually during our sojourn here on earth. Such inspiration and guidance are gifts He freely gives if we but seek them. What treasures they are!

I am always humbled and grateful when my Heavenly Father communicates with me through His inspiration. I have learned to recognize it, to trust it, and to follow it. Time and time again I have been the recipient of such inspiration. One rather dramatic experience took place in August of 1987 during the dedication of the Frankfurt Germany Temple. President Ezra Taft Benson had been with us for the first day or two of the dedication but had returned home, and so it became my opportunity to conduct the remaining sessions.

On Saturday we had a session for our Dutch members who were in the Frankfurt Temple district. I was well acquainted with one of our outstanding leaders from the Netherlands, Brother Peter Mourik. Just prior to the session, I had the distinct impression that Brother Mourik should be called upon to speak to his fellow Dutch members during the session and that, in fact, he should be the first speaker. Not having seen him in the temple that morning, I passed a note to Elder Carlos E. Asay, our Area President, asking whether Peter Mourik was in attendance at the session. Just prior to standing up to begin the session, I received a note back from Elder Asay indicating that Brother Mourik was actually not in attendance, that he was involved elsewhere, and that he was planning to attend the dedicatory session in the temple the following day with the servicemen stakes.

As I stood at the pulpit to welcome the people and to outline the program, I received unmistakable inspiration once again that I was to announce Peter Mourik as the first speaker. This was counter to all my instincts, for I had just heard from Elder Asay that Brother Mourik was definitely not in the temple. Trusting in the inspiration, however, I announced the choir presentation and the prayer and then indicated that our first speaker would be Brother Peter Mourik.

As I returned to my seat, I glanced toward Elder Asay; I saw on his face a look of alarm. He later told me that when I had announced Brother Mourik as the first speaker, he couldn’t believe his ears. He said he knew that I had received his note and that I indeed had read it, and he couldn’t fathom why I would then announce Brother Mourik as a speaker, knowing he wasn’t anywhere in the temple.

During the time all of this was taking place, Peter Mourik was in a meeting at the area offices in Porthstrasse. As his meeting was going forward, he suddenly turned to Elder Thomas A. Hawkes Jr., who was then the regional representative, and asked, “How fast can you get me to the temple?”

Elder Hawkes, who was known to drive rather rapidly in his small sports car, answered, “I can have you there in 10 minutes! But why do you need to go to the temple?”

Brother Mourik admitted he did not know why he needed to go to the temple but that he knew he had to get there. The two of them set out for the temple immediately.

During the magnificent choir number, I glanced around, thinking that at any moment I would see Peter Mourik. I did not. Remarkably, however, I felt no alarm. I had a sweet, undeniable assurance that all would be well.

Brother Mourik entered the front door of the temple just as the opening prayer was concluding, still not knowing why he was there. As he hurried down the hall, he saw my image on the monitor and heard me announce, “We will now hear from Brother Peter Mourik.”

To the astonishment of Elder Asay, Peter Mourik immediately walked into the room and took his place at the podium.

Following the session, Brother Mourik and I discussed that which had taken place prior to his opportunity to speak. I have pondered the inspiration which came that day not only to me but also to Peter Mourik. That remarkable experience has provided an undeniable witness to me of the importance of being worthy to receive such inspiration and then trusting it—and following it—when it comes. I know without question that the Lord intended for those who were present at that session of the Frankfurt Temple dedication to hear the powerful, touching testimony of His servant Brother Peter Mourik.

My beloved brothers and sisters, communication with our Father in Heaven—including our prayers to Him and His inspiration to us—is necessary in order for us to weather the storms and trials of life. The Lord invites us, “Draw near unto me and I will draw near unto you; seek me diligently and ye shall find me.” As we do so, we will feel His Spirit in our lives, providing us the desire and the courage to stand strong and firm in righteousness—to “stand … in holy places, and be not moved.”

As the winds of change swirl around us and the moral fiber of society continues to disintegrate before our very eyes, may we remember the Lord’s precious promise to those who trust in Him: “Fear thou not; for I am with thee: be not dismayed; for I am thy God: I will strengthen thee; yea, I will help thee; yea, I will uphold thee with the right hand of my righteousness.”

What a promise! May such be our blessing, I sincerely pray in the sacred name of our Lord and Savior, Jesus Christ, amen.

\newpage
\settalk{Until We Meet Again}
\setspeaker{Thomas S. Monson}
\settalkdate{October 2011}
\talkheading{Until We Meet Again}{Thomas S. Monson}

My brothers and sisters, I know you will agree with me that this has been a most inspiring conference. We have felt the Spirit of the Lord in rich abundance these past two days as our hearts have been touched and our testimonies of this divine work have been strengthened. We express thanks to each one who has participated, including those Brethren offering prayers.

We are all here because we love the Lord and want to serve Him. I testify to you that our Heavenly Father is mindful of us. I acknowledge His hand in all things.

Once again the music has been wonderful, and I express my personal gratitude and that of the entire Church to those willing to share with us their talents in this regard.

We express our deep appreciation to those Brethren who have been released during this conference. They have served faithfully and well and have made significant contributions to the work of the Lord.

I express profound appreciation to my faithful and dedicated counselors and thank them publicly for the support and assistance they provide to me. They are truly men of wisdom and understanding, and their service is invaluable.

I thank my brethren of the Quorum of the Twelve for their most able and untiring service in the work of the Lord. Likewise I express my gratitude to the members of the Quorums of the Seventy and to the Presiding Bishopric for their selfless and effective service. I similarly express my appreciation for the women and men who serve as general auxiliary officers.

Brothers and sisters, I assure you that our Heavenly Father is mindful of the challenges we face in the world today. He loves each of us and will bless us as we strive to keep His commandments and seek Him through prayer.

How blessed we are to have the restored gospel of Jesus Christ. It provides answers to the questions concerning where we came from, why we are here, and where we will go when we depart from this life. It gives meaning and purpose and hope to our lives.

I thank you for the service you so willingly give to one another. We are God’s hands here on this earth, with a mandate to love and to serve His children.

I thank you for all that you do in your wards and your branches. I express my gratitude for your willingness to serve in the positions to which you are called, whatever they may be. Each is important in furthering the work of the Lord.

Conference is now over. As we return to our homes, may we do so safely. May we find all has been well during our absence. May the spirit we have felt here be and abide with us as we go about those things which occupy us each day. May we show increased kindness one toward another. May we ever be found doing the work of the Lord.

May heaven’s blessings be with you. May your homes be filled with harmony and love. May you constantly nourish your testimonies, that they might be a protection for you against the adversary.

As your humble servant, I desire with all my heart to do God’s will and to serve Him and to serve you.

I love you; I pray for you. I would ask once again that you would remember me and all the General Authorities in your prayers. We are one with you in moving forward this marvelous work. I testify to you that we are all in this together and that every man, woman, and child has a part to play. May God give us the strength and the ability and the determination to play our part well.

I bear my testimony to you that this work is true, that our Savior lives, and that He guides and directs His Church here upon the earth. I leave with you my witness and my testimony that God our Eternal Father lives and loves us. He is indeed our Father, and He is personal and real. May we realize and understand how close to us He is willing to come, how far He is willing to go to help us, how much He loves us, and how much He does and is willing to do for us.

May He bless you. May His promised peace be with you now and always.

I say farewell to you until we meet again in six months’ time, and I do so in the name of Jesus Christ, our Savior and Redeemer, amen.

\newpage
\settalk{As We Close This Conference}
\setspeaker{Thomas S. Monson}
\settalkdate{April 2012}
\talkheading{As We Close This Conference}{Thomas S. Monson}

My heart is full as we come to the close of this glorious conference. We have been so richly blessed as we have listened to the counsel and testimonies of those who have spoken to us. I think you will agree with me that we have felt the Spirit of the Lord as our hearts have been touched and our testimonies strengthened.

Once again we have enjoyed beautiful music, which has enhanced and enriched each session of conference. I express my gratitude to all who have shared with us their talents in this regard.

My heartfelt thanks go to each who has spoken to us as well as to those who have offered prayers at each of the sessions.

There are countless individuals who work either behind the scenes or in less visible positions each conference. It would not be possible for us to hold these sessions without their assistance. My thanks go to all of them as well.

I know you join with me in expressing profound gratitude to those brethren and sisters who have been released during this conference. We will miss them. Their contributions to the work of the Lord have been enormous and will be felt throughout generations to come.

We have also sustained, through uplifted hands, brethren and sisters who have been called to new positions during this conference. We welcome them and want them to know that we look forward to serving with them in the cause of the Master. They have been called by inspiration from on high.

We have had unprecedented coverage of this conference, reaching across the continents and oceans to people everywhere. Though we are far removed from many of you, we feel of your spirit and your dedication, and we send our love and appreciation to you wherever you are.

How blessed we are, my brothers and sisters, to have the restored gospel of Jesus Christ in our lives and in our hearts. It provides answers to life’s greatest questions. It provides meaning and purpose and hope to our lives.

We live in troubled times. I assure you that our Heavenly Father is mindful of the challenges we face. He loves each of us and desires to bless us and to help us. May we call upon Him in prayer, as He admonished when He said, “Pray always, and I will pour out my Spirit upon you, and great shall be your blessing—yea, even more than if you should obtain treasures of earth.”

My dear brothers and sisters, may your homes be filled with love and courtesy and with the Spirit of the Lord. Love your families. If there are disagreements or contentions among you, I urge you to settle them now. Said the Savior:

“There shall be no disputations among you. …

“For verily, verily I say unto you, he that hath the spirit of contention is not of me, but is of the devil, who is the father of contention, and he stirreth up the hearts of men to contend with anger, one with another.

“[But] behold, this is not my doctrine … ; but this is my doctrine, that such things should be done away.”

As your humble servant, I echo the words of King Benjamin in his address to his people when he said:

“I have not commanded you to … think that I of myself am more than a mortal man.

“But I am like as yourselves, subject to all manner of infirmities in body and mind; yet I have been chosen … by the hand of the Lord … and have been kept and preserved by his matchless power, to serve you with all the might, mind and strength which the Lord hath granted unto me.”

My beloved brothers and sisters, I desire with all my heart to do God’s will and to serve Him and to serve you.

Now as we leave this conference, I invoke the blessings of heaven upon each of you. May you who are away from your homes return to them safely. May you ponder the truths you have heard, and may they help you to become even better than you were when conference began two days ago.

Until we meet again in six months’ time, I ask the Lord’s blessings to be upon you and, indeed, upon all of us, and I do so in His holy name—even Jesus Christ, our Lord and Savior—amen.

\newpage
\settalk{As We Gather Once Again}
\setspeaker{Thomas S. Monson}
\settalkdate{April 2012}
\talkheading{As We Gather Once Again}{Thomas S. Monson}

My beloved brothers and sisters, as we gather once again in a general conference of the Church, I welcome you and express my love to you. We meet each six months to strengthen one another, to extend encouragement, to provide comfort, to build faith. We are here to learn. Some of you may be seeking answers to questions and challenges you are experiencing in your life. Some are struggling with disappointments or losses. Each can be enlightened and uplifted and comforted as the Spirit of the Lord is felt.

Should there be changes which need to be made in your life, may you find the incentive and the courage to do so as you listen to the inspired words which will be spoken. May each of us resolve anew to live so that we are worthy sons and daughters of our Heavenly Father. May we continue to oppose evil wherever it is found.

How blessed we are to have come to earth at such a time as this—a marvelous time in the long history of the world. We can’t all be together under one roof, but we now have the ability to partake of the proceedings of this conference through the wonders of television, radio, cable, satellite transmission, and the Internet—even on mobile devices. We come together as one, speaking many languages, living in many lands, but all of one faith and one doctrine and one purpose.

From a small beginning 182 years ago, our presence is now felt throughout the world. This great cause in which we are engaged will continue to go forth, changing and blessing lives as it does so. No cause, no force in the entire world can stop the work of God. Despite what comes, this great cause will go forward. You recall the prophetic words of the Prophet Joseph Smith: “No unhallowed hand can stop the work from progressing; persecutions may rage, mobs may combine, armies may assemble, calumny may defame, but the truth of God will go forth boldly, nobly, and independent, till it has penetrated every continent, visited every clime, swept every country, and sounded in every ear, till the purposes of God shall be accomplished, and the Great Jehovah shall say the work is done.”

There is much that is difficult and challenging in the world today, my brothers and sisters, but there is also much that is good and uplifting. As we declare in our thirteenth article of faith, “If there is anything virtuous, lovely, or of good report or praiseworthy, we seek after these things.” May we ever continue to do so.

I thank you for your faith and devotion to the gospel. I thank you for the love and care you show one to another. I thank you for the service you provide in your wards and branches and in your stakes and districts. It is such service that enables the Lord to accomplish many of His purposes here upon the earth.

I express my thanks to you for your kindnesses to me wherever I go. I thank you for your prayers in my behalf. I have felt those prayers and am most grateful for them.

Now, my brothers and sisters, we have come to be instructed and inspired. Many messages will be shared during the next two days. I can assure you that those men and women who will address you have sought heaven’s help and direction as they have prepared their messages. They have been inspired concerning that which they will share with us.

Our Heavenly Father is mindful of each of us and our needs. May we be filled with His Spirit as we partake of the proceedings of this conference. This is my sincere prayer in the sacred name of our Lord and Savior, Jesus Christ, amen.

\newpage
\settalk{Believe, Obey, and Endure}
\setspeaker{Thomas S. Monson}
\settalkdate{April 2012}
\talkheading{Believe, Obey, and Endure}{Thomas S. Monson}

My dear young sisters, the responsibility to address you is humbling. I pray for divine help, that I may be made equal to such an opportunity.

A mere 20 years ago you had not yet commenced your journey through mortality. You were still in your heavenly home. There you were among those who loved you and were concerned for your eternal well-being. Eventually, earth life became essential to your progress. Farewells were no doubt spoken, and expressions of confidence given. You gained bodies and became mortal, cut off from the presence of your Heavenly Father.

A joyous welcome, however, awaited you here on earth. Those first years were precious, special years. Satan had no power to tempt you, for you had not yet become accountable. You were innocent before God.

Soon you entered that period some have labeled “the terrible teens.” I prefer “the terrific teens.” What a time of opportunity, a season of growth, a semester of development—marked by the acquisition of knowledge and the quest for truth.

No one has described the teenage years as being easy. They are often years of insecurity, of feeling as though you just don’t measure up, of trying to find your place with your peers, of trying to fit in. This is a time when you are becoming more independent—and perhaps desire more freedom than your parents are willing to give you right now. They are also prime years when Satan will tempt you and will do his utmost to entice you from the path which will lead you back to that heavenly home from which you came and back to your loved ones there and back to your Heavenly Father.

The world around you is not equipped to provide the help you need to make it through this often-treacherous journey. So many in our society today seem to have slipped from the moorings of safety and drifted from the harbor of peace.

Permissiveness, immorality, pornography, drugs, the power of peer pressure—all these and more—cause many to be tossed about on a sea of sin and crushed on the jagged reefs of lost opportunities, forfeited blessings, and shattered dreams.

Is there a way to safety? Is there an escape from threatened destruction? The answer is a resounding yes! I counsel you to look to the lighthouse of the Lord. I have said it before; I will say it again: there is no fog so dense, no night so dark, no gale so strong, no mariner so lost but what the lighthouse of the Lord can rescue. It beckons through the storms of life. It calls, “This way to safety. This way to home.” It sends forth signals of light easily seen and never failing. If followed, those signals will guide you back to your heavenly home.

I wish to talk with you tonight about three essential signals from the Lord’s lighthouse which will help you to return to that Father who eagerly awaits your triumphant homecoming. Those three signals are believe, obey, and endure.

First, I mention a signal which is basic and essential: believe. Believe that you are a daughter of Heavenly Father, that He loves you, and that you are here for a glorious purpose—to gain your eternal salvation. Believe that remaining strong and faithful to the truths of the gospel is of utmost importance. I testify that it is!

My young friends, believe in the words you say each week as you recite the Young Women theme. Think about the meaning of those words. There is truth there. Strive always to live the values which are set forth. Believe, as your theme states, that if you accept and act upon those values, you will be prepared to strengthen your home and your family, to make and keep sacred covenants, to receive the ordinances of the temple, and to eventually enjoy the blessings of exaltation. These are beautiful gospel truths, and by following them, you will be happier throughout your life here and hereafter than you will be if you disregard them.

Most of you were taught the truths of the gospel from the time you were a toddler. You were taught by loving parents and caring teachers. The truths they imparted to you helped you gain a testimony; you believed what you were taught. Although that testimony can continue to be fed spiritually and to grow as you study, as you pray for guidance, and as you attend your Church meetings each week, it is up to you to keep that testimony alive. Satan will try with all his might to destroy it. Throughout your entire life you will need to nurture it. As with the flame of a brightly burning fire, your testimony—if not continually fed—will fade to glowing embers and then cool completely. You must not let this happen.

Besides attending your Sunday meetings and your weeknight activities, when you have the chance to be involved in seminary, whether in the early morning or in released-time classes, take advantage of that opportunity. Many of you are attending seminary now. As with anything in life, much of what you take from your seminary experience depends on your attitude and your willingness to be taught. May your attitude be one of humility and a desire to learn. How grateful I am for the opportunity I had as a teenager to attend early-morning seminary, for it played a vital role in my development and the development of my testimony. Seminary can change lives.

Some years ago I was on a board of directors with a fine man who had been extremely successful in life. I was impressed with his integrity and his loyalty to the Church. I learned that he had gained a testimony and had joined the Church because of seminary. When he married, his wife had been a lifelong member of the Church. He belonged to no church. Through the years and despite her efforts, he showed no interest in attending church with his wife and children. And then he began driving two of his daughters to early-morning seminary. He would remain in the car while they had their class, and then he would drive them to school. One day it was raining, and one of his daughters said, “Come in, Dad. You can sit in the hall.” He accepted the invitation. The door to the classroom was open, and he began to listen. His heart was touched. For the rest of that school year, he attended seminary with his daughters, which led eventually to his membership and a lifetime of activity in the Church. Let seminary help build and strengthen your testimony.

There will be times when you will face challenges which might jeopardize your testimony, or you may neglect it as you pursue other interests. I plead with you to keep it strong. It is your responsibility, and yours alone, to keep its flame burning brightly. Effort is required, but it is effort you will never, ever regret. I’m reminded of the words of a song written by Julie de Azevedo Hanks. Referring to her testimony, she wrote:

\begin{verse}
Through the winds of change\\
Encircled by the clouds of pain\\
I guard it with my life\\
I need the warmth—I need the light\\
Though the storm will rage\\
I stand against the pounding rain\\
I remain\\
A keeper of the flame.
\end{verse}

May you believe and then may you keep the flame of your testimony burning brightly, come what may.

Next, young women, may you obey. Obey your parents. Obey the laws of God. They are given to us by a loving Heavenly Father. When they are obeyed, our lives will be more fulfilling, less complicated. Our challenges and problems will be easier to bear. We will receive the Lord’s promised blessings. He has said, “The Lord requireth the heart and a willing mind; and the willing and obedient shall eat the good of the land of Zion in these last days.”

You have but one life to live. Keep it as free from trouble as you can. You will be tempted, sometimes by individuals you had thought friends.

Some years ago I spoke to a Mia Maid adviser who told me of an experience she had with one of the young women in her class. This young woman had been tempted time and time again to leave the pathway of truth and follow the detour of sin. Through the constant persuasion of some of her friends at school, she had finally agreed to follow such a detour. The plan was set: she would tell her parents she was going to her activity night for Young Women. She planned, however, to be there only long enough for her girlfriends and their dates to pick her up. They would then attend a party where alcoholic beverages would be consumed and where the behavior would be in complete violation of what this young woman knew was right.

The teacher had prayed for inspiration in helping all her girls but especially this particular young woman, who seemed so uncertain about her commitment to the gospel. The teacher had received inspiration that night to abandon what she had previously planned and to speak to the girls about remaining morally clean. As she began sharing her thoughts and feelings, the young woman in question checked her watch often to make sure she didn’t miss her rendezvous with her friends. However, as the discussion progressed, her heart was touched, her conscience awakened, and her determination renewed. When it came, she ignored the repeated sound of the automobile horn summoning her. She remained throughout the evening with her teacher and the other girls in the class. The temptation to detour from God’s approved way had been averted. Satan had been frustrated. The young woman remained after the others had left in order to thank her teacher for the lesson and to let her know how it had helped her avoid what might have been a tragic outcome. A teacher’s prayer had been answered.

I subsequently learned that because she had made her decision not to go with her friends that night—some of the most popular girls and boys at school—the young woman was shunned by them and for many months had no friends at school. They couldn’t accept that she was unwilling to do the things they did. It was an extremely difficult and lonely period for her, but she remained steadfast and eventually gained friends who shared her standards. Now, several years later, she has a temple marriage and four beautiful children. How different her life could have been. Our decisions determine our destiny.

Precious young women, make every decision you contemplate pass this test: “What does it do to me? What does it do for me?” And let your code of conduct emphasize not “What will others think?” but rather “What will I think of myself?” Be influenced by that still, small voice. Remember that one with authority placed his hands on your head at the time of your confirmation and said, “Receive the Holy Ghost.” Open your hearts, even your very souls, to the sound of that special voice which testifies of truth. As the prophet Isaiah promised, “Thine ears shall hear a word … saying, This is the way, walk ye in it.”

The tenor of our times is permissiveness. Magazines and television shows portray the stars of the movie screen, the heroes of the athletic field—those whom many young people long to emulate—as disregarding the laws of God and flaunting sinful practices, seemingly with no ill effect. Don’t you believe it! There is a time of reckoning—even a balancing of the ledger. Every Cinderella has her midnight—if not in this life, then in the next. Judgment Day will come for all. Are you prepared? Are you pleased with your own performance?

If any has stumbled in her journey, I promise you that there is a way back. The process is called repentance. Our Savior died to provide you and me that blessed gift. Though the path is difficult, the promise is real. Said the Lord: “Though your sins be as scarlet, they shall be as white as snow.” “And I will remember [them] no more.”

My beloved young sisters, you have the precious gift of agency. I plead with you to choose to obey.

Finally, may you endure. What does it mean to endure? I love this definition: to withstand with courage. Courage may be necessary for you to believe; it will at times be necessary as you obey. It will most certainly be required as you endure until that day when you will leave this mortal existence.

I have spoken over the years with many individuals who have told me, “I have so many problems, such real concerns. I’m overwhelmed with the challenges of life. What can I do?” I have offered to them, and I now offer to you, this specific suggestion: seek heavenly guidance one day at a time. Life by the yard is hard; by the inch it’s a cinch. Each of us can be true for just one day—and then one more and then one more after that—until we’ve lived a lifetime guided by the Spirit, a lifetime close to the Lord, a lifetime of good deeds and righteousness. The Savior promised, “Look unto me, and endure to the end, and ye shall live; for unto him that endureth to the end will I give eternal life.”

For this purpose have you come into mortality, my young friends. There is nothing more important than the goal you strive to attain—even eternal life in the kingdom of your Father.

You are precious, precious daughters of our Heavenly Father sent to earth at this day and time for a purpose. You have been withheld until this very hour. Wonderful, glorious things are in store for you if you will only believe, obey, and endure. May this be your blessing, I pray in the name of Jesus Christ, our Savior, amen.

\newpage
\settalk{The Race of Life}
\setspeaker{Thomas S. Monson}
\settalkdate{April 2012}
\talkheading{The Race of Life}{Thomas S. Monson}

My beloved brothers and sisters, this morning I wish to speak to you of eternal truths—those truths which will enrich our lives and see us safely home.

Everywhere people are in a hurry. Jet-powered aircraft speed their precious human cargo across broad continents and vast oceans so that business meetings might be attended, obligations met, vacations enjoyed, or families visited. Roadways everywhere—including freeways, thruways, and motorways—carry millions of automobiles, occupied by more millions of people, in a seemingly endless stream and for a multitude of reasons as we rush about the business of each day.

In this fast-paced life, do we ever pause for moments of meditation—even thoughts of timeless truths?

When compared to eternal verities, most of the questions and concerns of daily living are really rather trivial. What should we have for dinner? What color should we paint the living room? Should we sign Johnny up for soccer? These questions and countless others like them lose their significance when times of crisis arise, when loved ones are hurt or injured, when sickness enters the house of good health, when life’s candle dims and darkness threatens. Our thoughts become focused, and we are easily able to determine what is really important and what is merely trivial.

I recently visited with a woman who has been battling a life-threatening disease for over two years. She indicated that prior to her illness, her days were filled with activities such as cleaning her house to perfection and filling it with beautiful furnishings. She visited her hairdresser twice a week and spent money and time each month adding to her wardrobe. Her grandchildren were invited to visit infrequently, for she was always concerned that what she considered her precious possessions might be broken or otherwise ruined by tiny and careless hands.

And then she received the shocking news that her mortal life was in jeopardy and that she might have very limited time left here. She said that at the moment she heard the doctor’s diagnosis, she knew immediately that she would spend whatever time she had remaining with her family and friends and with the gospel at the center of her life, for these represented what was most precious to her.

Such moments of clarity come to all of us at one time or another, although not always through so dramatic a circumstance. We see clearly what it is that really matters in our lives and how we should be living.

Said the Savior:

“Lay not up for yourselves treasures upon earth, where moth and rust doth corrupt, and where thieves break through and steal:

“But lay up for yourselves treasures in heaven, where neither moth nor rust doth corrupt, and where thieves do not break through nor steal:

“For where your treasure is, there will your heart be also.”

In our times of deepest reflection or greatest need, the soul of man reaches heavenward, seeking a divine response to life’s greatest questions: Where did we come from? Why are we here? Where do we go after we leave this life?

Answers to these questions are not discovered within the covers of academia’s textbooks or by checking the Internet. These questions transcend mortality. They embrace eternity.

Where did we come from? This query is inevitably thought, if not spoken, by every human being.

The Apostle Paul told the Athenians on Mars’ Hill that “we are the offspring of God.” Since we know that our physical bodies are the offspring of our mortal parents, we must probe for the meaning of Paul’s statement. The Lord has declared that “the spirit and the body are the soul of man.” Thus it is the spirit which is the offspring of God. The writer of Hebrews refers to Him as “the Father of spirits.” The spirits of all men are literally His “begotten sons and daughters.”

We note that inspired poets have, for our contemplation of this subject, written moving messages and recorded transcendent thoughts. William Wordsworth penned the truth:

\begin{verse}
Our birth is but a sleep and a forgetting:\\
The soul that rises with us, our life’s Star,\\
Hath had elsewhere its setting,\\
And cometh from afar:\\
Not in entire forgetfulness,\\
And not in utter nakedness,\\
But trailing clouds of glory do we come\\
From God, who is our home:\\
Heaven lies about us in our infancy!
\end{verse}

Parents ponder their responsibility to teach, to inspire, and to provide guidance, direction, and example. And while parents ponder, children—and particularly youth—ask the penetrating question, “Why are we here?” Usually it is spoken silently to the soul and phrased, “Why am I here?”

How grateful we should be that a wise Creator fashioned an earth and placed us here, with a veil of forgetfulness of our previous existence so that we might experience a time of testing, an opportunity to prove ourselves in order to qualify for all that God has prepared for us to receive.

Clearly, one primary purpose of our existence upon the earth is to obtain a body of flesh and bones. We have also been given the gift of agency. In a thousand ways we are privileged to choose for ourselves. Here we learn from the hard taskmaster of experience. We discern between good and evil. We differentiate as to the bitter and the sweet. We discover that there are consequences attached to our actions.

By obedience to God’s commandments, we can qualify for that “house” spoken of by Jesus when He declared: “In my Father’s house are many mansions. … I go to prepare a place for you … that where I am, there ye may be also.”

Although we come into mortality “trailing clouds of glory,” life moves relentlessly forward. Youth follows childhood, and maturity comes ever so imperceptibly. From experience we learn the need to reach heavenward for assistance as we make our way along life’s pathway.

God, our Father, and Jesus Christ, our Lord, have marked the way to perfection. They beckon us to follow eternal verities and to become perfect, as They are perfect.

The Apostle Paul likened life to a race. To the Hebrews he urged, “Let us lay aside … the sin which doth so easily beset us, and let us run with patience the race that is set before us.”

In our zeal, let us not overlook the sage counsel from Ecclesiastes: “The race is not to the swift, nor the battle to the strong.” Actually, the prize belongs to him or her who endures to the end.

When I reflect on the race of life, I remember another type of race, even from childhood days. My friends and I would take pocketknives in hand and, from the soft wood of a willow tree, fashion small toy boats. With a triangular-shaped cotton sail in place, each would launch his crude craft in the race down the relatively turbulent waters of Utah’s Provo River. We would run along the river’s bank and watch the tiny vessels sometimes bobbing violently in the swift current and at other times sailing serenely as the water deepened.

During a particular race we noted that one boat led all the rest toward the appointed finish line. Suddenly, the current carried it too close to a large whirlpool, and the boat heaved to its side and capsized. Around and around it was carried, unable to make its way back into the main current. At last it came to an uneasy rest amid the flotsam and jetsam that surrounded it, held fast by the tentacles of the grasping green moss.

The toy boats of childhood had no keel for stability, no rudder to provide direction, and no source of power. Inevitably, their destination was downstream—the path of least resistance.

Unlike toy boats, we have been provided divine attributes to guide our journey. We enter mortality not to float with the moving currents of life but with the power to think, to reason, and to achieve.

Our Heavenly Father did not launch us on our eternal voyage without providing the means whereby we could receive from Him guidance to ensure our safe return. I speak of prayer. I speak too of the whisperings from that still, small voice; and I do not overlook the holy scriptures, which contain the word of the Lord and the words of the prophets—provided to us to help us successfully cross the finish line.

At some period in our mortal mission, there appears the faltering step, the wan smile, the pain of sickness—even the fading of summer, the approach of autumn, the chill of winter, and the experience we call death.

Every thoughtful person has asked himself the question best phrased by Job of old: “If a man die, shall he live again?” Try as we might to put the question out of our thoughts, it always returns. Death comes to all mankind. It comes to the aged as they walk on faltering feet. Its summons is heard by those who have scarcely reached midway in life’s journey. At times it hushes the laughter of little children.

But what of an existence beyond death? Is death the end of all? Robert Blatchford, in his book God and My Neighbor, attacked with vigor accepted Christian beliefs such as God, Christ, prayer, and particularly immortality. He boldly asserted that death was the end of our existence and that no one could prove otherwise. Then a surprising thing happened. His wall of skepticism suddenly crumbled to dust. He was left exposed and undefended. Slowly he began to feel his way back to the faith he had ridiculed and abandoned. What had caused this profound change in his outlook? His wife died. With a broken heart he went into the room where lay all that was mortal of her. He looked again at the face he loved so well. Coming out, he said to a friend: “It is she, and yet it is not she. Everything is changed. Something that was there before is taken away. She is not the same. What can be gone if it be not the soul?”

Later he wrote: “Death is not what some people imagine. It is only like going into another room. In that other room we shall find … the dear women and men and the sweet children we have loved and lost.”

My brothers and sisters, we know that death is not the end. This truth has been taught by living prophets throughout the ages. It is also found in our holy scriptures. In the Book of Mormon we read specific and comforting words:

“Now, concerning the state of the soul between death and the resurrection—Behold, it has been made known unto me by an angel, that the spirits of all men, as soon as they are departed from this mortal body, yea, the spirits of all men, whether they be good or evil, are taken home to that God who gave them life.

“And then shall it come to pass, that the spirits of those who are righteous are received into a state of happiness, which is called paradise, a state of rest, a state of peace, where they shall rest from all their troubles and from all care, and sorrow.”

After the Savior was crucified and His body had lain in the tomb for three days, the spirit again entered. The stone was rolled away, and the resurrected Redeemer walked forth, clothed with an immortal body of flesh and bones.

The answer to Job’s question, “If a man die, shall he live again?” came when Mary and others approached the tomb and saw two men in shining garments who spoke to them: “Why seek ye the living among the dead? He is not here, but is risen.”

As the result of Christ’s victory over the grave, we shall all be resurrected. This is the redemption of the soul. Paul wrote: “There are … celestial bodies, and bodies terrestrial: but the glory of the celestial is one, and the glory of the terrestrial is another.”

It is the celestial glory which we seek. It is in the presence of God we desire to dwell. It is a forever family in which we want membership. Such blessings are earned through a lifetime of striving, seeking, repenting, and finally succeeding.

Where did we come from? Why are we here? Where do we go after this life? No longer need these universal questions remain unanswered. From the very depths of my soul and in all humility, I testify that those things of which I have spoken are true.

Our Heavenly Father rejoices for those who keep His commandments. He is concerned also for the lost child, the tardy teenager, the wayward youth, the delinquent parent. Tenderly the Master speaks to these and indeed to all: “Come back. Come up. Come in. Come home. Come unto me.”

In one week we will celebrate Easter. Our thoughts will turn to the Savior’s life, His death, and His Resurrection. As His special witness, I testify to you that He lives and that He awaits our triumphant return. That such a return will be ours, I pray humbly in His holy name—even Jesus Christ, our Savior and our Redeemer, amen.

\newpage
\settalk{Willing and Worthy to Serve}
\setspeaker{Thomas S. Monson}
\settalkdate{April 2012}
\talkheading{Willing and Worthy to Serve}{Thomas S. Monson}

My beloved brethren, how good it is to meet with you once again. Whenever I attend the general priesthood meeting, I reflect on the teachings of some of God’s noble leaders who have spoken in the general priesthood meetings of the Church. Many have passed to their eternal reward, and yet from the brilliance of their minds, from the depths of their souls, and from the warmth of their hearts, they have given us inspired direction. I share with you tonight some of their teachings concerning the priesthood.

From the Prophet Joseph Smith: “Priesthood is an everlasting principle, and existed with God from eternity, and will to eternity, without beginning of days or end of years.”

From the words of President Wilford Woodruff, we learn: “The Holy Priesthood is the channel through which God communicates and deals with man upon the earth; and the heavenly messengers that have visited the earth to communicate with man are men who held and honored the priesthood while in the flesh; and everything that God has caused to be done for the salvation of man, from the coming of man upon the earth to the redemption of the world, has been and will be by virtue of the everlasting priesthood.”

President Joseph F. Smith further clarified: “The Priesthood … is … the power of God delegated to man by which man can act in the earth for the salvation of the human family, in the name of the Father and the Son and the Holy Ghost, and act legitimately; not assuming that authority, nor borrowing it from generations that are dead and gone, but authority that has been given in this day in which we live by ministering angels and spirits from above, direct from the presence of Almighty God.”

And finally from President John Taylor: “What is priesthood? … It is the government of God, whether on the earth or in the heavens, for it is by that power, agency, or principle that all things are governed on the earth and in the heavens, and by that power that all things are upheld and sustained. It governs all things—it directs all things—it sustains all things—and has to do with all things that God and truth are associated with.”

How blessed we are to be here in these last days, when the priesthood of God is upon the earth. How privileged we are to bear that priesthood. The priesthood is not so much a gift as it is a commission to serve, a privilege to lift, and an opportunity to bless the lives of others.

With these opportunities come responsibilities and duties. I love and cherish the noble word duty and all that it implies.

In one capacity or another, in one setting or another, I have been attending priesthood meetings for the past 72 years—since I was ordained a deacon at the age of 12. Time certainly marches on. Duty keeps cadence with that march. Duty does not dim nor diminish. Catastrophic conflicts come and go, but the war waged for the souls of men continues without abatement. Like a clarion call comes the word of the Lord to you, to me, and to priesthood holders everywhere: “Wherefore, now let every man learn his duty, and to act in the office in which he is appointed, in all diligence.”

The call of duty came to Adam, to Noah, to Abraham, to Moses, to Samuel, to David. It came to the Prophet Joseph Smith and to each of his successors. The call of duty came to the boy Nephi when he was instructed by the Lord, through his father Lehi, to return to Jerusalem with his brothers to obtain the brass plates from Laban. Nephi’s brothers murmured, saying it was a hard thing which had been asked of them. What was Nephi’s response? Said he, “I will go and do the things which the Lord hath commanded, for I know that the Lord giveth no commandments unto the children of men, save he shall prepare a way for them that they may accomplish the thing which he commandeth them.”

When that same call comes to you and to me, what will be our response? Will we murmur, as did Laman and Lemuel, and say, “This is a hard thing required of us”? Or will we, with Nephi, individually declare, “I will go. I will do”? Will we be willing to serve and to obey?

At times the wisdom of God appears as being foolish or just too difficult, but one of the greatest and most valuable lessons we can learn in mortality is that when God speaks and a man obeys, that man will always be right.

When I think of the word duty and how performing our duty can enrich our lives and the lives of others, I recall the words penned by a renowned poet and author:

\begin{verse}
I slept and dreamt\\
That life was joy\\
I awoke and saw\\
That life was duty\\
I acted and behold\\
Duty was joy.
\end{verse}

Robert Louis Stevenson put it another way. Said he, “I know what pleasure is, for I have done good work.”

As we perform our duties and exercise our priesthood, we will find true joy. We will experience the satisfaction of having completed our tasks.

We have been taught the specific duties of the priesthood which we hold, whether it be the Aaronic or the Melchizedek Priesthood. I urge you to contemplate those duties and then do all within your power to fulfill them. In order to do so, each must be worthy. Let us have ready hands, clean hands, and willing hands, that we may participate in providing what our Heavenly Father would have others receive from Him. If we are not worthy, it is possible to lose the power of the priesthood; and if we lose it, we have lost the essence of exaltation. Let us be worthy to serve.

President Harold B. Lee, one of the great teachers in the Church, said: “When one becomes a holder of the priesthood, he becomes an agent of the Lord. He should think of his calling as though he were on the Lord’s errand.”

During World War II, in the early part of 1944, an experience involving the priesthood took place as United States marines were taking Kwajalein Atoll, part of the Marshall Islands and located in the Pacific Ocean about midway between Australia and Hawaii. What took place in this regard was related by a correspondent—not a member of the Church—who worked for a newspaper in Hawaii. In the 1944 newspaper article he wrote following the experience, he explained that he and other correspondents were in the second wave behind the marines at Kwajalein Atoll. As they advanced, they noticed a young marine floating facedown in the water, obviously badly wounded. The shallow water around him was red with his blood. And then they noticed another marine moving toward his wounded comrade. The second marine was also wounded, with his left arm hanging helplessly by his side. He lifted up the head of the one who was floating in the water in order to keep him from drowning. In a panicky voice he called for help. The correspondents looked again at the boy he was supporting and called back, “Son, there is nothing we can do for this boy.”

“Then,” wrote the correspondent, “I saw something that I had never seen before.” This boy, badly wounded himself, made his way to the shore with the seemingly lifeless body of his fellow marine. He “put the head of his companion on his knee. … What a picture that was—these two mortally wounded boys—both … clean, wonderful-looking young men, even in their distressing situation. And the one boy bowed his head over the other and said, ‘I command you, in the name of Jesus Christ and by the power of the priesthood, to remain alive until I can get medical help.’” The correspondent concluded his article: “The three of us [the two marines and I] are here in the hospital. The doctors don’t know [how they made it alive], but I know.”

Miracles are everywhere to be found when the priesthood is understood, its power is honored and used properly, and faith is exerted. When faith replaces doubt, when selfless service eliminates selfish striving, the power of God brings to pass His purposes.

The call of duty can come quietly as we who hold the priesthood respond to the assignments we receive. President George Albert Smith, that modest but effective leader, declared, “It is your duty first of all to learn what the Lord wants and then by the power and strength of His holy Priesthood to [so] magnify your calling in the presence of your fellows … that the people will be glad to follow you.”

Such a call of duty—a much less dramatic call but one which nonetheless helped to save a soul—came to me in 1950 when I was a newly called bishop. My responsibilities as a bishop were many and varied, and I tried to the best of my ability to do all that was required of me. The United States was engaged in a different war by then. Because many of our members were serving in the armed services, an assignment came from Church headquarters for all bishops to provide each serviceman a subscription to the Church News and the Improvement Era, the Church’s magazine at that time. In addition, each bishop was asked to write a personal, monthly letter to each serviceman from his ward. Our ward had 23 men in uniform. The priesthood quorums, with effort, supplied the funds for the subscriptions to the publications. I undertook the task, even the duty, to write 23 personal letters each month. After all these years I still have copies of many of my letters and the responses received. Tears come easily when these letters are reread. It is a joy to learn again of a soldier’s pledge to live the gospel, a sailor’s decision to keep faith with his family.

One evening I handed to a sister in the ward the stack of 23 letters for the current month. Her assignment was to handle the mailing and to maintain the constantly changing address list. She glanced at one envelope and, with a smile, asked, “Bishop, don’t you ever get discouraged? Here is another letter to Brother Bryson. This is the 17th letter you have sent to him without a reply.”

I responded, “Well, maybe this will be the month.” As it turned out, that was the month. For the first time, he responded to my letter. His reply is a keepsake, a treasure. He was serving far away on a distant shore, isolated, homesick, alone. He wrote, “Dear Bishop, I ain’t much at writin’ letters.” (I could have told him that several months earlier.) His letter continued, “Thank you for the Church News and magazines, but most of all thank you for the personal letters. I have turned over a new leaf. I have been ordained a priest in the Aaronic Priesthood. My heart is full. I am a happy man.”

Brother Bryson was no happier than was his bishop. I had learned the practical application of the adage “Do [your] duty; that is best; leave unto [the] Lord the rest.”

Years later, while attending the Salt Lake Cottonwood Stake when James E. Faust served as its president, I related that account in an effort to encourage attention to our servicemen. After the meeting, a fine-looking young man came forward. He took my hand in his and asked, “Bishop Monson, do you remember me?”

I suddenly realized who he was. “Brother Bryson!” I exclaimed. “How are you? What are you doing in the Church?”

With warmth and obvious pride, he responded, “I’m fine. I serve in the presidency of my elders quorum. Thank you again for your concern for me and the personal letters which you sent and which I treasure.”

Brethren, the world is in need of our help. Are we doing all we should? Do we remember the words of President John Taylor: “If you do not magnify your callings, God will hold you responsible for those whom you might have saved had you done your duty”? There are feet to steady, hands to grasp, minds to encourage, hearts to inspire, and souls to save. The blessings of eternity await you. Yours is the privilege to be not spectators but participants on the stage of priesthood service. Let us hearken to the stirring reminder found in the Epistle of James: “Be ye doers of the word, and not hearers only, deceiving your own selves.”

Let us learn and contemplate our duty. Let us be willing and worthy to serve. Let us in the performance of our duty follow in the footsteps of the Master. As you and I walk the pathway Jesus walked, we will discover He is more than the babe in Bethlehem, more than the carpenter’s son, more than the greatest teacher ever to live. We will come to know Him as the Son of God, our Savior and our Redeemer. When to Him came the call of duty, He answered, “Father, thy will be done, and the glory be thine forever.” May each of us do likewise, I pray in His holy name, the name of Jesus Christ, the Lord, amen.

\newpage
\settalk{Consider the Blessings}
\setspeaker{Thomas S. Monson}
\settalkdate{October 2012}
\talkheading{Consider the Blessings}{Thomas S. Monson}

My beloved brothers and sisters, this conference marks 49 years since I was sustained, on October 4, 1963, as a member of the Quorum of the Twelve Apostles. Forty-nine years is a long time. In many ways, however, the time seems very short since I stood at the pulpit in the Tabernacle and gave my very first general conference address.

Much has changed since October 4, 1963. We live in a unique time in the world’s history. We are blessed with so very much. And yet it is sometimes difficult to view the problems and permissiveness around us and not become discouraged. I have found that, rather than dwelling on the negative, if we will take a step back and consider the blessings in our lives, including seemingly small, sometimes overlooked blessings, we can find greater happiness.

As I have reviewed the past 49 years, I have made some discoveries. One is that countless experiences I have had were not necessarily those one would consider extraordinary. In fact, at the time they transpired, they often seemed unremarkable and even ordinary. And yet, in retrospect, they enriched and blessed lives—not the least of which was my own. I would recommend this same exercise to you—namely, that you take an inventory of your life and look specifically for the blessings, large and small, you have received.

Reinforced constantly during my own review of the years has been my knowledge that our prayers are heard and answered. We are familiar with the truth found in 2 Nephi in the Book of Mormon: “Men are, that they might have joy.” I testify that much of that joy comes as we recognize that we can communicate with our Heavenly Father through prayer and that those prayers will be heard and answered—perhaps not how and when we expected they would be answered, but they will be answered and by a Heavenly Father who knows and loves us perfectly and who desires our happiness. Hasn’t He promised us, “Be thou humble; and the Lord thy God shall lead thee by the hand, and give thee answer to thy prayers”?

For the next few minutes allotted to me, I would like to share with you just a tiny sampling of the experiences I have had wherein prayers were heard and answered and which, in retrospect, brought blessings into my life as well as the lives of others. My daily journal, kept over all these years, has helped provide some specifics which I most likely would not otherwise be able to recount.

In early 1965, I was assigned to attend stake conferences and to hold other meetings throughout the South Pacific area. This was my first visit to that part of the world, and it was a time never to be forgotten. Much that was spiritual in nature occurred during this assignment as I met with leaders, members, and missionaries.

On the weekend of Saturday and Sunday, February 20 and 21, we were in Brisbane, Australia, to hold regular conference sessions of the Brisbane Stake. During meetings on Saturday, I was introduced to the district president from an adjoining area. As I shook his hand, I had a strong impression that I needed to speak with him and to provide counsel, and so I asked him if he would accompany me to the Sunday morning session the following day so that this could be accomplished.

Following the Sunday session, we had an opportunity to visit together. We talked of his many responsibilities as district president. As we did so, I felt impressed to offer him specific suggestions concerning missionary work and how he and his members could help the full-time missionaries in their labors in his area. I later learned that this man had been praying for guidance in this regard. To him our visit was a special witness that his prayers were heard and answered. This was a seemingly unremarkable meeting but one which I am convinced was guided by the Spirit and which made a difference in that district president’s life and administration, in the lives of his members, and in the success of the missionaries there.

My brothers and sisters, the Lord’s purposes are often accomplished as we pay heed to the guidance of the Spirit. I believe that the more we act upon the inspiration and impressions which come to us, the more the Lord will entrust to us His errands.

I have learned, as I have mentioned in previous messages, never to postpone a prompting. On one occasion many years ago, I was swimming laps at the old Deseret Gym in Salt Lake City when I felt the inspiration to go to the University Hospital to visit a good friend of mine who had lost the use of his lower limbs because of a malignancy and the surgery which followed. I immediately left the pool, dressed, and was soon on my way to see this good man.

When I arrived at his room, I found that it was empty. Upon inquiry I learned I would probably find him in the swimming pool area of the hospital, an area which was used for physical therapy. Such turned out to be the case. He had guided himself there in his wheelchair and was the only occupant of the room. He was on the far side of the pool, near the deep end. I called to him, and he maneuvered his wheelchair over to greet me. We had an enjoyable visit, and I accompanied him back to his hospital room, where I gave him a blessing.

I learned later from my friend that he had been utterly despondent that day and had been contemplating taking his own life. He had prayed for relief but began to feel that his prayers had gone unanswered. He went to the pool with the thought that this would be a way to end his misery—by guiding his wheelchair into the deep end of the pool. I had arrived at a critical moment, in response to what I know was inspiration from on high.

My friend was able to live many more years—years filled with happiness and gratitude. How pleased I am to have been an instrument in the Lord’s hands on that critical day at the swimming pool.

On another occasion, as Sister Monson and I were driving home after visiting friends, I felt impressed that we should go into town—a drive of many miles—to pay a visit to an elderly widow who had once lived in our ward. Her name was Zella Thomas. At the time, she was a resident in a care center. That early afternoon we found her to be extremely frail but lying peacefully on her bed.

Zella had long been blind, but she recognized our voices immediately. She asked if I might give her a blessing, adding that she was prepared to die if the Lord wanted her to return home. There was a sweet, peaceful spirit in the room, and all of us knew that her remaining time in mortality would be brief. Zella took me by the hand and said that she had prayed fervently that I would come to see her and provide her a blessing. I told her that we had come because of direct inspiration from our Heavenly Father. I kissed her on the forehead, knowing that I perhaps would not again see her in mortality. Such proved to be the case, for she passed away the following day. To have been able to provide some comfort and peace to our sweet Zella was a blessing to her and to me.

The opportunity to be a blessing in the life of another often comes unexpectedly. On one extremely cold Saturday night during the winter of 1983–84, Sister Monson and I drove several miles to the mountain valley of Midway, Utah, where we have a home. The temperature that night was minus 24 degrees Fahrenheit (–31°C), and we wanted to make certain all was well at our home there. We checked and found that it was fine, so we left to return to Salt Lake City. We barely made it the few miles to the highway before our car stopped working. We were completely stranded. I have seldom, if ever, been as cold as we were that night.

Reluctantly we began walking toward the nearest town, the cars whizzing past us. Finally one car stopped, and a young man offered to help. We eventually found that the diesel fuel in our gas tank had thickened because of the cold, making it impossible for us to drive the car. This kind young man drove us back to our Midway home. I attempted to reimburse him for his services, but he graciously declined. He indicated that he was a Boy Scout and wanted to do a good turn. I identified myself to him, and he expressed his appreciation for the privilege to be of help. Assuming that he was about missionary age, I asked him if he had plans to serve a mission. He indicated he was not certain just what he wanted to do.

On the following Monday morning, I wrote a letter to this young man and thanked him for his kindness. In the letter I encouraged him to serve a full-time mission. I enclosed a copy of one of my books and underscored the chapters on missionary service.

About a week later the young man’s mother telephoned and advised that her son was an outstanding young man but that because of certain influences in his life, his long-held desire to serve a mission had diminished. She indicated she and his father had fasted and prayed that his heart would be changed. They had placed his name on the prayer roll of the Provo Utah Temple. They hoped that somehow, in some way, his heart would be touched for good and he would return to his desire to fill a mission and to serve the Lord faithfully. The mother wanted me to know that she looked upon the events of that cold evening as an answer to their prayers in his behalf. I said, “I agree with you.”

After several months and more communication with this young man, Sister Monson and I were overjoyed to attend his missionary farewell prior to his departure for the Canada Vancouver Mission.

Was it chance that our paths crossed on that cold December night? I do not for one moment believe so. Rather, I believe our meeting was an answer to a mother’s and father’s heartfelt prayers for the son they cherished.

Again, my brothers and sisters, our Heavenly Father is aware of our needs and will help us as we call upon Him for assistance. I believe that no concern of ours is too small or insignificant. The Lord is in the details of our lives.

I should like to conclude by relating one recent experience which had an impact on hundreds. It occurred at the cultural celebration for the Kansas City Temple, just five months ago. As with so much that happens in our lives, at the time it seemed to be just another experience where everything worked out. However, as I learned of the circumstances associated with the cultural celebration the evening before the temple was dedicated, I realized that the performance that night was not ordinary. Rather, it was quite remarkable.

As with all cultural events held in conjunction with temple dedications, the youth in the Kansas City Missouri Temple District had rehearsed the performance in separate groups in their own areas. The plan was that they would meet all together in the large rented municipal center on the Saturday morning of the performance so that they could learn when and where to enter, where they were to stand, how much space should be between them and the person next to them, how to exit the main floor, and so forth—many details which they would have to grasp during the day as those in charge put the various scenes together so that the final performance would be polished and professional.

There was just one major problem that day. The entire production was dependent on prerecorded segments that would be shown on the large screen known as a Jumbotron. These recorded segments were critical to the entire production. They not only tied it all together, but each televised segment would introduce the next performance. The video segments provided the framework on which the entire production depended. And the Jumbotron was not working.

Technicians worked frantically to solve the problem while the youth waited, hundreds of them, losing precious rehearsal time. The situation began to look impossible.

The writer and director of the celebration, Susan Cooper, later explained: “As we moved from plan A to B to Z, we knew that it wasn’t working. … As we were looking at the schedule, we knew that it was going to be beyond us, but we knew that we had one of the greatest strengths on the floor below—3,000 youth. We needed to go down and tell [them] what was happening and draw upon their faith.”

Just an hour before the audience would begin to enter the center, 3,000 youth knelt on the floor and prayed together. They prayed that those working on the Jumbotron would be inspired to know what to do to repair it; they asked their Heavenly Father to make up for what they themselves could not do because of the shortage of time.

Said one who wrote about it afterward, “It was a prayer the youth will never forget, not because the floor was hard, but because the Spirit melted their bones.”

It was not long before one of the technicians came to tell them that the problem had been discovered and corrected. He attributed the solution to luck, but all those youth knew better.

When we entered the municipal center that evening, we had no idea of the difficulties of the day. Only later did we learn of them. What we witnessed, however, was a beautiful, polished performance—one of the best I have seen. The youth radiated a glorious, powerful spirit which was felt by all who were present. They seemed to know just where to enter, where to stand, and how to interact with all the other performers around them. When I learned that their rehearsals had been cut short and that many of the numbers had not been rehearsed by the entire group, I was astonished. No one would have known. The Lord had indeed made up the difference.

I never cease to be amazed by how the Lord can motivate and direct the length and breadth of His kingdom and yet have time to provide inspiration concerning one individual—or one cultural celebration or one Jumbotron. The fact that He can, that He does, is a testimony to me.

My brothers and sisters, the Lord is in all of our lives. He loves us. He wants to bless us. He wants us to seek His help. As He guides us and directs us and as He hears and answers our prayers, we will find the happiness here and now that He desires for us. May we be aware of His blessings in our lives, I pray in the name of Jesus Christ, our Savior, amen.

\newpage
\settalk{God Be with You Till We Meet Again}
\setspeaker{Thomas S. Monson}
\settalkdate{October 2012}
\talkheading{God Be with You Till We Meet Again}{Thomas S. Monson}

My dear brothers and sisters, we have come to the close of another inspiring general conference. I personally have been spiritually fed and uplifted and know that you too have felt the special spirit of this conference.

We offer our heartfelt gratitude to all who have participated in any way. The truths of the gospel have been beautifully taught and reemphasized. As we take the messages of the past two days into our hearts and into our lives, we will be blessed.

As always, the proceedings of this conference will be available in the coming issues of the Ensign and the Liahona magazines. I encourage you to read the talks once again and to ponder the messages contained therein. I have found in my own life that I gain even more from these inspired sermons when I study them in greater depth.

We have had unprecedented coverage of the conference, reaching across the continents and the oceans to people everywhere. Though we are far removed from many of you, we feel of your spirit and send our love and appreciation to you.

To our Brethren who have been released at this conference, may I express the heartfelt gratitude of all of us for your many years of devoted service. Countless are those who have been blessed by your contributions to the work of the Lord.

Brothers and sisters, I have just recently celebrated my 85th birthday, and I am grateful for each year the Lord has granted me. As I reflect upon my life’s experiences, I thank Him for His many blessings to me. As I mentioned in my message this morning, I have felt His hand directing my efforts as I have tried earnestly to serve Him and to serve all of you.

The office of the President of the Church is a demanding one. How grateful I am for my two faithful counselors, who serve by my side and who are always willing and exceptionally able to assist in the work which comes to the First Presidency. I express my gratitude as well for the noble men who comprise the Quorum of the Twelve Apostles. They work tirelessly in the cause of the Master, with the members of the Quorums of the Seventy providing inspired assistance to them.

I wish also to commend you, my brothers and sisters, wherever you are throughout the world, for all that you do in your wards and branches, your stakes and districts. As you willingly fulfill callings when you are asked, you are helping to build the kingdom of God on earth.

May we ever watch over one another, assisting in times of need. Let us not be critical and judgmental but let us be tolerant, ever emulating the Savior’s example of loving-kindness. In that vein, may we willingly serve one another. May we pray for the inspiration to know of the needs of those around us, and then may we go forward and provide assistance.

Let us be of good cheer as we go about our lives. Although we live in increasingly perilous times, the Lord loves us and is mindful of us. He is always on our side as we do what is right. He will help us in time of need. Difficulties come into our lives, problems we do not anticipate and which we would never choose. None of us is immune. The purpose of mortality is to learn and to grow to be more like our Father, and it is often during the difficult times that we learn the most, as painful as the lessons may be. Our lives can also be filled with joy as we follow the teachings of the gospel of Jesus Christ.

The Lord admonished, “Be of good cheer; I have overcome the world.” What great happiness this knowledge should bring to us. He lived for us and He died for us. He paid the price for our sins. May we emulate His example. May we show our great gratitude to Him by accepting His sacrifice and living lives that will qualify us to return and one day live with Him.

As I have mentioned at previous conferences, I thank you for your prayers in my behalf. I need them; I feel them. We as General Authorities also remember all of you and pray for our Heavenly Father’s choicest blessings to be with you.

Now, my beloved brothers and sisters, we adjourn for six months. May God be with you until we meet again at that time. In the name of our Savior and Redeemer, even Jesus Christ the Lord, amen.

\newpage
\settalk{See Others as They May Become}
\setspeaker{Thomas S. Monson}
\settalkdate{October 2012}
\talkheading{See Others as They May Become}{Thomas S. Monson}

My dear brethren, twice each year this magnificent Conference Center is filled to capacity with the priesthood of God as we gather to hear messages of inspiration. There is a marvelous spirit which permeates the general priesthood meeting of the Church. This spirit emanates from the Conference Center and enters every building where the sons of God assemble. We have surely felt that spirit tonight.

Some years ago, before this beautiful Conference Center was built, a visitor to Temple Square in Salt Lake City attended a general conference session in the Tabernacle. He listened to the messages of the Brethren. He paid attention to the prayers. He heard the beautiful music by the Tabernacle Choir. He marveled at the grandeur of the magnificent Tabernacle organ. When the meeting had ended he was heard to say, “I would give everything I possess if I knew that what those speakers said today was true.” In essence he was saying, “I wish that I had a testimony of the gospel.”

There is absolutely nothing in this world that will provide more comfort and happiness than a testimony of the truth. Although to varying degrees, I believe every man or young man here tonight has a testimony. If you feel that you do not yet have the depth of testimony you would wish, I admonish you to work to achieve such a testimony. If it is strong and deep, labor to keep it that way. How blessed we are to have a knowledge of the truth.

My message tonight, brethren, is that there are countless individuals who have little or no testimony right now, those who could and would receive such a testimony if we would be willing to make the effort to share ours and to help them change. In some instances we can provide the incentive for change. I mention first those who are members but who are not at present fully committed to the gospel.

Many years ago, at an area conference held in Helsinki, Finland, I heard a powerful, memorable, and motivating message given in a mothers and daughters’ session. I have not forgotten that message, though nearly 40 years have passed since I heard it. Among many truths the speaker discussed, she said that a woman needs to be told she is beautiful. She needs to be told she is valued. She needs to be told she is worthwhile.

Brethren, I know that men are very much like women in this regard. We need to be told that we amount to something, that we are capable and worthwhile. We need to be given a chance to serve. For those members who have slipped from activity or who hold back and remain noncommittal, we can prayerfully seek for some way to reach them. Asking them to serve in some capacity may just be the incentive they need to return to full activity. But those leaders who could help in this regard are sometimes reluctant to do so. We need to bear in mind that people can change. They can put behind them bad habits. They can repent from transgressions. They can bear the priesthood worthily. And they can serve the Lord diligently. May I provide a few illustrations.

When I first became a member of the Quorum of the Twelve Apostles, I had the opportunity to accompany President N. Eldon Tanner, a counselor to President David O. McKay, to a stake conference in Alberta, Canada. During the meeting, the stake president read the names of four brethren who had qualified to be ordained elders. These were men whom President Tanner knew, for at one time he had lived in that area. But President Tanner knew and remembered them as they once were and did not know that they had turned their lives around and had fully qualified to become elders.

The stake president read the name of the first man and asked him to stand. President Tanner whispered to me, “Look at him. I never thought he would make it.” The stake president read the name of the second man, and he stood. President Tanner nudged me again and reported his astonishment. And so it was with all four of the brethren.

After the meeting, President Tanner and I had the opportunity to congratulate these four brethren. They had demonstrated that men can change.

During the 1940s and 1950s, an American prison warden, Clinton Duffy, was well known for his efforts to rehabilitate the men in his prison. Said one critic, “You should know that leopards don’t change their spots!”

Replied Warden Duffy, “You should know I don’t work with leopards. I work with men, and men change every day.”

Many years ago it was my opportunity to serve as president of the Canadian Mission. There we had a branch with very limited priesthood. We always had a missionary presiding over the branch. I received a strong impression that we needed to have a member of the branch preside there.

We had one adult member in the branch who was a deacon in the Aaronic Priesthood but who didn’t attend or participate enough to be advanced in the priesthood. I felt inspired to call him as the branch president. I shall always remember the day that I had an interview with him. I told him that the Lord had inspired me to call him to be the president of the branch. After much protest on his part, and much encouragement on the part of his wife, he indicated that he would serve. I ordained him a priest.

It was the beginning of a new day for that man. His life was quickly put in order, and he assured me that he would live the commandments as he was expected to live them. In a few months he was ordained an elder. He and his wife and family eventually went to the temple and were sealed. Their children served missions and married in the house of the Lord.

Sometimes letting our brethren know they are needed and valued can help them take that step into commitment and full activity. This can be true of priesthood holders regardless of age. It is our responsibility to give them opportunities to live as they should. We can help them to overcome their shortcomings. We must develop the capacity to see men not as they are at present but as they may become when they receive testimonies of the gospel of Christ.

I once attended a meeting in Leadville, Colorado. Leadville is situated at an altitude of over 10,000 feet (3,000 m). I remember that particular meeting because of the high altitude, but I also remember it for what took place that evening. There were just a small number of priesthood holders present. As with the branch in the Canadian Mission, that branch was presided over by a missionary and always had been.

That night we had a lovely meeting, but as we were singing the closing song, the inspiration came to me that there ought to be a local branch president presiding. I turned to the mission president and asked, “Isn’t there someone here who could preside—a local man?”

He replied, “I don’t know of one.”

During the singing of that song, I looked carefully at the men who were seated on the first three rows. My attention seemed to be focused on one of the brethren. I said to the mission president, “Could he serve as the branch president?”

He replied, “I don’t know. Perhaps he could.”

I said, “President, I’ll take him into the other room and interview him. You speak after the closing song until we return.”

When the two of us walked back in the room, the mission president concluded his testimony. I presented the name of the brother to be the new branch president. From that day forward, Leadville, Colorado, had a local member leading the unit there.

The same principle, brethren, applies to those who are not yet members. We should develop the capacity to see men not as they are but as they can become when they are members of the Church, when they have a testimony of the gospel, and when their lives are in harmony with its teachings.

Back in the year 1961, a worldwide conference was held for mission presidents, and every mission president in the Church was brought to Salt Lake City for those meetings. I came to Salt Lake City from my mission in Toronto, Canada.

In one particular meeting, N. Eldon Tanner, who was then an Assistant to the Quorum of the Twelve, had just returned from his initial experience of presiding over the missions in Great Britain and western Europe. He told of a missionary who had been the most successful missionary whom he had met in all of the interviews he had conducted. He said that as he interviewed that missionary, he said to him, “I suppose that all of the people whom you baptized came into the Church by way of referrals.”

The young man answered, “No, we found them all by tracting.”

Brother Tanner asked him what was different about his approach—why he had such phenomenal success when others didn’t. The young man said that he attempted to baptize every person whom he met. He said that if he knocked on the door and saw a man smoking a cigar and dressed in old clothes and seemingly uninterested in anything—particularly religion—the missionary would picture in his own mind what that man would look like under a different set of circumstances. In his mind he would look at him as clean-shaven and wearing a white shirt and white trousers. And the missionary could see himself leading that man into the waters of baptism. He said, “When I look at someone that way, I have the capacity to bear my testimony to him in a way that can touch his heart.”

We have the responsibility to look at our friends, our associates, our neighbors this way. Again, we have the responsibility to see individuals not as they are but rather as they can become. I would plead with you to think of them in this way.

Brethren, the Lord told us something about the importance of this priesthood that we hold. He told us that we receive it with an oath and a covenant. He gave unto us the instruction that we must be faithful and true in all that we receive, and that we have the responsibility to keep this covenant even unto the end. And then all that the Father has shall be given unto us.

Courage is the word we need to hear and hold near our hearts—courage to turn our backs on temptation, courage to lift up our voices in testimony to all whom we meet, remembering that everyone must have an opportunity to hear the message. It is not an easy thing for most to do this. But we can come to believe in the words of Paul to Timothy:

“For God hath not given us the spirit of fear; but of power, and of love, and of a sound mind.

“Be not thou therefore ashamed of the testimony of our Lord.”

In May of 1974, I was with Brother John H. Groberg in the Tongan islands. We had an appointment to visit the king of Tonga, and we met with him in a formal session. We exchanged the normal pleasantries. However, before we left, John Groberg said something that was out of the ordinary. He said, “Your Majesty, you should really become a Mormon and your subjects as well, for then your problems and their problems would largely be solved.”

The king smiled broadly and answered, “John Groberg, perhaps you’re right.”

I thought of the Apostle Paul before Agrippa. I thought of Agrippa’s response to Paul’s testimony: “Almost thou persuadest me to be a Christian.” Brother Groberg had the courage to bear his testimony to a king.

Tonight there are many thousands of our number who are serving the Lord full-time as His missionaries. In response to a call, they have left behind home, family, friends, and school and have gone forward to serve. Those who don’t understand ask the question, “Why do they respond so readily and willingly give so much?”

Our missionaries could well answer in the words of Paul, that peerless missionary of an earlier day: “For though I preach the gospel, I have nothing to glory of: for necessity is laid upon me; yea, woe is unto me, if I preach not the gospel!”

The holy scriptures contain no proclamation more relevant, no responsibility more binding, no instruction more direct than the injunction given by the resurrected Lord as He appeared in Galilee to the eleven disciples. Said He:

“All power is given unto me in heaven and in earth.

“Go ye therefore, and teach all nations, baptizing them in the name of the Father, and of the Son, and of the Holy Ghost:

“Teaching them to observe all things whatsoever I have commanded you: and, lo, I am with you alway, even unto the end of the world.”

This divine command, coupled with its glorious promise, is our watchword today, as it was in the meridian of time. Missionary work is an identifying feature of The Church of Jesus Christ of Latter-day Saints. Always has it been; ever shall it be. As the Prophet Joseph Smith declared, “After all that has been said, the greatest and most important duty is to preach the Gospel.”

Within two short years, all of the full-time missionaries currently serving in this royal army of God will have concluded their full-time labors and will have returned to their homes and loved ones. Their replacements are found tonight in the ranks of the Aaronic Priesthood of the Church. Young men, are you ready to respond? Are you willing to work? Are you prepared to serve?

President John Taylor summed up the requirements: “The kind of men we want as bearers of this gospel message are men who have faith in God; men who have faith in their religion; men who honor their priesthood; … men full of the Holy Ghost and the power of God[;] … men of honor, integrity, virtue and purity.”

Brethren, to each of us comes the mandate to share the gospel of Christ. When our lives comply with God’s own standard, those within our sphere of influence will never speak the lament, “The harvest is past, the summer is ended, and we are not saved.”

The perfect Shepherd of souls, the missionary who redeemed mankind, gave us His divine assurance:

“If it so be that you should labor all your days in crying repentance unto this people, and bring, save it be one soul unto me, how great shall be your joy with him in the kingdom of my Father!

“And now, if your joy will be great with one soul that you have brought unto me into the kingdom of my Father, how great will be your joy if you should bring many souls unto me!”

Of Him who spoke these words, I declare my personal witness. He is the Son of God, our Redeemer, and our Savior.

I pray that we will have the courage to extend the hand of fellowship, the tenacity to try and try again, and the humility needed to seek guidance from our Father as we fulfill our mandate to share the gospel. The responsibility is upon us, brethren. In the name of Jesus Christ, amen.

\newpage
\settalk{Welcome to Conference}
\setspeaker{Thomas S. Monson}
\settalkdate{October 2012}
\talkheading{Welcome to Conference}{Thomas S. Monson}

As far as I can see, every seat is filled—except for a few right there in the back. There is room for improvement. This is a courtesy to those who might be just a bit tardy, because of the traffic, to find a seat when they come.

This is a great day—conference day. We have heard a beautiful choir sing magnificent music. Every time I hear the choir or hear the organ or hear the piano, I think of my mother, who said, “I love all the acclaim that has been given you, all the degrees you have obtained, and all the work you have done. My only regret is that you did not stay with the piano.” Thanks, Mother. I wish I had.

How good it is, my brothers and sisters, to welcome you to the 182nd Semiannual General Conference of The Church of Jesus Christ of Latter-day Saints.

Since we met six months ago, three new temples have been dedicated, and one temple has been rededicated. In May, it was my privilege to dedicate the beautiful Kansas City Missouri Temple and to attend the cultural celebration associated with it. I will mention that celebration in greater detail in my remarks tomorrow morning.

In June, President Dieter F. Uchtdorf dedicated the long-awaited temple in Manaus, Brazil, and in early September, President Henry B. Eyring rededicated the newly refurbished temple in Buenos Aires, Argentina, a temple which I had the privilege to dedicate nearly 27 years ago. Just two weeks ago, President Boyd K. Packer dedicated the lovely Brigham City Temple in the hometown where he was born and raised.

As I have indicated previously, no Church-built facility is more important than a temple, and we are pleased to have 139 temples in operation throughout the world, with 27 more announced or under construction. We are grateful for these sacred edifices and the blessings they bring into our lives.

This morning I am pleased to announce two additional temples, which in coming months and years will be built in the following locations: Tucson, Arizona, and Arequipa, Peru. Details concerning these temples will be provided in the future as necessary permits and approvals are obtained.

Brothers and sisters, I now turn to another matter—namely, missionary service.

For some time the First Presidency and the Quorum of the Twelve Apostles have allowed young men from certain countries to serve at the age of 18 when they are worthy, able, have graduated from high school, and have expressed a sincere desire to serve. This has been a country-specific policy and has allowed thousands of young men to serve honorable missions and also fulfill required military obligations and educational opportunities.

Our experience with these 18-year-old missionaries has been positive. Their mission presidents report that they are obedient, faithful, mature, and serve just as competently as do the older missionaries who serve in the same missions. Their faithfulness, obedience, and maturity have caused us to desire the same option of earlier missionary service for all young men, regardless of the country from which they come.

I am pleased to announce that effective immediately all worthy and able young men who have graduated from high school or its equivalent, regardless of where they live, will have the option of being recommended for missionary service beginning at the age of 18, instead of age 19. I am not suggesting that all young men will—or should—serve at this earlier age. Rather, based on individual circumstances as well as upon a determination by priesthood leaders, this option is now available.

As we have prayerfully pondered the age at which young men may begin their missionary service, we have also given consideration to the age at which a young woman might serve. Today I am pleased to announce that able, worthy young women who have the desire to serve may be recommended for missionary service beginning at age 19, instead of age 21.

We affirm that missionary work is a priesthood duty—and we encourage all young men who are worthy and who are physically able and mentally capable to respond to the call to serve. Many young women also serve, but they are not under the same mandate to serve as are the young men. We assure the young sisters of the Church, however, that they make a valuable contribution as missionaries, and we welcome their service.

We continue to need many more senior couples. As your circumstances allow, as you are eligible for retirement, and as your health permits, I encourage you to make yourselves available for full-time missionary service. Both husband and wife will have a greater joy as they together serve our Father’s children.

Now, my brothers and sisters, may we listen attentively to the messages which will be presented during the next two days, that we may feel the Spirit of the Lord and gain the knowledge He would desire for us. That this may be our experience I pray in the name of Jesus Christ, amen.

\newpage
\settalk{Come, All Ye Sons of God}
\setspeaker{Thomas S. Monson}
\settalkdate{April 2013}
\talkheading{Come, All Ye Sons of God}{Thomas S. Monson}

Twice each year this magnificent Conference Center seems to say to us, with its persuasive voice, “Come, all ye sons of God who have received the priesthood.” There is a characteristic spirit which pervades the general priesthood meeting of the Church.

Tonight there are many thousands of our number throughout the world who are serving the Lord as His missionaries. As I mentioned in my message this morning, we currently have over 65,000 missionaries in the field, with thousands more who are waiting to enter the missionary training center or whose applications are currently being processed. We love and commend those who are willing and anxious to serve.

The holy scriptures contain no proclamation more relevant, no responsibility more binding, no instruction more direct than the injunction given by the resurrected Lord as He appeared in Galilee to the eleven disciples. Said He:

“Go ye therefore, and teach all nations, baptizing them in the name of the Father, and of the Son, and of the Holy Ghost:

“Teaching them to observe all things whatsoever I have commanded you: and, lo, I am with you alway, even unto the end of the world.”

This divine command, coupled with its glorious promise, is our watchword today as it was in the meridian of time. Missionary work is an identifying feature of The Church of Jesus Christ of Latter-day Saints. Always has it been; ever shall it be. As the Prophet Joseph Smith declared, “After all that has been said, the greatest and most important duty is to preach the Gospel.”

Within two short years, all of the full-time missionaries currently serving in this royal army of God will have concluded their labors and will have returned to their homes and loved ones. For the elders, their replacements are found tonight in the ranks of the Aaronic Priesthood of the Church. Young men, are you ready to respond? Are you willing to work? Are you prepared to serve?

At best, missionary work necessitates drastic adjustment to one’s pattern of living. It requires long hours and great devotion, selfless sacrifice and fervent prayer. As a result, dedicated missionary service returns a dividend of eternal joy which extends throughout mortality and into eternity.

The challenge is to be more profitable servants in the Lord’s vineyard. This applies to all of us, whatever our age, and not alone to those who are preparing to serve as full-time missionaries, for to each of us comes the mandate to share the gospel of Christ.

May I suggest a formula that will ensure our success: first, search the scriptures with diligence; second, plan your life with purpose (and, I might add, plan your life regardless of your age); third, teach the truth with testimony; and fourth, serve the Lord with love.

Let us consider each of the four parts of the formula.

First, search the scriptures with diligence.

The scriptures testify of God and contain the words of eternal life. They become the foundation of our message.

The emphasis of the Church curricula is the holy scriptures, programmed and coordinated through the correlation effort. We are encouraged, as well, to study the scriptures each day both individually and with our families.

Let me provide but one reference which has immediate application to our lives. In the Book of Mormon, the 17th chapter of Alma, we read the account of Alma’s joy as he once more saw the sons of Mosiah and noted their steadfastness in the cause of truth. The record tells us, “They had waxed strong in the knowledge of the truth; for they were men of a sound understanding and they had searched the scriptures diligently, that they might know the word of God.

“But this is not all; they had given themselves to much prayer, and fasting; therefore they had the spirit of prophecy, and the spirit of revelation, and when they taught, they taught with power and authority of God.”

Brethren, search the scriptures with diligence.

Second in our formula: plan your life with purpose.

Perhaps no generation of youth has faced such far-reaching decisions as the youth of today. Provision must be made for school, mission, and marriage. For some, military service will be included.

Preparation for a mission begins early. In addition to spiritual preparation, a wise parent will provide the means whereby a young son might commence his personal missionary fund. He may well be encouraged as the years go by to study a foreign language so that, if necessary, his language skills could be utilized. Eventually there comes that glorious day when the bishop and stake president invite the young man in for a visit. Worthiness is ascertained; a missionary recommendation is completed.

During no other time does the entire family so anxiously watch and wait for the mailman and the letter which contains the return address 47 East South Temple, Salt Lake City, Utah. The letter arrives; the suspense is overwhelming; the call is read. Often the assigned field of labor is far away from home. Regardless of the location, however, the response of the prepared and obedient missionary is the same: “I will serve.”

Preparations for departure begin. Young men, I hope you appreciate the sacrifices which your parents so willingly make in order for you to serve. Their labors will sustain you, their faith encourage you, their prayers uphold you. A mission is a family affair. Though the expanse of continents or oceans may separate, hearts are as one.

Brethren, as you plan with purpose your lives, remember that your missionary opportunities are not restricted to the period of a formal call. For those of you who serve in the military, such time can and should be profitable. Each year our young men in uniform bring many souls into the kingdom of God by honoring their priesthood, living the commandments of God, and teaching to others the Lord’s divine word.

Do not overlook your privilege to be missionaries while you are pursuing your formal education. Your example as a Latter-day Saint will be observed, weighed, and ofttimes emulated.

Brethren, whatever your age, whatever your circumstance, I admonish you to plan your life with purpose.

Now to the third point in our formula: teach the truth with testimony.

Obey the counsel of the Apostle Peter, who urged, “Be ready always to give an answer to every man that asketh you a reason of the hope that is in you.” Lift up your voices and testify to the true nature of the Godhead. Declare your witness concerning the Book of Mormon. Convey the glorious and beautiful truths contained in the plan of salvation.

When I served as a mission president in Canada more than 50 years ago, one young missionary who came from a small, rural community marveled at the size of Toronto. He was short in stature but tall in testimony. Not long after his arrival, together with his companion, he called at the home of Elmer Pollard in Oshawa, Ontario, Canada. Feeling sorry for the young men who, during a blinding blizzard, were going house to house, Mr. Pollard invited the missionaries into his home. They presented to him their message. He did not catch the spirit. In due time he asked that they leave and not return. His last words to the elders as they departed his front porch were spoken in derision: “You can’t tell me you actually believe Joseph Smith was a prophet of God!”

The door was shut. The elders walked down the path. Our country boy spoke to his companion: “Elder, we didn’t respond to Mr. Pollard. He said we didn’t believe Joseph Smith was a true prophet. Let’s return and bear our testimonies to him.” At first the more experienced missionary hesitated but finally agreed to accompany his companion. Fear struck their hearts as they approached the door from which they had just been ejected. They knocked, confronted Mr. Pollard, spent an agonizing moment, and then with power borne of the Spirit, our inexperienced missionary spoke: “Mr. Pollard, you said we didn’t really believe Joseph Smith was a prophet of God. I testify to you that Joseph was a prophet. He did translate the Book of Mormon. He saw God the Father and Jesus the Son. I know it.”

Some time later, Mr. Pollard, now Brother Pollard, stood in a priesthood meeting and declared, “That night I could not sleep. Resounding in my ears I heard the words ‘Joseph Smith was a prophet of God. I know it. I know it. I know it.’ The next day I telephoned the missionaries and asked them to return. Their message, coupled with their testimonies, changed my life and the lives of my family.” Brethren, teach the truth with testimony.

The final point in our formula is to serve the Lord with love. There is no substitute for love. Successful missionaries love their companions, their mission leaders, and the precious persons whom they teach. In the fourth section of the Doctrine and Covenants, the Lord established the qualifications for the labors of the ministry. Let us consider but a few verses:

“O ye that embark in the service of God, see that ye serve him with all your heart, might, mind and strength, that ye may stand blameless before God at the last day. …

“And faith, hope, charity and love, with an eye single to the glory of God, qualify him for the work.

“Remember faith, virtue, knowledge, temperance, patience, brotherly kindness, godliness, charity, humility, diligence.”

Well might each of you within the sound of my voice ask himself the question “Today, have I increased in faith, in virtue, in knowledge, in godliness, in love?”

Through your dedicated devotion at home or abroad, those souls whom you help to save may well be those whom you love the most.

Many years ago dear friends of mine, Craig Sudbury and his mother, Pearl, came to my office prior to Craig’s departure for the Australia Melbourne Mission. Fred Sudbury, Craig’s father, was noticeably absent. Twenty-five years earlier, Craig’s mother had married Fred, who did not share her love for the Church and, indeed, was not a member.

Craig confided to me his deep and abiding love for his parents and his hope that somehow, in some way, his father would be touched by the Spirit and open his heart to the gospel of Jesus Christ. I prayed for inspiration concerning how such a desire might be fulfilled. The inspiration came, and I said to Craig, “Serve the Lord with all your heart. Be obedient to your sacred calling. Each week write a letter to your parents, and on occasion, write to Dad personally, and let him know how much you love him, and tell him why you’re grateful to be his son.” He thanked me and, with his mother, departed the office.

I was not to see Craig’s mother for some 18 months, when she came to my office and, in sentences punctuated by tears, said to me, “It has been almost two years since Craig left for his mission. He has never failed in writing a letter to us each week. Recently, my husband, Fred, stood for the first time in a testimony meeting and surprised me and shocked everyone who was there by announcing that he had made the decision to become a member of the Church. He indicated that he and I would go to Australia to meet Craig at the conclusion of his mission so that Fred could be Craig’s final baptism as a full-time missionary.”

No missionary stood so tall as did Craig Sudbury when, in far-off Australia, he helped his father into water waist-deep and, raising his right arm to the square, repeated those sacred words: “Frederick Charles Sudbury, having been commissioned of Jesus Christ, I baptize you in the name of the Father, and of the Son, and of the Holy Ghost.”

Love had won its victory. Serve the Lord with love.

Brethren, may each one of us search the scriptures with diligence, plan his life with purpose, teach the truth with testimony, and serve the Lord with love.

The perfect Shepherd of our souls, the missionary who redeemed mankind, gave us His divine assurance:

“If it so be that you should labor all your days in crying repentance unto this people, and bring, save it be one soul unto me, how great shall be your joy with him in the kingdom of my Father!

“And now, if your joy will be great with one soul that you have brought unto me into the kingdom of my Father, how great will be your joy if you should bring many souls unto me!”

Of Him who spoke these words I declare my witness: He is the Son of God, our Redeemer, and our Savior.

I pray that we may ever respond to His gentle invitation, “Follow thou me.” In His holy name—even the name of Jesus Christ the Lord—amen.

\newpage
\settalk{Obedience Brings Blessings}
\setspeaker{Thomas S. Monson}
\settalkdate{April 2013}
\talkheading{Obedience Brings Blessings}{Thomas S. Monson}

My beloved brothers and sisters, how grateful I am to be with you this morning. I seek an interest in your faith and prayers as I respond to the privilege to address you.

Throughout the ages, men and women have sought for knowledge and understanding concerning this mortal existence and their place and purpose in it, as well as for the way to peace and happiness. Such a search is undertaken by each of us.

This knowledge and understanding are available to all mankind. They are contained in truths which are eternal. In Doctrine and Covenants section 1, verse 39, we read, “Behold, and lo, the Lord is God, and the Spirit beareth record, and the record is true, and the truth abideth forever and ever.”

The poet wrote:

\begin{verse}
Tho the heavens depart and the earth’s fountains burst,\\
Truth, the sum of existence, will weather the worst,\\
Eternal, unchanged, evermore.
\end{verse}

Some would ask, “Where is such truth to be found, and how are we to recognize it?” In a revelation given through the Prophet Joseph Smith at Kirtland, Ohio, in May of 1833, the Lord declared:

“Truth is knowledge of things as they are, and as they were, and as they are to come. …

“The Spirit of truth is of God. …

“And no man receiveth a fulness unless he keepeth his commandments.

“He that keepeth [God’s] commandments receiveth truth and light, until he is glorified in truth and knoweth all things.”

What a glorious promise! “He that keepeth [God’s] commandments receiveth truth and light, until he is glorified in truth and knoweth all things.”

There is no need for you or for me, in this enlightened age when the fulness of the gospel has been restored, to sail uncharted seas or to travel unmarked roads in search of truth. A loving Heavenly Father has plotted our course and provided an unfailing guide—even obedience. A knowledge of truth and the answers to our greatest questions come to us as we are obedient to the commandments of God.

We learn obedience throughout our lives. Beginning when we are very young, those responsible for our care set forth guidelines and rules to ensure our safety. Life would be simpler for all of us if we would obey such rules completely. Many of us, however, learn through experience the wisdom of being obedient.

When I was growing up, each summer from early July until early September, my family stayed at our cabin at Vivian Park in Provo Canyon in Utah.

One of my best friends during those carefree days in the canyon was Danny Larsen, whose family also owned a cabin at Vivian Park. Each day he and I roamed this boy’s paradise, fishing in the stream and the river, collecting rocks and other treasures, hiking, climbing, and simply enjoying each minute of each hour of each day.

One morning Danny and I decided we wanted to have a campfire that evening with all our canyon friends. We just needed to clear an area in a nearby field where we could all gather. The June grass which covered the field had become dry and prickly, making the field unsuitable for our purposes. We began to pull at the tall grass, planning to clear a large, circular area. We tugged and yanked with all our might, but all we could get were small handfuls of the stubborn weeds. We knew this task would take the entire day, and already our energy and enthusiasm were waning.

And then what I thought was the perfect solution came into my eight-year-old mind. I said to Danny, “All we need is to set these weeds on fire. We’ll just burn a circle in the weeds!” He readily agreed, and I ran to our cabin to get a few matches.

Lest any of you think that at the tender age of eight we were permitted to use matches, I want to make it clear that both Danny and I were forbidden to use them without adult supervision. Both of us had been warned repeatedly of the dangers of fire. However, I knew where my family kept the matches, and we needed to clear that field. Without so much as a second thought, I ran to our cabin and grabbed a few matchsticks, making certain no one was watching. I hid them quickly in one of my pockets.

Back to Danny I ran, excited that in my pocket I had the solution to our problem. I recall thinking that the fire would burn only as far as we wanted and then would somehow magically extinguish itself.

I struck a match on a rock and set the parched June grass ablaze. It ignited as though it had been drenched in gasoline. At first Danny and I were thrilled as we watched the weeds disappear, but it soon became apparent that the fire was not about to go out on its own. We panicked as we realized there was nothing we could do to stop it. The menacing flames began to follow the wild grass up the mountainside, endangering the pine trees and everything else in their path.

Finally we had no option but to run for help. Soon all available men and women at Vivian Park were dashing back and forth with wet burlap bags, beating at the flames in an attempt to extinguish them. After several hours the last remaining embers were smothered. The ages-old pine trees had been saved, as were the homes the flames would eventually have reached.

Danny and I learned several difficult but important lessons that day—not the least of which was the importance of obedience.

There are rules and laws to help ensure our physical safety. Likewise, the Lord has provided guidelines and commandments to help ensure our spiritual safety so that we might successfully navigate this often-treacherous mortal existence and return eventually to our Heavenly Father.

Centuries ago, to a generation steeped in the tradition of animal sacrifice, Samuel boldly declared, “To obey is better than sacrifice, and to hearken than the fat of rams.”

In this dispensation, the Lord revealed to the Prophet Joseph Smith that He requires “the heart and a willing mind; and the willing and obedient shall eat the good of the land of Zion in these last days.”

All prophets, ancient and modern, have known that obedience is essential to our salvation. Nephi declared, “I will go and do the things which the Lord hath commanded.” Though others faltered in their faith and their obedience, never once did Nephi fail to do that which the Lord asked of him. Untold generations have been blessed as a result.

A soul-stirring account of obedience is that of Abraham and Isaac. How painfully difficult it must have been for Abraham, in obedience to God’s command, to take his beloved Isaac into the land of Moriah to offer him as a sacrifice. Can we imagine the heaviness of Abraham’s heart as he journeyed to the appointed place? Surely anguish must have racked his body and tortured his mind as he bound Isaac, laid him on the altar, and took the knife to slay him. With unwavering faith and implicit trust in the Lord, he responded to the Lord’s command. How glorious was the pronouncement, and with what wondered welcome did it come: “Lay not thine hand upon the lad, neither do thou any thing unto him: for now I know that thou fearest God, seeing thou hast not withheld thy son, thine only son from me.”

Abraham had been tried and tested, and for his faithfulness and obedience the Lord gave him this glorious promise: “In thy seed shall all the nations of the earth be blessed; because thou hast obeyed my voice.”

Although we are not asked to prove our obedience in such a dramatic and heart-wrenching way, obedience is required of us as well.

Declared President Joseph F. Smith in October 1873, “Obedience is the first law of heaven.”

Said President Gordon B. Hinckley, “The happiness of the Latter-day Saints, the peace of the Latter-day Saints, the progress of the Latter-day Saints, the prosperity of the Latter-day Saints, and the eternal salvation and exaltation of this people lie in walking in obedience to the counsels of … God.”

Obedience is a hallmark of prophets; it has provided strength and knowledge to them throughout the ages. It is essential for us to realize that we, as well, are entitled to this source of strength and knowledge. It is readily available to each of us today as we obey God’s commandments.

Throughout the years, I have known countless individuals who have been particularly faithful and obedient. I have been blessed and inspired by them. May I share with you an account of two such individuals.

Walter Krause was a steadfast member of the Church who, with his family, lived in what became known as East Germany following the Second World War. Despite the hardships he faced because of the lack of freedom in that area of the world at the time, Brother Krause was a man who loved and served the Lord. He faithfully and conscientiously fulfilled each assignment given to him.

The other man, Johann Denndorfer, a native of Hungary, was converted to the Church in Germany and was baptized there in 1911 at the age of 17. Not too long afterward he returned to Hungary. Following the Second World War, he found himself virtually a prisoner in his native land, in the city of Debrecen. Freedom had also been taken from the people of Hungary.

Brother Walter Krause, who did not know Brother Denndorfer, received the assignment to be his home teacher and to visit him on a regular basis. Brother Krause called his home teaching companion and said to him, “We have received an assignment to visit Brother Johann Denndorfer. Would you be available to go with me this week to see him and give him a gospel message?” And then he added, “Brother Denndorfer lives in Hungary.”

His startled companion asked, “When will we leave?”

“Tomorrow,” came the reply from Brother Krause.

“When will we return home?” asked the companion.

Brother Krause responded, “Oh, in about a week—if we get back.”

Away the two home teaching companions went to visit Brother Denndorfer, traveling by train and bus from the northeastern area of Germany to Debrecen, Hungary—a substantial journey. Brother Denndorfer had not had home teachers since before the war. Now, when he saw these servants of the Lord, he was overwhelmed with gratitude that they had come. At first he declined to shake hands with them. Rather, he went to his bedroom and took from a small cabinet a box containing his tithing that he had saved for years. He presented the tithing to his home teachers and said, “Now I am current with the Lord. Now I feel worthy to shake the hands of servants of the Lord!” Brother Krause told me later that he had been touched beyond words to think that this faithful brother, who had no contact with the Church for many years, had obediently and consistently taken from his meager earnings 10 percent with which to pay his tithing. He had saved it not knowing when or if he might have the privilege of paying it.

Brother Walter Krause passed away nine years ago at the age of 94. He served faithfully and obediently throughout his life and was an inspiration to me and to all who knew him. When asked to fulfill assignments, he never questioned, he never murmured, and he never made excuses.

My brothers and sisters, the great test of this life is obedience. “We will prove them herewith,” said the Lord, “to see if they will do all things whatsoever the Lord their God shall command them.”

Declared the Savior, “For all who will have a blessing at my hands shall abide the law which was appointed for that blessing, and the conditions thereof, as were instituted from before the foundation of the world.”

No greater example of obedience exists than that of our Savior. Of Him, Paul observed:

“Though he were a Son, yet learned he obedience by the things which he suffered;

“And being made perfect, he became the author of eternal salvation unto all them that obey him.”

The Savior demonstrated genuine love of God by living the perfect life, by honoring the sacred mission that was His. Never was He haughty. Never was He puffed up with pride. Never was He disloyal. Ever was He humble. Ever was He sincere. Ever was He obedient.

Though He was tempted by that master of deceit, even the devil, though He was physically weakened from fasting 40 days and 40 nights and was an hungered, yet when the evil one proffered Jesus the most alluring and tempting proposals, He gave to us a divine example of obedience by refusing to deviate from what He knew was right.

When faced with the agony of Gethsemane, where He endured such pain that “his sweat was as it were great drops of blood falling down to the ground,” He exemplified the obedient Son by saying, “Father, if thou be willing, remove this cup from me: nevertheless not my will, but thine, be done.”

As the Savior instructed His early Apostles, so He instructs you and me, “Follow thou me.” Are we willing to obey?

The knowledge which we seek, the answers for which we yearn, and the strength which we desire today to meet the challenges of a complex and changing world can be ours when we willingly obey the Lord’s commandments. I quote once again the words of the Lord: “He that keepeth [God’s] commandments receiveth truth and light, until he is glorified in truth and knoweth all things.”

It is my humble prayer that we may be blessed with the rich rewards promised to the obedient. In the name of Jesus Christ, our Lord and Savior, amen.

\newpage
\settalk{Until We Meet Again}
\setspeaker{Thomas S. Monson}
\settalkdate{April 2013}
\talkheading{Until We Meet Again}{Thomas S. Monson}

My brothers and sisters, what a glorious conference we have had. I know you will agree with me that the messages have been inspiring. Our hearts have been touched, and our testimonies of this divine work have been strengthened as we have felt the Spirit of the Lord. May we long remember what we have heard these past two days. I urge you to study the messages further when they are printed in coming issues of the Ensign and Liahona magazines.

We express our gratitude to each one who has spoken to us, as well as to those who have offered prayers. In addition, the music has been uplifting and inspiring. We love our wonderful Tabernacle Choir and thank all others who provided music as well.

We join together in expressing our gratitude to those of the presidency and board of the general Young Women, who were released yesterday. Their service has been outstanding and their dedication complete.

We have sustained, by uplifted hands, brethren and sisters who have been called to new positions during this conference. We want all of them to know that we look forward to serving with them in the cause of the Master.

We are a worldwide Church, brothers and sisters. Our membership is found across the globe. I admonish you to be good citizens of the nations in which you live and good neighbors in your communities, reaching out to those of other faiths as well as to our own. May we be tolerant of, as well as kind and loving to, those who do not share our beliefs and our standards. The Savior brought to this earth a message of love and goodwill to all men and women. May we ever follow His example.

I pray that we may be aware of the needs of those around us. There are some, particularly among the young, who are tragically involved in drugs, immorality, pornography, and so on. There are those who are lonely, including widows and widowers, who long for the company and concern of others. May we ever be ready to extend to them a helping hand and a loving heart.

We live at a time in the world’s history when there are many difficult challenges but also great opportunities and reasons for rejoicing. There are, of course, those times when we experience disappointments, heartaches, and even tragedies in our lives. However, if we will put our trust in the Lord, He will help us through our difficulties, whatever they may be. The Psalmist provided this assurance: “Weeping may endure for a night, but joy cometh in the morning.”

My brothers and sisters, I want you to know how grateful I am for the gospel of Jesus Christ, restored in these latter days through the Prophet Joseph Smith. It is the key to our happiness. May we be humble and prayerful, having the faith that our Heavenly Father can guide and bless us in our lives.

I bear my personal witness and testimony to you that God lives, that He hears the prayers of humble hearts. His Son, our Savior and Redeemer, speaks to each of us: “Behold, I stand at the door, and knock: if any man hear my voice, and open the door, I will come in to him.” May we believe these words and take advantage of this promise.

As this conference now concludes, I invoke the blessings of heaven upon each of you. May your homes be filled with peace, harmony, courtesy, and love. May they be filled with the Spirit of the Lord. May you nurture and nourish your testimonies of the gospel, that they will be a protection to you against the buffetings of Satan.

Until we meet again in six months, I pray that the Lord will bless and keep you, my brothers and sisters. May His promised peace be with you now and always. Thank you for your prayers in my behalf and in behalf of all of the General Authorities. We are deeply grateful for you. In the name of our Savior and Redeemer, whom we serve, even Jesus Christ, the Lord, amen.

\newpage
\settalk{Welcome to Conference}
\setspeaker{Thomas S. Monson}
\settalkdate{April 2013}
\talkheading{Welcome to Conference}{Thomas S. Monson}

My beloved brothers and sisters, how pleased I am to welcome you to the 183rd Annual General Conference of the Church.

During the six months since last we met, it has been my opportunity to travel a bit and to meet with some of you in your own areas. Following general conference in October, I traveled to Germany, where it was my privilege to meet with our members at several locations in that country as well as in parts of Austria.

At the end of October I dedicated the Calgary Alberta Temple in Canada, with the assistance of Elder and Sister M. Russell Ballard, Elder and Sister Craig C. Christensen, and Elder and Sister William R. Walker. In November I rededicated the Boise Idaho Temple. Also traveling with me and participating in the dedication were Elder and Sister David A. Bednar, Elder and Sister Craig C. Christensen, and Elder and Sister William R. Walker.

The cultural celebrations held in conjunction with both of these dedications were outstanding. I did not personally attend the cultural celebration in Calgary, inasmuch as it was Sister Monson’s 85th birthday and I felt I should be with her. However, she and I were privileged to watch the celebration in our living room over closed-circuit television, and then I flew to Calgary the following morning for the dedication. In Boise over 9,000 youth from the temple district participated in the cultural celebration. There were so many young people involved that there was not room for family members to attend in the arena in which they performed.

Just last month President Dieter F. Uchtdorf, accompanied by Sister Uchtdorf, Elder and Sister Jeffrey R. Holland, and Elder and Sister Gregory A. Schwitzer, traveled to Tegucigalpa, Honduras, to dedicate our newly completed temple there. A magnificent youth celebration took place the evening prior to the dedication.

There are other temples which have been announced and which are at various stages in the preliminary process or which are under construction.

It is my privilege this morning to announce two additional temples, which in coming months and years will be built in the following locations: Cedar City, Utah, and Rio de Janeiro, Brazil. Brothers and sisters, temple building continues unabated.

As you know, in the October general conference I announced changes in the ages at which young men and young women might serve as full-time missionaries, with the young men now being able to serve at age 18 and the young women at 19.

The response of our young people has been remarkable and inspiring. As of April 4—two days ago—we have 65,634 full-time missionaries serving, with over 20,000 more who have received their calls but who have not yet entered a missionary training center and over 6,000 more in the interview process with their bishops and stake presidents. It has been necessary for us to create 58 new missions to accommodate the increased numbers of missionaries.

To help maintain this missionary force, and because many of our missionaries come from modest circumstances, we invite you, as you are able, to contribute generously to the General Missionary Fund of the Church.

Now, brothers and sisters, we will hear inspired messages today and tomorrow. Those who will address us have sought prayerfully to know that which the Lord would have us hear at this time.

I urge you to be attentive and receptive to the messages which we will hear. That we may do so is my prayer in the name of Jesus Christ, the Lord, amen.

\newpage
\settalk{“I Will Not Fail Thee, nor Forsake Thee”}
\setspeaker{Thomas S. Monson}
\settalkdate{October 2013}
\talkheading{“I Will Not Fail Thee, nor Forsake Thee”}{Thomas S. Monson}

In my journal tonight, I shall write, “This has been one of the most inspiring sessions of any general conference I’ve attended. Everything has been of the greatest and most spiritual nature.”

Brothers and sisters, six months ago as we met together in our general conference, my sweet wife, Frances, lay in the hospital, having suffered a devastating fall just a few days earlier. In May, after weeks of valiantly struggling to overcome her injuries, she slipped into eternity. Her loss has been profound. She and I were married in the Salt Lake Temple on October 7, 1948. Tomorrow would have been our 65th wedding anniversary. She was the love of my life, my trusted confidant, and my closest friend. To say that I miss her does not begin to convey the depth of my feelings.

This conference marks 50 years since I was called to the Quorum of the Twelve Apostles by President David O. McKay. Through all these years I have felt nothing but the full and complete support of my sweet companion. Countless are the sacrifices she made so that I could fulfill my calling. Never did I hear a word of complaint from her as I was often required to spend days and sometimes weeks away from her and from our children. She was an angel, indeed.

I wish to express my thanks, as well as those of my family, for the tremendous outpouring of love which has come to us since Frances’s passing. Hundreds of cards and letters were sent from around the world expressing admiration for her and condolences to our family. We received dozens of beautiful floral arrangements. We are grateful for the numerous contributions which have been offered in her name to the General Missionary Fund of the Church. On behalf of those of us whom she left behind, I express deep gratitude for your kind and heartfelt expressions.

Of utmost comfort to me during this tender time of parting have been my testimony of the gospel of Jesus Christ and the knowledge I have that my dear Frances lives still. I know that our separation is temporary. We were sealed in the house of God by one having authority to bind on earth and in heaven. I know that we will be reunited one day and will never again be separated. This is the knowledge that sustains me.

Brothers and sisters, it may be safely assumed that no person has ever lived entirely free of suffering and sorrow, nor has there ever been a period in human history that did not have its full share of turmoil and misery.

When the pathway of life takes a cruel turn, there is the temptation to ask the question “Why me?” At times there appears to be no light at the end of the tunnel, no sunrise to end the night’s darkness. We feel encompassed by the disappointment of shattered dreams and the despair of vanished hopes. We join in uttering the biblical plea, “Is there no balm in Gilead?” We feel abandoned, heartbroken, alone. We are inclined to view our own personal misfortunes through the distorted prism of pessimism. We become impatient for a solution to our problems, forgetting that frequently the heavenly virtue of patience is required.

The difficulties which come to us present us with the real test of our ability to endure. A fundamental question remains to be answered by each of us: Shall I falter, or shall I finish? Some do falter as they find themselves unable to rise above their challenges. To finish involves enduring to the very end of life itself.

As we ponder the events that can befall all of us, we can say with Job of old, “Man is born unto trouble.” Job was a “perfect and upright” man who “feared God, and eschewed evil.” Pious in his conduct, prosperous in his fortune, Job was to face a test which could have destroyed anyone. Shorn of his possessions, scorned by his friends, afflicted by his suffering, shattered by the loss of his family, he was urged to “curse God, and die.” He resisted this temptation and declared from the depths of his noble soul:

“Behold, my witness is in heaven, and my record is on high.”

“I know that my redeemer liveth.”

Job kept the faith. Will we do likewise as we face those challenges which will be ours?

Whenever we are inclined to feel burdened down with the blows of life, let us remember that others have passed the same way, have endured, and then have overcome.

The history of the Church in this, the dispensation of the fulness of times, is replete with the experiences of those who have struggled and yet who have remained steadfast and of good cheer. The reason? They have made the gospel of Jesus Christ the center of their lives. This is what will pull us through whatever comes our way. We will still experience difficult challenges, but we will be able to face them, to meet them head-on, and to emerge victorious.

From the bed of pain, from the pillow wet with tears, we are lifted heavenward by that divine assurance and precious promise: “I will not fail thee, nor forsake thee.” Such comfort is priceless.

As I have traveled far and wide throughout the world fulfilling the responsibilities of my calling, I have come to know many things—not the least of which is that sadness and suffering are universal. I cannot begin to measure all of the heartache and sorrow I have witnessed as I have visited with those who are dealing with grief, experiencing illness, facing divorce, struggling with a wayward son or daughter, or suffering the consequences of sin. The list could go on and on, for there are countless problems which can befall us. To single out one example is difficult, and yet whenever I think of challenges, my thoughts turn to Brother Brems, one of my boyhood Sunday School teachers. He was a faithful member of the Church, a man with a heart of gold. He and his wife, Sadie, had eight children, many of whom were the same ages as those in our family.

After Frances and I were married and moved from the ward, we saw Brother and Sister Brems and members of their family at weddings and funerals, as well as at ward reunions.

In 1968, Brother Brems lost his wife, Sadie. Two of his eight children also passed away as the years went by.

One day nearly 13 years ago, Brother Brems’s oldest granddaughter telephoned me. She explained that her grandfather had reached his 105th birthday. She said, “He lives in a small care center but meets with his entire family each Sunday, where he delivers a gospel lesson.” She continued, “This past Sunday, Grandpa announced to us, ‘My dears, I am going to die this week. Will you please call Tommy Monson. He will know what to do.’”

I visited Brother Brems the very next evening. I had not seen him for a while. I could not speak to him, for he had lost his hearing. I could not write a message for him to read, because he had lost his sight. I was told that the family communicated with him by taking the finger of his right hand and then tracing on the palm of his left hand the name of the person visiting. Any message had to be conveyed in this same way. I followed the procedure by taking his finger and spelling T-O-M-M-Y M-O-N-S-O-N, the name by which he had always known me. Brother Brems became excited and, taking my hands, placed them on his head. I knew his desire was to receive a priesthood blessing. The driver who had taken me to the care center joined me as we placed our hands on the head of Brother Brems and provided the desired blessing. Afterward, tears streamed from his sightless eyes. He grasped our hands in gratitude. Although he had not heard the blessing we had given him, the Spirit was strong, and I believe he was inspired to know we had provided the blessing which he needed. This sweet man could no longer see. He could no longer hear. He was confined night and day to a small room in a care center. And yet the smile on his face and the words he spoke touched my heart. “Thank you,” he said. “My Heavenly Father has been so good to me.”

Within a week, just as Brother Brems had predicted, he passed away. Never did he dwell on what he was lacking; rather, he was always deeply grateful for his many blessings.

Our Heavenly Father, who gives us so much to delight in, also knows that we learn and grow and become stronger as we face and survive the trials through which we must pass. We know that there are times when we will experience heartbreaking sorrow, when we will grieve, and when we may be tested to our limits. However, such difficulties allow us to change for the better, to rebuild our lives in the way our Heavenly Father teaches us, and to become something different from what we were—better than we were, more understanding than we were, more empathetic than we were, with stronger testimonies than we had before.

This should be our purpose—to persevere and endure, yes, but also to become more spiritually refined as we make our way through sunshine and sorrow. Were it not for challenges to overcome and problems to solve, we would remain much as we are, with little or no progress toward our goal of eternal life. The poet expressed much the same thought in these words:

\begin{verse}
Good timber does not grow with ease,\\
The stronger wind, the stronger trees.\\
The further sky, the greater length.\\
The more the storm, the more the strength.\\
By sun and cold, by rain and snow,\\
In trees and men good timbers grow.
\end{verse}

Only the Master knows the depths of our trials, our pain, and our suffering. He alone offers us eternal peace in times of adversity. He alone touches our tortured souls with His comforting words:

“Come unto me, all ye that labour and are heavy laden, and I will give you rest.

“Take my yoke upon you, and learn of me; for I am meek and lowly in heart: and ye shall find rest unto your souls.

“For my yoke is easy, and my burden is light.”

Whether it is the best of times or the worst of times, He is with us. He has promised that this will never change.

My brothers and sisters, may we have a commitment to our Heavenly Father that does not ebb and flow with the years or the crises of our lives. We should not need to experience difficulties for us to remember Him, and we should not be driven to humility before giving Him our faith and trust.

May we ever strive to be close to our Heavenly Father. To do so, we must pray to Him and listen to Him every day. We truly need Him every hour, whether they be hours of sunshine or of rain. May His promise ever be our watchword: “I will not fail thee, nor forsake thee.”

With all the strength of my soul, I testify that God lives and loves us, that His Only Begotten Son lived and died for us, and that the gospel of Jesus Christ is that penetrating light which shines through the darkness of our lives. May it ever be so, I pray in the sacred name of Jesus Christ, amen.

\newpage
\settalk{Till We Meet Again}
\setspeaker{Thomas S. Monson}
\settalkdate{October 2013}
\talkheading{Till We Meet Again}{Thomas S. Monson}

My brothers and sisters, my heart is full as we bring to a close this wonderful general conference of the Church. We have been spiritually fed as we have listened to the counsel and testimonies of those who have participated in each session.

We have been blessed to meet here in the magnificent Conference Center in peace and safety. We have had unprecedented coverage of the conference, reaching across the continents to people everywhere. Though we are physically far removed from many of you, we feel of your spirit.

To our Brethren who have been released at this conference, may I express the heartfelt thanks of the entire Church for your years of devoted service. Countless are those who have been blessed by your contributions to the work of the Lord.

I express gratitude to the Tabernacle Choir and to the other choirs which participated in this conference. The music has been beautiful and has added greatly to the Spirit we have felt at each session.

I thank you for your prayers in my behalf and in behalf of all the General Authorities and general officers of the Church. We are strengthened by them.

May heaven’s blessings be with you. May your homes be filled with love and courtesy and with the Spirit of the Lord. May you constantly nourish your testimonies of the gospel that they will be a protection to you against the buffetings of the adversary.

Conference is now over. As we return to our homes, may we do so safely. May the Spirit we have felt here be and abide with us as we go about those things which occupy us each day. May we show increased kindness toward one another, and may we ever be found doing the work of the Lord.

My brothers and sisters, may God bless you. May His promised peace be with you now and always. I bid you farewell until we meet again in six months’ time. In the name of our Savior, even Jesus Christ the Lord, amen.

\newpage
\settalk{True Shepherds}
\setspeaker{Thomas S. Monson}
\settalkdate{October 2013}
\talkheading{True Shepherds}{Thomas S. Monson}

Tonight in the Conference Center in Salt Lake City and in locations far and near are assembled those who bear the priesthood of God. Truly you are “a royal priesthood”—even “a chosen generation,” as the Apostle Peter declared. I am honored to have the privilege to address you.

When I was growing up, each summer our family would drive to Provo Canyon, about 45 miles (72 km) south and a little east of Salt Lake City, where we would stay in the family cabin for several weeks. We boys were always anxious to get on the fishing stream or into the swimming hole, and we would try to push the car a little faster. In those days, the automobile my father drove was a 1928 Oldsmobile. If he went over 35 miles (56 km) an hour, my mother would say, “Keep it down! Keep it down!” I would say, “Put the accelerator down, Dad! Put it down!”

Dad would drive about 35 miles an hour all the way up to Provo Canyon or until we would come around a bend in the road and our journey would be halted by a herd of sheep. We would watch as hundreds of sheep filed past us, seemingly without a shepherd, a few dogs yapping at their heels as they moved along. Way back in the rear we could see the sheepherder on his horse—not a bridle on it but a halter. He was occasionally slouched down in the saddle dozing, since the horse knew which way to go and the yapping dogs did the work.

Contrast that to the scene which I viewed in Munich, Germany, many years ago. It was a Sunday morning, and we were en route to a missionary conference. As I looked out the window of the mission president’s automobile, I saw a shepherd with a staff in his hand, leading the sheep. They followed him wherever he went. If he moved to the left, they followed him to the left. If he moved to the right, they followed him in that direction. I made the comparison between the true shepherd who led his sheep and the sheepherder who rode casually behind his sheep.

Jesus said, “I am the good shepherd, and know my sheep.” He provides for us the perfect example of what a true shepherd should be.

Brethren, as the priesthood of God we have a shepherding responsibility. The wisdom of the Lord has provided guidelines whereby we might be shepherds to the families of the Church, where we can serve, we can teach, and we can testify to them. Such is called home teaching, and it is about this that I wish to speak to you tonight.

The bishop of each ward in the Church oversees the assigning of priesthood holders as home teachers to visit the homes of members every month. They go in pairs. Where possible, a young man who is a priest or a teacher in the Aaronic Priesthood accompanies an adult holding the Melchizedek Priesthood. As they go into the homes of those for whom they are responsible, the Aaronic Priesthood holder should take part in the teaching which takes place. Such an assignment will help to prepare these young men for missions as well as for a lifetime of priesthood service.

The home teaching program is a response to modern revelation commissioning those ordained to the priesthood “to teach, expound, exhort, baptize, … and visit the house of each member, and exhort them to pray vocally and in secret and attend to all family duties, … to watch over the church always, and be with and strengthen them; and see that there is no iniquity in the church, neither hardness with each other, neither lying, backbiting, nor evil speaking.”

President David O. McKay admonished: “Home teaching is one of our most urgent and most rewarding opportunities to nurture and inspire, to counsel and direct our Father’s children. … [It] is a divine service, a divine call. It is our duty as Home Teachers to carry the … spirit into every home and heart. To love the work and do our best will bring unbounded peace, joy and satisfaction to [a noble,] dedicated [teacher] of God’s children.”

From the Book of Mormon we read that Alma “consecrated all their priests and all their teachers; and none were consecrated except they were just men.

“Therefore they did watch over their people, and did nourish them with things pertaining to righteousness.”

In performing our home teaching responsibilities, we are wise if we learn and understand the challenges of the members of each family, that we might be effective in teaching and in providing needed assistance.

A home teaching visit is also more likely to be successful if an appointment is made in advance. To illustrate this point, let me share with you an experience I had some years ago. At that time the Missionary Executive Committee was comprised of Spencer W. Kimball, Gordon B. Hinckley, and Thomas S. Monson. One evening Brother and Sister Hinckley hosted a dinner in their home for the committee members and our wives. We had just finished a lovely meal when there was a knock at the door. President Hinckley opened the door and found one of his home teachers standing there. The home teacher said, “I know I didn’t make an appointment to come, and I don’t have with me my companion, but I felt I should come tonight. I didn’t know you would be entertaining company.”

President Hinckley graciously invited the home teacher to come in and sit down and to instruct three Apostles and our wives concerning our duty as members. With a bit of trepidation, the home teacher did his best. President Hinckley thanked him for coming, after which he made a hurried exit.

I mention one more example of the incorrect way to accomplish home teaching. President Marion G. Romney, who was a counselor in the First Presidency some years ago, used to tell about his home teacher who once went to the Romney home on a cold winter night. He kept his hat in his hand and shifted nervously when invited to sit down and give his message. As he remained standing, he said, “Well, I’ll tell you, Brother Romney, it’s cold outside, and I left my car engine running so it wouldn’t stop. I just came by so I could tell the bishop I had made my visits.”

President Ezra Taft Benson, after relating President Romney’s experience in a meeting of priesthood holders, then said, “We can do better than that, brethren—much better!” I agree.

Home teaching is more than a mechanical visit once per month. Ours is the responsibility to teach, to inspire, to motivate, and where we visit those who are not active, to bring to activity and to eventual exaltation the sons and daughters of God.

To assist in our efforts, I share this wise counsel which surely applies to home teachers. It comes from Abraham Lincoln, who said, “If you would win a man to your cause, first convince him that you are his sincere friend.” President Ezra Taft Benson urged: “Above all, be a genuine friend to the individuals and families you teach. … A friend makes more than a dutiful visit each month. A friend is more concerned about helping people than getting credit. A friend cares. A friend [shows love]. A friend listens, and a friend reaches out.”

Home teaching answers many prayers and permits us to see the transformations which can take place in people’s lives.

An example of this would be Dick Hammer, who came to Utah with the Civilian Conservation Corps during the Depression. He met and married a Latter-day Saint young woman. He opened Dick’s Café in St. George, Utah, which became a popular meeting spot.

Assigned as home teacher to the Hammer family was Willard Milne, a friend of mine. Since I knew Dick Hammer as well, having printed the menus for his café, I would ask my friend Brother Milne when I visited St. George, “How is our friend Dick Hammer coming?”

The reply would generally be, “He’s coming, but slowly.”

When Willard Milne and his companion visited the Hammer home each month, they always managed to present a gospel message and to share their testimonies with Dick and the family.

The years passed by, and then one day Willard phoned me with good news. “Brother Monson,” he began, “Dick Hammer is converted and is going to be baptized. He is in his 90th year, and we have been friends all our adult lives. His decision warms my heart. I’ve been his home teacher for many years.” There was a catch in Willard’s voice as he conveyed his welcome message.

Brother Hammer was indeed baptized and a year later entered that beautiful St. George Temple and there received his endowment and sealing blessings.

I asked Willard, “Did you ever become discouraged as his home teacher for such a long time?”

He replied, “No, it was worth every effort. As I witness the joy which has come to the members of the Hammer family, my heart fills with gratitude for the blessings the gospel has brought into their lives and for the privilege I have had to help in some way. I am a happy man.”

Brethren, it will be our privilege through the years to visit and teach many individuals—those who are less active as well as those who are fully committed. If we are conscientious in our calling, we will have many opportunities to bless lives. Our visits to those who have distanced themselves from Church activity can be the key which will eventually open the doors to their return.

With this thought in mind, let us reach out to those for whom we are responsible and bring them to the table of the Lord to feast on His word and to enjoy the companionship of His Spirit and be “no more strangers and foreigners, but fellowcitizens with the saints, and of the household of God.”

If any of you has slipped into complacency concerning your home teaching visits, may I say that there is no time like the present to rededicate yourself to fulfilling your home teaching duties. Decide now to make whatever effort is necessary to reach those for whom you have been given responsibility. There are times when a little extra prodding may be needed, as well, to help your home teaching companion find the time to go with you, but if you are persistent, you will succeed.

Brethren, our efforts in home teaching are ongoing. The work will never be concluded until our Lord and Master says, “It is enough.” There are lives to brighten. There are hearts to touch. There are souls to save. Ours is the sacred privilege to brighten, to touch, and to save those precious souls entrusted to our care. We should do so faithfully and with hearts filled with gladness.

In closing I turn to one particular example to describe the type of home teachers we should be. There is one Teacher whose life overshadows all others. He taught of life and death, of duty and destiny. He lived not to be served but to serve, not to receive but to give, not to save His life but to sacrifice it for others. He described a love more beautiful than lust, a poverty richer than treasure. It was said of this Teacher that He taught with authority and not as did the scribes. His laws were not inscribed upon stone but upon human hearts.

I speak of the Master Teacher, even Jesus Christ, the Son of God, the Savior and Redeemer of all mankind. The biblical account says of Him, He “went about doing good.” With Him as our unfailing guide and exemplar, we shall qualify for His divine help in our home teaching. Lives will be blessed. Hearts will be comforted. Souls will be saved. We will become true shepherds. That this may be so, I pray in the name of that great Shepherd, Jesus Christ, amen.

\newpage
\settalk{We Never Walk Alone}
\setspeaker{Thomas S. Monson}
\settalkdate{October 2013}
\talkheading{We Never Walk Alone}{Thomas S. Monson}

My dear sisters, the spirit we feel this evening is a reflection of your strength, your devotion, and your goodness. To quote the Master: “Ye are the salt of the earth. … Ye are the light of the world.”

As I have contemplated my opportunity to address you, I have been reminded of the love my dear wife, Frances, had for Relief Society. During her lifetime she served in many positions in Relief Society. When she and I were both just 31 years of age, I was called to be president of the Canadian Mission. During the three years of that assignment, Frances presided over all of the Relief Societies in that vast area, which encompassed the provinces of Ontario and Quebec. Some of her closest friendships came as a result of that assignment, as well as from the many callings she later filled in our own ward Relief Society. She was a faithful daughter of our Heavenly Father, my beloved companion, and my dearest friend. I miss her more than words can express.

I too love Relief Society. I testify to you that it was organized by inspiration and is a vital part of the Lord’s Church here upon the earth. It would be impossible to calculate all the good which has come from this organization and all the lives which have been blessed because of it.

Relief Society is made up of a variety of women. There are those of you who are single—perhaps in school, perhaps working—yet forging a full and rich life. Some of you are busy mothers of growing children. Still others of you have lost your husbands because of divorce or death and are struggling to raise your children without the help of a husband and father. Some of you have raised your children but have realized that their need for your help is ongoing. There are many of you who have aging parents who require the loving care only you can give.

Wherever we are in life, there are times when all of us have challenges and struggles. Although they are different for each, they are common to all.

Many of the challenges we face exist because we live in this mortal world, populated by all manner of individuals. At times we ask in desperation, “How can I keep my sights firmly fixed on the celestial as I navigate through this telestial world?”

There will be times when you will walk a path strewn with thorns and marked by struggle. There may be times when you feel detached—even isolated—from the Giver of every good gift. You worry that you walk alone. Fear replaces faith.

When you find yourself in such circumstances, I plead with you to remember prayer. I love the words of President Ezra Taft Benson concerning prayer. Said he:

“All through my life the counsel to depend on prayer has been prized above almost any other advice I have … received. It has become an integral part of me—an anchor, a constant source of strength, and the basis of my knowledge of things divine. …

“… Though reverses come, in prayer we can find reassurance, for God will speak peace to the soul. That peace, that spirit of serenity, is life’s greatest blessing.”

The Apostle Paul admonished:

“Let your requests be made known unto God.

“And the peace of God, which passeth all understanding, shall keep your hearts and minds through Christ Jesus.”

What a glorious promise! Peace is that which we seek, that for which we yearn.

We were not placed on this earth to walk alone. What an amazing source of power, of strength, and of comfort is available to each of us. He who knows us better than we know ourselves, He who sees the larger picture and who knows the end from the beginning, has assured us that He will be there for us to provide help if we but ask. We have the promise: “Pray always, and be believing, and all things shall work together for your good.”

As our prayers ascend heavenward, let us not forget the words taught to us by the Savior. When He faced the excruciating agony of Gethsemane and the cross, He prayed to the Father, “Not my will, but thine, be done.” Difficult as it may at times be, it is for us, as well, to trust our Heavenly Father to know best how and when and in what manner to provide the help we seek.

I cherish the words of the poet:

\begin{verse}
I know not by what methods rare,\\
But this I know, God answers prayer.\\
I know that He has given His Word,\\
Which tells me prayer is always heard,\\
And will be answered, soon or late.\\
And so I pray and calmly wait.\\
I know not if the blessing sought\\
Will come in just the way I thought;\\
But leave my prayers with Him alone,\\
Whose will is wiser than my own,\\
Assured that He will grant my quest,\\
Or send some answer far more blest.
\end{verse}

Of course, prayer is not just for times of trouble. We are told repeatedly in the scriptures to “pray always” and to keep a prayer in our hearts. The words of a favorite and familiar hymn pose a question which we would do well to ask ourselves daily: “Did you think to pray?”

Allied with prayer in helping us cope in our often difficult world is scripture study. The words of truth and inspiration found in our four standard works are prized possessions to me. I never tire of reading them. I am lifted spiritually whenever I search the scriptures. These holy words of truth and love give guidance to my life and point the way to eternal perfection.

As we read and ponder the scriptures, we will experience the sweet whisperings of the Spirit to our souls. We can find answers to our questions. We learn of the blessings which come through keeping God’s commandments. We gain a sure testimony of our Heavenly Father and our Savior, Jesus Christ, and of Their love for us. When scripture study is combined with our prayers, we can of a certainty know that the gospel of Jesus Christ is true.

Said President Gordon B. Hinckley, “May the Lord bless each of us to feast upon his holy [words] and to draw from [them] that strength, that peace, [and] that knowledge ‘which passeth all understanding’ (Philip. 4:7).”

As we remember prayer and take time to turn to the scriptures, our lives will be infinitely more blessed and our burdens will be made lighter.

May I share with you the account of how our Heavenly Father answered the prayers and pleadings of one woman and provided her the peace and assurance she so desperately sought?

Tiffany’s difficulties began last year when she had guests at her home for Thanksgiving and then again for Christmas. Her husband had been in medical school and was now in the second year of his medical residency. Because of the long work hours required of him, he was not able to help her as much as they both would have liked, and so most of that which needed to be accomplished during this holiday season, in addition to the care of their four young children, fell to Tiffany. She was becoming overwhelmed, and then she learned that one who was dear to her had been diagnosed with cancer. The stress and worry began to take a heavy toll on her, and she slipped into a period of discouragement and depression. She sought medical help, and yet nothing changed. Her appetite disappeared, and she began to lose weight, which her tiny frame could ill afford. She sought peace through the scriptures and prayed for deliverance from the gloom which was overtaking her. When neither peace nor help seemed to come, she began to feel abandoned by God. Her family and friends prayed for her and tried desperately to help. They delivered her favorite foods in an attempt to keep her physically healthy, but she could take only a few bites and then would be unable to finish.

On one particularly trying day, a friend attempted in vain to entice her with foods she had always loved. When nothing worked, the friend said, “There must be something that sounds good to you.”

Tiffany thought for a moment and said, “The only thing I can think of that sounds good is homemade bread.”

But there was none on hand.

The following afternoon Tiffany’s doorbell rang. Her husband happened to be home and answered it. When he returned, he was carrying a loaf of homemade bread. Tiffany was astonished when he told her it had come from a woman named Sherrie, whom they barely knew. She was a friend of Tiffany’s sister Nicole, who lived in Denver, Colorado. Sherrie had been introduced to Tiffany and her husband briefly several months earlier when Nicole and her family were staying with Tiffany for Thanksgiving. Sherrie, who lived in Omaha, had come to Tiffany’s home to visit with Nicole.

Now, months later, with the delicious bread in hand, Tiffany called her sister Nicole to thank her for sending Sherrie on an errand of mercy. Instead, she learned Nicole had not instigated the visit and had no knowledge of it.

The rest of the story unfolded as Nicole checked with her friend Sherrie to find out what had prompted her to deliver that loaf of bread. What she learned was an inspiration to her, to Tiffany, to Sherrie—and it is an inspiration to me.

On that particular morning of the bread delivery, Sherrie had been prompted to make two loaves of bread instead of the one she had planned to make. She said she felt impressed to take the second loaf with her in her car that day, although she didn’t know why. After lunch at a friend’s home, her one-year-old daughter began to cry and needed to be taken home for a nap. Sherrie hesitated when the unmistakable feeling came to her that she needed to deliver that extra loaf of bread to Nicole’s sister Tiffany, who lived 30 minutes away on the other side of town and whom she barely knew. She tried to rationalize away the thought, wanting to get her very tired daughter home and feeling sheepish about delivering a loaf of bread to people who were almost strangers. However, the impression to go to Tiffany’s home was strong, so she heeded the prompting.

When she arrived, Tiffany’s husband answered the door. Sherrie reminded him that she was Nicole’s friend whom he’d met briefly at Thanksgiving, handed him the loaf of bread, and left.

And so it happened that the Lord sent a virtual stranger across town to deliver not just the desired homemade bread but also a clear message of love to Tiffany. What happened to her cannot be explained in any other way. She had an urgent need to feel that she wasn’t alone—that God was aware of her and had not abandoned her. That bread—the very thing she wanted—was delivered to her by someone she barely knew, someone who had no knowledge of her need but who listened to the prompting of the Spirit and followed that prompting. It became an obvious sign to Tiffany that her Heavenly Father was aware of her needs and loved her enough to send help. He had responded to her cries for relief.

My dear sisters, your Heavenly Father loves you—each of you. That love never changes. It is not influenced by your appearance, by your possessions, or by the amount of money you have in your bank account. It is not changed by your talents and abilities. It is simply there. It is there for you when you are sad or happy, discouraged or hopeful. God’s love is there for you whether or not you feel you deserve love. It is simply always there.

As we seek our Heavenly Father through fervent, sincere prayer and earnest, dedicated scripture study, our testimonies will become strong and deeply rooted. We will know of God’s love for us. We will understand that we do not ever walk alone. I promise you that you will one day stand aside and look at your difficult times, and you will realize that He was always there beside you. I know this to be true in the passing of my eternal companion—Frances Beverly Johnson Monson.

I leave with you my blessing. I leave with you my gratitude for all the good you do and for the lives you lead. That you may be blessed with every good gift is my prayer in the name of our Savior and Redeemer, even Jesus Christ the Lord, amen.

\newpage
\settalk{Welcome to Conference}
\setspeaker{Thomas S. Monson}
\settalkdate{October 2013}
\talkheading{Welcome to Conference}{Thomas S. Monson}

How good it is, my beloved brothers and sisters, to meet together once again. It has been just over 183 years since the Church was organized by the Prophet Joseph Smith, under the direction of the Lord. At that meeting on April 6, 1830, there were six members of the Church present.

I am happy to announce that two weeks ago, the membership of the Church reached 15 million. The Church continues to grow steadily and to change the lives of more and more people every year. It is spreading across the earth as our missionary force seeks out those who are searching for the truth.

It has scarcely been one year since I announced the lowering of the age of missionary service. Since that time the number of full-time missionaries serving has increased from 58,500 in October 2012 to 80,333 today. What a tremendous and inspiring response we have witnessed!

The holy scriptures contain no proclamation more relevant, no responsibility more binding, no instruction more direct than the injunction given by the resurrected Lord as He appeared in Galilee to the eleven disciples. Said He, “Go ye therefore, and teach all nations, baptizing them in the name of the Father, and of the Son, and of the Holy Ghost.” The Prophet Joseph Smith declared, “After all that has been said, the greatest and most important duty is to preach the Gospel.” Some of you here today will yet remember the words of President David O. McKay, who phrased the familiar “Every member a missionary!”

To their words I add my own. Now is the time for members and missionaries to come together, to work together, to labor in the Lord’s vineyard to bring souls unto Him. He has prepared the means for us to share the gospel in a multitude of ways, and He will assist us in our labors if we will act in faith to fulfill His work.

To help maintain our ever-increasing missionary force, I have asked our members in the past to contribute, as they are able, to their ward missionary fund or to the General Missionary Fund of the Church. The response to that request has been gratifying and has helped support thousands of missionaries whose circumstances do not allow them to support themselves. I thank you for your generous contributions. The need for help is ongoing, that we might continue to assist those whose desire to serve is great but who do not, by themselves, have the means to do so.

Now, brothers and sisters, we have come here to be instructed and inspired. Many messages, covering a variety of gospel topics, will be given during the next two days. Those men and women who will speak to you have sought heaven’s help concerning the messages they will give.

It is my prayer that we may be filled with the Spirit of the Lord as we listen and learn. In the name of our Savior, Jesus Christ, amen.

\newpage
\settalk{Be Strong and of a Good Courage}
\setspeaker{Thomas S. Monson}
\settalkdate{April 2014}
\talkheading{Be Strong and of a Good Courage}{Thomas S. Monson}

My beloved brethren, how good it is to be with you once again. I pray for heavenly help as I respond to the opportunity to address you.

Beyond this Conference Center are additional thousands assembled in chapels and in other settings throughout much of the world. A common thread binds all of us together, for we have been entrusted to bear the priesthood of God.

We are here upon the earth at a remarkable period in its history. Our opportunities are almost limitless, and yet we also face a multitude of challenges, some of them unique to our time.

We live in a world where moral values have, in great measure, been tossed aside, where sin is flagrantly on display, and where temptations to stray from the strait and narrow path surround us. We are faced with persistent pressures and insidious influences tearing down what is decent and attempting to substitute the shallow philosophies and practices of a secular society.

Because of these and other challenges, decisions are constantly before us which can determine our destiny. In order for us to make the correct decisions, courage is needed—the courage to say no when we should, the courage to say yes when that is appropriate, the courage to do the right thing because it is right.

Inasmuch as the trend in society today is rapidly moving away from the values and principles the Lord has given us, we will almost certainly be called upon to defend that which we believe. Will we have the courage to do so?

Said President J. Reuben Clark Jr., who for many years was a member of the First Presidency: “Not unknown are cases where [those] of presumed faith … have felt that, since by affirming their full faith they might call down upon themselves the ridicule of their unbelieving colleagues, they must either modify or explain away their faith, or destructively dilute it, or even pretend to cast it away. Such are hypocrites.” None of us would wish to wear such a label, and yet are we reluctant to declare our faith in some circumstances?

We can help ourselves in our desire to do what is right if we put ourselves in places and participate in activities where our thoughts are influenced for good and where the Spirit of the Lord will be comfortable.

I recall reading some time ago the counsel a father gave to his son when he went away to school: “If you ever find yourself where you shouldn’t ought to be, get out!” I offer to each of you the same advice: “If you ever find yourself where you shouldn’t ought to be, get out!”

The call for courage comes constantly to each of us. Every day of our lives courage is needed—not just for the momentous events but more often as we make decisions or respond to circumstances around us. Said Scottish poet and novelist Robert Louis Stevenson: “Everyday courage has few witnesses. But yours is no less noble because no drum beats for you and no crowds shout your name.”

Courage comes in many forms. Wrote the Christian author Charles Swindoll: “Courage is not limited to the battlefield … or bravely catching a thief in your house. The real tests of courage are much quieter. They are inner tests, like remaining faithful when no one’s looking, … like standing alone when you’re misunderstood.” I would add that this inner courage also includes doing the right thing even though we may be afraid, defending our beliefs at the risk of being ridiculed, and maintaining those beliefs even when threatened with a loss of friends or of social status. He who stands steadfastly for that which is right must risk becoming at times disapproved and unpopular.

While serving in the United States Navy in World War II, I learned of brave deeds, instances of valor, and examples of courage. One which I shall never forget was the quiet courage of an 18-year-old seaman—not of our faith—who was not too proud to pray. Of 250 men in the company, he was the only one who each night knelt down by the side of his bunk, at times amidst the jeers of bullies and the jests of unbelievers. With bowed head, he prayed to God. He never wavered. He never faltered. He had courage.

I listened not long ago to an example of one who surely seemed to lack this inner courage. A friend told of a spiritual and faith-promoting sacrament meeting she and her husband had attended in their ward. A young man who held the office of priest in the Aaronic Priesthood touched the hearts of the entire congregation as he spoke of gospel truths and of the joys of keeping the commandments. He bore a fervent, touching testimony as he stood at the pulpit, appearing clean and neat in his white shirt and tie.

Later that same day, as this woman and her husband drove out of their neighborhood, they saw this same young man who had so inspired them just a few hours earlier. Now, however, he presented a completely different picture as he walked down the sidewalk dressed in scruffy clothes—and smoking a cigarette. My friend and her husband were not only greatly disappointed and saddened, but they were also confused by how he could so convincingly seem to be one person in sacrament meeting and then so quickly seem to be someone else entirely.

Brethren, are you the same person wherever you are and whatever you are doing—the person our Heavenly Father wants you to be and the person you know you should be?

In an interview published in a national magazine, well-known American NCAA basketball player Jabari Parker, a member of the Church, was asked to share the best advice he had received from his father. Replied Jabari, “[My father] said, Just be the same person you are in the dark that you are in the light.” Important advice, brethren, for all of us.

Our scriptures are filled with examples of the type of courage needed by each of us today. The prophet Daniel exhibited supreme courage by standing up for that which he knew to be right and by demonstrating the courage to pray, though threatened with death were he to do so.

Courage characterized the life of Abinadi, as shown by his willingness to offer his life rather than to deny the truth.

Who can help but be inspired by the lives of the 2,000 stripling sons of Helaman, who taught and demonstrated the need for courage to follow the teachings of parents, to be chaste and pure?

Perhaps each of these scriptural accounts is crowned by the example of Moroni, who had the courage to persevere in righteousness to the very end.

Throughout his life, the Prophet Joseph Smith provided countless examples of courage. One of the most dramatic occurred as he and other brethren were chained together—imagine, chained together—and held in an unfinished cabin next to the courthouse in Richmond, Missouri. Parley P. Pratt, who was among those held captive, wrote of one particular night: “We had lain as if in sleep till the hour of midnight had passed, and our ears and hearts had been pained, while we had listened for hours to the obscene jests, the horrid oaths, the dreadful blasphemies and filthy language of our guards.”

Continued Elder Pratt:

“I had listened till I became so disgusted, shocked, horrified, and so filled with the spirit of indignant justice that I could scarcely refrain from rising upon my feet and rebuking the guards; but [I] had said nothing to Joseph, or any one else, although I lay next to him and knew he was awake. On a sudden he arose to his feet, and spoke in a voice of thunder, or as the roaring lion, uttering, as near as I can recollect, the following words:

“‘SILENCE. … In the name of Jesus Christ I rebuke you, and command you to be still; I will not live another minute and hear such language. Cease such talk, or you or I die THIS INSTANT!’”

Joseph “stood erect in terrible majesty,” as described by Elder Pratt. He was chained, without a weapon, and yet he was calm and dignified. He looked down upon the quailing guards, who were shrinking into a corner or crouching at his feet. These seemingly incorrigible men begged his pardon and remained quiet.

Not all acts of courage bring such spectacular or immediate results, and yet all of them do bring peace of mind and a knowledge that right and truth have been defended.

It is impossible to stand upright when one plants his roots in the shifting sands of popular opinion and approval. Needed is the courage of a Daniel, an Abinadi, a Moroni, or a Joseph Smith in order for us to hold strong and fast to that which we know is right. They had the courage to do not that which was easy but that which was right.

We will all face fear, experience ridicule, and meet opposition. Let us—all of us—have the courage to defy the consensus, the courage to stand for principle. Courage, not compromise, brings the smile of God’s approval. Courage becomes a living and an attractive virtue when it is regarded not only as a willingness to die manfully but also as the determination to live decently. As we move forward, striving to live as we should, we will surely receive help from the Lord and can find comfort in His words. I love His promise recorded in the book of Joshua:

“I will not fail thee, nor forsake thee. …

“… Be strong and of a good courage; be not afraid, neither be thou dismayed: for the Lord thy God is with thee whithersoever thou goest.”

My beloved brethren, with the courage of our convictions, may we declare, with the Apostle Paul, “I am not ashamed of the gospel of Christ.” And then, with that same courage, may we follow Paul’s counsel: “Be thou an example of the believers, in word, in conversation, in charity, in spirit, in faith, in purity.”

Catastrophic conflicts come and go, but the war waged for the souls of men continues without abatement. Like a clarion call comes the word of the Lord to you, to me, and to priesthood holders everywhere: “Wherefore, now let every man learn his duty, and to act in the office in which he is appointed, in all diligence.” Then we will be, as the Apostle Peter declared, even “a royal priesthood,” united in purpose and endowed with power from on high.

May each one leave here tonight with the determination and the courage to say, with Job of old, “While my breath is in me, … I will not remove mine integrity from me.” That this may be so is my humble prayer in the name of Jesus Christ, our Lord, amen.

\newpage
\settalk{Love—the Essence of the Gospel}
\setspeaker{Thomas S. Monson}
\settalkdate{April 2014}
\talkheading{Love—the Essence of the Gospel}{Thomas S. Monson}

My beloved brothers and sisters, when our Savior ministered among men, He was asked by the inquiring lawyer, “Master, which is the great commandment in the law?”

Matthew records that Jesus responded:

“Thou shalt love the Lord thy God with all thy heart, and with all thy soul, and with all thy mind.

“This is the first and great commandment.

“And the second is like unto it, Thou shalt love thy neighbour as thyself.”

Mark concludes the account with the Savior’s statement: “There is none other commandment greater than these.”

We cannot truly love God if we do not love our fellow travelers on this mortal journey. Likewise, we cannot fully love our fellowmen if we do not love God, the Father of us all. The Apostle John tells us, “This commandment have we from him, That he who loveth God love his brother also.” We are all spirit children of our Heavenly Father and, as such, are brothers and sisters. As we keep this truth in mind, loving all of God’s children will become easier.

Actually, love is the very essence of the gospel, and Jesus Christ is our Exemplar. His life was a legacy of love. The sick He healed; the downtrodden He lifted; the sinner He saved. At the end the angry mob took His life. And yet there rings from Golgotha’s hill the words: “Father, forgive them; for they know not what they do”—a crowning expression in mortality of compassion and love.

There are many attributes which are manifestations of love, such as kindness, patience, selflessness, understanding, and forgiveness. In all our associations, these and other such attributes will help make evident the love in our hearts.

Usually our love will be shown in our day-to-day interactions one with another. All important will be our ability to recognize someone’s need and then to respond. I have always cherished the sentiment expressed in the short poem:

\begin{verse}
I have wept in the night\\
For the shortness of sight\\
That to somebody’s need made me blind;\\
But I never have yet\\
Felt a tinge of regret\\
For being a little too kind.
\end{verse}

I recently was made aware of a touching example of loving kindness—one that had unforeseen results. The year was 1933, when because of the Great Depression, employment opportunities were scarce. The location was the eastern part of the United States. Arlene Biesecker had just graduated from high school. After a lengthy search for employment, she was finally able to obtain work at a clothing mill as a seamstress. The mill workers were paid only for each of the correctly completed pieces they sewed together daily. The more pieces they produced, the more they were paid.

One day shortly after starting at the mill, Arlene was faced with a procedure that had her confused and frustrated. She sat at her sewing machine trying to unpick her unsuccessful attempt to complete the piece on which she was working. There seemed to be no one to help her, for all of the other seamstresses were hurrying to complete as many pieces as they could. Arlene felt helpless and hopeless. Quietly, she began to cry.

Across from Arlene sat Bernice Rock. She was older and more experienced as a seamstress. Observing Arlene’s distress, Bernice left her own work and went to Arlene’s side, kindly giving her instruction and help. She stayed until Arlene gained confidence and was able to successfully complete the piece. Bernice then went back to her own machine, having missed the opportunity to complete as many pieces as she could have, had she not helped.

With this one act of loving kindness, Bernice and Arlene became lifelong friends. Each eventually married and had children. Sometime in the 1950s, Bernice, who was a member of the Church, gave Arlene and her family a copy of the Book of Mormon. In 1960, Arlene and her husband and children were baptized members of the Church. Later they were sealed in a holy temple of God.

As a result of the compassion shown by Bernice as she went out of her way to help one whom she didn’t know but who was in distress and needed assistance, countless individuals, both living and dead, now enjoy the saving ordinances of the gospel.

Every day of our lives we are given opportunities to show love and kindness to those around us. Said President Spencer W. Kimball: “We must remember that those mortals we meet in parking lots, offices, elevators, and elsewhere are that portion of mankind God has given us to love and to serve. It will do us little good to speak of the general brotherhood of mankind if we cannot regard those who are all around us as our brothers and sisters.”

Often our opportunities to show our love come unexpectedly. An example of such an opportunity appeared in a newspaper article in October 1981. So impressed was I with the love and compassion related therein that I have kept the clipping in my files for over 30 years.

The article indicates that an Alaska Airlines nonstop flight from Anchorage, Alaska, to Seattle, Washington—a flight carrying 150 passengers—was diverted to a remote Alaskan town in order to transport a gravely injured child. The two-year-old boy had severed an artery in his arm when he fell on a piece of glass while playing near his home. The town was 450 miles (725 km) south of Anchorage and was certainly not on the flight path. However, medics at the scene had sent out a frantic request for help, and so the flight was diverted to pick up the child and take him to Seattle so that he could be treated in a hospital.

When the flight touched down near the remote town, medics informed the pilot that the boy was bleeding so badly he could not survive the flight to Seattle. A decision was made to fly another 200 miles (320 km) out of the way to Juneau, Alaska, the nearest city with a hospital.

After transporting the boy to Juneau, the flight headed for Seattle, now hours behind schedule. Not one passenger complained, even though most of them would miss appointments and connecting flights. In fact, as the minutes and hours ticked by, they took up a collection, raising a considerable sum for the boy and his family.

As the flight was about to land in Seattle, the passengers broke into a cheer when the pilot announced that he had received word by radio that the boy was going to be all right.

To my mind come the words of the scripture: “Charity is the pure love of Christ, … and whoso is found possessed of it at the last day, it shall be well with him.”

Brothers and sisters, some of our greatest opportunities to demonstrate our love will be within the walls of our own homes. Love should be the very heart of family life, and yet sometimes it is not. There can be too much impatience, too much arguing, too many fights, too many tears. Lamented President Gordon B. Hinckley: “Why is it that the [ones] we love [most] become so frequently the targets of our harsh words? Why is it that [we] sometimes speak as if with daggers that cut to the quick?” The answers to these questions may be different for each of us, and yet the bottom line is that the reasons do not matter. If we would keep the commandment to love one another, we must treat each other with kindness and respect.

Of course there will be times when discipline needs to be meted out. Let us remember, however, the counsel found in the Doctrine and Covenants—namely, that when it is necessary for us to reprove another, we afterward show forth an increase of love.

I would hope that we would strive always to be considerate and to be sensitive to the thoughts and feelings and circumstances of those around us. Let us not demean or belittle. Rather, let us be compassionate and encouraging. We must be careful that we do not destroy another person’s confidence through careless words or actions.

Forgiveness should go hand in hand with love. In our families, as well as with our friends, there can be hurt feelings and disagreements. Again, it doesn’t really matter how small the issue was. It cannot and should not be left to canker, to fester, and ultimately to destroy. Blame keeps wounds open. Only forgiveness heals.

A lovely lady who has since passed away visited with me one day and unexpectedly recounted some regrets. She spoke of an incident which had taken place many years earlier and involved a neighboring farmer, once a good friend but with whom she and her husband had disagreed on multiple occasions. One day the farmer asked if he could take a shortcut across her property to reach his own acreage. At this point she paused in her narrative to me and, with a tremor in her voice, said, “Brother Monson, I didn’t let him cross our property then or ever but required him to take the long way around on foot to reach his property. I was wrong, and I regret it. He’s gone now, but oh, I wish I could say to him, ‘I’m so sorry.’ How I wish I had a second chance to be kind.”

As I listened to her, there came to my mind the doleful observation of John Greenleaf Whittier: “Of all sad words of tongue or pen, the saddest are these: ‘It might have been!’” Brothers and sisters, as we treat others with love and kind consideration, we will avoid such regrets.

Love is expressed in many recognizable ways: a smile, a wave, a kind comment, a compliment. Other expressions may be more subtle, such as showing interest in another’s activities, teaching a principle with kindness and patience, visiting one who is ill or homebound. These words and actions and many others can communicate love.

Dale Carnegie, a well-known American author and lecturer, believed that each person has within himself or herself the “power to increase the sum total of [the] world’s happiness … by giving a few words of sincere appreciation to someone who is lonely or discouraged.” Said he, “Perhaps you will forget tomorrow the kind words you say today, but the recipient may cherish them over a lifetime.”

May we begin now, this very day, to express love to all of God’s children, whether they be our family members, our friends, mere acquaintances, or total strangers. As we arise each morning, let us determine to respond with love and kindness to whatever might come our way.

Beyond comprehension, my brothers and sisters, is the love of God for us. Because of this love, He sent His Son, who loved us enough to give His life for us, that we might have eternal life. As we come to understand this incomparable gift, our hearts will be filled with love for our Eternal Father, for our Savior, and for all mankind. That such may be so is my earnest prayer in the sacred name of Jesus Christ, amen.

\newpage
\settalk{Until We Meet Again}
\setspeaker{Thomas S. Monson}
\settalkdate{April 2014}
\talkheading{Until We Meet Again}{Thomas S. Monson}

My brothers and sisters, what a wonderful conference this has been. We have been fed spiritually as we have listened to the inspired words of the men and women who have addressed us. The music has been superb, the messages have been prepared and delivered under the promptings of the Holy Spirit, and the prayers have drawn us nearer to heaven. We have been uplifted in every way as we have participated together.

I hope that we will take the time to read the conference messages when they become available on LDS.org within the next few days and when they are printed in coming issues of the Ensign and Liahona magazines, for they are deserving of our careful review and study.

I know you join with me in expressing our sincere gratitude to those brethren and sisters who were released during this conference. They have served well and have made significant contributions to the work of the Lord. Their dedication has been complete.

We have also sustained, by uplifted hands, brethren who have been called to new positions of responsibility. We welcome them and want them to know that we look forward to serving with them in the cause of the Master.

As we ponder the messages we have heard, may we resolve to do a little better than we have done in the past. May we be kind and loving to those who do not share our beliefs and our standards. The Savior brought to this earth a message of love and goodwill to all men and women. May we ever follow His example.

We face many serious challenges in the world today, but I assure you that our Heavenly Father is mindful of us. He will guide and bless us as we put our faith and trust in Him and will see us through whatever difficulties come our way.

May heaven’s blessings be with each of us. May our homes be filled with love and courtesy and with the Spirit of the Lord. May we constantly nourish our testimonies of the gospel, that they will be a protection for us against the buffetings of the adversary. May the Spirit we have felt during these past two days be and abide with us as we go about those things which occupy us each day, and may we ever be found doing the work of the Lord.

I bear testimony that this work is true, that our Savior lives, and that He guides and directs His Church here upon the earth. I leave with you my witness and my testimony that God our Eternal Father lives and loves us. He is indeed our Father, and He is personal and real. May we realize how close to us He is willing to come, how far He is willing to go to help us, and how much He loves us.

My brothers and sisters, may God bless you. May His promised peace be with you now and always.

I bid you farewell until we meet again in six months’ time, and I do so in the name of Jesus Christ, our Lord and Savior, amen.

\newpage
\settalk{Welcome to Conference}
\setspeaker{Thomas S. Monson}
\settalkdate{April 2014}
\talkheading{Welcome to Conference}{Thomas S. Monson}

My beloved brothers and sisters, how pleased I am to welcome you to this worldwide conference of The Church of Jesus Christ of Latter-day Saints. We are gathered together as a great family, more than 15 million strong, united in our faith and in our desire to listen to and learn from the messages which will be presented to us.

The past six months have gone by quickly as the work of the Church has moved forward unhindered. It was my privilege just over a month ago to dedicate the Gilbert Arizona Temple, a magnificent structure. The evening before the dedication, a cultural event was held at the nearby Discovery Park. Twelve thousand young people performed a 90-minute program. The dancing, the singing, and the musical performances were outstanding.

This area had been experiencing an especially dry season, and I believe many prayers had been sent heavenward over the preceding several weeks for much-needed rain. Unfortunately, it came just before the performance and stayed for the entire production! Despite the fact that the youth were soaked through with the rain and chilled from the cool temperature, we all felt the Spirit of the Lord. The theme of the program, “Live True to the Faith”—think about that: “Live True to the Faith”—was portrayed magnificently by smiling and enthusiastic young men and young women. Despite the cold and the rain, this was a faith-filled and inspiring experience these young people will ever treasure and will be relating to their children and grandchildren in the years to come.

The following day, the dedication of the Gilbert Arizona Temple took place. It became the 142nd operating temple in the Church. Unlike the evening before, the day was beautiful and filled with sunshine. The sessions were truly inspiring. Attending with me were President Henry B. Eyring, Elder and Sister Tad R. Callister, Elder and Sister William R. Walker, and Elder and Sister Kent F. Richards.

In May the Fort Lauderdale Florida Temple will be dedicated. Other temples are scheduled to be completed and dedicated later this year. In 2015 we anticipate completing and dedicating new temples in many parts of the world. This process will continue. When all the previously announced temples are completed, we will have 170 operating temples throughout the world.

Although we are currently concentrating our efforts on completing the previously announced temples and will not be announcing any new temples in the immediate future, we will continue the process of determining needs and of finding locations for temples yet to come. Announcements will then be made in future general conferences. We are a temple-building and a temple-attending people.

Now, brothers and sisters, we are anxious to listen to the messages which will be presented to us today and tomorrow. Those who will address us have sought heaven’s help and direction as they have prepared their messages.

May we—all of us, here and elsewhere—be filled with the Spirit of the Lord and be uplifted and inspired as we listen and learn. In the name of Jesus Christ, our Savior, amen.

\newpage
\settalk{Guided Safely Home}
\setspeaker{Thomas S. Monson}
\settalkdate{October 2014}
\talkheading{Guided Safely Home}{Thomas S. Monson}

Brethren, we are assembled as a mighty body of the priesthood, both here in the Conference Center and in locations throughout the world. I am honored yet humbled by the responsibility which is mine to address a few remarks to you. I pray for the Spirit of the Lord to attend me as I do so.

Seventy-five years ago, on February 14, 1939, in Hamburg, Germany, a public holiday was celebrated. Amid fervent speeches, cheering throngs, and the playing of patriotic anthems, the new battleship Bismarck was put to sea via the River Elbe. This, the most powerful vessel afloat, was a breathtaking spectacle of armor and machinery. Construction required more than 57,000 blueprints for the 380-millimeter, radar-controlled, double-gun turrets. The vessel featured 28,000 miles (45,000 km) of electrical circuits. It weighed over 35,000 tons, and armor plate provided maximum safety. Majestic in appearance, gigantic in size, awesome in firepower, the mighty colossus was considered unsinkable.

The Bismarck’s appointment with destiny came more than two years later, when on May 24, 1941, the two most powerful warships in the British Navy, the Prince of Wales and the Hood, engaged in battle the Bismarck and the German cruiser Prinz Eugen. Within five minutes the Bismarck had sent to the depths of the Atlantic the Hood and all but three men of a crew of over 1,400. The other British battleship, the Prince of Wales, had suffered heavy damage and turned away.

Over the next three days the Bismarck was engaged again and again by British warships and aircraft. In all, the British concentrated the strength of five battleships, two aircraft carriers, 11 cruisers, and 21 destroyers in an effort to find and to sink the mighty Bismarck.

During these battles, shell after shell inflicted only superficial damage on the Bismarck. Was it unsinkable after all? Then a torpedo scored a lucky hit, which jammed the Bismarck’s rudder. Repair efforts proved fruitless. With guns primed and the crews at ready, the Bismarck could only steer a slow circle. Just beyond reach was the powerful German air force. The Bismarck could not reach the safety of home port. Neither could provide the needed haven, for the Bismarck had lost the ability to steer a charted course. No rudder, no help, no port. The end drew near. British guns blazed as the German crew scuttled and sank the once seemingly indestructible vessel. The hungry waves of the Atlantic first lapped at the sides and then swallowed the pride of the German navy. The Bismarck was no more.

Like the Bismarck, each of us is a miracle of engineering. Our creation, however, was not limited by human genius. Man can devise the most complex machines but cannot give them life or bestow upon them the powers of reason and judgment. These are divine gifts, bestowed only by God.

Like the vital rudder of a ship, brethren, we have been provided a way to determine the direction we travel. The lighthouse of the Lord beckons to all as we sail the seas of life. Our purpose is to steer an undeviating course toward our desired goal—even the celestial kingdom of God. A man without a purpose is like a ship without a rudder, never likely to reach home port. To us comes the signal: chart your course, set your sail, position your rudder, and proceed.

As with the mighty Bismarck, so it is with man. The thrust of the turbines and the power of the propellers are useless without that sense of direction, that harnessing of the energy, that directing of the power provided by the rudder, hidden from view, relatively small in size but absolutely essential in function.

Our Father provided the sun, the moon, and the stars—heavenly galaxies to guide mariners who sail the lanes of the sea. To us, as we walk the pathway of life, He provides a clear map and points the way toward our desired destination. He cautions: beware the detours, the pitfalls, the traps. We cannot be deceived by those who would lead us astray, those clever pied pipers of sin beckoning here or there. Instead, we pause to pray; we listen to that still, small voice which speaks to the depths of our souls the Master’s gentle invitation, “Come, follow me.”

Yet there are those who do not hear, who will not obey, who prefer to walk a path of their own making. Too often they succumb to the temptations which surround all of us and which can appear so enticing.

As bearers of the priesthood, we have been placed on earth in troubled times. We live in a complex world with currents of conflict everywhere to be found. Political schemes ruin the stability of nations, despots grasp for power, and segments of society seem forever downtrodden, deprived of opportunity and left with a feeling of failure. The sophistries of men ring in our ears, and sin surrounds us.

Ours is the responsibility to be worthy of all the glorious blessings our Father in Heaven has in store for us. Wherever we go, our priesthood goes with us. Are we standing in holy places? Please, before you put yourself and your priesthood in jeopardy by venturing into places or participating in activities which are not worthy of you or of that priesthood, pause to consider the consequences.

We who have been ordained to the priesthood of God can make a difference. When we maintain our personal purity and honor our priesthood, we become righteous examples for others to follow. The Apostle Paul admonished, “Be thou an example of the believers, in word, in conversation, in charity, in spirit, in faith, in purity.” He also wrote that the followers of Christ should be “as lights in the world.” Providing an example of righteousness can help to illuminate an increasingly dark world.

Many of you will remember President N. Eldon Tanner, who served as a counselor to four Presidents of the Church. He provided an undeviating example of righteousness throughout his career in industry, during service in the government in Canada, and as an Apostle of Jesus Christ. He gave us this inspired counsel: “Nothing will bring greater joy and success than to live according to the teachings of the gospel. Be an example; be an influence for good.”

He continued: “Every one of us has been foreordained for some work as [God’s] chosen servant on whom he has seen fit to confer the priesthood and power to act in his name. Always remember that people are looking to you for leadership and you are influencing the lives of individuals either for good or for bad, which influence will be felt for generations to come.”

We are strengthened by the truth that the greatest force in the world today is the power of God as it works through man. To sail safely the seas of mortality, we need the guidance of that Eternal Mariner—even the great Jehovah. We reach out, we reach up to obtain heavenly help.

A well-known example of one who did not reach upward is that of Cain, son of Adam and Eve. Powerful in potential but weak of will, Cain permitted greed, envy, disobedience, and even murder to jam that personal rudder which would have guided him to safety and exaltation. The downward gaze replaced the upward look; Cain fell.

In another time and by a wicked king, a servant of God was tested. Aided by the inspiration of heaven, Daniel interpreted for the king the writing on the wall. Concerning the proffered rewards—even a royal robe, a necklace of gold, and political power—Daniel said, “Let thy gifts be to thyself, and give thy rewards to another.” Great riches and power had been offered to Daniel, rewards representing the things of the world and not of God. Daniel resisted and remained faithful.

Later, when Daniel worshipped God despite a decree declaring such to be forbidden, he was thrown into a den of lions. The biblical account tells us that the following morning, “Daniel was taken up out of the den, and no manner of hurt was found upon him, because he believed in … God.” In a time of critical need, Daniel’s determination to steer a steady course yielded divine protection and provided a sanctuary of safety. Such protection and safety can be ours as we also steer that steady course toward our eternal home.

The clock of history, like the sands of the hourglass, marks the passage of time. A new cast occupies the stage of life. The problems of our day loom ominously before us. Throughout the history of the world, Satan has worked tirelessly for the destruction of the followers of the Savior. If we succumb to his enticings, we—like the mighty Bismarck—will lose that rudder which can guide us to safety. Instead, surrounded by the sophistication of modern living, we look heavenward for that unfailing sense of direction, that we might chart and follow a wise and proper course. Our Heavenly Father will not leave our sincere petition unanswered. As we seek heavenly help, our rudder, unlike that of the Bismarck, will not fail.

As we venture forth on our individual voyages, may we sail safely the seas of life. May we have the courage of a Daniel, that we might remain true and faithful despite the sin and temptation which surround us. May our testimonies be as deep and as strong as that of Jacob, the brother of Nephi, who, when confronted by one who sought in every way possible to destroy his faith, declared, “I could not be shaken.”

With the rudder of faith guiding our passage, brethren, we too will find our way safely home—home to God, to dwell with Him eternally. That such may be so for each of us, I pray in the sacred name of Jesus Christ, our Savior and Redeemer, amen.

\newpage
\settalk{Ponder the Path of Thy Feet}
\setspeaker{Thomas S. Monson}
\settalkdate{October 2014}
\talkheading{Ponder the Path of Thy Feet}{Thomas S. Monson}

My beloved brothers and sisters, I am humbled as I stand before you this morning. I ask for your faith and prayers in my behalf as I share with you my message.

All of us commenced a wonderful and essential journey when we left the spirit world and entered this often-challenging stage called mortality. The primary purposes of our existence upon the earth are to obtain a body of flesh and bones, to gain experience that could come only through separation from our heavenly parents, and to see if we would keep the commandments. In the book of Abraham chapter 3 we read: “And we will prove them herewith, to see if they will do all things whatsoever the Lord their God shall command them.”

When we came to the earth, we brought with us that great gift from God—even our agency. In thousands of ways we are privileged to choose for ourselves. Here we learn from the hard taskmaster of experience. We discern between good and evil. We differentiate as to the bitter and the sweet. We learn that decisions determine destiny.

I am certain we left our Father with an overwhelming desire to return to Him, that we might gain the exaltation He planned for us and which we ourselves so much wanted. Although we are left to find and follow that path which will lead us back to our Father in Heaven, He did not send us here without direction and guidance. Rather, He has given us the tools we need, and He will assist us as we seek His help and strive to do all in our power to endure to the end and gain eternal life.

To help guide us we have the words of God and of His Son found in our holy scriptures. We have the counsel and teachings of God’s prophets. Of paramount importance, we have been provided with a perfect example to follow—even the example of our Lord and Savior, Jesus Christ—and we have been instructed to follow that example. Said the Savior Himself: “Come, follow me.” “The works which ye have seen me do that shall ye also do.” He posed the question, “What manner of men ought ye to be?” And then He answered, “Verily I say unto you, even as I am.” “He marked the path and led the way.”

As we look to Jesus as our Exemplar and as we follow in His footsteps, we can return safely to our Heavenly Father to live with Him forever. Said the prophet Nephi, “Unless a man shall endure to the end, in following the example of the Son of the living God, he cannot be saved.”

One woman, each time she related experiences she had during a visit to the Holy Land, would exclaim, “I walked where Jesus walked!”

She had been in the vicinity where Jesus lived and taught. Perhaps she stood on a rock on which He had once stood or looked at a mountain range He had once gazed upon. The experiences, in and of themselves, were thrilling to her; but physically walking where Jesus walked is less important than walking as He walked. Emulating His actions and following His example are far more important than trying to retrace the remnants of the trails He traversed in mortality.

When Jesus extended to a certain rich man the invitation, “Come, follow me,” He did not intend merely that the rich man follow Him up and down the hills and valleys of the countryside.

We need not walk by the shores of Galilee or among the Judean hills to walk where Jesus walked. All of us can walk the path He walked when, with His words ringing in our ears, His Spirit filling our hearts, and His teachings guiding our lives, we choose to follow Him as we journey through mortality. His example lights the way. Said He, “I am the way, the truth, and the life.”

As we examine the path Jesus walked, we will see that it took Him through many of the same challenges we ourselves will face in life.

For example, Jesus walked the path of disappointment. Although He experienced many disappointments, one of the most poignant was depicted in His lament over Jerusalem as He closed His public ministry. The children of Israel had rejected the safety of the protecting wing which He had offered them. As He looked out over the city soon to be abandoned to destruction, He was overcome by emotions of deep sorrow. In anguish He cried out, “O Jerusalem, Jerusalem, which killest the prophets, and stonest them that are sent unto thee; how often would I have gathered thy children together, as a hen doth gather her brood under her wings, and ye would not!”

Jesus walked the path of temptation. Lucifer, that evil one, amassing his greatest strength, his most inviting sophistry, tempted Him who had fasted for 40 days and 40 nights. Jesus did not succumb; rather, He resisted each temptation. His parting words: “Get thee hence, Satan.”

Jesus walked the path of pain. Consider Gethsemane, where He was “in an agony … and his sweat was as it were great drops of blood falling down to the ground.” And none can forget His suffering on the cruel cross.

Each of us will walk the path of disappointment, perhaps because of an opportunity lost, a power misused, a loved one’s choices, or a choice we ourselves make. The path of temptation too will be the path of each. We read in the 29th section of the Doctrine and Covenants: “And it must needs be that the devil should tempt the children of men, or they could not be agents unto themselves.”

Likewise shall we walk the path of pain. We, as servants, can expect no more than the Master, who left mortality only after great pain and suffering.

While we will find on our path bitter sorrow, we can also find great happiness.

We, with Jesus, can walk the path of obedience. It will not always be easy, but let our watchword be the heritage bequeathed us by Samuel: “Behold, to obey is better than sacrifice, and to hearken than the fat of rams.” Let us remember that the end result of disobedience is captivity and death, while the reward for obedience is liberty and eternal life.

We, like Jesus, can walk the path of service. As a glowing searchlight of goodness is the life of Jesus as He ministered among men. He brought strength to the limbs of the cripple, sight to the eyes of the blind, hearing to the ears of the deaf.

Jesus walked the path of prayer. He taught us how to pray by giving us the beautiful prayer we know as the Lord’s Prayer. And who can forget His prayer in Gethsemane, “Not my will, but thine, be done”?

Other instructions given to us by the Savior are at our fingertips, found in the holy scriptures. In His Sermon on the Mount, He tells us to be merciful, to be humble, to be righteous, to be pure in heart, to be peacemakers. He instructs us to stand up bravely for our beliefs, even when we are ridiculed and persecuted. He asks us to let our lights shine so that others may see them and may desire to glorify our Father in Heaven. He teaches us to be morally clean in both our thoughts and our actions. He tells us it is far more important to lay up treasures in heaven than on earth.

His parables teach with power and authority. With the account of the good Samaritan, He teaches us to love and to serve our neighbors. In His parable of the talents, He teaches us to improve ourselves and to strive for perfection. With the parable of the lost sheep, He instructs us to go to the rescue of those who have left the path and have lost their way.

As we strive to place Christ at the center of our lives by learning His words, by following His teachings, and by walking in His path, He has promised to share with us the eternal life that He died to gain. There is no higher end than this, that we should choose to accept His discipline and become His disciples and do His work throughout our lives. Nothing else, no other choice we make, can make of us what He can.

As I think of those who have truly tried to follow the example of the Savior and who have walked in His path, there comes readily to my mind the names of Gustav and Margarete Wacker—two of the most Christlike individuals I have ever known. They were native Germans who had immigrated to eastern Canada, and I met them when I served as a mission president there. Brother Wacker earned his living as a barber. Though their means were limited, they shared all they had. They were not blessed with children, but they nurtured all who entered their home. Men and women of learning and sophistication sought out these humble, unlettered servants of God and counted themselves fortunate if they could spend an hour in their presence.

Their appearance was ordinary, their English halting and somewhat difficult to understand, their home unpretentious. They didn’t own a car or a television, nor did they do any of the things to which the world usually pays attention. Yet the faithful beat a path to their door in order to partake of the spirit that was there. Their home was a heaven on earth, and the spirit they radiated was of pure peace and goodness.

We too can have that spirit and can share it with the world as we walk the path of our Savior and follow His perfect example.

We read in Proverbs the admonition, “Ponder the path of thy feet.” As we do, we will have the faith, even the desire, to walk the path which Jesus walked. We will have no doubt that we are on a path which our Father would have us follow. The Savior’s example provides a framework for everything that we do, and His words provide an unfailing guide. His path will take us safely home. May this be our blessing, I pray in the name of Jesus Christ, whom I love, whom I serve, and of whom I testify, amen.

\newpage
\settalk{Until We Meet Again}
\setspeaker{Thomas S. Monson}
\settalkdate{October 2014}
\talkheading{Until We Meet Again}{Thomas S. Monson}

My brothers and sisters, we have experienced two glorious days of inspired messages. Our hearts have been touched and our faith strengthened as we have partaken of the spirit which has been present during these conference sessions. As we conclude, we thank our Heavenly Father for His many blessings to us.

We have been lifted and inspired by the beautiful music that has been provided during the sessions. The prayers which have been given have drawn us nearer to heaven.

May I express the heartfelt thanks of the entire Church to our Brethren who have been released at this conference. We will miss them. Their contributions to the work of the Lord have been enormous and will be felt throughout generations to come.

May we return to our homes with a resolve in our hearts to be a little better than we have been in the past. May we be a little kinder and more thoughtful. May we reach out in helpfulness, not only to our fellow members but also to those who are not of our faith. As we associate with them, may we show our respect for them.

There are those who struggle every day with challenges. Let us extend to them our concern, as well as a helping hand. As we care for each other, we will be blessed.

May we remember the elderly and those who are homebound. As we take time to visit them, they will know that they are loved and valued. May we follow the mandate to “succor the weak, lift up the hands which hang down, and strengthen the feeble knees.”

May we be people of honesty and integrity, trying to do the right thing at all times and in all circumstances. May we be faithful followers of Christ, examples of righteousness, thus becoming “lights in the world.”

My brothers and sisters, I thank you for your prayers in my behalf. They strengthen me and lift me as I strive with all my heart and strength to do God’s will and to serve Him and to serve you.

As we leave this conference, I invoke the blessings of heaven upon each of you. May you who are away from your homes return to them safely and find all in order. May we all ponder the truths we have heard, and may they help us to become even more valiant disciples than we were when this conference began.

Until we meet again in six months’ time, I ask the Lord’s blessings to be upon you and, indeed, upon all of us, and I do so in His holy name—even Jesus Christ, our Lord and Savior—amen.

\newpage
\settalk{Welcome to Conference}
\setspeaker{Thomas S. Monson}
\settalkdate{October 2014}
\talkheading{Welcome to Conference}{Thomas S. Monson}

My brothers and sisters, how pleased I am to welcome you to this great world conference. We are gathered together in locations around the world to listen to and learn from the brethren and sisters whom we have sustained as General Authorities and general officers of the Church. They have sought heaven’s help concerning the messages which they will present, and they have felt inspiration regarding what will be said.

This conference marks the 90-year anniversary of radio broadcasts of general conference. During the October conference of 1924, the sessions were broadcast on the radio for the first time through Church-owned KSL. This conference also marks the 65-year anniversary of television broadcasts of conference. At the general conference held in October 1949, the sessions were first televised throughout the Salt Lake area over KSL television.

We acknowledge the blessings of modern media in allowing millions of members of the Church to watch or listen to general conference. The sessions of this weekend are being broadcast via television, radio, cable, satellite transmission, and the Internet, including on mobile devices.

During the past six months since we last met, one new temple has been dedicated and one rededicated. In May, President Dieter F. Uchtdorf dedicated the Fort Lauderdale Florida Temple. A wonderful youth cultural celebration was presented the day prior to the dedication. The following day, on Sunday, May 4, the temple was dedicated in three sessions.

Just two weeks ago it was my privilege to rededicate the Ogden Utah Temple, originally dedicated in 1972 by President Joseph Fielding Smith. A grand cultural celebration took place the day before the rededication, with so many youth participating that two separate performances were presented, with a different cast for each. In all, 16,000 youth participated. The rededication services took place the following day, with many of the Brethren participating, along with the auxiliary leaders and the temple president, his counselors, and their wives.

Our temple building continues in earnest. Next month the new Phoenix Arizona Temple will be dedicated, and next year, in 2015, we anticipate dedicating or rededicating at least five temples, with more possible, depending on completion.

As I mentioned in April, when all the previously announced temples are constructed and dedicated, we will have 170 operating temples throughout the world. Because we are concentrating our efforts on completing temples which were previously announced, we are not at the present time announcing any new temples. However, in the future, as we identify needs and locate properties, announcements of additional temples will be made.

The Church continues to grow. We are now more than 15 million strong and increasing in numbers. Our missionary efforts are going forward unhindered. We have over 88,000 missionaries serving, sharing the gospel message the world over. We reaffirm that missionary work is a priesthood duty, and we encourage all worthy and able young men to serve. We are very grateful for the young women who also serve. They make a significant contribution, although they are not under the same mandate to serve as are the young men.

Now I invite you to give your attention to the brethren and sisters who will participate today and tomorrow in our conference sessions. All who have been asked to speak feel a great responsibility in doing so. As we listen, may our hearts be touched and our faith increased, I humbly pray in the name of Jesus Christ, amen.

\newpage
\settalk{Blessings of the Temple}
\setspeaker{Thomas S. Monson}
\settalkdate{April 2015}
\talkheading{Blessings of the Temple}{Thomas S. Monson}

My beloved brothers and sisters, how grateful I am to be with you this beautiful Easter morning when our thoughts turn to the Savior of the world. I extend my love and greetings to each of you and pray that our Heavenly Father will inspire my words.

This conference marks seven years since I was sustained as President of the Church. They have been busy years, filled not only with a few challenges but also with countless blessings. Among the most enjoyable and sacred of these blessings has been my opportunity to dedicate and rededicate temples.

Most recently, this past November it was my privilege to dedicate the beautiful new Phoenix Arizona Temple. I was joined by President Dieter F. Uchtdorf, Elder Dallin H. Oaks, Elder Richard J. Maynes, Elder Lynn G. Robbins, and Elder Kent F. Richards. On the evening prior to the dedication, a marvelous cultural celebration was held where over 4,000 of our youth from the temple district performed beautifully. The following day the temple was dedicated in three sacred and inspiring sessions.

The building of temples is a very clear indication of the growth of the Church. We currently have 144 temples in operation worldwide, with 5 being renovated and 13 more under construction. In addition, 13 temples which were previously announced are in various stages of preparation before construction begins. This year we anticipate rededicating 2 temples and dedicating 5 new temples which are scheduled for completion.

For the past two years, as we have concentrated our efforts on completing previously announced temples, we have held in abeyance plans for any additional temples. This morning, however, I am very pleased to announce three new temples which will be built in the following locations: Abidjan, Ivory Coast; Port-au-Prince, Haiti; and Bangkok, Thailand. What marvelous blessings are in store for our faithful members in these areas and, indeed, wherever temples are located throughout the world.

The process of determining needs and finding locations for additional temples is ongoing, for we desire that as many members as possible have an opportunity to attend the temple without great sacrifices of time and resources. As we have done in the past, we will keep you informed as decisions are made in this regard.

As I think of temples, my thoughts turn to the many blessings we receive therein. As we enter through the doors of the temple, we leave behind us the distractions and confusion of the world. Inside this sacred sanctuary, we find beauty and order. There is rest for our souls and a respite from the cares of our lives.

As we attend the temple, there can come to us a dimension of spirituality and a feeling of peace which will transcend any other feeling which could come into the human heart. We will grasp the true meaning of the words of the Savior when He said: “Peace I leave with you, my peace I give unto you. … Let not your heart be troubled, neither let it be afraid.”

Such peace can permeate any heart—hearts that are troubled, hearts that are burdened down with grief, hearts that feel confusion, hearts that plead for help.

I recently learned firsthand of a young man who attended the temple with a heart pleading for help. Many months earlier he had received his call to serve in a mission in South America. However, his visa was delayed for such a lengthy period that he was reassigned to a mission in the United States. Although disappointed that he could not serve in the area of his original call, he nonetheless worked hard in his new assignment, determined to serve to the best of his ability. He became discouraged, however, because of negative experiences he had with missionaries who seemed to him to be more interested in having a good time than in sharing the gospel.

A few short months later this young man suffered a very serious health challenge which left him partially paralyzed, and so he was sent home on a medical leave.

Some months later the young man had healed completely, and his paralysis had disappeared. He was informed that he would once again be able to serve as a missionary, a blessing for which he had prayed daily. The only disappointing news was that he would return to the same mission which he had left, where he felt the behaviors and attitudes of some missionaries were less than they should be.

He had come to the temple to seek comfort and a confirmation that he could have a good experience as a missionary. His parents also had prayed that this temple visit would provide the help their son needed.

As the young man entered the celestial room following the session, he sat in a chair and began to pray for guidance from his Heavenly Father.

Another who entered the celestial room shortly afterward was a young man whose name is Landon. As he walked into the room, his gaze was immediately drawn to the young man sitting on the chair, eyes closed and obviously praying. Landon received an unmistakable prompting that he should speak with the young man. Hesitant to interrupt, however, he decided to wait. After several minutes had gone by, the young man was still praying. Landon knew he could no longer postpone the prompting. He approached the young man and gently touched his shoulder. The young man opened his eyes, startled that he had been disturbed. Landon said quietly, “I have felt impressed that I need to talk with you, although I am not certain why.”

As they began to converse, the young man poured out his heart to Landon, explaining his circumstances and ending with his desire to receive some comfort and encouragement concerning his mission. Landon, who had returned from a successful mission just a year earlier, told of his own mission experiences, the challenges and concerns he had faced, the manner in which he had turned to the Lord for help, and the blessings he had received. His words were comforting and reassuring, and his enthusiasm for his mission was contagious. Eventually, as the young man’s fears subsided, a feeling of peace came to him. He felt deep gratitude as he realized his prayer had been answered.

The two young men prayed together, and then Landon prepared to leave, happy that he had listened to the inspiration which had come to him. As he stood to go, the young man asked Landon, “Where did you serve your mission?” To this point, neither of them had mentioned to the other the name of the mission in which he had served. When Landon replied with the name of his mission, tears welled up in the eyes of the young man. Landon had served in the very mission to which the young man would be returning!

In a recent letter to me, Landon shared with me the young man’s parting words to him: “I had faith Heavenly Father would bless me, but I never could have imagined that He would send someone to help me who had served in my own mission. I know now that all will be well.” The humble prayer of a sincere heart had been heard and answered.

My brothers and sisters, in our lives we will have temptations; we will have trials and challenges. As we go to the temple, as we remember the covenants we make there, we will be better able to overcome those temptations and to bear our trials. In the temple we can find peace.

The blessings of the temple are priceless. One for which I am grateful every day of my life is that which my beloved wife, Frances, and I received as we knelt at a sacred altar and made covenants binding us together for all eternity. There is no blessing more precious to me than the peace and comfort I receive from the knowledge I have that she and I will be together again.

May our Heavenly Father bless us that we may have the spirit of temple worship, that we may be obedient to His commandments, and that we may follow carefully the steps of our Lord and Savior, Jesus Christ. I testify that He is our Redeemer. He is the Son of God. He it is who came forth from the grave that first Easter morning, bringing with Him the gift of everlasting life for all of God’s children. On this beautiful day, as we celebrate that momentous event, may we offer prayers of gratitude for His great and marvelous gifts to us. That this may be so, I pray humbly in His holy name, amen.

\newpage
\settalk{The Priesthood—a Sacred Gift}
\setspeaker{Thomas S. Monson}
\settalkdate{April 2015}
\talkheading{The Priesthood—a Sacred Gift}{Thomas S. Monson}

One of my most vivid memories is attending priesthood meeting as a newly ordained deacon and singing during the opening hymn, “Come, all ye sons of God who have received the priesthood.” Tonight, to all assembled here in the Conference Center and, indeed, throughout the world, I echo the spirit of that special hymn and say to you: Come, all ye sons of God who have received the priesthood, let us consider our callings; let us reflect on our responsibilities; let us determine our duty; and let us follow Jesus Christ, our Lord. While we may differ in age, in custom, or in nationality, we are united as one in our priesthood callings.

To each of us, the restoration of the Aaronic Priesthood to Oliver Cowdery and Joseph Smith by John the Baptist is most significant. Likewise, the restoration of the Melchizedek Priesthood to Joseph and Oliver by Peter, James, and John is a cherished event.

Let us take most seriously the callings, the responsibilities, and the duties which come with the priesthood we hold.

I felt a great responsibility when I was called to be secretary of my deacons quorum. I prepared most conscientiously the records I kept, for I wanted to do the very best I knew how to do in that calling. I took great pride in my work. Doing all I can, to the very best of my ability, has been my goal in any position I have ever held.

I hope each young man who has been ordained to the Aaronic Priesthood is given a spiritual awareness of the sacredness of his ordained calling, as well as opportunities to magnify that calling. I received such an opportunity as a deacon when the bishopric asked that I take the sacrament to a shut-in who lived about a mile from our chapel. That special Sunday morning, as I knocked on Brother Wright’s door and heard his feeble voice call, “Come in,” I entered not only his humble cottage but also a room filled with the Spirit of the Lord. I approached Brother Wright’s bedside and carefully placed a piece of the bread to his lips. I then held the cup of water, that he might drink. As I departed, I saw tears in his eyes as he said, “God bless you, my boy.” And God did bless me—with an appreciation for the sacred emblems of the sacrament and for the priesthood which I held.

No deacon, teacher, or priest from our ward will ever forget the memorable visits we made to Clarkston, Utah, to the gravesite of Martin Harris, one of the Three Witnesses of the Book of Mormon. As we surrounded the tall granite shaft which marks his grave, and as one of the quorum leaders read to us those penetrating words from “The Testimony of Three Witnesses,” found at the beginning of the Book of Mormon, we developed a love for that sacred record and for the truths found therein.

During those years our objective was to become as the sons of Mosiah. Of them it was said:

“They had waxed strong in the knowledge of the truth; for they were men of a sound understanding and they had searched the scriptures diligently, that they might know the word of God.

“But this is not all; they had given themselves to much prayer, and fasting; therefore they had the spirit of prophecy, and the spirit of revelation, and when they taught, they taught with power and authority of God.”

I cannot think of a more worthy goal for a young man to have than to be described as were the valiant and righteous sons of Mosiah.

As I approached my 18th birthday and prepared to enter the mandatory military service required of young men during World War II, I was recommended to receive the Melchizedek Priesthood, but first I needed to telephone my stake president, Paul C. Child, for an interview. He was one who loved and understood the holy scriptures, and it was his intent that all others should similarly love and understand them. Having heard from some of my friends of his rather detailed and searching interviews, I desired minimum exposure of my scriptural knowledge; therefore, when I called him I suggested we meet the following Sunday at a time I knew was just an hour before his sacrament meeting time.

His response: “Oh, Brother Monson, that would not provide us sufficient time to peruse the scriptures.” He then suggested a time three hours before his sacrament meeting, and he instructed me to bring with me my personally marked and referenced set of scriptures.

When I arrived at his home on Sunday, I was greeted warmly, and then the interview began. President Child said, “Brother Monson, you hold the Aaronic Priesthood. Have you ever had angels minister to you?” I replied that I had not. When he asked if I knew I was entitled to such, I again replied that I had not known.

He instructed, “Brother Monson, repeat from memory the 13th section of the Doctrine and Covenants.”

I began, “‘Upon you my fellow servants, in the name of Messiah I confer the Priesthood of Aaron, which holds the keys of the ministering of angels—’”

“Stop,” President Child directed. Then, in a calm, kindly tone, he counseled, “Brother Monson, never forget that as a holder of the Aaronic Priesthood you are entitled to the ministering of angels.”

It was almost as though an angel were in the room that day. I have never forgotten the interview. I yet feel the spirit of that solemn occasion as we together read of the responsibilities, the duties, and the blessings of the Aaronic Priesthood and the Melchizedek Priesthood—blessings which come not only to us but also to our families and to others we will have the privilege to serve.

I was ordained an elder, and on the day of my departure for active duty with the navy, a member of my ward bishopric joined my family and friends at the train station to bid me farewell. Just before train time, he placed in my hand a small volume titled Missionary Handbook. I laughed and commented that I wasn’t going on a mission.

He answered, “Take it anyway. It may come in handy.”

It did. I needed a hard, rectangular object to place in the bottom of my seabag so that my clothing would stay more firm and would thus be less wrinkled. The Missionary Handbook was just what I needed, and it served well in my seabag for 12 weeks.

The night before our Christmas leave, our thoughts were of home. The barracks were quiet, but then the silence was broken by my buddy in the adjoining bunk—a Mormon boy, Leland Merrill—who began to moan in pain. I inquired concerning the reason, and he said he felt really sick. He did not want to go to the base dispensary, for he knew that doing such would prevent his going home the following day.

He seemed to grow worse as the hours passed. Finally, knowing that I was an elder, he asked me to give him a priesthood blessing.

I had never before given a priesthood blessing, I had never received a blessing, and I had never witnessed a blessing being given. As I prayed silently for help, I remembered the Missionary Handbook in the bottom of my seabag. I quickly emptied the bag and took the book to the night-light. There I read how one blesses the sick. With many curious sailors looking on, I proceeded with the blessing. Before I could put everything back into my bag, Leland Merrill was sleeping like a child. He awakened the following morning feeling fine. The gratitude each of us felt for the power of the priesthood was immense.

The years have brought me more opportunities to provide blessings to those in need than I could possibly count. Each opportunity has found me deeply grateful that God has entrusted to me this sacred gift. I revere the priesthood. I have witnessed its power time and time again. I have seen its strength. I have marveled at the miracles it has wrought.

Brethren, each of us has been entrusted with one of the most precious gifts ever bestowed upon mankind. As we honor our priesthood and live our lives so that we are at all times worthy, the blessings of the priesthood will flow through us. I love the words found in the Doctrine and Covenants, section 121, verse 45, which tell us what we must do to be worthy: “Let thy bowels … be full of charity towards all men, and to the household of faith, and let virtue garnish thy thoughts unceasingly; then shall thy confidence wax strong in the presence of God; and the doctrine of the priesthood shall distil upon thy soul as the dews from heaven.”

As bearers of the priesthood of God, we are engaged in the work of the Lord Jesus Christ. We have answered His call; we are on His errand. Let us learn of Him. Let us follow in His footsteps. Let us live by His precepts. By so doing, we will be prepared for any service He calls us to perform. This is His work. This is His Church. Indeed, He is our captain, the King of Glory, even the Son of God. I testify that He lives and bear this witness in His holy name, the name of Jesus Christ, amen.

\newpage
\settalk{Be an Example and a Light}
\setspeaker{Thomas S. Monson}
\settalkdate{October 2015}
\talkheading{Be an Example and a Light}{Thomas S. Monson}

Brothers and sisters, how good it is to be with you once again. As you know, since we were together in April, we have been saddened by the loss of three of our beloved Apostles: President Boyd K. Packer, Elder L. Tom Perry, and Elder Richard G. Scott. They have returned to their heavenly home. We miss them. How grateful we are for their examples of Christlike love and for the inspired teachings they have left to all of us.

We extend a heartfelt welcome to our newest Apostles, Elder Ronald A. Rasband, Elder Gary E. Stevenson, and Elder Dale G. Renlund. These are men dedicated to the work of the Lord. They are well qualified to fill the important positions to which they have been called.

Recently, as I have been reading and pondering the scriptures, two passages in particular have stayed with me. Both are familiar to us. The first is from the Sermon on the Mount: “Let your light so shine before men, that they may see your good works, and glorify your Father which is in heaven.” The second scripture is one which came to my mind as I pondered the meaning of the first. It is from the Apostle Paul’s Epistle to Timothy: “Be thou an example of the believers, in word, in conversation, in charity, in spirit, in faith, in purity.”

I believe the second scripture explains, in great part, how we can accomplish the first. We become examples of the believers by living the gospel of Jesus Christ in word, in conversation, in charity, in spirit, in faith, and in purity. As we do so, our lights will shine for others to see.

Each of us came to earth having been given the Light of Christ. As we follow the example of the Savior and live as He lived and as He taught, that light will burn within us and will light the way for others.

The Apostle Paul lists six attributes of a believer, attributes that will allow our lights to shine. Let us look at each one.

I mention the first two attributes together—being an example in word and in conversation. The words we use can lift and inspire, or they can harm and demean. In the world today there is a profusion of profanity with which we seem to be surrounded at nearly every turn. It is difficult to avoid hearing the names of Deity being used casually and thoughtlessly. Coarse comments seem to have become a staple of television, movies, books, and music. Bandied about are slanderous remarks and angry rhetoric. Let us speak to others with love and respect, ever keeping our language clean and avoiding words or comments that would wound or offend. May we follow the example of the Savior, who spoke with tolerance and kindness throughout His ministry.

The next attribute mentioned by Paul is charity, which has been defined as “the pure love of Christ.” I am confident there are within our sphere of influence those who are lonely, those who are ill, and those who feel discouraged. Ours is the opportunity to help them and to lift their spirits. The Savior brought hope to the hopeless and strength to the weak. He healed the sick; He caused the lame to walk, the blind to see, the deaf to hear. He even raised the dead to life. Throughout His ministry He reached out in charity to any in need. As we emulate His example, we will bless lives, including our own.

Next, we are to be an example in spirit. To me that means we strive to have in our lives kindness, gratitude, forgiveness, and goodwill. These qualities will provide for us a spirit which will touch the lives of those around us. It has been my opportunity through the years to associate with countless individuals who possess such a spirit. We experience a special feeling when we are with them, a feeling that makes us want to associate with them and to follow their example. They radiate the Light of Christ and help us feel His love for us.

To illustrate that the light which comes from a pure and loving spirit is recognized by others, I share with you an experience of many years ago.

At that time, leaders of the Church met with officials in Jerusalem to work out a lease agreement for land on which the Church’s Jerusalem Center would be built. In order to obtain the permissions needed, the Church had to agree that no proselyting would be undertaken by our members who would occupy the center. After that agreement had been made, one of the Israeli officials, who was well acquainted with the Church and its members, remarked that he knew the Church would honor the no-proselyting agreement. “But,” he said, referring to the students who would attend there, “what are [we] going to do about the light that is in their eyes?” May that special light ever shine within us, that it might be recognized and appreciated by others.

To be an example of faith means that we trust in the Lord and in His word. It means that we possess and that we nourish the beliefs that will guide our thoughts and our actions. Our faith in the Lord Jesus Christ and in our Heavenly Father will influence all that we do. Amidst the confusion of our age, the conflicts of conscience, and the turmoil of daily living, an abiding faith becomes an anchor to our lives. Remember that faith and doubt cannot exist in the same mind at the same time, for one will dispel the other. I reiterate what we have been told repeatedly—that in order to gain and to keep the faith we need, it is essential that we read and study and ponder the scriptures. Communication with our Heavenly Father through prayer is vital. We cannot afford to neglect these things, for the adversary and his hosts are relentlessly seeking for a chink in our armor, a lapse in our faithfulness. Said the Lord, “Search diligently, pray always, and be believing, and all things shall work together for your good.”

Finally, we are to be pure, which means that we are clean in body, mind, and spirit. We know that our body is a temple, to be treated with reverence and respect. Our minds should be filled with uplifting and ennobling thoughts and kept free from those things which will pollute. In order to have the Holy Ghost as our constant companion, we must be worthy. Brothers and sisters, purity will bring us peace of mind and will qualify us to receive the Savior’s promises. Said He, “Blessed are the pure in heart: for they shall see God.”

As we prove to be examples in word, in conversation, in charity, in spirit, in faith, and in purity, we will qualify to be lights to the world.

May I say to all of you, and particularly to you young people, that as the world moves further and further away from the principles and guidelines given to us by a loving Heavenly Father, we will stand out from the crowd because we are different. We will stand out because we dress modestly. We will be different because we do not use profanity and because we do not partake of substances which are harmful to our bodies. We will be different because we avoid off-color humor and degrading remarks. We will be different as we decide not to fill our minds with media choices that are base and demeaning and that will remove the Spirit from our homes and our lives. We will certainly stand out as we make choices regarding morality—choices which adhere to gospel principles and standards. Those things which make us different from most of the world also provide us with that light and that spirit which will shine in an increasingly dark world.

It is often difficult to be different and to stand alone in a crowd. It is natural to fear what others might think or say. Comforting are the words of the psalm: “The Lord is my light and my salvation; whom shall I fear? the Lord is the strength of my life; of whom shall I be afraid?” As we make Christ the center of our lives, our fears will be replaced by the courage of our convictions.

Life is perfect for none of us, and at times the challenges and difficulties we face may become overwhelming, causing our light to dim. However, with help from our Heavenly Father, coupled with support from others, we can regain that light which will illuminate our own path once again and provide the light others may need.

To illustrate, I share with you the touching words of a favorite poem I first read many years ago:

\begin{verse}
I met a stranger in the night\\
Whose lamp had ceased to shine.\\
I paused and let him light\\
His lamp from mine.
\end{verse}

My brothers and sisters, our opportunities to shine surround us each day, in whatever circumstance we find ourselves. As we follow the example of the Savior, ours will be the opportunity to be a light in the lives of others, whether they be our own family members and friends, our co-workers, mere acquaintances, or total strangers.

To each of you, I say that you are a son or daughter of our Heavenly Father. You have come from His presence to live on this earth for a season, to reflect the Savior’s love and teachings, and to bravely let your light shine for all to see. When that season on earth has ended, if you have done your part, yours will be the glorious blessing of returning to live with Him forever.

How reassuring are the Savior’s words: “I am the light of the world: he that followeth me shall not walk in darkness, but shall have the light of life.” Of Him I testify. He is our Savior and Redeemer, our Advocate with the Father. He is our Exemplar and our strength. He is “the light which shineth in darkness.” That each of us within the sound of my voice may pledge to follow Him, thus becoming a shining light to the world, is my prayer in His holy name, even Jesus Christ the Lord, amen.

\newpage
\settalk{Keep the Commandments}
\setspeaker{Thomas S. Monson}
\settalkdate{October 2015}
\talkheading{Keep the Commandments}{Thomas S. Monson}

My beloved brethren, how good it is to be with you once again. We have been inspired this evening by the words which we have heard. I pray that I too will be guided in what I say.

My message to you tonight is straightforward. It is this: keep the commandments.

God’s commandments are not given to frustrate us or to become obstacles to our happiness. Just the opposite is true. He who created us and who loves us perfectly knows just how we need to live our lives in order to obtain the greatest happiness possible. He has provided us with guidelines which, if we follow them, will see us safely through this often treacherous mortal journey. We remember the words of the familiar hymn: “Keep the commandments! In this there is safety; in this there is peace.”

Our Heavenly Father loves us enough to say: Thou shalt not lie; thou shalt not steal; thou shalt not commit adultery; thou shalt love thy neighbor as thyself; and so on. We know the commandments. He understands that when we keep the commandments, our lives will be happier, more fulfilling, and less complicated. Our challenges and problems will be easier to bear, and we will receive His promised blessings. But while He gives us laws and commandments, He also allows us to choose whether to accept them or to reject them. Our decisions in this regard will determine our destiny.

I am confident that each of us has as his ultimate goal life everlasting in the presence of our Heavenly Father and His Son, Jesus Christ. It is imperative, therefore, for us to make choices throughout our lives that will lead us to this great goal. We know, however, that the adversary is committed to our failure. He and his hosts are relentless in their efforts to thwart our righteous desires. They represent a grave and constant threat to our eternal salvation unless we are also relentless in our determination and efforts to achieve our goal. The Apostle Peter warns us, “Be vigilant; because your adversary the devil, as a roaring lion, walketh about, seeking whom he may devour.”

Although there is no time in our lives when we are exempt from temptation, you young men are at an age when you may be particularly vulnerable. Teenage years are often years of insecurity, of feeling as though you don’t measure up, of trying to find your place with your peers, and of trying to fit in. You may be tempted to lower your standards and to follow the crowd in order to be accepted by those you desire to have as friends. Please be strong, and be alert to anything that would rob you of the blessings of eternity. The choices you make here and now are forever important.

We read in 1 Corinthians: “There are … so many kinds of voices in the world.” We are surrounded by persuasive voices, beguiling voices, belittling voices, sophisticated voices, and confusing voices. I might add that these are loud voices. I admonish you to turn the volume down and to be influenced instead by that still, small voice which will guide you to safety. Remember that one with authority placed his hands on your head after you were baptized, confirming you a member of the Church and saying, “Receive the Holy Ghost.” Open your hearts, even your very souls, to the sound of that special voice which testifies of truth. As the prophet Isaiah promised, “Thine ears shall hear a word … , saying, This is the way, walk ye in it.” May we ever be in tune, that we might hear this comforting, guiding voice which will keep us safe.

Disregard for the commandments has opened the way for what I consider to be the plagues of our day. They include the plague of permissiveness, the plague of pornography, the plague of drugs, the plague of immorality, and the plague of abortion, to name just a few. The scriptures tell us that the adversary is “the founder of all these things.” We know that he is “the father of all lies, to deceive and to blind men.”

I plead with you to avoid anything that will deprive you of your happiness here in mortality and eternal life in the world to come. With his deceptions and lies, the adversary will lead you down a slippery slope to your destruction if you allow him to do so. You will likely be on that slippery slope before you even realize that there is no way to stop. You have heard the messages of the adversary. He cunningly calls: Just this once won’t matter; everyone is doing it; don’t be old-fashioned; times have changed; it can’t hurt anyone; your life is yours to live. The adversary knows us, and he knows the temptations which will be difficult for us to ignore. How vital it is that we exercise constant vigilance in order to avoid giving in to such lies and temptations.

Great courage will be required as we remain faithful and true amid the ever-increasing pressures and insidious influences with which we are surrounded and which distort the truth, tear down the good and the decent, and attempt to substitute the man-made philosophies of the world. If the commandments had been written by man, then to change them by inclination or legislation or by any other means would be the prerogative of man. The commandments, however, were God-given. Using our agency, we can set them aside. We cannot, however, change them, just as we cannot change the consequences which come from disobeying and breaking them.

May we realize that our greatest happiness in this life will come as we follow God’s commandments and obey His laws! I love the words found in Isaiah chapter 32, verse 17: “The work of righteousness shall be peace; and the effect of righteousness quietness and assurance for ever.” Such peace, such assurance can come only through righteousness.

We cannot allow ourselves the slightest bit of leeway in dealing with sin. We cannot allow ourselves to believe that we can participate “just a little” in disobeying the commandments of God, for the sin can grab us with an iron hand from which it is excruciatingly painful to free ourselves. The addictions which can come from drugs, alcohol, pornography, and immorality are real and are nearly impossible to break without great struggle and much help.

If any of you has stumbled in his journey, I assure you that there is a way back. The process is called repentance. Although the path is difficult, your eternal salvation depends on it. What could be more worthy of your efforts? I plead with you to determine right here and now to take the steps necessary to fully repent. The sooner you do so, the sooner you will be able to experience the peace and the quietness and the assurance spoken of by Isaiah.

A short while ago I heard the testimony of a woman who, with her husband, strayed from the path of safety, breaking commandments and, in the process, nearly destroying their family. When each of them could finally see through the thick haze of addiction and recognize how unhappy their lives had become, as well as how much they were hurting their loved ones, they began to change. The repentance process felt slow and was, at times, painful, but with the help of priesthood leaders, along with help from family and loyal friends, they made their way back.

I share with you a portion of this sister’s testimony of the healing power of repentance: “How does someone go from being one of the lost sheep and gripped by [sin], to this peace and happiness we now feel? How does that happen? The answer … is because of a perfect gospel, a perfect Son and His sacrifice for me. … Where there was darkness, there is now light. Where there was despair and pain, there is joy and hope. We have been infinitely blessed by the change that can only come through repentance made possible by the Atonement of Jesus Christ.”

Our Savior died to provide you and me that blessed gift. Despite the fact that the path is difficult, the promise is real. Said the Lord to those who repent:

“Though your sins be as scarlet, they shall be as white as snow.”

“And I will remember [them] no more.”

Throughout our lives we will need to nurture strong testimonies by studying the scriptures and by praying and by pondering the truths of the gospel of Jesus Christ. When firmly planted, our testimonies of the gospel, of the Savior, and of our Heavenly Father will influence all that we do.

I testify that all of us are beloved sons of our Father in Heaven, sent to earth at this day and time for a purpose, and given the priesthood of God so that we can serve others and perform God’s work here upon the earth. We have been commanded to live our lives so that we remain worthy to possess that priesthood.

My brethren, may we keep the commandments! Wonderful and glorious are the rewards which are in store for us if we do. May this be our blessing, I pray in the name of Jesus Christ, our Savior and our Redeemer, amen.

\newpage
\settalk{A Sacred Trust}
\setspeaker{Thomas S. Monson}
\settalkdate{April 2016}
\talkheading{A Sacred Trust}{Thomas S. Monson}

My beloved brethren, I pray for the Spirit to guide my remarks this evening. A common thread binds us together. We have been entrusted to bear the priesthood of God and to act in His name. We are the recipients of a sacred trust. Much is expected of us.

We read in the Doctrine and Covenants, section 121, verse 36, “The rights of the priesthood are inseparably connected with the powers of heaven.” What a wonderful gift we have been given. Ours is the responsibility to guard and protect that priesthood and to be worthy of all the glorious blessings our Father in Heaven has in store for us—and for others through us.

Wherever you go, your priesthood goes with you. Are you standing in holy places? Before you put yourself and your priesthood in jeopardy by venturing into places or participating in activities which are not worthy of you or of that priesthood, pause to consider the consequences. Remember who you are and what God expects you to become. You are a child of promise. You are a man of might. You are a son of God.

This precious gift of priesthood power brings with it not only solemn responsibilities but also special blessings for ourselves and for others. May we, in whatever place we may find ourselves, always be worthy to call upon its power, for we never know when our need and our opportunity to do so may come.

During World War II, a friend of mine was serving in the South Pacific when his plane was shot down over the ocean. He and the other crew members successfully parachuted from the burning plane, inflated their life rafts, and clung to those rafts for three days.

On the third day they spotted what they knew to be a rescue vessel. It passed them by. The next morning it passed them by again. They began to despair as they realized that this was the last day the rescue vessel would be in the area.

Then the Holy Spirit spoke to my friend: “You have the priesthood. Command the rescuers to pick you up.”

He did as prompted: “In the name of Jesus Christ and by the power of the priesthood, turn about and pick us up.”

Within a few minutes the vessel was beside them, helping them on deck. A faithful and worthy bearer of the priesthood, in his extremity, had exercised that priesthood, blessing his life and the lives of others.

May we determine, here and now, ever to be prepared for our time of need, our time of service, our time of blessing.

As we now conclude this general priesthood session, I say to you that you “are a chosen generation, a royal priesthood” (1 Peter 2:9). May we ever be worthy of these divine accolades, I pray with all my heart in the name of Jesus Christ, our Savior, amen.

\newpage
\settalk{Choices}
\setspeaker{Thomas S. Monson}
\settalkdate{April 2016}
\talkheading{Choices}{Thomas S. Monson}

Brothers and sisters, before I begin my formal message today, I would like to announce four new temples which, in coming months and years, will be built in the following locations: Quito, Ecuador; Harare, Zimbabwe; Belém, Brazil; and a second temple in Lima, Peru.

When I became a member of the Quorum of the Twelve Apostles in 1963, there were 12 operating temples in the entire Church. With the dedication of the Provo City Center Temple two weeks ago, there are now 150 temples in operation throughout the world. How grateful we are for the blessings we receive in these holy houses.

Now, brothers and sisters, I wish to express my gratitude for the opportunity to share a few thoughts with you this morning.

I have been thinking recently about choices. It has been said that the door of history turns on small hinges, and so do people’s lives. The choices we make determine our destiny.

When we left our premortal existence and entered mortality, we brought with us the gift of agency. Our goal is to obtain celestial glory, and the choices we make will, in large part, determine whether or not we reach our goal.

Most of you are familiar with Alice in Lewis Carroll’s classic novel Alice’s Adventures in Wonderland. You will remember that she comes to a crossroads with two paths before her, each stretching onward but in opposite directions. As she contemplates which way to turn, she is confronted by the Cheshire Cat, of whom Alice asks, “Which path shall I follow?”

The cat answers, “That depends where you want to go. If you do not know where you want to go, it doesn’t matter which path you take.”

Unlike Alice, we know where we want to go, and it does matter which way we go, for the path we follow in this life leads to our destination in the next life.

May we choose to build up within ourselves a great and powerful faith which will be our most effective defense against the designs of the adversary—a real faith, the kind of faith which will sustain us and will bolster our desire to choose the right. Without such faith, we go nowhere. With it, we can accomplish our goals.

Although it is imperative that we choose wisely, there are times when we will make foolish choices. The gift of repentance, provided by our Savior, enables us to correct our course settings, that we might return to the path which will lead us to that celestial glory we seek.

May we maintain the courage to defy the consensus. May we ever choose the harder right instead of the easier wrong.

As we contemplate the decisions we make in our lives each day—whether to make this choice or that choice—if we choose Christ, we will have made the correct choice.

That this may ever be so is my heartfelt and humble prayer in the name of Jesus Christ, our Lord and Savior, amen.

\newpage
\settalk{Principles and Promises}
\setspeaker{Thomas S. Monson}
\settalkdate{October 2016}
\talkheading{Principles and Promises}{Thomas S. Monson}

Tonight, brethren, I pray for the guidance of our Heavenly Father as I share my message with you.

In 1833 the Lord revealed to the Prophet Joseph Smith a plan for healthy living. That plan is found in the 89th section of the Doctrine and Covenants and is known as the Word of Wisdom. It gives specific direction regarding the food we eat, and it prohibits the use of substances which are harmful to our bodies.

Those who are obedient to the Lord’s commandments and who faithfully observe the Word of Wisdom are promised particular blessings, among which are good health and added physical stamina.

Recently I read the true account of a dramatic manifestation concerning these promises. A faithful member of the Church, John A. Larsen, served during World War II in the United States Coast Guard on the ship USS Cambria. During a battle in the Philippines, word came of an approaching squadron of bombers and kamikaze fighter planes. Orders were given for immediate evacuation. Since the USS Cambria was already gone, John and three companions gathered their gear and hurried to the beach, hoping for a lift out to one of the departing ships. Fortunately, a landing craft picked them up and sped toward the last ship leaving the bay. The men on that departing ship, in an effort to evacuate as quickly as possible, were busy on deck and had time only to throw ropes to the four men, that they might hopefully be able to climb to the deck.

John, with a heavy radio strapped to his back, found himself dangling at the end of a 40-foot (12 m) rope, at the side of a ship headed out to the open sea. He began pulling himself up, hand over hand, knowing that if he lost his grip, he would almost certainly perish. After climbing only a third of the way, he felt his arms burning with pain. He had become so weak that he felt he could no longer hold on.

With his strength depleted, as he grimly contemplated his fate, John silently cried unto God, telling Him that he had always kept the Word of Wisdom and had lived a clean life—and he now desperately needed the promised blessings.

John later said that as he finished his prayer, he felt a great surge of strength. He began climbing once again and fairly flew up the rope. When he reached the deck, his breathing was normal and not the least bit labored. The blessings of added health and stamina promised in the Word of Wisdom had been his. He gave thanks to his Heavenly Father then, and throughout the remainder of his life, for the answer to his desperate prayer for help.

Brethren, may we care for our bodies and our minds by observing the principles set forth in the Word of Wisdom, a divinely provided plan. With all my heart and soul, I testify of the glorious blessings which await us as we do. That this may be so, I pray in the name of our Lord and Savior, Jesus Christ, amen.

\newpage
\settalk{The Perfect Path to Happiness}
\setspeaker{Thomas S. Monson}
\settalkdate{October 2016}
\talkheading{The Perfect Path to Happiness}{Thomas S. Monson}

My beloved brothers and sisters, both here in the Conference Center and throughout the world, how grateful I am for the opportunity to share my thoughts with you this morning.

Fifty-two years ago, in July 1964, I had an assignment in New York City during the time the World’s Fair was hosted there. Early one morning I visited the Mormon Pavilion at the fair. I arrived just prior to a showing of the Church’s film Man’s Search for Happiness, a portrayal of the plan of salvation which has since become a Church classic. I sat next to a young man who was perhaps 35 years of age. We spoke briefly. He was not a member of our Church. Then the lights dimmed, and the show commenced.

We listened to the voice of the narrator as he posed the poignant and universal questions: Where did I come from? Why am I here? Where do I go when I leave this life? All ears strained to hear the answers, and all eyes were fixed on the images portrayed. A description of our premortal life was given, along with an explanation of our purpose on earth. We witnessed a touching depiction of the passing from this life of an elderly grandfather and of his glorious reunion with loved ones who had preceded him to the spirit world.

At the conclusion of this beautiful portrayal of our Heavenly Father’s plan for us, the crowd silently filed out, many visibly touched by the message of the film. The young visitor next to me did not arise. I asked if he had enjoyed the presentation. His emphatic response: “This is the truth!”

Our Father’s plan for our happiness and our salvation is shared by our missionaries throughout the world. Not all who hear this divine message accept and embrace it. However, men and women everywhere, just like my young friend at the New York World’s Fair, recognize its truths, and they plant their feet on the path that will lead them safely home. Their lives are forever changed.

Essential to the plan is our Savior, Jesus Christ. Without His atoning sacrifice, all would be lost. It is not enough, however, merely to believe in Him and His mission. We need to work and learn, search and pray, repent and improve. We need to know God’s laws and live them. We need to receive His saving ordinances. Only by so doing will we obtain true, eternal happiness.

We are blessed to have the truth. We have a mandate to share the truth. Let us live the truth, that we might merit all that the Father has for us. He does nothing save it be for our benefit. He has told us, “This is my work and my glory—to bring to pass the immortality and eternal life of man.”

From the depths of my soul and in all humility, I testify of the great gift which is our Father’s plan for us. It is the one perfect path to peace and happiness both here and in the world to come.

My brothers and sisters, I leave with you my love and my blessing as I close, and I do so in the name of our Savior and Redeemer, even Jesus Christ, amen.

\newpage
\settalk{Kindness, Charity, and Love}
\setspeaker{Thomas S. Monson}
\settalkdate{April 2017}
\talkheading{Kindness, Charity, and Love}{Thomas S. Monson}

My dear brethren, I am honored by the privilege to address you in this worldwide gathering of faithful holders of the priesthood of God. This evening I reference a topic about which I have spoken before.

The prophet Mormon outlined one of the key characteristics of the Savior and which is to be emulated by His disciples. Said he:

“And if a man be meek and lowly in heart, and confesses by the power of the Holy Ghost that Jesus is the Christ, he must needs have charity; for if he have not charity he is nothing; wherefore he must needs have charity.

“And charity suffereth long, and is kind, and envieth not, and is not puffed up, seeketh not her own, is not easily provoked. …

“Wherefore, my beloved brethren, if ye have not charity, ye are nothing, for charity never faileth. Wherefore, cleave unto charity, which is the greatest of all, for all things must fail—

“But charity is the pure love of Christ, and it endureth forever; and whoso is found possessed of it at the last day, it shall be well with him.”

Brethren, we do not honor the priesthood of God if we are not kind to others.

My dear friend and associate Elder Joseph B. Wirthlin was truly a kind man. Said he:

“Kindness is the essence of a celestial life. Kindness is how a Christlike person treats others. Kindness should permeate all of our words and actions at work, at school, at church, and especially in our homes.

“Jesus, our Savior, was the epitome of kindness and compassion.”

The scriptures teach us that the righteous exercise of the priesthood is dependent upon our living the principles of kindness, charity, and love. In the Doctrine and Covenants we read:

“No power or influence can or ought to be maintained by virtue of the priesthood, only by persuasion, … by gentleness and meekness, and by love unfeigned;

“By kindness, and pure knowledge, which shall greatly enlarge the soul without hypocrisy, and without guile.”

Brethren, let us examine our lives and determine to follow the Savior’s example by being kind, loving, and charitable. And as we do so, we will be in a better position to call down the powers of heaven for ourselves, for our families, and for our fellow travelers in this sometimes difficult journey back to our heavenly home. I so pray in the name of Jesus Christ the Lord, amen.

\newpage
\settalk{The Power of the Book of Mormon}
\setspeaker{Thomas S. Monson}
\settalkdate{April 2017}
\talkheading{The Power of the Book of Mormon}{Thomas S. Monson}

My dear brothers and sisters, I greet you most warmly as we are met again in a great general conference of The Church of Jesus Christ of Latter-day Saints. Before I begin my formal message today, I would like to announce five new temples which will be built in the following locations: Brasília, Brazil; greater Manila, Philippines, area; Nairobi, Kenya; Pocatello, Idaho, USA; and Saratoga Springs, Utah, USA.

This morning I speak about the power of the Book of Mormon and the critical need we have as members of this Church to study, ponder, and apply its teachings in our lives. The importance of having a firm and sure testimony of the Book of Mormon cannot be overstated.

We live in a time of great trouble and wickedness. What will protect us from the sin and evil so prevalent in the world today? I maintain that a strong testimony of our Savior, Jesus Christ, and of His gospel will help see us through to safety. If you are not reading the Book of Mormon each day, please do so. If you will read it prayerfully and with a sincere desire to know the truth, the Holy Ghost will manifest its truth to you. If it is true—and I solemnly testify that it is—then Joseph Smith was a prophet who saw God the Father and His Son, Jesus Christ.

Because the Book of Mormon is true, The Church of Jesus Christ of Latter-day Saints is the Lord’s Church on the earth, and the holy priesthood of God has been restored for the benefit and blessing of His children.

If you do not have a firm testimony of these things, do that which is necessary to obtain one. It is essential for you to have your own testimony in these difficult times, for the testimonies of others will carry you only so far. However, once obtained, a testimony needs to be kept vital and alive through continued obedience to the commandments of God and through daily prayer and scripture study.

My dear associates in the work of the Lord, I implore each of us to prayerfully study and ponder the Book of Mormon each day. As we do so, we will be in a position to hear the voice of the Spirit, to resist temptation, to overcome doubt and fear, and to receive heaven’s help in our lives. I so testify with all my heart in the name of Jesus Christ, amen.

\end{document}
